% Auto-generated from data/segments.json + layout.json (+ transforms) — do not edit by hand.
% volume: 39Abhi11
% mode: reading
% segments: 7471
% transform_rules: 0
% toc_marks: 310

% Volume layout defaults (layout.json ``layout``).
\csromanlayoutapply{1}{1.5}{21.6pt}{5pt}{6.3pt}{65pt}{2.5em}%

% Reading physical-page layout (layout.json ``page_layout_reading_mode``).
\csromanreadingpagelayoutvolume{1}{1.5}{21.6pt}{5pt}{6.3pt}{65pt}{2.5em}%
\csromanreadingpagelayoutenable

\pitaka{\textbf{อภิธมฺมปิฏก}}
\csromanheader{2}{\textbf{ธมฺมานุโลม}}
\csromanmatikaanchor{mtk.0}
\csromantocmarknopage{chapter}{ปฏฺฐานปาฬิ}{mtk.0}
\bookmark[dest={mtk.0},level=0]{ปฏฺฐานปาฬิ}
\gambhira{\textbf{ทุกปฏฺฐานปาฬิ}}
\namakkaram{นโม ตสฺส ภควโต อรหโต สมฺมาสมฺพุทฺธสฺส.}
\csromanmatikaanchor{mtk.1}
\csromanmatikaanchor{mtk.2}
\csromantocmarknopage{section}{จตุตฺถภาค}{mtk.1}
\bookmark[dest={mtk.1},level=1]{จตุตฺถภาค}
\csromantocmarknopage{chapter}{12. กิเลสโคจฺฉก}{mtk.2}
\bookmark[dest={mtk.2},level=0]{12. กิเลสโคจฺฉก}
\chapterhead{12. กิเลสโคจฺฉก}
\csromanmatikaanchor{mtk.3}
\csromantocmarkref{section}{75. กิเลสทุก}{mtk.3}
\bookmark[dest={mtk.3},level=1]{75. กิเลสทุก}
\csromanheader{1}{75. กิเลสทุก 1. ปฏิจฺจวาร}
\proseitem{1}{\csromanfolio{1}กิเลสํ ธมฺมํ ปฏิจฺจ กิเลโส ธมฺโม อุปฺปชฺชติ เหตุปจฺจยา.\csromanspacer{} โลภํ ปฏิจฺจ โมโห, ทิฏฺฐิ, ถินํ\footnote{ถีนํ (สี, สฺยา)}, อุทฺธจฺจํ, อหิริกํ, อโนตฺตปฺปํ.\csromanspacer{} โลภํ ปฏิจฺจ โมโห, ทิฏฺฐิ, อุทฺธจฺจํ, อหิริกํ, อโนตฺตปฺปํ.\csromanspacer{} โลภํ ปฏิจฺจ โมโห, มาโน, ถินํ, อุทฺธจฺจํ, อหิริกํ, อโนตฺตปฺปํ.\csromanspacer{} โลภํ ปฏิจฺจ โมโห, มาโน, อุทฺธจฺจํ, อหิริกํ, อโนตฺตปฺปํ.\csromanspacer{} โลภํ ปฏิจฺจ โมโห, ถินํ, อุทฺธจฺจํ, อหิริกํ, อโนตฺตปฺปํ.\csromanspacer{} โลภํ ปฏิจฺจ โมโห, อุทฺธจฺจํ, อหิริกํ, อโนตฺตปฺปํ.\csromanspacer{} โทสํ ปฏิจฺจ โมโห, ถินํ, อุทฺธจฺจํ, อหิริกํ, อโนตฺตปฺปํ.\csromanspacer{} โทสํ ปฏิจฺจ โมโห, อุทฺธจฺจํ, อหิริกํ, อโนตฺตปฺปํ.\csromanspacer{} วิจิกิจฺฉํ ปฏิจฺจ โมโห, อุทฺธจฺจํ, อหิริกํ, อโนตฺตปฺปํ.\csromanspacer{} อุทฺธจฺจํ ปฏิจฺจ โมโห, อหิริกํ, อโนตฺตปฺปํ. (1)}

\prose{\csromanfolio{2}กิเลสํ ธมฺมํ ปฏิจฺจ โนกิเลโส ธมฺโม อุปฺปชฺชติ เหตุปจฺจยา.\csromanspacer{} กิเลสํ ปฏิจฺจ สมฺปยุตฺตกา ขนฺธา จิตฺตสมุฏฺฐานญฺจ รูปํ. (2)}

\prose{กิเลสํ ธมฺมํ ปฏิจฺจ กิเลโส จ โนกิเลโส จ ธมฺมา อุปฺปชฺชนฺติ เหตุปจฺจยา.\csromanspacer{} โลภํ ปฏิจฺจ โมโห, ทิฏฺฐิ, ถินํ, อุทฺธจฺจํ, อหิริกํ, อโนตฺตปฺปํ สมฺปยุตฺตกา จ ขนฺธา จิตฺตสมุฏฺฐานญฺจ รูปํ. (จกฺกํ.) (3)}

\proseitem{2}{โนกิเลสํ ธมฺมํ ปฏิจฺจ โนกิเลโส ธมฺโม อุปฺปชฺชติ เหตุปจฺจยา.\csromanspacer{} โนกิเลสํ เอกํ ขนฺธํ ปฏิจฺจ ตโย ขนฺธา จิตฺตสมุฏฺฐานญฺจ รูปํ ฯเปฯ เทฺว ขนฺเธ ฯเปฯ.\csromanspacer{} ปฏิสนฺธิกฺขเณ ฯเปฯ.\csromanspacer{} ขนฺเธ ปฏิจฺจ วตฺถุ, วตฺถุํ ปฏิจฺจ ขนฺธา.\csromanspacer{} เอกํ มหาภูตํ ฯเปฯ. (1)}

\prose{โนกิเลสํ ธมฺมํ ปฏิจฺจ กิเลโส ธมฺโม อุปฺปชฺชติ เหตุปจฺจยา.\csromanspacer{} โนกิเลเส ขนฺเธ ปฏิจฺจ กิเลสา. (2)}

\prose{โนกิเลสํ ธมฺมํ ปฏิจฺจ กิเลโส จ โนกิเลโส จ ธมฺมา อุปฺปชฺชนฺติ เหตุปจฺจยา.\csromanspacer{} โนกิเลสํ เอกํ ขนฺธํ ปฏิจฺจ ตโย ขนฺธา กิเลสา จ จิตฺตสมุฏฺฐานญฺจ รูปํ ฯเปฯ เทฺว ขนฺเธ ฯเปฯ. (3)}

\proseitem{3}{กิเลสญฺจ โนกิเลสญฺจ ธมฺมํ ปฏิจฺจ กิเลโส ธมฺโม อุปฺปชฺชติ เหตุปจฺจยา.\csromanspacer{} โลภญฺจ สมฺปยุตฺตเก จ ขนฺเธ ปฏิจฺจ โมโห, ทิฏฺฐิ, ถินํ, อุทฺธจฺจํ, อหิริกํ, อโนตฺตปฺปํ. (จกฺกํ.) (1)}

\prose{กิเลสญฺจ โนกิเลสญฺจ ธมฺมํ ปฏิจฺจ โนกิเลโส ธมฺโม อุปฺปชฺชติ เหตุปจฺจยา.\csromanspacer{} โนกิเลสํ เอกํ ขนฺธญฺจ กิเลเส จ ปฏิจฺจ ตโย ขนฺธา จิตฺตสมุฏฺฐานญฺจ รูปํ ฯเปฯ เทฺว ขนฺเธ จ ฯเปฯ.\csromanspacer{} กิเลเส จ มหาภูเต จ ปฏิจฺจ จิตฺตสมุฏฺฐานํ รูปํ. (2)}

\prose{กิเลสญฺจ โนกิเลสญฺจ ธมฺมํ ปฏิจฺจ กิเลโส จ โนกิเลโส จ ธมฺมา อุปฺปชฺชนฺติ เหตุปจฺจยา.\csromanspacer{} โนกิเลสํ เอกํ ขนฺธญฺจ โลภญฺจ ปฏิจฺจ ตโย ขนฺธา โมโห, ทิฏฺฐิ, ถินํ, อุทฺธจฺจํ, อหิริกํ, อโนตฺตปฺปํ จิตฺตสมุฏฺฐานญฺจ รูปํ ฯเปฯ เทฺว ขนฺเธ จ ฯเปฯ. (จกฺกํ, สํขิตฺตํ.) (3)}

\csromanheader{2}{1. ปจฺจยานุโลม 2. สงฺขฺยาวาร}
\proseitem{4}{เหตุยา นว, อารมฺมเณ นว, อธิปติยา นว, อนนฺตเร นว, สมนนฺตเร นว, (สพฺพตฺถ นว,) วิปาเก เอกํ ฯเปฯ อวิคเต นว.}

\csromanheader{2}{2. ปจฺจยปจฺจนีย 1. วิภงฺควาร}
\csromanheader{2}{\textbf{นเหตุ}}
\proseitem{5}{\csromanfolio{3}กิเลสํ ธมฺมํ ปฏิจฺจ กิเลโส ธมฺโม อุปฺปชฺชติ นเหตุปจฺจยา.\csromanspacer{} วิจิกิจฺฉํ ปฏิจฺจ โมโห.\csromanspacer{} อุทฺธจฺจํ ปฏิจฺจ โมโห. (1)}

\prose{โนกิเลสํ ธมฺมํ ปฏิจฺจ โนกิเลโส ธมฺโม อุปฺปชฺชติ นเหตุปจฺจยา.\csromanspacer{} อเหตุกํ โนกิเลสํ เอกํ ขนฺธํ ปฏิจฺจ ตโย ขนฺธา จิตฺตสมุฏฺฐานญฺจ รูปํ ฯเปฯ เทฺว ขนฺเธ ฯเปฯ.\csromanspacer{} อเหตุกปฏิสนฺธิกฺขเณ ฯเปฯ. (ยาว อสญฺญสตฺตา.) (1)}

\prose{โนกิเลสํ ธมฺมํ ปฏิจฺจ กิเลโส ธมฺโม อุปฺปชฺชติ นเหตุปจฺจยา.\csromanspacer{} วิจิกิจฺฉาสหคเต อุทฺธจฺจสหคเต ขนฺเธ ปฏิจฺจ วิจิกิจฺฉาสหคโต อุทฺธจฺจสหคโต โมโห. (2)}

\prose{กิเลสญฺจ โนกิเลสญฺจ ธมฺมํ ปฏิจฺจ กิเลโส ธมฺโม อุปฺปชฺชติ นเหตุปจฺจยา.\csromanspacer{} วิจิกิจฺฉาสหคเต ขนฺเธ จ วิจิกิจฺฉญฺจ ปฏิจฺจ วิจิกิจฺฉาสหคโต โมโห.\csromanspacer{} อุทฺธจฺจสหคเต ขนฺเธ จ อุทฺธจฺจญฺจ ปฏิจฺจ อุทฺธจฺจสหคโต โมโห. (1)}

\csromanheader{2}{\textbf{น-อารมฺมณาทิ}}
\proseitem{6}{กิเลสํ ธมฺมํ ปฏิจฺจ โนกิเลโส ธมฺโม อุปฺปชฺชติ น-อารมฺมณปจฺจยา.\csromanspacer{} กิเลเส ปฏิจฺจ จิตฺตสมุฏฺฐานํ รูปํ. (1)}

\prose{โนกิเลสํ ธมฺมํ ปฏิจฺจ โนกิเลโส ธมฺโม อุปฺปชฺชติ น-อารมฺมณปจฺจยา.\csromanspacer{} โนกิเลเส ขนฺเธ ปฏิจฺจ จิตฺตสมุฏฺฐานํ รูปํ.\csromanspacer{} ปฏิสนฺธิกฺขเณ ฯเปฯ.\csromanspacer{} ขนฺเธ ปฏิจฺจ วตฺถุ. (ยาว อสญฺญสตฺตา.) (1)}

\prose{กิเลสญฺจ โนกิเลสญฺจ ธมฺมํ ปฏิจฺจ โนกิเลโส ธมฺโม อุปฺปชฺชติ น-อารมฺมณปจฺจยา.\csromanspacer{} กิเลเส จ สมฺปยุตฺตเก จ ขนฺเธ ปฏิจฺจ จิตฺตสมุฏฺฐานํ รูปํ.\csromanspacer{} กิเลเส จ มหาภูเต จ ปฏิจฺจ จิตฺตสมุฏฺฐานํ รูปํ. (1)}

\prose{น-อธิปติปจฺจยา.\csromanspacer{} น-อนนฺตรปจฺจยา.\csromanspacer{} นสมนนฺตรปจฺจยา.\csromanspacer{} น-อญฺญมญฺญปจฺจยา.\csromanspacer{} น-อุปนิสฺสยปจฺจยา.}

\csromanheader{2}{\textbf{นปุเรชาตาทิ}}
\proseitem{7}{\csromanfolio{4}กิเลสํ ธมฺมํ ปฏิจฺจ กิเลโส ธมฺโม อุปฺปชฺชติ นปุเรชาตปจฺจยา.\csromanspacer{} อรูเป โลภํ ปฏิจฺจ โมโห, ทิฏฺฐิ, ถินํ, อุทฺธจฺจํ, อหิริกํ, อโนตฺตปฺปํ.\csromanspacer{} โลภํ ปฏิจฺจ โมโห, ทิฏฺฐิ, อุทฺธจฺจํ, อหิริกํ, อโนตฺตปฺปํ.\csromanspacer{} โลภํ ปฏิจฺจ โมโห, มาโน, ถินํ, อุทฺธจฺจํ, อหิริกํ, อโนตฺตปฺปํ.\csromanspacer{} โลภํ ปฏิจฺจ โมโห, มาโน, อุทฺธจฺจํ, อหิริกํ, อโนตฺตปฺปํ.\csromanspacer{} โลภํ ปฏิจฺจ โมโห, ถินํ, อุทฺธจฺจํ, อหิริกํ, อโนตฺตปฺปํ.\csromanspacer{} โลภํ ปฏิจฺจ โมโห, อุทฺธจฺจํ, อหิริกํ, อโนตฺตปฺปํ.\csromanspacer{} วิจิกิจฺฉํ ปฏิจฺจ โมโห, อุทฺธจฺจํ, อหิริกํ, อโนตฺตปฺปํ.\csromanspacer{} อุทฺธจฺจํ ปฏิจฺจ โมโห, อหิริกํ, อโนตฺตปฺปํ. (อรูเป โทสมูลกํ นตฺถิ.) (1)}

\prose{กิเลสํ ธมฺมํ ปฏิจฺจ โนกิเลโส ธมฺโม อุปฺปชฺชติ นปุเรชาตปจฺจยา.\csromanspacer{} อรูเป กิเลเส ปฏิจฺจ สมฺปยุตฺตกา ขนฺธา.\csromanspacer{} กิเลเส ปฏิจฺจ จิตฺตสมุฏฺฐานํ รูปํ. (เอวํ นวปิ ปญฺหา กาตพฺพา.) นปจฺฉาชาตปจฺจยา.\csromanspacer{} น-อาเสวนปจฺจยา.}

\csromanheader{2}{\textbf{นกมฺม}}
\proseitem{8}{กิเลสํ ธมฺมํ ปฏิจฺจ โนกิเลโส ธมฺโม อุปฺปชฺชติ นกมฺมปจฺจยา.\csromanspacer{} กิเลเส ปฏิจฺจ สมฺปยุตฺตกา เจตนา. (1)}

\prose{โนกิเลสํ ธมฺมํ ปฏิจฺจ โนกิเลโส ธมฺโม อุปฺปชฺชติ นกมฺมปจฺจยา.\csromanspacer{} โนกิเลเส ขนฺเธ ปฏิจฺจ สมฺปยุตฺตกา เจตนา.\csromanspacer{} พาหิรํ.\csromanspacer{} อาหารสมุฏฺฐานํ.\csromanspacer{} อุตุสมุฏฺฐานํ ฯเปฯ. (1)}

\prose{กิเลสญฺจ โนกิเลสญฺจ ธมฺมํ ปฏิจฺจ โนกิเลโส ธมฺโม อุปฺปชฺชติ นกมฺมปจฺจยา.\csromanspacer{} กิเลเส จ สมฺปยุตฺตเก จ ขนฺเธ ปฏิจฺจ สมฺปยุตฺตกา เจตนา. (1) (เอวํ สพฺเพ ปจฺจยา กาตพฺพา.)}

\csromanheader{2}{2. ปจฺจยปจฺจนีย 2. สงฺขฺยาวาร}
\csromanheader{2}{\textbf{สุทฺธ}}
\proseitem{9}{นเหตุยา จตฺตาริ, น-อารมฺมเณ ตีณิ, น-อธิปติยา นว, น-อนนฺตเร ตีณิ, นสมนนฺตเร ตีณิ, น-อญฺญมญฺเญ ตีณิ, น-อุปนิสฺสเย ตีณิ, นปุเรชาเต นว, นปจฺฉาชาเต นว, น-อาเสวเน นว, นกมฺเม ตีณิ, \csromanfolio{5}นวิปาเก นว, น-อาหาเร เอกํ, น-อินฺทฺริเย เอกํ, นฌาเน เอกํ, นมคฺเค เอกํ, นสมฺปยุตฺเต ตีณิ, นวิปฺปยุตฺเต นว, โนนตฺถิยา ตีณิ, โนวิคเต ตีณิ.}

\vspace{\tipitakaparskip}
\csromancenter{(เอวํ อิตเร เทฺว คณนาปิ สหชาตวาโรปิ กาตพฺโพ.)}
\csromansectionrule
\csromanheader{2}{75. กิเลสทุก 3. ปจฺจยวาร}
\proseitem{10}{กิเลสํ ธมฺมํ ปจฺจยา กิเลโส ธมฺโม อุปฺปชฺชติ เหตุปจฺจยา.\csromanspacer{} ตีณิ. (ปฏิจฺจสทิสา.)}

\prose{โนกิเลสํ ธมฺมํ ปจฺจยา โนกิเลโส ธมฺโม อุปฺปชฺชติ เหตุปจฺจยา.\csromanspacer{} โนกิเลสํ เอกํ ขนฺธํ ปจฺจยา ตโย ขนฺธา ฯเปฯ. (ยาว อชฺฌตฺติกา มหาภูตา.) วตฺถุํ ปจฺจยา โนกิเลสา ขนฺธา. (1)}

\prose{โนกิเลสํ ธมฺมํ ปจฺจยา กิเลโส ธมฺโม อุปฺปชฺชติ เหตุปจฺจยา.\csromanspacer{} โนกิเลเส ขนฺเธ ปจฺจยา กิเลสา.\csromanspacer{} วตฺถุํ ปจฺจยา กิเลสา. (2)}

\prose{โนกิเลสํ ธมฺมํ ปจฺจยา กิเลโส จ โนกิเลโส จ ธมฺมา อุปฺปชฺชนฺติ เหตุปจฺจยา.\csromanspacer{} โนกิเลสํ เอกํ ขนฺธํ ปจฺจยา ตโย ขนฺธา กิเลสา จ จิตฺตสมุฏฺฐานญฺจ รูปํ ฯเปฯ เทฺว ขนฺเธ ฯเปฯ.\csromanspacer{} วตฺถุํ ปจฺจยา กิเลสา.\csromanspacer{} มหาภูเต ปจฺจยา จิตฺตสมุฏฺฐานํ รูปํ.\csromanspacer{} วตฺถุํ ปจฺจยา กิเลสา จ สมฺปยุตฺตกา จ ขนฺธา. (3)}

\proseitem{11}{กิเลสญฺจ โนกิเลสญฺจ ธมฺมํ ปจฺจยา กิเลโส ธมฺโม อุปฺปชฺชติ เหตุปจฺจยา.\csromanspacer{} โลภญฺจ สมฺปยุตฺตเก จ ขนฺเธ ปจฺจยา โมโห.\csromanspacer{} ทิฏฺฐิ, ถินํ, อุทฺธจฺจํ, อหิริกํ, อโนตฺตปฺปํ. (จกฺกํ.) โลภญฺจ วตฺถุญฺจ ปจฺจยา กิเลสา. (1)}

\prose{กิเลสญฺจ โนกิเลสญฺจ ธมฺมํ ปจฺจยา โนกิเลโส ธมฺโม อุปฺปชฺชติ เหตุปจฺจยา.\csromanspacer{} โนกิเลสํ เอกํ ขนฺธญฺจ กิเลสญฺจ ปจฺจยา ตโย ขนฺธา จิตฺตสมุฏฺฐานญฺจ รูปํ ฯเปฯ เทฺว ขนฺเธ จ ฯเปฯ.\csromanspacer{} กิเลเส จ มหาภูเต จ ปจฺจยา จิตฺตสมุฏฺฐานํ รูปํ.\csromanspacer{} กิเลเส จ วตฺถุญฺจ ปจฺจยา โนกิเลสา ขนฺธา. (2)}

\prose{\csromanfolio{6}กิเลสญฺจ โนกิเลสญฺจ ธมฺมํ ปจฺจยา กิเลโส จ โนกิเลโส จ ธมฺมา อุปฺปชฺชนฺติ เหตุปจฺจยา.\csromanspacer{} โนกิเลสํ เอกํ ขนฺธญฺจ โลภญฺจ ปจฺจยา ตโย ขนฺธา โมโห, ทิฏฺฐิ, ถินํ, อุทฺธจฺจํ, อหิริกํ, อโนตฺตปฺปํ, จิตฺตสมุฏฺฐานญฺจ รูปํ ฯเปฯ เทฺว ขนฺเธ จ ฯเปฯ. (จกฺกํ.) โลภญฺจ วตฺถุญฺจ ปจฺจยา โมโห, ทิฏฺฐิ, ถินํ, อุทฺธจฺจํ, อหิริกํ, อโนตฺตปฺปํ สมฺปยุตฺตกา จ ขนฺธา. (จกฺกํ.) (3) (อารมฺมณปจฺจเย โนกิเลสมูเล ปญฺจ วิญฺญาณา กาตพฺพา.)}

\csromanheader{2}{1. ปจฺจยานุโลม 2. สงฺขฺยาวาร}
\proseitem{12}{เหตุยา นว, อารมฺมเณ นว, อธิปติยา นว, (สพฺพตฺถ นว,) วิปาเก เอกํ ฯเปฯ อวิคเต นว.}

\csromanheader{2}{2. ปจฺจยปจฺจนีย 1. วิภงฺควาร}
\csromanheader{2}{\textbf{นเหตุ}}
\proseitem{13}{กิเลสํ ธมฺมํ ปจฺจยา กิเลโส ธมฺโม อุปฺปชฺชติ นเหตุปจฺจยา.\csromanspacer{} วิจิกิจฺฉํ ปจฺจยา วิจิกิจฺฉาสหคโต โมโห.\csromanspacer{} อุทฺธจฺจํ ปจฺจยา อุทฺธจฺจหคโต โมโห. (1)}

\prose{โนกิเลสํ ธมฺมํ ปจฺจยา โนกิเลโส ธมฺโม อุปฺปชฺชติ นเหตุปจฺจยา. (ยาว อสญฺญสตฺตา.) จกฺขายตนํ ปจฺจยา จกฺขุวิญฺญาณํ ฯเปฯ.\csromanspacer{} กายายตนํ ฯเปฯ.\csromanspacer{} วตฺถุํ ปจฺจยา อเหตุกา โนกิเลสา ขนฺธา. (1)}

\prose{โนกิเลสํ ธมฺมํ ปจฺจยา กิเลโส ธมฺโม อุปฺปชฺชติ นเหตุปจฺจยา.\csromanspacer{} วิจิกิจฺฉาคเต อุทฺธจฺจสหคเต ขนฺเธ จ วตฺถุญฺจ ปจฺจยา วิจิกิจฺฉาสหคโต อุทฺธจฺจสหคโต โมโห. (2)}

\prose{กิเลสญฺจ โนกิเลสญฺจ ธมฺมํ ปจฺจยา กิเลโส ธมฺโม อุปฺปชฺชติ นเหตุปจฺจยา.\csromanspacer{} วิจิกิจฺฉญฺจ สมฺปยุตฺตเก จ ขนฺเธ วตฺถุญฺจ ปจฺจยา วิจิกิจฺฉาสหคโต โมโห.\csromanspacer{} อุทฺธจฺจญฺจ สมฺปยุตฺตเก จ ขนฺเธ วตฺถุญฺจ ปจฺจยา อุทฺธจฺจสหคโต โมโห. (สํขิตฺตํ.) (1)}

\csromanheader{2}{2. ปจฺจยปจฺจนีย 2. สงฺขฺยาวาร}
\proseitem{14}{\csromanfolio{7}นเหตุยา จตฺตาริ, น-อารมฺมเณ ตีณิ, น-อธิปติยา นว ฯเปฯ นกมฺเม ตีณิ, นวิปาเก นว, น-อาหาเร เอกํ ฯเปฯ โนวิคเต ตีณิ.}

\vspace{\tipitakaparskip}
\csromancenter{(เอวํ อิตเร เทฺว คณนาปิ นิสฺสยวาโรปิ กาตพฺโพ.)}
\csromanheader{2}{75. กิเลสทุก 5. สํสฏฺฐวาร}
\csromanheader{2}{1. ปจฺจยานุโลมาทิ}
\proseitem{15}{กิเลสํ ธมฺมํ สํสฏฺโฐ กิเลโส ธมฺโม อุปฺปชฺชติ เหตุปจฺจยา.\csromanspacer{} โลภํ สํสฏฺโฐ โมโห, ทิฏฺฐิ, ถินํ, อุทฺธจฺจํ, อหิริกํ, อโนตฺตปฺปํ. (จกฺกํ.) (เอวํ นว ปญฺหา กาตพฺพา.)}

\prose{\textbf{เหตุ}ยา นว, อารมฺมเณ นว, (สพฺพตฺถ นว,) วิปาเก เอกํ ฯเปฯ อวิคเต นว.\csromanspacer{} อนุโลมํ.}

\prose{กิเลสํ ธมฺมํ สํสฏฺโฐ กิเลโส ธมฺโม อุปฺปชฺชติ นเหตุปจฺจยา. (เอวํ นเหตุปญฺหา จตฺตาริ กาตพฺพา.)}

\prose{นเหตุยา จตฺตาริ, น-อธิปติยา นว, นปุเรชาเต นว, นปจฺฉาชาเต นว, น-อาเสวเน นว, นกมฺเม ตีณิ, นวิปาเก นว, นฌาเน เอกํ, นมคฺเค เอกํ, นวิปฺปยุตฺเต นว.}

\vspace{\tipitakaparskip}
\csromancenter{ปจฺจนียํ.}
\csromancenter{(เอวํ อิตเร เทฺว คณนาปิ สมฺปยุตฺตวาโรปิ กาตพฺโพ.)}
\csromanheader{2}{75. กิเลสทุก 7. ปญฺหาวาร}
\proseitem{16}{กิเลโส ธมฺโม กิเลสสฺส ธมฺมสฺส เหตุปจฺจเยน ปจฺจโย.\csromanspacer{} กิเลสา เหตู สมฺปยุตฺตกานํ กิเลสานํ เหตุปจฺจเยน ปจฺจโย. \csromanfolio{8}(มูลํ ปุจฺฉิตพฺพํ.) กิเลสา เหตู สมฺปยุตฺตกานํ ขนฺธานํ จิตฺตสมุฏฺฐานานญฺจ รูปานํ เหตุปจฺจเยน ปจฺจโย. (มูลํ ปุจฺฉิตพฺพํ.) กิเลสา เหตู สมฺปยุตฺตกานํ ขนฺธานํ กิเลสานํ จิตฺตสมุฏฺฐานานญฺจ รูปานํ เหตุปจฺจเยน ปจฺจโย. (3)}

\prose{โนกิเลโส ธมฺโม โนกิเลสสฺส ธมฺมสฺส เหตุปจฺจเยน ปจฺจโย.\csromanspacer{} โนกิเลสา เหตู สมฺปยุตฺตกานํ ขนฺธานํ จิตฺตสมุฏฺฐานานญฺจ รูปานํ เหตุปจฺจเยน ปจฺจโย.\csromanspacer{} ปฏิสนฺธิกฺขเณ ฯเปฯ. (1)}

\csromanheader{2}{\textbf{อารมฺมณ}}
\proseitem{17}{กิเลโส ธมฺโม กิเลสสฺส ธมฺมสฺส อารมฺมณปจฺจเยน ปจฺจโย.\csromanspacer{} กิเลเส อารพฺภ กิเลสา อุปฺปชฺชนฺติ. (มูลํ ปุจฺฉิตพฺพํ.) กิเลเส อารพฺภ โนกิเลสา ขนฺธา อุปฺปชฺชนฺติ. (มูลํ ปุจฺฉิตพฺพํ.) กิเลเส อารพฺภ กิเลสา จ สมฺปยุตฺตกา จ ขนฺธา อุปฺปชฺชนฺติ. (3)}

\proseitem{18}{โนกิเลโส ธมฺโม โนกิเลสสฺส ธมฺมสฺส อารมฺมณปจฺจเยน ปจฺจโย.\csromanspacer{} ทานํ ฯเปฯ สีลํ ฯเปฯ อุโปสถกมฺมํ ฯเปฯ.\csromanspacer{} ปุพฺเพ สุจิณฺณานิ ฯเปฯ.\csromanspacer{} ฌานา วุฏฺฐหิตฺวา ฌานํ ปจฺจเวกฺขติ, อสฺสาเทติ อภินนฺทติ, ตํ อารพฺภ ราโค ฯเปฯ ทิฏฺฐิ ฯเปฯ วิจิกิจฺฉา ฯเปฯ อุทฺธจฺจํ ฯเปฯ โทมนสฺสํ อุปฺปชฺชติ.\csromanspacer{} อริยา มคฺคา วุฏฺฐหิตฺวา ฯเปฯ ผลสฺส, อาวชฺชนาย อารมฺมณปจฺจเยน ปจฺจโย.\csromanspacer{} จกฺขุํ ฯเปฯ.\csromanspacer{} วตฺถุํ โนกิเลเส ขนฺเธ อนิจฺจโต ฯเปฯ โทมนสฺสํ อุปฺปชฺชติ.\csromanspacer{} ทิพฺเพน จกฺขุนา รูปํ ปสฺสติ.\csromanspacer{} ทิพฺพาย โสตธาตุยา สทฺทํ สุณาติ ฯเปฯ.\csromanspacer{} อนาคตํสญาณสฺส, อาวชฺชนาย อารมฺมณปจฺจเยน ปจฺจโย. (1)}

\prose{โนกิเลโส ธมฺโม กิเลสสฺส ธมฺมสฺส อารมฺมณปจฺจเยน ปจฺจโย.\csromanspacer{} ทานํ ฯเปฯ.\csromanspacer{} ฌานา วุฏฺฐหิตฺวา ฌานํ อสฺสาเทติ อภินนฺทติ, ตํ อารพฺภ ราโค ฯเปฯ ทิฏฺฐิ ฯเปฯ วิจิกิจฺฉา ฯเปฯ อุทฺธจฺจํ อุปฺปชฺชติ.\csromanspacer{} ฌาเน ปริหีเน วิปฺปฏิสาริสฺส โทมนสฺสํ อุปฺปชฺชติ.\csromanspacer{} จกฺขุํ ฯเปฯ.\csromanspacer{} วตฺถุํ โนกิเลเส ขนฺเธ อสฺสาเทติ อภินนฺทติ, ตํ อารพฺภ ราโค ฯเปฯ โทมนสฺสํ อุปฺปชฺชติ. (2)}

\prose{โนกิเลโส ธมฺโม กิเลสสฺส จ โนกิเลสสฺส จ ธมฺมสฺส อารมฺมณปจฺจเยน ปจฺจโย.\csromanspacer{} ทานํ ฯเปฯ.\csromanspacer{} ฌานา วุฏฺฐหิตฺวา ฯเปฯ.\csromanspacer{} จกฺขุํ ฯเปฯ.\csromanspacer{} วตฺถุํ โนกิเลเส ขนฺเธ อสฺสาเทติ อภินนฺทติ, ตํ อารพฺภ กิเลสา จ สมฺปยุตฺตกา จ ขนฺธา อุปฺปชฺชนฺติ. (3)}

\prose{\csromanfolio{9}กิเลโส จ โนกิเลโส จ ธมฺมา กิเลสสฺส ธมฺมสฺส อารมฺมณปจฺจเยน ปจฺจโย.\csromanspacer{} ตีณิ.}

\csromanheader{2}{\textbf{อธิปติ}}
\proseitem{19}{กิเลโส ธมฺโม กิเลสสฺส ธมฺมสฺส อธิปติปจฺจเยน ปจฺจโย.\csromanspacer{} \textbf{อารมฺมณาธิปติ\csromandash{}}กิเลเส ครุํ กตฺวา กิเลสา อุปฺปชฺชนฺติ.\csromanspacer{} ตีณิ. (อารมฺมณาธิปติเยว.) (3)}

\prose{โนกิเลโส ธมฺโม โนกิเลสสฺส ธมฺมสฺส อธิปติปจฺจเยน ปจฺจโย.\csromanspacer{} อารมฺมณาธิปติ, สหชาตาธิปติ.\csromanspacer{}\csromanspacer{}\textbf{อารมฺมณาธิปติ\csromandash{}}ทานํ ฯเปฯ สีลํ ฯเปฯ อุโปสถกมฺมํ กตฺวา ตํ ครุํ กตฺวา ปจฺจเวกฺขติ, อสฺสาเทติ อภินนฺทติ.\csromanspacer{} ตํ ครุํ กตฺวา ราโค อุปฺปชฺชติ.\csromanspacer{} ทิฏฺฐิ อุปฺปชฺชติ.\csromanspacer{} ปุพฺเพ ฯเปฯ.\csromanspacer{} ฌานา ฯเปฯ.\csromanspacer{} อริยา มคฺคา วุฏฺฐหิตฺวา มคฺคํ ครุํ กตฺวา ปจฺจเวกฺขนฺติ ฯเปฯ ผลสฺส อธิปติปจฺจเยน ปจฺจโย.\csromanspacer{} จกฺขุํ ฯเปฯ.\csromanspacer{} วตฺถุํ โนกิเลเส ขนฺเธ ครุํ กตฺวา อสฺสาเทติ อภินนฺทติ, ตํ ครุํ กตฺวา ราโค อุปฺปชฺชติ.\csromanspacer{} ทิฏฺฐิ อุปฺปชฺชติ.\csromanspacer{} \textbf{สหชาตาธิปติ\csromandash{}}โนกิเลสาธิปติ สมฺปยุตฺตกานํ ขนฺธานํ จิตฺตสมุฏฺฐานานญฺจ รูปานํ อธิปติปจฺจเยน ปจฺจโย. (1)}

\prose{โนกิเลโส ธมฺโม กิเลสสฺส ธมฺมสฺส อธิปติปจฺจเยน ปจฺจโย.\csromanspacer{} อารมฺมณาธิปติ, สหชาตาธิปติ.\csromanspacer{}\csromanspacer{}\textbf{อารมฺมนาธิปติ\csromandash{}}ทานํ ฯเปฯ.\csromanspacer{} ฌานํ ฯเปฯ.\csromanspacer{} จกฺขุํ ฯเปฯ.\csromanspacer{} วตฺถุํ โนกิเลเส ขนฺเธ ครุํ กตฺวา อสฺสาเทติ อภินนฺทติ, ตํ ครุํ กตฺวา ราโค อุปฺปชฺชติ.\csromanspacer{} ทิฏฺฐิ อุปฺปชฺชติ.\csromanspacer{} \textbf{สหชาตาธิปติ\csromandash{}}โนกิเลสาธิปติ สมฺปยุตฺตกานํ กิเลสานํ อธิปติปจฺจเยน ปจฺจโย. (2)}

\prose{โนกิเลโส ธมฺโม กิเลสสฺส จ โนกิเลสสฺส จ ธมฺมสฺส อธิปติปจฺจเยน ปจฺจโย.\csromanspacer{} อารมฺมณาธิปติ, สหชาตาธิปติ.\csromanspacer{}\csromanspacer{}\textbf{อารมฺมณาธิปติ\csromandash{}}ทานํ ฯเปฯ.\csromanspacer{} ฌานํ ฯเปฯ.\csromanspacer{} จกฺขุํ ฯเปฯ.\csromanspacer{} วตฺถุํ โนกิเลเส ขนฺเธ ครุํ กตฺวา อสฺสาเทติ อภินนฺทติ, ตํ ครุํ กตฺวา กิเลสา จ สมฺปยุตฺตกา จ ขนฺธา อุปฺปชฺชนฺติ.\csromanspacer{} \textbf{สหชาตาธิปติ\csromandash{}}โนกิเลสาธิปติ สมฺปยุตฺตกานํ ขนฺธานํ กิเลสานญฺจ จิตฺตสมุฏฺฐานานญฺจ รูปานํ อธิปติปจฺจเยน ปจฺจโย. (3)}

\prose{กิเลโส จ โนกิเลโส จ ธมฺมา กิเลสสฺส ธมฺมสฺส อธิปติปจฺจเยน ปจฺจโย.\csromanspacer{} ตีณิ. (อารมฺมณาธิปติเยว.) (3)}

\csromanheader{2}{\textbf{อนนฺตราทิ}}
\proseitem{20}{\csromanfolio{10}กิเลโส ธมฺโม กิเลสสฺส ธมฺมสฺส อนนฺตรปจฺจเยน ปจฺจโย.\csromanspacer{} ปุริมา ปุริมา กิเลสา ปจฺฉิมานํ ปจฺฉิมานํ กิเลสานํ อนนฺตรปจฺจเยน ปจฺจโย. (มูลํ กาตพฺพํ.) ปุริมา ปุริมา กิเลสา ปจฺฉิมานํ ปจฺฉิมานํ โนกิเลสานํ ขนฺธานํ อนนฺตรปจฺจเยน ปจฺจโย.\csromanspacer{} กิเลสา วุฏฺฐานสฺส อนนฺตรปจฺจเยน ปจฺจโย. (มูลํ กาตพฺพํ.) ปุริมา ปุริมา กิเลสา ปจฺฉิมานํ ปจฺฉิมานํ กิเลสานํ สมฺปยุตฺตกานญฺจ ขนฺธานํ อนนฺตรปจฺจเยน ปจฺจโย. (3)}

\proseitem{21}{โนกิเลโส ธมฺโม โนกิเลสสฺส ธมฺมสฺส อนนฺตรปจฺจเยน ปจฺจโย.\csromanspacer{} ปุริมา ปุริมา โนกิเลสา ขนฺธา ปจฺฉิมานํ ปจฺฉิมานํ โนกิเลสานํ ขนฺธานํ อนนฺตรปจฺจเยน ปจฺจโย ฯเปฯ ผลสมาปตฺติยา อนนฺตรปจฺจเยน ปจฺจโย. (มูลํ กาตพฺพํ.) ปุริมา ปุริมา โนกิเลสา ขนฺธา ปจฺฉิมานํ ปจฺฉิมานํ กิเลสานํ อนนฺตรปจฺจเยน ปจฺจโย.\csromanspacer{} อาวชฺชนา กิเลสานํ อนนฺตรปจฺจเยน ปจฺจโย. (มูลํ กาตพฺพํ.) ปุริมา ปุริมา โนกิเลสา ขนฺธา ปจฺฉิมานํ ปจฺฉิมานํ กิเลสานํ สมฺปยุตฺตกานญฺจ ขนฺธานํ อนนฺตรปจฺจเยน ปจฺจโย.\csromanspacer{} อาวชฺชนา กิเลสานํ สมฺปยุตฺตกานญฺจ ขนฺธานํ อนนฺตรปจฺจเยน ปจฺจโย. (3)}

\proseitem{22}{กิเลโส จ โนกิเลโส จ ธมฺมา กิเลสสฺส ธมฺมสฺส อนนฺตรปจฺจเยน ปจฺจโย.\csromanspacer{} ปุริมา ปุริมา กิเลสา จ สมฺปยุตฺตกา จ ขนฺธา ปจฺฉิมานํ ปจฺฉิมานํ กิเลสานํ อนนฺตรปจฺจเยน ปจฺจโย. (มูลํ กาตพฺพํ.) ปุริมา ปุริมา กิเลสา จ สมฺปยุตฺตกา จ ขนฺธา ปจฺฉิมานํ ปจฺฉิมานํ โนกิเลสานํ ขนฺธานํ อนนฺตรปจฺจเยน ปจฺจโย.\csromanspacer{} กิเลสา จ สมฺปยุตฺตกา จ ขนฺธา วุฏฺฐานสฺส อนนฺตรปจฺจเยน ปจฺจโย. (มูลํ กาตพฺพํ.) ปุริมา ปุริมา กิเลสา จ สมฺปยุตฺตกา จ ขนฺธา ปจฺฉิมานํ ปจฺฉิมานํ กิเลสานํ สมฺปยุตฺตกานญฺจ ขนฺธานํ อนนฺตรปจฺจเยน ปจฺจโย. (3)}

\prose{สมนนฺตรปจฺจเยน ปจฺจโย.\csromanspacer{} สหชาตปจฺจเยน ปจฺจโย.\csromanspacer{} อญฺญมญฺญปจฺจเยน ปจฺจโย.\csromanspacer{} นิสฺสยปจฺจเยน ปจฺจโย.}

\csromanheader{2}{\textbf{อุปนิสฺสย}}
\proseitem{23}{\csromanfolio{11}กิเลโส ธมฺโม กิเลสสฺส ธมฺมสฺส อุปนิสฺสยปจฺจเยน ปจฺจโย.\csromanspacer{} \textbf{อารมฺมณูปนิสฺสโย, อนนฺตรูปนิสฺสโย, ปกตูปนิสฺสโย ฯเปฯ.}\csromanspacer{} ปกตูปนิสฺสโย\csromandash{}กิเลสา กิเลสานํ อุปนิสฺสยปจฺจเยน ปจฺจโย.\csromanspacer{} ตีณิ.}

\proseitem{24}{โนกิเลโส ธมฺโม โนกิเลสสฺส ธมฺมสฺส อุปนิสฺสยปจฺจเยน ปจฺจโย.\csromanspacer{} \textbf{อารมฺมณูปนิสฺสโย, อนนฺตรูปนิสฺสโย, ปกตูปนิสฺสโย ฯเปฯ.}\csromanspacer{} ปกตูปนิสฺสโย\csromandash{}สทฺธํ อุปนิสฺสาย ทานํ เทติ ฯเปฯ.\csromanspacer{} มานํ ชปฺเปติ.\csromanspacer{} ทิฏฺฐิํ คณฺหาติ.\csromanspacer{} สีลํ ฯเปฯ.\csromanspacer{} ปญฺญํ.\csromanspacer{} ราคํ.\csromanspacer{} โทสํ.\csromanspacer{} โมหํ.\csromanspacer{} มานํ.\csromanspacer{} ทิฏฺฐิํ.\csromanspacer{} ปตฺถนํ.\csromanspacer{} กายิกํ สุขํ ฯเปฯ.\csromanspacer{} เสนาสนํ อุปนิสฺสาย ทานํ เทติ ฯเปฯ สํฆํ ภินฺทติ.\csromanspacer{} สทฺธา ฯเปฯ.\csromanspacer{} เสนาสนํ สทฺธาย ฯเปฯ ผลสมาปตฺติยา อุปนิสฺสยปจฺจเยน ปจฺจโย. (1)}

\prose{โนกิเลโส ธมฺโม กิเลสสฺส ธมฺมสฺส อุปนิสฺสยปจฺจเยน ปจฺจโย.\csromanspacer{} \textbf{อารมฺมณูปนิสฺสโย, อนนฺตรูปนิสฺสโย, ปกตูปนิสฺสโย ฯเปฯ.}\csromanspacer{} ปกตูปนิสฺสโย\csromandash{}สทฺธํ อุปนิสฺสาย มานํ ชปฺเปติ.\csromanspacer{} ทิฏฺฐิํ คณฺหาติ.\csromanspacer{} สีลํ ฯเปฯ.\csromanspacer{} เสนาสนํ อุปนิสฺสาย ปาณํ หนติ ฯเปฯ สํฆํ ภินฺทติ, สทฺธา ฯเปฯ.\csromanspacer{} เสนาสนํ กิเลสานํ อุปนิสฺสยปจฺจเยน ปจฺจโย. (2)}

\prose{โนกิเลโส ธมฺโม กิเลสสฺส จ โนกิเลสสฺส จ ธมฺมสฺส อุปนิสฺสยปจฺจเยน ปจฺจโย.\csromanspacer{} \textbf{อารมฺมณูปนิสฺสโย, อนนฺตรูปนิสฺสโย, ปกตูปนิสฺสโย ฯเปฯ.}\csromanspacer{} ปกตูปนิสฺสโย\csromandash{}สทฺธํ อุปนิสฺสาย มานํ ชปฺเปติ.\csromanspacer{} ทิฏฺฐิํ คณฺหาติ.\csromanspacer{} สีลํ ฯเปฯ.\csromanspacer{} เสนาสนํ อุปนิสฺสาย ปาณํ หนติ ฯเปฯ สํฆํ ภินฺทติ.\csromanspacer{} สทฺธา ฯเปฯ.\csromanspacer{} เสนาสนํ กิเลสานํ สมฺปยุตฺตกานญฺจ ขนฺธานํ อุปนิสฺสยปจฺจเยน ปจฺจโย. (3)}

\prose{กิเลโส จ โนกิเลโส จ ธมฺมา กิเลสสฺส ธมฺมสฺส อุปนิสฺสยปจฺจเยน ปจฺจโย.\csromanspacer{} ตีณิ.}

\csromanheader{2}{\textbf{ปุเรชาต}}
\proseitem{25}{โนกิเลโส ธมฺโม โนกิเลสสฺส ธมฺมสฺส ปุเรชาตปจฺจเยน ปจฺจโย.\csromanspacer{} อารมฺมณปุเรชาตํ, วตฺถุปุเรชาตํ.\csromanspacer{}\csromanspacer{}\textbf{อารมฺมณปุเรชาตํ\csromandash{}} จกฺขุํ ฯเปฯ.\csromanspacer{} วตฺถุํ อนิจฺจโต ฯเปฯ โทมนสฺสํ อุปฺปชฺชติ.\csromanspacer{} ทิพฺเพน \csromanfolio{12}จกฺขุนา รูปํ ปสฺสติ.\csromanspacer{} ทิพฺพาย โสตธาตุยา สทฺทํ สุณาติ.\csromanspacer{} รูปายตนํ จกฺขุวิญฺญาณสฺส ฯเปฯ.\csromanspacer{} โผฏฺฐพฺพายตนํ กายวิญฺญาณสฺส ฯเปฯ.\csromanspacer{} \textbf{วตฺถุปุเรชาตํ\csromandash{}}จกฺขายตนํ จกฺขุวิญฺญาณสฺส ฯเปฯ.\csromanspacer{} กายายตนํ กายวิญฺญาณสฺส ฯเปฯ.\csromanspacer{} วตฺถุ โนกิเลสานํ ขนฺธานํ ปุเรชาตปจฺจเยน ปจฺจโย. (1)}

\prose{โนกิเลโส ธมฺโม กิเลสสฺส ธมฺมสฺส ปุเรชาตปจฺจเยน ปจฺจโย.\csromanspacer{} อารมฺมณปุเรชาตํ, วตฺถุปุเรชาตํ.\csromanspacer{}\csromanspacer{}\textbf{อารมฺมณปุเรชาตํ\csromandash{}}จกฺขุํ ฯเปฯ.\csromanspacer{} วตฺถุํ อสฺสาเทติ อภินนฺทติ, ตํ อารพฺภ ราโค ฯเปฯ โทมนสฺสํ อุปฺปชฺชติ.\csromanspacer{} \textbf{วตฺถุปุเรชาตํ\csromandash{}}วตฺถุ กิเลสานํ ปุเรชาตปจฺจเยน ปจฺจโย. (2)}

\prose{โนกิเลโส ธมฺโม กิเลสสฺส จ โนกิเลสสฺส จ ธมฺมสฺส ปุเรชาตปจฺจเยน ปจฺจโย.\csromanspacer{} อารมฺมณปุเรชาตํ, วตฺถุปุเรชาตํ.\csromanspacer{}\csromanspacer{}\textbf{อารมฺมณปุเรชาตํ\csromandash{}}จกฺขุํ ฯเปฯ.\csromanspacer{} วตฺถุํ อสฺสาเทติ อภินนฺทติ, ตํ อารพฺภ กิเลสา จ สมฺปยุตฺตกา จ ขนฺธา อุปฺปชฺชนฺติ.\csromanspacer{} \textbf{วตฺถุปุเรชาตํ\csromandash{}} วตฺถุ กิเลสานํ สมฺปยุตฺตกานญฺจ ขนฺธานํ ปุเรชาตปจฺจเยน ปจฺจโย. (3)}

\csromanheader{2}{\textbf{ปจฺฉาชาตาเสวน}}
\proseitem{26}{กิเลโส ธมฺโม โนกิเลสสฺส ธมฺมสฺส ปจฺฉาชาตปจฺจเยน ปจฺจโย. (สํขิตฺตํ.) (1)}

\prose{โนกิเลโส ธมฺโม โนกิเลสสฺส ธมฺมสฺส ปจฺฉาชาตปจฺจเยน ปจฺจโย. (สํขิตฺตํ.) (1)}

\prose{กิเลโส จ โนกิเลโส จ ธมฺมา โนกิเลสสฺส ธมฺมสฺส ปจฺฉาชาตปจฺจเยน ปจฺจโย. (สํขิตฺตํ.) อาเสวนปจฺจเยน ปจฺจโย.\csromanspacer{} นว.}

\csromanheader{2}{\textbf{กมฺม}}
\proseitem{27}{โนกิเลโส ธมฺโม โนกิเลสสฺส ธมฺมสฺส กมฺมปจฺจเยน ปจฺจโย.\csromanspacer{} \textbf{สหชาตา}, นานากฺขณิกา.\csromanspacer{}\csromanspacer{}\textbf{สหชาตา} โนกิเลสา เจตนา สมฺปยุตฺตกานํ ขนฺธานํ จิตฺตสมุฏฺฐานานญฺจ รูปานํ กมฺมปจฺจเยน ปจฺจโย.\csromanspacer{} ปฏิสนฺธิกฺขเณ ฯเปฯ.\csromanspacer{} \textbf{นานากฺขณิกา} โนกิเลสา เจตนา วิปากานํ ขนฺธานํ กฏตฺตา จ รูปานํ กมฺมปจฺจเยน ปจฺจโย. (1)}

\prose{\csromanfolio{13}โนกิเลโส ธมฺโม กิเลสสฺส ธมฺมสฺส กมฺมปจฺจเยน ปจฺจโย.\csromanspacer{} โนกิเลสา เจตนา กิเลสานํ กมฺมปจฺจเยน ปจฺจโย. (2)}

\prose{โนกิเลโส ธมฺโม กิเลสสฺส จ โนกิเลสสฺส จ ธมฺมสฺส กมฺมปจฺจเยน ปจฺจโย.\csromanspacer{} โนกิเลสา เจตนา สมฺปยุตฺตกานํ ขนฺธานํ กิเลสานํ จิตฺตสมุฏฺฐานานญฺจ รูปานํ กมฺมปจฺจเยน ปจฺจโย. (3)}

\csromanheader{2}{\textbf{วิปากาทิ}}
\proseitem{28}{โนกิเลโส ธมฺโม โนกิเลสสฺส ธมฺมสฺส วิปากปจฺจเยน ปจฺจโย.\csromanspacer{} เอกํ.\csromanspacer{} อาหารปจฺจเยน ปจฺจโย.\csromanspacer{} ตีณิ.\csromanspacer{} อินฺทฺริยปจฺจเยน ปจฺจโย.\csromanspacer{} ตีณิ.\csromanspacer{} ฌานปจฺจเยน ปจฺจโย.\csromanspacer{} ตีณิ.\csromanspacer{} มคฺคปจฺจเยน ปจฺจโย.\csromanspacer{} นว.\csromanspacer{} สมฺปยุตฺตปจฺจเยน ปจฺจโย.\csromanspacer{} นว.}

\csromanheader{2}{\textbf{วิปยุตฺต}}
\proseitem{29}{กิเลโส ธมฺโม โนกิเลสสฺส ธมฺมสฺส วิปฺปยุตฺตปจฺจเยน ปจฺจโย.\csromanspacer{} สหชาตํ, ปจฺฉาชาตํ. (สํขิตฺตํ.) (1)}

\prose{โนกิเลโส ธมฺโม โนกิเลสสฺส ธมฺมสฺส วิปฺปยุตฺตปจฺจเยน ปจฺจโย.\csromanspacer{} สหชาตํ, ปุเรชาตํ, ปจฺฉาชาตํ. (สํขิตฺตํ.) (1)}

\prose{โนกิเลโส ธมฺโม กิเลสสฺส ธมฺมสฺส วิปฺปยุตฺตปจฺจเยน ปจฺจโย.\csromanspacer{} \textbf{ปุเรชาตํ} วตฺถุ กิเลสานํ วิปฺปยุตฺตปจฺจเยน ปจฺจโย. (2)}

\prose{โนกิเลโส ธมฺโม กิเลสสฺส จ โนกิเลสสฺส จ ธมฺมสฺส วิปฺปยุตฺตปจฺจเยน ปจฺจโย.\csromanspacer{} \textbf{ปุเรชาตํ} วตฺถุ กิเลสานํ สมฺปยุตฺตกานญฺจ ขนฺธานํ วิปฺปยุตฺตปจฺจเยน ปจฺจโย. (3)}

\prose{กิเลโส จ โนกิเลโส จ ธมฺมา โนกิเลสสฺส ธมฺมสฺส วิปฺปยุตฺตปจเยน ปจฺจโย.\csromanspacer{} สหชาตํ, ปจฺฉาชตํ. (สํขิตฺตํ.\csromanspacer{} วิตฺถาเรตพฺพํ.)}

\csromanheader{2}{\textbf{อตฺถิ-อาทิ}}
\proseitem{30}{กิเลโส ธมฺโม กิเลสสฺส ธมฺมสฺส อตฺถิปจฺจเยน ปจฺจโย.\csromanspacer{} เอกํ. (ปฏิจฺจสทิสํ.) (1)}

\prose{กิเลโส ธมฺโม โนกิเลสสฺส ธมฺมสฺส อตฺถิปจฺจเยน ปจฺจโย.\csromanspacer{} สหชาตํ, ปจฺฉาชาตํ. (สํขิตฺตํ.) (2)}

\prose{\csromanfolio{14}กิเลโส ธมฺโม กิเลสสฺส จ โนกิเลสสฺส จ ธมฺมสฺส อตฺถิปจฺจเยน ปจฺจโย. (ปฏิจฺจสทิสํ.) (3)}

\prose{โนกิเลโส ธมฺโม โนกิเลสสฺส ธมฺมสฺส อตฺถิปจฺจเยน ปจฺจโย.\csromanspacer{} สหชาตํ, ปุเรชาตํ, ปจฺฉาชาตํ, อาหารํ, อินฺทฺริยํ. (สํขิตฺตํ.) (1)}

\prose{โนกิเลโส ธมฺโม กิเลสสฺส ธมฺมสฺส อตฺถิปจฺจเยน ปจฺจโย.\csromanspacer{} สหชาตํ, ปุเรชาตํ. (สํขิตฺตํ.\csromanspacer{} สหชาตํ สหชาตสทิสํ.\csromanspacer{} ปุเรชาตํ ปุเรชาตสทิสํ).\csromanspacer{}\csromanspacer{}โนกิเลโส ธมฺโม กิเลสสฺส จ โนกิเลสสฺส จ ธมฺมสฺส อตฺถิปจฺจเยน ปจฺจโย.\csromanspacer{} สหชาตํ, ปุเรชาตํ. (สหชาตํ สหชาตสทิสํ.\csromanspacer{} ปุเรชาตํ ปุเรชาตสทิสํ.) (3)}

\proseitem{31}{กิเลโส จ โนกิเลโส จ ธมฺมา กิเลสสฺส ธมฺมสฺส อตฺถิปจฺจเยน ปจฺจโย.\csromanspacer{} สหชาตํ, ปุเรชาตํ.\csromanspacer{}\csromanspacer{}\textbf{สหชาโต} โลโภ จ สมฺปยุตฺตกา จ ขนฺธา โมหสฺส, ทิฏฺฐิยา, ถินสฺส, อุทฺธจฺจสฺส, อหิริกสฺส, อโนตฺตปฺปสฺส อตฺถิปจฺจเยน ปจฺจโย. (จกฺกํ.) \textbf{สหชาโต} โลโภ จ วตฺถุ จ โมหสฺส, ทิฏฺฐิยา, ถินสฺส, อุทฺธจฺจสฺส, อหิริกสฺส, อโนตฺตปฺปสฺส อตฺถิปจฺจเยน ปจฺจโย. (จกฺกํ.) (1)}

\prose{กิเลโส จ โนกิเลโส จ ธมฺมา โนกิเลสสฺส ธมฺมสฺส อตฺถิปจฺจเยน ปจฺจโย.\csromanspacer{} สหชาตํ, ปุเรชาตํ, ปจฺฉาชาตํ, อาหารํ, อินฺทฺริยํ.\csromanspacer{}\csromanspacer{}\textbf{สหชาโต} โนกิเลโส เอโก ขนฺโธ จ กิเลโส จ ติณฺณนฺนํ ขนฺธานํ จิตฺตสมุฏฺฐานานญฺจ รูปานํ อตฺถิปจฺจเยน ปจฺจโย ฯเปฯ เทฺว ขนฺธา จ ฯเปฯ.\csromanspacer{} \textbf{สหชาตา} กิเลสา จ มหาภูตา จ จิตฺตสมุฏฺฐานานํ รูปานํ อตฺถิปจฺจเยน ปจฺจโย.\csromanspacer{} \textbf{สหชาตา} กิเลสา จ วตฺถุ จ โนกิเลสานํ ขนฺธานํ อตฺถิปจฺจเยน ปจฺจโย.\csromanspacer{} \textbf{ปจฺฉาชาตา} กิเลสา จ สมฺปยุตฺตกา ขนฺธา จ ปุเรชาตสฺส อิมสฺส กายสฺส อตฺถิปจฺจเยน ปจฺจโย.\csromanspacer{} \textbf{ปจฺฉาชาตา} กิเลสา จ สมฺปยุตฺตกา จ ขนฺธา กพฬีกาโร อาหาโร จ อิมสฺส กายสฺส อตฺถิปจฺจเยน ปจฺจโย.\csromanspacer{} \textbf{ปจฺฉาชาตา} กิเลสา จ สมฺปยุตฺตกา จ ขนฺธา รูปชีวิตินฺทฺริยญฺจ กฏตฺตารูปานํ อตฺถิปจฺจเยน ปจฺจโย. (2)}

\prose{กิเลโส จ โนกิเลโส จ ธมฺมา กิเลสสฺส จ โนกิเลสสฺส จ ธมฺมสฺส อตฺถิปจฺจเยน ปจฺจโย.\csromanspacer{} สหชาตํ, ปุเรชาตํ.\csromanspacer{}\csromanspacer{} \csromanfolio{15}\textbf{สหชาโต} โนกิเลโส เอโก ขนฺโธ จ โลโภ จ ติณฺณนฺนํ ขนฺธานํ จิตฺตสมุฏฺฐานานญฺจ รูปานํ โมหสฺส จ, ทิฏฺฐิยา, ถินสฺส, อุทฺธจฺจสฺส, อหิริกสฺส, อโนตฺตปฺปสฺส อตฺถิปจฺจเยน ปจฺจโย ฯเปฯ เทฺว ขนฺธา จ ฯเปฯ.\csromanspacer{} \textbf{สหชาโต} โลโภ จ วตฺถุ จ โมหสฺส, ทิฏฺฐิยา, ถินสฺส, อุทฺธจฺจสฺส, อหิริกสฺส, อโนตฺตปฺปสฺส สมฺปยุตฺตกานญฺจ ขนฺธานํ อตฺถิปจฺจเยน ปจฺจโย. (จกฺกํ.) (3)}

\prose{นตฺถิปจฺจเยน ปจฺจโย.\csromanspacer{} วิคตปจฺจเยน ปจฺจโย.\csromanspacer{} อวิคตปจฺจเยน ปจฺจโย.}

\csromanheader{2}{1. ปจฺจยานุโลม 2. สงฺขฺยาวาร}
\csromanheader{2}{\textbf{สุทฺธ}}
\proseitem{32}{เหตุยา จตฺตาริ, อารมฺมเณ นว, อธิปติยา นว, อนนฺตเร นว, สมนนฺตเร นว, สหชาเต นว, อญฺญมญฺเญ นว, นิสฺสเย นว, อุปนิสฺสเย นว, ปุเรชาเต ตีณิ, ปจฺฉาชาเต ตีณิ, อาเสวเน นว, กมฺเม ตีณิ, วิปาเก เอกํ, อาหาเร ตีณิ, อินฺทฺริเย ตีณิ, ฌาเน ตีณิ, มคฺเค นว, สมฺปยุตฺเต นว, วิปฺปยุตฺเต ปญฺจ, อตฺถิยา นว, นตฺถิยา นว, วิคเต นว, อวิคเต นว.}

\csromanheader{2}{2. ปจฺจนียุทฺธาร}
\proseitem{33}{กิเลโส ธมฺโม กิเลสสฺส ธมฺมสฺส อารมฺมณปจฺจเยน ปจฺจโย.\csromanspacer{} สหชตปจฺจเยน ปจฺจโย.\csromanspacer{} อุปนิสฺสยปจฺจเยน ปจฺจโย. (1)}

\prose{กิเลโส ธมฺโม โนกิเลสสฺส ธมฺมสฺส อารมฺมณปจฺจเยน ปจฺจโย.\csromanspacer{} สหชาตปจฺจเยน ปจฺจโย.\csromanspacer{} อุปนิสฺสยปจฺจเยน ปจฺจโย.\csromanspacer{} ปจฺฉาชาตปจฺจเยน ปจฺจโย. (2)}

\prose{กิเลโส ธมฺโม กิเลสสฺส จ โนกิเลสสฺส จ ธมฺมสฺส อารมฺมณปจฺจเยน ปจฺจโย.\csromanspacer{} สหชาตปจฺจเยน ปจฺจโย.\csromanspacer{} อุปนิสฺสยปจฺจเยน ปจฺจโย. (3)}

\proseitem{34}{โนกิเลโส ธมฺโม โนกิเลสสฺส ธมฺมสฺส อารมฺมณปจฺจเยน ปจฺจโย.\csromanspacer{} สหชาตปจฺจเยน ปจฺจโย.\csromanspacer{} อุปนิสฺสยปจฺจเยน \csromanfolio{16}ปจฺจโย.\csromanspacer{} ปุเรชาตปจฺจเยน ปจฺจโย.\csromanspacer{} ปจฺฉาชาตปจฺจเยน ปจฺจโย.\csromanspacer{} กมฺมปจฺจเยน ปจฺจโย.\csromanspacer{} อาหารปจฺจเยน ปจฺจโย.\csromanspacer{} อินฺทฺริยปจฺจเยน ปจฺจโย. (1)}

\prose{โนกิเลโส ธมฺโม กิเลสสฺส ธมฺมสฺส อารมฺมณปจฺจเยน ปจฺจโย.\csromanspacer{} สหชาตปจฺจเยน ปจฺจโย.\csromanspacer{} อุปนิสฺสยปจฺจเยน ปจฺจโย.\csromanspacer{} ปุเรชาตปจฺจเยน ปจฺจโย. (2)}

\prose{โนกิเลโส ธมฺโม กิเลสสฺส จ โนกิเลสสฺส จ ธมฺมสฺส อารมฺมณปจฺจเยน ปจฺจโย.\csromanspacer{} สหชาตปจฺจเยน ปจฺจโย.\csromanspacer{} อุปนิสฺสยปจฺจเยน ปจฺจโย.\csromanspacer{} ปุเรชาตปจฺจเยน ปจฺจโย. (3)}

\proseitem{35}{กิเลโส จ โนกิเลโส จ ธมฺมา กิเลสสฺส ธมฺมสฺส อารมฺมณปจฺจเยน ปจฺจโย.\csromanspacer{} สหชาตปจฺจเยน ปจฺจโย.\csromanspacer{} อุปนิสฺสยปจฺจเยน ปจฺจโย. (1)}

\prose{กิเลโส จ โนกิเลโส จ ธมฺมา โนกิเลสสฺส ธมฺมสฺส อารมฺมณปจฺจเยน ปจฺจโย.\csromanspacer{} สหชาตปจฺจเยน ปจฺจโย.\csromanspacer{} อุปนิสฺสยปจฺจเยน ปจฺจโย.\csromanspacer{} ปจฺฉาชาตปจฺจเยน ปจฺจโย. (2)}

\prose{กิเลโส จ โนกิเลโส จ ธมฺมา กิเลสสฺส จ โนกิเลสสฺส จ ธมฺมสฺส อารมฺมณปจฺจเยน ปจฺจโย.\csromanspacer{} สหชาตปจฺจเยน ปจฺจโย.\csromanspacer{} อุปนิสฺสยปจฺจเยน ปจฺจโย. (3)}

\csromanheader{2}{2. ปจฺจยปจฺจนีย 2. สงฺขฺยาวาร}
\proseitem{36}{นเหตุยา นว, น-อารมฺมเณ นว, น-อธิปติยา นว, (สพฺพตฺถ นว,) โน-อวิคเต นว.}

\csromanheader{2}{3. ปจฺจยานุโลมปจฺจนีย}
\proseitem{37}{เหตุปจฺจยา น-อารมฺมเณ จตฺตาริ, น-อธิปติยา จตฺตาริ, น-อนนฺตเร จตฺตาริ, นสมนนฺตเร จตฺตาริ, น-อญฺญมญฺเญ เทฺว, น-อุปนิสฺสเย จตฺตาริ, (สพฺพตฺถ จตฺตาริ,) นสมฺปยุตฺเต เทฺว, นวิปฺปยุตฺเต จตฺตาริ, โนนตฺถิยา จตฺตาริ, โนวิคเต จตฺตาริ.}

\csromanmatikaanchor{mtk.4}
\csromanmatikaanchor{mtk.5}
\csromantocmarkref{section}{76. สํกิเลสทุก}{mtk.4}
\bookmark[dest={mtk.4},level=1]{76. สํกิเลสทุก}
\csromantocmarkref{section}{77. สํกิลิฏฺฐทุก}{mtk.5}
\bookmark[dest={mtk.5},level=1]{77. สํกิลิฏฺฐทุก}
\csromanheader{1}{77. สํกิลิฏฺฐทุก \textbf{4. ปจฺจยปจฺจนียานุโลม}}
\proseitem{38}{\csromanfolio{17}นเหตุปจฺจยา อารมฺมเณ นว, อธิปติยา นว, (อนุโลมมาติกา กาตพฺพา.) ฯเปฯ อวิคเต นว.}

\nitthitam{กิเลสทุกํ นิฏฺฐิตํ.}
\csromanheader{2}{76. สํกิเลสิกทุก 1. ปฏิจฺจวาร}
\vspace{\tipitakaparskip}
\csromancenter{สํกิเลสิกํ ธมฺมํ ปฏิจฺจ สํกิเลสิโก ธมฺโม อุปฺปชฺชติ เหตุปจฺจยา.}
\csromancenter{(ยถา โลกิยทุกํ, เอวํ นินฺนานากรณํ.)}
\nitthitam{สํกิเลสิกทุกํ นิฏฺฐิตํ.}
\csromanheader{2}{77. สํกิลิฏฺฐทุก 1. ปฏิจฺจวาร}
\proseitem{40}{สํกิลิฏฺฐํ ธมฺมํ ปฏิจฺจ สํกิลิฏฺโฐ ธมฺโม อุปฺปชฺชติ เหตุปจฺจยา.\csromanspacer{} สํกิลิฏฺฐํ เอกํ ขนฺธํ ปฏิจฺจ ตโย ขนฺธา ฯเปฯ เทฺว ขนฺเธ ฯเปฯ. (1)}

\prose{สํกิลิฏฺฐํ ธมฺมํ ปฏิจฺจ อสํกิลิฏฺโฐ ธมฺโม อุปฺปชฺชติ เหตุปจฺจยา.\csromanspacer{} สํกิลิฏฺเฐ ขนฺเธ ปฏิจฺจ จิตฺตสมุฏฺฐานํ รูปํ. (2)}

\prose{สํกิลิฏฺฐํ ธมฺมํ ปฏิจฺจ สํกิลิฏฺโฐ จ อสํกิลิฏฺโฐ จ ธมฺมา อุปฺปชฺชนฺติ เหตุปจฺจยา.\csromanspacer{} สํกิลิฏฺฐํ เอกํ ขนฺธํ ปฏิจฺจ ตโย ขนฺธา จิตฺตสมุฏฺฐานญฺจ รูปํ ฯเปฯ เทฺว ขนฺเธ ฯเปฯ. (3)}

\prose{อสํกิลิฏฺฐํ ธมฺมํ ปฏิจฺจ อสํกิลิฏฺโฐ ธมฺโม อุปฺปชฺชติ เหตุปจฺจยา.\csromanspacer{} อสํกิลิฏฺฐํ เอกํ ขนฺธํ ปฏิจฺจ ตโย ขนฺธา จิตฺตสมุฏฺฐานญฺจ รูปํ ฯเปฯ เทฺว ขนฺเธ ฯเปฯ.\csromanspacer{} ปฏิสนฺธิกฺขเณ ฯเปฯ.\csromanspacer{} ขนฺเธ ปฏิจฺจ วตฺถุ, ปตฺถุํ ปฏิจฺจ ขนฺธา.\csromanspacer{} เอกํ มหาภูตํ ฯเปฯ. (1)}

\prose{\csromanfolio{18}สํกิลิฏฺฐญฺจ อสํกิลิฏฺฐญฺจ ธมฺมํ ปฏิจฺจ อสํกิลิฏฺโฐ ธมฺโม อุปฺปชฺชติ เหตุปจฺจยา.\csromanspacer{} สํกิลิฏฺเฐ ขนฺเธ จ มหาภูเต จ ปฏิจฺจ จิตฺตสมุฏฺฐานํ รูปํ. (สํขิตฺตํ.) (1)}

\csromanheader{2}{1. ปจฺจยานุโลม 2. สงฺขฺยาวาร}
\proseitem{41}{เหตุยา ปญฺจ, อารมฺมเณ เทฺว, อธิปติยา ปญฺจ, อนนฺตเร เทฺว, สมนนฺตเร เทฺว, สหชาเต ปญฺจ, อญฺญมญฺเญ เทฺว, นิสฺสเย ปญฺจ, อุปนิสฺสเย เทฺว, ปุเรชาเต เทฺว, อาเสวเน เทฺว, กมฺเม ปญฺจ, วิปาเก เอกํ อาหาเร ปญฺจ, (สํขิตฺตํ.) อวิคเต ปญฺจ.}

\csromanheader{2}{2. ปจฺจยปจฺจนีย 1. วิภงฺควาร}
\csromanheader{2}{\textbf{นเหตุ}}
\proseitem{42}{สํกิลิฏฺฐํ ธมฺมํ ปฏิจฺจ สํกิลิฏฺโฐ ธมฺโม อุปฺปชฺชติ นเหตุปจฺจยา.\csromanspacer{} วิจิกิจฺฉาสหคเต อุทฺธจฺจสหคเต ขนฺเธ ปฏิจฺจ วิจิกิจฺฉาสหคโต อุทฺธจฺจสหคโต โมโห. (1)}

\prose{อสํกิลิฏฺฐํ ธมฺมํ ปฏิจฺจ อสํกิลิฏฺโฐ ธมฺโม อุปฺปชฺชติ นเหตุปจฺจยา.\csromanspacer{} อเหตุกํ อสํกิลิฏฺฐํ เอกํ ขนฺธํ ปฏิจฺจ ตโย ขนฺธา จิตฺตสมุฏฺฐานญฺจ รูปํ ฯเปฯ เทฺว ขนฺเธ ฯเปฯ.\csromanspacer{} อเหตุกปฏิสนฺธิกฺขเณ ฯเปฯ. (ยาว อสญฺญสตฺตา.) (1)}

\csromanheader{2}{2. ปจฺจยปจฺจนีย 2. สงฺขฺยาวาร}
\vspace{\tipitakaparskip}
\csromancenter{\textbf{นเหตุ}ยา เทฺว, น-อารมฺมเณ ตีณิ, น-อธิปติยา ปญฺจ, น-อนนฺตเร ตีณิ, นสมนนฺตเร ตีณิ, น-อญฺญมญฺเญ ตีณิ, น-อุปนิสฺสเย ตีณิ, นปุเรชาเต จตฺตาริ, นปจฺฉาชาเต ปญฺจ, น-อาเสวเน ปญฺจ, นกมฺเม เทฺว, นวิปาเก ปญฺจ, น-อาหาเร เอกํ, น-อินฺทฺริเย เอกํ, นฌาเน เอกํ, นมคฺเค เอกํ, นสมฺปยุตฺเต ตีณิ, นวิปฺปยุตฺเต เทฺว, โนนตฺถิยา ตีณิ, โนวิคเต ตีณิ.}
\csromancenter{(เอวํ อิตเร เทฺว คณนาปิ สหชาตวาโรปิ กาตพฺโพ.)}
\csromanheader{2}{77. สํกิลิฏฺฐทุก \textbf{77. สํกิลิฏฺฐทุก 3. ปจฺจยวาร}}
\proseitem{44}{\csromanfolio{19}สํกิลิฏฺฐิ ธมฺมํ ปจฺจยา สํกิลิฏฺโฐ ธมฺโม อุปฺปชฺชติ เหตุปจฺจยา.\csromanspacer{} ตีณิ. (ปฏิจฺจสทิสํ.)}

\prose{อสํกิลิฏฺฐํ ธมฺมํ ปจฺจยา อสํกิลิฏฺโฐ ธมฺโม อุปฺปชฺชติ เหตุปจฺจยา.\csromanspacer{} อสํกิลิฏฺฐํ เอกํ ขนฺธํ ปจฺจยา ตโย ขนฺธา จิตฺตสมุฏฺฐานญฺจ รูปํ ฯเปฯ เทฺว ขนฺเธ ฯเปฯ.\csromanspacer{} ปฏิสนฺธิกฺขเณ ฯเปฯ. (ยาว อชฺฌตฺติกา มหาภูตา.) วตฺถุํ ปจฺจยา อสํกิลิฏฺฐา ขนฺธา. (1)}

\prose{อสํกิลิฏฺฐํ ธมฺมํ ปจฺจยา สํกิลิฏฺโฐ ธมฺโม อุปฺปชฺชติ เหตุปจฺจยา.\csromanspacer{} วตฺถุํ ปจฺจยา สํกิลิฏฺฐา ขนฺธา. (2)}

\prose{อสํกิลิฏฺฐํ ธมฺมํ ปจฺจยา สํกิลิฏฺโฐ จ อสํกิลิฏฺโฐ จ ธมฺมา อุปฺปชฺชนฺติ เหตุปจฺจยา.\csromanspacer{} วตฺถุํ ปจฺจยา สํกิลิฏฺฐา ขนฺธา.\csromanspacer{} มหาภูเต ปจฺจยา จิตฺตสมุฏฺฐานํ รูปํ. (3)}

\proseitem{45}{สํกิลิฏฺฐญฺจ อสํกิลิฏฺฐญฺจ ธมฺมํ ปจฺจยา สํกิลิฏฺโฐ ธมฺโม อุปฺปชฺชติ เหตุปจฺจยา.\csromanspacer{} สํกิลิฏฺฐํ เอกํ ขนฺธญฺจ วตฺถุญฺจ ปจฺจยา ตโย ขนฺธา ฯเปฯ เทฺว ขนฺเธ จ ฯเปฯ. (1)}

\prose{สํกิลิฏฺฐญฺจ อสํกิลิฏฺฐญฺจ ธมฺมํ ปจฺจยา อสํกิลิฏฺโฐ ธมฺโม อุปฺปชฺชติ เหตุปจฺจยา.\csromanspacer{} สํกิลิฏฺเฐ ขนฺเธ จ มหาภูเต จ ปจฺจยา จิตฺตสมุฏฺฐานํ รูปํ. (2)}

\prose{สํกิลิฏฺฐญฺจ อสํกิลิฏฺฐญฺจ ธมฺมํ ปจฺจยา สํกิลิฏฺโฐ จ อสํกิลิฏฺโฐ จ ธมฺมา อุปฺปชฺชนฺติ เหตุปจฺจยา.\csromanspacer{} สํกิลิฏฺฐํ เอกํ ขนฺธญฺจ วตฺถุญฺจ ปจฺจยา ตโย ขนฺธา ฯเปฯ เทฺว ขนฺเธ จ ฯเปฯ.\csromanspacer{} สํกิลิฏฺเฐ ขนฺเธ จ มหาภูเต จ ปจฺจยา จิตฺตสมุฏฺฐานํ รูปํ. (สํขิตฺตํ.) (3)}

\csromanheader{2}{1. ปจฺจยานุโลม 2. สงฺขฺยาวาร}
\vspace{\tipitakaparskip}
\csromancenter{\textbf{เหตุ}ยา นว, อารมฺมเณ จตฺตาริ, อธิปติยา นว, อนนฺตเร จตฺตาริ, สมนนฺตเร จตฺตาริ, สหชาเต นว, อญฺญมญฺเญ จตฺตาริ, นิสฺสเย นว, อุปนิสฺสเย จตฺตาริ, ปุเรชาเต จตฺตาริ, อาเสวเน}
\prosecont{\csromanfolio{20}จตฺตาริ, กมฺเม นว, วิปาเก เอกํ, อาหาเร นว, อินฺทฺริเย นว ฯเปฯ วิคเต จตฺตาริ, อวิคเต นว.}

\csromanheader{2}{2. ปจฺจยปจฺจนีย 1. วิภงฺควาร}
\csromanheader{2}{\textbf{นเหตุ}}
\proseitem{47}{สํกิลิฏฺฐํ ธมฺมํ ปจฺจยา สํกิลิฏฺโฐ ธมฺโม อุปฺปชฺชติ นเหตุปจฺจยา.\csromanspacer{} วิจิกิจฺฉาสหคเต อุทฺธจฺจสหคเต ขนฺเธ ปจฺจยา วิจิกิจฺฉาสหคโต อุทฺธจฺจสหคโต โมโห. (1)}

\prose{อสํกิลิฏฺฐํ ธมฺมํ ปจฺจยา อสํกิลิฏฺโฐ ธมฺโม อุปฺปชฺชติ นเหตุปจฺจยา.\csromanspacer{} อเหตุกํ อสํกิลิฏฺฐํ ฯเปฯ. (ยาว อสญฺญสตฺตา.) จกฺขายตนํ ปจฺจยา จกฺขุวิญฺญาณํ ฯเปฯ.\csromanspacer{} กายายตนํ ปจฺจยา กายวิญฺญาณํ.\csromanspacer{} วตฺถุํ ปจฺจยา อเหตุกา อสํกิลิฏฺฐา ขนฺธา. (1)}

\prose{อสํกิลิฏฺฐํ ธมฺมํ ปจฺจยา สํกิลิฏฺโฐ ธมฺโม อุปฺปชฺชติ นเหตุปจฺจยา.\csromanspacer{} วตฺถุํ ปจฺจยา วิจิกิจฺฉาสหคโต อุทฺธจฺจสหคโต โมโห. (2)}

\prose{สํกิลิฏฺฐญฺจ อสํกิลิฏฺฐญฺจ ธมฺมํ ปจฺจยา สํกิลิฏฺโฐ ธมฺโม อุปฺปชฺชติ นเหตุปจฺจยา.\csromanspacer{} วิจิกิจฺฉาสหคเต อุทฺธจฺจสหคเต ขนฺเธ จ วตฺถุญฺจ ปจฺจยา วิจิกิจฺฉาสหคโต อุทฺธจฺจสคโต โมโห. (1)}

\csromanheader{2}{2. ปจฺจยปจฺจนีย 2. สงฺขฺยาวาร}
\vspace{\tipitakaparskip}
\csromancenter{\textbf{นเหตุ}ยา จตฺตาริ, น-อารมฺมเณ ตีณิ, น-อธิปติยา นว, น-อนนฺตเร ตีณิ, นสมนนฺตเร ตีณิ, น-อญฺญมญฺเญ ตีณิ, น-อุปนิสฺสเย ตีณิ, นปุเรชาเต จตฺตาริ, นปจฺฉาชาเต นว, น-อาเสวเน นว, นกมฺเม จตฺตาริ, นวิปาเก นว, น-อาหาเร เอกํ, น-อินฺทฺริเย เอกํ, นฌาเน เอกํ, นมคฺเค เอกํ, นสมฺปยุตฺเต ตีณิ, นวิปฺปยุตฺเต เทฺว, โนนตฺถิยา ตีณิ, โนวิคเต ตีณิ.}
\csromancenter{(เอวํ อิตเร เทฺว คณนาปิ นิสฺสยวาโรปิ กาตพฺโพ.)}
\csromanheader{2}{77. สํกิลิฏฺฐทุก \textbf{77. สํกิลิฏฺฐทุก 5. สํสฏฺฐวาร}}
\csromanheader{2}{1. ปจฺจยานุโลมาทิ-เหตุ}
\proseitem{49}{\csromanfolio{21}สํกิลิฏฺฐํ ธมฺมํ สํสฏฺโฐ สํกิลิฏฺโฐ ธมฺโม อุปฺปชฺชติ เหตุปจฺจยา.\csromanspacer{} สํกิลิฏฺฐํ เอกํ ขนฺธํ สํสฏฺฐา ตโย ขนฺธา ฯเปฯ เทฺว ขนฺเธ ฯเปฯ. (1)}

\prose{อสํกิลิฏฺฐํ ธมฺมํ สํสฏฺโฐ อสํกิลิฏฺโฐ ธมฺโม อุปฺปชฺชติ เหตุปจฺจยา.\csromanspacer{} อสํกิลิฏฺฐํ เอกํ ขนฺธํ สํสฏฺฐา ตโย ขนฺธา ฯเปฯ เทฺว ขนฺเธ ฯเปฯ.\csromanspacer{} ปฏิสนฺธิกฺขเณ ฯเปฯ. (1)}

\prose{\textbf{เหตุ}ยา เทฺว, อารมฺมเณ เทฺว, อธิปติยา เทฺว, (สพฺพตฺถ เทฺว.) วิปาเก เอกํ ฯเปฯ อวิคเต เทฺว.\csromanspacer{} อนุโลมํ.}

\proseitem{50}{สํกิลิฏฺฐํ ธมฺมํ สํสฏฺโฐ สํกิลิฏฺโฐ ธมฺโม อุปฺปชฺชติ นเหตุปจฺจยา.\csromanspacer{} วิจิกิจฺฉาสหคเต อุทฺธจฺจสหคเต ขนฺเธ สํสฏฺโฐ วิจิกิจฺฉาสหคโต อุทฺธจฺจสหคโต โมโห. (1)}

\prose{อสํกิลิฏฺฐํ ธมฺมํ สํสฏฺโฐ อสํกิลิฏฺโฐ ธมฺโม อุปฺปชฺชติ นเหตุปจฺจยา.\csromanspacer{} อเหตุกํ อสํกิลิฏฺฐํ เอกํ ขนฺธํ สํสฏฺฐา ตโย ขนฺธา ฯเปฯ เทฺว ขนฺเธ ฯเปฯ.\csromanspacer{} อเหตุกปฏิสนฺธิกฺขเณ ฯเปฯ. (1)}

\prose{นเหตุยา เทฺว, น-อธิปติยา เทฺว, นปุเรชาเต เทฺว, นปจฺฉาชาเต เทฺว, น-อาเสวเน เทฺว, นกมฺเม เทฺว, นวิปาเก เทฺว, นฌาเน เอกํ, นมคฺเค เอกํ, นวิปฺปยุตฺเต เทฺว.}

\vspace{\tipitakaparskip}
\csromancenter{ปจฺจนียํ.}
\csromancenter{(เอวํ อิตเร เทฺว คณนาปิ สมฺปยุตฺตวาโรปิ กาตพฺโพ.)}
\csromanheader{2}{77. สํกิลิฏฺฐทุก 7. ปญฺหาวาร}
\proseitem{51}{สํกิลิฏฺโฐ ธมฺโม สํกิลิฏฺฐสฺส ธมฺมสฺส เหตุปจฺจเยน ปจฺจโย.\csromanspacer{} สํกิลิฏฺฐา เหตู สมฺปยุตฺตกานํ ขนฺธานํ เหตุปจฺจเยน ปจฺจโย. (1)}

\prose{\csromanfolio{22}สํกิลิฏฺโฐ ธมฺโม อสํกิลิฏฺฐสฺส ธมฺมสฺส เหตุปจฺจเยน ปจฺจโย.\csromanspacer{} สํกิลิฏฺฐา เหตู จิตฺตสมุฏฺฐานานํ รูปานํ เหตุปจฺจเยน ปจฺจโย. (2)}

\prose{สํกิลิฏฺโฐ ธมฺโม สํกิลิฏฺฐสฺส จ อสํกิลิฏฺฐสฺส จ ธมฺมสฺส เหตุปจฺจเยน ปจฺจโย.\csromanspacer{} สํกิลิฏฺฐา เหตู สมฺปยุตฺตกานํ ขนฺธานํ จิตฺตสมุฏฺฐานานญฺจ รูปานํ เหตุปจฺจเยน ปจฺจโย. (3)}

\prose{อสํกิลิฏฺโฐ ธมฺโม อสํกิลิฏฺฐสฺส ธมฺมสฺส เหตุปจฺจเยน ปจฺจโย.\csromanspacer{} อสํกิลิฏฺฐา เหตู สมฺปยุตฺตกานํ ขนฺธานํ จิตฺตสมุฏฺฐานานญฺจ รูปานํ เหตุปจฺจเยน ปจฺจโย.\csromanspacer{} ปฏิสนฺธิกฺขเณ ฯเปฯ. (1)}

\csromanheader{2}{\textbf{อารมฺมณ}}
\proseitem{52}{สํกิลิฏฺโฐ ธมฺโม สํกิลิฏฺฐสฺส ธมฺมสฺส อารมฺมณปจฺจเยน ปจฺจโย.\csromanspacer{} ราคํ อสฺสาเทติ อภินนฺทติ, ตํ อารพฺภ ราโค อุปฺปชฺชติ.\csromanspacer{} ทิฏฺฐิ อุปฺปชฺชติ.\csromanspacer{} วิจิกิจฺฉา ฯเปฯ อุทฺธจฺจํ ฯเปฯ โทมนสฺสํ อุปฺปชฺชติ.\csromanspacer{} ทิฏฺฐิํ อสฺสาเทติ ฯเปฯ. (กุสลตฺติกสทิสํ.) วิจิกิจฺฉํ อารพฺภ ฯเปฯ อุทฺธจฺจํ อารพฺภ ฯเปฯ โทมนสฺสํ อุปฺปชฺชติ.\csromanspacer{} ทิฏฺฐิ ฯเปฯ วิจิกิจฺฉา ฯเปฯ อุทฺธจฺจํ อุปฺปชฺชติ. (1)}

\prose{สํกิลิฏฺโฐ ธมฺโม อสํกิลิฏฺฐสฺส ธมฺมสฺส อารมฺมณปจฺจเยน ปจฺจโย.\csromanspacer{} อริยา ปหีเน กิเลเส ปจฺจเวกฺขนฺติ, วิกฺขมฺภิเต กิเลเส ฯเปฯ.\csromanspacer{} ปุพฺเพ สมุทาจิณฺเณ ฯเปฯ.\csromanspacer{} สํกิลิฏฺเฐ ขนฺเธ อนิจฺจโต ฯเปฯ วิปสฺสติ.\csromanspacer{} เจโตปริยญาเณน สํกิลิฏฺฐจิตฺตสมงฺคิสฺส จิตฺตํ ชานาติ.\csromanspacer{} สํกิลิฏฺฐา ขนฺธา เจโตปริยญาณสฺส, ปุพฺเพนิวาสานุสฺสติญาณสฺส, ยถากมฺมูปคญาณสฺส, อนาคตํสญาณสฺส, อาวชฺชนาย อารมฺมณปจฺจเยน ปจฺจโย. (2)}

\proseitem{53}{อสํกิลิฏฺโฐ ธมฺโม อสํกิลิฏฺฐสฺส ธมฺมสฺส อารมฺมณปจฺจเยน ปจฺจโย.\csromanspacer{} ทานํ ฯเปฯ สีลํ ฯเปฯ อุโปสถกมฺมํ ฯเปฯ.\csromanspacer{} ปุพฺเพ ฯเปฯ.\csromanspacer{} ฌานา วุฏฺฐหิตฺวา ฌานํ ปจฺจเวกฺขติ.\csromanspacer{} อริยา มคฺคา วุฏฺฐหิตฺวา มคฺคํ ปจฺจเวกฺขนฺติ ฯเปฯ อาวชฺชนาย อารมฺมณปจฺจเยน ปจฺจโย.\csromanspacer{} จกฺขุํ ฯเปฯ.\csromanspacer{} วตฺถุํ อสํกิลิฏฺเฐ ขนฺเธ อนิจฺจโต ฯเปฯ วิปสฺสติ.\csromanspacer{} ทิพฺเพน จกฺขุนา รูปํ ปสฺสติ ฯเปฯ.\csromanspacer{} อนาคตํสญาณสฺส, อาวชฺชนาย อารมฺมณปจฺจเยน ปจฺจโย. (1)}

\prose{อสํกิลิฏฺโฐ ธมฺโม สํกิลิฏฺฐสฺส ธมฺมสฺส อารมฺมณปจฺจเยน ปจฺจโย.\csromanspacer{} ทานํ ฯเปฯ สีลํ ฯเปฯ.\csromanspacer{} ฌานา วุฏฺฐหิตฺวา ฯเปฯ.\csromanspacer{} จกฺขุํ ฯเปฯ.\csromanspacer{} วตฺถุํ อสํกิลิฏฺเฐ}

\vspace{\tipitakaparskip}
\csromancenter{สํกิลิฏฺฐทุก ขนฺเธ อสฺสาเทติ อภินนฺทติ, ตํ อารพฺภ ราโค ฯเปฯ โทมนสฺสํ อุปฺปชฺชติ. (2)}
\csromanheader{2}{\textbf{อธิปติ}}
\proseitem{54}{\csromanfolio{23}สํกิลิฏฺโฐ ธมฺโม สํกิลิฏฺฐสฺส ธมฺมสฺส อธิปติปจฺจเยน ปจฺจโย.\csromanspacer{} อารมฺมณาธิปติ, สหชาตาธิปติ.\csromanspacer{}\csromanspacer{}\textbf{อารมฺมณาธิปติ\csromandash{}}ราคํ ครุํ กตฺวา อสฺสาเทติ อภินนฺทติ, ตํ ครุํ กตฺวา ราโค อุปฺปชฺชติ.\csromanspacer{} ทิฏฺฐิ อุปฺปชฺชติ.\csromanspacer{} ทิฏฺฐิํ ครุํ กตฺวา อสฺสาเทติ อภินนฺทติ, ตํ ครุํ กตฺวา ราโค อุปฺปชฺชติ.\csromanspacer{} ทิฏฺฐิ อุปฺปชฺชติ.\csromanspacer{} \textbf{สหชาตาธิปติ\csromandash{}}สํกิลิฏฺฐาธิปติ สมฺปยุตฺตกานํ ขนฺธานํ อธิปติปจฺจเยน ปจฺจโย. (1)}

\prose{สํกิลิฏฺโฐ ธมฺโม อสํกิลิฏฺฐสฺส ธมฺมสฺส อธิปติปจฺจเยน ปจฺจโย.\csromanspacer{} \textbf{สหชาตาธิปติ\csromandash{}}สํกิลิฏฺฐาธิปติ จิตฺตสมุฏฺฐานานํ รูปานํ อธิปติปจฺจเยน ปจฺจโย. (2)}

\prose{สํกิลิฏฺโฐ ธมฺโม สํกิลิฏฺฐสฺส จ อสํกิลิฏฺฐสฺส จ ธมฺมสฺส อธิปติปจฺจเยน ปจฺจโย.\csromanspacer{} \textbf{สหชาตาธิปติ\csromandash{}}สํกิลิฏฺฐาธิปติ สมฺปยุตฺตกานํ ขนฺธานํ จิตฺตสมุฏฺฐานานญฺจ รูปานํ อธิปติปจฺจเยน ปจฺจโย. (3)}

\proseitem{55}{อสํกิลิฏฺโฐ ธมฺโม อสํกิลิฏฺฐสฺส ธมฺมสฺส อธิปติปจฺจเยน ปจฺจโย.\csromanspacer{} อารมฺมณาธิปติ, สหชาตาธิปติ.\csromanspacer{}\csromanspacer{}\textbf{อารมฺมณาธิปติ\csromandash{}}ทานํ ฯเปฯ สีลํ ฯเปฯ อุโปสถกมฺมํ ฯเปฯ.\csromanspacer{} ปุพฺเพ ฯเปฯ.\csromanspacer{} ฌานา วุฏฺฐหิตฺวา ฌานํ ครุํ กตฺวา ปจฺจเวกฺขติ.\csromanspacer{} อริยา มคฺคา วุฏฺฐหิตฺวา มคฺคํ ครุํ กตฺวา ปจฺจเวกฺขนฺติ ฯเปฯ.\csromanspacer{} นิพฺพานํ ผลสฺส อธิปติปจฺจเยน ปจฺจโย.\csromanspacer{} \textbf{สหชาตาธิปติ\csromandash{}}อสํกิลิฏฺฐาธิปติ สมฺปยุตฺตกานํ ขนฺธานํ จิตฺตสมุฏฺฐานานญฺจ รูปานํ อธิปติปจฺจเยน ปจฺจโย. (1)}

\prose{อสํกิลิฏฺโฐ ธมฺโม สํกิลิฏฺฐสฺส ธมฺมสฺส อธิปติปจฺจเยน ปจฺจโย.\csromanspacer{} \textbf{อารมฺมณาธิปติ\csromandash{}}ทานํ ฯเปฯ สีลํ ฯเปฯ อุโปสถกมฺมํ ฯเปฯ.\csromanspacer{} ปุพฺเพ ฯเปฯ.\csromanspacer{} ฌานา วุฏฺฐหิตฺวา ฯเปฯ.\csromanspacer{} จกฺขุํ ฯเปฯ.\csromanspacer{} วตฺถุํ อสํกิลิฏฺเฐ ขนฺเธ ครุํ กตฺวา อสฺสาเทติ อภินนฺทติ, ตํ ครุํ กตฺวา ราโค อุปฺปชฺชติ.\csromanspacer{} ทิฏฺฐิ อุปฺปชฺชติ. (2)}

\csromanheader{2}{\textbf{อนนฺตราทิ}}
\proseitem{56}{\csromanfolio{24}สํกิลิฏฺโฐ ธมฺโม สํกิลิฏฺฐสฺส ธมฺมสฺส อนนฺตรปจฺจเยน ปจฺจโย.\csromanspacer{} ปุริมา ปุริมา สํกิลิฏฺฐา ขนฺธา ปจฺฉิมานํ ปจฺฉิมานํ สํกิลิฏฺฐานํ ขนฺธานํ อนนฺตรปจฺจเยน ปจฺจโย. (1)}

\prose{สํกิลิฏฺโฐ ธมฺโม อสํกิลิฏฺฐสฺส ธมฺมสฺส อนนฺตรปจฺจเยน ปจฺจโย.\csromanspacer{} สํกิลิฏฺฐา ขนฺธา วุฏฺฐานสฺส อนนฺตรปจฺจเยน ปจฺจโย. (2)}

\proseitem{57}{อสํกิลิฏฺโฐ ธมฺโม อสํกิลิฏฺฐสฺส ธมฺมสฺส อนนฺตรปจฺจเยน ปจฺจโย.\csromanspacer{} ปุริมา ปุริมา อสํกิลิฏฺฐา ขนฺธา ปจฺฉิมานํ ปจฺฉิมานํ อสํกิลิฏฺฐานํ ขนฺธานํ ฯเปฯ ผลสมาปตฺติยา อนนฺตรปจฺจเยน ปจฺจโย. (1)}

\prose{อสํกิลิฏฺโฐ ธมฺโม สํกิลิฏฺฐสฺส ธมฺมสฺส อนนฺตรปจฺจเยน ปจฺจโย.\csromanspacer{} อาวชฺชนา สํกิลิฏฺฐานํ ขนฺธานํ อนนฺตรปจฺจเยน ปจฺจโย. (2)}

\prose{สมนนฺตรปจฺจเยน ปจฺจโย.\csromanspacer{} จตฺตาริ.\csromanspacer{} สหชาตปจฺจเยน ปจฺจโย.\csromanspacer{} ปญฺจ.\csromanspacer{} อญฺญมญฺญปจฺจเยน ปจฺจโย.\csromanspacer{} เทฺว.\csromanspacer{} นิสฺสยปจฺจเยน ปจฺจโย.\csromanspacer{} สตฺต.}

\csromanheader{2}{\textbf{อุปนิสฺสย}}
\proseitem{58}{สํกิลิฏฺโฐ ธมฺโม สํกิลิฏฺฐสฺส ธมฺมสฺส อุปนิสฺสยปจฺจเยน ปจฺจโย.\csromanspacer{} \textbf{อารมฺมณูปนิสฺสโย, อนนฺตรูปนิสฺสโย, ปกตูปนิสฺสโย ฯเปฯ.}\csromanspacer{} ปกตูปนิสฺสโย\csromandash{}ราคํ อุปนิสฺสาย ปาณํ หนติ ฯเปฯ สํฆํ ภินฺทติ.\csromanspacer{} โทสํ ฯเปฯ.\csromanspacer{} ปตฺถนํ อุปนิสฺสาย ปาณํ หนติ ฯเปฯ สํฆํ ภินฺทติ.\csromanspacer{} ราโค ฯเปฯ.\csromanspacer{} ปตฺถนา ราคสฺส ฯเปฯ ปตฺถนาย อุปนิสฺสยปจฺจเยน ปจฺจโย. (1)}

\prose{สํกิลิฏฺโฐ ธมฺโม อสํกิลิฏฺฐสฺส ธมฺมสฺส อุปนิสฺสยปจฺจเยน ปจฺจโย.\csromanspacer{} \textbf{อนนฺตรูปนิสฺสโย, ปกตูปนิสฺสโย ฯเปฯ.}\csromanspacer{} ปกตูปนิสฺสโย\csromandash{} ราคํ อุปนิสฺสาย ทานํ เทติ ฯเปฯ สมาปตฺติํ อุปฺปาเทติ.\csromanspacer{} โทสํ ฯเปฯ.\csromanspacer{} ปตฺถนํ อุปนิสฺสาย ทานํ เทติ ฯเปฯ สมาปตฺติํ อุปฺปาเทติ.\csromanspacer{} ราโค ฯเปฯ.\csromanspacer{} ปตฺถนา สทฺธาย ฯเปฯ ปญฺญาย.\csromanspacer{} กายิกสฺส สุขสฺส.\csromanspacer{} กายิกสฺส ทุกฺขสฺส.\csromanspacer{} มคฺคสฺส, ผลสมาปตฺติยา อุปนิสฺสยปจฺจเยน ปจฺจโย. (2)}

\proseitem{59}{อสํกิลิฏฺโฐ ธมฺโม อสํกิลิฏฺฐสฺส ธมฺมสฺส อุปนิสฺสยปจฺจเยน ปจฺจโย.\csromanspacer{} \textbf{อารมฺมณูปนิสฺสโย, อนนฺตรูปนิสฺสโย, ปกตูปนิสฺสโย ฯเปฯ.}\csromanspacer{} ปกตูปนิสฺสโย\csromandash{}สทฺธํ อุปนิสฺสาย ทานํ เทติ ฯเปฯ}

\vspace{\tipitakaparskip}
\csromancenter{สํกิลิฏฺฐทุก สมาปตฺติํ อุปฺปาเทติ.\csromanspacer{} สีลํ ฯเปฯ.\csromanspacer{} ปญฺญํ.\csromanspacer{} กายิกํ สุขํ.\csromanspacer{} กายิกํ ทุกฺขํ.\csromanspacer{} อุตุํ.\csromanspacer{} โภชนํ.\csromanspacer{} เสนาสนํ อุปนิสฺสาย ทานํ เทติ ฯเปฯ สมาปตฺติํ อุปฺปาเทติ.\csromanspacer{} สทฺธา ฯเปฯ.\csromanspacer{} เสนาสนํ สทฺธาย ฯเปฯ ผลสมาปตฺติยา อุปานิสฺสยปจฺจเยน ปจฺจโย. (1)}
\prose{\csromanfolio{25}อสํกิลิฏฺโฐ ธมฺโม สํกิลิฏฺฐสฺส ธมฺมสฺส อุปนิสฺสยปจฺจเยน ปจฺจโย.\csromanspacer{} \textbf{อารมฺมณูปนิสฺสโย, อนนฺตรูปนิสฺสโย, ปกตูปนิสฺสโย ฯเปฯ.}\csromanspacer{} ปกตูปนิสฺสโย\csromandash{}สทฺธํ อุปนิสฺสาย มานํ ชปฺเปติ.\csromanspacer{} ทิฏฺฐิํ คณฺหาติ.\csromanspacer{} สีลํ ฯเปฯ.\csromanspacer{} ปญฺญํ.\csromanspacer{} กายิกํ สุขํ.\csromanspacer{} กายิกํ ทุกฺขํ.\csromanspacer{} อุตุํ.\csromanspacer{} โภชนํ.\csromanspacer{} เสนาสนํ อุปนิสฺสาย ปาณํ หนติ ฯเปฯ สํฆํ ภินฺทติ.\csromanspacer{} สทฺธา ฯเปฯ.\csromanspacer{} เสนาสนํ ราคสฺส ฯเปฯ ปตฺถนาย อุปนิสฺสยปจฺจเยน ปจฺจโย. (2)}

\csromanheader{2}{\textbf{ปุเรชาต}}
\proseitem{60}{อสํกิลิฏฺโฐ ธมฺโม อสํกิลิฏฺฐสฺส ธมฺมสฺส ปุเรชาตปจฺจเยน ปจฺจโย.\csromanspacer{} อารมฺมณปุเรชาตํ, วตฺถุปุเรชาตํ. (สํขิตฺตํ.) อสํกิลิฏฺโฐ ธมฺโม สํกิลิฏฺฐสฺส ธมฺมสฺส ปุเรชาตปจฺจเยน ปจฺจโย.\csromanspacer{} อารมฺมณปุเรชาตํ, วตฺถุปุเรชาตํ. (สํขิตฺตํ.) (2)}

\csromanheader{2}{\textbf{ปจฺฉาชาตาเสวน}}
\proseitem{61}{สํกิลิฏฺโฐ ธมฺโม อสํกิลิฏฺฐสฺส ธมฺมสฺส ปจฺฉาชาตปจฺจเยน ปจฺจโย. (สํขิตฺตํ.) อสํกิลิฏฺโฐ ธมฺโม อสํกิลิฏฺฐสฺส ธมฺมสฺส ปจฺฉาชาตปจฺจเยน ปจฺจโย. (สํขิตฺตํ.) อาเสวนปจฺจเยน ปจฺจโย เทฺว.}

\csromanheader{2}{\textbf{กมฺม}}
\proseitem{62}{สํกิลิฏฺโฐ ธมฺโม สํกิลิฏฺฐสฺส ธมฺมสฺส กมฺมปจฺจเยน ปจฺจโย.\csromanspacer{} สํกิลิฏฺฐา เจตนา สมฺปยุตฺตกานํ ขนฺธานํ กมฺมปจฺจเยน ปจฺจโย. (1)}

\prose{สํกิลิฏฺโฐ ธมฺโม อสํกิลิฏฺฐสฺส ธมฺมสฺส กมฺมปจฺจเยน ปจฺจโย.\csromanspacer{} \textbf{สหชาตา}, นานากฺขณิกา.\csromanspacer{}\csromanspacer{}\textbf{สหชาตา} สํกิลิฏฺฐา เจตนา จิตฺตสมุฏฺฐานานํ รูปานํ กมฺมปจฺจเยน ปจฺจโย.\csromanspacer{} \textbf{นานากฺขณิกา} สํกิลิฏฺฐา เจตนา วิปากานํ ขนฺธานํ กฏตฺตา จ รูปานํ กมฺมปจฺจเยน ปจฺจโย. (มูลํ กาตพฺพํ.) สํกิลิฏฺฐา เจตนา สมฺปยุตฺตกานํ ขนฺธานํ จิตฺตสมุฏฺฐานานญฺจ รูปานํ กมฺมปจฺจเยน ปจฺจโย. (3)}

\prose{\csromanfolio{26}อสํกิลิฏฺโฐ ธมฺโม อสํกิลิฏฺฐสฺส ธมฺมสฺส กมฺมปจฺจเยน ปจฺจโย.\csromanspacer{} \textbf{สหชาตา}, นานากฺขณิกา.\csromanspacer{}\csromanspacer{}\textbf{สหชาตา} อสํกิลิฏฺฐา เจตนา สมฺปยุตฺตกานํ ขนฺธานํ จิตฺตสมุฏฺฐานานญฺจ รูปานํ กมฺมปจฺจเยน ปจฺจโย.\csromanspacer{} ปฏิสนฺธิกฺขเณ ฯเปฯ.\csromanspacer{} \textbf{นานากฺขณิกา} อสํกิลิฏฺฐา เจตนา วิปากานํ ขนฺธานํ กฏตฺตา จ รูปานํ กมฺมปจฺจเยน ปจฺจโย. (1)}

\csromanheader{2}{\textbf{วิปากาทิ}}
\proseitem{63}{อสํกิลิฏฺโฐ ธมฺโม อสํกิลิฏฺฐสฺส ธมฺมสฺส วิปากปจฺจเยน ปจฺจโย.\csromanspacer{} เอกํ.}

\prose{สํกิลิฏฺโฐ ธมฺโม สํกิลิฏฺฐสฺส ธมฺมสฺส อาหารปจฺจเยน ปจฺจโย.\csromanspacer{} อินฺทฺริยปจฺจเยน ปจฺจโย.\csromanspacer{} ฌานปจฺจเยน ปจฺจโย.\csromanspacer{} มคฺคปจฺจเยน ปจฺจโย.\csromanspacer{} สมฺปยุตฺตปจฺจเยน ปจฺจโย.}

\csromanheader{2}{\textbf{วิปฺปยุตฺต}}
\proseitem{64}{สํกิลิฏฺโฐ ธมฺโม อสํกิลิฏฺฐสฺส ธมฺมสฺส วิปฺปยุตฺตปจฺจเยน ปจฺจโย.\csromanspacer{} สหชาตํ, ปจฺฉาชาตํ. (สํขิตฺตํ.) (1)}

\prose{อสํกิลิฏฺโฐ ธมฺโม อสํกิลิฏฺฐสฺส ธมฺมสฺส วิปฺปยุตฺตปจฺจเยน ปจฺจโย.\csromanspacer{} สหชาตํ, ปุเรชาตํ, ปจฺฉาชาตํ. (สํขิตฺตํ.) (1)}

\prose{อสํกิลิฏฺโฐ ธมฺโม สํกิลิฏฺฐสฺส ธมฺมสฺส วิปฺปยุตฺตปจฺจเยน ปจฺจโย.\csromanspacer{} \textbf{ปุเรชาตํ} วตฺถุ สํกิลิฏฺฐานํ ขนฺธานํ วิปฺปยุตฺตปจฺจเยน ปจฺจโย. (2)}

\csromanheader{2}{\textbf{อตฺถิ-อาทิ}}
\proseitem{65}{สํกิลิฏฺโฐ ธมฺโม สํกิลิฏฺฐสฺส ธมฺมสฺส อตฺถิปจฺจเยน ปจฺจโย.\csromanspacer{} เอกํ. (ปฏิจฺจวารสทิสํ.) สํกิลิฏฺโฐ ธมฺโม อสํกิลิฏฺฐสฺส ธมฺมสฺส อตฺถิปจฺจเยน ปจฺจโย.\csromanspacer{} สหชาตํ, ปจฺฉาชาตํ. (สํขิตฺตํ.) สํกิลิฏฺโฐ ธมฺโม สํกิลิฏฺฐสฺส จ อสํกิลิฏฺฐสฺส จ ธมฺมสฺส อตฺถิปจฺจเยน ปจฺจโย. (ปฏิจฺจสทิสํ.) (3)}

\prose{อสํกิลิฏฺโฐ ธมฺโม อสํกิลิฏฺฐสฺส ธมฺมสฺส อตฺถิปจฺจเยน ปจฺจโย.\csromanspacer{} สหชาตํ, ปุเรชาตํ, ปจฺฉาชาตํ, อาหารํ, อินฺทฺริยํ. (สํขิตฺตํ.) อสํกิลิฏฺโฐ ธมฺโม สํกิลิฏฺฐสฺส ธมฺมสฺส อตฺถิปจฺจเยน ปจฺจโย.\csromanspacer{} \textbf{ปุเรชาตํ}. (สํขิตฺตํ.) (2)}

\csromanheader{2}{77. สํกิลิฏฺฐทุก}
\proseitem{66}{\csromanfolio{27}สํกิลิฏฺโฐ จ อสํกิลิฏฺโฐ จ ธมฺมา สํกิลิฏฺฐสฺส ธมฺมสฺส อตฺถิปจฺจเยน ปจฺจโย.\csromanspacer{} สหชาตํ, ปุเรชาตํ.\csromanspacer{}\csromanspacer{}สหชาโต สํกิลิฏฺโฐ เอโก ขนฺโธ จ วตฺถุ จ ติณฺณนฺนํ ขนฺธานํ อตฺถิปจฺจเยน ปจฺจโย ฯเปฯ เทฺว ขนฺธา จ ฯเปฯ. (สํขิตฺตํ.) (1)}

\prose{สํกิลิฏฺโฐ จ อสํกิลิฏฺโฐ จ ธมฺมา อสํกิลิฏฺฐสฺส ธมฺมสฺส อตฺถิปจฺจเยน ปจฺจโย.\csromanspacer{} สหชาตํ, ปจฺฉาชาตํ, อาหารํ, อินฺทฺริยํ.\csromanspacer{}\csromanspacer{}\textbf{สหชาตา} สํกิลิฏฺฐา ขนฺธา จ มหาภูตา จ จิตฺตสมุฏฺฐานานํ รูปานํ อตฺถิปจฺจเยน ปจฺจโย.\csromanspacer{} \textbf{ปจฺฉาชาตา} สํกิลิฏฺฐา ขนฺธา จ กพฬีกาโร อาหาโร จ อิมสฺส กายสฺส อตฺถิปจฺจเยน ปจฺจโย.\csromanspacer{} \textbf{ปจฺฉาชาตา} สํกิลิฏฺฐา ขนฺธา จ รูปชีวิตินฺทฺริยญฺจ กฏตฺตารูปานํ อตฺถิปจฺจเยน ปจฺจโย. (2)}

\prose{นตฺถิปจฺจเยน ปจฺจโย.\csromanspacer{} วิคตปจฺจเยน ปจฺจโย.\csromanspacer{} อวิคตปจฺจเยน ปจฺจโย.}

\csromanheader{2}{1. ปจฺจยานุโลม 2. สงฺขฺยาวาร}
\csromanheader{2}{\textbf{สุทฺธ}}
\proseitem{67}{เหตุยา จตฺตาริ, อารมฺมเณ จตฺตาริ, อธิปติยา ปญฺจ, อนนฺตเร จตฺตาริ, สมนนฺตเร จตฺตาริ, สหชาเต ปญฺจ, อญฺญมญฺเญ เทฺว, นิสฺสเย สตฺต, อุปนิสฺสเย จตฺตาริ, ปุเรชาเต เทฺว, ปจฺฉาชาเต เทฺว, อาเสวเน เทฺว, กมฺเม จตฺตาริ, วิปาเก เอกํ, อาหาเร จตฺตาริ, อินฺทฺริเย จตฺตาริ, ฌาเน จตฺตาริ, มคฺเค จตฺตาริ, สมฺปยุตฺเต เทฺว, วิปฺปยุตฺเต ตีณิ, อตฺถิยา สตฺต, นตฺถิยา จตฺตาริ, วิคเต จตฺตาริ, อวิคเต สตฺต.}

\csromanheader{2}{2. ปจฺจนียุทฺธาร}
\proseitem{68}{สํกิลิฏฺโฐ ธมฺโม สํกิลิฏฺฐสฺส ธมฺมสฺส อารมฺมณปจฺจเยน ปจฺจโย.\csromanspacer{} สหชาตปจฺจเยน ปจฺจโย.\csromanspacer{} อุปนิสฺสยปจฺจเยน ปจฺจโย. (1)}

\prose{สํกิลิฏฺโฐ ธมฺโม อสํกิลิฏฺฐสฺส ธมฺมสฺส อารมฺมณปจฺจเยน ปจฺจโย.\csromanspacer{} สหชาตปจฺจเยน ปจฺจโย.\csromanspacer{} อุปนิสฺสยปจฺจเยน ปจฺจโย.\csromanspacer{} ปจฺฉาชาตปจฺจเยน ปจฺจโย.\csromanspacer{} กมฺมปจฺจเยน ปจฺจโย. (2)}

\prose{\csromanfolio{28}สํกิลิฏฺโฐ ธมฺโม สํกิลิฏฺฐสฺส จ อสํกิลิฏฺฐสฺส จ ธมฺมสฺส สหชาตปจฺจเยน ปจฺจโย. (3)}

\proseitem{69}{อสํกิลิฏฺโฐ ธมฺโม อสํกิลิฏฺฐสฺส ธมฺมสฺส อารมฺมณปจฺจเยน ปจฺจโย.\csromanspacer{} สหชาตปจฺจเยน ปจฺจโย.\csromanspacer{} อุปนิสฺสยปจฺจเยน ปจฺจโย.\csromanspacer{} ปุเรชาตปจฺจเยน ปจฺจโย.\csromanspacer{} ปจฺฉาชาตปจฺจเยน ปจฺจโย.\csromanspacer{} กมฺมปจฺจเยน ปจฺจโย.\csromanspacer{} อาหารปจฺจเยน ปจฺจโย.\csromanspacer{} อินฺทฺริยปจฺจเยน ปจฺจโย. (1)}

\prose{อสํกิลิฏฺโฐ ธมฺโม สํกิลิฏฺฐสฺส ธมฺมสฺส อารมฺมณปจฺจเยน ปจฺจโย.\csromanspacer{} อุปนิสฺสยปจฺจเยน ปจฺจโย.\csromanspacer{} ปุเรชาตปจฺจเยน ปจฺจโย. (2)}

\prose{สํกิลิฏฺโฐ จ อสํกิลิฏฺโฐ จ ธมฺมา สํกิลิฏฺฐสฺส ธมฺมสฺส สหชาตํ.\csromanspacer{} ปุเรชาตํ. (1)}

\prose{สํกิลิฏฺโฐ จ อสํกิลิฏฺโฐ จ ธมฺมา อสํกิลิฏฺฐสฺส ธมฺมสฺส สหชาตํ.\csromanspacer{} ปจฺฉาชาตํ.\csromanspacer{} อาหารํ.\csromanspacer{} อินฺทฺริยํ. (2)}

\csromanheader{2}{2. ปจฺจยปจฺจนีย 2. สงฺขฺยาวาร}
\csromanheader{2}{\textbf{สุทฺธ}}
\proseitem{70}{นเหตุยา สตฺต, น-อารมฺมเณ สตฺต, น-อธิปติยา สตฺต, น-อนนฺตเร สตฺต, นสมนนฺตเร สตฺต, นสหชาเต ปญฺจ, น-อญฺญมญฺเญ ปญฺจ, นนิสฺสเย ปญฺจ, น-อุปนิสฺสเย สตฺต, นปุเรชาเต ฉ, นปจฺฉาชาเต สตฺต ฯเปฯ นมคฺเค สตฺต, นสมฺปยุตฺเต ปญฺจ, นวิปฺปยุตฺเต จตฺตาริ, โน-อตฺถิยา จตฺตาริ, โนนตฺถิยา สตฺต, โนวิคเต สตฺต, โน-อวิคเต จตฺตาริ.}

\csromanheader{2}{3. ปจฺจยานุโลมปจฺจนีย}
\proseitem{71}{เหตุปจฺจยา น-อารมฺมเณ จตฺตาริ, น-อธิปติยา จตฺตาริ, น-อนนฺตเร จตฺตาริ, นสมนนฺตเร จตฺตาริ, น-อญฺญมญฺเญ เทฺว, น-อุปนิสฺสเย จตฺตาริ ฯเปฯ นสมฺปยุตฺเต เทฺว, นวิปฺปยุตฺเต เทฺว, โนนตฺถิยา จตฺตาริ, โนวิคเต จตฺตาริ.}

\csromanmatikaanchor{mtk.6}
\csromanmatikaanchor{mtk.7}
\csromantocmarkref{section}{78. กิเลสสมฺปยุตฺตทุก}{mtk.6}
\bookmark[dest={mtk.6},level=1]{78. กิเลสสมฺปยุตฺตทุก}
\csromantocmarkref{section}{79. กิเลสสํกิเลสิกทุก}{mtk.7}
\bookmark[dest={mtk.7},level=1]{79. กิเลสสํกิเลสิกทุก}
\csromanheader{1}{79. กิเลสสํกิเลสิกทุก \textbf{4. ปจฺจยปจฺจนียานุโลม}}
\proseitem{72}{\csromanfolio{29}นเหตุปจฺจยา อารมฺมเณ จตฺตาริ, อธิปติยา ปญฺจ, (อนุโลมมาติกา.) ฯเปฯ อวิคเต สตฺต.}

\nitthitam{สํกิลิฏฺฐทุกํ นิฏฺฐิตํ.}
\csromanheader{1}{78. กิเลสสมฺปยุตฺตทุก 1. ปฏิจฺจวาร}
\proseitem{73}{กิเลสสมฺปยุตฺตํ ธมฺมํ ปฏิจฺจ กิเลสสมฺปยุตฺโต ธมฺโม อุปฺปชฺชติ เหตุปจฺจยา.\csromanspacer{} กิเลสสมฺปยุตฺตํ เอกํ ขนฺธํ ปฏิจฺจ ตโย ธนฺธา ฯเปฯ เทฺว ขนฺเธ ฯเปฯ. (1)}

\prose{กิเลสสมฺปยุตฺตํ ธมฺมํ ปฏิจฺจ กิเลสวิปฺปยุตฺโต ธมฺโม อุปฺปชฺชติ เหตุปจฺจยา.\csromanspacer{} กิเลสสมฺปยุตฺเต ขนฺเธ ปฏิจฺจ จิตฺตสมุฏฺฐานํ รูปํ. (2)}

\prose{(กิเลสสมฺปยุตฺตทุกํ สํกิลิฏฺฐทุกสทิสํ, นินฺนานากรณํ.)}

\nitthitam{กิเลสสมฺปยุตฺตทุกํ นิฏฺฐิตํ.}
\csromanheader{2}{79. กิเลสสํกิเลสิกทุก 1. ปฏิจฺจวาร}
\proseitem{74}{กิเลสญฺเจว สํกิเลสิกญฺจ ธมฺมํ ปฏิจฺจ กิเลโส เจว สํกิเลสิโก จ ธมฺโม อุปฺปชฺชติ เหตุปจฺจยา.\csromanspacer{} โลภํ ปฏิจฺจ โมโห, ทิฏฺฐิ, ถินํ, อุทฺธจฺจํ, อหิริกํ, อโนตฺตปฺปํ. (จกฺกํ.) (1)}

\prose{กิเลสญฺเจว สํกิเลสิกญฺจ ธมฺมํ ปฏิจฺจ สํกิเลสิโก เจว โน จ กิเลโส ธมฺโม อุปฺปชฺชติ เหตุปจฺจยา.\csromanspacer{} กิเลเส ปฏิจฺจ สมฺปยุตฺตกา ขนฺธา จิตฺตสมุฏฺฐานญฺจ รูปํ. (2)}

\prose{กิเลสญฺเจว สํกิเลสิกญฺจ ธมฺมํ ปฏิจฺจ กิเลโส เจว สํกิเลสิโก จ สํกิเลสิโก เจว โน จ กิเลโส จ ธมฺมา อุปฺปชฺชนฺติ เหตุปจฺจยา.\csromanspacer{} โลภํ ปฏิจฺจ โมโห, ทิฏฺฐิ, ถินํ, อุทฺธจฺจํ, อหิริกํ, อโนตฺตปฺปํ สมฺปยุตฺตกา จ ขนฺธา จิตฺตสมุฏฺฐานญฺจ รูปํ. (3)}

\prose{\csromanfolio{30}(เอวํ ปฏิจฺจวาโรปิ สหชาตวาโรปิ ปจฺจยวาโรปิ นิสฺสยวาโรปิ สํสฏฺฐวาโรปิ สมฺปยุตฺตวาโรปิ กิเลสทุกสทิสา.\csromanspacer{} นินฺนานากรณํ.\csromanspacer{} อามสนํ นานํ.)}

\csromanheader{2}{79. กิเลสสํกิเลสิกทุก 7. ปญฺหาวาร}
\proseitem{75}{กิเลโส เจว สํกิเลสิโก จ ธมฺโม กิเลสสฺส เจว สํกิเลสิกสฺส จ ธมฺมสฺส เหตุปจฺจเยน ปจฺจโย.\csromanspacer{} กิเลสา เจว สํกิเลสิกา จ เหตู สมฺปยุตฺตกานํ กิเลสานํ เหตุปจฺจเยน ปจฺจโย. (เอวํ จตฺตาริ, กิเลสทุกสทิสํ.) (4)}

\csromanheader{2}{\textbf{อารมฺมณ}}
\proseitem{76}{กิเลโส เจว สํกิเลสิโก จ ธมฺโม กิเลสสฺส เจว สํกิเลสิกสฺส จ ธมฺมสฺส อารมฺมณปจฺจเยน ปจฺจโย.\csromanspacer{} กิเลเส อารพฺภ กิเลสา อุปฺปชฺชนฺติ. (มูลํ กาตพฺพํ.) กิเลเส อารพฺภ สํกิเลสิกา เจว โน จ กิเลสา ขนฺธา อุปฺปชฺชนฺติ. (มูลํ กาตพฺพํ.) กิเลเส อารพฺภ กิเลสา จ สมฺปยุตฺตกา จ ขนฺธา อุปฺปชฺชนฺติ. (3)}

\proseitem{77}{สํกิเลสิโก เจว โน จ กิเลโส ธมฺโม สํกิเลสิกสฺส เจว โน จ กิเลสสฺส ธมฺมสฺส อารมฺมณปจฺจเยน ปจฺจโย.\csromanspacer{} ทานํ ฯเปฯ สีลํ ฯเปฯ อุโปสถกมฺมํ ฯเปฯ.\csromanspacer{} ปุพฺเพ สุจิณฺณานิ ฯเปฯ.\csromanspacer{} ฌานา วุฏฺฐหิตฺวา ฌานํ ปจฺจเวกฺขติ, อสฺสาเทติ อภินนฺทติ, ตํ อารพฺภ ราโค อุปฺปชฺชติ ฯเปฯ อุทฺธจฺจํ อุปฺปชฺชติ.\csromanspacer{} ฌาเน ปริหีเน วิปฺปฏิสาริสฺส โทมนสฺสํ อุปฺปชฺชติ.\csromanspacer{} อริยา โคตฺรภุํ ปจฺจเวกฺขนฺติ, โวทานํ ปจฺจเวกฺขนฺติ.\csromanspacer{} จกฺขุํ ฯเปฯ.\csromanspacer{} วตฺถุํ สํกิเลสิเก เจว โน จ กิเลเส ขนฺเธ อนิจฺจโต ฯเปฯ โทมนสฺสํ อุปฺปชฺชติ.\csromanspacer{} ทิพฺเพน จกฺขุนา รูปํ ปสฺสติ ฯเปฯ อาวชฺชนาย อารมฺมณปจฺจเยน ปจฺจโย. (อิตเร เทฺว กิเลสทุกสทิสา, ฆฏนารมฺมณาปิ กิเลสทุกสทิสา.)}

\csromanheader{2}{79. กิเลสสํกิเลสิกทุก \textbf{อธิปติ}}
\proseitem{78}{\csromanfolio{31}กิเลโส เจว สํกิเลสิโก จ ธมฺโม กิเลสสฺส เจว สํกิเลสิกสฺส จ ธมฺมสฺส อธิปติปจฺจเยน ปจฺจโย.\csromanspacer{} อารมฺมณาธิปติ.\csromanspacer{} ตีณิ.}

\prose{สํกิเลสิโก เจว โน จ กิเลโส ธมฺโม สํกิเลสิกสฺส เจว โน จ กิเลสสฺส ธมฺมสฺส อธิปติปจฺจเยน ปจฺจโย.\csromanspacer{} อารมฺมณาธิปติ, สหชาตาธิปติ.\csromanspacer{}\csromanspacer{}\textbf{อารมฺมณาธิปติ\csromandash{}}ทานํ ฯเปฯ สีลํ ฯเปฯ อุโปสถกมฺมํ ฯเปฯ.\csromanspacer{} ปุพฺเพ สุจิณฺณานิ ฯเปฯ.\csromanspacer{} ฌานา วุฏฺฐหิตฺวา ฌานํ ครุํ กตฺวา ปจฺจเวกฺขติ.\csromanspacer{} เสกฺขา โคตฺรภุํ ครุํ กตฺวา ปจฺจเวกฺขนฺติ, โวทานํ ครุํ กตฺวา ปจฺจเวกฺขนฺติ.\csromanspacer{} จกฺขุํ ฯเปฯ.\csromanspacer{} วตฺถุํ สํกิเลสิเก เจว โน จ กิเลเส ขนฺเธ ครุํ กตฺวา สํกิเลสิกา เจว โน จ กิเลสา ขนฺธา อุปฺปชฺชนฺติ.\csromanspacer{} \textbf{สหชาตาธิปติ\csromandash{}}สํกิเลสิกา เจว โน จ กิเลสาธิปติ สมฺปยุตฺตกานํ ขนฺธานํ จิตฺตสมุฏฺฐานานญฺจ รูปานํ อธิปติปจฺจเยน ปจฺจโย. (อิตเร เทฺวปิ กิเลสทุกสทิสา.\csromanspacer{} ฆฏนาธิปติปิ.)}

\csromanheader{2}{\textbf{อนนฺตราทิ}}
\proseitem{79}{กิเลโส เจว สํกิเลสิโก จ ธมฺโม กิเลสสฺส เจว สํกิเลสิกสฺส จ ธมฺมสฺส อนนฺตรปจฺจเยน ปจฺจโย.\csromanspacer{} ตีณิ. (กิเลสทุกสทิสา.)}

\prose{สํกิเลสิโก เจว โน จ กิเลโส ธมฺโม สํกิเลสิกสฺส เจว โน จ กิเลสสฺส ธมฺมสฺส อนนฺตรปจฺจเยน ปจฺจโย.\csromanspacer{} ปุริมา ปุริมา สํกิเลสิกา เจว โน จ กิเลสา ขนฺธา ปจฺฉิมานํ ปจฺฉิมานํ สํกิเลสิกานญฺเจว โน จ กิเลสานํ ขนฺธานํ อนนฺตรปจฺจเยน ปจฺจโย.\csromanspacer{} อนุโลมํ โคตฺรภุสฺส.\csromanspacer{} อนุโลมํ โวทานสฺส.\csromanspacer{} อาวชฺชนา สํกิเลสิกานญฺเจว โน จ กิเลสานํ ขนฺธานํ อนนฺตรปจฺจเยน ปจฺจโย.}

\prose{(อิตเร เทฺว อนนฺตรา กิเลสทุกสทิสา, นินฺนานากรณา.\csromanspacer{} ฆฏนานนฺตรมฺปิ สพฺเพ ปจฺจยา กิเลสทุกสทิสา, นินฺนานากรณา.\csromanspacer{} อุปนิสฺสเย โลกุตฺตรํ นตฺถิ, อิทํ ทุกํ กิเลสทุกสทิสํ, นินฺนานากรณํ.)}

\nitthitam{กิเลสสํกิเลสิกทุกํ นิฏฺฐิตํ.}
\csromanmatikaanchor{mtk.8}
\csromantocmarkref{section}{80. กิเลสสํกิลิฏฺฐทุก}{mtk.8}
\bookmark[dest={mtk.8},level=1]{80. กิเลสสํกิลิฏฺฐทุก}
\csromanheader{1}{80. กิเลสสํกิลิฏฺฐทุก 1. ปฏิจฺจวาร}
\proseitem{80}{\csromanfolio{32}กิเลสญฺเจว สํกิลิฏฺฐญฺจ ธมฺมํ ปฏิจฺจ กิเลโส เจว สํกิลิฏฺโฐ จ ธมฺโม อุปฺปชฺชติ เหตุปจฺจยา.\csromanspacer{} โลภํ ปฏิจฺจ โมโห, ทิฏฺฐิ, ถินํ, อุทฺธจฺจํ, อหิริกํ, อโนตฺตปฺปํ. (จกฺกํ.) (1)}

\prose{กิเลสญฺจว สํกิลิฏฺฐญฺจ ธมฺมํ ปฏิจฺจ สํกิลิฏฺโฐ เจว โน จ กิเลโส ธมฺโม อุปฺปชฺชติ เหตุปจฺจยา.\csromanspacer{} กิเลเส ปฏิจฺจ สมฺปยุตฺตกา ขนฺธา. (2)}

\prose{กิเลสญฺเจว สํกิลิฏฺฐญฺจ ธมฺมํ ปฏิจฺจ กิเลโส เจว สํกิลิฏฺโฐ จ สํกิลิฏฺโฐ เจว โน จ กิเลโส จ ธมฺมา อุปฺปชฺชติ เหตุปจฺจยา.\csromanspacer{} โลภํ ปฏิจฺจ โมโห, ทิฏฺฐิ, ถินํ, อุทฺธจฺจํ, อหิริกํ, อโนตฺตปฺปํ สมฺปยุตฺตกา จ ขนฺธา. (จกฺกํ.) (3)}

\proseitem{81}{สํกิลิฏฺฐญฺเจว โน จ กิเลสํ ธมฺมํ ปฏิจฺจ สํกิลิฏฺโฐ เจว โน จ กิเลโส ธมฺโม อุปฺปชฺชติ เหตุปจฺจยา.\csromanspacer{} สํกิลิฏฺฐญฺเจว โน จ กิเลสํ เอกํ ขนฺธํ ปฏิจฺจ ตโย ขนฺธา ฯเปฯ เทฺว ขนฺเธ ฯเปฯ. (1)}

\prose{สํกิลิฏฺฐญฺเจว โน จ กิเลสํ ธมฺมํ ปฏิจฺจ กิเลโส เจว สํกิลิฏฺโฐ จ ธมฺโม อุปฺปชฺชติ เหตุปจฺจยา.\csromanspacer{} สํกิลิฏฺเฐ เจว โน จ กิเลเส ขนฺเธ ปฏิจฺจ กิเลสา. (2)}

\prose{สํกิลิฏฺฐญฺเจว โน จ กิเลสํ ธมฺมํ ปฏิจฺจ กิเลโส เจว สํกิลิฏฺโฐ จ สํกิลิฏฺโฐ เจว โน จ กิเลโส จ ธมฺมา อุปฺปชฺชนฺติ เหตุปจฺจยา.\csromanspacer{} สํกิลิฏฺฐญฺเจว โน จ กิเลสํ เอกํ ขนฺธํ ปฏิจฺจ ตโย ขนฺธา กิเลสา จ ฯเปฯ เทฺว ขนฺเธ ฯเปฯ. (3)}

\proseitem{82}{กิเลสญฺเจว สํกิลิฏฺฐญฺจ สํกิลิฏฺฐญฺเจว โน จ กิเลสญฺจ ธมฺมํ ปฏิจฺจ กิเลโส เจว สํกิลิฏฺโฐ จ ธมฺโม อุปฺปชฺชติ เหตุปจฺจยา.\csromanspacer{} โลภญฺจ สมฺปยุตฺตเก จ ขนฺเธ ปฏิจฺจ โมโห, ทิฏฺฐิ, ถินํ, อุทฺธจฺจํ, อหิริกํ, อโนตฺตปฺปํ. (จกฺกํ.) (1)}

\csromanheader{2}{80. กิเลสสํกิลิฏฺฐทุก}
\prose{\csromanfolio{33}กิเลสญฺเจว สํกิลิฏฺฐญฺจ สํกิลิฏฺฐญฺเจว โน จ กิเลสญฺจ ธมฺมํ ปฏิจฺจ สํกิลิฏฺโฐ เจว โน จ กิเลโส ธมฺโม อุปฺปชฺชติ เหตุปจฺจยา.\csromanspacer{} สํกิลิฏฺฐญฺเจว โน จ กิเลสํ เอกํ ขนฺธญฺจ กิเลเส จ ปฏิจฺจ ตโย ขนฺธา ฯเปฯ เทฺว ขนฺเธ จ ฯเปฯ. (2)}

\prose{กิเลสญฺเจว สํกิลิฏฺฐญฺจ สํกิลิฏฺฐญฺเจว โน จ กิเลสญฺจ ธมฺมํ ปฏิจฺจ กิเลโส เจว สํกิลิฏฺโฐ จ สํกิลิฏฺโฐ เจว โน จ กิเลโส จ ธมฺมา อุปฺปชฺชนฺติ เหตุปจฺจยา.\csromanspacer{} สํกิลิฏฺฐญฺเจว โน จ กิเลสํ เอกํ ขนฺธญฺจ โลภญฺจ ปฏิจฺจ ตโย ขนฺธา โมโห, ทิฏฺฐิ, ถินํ, อุทฺธจฺจํ, อหิริกํ, อโนตฺตปฺปํ ฯเปฯ เทฺว ขนฺเธ จ ฯเปฯ. (จกฺกํ.) (3)}

\csromanheader{2}{1. ปจฺจยานุโลม 2. สงฺขฺยาวาร}
\proseitem{83}{เหตุยา นว, อารมฺมเณ นว, (สพฺพตฺถ นว,) กมฺเม นว, อาหาเร นว ฯเปฯ อวิคเต นว.}

\csromanheader{2}{2. ปจฺจยปจฺจนีย 1. วิภงฺควาร}
\csromanheader{2}{\textbf{นเหตุ}}
\proseitem{84}{กิเลสญฺเจว สํกิลิฏฺฐญฺจ ธมฺมํ ปฏิจฺจ กิเลโส เจว สํกิลิฏฺโฐ จ ธมฺโม อุปฺปชฺชติ นเหตุปจฺจยา.\csromanspacer{} วิจิกิจฺฉํ ปฏิจฺจ วิจิกิจฺฉาสหคโต โมโห.\csromanspacer{} อุทฺธจฺจํ ปฏิจฺจ อุทฺธจฺจสหคโต โมโห. (1)}

\prose{สํกิลิฏฺฐญฺเจว โน จ กิเลสํ ธมฺมํ ปฏิจฺจ กิเลโส เจว สํกิลิฏฺโฐ จ ธมฺโม อุปฺปชฺชติ นเหตุปจฺจยา.\csromanspacer{} วิจิกิจฺฉาสหคเต อุทฺธจฺจสหคเต ขนฺเธ ปฏิจฺจ วิจิกิจฺฉาสหคโต อุทฺธจฺจสหคโต โมโห. (1)}

\prose{กิลิสญฺเจว สํกิลิฏฺฐญฺจ สํกิลิฏฺฐญฺเจว โน จ กิเลสญฺจ ธมฺมํ ปฏิจฺจ กิเลโส เจว สํกิลิฏฺโฐ จ ธมฺโม อุปฺปชฺชติ นเหตุปจฺจยา.\csromanspacer{} วิจิกิจฺฉาสหคเต อุทฺธจฺจสหคเต ขนฺเธ จ วิจิกิจฺฉญฺจ อุทฺธจฺจญฺจ ปฏิจฺจ วิจิกิจฺฉาสหคโต อุทฺธจฺจสหคโต โมโห. (1)}

\csromanheader{2}{2. ปจฺจยปจฺจนีย 2. สงฺขฺยาวาร}
\proseitem{85}{\csromanfolio{34}นเหตุยา ตีณิ, น-อธิปติยา นว, นปุเรชาเต นว, นปจฺฉาชาเต นว, น-อาเสวเน นว, นกมฺเม ตีณิ, นวิปาเก นว, นวิปฺปยุตฺเต นว.}

\prose{(เอวํ อิตเร เทฺว คณนาปิ สหชาตวาโรปิ ปจฺจยวาโรปิ นิสฺสยวาโรปิ สํสฏฺฐวาโรปิ สมฺปยุตฺตวาโรปิ ปฏิจฺจวารสทิสา.)}

\csromanheader{2}{80. กิเลสสํกิลิฏฺฐทุก 7. ปญฺหาวาร}
\proseitem{86}{กิเลโส เจว สํกิลิฏฺโฐ จ ธมฺโม กิเลสสฺส เจว สํกิลิฏฺฐสฺส จ ธมฺมสฺส เหตุปจฺจเยน ปจฺจโย.\csromanspacer{} กิเลสา เจว สํกิลิฏฺฐา จ เหตู สมฺปยุตฺตกานํ กิเลสานํ เหตุปจฺจเยน ปจฺจโย. (1)}

\prose{กิเลโส เจว สํกิลิฏฺโฐ จ ธมฺโม สํกิลิฏฺฐสฺส เจว โน จ กิเลสสฺส ธมฺมสฺส เหตุปจฺจเยน ปจฺจโย.\csromanspacer{} กิเลสา เจว สํกิลิฏฺฐา จ เหตู สมฺปยุตฺตกานํ ขนฺธานํ เหตุปจฺจเยน ปจฺจโย. (2)}

\prose{กิเลโส เจว สํกิลิฏฺโฐ จ ธมฺโม กิเลสสฺส เจว สํกิลิฏฺฐสฺส จ สํกิลิฏฺฐสฺส เจว โน จ กิเลสสฺส จ ธมฺมสฺส เหตุปจฺจเยน ปจฺจโย.\csromanspacer{} กิเลสา เจว สํกิลิฏฺฐา จ เหตู สมฺปยุตฺตกานํ ขนฺธานํ กิเลสานญฺจ เหตุปจฺจเยน ปจฺจโย. (3)}

\csromanheader{2}{\textbf{อารมฺมณ}}
\proseitem{87}{กิเลโส เจว สํกิลิฏฺโฐ จ ธมฺโม กิเลสสฺส เจว สํกิลิฏฺฐสฺส จ ธมฺมสฺส อารมฺมณปจฺจเยน ปจฺจโย.\csromanspacer{} กิเลเส อารพฺภ กิเลสา อุปฺปชฺชนฺติ. (มูลํ กาตพฺพํ.) กิเลเส อารพฺภ สํกิลิฏฺฐา เจว โน จ กิเลสา ขนฺธา อุปฺปชฺชนฺติ. (มูลํ กาตพฺพํ.) กิเลเส อารพฺภ กิเลสา จ สมฺปยุตฺตกา ขนฺธา จ อุปฺปชฺชนฺติ. (3)}

\csromanheader{2}{80. กิเลสสํกิลิฏฺฐทุก}
\proseitem{88}{\csromanfolio{35}สํกิลิฏฺโฐ เจว โน จ กิเลโส ธมฺโม สํกิลิฏฺฐสฺส เจว โน จ กิเลสสฺส ธมฺมสฺส อารมฺมณปจฺจเยน ปจฺจโย.\csromanspacer{} สํกิลิฏฺเฐ เจว โน จ กิเลเส ขนฺเธ อารพฺภ สํกิลิฏฺฐา เจว โน จ กิเลสา ขนฺธา อุปฺปชฺชนฺติ. (มูลํ กาตพฺพํ.) สํกิลิฏฺเฐ เจว โน จ กิเลเส ขนฺเธ อารพฺภ กิเลสา อุปฺปชฺชนฺติ. (มูลํ กาตพฺพํ.) สํกิลิฏฺเฐ เจว โน จ กิเลเส ขนฺเธ อารพฺภ กิเลสา จ สมฺปยุตฺตกา ขนฺธา จ อุปฺปชฺชนฺติ. (3) (อิตเรปิ ตีณิ กาตพฺพา.)}

\csromanheader{2}{\textbf{อธิปติ}}
\proseitem{89}{กิเลโส เจว สํกิลิฏฺโฐ จ ธมฺโม กิเลสสฺส เจว สํกิลิฏฺฐสฺส จ ธมฺมสฺส อธิปติปจฺจเยน ปจฺจโย.\csromanspacer{} อารมฺมณาธิปติ.\csromanspacer{} ตีณิ.}

\prose{สํกิลิฏฺโฐ เจว โน จ กิเลโส ธมฺโม สํกิลิฏฺฐสฺส เจว โน จ กิเลสสฺส ธมฺมสฺส อธิปติปจฺจเยน ปจฺจโย.\csromanspacer{} อารมฺมณาธิปติ, สหชาตาธิปติ.\csromanspacer{}\csromanspacer{}\textbf{อารมฺมณาธิปติ\csromandash{}}สํกิลิฏฺเฐ เจว โน จ กิเลเส ขนฺเธ ครุํ กตฺวา ฯเปฯ.\csromanspacer{} ตีณิ. (เทฺว อธิปติ.\csromanspacer{} ตีณิปิ กาตพฺพา.\csromanspacer{} อิตเร เทฺวปิ ตีณิ กาตพฺพา.)}

\csromanheader{2}{\textbf{อนนฺตราทิ}}
\proseitem{90}{กิเลโส เจว สํกิลิฏฺโฐ จ ธมฺโม กิเลสสฺส เจว สํกิลิฏฺฐสฺส จ ธมฺมสฺส อนนฺตรปจฺจเยน ปจฺจโย. (นวปิ กาตพฺพา.\csromanspacer{} อาวชฺชนาปิ วุฏฺฐานมฺปิ นตฺถิ.) สมนนฺตรปจฺจเยน ปจฺจโย.\csromanspacer{} สหชาตปจฺจเยน ปจฺจโย.\csromanspacer{} อญฺญมญฺญปจฺจเยน ปจฺจโย.\csromanspacer{} นิสฺสยปจฺจเยน ปจฺจโย.\csromanspacer{} อุปนิสฺสยปจฺจเยน ปจฺจโย.\csromanspacer{} นว ปญฺหา. (ปุเรชาตปจฺจโย ปจฺฉาชาตปจฺจโยปิ นตฺถิ.) อาเสวนปจฺจเยน ปจฺจโย.}

\csromanheader{2}{\textbf{กมฺมาทิ}}
\proseitem{91}{สํกิลิฏฺโฐ เจว โน จ กิเลโส ธมฺโม สํกิลิฏฺฐสฺส เจว โน จ กิเลสสฺส ธมฺมสฺส กมฺมปจฺจเยน ปจฺจโย.\csromanspacer{} สํกิลิฏฺฐา เจว โน จ กิเลสา เจตนา สมฺปยุตฺตกานํ ขนฺธานํ กมฺมปจฺจเยน ปจฺจโย. (1)}

\prose{สํกิลิฏฺโฐ เจว โน จ กิเลโส ธมฺโม กิเลสสฺส เจว สํกิลิฏฺฐสฺส จ ธมฺมสฺส กมฺมปจฺจเยน ปจฺจโย.\csromanspacer{} สํกิลิฏฺฐา เจว, \csromanfolio{36}โน จ กิเลสา เจตนา สมฺปยุตฺตกานํ กิเลสานํ กมฺมปจฺจเยน ปจฺจโย. (2)}

\prose{สํกิลิฏฺโฐ เจว โน จ กิเลโส ธมฺโม กิเลสสฺส เจว สํกิลิฏฺฐสฺส จ สํกิลิฏฺฐสฺส เจว โน จ กิเลสสฺส ธมฺมสฺส กมฺมปจฺจเยน ปจฺจโย.\csromanspacer{} สํกิลิฏฺฐา เจว โน จ กิเลสา เจตนา สมฺปยุตฺตกานํ ขนฺธานํ กิเลสานญฺจ กมฺมปจฺจเยน ปจฺจโย. (3)}

\prose{อาหารปจฺจเยน ปจฺจโย.\csromanspacer{} ตีณิ.\csromanspacer{} อินฺทฺริยปจฺจเยน ปจฺจโย.\csromanspacer{} ตีณิ.\csromanspacer{} ฌานปจฺจเยน ปจฺจโย.\csromanspacer{} ตีณิ.\csromanspacer{} มคฺคปจฺจเยน ปจฺจโย.\csromanspacer{} นว.\csromanspacer{} สมฺปยุตฺตปจฺจเยน ปจฺจโย.\csromanspacer{} นว.\csromanspacer{} อตฺถิปจฺจเยน ปจฺจโย.\csromanspacer{} นว.\csromanspacer{} นตฺถิปจฺจเยน ปจฺจโย.\csromanspacer{} วิคตปจฺจเยน ปจฺจโย.\csromanspacer{} อวิคตปจฺจเยน ปจฺจโย.\csromanspacer{} นว.}

\csromanheader{2}{1. ปจฺจยานุโลม 2. สงฺขฺยาวาร}
\csromanheader{2}{\textbf{สุทฺธ}}
\proseitem{92}{เหตุยา ตีณิ, อารมฺมเณ นว, อธิปติยา นว, อนนฺตเร นว, สมนนฺตเร นว, สหชาเต นว, อญฺญมญฺเญ นว, นิสฺสเย นว, อุปนิสฺสเย นว, อาเสวเน นว, กมฺเม ตีณิ, อาหาเร ตีณิ, อินฺทฺริเย ตีณิ, ฌาเน ตีณิ, มคฺเค นว, สมฺปยุตฺเต นว, อตฺถิยา นว, นตฺถิยา นว, วิคเต นว, อวิคเต นว.}

\csromanheader{2}{2. ปจฺจนียุทฺธาร}
\proseitem{93}{กิเลโส เจว สํกิลิฏฺโฐ จ ธมฺโม กิเลสสฺส เจว สํกิลิฏฺฐสฺส จ ธมฺมสฺส อารมฺมณปจฺจเยน ปจฺจโย.\csromanspacer{} สหชาตปจฺจเยน ปจฺจโย.\csromanspacer{} อุปนิสฺสยปจฺจเยน ปจฺจโย. (นวปิ, ตีณิเยว ปทา กาตพฺพา.)}

\csromanheader{2}{2. ปจฺจยปจฺจนีย 2. สงฺขฺยาวาร}
\vspace{\tipitakaparskip}
\csromancenter{นเหตุยา นว, น-อารมฺมเณ นว, (สพฺพตฺถ นว,) โน-อวิคเต นว.}
\csromanmatikaanchor{mtk.9}
\csromanmatikaanchor{mtk.10}
\csromantocmarkref{section}{81. กิเลสกิเลสสมฺปยุตฺตทุก}{mtk.9}
\bookmark[dest={mtk.9},level=1]{81. กิเลสกิเลสสมฺปยุตฺตทุก}
\csromantocmarkref{section}{82. กิเลสวิปฺปยุตฺตสํกิเลสิกทุก}{mtk.10}
\bookmark[dest={mtk.10},level=1]{82. กิเลสวิปฺปยุตฺตสํกิเลสิกทุก}
\csromanheader{1}{82. กิเลสวิปฺปยุตฺตสํกิเลสิกทุก \textbf{3. ปจฺจยานุโลมปจฺจนีย}}
\proseitem{95}{\csromanfolio{37}เหตุปจฺจยา น-อารมฺมเณ ตีณิ, น-อธิปติยา ตีณิ, น-อนนฺตเร ตีณิ, นสมนนฺตเร ตีณิ, น-อุปนิสฺสเย ตีณิ ฯเปฯ นมคฺเค ตีณิ ฯเปฯ นวิปฺปยุตฺเต ตีณิ, โนนตฺถิยา ตีณิ, โนวิคเต ตีณิ.}

\csromanheader{2}{4. ปจฺจยปจฺจนียานุโลม}
\csromanheader{2}{96. นเหตุปจฺจยา อารมฺมเณ นว, (อนุโลมมาติกา กาตพฺพา) ฯเปฯ อวิคเต นว.}
\nitthitam{กิเลสสํกิลิฏฺฐทุกํ นิฏฺฐิตํ.}
\csromanheader{1}{81. กิเลสกิเลสสมฺปยุตฺตทุก 1. ปฏิจฺจวาร}
\proseitem{97}{กิเลสญฺเจว กิเลสสมฺปยุตฺตญฺจ ธมฺมํ ปฏิจฺจ กิเลโส เจว กิเลสสมฺปยุตฺโต จ ธมฺโม อุปฺปชฺชติ เหตุปจฺจยา.\csromanspacer{} โลภํ ปฏิจฺจ โมโห, ทิฏฺฐิ, ถินํ, อุทฺธจฺจํ, อหิริกํ, อโนตฺตปฺปํ. (กิเลสสํกิลิฏฺฐทุกสทิสํ, นินฺนานากรณํ, สพฺเพ วารา.)}

\nitthitam{กิเลสกิเลสสมฺปยุตฺตทุกํ นิฏฺฐิตํ.}
\csromanheader{2}{82. กิเลสวิปฺปยุตฺตสํกิเลสิกทุก 1. ปฏิจฺจวาร}
\proseitem{98}{กิเลสวิปฺปยุตฺตํ สํกิเลสิกํ ธมฺมํ ปฏิจฺจ กิเลสวิปฺปยุตฺโต สํกิเลสิโก ธมฺโม อุปฺปชฺชติ เหตุปจฺจยา.\csromanspacer{} กิเลสวิปฺปยุตฺตํ สํกิเลสิกํ เอกํ ขนฺธํ ปฏิจฺจ ตโย ขนฺธา จิตฺตสมุฏฺฐานญฺจ รูปํ ฯเปฯ เทฺว ขนฺเธ ฯเปฯ. (ยถา โลกิยทุกํ, เอวํ นินฺนานากรณํ.) กิเลสวิปฺปยุตฺตสํกิเลสิกทุกํ นิฏฺฐิตํ.\csromanspacer{} กิเลสโคจฺฉกํ นิฏฺฐิตํ.}

\csromanmatikaanchor{mtk.11}
\csromantocmarknopage{chapter}{13. ปิฏฺฐิทุก}{mtk.11}
\bookmark[dest={mtk.11},level=0]{13. ปิฏฺฐิทุก}
\chapterheadpageread{13. ปิฏฺฐิทุก}
\csromanmatikaanchor{mtk.12}
\csromantocmarkref{section}{83. ทสฺสเนนปหาตพฺพทุก}{mtk.12}
\bookmark[dest={mtk.12},level=1]{83. ทสฺสเนนปหาตพฺพทุก}
\csromanheader{1}{83. ทสฺสเนนปหาตพฺพทุก 1. ปฏิจฺจวาร}
\proseitem{1}{\csromanfolio{38}ทสฺสเนน ปหาตพฺพํ ธมฺมํ ปฏิจฺจ ทสฺสเนน ปหาตพฺโพ ธมฺโม อุปฺปชฺชติ เหตุปจฺจยา.\csromanspacer{} ทสฺสเนน ปหาตพฺพํ เอกํ ขนฺธํ ปฏิจฺจ ตโย ขนฺธา ฯเปฯ เทฺว ขนฺเธ ฯเปฯ. (1)}

\prose{ทสฺสเนน ปหาตพฺพํ ธมฺมํ ปฏิจฺจ นทสฺสเนน ปหาตพฺโพ ธมฺโม อุปฺปชฺชติ เหตุปจฺจยา.\csromanspacer{} ทสฺสเนน ปหาตพฺเพ ขนฺเธ ปฏิจฺจ จิตฺตสมุฏฺฐานํ รูปํ. (2)}

\prose{ทสฺสเนน ปหาตพฺพํ ธมฺมํ ปฏิจฺจ ทสฺสเนน ปหาตพฺโพ จ นทสฺสเนน ปหาตพฺโพ จ ธมฺมา อุปฺปชฺชนฺติ เหตุปจฺจยา.\csromanspacer{} ทสฺสเนน ปหาตพฺพํ เอกํ ขนฺธํ ปฏิจฺจ ตโย ขนฺธา จิตฺตสมุฏฺฐานญฺจ ฯเปฯ เทฺว ขนฺเธ ฯเปฯ. (3)}

\prose{นทสฺสเนน ปหาตพฺพํ ธมฺมํ ปฏิจฺจ นทสฺสเนน ปหาตพฺโพ ธมฺโม อุปฺปชฺชติ เหตุปจฺจยา.\csromanspacer{} นทสฺสเนน ปหาตพฺพํ เอกํ ขนฺธํ ปฏิจฺจ ตโย ขนฺธา จิตฺตสมุฏฺฐานญฺจ รูปํ ฯเปฯ เทฺว ขนฺเธ ฯเปฯ.\csromanspacer{} ปฏิสนฺธิกฺขเณ ฯเปฯ.\csromanspacer{} ขนฺเธ ปฏิจฺจ วตฺถุ, วตฺถุํ ปฏิจฺจ ขนฺธา.\csromanspacer{} เอกํ มหาภูตํ ฯเปฯ. (1)}

\prose{ทสฺสเนน ปหาตพฺพญฺจ นทสฺสเนน ปหาตพฺพญฺจ ธมฺมํ ปฏิจฺจ นทสฺสเนน ปหาตพฺโพ ธมฺโม อุปฺปชฺชติ เหตุปจฺจยา.\csromanspacer{} ทสฺสเนน ปหาตพฺเพ ขนฺเธ จ มหาภูเต จ ปฏิจฺจ จิตฺตสมุฏฺฐานํ รูปํ. (สํขิตฺตํ.) (1)}

\csromanheader{2}{1. ปจฺจยานุโลม 2. สงฺขฺยาวาร}
\vspace{\tipitakaparskip}
\csromancenter{\textbf{เหตุ}ยา ปญฺจ, อารมฺมเณ เทฺว, อธิปติยา ปญฺจ, อนนฺตเร เทฺว, สมนนฺตเร เทฺว, สหชาเต ปญฺจ, อญฺญมญฺเญ เทฺว, นิสฺสเย ปญฺจ, อุปนิสฺสเย เทฺว, ปุเรชาเต เทฺว, อาเสวเน เทฺว, กมฺเม ปญฺจ, วิปาเก เอกํ, อาหาเร ปญฺจ ฯเปฯ อวิคเต ปญฺจ.}
\csromanheader{2}{83. ทสฺสเนนปหาตพฺพทุก \textbf{2. ปจฺจยปจฺจนีย 1. วิภงฺควาร}}
\csromanheader{2}{\textbf{นเหตุ}}
\proseitem{3}{\csromanfolio{39}ทสฺสเนน ปหาตพฺพํ ธมฺมํ ปฏิจฺจ ทสฺสเนน ปหาตพฺโพ ธมฺโม อุปฺปชฺชติ นเหตุปจฺจยา.\csromanspacer{} วิจิกิจฺฉาสหคเต ขนฺเธ ปฏิจฺจ วิจิกิจฺฉาสหคโต โมโห. (1)}

\prose{นทสฺสเนน ปหาตพฺพํ ธมฺมํ ปฏิจฺจ นทสฺสเนน ปหาตพฺโพ ธมฺโม อุปชฺชติ นเหตุปจฺจยา.\csromanspacer{} อเหตุกํ นทสฺสเนน ปหาตพฺพํ เอกํ ขนฺธํ ฯเปฯ.\csromanspacer{} อเหตุกปฏิสนฺธิกฺขเณ ฯเปฯ. (ยาว อสญฺญสตฺตา.) อุทฺธจฺจสหคเต ขนฺเธ ปฏิจฺจ อุทฺธจฺจสหคโต โมโห. (1)}

\csromanheader{2}{2. ปจฺจยปจฺจนีย 2. สงฺขฺยาวาร}
\vspace{\tipitakaparskip}
\csromancenter{\textbf{นเหตุ}ยา เทฺว, น-อารมฺมเณ ตีณิ, น-อธิปติยา ปญฺจ, น-อนนฺตเร ตีณิ, นสมนนฺตเร ตีณิ, น-อญฺญมญฺเญ ตีณิ, น-อุปนิสฺสเย ตีณิ, นปุเรชาเต จตฺตาริ, นปจฺฉาชาเต ปญฺจ, น-อาเสวเน ปญฺจ, นกมฺเม เทฺว, นวิปาเก ปญฺจ, น-อาหาเร เอกํ, น-อินฺทฺริเย เอกํ, นฌาเน เอกํ, นมคฺเค เอกํ, นสมฺปยุตฺเต ตีณิ, นวิปฺปยุตฺเต เทฺว, โนนตฺถิยา ตีณิ, โนวิคเต ตีณิ.}
\csromancenter{(อิตเร เทฺว คณนาปิ สหชาตวาโรปิ กาตพฺโพ.)}
\csromanheader{2}{83. ทสฺสเนนปหาตพฺพทุก 3. ปจฺจยวาร}
\proseitem{5}{ทสฺสเนน ปหาตพฺพํ ธมฺมํ ปจฺจยา ทสฺสเนน ปหาตพฺโพ ธมฺโม อุปฺปชฺชติ เหตุปจฺจยา.\csromanspacer{} ตีณิ. (ปฏิจฺจสทิสา.)}

\proseitem{6}{นทสฺสเนน ปหาตพฺพํ ธมฺมํ ปจฺจยา นทสฺสเนน ปหาตพฺโพ ธมฺโม อุปฺปชฺชติ เหตุปจฺจยา.\csromanspacer{} นทสฺสเนน ปหาตพฺพํ เอกํ ขนฺธํ ปจฺจยา ตโย ขนฺธา \csromanfolio{40}จิตฺตสมุฏฺฐานญฺจ รูปํ ฯเปฯ เทฺว ขนฺเธ ฯเปฯ.\csromanspacer{} ปฏิสนฺธิกฺขเณ ฯเปฯ. (ยาว อชฺฌตฺติกา มหาภูตา.) วตฺถุํ ปจฺจยา นทสฺสเนน ปหาตพฺพา ขนฺธา. (1)}

\prose{นทสฺสเนน ปหาตพฺพํ ธมฺมํ ปจฺจยา ทสฺสเนน ปหาตพฺโพ ธมฺโม อุปฺปชฺชติ เหตุปจฺจยา.\csromanspacer{} วตฺถุํ ปจฺจยา ทสฺสเนน ปหาตพฺพา ขนฺธา. (2)}

\prose{นทสฺสเนน ปหาตพฺพํ ธมฺมํ ปจฺจยา ทสฺสเนน ปหาตพฺโพ จ นทสฺสเนน ปหาตพฺโพ จ ธมฺมา อุปฺปชฺชนฺติ เหตุปจฺจยา.\csromanspacer{} วตฺถุํ ปจฺจยา ทสฺสเนน ปหาตพฺพา ขนฺธา.\csromanspacer{} มหาภูเต ปจฺจยา จิตฺตสมุฏฺฐานํ. (3)}

\proseitem{7}{ทสฺสเนน ปหาตพฺพญฺจ นทสฺสเนน ปหาตพฺพญฺจ ธมฺมํ ปจฺจยา ทสฺสเนน ปหาตพฺโพ ธมฺโม อุปฺปชฺชติ เหตุปจฺจยา.\csromanspacer{} ทสฺสเนน ปหาตพฺพํ เอกํ ขนฺธญฺจ วตฺถุญฺจ ปจฺจยา ตโย ขนฺธา ฯเปฯ เทฺว ขนฺเธ จ ฯเปฯ. (1)}

\prose{ทสฺสเนน ปหาตพฺพญฺจ นทสฺสเนน ปหาตพฺพญฺจ ธมฺมํ ปจฺจยา นทสฺสเนน ปหาตพฺโพ ธมฺโม อุปฺปชฺชติ เหตุปจฺจยา.\csromanspacer{} ทสฺสเนน ปหาตพฺเพ ขนฺเธ จ มหาภูเต จ ปจฺจยา จิตฺตสมุฏฺฐานํ รูปํ. (2)}

\prose{ทสฺสเนน ปหาตพฺพญฺจ นทสฺสเนน ปหาตพฺพญฺจ ธมฺมํ ปจฺจยา ทสฺสเนน ปหาตพฺโพ จ นทสฺสเนน ปหาตพฺโพ จ ธมฺมา อุปฺปชฺชนฺติ เหตุปจฺจยา.\csromanspacer{} ทสฺสเนน ปหาตพฺพํ เอกํ ขนฺธญฺจ วตฺถุญฺจ ปจฺจยา ตโย ขนฺธา ฯเปฯ เทฺว ขนฺเธ จ ฯเปฯ.\csromanspacer{} ทสฺสเนน ปหาตพฺเพ ขนฺเธ จ มหาภูเต จ ปจฺจยา จิตฺตสมุฏฺฐานํ รูปํ. (3)}

\csromanheader{2}{1. ปจฺจยานุโลม 2. สงฺขฺยาวาร}
\vspace{\tipitakaparskip}
\csromancenter{เหตุยา นว, อารมฺมเณ จตฺตาริ, อธิปติยา นว ฯเปฯ อวิคเต นว.}
\csromanheader{2}{2. ปจฺจยปจฺจนีย 1. วิภงฺควาร}
\csromanheader{2}{\textbf{นเหตุ}}
\proseitem{9}{ทสฺสเนน ปหาตพฺพํ ธมฺมํ ปจฺจยา ทสฺสเนน ปหาตพฺโพ ธมฺโม อุปฺปชฺชติ นเหตุปจฺจยา.\csromanspacer{} วิจิกิจฺฉาสหคเต ขนฺเธ ปจฺจยา วิจิกิจฺฉาสหคโต โมโห. (1)}

\csromanheader{2}{83. ทสฺสเนนปหาตพฺพทุก}
\prose{\csromanfolio{41}นทสฺสเนน ปหาตพฺพํ ธมฺมํ ปจฺจยา นทสฺสเนน ปหาตพฺโพ ธมฺโม อุปฺปชฺชติ นเหตุปจฺจยา.\csromanspacer{} อเหตุกํ นทสฺสเนน ปหาตพฺพํ เอกํ ขนฺธํ ปจฺจยา ตโย ขนฺธา จิตฺตสมุฏฺฐานญฺจ รูปํ ฯเปฯ เทฺว ขนฺเธ ฯเปฯ. (ยาว อสญฺญสตฺตา.) จกฺขายตนํ ปจฺจยา จกฺขุวิญฺญาณํ ฯเปฯ.\csromanspacer{} กายายตนํ ปจฺจยา กายวิญฺญาณํ.\csromanspacer{} วตฺถุํ ปจฺจยา อเหตุกา นทสฺสเนน ปหาตพฺพา ขนฺธา.\csromanspacer{} อุทฺธจฺจสหคเต ขนฺเธ ปจฺจยา อุทฺธสหคโต โมโห.\csromanspacer{} วตฺถุํ ปจฺจยา อุทฺธจฺจสหคโต โมโห. (1)}

\prose{นทสฺสเนน ปหาตพฺพํ ธมฺมํ ปจฺจยา ทสฺสเนน ปหาตพฺโพ ธมฺโม อุปฺปชฺชติ นเหตุปจฺจยา.\csromanspacer{} วตฺถุํ ปจฺจยา วิจิกิจฺฉาสหคโต โมโห. (2)}

\prose{ทสฺสเนน ปหาตพฺพญฺจ นทสฺสเนน ปหาตพฺพญฺจ ธมฺมํ ปจฺจยา ทสฺสเนน ปหาตพฺโพ ธมฺโม อุปฺปชฺชติ นเหตุปจฺจยา.\csromanspacer{} วิจิกิจฺฉาสหคเต ขนฺเธ จ วตฺถุญฺจ ปจฺจยา วิจิกิจฺฉาสหคโต โมโห. (1)}

\csromanheader{2}{2. ปจฺจยปจฺจนีย 2. สงฺขฺยาวาร}
\csromanheader{2}{\textbf{สุทฺธ}}
\proseitem{10}{นเหตุยา จตฺตาริ, น-อารมฺมเณ ตีณิ, น-อธิปติยา นว, น-อนนฺตเร ตีณิ, นสมนนฺตเร ตีณิ, น-อญฺญมญฺเญ ตีณิ, น-อุปนิสฺสเย ตีณิ, นปุเรชาเต จตฺตาริ, นปจฺฉาชาเต นว, น-อาเสวเน นว, นกมฺเม จตฺตาริ, นวิปาเก นว, น-อาหาเร เอกํ, น-อินฺทฺริเย เอกํ, นฌาเน เอกํ, นมคฺเค เอกํ, นสมฺปยุตฺเต ตีณิ, นวิปฺปยุตฺเต เทฺว, โนนตฺถิยา ตีณิ, โนวิคเต ตีณิ.}

\vspace{\tipitakaparskip}
\csromancenter{(อิตเร เทฺว คณนาปิ นิสฺสยวาโรปิ กาตพฺโพ.)}
\csromanheader{2}{83. ทสฺสเนนปหาตพฺพทุก 5. สํสฏฺฐวาร}
\csromanheader{2}{1. ปจฺจยานุโลมาทิ}
\proseitem{11}{ทสฺสเนน ปหาตพฺพํ ธมฺมํ สํสฏฺโฐ ทสฺสเนน ปหาตพฺโพ ธมฺโม อุปฺปชฺชติ เหตุปจฺจยา. (สํขิตฺตํ.) (1)}

\prose{\csromanfolio{42}นทสฺสเนน ปหาตพฺพํ ธมฺมํ สํสฏฺโฐ นทสฺสเนน ปหาตพฺโพ ธมฺโม อุปฺปชฺชติ เหตุปจฺจยา. (สํขิตฺตํ.) (1)}

\prose{\textbf{เหตุ}ยา เทฺว, อารมฺมเณ เทฺว, อธิปติยา เทฺว, อวิคเต เทฺว.\csromanspacer{} อนุโลมํ.}

\proseitem{12}{ทสฺสเนน ปหาตพฺพํ ธมฺมํ สํสฏฺโฐ ทสฺสเนน ปหาตพฺโพ ธมฺโม อุปฺปชฺชติ นเหตุปจฺจยา.\csromanspacer{} วิจิกิจฺฉาสหคเต ขนฺเธ สํสฏฺโฐ วิจิกิจฺฉาสหคโต โมโห. (1)}

\prose{นทสฺสเนน ปหาตพฺพํ ธมฺมํ สํสฏฺโฐ นทสฺสเนน ปหาตพฺโพ ธมฺโม อุปฺปชฺชติ นเหตุปจฺจยา. (สํขิตฺตํ.) (1)}

\prose{นเหตุยา เทฺว, น-อธิปติยา เทฺว, นปุเรชาเต เทฺว, นปจฺฉาชาเต เทฺว, น-อาเสวเน เทฺว, นกมฺเม เทฺว, นวิปาเก เทฺว, นฌาเน เอกํ, นมคฺเค เอกํ, นวิปฺปยุตฺเต เทฺว.}

\vspace{\tipitakaparskip}
\csromancenter{ปจฺจนียํ.}
\csromancenter{(เอวํ อิตเร เทฺว คณนาปิ สมฺปยุตฺตวาโรปิ กาตพฺโพ.)}
\csromanheader{2}{83. ทสฺสเนนปหาตพฺพทุก 7. ปญฺหาวาร}
\proseitem{13}{ทสฺสเนน ปหาตพฺโพ ธมฺโม ทสฺสเนน ปหาตพฺพสฺส ธมฺมสฺส เหตุปจฺจเยน ปจฺจโย.\csromanspacer{} ตีณิ.}

\prose{นทสฺสเนน ปหาตพฺโพ ธมฺโม นทสฺสเนน ปหาตพฺพสฺส ธมฺมสฺส เหตุปจฺจเยน ปจฺจโย.\csromanspacer{} นทสฺสเนน ปหาตพฺพา เหตู สมฺปยุตฺตกานํ ขนฺธานํ จิตฺตสมุฏฺฐานญฺจ รูปานํ เหตุปจฺจเยน ปจฺจโย.\csromanspacer{} ปฏิสนฺธิกฺขเณ ฯเปฯ. (1)}

\csromanheader{2}{83. ทสฺสเนนปหาตพฺพทุก \textbf{อารมฺมณ}}
\proseitem{14}{\csromanfolio{43}ทสฺสเนน ปหาตพฺโพ ธมฺโม ทสฺสเนน ปหาตพฺพสฺส ธมฺมสฺส อารมฺมณปจฺจเยน ปจฺจโย.\csromanspacer{} ทสฺสเนน ปหาตพฺพํ ราคํ อสฺสาเทติ อภินนฺทติ, ตํ อารพฺภ ทสฺสเนน ปหาตพฺโพ ราโค อุปฺปชฺชติ.\csromanspacer{} ทิฏฺฐิ ฯเปฯ วิจิกิจฺฉา อุปฺปชฺชติ.\csromanspacer{} ทสฺสเนน ปหาตพฺพํ โทมนสฺสํ อุปฺปชฺชติ.\csromanspacer{} ทสฺสเนน ปหาตพฺพํ ทิฏฺฐิํ อสฺสาเทติ อภินนฺทติ, ตํ อารพฺภ ทสฺสเนน ปหาตพฺโพ ราโค อุปฺปชฺชติ.\csromanspacer{} ทิฏฺฐิ ฯเปฯ วิจิกิจฺฉา ฯเปฯ โทมนสฺสํ อุปฺปชฺชติ.\csromanspacer{} วิจิกิจฺฉํ อารพฺภ วิจิกิจฺฉา อุปฺปชฺชติ.\csromanspacer{} ทิฏฺฐิ ฯเปฯ โทมนสฺสํ อุปฺปชฺชติ.\csromanspacer{} ทสฺสเนน ปหาตพฺพํ โทมนสฺสํ อารพฺภ ทสฺสเนน ปหาตพฺพํ โทมนสฺสํ อุปฺปชฺชติ.\csromanspacer{} ทิฏฺฐิ อุปฺปชฺชติ.\csromanspacer{} วิจิกิจฺฉา อุปฺปชฺชติ. (1)}

\prose{ทสฺสเนน ปหาตพฺโพ ธมฺโม นทสฺสเนน ปหาตพฺพสฺส ธมฺมสฺส อารมฺมณปจฺจเยน ปจฺจโย.\csromanspacer{} อริยา ทสฺสเนน ปหาตพฺเพ ปหีเน กิเลเส ปจฺจเวกฺขนฺติ.\csromanspacer{} ปุพฺเพ สมุทาจิณฺเณ ฯเปฯ.\csromanspacer{} ทสฺสเนน ปหาตพฺเพ ขนฺเธ อนิจฺจโต ฯเปฯ วิปสฺสติ.\csromanspacer{} เจโตปริยญาเณน ทสฺสเนน ปหาตพฺพจิตฺตสมงฺคิสฺส จิตฺตํ ชานาติ.\csromanspacer{} ทสฺสเนน ปหาตพฺพา ขนฺธา เจโตปริยญาณสฺส, ปุพฺเพนิวาสานุสฺสติญาณสฺส, ยถากมฺมูปคญาณสฺส, อนาคตํสญาณสฺส, อาวชฺชนาย อารมฺมณปจฺจเยน ปจฺจโย. (2)}

\proseitem{15}{นทสฺสเนน ปหาตพฺโพ ธมฺโม นทสฺสเนน ปหาตพฺพสฺส ธมฺมสฺส อารมฺมณปจฺจเยน ปจฺจโย.\csromanspacer{} ทานํ ฯเปฯ สีลํ ฯเปฯ อุโปสถกมฺมํ กตฺวา ตํ ปจฺจเวกฺขติ, อสฺสาเทติ อภินนฺทติ, ตํ อารพฺภ นทสฺสเนน ปหาตพฺโพ ราโค อุปฺปชฺชติ.\csromanspacer{} อุทฺธจฺจํ ฯเปฯ.\csromanspacer{} นทสฺสเนน ปหาตพฺพํ โทมนสฺสํ อุปฺปชฺชติ.\csromanspacer{} ปุพฺเพ สุจิณฺณานิ ฯเปฯ.\csromanspacer{} ฌานา ฯเปฯ.\csromanspacer{} อริยา มคฺคา วุฏฺฐหิตฺวา มคฺคํ ปจฺจเวกฺขนฺติ ฯเปฯ ผลสฺส, อาวชฺชนาย อารมฺมณปจฺจเยน ปจฺจโย.\csromanspacer{} อริยา นทสฺสเนน ปหาตพฺเพ ปหีเน กิเลเส ปจฺจเวกฺขนฺติ.\csromanspacer{} วิกฺขมฺภิเต กิเลเส ฯเปฯ.\csromanspacer{} ปุพฺเพ ฯเปฯ.\csromanspacer{} จกฺขุํ ฯเปฯ.\csromanspacer{} วตฺถุํ นทสฺสเนน ปหาตพฺเพ ขนฺเธ อนิจฺจโต ฯเปฯ วิปสฺสติ, อสฺสาเทติ อภินนฺทติ, ตํ อารพฺภ นทสฺสเนน ปหาตพฺโพ ราโค อุปฺปชฺชติ.\csromanspacer{} อุทฺธจฺจํ ฯเปฯ.\csromanspacer{} นทสฺสเนน ปหาตพฺพํ โทมนสฺสํ อุปฺปชฺชติ.\csromanspacer{} ทิพฺเพน จกฺขุนา รูปํ ปสฺสติ ฯเปฯ อนาคตํสญาณสฺส, อาวชฺชนาย อารมฺมณปจฺจเยน ปจฺจโย. (1)}

\prose{\csromanfolio{44}นทสฺสเนน ปหาตพฺโพ ธมฺโม ทสฺสเนน ปหาตพฺพสฺส ธมฺมสฺส อารมฺมณปจฺจเยน ปจฺจโย.\csromanspacer{} ทานํ ฯเปฯ สีลํ ฯเปฯ อุโปสถกมฺมํ ฯเปฯ.\csromanspacer{} ปุพฺเพ ฯเปฯ.\csromanspacer{} ฌานา ฯเปฯ.\csromanspacer{} จกฺขุํ ฯเปฯ.\csromanspacer{} วตฺถุํ นทสฺสเนน ปหาตพฺเพ ขนฺเธ อสฺสาเทติ อภินนฺทติ, ตํ อารพฺภ ทสฺสเนน ปหาตพฺโพ ราโค อุปฺปชฺชติ.\csromanspacer{} ทิฏฺฐิ ฯเปฯ วิจิกิจฺฉา ฯเปฯ ทสฺสเนน ปหาตพฺพํ โทมนสฺสํ อุปฺปชฺชติ. (2)}

\csromanheader{2}{\textbf{อธิปติ}}
\proseitem{16}{ทสฺสเนน ปหาตพฺโพ ธมฺโม ทสฺสเนน ปหาตพฺพสฺส ธมฺมสฺส อธิปติปจฺจเยน ปจฺจโย.\csromanspacer{} อารมฺมณาธิปติ, สหชาตาธิปติ.\csromanspacer{}\csromanspacer{}\textbf{อารมฺมณาธิปติ\csromandash{}}ทสฺสเนน ปหาตพฺพํ ราคํ ครุํ กตฺวา อสฺสาเทติ อภินนฺทติ, ตํ ครุํ กตฺวา ทสฺสเนน ปหาตพฺโพ ราโค อุปฺปชฺชติ.\csromanspacer{} ทิฏฺฐิ อุปฺปชฺชติ.\csromanspacer{} ทิฏฺฐิํ ครุํ กตฺวา ฯเปฯ.\csromanspacer{} \textbf{สหชาตาธิปติ\csromandash{}}ทสฺสเนน ปหาตพฺพาธิปติ สมฺปยุตฺตกานํ ขนฺธานํ อธิปติปจฺจเยน ปจฺจโย. (1)}

\prose{ทสฺสเนน ปหาตพฺโพ ธมฺโม นทสฺสเนน ปหาตพฺพสฺส ธมฺมสฺส อธิปติปจฺจเยน ปจฺจโย.\csromanspacer{} \textbf{สหชาตาธิปติ\csromandash{}}ทสฺสเนน ปหาตพฺพาธิปติ จิตฺตสมุฏฺฐานานํ รูปานํ อธิปติปจฺจเยน ปจฺจโย. (2)}

\prose{ทสฺสเนน ปหาตพฺโพ ธมฺโม ทสฺสเนน ปหาตพฺพสฺส จ นทสฺสเนน ปหาตพฺพสฺส จ ธมฺมสฺส อธิปติปจฺจเยน ปจฺจโย.\csromanspacer{} \textbf{สหชาตาธิปติ\csromandash{}} ทสฺสเนน ปหาตพฺพาธิปติ สมฺปยุตฺตกานํ ขนฺธานํ จิตฺตสมุฏฺฐานานญฺจ รูปานํ อธิปติปจฺจเยน ปจฺจโย. (3)}

\proseitem{17}{นทสฺสเนน ปหาตพฺโพ ธมฺโม นทสฺสเนน ปหาตพฺพสฺส ธมฺมสฺส อธิปติปจฺจเยน ปจฺจโย.\csromanspacer{} อารมฺมณาธิปติ, สหชาตาธิปติ.\csromanspacer{}\csromanspacer{}\textbf{อารมฺมณาธิปติ\csromandash{}}ทานํ ฯเปฯ สีลํ ฯเปฯ อุโปสถกมฺมํ ฯเปฯ.\csromanspacer{} ปุพฺเพ ฯเปฯ.\csromanspacer{} ฌานา วุฏฺฐหิตฺวา ฌานํ ครุํ กตฺวา ปจฺจเวกฺขติ, อสฺสาเทติ อภินนฺทติ, ตํ ครุํ กตฺวา นทสฺสเนน ปหาตพฺโพ ราโค อุปฺปชฺชติ อริยา มคฺคา วุฏฺฐหิตฺวา มคฺคํ ครุํ กตฺวา ฯเปฯ ผลสฺส อธิปติปจฺจเยน ปจฺจโย.\csromanspacer{} จกฺขุํ ฯเปฯ.\csromanspacer{} วตฺถุํ นทสฺสเนน ปหาตพฺเพ ขนฺเธ ครุํ กตฺวา อสฺสาเทติ อภินนฺทติ, ตํ ครุํ กตฺวา นทสฺสเนน ปหาตพฺโพ ราโค อุปฺปชฺชติ.\csromanspacer{}\csromanspacer{}\textbf{สหชาตาธิปติ\csromandash{}}นทสฺสเนน ปหาตพฺพาธิปติ สมฺปยุตฺตกานํ ขนฺธานํ จิตฺตสมุฏฺฐานานญฺจ รูปานํ อธิปติปจฺจเยน ปจฺจโย. (1)}

\csromanheader{2}{83. ทสฺสเนนปหาตพฺพทุก}
\prose{\csromanfolio{45}นทสฺสเนน ปหาตพฺโพ ธมฺโม ทสฺสเนน ปหาตพฺพสฺส ธมฺมสฺส อธิปติปจฺจเยน ปจฺจโย.\csromanspacer{} \textbf{อารมฺมณาธิปติ\csromandash{}}ทานํ ทตฺวา ฯเปฯ ฌานํ ฯเปฯ.\csromanspacer{} จกฺขุํ ฯเปฯ.\csromanspacer{} วตฺถุํ นทสฺสเนน ปหาตพฺเพ ขนฺเธ ครุํ กตฺวา อสฺสาเทติ อภินนฺทติ, ตํ ครุํ กตฺวา ทสฺสเนน ปหาตพฺโพ ราโค อุปฺปชฺชติ.\csromanspacer{} ทิฏฺฐิ อุปฺปชฺชติ. (2)}

\csromanheader{2}{\textbf{อนนฺตร}}
\proseitem{18}{ทสฺสเนน ปหาตพฺโพ ธมฺโม ทสฺสเนน ปหาตพฺพสฺส ธมฺมสฺส อนนฺตรปจฺจเยน ปจฺจโย.\csromanspacer{} ปุริมา ปุริมา ทสฺสเนน ปหาตพฺพา ขนฺธา ปจฺฉิมานํ ปจฺฉิมานํ ทสฺสเนน ปหาตพฺพานํ ขนฺธานํ อนนฺตรปจฺจเยน ปจฺจโย. (1)}

\prose{ทสฺสเนน ปหาตพฺโพ ธมฺโม นทสฺสเนน ปหาตพฺพสฺส ธมฺมสฺส อนนฺตรปจฺจเยน ปจฺจโย.\csromanspacer{} ทสฺสเนน ปหาตพฺพา ขนฺธา วุฏฺฐานสฺส อนนฺตรปจฺจเยน ปจฺจโย. (มูลํ กาตพฺพํ.) ปุริมา ปุริมา นทสฺสเนน ปหาตพฺพา ขนฺธา ฯเปฯ ผลสมาปตฺติยา อนนฺตรปจฺจเยน ปจฺจโย. (มูลํ กาตพฺพํ.) อาวชฺชานา ทสฺสเนน ปหาตพฺพานํ ขนฺธานํ อนนฺตรปจฺจเยน ปจฺจโย. (2)}

\prose{สมนนฺตรปจฺจเยน ปจฺจโย.\csromanspacer{} สหชาตปจฺจเยน ปจฺจโย.\csromanspacer{} ปญฺจ.\csromanspacer{} อญฺญมญฺญปจฺจเยน ปจฺจโย.\csromanspacer{} เทฺว.\csromanspacer{} นิสฺสยปจฺจเยน ปจฺจโย.\csromanspacer{} สตฺต.}

\csromanheader{2}{\textbf{อุปนิสฺสย}}
\proseitem{19}{ทสฺสเนน ปหาตพฺโพ ธมฺโม ทสฺสเนน ปหาตพฺพสฺส ธมฺมสฺส อุปนิสฺสยปจฺจเยน ปจฺจโย.\csromanspacer{} \textbf{อารมฺมณูปนิสฺสโย, อนนฺตรูปนิสฺสโย}, \textbf{ปกตูปนิสฺสโย ฯเปฯ.}\csromanspacer{} ปกตูปนิสฺสโย\csromandash{}ทสฺสเนน ปหาตพฺพํ ราคํ อุปนิสฺสาย ปาณํ หนติ ฯเปฯ สํฆํ ภินฺทติ.\csromanspacer{} ทสฺสเนน ปหาตพฺพํ โทสํ.\csromanspacer{} โมหํ.\csromanspacer{} ทิฏฺฐิํ.\csromanspacer{} ปตฺถนํ อุปนิสฺสาย ปาณํ หนติ ฯเปฯ สํฆํ ภินฺทติ.\csromanspacer{} ทสฺสเนน ปหาตพฺโพ ราโค ฯเปฯ ปตฺถนา ทสฺสเนน ปหาตพฺพสฺส ราคสฺส ฯเปฯ ปตฺถนาย อุปนิสฺสยปจฺจเยน ปจฺจโย. (1)}

\prose{ทสฺสเนน ปหาตพฺโพ ธมฺโม นทสฺสเนน ปหาตพฺพสฺส ธมฺมสฺส อุปนิสฺสยปจฺจเยน ปจฺจโย.\csromanspacer{} \textbf{อนนฺตรูปนิสฺสโย, ปกตูปนิสฺสโย ฯเปฯ.}\csromanspacer{} ปกตูปนิสฺสโย\csromandash{}ทสฺสเนน ปหาตพฺพํ ราคํ อุปนิสฺสาย ทานํ เทติ ฯเปฯ สมาปตฺติํ อุปฺปาเทติ.\csromanspacer{} ทสฺสเนน ปหาตพฺพํ โทสํ.\csromanspacer{} โมหํ.\csromanspacer{} ทิฏฺฐิํ.\csromanspacer{} ปตฺถนํ \csromanfolio{46}อุปนิสฺสาย ทานํ เทติ ฯเปฯ สมาปตฺติํ อุปฺปาเทติ.\csromanspacer{} ทสฺสเนน ปหาตพฺโพ ราโค ฯเปฯ ปตฺถนา สทฺธาย ฯเปฯ ปญฺญาย.\csromanspacer{} นทสฺสเนน ปหาตพฺพสฺส ราคสฺส.\csromanspacer{} โทสสฺส.\csromanspacer{} โมหสฺส.\csromanspacer{} มานสฺส.\csromanspacer{} ปตฺถนาย.\csromanspacer{} กายิกสฺส สุขสฺส ฯเปฯ ผลสมาปตฺติยา อุปนิสฺสยปจฺจเยน ปจฺจโย. (2)}

\proseitem{20}{นทสฺสเนน ปหาตพฺโพ ธมฺโม นทสฺสเนน ปหาตพฺพสฺส ธมฺมสฺส อุปนิสฺสยปจฺจเยน ปจฺจโย.\csromanspacer{} \textbf{อารมฺมณูปนิสฺสโย, อนนฺตรูปนิสฺสโย,} \textbf{ปกตูปนิสฺสโย ฯเปฯ.}\csromanspacer{} ปกตูปนิสฺสโย\csromandash{}สทฺธํ อุปนิสฺสาย ทานํ เทติ ฯเปฯ สมาปตฺติํ อุปฺปาเทติ.\csromanspacer{} มานํ ชปฺเปติ.\csromanspacer{} สีลํ ฯเปฯ.\csromanspacer{} ปญฺญํ.\csromanspacer{} นทสฺสเนน ปหาตพฺพํ ราคํ.\csromanspacer{} โทสํ.\csromanspacer{} โมหํ.\csromanspacer{} มานํ.\csromanspacer{} ปตฺถนํ.\csromanspacer{} กายิกํ สุขํ ฯเปฯ.\csromanspacer{} เสนาสนํ อุปนิสฺสาย ทานํ เทติ ฯเปฯ สมาปตฺติํ อุปฺปาเทติ.\csromanspacer{} มานํ ชปฺเปติ.\csromanspacer{} สทฺธา ฯเปฯ.\csromanspacer{} ปญฺญา.\csromanspacer{} นทสฺสเนน ปหาตพฺโพ ราโค ฯเปฯ ปตฺตนา.\csromanspacer{} กายิกํ สุขํ ฯเปฯ.\csromanspacer{} เสนาสนํ สทฺธาย ฯเปฯ ปญฺญาย.\csromanspacer{} นทสฺสเนน ปหาตพฺพสฺส ราคสฺส ฯเปฯ ปตฺถนาย.\csromanspacer{} กายิกสฺส สุขสฺส.\csromanspacer{} กายิกสฺส ทุกฺขสฺส.\csromanspacer{} มคฺคสฺส, ผลสมาปตฺติยา อุปนิสฺสยปจฺจเยน ปจฺจโย. (1)}

\prose{นทสฺสเนน ปหาตพฺโพ ธมฺโม ทสฺสเนน ปหาตพฺพสฺส ธมฺมสฺส อุปนิสฺสยปจฺจเยน ปจฺจโย.\csromanspacer{} \textbf{อารมฺมณูปนิสฺสโย, อนนฺตรูปนิสฺสโย,} \textbf{ปกตูปนิสฺสโย ฯเปฯ.}\csromanspacer{} ปกตูปนิสฺสโย\csromandash{}สทฺธํ อุปนิสฺสาย ทิฏฺฐิํ คณฺหาติ.\csromanspacer{} สีลํ ฯเปฯ.\csromanspacer{} ปญฺญํ อุปนิสฺสาย ทิฏฺฐิํ คณฺหาติ.\csromanspacer{} นทสฺสเนน ปหาตพฺพํ ราคํ.\csromanspacer{} โทสํ.\csromanspacer{} โมหํ.\csromanspacer{} มานํ.\csromanspacer{} ปตฺถนํ.\csromanspacer{} กายิกํ สุขํ ฯเปฯ.\csromanspacer{} เสนาสนํ อุปนิสฺสาย ปาณํ หนติ ฯเปฯ สํฆํ ภินฺทติ.\csromanspacer{} สทฺธา ฯเปฯ.\csromanspacer{} เสนาสนํ ทสฺสเนน ปหาตพฺพสฺส ราคสฺส.\csromanspacer{} โทสสฺส.\csromanspacer{} โมหสฺส.\csromanspacer{} ทิฏฺฐิยา.\csromanspacer{} ปตฺถนาย อุปนิสฺสยปจฺจเยน ปจฺจโย. (2)}

\csromanheader{2}{\textbf{ปุเรชาต}}
\proseitem{21}{นทสฺสเนน ปหาตพฺโพ ธมฺโม นทสฺสเนน ปหาตพฺพสฺส ธมฺมสฺส ปุเรชาตปจฺจเยน ปจฺจโย.\csromanspacer{} อารมฺมณปุเรชาตํ, วตฺถุปุเรชาตํ.\csromanspacer{}\csromanspacer{}\textbf{อารมฺมณปุเรชาตํ\csromandash{}}จกฺขุํ ฯเปฯ.\csromanspacer{} วตฺถุํ อนิจฺจโต ฯเปฯ วิปสฺสติ, อสฺสาเทติ อภินนฺทติ, ตํ อารพฺภ นทสฺสเนน ปหาตพฺโพ ราโค ฯเปฯ อุทฺธจฺจํ ฯเปฯ.\csromanspacer{} นทสฺสเนน ปหาตพฺพํ โทมนสฺสํ อุปฺปชฺชติ.\csromanspacer{} ทิพฺเพน จกฺขุนา รูปํ ปสฺสติ.\csromanspacer{} ทิพฺพาย โสตธาตุยา สทฺทํ สุณาติ.\csromanspacer{} รูปายตนํ จกฺขุวิญฺญาณสฺส ฯเปฯ.\csromanspacer{} โผฏฺฐพฺพายตนํ กายวิญฺญาณสฺส ฯเปฯ.\csromanspacer{} \textbf{วตฺถุปุเรชาตํ\csromandash{}}จกฺขายตนํ}

\vspace{\tipitakaparskip}
\csromancenter{ทสฺสเนนปหาตพฺพทุก จกฺขุวิญฺญาณสฺส ฯเปฯ.\csromanspacer{} กายายตนํ กายวิญฺญาณสฺส ฯเปฯ.\csromanspacer{} วตฺถุ นทสฺสเนน ปหาตพฺพานํ ขนฺธานํ ปุเรชาตปจฺจเยน ปจฺจโย. (1)}
\prose{\csromanfolio{47}นทสฺสเนน ปหาตพฺโพ ธมฺโม ทสฺสเนน ปหาตพฺพสฺส ธมฺมสฺส ปุเรชาตปจฺจเยน ปจฺจโย.\csromanspacer{} อารมฺมณปุเรชาตํ, วตฺถุปุเรชาตํ.\csromanspacer{}\csromanspacer{}\textbf{อารมฺมณปุเรชาตํ\csromandash{}}จกฺขุํ ฯเปฯ.\csromanspacer{} วตฺถุํ อสฺสาเทติ อภินนฺทติ, ตํ อารพฺภ ทสฺสเนน ปหาตพฺโพ ราโค ฯเปฯ ทิฏฺฐิ ฯเปฯ วิจิกิจฺฉา ฯเปฯ.\csromanspacer{} ทสฺสเนน ปหาตพฺพํ โทมนสฺสํ อุปฺปชฺชติ.\csromanspacer{} \textbf{วตฺถุปุเรชาตํ\csromandash{}}วตฺถุ ทสฺสเนน ปหาตพฺพานํ ขนฺธานํ ปุเรชาตปจฺจเยน ปจฺจโย. (2)}

\csromanheader{2}{\textbf{ปจฺฉาชาตาเสวน}}
\proseitem{22}{ทสฺสเนน ปหาตพฺโพ ธมฺโม นทสฺสเนน ปหาตพฺพสฺส ธมฺมสฺส ปจฺฉาชาตปจฺจเยน ปจฺจโย. (สํขิตฺตํ.) (1)}

\prose{นทสฺสเนน ปหาตพฺโพ ธมฺโม นทสฺสเนน ปหาตพฺพสฺส ธมฺมสฺส ปจฺฉาชาตปจฺจเยน ปจฺจโย. (สํขิตฺตํ.) อาเสวนปจฺจเยน ปจฺจโย.\csromanspacer{} เทฺว.}

\csromanheader{2}{\textbf{กมฺม}}
\proseitem{23}{ทสฺสเนน ปหาตพฺโพ ธมฺโม ทสฺสเนน ปหาตพฺพสฺส ธมฺมสฺส กมฺมปจฺจเยน ปจฺจโย.\csromanspacer{} ทสฺสเนน ปหาตพฺพา เจตนา สมฺปยุตฺตกานํ ขนฺธานํ กมฺมปจฺจเยน ปจฺจโย. (1)}

\prose{ทสฺสเนน ปหาตพฺโพ ธมฺโม นทสฺสเนน ปหาตพฺพสฺส ธมฺมสฺส กมฺมปจฺจเยน ปจฺจโย.\csromanspacer{} \textbf{สหชาตา}, นานากฺขณิกา.\csromanspacer{}\csromanspacer{}\textbf{สหชาตา} ทสฺสเนน ปหาตพฺพา เจตนา จิตฺตสมุฏฺฐานานํ รูปานํ กมฺมปจฺจเยน ปจฺจโย.\csromanspacer{} \textbf{นานากฺขณิกา} ทสฺสเนน ปหาตพฺพา เจตนา วิปากานํ ขนฺธานํ กฏตฺตา จ รูปานํ กมฺมปจฺจเยน ปจฺจโย. (2)}

\prose{ทสฺสเนน ปหาตพฺโพ ธมฺโม ทสฺสเนน ปหาตพฺพสฺส จ นทสฺสเนน ปหาตพฺพสฺส จ ธมฺมสฺส กมฺมปจฺจเยน ปจฺจโย.\csromanspacer{} ทสฺสเนน ปหาตพฺพา เจตนา สมฺปยุตฺตกานํ ขนฺธานํ จิตฺตสมุฏฺฐานานญฺจ รูปานํ กมฺมปจฺจเยน ปจฺจโย. (3)}

\proseitem{24}{\csromanfolio{48}นทสฺสเนน ปหาตพฺโพ ธมฺโม นทสฺสเนน ปหาตพฺพสฺส ธมฺมสฺส กมฺมปจฺจเยน ปจฺจโย.\csromanspacer{} \textbf{สหชาตา}, นานากฺขณิกา.\csromanspacer{}\csromanspacer{}\textbf{สหชาตา} นทสฺสเนน ปหาตพฺพา เจตนา สมฺปยุตฺตกานํ ขนฺธานํ จิตฺตสมุฏฺฐานานญฺจ รูปานํ กมฺมปจฺจเยน ปจฺจโย.\csromanspacer{} ปฏิสนฺธิกฺขเณ ฯเปฯ.\csromanspacer{} \textbf{นานากฺขณิกา} นทสฺสเนน ปหาตพฺพา เจตนา วิปากานํ ขนฺธานํ กฏตฺตา จ รูปานํ กมฺมปจฺจเยน ปจฺจโย.}

\csromanheader{2}{\textbf{วิปากาทิ}}
\proseitem{25}{นทสฺสเนน ปหาตพฺโพ ธมฺโม นทสฺสเนน ปหาตพฺพสฺส ธมฺมสฺส วิปากปจฺจเยน ปจฺจโย.\csromanspacer{} เอกํ.\csromanspacer{} อาหารปจฺจเยน ปจฺจโย.\csromanspacer{} จตฺตาริ.\csromanspacer{} อินฺทฺริยปจฺจเยน ปจฺจโย.\csromanspacer{} จตฺตาริ.\csromanspacer{} ฌานปจฺจเยน ปจฺจโย.\csromanspacer{} จตฺตาริ.\csromanspacer{} มคฺคปจฺจเยน ปจฺจโย.\csromanspacer{} จตฺตาริ.\csromanspacer{} สมฺปยุตฺตปจฺจเยน ปจฺจโย.\csromanspacer{} เทฺว.}

\csromanheader{2}{\textbf{วิปฺปยุตฺต}}
\proseitem{26}{ทสฺสเนน ปหาตพฺโพ ธมฺโม นทสฺสเนน ปหาตพฺพสฺส ธมฺมสฺส วิปฺปยุตฺตปจฺจเยน ปจฺจโย.\csromanspacer{} สหชาตํ, ปจฺฉาชาตํ. (สํขิตฺตํ.) (1)}

\prose{นทสฺสเนน ปหาตพฺโพ ธมฺโม นทสฺสเนน ปหาตพฺพสฺส ธมฺมสฺส วิปฺปยุตฺตปจฺจเยน ปจฺจโย.\csromanspacer{} สหชาตํ, ปุเรชาตํ, ปจฺฉาชาตํ. (สํขิตฺตํ.)}

\prose{นทสฺสเนน ปหาตพฺโพ ธมฺโม ทสฺสเนน ปหาตพฺพสฺส ธมฺมสฺส วิปฺปยุตฺตปจฺจเยน ปจฺจโย.\csromanspacer{} \textbf{ปุเรชาตํ} วตฺถุ ทสฺสเนน ปหาตพฺพานํ ขนฺธานํ วิปฺปยุตฺตปจฺจเยน ปจฺจโย. (2)}

\csromanheader{2}{\textbf{อตฺถิ-อาทิ}}
\proseitem{27}{ทสฺสเนน ปหาตพฺโพ ธมฺโม ทสฺสเนน ปหาตพฺพสฺส ธมฺมสฺส อตฺถิปจฺจเยน ปจฺจโย. (ปฏิจฺจสทิสํ.) ทสฺสเนน ปหาตพฺโพ ธมฺโม นทสฺสเนน ปหาตพฺพสฺส ธมฺมสฺส อตฺถิปจฺจเยน ปจฺจโย.\csromanspacer{} สหชาตํ, ปจฺฉาชาตํ. (สํขิตฺตํ.) ทสฺสเนน ปหาตพฺโพ ธมฺโม ทสฺสเนน ปหาตพฺพสฺส จ นทสฺสเนน ปหาตพฺพสฺส จ ธมฺมสฺส อตฺถิปจฺจเยน ปจฺจโย. (ปฏิจฺจวารสทิสา.) (3)}

\csromanheader{2}{83. ทสฺสเนนปหาตพฺพทุก}
\prose{\csromanfolio{49}นทสฺสเนน ปหาตพฺโพ ธมฺโม นทสฺสเนน ปหาตพฺพสฺส ธมฺมสฺส อตฺถิปจฺจเยน ปจฺจโย.\csromanspacer{} \textbf{สหชาตํ,} ปุเรชาตํ, ปจฺฉาชาตํ, อาหารํ, อินฺทฺริยํ. (สํขิตฺตํ.) นทสฺสเนน ปหาตพฺโพ ธมฺโม ทสฺสเนน ปหาตพฺพสฺส ธมฺมสฺส อตฺถิปจฺจเยน ปจฺจโย.\csromanspacer{} ปุเรชาตํ\csromandash{}จกฺขุํ ฯเปฯ. (สํขิตฺตํ.\csromanspacer{} ปุเรชาตสทิสํ.) (2)}

\proseitem{28}{ทสฺสเนน ปหาตพฺโพ จ นทสฺสเนน ปหาตพฺโพ จ ธมฺมา ทสฺสเนน ปหาตพฺพสฺส ธมฺมสฺส อตฺถิปจฺจเยน ปจฺจโย.\csromanspacer{} \textbf{สหชาตํ,} ปุเรชาตํ.\csromanspacer{}\csromanspacer{}สหชาโต ทสฺสเนน ปหาตพฺโพ เอโก ขนฺโธ จ วตฺถุ จ ติณฺณนฺนํ ขนฺธานํ อตฺถิปจฺจเยน ปจฺจโย ฯเปฯ เทฺว ขนฺธา จ ฯเปฯ. (1)}

\prose{ทสฺสเนน ปหาตพฺโพ จ นทสฺสเนน ปหาตพฺโพ จ ธมฺมา นทสฺสเนน ปหาตพฺพสฺส ธมฺมสฺส อตฺถิปจฺจเยน ปจฺจโย.\csromanspacer{} \textbf{สหชาตํ,} ปจฺฉาชาตํ, อาหารํ, อินฺทฺริยํ.\csromanspacer{}\csromanspacer{}\textbf{สหชาตา} ทสฺสเนน ปหาตพฺพา ขนฺธา จ มหาภูตา จ จิตฺตสมุฏฺฐานานํ รูปานํ อตฺถิปจฺจเยน ปจฺจโย.\csromanspacer{} \textbf{ปจฺฉาชาตา} ทสฺสเนน ปหาตพฺพา ขนฺธา จ กพฬีกาโร อาหาโร จ อิมสฺส กายสฺส อตฺถิปจฺจเยน ปจฺจโย.\csromanspacer{} \textbf{ปจฺฉาชาตา} ทสฺสเนน ปหาตพฺพา ขนฺธา จ รูปชีวิตินฺทฺริยญฺจ กฏตฺตารูปานํ อตฺถิปจฺจเยน ปจฺจโย. (2)}

\prose{นตฺถิปจฺจเยน ปจฺจโย.\csromanspacer{} วิคตปจฺจเยน ปจฺจโย.\csromanspacer{} อวิคตปจฺจเยน ปจฺจโย.}

\csromanheader{2}{1. ปจฺจยานุโลม 2. สงฺขฺยาวาร}
\csromanheader{2}{\textbf{สุทฺธ}}
\proseitem{29}{เหตุยา จตฺตาริ, อารมฺมเณ จตฺตาริ, อธิปติยา ปญฺจ, อนนฺตเร จตฺตาริ, สมนนฺตเร จตฺตาริ, สหชาเต ปญฺจ, อญฺญมญฺเญ เทฺว, นิสฺสเย สตฺต, อุปนิสฺสเย จตฺตาริ, ปุเรชาเต เทฺว, ปจฺฉาชาเต เทฺว, อาเสวเน เทฺว, กมฺเม จตฺตาริ, วิปาเก เอกํ, อาหาเร จตฺตาริ, อินฺทฺริเย จตฺตาริ, ฌาเน จตฺตาริ, มคฺเค จตฺตาริ, สมฺปยุตฺเต เทฺว, วิปฺปยุตฺเต ตีณิ, อตฺถิยา สตฺต, นตฺถิยา จตฺตาริ, วิคเต จตฺตาริ, อวิคเต สตฺต.}

\csromanheader{2}{2. ปจฺจนียุทฺธาร}
\proseitem{30}{\csromanfolio{50}ทสฺสเนน ปหาตพฺโพ ธมฺโม ทสฺสเนน ปหาตพฺพสฺส ธมฺมสฺส อารมฺมณปจฺจเยน ปจฺจโย.\csromanspacer{} สหชาตปจฺจเยน ปจฺจโย.\csromanspacer{} อุปนิสฺสยปจฺจเยน ปจฺจโย. (1)}

\prose{ทสฺสเนน ปหาตพฺโพ ธมฺโม นทสฺสเนน ปหาตพฺพสฺส ธมฺมสฺส อารมฺมณปจฺจเยน ปจฺจโย.\csromanspacer{} สหชาตปจฺจเยน ปจฺจโย.\csromanspacer{} อุปนิสฺสยปจฺจเยน ปจฺจโย.\csromanspacer{} ปจฺฉาชาตปจฺจเยน ปจฺจโย.\csromanspacer{} กมฺมปจฺจเยน ปจฺจโย. (2)}

\prose{ทสฺสเนน ปหาตพฺโพ ธมฺโม ทสฺสเนน ปหาตพฺพสฺส จ นทสฺสเนน ปหาตพฺพสฺส จ ธมฺมสฺส สหชาตปจฺจเยน ปจฺจโย. (3)}

\proseitem{31}{นทสฺสเนน ปหาตพฺโพ ธมฺโม นทสฺสเนน ปหาตพฺพสฺส ธมฺมสฺส อารมฺมณปจฺจเยน ปจฺจโย.\csromanspacer{} สหชาตปจฺจเยน ปจฺจโย.\csromanspacer{} อุปนิสฺสยปจฺจเยน ปจฺจโย.\csromanspacer{} ปุเรชาตปจฺจเยน ปจฺจโย.\csromanspacer{} ปจฺฉาชาตปจฺจเยน ปจฺจโย.\csromanspacer{} กมฺมปจฺจเยน ปจฺจโย.\csromanspacer{} อาหารปจฺจเยน ปจฺจโย.\csromanspacer{} อินฺทฺริยปจฺจเยน ปจฺจโย. (1)}

\prose{นทสฺสเนน ปหาตพฺโพ ธมฺโม ทสฺสเนน ปหาตพฺพสฺส ธมฺมสฺส อารมฺมณปจฺจเยน ปจฺจโย.\csromanspacer{} อุปนิสฺสยปจฺจเยน ปจฺจโย.\csromanspacer{} ปุเรชาตปจฺจเยน ปจฺจโย. (2)}

\prose{ทสฺสเนน ปหาตพฺโพ จ นทสฺสเนน ปหาตพฺโพ จ ธมฺมา ทสฺสเนน ปหาตพฺพสฺส ธมฺมสฺส สหชาตํ, ปุเรชาตํ. (1)}

\prose{ทสฺสเนน ปหาตพฺโพ จ นทสฺสเนน ปหาตพฺโพ จ ธมฺมา นทสฺสเนน ปหาตพฺพสฺส ธมฺมสฺส สหชาตํ, ปจฺฉาชาตํ, อาหารํ, อินฺทฺริยํ. (2)}

\csromanheader{2}{2. ปจฺจยปจฺจนีย 2. สงฺขฺยาวาร}
\csromanheader{2}{\textbf{สุทฺธ}}
\proseitem{32}{นเหตุยา สตฺต, น-อารมฺมเณ สตฺต, น-อธิปติยา สตฺต, น-อนนฺตเร สตฺต, นสมนนฺตเร สตฺต, นสหชาเต ปญฺจ, น-อญฺญมญฺเญ ปญฺจ, นนิสฺสเย ปญฺจ, น-อุปนิสฺสเย สตฺต, นปุเรชาเต ฉ, นปจฺฉาชาเต สตฺต ฯเปฯ นมคฺเค สตฺต, นสมฺปยุตฺเต ปญฺจ, นวิปฺปยุตฺเต จตฺตาริ, โน-อตฺถิยา จตฺตาริ, โนนตฺถิยา สตฺต, โนวิคเต สตฺต, โน-อวิคเต จตฺตาริ.}

\csromanmatikaanchor{mtk.13}
\csromantocmarkref{section}{84. ภาวนายปหาตพฺพทุก}{mtk.13}
\bookmark[dest={mtk.13},level=1]{84. ภาวนายปหาตพฺพทุก}
\csromanheader{1}{84. ภาวนายปหาตพฺพทุก \textbf{3. ปจฺจยานุโลมปจฺจนีย}}
\proseitem{33}{\csromanfolio{51}เหตุปจฺจยา น-อารมฺมเณ จตฺตาริ, น-อธิปติยา นว, น-อนนฺตเร จตฺตาริ, นสมนนฺตเร จตฺตาริ, น-อญฺญมญฺเญ เทฺว, น-อุปนิสฺสเย จตฺตาริ, (สพฺพตฺถ จตฺตาริ.) นสมฺปยุตฺเต เทฺว, นวิปฺปยุตฺเต เทฺว, โนนตฺถิยา จตฺตาริ, โนวิคเต จตฺตาริ.}

\csromanheader{2}{4. ปจฺจยปจฺจนียานุโลม}
\proseitem{34}{นเหตุปจฺจยา อารมฺมเณ จตฺตาริ, อธิปติยา ปญฺจ, (อนุโลมมาติกา กาตพฺพา.) ฯเปฯ อวิคเต สตฺต.}

\nitthitam{ทสฺสเนน ปหาตพฺพทุกํ นิฏฺฐิตํ.}
\csromanheader{2}{84. ภาวนายปหาตพฺพทุก 1. ปฏิจฺจวาร}
\csromanheader{2}{1. ปจฺจยานุโลมาทิ}
\proseitem{35}{ภาวนาย ปหาตพฺพํ ธมฺมํ ปฏิจฺจ ภาวนาย ปหาตพฺโพ ธมฺโม อุปฺปชฺชติ เหตุปจฺจยา. (ยถา ทสฺสนทุกํ, เอวํ วิตฺถาเรตพฺพํ, นินฺนานากรณํ.)}

\prose{เหตุยา ปญฺจ ฯเปฯ อวิคเต ปญฺจ.\csromanspacer{} อนุโลมํ.}

\prose{ภาวนาย ปหาตพฺพํ ธมฺมํ ปฏิจฺจ ภาวนาย ปหาตพฺโพ ธมฺโม อุปฺปชฺชติ นเหตุปจฺจยา.\csromanspacer{} อุทฺธจฺจสหคเต ขนฺเธ ปฏิจฺจ อุทฺธจฺจสหคโต โมโห.}

\prose{นภาวนาย ปหาตพฺพํ ธมฺมํ ปฏิจฺจ นภาวนาย ปหาตพฺโพ ธมฺโม อุปฺปชฺชติ นเหตุปจฺจยา.\csromanspacer{} อเหตุกํ นภาวนาย ปหาตพฺพํ เอกํ ขนฺธํ ฯเปฯ. (ยาว อสญฺญสตฺตา.) วิจิกิจฺฉาสหคเต ขนฺเธ ปฏิจฺจ วิจิกิจฺฉาสหคโต โมโห. (สํขิตฺตํ.)}

\prose{นเหตุยา เทฺว ฯเปฯ โนวิคเต ตีณิ.}

\vspace{\tipitakaparskip}
\csromancenter{ปจฺจนียํ.}
\prose{\csromanfolio{52}(ปจฺจยวารปจฺจนีเย นเหตุปจฺจเย อุทฺธจฺจสหคเต ตีณิ, โมโห อุทฺธริตพฺโพ, สพฺเพปิ วารา ทสฺสนทุกสทิสา, อุทฺธจฺจปจฺจนียมฺปิ นานํ.)}

\csromanheader{2}{84. ภาวนายปหาตพฺพทุก 7. ปญฺหาวาร}
\proseitem{36}{ภาวนาย ปหาตพฺโพ ธมฺโม ภาวนาย ปหาตพฺพสฺส ธมฺมสฺส เหตุปจฺจเยน ปจฺจโย.\csromanspacer{} ภาวนาย ปหาตพฺพา เหตู สมฺปยุตฺตกานํ ขนฺธานํ เหตุปจฺจเยน ปจฺจโย.\csromanspacer{} ตีณิ.}

\prose{นภาวนาย ปหาตพฺโพ ธมฺโม นภาวนาย ปหาตพฺพสฺส ธมฺมสฺส เหตุปจฺจเยน ปจฺจโย. (สํขิตฺตํ.) (1)}

\csromanheader{2}{\textbf{อารมฺมณ}}
\proseitem{37}{ภาวนาย ปหาตพฺโพ ธมฺโม ภาวนาย ปหาตพฺพสฺส ธมฺมสฺส อารมฺมณปจฺจเยน ปจฺจโย.\csromanspacer{} ภาวนาย ปหาตพฺพํ ราคํ อสฺสาเทติ อภินนฺทติ, ตํ อารพฺภ ภาวนาย ปหาตพฺโพ ราโค อุปฺปชฺชติ.\csromanspacer{} อุทฺธจฺจํ อุปฺปชฺชติ.\csromanspacer{} ภาวนาย ปหาตพฺพํ โทมนสฺสํ อุปฺปชฺชติ.\csromanspacer{} อุทฺธจฺจํ อารพฺภ อุทฺธจฺจํ อุปฺปชฺชติ.\csromanspacer{} ภาวนาย ปหาตพฺพํ โทมนสฺสํ อุปฺปชฺชติ.\csromanspacer{} ภาวนาย ปหาตพฺพํ โทมนสฺสํ อารพฺภ ภาวนาย ปหาตพฺพํ โทมนสฺสํ อุปฺปชฺชติ.\csromanspacer{} อุทฺธจฺจํ อุปฺปชฺชติ. (1)}

\prose{ภาวนาย ปหาตพฺโพ ธมฺโม นภาวนาย ปหาตพฺพสฺส ธมฺมสฺส อารมฺมณปจฺจเยน ปจฺจโย.\csromanspacer{} อริยา ภาวนาย ปหาตพฺเพ ปหีเน กิเลเส ฯเปฯ.\csromanspacer{} วิกฺขมฺภิเต กิเลเส ฯเปฯ.\csromanspacer{} ปุพฺเพ สมุทาจิณฺเณ ฯเปฯ.\csromanspacer{} ภาวนาย ปหาตพฺเพ ขนฺเธ อนิจฺจโต ฯเปฯ วิปสฺสติ, อสฺสาเทติ อภินนฺทติ, ตํ อารพฺภ นภาวนาย ปหาตพฺโพ ราโค อุปฺปชฺชติ.\csromanspacer{} ทิฏฺฐิ ฯเปฯ วิจิกิจฺฉา ฯเปฯ.\csromanspacer{} นภาวนาย ปหาตพฺพํ โทมนสฺสํ อุปฺปชฺชติ.\csromanspacer{} เจโตปริยญาเณน ภาวนาย ปหาตพฺพจิตฺตสมงฺคิสฺส จิตฺตํ ชานาติ.\csromanspacer{} ภาวนาย ปหาตพฺพา ขนฺธา เจโตปริยญาณสฺส, ปุพฺเพ นิวาสานุสฺสติญาณสฺส,}

\vspace{\tipitakaparskip}
\csromancenter{ภาวนายปหาตพฺพทุก ยถากมฺมูปคญาณสฺส, อนาคตํสญาณสฺส, อาวชฺชนาย อารมฺมณปจฺจเยน ปจฺจโย. (2)}
\proseitem{38}{\csromanfolio{53}นภาวนาย ปหาตพฺโพ ธมฺโม นภาวนาย ปหาตพฺพสฺส ธมฺมสฺส อารมฺมณปจฺจเยน ปจฺจโย.\csromanspacer{} ทานํ ฯเปฯ สีลํ ฯเปฯ อุโปสถกมฺมํ กตฺวา ตํ ปจฺจเวกฺขติ, อสฺสาเทติ อภินนฺทติ, ตํ อารพฺภ นภาวนาย ปหาตพฺโพ ราโค อุปฺปชฺชติ.\csromanspacer{} ทิฏฺฐิ ฯเปฯ วิจิกิจฺฉา ฯเปฯ.\csromanspacer{} นภาวนาย ปหาตพฺพํ โทมนสฺสํ อุปฺปชฺชติ.\csromanspacer{} ปุพฺเพ ฯเปฯ.\csromanspacer{} ฌานา ฯเปฯ.\csromanspacer{} อริยา มคฺคา วุฏฺฐหิตฺวา ฯเปฯ ผลสฺส, อาวชฺชนาย อารมฺมณปจฺจเยน ปจฺจโย.\csromanspacer{} อริยา นภาวนาย ปหาตพฺเพ ปหีเน กิเลเส ฯเปฯ.\csromanspacer{} จกฺขุํ ฯเปฯ.\csromanspacer{} วตฺถุํ นภาวนาย ปหาตพฺเพ ขนฺเธ อนิจฺจโต ฯเปฯ วิปสฺสติ, อสฺสาเทติ อภินนฺทติ, ตํ อารพฺภ นภาวนาย ปหาตพฺโพ ราโค อุปฺปชฺชติ.\csromanspacer{} ทิฏฺฐิ ฯเปฯ วิจิกิจฺฉา ฯเปฯ.\csromanspacer{} นภาวนาย ปหาตพฺพํ โทมนสฺสํ อุปฺปชฺชติ.\csromanspacer{} ทิพฺเพน จกฺขุนา รูปํ ปสฺสติ ฯเปฯ ยถากมฺมูปคญาณสฺส, อนาคตํสญาณสฺส, อาวชฺชนาย อารมฺมณปจฺจเยน ปจฺจโย. (1)}

\prose{นภาวนาย ปหาตพฺโพ ธมฺโม ภาวนาย ปหาตพฺพสฺส ธมฺมสฺส อารมฺมณปจฺจเยน ปจฺจโย.\csromanspacer{} ทานํ ฯเปฯ สีลํ ฯเปฯ ฌานํ ฯเปฯ.\csromanspacer{} จกฺขุํ ฯเปฯ.\csromanspacer{} วตฺถุํ นภาวนาย ปหาตพฺเพ ขนฺเธ อสฺสาเทติ อภินนฺทติ, ตํ อารพฺภ ภาวนาย ปหาตพฺโพ ราโค อุปฺปชฺชติ.\csromanspacer{} อุทฺธจฺจํ อุปฺปชฺชติ.\csromanspacer{} ภาวนาย ปหาตพฺพํ โทมนสฺสํ อุปฺปชฺชติ. (2)}

\csromanheader{2}{\textbf{อธิปติ}}
\proseitem{39}{ภาวนาย ปหาตพฺโพ ธมฺโม ภาวนาย ปหาตพฺพสฺส ธมฺมสฺส อธิปติปจฺจเยน ปจฺจโย.\csromanspacer{} อารมฺมณาธิปติ, สหชาตาธิปติ.\csromanspacer{}\csromanspacer{}\textbf{อารมฺมณาธิปติ\csromandash{}}ภาวนาย ปหาตพฺพํ ราคํ ครุํ กตฺวา อสฺสาเทติ อภินนฺทติ, ตํ ครุํ กตฺวา ภาวนาย ปหาตพฺโพ ราโค อุปฺปชฺชติ.\csromanspacer{} \textbf{สหชาตาธิปติ\csromandash{}}ภาวนาย ปหาตพฺพาธิปติ สมฺปยุตฺตกานํ ขนฺธานํ อธิปติปจฺจเยน ปจฺจโย. (1)}

\prose{ภาวนาย ปหาตพฺโพ ธมฺโม นภาวนาย ปหาตพฺพสฺส ธมฺมสฺส อธิปติปจฺจเยน ปจฺจโย.\csromanspacer{} อารมฺมณาธิปติ, สหชาตาธิปติ.\csromanspacer{}\csromanspacer{}\textbf{อารมฺมณาธิปติ\csromandash{}}ภาวนาย ปหาตพฺพํ ราคํ ครุํ กตฺวา อสฺสาเทติ อภินนฺทติ, ตํ ครุํ กตฺวา นภาวนาย ปหาตพฺโพ ราโค อุปฺปชฺชติ. \csromanfolio{54}ทิฏฺฐิ อุปฺปชฺชติ.\csromanspacer{} \textbf{สหชาตาธิปติ\csromandash{}}ภาวนาย ปหาตพฺพาธิปติ จิตฺตสมุฏฺฐานานํ รูปานํ อธิปติปจฺจเยน ปจฺจโย. (2)}

\prose{ภาวนาย ปหาตพฺโพ ธมฺโม ภาวนาย ปหาตพฺพสฺส จ นภาวนาย ปหาตพฺพสฺส จ ธมฺมสฺส อธิปติปจฺจเยน ปจฺจโย.\csromanspacer{} \textbf{สหชาตาธิปติ\csromandash{}}ภาวนาย ปหาตพฺพาธิปติ สมฺปยุตฺตกานํ ขนฺธานํ จิตฺตสมุฏฺฐานานญฺจ รูปานํ อธิปติปจฺจเยน ปจฺจโย. (3)}

\proseitem{40}{นภาวนาย ปหาตพฺโพ ธมฺโม นภาวนาย ปหาตพฺพสฺส ธมฺมสฺส อธิปติปจฺจเยน ปจฺจโย.\csromanspacer{} อารมฺมณาธิปติ, สหชาตาธิปติ.\csromanspacer{}\csromanspacer{}\textbf{อารมฺมณาธิปติ\csromandash{}}ทานํ ทตฺวา, สีลํ ฯเปฯ อุโปสถกมฺมํ กตฺวา ตํ ครุํ กตฺวา ปจฺจเวกฺขติ, อสฺสาเทติ อภินนฺทติ, ตํ ครุํ กตฺวา นภาวนาย ปหาตพฺโพ ราโค อุปฺปชฺชติ.\csromanspacer{} \textbf{สหชาตาธิปติ\csromandash{}}นภาวนาย ปหาตพฺพาธิปติ สมฺปยุตฺตกานํ ขนฺธานํ จิตฺตสมุฏฺฐานานญฺจ รูปานํ อธิปติปจฺจเยน ปจฺจโย. (1)}

\prose{นภาวนาย ปหาตพฺโพ ธมฺโม ภาวนาย ปหาตพฺพสฺส ธมฺมสฺส อธิปติปจฺจเยน ปจฺจโย.\csromanspacer{} \textbf{อารมฺมณาธิปติ\csromandash{}}ทานํ ฯเปฯ ฌานํ ฯเปฯ.\csromanspacer{} จกฺขุํ ฯเปฯ.\csromanspacer{} วตฺถุํ นภาวนาย ปหาตพฺเพ ขนฺเธ ครุํ กตฺวา อสฺสาเทติ อภินนฺทติ, ตํ ครุํ กตฺวา ภาวนาย ปหาตพฺโพ ราโค อุปฺปชฺชติ. (2)}

\csromanheader{2}{\textbf{อนนฺตราทิ}}
\proseitem{41}{ภาวนาย ปหาตพฺโพ ธมฺโม ภาวนาย ปหาตพฺพสฺส ธมฺมสฺส อนนฺตรปจฺจเยน ปจฺจโย.\csromanspacer{} จตฺตาริ. (ทสฺสนทุกสทิสา ภาวนา นินฺนานากรณา.) สมนนฺตรปจฺจเยน ปจฺจโย.\csromanspacer{} จตฺตาริ.\csromanspacer{} สหชาตปจฺจเยน ปจฺจโย.\csromanspacer{} ปญฺจ.\csromanspacer{} อญฺญมญฺญปจฺจเยน ปจฺจโย.\csromanspacer{} เทฺว.\csromanspacer{} นิสฺสยปจฺจเยน ปจฺจโย.\csromanspacer{} สตฺต.}

\csromanheader{2}{\textbf{อุปนิสฺสย}}
\proseitem{42}{ภาวนาย ปหาตพฺโพ ธมฺโม ภาวนาย ปหาตพฺพสฺส ธมฺมสฺส อุปนิสฺสยปจฺจเยน ปจฺจโย.\csromanspacer{} อารมฺมณูปนิสฺสโย, อนนฺตรูปนิสฺสโย, ปกตูปนิสฺสโย ฯเปฯ.\csromanspacer{} ปกตูปนิสฺสโย\csromandash{}ภาวนาย ปหาตพฺโพ ราโค.\csromanspacer{} โทโส.\csromanspacer{} โมโห.\csromanspacer{} มาโน.\csromanspacer{} ปตฺถนา ภาวนาย ปหาตพฺพสฺส}

\vspace{\tipitakaparskip}
\csromancenter{ภาวนายปหาตพฺพทุก ราคสฺส.\csromanspacer{} โทสสฺส.\csromanspacer{} โมหสฺส.\csromanspacer{} มานสฺส.\csromanspacer{} ปตฺถนาย อุปนิสฺสยปจฺจเยน ปจฺจโย. (1)}
\prose{\csromanfolio{55}ภาวนาย ปหาตพฺโพ ธมฺโม นภาวนาย ปหาตพฺพสฺส ธมฺมสฺส อุปนิสฺสยปจฺจเยน ปจฺจโย.\csromanspacer{} อารมฺมณูปนิสฺสโย, อนนฺตรูปนิสฺสโย, ปกตูปนิสฺสโย ฯเปฯ.\csromanspacer{} ปกตูปนิสฺสโย\csromandash{}ภาวนาย ปหาตพฺพํ ราคํ อุปนิสฺสาย ทานํ เทติ ฯเปฯ สมาปตฺติํ อุปฺปาเทติ.\csromanspacer{} ปาณํ หนติ ฯเปฯ สํฆํ ภินฺทติ.\csromanspacer{} ภาวนาย ปหาตพฺพํ โทสํ.\csromanspacer{} โมหํ.\csromanspacer{} มานํ.\csromanspacer{} ปตฺถนํ อุปนิสฺสาย ทานํ เทติ ฯเปฯ สมาปตฺติํ อุปฺปาเทติ.\csromanspacer{} ปาณํ หนติ ฯเปฯ สํฆํ ภินฺทติ.\csromanspacer{} ภาวนาย ปหาตพฺโพ ราโค ฯเปฯ ปตฺถนา สทฺธาย ฯเปฯ ปญฺญาย.\csromanspacer{} นภาวนาย ปหาตพฺพสฺส ราคสฺส.\csromanspacer{} โทสสฺส.\csromanspacer{} โมหสฺส.\csromanspacer{} ทิฏฺฐิยา.\csromanspacer{} ปตฺถนาย.\csromanspacer{} กายิกสฺส สุขสฺส.\csromanspacer{} กายิกสฺส ทุกฺขสฺส.\csromanspacer{} มคฺคสฺส, ผลสมาปตฺติยา อุปนิสฺสยปจฺจเยน ปจฺจโย. (2)}

\proseitem{43}{นภาวนาย ปหาตพฺโพ ธมฺโม นภาวนาย ปหาตพฺพสฺส ธมฺมสฺส อุปนิสฺสยปจฺจเยน ปจฺจโย.\csromanspacer{} อารมฺมณูปนิสฺสโย, อนนฺตรูปนิสฺสโย, ปกตูปนิสฺสโย ฯเปฯ.\csromanspacer{} ปกตูปนิสฺสโย\csromandash{}สทฺธํ อุปนิสฺสาย ทานํ เทติ ฯเปฯ สมาปตฺติํ อุปฺปาเทติ.\csromanspacer{} ทิฏฺฐิํ คณฺหาติ.\csromanspacer{} สีลํ ฯเปฯ.\csromanspacer{} ปญฺญํ.\csromanspacer{} นภาวนาย ปหาตพฺพํ ราคํ.\csromanspacer{} โทสํ.\csromanspacer{} โมหํ.\csromanspacer{} ทิฏฺฐิํ.\csromanspacer{} ปตฺถนํ.\csromanspacer{} กายิกํ สุขํ.\csromanspacer{} กายิกํ ทุกฺขํ ฯเปฯ.\csromanspacer{} เสนาสนํ อุปนิสฺสาย ทานํ เทติ ฯเปฯ สมาปตฺติํ อุปฺปาเทติ.\csromanspacer{} ปาณํ หนติ ฯเปฯ สํฆํ ภินฺทติ.\csromanspacer{} สทฺธา ฯเปฯ.\csromanspacer{} เสนาสนํ สทฺธาย ฯเปฯ ปญฺญาย.\csromanspacer{} นภาวนาย ปหาตพฺพสฺส ราคสฺส.\csromanspacer{} โทสสฺส.\csromanspacer{} โมหสฺส.\csromanspacer{} ทิฏฺฐิยา.\csromanspacer{} ปตฺถนาย.\csromanspacer{} กายิกสฺส สุขสฺส.\csromanspacer{} กายิกสฺส ทุกฺขสฺส.\csromanspacer{} มคฺคสฺส, ผลสมาปตฺติยา อุปนิสฺสยปจฺจเยน ปจฺจโย. (1)}

\prose{นภาวนาย ปหาตพฺโพ ธมฺโม ภาวนาย ปหาตพฺพสฺส ธมฺมสฺส อุปนิสฺสยปจฺจเยน ปจฺจโย.\csromanspacer{} อารมฺมณูปนิสฺสโย, อนนฺตรูปนิสฺสโย, ปกตูปนิสฺสโย ฯเปฯ.\csromanspacer{} ปกตูปนิสฺสโย\csromandash{}สทฺธํ อุปนิสฺสาย มานํ ชปฺเปติ ฯเปฯ.\csromanspacer{} สีลํ ฯเปฯ.\csromanspacer{} ปญฺญํ.\csromanspacer{} ราคํ ฯเปฯ.\csromanspacer{} กายิกํ สุขํ.\csromanspacer{} กายิกํ ทุกฺขํ.\csromanspacer{} เสนาสนํ อุปนิสฺสาย มานํ ชปฺเปติ.\csromanspacer{} สทฺธา ฯเปฯ.\csromanspacer{} เสนาสนํ ภาวนาย ปหาตพฺพสฺส ราคสฺส.\csromanspacer{} โทสสฺส.\csromanspacer{} โมหสฺส.\csromanspacer{} มานสฺส.\csromanspacer{} ปตฺถนาย อุปนิสฺสยปจฺจเยน ปจฺจโย. (2)}

\csromanheader{2}{\textbf{ปุเรชาตาทิ}}
\proseitem{44}{\csromanfolio{56}นภาวนาย ปหาตพฺโพ ธมฺโม นภาวนาย ปหาตพฺพสฺส ธมฺมสฺส ปุเรชาตปจฺจเยน ปจฺจโย.\csromanspacer{} อารมฺมณปุเรชาตํ, วตฺถุปุเรชาตํ.\csromanspacer{}\csromanspacer{}\textbf{อารมฺมณปุเรชาตํ\csromandash{}}จกฺขุํ ฯเปฯ.\csromanspacer{} วตฺถุํ อนิจฺจโต ฯเปฯ วิปสฺสติ, อสฺสาเทติ อภินนฺทติ, ตํ อารพฺภ นภาวนาย ปหาตพฺโพ ราโค อุปฺปชฺชติ.\csromanspacer{} ทิฏฺฐิ อุปฺปชฺชติ.\csromanspacer{} วิจิกิจฺฉา อุปฺปชฺชติ.\csromanspacer{} นภาวนาย ปหาตพฺพํ โทมนสฺสํ อุปฺปชฺชติ.\csromanspacer{} ทิพฺเพน จกฺขุนา รูปํ ปสฺสติ ฯเปฯ.\csromanspacer{} โผฏฺฐพฺพายตนํ กายวิญฺญาณสฺส ฯเปฯ.\csromanspacer{} \textbf{วตฺถุปุเรชาตํ\csromandash{}}จกฺขายตนํ จกฺขุวิญฺญาณสฺส ฯเปฯ กายายตนํ กายวิญฺญาณสฺส ฯเปฯ.\csromanspacer{} วตฺถุ นภาวนาย ปหาตพฺพานํ ขนฺธานํ ปุเรชาตปจฺจเยน ปจฺจโย. (1)}

\prose{นภาวนาย ปหาตพฺโพ ธมฺโม ภาวนาย ปหาตพฺพสฺส ธมฺมสฺส ปุเรชาตปจฺจเยน ปจฺจโย.\csromanspacer{} อารมฺมณปุเรชาตํ, วตฺถุปุเรชาตํ.\csromanspacer{}\csromanspacer{}\textbf{อารมฺมณปุเรชาตํ\csromandash{}}จกฺขุํ ฯเปฯ.\csromanspacer{} วตฺถุํ อสฺสาเทติ อภินนฺทติ, ตํ อารพฺภ ภาวนาย ปหาตพฺโพ ราโค อุปฺปชฺชติ.\csromanspacer{} อุทฺธจฺจํ อุปฺปชฺชติ.\csromanspacer{} ภาวนาย ปหาตพฺพํ โทมนสฺสํ อุปฺปชฺชติ.\csromanspacer{} \textbf{วตฺถุปุเรชาตํ\csromandash{}}วตฺถุ ภาวนาย ปหาตพฺพานํ ขนฺธานํ ปุเรชาตปจฺจเยน ปจฺจโย. (2)}

\proseitem{45}{ปจฺฉาชาตปจฺจเยน ปจฺจโย.\csromanspacer{} เทฺว.\csromanspacer{} อาเสวนปจฺจเยน ปจฺจโย.\csromanspacer{} เทฺว.\csromanspacer{} กมฺมปจฺจเยน ปจฺจโย.\csromanspacer{} ภาวนาย ปหาตพฺพา เจตนา สมฺปยุตฺตกานํ ขนฺธานํ กมฺมปจฺจเยน ปจฺจโย. (มูลํ กาตพฺพํ.) ภาวนาย ปหาตพฺพา เจตนา จิตฺตสมุฏฺฐานานํ รูปานํ กมฺมปจฺจเยน ปจฺจโย. (มูลํ กาตพฺพํ.) ภาวนาย ปหาตพฺพา เจตนา สมฺปยุตฺตกานํ ขนฺธานํ จิตฺตสมุฏฺฐานานญฺจ รูปานํ กมฺมปจฺจเยน ปจฺจโย. (3)}

\prose{นภาวนาย ปหาตพฺโพ ธมฺโม นภาวนาย ปหาตพฺพสฺส ธมฺมสฺส กมฺมปจฺจเยน ปจฺจโย.\csromanspacer{} \textbf{สหชาตา}, นานากฺขณิกา.\csromanspacer{}\csromanspacer{}\textbf{สหชาตา} นภาวนาย ปหาตพฺพา เจตนา สมฺปยุตฺตกานํ ขนฺธานํ จิตฺตสมุฏฺฐานานญฺจ รูปานํ กมฺมปจฺจเยน ปจฺจโย.\csromanspacer{} ปฏิสนฺธิกฺขเณ ฯเปฯ.\csromanspacer{} \textbf{นานากฺขณิกา} นภาวนาย ปหาตพฺพา เจตนา วิปากานํ ขนฺธานํ กฏตฺตา จ รูปานํ กมฺมปจฺจเยน ปจฺจโย.\csromanspacer{} วิปากปจฺจเยน ปจฺจโย.\csromanspacer{} เอกํ ฯเปฯ.\csromanspacer{} อวิคตปจฺจเยน ปจฺจโย. (สพฺพปจฺจยา ทสฺสนทุกสทิสา.\csromanspacer{} ภาวนา นินฺนานากรณา.)}

\csromanmatikaanchor{mtk.14}
\csromantocmarkref{section}{85. ทสฺสเนนปหาตพฺพเหตุกทุก}{mtk.14}
\bookmark[dest={mtk.14},level=1]{85. ทสฺสเนนปหาตพฺพเหตุกทุก}
\csromanheader{1}{85. ทสฺสเนนปหาตพฺพเหตุกทุก \textbf{1. ปจฺจยานุโลม 2. สงฺขฺยาวาร}}
\csromanheader{2}{\textbf{สุทฺธ}}
\proseitem{46}{\csromanfolio{57}\textbf{เหตุ}ยา จตฺตาริ, อารมฺมเณ จตฺตาริ, อธิปติยา ปญฺจ, อนนฺตเร จตฺตาริ, สมนนฺตเร จตฺตาริ, สหชาเต ปญฺจ, อญฺญมญฺเญ เทฺว, นิสฺสเย สตฺต, อุปนิสฺสเย จตฺตาริ, ปุเรชาเต เทฺว, ปจฺฉาชาเต เทฺว, อาเสวเน เทฺว, กมฺเม จตฺตาริ, วิปาเก เอกํ, อาหาเร จตฺตาริ, อินฺทฺริเย จตฺตาริ, ฌาเน จตฺตาริ, มคฺเค จตฺตาริ, สมฺปยุตฺเต เทฺว, วิปฺปยุตฺเต ตีณิ, อตฺถิยา สตฺต, นตฺถิยา จตฺตาริ, วิคเต จตฺตาริ, อวิคเต สตฺต.}

\prose{(ปจฺจนียวิภงฺโค ทสฺสนทุกสทิโส วิภชิตพฺโพ.\csromanspacer{} เอวํ ตีณิ คณนาปิ คเณตพฺพา.)}

\nitthitam{ภาวนาย ปหาตพฺพทุกํ นิฏฺฐิตํ.}
\csromanheader{2}{85. ทสฺสเนนปหาตพฺพเหตุกทุก 1. ปฏิจฺจวาร}
\proseitem{47}{ทสฺสเนน ปหาตพฺพเหตุกํ ธมฺมํ ปฏิจฺจ ทสฺสเนน ปหาตพฺพเหตุโก ธมฺโม อุปฺปชฺชติ เหตุปจฺจยา.\csromanspacer{} ทสฺสเนน ปหาตพฺพเหตุกํ เอกํ ขนฺธํ ปฏิจฺจ ตโย ขนฺธา ฯเปฯ เทฺว ขนฺเธ ฯเปฯ. (1)}

\prose{ทสฺสเนน ปหาตพฺพเหตุกํ ธมฺมํ ปฏิจฺจ นทสฺสเนน ปหาตพฺพเหตุโก ธมฺโม อุปฺปชฺชติ เหตุปจฺจยา.\csromanspacer{} ทสฺสเนน ปหาตพฺพเหตุเก ขนฺเธ ปฏิจฺจ จิตฺตสมุฏฺฐานํ รูปํ. (2)}

\prose{ทสฺสเนน ปหาตพฺพเหตุกํ ธมฺมํ ปฏิจฺจ ทสฺสเนน ปหาตพฺพเหตุโก จ นทสฺสเนน ปหาตพฺพเหตุโก จ ธมฺมา อุปฺปชฺชนฺติ เหตุปจฺจยา.\csromanspacer{} ทสฺสเนน ปหาตพฺพเหตุกํ เอกํ ขนฺธํ ปฏิจฺจ ตโย ขนฺธา จิตฺตสมุฏฺฐานญฺจ รูปํ ฯเปฯ เทฺว ขนฺเธ ฯเปฯ. (3)}

\proseitem{48}{\csromanfolio{58}นทสฺสเนน ปหาตพฺพเหตุกํ ธมฺมํ ปฏิจฺจ นทสฺสเนน ปหาตพฺพเหตุโก ธมฺโม อุปฺปชฺชติ เหตุปจฺจยา.\csromanspacer{} นทสฺสเนน ปหาตพฺพเหตุกํ เอกํ ขนฺธํ ปฏิจฺจ ตโย ขนฺธา จิตฺตสมุฏฺฐานญฺจ รูปํ ฯเปฯ เทฺว ขนฺเธ ฯเปฯ.\csromanspacer{} วิจิกิจฺฉาสหคตํ โมหํ ปฏิจฺจ จิตฺตสมุฏฺฐานํ รูปํ.\csromanspacer{} ปฏิสนฺธิกฺขเณ ฯเปฯ. (ยาว อชฺฌตฺติกา มหาภูตา.) (1)}

\prose{นทสฺสเนน ปหาตพฺพเหตุกํ ธมฺมํ ปฏิจฺจ ทสฺสเนน ปหาตพฺพเหตุโก ธมฺโม อุปฺปชฺชติ เหตุปจฺจยา.\csromanspacer{} วิจิกิจฺฉาสหคตํ โมหํ ปฏิจฺจ สมฺปยุตฺตกา ขนฺธา. (2)}

\prose{นทสฺสเนน ปหาตพฺพเหตุกํ ธมฺมํ ปฏิจฺจ ทสฺสเนน ปหาตพฺพเหตุโก จ นทสฺสเนน ปหาตพฺพเหตุโก จ ธมฺมา อุปฺปชฺชนฺติ เหตุปจฺจยา.\csromanspacer{} วิจิกิจฺฉาสหคตํ โมหํ ปฏิจฺจ สมฺปยุตฺตกา ขนฺธา จิตฺตสมุฏฺฐานญฺจ รูปํ. (3)}

\proseitem{49}{ทสฺสเนน ปหาตพฺพเหตุกญฺจ นทสฺสเนน ปหาตพฺพเหตุกญฺจ ธมฺมํ ปฏิจฺจ ทสฺสเนน ปหาตพฺพเหตุโก ธมฺโม อุปฺปชฺชติ เหตุปจฺจยา.\csromanspacer{} วิจิกิจฺฉาสหคตํ เอกํ ขนฺธญฺจ โมหญฺจ ปฏิจฺจ ตโย ขนฺธา ฯเปฯ เทฺว ขนฺเธ จ ฯเปฯ. (1)}

\prose{ทสฺสเนน ปหาตพฺพเหตุกญฺจ นทสฺสเนน ปหาตพฺพเหตุกญฺจ ธมฺมํ ปฏิจฺจ นทสฺสเนน ปหาตพฺพเหตุโก ธมฺโม อุปฺปชฺชติ เหตุปจฺจยา.\csromanspacer{} ทสฺสเนน ปหาตพฺพเหตุเก ขนฺเธ จ มหาภูเต จ ปฏิจฺจ จิตฺตสมุฏฺฐานํ รูปํ.\csromanspacer{} วิจิกิจฺฉาสหคเต ขนฺเธ จ โมหญฺจ ปฏิจฺจ จิตฺตสมุฏฺฐานํ รูปํ. (2)}

\prose{ทสฺสเนน ปหาตพฺพเหตุกญฺจ นทสฺสเนน ปหาตพฺพเหตุกญฺจ ธมฺมํ ปฏิจฺจ ทสฺสเนน ปหาตพฺพเหตุโก จ นทสฺสเนน ปหาตพฺพเหตุโก จ ธมฺมา อุปฺปชฺชนฺติ เหตุปจฺจยา.\csromanspacer{} วิจิกิจฺฉาสหคตํ เอกํ ขนฺธญฺจ โมหญฺจ ปฏิจฺจ ตโย ขนฺธา จิตฺตสมุฏฺฐานญฺจ รูปํ ฯเปฯ เทฺว ขนฺเธ จ ฯเปฯ. (3)}

\csromanheader{2}{\textbf{อารมฺมณ}}
\proseitem{50}{ทสฺสเนน ปหาตพฺพเหตุกํ ธมฺมํ ปฏิจฺจ ทสฺสเนน ปหาตพฺพเหตุโก ธมฺโม อุปฺปชฺชติ อารมฺมณปจฺจยา.\csromanspacer{} ทสฺสเนน ปหาตพฺพเหตุกํ เอกํ ขนฺธํ ปฏิจฺจ ตโย ขนฺธา ฯเปฯ เทฺว ขนฺเธ ฯเปฯ. (1)}

\csromanheader{2}{85. ทสฺสเนนปหาตพฺพเหตุกทุก}
\prose{\csromanfolio{59}ทสฺสเนน ปหาตพฺพเหตุกํ ธมฺมํ ปฏิจฺจ นทสฺสเนน ปหาตพฺพเหตุโก ธมฺโม อุปฺปชฺชติ อารมฺมณปจฺจยา.\csromanspacer{} วิจิกิจฺฉาสหคเต ขนฺเธ ปฏิจฺจ วิจิกิจฺฉาสหคโต โมโห. (2)}

\prose{ทสฺสเนน ปหาตพฺพเหตุกํ ธมฺมํ ปฏิจฺจ ทสฺสเนน ปหาตพฺพเหตุโก จ นทสฺสเนน ปหาตพฺพเหตุโก จ ธมฺมา อุปฺปชฺชนฺติ อารมฺมณปจฺจยา.\csromanspacer{} วิจิกิจฺฉาสหคตํ เอกํ ขนฺธํ ปฏิจฺจ ตโย ขนฺธา โมโห จ ฯเปฯ เทฺว ขนฺเธ ฯเปฯ. (3)}

\proseitem{51}{นทสฺสเนน ปหาตพฺพเหตุกํ ธมฺมํ ปฏิจฺจ นทสฺสเนน ปหาตพฺพเหตุโก ธมฺโม อุปฺปชฺชติ อารมฺมณปจฺจยา.\csromanspacer{} นทสฺสเนน ปหาตพฺพเหตุกํ เอกํ ขนฺธํ ปฏิจฺจ ตโย ขนฺธา ฯเปฯ เทฺว ขนฺเธ ฯเปฯ.\csromanspacer{} ปฏิสนฺธิกฺขเณ ฯเปฯ. (1)}

\prose{นทสฺสเนน ปหาตพฺพเหตุกํ ธมฺมํ ปฏิจฺจ ทสฺสเนน ปหาตพฺพเหตุโก ธมฺโม อุปฺปชฺชติ อารมฺมณปจฺจยา.\csromanspacer{} วิจิกิจฺฉาสหคตํ โมหํ ปฏิจฺจ สมฺปยุตฺตกา ขนฺธา. (2)}

\prose{ทสฺสเนน ปหาตพฺพเหตุกญฺจ นทสฺสเนน ปหาตพฺพเหตุกญฺจ ธมฺมํ ปฏิจฺจ ทสฺสเนน ปหาตพฺพเหตุโก ธมฺโม อุปฺปชฺชติ อารมฺมณปจฺจยา.\csromanspacer{} วิจิกิจฺฉาสหคตํ เอกํ ขนฺธญฺจ โมหญฺจ ปฏิจฺจ ตโย ขนฺธา ฯเปฯ เทฺว ขนฺเธ จ ฯเปฯ. (1) (สํขิตฺตํ.)}

\csromanheader{2}{1. ปจฺจยานุโลม 2. สงฺขฺยาวาร}
\proseitem{52}{เหตุยา นว, อารมฺมเณ ฉ, อธิปติยา ปญฺจ, อนนฺตเร ฉ, สมนนฺตเร ฉ, สหชาเต นว, อญฺญมญฺเญ ฉ, นิสฺสเย นว, อุปนิสฺสเย ฉ, ปุเรชาเต ฉ, อาเสวเน ฉ, กมฺเม นว, วิปาเก เอกํ, อาหาเร นว, อินฺทฺริเย นว, ฌาเน นว, มคฺเค นว, สมฺปยุตฺเต ฉ, วิปฺปยุตฺเต นว, อตฺถิยา นว, นตฺถิยา ฉ, วิคเต ฉ, อวิคเต นว.}

\csromanheader{2}{2. ปจฺจยปจฺจนีย 1. วิภงฺควาร}
\csromanheader{2}{\textbf{นเหตุ}}
\proseitem{53}{\csromanfolio{60}ทสฺสเนน ปหาตพฺพเหตุกํ ธมฺมํ ปฏิจฺจ นทสฺสเนน ปหาตพฺพเหตุโก ธมฺโม อุปฺปชฺชติ นเหตุปจฺจยา.\csromanspacer{} วิจิกิจฺฉาสหคเต ขนฺเธ ปฏิจฺจ วิจิกิจฺฉาสหคโต โมโห. (1)}

\prose{นทสฺสเนน ปหาตพฺพเหตุกํ ธมฺมํ ปฏิจฺจ นทสฺสเนน ปหาตพฺพเหตุโก ธมฺโม อุปฺปชฺชติ นเหตุปจฺจยา.\csromanspacer{} อเหตุกํ นทสฺสเนน ปหาตพฺพเหตุกํ เอกํ ขนฺธํ ปฏิจฺจ ตโย ขนฺธา จิตฺตสมุฏฺฐานญฺจ รูปํ ฯเปฯ. (ยาว อสญฺญสตฺตา.) อุทฺธจฺจสหคเต ขนฺเธ ปฏิจฺจ อุทฺธจฺจสหคโต โมโห. (1)}

\csromanheader{2}{2. ปจฺจยปจฺจนีย 2. สงฺขฺยาวาร}
\vspace{\tipitakaparskip}
\csromancenter{\textbf{นเหตุ}ยา เทฺว, น-อารมฺมเณ ตีณิ, น-อธิปติยา นว, น-อนนฺตเร ตีณิ, นสมนนฺตเร ตีณิ, น-อญฺญมญฺเญ ตีณิ, น-อุปนิสฺสเย ตีณิ, นปุเรชาเต สตฺต, นปจฺฉาชาเต นว, น-อาเสวเน นว, นกมฺเม จตฺตาริ, นวิปาเก นว, น-อาหาเร เอกํ, น-อินฺทฺริเย เอกํ, นฌาเน เอกํ, นมคฺเค เอกํ, นสมฺปยุตฺเต ตีณิ, นวิปฺปยุตฺเต ฉ, โนนตฺถิยา ตีณิ, โนวิคเต ตีณิ.}
\csromanheader{2}{3. ปจฺจยานุโลมปจฺจนีย}
\proseitem{55}{เหตุปจฺจยา น-อารมฺมเณ ตีณิ, น-อธิปติยา นว ฯเปฯ นปุเรชาเต สตฺต ฯเปฯ นวิปฺปยุตฺเต จตฺตาริ, โนนตฺถิยา ตีณิ, โนวิคเต ตีณิ.}

\csromanheader{2}{4. ปจฺจยปจฺจนียานุโลม}
\vspace{\tipitakaparskip}
\csromancenter{\textbf{นเหตุ}ปจฺจยา อารมฺมเณ เทฺว, อนนฺตเร เทฺว ฯเปฯ วิปาเก เอกํ ฯเปฯ มคฺเค เทฺว ฯเปฯ อวิคเต เทฺว.}
\csromancenter{(สหชาตวาโร ปฏิจฺจวารสทิโส.)}
\csromancenter{ทสฺสเนนปหาตพฺพเหตุกทุก \textbf{85.}\csromanspacer{}\textbf{ ทสฺสเนนปหาตพฺพเหตุกทุก 3.}\csromanspacer{}\textbf{ ปจฺจยวาร}}
\proseitem{57}{\csromanfolio{61}ทสฺสเนน ปหาตพฺพเหตุกํ ธมฺมํ ปจฺจยา ทสฺสเนน ปหาตพฺพเหตุโก ธมฺโม อุปฺปชฺชติ เหตุปจฺจยา.\csromanspacer{} ตีณิ.}

\prose{นทสฺสเนน ปหาตพฺพเหตุกํ ธมฺมํ ปจฺจยา นทสฺสเนน ปหาตพฺพเหตุโก ธมฺโม อุปฺปชฺชติ เหตุปจฺจยา.\csromanspacer{} นทสฺสเนน ปหาตพฺพเหตุกํ เอกํ ขนฺธํ ปจฺจยา ตโย ขนฺธา จิตฺตสมุฏฺฐานญฺจ รูปํ ฯเปฯ เทฺว ขนฺเธ ฯเปฯ.\csromanspacer{} ปฏิสนฺธิกฺขเณ ฯเปฯ. (ยาว อชฺฌตฺติกา มหาภูตา.) วตฺถุํ ปจฺจยา นทสฺสเนน ปหาตพฺพเหตุกา ขนฺธา. (1)}

\prose{นทสฺสเนน ปหาตพฺพเหตุกํ ธมฺมํ ปจฺจยา ทสฺสเนน ปหาตพฺพเหตุโก ธมฺโม อุปฺปชฺชติ เหตุปจฺจยา.\csromanspacer{} วตฺถุํ ปจฺจยา ทสฺสเนน ปหาตพฺพเหตุกา ขนฺธา.\csromanspacer{} วิจิกิจฺฉาสหคตํ โมหํ ปจฺจยา สมฺปยุตฺตกา ขนฺธา. (2)}

\prose{นทสฺสเนน ปหาตพฺพเหตุกํ ธมฺมํ ปจฺจยา ทสฺสเนน ปหาตพฺพเหตุโก จ นทสฺสเนน ปหาตพฺพเหตุโก จ ธมฺมา อุปฺปชฺชนฺติ เหตุปจฺจยา.\csromanspacer{} วตฺถุํ ปจฺจยา ทสฺสเนน ปหาตพฺพเหตุกา ขนฺธา.\csromanspacer{} มหาภูเต ปจฺจยา จิตฺตสมุฏฺฐานํ รูปํ.\csromanspacer{} วิจิกิจฺฉาสหคตํ โมหํ ปจฺจยา สมฺปยุตฺตกา ขนฺธา จิตฺตสมุฏฺฐานญฺจ รูปํ. (3)}

\proseitem{58}{ทสฺสเนน ปหาตพฺพเหตุกญฺจ นทสฺสเนน ปหาตพฺพเหตุกญฺจ ธมฺมํ ปจฺจยา ทสฺสเนน ปหาตพฺพเหตุโก ธมฺโม อุปฺปชฺชติ เหตุปจฺจยา.\csromanspacer{} ทสฺสเนน ปหาตพฺพเหตุกํ เอกํ ขนฺธญฺจ วตฺถุญฺจ ปจฺจยา ตโย ขนฺธา ฯเปฯ เทฺว ขนฺเธ จ ฯเปฯ.\csromanspacer{} วิจิกิจฺฉาสหคตํ เอกํ ขนฺธญฺจ วตฺถุญฺจ ปจฺจยา ตโย ขนฺธา ฯเปฯ เทฺว ขนฺเธ จ ฯเปฯ. (1)}

\prose{ทสฺสเนน ปหาตพฺพเหตุกญฺจ นทสฺสเนน ปหาตพฺพเหตุกญฺจ ธมฺมํ ปจฺจยา นทสฺสเนน ปหาตพฺพเหตุโก ธมฺโม อุปฺปชฺชติ เหตุปจฺจยา.\csromanspacer{} ทสฺสเนน ปหาตพฺพเหตุเก ขนฺเธ จ มหาภูเต จ ปจฺจยา จิตฺตสมุฏฺฐานํ รูปํ.\csromanspacer{} วิจิกิจฺฉาสหคเต ขนฺเธ จ โมหญฺจ ปจฺจยา จิตฺตสมุฏฺฐานํ รูปํ. (2)}

\prose{\csromanfolio{62}ทสฺสเนน ปหาตพฺพเหตุกญฺจ นทสฺสเนน ปหาตพฺพเหตุกญฺจ ธมฺมํ ปจฺจยา ทสฺสเนน ปหาตพฺพเหตุโก จ นทสฺสเนน ปหาตพฺพเหตุโก จ ธมฺมา อุปฺปชฺชนฺติ เหตุปจฺจยา.\csromanspacer{} ทสฺสเนน ปหาตพฺพเหตุกํ เอกํ ขนฺธญฺจ วตฺถุญฺจ ปจฺจยา ตโย ขนฺธา ฯเปฯ เทฺว ขนฺเธ จ ฯเปฯ.\csromanspacer{} ทสฺสเนน ปหาตพฺพเหตุเก ขนฺเธ จ มหาภูเต จ ปจฺจยา จิตฺตสมุฏฺฐานํ รูปํ.\csromanspacer{} วิจิกิจฺฉาสหคตํ เอกํ ขนฺธญฺจ โมหญฺจ ปจฺจยา ตโย ขนฺธา จิตฺตสมุฏฺฐานญฺจ รูปํ ฯเปฯ เทฺว ขนฺเธ จ ฯเปฯ. (3)}

\csromanheader{2}{\textbf{อารมฺมณ}}
\proseitem{59}{ทสฺสเนน ปหาตพฺพเหตุกํ ธมฺมํ ปจฺจยา ทสฺสเนน ปหาตพฺพเหตุโก ธมฺโม อุปฺปชฺชติ อารมฺมณปจฺจยา.\csromanspacer{} ตีณิ. (ปฏิจฺจสทิสา.)}

\prose{นทสฺสเนน ปหาตพฺพเหตุกํ ธมฺมํ ปจฺจยา นทสฺสเนน ปหาตพฺพเหตุโก ธมฺโม อุปฺปชฺชติ อารมฺมณปจฺจยา.\csromanspacer{} นทสฺสเนน ปหาตพฺพเหตุกํ เอกํ ขนฺธํ ปจฺจยา ตโย ขนฺธา ฯเปฯ เทฺว ขนฺเธ ฯเปฯ.\csromanspacer{} ปฏิสนฺธิกฺขเณ ฯเปฯ.\csromanspacer{} จกฺขายตนํ ปจฺจยา จกฺขุวิญฺญาณํ ฯเปฯ.\csromanspacer{} วตฺถุํ ปจฺจยา นทสฺสเนน ปหาตพฺพเหตุกา ขนฺธา.\csromanspacer{} วตฺถุํ ปจฺจยา วิจิกิจฺฉาสหคโต โมโห.\csromanspacer{}\csromanspacer{}นทสฺสเนน ปหาตพฺพเหตุกํ ธมฺมํ ปจฺจยา ทสฺสเนน ปหาตพฺพเหตุโก ธมฺโม อุปฺปชฺชติ อารมฺมณปจฺจยา.\csromanspacer{} วตฺถุํ ปจฺจยา ทสฺสเนน ปหาตพฺพเหตุกา ขนฺธา.\csromanspacer{} วิจิกิจฺฉาสหคตํ โมหํ ปจฺจยา สมฺปยุตฺตกา ขนฺธา.\csromanspacer{}\csromanspacer{}นทสฺสเนน ปหาตพฺพเหตุกํ ธมฺมํ ปจฺจยา ทสฺสเนน ปหาตพฺพเหตุโก จ นทสฺสเนน ปหาตพฺพเหตุโก จ ธมฺมา อุปฺปชฺชนฺติ อารมฺมณปจฺจยา.\csromanspacer{} วตฺถุํ ปจฺจยา วิจิกิจฺฉาสหคตา ขนฺธา จ โมโห จ. (3)}

\proseitem{60}{ทสฺสเนน ปหาตพฺพเหตุกญฺจ นทสฺสเนน ปหาตพฺพเหตุกญฺจ ธมฺมํ ปจฺจยา ทสฺสเนน ปหาตพฺพเหตุโก ธมฺโม อุปฺปชฺชติ อารมฺมณปจฺจยา.\csromanspacer{} ทสฺสเนน ปหาตพฺพเหตุกํ เอกํ ขนฺธญฺจ วตฺถุญฺจ ปจฺจยา ตโย ขนฺธา ฯเปฯ เทฺว ขนฺเธ จ ฯเปฯ.\csromanspacer{} วิจิกิจฺฉาสหคตํ เอกํ ขนฺธญฺจ โมหญฺจ ปจฺจยา ตโย ขนฺธา ฯเปฯ เทฺว ขนฺเธ จ ฯเปฯ. (1)}

\prose{ทสฺสเนน ปหาตพฺพเหตุกญฺจ นทสฺสเนน ปหาตพฺพเหตุกญฺจ ธมฺมํ ปจฺจยา นทสฺสเนน ปหาตพฺพเหตุโก ธมฺโม อุปฺปชฺชติ อารมฺมณปจฺจยา.}

\vspace{\tipitakaparskip}
\csromancenter{ทสฺสเนนปหาตพฺพเหตุกทุก วิจิกิจฺฉาสหคเต ขนฺเธ จ วตฺถุญฺจ ปจฺจยา วิจิกิจฺฉาสหคโต โมโห. (2)}
\prose{\csromanfolio{63}ทสฺสเนน ปหาตพฺพเหตุกญฺจ นทสฺสเนน ปหาตพฺพเหตุกญฺจ ธมฺมํ ปจฺจยา ทสฺสเนน ปหาตพฺพเหตุโก จ นทสฺสเนน ปหาตพฺพเหตุโก จ ธมฺมา อุปฺปชฺชนฺติ อารมฺมณปจฺจยา.\csromanspacer{} วิจิกิจฺฉาสหคตํ เอกํ ขนฺธญฺจ วตฺถุญฺจ ปจฺจยา ตโย ขนฺธา โมโห จ ฯเปฯ เทฺว ขนฺเธ จ ฯเปฯ. (3)}

\csromanheader{2}{1. ปจฺจยานุโลม 2. สงฺขฺยาวาร}
\proseitem{61}{เหตุยา นว, อารมฺมเณ นว, อธิปติยา นว, (สพฺพตฺถ นว.) วิปาเก เอกํ ฯเปฯ อวิคเต นว.}

\csromanheader{2}{2. ปจฺจยปจฺจนีย 1. วิภงฺควาร}
\csromanheader{2}{\textbf{นเหตุ}}
\proseitem{62}{ทสฺสเนน ปหาตพฺพเหตุกํ ธมฺมํ ปจฺจยา นทสฺสเนน ปหาตพฺพเหตุโก ธมฺโม อุปฺปชฺชติ นเหตุปจฺจยา.\csromanspacer{} วิจิกิจฺฉาสหคเต ขนฺเธ ปจฺจยา วิจิกิจฺฉาสหคโต โมโห. (1)}

\prose{นทสฺสเนน ปหาตพฺพเหตุกํ ธมฺมํ ปจฺจยา นทสฺสเนน ปหาตพฺพเหตุโก ธมฺโม อุปฺปชฺชติ นเหตุปจฺจยา.\csromanspacer{} อเหตุกํ นทสฺสเนน ปหาตพฺพเหตุกํ ฯเปฯ. (ยาว อสญฺญสตฺตา.) จกฺขายตนํ ปจฺจยา จกฺขุวิญฺญาณํ ฯเปฯ.\csromanspacer{} กายายตนํ ปจฺจยา กายวิญฺญาณํ.\csromanspacer{} วตฺถุํ ปจฺจยา อเหตุกา นทสฺสเนน ปหาตพฺพเหตุกา ขนฺธา.\csromanspacer{} อุทฺธจฺจสหคเต ขนฺเธ จ วตฺถุญฺจ ปจฺจยา อุทฺธจฺจสหคโต โมโห. (1)}

\prose{ทสฺสเนน ปหาตพฺพเหตุกญฺจ นทสฺสเนน ปหาตพฺพเหตุกญฺจ ธมฺมํ ปจฺจยา นทสฺสเนน ปหาตพฺพเหตุโก ธมฺโม อุปฺปชฺชติ นเหตุปจฺจยา.\csromanspacer{} วิจิกิจฺฉาสหคเต ขนฺเธ จ วตฺถุญฺจ ปจฺจยา วิจิกิจฺฉาสหคโต โมโห. (สํขิตฺตํ.) (1)}

\csromanheader{2}{2. ปจฺจยปจฺจนีย 2. สงฺขฺยาวาร}
\proseitem{63}{\csromanfolio{64}นเหตุยา ตีณิ, น-อารมฺมเณ ตีณิ, น-อธิปติยา นว, น-อนนฺตเร ตีณิ, นสมนนฺตเร ตีณิ, น-อญฺญมญฺเญ ตีณิ, น-อุปนิสฺสเย ตีณิ, นปุเรชาเต สตฺต, นปจฺฉาชาเต นว, น-อาเสวเน นว, นกมฺเม จตฺตาริ, นวิปาเก นว, น-อาหาเร เอกํ, น-อินฺทฺริเย เอกํ, นฌาเน เอกํ, นมคฺเค เอกํ, นสมฺปยุตฺเต ตีณิ, นวิปฺปยุตฺเต ฉ, โนนตฺถิยา ตีณิ, โนวิคเต ตีณิ.}

\vspace{\tipitakaparskip}
\csromancenter{(เอวํ อิตเร เทฺว คณนาปิ นิสฺสยวาโรปิ กาตพฺโพ.)}
\csromanheader{2}{85. ทสฺสเนนปหาตพฺพเหตุกทุก 5. สํสฏฺฐวาร}
\csromanheader{2}{\textbf{ปจฺจยจตุกฺก}}
\proseitem{64}{ทสฺสเนน ปหาตพฺพเหตุกํ ธมฺมํ สํสฏฺโฐ ทสฺสเนน ปหาตพฺพเหตุโก ธมฺโม อุปฺปชฺชติ เหตุปจฺจยา.\csromanspacer{} ทสฺสเนน ปหาตพฺพเหตุกํ เอกํ ขนฺธํ สํสฏฺฐา ตโย ขนฺธา ฯเปฯ เทฺว ขนฺเธ ฯเปฯ. (1)}

\prose{นทสฺสเนน ปหาตพฺพเหตุกํ ธมฺมํ สํสฏฺโฐ นทสฺสเนน ปหาตพฺพเหตุโก ธมฺโม อุปฺปชฺชติ เหตุปจฺจยา.\csromanspacer{} นทสฺสเนน ปหาตพฺพเหตุกํ เอกํ ขนฺธํ สํสฏฺฐา ตโย ขนฺธา ฯเปฯ เทฺว ขนฺเธ ฯเปฯ.\csromanspacer{} ปฏิสนฺธิกฺขเณ ฯเปฯ. (1)}

\prose{นทสฺสเนน ปหาตพฺพเหตุกํ ธมฺมํ สํสฏฺโฐ ทสฺสเนน ปหาตพฺพเหตุโก ธมฺโม อุปฺปชฺชติ เหตุปจฺจยา.\csromanspacer{} วิจิกิจฺฉาสหคตํ โมหํ สํสฏฺฐา สมฺปยุตฺตกา ขนฺธา. (2)}

\prose{ทสฺสเนน ปหาตพฺพเหตุกญฺจ นทสฺสเนน ปหาตพฺพเหตุกญฺจ ธมฺมํ สํสฏฺโฐ ทสฺสเนน ปหาตพฺพเหตุโก ธมฺโม อุปฺปชฺชติ เหตุปจฺจยา.\csromanspacer{} วิจิกิจฺฉาสหคตํ เอกํ ขนฺธญฺจ โมหญฺจ สํสฏฺฐา ตโย ขนฺธา ฯเปฯ เทฺว ขนฺเธ จ ฯเปฯ. (1)}

\csromanheader{2}{85. ทสฺสเนนปหาตพฺพเหตุกทุก \textbf{อารมฺมณ}}
\proseitem{65}{\csromanfolio{65}ทสฺสเนน ปหาตพฺพเหตุกํ ธมฺมํ สํสฏฺโฐ ทสฺสเนน ปหาตพฺพเหตุโก ธมฺโม อุปฺปชฺชติ อารมฺมณปจฺจยา.\csromanspacer{} ตีณิ. (ปฏิจฺจสทิสา.)}

\prose{นทสฺสเนน ปหาตพฺพเหตุกํ ธมฺมํ สํสฏฺโฐ นทสฺสเนน ปหาตพฺพเหตุโก ธมฺโม อุปฺปชฺชติ อารมฺมณปจฺจยา.\csromanspacer{} นทสฺสเนน ปหาตพฺพเหตุกํ เอกํ ขนฺธํ สํสฏฺฐา ตโย ขนฺธา ฯเปฯ เทฺว ขนฺเธ ฯเปฯ.\csromanspacer{} ปฏิสนฺธิกฺขเณ ฯเปฯ. (1)}

\prose{นทสฺสเนน ปหาตพฺพเหตุกํ ธมฺมํ สํสฏฺโฐ ทสฺสเนน ปหาตพฺพเหตุโก ธมฺโม อุปฺปชฺชติ อารมฺมณปจฺจยา. (ปฏิจฺจสทิสํ.) (2)}

\prose{ทสฺสเนน ปหาตพฺพเหตุกญฺจ นทสฺสเนน ปหาตพฺพเหตุกญฺจ ธมฺมํ สํสฏฺโฐ ทสฺสเนน ปหาตพฺพเหตุโก ธมฺโม อุปฺปชฺชติ อารมฺมณปจฺจยา. (ปฏิจฺจสทิสํ.\csromanspacer{} สํขิตฺตํ.) (1)}

\proseitem{66}{เหตุยา จตฺตาริ, อารมฺมเณ ฉ, อธิปติยา เทฺว, อนนฺตเร ฉ, (สพฺพตฺถ ฉ,) วิปาเก เอกํ ฯเปฯ อวิคเต ฉ.\csromanspacer{} อนุโลมํ.}

\csromanheader{2}{\textbf{นเหตุ}}
\proseitem{67}{ทสฺสเนน ปหาตพฺพเหตุกํ ธมฺมํ สํสฏฺโฐ นทสฺสเนน ปหาตพฺพเหตุโก ธมฺโม อุปฺปชฺชติ นเหตุปจฺจยา.\csromanspacer{} วิจิกิจฺฉาสหคเต ขนฺเธ สํสฏฺโฐ วิจิกิจฺฉาสหคโต โมโห. (1)}

\prose{นทสฺสเนน ปหาตพฺพเหตุกํ ธมฺมํ สํสฏฺโฐ นทสฺสเนน ปหาตพฺพเหตุโก ธมฺโม อุปฺปชฺชติ นเหตุปจฺจยา.\csromanspacer{} อเหตุกํ นทสฺสเนน ปหาตพฺพเหตุกํ เอกํ ขนฺธํ สํสฏฺฐา ตโย ขนฺธา ฯเปฯ เทฺว ขนฺเธ ฯเปฯ.\csromanspacer{} อเหตุกปฏิสนฺธิกฺขเณ ฯเปฯ.\csromanspacer{} อุทฺธจฺจสหคเต ขนฺเธ สํสฏฺโฐ อุทฺธจฺจสหคโต โมโห. (1)}

\vspace{\tipitakaparskip}
\csromancenter{\textbf{นเหตุ}ยา เทฺว, น-อธิปติยา ฉ, นปุเรชาเต ฉ, นปจฺฉาชาเต ฉ, น-อาเสวเน ฉ, นกมฺเม จตฺตาริ, นวิปาเก ฉ, นฌาเน เอกํ, นมคฺเค เอกํ, นวิปฺปยุตฺเต ฉ.}
\csromancenter{ปจฺจนียํ.}
\csromanheader{2}{\textbf{เหตุทุก}}
\vspace{\tipitakaparskip}
\csromancenter{\textbf{เหตุ}ปจฺจยา น-อธิปติยา จตฺตาริ, นปุเรชาเต จตฺตาริ ฯเปฯ นวิปฺปยุตฺเต จตฺตาริ.}
\csromanheader{2}{\textbf{นเหตุทุก}}
\vspace{\tipitakaparskip}
\csromancenter{นเหตุปจฺจยา อารมฺมเณ เทฺว, อนนฺตเร เทฺว ฯเปฯ วิปาเก เอกํ ฯเปฯ อวิคเต เทฺว.}
\csromancenter{(สมฺปยุตฺตวาโร สํสฏฺฐวารสทิโส.)}
\csromanheader{2}{85. ทสฺสเนนปหาตพฺพเหตุกทุก 7. ปญฺหาวาร}
\proseitem{68}{\csromanfolio{66}ทสฺสเนน ปหาตพฺพเหตุโก ธมฺโม ทสฺสเนน ปหาตพฺพเหตุกสฺส ธมฺมสฺส เหตุปจฺจเยน ปจฺจโย.\csromanspacer{} ทสฺสเนน ปหาตพฺพเหตุกา เหตู สมฺปยุตฺตกานํ ขนฺธานํ เหตุปจฺจเยน ปจฺจโย. (มูลํ กาตพฺพํ.) ทสฺสเนน ปหาตพฺพเหตุกา เหตู จิตฺตสมุฏฺฐานานํ รูปานํ เหตุปจฺจเยน ปจฺจโย. (มูลํ กาตพฺพํ.) ทสฺสเนน ปหาตพฺพเหตุกา เหตู สมฺปยุตฺตกานํ ขนฺธานํ จิตฺตสมุฏฺฐานานญฺจ รูปานํ เหตุปจฺจเยน ปจฺจโย. (3)}

\proseitem{69}{นทสฺสเนน ปหาตพฺพเหตุโก ธมฺโม นทสฺสเนน ปหาตพฺพเหตุกสฺส ธมฺมสฺส เหตุปจฺจเยน ปจฺจโย.\csromanspacer{} นทสฺสเนน ปหาตพฺพเหตุกา เหตู สมฺปยุตฺตกานํ ขนฺธานํ จิตฺตสมุฏฺฐานานญฺจ รูปานํ เหตุปจฺจเยน ปจฺจโย.\csromanspacer{} วิจิกิจฺฉาสหคโต โมโห จิตฺตสมุฏฺฐานานํ รูปานํ เหตุปจฺจเยน ปจฺจโย.\csromanspacer{} ปฏิสนฺธิกฺขเณ ฯเปฯ. (มูลํ กาตพฺพํ.) วิจิกิจฺฉาสหคโต โมโห สมฺปยุตฺตกานํ ขนฺธานํ เหตุปจฺจเยน ปจฺจโย. (มูลํ กาตพฺพํ.) วิจิกิจฺฉาสหคโต โมโห สมฺปยุตฺตกานํ ขนฺธานํ จิตฺตสมุฏฺฐานานญฺจ รูปานํ เหตุปจฺจเยน ปจฺจโย. (3)}

\csromanheader{2}{85. ทสฺสเนนปหาตพฺพเหตุกทุก \textbf{อารมฺมณ}}
\proseitem{70}{\csromanfolio{67}ทสฺสเนน ปหาตพฺพเหตุโก ธมฺโม ทสฺสเนน ปหาตพฺพเหตุกสฺส ธมฺมสฺส อารมฺมณปจฺจเยน ปจฺจโย.\csromanspacer{} ทสฺสเนน ปหาตพฺพเหตุเก ขนฺเธ อารพฺภ ทสฺสเนน ปหาตพฺพเหตุกา ขนฺธา อุปฺปชฺชนฺติ. (มูลํ กาตพฺพํ.) ทสฺสเนน ปหาตพฺพเหตุเก ขนฺเธ อารพฺภ นทสฺสเนน ปหาตพฺพเหตุกา ขนฺธา จ โมโห จ อุปฺปชฺชนฺติ. (มูลํ กาตพฺพํ.) ทสฺสเนน ปหาตพฺพเหตุเก ขนฺเธ อารพฺภ วิจิกิจฺฉาสหคตา ขนฺธา จ โมโห จ อุปฺปชฺชนฺติ. (3)}

\proseitem{71}{นทสฺสเนน ปหาตพฺพเหตุโก ธมฺโม นทสฺสเนน ปหาตพฺพเหตุกสฺส ธมฺมสฺส อารมฺมณปจฺจเยน ปจฺจโย.\csromanspacer{} ทานํ ฯเปฯ สีลํ ฯเปฯ อุโปสถกมฺมํ กตฺวา ตํ ปจฺจเวกฺขติ, อสฺสาเทติ อภินนฺทติ, ตํ อารพฺภ นทสฺสเนน ปหาตพฺพเหตุโก ราโค อุปฺปชฺชติ.\csromanspacer{} อุทฺธจฺจํ อุปฺปชฺชติ.\csromanspacer{} นทสฺสเนน ปหาตพฺพเหตุกํ โทมนสฺสํ อุปฺปชฺชติ.\csromanspacer{} ปุพฺเพ สุจิณฺณานิ ฯเปฯ.\csromanspacer{} ฌานา วุฏฺฐหิตฺวา ฯเปฯ.\csromanspacer{} อริยา มคฺคา วุฏฺฐหิตฺวา มคฺคํ ปจฺจเวกฺขนฺติ ฯเปฯ.\csromanspacer{} อริยา นทสฺสเนน ปหาตพฺพเหตุเก ปหีเน กิเลเส ปจฺจเวกฺขนฺติ.\csromanspacer{} วิกฺขมฺภิเต กิเลเส ฯเปฯ.\csromanspacer{} ปุพฺเพ สมุทาจิณฺเณ กิเลเส ชานนฺติ.\csromanspacer{} จกฺขุํ ฯเปฯ.\csromanspacer{} วตฺถุํ นทสฺสเนน ปหาตพฺพเหตุเก ขนฺเธ จ โมหญฺจ อนิจฺจโต ฯเปฯ วิปสฺสติ, อสฺสาเทติ อภินนฺทติ, ตํ อารพฺภ นทสฺสเนน ปหาตพฺพเหตุโก ราโค อุปฺปชฺชติ.\csromanspacer{} อุทฺธจฺจํ อุปฺปชฺชติ.\csromanspacer{} นทสฺสเนน ปหาตพฺพเหตุกํ โทมนสฺสํ อุปฺปชฺชติ.\csromanspacer{} ทิพฺเพน จกฺขุนา รูปํ ปสฺสติ ฯเปฯ อนาคตํสญาณสฺส, อาวชฺชนาย, โมหสฺส จ อารมฺมณปจฺจเยน ปจฺจโย. (1)}

\prose{นทสฺสเนน ปหาตพฺพเหตุโก ธมฺโม ทสฺสเนน ปหาตพฺพเหตุกสฺส ธมฺมสฺส อารมฺมณปจฺจเยน ปจฺจโย.\csromanspacer{} ทานํ ฯเปฯ.\csromanspacer{} ฌานํ ฯเปฯ.\csromanspacer{} จกฺขุํ ฯเปฯ.\csromanspacer{} วตฺถุํ นทสฺสเนน ปหาตพฺพเหตุเก ขนฺเธ จ โมหญฺจ อสฺสาเทติ อภินนฺทติ, ตํ อารพฺภ ทสฺสเนน ปหาตพฺพเหตุโก ราโค อุปฺปชฺชติ.\csromanspacer{} ทิฏฺฐิ ฯเปฯ วิจิกิจฺฉา ฯเปฯ.\csromanspacer{} ทสฺสเนน ปหาตพฺพเหตุกํ โทมนสฺสํ อุปฺปชฺชติ. (2)}

\prose{นทสฺสเนน ปหาตพฺพเหตุโก ธมฺโม ทสฺสเนน ปหาตพฺพเหตุกสฺส จ นทสฺสเนน ปหาตพฺพเหตุกสฺส จ ธมฺมสฺส อารมฺมณปจฺจเยน ปจฺจโย.\csromanspacer{} จกฺขุํ ฯเปฯ.\csromanspacer{} วตฺถุํ นทสฺสเนน ปหาตพฺพเหตุเก ขนฺเธ จ โมหญฺจ อารพฺภ วิจิกิจฺฉาสหคตา ขนฺธา จ โมโห จ อุปฺปชฺชนฺติ. (3)}

\proseitem{72}{\csromanfolio{68}ทสฺสเนน ปหาตพฺพเหตุโก จ นทสฺสเนน ปหาตพฺพเหตุโก จ ธมฺมา ทสฺสเนน ปหาตพฺพเหตุกสฺส ธมฺมสฺส อารมฺมณปจฺจเยน ปจฺจโย.\csromanspacer{} วิจิกิจฺฉาสหคเต ขนฺเธ จ โมหญฺจ อารพฺภ ทสฺสเนน ปหาตพฺพเหตุกา ขนฺธา อุปฺปชฺชนฺติ. (มูลํ กาตพฺพํ.) วิจิกิจฺฉาสหคเต ขนฺเธ จ โมหญฺจ อารพฺภ นทสฺสเนน ปหาตพฺพเหตุกา ขนฺธา จ โมโห จ อุปฺปชฺชนฺติ. (มูลํ กาตพฺพํ.) วิจิกิจฺฉาสหคเต ขนฺเธ จ โมหญฺจ อารพฺภ วิจิกิจฺฉาสหคตา ขนฺธา จ โมโห จ อุปฺปชฺชนฺติ. (3)}

\csromanheader{2}{\textbf{อธิปติ}}
\proseitem{73}{ทสฺสเนน ปหาตพฺพเหตุโก ธมฺโม ทสฺสเนน ปหาตพฺพเหตุกสฺส ธมฺมสฺส อธิปติปจฺจเยน ปจฺจโย.\csromanspacer{} อารมฺมณาธิปติ, สหชาตาธิปติ.\csromanspacer{}\csromanspacer{}\textbf{อารมฺมณาธิปติ\csromandash{}}ทสฺสเนน ปหาตพฺพเหตุเก ขนฺเธ ครุํ กตฺวา ทสฺสเนน ปหาตพฺพเหตุกา ขนฺธา อุปฺปชฺชนฺติ.\csromanspacer{} \textbf{สหชาตาธิปติ\csromandash{}}ทสฺสเนน ปหาตพฺพเหตุกาธิปติ สมฺปยุตฺตกานํ ขนฺธานํ อธิปติปจฺจเยน ปจฺจโย. (1)}

\prose{ทสฺสเนน ปหาตพฺพเหตุโก ธมฺโม นทสฺสเนน ปหาตพฺพเหตุกสฺส ธมฺมสฺส อธิปติปจฺจเยน ปจฺจโย.\csromanspacer{} \textbf{สหชาตาธิปติ\csromandash{}}ทสฺสเนน ปหาตพฺพเหตุกาธิปติ จิตฺตสมุฏฺฐานานํ รูปานํ อธิปติปจฺจเยน ปจฺจโย. (2)}

\prose{ทสฺสเนน ปหาตพฺพเหตุโก ธมฺโม ทสฺสเนน ปหาตพฺพเหตุกสฺส จ นทสฺสเนน ปหาตพฺพเหตุกสฺส จ ธมฺมสฺส อธิปติปจฺจเยน ปจฺจโย.\csromanspacer{} \textbf{สหชาตาธิปติ\csromandash{}}ทสฺสเนน ปหาตพฺพเหตุกาธิปติ สมฺปยุตฺตกานํ ขนฺธานํ จิตฺตสมุฏฺฐานานญฺจ รูปานํ อธิปติปจฺจเยน ปจฺจโย. (3)}

\proseitem{74}{นทสฺสเนน ปหาตพฺพเหตุโก ธมฺโม นทสฺสเนน ปหาตพฺพเหตุกสฺส ธมฺมสฺส อธิปติปจฺจเยน ปจฺจโย.\csromanspacer{} อารมฺมณาธิปติ, สหชาตาธิปติ.\csromanspacer{}\csromanspacer{}\textbf{อารมฺมณาธิปติ\csromandash{}}ทานํ ฯเปฯ สีลํ ฯเปฯ อุโปสถกมฺมํ กตฺวา ตํ ครุํ กตฺวา ปจฺจเวกฺขติ, อสฺสาเทติ อภินนฺทติ, ตํ ครุํ กตฺวา นทสฺสเนน ปหาตพฺพเหตุโก ราโค อุปฺปชฺชติ.\csromanspacer{} ปุพฺเพ สุจิณฺณานิ ฯเปฯ.\csromanspacer{} ฌานา ฯเปฯ.\csromanspacer{} อริยา มคฺคา วุฏฺฐหิตฺวา มคฺคํ ครุํ กตฺวา ฯเปฯ}

\vspace{\tipitakaparskip}
\csromancenter{ทสฺสเนนปหาตพฺพเหตุกทุก ผลสฺส อธิปติปจฺจเยน ปจฺจโย.\csromanspacer{} จกฺขุํ ฯเปฯ.\csromanspacer{} วตฺถุํ นทสฺสเนน ปหาตพฺพเหตุเก ขนฺเธ ครุํ กตฺวา อสฺสาเทติ อภินนฺทติ, ตํ ครุํ กตฺวา นทสฺสเนน ปหาตพฺพเหตุโก ราโค อุปฺปชฺชติ.\csromanspacer{} \textbf{สหชาตาธิปติ\csromandash{}} นทสฺสเนน ปหาตพฺพเหตุกาธิปติ สมฺปยุตฺตกานํ ขนฺธานํ จิตฺตสมุฏฺฐานานญฺจ รูปานํ อธิปติปจฺจเยน ปจฺจโย. (1)}
\prose{\csromanfolio{69}นทสฺสเนน ปหาตพฺพเหตุโก ธมฺโม ทสฺสเนน ปหาตพฺพเหตุกสฺส ธมฺมสฺส อธิปติปจฺจเยน ปจฺจโย.\csromanspacer{} \textbf{อารมฺมณาธิปติ\csromandash{}}ทานํ ฯเปฯ สีลํ ฯเปฯ อุโปสถกมฺมํ กตฺวา ฯเปฯ.\csromanspacer{} ฌานา ฯเปฯ.\csromanspacer{} จกฺขุํ ฯเปฯ.\csromanspacer{} วตฺถุํ นทสฺสเนน ปหาตพฺพเหตุเก ขนฺเธ ครุํ กตฺวา อสฺสาเทติ อภินนฺทติ, ตํ ครุํ กตฺวา ทสฺสเนน ปหาตพฺพเหตุโก ราโค อุปฺปชฺชติ.\csromanspacer{} ทิฏฺฐิ อุปฺปชฺชติ. (2)}

\csromanheader{2}{\textbf{อนนฺตร}}
\proseitem{75}{ทสฺสเนน ปหาตพฺพเหตุโก ธมฺโม ทสฺสเนน ปหาตพฺพเหตุกสฺส ธมฺมสฺส อนนฺตรปจฺจเยน ปจฺจโย.\csromanspacer{} ปุริมา ปุริมา ทสฺสเนน ปหาตพฺพเหตุกา ขนฺธา ปจฺฉิมานํ ปจฺฉิมานํ ทสฺสเนน ปหาตพฺพเหตุกานํ ขนฺธานํ อนนฺตรปจฺจเยน ปจฺจโย. (มูลํ กาตพฺพํ.) ปุริมา ปุริมา วิจิกิจฺฉาสหคตา ขนฺธา ปจฺฉิมสฺส ปจฺฉิมสฺส วิจิกิจฺฉาสหคตสฺส โมหสฺส อนนฺตรปจฺจเยน ปจฺจโย.\csromanspacer{} ทสฺสเนน ปหาตพฺพเหตุกา ขนฺธา วุฏฺฐานสฺส อนนฺตรปจฺจเยน ปจฺจโย. (มูลํ กาตพฺพํ.) ปุริมา ปุริมา วิจิกิจฺฉาสหคตา ขนฺธา ปจฺฉิมานํ ปจฺฉิมานํ วิจิกิจฺฉาสหคตานํ ขนฺธานํ โมหสฺส จ อนนฺตรปจฺจเยน ปจฺจโย. (3)}

\proseitem{76}{นทสฺสเนน ปหาตพฺพเหตุโก ธมฺโม นทสฺสเนน ปหาตพฺพเหตุกสฺส ธมฺมสฺส อนนฺตรปจฺจเยน ปจฺจโย.\csromanspacer{} ปุริโม ปุริโม วิจิกิจฺฉาสหคโต โมโห ปจฺฉิมสฺส ปจฺฉิมสฺส วิจิกิจฺฉาสหคตสฺส โมหสฺส อนนฺตรปจฺจเยน ปจฺจโย.\csromanspacer{} ปุริมา ปุริมา นทสฺสเนน ปหาตพฺพเหตุกา ขนฺธา ฯเปฯ ผลสมาปตฺติยา อนนฺตรปจฺจเยน ปจฺจโย. (มูลํ กาตพฺพํ.) ปุริโม ปุริโม วิจิกิจฺฉาสหคโต โมโห ปจฺฉิมานํ ปจฺฉิมานํ วิจิกิจฺฉาสหคตานํ ขนฺธานํ อนนฺตรปจฺจเยน ปจฺจโย.\csromanspacer{} อาวชฺชนา ทสฺสเนน ปหาตพฺพเหตุกานํ ขนฺธานํ อนนฺตรปจฺจเยน ปจฺจโย. (มูลํ กาตพฺพํ.) ปุริโม ปุริโม วิจิกิจฺฉาสหคโต โมโห ปจฺฉิมานํ \csromanfolio{70}ปจฺฉิมานํ วิจิกิจฺฉาสหคตานํ ขนฺธานํ โมหสฺส จ อนนฺตรปจฺจเยน ปจฺจโย.\csromanspacer{} อาวชฺชนา วิจิกิจฺฉาสหคตานํ ขนฺธานํ โมหสฺส จ อนนฺตรปจฺจเยน ปจฺจโย. (3)}

\proseitem{77}{ทสฺสเนน ปหาตพฺพเหตุโก จ นทสฺสเนน ปหาตพฺพเหตุโก จ ธมฺมา ทสฺสเนน ปหาตพฺพเหตุกสฺส ธมฺมสฺส อนนฺตรปจฺจเยน ปจฺจโย.\csromanspacer{} ปุริมา ปุริมา วิจิกิจฺฉาสหคตา ขนฺธา จ โมโห จ ปจฺฉิมานํ ปจฺฉิมานํ วิจิกิจฺฉาสหคตานํ ขนฺธานํ อนนฺตรปจฺจเยน ปจฺจโย. (มูลํ กาตพฺพํ.) ปุริมา ปุริมา วิจิกิจฺฉาสหคตา ขนฺธา จ โมโห จ ปจฺฉิมสฺส ปจฺฉิมสฺส วิจิกิจฺฉาสหคตสฺส โมหสฺส อนนฺตรปจฺจเยน ปจฺจโย.\csromanspacer{} วิจิกิจฺฉาสหคตา ขนฺธา จ โมโห จ วุฏฺฐานสฺส อนนฺตรปจฺจเยน ปจฺจโย. (มูลํ กาตพฺพํ.) ปุริมา ปุริมา วิจิกิจฺฉาสหคตา ขนฺธา จ โมโห จ ปจฺฉิมานํ ปจฺฉิมานํ วิจิกิจฺฉาสหคตานํ ขนฺธานํ โมหสฺส จ อนนฺตรปจฺจเยน ปจฺจโย. (3)}

\prose{สมนนฺตรปจฺจเยน ปจฺจโย.\csromanspacer{} นว.\csromanspacer{} สหชาตปจฺจเยน ปจฺจโย.\csromanspacer{} นว.\csromanspacer{} อญฺญมญฺญปจฺจเยน ปจฺจโย.\csromanspacer{} ฉ.\csromanspacer{} นิสฺสยปจฺจเยน ปจฺจโย.\csromanspacer{} นว.}

\csromanheader{2}{\textbf{อุปนิสฺสย}}
\proseitem{78}{ทสฺสเนน ปหาตพฺพเหตุโก ธมฺโม ทสฺสเนน ปหาตพฺพเหตุกสฺส ธมฺมสฺส อุปนิสฺสยปจฺจเยน ปจฺจโย.\csromanspacer{} \textbf{อารมฺมณูปนิสฺสโย, อนนฺตรูปนิสฺสโย,} \textbf{ปกตูปนิสฺสโย ฯเปฯ.}\csromanspacer{} ปกตูปนิสฺสโย\csromandash{}ทสฺสเนน ปหาตพฺพเหตุกา ขนฺธา ทสฺสเนน ปหาตพฺพเหตุกานํ ขนฺธานํ อุปนิสฺสยปจฺจเยน ปจฺจโย. (อวเสเสสุ ทฺวีสุ \textbf{อนนฺตรูปนิสฺสโย,} ปกตูปนิสฺสโย.) (มูลํ กาตพฺพํ.) ทสฺสเนน ปหาตพฺพเหตุกา ขนฺธา นทสฺสเนน ปหาตพฺพเหตุกานํ ขนฺธานํ โมหสฺส จ อุปนิสฺสยปจฺจเยน ปจฺจโย. (มูลํ กาตพฺพํ.) ทสฺสเนน ปหาตพฺพเหตุกา ขนฺธา วิจิกิจฺฉาสหคตานํ ขนฺธานํ โมหสฺส จ อุปนิสฺสยปจฺจเยน ปจฺจโย. (3)}

\proseitem{79}{นทสฺสเนน ปหาตพฺพเหตุโก ธมฺโม นทสฺสเนน ปหาตพฺพเหตุกสฺส ธมฺมสฺส อุปนิสฺสยปจฺจเยน ปจฺจโย.\csromanspacer{} \textbf{อารมฺมณูปนิสฺสโย, อนนฺตรูปนิสฺสโย,} \textbf{ปกตูปนิสฺสโย ฯเปฯ.}\csromanspacer{} ปกตูปนิสฺสโย\csromandash{}สทฺธํ อุปนิสฺสาย ทานํ เทติ ฯเปฯ สมาปตฺติํ อุปฺปาเทติ.\csromanspacer{} มานํ ชปฺเปติ.}

\vspace{\tipitakaparskip}
\csromancenter{ทสฺสเนนปหาตพฺพเหตุกทุก สีลํ ฯเปฯ.\csromanspacer{} ปญฺญํ.\csromanspacer{} นทสฺสเนน ปหาตพฺพเหตุกํ ราคํ.\csromanspacer{} โทสํ.\csromanspacer{} โมหํ.\csromanspacer{} มานํ.\csromanspacer{} ปตฺถนํ.\csromanspacer{} กายิกํ สุขํ.\csromanspacer{} กายิกํ ทุกฺขํ.\csromanspacer{} อุตุํ.\csromanspacer{} โภชนํ.\csromanspacer{} เสนาสนํ อุปนิสฺสาย ทานํ เทติ ฯเปฯ สมาปตฺติํ อุปฺปาเทติ.\csromanspacer{} สทฺธา ฯเปฯ.\csromanspacer{} เสนาสนํ สทฺธาย ฯเปฯ ปญฺญาย.\csromanspacer{} นทสฺสเนน ปหาตพฺพเหตุกสฺส ราคสฺส ฯเปฯ ปตฺถนาย ผลสมาปตฺติยา อุปนิสฺสยปจฺจเยน ปจฺจโย. (1)}
\prose{\csromanfolio{71}นทสฺสเนน ปหาตพฺพเหตุโก ธมฺโม ทสฺสเนน ปหาตพฺพเหตุกสฺส ธมฺมสฺส อุปนิสฺสยปจฺจเยน ปจฺจโย.\csromanspacer{} \textbf{อารมฺมณูปนิสฺสโย, อนนฺตรูปนิสฺสโย,} \textbf{ปกตูปนิสฺสโย ฯเปฯ.}\csromanspacer{} ปกตูปนิสฺสโย\csromandash{}สทฺธํ อุปนิสฺสาย ทิฏฺฐิํ คณฺหาติ.\csromanspacer{} สีลํ ฯเปฯ.\csromanspacer{} ปญฺญํ.\csromanspacer{} นทสฺสเนน ปหาตพฺพเหตุกํ ราคํ.\csromanspacer{} โทสํ.\csromanspacer{} โมหํ.\csromanspacer{} มานํ.\csromanspacer{} ปตฺถนํ.\csromanspacer{} กายิกํ สุขํ.\csromanspacer{} กายิกํ ทุกฺขํ.\csromanspacer{} อุตุํ.\csromanspacer{} โภชนํ.\csromanspacer{} เสนาสนํ อุปนิสฺสาย ปาณํ หนติ ฯเปฯ สํฆํ ภินฺทติ.\csromanspacer{} สทฺธา ฯเปฯ.\csromanspacer{} เสนาสนํ ทสฺสเนน ปหาตพฺพเหตุกสฺส ราคสฺส.\csromanspacer{} โทสสฺส.\csromanspacer{} โมหสฺส.\csromanspacer{} ทิฏฺฐิยา.\csromanspacer{} ปตฺถนาย อุปนิสฺสยปจฺจเยน ปจฺจโย. (2)}

\prose{นทสฺสเนน ปหาตพฺพเหตุโก ธมฺโม ทสฺสเนน ปหาตพฺพเหตุกสฺส จ นทสฺสเนน ปหาตพฺพเหตุกสฺส จ ธมฺมสฺส อุปนิสฺสยปจฺจเยน ปจฺจโย.\csromanspacer{} \textbf{อนนฺตรูปนิสฺสโย, ปกตูปนิสฺสโย ฯเปฯ.}\csromanspacer{} ปกตูปนิสฺสโย\csromandash{} สทฺธา ฯเปฯ ปญฺญา.\csromanspacer{} นทสฺสเนน ปหาตพฺพเหตุโก ราโค.\csromanspacer{} โทโส.\csromanspacer{} โมโห.\csromanspacer{} มาโน.\csromanspacer{} ปตฺถนา.\csromanspacer{} กายิกํ สุขํ ฯเปฯ.\csromanspacer{} เสนาสนํ.\csromanspacer{} วิจิกิจฺฉาสหคตานํ ขนฺธานํ โมหสฺส จ อุปนิสฺสยปจฺจเยน ปจฺจโย. (3)}

\proseitem{80}{ทสฺสเนน ปหาตพฺพเหตุโก จ นทสฺสเนน ปหาตพฺพเหตุโก จ ธมฺมา ทสฺสเนน ปหาตพฺพเหตุกสฺส ธมฺมสฺส อุปนิสฺสยปจฺจเยน ปจฺจโย.\csromanspacer{} \textbf{อนนฺตรูปนิสฺสโย, ปกตูปนิสฺสโย ฯเปฯ.}\csromanspacer{} ปกตูปนิสฺสโย\csromandash{} วิจิกิจฺฉาสหคตา ขนฺธา จ โมโห จ ทสฺสเนน ปหาตพฺพเหตุกานํ ขนฺธานํ อุปนิสฺสยปจฺจเยน ปจฺจโย. (มูลํ กาตพฺพํ.) วิจิกิจฺฉาสหคตา ขนฺธา จ โมโห จ นทสฺสเนน ปหาตพฺพเหตุกานํ ขนฺธานํ โมหสฺส จ อุปนิสฺสยปจฺจเยน ปจฺจโย. (มูลํ กาตพฺพํ.) วิจิกิจฺฉาสหคตา ขนฺธา จ โมโห จ วิจิกิจฺฉาสหคตานํ ขนฺธานํ โมหสฺส จ อุปนิสฺสยปจฺจเยน ปจฺจโย. (3)}

\csromanheader{2}{\textbf{ปุเรชาต}}
\proseitem{81}{\csromanfolio{72}นทสฺสเนน ปหาตพฺพเหตุโก ธมฺโม นทสฺสเนน ปหาตพฺพเหตุกสฺส ธมฺมสฺส ปุเรชาตปจฺจเยน ปจฺจโย.\csromanspacer{} อารมฺมณปุเรชาตํ, วตฺถุปุเรชาตํ.\csromanspacer{}\csromanspacer{}\textbf{อารมฺมณปุเรชาตํ\csromandash{}}จกฺขุํ ฯเปฯ.\csromanspacer{} วตฺถุํ อนิจฺจโต ฯเปฯ วฺปสฺสติ, อสฺสาเทตุ อภินนฺทติ, ตํ อารพฺภ นทสฺสเนน ปหาตพฺพเหตุโก ราโค อุปฺปชฺชติ.\csromanspacer{} อุทฺธจฺจํ อุปฺปชฺชติ.\csromanspacer{} นทสฺสเนน ปหาตพฺพเหตุกํ โทมนสฺสํ อุปฺปชฺชติ.\csromanspacer{} ทิพฺเพน ฯเปฯ. (สํขิตฺตํ.\csromanspacer{} วตฺถุปุเรชาตํ.\csromanspacer{} สํขิตฺตํ.) (1)}

\prose{นทสฺสเนน ปหาตพฺพเหตุโก ธมฺโม ทสฺสเนน ปหาตพฺพเหตุกสฺส ธมฺมสฺส ปุเรชาตปจฺจเยน ปจฺจโย.\csromanspacer{} อารมฺมณปุเรชาตํ, วตฺถุปุเรชาตํ.\csromanspacer{}\csromanspacer{}\textbf{อารมฺมณปุเรชาตํ\csromandash{}}จกฺขุํ ฯเปฯ.\csromanspacer{} วตฺถุํ อสฺสาเทติ อภินนฺทติ, ตํ อารพฺภ ทสฺสเนน ปหาตพฺพเหตุโก ราโค อุปฺปชฺชติ.\csromanspacer{} ทิฏฺฐิ ฯเปฯ วิจิกิจฺฉา ฯเปฯ.\csromanspacer{} ทสฺสเนน ปหาตพฺพเหตุกํ โทมนสฺสํ อุปฺปชฺชติ. (วตฺถุปุเรชาตํ.\csromanspacer{} สํขิตฺตํ.) (2)}

\prose{นทสฺสเนน ปหาตพฺพเหตุโก ธมฺโม ทสฺสเนน ปหาตพฺพเหตุกสฺส จ นทสฺสเนน ปหาตพฺพเหตุกสฺส จ ธมฺมสฺส ปุเรชาตปจฺจเยน ปจฺจโย.\csromanspacer{} อารมฺมณปุเรชาตํ, วตฺถุปุเรชาตํ.\csromanspacer{}\csromanspacer{}\textbf{อารมฺมณปุเรชาตํ\csromandash{}} จกฺขุํ ฯเปฯ.\csromanspacer{} วตฺถุํ อารพฺภ วิจิกิจฺฉาสหคตา ขนฺธา จ โมโห จ อุปฺปชฺชนฺติ. (วตฺถุปุเรชาตํ.\csromanspacer{} สํขิตฺตํ.) (3)}

\csromanheader{2}{\textbf{ปจฺฉาชาตาเสวน}}
\proseitem{82}{ทสฺสเนน ปหาตพฺพเหตุโก ธมฺโม นทสฺสเนน ปหาตพฺพเหตุกสฺส ธมฺมสฺส ปจฺฉาชาตปจฺจเยน ปจฺจโย. (สํกฺขิตฺตํ.).\csromanspacer{}\csromanspacer{}นทสฺสเนน ปหาตพฺพเหตุโก ธมฺโม นทสฺสเนน ปหาตพฺพเหตุกสฺส ธมฺมสฺส ปจฺฉาชาตปจฺจเยน ปจฺจโย. (สํขิตฺตํ.).\csromanspacer{}\csromanspacer{}ทสฺสเนน ปหาตพฺพเหตุโก จ นทสฺสเนน ปหาตพฺพเหตุโก จ ธมฺมา นทสฺสเนน ปหาตพฺพเหตุกสฺส ธมฺมสฺส ปจฺฉาชาตปจฺจเยน ปจฺจโย. (สํขิตฺตํ.) อาเสวนปจฺจเยน ปจฺจโย.}

\csromanheader{2}{\textbf{กมฺมาทิ}}
\proseitem{83}{ทสฺสเนน ปหาตพฺพเหตุโก ธมฺโม ทสฺสเนน ปหาตพฺพเหตุกสฺส ธมฺมสฺส กมฺมปจฺจเยน ปจฺจโย.\csromanspacer{} ทสฺสเนน ปหาตพฺพเหตุกา}

\vspace{\tipitakaparskip}
\csromancenter{ทสฺสเนนปหาตพฺพเหตุกทุก เจตนา สมฺปยุตฺตกานํ ขนฺธานํ กมฺมปจฺจเยน ปจฺจโย. (มูลํ กาตพฺพํ.) \textbf{สหชาตา}, นานากฺขณิกา.\csromanspacer{}\csromanspacer{}\textbf{สหชาตา} ทสฺสเนน ปหาตพฺพเหตุกา เจตนา จิตฺตสมุฏฺฐานานํ รูปานํ กมฺมปจฺจเยน ปจฺจโย.\csromanspacer{} \textbf{นานากฺขณิกา} ทสฺสเนน ปหาตพฺพเหตุกา เจตนา วิปากานํ ขนฺธานํ กฏตฺตา จ รูปานํ กมฺมปจฺจเยน ปจฺจโย.\csromanspacer{}\csromanspacer{}ทสฺสเนน ปหาตพฺพเหตุโก ธมฺโม ทสฺสเนน ปหาตพฺพเหตุกสฺส จ นทสฺสเนน ปหาตพฺพเหตุกสฺส จ ธมฺมสฺส กมฺมปจฺจเยน ปจฺจโย.\csromanspacer{} ทสฺสเนน ปหาตพฺพเหตุกา เจตนา สมฺปยุตฺตกานํ ขนฺธานํ จิตฺตสมุฏฺฐานานญฺจ รูปานํ กมฺมปจฺจเยน ปจฺจโย. (3)}
\prose{\csromanfolio{73}นทสฺสเนน ปหาตพฺพเหตุโก ธมฺโม นทสฺสเนน ปหาตพฺพเหตุกสฺส ธมฺมสฺส กมฺมปจฺจเยน ปจฺจโย.\csromanspacer{} \textbf{สหชาตา}, นานากฺขณิกา.\csromanspacer{}\csromanspacer{}\textbf{สหชาตา} นทสฺสเนน ปหาตพฺพเหตุกา เจตนา สมฺปยุตฺตกานํ ขนฺธานํ จิตฺตสมุฏฺฐานานญฺจ รูปานํ กมฺมปจฺจเยน ปจฺจโย.\csromanspacer{} ปฏิสนฺธิกฺขเณ ฯเปฯ.\csromanspacer{} \textbf{นานากฺขณิกา} นทสฺสเนน ปหาตพฺพเหตุกา เจตนา วิปากานํ ขนฺธานํ กฏตฺตา จ รูปานํ กมฺมปจฺจเยน ปจฺจโย. (1)}

\prose{วิปากปจฺจเยน ปจฺจโย.\csromanspacer{} อาหารปจฺจเยน ปจฺจโย.\csromanspacer{} อินฺทฺริยปจฺจเยน ปจฺจโย.\csromanspacer{} ฌานปจฺจเยน ปจฺจโย.\csromanspacer{} มคฺคปจฺจเยน ปจฺจโย.\csromanspacer{} สมฺปยุตฺตปจฺจเยน ปจฺจโย.\csromanspacer{} ฉ.\csromanspacer{} วิปฺปยุตฺตปจฺจเยน ปจฺจโย.\csromanspacer{} ปญฺจ.}

\csromanheader{2}{\textbf{อตฺถิ-อาทิ}}
\proseitem{84}{ทสฺสเนน ปหาตพฺพเหตุโก ธมฺโม ทสฺสเนน ปหาตพฺพเหตุกสฺส ธมฺมสฺส อตฺถิปจฺจเยน ปจฺจโย.\csromanspacer{} ตีณิ.}

\prose{นทสฺสเนน ปหาตพฺพเหตุโก ธมฺโม นทสฺสเนน ปหาตพฺพเหตุกสฺส ธมฺมสฺส อตฺถิปจฺจเยน ปจฺจโย.\csromanspacer{} สหชาตํ, ปุเรชาตํ, ปจฺฉาชาตํ, อาหารํ, อินฺทฺริยํ. (สํขิตฺตํ.).\csromanspacer{}\csromanspacer{}นทสฺสเนน ปหาตพฺพเหตุโก ธมฺโม ทสฺสเนน ปหาตพฺพเหตุกสฺส ธมฺมสฺส อตฺถิปจฺจเยน ปจฺจโย.\csromanspacer{} สหชาตํ, ปุเรชาตํ. (สํขิตฺตํ.).\csromanspacer{}\csromanspacer{}นทสฺสเนน ปหาตพฺพเหตุโก ธมฺโม ทสฺสเนน ปหาตพฺพเหตุกสฺส จ นทสฺสเนน ปหาตพฺพเหตุกสฺส จ ธมฺมสฺส อตฺถิปจฺจเยน ปจฺจโย.\csromanspacer{} สหชาตํ, ปุเรชาตํ. (สํขิตฺตํ.) (3)}

\proseitem{85}{\csromanfolio{74}ทสฺสเนน ปหาตพฺพเหตุโก จ นทสฺสเนน ปหาตพฺพเหตุโก จ ธมฺมา ทสฺสเนน ปหาตพฺพเหตุกสฺส ธมฺมสฺส อตฺถิปจฺจเยน ปจฺจโย.\csromanspacer{} \textbf{สหชาตํ,} ปุเรชาตํ.\csromanspacer{}\csromanspacer{}สหชาโต ทสฺสเนน ปหาตพฺพเหตุโก เอโก ขนฺโธ จ วตฺถุ จ ติณฺณนฺนํ ขนฺธานํ อตฺถิปจฺจเยน ปจฺจโย ฯเปฯ เทฺว ขนฺธา จ ฯเปฯ.\csromanspacer{} วิจิกิจฺฉาสหคโต เอโก ขนฺโธ จ โมโห จ ติณฺณนฺนํ ขนฺธานํ อตฺถิปจฺจเยน ปจฺจโย ฯเปฯ เทฺว ขนฺธา จ ฯเปฯ. (1)}

\prose{ทสฺสเนน ปหาตพฺพเหตุโก จ นทสฺสเนน ปหาตพฺพเหตุโก จ ธมฺมา นทสฺสเนน ปหาตพฺพเหตุกสฺส ธมฺมสฺส อตฺถิปจฺจเยน ปจฺจโย.\csromanspacer{} \textbf{สหชาตํ,} ปุเรชาตํ, ปจฺฉาชาตํ, อาหารํ, อินฺทฺริยํ. (สํขิตฺตํ.) (2)}

\prose{ทสฺสเนน ปหาตพฺพเหตุโก จ นทสฺสเนน ปหาตพฺพเหตุโก จ ธมฺมา ทสฺสเนน ปหาตพฺพเหตุกสฺส จ นทสฺสเนน ปหาตพฺพเหตุกสฺส จ ธมฺมสฺส อตฺถิปจฺจเยน ปจฺจโย.\csromanspacer{} \textbf{สหชาตํ,} ปุเรชาตํ. (สํขิตฺตํ.) (3)}

\prose{นตฺถิปจฺจเยน ปจฺจโย.\csromanspacer{} วิคตปจฺจเยน ปจฺจโย.\csromanspacer{} อวิคตปจฺจเยน ปจฺจโย.}

\csromanheader{2}{1. ปจฺจยานุโลม 2. สงฺขฺยาวาร}
\proseitem{86}{เหตุยา ฉ, อารมฺมเณ นว, อธิปติยา ปญฺจ, อนนฺตเร นว, สมนนฺตเร นว, สหชาเต นว, อญฺญมญฺเญ ฉ, นิสฺสเย นว, อุปนิสฺสเย นว, ปุเรชาเต ตีณิ, ปจฺฉาชาเต ตีณิ, อาเสวเน นว, กมฺเม จตฺตาริ, วิปาเก เอกํ, อาหาเร จตฺตาริ, อินฺทฺริเย จตฺตาริ, ฌาเน จตฺตาริ, มคฺเค จตฺตาริ, สมฺปยุตฺเต ฉ, วิปฺปยุตฺเต ปญฺจ, อตฺถิยา นว, นตฺถิยา นว, วิคเต นว, อวิคเต นว.}

\csromanheader{2}{2. ปจฺจนียุทฺธาร}
\proseitem{87}{ทสฺสเนน ปหาตพฺพเหตุโก ธมฺโม ทสฺสเนน ปหาตพฺพเหตุกสฺส ธมฺมสฺส อารมฺมณปจฺจเยน ปจฺจโย.\csromanspacer{} สหชาตปจฺจเยน ปจฺจโย.\csromanspacer{} อุปนิสฺสยปจฺจเยน ปจฺจโย.\csromanspacer{}\csromanspacer{}ทสฺสเนน ปหาตพฺพเหตุโก ธมฺโม นทสฺสเนน ปหาตพฺพเหตุกสฺส ธมฺมสฺส อารมฺมณปจฺจเยน ปจฺจโย.\csromanspacer{} สหชาตปจฺจเยน ปจฺจโย.\csromanspacer{} อุปนิสฺสยปจฺจเยน ปจฺจโย.}

\vspace{\tipitakaparskip}
\csromancenter{ทสฺสเนนปหาตพฺพเหตุกทุก ปจฺฉาชาตปจฺจเยน ปจฺจโย.\csromanspacer{} กมฺมปจฺจเยน ปจฺจโย.\csromanspacer{}\csromanspacer{}ทสฺสเนน ปหาตพฺพเหตุโก ธมฺโม ทสฺสเนน ปหาตพฺพเหตุกสฺส จ นทสฺสเนน ปหาตพฺพเหตุกสฺส จ ธมฺมสฺส อารมฺมณปจฺจเยน ปจฺจโย.\csromanspacer{} สหชาตปจฺจเยน ปจฺจโย.\csromanspacer{} อุปนิสฺสยปจฺจเยน ปจฺจโย. (3)}
\proseitem{88}{\csromanfolio{75}นทสฺสเนน ปหาตพฺพเหตุโก ธมฺโม นทสฺสเนน ปหาตพฺพเหตุกสฺส ธมฺมสฺส อารมฺมณปจฺจเยน ปจฺจโย.\csromanspacer{} สหชาตปจฺจเยน ปจฺจโย.\csromanspacer{} อุปนิสฺสยปจฺจเยน ปจฺจโย.\csromanspacer{} ปุเรชาตปจฺจเยน ปจฺจโย.\csromanspacer{} ปจฺฉาชาตปจฺจเยน ปจฺจโย.\csromanspacer{} กมฺมปจฺจเยน ปจฺจโย.\csromanspacer{} อาหารปจฺจเยน ปจฺจโย.\csromanspacer{} อินฺทฺริยปจฺจเยน ปจฺจโย.\csromanspacer{}\csromanspacer{}นทสฺสเนน ปหาตพฺพเหตุโก ธมฺโม ทสฺสเนน ปหาตพฺพเหตุกสฺส ธมฺมสฺส อารมฺมณปจฺจเยน ปจฺจโย.\csromanspacer{} สหชาตปจฺจเยน ปจฺจโย.\csromanspacer{} อุปนิสฺสยปจฺจเยน ปจฺจโย.\csromanspacer{} ปุเรชาตปจฺจเยน ปจฺจโย.\csromanspacer{}\csromanspacer{}นทสฺสเนน ปหาตพฺพเหตุโก ธมฺโม ทสฺสเนน ปหาตพฺพเหตุกสฺส จ นทสฺสเนน ปหาตพฺพเหตุกสฺส จ ธมฺมสฺส อารมฺมณปจฺจเยน ปจฺจโย.\csromanspacer{} สหชาตปจฺจเยน ปจฺจโย.\csromanspacer{} อุปนิสฺสยปจฺจเยน ปจฺจโย.\csromanspacer{} ปุเรชาตปจฺจเยน ปจฺจโย. (3)}

\proseitem{89}{ทสฺสเนน ปหาตพฺพเหตุโก จ นทสฺสเนน ปหาตพฺพเหตุโก จ ธมฺมา ทสฺสเนน ปหาตพฺพเหตุกสฺส ธมฺมสฺส อารมฺมณปจฺจเยน ปจฺจโย.\csromanspacer{} สหชาตปจฺจเยน ปจฺจโย.\csromanspacer{} อุปนิสฺสยปจฺจเยน ปจฺจโย.\csromanspacer{}\csromanspacer{}ทสฺสเนน ปหาตพฺพเหตุโก จ นทสฺสเนน ปหาตพฺพเหตุโก จ ธมฺมา นทสฺสเนน ปหาตพฺพเหตุกสฺส ธมฺมสฺส อารมฺมณปจฺจเยน ปจฺจโย.\csromanspacer{} สหชาตปจฺจเยน ปจฺจโย.\csromanspacer{} อุปนิสฺสยปจฺจเยน ปจฺจโย.\csromanspacer{} ปจฺฉาชาตปจฺจเยน ปจฺจโย.\csromanspacer{}\csromanspacer{}ทสฺสเนน ปหาตพฺพเหตุโก จ นทสฺสเนน ปหาตพฺพเหตุโก จ ธมฺมา ทสฺสเนน ปหาตพฺพเหตุกสฺส จ นทสฺสเนน ปหาตพฺพเหตุกสฺส จ ธมฺมสฺส อรมฺมณปจฺจเยน ปจฺจโย.\csromanspacer{} สหชาตปจฺจเยน ปจฺจโย.\csromanspacer{} อุปนิสฺสยปจฺจเยน ปจฺจโย. (3)}

\csromanheader{2}{2. ปจฺจยปจฺจนีย 2. สงฺขฺยาวาร}
\vspace{\tipitakaparskip}
\csromancenter{นเหตุยา นว, น-อารมฺมเณ นว, (สพฺพตฺถ นว,) โน-อวิคเต นว.}
\csromanheader{2}{3. ปจฺจยานุโลมปจฺจนีย}
\proseitem{91}{\csromanfolio{76}\textbf{เหตุ}ปจฺจยา น-อารมฺมเณ ฉ, น-อธิปติยา ฉ, น-อนนฺตเร ฉ, นสมนนฺตเร ฉ, น-อญฺญมญฺเญ เทฺว, น-อุปนิสฺสเย ฉ ฯเปฯ นฺสมฺปยุตฺเต เทฺว, นวิปฺปยุตฺเต ตีณิ, โนนตฺถิยา ฉ, โนวิคเต ฉ.}

\csromanheader{2}{4. ปจฺจยปจฺจนียานุโลม}
\proseitem{92}{นเหตุปจฺจยา อารมฺมเณ นว, อธิปติยา ปญฺจ, (อนุโลมมาติกา.) ฯเปฯ อวิคเต นว.}

\nitthitam{ทสฺสเนน ปหาตพฺพเหตุกทุกํ นิฏฺฐิตํ.}
\csromanmatikaanchor{mtk.15}
\csromantocmarkref{section}{86. ภาวนายปหาตพฺพเหตุกทุก}{mtk.15}
\bookmark[dest={mtk.15},level=1]{86. ภาวนายปหาตพฺพเหตุกทุก}
\csromanheader{1}{86. ภาวนายปหาตพฺพเหตุกทุก 1. ปฏิจฺจวาร}
\proseitem{93}{ภาวนาย ปหาตพฺพเหตุกํ ธมฺมํ ปฏิจฺจ ภาวนาย ปหาตพฺพเหตุโก ธมฺโม อุปฺปชฺชติ เหตุปจฺจยา.\csromanspacer{} ภาวนาย ปหาตพฺพเหตุกํ เอกํ ขนฺธํ ปฏิจฺจ ตโย ขนฺธา ฯเปฯ เทฺว ขนฺเธ ฯเปฯ. (เอวํ ปฏิจฺจวาโรปิ สหชาตวาโรปิ ปจฺจยวาโรปิ นิสฺสยวาโรปิ สํสฏฺฐวาโรปิ สมฺปยุตฺตวาโรปิ ทสฺสเนนปหาตพฺพเหตุกทุก- สทิสา.\csromanspacer{} อุทฺธจฺจสหคโต โมโห วิจิกิจฺฉาสหคตโมหฏฺฐาเน ฐเปตพฺโพ.)}

\csromanheader{2}{86. ภาวนายปหาตพฺพเหตุกทุก 7. ปญฺหาวาร}
\proseitem{94}{ภาวนาย ปหาตพฺพเหตุโก ธมฺโม ภาวนาย ปหาตพฺพเหตุกสฺส ธมฺมสฺส เหตุปจฺจเยน ปจฺจโย.\csromanspacer{} ฉ ฯเปฯ. (ทสฺสเนน ปหาตพฺพเหตุกทุกสทิสา.)}

\csromanheader{2}{86. ภาวนายปหาตพฺพเหตุกทุก \textbf{อารมฺมณ}}
\proseitem{95}{\csromanfolio{77}ภาวนาย ปหาตพฺพเหตุโก ธมฺโม ภาวนาย ปหาตพฺพเหตุกสฺส ธมฺมสฺส อารมฺมณปจฺจเยน ปจฺจโย.\csromanspacer{} ตีณิ. (อารพฺภ.\csromanspacer{} ทสฺสเนน ปหาตพฺพเหตุกทุกสทิสา.)}

\proseitem{96}{นภาวนาย ปหาตพฺพเหตุโก ธมฺโม นภาวนาย ปหาตพฺพเหตุกสฺส ธมฺมสฺส อารมฺมณปจฺจเยน ปจฺจโย.\csromanspacer{} ทานํ ฯเปฯ สีลํ ฯเปฯ อุโปสถกมฺมํ กตฺวา ตํ ปจฺจเวกฺขติ, อสฺสาเทติ อภินนฺทติ, ตํ อารพฺภ นภาวนาย ปหาตพฺพเหตุโก ราโค อุปฺปชฺชติ.\csromanspacer{} ทิฏฺฐิ ฯเปฯ วิจิกิจฺฉา ฯเปฯ.\csromanspacer{} นภาวนาย ปหาตพฺพเหตุกํ โทมนสฺสํ อุปฺปชฺชติ.\csromanspacer{} ปุพฺเพ สุจิณฺณานิ ฯเปฯ.\csromanspacer{} ฌานา ฯเปฯ.\csromanspacer{} อริยา มคฺคา ฯเปฯ ผลสฺส, อาวชฺชนาย อารมฺมณปจฺจเยน ปจฺจโย.\csromanspacer{} อริยา นภาวนาย ปหาตพฺพเหตุเก ปหีเน กิเลเส ฯเปฯ.\csromanspacer{} ปุพฺเพ สมุทาจิณฺเณ ฯเปฯ.\csromanspacer{} จกฺขุํ ฯเปฯ.\csromanspacer{} วตฺถุํ นภาวนาย ปหาตพฺพเหตุเก ขนฺเธ จ โมหญฺจ อนิจฺจโต ฯเปฯ วิปสฺสติ, อสฺสาเทติ อภินนฺทติ, ตํ อารพฺภ นภาวนาย ปหาตพฺพเหตุโก ราโค อุปฺปชฺชติ.\csromanspacer{} ทิฏฺฐิ ฯเปฯ.\csromanspacer{} วิจิกิจฺฉา ฯเปฯ.\csromanspacer{} นภาวนาย ปหาตพฺพเหตุกํ โทมนสฺสํ อุปฺปชฺชติ.\csromanspacer{} ทิพฺเพน จกฺขุนา รูปํ ปสฺสติ ฯเปฯ อนาคตํสญาณสฺส, อาวชฺชนาย, โมหสฺส จ อารมฺมณปจฺจเยน ปจฺจโย. (1)}

\prose{นภาวนาย ปหาตพฺพเหตุโก ธมฺโม ภาวนาย ปหาตพฺพเหตุกสฺส ธมฺมสฺส อารมฺมณปจฺจเยน ปจฺจโย.\csromanspacer{} ทานํ ฯเปฯ.\csromanspacer{} สีลํ ฯเปฯ.\csromanspacer{} ฌานํ ฯเปฯ.\csromanspacer{} จกฺขุํ ฯเปฯ.\csromanspacer{} วตฺถุํ นภาวนาย ปหาตพฺพเหตุเก ขนฺเธ จ โมหญฺจ อสฺสาเทติ อภินนฺทติ, ตํ อารพฺภ ภาวนาย ปหาตพฺพเหตุโก ราโค อุปฺปชฺชติ.\csromanspacer{} อุทฺธจฺจํ ฯเปฯ.\csromanspacer{} ภาวนาย ปหาตพฺพเหตุกํ โทมนสฺสํ อุปฺปชฺชติ. (2)}

\prose{นภาวนาย ปหาตพฺพเหตุโก ธมฺโม ภาวนาย ปหาตพฺพเหตุกสฺส จ นภาวนาย ปหาตพฺพเหตุกสฺส จ ธมฺมสฺส อารมฺมณปจฺจเยน ปจฺจโย.\csromanspacer{} จกฺขุํ ฯเปฯ.\csromanspacer{} วตฺถุํ นภาวนาย ปหาตพฺพเหตุเก ขนฺเธ จ โมหญฺจ อารพฺภ อุทฺธจฺจสหคตา ขนฺธา จ โมโห จ อุปฺปชฺชนฺติ. (ฆฏนารมฺมณา ตีณิปิ กาตพฺพา.) (3)}

\csromanheader{2}{\textbf{อธิปตฺยาทิ}}
\proseitem{97}{\csromanfolio{78}ภาวนาย ปหาตพฺพเหตุโก ธมฺโม ภาวนาย ปหาตพฺพเหตุกสฺส ธมฺมสฺส อธิปติปจฺจเยน ปจฺจโย.\csromanspacer{} อารมฺมณาธิปติ, สหชาตาธิปติ.\csromanspacer{}\csromanspacer{}\textbf{อารมฺมณาธิปติ\csromandash{}}ภาวนาย ปหาตพฺพเหตุกํ ราคํ ครุํ กตฺวา อสฺสาเทติ อภินนฺทติ, ตํ ครุํ กตฺวา ภาวนาย ปหาตพฺพเหตุโก ราโค อุปฺปชฺชติ.\csromanspacer{} \textbf{สหชาตาธิปติ\csromandash{}} ภาวนาย ปหาตพฺพเหตุกาธิปติ สมฺปยุตฺตกานํ ขนฺธานํ อธิปติปจฺจเยน ปจฺจโย. (1)}

\prose{ภาวนาย ปหาตพฺพเหตุโก ธมฺโม นภาวนาย ปหาตพฺพเหตุกสฺส ธมฺมสฺส อธิปติปจฺจเยน ปจฺจโย.\csromanspacer{} อารมฺมณาธิปติ, สหชาตาธิปติ.\csromanspacer{}\csromanspacer{}\textbf{อารมฺมณาธิปติ\csromandash{}}ภาวนาย ปหาตพฺพเหตุกํ ราคํ ครุํ กตฺวา อสฺสาเทติ อภินนฺทติ, ตํ ครุํ กตฺวา นภาวนาย ปหาตพฺพเหตุโก ราโค อุปฺปชฺชติ.\csromanspacer{} ทิฏฺฐิ อุปฺปชฺชติ.\csromanspacer{} \textbf{สหชาตาธิปติ\csromandash{}}ภาวนาย ปหาตพฺพเหตุกาธิปติ จิตฺตสมุฏฺฐานานํ รูปานํ อธิปติปจฺจเยน ปจฺจโย. (2)}

\prose{ภาวนาย ปหาตพฺพเหตุโก ธมฺโม ภาวนาย ปหาตพฺพเหตุกสฺส จ นภาวนาย ปหาตพฺพเหตุกสฺส จ ธมฺมสฺส อธิปติปจฺจเยน ปจฺจโย.\csromanspacer{} \textbf{สหชาตาธิปติ\csromandash{}}ภาวนาย ปหาตพฺพเหตุกาธิปติ สมฺปยุตฺตกานํ ขนฺธานํ จิตฺตสมุฏฺฐานานญฺจ รูปานํ อธิปติปจฺจเยน ปจฺจโย. (3)}

\proseitem{98}{นภาวนาย ปหาตพฺพเหตุโก ธมฺโม นภาวนาย ปหาตพฺพเหตุกสฺส ธมฺมสฺส อธิปติปจฺจเยน ปจฺจโย.\csromanspacer{} อารมฺมณาธิปติ, สหชาตาธิปติ.\csromanspacer{}\csromanspacer{}\textbf{อารมฺมณาธิปติ\csromandash{}}ทานํ ฯเปฯ สีลํ ฯเปฯ อุโปสถกมฺมํ กตฺวา ตํ ครุํ กตฺวา ปจฺจเวกฺขติ, อสฺสาเทติ อภินนฺทติ, ตํ ครุํ กตฺวา นภาวนาย ปหาตพฺพเหตุโก ราโค อุปฺปชฺชติ.\csromanspacer{} ทิฏฺฐิ อุปฺปชฺชติ.\csromanspacer{} ปุพฺเพ ฯเปฯ.\csromanspacer{} ฌานา ฯเปฯ.\csromanspacer{} อริยา มคฺคา วุฏฺฐหิตฺวา ฯเปฯ ผลสฺส อธิปติปจฺจเยน ปจฺจโย.\csromanspacer{} จกฺขุํ ฯเปฯ.\csromanspacer{} วตฺถุํ นภาวนาย ปหาตพฺพเหตุเก ขนฺเธ ครุํ กตฺวา อสฺสาเทติ อภินนฺทติ, ตํ ครุํ กตฺวา นภาวนาย ปหาตพฺพเหตุโก ราโค อุปฺปชฺชติ.\csromanspacer{} ทิฏฺฐิ อุปฺปชฺชติ.\csromanspacer{} \textbf{สหชาตาธิปติ\csromandash{}}นภาวนาย ปหาตพฺพเหตุกาธิปติ สมฺปยุตฺตกานํ ขนฺธานํ จิตฺตสมุฏฺฐานานญฺจ รูปานํ อธิปติปจฺจเยน ปจฺจโย. (1)}

\csromanheader{2}{86. ภาวนายปหาตพฺพเหตุกทุก}
\prose{\csromanfolio{79}นภาวนาย ปหาตพฺพเหตุโก ธมฺโม ภาวนาย ปหาตพฺพเหตุกสฺส ธมฺมสฺส อธิปติปจฺจเยน ปจฺจโย.\csromanspacer{} \textbf{อารมฺมณาธิปติ\csromandash{}}ทานํ ฯเปฯ ฌานํ ฯเปฯ.\csromanspacer{} จกฺขุํ ฯเปฯ.\csromanspacer{} วตฺถุํ นภาวนาย ปหาตพฺพเหตุเก ขนฺเธ ครุํ กตฺวา อสฺสาเทติ อภินนฺทติ, ตํ ครุํ กตฺวา ภาวนาย ปหาตพฺพเหตุโก ราโค อุปฺปชฺชติ. (2)}

\prose{(อนนฺตรปจฺจเย นภาวนาย ปหาตพฺพเหตุกการณา วิจิกิจฺฉาสหคโต โมโห น กาตพฺโพ, อุทฺธจฺจสหคโต โมโห กาตพฺโพ.) สมนนฺตรปจฺจเยน ปจฺจโย.\csromanspacer{} สหชาตปจฺจเยน ปจฺจโย.\csromanspacer{} นว.\csromanspacer{} อญฺญมญฺญปจฺจเยน ปจฺจโย ฉ.\csromanspacer{} นิสฺสยปจฺจเยน ปจฺจโย.\csromanspacer{} นว.}

\csromanheader{2}{\textbf{อุปนิสฺสย}}
\proseitem{99}{ภาวนาย ปหาตพฺพเหตุโก ธมฺโม ภาวนาย ปหาตพฺพเหตุกสฺส ธมฺมสฺส อุปนิสฺสยปจฺจเยน ปจฺจโย.\csromanspacer{} อารมฺมณูปนิสฺสโย, อนนฺตรูปนิสฺสโย, ปกตูปนิสฺสโย ฯเปฯ.\csromanspacer{} ปกตูปนิสฺสโย\csromandash{}ภาวนาย ปหาตพฺพเหตุกา ขนฺธา ภาวนาย ปหาตพฺพเหตุกานํ ขนฺธานํ อุปนิสยปจฺจเยน ปจฺจโย. (มูลํ กาตพฺพํ.) ภาวนาย ปหาตพฺพเหตุกา ขนฺธา นภาวนาย ปหาตพฺพเหตุกานํ ขนฺธานํ โมหสฺส จ อุปนิสฺสยปจฺจเยน ปจฺจโย.\csromanspacer{} สกภณฺเฑ ฉนฺทราโค ปรภณฺเฑ ฉนฺทราคสฺส อุปนิสฺสยปจฺจเยน ปจฺจโย.\csromanspacer{} สกปริคฺคเห ฉนฺทราโค ปรปริคฺคเห ฉนฺทราคสฺส อุปนิสฺสยปจฺจเยน ปจฺจโย. (มูลํ กาตพฺพํ.) ภาวนาย ปหาตพฺพเหตุกา ขนฺธา อุทฺธจฺจสหคตานํ ขนฺธานํ โมหสฺส จ อุปนิสฺสยปจฺจเยน ปจฺจโย. (3)}

\proseitem{100}{นภาวนาย ปหาตพฺพเหตุโก ธมฺโม นภาวนาย ปหาตพฺพเหตุกสฺส ธมฺมสฺส อุปนิสฺสยปจฺจเยน ปจฺจโย.\csromanspacer{} อารมฺมณูปนิสฺสโย, อนนฺตรูปนิสฺสโย, ปกตูปนิสฺสโย ฯเปฯ.\csromanspacer{} ปกตูปนิสฺสโย\csromandash{}สทฺธํ อุปนิสฺสาย ทานํ เทติ ฯเปฯ.\csromanspacer{} สมาปตฺติํ อุปฺปาเทติ.\csromanspacer{} ทิฏฺฐิํ คณฺหาติ.\csromanspacer{} สีลํ ฯเปฯ.\csromanspacer{} ปญฺญํ.\csromanspacer{} นภาวนาย ปหาตพฺพเหตุกํ ราคํ.\csromanspacer{} โทสํ.\csromanspacer{} โมหํ.\csromanspacer{} ทิฏฺฐิํ.\csromanspacer{} ปตฺถนํ.\csromanspacer{} กายิกํ สุขํ.\csromanspacer{} กายิกํ ทุกฺขํ ฯเปฯ.\csromanspacer{} เสนาสนํ อุปนิสฺสาย ทานํ เทติ ฯเปฯ ปาณํ หนติ ฯเปฯ สํฆํ ภินฺทติ.\csromanspacer{} สทฺธา ฯเปฯ.\csromanspacer{} เสนาสนํ สทฺธาย ฯเปฯ ปญฺญาย.\csromanspacer{} นภาวนาย ปหาตพฺพเหตุกสฺส \csromanfolio{80}ราคสฺส.\csromanspacer{} โทสสฺส.\csromanspacer{} โมหสฺส.\csromanspacer{} ทิฏฺฐิยา.\csromanspacer{} ปตฺถนาย กายิกสฺส สุขสฺส ฯเปฯ ผลสมาปตฺติยา อุปนิสฺสยปจฺจเยน ปจฺจโย. (1)}

\prose{นภาวนาย ปหาตพฺพเหตุโก ธมฺโม ภาวนาย ปหาตพฺพเหตุกสฺส ธมฺมสฺส อุปนิสฺสยปจฺจเยน ปจฺจโย.\csromanspacer{} อารมฺมณูปนิสฺสโย, อนนฺตรูปนิสฺสโย, ปกตูปนิสฺสโย ฯเปฯ.\csromanspacer{} ปกตูปนิสฺสโย\csromandash{}สทฺธํ อุปนิสฺสาย มานํ ชปฺเปติ ฯเปฯ.\csromanspacer{} สทฺธา ฯเปฯ.\csromanspacer{} เสนาสนํ ภาวนาย ปหาตพฺพเหตุกสฺส ราคสฺส.\csromanspacer{} โทสสฺส.\csromanspacer{} โมหสฺส.\csromanspacer{} มานสฺส.\csromanspacer{} ปตฺถนาย อุปนิสฺสยปจฺจเยน ปจฺจโย. (2)}

\prose{นภาวนาย ปหาตพฺพเหตุโก ธมฺโม ภาวนาย ปหาตพฺพเหตุกสฺส จ นภาวนาย ปหาตพฺพเหตุกสฺส จ ธมฺมสฺส อุปนิสฺสยปจฺจเยน ปจฺจโย.\csromanspacer{} อารมฺมณูปนิสฺสโย, อนนฺตรูปนิสฺสโย, ปกตูปนิสฺสโย ฯเปฯ.\csromanspacer{} ปกตูปนิสฺสโย\csromandash{}สทฺธา ฯเปฯ ปญฺญา.\csromanspacer{} กายิกํ สุขํ ฯเปฯ.\csromanspacer{} เสนาสนํ อุทฺธจฺจสหคตานํ ขนฺธานํ โมหสฺส จ อุปนิสฺสยปจฺจเยน ปจฺจโย. (ฆฏนูปนิสฺสยาปิ ตีณิปิ กาตพฺพา.) (3)}

\csromanheader{2}{\textbf{ปุเรชาตาทิ}}
\proseitem{101}{นภาวนาย ปหาตพฺพเหตุโก ธมฺโม นภาวนาย ปหาตพฺพเหตุกสฺส ธมฺมสฺส ปุเรชาตปจฺจเยน ปจฺจโย.\csromanspacer{} ตีณิ.\csromanspacer{} ปจฺฉาชาตปจฺจเยน ปจฺจโย.\csromanspacer{} ตีณิ.\csromanspacer{} อาเสวนปจฺจเยน ปจฺจโย.\csromanspacer{} นว.\csromanspacer{} กมฺมปจฺจเยน ปจฺจโย. (นภาวนาย ปหาตพฺพภาชนการเณ นานากฺขณิกา ลพฺภติ.) ฯเปฯ.\csromanspacer{} โนวิคตปจฺจเยน ปจฺจโย. (สํขิตฺตํ.\csromanspacer{} ยถา ทสฺสเนน ปหาตพฺพเหตุกทุกํ, เอวํ ภาวนาย ปหาตพฺพเหตุกปจฺจยาปิ ปจฺจนียาปิ วิภาโคปิ คณนาปิ นินฺนานากรณา.)}

\prose{นทสฺสเนน ปหาตพฺโพ ธมฺโม นทสฺสเนน ปหาตพฺพสฺส ธมฺมสฺส ฯเปฯ. (ปรนฺเตน สกภณฺฑฉนฺทราโคปิ กาตพฺโพ.)}

\prose{ภาวนาย ปหาตพฺโพ ธมฺโม นภาวนาย ปหาตพฺพสฺส ธมฺมสฺส ฯเปฯ. (ปรนฺเตน “สกภณฺฑฉนฺทราโค”ติ กาตพฺพํ.)}

\nitthitam{ภาวนาย ปหาตพฺพเหตุกทุกํ นิฏฺฐิตํ.}
\csromanmatikaanchor{mtk.16}
\csromantocmarkref{section}{87. สวิตกฺกทุก}{mtk.16}
\bookmark[dest={mtk.16},level=1]{87. สวิตกฺกทุก}
\csromanheader{1}{87. สวิตกฺกทุก \textbf{87. สวิตกฺกทุก 1. ปฏิจฺจวาร}}
\proseitem{102}{\csromanfolio{81}สวิตกฺกํ ธมฺมํ ปฏิจฺจ สวิตกฺโก ธมฺโม อุปฺปชฺชติ เหตุปจฺจยา.\csromanspacer{} สวิตกฺกํ เอกํ ขนฺธํ ปฏิจฺจ ตโย ขนฺธา ฯเปฯ เทฺว ขนฺเธ ฯเปฯ.\csromanspacer{} ปฏิสนฺธิกฺขเณ ฯเปฯ.\csromanspacer{} สวิตกฺกํ ธมฺมํ ปฏิจฺจ อวิตกฺโก ธมฺโม อุปฺปชฺชติ เหตุปจฺจยา.\csromanspacer{} สวิตกฺเก ขนฺเธ ปฏิจฺจ วิตกฺโก จิตฺตสมุฏฺฐานญฺจ รูปํ.\csromanspacer{} ปฏิสนฺธิกฺขเณ ฯเปฯ.\csromanspacer{} สวิตกฺกํ ธมฺมํ ปฏิจฺจ สวิตกฺโก จ อวิตกฺโก จ ธมฺมา อุปฺปชฺชนฺติ เหตุปจฺจยา.\csromanspacer{} สวิตกฺกํ เอกํ ขนฺธํ ปฏิจฺจ ตโย ขนฺธา วิตกฺโก จ จิตฺตสมุฏฺฐานญฺจ รูปํ ฯเปฯ เทฺว ขนฺเธ ฯเปฯ.\csromanspacer{} ปฏิสนฺธิกฺขเณ ฯเปฯ. (3)}

\proseitem{103}{อวิตกฺกํ ธมฺมํ ปฏิจฺจ อวิตกฺโก ธมฺโม อุปฺปชฺชติ เหตุปจฺจยา.\csromanspacer{} อวิตกฺกํ เอกํ ขนฺธํ ปฏิจฺจ ตโย ขนฺธา จิตฺตสมุฏฺฐานญฺจ รูปํ ฯเปฯ เทฺว ขนฺเธ ฯเปฯ.\csromanspacer{} วิตกฺกํ ปฏิจฺจ จิตฺตสมุฏฺฐานํ รูปํ.\csromanspacer{} ปฏิสนฺธิกฺขเณ อวิตกฺกํ เอกํ ขนฺธํ ปฏิจฺจ ตโย ขนฺธา กฏตฺตา จ รูปํ ฯเปฯ เทฺว ขนฺเธ ฯเปฯ.\csromanspacer{} วิตกฺกํ ปฏิจฺจ กฏตฺตารูปํ.\csromanspacer{} ขนฺเธ ปฏิจฺจ วตฺถุ, วตฺถุํ ปฏิจฺจ ขนฺธา.\csromanspacer{} วิตกฺกํ ปฏิจฺจ วตฺถุ, วตฺถุํ ปฏิจฺจ วิตกฺโก.\csromanspacer{} เอกํ มหาภูตํ ฯเปฯ. (1)}

\prose{อวิตกฺกํ ธมฺมํ ปฏิจฺจ สวิตกฺโก ธมฺโม อุปฺปชฺชติ เหตุปจฺจยา.\csromanspacer{} วิตกฺกํ ปฏิจฺจ สมฺปยุตฺตกา ขนฺธา.\csromanspacer{} ปฏิสนฺธิกฺขเณ วิตกฺกํ ปฏิจฺจ สมฺปยุตฺตกา ขนฺธา.\csromanspacer{} ปฏิสนฺธิกฺขเณ วตฺถุํ ปฏิจฺจ สวิตกฺกา ขนฺธา. (2)}

\prose{อวิตกฺกํ ธมฺมํ ปฏิจฺจ สวิตกฺโก จ อวิตกฺโก จ ธมฺมา อุปฺปชฺชนฺติ เหตุปจฺจยา.\csromanspacer{} วิตกฺกํ ปฏิจฺจ สมฺปยุตฺตกา ขนฺธา จิตฺตสมุฏฺฐานญฺจ รูปํ.\csromanspacer{} วิตกฺกํ ปฏิจฺจ สวิตกฺกา ขนฺธา.\csromanspacer{} มหาภูเต ปฏิจฺจ จิตฺตสมุฏฺฐานํ รูปํ.\csromanspacer{} ปฏิสนฺธิกฺขเณ วิตกฺกํ ปฏิจฺจ สมฺปยุตฺตกา ขนฺธา กฏตฺตา จ รูปํ.\csromanspacer{} ปฏิสนฺธิกฺขเณ วิตกฺกํ ปฏิจฺจ สวิตกฺกา ขนฺธา.\csromanspacer{} มหาภูเต ปฏิจฺจ กฏตฺตารูปํ.\csromanspacer{} ปฏิสนฺธิกฺขเณ วตฺถุํ ปฏิจฺจ สวิตกฺกา ขนฺธา.\csromanspacer{} มหาภูเต ปฏิจฺจ กฏตฺตารูปํ.\csromanspacer{} ปฏิสนฺธิกฺขเณ วตฺถุํ ปฏิจฺจ วิตกฺโก สมฺปยุตฺตกา จ ขนฺธา. (3)}

\proseitem{104}{สวิตกฺกญฺจ อวิตกฺกญฺจ ธมฺมํ ปฏิจฺจ สวิตกฺโก ธมฺโม อุปฺปชฺชติ เหตุปจฺจยา.\csromanspacer{} สวิตกฺกํ เอกํ ขนฺธญฺจ วิตกฺกญฺจ ปฏิจฺจ ตโย ขนฺธา ฯเปฯ เทฺว \csromanfolio{82}ขนฺเธ ฯเปฯ.\csromanspacer{} ปฏิสนฺธิกฺขเณ สวิตกฺกํ เอกํ ขนฺธญฺจ วิตกฺกญฺจ ปฏิจฺจ ตโย ขนฺธา ฯเปฯ เทฺว ขนฺเธ ฯเปฯ.\csromanspacer{} ปฏิสนฺธิกฺขเณ สวิตกฺกํ เอกํ ขนฺธญฺจ วตฺถุญฺจ ปฏิจฺจ ตโย ขนฺธา ฯเปฯ เทฺว ขนฺเธ ฯเปฯ. (1)}

\prose{สวิตกฺกญฺจ อวิตกฺกญฺจ ธมฺมํ ปฏิจฺจ อวิตกฺโก ธมฺโม อุปฺปชฺชติ เหตุปจฺจยา.\csromanspacer{} สวิตกฺเก ขนฺเธ จ วิตกฺกญฺจ มหาภูเต จ ปฏิจฺจ จิตฺตสมุฏฺฐานํ รูปํ.\csromanspacer{} ปฏิสนฺธิกฺขเณ สวิตกฺเก ขนฺเธ จ วิตกฺกญฺจ มหาภูเต จ ปฏิจฺจ กฏตฺตารูปํ.\csromanspacer{} ปฏิสนฺธิกฺขเณ สวิตกฺเก ขนฺเธ จ วตฺถุญฺจ ปฏิจฺจ วิตกฺโก. (2)}

\prose{สวิตกฺกญฺจ อวิตกฺกญฺจ ธมฺมํ ปฏิจฺจ สวิตกฺโก จ อวิตกฺโก จ ธมฺมา อุปฺปชฺชนฺติ เหตุปจฺจยา.\csromanspacer{} สวิตกฺกํ เอกํ ขนฺธญฺจ วิตกฺกญฺจ ปฏิจฺจ ตโย ขนฺธา จิตฺตสมุฏฺฐานญฺจ รูปํ ฯเปฯ เทฺว ขนฺเธ ฯเปฯ.\csromanspacer{} สวิตกฺกํ เอกํ ขนฺธญฺจ วิตกฺกญฺจ ปฏิจฺจ ตโย ขนฺธา ฯเปฯ เทฺว ขนฺเธ ฯเปฯ.\csromanspacer{} สวิตกฺเก ขนฺเธ จ วิตกฺกญฺจ มหาภูเต จ ปฏิจฺจ จิตฺตสมุฏฺฐานํ รูปํ.\csromanspacer{} ปฏิสนฺธิกฺขเณ สวิตกฺกํ เอกํ ขนฺธญฺจ วิตกฺกญฺจ ปฏิจฺจ ตโย ขนฺธา กฏตฺตา จ รูปํ ฯเปฯ เทฺว ขนฺเธ ฯเปฯ.\csromanspacer{} ปฏิสนฺธิกฺขเณ สวิตกฺกํ เอกํ ขนฺธญฺจ วิตกฺกญฺจ ปฏิจฺจ ตโย ขนฺธา ฯเปฯ เทฺว ขนฺเธ ฯเปฯ.\csromanspacer{} สวิตกฺเก ขนฺเธ จ วิตกฺกญฺจ มหาภูเต จ ปฏิจฺจ กฏตฺตารูปํ.\csromanspacer{} ปฏิสนฺธิกฺขเณ สวิตกฺกํ เอกํ ขนฺธญฺจ วตฺถุญฺจ ปฏิจฺจ ตโย ขนฺธา วิตกฺโก จ ฯเปฯ เทฺว ขนฺเธ ฯเปฯ. (3) (สํขิตฺตํ.)}

\csromanheader{2}{1. ปจฺจยานุโลม 2. สงฺขฺยาวาร}
\proseitem{105}{เหตุยา นว, อารมฺมเณ นว, อธิปติยา นว ฯเปฯ อุปนิสฺสเย นว, ปุเรชาเต ฉ, อาเสวเน ฉ, กมฺเม นว, วิปาเก นว, (สพฺพตฺถ นว.) อวิคเต นว.}
\csromansectionrule

\csromanheader{2}{2. ปจฺจยปจฺจนีย 1. วิภงฺควาร}
\csromanheader{2}{\textbf{นเหตุ}}
\proseitem{106}{สวิตกฺกํ ธมฺมํ ปฏิจฺจ สวิตกฺโก ธมฺโม อุปฺปชฺชติ นเหตุปจฺจยา.\csromanspacer{} อเหตุกํ สวิตกฺกํ เอกํ ขนฺธํ ปฏิจฺจ ตโย ขนฺธา ฯเปฯ เทฺว ขนฺเธ ฯเปฯ.\csromanspacer{} อเหตุกปฏิสนฺธิกฺขเณ ฯเปฯ.\csromanspacer{} วิจิกิจฺฉาสหคเต อุทฺธจฺจสหคเต ขนฺเธ}

\vspace{\tipitakaparskip}
\csromancenter{สวิตกฺกทุก ปฏิจฺจ วิจิกิจฺฉาสหคโต อุทฺธจฺจสหคโต โมโห. (สวิตกฺกมูลกา อวเสสา เทฺว ปญฺหา กาตพฺพา.\csromanspacer{} อเหตุกํ นินฺนานํ.) (3)}
\prose{\csromanfolio{83}อวิตกฺกํ ธมฺมํ ปฏิจฺจ อวิตกฺโก ธมฺโม อุปฺปชฺชติ นเหตุปจฺจยา.\csromanspacer{} อเหตุกํ อวิตกฺกํ เอกํ ขนฺธํ ปฏิจฺจ ตโย ขนฺธา ฯเปฯ เทฺว ขนฺเธ ฯเปฯ.\csromanspacer{} อเหตุกํ วิตกฺกํ ปฏิจฺจ จิตฺตสมุฏฺฐานํ รูปํ.\csromanspacer{} อเหตุกปฏิสนฺธิกฺขเณ วิตกฺกํ ปฏิจฺจ กฏตฺตารูปํ.\csromanspacer{} วิตกฺกํ ปฏิจฺจ วตฺถุ, วตฺถุํ ปฏิจฺจ วิตกฺโก.\csromanspacer{} เอกํ มหาภูตํ ฯเปฯ. (1)}

\prose{อวิตกฺกํ ธมฺมํ ปฏิจฺจ สวิตกฺโก ธมฺโม อุปฺปชฺชติ นเหตุปจฺจยา.\csromanspacer{} อเหตุกํ วิตกฺกํ ปฏิจฺจ สมฺปยุตฺตกา ขนฺธา.\csromanspacer{} อเหตุกปฏิสนฺธิกฺขเณ วิตกฺกํ ปฏิจฺจ สมฺปยุตฺตกา ขนฺธา.\csromanspacer{} ปฏิสนฺธิกฺขเณ วตฺถุํ ปฏิจฺจ อเหตุกา สวิตกฺกา ขนฺธา.\csromanspacer{} วิตกฺกํ ปฏิจฺจ วิจิกิจฺฉาสหคโต อุทฺธจฺจสหคโต โมโห. (2)}

\prose{อวิตกฺกํ ธมฺมํ ปฏิจฺจ สวิตกฺโก จ อวิตกฺโก จ ธมฺมา อุปฺปชฺชนฺติ นเหตุปจฺจยา. (สํขิตฺตํ.\csromanspacer{} เหตุปจฺจยสทิสํ. “อเหตุกนฺ”ติ นิยาเมตพฺพํ.) (3)}

\prose{สวิตกฺกญฺจ อวิตกฺกญฺจ ธมฺมํ ปฏิจฺจ สวิตกฺโก ธมฺโม อุปฺปชฺชติ นเหตุปจฺจยา.\csromanspacer{} อเหตุกํ สวิตกฺกํ เอกํ ขนฺธญฺจ วิตกฺกญฺจ ปฏิจฺจ ตโย ขนฺธา ฯเปฯ เทฺว ขนฺเธ ฯเปฯ.\csromanspacer{} อเหตุกปฏิสนฺธิกฺขเณ สวิตกฺกํ เอกํ ขนฺธญฺจ วิตกฺกญฺจ ปฏิจฺจ ตโย ขนฺธา ฯเปฯ เทฺว ขนฺเธ ฯเปฯ.\csromanspacer{} อเหตุกปฏิสนฺธิกฺขเณ สวิตกฺกํ เอกํ ขนฺธญฺจ วตฺถุญฺจ วิตกฺกญฺจ ปฏิจฺจ ตโย ขนฺธา ฯเปฯ เทฺว ขนฺเธ ฯเปฯ.\csromanspacer{} วิจิกิจฺฉาสหคเต อุทฺธจฺจสหคเต ขนฺเธ จ วิตกฺกญฺจ ปฏิจฺจ วิจิกิจฺฉาสหคโต อุทฺธจฺจสหคโต โมโห. (อวเสสา เทฺว ปญฺหา เหตุปจฺจยสทิสา นินฺนานา, “อเหตุกนฺ”ติ นิยาเมตพฺพํ.) (3)}

\csromanheader{2}{2. ปจฺจยปจฺจนีย 2. สงฺขฺยาวาร}
\proseitem{107}{นเหตุยา นว, น-อารมฺมเณ ตีณิ, น-อธิปติยา นว, น-อนนฺตเร ตีณิ ฯเปฯ น-อุปนิสฺสเย ตีณิ, นปุเรชาเต นว, นปจฺฉาชาเต นว, น-อาเสวเน นว, นกมฺเม จตฺตาริ, นวิปาเก นว, น-อาหาเร เอกํ, น-อินฺทฺริเย เอกํ, นฌาเน เอกํ, นมคฺเค นว, นสมฺปยุตฺเต ตีณิ, นวิปฺปยุตฺเต ฉ, โนนตฺถิยา ตีณิ, โนวิคเต ตีณิ.}

\csromanheader{2}{3. ปจฺจยานุโลมปจฺจนีย}
\proseitem{108}{\csromanfolio{84}\textbf{เหตุ}ปจฺจยา น-อารมฺมเณ ตีณิ, น-อธิปติยา นว, น-อนนฺตเร ตีณิ ฯเปฯ นกมฺเม จตฺตาริ, นวิปาเก นว, นสมฺปยุตฺเต ตีณิ, นวิปฺปยุตฺเต ฉ, โนนตฺถิยา ตีณิ, โนวิคเต ตีณิ.}

\csromanheader{2}{4. ปจฺจยปจฺจนียานุโลม}
\proseitem{109}{นเหตุปจฺจยา อารมฺมเณ นว ฯเปฯ อนนฺตเร นว ฯเปฯ ปุเรชาเต ฉ, อาเสวเน ปญฺจ, กมฺเม นว ฯเปฯ มคฺเค ตีณิ, สมฺปยุตฺเต นว. (สพฺพตฺถ นว.)}

\vspace{\tipitakaparskip}
\csromancenter{(สหชาตวาโร ปฏิจฺจวารสทิโส.)}
\csromanheader{2}{87. สวิตกฺกทุก 3. ปจฺจยวาร}
\proseitem{110}{สวิตกฺกํ ธมฺมํ ปจฺจยา สวิตกฺโก ธมฺโม อุปฺปชฺชติ เหตุปจฺจยา.\csromanspacer{} ตีณิ. (ปฏิจฺจวารสทิสา.)}

\prose{อวิตกฺกํ ธมฺมํ ปจฺจยา อวิตกฺโก ธมฺโม อุปฺปชฺชติ เหตุปจฺจยา.\csromanspacer{} อวิตกฺกํ เอกํ ขนฺธํ ปจฺจยา ตโย ขนฺธา จิตฺตสมุฏฺฐานญฺจ รูปํ ฯเปฯ เทฺว ขนฺเธ ฯเปฯ.\csromanspacer{} วิตฺตกฺกํ ปจฺจยา จิตฺตสมุฏฺฐานํ รูปํ.\csromanspacer{} ปฏิสนฺธิกฺขเณ อวิตกฺกํ เอกํ ขนฺธํ ปจฺจยา ตโย ขนฺธา กฏตฺตา จ รูปํ ฯเปฯ เทฺว ขนฺเธ ฯเปฯ.\csromanspacer{} ปฏิสนฺธิกฺขเณ วิตกฺกํ ปจฺจยา กฏตฺตารูปํ.\csromanspacer{} ขนฺเธ ปจฺจยา วตฺถุ, วตฺถุํ ปจฺจยา ขนฺธา.\csromanspacer{} วิตกฺกํ ปจฺจยา วตฺถุ, วตฺถุํ, ปจฺจยา วิตกฺโก.\csromanspacer{} เอกํ มหาภูตํ ปจฺจยา ตโย มหาภูตา ฯเปฯ.\csromanspacer{} วตฺถุํ ปจฺจยา อวิตกฺกา ขนฺธา.\csromanspacer{} วตฺถุํ ปจฺจยา วิตกฺโก. (1)}

\prose{อวิตกฺกํ ธมฺมํ ปจฺจยา สวิตกฺโก ธมฺโม อุปฺปชฺชติ เหตุปจฺจยา.\csromanspacer{} วิตกฺกํ ปจฺจยา สมฺปยุตฺตกา ขนฺธา.\csromanspacer{} วตฺถุํ ปจฺจยา สวิตกฺกา ขนฺธา. (ปฏิสนฺธิยาปิ เทฺว.) (2)}

\prose{อวิตกฺกํ ธมฺมํ ปจฺจยา สวิตกฺโก จ อวิตกฺโก จ ธมฺมา อุปฺปชฺชนฺติ เหตุปจฺจยา.\csromanspacer{} วิตกฺกํ ปจฺจยา สมฺปยุตฺตกา ขนฺธา จิตฺตสมุฏฺฐานญฺจ รูปํ.\csromanspacer{} วิตกฺกํ}

\vspace{\tipitakaparskip}
\csromancenter{สวิตกฺกทุก ปจฺจยา สมฺปยุตฺตกา ขนฺธา.\csromanspacer{} มหาภูเต ปจฺจยา จิตฺตสมุฏฺฐานํ รูปํ.\csromanspacer{} วตฺถุํ ปจฺจยา สวิตกฺกา ขนฺธา.\csromanspacer{} มหาภูเต ปจฺจยา จิตฺตสมุฏฺฐานํ รูปํ.\csromanspacer{} วตฺถุํ ปจฺจยา วิตกฺโก สมฺปยุตฺตกา จ ขนฺธา.\csromanspacer{} ปฏิสนฺธิกฺขเณ ฯเปฯ. (ปฏิสนฺธิยาปิ ปวตฺติสทิสาเยว.) (3)}
\proseitem{111}{\csromanfolio{85}สวิตกฺกญฺจ อวิตกฺกญฺจ ธมฺมํ ปจฺจยา สวิตกฺโก ธมฺโม อุปฺปชฺชติ เหตุปจฺจยา.\csromanspacer{} สวิตกฺกํ เอกํ ขนฺธญฺจ วิตกฺกญฺจ ปจฺจยา ตโย ขนฺธา ฯเปฯ เทฺว ขนฺเธ ฯเปฯ.\csromanspacer{} สวิตกฺกํ เอกํ ขนฺธญฺจ วตฺถุญฺจ ปจฺจยา ตโย ขนฺธา ฯเปฯ เทฺว ขนฺเธ ฯเปฯ. (ปฏิสนฺธิกฺขเณ.\csromanspacer{} เทฺว กาตพฺพา.) (1)}

\prose{สวิตกฺกญฺจ อวิตกฺกญฺจ ธมฺมํ ปจฺจยา อวิตกฺโก ธมฺโม อุปฺปชฺชติ เหตุปจฺจยา.\csromanspacer{} สวิตกฺเก ขนฺเธ จ วิตกฺกญฺจ ปจฺจยา จิตฺตสมุฏฺฐานํ รูปํ.\csromanspacer{} สวิตกฺเก ขนฺเธ จ วิตกฺกญฺจ มหาภูเต จ ปจฺจยา จิตฺตสมุฏฺฐานํ รูปํ.\csromanspacer{} สวิตกฺเก ขนฺเธ จ วตฺถุญฺจ ปจฺจยา วิตกฺโก.\csromanspacer{} ปฏิสนฺธิกฺขเณ ฯเปฯ. (ตีณิ.\csromanspacer{} ปฏิสนฺธิยาปิ.) (2)}

\prose{สวิตกฺกญฺจ อวิตกฺกญฺจ ธมฺมํ ปจฺจยา สวิตกฺโก จ อวิตกฺโก จ ธมฺมา อุปฺปชฺชนฺติ เหตุปจฺจยา.\csromanspacer{} สวิตกฺกํ เอกํ ขนฺธญฺจ วิตกฺกญฺจ ปจฺจยา ตโย ขนฺธา จิตฺตสมุฏฺฐานญฺจ รูปํ ฯเปฯ เทฺว ขนฺเธ ฯเปฯ.\csromanspacer{} สวิตกฺกํ เอกํ ขนฺธญฺจ วิตกฺกญฺจ วตฺถุญฺจ ปจฺจยา ตโย ขนฺธา ฯเปฯ เทฺว ขนฺเธ ฯเปฯ.\csromanspacer{} สวิตกฺเก ขนฺเธ จ วิตกฺกญฺจ มหาภูเต จ ปจฺจยา จิตฺตสมุฏฺฐานํ รูปํ.\csromanspacer{} สวิตกฺกํ เอกํ ขนฺธญฺจ วตฺถุญฺจ ปจฺจยา ตโย ขนฺธา วิตกฺโก จ ฯเปฯ เทฺว ขนฺเธ ฯเปฯ.\csromanspacer{} ปฏิสนฺธิกฺขเณ ฯเปฯ. (3)}

\csromanheader{2}{1. ปจฺจยานุโลม 2. สงฺขฺยาวาร}
\vspace{\tipitakaparskip}
\csromancenter{เหตุยา นว, อารมฺมเณ นว, อธิปติยา นว, (สพฺพตฺถ นว.) อวิคเต นว.}
\csromanheader{2}{2. ปจฺจยปจฺจนีย 1. วิภงฺควาร}
\proseitem{113}{สวิตกฺกํ ธมฺมํ ปจฺจยา สวิตกฺโก ธมฺโม อุปฺปชฺชติ นเหตุปจฺจยา. (นว ปญฺหา กาตพฺพา. “อเหตุกา”ติ นิยาเมตพฺพา ตีณิเยว.\csromanspacer{} โมโห อุทฺธริตพฺโพ, ยถา ปฏิจฺจวาเร เหตุปจฺจยสทิสาเยว ปญฺหา ปญฺจวิญฺญาณา อติเรกา โมโห วิตกฺกํ.)}

\csromanheader{2}{2. ปจฺจยปจฺจนีย 2. สงฺขฺยาวาร}
\proseitem{114}{\csromanfolio{86}นเหตุยา นว, น-อารมฺมเณ ตีณิ, น-อธิปติยา นว, น-อนนฺตเร ตีณิ ฯเปฯ น-อุปนิสฺสเย ตีณิ, นปุเรชาเต นว, นปจฺฉาชาเต นว, น-อาเสวเน นว, นกมฺเม จตฺตาริ, นวิปาเก นว, น-อาหาเร เอกํ, น-อินฺทฺริเย เอกํ, นฌาเน เอกํ, นมคฺเค นว, นสมฺปยุตฺเต ตีณิ, นวิปฺปยุตฺเต ฉ, โนนตฺถิยา ตีณิ, โนวิคเต ตีณิ. (เอวํ อิตเร เทฺว คณนาปิ นิสฺสยวาโรปิ กาตพฺโพ.)}

\csromanheader{2}{87. สวิตกฺกทุก 5. สํสฏฺฐวาร \textbf{1. ปจฺจยานุโลมาทิ}}
\proseitem{115}{สวิตกฺกํ ธมฺมํ สํสฏฺโฐ สวิตกฺโก ธมฺโม อุปฺปชฺชติ เหตุปจฺจยา.\csromanspacer{} สวิตกฺกํ เอกํ ขนฺธํ สํสฏฺฐา ตโย ขนฺธา ฯเปฯ เทฺว ขนฺเธ ฯเปฯ.\csromanspacer{} ปฏิสนฺธิกฺขเณ ฯเปฯ.\csromanspacer{} สวิตกฺกํ ธมฺมํ สํสฏฺโฐ อวิตกฺโก ธมฺโม อุปฺปชฺชติ เหตุปจฺจยา.\csromanspacer{} สวิตกฺเก ขนฺเธ สํสฏฺโฐ วิตกฺโก.\csromanspacer{} ปฏิสนฺธิกฺขเณ ฯเปฯ.\csromanspacer{} สวิตกฺกํ ธมฺมํ สํสฏฺโฐ สวิตกฺโก จ อวิตกฺโก จ ธมฺมา อุปฺปชฺชนฺติ เหตุปจฺจยา.\csromanspacer{} สวิตกฺกํ เอกํ ขนฺธํ สํสฏฺฐา ตโย ขนฺธา วิตกฺโก จ ฯเปฯ เทฺว ขนฺเธ ฯเปฯ.\csromanspacer{} ปฏิสนฺธิกฺขเณ ฯเปฯ. (3)}

\prose{อวิตกฺกํ ธมฺมํ สํสฏฺโฐ อวิตกฺโก ธมฺโม อุปฺปชฺชติ เหตุปจฺจยา.\csromanspacer{} อวิตกฺกํ เอกํ ขนฺธํ สํสฏฺฐา ตโย ขนฺธา ฯเปฯ เทฺว ขนฺเธ ฯเปฯ.\csromanspacer{} ปฏิสนฺธิกฺขเณ ฯเปฯ.\csromanspacer{} อวิตกฺกํ ธมฺมํ สํสฏฺโฐ สวิตกฺโก ธมฺโม อุปฺปชฺชติ เหตุปจฺจยา.\csromanspacer{} วิตกฺกํ สํสฏฺฐา สมฺปยุตฺตกา ขนฺธา.\csromanspacer{} ปฏิสนฺธิกฺขเณ ฯเปฯ. (2)}

\prose{สวิตกฺกญฺจ อวิตกฺกญฺจ ธมฺมํ สํสฏฺโฐ สวิตกฺโก ธมฺโม อุปฺปชฺชติ เหตุปจฺจยา.\csromanspacer{} สวิตกฺกํ เอกํ ขนฺธญฺจ วิตกฺกญฺจ สํสฏฺฐา ตโย ขนฺธา ฯเปฯ เทฺว ขนฺเธ ฯเปฯ.\csromanspacer{} ปฏิสนฺธิกฺขเณ ฯเปฯ. (1) (สํขิตฺตํ.)}

\vspace{\tipitakaparskip}
\csromancenter{\textbf{เหตุ}ยา ฉ, อารมฺมเณ ฉ, อธิปติยา ฉ, (สพฺพตฺถ ฉ.) อวิคเต ฉ.}
\csromancenter{อนุโลมํ.}
\csromanheader{2}{87. สวิตกฺกทุก}
\prose{\csromanfolio{87}สวิตกฺกํ ธมฺมํ สํสฏฺโฐ สวิตกฺโก ธมฺโม อุปฺปชฺชติ นเหตุปจฺจยา. (เอวํ ฉ ปญฺหา กาตพฺพา, อนุโลมสทิสา, “อเหตุกา”ติ นิยาเมตพฺพา.\csromanspacer{} ตีณิเยว, โมโห อุทฺธริตพฺโพ.)}

\prose{นเหตุยา ฉ, น-อธิปติยา ฉ, นปุเรชาเต ฉ, นปจฺฉาชาเต ฉ, น-อาเสวเน ฉ, นกมฺเม จตฺตาริ, นวิปาเก ฉ, นฌาเน เอกํ, นมคฺเค ฉ, นวิปฺปยุตฺเต ฉ.}

\vspace{\tipitakaparskip}
\csromancenter{ปจฺจนียํ.}
\csromancenter{(เอวํ อิตเร เทฺว คณนาปิ สมฺปยุตฺตวาโรปิ กาตพฺโพ.)}
\csromanheader{2}{87. สวิตกฺกทุก 7. ปญฺหาวาร}
\proseitem{116}{สวิตกฺโก ธมฺโม สวิตกฺกสฺส ธมฺมสฺส เหตุปจฺจเยน ปจฺจโย.\csromanspacer{} สวิตกฺกา เหตู สมฺปยุตฺตกานํ ขนฺธานํ เหตุปจฺจเยน ปจฺจโย.\csromanspacer{} ปฏิสนฺธิกฺขเณ ฯเปฯ.\csromanspacer{} สวิตกฺโก ธมฺโม อวิตกฺกสฺส ธมฺมสฺส เหตุปจฺจเยน ปจฺจโย.\csromanspacer{} สวิตกฺกา เหตู วิตกฺกสฺส จิตฺตสมุฏฺฐานานญฺจ รูปานํ เหตุปจฺจเยน ปจฺจโย.\csromanspacer{} ปฏิสนฺธิกฺขเณ ฯเปฯ. (มูลํ กาตพฺพํ.) สวิตกฺกา เหตู สมฺปยุตฺตกานํ ขนฺธานํ วิตกฺกสฺส จ จิตฺตสมุฏฺฐานานญฺจ รูปานํ เหตุปจฺจเยน ปจฺจโย.\csromanspacer{} ปฏิสนฺธิกฺขเณ ฯเปฯ. (3)}

\prose{อวิตกฺโก ธมฺโม อวิตกฺกสฺส ธมฺมสฺส เหตุปจฺจเยน ปจฺจโย.\csromanspacer{} อวิตกฺกา เหตู สมฺปยุตฺตกานํ ขนฺธานํ จิตฺตสมุฏฺฐานานญฺจ รูปานํ เหตุปจฺจเยน ปจฺจโย.\csromanspacer{} ปฏิสนฺธิกฺขเณ ฯเปฯ. (1)}

\csromanheader{2}{\textbf{อารมฺมณ}}
\proseitem{117}{สวิตกฺโก ธมฺโม สวิตกฺกสฺส ธมฺมสฺส อารมฺมณปจฺจเยน ปจฺจโย.\csromanspacer{} สวิตกฺเก ขนฺเธ อารพฺภ สวิตกฺกา ขนฺธา อุปฺปชฺชนฺติ. (มูลํ กาตพฺพํ.) สวิตกฺเก ขนฺเธ อารพฺภ อวิตกฺกา ขนฺธา จ วิตกฺโก จ \csromanfolio{88}อุปฺปชฺชนฺติ. (มูลํ กาตพฺพํ.) สวิตกฺเก ขนฺเธ อารพฺภ สวิตฺตกฺกา ขนฺธา จ วิตกฺโก จ อุปฺปชฺชนฺติ. (3)}

\proseitem{118}{อวิตกฺโก ธมฺโม อวิตกฺกสฺส ธมฺมสฺส อารมฺมณปจฺจเยน ปจฺจโย.\csromanspacer{} อริยา อวิตกฺกา ฌานา วิฏฺฐหิตฺวา อวิตกฺกํ ฌานํ ปจฺจเวกฺขนฺติ, มคฺคา วุฏฺฐหิตฺวา มคฺคํ ปจฺจเวกฺขนฺติ, ผลา วุฏฺฐหิตฺวา ผลํ ปจฺจเวกฺขนฺติ, นิพฺพานํ ปจฺจเวกฺขนฺติ.\csromanspacer{} นิพฺพานํ อวิตกฺกสฺส มคฺคสฺส, ผลสฺส, วิตกฺกสฺส จ อารมฺมณปจฺจเยน ปจฺจโย.\csromanspacer{} จกฺขุํ ฯเปฯ.\csromanspacer{} วตฺถุํ อวิตกฺเก ขนฺเธ จ วิตกฺกญฺจ อนิจฺจโต ฯเปฯ วิปสฺสติ, อสฺสาเทติ อภินนฺทติ, ตํ อารพฺภ วิตกฺโก อุปฺปชฺชติ.\csromanspacer{} ทิพฺเพน จกฺขุนา รูปํ ปสฺสติ.\csromanspacer{} ทิพฺพาย โสตธาตุยา สทฺทํ สุณาติ.\csromanspacer{} เจโตปริยญาเณน อวิตกฺกจิตฺตสมงฺคีสฺส จิตฺตํ ชานาติ.\csromanspacer{} อากาสานญฺจายตนํ ฯเปฯ.\csromanspacer{} อากิญฺจญฺญายตนํ ฯเปฯ.\csromanspacer{} รูปายตนํ จกฺขุวิญฺญาณสฺส ฯเปฯ.\csromanspacer{} โผฏฺฐพฺพายตนํ ฯเปฯ.\csromanspacer{} อวิตกฺกา ขนฺธา อิทฺธิวิธญาณสฺส, เจโตปริยญาณสฺส, ปุพฺเพนิวาสานุสฺสติญาณสฺส, ยถากมฺมูปคญาณสฺส, อนาคตํสญาณสฺส, วิตกฺกสฺส จ อารมฺมณปจฺจเยน ปจฺจโย.\csromanspacer{} อวิตกฺเก ขนฺเธ จ วิตกฺกญฺจ อารพฺภ อวิตกฺกา ขนฺธา จ วิตกฺโก จ อุปฺปชฺชนฺติ. (1)}

\prose{อวิตกฺโก ธมฺโม สวิตกฺกสฺส ธมฺมสฺส อารมฺมณปจฺจเยน ปจฺจโย.\csromanspacer{} อริยา อวิตกฺกา ฌานา วุฏฺฐหิตฺวา ฯเปฯ มคฺคา วุฏฺฐหิตฺวา ฯเปฯ ผลา วุฏฺฐหิตฺวา ผลํ ปจฺจเวกฺขนฺติ, นิพฺพานํ ปจฺจเวกฺขนฺติ.\csromanspacer{} นิพฺพานํ โคตฺรภุสฺส, โวทานสฺส, สวิตกฺกสฺส มคฺคสฺส, ผลสฺส, อาวชฺชนาย อารมฺมณปจฺจเยน ปจฺจโย.\csromanspacer{} จกฺขุํ ฯเปฯ.\csromanspacer{} วตฺถุํ อวิตกฺเก ขนฺเธ จ วิตกฺกญฺจ อนิจฺจโต ฯเปฯ วิปสฺสติ, อสฺสาเทติ อภินนฺทติ, ตํ อารพฺภ ราโค อุปฺปชฺชติ ฯเปฯ โทมนสฺสํ อุปฺปชฺชติ.\csromanspacer{} อวิตกฺเก ขนฺเธ จ วิตกฺกญฺจ อารพฺภ สวิตกฺกา ขนฺธา อุปฺปชฺชนฺติ. (2)}

\prose{อวิตกฺโก ธมฺโม สวิตกฺกสฺส จ อวิตกฺกสฺส จ ธมฺมสฺส อารมฺมณปจฺจเยน ปจฺจโย.\csromanspacer{} อริยา อวิตกฺกา ฌานา วุฏฺฐหิตฺวา ฯเปฯ มคฺคา วุฏฺฐหิตฺวา ฯเปฯ ผลา วุฏฺฐหิตฺวา ผลํ ปจฺจเวกฺขนฺติ, นิพฺพานํ ปจฺจเวกฺขนฺติ.\csromanspacer{} นิพฺพานํ โคตฺรภุสฺส, โวทานสฺส, สวิตกฺกสฺส มคฺคสฺส, ผลสฺส, อาวชฺชนาย, วิตกฺกสฺส จ อารมฺมณปจฺจเยน ปจฺจโย.\csromanspacer{} จกฺขุํ ฯเปฯ.\csromanspacer{} วตฺถุํ อวิตกฺเก ขนฺเธ จ วิตกฺกญฺจ อนิจฺจโต ฯเปฯ วิปสฺสติ, อสฺสาเทติ อภินนฺทติ, ตํ อารพฺภ}

\vspace{\tipitakaparskip}
\csromancenter{สวิตกฺกทุก สวิตกฺกา ขนฺธา จ วิตกฺโก จ อุปฺปชฺชนฺติ.\csromanspacer{} อวิตกฺเก ขนฺเธ จ วิตกฺกญฺจ อารพฺภ สวิตกฺกา ขนฺธา จ วิตกฺโก จ อุปฺปชฺชนฺติ. (3)}
\proseitem{119}{\csromanfolio{89}สวิตกฺโก จ อวิตกฺโก จ ธมฺมา สวิตกฺกสฺส ธมฺมสฺส อารมฺมณปจฺจเยน ปจฺจโย.\csromanspacer{} สวิตกฺเก ขนฺเธ จ วิตกฺกญฺจ อารพฺภ สวิตกฺกา ขนฺธา อุปฺปชฺชนฺติ. (มูลํ กาตพฺพํ.) สวิตกฺเก ขนฺเธ จ วิตกฺกญฺจ อารพฺภ อวิตกฺกา ขนฺธา จ วิตกฺโก จ อุปฺปชฺชนฺติ. (มูลํ กาตพฺพํ.) สวิตกฺเก ขนฺเธ จ วิตกฺกญฺจ อารพฺภ สวิตกฺกา ขนฺธา จ วิตกฺโก จ อุปฺปชฺชนฺติ. (3)}

\csromanheader{2}{\textbf{อธิปติ}}
\proseitem{120}{สวิตกฺโก ธมฺโม สวิตกฺกสฺส ธมฺมสฺส อธิปติปจฺจเยน ปจฺจโย.\csromanspacer{} \textbf{อารมฺมณาธิปติ}, สหชาตาธิปติ.\csromanspacer{}\csromanspacer{}\textbf{อารมฺมณาธิปติ\csromandash{}}สวิตกฺเก ขนฺเธ ครุํ กตฺวา สวิตกฺกา ขนฺธา อุปฺปชฺชนฺติ.\csromanspacer{} \textbf{สหชาตาธิปติ\csromandash{}} สวิตกฺกาธิปติ สมฺปยุตฺตกานํ ขนฺธานํ อธิปติปจฺจเยน ปจฺจโย. (1)}

\prose{สวิตกฺโก ธมฺโม อวิตกฺกสฺส ธมฺมสฺส อธิปติปจฺจเยน ปจฺจโย.\csromanspacer{} \textbf{อารมฺมณาธิปติ}, สหชาตาธิปติ.\csromanspacer{}\csromanspacer{}\textbf{อารมฺมณาธิปติ\csromandash{}}สวิตกฺเก ขนฺเธ ครุํ กตฺวา วิตกฺโก อุปฺปชฺชติ.\csromanspacer{} \textbf{สหชาตาธิปติ\csromandash{}}สวิตกฺกาธิปติ วิตกฺกสฺส จ จิตฺตสมุฏฺฐานานญฺจ รูปานํ อธิปติปจฺจเยน ปจฺจโย. (2)}

\prose{สวิตกฺโก ธมฺโม สวิตกฺกสฺส จ อวิตกฺกสฺส จ ธมฺมสฺส อธิปติปจฺจเยน ปจฺจโย.\csromanspacer{} \textbf{อารมฺมณาธิปติ}, สหชาตาธิปติ.\csromanspacer{}\csromanspacer{}\textbf{อารมฺมณาธิปติ\csromandash{}}สวิตกฺเก ขนฺเธ ครุํ กตฺวา สวิตกฺกา ขนฺธา จ วิตกฺโก จ อุปฺปชฺชนฺติ.\csromanspacer{} \textbf{สหชาตาธิปติ\csromandash{}}สวิตกฺกาธิปติ สมฺปยุตฺตกานํ ขนฺธานํ วิตกฺกสฺส จ จิตฺตสมุฏฺฐานานญฺจ รูปานํ อธิปติปจฺจเยน ปจฺจโย. (3)}

\proseitem{121}{อวิตกฺโก ธมฺโม อวิตกฺกสฺส ธมฺมสฺส อธิปติปจฺจเยน ปจฺจโย.\csromanspacer{} \textbf{อารมฺมณาธิปติ}, สหชาตาธิปติ.\csromanspacer{}\csromanspacer{}\textbf{อารมฺมณาธิปติ\csromandash{}}อริยา อวิตกฺกา ฌานา วุฏฺฐหิตฺวา ฯเปฯ มคฺคา วุฏฺฐหิตฺวา ฯเปฯ ผลา วุฏฺฐหิตฺวา ผลํ ครุํ กตฺวา ปจฺจเวกฺขนฺติ, นิพฺพานํ ครุํ กตฺวา ปจฺจเวกฺขนฺติ.\csromanspacer{} นิพฺพานํ อวิตกฺกสฺส มคฺคสฺส, ผลสฺส, วิตกฺกสฺส จ อธิปติปจฺจเยน ปจฺจโย.\csromanspacer{} จกฺขุํ ฯเปฯ.\csromanspacer{} วตฺถุํ อวิตกฺเก ขนฺเธ จ วิตกฺกญฺจ ครุํ กตฺวา อสฺสาเทติ อภินนฺทติ, ตํ ครุํ กตฺวา วิตกฺโก อุปฺปชฺชติ.\csromanspacer{} \textbf{สหชาตาธิปติ\csromandash{}} \csromanfolio{90}อวิตกฺกาธิปติ สมฺปยุตฺตกานํ ขนฺธานํ จิตฺตสมุฏฺฐานานญฺจ รูปานํ อธิปติปจฺจเยน ปจฺจโย. (1)}

\prose{อวิตกฺโก ธมฺโม สวิตกฺกสฺส ธมฺมสฺส อธิปติปจฺจเยน ปจฺจโย.\csromanspacer{} \textbf{อารมฺมณาธิปติ\csromandash{}}อริยา อวิตกฺกา ฌานา วุฏฺฐหิตฺวา ฯเปฯ มคฺคา วุฏฺฐหิตฺวา ฯเปฯ ผลา วุฏฺฐหิตฺวา ผลํ ครุํ กตฺวา ปจฺจเวกฺขนฺติ, นิพฺพานํ ครุํ กตฺวา ปจฺจเวกฺขนฺติ.\csromanspacer{} นิพฺพานํ โคตฺรภุสฺส, โวทานสฺส, สวิตกฺกสฺส มคฺคสฺส, ผลสฺส อธิปติปจฺจเยน ปจฺจโย.\csromanspacer{} จกฺขุํ ฯเปฯ.\csromanspacer{} วตฺถุํ อวิตกฺเก ขนฺเธ จ วิตกฺกญฺจ ครุํ กตฺวา อสฺสาเทติ อภินนฺทติ, ตํ ครุํ กตฺวา ราโค อุปฺปชฺชติ, ทิฏฺฐิ อุปฺปชฺชติ.\csromanspacer{} อวิตกฺเก ขนฺเธ จ วิตกฺกญฺจ ครุํ กตฺวา สวิตกฺกา ขนฺธา อุปฺปชฺชนฺติ. (2)}

\prose{อวิตกฺโก ธมฺโม สวิตกฺกสฺส จ อวิตกฺกสฺส จ ธมฺมสฺส อธิปติปจฺจเยน ปจฺจโย.\csromanspacer{} \textbf{อารมฺมณาธิปติ\csromandash{}}อริยา อวิตกฺกา ฌานา วุฏฺฐหิตฺวา ฯเปฯ มคฺคา วุฏฺฐหิตฺวา ฯเปฯ ผลา วุฏฺฐหิตฺวา ผลํ ครุํ กตฺวา ปจฺจเวกฺขนฺติ, นิพฺพานํ ครุํ กตฺวา ปจฺจเวกฺขนฺติ.\csromanspacer{} นิพฺพานํ โคตฺรภุสฺส, โวทานสฺส, สวิตกฺกสฺส มคฺคสฺส, ผลสฺส, วิตกฺกสฺส จ อธิปติปจฺจเยน ปจฺจโย.\csromanspacer{} จกฺขุํ ฯเปฯ.\csromanspacer{} วตฺถุํ อวิตกฺเก ขนฺเธ จ วิตกฺกญฺจ ครุํ กตฺวา อสฺสาเทติ อภินนฺทติ, ตํ ครุํ กตฺวา ราโค อุปฺปชฺชติ, ทิฏฺฐิ อุปฺปชฺชติ.\csromanspacer{} อวิตกฺเก ขนฺเธ จ วิตกฺกญฺจ ครุํ กตฺวา สวิตกฺกา ขนฺธา จ วิตกฺโก จ อุปฺปชฺชนฺติ. (3)}

\prose{สวิตกฺโก จ อวิตกฺโก จ ธมฺมา สวิตกฺกสฺส ธมฺมสฺส อธิปติปจฺจเยน ปจฺจโย.\csromanspacer{} \textbf{อารมฺมณาธิปติ\csromandash{}}สวิตกฺเก ขนฺเธ จ วิตกฺกญฺจ ครุํ กตฺวา สวิตกฺกา ขนฺธา อุปฺปชฺชนฺติ. (มูลํ กาตพฺพํ.) สวิตกฺเก ขนฺเธ จ วิตกฺกญฺจ ครุํ กตฺวา วิตกฺโก อุปฺปชฺชติ. (มูลํ กาตพฺพํ.) สวิตกฺเก ขนฺเธ จ วิตกฺกญฺจ ครุํ กตฺวา สวิตกฺกา ขนฺธา จ วิตกฺโก จ อุปฺปชฺชนฺติ. (3)}

\csromanheader{2}{\textbf{อนนฺตราทิ}}
\proseitem{122}{สวิตกฺโก ธมฺโม สวิตกฺกสฺส ธมฺมสฺส อนนฺตรปจฺจเยน ปจฺจโย.\csromanspacer{} ปุริมา ปุริมา สวิตกฺกา ขนฺธา ปจฺฉิมานํ ปจฺฉิมานํ สวิตกฺกานํ ขนฺธานํ อนนฺตรปจฺจเยน ปจฺจโย. (1)}

\csromanheader{2}{87. สวิตกฺกทุก}
\prose{\csromanfolio{91}สวิตกฺโก ธมฺโม อวิตกฺกสฺส ธมฺมสฺส อนนฺตรปจฺจเยน ปจฺจโย.\csromanspacer{} ปุริมา ปุริมา สวิตกฺกา ขนฺธา ปจฺฉิมสฺส ปจฺฉิมสฺส วิตกฺกสฺส อนนฺตรปจฺจเยน ปจฺจโย.\csromanspacer{} สวิตกฺกํ จุติจิตฺตํ อวิตกฺกสฺส อุปปตฺติจิตฺตสฺส อนนฺตรปจฺจเยน ปจฺจโย.\csromanspacer{} อาวชฺชนา ปญฺจนฺนํ วิญฺญาณานํ อนนฺตรปจฺจเยน ปจฺจโย.\csromanspacer{} สวิตกฺกา ขนฺธา อวิตกฺกสฺส วุฏฺฐานสฺส อนนฺตรปจฺจเยน ปจฺจโย.\csromanspacer{} ทุติยสฺส ฌานสฺส ปริกมฺมํ ทุติยสฺส ฌานสฺส อนนฺตรปจฺจเยน ปจฺจโย.\csromanspacer{} ตติยสฺส ฌานสฺส ปริกมฺมํ ฯเปฯ.\csromanspacer{} เนวสญฺญานาสญฺญายตนสฺส ปริกมฺมํ เนวสญฺญานาสญฺญายตนสฺส ฯเปฯ.\csromanspacer{} ทิพฺพสฺส จกฺขุสฺส ปริกมฺมํ ฯเปฯ.\csromanspacer{} ทิพฺพาย โสตธาตุยา ปริกมฺมํ ฯเปฯ.\csromanspacer{} อิทฺธิวิธญาณสฺส ปริกมฺมํ ฯเปฯ.\csromanspacer{} เจโตปริยญาณสฺส ปริกมฺมํ ฯเปฯ.\csromanspacer{} ปุพฺเพนิวาสานุสฺสติญาณสฺส ปริกมฺมํ ฯเปฯ.\csromanspacer{} ยถากมฺมูปคญาณสฺส ปริกมฺมํ ยถากมฺมูปคญาณสฺส ฯเปฯ.\csromanspacer{} อนาคตํสญาณสฺส ปริกมฺมํ อนาคตํสญาณสฺส อนนฺตรปจฺจเยน ปจฺจโย.\csromanspacer{} โคตฺรภุ อวิตกฺกสฺส มคฺคสฺส.\csromanspacer{} โวทานํ อวิตกฺกสฺส มคฺคสฺส.\csromanspacer{} อนุโลมํ อวิตกฺกาย ผลสมาปตฺติยา อนนฺตรปจฺจเยน ปจฺจโย. (2)}

\prose{สวิตกฺโก ธมฺโม สวิตกฺกสฺส จ อวิตกฺกสฺส จ ธมฺมสฺส อนนฺตรปจฺจเยน ปจฺจโย.\csromanspacer{} ปุริมา ปุริมา สวิตกฺกา ขนฺธา ปจฺฉิมานํ ปจฺฉิมานํ สวิตกฺกานํ ขนฺธานํ วิตกฺกสฺส จ อนนฺตรปจฺจเยน ปจฺจโย. (3)}

\proseitem{123}{อวิตกฺโก ธมฺโม อวิตกฺกสฺส ธมฺมสฺส อนนฺตรปจฺจเยน ปจฺจโย.\csromanspacer{} ปุริโม ปุริโม วิตกฺโก ปจฺฉิมสฺส ปจฺฉิมสฺส วิตกฺกสฺส อนนฺตรปจฺจเยน ปจฺจโย.\csromanspacer{} ปุริมา ปุริมา อวิตกฺกา ขนฺธา ปจฺฉิมานํ ปจฺฉิมานํ อวิตกฺกานํ ขนฺธานํ อนนฺตรปจฺจเยน ปจฺจโย.\csromanspacer{} อวิตกฺโก มคฺโค อวิตกฺกสฺส ผลสฺส ฯเปฯ.\csromanspacer{} อวิตกฺกํ ผลํ อวิตกฺกสฺส ผลสฺส ฯเปฯ.\csromanspacer{} นิโรธา วุฏฺฐหนฺตสฺส เนวสญฺญานาสญฺญายตนํ อวิตกฺกาย ผลสมาปตฺติยา อนนฺตรปจฺจเยน ปจฺจโย. (1)}

\prose{อวิตกฺโก ธมฺโม สวิตกฺกสฺส ธมฺมสฺส อนนฺตรปจฺจเยน ปจฺจโย.\csromanspacer{} ปุริโม ปุริโม วิตกฺโก ปจฺฉิมานํ ปจฺฉิมานํ สวิตกฺกานํ ขนฺธานํ อนนฺตรปจฺจเยน ปจฺจโย.\csromanspacer{} อวิตกฺกํ จุติจิตฺตํ สวิตกฺกสฺส อุปปตฺติจิตฺตสฺส.\csromanspacer{} อวิตกฺกํ ภวงฺคํ อาวชฺชนาย.\csromanspacer{} อวิตกฺกา ขนฺธา สวิตกฺกสฺส วุฏฺฐานสฺส อนนฺตรปจฺจเยน ปจฺจโย. (2)}

\prose{\csromanfolio{92}อวิตกฺโก ธมฺโม สวิตกฺกสฺส จ อวิตกฺกสฺส จ ธมฺมสฺส อนนฺตรปจฺจเยน ปจฺจโย.\csromanspacer{} ปุริโม ปุริโม วิตกฺโก ปจฺฉิมานํ ปจฺฉิมานํ สวิตกฺกานํ ขนฺธานํ วิตกฺกสฺส จ อนนฺตรปจฺจเยน ปจฺจโย. (3)}

\proseitem{124}{สวิตกฺโก จ อวิตกฺโก จ ธมฺมา สวิตกฺกสฺส ธมฺมสฺส อนนฺตรปจฺจเยน ปจฺจโย.\csromanspacer{} ปุริมา ปุริมา สวิตกฺกา ขนฺธา จ วิตกฺโก จ ปจฺฉิมานํ ปจฺฉิมานํ สวิตกฺกานํ ขนฺธานํ อนนฺตรปจฺจเยน ปจฺจโย. (1)}

\prose{สวิตกฺโก จ อวิตกฺโก จ ธมฺมา อวิตกฺกสฺส ธมฺมสฺส อนนฺตรปจฺจเยน ปจฺจโย.\csromanspacer{} ปุริมา ปุริมา สวิตกฺกา ขนฺธา จ วิตกฺโก จ ปจฺฉิมสฺส ปจฺฉิมสฺส วิตกฺกสฺส อนนฺตรปจฺจเยน ปจฺจโย.\csromanspacer{} สวิตกฺกํ จุติจิตฺตญฺจ วิตกฺโก จ อวิตกฺกสฺส อุปปตฺติจิตฺตสฺส ฯเปฯ.\csromanspacer{} อาวชฺชนา จ วิตกฺโก จ ปญฺจนฺนํ วิญฺญาณานํ ฯเปฯ.\csromanspacer{} สวิตกฺกา ขนฺธา จ วิตกฺโก จ อวิตกฺกสฺส วุฏฺฐานสฺส อนนฺตรปจฺจเยน ปจฺจโย.\csromanspacer{} ทุติยสฺส ฌานสฺส ปริกมฺมญฺจ วิตกฺโก จ ฯเปฯ. (เหฏฺฐา ลิขิตํ เลขํ อิมินา การเณน ทฏฺฐพฺพํ.) อนุโลมญฺจ วิตกฺโก จ อวิตกฺกาย ผลสมาปตฺติยา อนนฺตรปจฺจเยน ปจฺจโย. (2)}

\prose{สวิตกฺโก จ อวิตกฺโก จ ธมฺมา สวิตกฺกสฺส จ อวิตกฺกสฺส จ ธมฺมสฺส อนนฺตรปจฺจเยน ปจฺจโย.\csromanspacer{} ปุริมา ปุริมา สวิตกฺกา ขนฺธา จ วิตกฺโก จ ปจฺฉิมานํ ปจฺฉิมานํ สวิตกฺกานํ ขนฺธานํ วิตกฺกสฺส จ อนนฺตรปจฺจเยน ปจฺจโย. (3)}

\prose{สมนนฺตรปจฺจเยน ปจฺจโย.\csromanspacer{} สหชาตปจฺจเยน ปจฺจโย.\csromanspacer{} นว.\csromanspacer{} อญฺญมญฺญปจฺจเยน ปจฺจโย.\csromanspacer{} นว.\csromanspacer{} นิสฺสยปจฺจเยน ปจฺจโย.\csromanspacer{} นว.}

\csromanheader{2}{\textbf{อุปนิสฺสย}}
\proseitem{125}{สวิตกฺโก ธมฺโม สวิตกฺกสฺส ธมฺมสฺส อุปนิสฺสยปจฺจเยน ปจฺจโย.\csromanspacer{} อารมฺมณูปนิสฺสโย, อนนฺตรูปนิสฺสโย, ปกตูปนิสฺสโย ฯเปฯ.\csromanspacer{} ปกตูปนิสฺสโย\csromandash{}สวิตกฺกา ขนฺธา สวิตกฺกานํ ขนฺธานํ อุปนิสฺสยปจฺจเยน ปจฺจโย. (มูลํ กาตพฺพํ.) สวิตกฺกา ขนฺธา อวิตกฺกานํ ขนฺธานํ วิตกฺกสฺส จ อุปนิสฺสยปจฺจเยน ปจฺจโย. (มูลํ กาตพฺพํ.) สวิตกฺกา ขนฺธา สวิตกฺกานํ ขนฺธานํ วิตกฺกสฺส จ อุปนิสฺสยปจฺจเยน ปจฺจโย. (3)}

\csromanheader{2}{87. สวิตกฺกทุก}
\proseitem{126}{\csromanfolio{93}อวิตกฺโก ธมฺโม อวิตกฺกสฺส ธมฺมสฺส อุปนิสฺสยปจฺจเยน ปจฺจโย.\csromanspacer{} อารมฺมณูปนิสฺสโย, อนนฺตรูปนิสฺสโย, ปกตูปนิสฺสโย ฯเปฯ.\csromanspacer{} ปกตูปนิสฺสโย\csromandash{}อวิตกฺกํ สทฺธํ อุปนิสฺสาย อวิตกฺกํ ฌานํ อุปฺปาเทติ, มคฺคํ อุปฺปาเทติ, อภิญฺญํ อุปฺปาเทติ, สมาปตฺติํ อุปฺปาเทติ.\csromanspacer{} อวิตกฺกํ สีลํ ฯเปฯ ปญฺญํ.\csromanspacer{} กายิกํ สุขํ.\csromanspacer{} กายิกํ ทุกฺขํ.\csromanspacer{} อุตุํ.\csromanspacer{} โภชนํ.\csromanspacer{} เสนาสนํ วิตกฺกํ อุปนิสฺสาย อวิตกฺกํ ฌานํ อุปฺปาเทติ.\csromanspacer{} มคฺคํ ฯเปฯ.\csromanspacer{} อภิญฺญํ ฯเปฯ สมาปตฺติํ อุปฺปาเทติ.\csromanspacer{} อวิตกฺกา สทฺธา ฯเปฯ.\csromanspacer{} เสนาสนํ วิตกฺโก จ อวิตกฺกาย สทฺธาย ฯเปฯ ปญฺญาย.\csromanspacer{} กายิกสฺส สุขสฺส.\csromanspacer{} กายิกสฺส ทุกฺขสฺส.\csromanspacer{} อวิตกฺกสฺส มคฺคสฺส, ผลสมาปตฺติยา อุปนิสฺสยปจฺจเยน ปจฺจโย. (1)}

\prose{อวิตกฺโก ธมฺโม สวิตกฺกสฺส ธมฺมสฺส อุปนิสฺสยปจฺจเยน ปจฺจโย. (ตีณิปิ อุปนิสฺสยา สพฺพตฺถ กาตพฺพา.) อวิตกฺกํ สทฺธํ อุปนิสฺสาย ทานํ เทติ, สีลํ สมาทิยติ, อุโปสถกมฺมํ กโรติ.\csromanspacer{} สวิตกฺกํ ฌานํ อุปฺปาเทติ, วิปสฺสนํ ฯเปฯ มคฺคํ ฯเปฯ สมาปตฺติํ อุปฺปาเทติ, มานํ ชปฺเปติ, ทิฏฺฐิํ คณฺหาติ.\csromanspacer{} อวิตกฺกํ สีลํ ฯเปฯ.\csromanspacer{} เสนาสนํ วิตกฺกํ อุปนิสฺสาย ทานํ เทติ ฯเปฯ สมาปตฺติํ อุปฺปาเทติ, ปาณํ หนติ ฯเปฯ สํฆํ ภินฺทติ.\csromanspacer{} อวิตกฺกา สทฺธา ฯเปฯ.\csromanspacer{} เสนาสนํ วิตกฺโก จ สวิตกฺกาย สทฺธาย ฯเปฯ ปญฺญาย.\csromanspacer{} ราคสฺส ฯเปฯ.\csromanspacer{} ปตฺถนาย สวิตกฺกสฺส มคฺคสฺส, ผลสมาปตฺติยา อุปนิสฺสยปจฺจเยน ปจฺจโย. (2)}

\prose{อวิตกฺโก ธมฺโม สวิตกฺกสฺส จ อวิตกฺกสฺส จ ธมฺมสฺส อุปนิสฺสยปจฺจเยน ปจฺจโย.\csromanspacer{} อวิตกฺกํ สทฺธํ อุปนิสฺสาย ทานํ เทติ ฯเปฯ. (ทุติยวาเร ลิขิตปทา สพฺเพ กาตพฺพา.) สมาปตฺติํ อุปฺปาเทติ, มานํ ชปฺเปติ, ทิฏฺฐิํ คณฺหาติ.\csromanspacer{} สีลํ ฯเปฯ.\csromanspacer{} ปญฺญํ ฯเปฯ.\csromanspacer{} เสนาสนํ วิตกฺกํ อุปนิสฺสาย ทานํ เทติ ฯเปฯ ปาณํ หนติ ฯเปฯ สํฆํ ภินฺทติ.\csromanspacer{} อวิตกฺกา สทฺธา ฯเปฯ.\csromanspacer{} เสนาสนํ วิตกฺโก จ สวิตกฺกาย สทฺธาย ฯเปฯ ปญฺญาย.\csromanspacer{} ราคสฺส ฯเปฯ.\csromanspacer{} ปตฺถนาย สวิตกฺกสฺส มคฺคสฺส, ผลสมาปตฺติยา วิตกฺกสฺส จ อุปนิสฺสยปจฺจเยน ปจฺจโย. (3)}

\proseitem{127}{สวิตกฺโก จ อวิตกฺโก จ ธมฺมา สวิตกฺกสฺส ธมฺมสฺส อุปนิสฺสยปจฺจเยน ปจฺจโย.\csromanspacer{} สวิตกฺกา ขนฺธา จ วิตกฺโก จ สวิตกฺกานํ ขนฺธานํ อุปนิสฺสยปจฺจเยน ปจฺจโย. (มูลํ กาตพฺพํ.) สวิตกฺกา ขนฺธา จ วิตกฺโก จ อวิตกฺกานํ ขนฺธานํ วิตกฺกสฺส จ อุปนิสฺสยปจฺจเยน ปจฺจโย. \csromanfolio{94}(มูลํ กาตพฺพํ.) สวิตกฺกา ขนฺธา จ วิตกฺโก จ สวิตกฺกานํ ขนฺธานํ วิตกฺกสฺส จ อุปนิสฺสยปจฺจเยน ปจฺจโย. (3)}

\csromanheader{2}{\textbf{ปุเรชาต}}
\proseitem{128}{อวิตกฺโก ธมฺโม อวิตกฺกสฺส ธมฺมสฺส ปุเรชาตปจฺจเยน ปจฺจโย.\csromanspacer{} อารมฺมณปุเรชาตํ, วตฺถุปุเรชาตํ.\csromanspacer{}\csromanspacer{}\textbf{อารมฺมณปุเรชาตํ\csromandash{}} จกฺขุํ ฯเปฯ.\csromanspacer{} วตฺถุํ อนิจฺจโต ฯเปฯ วิปสฺสติ, อสฺสาเทติ อภินนฺทติ, ตํ อารพฺภ วิตกฺโก อุปฺปชฺชติ.\csromanspacer{} ทิพฺเพน จกฺขุนา รูปํ ปสฺสติ ฯเปฯ.\csromanspacer{} โผฏฺฐพฺพายตนํ กายวิญฺญาณสฺส ฯเปฯ.\csromanspacer{} \textbf{วตฺถุปุเรชาตํ\csromandash{}}จกฺขายตนํ จกฺขุวิญฺญาณสฺส ฯเปฯ.\csromanspacer{} กายายตนํ กายวิญฺญาณสฺส ฯเปฯ.\csromanspacer{} วตฺถุ อวิตกฺกานํ ขนฺธานํ วิตกฺกสฺส จ ปุเรชาตปจฺจเยน ปจฺจโย. (1)}

\prose{อวิตกฺโก ธมฺโม สวิตกฺกสฺส ธมฺมสฺส ปุเรชาตปจฺจเยน ปจฺจโย.\csromanspacer{} อารมฺมณปุเรชาตํ, วตฺถุปุเรชาตํ.\csromanspacer{}\csromanspacer{}\textbf{อารมฺมณปุเรชาตํ\csromandash{}}จกฺขุํ ฯเปฯ.\csromanspacer{} วตฺถุํ อนิจฺจโต ฯเปฯ วิปสฺสติ, อสฺสาเทติ อภินนฺทติ, ตํ อารพฺภ สวิตกฺกา ขนฺธา อุปฺปชฺชนฺติ.\csromanspacer{} \textbf{วตฺถุปุเรชาตํ\csromandash{}}วตฺถุ สวิตกฺกานํ ขนฺธานํ ปุเรชาตปจฺจเยน ปจฺจโย. (2)}

\prose{อวิตกฺโก ธมฺโม สวิตกฺกสฺส จ อวิตกฺกสฺส จ ธมฺมสฺส ปุเรชาตปจฺจเยน ปจฺจโย.\csromanspacer{} อารมฺมณปุเรชาตํ, วตฺถุปุเรชาตํ.\csromanspacer{}\csromanspacer{}\textbf{อารมฺมณปุเรชาตํ\csromandash{}}จกฺขุํ ฯเปฯ.\csromanspacer{} วตฺถุํ อนิจฺจโต ฯเปฯ วิปสฺสติ, อสฺสาเทติ อภินนฺทติ, ตํ อารพฺภ สวิตกฺกา ขนฺธา จ วิตกฺโก จ อุปฺปชฺชนฺติ.\csromanspacer{} \textbf{วตฺถุปุเรชาตํ\csromandash{}}วตฺถุ สวิตกฺกานํ ขนฺธานํ วิตกฺกสฺส จ ปุเรชาตปจฺจเยน ปจฺจโย. (3)}

\csromanheader{2}{\textbf{ปจฺฉาชาตาเสวน}}
\proseitem{129}{สวิตกฺโก ธมฺโม อวิตกฺกสฺส ธมฺมสฺส ปจฺฉาชาตปจฺจเยน ปจฺจโย. (ตีณิ, ปจฺฉาชาตา.) อาเสวนปจฺจเยน ปจฺจโย.\csromanspacer{} นว.}

\csromanheader{2}{\textbf{กมฺมาทิ}}
\proseitem{130}{สวิตกฺโก ธมฺโม สวิตกฺกสฺส ธมฺมสฺส กมฺมปจฺจเยน ปจฺจโย.\csromanspacer{} \textbf{สหชาตา}, นานากฺขณิกา.\csromanspacer{}\csromanspacer{}\textbf{สหชาตา} สวิตกฺกา เจตนา สมฺปยุตฺตกานํ ขนฺธานํ กมฺมปจฺจเยน ปจฺจโย.\csromanspacer{} ปฏิสนฺธิกฺขเณ ฯเปฯ.\csromanspacer{} \textbf{นานากฺขณิกา} สวิตกฺกา เจตนา วิปากานํ สวิตกฺกานํ ขนฺธานํ กมฺมปจฺจเยน ปจฺจโย. (เอวํ จตฺตาริ, สหชาตาปิ นานากฺขณิกาปิ กาตพฺพา.)}

\csromanheader{2}{87. สวิตกฺกทุก}
\prose{\csromanfolio{95}วิปากปจฺจเยน ปจฺจโย.\csromanspacer{} นว.\csromanspacer{} อาหารปจฺจเยน ปจฺจโย.\csromanspacer{} จตฺตาริ.\csromanspacer{} อินฺทฺริยปจฺจเยน ปจฺจโย.\csromanspacer{} จตฺตาริ.\csromanspacer{} ฌานปจฺจเยน ปจฺจโย.\csromanspacer{} นว.\csromanspacer{} มคฺคปจฺจเยน ปจฺจโย.\csromanspacer{} นว.\csromanspacer{} สมฺปยุตฺตปจฺจเยน ปจฺจโย.\csromanspacer{} ฉ.}

\proseitem{131}{สวิตกฺโก ธมฺโม สวิตกฺกสฺส ธมฺมสฺส วิปฺปยุตฺตปจฺจเยน ปจฺจโย.\csromanspacer{} \textbf{สหชาตํ}, ปจฺฉาชาตํ. (สํขิตฺตํ.) (1)}

\prose{อวิตกฺโก ธมฺโม อวิตกฺกสฺส ธมฺมสฺส วิปฺปยุตฺตปจฺจเยน ปจฺจโย.\csromanspacer{} \textbf{สหชาตํ}, ปุเรชาตํ, ปจฺฉาชาตํ. (สํขิตฺตํ.).\csromanspacer{} อวิตกฺโก ธมฺโม สวิตกฺกสฺส ธมฺมสฺส วิปฺปยุตฺตปจฺจเยน ปจฺจโย.\csromanspacer{} \textbf{สหชาตํ}, ปุเรชาตํ. (สํขิตฺตํ.).\csromanspacer{} อวิตกฺโก ธมฺโม สวิตกฺกสฺส จ อวิตกฺกสฺส จ ธมฺมสฺส วิปฺปยุตฺตปจฺจเยน ปจฺจโย.\csromanspacer{} \textbf{สหชาตํ}, ปุเรชาตํ.\csromanspacer{}\csromanspacer{}\textbf{สหชาตํ} ปฏิสนฺธิกฺขเณ วตฺถุ วิตกฺกสฺส สมฺปยุตฺตกานญฺจ ขนฺธานํ วิปฺปยุตฺตปจฺจเยน ปจฺจโย.\csromanspacer{} \textbf{ปุเรชาตํ} วตฺถุ วิตกฺกสฺส สมฺปยุตฺตกานญฺจ ขนฺธานํ วิปฺปยุตฺตปจฺจเยน ปจฺจโย. (3)}

\prose{สวิตกฺโก จ อวิตกฺโก จ ธมฺมา อวิตกฺกสฺส ธมฺมสฺส วิปฺปยุตฺตปจฺจเยน ปจฺจโย.\csromanspacer{} \textbf{สหชาตํ}, ปจฺฉาชาตํ. (สํขิตฺตํ.)}

\csromanheader{2}{\textbf{อตฺถิ-อาทิ}}
\proseitem{132}{สวิตกฺโก ธมฺโม สวิตกฺกสฺส ธมฺมสฺส อตฺถิปจฺจเยน ปจฺจโย. (เอกํ, ปฏิจฺจสทิสํ.).\csromanspacer{} สวิตกฺโก ธมฺโม อวิตกฺกสฺส ธมฺมสฺส อตฺถิปจฺจเยน ปจฺจโย.\csromanspacer{} \textbf{สหชาตํ}, ปจฺฉาชาตํ. (สํขิตฺตํ.).\csromanspacer{} สวิตกฺโก ธมฺโม สวิตกฺกสฺส จ อวิตกฺกสฺส จ ธมฺมสฺส อตฺถิปจฺจเยน ปจฺจโย. (ปฏิจฺจสทิสํ.) (3)}

\prose{อวิตกฺโก ธมฺโม อวิตกฺกสฺส ธมฺมสฺส อตฺถิปจฺจเยน ปจฺจโย.\csromanspacer{} \textbf{สหชาตํ}, ปุเรชาตํ, ปจฺฉาชาตํ, อาหารํ, อินฺทฺริยํ. (สํขิตฺตํ.).\csromanspacer{} อวิตกฺโก ธมฺโม สวิตกฺกสฺส ธมฺมสฺส อตฺถิปจฺจเยน ปจฺจโย.\csromanspacer{} \textbf{สหชาตํ}, ปุเรชาตํ. (สํขิตฺตํ.).\csromanspacer{} อวิตกฺโก ธมฺโม สวิตกฺกสฺส จ อวิตกฺกสฺส จ ธมฺมสฺส อตฺถิปจฺจเยน ปจฺจโย.\csromanspacer{} \textbf{สหชาตํ}, ปุเรชาตํ.\csromanspacer{}\csromanspacer{}\textbf{สหชาโต} วิตกฺโก สมฺปยุตฺตกานํ ขนฺธานํ จิตฺตสมุฏฺฐานานญฺจ รูปานํ อตฺถิปจฺจเยน ปจฺจโย.\csromanspacer{} ปฏิสนฺธิกฺขเณ วิตกฺโก สมฺปยุตฺตกานํ ขนฺธานํ กฏตฺตา \csromanfolio{96}จ รูปานํ อตฺถิปจฺจเยน ปจฺจโย.\csromanspacer{} ปฏิสนฺธิกฺขเณ วตฺถุ วิตกฺกสฺส สมฺปยุตฺตกานญฺจ ขนฺธานํ อตฺถิปจฺจเยน ปจฺจโย.\csromanspacer{} \textbf{ปุเรชาตํ} จกฺขุํ ฯเปฯ.\csromanspacer{} วตฺถุํ อนิจฺจโต ฯเปฯ วิปสฺสติ, อสฺสาเทติ อภินนฺทติ, ตํ อารพฺภ วิตกฺโก จ สมฺปยุตฺตกา จ ขนฺธา อุปฺปชฺชนฺติ.\csromanspacer{} วตฺถุ วิตกฺกสฺส สมฺปยุตฺตกานญฺจ ขนฺธานํ อตฺถิปจฺจเยน ปจฺจโย. (3)}

\proseitem{133}{สวิตกฺโก จ อวิตกฺโก จ ธมฺมา สวิตกฺกสฺส ธมฺมสฺส อตฺถิปจฺจเยน ปจฺจโย.\csromanspacer{} สหชาตํ, ปุเรชาตํ.\csromanspacer{}\csromanspacer{}\textbf{สหชาโต} สวิตกฺโก เอโก ขนฺโธ จ วิตกฺโก จ ติณฺณนฺนํ ขนฺธานํ อตฺถิปจฺจเยน ปจฺจโย ฯเปฯ เทฺว ขนฺธา จ ฯเปฯ.\csromanspacer{} \textbf{สหชาโต} สวิตกฺโก เอโก ขนฺโธ จ วตฺถุ จ ติณฺณนฺนํ ขนฺธานํ อตฺถิปจฺจเยน ปจฺจโย ฯเปฯ เทฺว ขนฺธา จ ฯเปฯ. (ปฏิสนฺธิกฺขเณ สหชาตาปิ เทฺวปิ กาตพฺพา.) (1)}

\prose{สวิตกฺโก จ อวิตกฺโก จ ธมฺมา อวิตกฺกสฺส ธมฺมสฺส อตฺถิปจฺจเยน ปจฺจโย.\csromanspacer{} สหชาตํ, ปุเรชาตํ, ปจฺฉาชาตํ, อาหารํ, อินฺทฺริยํ.\csromanspacer{}\csromanspacer{}\textbf{สหชาตา} สวิตกฺกา ขนฺธา จ วิตกฺโก จ จิตฺตสมุฏฺฐานานํ รูปานํ อตฺถิปจฺจเยน ปจฺจโย.\csromanspacer{} \textbf{สหชาตา} สวิตกฺกา ขนฺธา จ วตฺถุ จ วิตกฺกสฺส อตฺถิปจฺจเยน ปจฺจโย. (ปฏิสนฺธิกฺขเณ.\csromanspacer{} ตีณิ.) \textbf{ปจฺฉาชาตา} สวิตกฺกา ขนฺธา จ วิตกฺโก จ ปุเรชาตสฺส อิมสฺส กายสฺส อตฺถิปจฺจเยน ปจฺจโย.\csromanspacer{} \textbf{ปจฺฉาชาตา} สวิตกฺกา ขนฺธา จ วิตกฺโก จ กพฬีกาโร อาหาโร จ อิมสฺส กายสฺส อตฺถิปจฺจเยน ปจฺจโย.\csromanspacer{} \textbf{ปจฺฉาชาตา} สวิตกฺกา ขนฺธา จ วิตกฺโก จ รูปชีวิตินฺทฺริยญฺจ กฏตฺตารูปานํ อตฺถิปจฺจเยน ปจฺจโย. (2)}

\prose{สวิตกฺโก จ อวิตกฺโก จ ธมฺมา สวิตกฺกสฺส จ อวิตกฺกสฺส จ ธมฺมสฺส อตฺถิปจฺจเยน ปจฺจโย.\csromanspacer{} สหชาตํ, ปุเรชาตํ.\csromanspacer{}\csromanspacer{}\textbf{สหชาโต} สวิตกฺโก เอโก ขนฺโธ จ วิตกฺโก จ ติณฺณนฺนํ ขนฺธานํ จิตฺตสมุฏฺฐานานญฺจ รูปานํ อตฺถิปจฺจเยน ปจฺจโย ฯเปฯ เทฺว ขนฺธา จ ฯเปฯ.\csromanspacer{} \textbf{สหชาโต} สวิตกฺโก เอโก ขนฺโธ จ วตฺถุ จ ติณฺณนฺนํ ขนฺธานํ วิตกฺกสฺส จ อตฺถิปจฺจเยน ปจฺจโย ฯเปฯ เทฺว ขนฺธา จ ฯเปฯ. (ปฏิสนฺธิยาปิ เทฺว.) (3)}

\prose{นตฺถิปจฺจเยน ปจฺจโย.\csromanspacer{} วิคตปจฺจเยน ปจฺจโย.\csromanspacer{} อวิคตปจฺจเยน ปจฺจโย.}

\csromanheader{2}{87. สวิตกฺกทุก \textbf{1. ปจฺจยานุโลม 2. สงฺขฺยาวาร}}
\csromanheader{2}{\textbf{สุทฺธ}}
\proseitem{134}{\csromanfolio{97}เหตุยา จตฺตาริ, อารมฺมเณ นว, อธิปติยา นว, อนนฺตเร นว, สมนนฺตเร นว, สหชาเต นว, อญฺญมญฺเญ นว, นิสฺสเย นว, อุปนิสฺสเย นว, ปุเรชาเต ตีณิ, ปจฺฉาชาเต ตีณิ, อาเสวเน นว, กมฺเม จตฺตาริ, วิปาเก นว, อาหาเร จตฺตาริ, อินฺทฺริเย จตฺตาริ, ฌาเน นว, มคฺเค นว, สมฺปยุตฺเต ฉ, วิปฺปยุตฺเต ปญฺจ, อตฺถิยา นว, นตฺถิยา นว, วิคเต นว, อวิคเต นว.}

\vspace{\tipitakaparskip}
\csromancenter{อนุโลมํ.}
\csromanheader{2}{2. ปจฺจนียุทฺธาร}
\proseitem{135}{สวิตกฺโก ธมฺโม สวิตกฺกสฺส ธมฺมสฺส อารมฺมณปจฺจเยน ปจฺจโย.\csromanspacer{} สหชาตปจฺจเยน ปจฺจโย.\csromanspacer{} อุปนิสฺสยปจฺจเยน ปจฺจโย.\csromanspacer{} กมฺมปจฺจเยน ปจฺจโย.\csromanspacer{}\csromanspacer{}สวิตกฺโก ธมฺโม อวิตกฺกสฺส ธมฺมสฺส อารมฺมณปจฺจเยน ปจฺจโย.\csromanspacer{} สหชาตปจฺจเยน ปจฺจโย.\csromanspacer{} อุปนิสฺสยปจฺจเยน ปจฺจโย.\csromanspacer{} ปจฺฉาชาตปจฺจเยน ปจฺจโย.\csromanspacer{} กมฺมปจฺจเยน ปจฺจโย.\csromanspacer{}\csromanspacer{}สวิตกฺโก ธมฺโม สวิตกฺกสฺส จ อวิตกฺกสฺส จ ธมฺมสฺส อารมฺมณปจฺจเยน ปจฺจโย.\csromanspacer{} สหชาตปจฺจเยน ปจฺจโย.\csromanspacer{} อุปนิสฺสยปจฺจเยน ปจฺจโย.\csromanspacer{} กมฺมปจฺจเยน ปจฺจโย. (3)}

\prose{อวิตกฺโก ธมฺโม อวิตกฺกสฺส ธมฺมสฺส อารมฺมณปจฺจเยน ปจฺจโย.\csromanspacer{} สหชาตปจฺจเยน ปจฺจโย.\csromanspacer{} อุปนิสฺสยปจฺจเยน ปจฺจโย.\csromanspacer{} ปุเรชาตปจฺจเยน ปจฺจโย.\csromanspacer{} ปจฺฉาชาตปจฺจเยน ปจฺจโย.\csromanspacer{} กมฺมปจฺจเยน ปจฺจโย.\csromanspacer{} อาหารปจฺจเยน ปจฺจโย.\csromanspacer{} อินฺทฺริยปจฺจเยน ปจฺจโย.\csromanspacer{}\csromanspacer{}อวิตกฺโก ธมฺโม สวิตกฺกสฺส ธมฺมสฺส อารมฺมณปจฺจเยน ปจฺจโย.\csromanspacer{} สหชาตปจฺจเยน ปจฺจโย.\csromanspacer{} อุปนิสฺสยปจฺจเยน ปจฺจโย.\csromanspacer{} ปุเรชาตปจฺจเยน ปจฺจโย.\csromanspacer{}\csromanspacer{}อวิตกฺโก ธมฺโม สวิตกฺกสฺส จ อวิตกฺกสฺส จ ธมฺมสฺส อารมฺมณปจฺจเยน ปจฺจโย.\csromanspacer{} สหชาตปจฺจเยน ปจฺจโย.\csromanspacer{} อุปนิสฺสยปจฺจเยน ปจฺจโย.\csromanspacer{} ปุเรชาตปจฺจเยน ปจฺจโย. (3)}

\proseitem{137}{\csromanfolio{98}สวิตกฺโก จ อวิตกฺโก จ ธมฺมา สวิตกฺกสฺส ธมฺมสฺส อารมฺมณปจฺจเยน ปจฺจโย.\csromanspacer{} สหชาตปจฺจเยน ปจฺจโย.\csromanspacer{} อุปนิสฺสยปจฺจเยน ปจฺจโย.\csromanspacer{}\csromanspacer{}สวิตกฺโก จ อวิตกฺโก จ ธมฺมา อวิตกฺกสฺส ธมฺมสฺส อารมฺมณปจฺจเยน ปจฺจโย.\csromanspacer{} สหชาตปจฺจเยน ปจฺจโย.\csromanspacer{} อุปนิสฺสยปจฺจเยน ปจฺจโย.\csromanspacer{} ปจฺฉาชาตปจฺจเยน ปจฺจโย.\csromanspacer{}\csromanspacer{}สวิตกฺโก จ อวิตกฺโก จ ธมฺมา สวิตกฺกสฺส จ อวิตกฺกสฺส จ ธมฺมสฺส อารมฺมณปจฺจเยน ปจฺจโย.\csromanspacer{} สหชาตปจฺจเยน ปจฺจโย.\csromanspacer{} อุปนิสฺสยปจฺจเยน ปจฺจโย. (3)}

\csromanheader{2}{2. ปจฺจยปจฺจนีย 2. สงฺขฺยาวาร}
\vspace{\tipitakaparskip}
\csromancenter{นเหตุยา นว, น-อารมฺมเณ นว, (สพฺพตฺถ นว,) โน-อวิคเต นว.}
\csromanheader{2}{3. ปจฺจยานุโลมปจฺจนีย}
\proseitem{139}{เหตุปจฺจยา น-อารมฺมเณ จตฺตาริ ฯเปฯ นสมนนฺตเร จตฺตาริ, น-อญฺญมญฺเญ เทฺว, น-อุปนิสฺสเย จตฺตาริ ฯเปฯ นสมฺปยุตฺเต เทฺว, นวิปฺปยุตฺเต จตฺตาริ, โนนตฺถิยา จตฺตาริ, โนวิคเต จตฺตาริ.}

\csromanheader{2}{4. ปจฺจยปจฺจนียานุโลม}
\proseitem{140}{นเหตุปจฺจยา อารมฺมเณ นว, อธิปติยา นว, (อนุโลมมาติกา วิตฺถาเรตพฺพา) ฯเปฯ อวิคเต นว.}

\nitthitam{สวิตกฺกทุกํ นิฏฺฐิตํ.}
\csromanmatikaanchor{mtk.17}
\csromantocmarkref{section}{88. สวิจารทุก}{mtk.17}
\bookmark[dest={mtk.17},level=1]{88. สวิจารทุก}
\csromanheader{1}{88. สวิจารทุก 1. ปฏิจฺจวาร}
\proseitem{141}{สวิจารํ ธมฺมํ ปฏิจฺจ สวิจาโร ธมฺโม อุปฺปชฺชติ เหตุปจฺจยา.\csromanspacer{} สวิจารํ เอกํ ขนฺธํ ปฏิจฺจ ตโย ขนฺธา ฯเปฯ เทฺว ขนฺเธ ฯเปฯ.\csromanspacer{} ปฏิสนฺธิกฺขเณ ฯเปฯ. (ยถา สวิตกฺกทุกํ, เอวํ กาตพฺพํ, นินฺนานากรณํ.\csromanspacer{} อิธ มคฺเค จตฺตาริ กาตพฺพานิ.\csromanspacer{} สวิจารทุเก อิมํ นานากรณํ.)}

\nitthitam{สวิจารทุกํ นิฏฺฐิตํ.}
\csromanmatikaanchor{mtk.18}
\csromantocmarkref{section}{89. สปฺปีติกทุก}{mtk.18}
\bookmark[dest={mtk.18},level=1]{89. สปฺปีติกทุก}
\csromanheader{1}{89. สปฺปีติกทุก \textbf{89. สปฺปีติกทุก 1. ปฏิจฺจวาร}}
\proseitem{142}{\csromanfolio{99}สปฺปีติกํ ธมฺมํ ปฏิจฺจ สปฺปีติโก ธมฺโม อุปฺปชฺชติ เหตุปจฺจยา.\csromanspacer{} สปฺปีติกํ เอกํ ขนฺธํ ปฏิจฺจ ตโย ขนฺธา ฯเปฯ เทฺว ขนฺเธ ฯเปฯ.\csromanspacer{} ปฏิสนฺธิกฺขเณ ฯเปฯ.\csromanspacer{} สปฺปีติกํ ธมฺมํ ปฏิจฺจ อปฺปีติโก ธมฺโม อุปฺปชฺชติ เหตุปจฺจยา.\csromanspacer{} สปฺปีติเก ขนฺเธ ปฏิจฺจ ปีติ จ จิตฺตสมุฏฺฐานญฺจ รูปํ.\csromanspacer{} ปฏิสนฺธิกฺขเณ ฯเปฯ.\csromanspacer{} สปฺปีติกํ ธมฺมํ ปฏิจฺจ สปฺปีติโก จ อปฺปีติโก จ ธมฺมา อุปฺปชฺชนฺติ เหตุปจฺจยา.\csromanspacer{} สปฺปีติกํ เอกํ ขนฺธํ ปฏิจฺจ ตโย ขนฺธา ปีติ จ จิตฺตสมุฏฺฐานญฺจ รูปํ ฯเปฯ เทฺว ขนฺเธ ฯเปฯ.\csromanspacer{} ปฏิสนฺธิกฺขเณ ฯเปฯ. (3)}

\prose{อปฺปีติกํ ธมฺมํ ปฏิจฺจ อปฺปีติโก ธมฺโม อุปฺปชฺชติ เหตุปจฺจยา.\csromanspacer{} อปฺปีติกํ เอกํ ขนฺธํ ปฏิจฺจ ตโย ขนฺธา จิตฺตสมุฏฺฐานญฺจ รูปํ ฯเปฯ เทฺว ขนฺเธ ฯเปฯ.\csromanspacer{} ปีติํ ปฏิจฺจ จิตฺตสมุฏฺฐานํ รูปํ.\csromanspacer{} ปฏิสนฺธิกฺขเณ อปฺปีติกํ เอกํ ขนฺธํ ปฏิจฺจ ตโย ขนฺธา กฏตฺตา จ รูปํ.\csromanspacer{} เทฺว ขนฺเธ ฯเปฯ.\csromanspacer{} ปีติํ ปฏิจฺจ กฏตฺตารูปํ.\csromanspacer{} ขนฺเธ ปฏิจฺจ วตฺถุ, วตฺถุํ ปฏิจฺจ ขนฺธา.\csromanspacer{} ปีติํ ปฏิจฺจ วตฺถุ, วตฺถุํ ปฏิจฺจ ปีติ.\csromanspacer{} เอกํ มหาภูตํ ฯเปฯ. (ยถา สวิตกฺกทุกํ สพฺพตฺถ, เอวํ สปฺปีติกทุกํ กาตพฺพํ.\csromanspacer{} สพฺพตฺถ ปวตฺติปฏิสนฺธิ นวปิ ปญฺหา.)}

\csromanheader{2}{1. ปจฺจยานุโลม 2. สงฺขฺยาวาร}
\vspace{\tipitakaparskip}
\csromancenter{\textbf{เหตุ}ยา นว, อารมฺมเณ นว, อธิปติยา นว ฯเปฯ ปุเรชาเต ฉ ฯเปฯ กมฺเม นว, วิปาเก นว ฯเปฯ อวิคเต นว.}
\csromansectionrule
\csromanheader{2}{2. ปจฺจยปจฺจนีย 1. วิภงฺควาร}
\csromanheader{2}{\textbf{นเหตุ}}
\proseitem{144}{สปฺปีติกํ ธมฺมํ ปฏิจฺจ สปฺปีติโก ธมฺโม อุปฺปชฺชติ นเหตุปจฺจยา.\csromanspacer{} อเหตุกํ สปฺปีติกํ เอกํ ขนฺธํ ปฏิจฺจ ตโย ขนฺธา ฯเปฯ เทฺว ขนฺเธ ฯเปฯ.\csromanspacer{} สปฺปีติกํ ธมฺมํ ปฏิจฺจ อปฺปีติโก ธมฺโม อุปฺปชฺชติ นเหตุปจฺจยา.\csromanspacer{} อเหตุเก สปฺปีติเก ขนฺเธ ปฏิจฺจ ปีติ จ จิตฺตสมุฏฺฐานญฺจ รูปํ.\csromanspacer{}\csromanspacer{} \csromanfolio{100}สปฺปีติกํ ธมฺมํ ปฏิจฺจ สปฺปีติโก จ อปฺปีติโก จ ธมฺมา อุปฺปชฺชนฺติ นเหตุปจฺจยา.\csromanspacer{} อเหตุกํ สปฺปีติกํ เอกํ ขนฺธํ ปฏิจฺจ ตโย ขนฺธา ปีติ จ จิตฺตสมุฏฺฐานญฺจ รูปํ ฯเปฯ เทฺว ขนฺเธ ฯเปฯ. (3)}

\prose{อปฺปีติกํ ธมฺมํ ปฏิจฺจ อปฺปีติโก ธมฺโม อุปฺปชฺชติ นเหตุปจฺจยา.\csromanspacer{} อเหตุกํ อปฺปีติกํ เอกํ ขนฺธํ ปฏิจฺจ ตโย ขนฺธา จิตฺตสมุฏฺฐานญฺจ รูปํ ฯเปฯ เทฺว ขนฺเธ ฯเปฯ.\csromanspacer{} ปีติํ ปฏิจฺจ จิตฺตสมุฏฺฐานํ รูปํ.\csromanspacer{} อเหตุกปฏิสนฺธิกฺขเณ อปฺปีติกํ เอกํ ขนฺธํ ปฏิจฺจ ตโย ขนฺธา กฏตฺตา จ รูปํ ฯเปฯ เทฺว ขนฺเธ ปฏิจฺจ เทฺว ขนฺธา กฏตฺตา จ รูปํ.\csromanspacer{} ขนฺเธ ปฏิจฺจ วตฺถุ, วตฺถุํ ปฏิจฺจ ขนฺธา.\csromanspacer{} เอกํ มหาภูตํ ฯเปฯ. (ยาว อสญฺญสตฺตาปิ.) วิจิกิจฺฉาสหคเต อุทฺธจฺจสหคเต ขนฺเธ ปฏิจฺจ วิจิกิจฺฉาสหคโต อุทฺธจฺจสหคโต โมโห.\csromanspacer{}\csromanspacer{}อปฺปีติกํ ธมฺมํ ปฏิจฺจ สปฺปีติโก ธมฺโม อุปฺปชฺชติ นเหตุปจฺจยา.\csromanspacer{} อเหตุกํ ปีติํ ปฏิจฺจ สปฺปีติกา ขนฺธา. (มูลํ กาตพฺพํ.) ปีติํ ปฏิจฺจ สปฺปีติกา ขนฺธา จิตฺตสมุฏฺฐานญฺจ รูปํ. (3)}

\proseitem{145}{สปฺปีติกญฺจ อปฺปีติกญฺจ ธมฺมํ ปฏิจฺจ สปฺปีติโก ธมฺโม อุปฺปชฺชติ นเหตุปจฺจยา.\csromanspacer{} อเหตุกํ สปฺปีติกํ เอกํ ขนฺธญฺจ ปีติญฺจ ปฏิจฺจ ตโย ขนฺธา ฯเปฯ เทฺว ขนฺเธ จ ฯเปฯ.\csromanspacer{} สปฺปีติกญฺจ อปฺปีติกญฺจ ธมฺมํ ปฏิจฺจ อปฺปีติโก ธมฺโม อุปฺปชฺชติ นเหตุปจฺจยา.\csromanspacer{} อเหตุเก สปฺปีติเก ขนฺเธ จ ปีติญฺจ ปฏิจฺจ จิตฺตสมุฏฺฐานํ รูปํ.\csromanspacer{} อเหตุเก สปฺปีติเก ขนฺเธ จ มหาภูเต จ ปฏิจฺจ จิตฺตสมุฏฺฐานํ รูปํ.\csromanspacer{}\csromanspacer{}สปฺปีติกญฺจ อปฺปีติกญฺจ ธมฺมํ ปฏิจฺจ สปฺปีติโก จ อปฺปีติโก จ ธมฺมา อุปฺปชฺชนฺติ นเหตุปจฺจยา.\csromanspacer{} อเหตุกํ สปฺปีติกํ เอกํ ขนฺธญฺจ ปีติญฺจ ปฏิจฺจ ตโย ขนฺธา จิตฺตสมุฏฺฐานญฺจ รูปํ ฯเปฯ เทฺว ขนฺเธ จ ฯเปฯ.\csromanspacer{} อเหตุกํ สปฺปีติกํ เอกํ ขนฺธญฺจ ปีติญฺจ ปฏิจฺจ ตโย ขนฺธา ฯเปฯ เทฺว ขนฺเธ จ ฯเปฯ.\csromanspacer{} อเหตุเก สปฺปีติเก ขนฺเธ จ ปีติญฺจ มหาภูเต จ ปฏิจฺจ จิตฺตสมุฏฺฐานํ รูปํ. (3)}

\csromanheader{2}{2. ปจฺจยปจฺจนีย 2. สงฺขฺยาวาร}
\csromanheader{2}{\textbf{สุทฺธ}}
\proseitem{146}{นเหตุยา นว, น-อารมฺมเณ ตีณิ, น-อธิปติยา นว, น-อนนฺตเร ตีณิ ฯเปฯ น-อุปนิสฺสเย ตีณิ, นปุเรชาเต นว, นปจฺฉาชาเต นว, น-อาเสวเน นว, นกมฺเม จตฺตาริ, นวิปาเก นว, น-อาหาเร เอกํ,}

\vspace{\tipitakaparskip}
\csromancenter{สปฺปีติกทุก น-อินฺทฺริเย เอกํ, นฌาเน เอกํ, นมคฺเค นว, นสมฺปยุตฺเต ตีณิ, นวิปฺปยุตฺเต ฉ, โนนตฺถิยา ตีณิ, โนวิคเต ตีณิ.}
\csromancenter{(เอวํ อิตเร เทฺว คณนาปิ สหชาตวาโรปิ กาตพฺโพ.)}
\csromanheader{2}{89. สปฺปีติกทุก 3. ปจฺจยวาร}
\csromanheader{2}{1. ปจฺจยานุโลมาทิ}
\proseitem{147}{\csromanfolio{101}สปฺปีติกํ ธมฺมํ ปจฺจยา สปฺปีติโก ธมฺโม อุปฺปชฺชติ เหตุปจฺจยา. (สํขิตฺตํ.\csromanspacer{} ยถา สวิตกฺกทุเก อนุโลมปจฺจยวารํ, เอวํ ปวตฺติปฏิสนฺธิ นว ปญฺหา ปริปุณฺณา ปีติ นินฺนานากรณา.)}

\prose{เหตุยา นว, อารมฺมเณ นว, อธิปติยา นว ฯเปฯ อวิคเต นว.\csromanspacer{} อนุโลมํ.}

\prose{สปฺปีติกํ ธมฺมํ ปจฺจยา สปฺปีติโก ธมฺโม อุปฺปชฺชติ นเหตุปจฺจยา.\csromanspacer{} ตีณิ. (ปฏิจฺจสทิสา.)}

\prose{อปฺปีติกํ ธมฺมํ ปจฺจยา อปฺปีติโก ธมฺโม อุปฺปชฺชติ นเหตุปจฺจยา. (ปวตฺติปฏิสนฺธิ กาตพฺพา.\csromanspacer{} ปฏิจฺจวารสทิสา.\csromanspacer{} ยาว อสญฺญสตฺตา.) จกฺขายตนํ ปจฺจยา จกฺขุวิญฺญาณํ ฯเปฯ.\csromanspacer{} กายายตนํ ปจฺจยา กายวิญฺญาณํ.\csromanspacer{} วตฺถุํ ปจฺจยา อเหตุกา อปฺปีติกา ขนฺธา ปีติ จ.\csromanspacer{} วิจิกิจฺฉาสหคเต อุทฺธจฺจสหคเต ขนฺเธ จ วตฺถุญฺจ ปจฺจยา วิจิกิจฺฉาสหคโต อุทฺธจฺจสหคโต โมโห. (อนุโลมสทิสา.\csromanspacer{} นว ปญฺหา ปวตฺติเยว.\csromanspacer{} ปฏิสนฺธิ นตฺถิ.\csromanspacer{} เอโกเยว โมโห.)}

\prose{นเหตุยา นว, น-อารมฺมเณ ตีณิ, น-อธิปติยา นว, น-อนนฺตเร ตีณิ ฯเปฯ น-อุปนิสฺสเย ตีณิ, นปุเรชาเต นว, นปจฺฉาชาเต นว, น-อาเสวเน นว, นกมฺเม นว, นวิปาเก นว, น-อาหาเร เอกํ, น-อินฺทฺริเย เอกํ, นฌาเน เอกํ, นมคฺเค นว, นสมฺปยุตฺเต ตีณิ, นวิปฺปยุตฺเต ฉ, โนนตฺถิยา ตีณิ, โนวิคเต ตีณิ.}

\vspace{\tipitakaparskip}
\csromancenter{ปจฺจนียํ.}
\csromancenter{(เอวํ อิตเร เทฺว คณนาปิ สิสฺสยวาโรปิ กาตพฺโพ.)}
\csromanheader{2}{89. สปฺปีติกทุก 5. สํสฏฺฐวาร}
\csromanheader{2}{1. ปจฺจยานุโลมาทิ}
\proseitem{148}{\csromanfolio{102}สปฺปีติกํ ธมฺมํ สํสฏฺโฐ สปฺปีติโก ธมฺโม อุปฺปชฺชติ เหตุปจฺจยา ฯเปฯ.}

\prose{\textbf{เหตุ}ยา ฉ, อารมฺมเณ ฉ, (สพฺพตฺถ ฉ,) อวิคเต ฉ.\csromanspacer{} อนุโลมํ.}

\prose{นเหตุยา ฉ, น-อธิปติยา ฉ, นปุเรชาเต ฉ, นปจฺฉาชาเต ฉ, น-อาเสวเน ฉ, นกมฺเม จตฺตาริ, นวิปาเก ฉ, นฌาเน เอกํ, นมคฺเค ฉ, นวิปฺปยุตฺเต ฉ.}

\vspace{\tipitakaparskip}
\csromancenter{ปจฺจนียํ.}
\csromancenter{(เอวํ อิตเร เทฺว คณนาปิ สมฺปยุตฺตวาโรปิ กาตพฺโพ.)}
\csromanheader{2}{89. สปฺปีติกทุก 7. ปญฺหาวาร}
\proseitem{149}{สปฺปีติโก ธมฺโม สปฺปีติกสฺส ธมฺมสฺส เหตุปจฺจเยน ปจฺจโย.\csromanspacer{} สปฺปีติกา เหตู สมฺปยุตฺตกานํ ขนฺธานํ เหตุปจฺจเยน ปจฺจโย.\csromanspacer{} ปฏิสนฺธิกฺขเณ ฯเปฯ.\csromanspacer{} สปฺปีติโก ธมฺโม อปฺปีติกสฺส ธมฺมสฺส เหตุปจฺจเยน ปจฺจโย.\csromanspacer{} สปฺปีติกา เหตู ปีติยา จ จิตฺตสมุฏฺฐานานญฺจ รูปานํ เหตุปจฺจเยน ปจฺจโย.\csromanspacer{} ปฏิสนฺธิกฺขเณ ฯเปฯ. (มูลํ กาตพฺพํ.) สปฺปีติกา เหตู สมฺปยุตฺตกานํ ขนฺธานํ ปีติยา จ จิตฺตสมุฏฺฐานานญฺจ รูปานํ เหตุปจฺจเยน ปจฺจโย.\csromanspacer{} ปฏิสนฺธิกฺขเณ ฯเปฯ. (3)}

\prose{อปฺปีติโก ธมฺโม อปฺปีติกสฺส ธมฺมสฺส เหตุปจฺจเยน ปจฺจโย.\csromanspacer{} อปฺปีติกา เหตู สมฺปยุตฺตกานํ ขนฺธานํ จิตฺตสมุฏฺฐานานญฺจ รูปานํ เหตุปจฺจเยน ปจฺจโย.\csromanspacer{} ปฏิสนฺธิกฺขเณ ฯเปฯ. (1)}

\csromanheader{2}{89. สปฺปีติกทุก \textbf{อารมฺมณ}}
\proseitem{150}{\csromanfolio{103}สปฺปีติโก ธมฺโม สปฺปีติกสฺส ธมฺมสฺส อารมฺมณปจฺจเยน ปจฺจโย.\csromanspacer{} สปฺปีติเก ขนฺเธ อารพฺภ สปฺปีติกา ขนฺธา อุปฺปชฺชนฺติ. (มูลํ กาตพฺพํ.) สปฺปีติเก ขนฺเธ อารพฺภ อปฺปีติกา ขนฺธา จ ปีติ จ อุปฺปชฺชนฺติ. (มูลํ กาตพฺพํ.) สปฺปีติเก ขนฺเธ อารพฺภ สปฺปีติกา ขนฺธา จ ปีติ จ อุปฺปชฺชนฺติ. (3)}

\proseitem{151}{อปฺปีติโก ธมฺโม อปฺปีติกสฺส ธมฺมสฺส อรมฺมณปจฺจเยน ปจฺจโย.\csromanspacer{} อปฺปีติเกน จิตฺเตน ทานํ ทตฺวา, สีลํ ฯเปฯ อุโปสถกมฺมํ กตฺวา อปฺปีติเกน จิตฺเตน ปจฺจเวกฺขติ, อสฺสาเทติ อภินนฺทติ, ตํ อารพฺภ อปฺปีติโก ราโค อุปฺปชฺชติ.\csromanspacer{} ทิฏฺฐิ อุปฺปชฺชติ, วิจิกิจฺฉา อุปฺปชฺชติ.\csromanspacer{} อุทฺธจฺจํ อุปฺปชฺชติ.\csromanspacer{} โทมนสฺสํ อุปฺปชฺชติ.\csromanspacer{} อปฺปีติกา ฌานา วุฏฺฐหิตฺวา ฯเปฯ มคฺคา วุฏฺฐหิตฺวา ฯเปฯ ผลา วุฏฺฐหิตฺวา อปฺปีติเกน จิตฺเตน ผลํ ปจฺจเวกฺขติ.\csromanspacer{} อริยา อปฺปีติเกน จิตฺเตน นิพฺพานํ ปจฺจเวกฺขนฺติ.\csromanspacer{} นิพฺพานํ อปฺปีติกสฺส โคตฺรภุสฺส, โวทานสฺส, มคฺคสฺส, ผลสฺส, อาวชฺชนาย ปีติยา จ อารมฺมณปจฺจเยน ปจฺจโย.\csromanspacer{} อริยา อปฺปีติเกน จิตฺเตน อปฺปีติเก ปหีเน กิเลเส ฯเปฯ.\csromanspacer{} วิกฺขมฺภิเต กิเลเส ฯเปฯ.\csromanspacer{} ปุพฺเพ ฯเปฯ.\csromanspacer{} จกฺขุํ ฯเปฯ.\csromanspacer{} วตฺถุํ อปฺปีติเก ขนฺเธ จ ปีติญฺจ อปฺปีติเกน จิตฺเตน อนิจฺจโต ฯเปฯ วิปสฺสติ, อสฺสาเทติ อภินนฺทติ, ตํ อารพฺภ อปฺปีติโก ราโค อุปฺปชฺชติ ฯเปฯ โทมนสฺสํ อุปฺปชฺชติ.\csromanspacer{} ทิพฺเพน จกฺขุนา รูปํ ปสฺสติ ฯเปฯ.\csromanspacer{} โผฏฺฐพฺพายตนํ กายวิญฺญาณสฺส อารมฺมณปจฺจเยน ปจฺจโย.\csromanspacer{} อปฺปีติกา ขนฺธา อิทฺธิวิธญาณสฺส, เจโตปริยญาณสฺส, ปุพฺเพนิวาสานุสฺสติญาณสฺส, ยถากมฺมูปคญาณสฺส, อนาคตํสญาณสฺส, อาวชฺชานาย, ปีติยา จ อารมฺมณปจฺจเยน ปจฺจโย.\csromanspacer{} อปฺปีติเก ขนฺเธ จ ปีติญฺจ อารพฺภ อปฺปีติกา ขนฺธา จ ปีติ จ อุปฺปชฺชนฺติ. (1)}

\prose{อปฺปีติโก ธมฺโม สปฺปีติกสฺส ธมฺมสฺส อารมฺมณปจฺจเยน ปจฺจโย.\csromanspacer{} อปฺปีติเกน จิตฺเตน ทานํ ทตฺวา, สีลํ ฯเปฯ อุโปสถกมฺมํ ฯเปฯ.\csromanspacer{} สปฺปีติเกน จิตฺเตน ปจฺจเวกฺขติ, อสฺสาเทติ อภินนฺทติ, ตํ อารพฺภ สปฺปีติโก ราโค อุปฺปชฺชติ.\csromanspacer{} ทิฏฺฐิ อุปฺปชฺชติ.\csromanspacer{} อปฺปีติกา ฌานา วุฏฺฐหิตฺวา ฯเปฯ มคฺคา วุฏฺฐหิตฺวา ฯเปฯ ผลา วุฏฺฐหิตฺวา สปฺปีติเกน จิตฺเตน ผลํ ปจฺจเวกฺขติ.\csromanspacer{} อริยา สปฺปีติเกน จิตฺเตน นิพฺพานํ ปจฺจเวกฺขนฺติ.\csromanspacer{} นิพฺพานํ สปฺปีติกสฺส โคตฺรภุสฺส, โวทานสฺส, มคฺคสฺส, ผลสฺส อารมฺมณปจฺจเยน ปจฺจโย.\csromanspacer{} อริยา สปฺปีติเกน จิตฺเตน อปฺปีติเก ปหีเน กิเลเส ฯเปฯ.\csromanspacer{} วิกฺขมฺภิเต \csromanfolio{104}กิเลเส ฯเปฯ.\csromanspacer{} ปุพฺเพ ฯเปฯ.\csromanspacer{} จกฺขุํ ฯเปฯ.\csromanspacer{} วตฺถุํ อปฺปีติเก ขนฺเธ จ ปีติญฺจ สปฺปีติเกน จิตฺเตน อนิจฺจโต ฯเปฯ วิปสฺสติ, อสฺสาเทติ อภินนฺทติ, ตํ อารพฺภ สปฺปีติโก ราโค อุปฺปชฺชติ.\csromanspacer{} ทิฏฺฐิ อุปฺปชฺชติ.\csromanspacer{} อปฺปีติเก ขนฺเธ จ ปีติญฺจ อารพฺภ สปฺปีติกา ขนฺธา อุปฺปชฺชนฺติ. (2)}

\prose{อปฺปีติโก ธมฺโม สปฺปีติกสฺส จ อปฺปีติกสฺส จ ธมฺมสฺส อารมฺมณปจฺจเยน ปจฺจโย.\csromanspacer{} อปฺปีติเกน จิตฺเตน ทานํ ทตฺวา, สีลํ ฯเปฯ อุโปสถกมฺมํ ฯเปฯ.\csromanspacer{} สปฺปีติเกน จิตฺเตน ปจฺจเวกฺขติ, อสฺสาเทติ อภินนฺทติ, ตํ อารพฺภ สปฺปีติกา ขนฺธา จ ปีติ จ อุปฺปชฺชนฺติ.\csromanspacer{} อปฺปีติกา ฌานา ฯเปฯ มคฺคา ฯเปฯ ผลา วุฏฺฐหิตฺวา สปฺปีติเกน จิตฺเตน ผลํ ปจฺจเวกฺขติ.\csromanspacer{} อริยา สปฺปีติเกน จิตฺเตน นิพฺพานํ ปจฺจเวกฺขนฺติ.\csromanspacer{} นิพฺพานํ สปฺปีติกสฺส โคตฺรภุสฺส, โวทานสฺส, มคฺคสฺส, ผลสฺส, ปีติยา จ อารมฺมณปจฺจเยน ปจฺจโย.\csromanspacer{} อริยา สปฺปีติเกน จิตฺเตน อปฺปีติเก ปหีเน กิเลเส ปจฺจเวกฺขนฺติ.\csromanspacer{} วิกฺขมฺภิเต กิเลเส ฯเปฯ.\csromanspacer{} ปุพฺเพ ฯเปฯ.\csromanspacer{} จกฺขุํ ฯเปฯ.\csromanspacer{} วตฺถุํ อปฺปีติเก ขนฺเธ จ ปีติญฺจ สปฺปีติเกน จิตฺเตน อนิจฺจโต ฯเปฯ วิปสฺสติ, อสฺสาเทติ อภินนฺทติ, ตํ อารพฺภ สปฺปีติกา ขนฺธา จ ปีติ จ อุปฺปชฺชนฺติ.\csromanspacer{} อปฺปีติเก ขนฺเธ จ ปิติญฺจ อารพฺภ สปฺปีติกา ขนฺธา จ ปีติ จ อุปฺปชฺชนฺติ. (3)}

\proseitem{152}{สปฺปีติโก จ อปฺปีติโก จ ธมฺมา สปฺปีติกสฺส ธมฺมสฺส อารมฺมณปจฺจเยน ปจฺจโย.\csromanspacer{} สปฺปีติเก ขนฺเธ จ ปีติญฺจ อารพฺภ สปฺปีติกา ขนฺธา อุปฺปชฺชนฺติ. (มูลํ กาตพฺพํ.) สปฺปีติเก ขนฺเธ จ ปีติญฺจ อารพฺภ อปฺปีติกา ขนฺธา จ ปีติ จ อุปฺปชฺชนฺติ. (มูลํ กาตพฺพํ.) สปฺปีติเก ขนฺเธ จ ปีติญฺจ อารพฺภ สปฺปีติกา ขนฺธา จ ปีติ จ อุปฺปชฺชนฺติ. (3)}

\csromanheader{2}{\textbf{อธิปติ}}
\proseitem{153}{สปฺปีติโก ธมฺโม สปฺปีติกสฺส ธมฺมสฺส อธิปติปจฺจเยน ปจฺจโย.\csromanspacer{} อารมฺมณาธิปติ, สหชาตาธิปติ.\csromanspacer{}\csromanspacer{}\textbf{อารมฺมณาธิปติ\csromandash{}}สปฺปีติเก ขนฺเธ ครุํ กตฺวา สปฺปีติกา ขนฺธา อุปฺปชฺชนฺติ.\csromanspacer{} \textbf{สหชาตาธิปติ\csromandash{}} สปฺปีติกาธิปติ สมฺปยุตฺตกานํ ขนฺธานํ อธิปติปจฺจเยน ปจฺจโย. (1)}

\prose{สปฺปีติโก ธมฺโม อปฺปีติกสฺส ธมฺมสฺส อธิปติปจฺจเยน ปจฺจโย.\csromanspacer{} อารมฺมณาธิปติ, สหชาตาธิปติ.\csromanspacer{}\csromanspacer{}\textbf{อารมฺมณาธิปติ\csromandash{}}สปฺปีติเก ขนฺเธ ครุํ กตฺวา อปฺปีติกา ขนฺธา จ ปีติ จ อุปฺปชฺชนฺติ.\csromanspacer{} \textbf{สหชาตาธิปติ\csromandash{}}}

\vspace{\tipitakaparskip}
\csromancenter{สปฺปีติกทุก สปฺปีติกาธิปติ ปีติยา จ จิตฺตสมุฏฺฐานานญฺจ รูปานํ อธิปติปจฺจเยน ปจฺจโย. (2)}
\prose{\csromanfolio{105}สปฺปีติโก ธมฺโม สปฺปีติกสฺส จ อปฺปีติกสฺส จ ธมฺมสฺส อธิปติปจฺจเยน ปจฺจโย.\csromanspacer{} อารมฺมณาธิปติ, สหชาตาธิปติ.\csromanspacer{}\csromanspacer{}\textbf{อารมฺมณาธิปติ\csromandash{}}สปฺปีติเก ขนฺเธ ครุํ กตฺวา สปฺปีติกา ขนฺธา จ ปีติ จ อุปฺปชฺชนฺติ.\csromanspacer{} \textbf{สหชาตาธิปติ\csromandash{}}สปฺปีติกาธิปติ สมฺปยุตฺตกานํ ขนฺธานํ ปีติยา จ จิตฺตสมุฏฺฐานานญฺจ รูปานํ อธิปติปจฺจเยน ปจฺจโย. (3)}

\proseitem{154}{อปฺปีติโก ธมฺโม อปฺปีติกสฺส ธมฺมสฺส อธิปติปจฺจเยน ปจฺจโย.\csromanspacer{} อารมฺมณาธิปติ, สหชาตาธิปติ.\csromanspacer{}\csromanspacer{}\textbf{อารมฺมณาธิปติ\csromandash{}} อปฺปีติเกน จิตฺเตน ทานํ ฯเปฯ สีลํ ฯเปฯ อุโปสถกมฺมํ ฯเปฯ.\csromanspacer{} อปฺปีติเกน จิตฺเตน ตํ ครุํ กตฺวา ปจฺจเวกฺขติ, อสฺสาเทติ อภินนฺทติ, ตํ ครุํ กตฺวา อปฺปีติโก ราโค อุปฺปชฺชติ, ทิฏฺฐิ อุปฺปชฺชติ.\csromanspacer{} อปฺปีติกา ฌานา ฯเปฯ มคฺคา ฯเปฯ ผลา วุฏฺฐหิตฺวา อปฺปีติเกน จิตฺเตน ผลํ ครุํ กตฺวา ปจฺจเวกฺขติ.\csromanspacer{} อริยา อปฺปีติเกน จิตฺเตน นิพฺพานํ ครุํ กตฺวา ปจฺจเวกฺขนฺติ.\csromanspacer{} นิพฺพานํ อปฺปีติกสฺส โคตฺรภุสฺส, โวทานสฺส, มคฺคสฺส, ผลสฺส อธิปติปจฺจเยน ปจฺจโย.\csromanspacer{} จกฺขุํ ฯเปฯ.\csromanspacer{} วตฺถุํ อปฺปีติเก ขนฺเธ จ ปีติญฺจ อปฺปีติเกน จิตฺเตน ครุํ กตฺวา อสฺสาเทติ อภินนฺทติ, ตํ ครุํ กตฺวา อปฺปีติโก ราโค อุปฺปชฺชติ.\csromanspacer{} ทิฏฺฐิ อุปฺปชฺชติ.\csromanspacer{} อปฺปีติเก ขนฺเธ จ ปีติญฺจ ครุํ กตฺวา อปฺปีติกา ขนฺธา จ ปีติ จ อุปฺปชฺชนฺติ.\csromanspacer{} \textbf{สหชาตาธิปติ\csromandash{}}อปฺปีติกาธิปติ สมฺปยุตฺตกานํ ขนฺธานํ จิตฺตสมุฏฺฐานานญฺจ รูปานํ อธิปติปจฺจเยน ปจฺจโย. (1)}

\prose{อปฺปีติโก ธมฺโม สปฺปีติกสฺส ธมฺมสฺส อธิปติปจฺจเยน ปจฺจโย.\csromanspacer{} \textbf{อารมฺมณาธิปติ\csromandash{}}อปฺปีติเกน จิตฺเตน ทานํ ฯเปฯ สีลํ ฯเปฯ อุโปสถกมฺมํ ฯเปฯ. (สํขิตฺตํ.) นิพฺพานํ สปฺปีติกสฺส โคตฺรภุสฺส, โวทานสฺส, มคฺคสฺส, ผลสฺส อธิปติปจฺจเยน ปจฺจโย.\csromanspacer{} จกฺขุํ ฯเปฯ.\csromanspacer{} วตฺถุํ อปฺปีติเก ขนฺเธ จ ปีติญฺจ สปฺปีติเกน จิตฺเตน ครุํ กตฺวา อสฺสาเทติ อภินนฺทติ, ตํ ครุํ กตฺวา สปฺปีติโก ราโค อุปฺปชฺชติ.\csromanspacer{} ทิฏฺฐิ อุปฺปชฺชติ.\csromanspacer{} อปฺปีติเก ขนฺเธ จ ปีติญฺจ ครุํ กตฺวา สปฺปีติกา ขนฺธา อุปฺปชฺชนฺติ. (2)}

\prose{อปฺปีติโก ธมฺโม สปฺปีติกสฺส จ อปฺปีติกสฺส จ ธมฺมสฺส อธิปติปจฺจเยน ปจฺจโย.\csromanspacer{} \textbf{อารมฺมณาธิปติ\csromandash{}}ทานํ ฯเปฯ (สํขิตฺตํ.) นิพฺพานํ สปฺปีติกสฺส โคตฺรภุสฺส, โวทานสฺส, มคฺคสฺส, ผลสฺส, ปีติยา จ \csromanfolio{106}อธิปติปจฺจเยน ปจฺจโย.\csromanspacer{} จกฺขุํ ฯเปฯ.\csromanspacer{} วตฺถุํ อปฺปีติเก ขนฺเธ จ ปีติญฺจ สปฺปีติเกน จิตฺเตน ครุํ กตฺวา อสฺสาเทติ อภินนฺทติ, ตํ ครุํ กตฺวา ราโค อุปฺปชฺชติ.\csromanspacer{} ทิฏฺฐิ อุปฺปชฺชติ, อปฺปีติเก ขนฺเธ จ ปีติญฺจ ครุํ กตฺวา สปฺปีติกา ขนฺธา จ ปีติ จ อุปฺปชฺชนฺติ. (3)}

\proseitem{155}{สปฺปีติโก จ อปฺปีติโก จ ธมฺมา สปฺปีติกสฺส ธมฺมสฺส อธิปติปจฺจเยน ปจฺจโย.\csromanspacer{} \textbf{อารมฺมณาธิปติ\csromandash{}}สปฺปีติเก ขนฺเธ จ ปีติญฺจ ครุํ กตฺวา สปฺปีติกา ขนฺธา อุปฺปชฺชนฺติ. (มูลํ กาตพฺพํ.) สปฺปีติเก ขนฺเธ จ ปีติญฺจ ครุํ กตฺวา อปฺปีติกา ขนฺธา จ ปีติ จ อุปฺปชฺชนฺติ. (มูลํ กาตพฺพํ.) สปฺปีติเก ขนฺเธ จ ปีติญฺจ ครุํ กตฺวา สปฺปีติกา ขนฺธา จ ปีติ จ อุปฺปชฺชนฺติ. (3)}

\csromanheader{2}{\textbf{อนนฺตราทิ}}
\proseitem{156}{สปฺปีติโก ธมฺโม สปฺปีติกสฺส ธมฺมสฺส อนนฺตรปจฺจเยน ปจฺจโย.\csromanspacer{} ปุริมา ปุริมา สปฺปีติกา ขนฺธา ปจฺฉิมานํ ปจฺฉิมานํ สปฺปีติกานํ ขนฺธานํ อนนฺตรปจฺจเยน ปจฺจโย. (มูลํ กาตพฺพํ.) ปุริมา ปุริมา สปฺปีติกา ขนฺธา ปจฺฉิมาย ปจฺฉิมาย ปีติยา อนนฺตรปจฺจเยน ปจฺจโย.\csromanspacer{} สปฺปีติกํ จุติจิตฺตํ อปฺปีติกสฺส อุปปตฺติจิตฺตสฺส.\csromanspacer{} สปฺปีติกํ ภวงฺคํ อาวชฺชนาย.\csromanspacer{} สปฺปีติกา ขนฺธา อปฺปีติกสฺส วุฏฺฐานสฺส.\csromanspacer{} ปีติสหคตา วิปากมโนวิญฺญาณธาตุ กิริยมโนวิญฺญาณธาตุยา.\csromanspacer{} สปฺปีติกํ ภวงฺคํ อปฺปีติกสฺส ภวงฺคสฺส.\csromanspacer{} สปฺปีติกํ กุสลากุสลํ อปฺปีติกสฺส วุฏฺฐานสฺส.\csromanspacer{} กิริยํ วุฏฺฐานสฺส.\csromanspacer{} ผลํ วุฏฺฐานสฺส อนนฺตรปจฺจเยน ปจฺจโย. (มูลํ กาตพฺพํ.) ปุริมา ปุริมา สปฺปีติกา ขนฺธา ปจฺฉิมานํ ปจฺฉิมานํ สปฺปีติกานํ ขนฺธานํ ปีติยา จ อนนฺตรปจฺจเยน ปจฺจโย. (3)}

\proseitem{157}{อปฺปีติโก ธมฺโม อปฺปีติกสฺส ธมฺมสฺส อนนฺตรปจฺจเยน ปจฺจโย.\csromanspacer{} ปุริมา ปุริมา ปีติ ปจฺฉิมาย ปจฺฉิมาย ปีติยา อนนฺตรปจฺจเยน ปจฺจโย.\csromanspacer{} ปุริมา ปุริมา อปฺปีติกา ขนฺธา ปจฺฉิมานํ ปจฺฉิมานํ อปฺปีติกานํ ขนฺธานํ อนนฺตรปจฺจเยน ปจฺจโย.\csromanspacer{} อนุโลมํ โคตฺรภุสฺส ฯเปฯ.\csromanspacer{} อปฺปีติกาย ผลสมาปตฺติยา ปีติยา จ อนนฺตรปจฺจเยน ปจฺจโย. (มูลํ กาตพฺพํ.) ปุริมา ปุริมา ปีติ ปจฺฉิมานํ ปจฺฉิมานํ สปฺปีติกานํ ขนฺธานํ อนนฺตรปจฺจเยน ปจฺจโย.\csromanspacer{} อปฺปีติกํ จุติจิตฺตํ สปฺปีติกสฺส อุปปตฺติจิตฺตสฺส.\csromanspacer{} อาวชฺชนา สปฺปีติกานํ ขนฺธานํ.\csromanspacer{} อปฺปีติกา ขนฺธา สปฺปีติกสฺส วุฏฺฐานสฺส.\csromanspacer{} วิปากมโนธาตุ สปฺปีติกาย วิปากมโนวิญฺญาณธาตุยา.\csromanspacer{} อปฺปีติกํ}

\vspace{\tipitakaparskip}
\csromancenter{สปฺปีติกทุก ภวงฺคํ สปฺปีติกสฺส ภวงฺคสฺส.\csromanspacer{} อปฺปีติกํ กุสลากุสลํ สปฺปีติกสฺส วุฏฺฐานสฺส.\csromanspacer{} กิริยํ วุฏฺฐานสฺส.\csromanspacer{} ผลํ วุฏฺฐานสฺส อนนฺตรปจฺจเยน ปจฺจโย.\csromanspacer{} นิโรธา วุฏฺฐหนฺตสฺส เนวสญฺญานาสญฺญายตนํ สปฺปีติกาย ผลสมาปตฺติยา อนนฺตรปจฺจเยน ปจฺจโย. (2)}
\prose{\csromanfolio{107}อปฺปีติโก ธมฺโม สปฺปีติกสฺส จ อปฺปีติกสฺส จ ธมฺมสฺส อนนฺตรปจฺจเยน ปจฺจโย.\csromanspacer{} ปุริมา ปุริมา ปีติ ปจฺฉิมานํ ปจฺฉิมานํ สปฺปีติกานํ ขนฺธานํ ปีติยา จ อนนฺตรปจฺจเยน ปจฺจโย. (3)}

\proseitem{158}{สปฺปีติโก จ อปฺปีติโก จ ธมฺมา สปฺปีติกสฺส ธมฺมสฺส อนนฺตรปจฺจเยน ปจฺจโย.\csromanspacer{} ปุริมา ปุริมา สปฺปีติกา ขนฺธา จ ปีติ จ ปจฺฉิมานํ ปจฺฉิมานํ สปฺปีติกานํ ขนฺธานํ อนนฺตรปจฺจเยน ปจฺจโย. (1)}

\prose{สปฺปีติโก จ อปฺปีติโก จ ธมฺมา อปฺปีติกสฺส ธมฺมสฺส อนนฺตรปจฺจเยน ปจฺจโย.\csromanspacer{} ปุริมา ปุริมา สปฺปีติกา ขนฺธา จ ปีติ จ ปจฺฉิมาย ปจฺฉิมาย ปีติยา อนนฺตรปจฺจเยน ปจฺจโย.\csromanspacer{} สปฺปีติกํ จุติจิตฺตญฺจ ปีติ จ อปฺปีติกสฺส อุปปตฺติจิตฺตสฺส.\csromanspacer{} สปฺปีติกํ ภวงฺคญฺจ ปีติ จ อาวชฺชนาย.\csromanspacer{} สปฺปีติกา ขนฺธา จ ปีติ จ อปฺปีติกสฺส วุฏฺฐานสฺส.\csromanspacer{} สปฺปีติกา วิปากมโนวิญฺญาณธาตุ จ ปีติ จ กิริยมโนวิญฺญาณธาตุยา.\csromanspacer{} สปฺปีติกํ ภวงฺคญฺจ ปีติ จ อปฺปีติกสฺส ภวงฺคสฺส.\csromanspacer{} สปฺปีติกํ กุสลากุสลญฺจ ปีติ จ อปฺปีติกสฺส วุฏฺฐานสฺส.\csromanspacer{} กิริยญฺจ ปีติ จ วุฏฺฐานสฺส.\csromanspacer{} ผลญฺจ ปีติ จ วุฏฺฐานสฺส อนนฺตรปจฺจเยน ปจฺจโย. (2)}

\prose{สปฺปีติโก จ อปฺปีติโก จ ธมฺมา สปฺปีติกสฺส จ อปฺปีติกสฺส จ ธมฺมสฺส อนนฺตรปจฺจเยน ปจฺจโย.\csromanspacer{} ปุริมา ปุริมา สปฺปีติกา ขนฺธา จ ปีติ จ ปจฺฉิมานํ ปจฺฉิมานํ สปฺปีติกานํ ขนฺธานํ ปีติยา จ อนนฺตรปจฺจเยน ปจฺจโย. (3)}

\prose{สมนนฺตรปจฺจเยน ปจฺจโย.\csromanspacer{} นว.\csromanspacer{} สหชาตปจฺจเยน ปจฺจโย.\csromanspacer{} นว.\csromanspacer{} อญฺญมญฺญปจฺจเยน ปจฺจโย.\csromanspacer{} นว.\csromanspacer{} นิสฺสยปจฺจเยน ปจฺจโย.\csromanspacer{} นว.}

\csromanheader{2}{\textbf{อุปนิสฺสย}}
\proseitem{159}{สปฺปีติโก ธมฺโม สปฺปีติกสฺส ธมฺมสฺส อุปนิสฺสยปจฺจเยน ปจฺจโย.\csromanspacer{} อารมฺมณูปนิสฺสโย, อนนฺตรูปนิสฺสโย, ปกตูปนิสฺสโย ฯเปฯ.\csromanspacer{} ปกตูปนิสฺสโย\csromandash{}สปฺปีติกา ขนฺธา สปฺปีติกานํ ขนฺธานํ \csromanfolio{108}อุปนิสฺสยปจฺจเยน ปจฺจโย. (มูลํ กาตพฺพํ.) สปฺปีติกา ขนฺธา อปฺปีติกานํ ขนฺธานํ ปีติยา จ อุปนิสฺสยปจฺจเยน ปจฺจโย. (มูลํ กาตพฺพํ.) สปฺปีติกา ขนฺธา สปฺปีติกานํ ขนฺธานํ ปีติยา จ อุปนิสฺสยปจฺจเยน ปจฺจโย. (3)}

\proseitem{160}{อปฺปีติโก ธมฺโม อปฺปีติกสฺส ธมฺมสฺส อุปนิสฺสยปจฺจเยน ปจฺจโย.\csromanspacer{} อารมฺมณูปนิสฺสโย, อนนฺตรูปนิสฺสโย, ปกตูปนิสฺสโย ฯเปฯ.\csromanspacer{} ปกตูปนิสฺสโย\csromandash{}อปฺปีติกํ สทฺธํ อุปนิสฺสาย อปฺปีติเกน จิตฺเตน ทานํ เทติ, สิลํ สมาทิยติ, อุโปสถกมฺมํ กโรติ.\csromanspacer{} อปฺปีติกํ ฌานํ ฯเปฯ วิปสฺสนํ.\csromanspacer{} มคฺคํ.\csromanspacer{} อภิญฺญํ.\csromanspacer{} สมาปตฺติํ อุปฺปาเทติ.\csromanspacer{} มานํ ชปฺเปติ.\csromanspacer{} ทิฏฺฐิํ คณฺหาติ.\csromanspacer{} อปฺปีติกํ สีลํ ฯเปฯ ปญฺญํ.\csromanspacer{} ราคํ.\csromanspacer{} โทสํ.\csromanspacer{} โมหํ.\csromanspacer{} มานํ.\csromanspacer{} ทิฏฺฐิํ.\csromanspacer{} ปตฺถนํ.\csromanspacer{} กายิกํ สุขํ.\csromanspacer{} กายิกํ ทุกฺขํ.\csromanspacer{} อุตุํ.\csromanspacer{} โภชนํ.\csromanspacer{} เสนาสนํ ปีติํ อุปนิสฺสาย อปฺปีติเกน จิตฺเตน ทานํ เทติ ฯเปฯ สมาปตฺติํ อุปฺปาเทติ.\csromanspacer{} ปาณํ หนติ ฯเปฯ สํฆํ ภินฺทติ.\csromanspacer{} อปฺปีติกา สทฺธา ฯเปฯ.\csromanspacer{} เสนาสนํ ปีติ จ อปฺปีติกาย สทฺธาย ฯเปฯ ปตฺถนาย.\csromanspacer{} กายิกสฺส สุขสฺส.\csromanspacer{} กายิกสฺส ทุกฺขสฺส.\csromanspacer{} มคฺคสฺส ผลสมาปตฺติยา ปีติยา จ อุปนิสฺสยปจฺจเยน ปจฺจโย. (1)}

\prose{อปฺปีติโก ธมฺโม สปฺปีติกสฺส ธมฺมสฺส อุปนิสฺสยปจฺจเยน ปจฺจโย. (ตีณิ อุปนิสฺสยา.) อปฺปีติกํ สทฺธํ อุปนิสฺสาย สปฺปีติเกน จิตฺเตน ทานํ เทติ ฯเปฯ.\csromanspacer{} อปฺปีติกา ฌานา ฯเปฯ.\csromanspacer{} มานํ ชปฺเปติ.\csromanspacer{} ทิฏฺฐิํ คณฺหาติ.\csromanspacer{} อปฺปีติกํ สีลํ ฯเปฯ.\csromanspacer{} เสนาสนํ ปีติํ อุปนิสฺสาย สปฺปีติเกน จิตฺเตน ทานํ เทติ ฯเปฯ สมาปตฺติํ อุปฺปาเทติ.\csromanspacer{} สปฺปีติเกน จิตฺเตน อทินฺนํ อาทิยติ.\csromanspacer{} มุสา ฯเปฯ ปิสุณํ ฯเปฯ สมฺผํ ฯเปฯ สนฺธิํ ฯเปฯ นิลฺโลปํ ฯเปฯ เอกาคาริกํ ฯเปฯ ปริปนฺเถ ฯเปฯ ปรทารํ ฯเปฯ คามฆาตํ ฯเปฯ นิคมฆาตํ กโรติ.\csromanspacer{} อปฺปีติกา สทฺธา ฯเปฯ.\csromanspacer{} เสนาสนํ ปีติ จ สปฺปีติกาย สทฺธาย ฯเปฯ ปญฺญาย.\csromanspacer{} ราคสฺส.\csromanspacer{} โมหสฺส.\csromanspacer{} มานสฺส.\csromanspacer{} ทิฏฺฐิยา.\csromanspacer{} ปตฺถนาย.\csromanspacer{} มคฺคสฺส, ผลสมาปตฺติยา อุปนิสฺสยปจฺจเยน ปจฺจโย. (2)}

\prose{อปฺปีติโก ธมฺโม สปฺปีติกสฺส จ อปฺปีติกสฺส จ ธมฺมสฺส อุปนิสฺสยปจฺจเยน ปจฺจโย. (ตีณิ อุปนิสฺสยา.) อปฺปีติกํ สทฺธํ อุปนิสฺสาย สปฺปีติเกน จิตฺเตน ทานํ เทติ ฯเปฯ สมาปตฺติํ อุปฺปาเทติ.\csromanspacer{} มานํ ชปฺเปติ.\csromanspacer{} ทิฏฺฐิํ คณฺหาติ.\csromanspacer{} อปฺปีติกํ สีลํ ฯเปฯ.\csromanspacer{} เสนาสนํ ปีติํ อุปนิสฺสาย}

\vspace{\tipitakaparskip}
\csromancenter{สปฺปีติกทุก สปฺปีติเกน จิตฺเตน ทานํ เทติ ฯเปฯ สมาปตฺติํ อุปฺปาเทติ.\csromanspacer{} สปฺปีติเกน จิตฺเตน อทินฺนํ อาทิยติ ฯเปฯ. (ทุติยวารสทิสํ.) นิคมฆาตํ กโรติ.\csromanspacer{} อปฺปีติกา สทฺธา ฯเปฯ.\csromanspacer{} เสนาสนํ ปีติ จ สปฺปีติกาย สทฺธาย ฯเปฯ ปญฺญาย.\csromanspacer{} ราคสฺส.\csromanspacer{} โมหสฺส.\csromanspacer{} มานสฺส.\csromanspacer{} ทิฏฺฐิยา.\csromanspacer{} ปตฺถนาย.\csromanspacer{} มคฺคสฺส, ผลสมาปตฺติยา ปีติยา จ อุปนิสฺสยปจฺจเยน ปจฺจโย. (3)}
\prose{\csromanfolio{109}สปฺปีติโก จ อปฺปีติโก จ ธมฺมา สปฺปีติกสฺส ธมฺมสฺส อุปนิสฺสยปจฺจเยน ปจฺจโย. (ตีณิปิ อุปนิสฺสยา.) สปฺปีติกา ขนฺธา จ ปีติ จ สปฺปีติกานํ ขนฺธานํ อุปนิสฺสยปจฺจเยน ปจฺจโย. (มูลํ กาตพฺพํ.) สปฺปีติกา ขนฺธา จ ปีติ จ อปฺปีติกานํ ขนฺธานํ ปีติยา จ อุปนิสฺสยปจฺจเยน ปจฺจโย. (มูลํ กาตพฺพํ.) สปฺปีติกา ขนฺธา จ ปีติ จ สปฺปีติกานํ ขนฺธานํ ปีติยา จ อุปนิสฺสยปจฺจเยน ปจฺจโย. (3)}

\csromanheader{2}{\textbf{ปุเรชาต}}
\proseitem{161}{อปฺปีติโก ธมฺโม อปฺปีติกสฺส ธมฺมสฺส ปุเรชาตปจฺจเยน ปจฺจโย.\csromanspacer{} อารมฺมณปุเรชาตํ, วตฺถุปุเรชาตํ.\csromanspacer{}\csromanspacer{}\textbf{อารมฺมณปุเรชาตํ\csromandash{}} จกฺขุํ ฯเปฯ.\csromanspacer{} วตฺถุํ อปฺปีติเกน จิตฺเตน อนิจฺจโต ฯเปฯ วิปสฺสติ, อสฺสาเทติ อภินนฺทติ, ตํ อารพฺภ อปฺปีติโก ราโค ฯเปฯ โทมนสฺสํ อุปฺปชฺชติ.\csromanspacer{} ปีติ อุปฺปชฺชติ.\csromanspacer{} ทิพฺเพน จกฺขุนา รูปํ ปสฺสาติ ฯเปฯ.\csromanspacer{} โผฏฺฐพฺพายตนํ กายวิญฺญาณสฺส ฯเปฯ.\csromanspacer{} \textbf{วตฺถุปุเรชาตํ\csromandash{}}จกฺขายตนํ จกฺขุวิญฺญาณสฺส ฯเปฯ กายายตนํ กายวิญฺญาณสฺส ฯเปฯ.\csromanspacer{} วตฺถุ อปฺปีติกานํ ขนฺธานํ ปีติยา จ ปุเรชาตปจฺจเยน ปจฺจโย. (1)}

\prose{อปฺปีติโก ธมฺโม สปฺปีติกสฺส ธมฺมสฺส ปุเรชาตปจฺจเยน ปจฺจโย.\csromanspacer{} อารมฺมณปุเรชาตํ, วตฺถุปุเรชาตํ.\csromanspacer{}\csromanspacer{}\textbf{อารมฺมณปุเรชาตํ\csromandash{}}จกฺขุํ ฯเปฯ.\csromanspacer{} วตฺถุํ สปฺปีติเกน จิตฺเตน อนิจฺจโต ฯเปฯ วิปสฺสติ, อสฺสาเทติ อภินนฺทติ, ตํ อารพฺภ สปฺปีติโก ราโค อุปฺปชฺชติ.\csromanspacer{} ทิฏฺฐิ อุปฺปชฺชติ.\csromanspacer{} \textbf{วตฺถุปุเรชาตํ\csromandash{}}วตฺถุ สปฺปีติกานํ ขนฺธานํ ปุเรชาตปจฺจเยน ปจฺจโย. (2)}

\prose{อปฺปีติโก ธมฺโม สปฺปีติกสฺส จ อปฺปีติกสฺส จ ธมฺมสฺส ปุเรชาตปจฺจเยน ปจฺจโย.\csromanspacer{} อารมฺมณปุเรชาตํ, วตฺถุปุเรชาตํ.\csromanspacer{}\csromanspacer{}\textbf{อารมฺมณปุเรชาตํ\csromandash{}}จกฺขุํ ฯเปฯ.\csromanspacer{} วตฺถุํ สปฺปีติเกน จิตฺเตน อนิจฺจโต ฯเปฯ วิปสฺสติ, อสฺสาเทติ อภินนฺทติ, ตํ อารพฺภ ปีติ จ สมฺปยุตฺตกา ขนฺธา จ อุปฺปชฺชนฺติ. \csromanfolio{110}\textbf{วตฺถุปุเรชาตํ\csromandash{}}วตฺถุ สปฺปีติกานํ ขนฺธานํ ปีติยา จ ปุเรชาตปจฺจเยน ปจฺจโย. (3)}

\csromanheader{2}{\textbf{ปจฺฉาชาตาทิ}}
\proseitem{162}{สปฺปีติโก ธมฺโม อปฺปีติกสฺส ธมฺมสฺส ปจฺฉาชาตปจฺจเยน ปจฺจโย.\csromanspacer{} ตีณิ.\csromanspacer{} อาเสวนปจฺจเยน ปจฺจโย.\csromanspacer{} นว.\csromanspacer{} กมฺมปจฺจเยน ปจฺจโย.\csromanspacer{} ฉ. (สหชาตาปิ นานากฺขณิกาปิ กาตพฺพา.\csromanspacer{} เทฺว นานากฺขณิกา.) วิปากปจฺจเยน ปจฺจโย.\csromanspacer{} นว.\csromanspacer{} อาหารปจฺจเยน ปจฺจโย.\csromanspacer{} จตฺตาริ.\csromanspacer{} อินฺทฺริยปจฺจเยน ปจฺจโย.\csromanspacer{} จตฺตาริ.\csromanspacer{} ฌานปจฺจเยน ปจฺจโย.\csromanspacer{} นว.\csromanspacer{} มคฺคปจฺจเยน ปจฺจโย.\csromanspacer{} จตฺตาริ.\csromanspacer{} สมฺปยุตฺตปจฺจเยน ปจฺจโย.\csromanspacer{} ฉ.\csromanspacer{} วิปฺปยุตฺตปจฺจเยน ปจฺจโย.\csromanspacer{} ปญฺจ.\csromanspacer{} อตฺถิปจฺจเยน ปจฺจโย.\csromanspacer{} นว. (สํขิตฺตํ.\csromanspacer{} สวิตกฺกทุกสทิสํ กาตพฺพํ.) นตฺถิปจฺจเยน ปจฺจโย.\csromanspacer{} วิคตปจฺจเยน ปจฺจโย.\csromanspacer{} อวิคตปจฺจเยน ปจฺจโย.\csromanspacer{} นว.}

\csromanheader{2}{1. ปจฺจยานุโลม 2. สงฺขฺยาวาร}
\proseitem{163}{เหตุยา จตฺตาริ, อารมฺมเณ นว, อธิปติยา นว, อนนฺตเร นว, สมนนฺตเร นว, สหชาเต นว, อญฺญมญฺเญ นว, นิสฺสเย นว, อุปนิสฺสเย นว, ปุเรชาเต ตีณิ, ปจฺฉาชาเต ตีณิ, อาเสวเน นว, กมฺเม ฉ, วิปาเก นว, อาหาเร จตฺตาริ, อินฺทฺริเย จตฺตาริ, ฌาเน นว, มคฺเค จตฺตาริ, สมฺปยุตฺเต ฉ, วิปฺปยุตฺเต ปญฺจ, อตฺถิยา นว, นตฺถิยา นว, วิคเต นว, อวิคเต นว.}

\vspace{\tipitakaparskip}
\csromancenter{อนุโลมํ.}
\csromanheader{2}{2. ปจฺจนียุทฺธาร}
\proseitem{164}{สปฺปีติโก ธมฺโม สปฺปีติกสฺส ธมฺมสฺส อารมฺมณปจฺจเยน ปจฺจโย.\csromanspacer{} สหชาตปจฺจเยน ปจฺจโย.\csromanspacer{} อุปนิสฺสยปจฺจเยน ปจฺจโย.\csromanspacer{} กมฺมปจฺจเยน ปจฺจโย.\csromanspacer{}\csromanspacer{}สปฺปีติโก ธมฺโม อปฺปีติกสฺส ธมฺมสฺส อารมฺมณปจฺจเยน ปจฺจโย.\csromanspacer{} สหชาตปจฺจเยน ปจฺจโย.\csromanspacer{} อุปนิสฺสยปจฺจเยน ปจฺจโย.\csromanspacer{} ปจฺฉาชาตปจฺจเยน ปจฺจโย.\csromanspacer{} กมฺมปจฺจเยน ปจฺจโย.\csromanspacer{}\csromanspacer{}}

\vspace{\tipitakaparskip}
\csromancenter{ปีติสหคตทุก สปฺปีติโก ธมฺโม สปฺปีติกสฺส จ อปฺปีติกสฺส จ ธมฺมสฺส อารมฺมณปจฺจเยน ปจฺจโย.\csromanspacer{} สหชาตปจฺจเยน ปจฺจโย.\csromanspacer{} อุปนิสฺสยปจฺจเยน ปจฺจโย.\csromanspacer{} กมฺมปจฺจเยน ปจฺจโย. (3)}
\proseitem{165}{\csromanfolio{111}อปฺปีติโก ธมฺโม อปฺปีติกสฺส ธมฺมสฺส อารมฺมณปจฺจเยน ปจฺจโย.\csromanspacer{} สหชาตปจฺจเยน ปจฺจโย.\csromanspacer{} อุปนิสฺสยปจฺจเยน ปจฺจโย.\csromanspacer{} ปุเรชาตปจฺจเยน ปจฺจโย.\csromanspacer{} ปจฺฉาชาตปจฺจเยน ปจฺจโย.\csromanspacer{} กมฺมปจฺจเยน ปจฺจโย.\csromanspacer{} อาหารปจฺจเยน ปจฺจโย.\csromanspacer{} อินฺทฺริยปจฺจเยน ปจฺจโย.\csromanspacer{}\csromanspacer{}อปฺปีติโก ธมฺโม สปฺปีติกสฺส ธมฺมสฺส อารมฺมณปจฺจเยน ปจฺจโย.\csromanspacer{} สหชาตปจฺจเยน ปจฺจโย.\csromanspacer{} อุปนิสฺสยปจฺจเยน ปจฺจโย.\csromanspacer{} ปุเรชาตปจฺจเยน ปจฺจโย.\csromanspacer{}\csromanspacer{}อปฺปีติโก ธมฺโม สปฺปีติกสฺส จ อปฺปีติกสฺส จ ธมฺมสฺส อารมฺมณปจฺจเยน ปจฺจโย.\csromanspacer{} สหชาตปจฺจเยน ปจฺจโย.\csromanspacer{} อุปนิสฺสยปจฺจเยน ปจฺจโย.\csromanspacer{} ปุเรชาตปจฺจเยน ปจฺจโย. (3)}

\proseitem{166}{สปฺปีติโก จ อปฺปีติโก จ ธมฺมา สปฺปีติกสฺส ธมฺมสฺส อารมฺมณปจฺจเยน ปจฺจโย.\csromanspacer{} สหชาตปจฺจเยน ปจฺจโย.\csromanspacer{} อุปนิสฺสยปจฺจเยน ปจฺจโย.\csromanspacer{}\csromanspacer{}สปฺปีติโก จ อปฺปีติโก จ ธมฺมา อปฺปีติกสฺส ธมฺมสฺส อารมฺมณปจฺจเยน ปจฺจโย.\csromanspacer{} สหชาตปจฺจเยน ปจฺจโย.\csromanspacer{} อุปนิสฺสยปจฺจเยน ปจฺจโย.\csromanspacer{} ปจฺฉาชาตปจฺจเยน ปจฺจโย.\csromanspacer{}\csromanspacer{}สปฺปีติโก จ อปฺปีติโก จ ธมฺมา สปฺปีติกสฺส จ อปฺปีติกสฺส จ ธมฺมสฺส อารมฺมณปจฺจเยน ปจฺจโย.\csromanspacer{} สหชาตปจฺจเยน ปจฺจโย.\csromanspacer{} อุปนิสฺสยปจฺจเยน ปจฺจโย. (3)}

\prose{(ปจฺจนียวิภงฺคคณนาปิ สวิตกฺกทุกสทิสา.\csromanspacer{} ยทิปิ น สเมติ อิมํ อนุโลมํ ปจฺจเวกฺขิตฺวา คเณตพฺพํ.\csromanspacer{} อิตเร เทฺว คณนา คเณตพฺพา.)}

\nitthitam{สปฺปีติกทุกํ นิฏฺฐิตํ.}
\csromanmatikaanchor{mtk.19}
\csromantocmarkref{section}{90. ปีติสหคตทุก}{mtk.19}
\bookmark[dest={mtk.19},level=1]{90. ปีติสหคตทุก}
\csromanheader{1}{90. ปีติสหคตทุก 1. ปฏิจฺจวาร}
\proseitem{167}{ปีติสหคตํ ธมฺมํ ปฏิจฺจ ปีติสหคโต ธมฺโม อุปฺปชฺชติ เหตุปจฺจยา.\csromanspacer{} ปีติสหคตํ เอกํ ขนฺธํ ปฏิจฺจ ตโย ขนฺธา ฯเปฯ เทฺว ขนฺเธ ฯเปฯ.}

\prose{\csromanfolio{112}(เอวํ ปีติสหคตทุกํ วิตฺถาเรตพฺพํ.\csromanspacer{} สปฺปีติกทุกสทิสํ, นินฺนานากรณํ, อามสนํ นินฺนานํ.)}

\nitthitam{ปีติสหคตทุกํ นิฏฺฐิตํ.}
\csromanmatikaanchor{mtk.20}
\csromantocmarkref{section}{91. สุขสหคตทุก}{mtk.20}
\bookmark[dest={mtk.20},level=1]{91. สุขสหคตทุก}
\csromanheader{1}{91. สุขสหคตทุก 1. ปฏิจฺจวาร}
\csromanheader{2}{1. ปจฺจยานุโลมาทิ}
\proseitem{168}{สุขสหคตํ ธมฺมํ ปฏิจฺจ สุขสหคโต ธมฺโม อุปฺปชฺชติ เหตุปจฺจยา.\csromanspacer{} สุขสหคตํ เอกํ ขนฺธํ ปฏิจฺจ เทฺว ขนฺธา, เทฺว ขนฺเธ ปฏิจฺจ เอโก ขนฺโธ.\csromanspacer{} ปฏิสนฺธิกฺขเณ ฯเปฯ.\csromanspacer{} สุขสหคตํ ธมฺมํ ปฏิจฺจ นสุขสหคโต ธมฺโม อุปฺปชฺชติ เหตุปจฺจยา.\csromanspacer{} สุขสหคเต ขนฺเธ ปฏิจฺจ สุขํ จิตฺตสมุฏฺฐานญฺจ รูปํ. (เอวํ สุขสหคตทุกํ วิตฺถาเรตพฺพํ.\csromanspacer{} ยถา สปฺปีติกทุกสฺส อนุโลมปฏิจฺจวาโร.)}

\vspace{\tipitakaparskip}
\csromancenter{\textbf{เหตุ}ยา นว, อารมฺมเณ นว ฯเปฯ ปุเรชาเต ฉ, อาเสวเน ฉ, กมฺเม นว ฯเปฯ อวิคเต นว.\csromanspacer{} อนุโลมํ.}
\proseitem{169}{สุขสหคตํ ธมฺมํ ปฏิจฺจ สุขสหคโต ธมฺโม อุปฺปชฺชติ นเหตุปจฺจยา.\csromanspacer{} อเหตุกํ สุขสหคตํ เอกํ ขนฺธํ ปฏิจฺจ เทฺว ขนฺธา, เทฺว ขนฺเธ ฯเปฯ. (มูลํ กาตพฺพํ.) อเหตุเก สุขสหคเต ขนฺเธ ปฏิจฺจ สุขญฺจ จิตฺตสมุฏฺฐานญฺจ รูปํ.\csromanspacer{}\csromanspacer{}สุขสหคตํ ธมฺมํ ปฏิจฺจ สุขสหคโต จ นสุขสหคโต จ ธมฺมา อุปฺปชฺชนฺติ นเหตุปจฺจยา.\csromanspacer{} อเหตุกํ สุขสหคตํ เอกํ ขนฺธํ ปฏิจฺจ เทฺว ขนฺธา สุขญฺจ จิตฺตสมุฏฺฐานญฺจ รูปํ.\csromanspacer{} เทฺว ขนฺเธ ฯเปฯ. (3)}

\prose{นสุขสหคตํ ธมฺมํ ปฏิจฺจ นสุขสหคโต ธมฺโม อุปฺปชฺชติ นเหตุปจฺจยา.\csromanspacer{} อเหตุกํ นสุขสหคตํ เอกํ ขนฺธํ ปฏิจฺจ ตโย ขนฺธา จิตฺตสมุฏฺฐานญฺจ รูปํ ฯเปฯ เทฺว ขนฺเธ ฯเปฯ.\csromanspacer{} อเหตุกํ สุขํ ปฏิจฺจ จิตฺตสมุฏฺฐานํ รูปํ.\csromanspacer{} อเหตุกปฏิสนฺธิกฺขเณ นสุขสหคตํ เอกํ ขนฺธํ ปฏิจฺจ ตโย ขนฺธา กฏตฺตา จ}

\vspace{\tipitakaparskip}
\csromancenter{สุขสหคตทุก รูปํ ฯเปฯ เทฺว ขนฺเธ ฯเปฯ.\csromanspacer{} ขนฺเธ ปฏิจฺจ วตฺถุ, วตฺถุํ ปฏิจฺจ ขนฺธา.\csromanspacer{} เอกํ มหาภูตํ ฯเปฯ. (ยาว อสญฺญสตฺตา.) วิจิกิจฺฉาสหคเต อุทฺธจฺจสหคเต ขนฺเธ ปฏิจฺจ วิจิกิจฺฉาสหคโต อุทฺธจฺจสหคโต โมโห. (ยถา สปฺปีติเก นเหตุปจฺจยสทิสํ, นินฺนานํ, สพฺพตฺถเมว นว ปญฺหา.)}
\prose{\csromanfolio{113}นเหตุยา นว, น-อารมฺมเณ ตีณิ ฯเปฯ น-อุปนิสฺสเย ตีณิ, นปุเรชาเต นว, นปจฺฉาชาเต นว, น-อาเสวเน นว, นกมฺเม จตฺตาริ, นวิปาเก นว, น-อาหาเร เอกํ, น-อินฺทฺริเย เอกํ, นฌาเน ฉ, นมคฺเค นว, นสมฺปยุตฺเต ตีณิ, นวิปฺปยุตฺเต ฉ, โนนตฺถิยา ตีณิ, โนวิคเต ตีณิ.\csromanspacer{} ปจฺจนียํ. (เอวํ อิตเร เทฺว คณนาปิ สหชาตวาโรปิ ปฏิจฺจวารสทิสา.\csromanspacer{} ปจฺจยวาเร ปวตฺติปิ ปฏิสนฺธิปิ วิตฺถาเรตพฺพา, ยถา สปฺปีติกทุกปจฺจวารปจฺจนีเยปิ ปวตฺเต วตฺถุ จ วิตฺถาเรตพฺพํ, ยถา สปฺปีติกทุเก เอโกเยว โมโห, เอวํ อิตเร เทฺว คณนาปิ นิสฺสยวาโรปิ สํสฏฺฐวาโรปิ สมฺปยุตฺตวาโรปิ ยถา สปฺปีติกทุกํ, เอวํ กาตพฺพํ.)}

\csromanheader{2}{91. สุขสหคตทุก 7. ปญฺหาวาร}
\proseitem{170}{สุขสหคโต ธมฺโม สุขสหคตสฺส ธมฺมสฺส เหตุปจฺจเยน ปจฺจโย.\csromanspacer{} จตฺตาริ. (อารมฺมเณปิ อธิปติยาปิ สปฺปีติกทุกสทิสา. “สุขนฺ”ติ นานากรณํ.)}

\csromanheader{2}{\textbf{อนนฺตราทิ}}
\proseitem{171}{สุขสหคโต ธมฺโม สุขสหคตสฺส ธมฺมสฺส อนนฺตรปจฺจเยน ปจฺจโย.\csromanspacer{} ปุริมา ปุริมา สุขสหคตา ขนฺธา ปจฺฉิมานํ ปจฺฉิมานํ สุขสหคตานํ ขนฺธานํ อนนฺตรปจฺจเยน ปจฺจโย. (มูลํ กาตพฺพํ.) ปุริมา ปุริมา \csromanfolio{114}สุขสหคตา ขนฺธา ปจฺฉิมสฺส ปจฺฉิมสฺส สุขสฺส อนนฺตรปจฺจเยน ปจฺจโย.\csromanspacer{} สุขสหคตํ จุติจิตฺตํ นสุขสหคตสฺส อุปปตฺติจิตฺตสฺส.\csromanspacer{} สุขสหคตํ ภวงฺคํ อาวชฺชนาย.\csromanspacer{} สุขสหคตํ กายวิญฺญาณํ วิปากมโนธาตุยา.\csromanspacer{} สุขสหคตา วิปากมโนวิญฺญาณธาตุ กิริยมโนวิญฺญาณธาตุยา.\csromanspacer{} สุขสหคตํ ภวงฺคํ นสุขสหคตสฺส ภวงฺคสฺส.\csromanspacer{} สุขสหคตํ กุสลากุสลํ นสุขสหคตสฺส วุฏฺฐานสฺส.\csromanspacer{} กิริยํ วุฏฺฐานสฺส.\csromanspacer{} ผลํ วุฏฺฐานสฺส อนนฺตรปจฺจเยน ปจฺจโย. (มูลํ กาตพฺพํ.) ปุริมา ปุริมา สุขสหคตา ขนฺธา ปจฺฉิมานํ ปจฺฉิมานํ สุขสหคตานํ ขนฺธานํ สุขสฺส จ อนนฺตรปจฺจเยน ปจฺจโย. (3)}

\prose{นสุขสหคโต ธมฺโม นสุขสหคตสฺส ธมฺมสฺส อนนฺตรปจฺจเยน ปจฺจโย.\csromanspacer{} ปุริมา ปุริมา นสุขสหคตา ฯเปฯ. (มูลํ.\csromanspacer{} ตีณิปิ สปฺปีติกทุกสทิสา.)}

\proseitem{172}{สุขสหคโต จ นสุขสหคโต จ ธมฺมา สุขสหคตสฺส ธมฺมสฺส อนนฺตรปจฺจเยน ปจฺจโย.\csromanspacer{} ปุริมา ปุริมา สุขสหคตา ขนฺธา จ สุขญฺจ ปจฺฉิมานํ ปจฺฉิมานํ สุขสหคตานํ ขนฺธานํ อนนฺตรปจฺจเยน ปจฺจโย. (มูลํ กาตพฺพํ.) ปุริมา ปุริมา สุขสหคตา ขนฺธา จ สุขญฺจ ปจฺฉิมสฺส ปจฺฉิมสฺส สุขสฺส อนนฺตรปจฺจเยน ปจฺจโย.\csromanspacer{} สุขสหคตํ จุติจิตฺตญฺจ สุขญฺจ นสุขสหคตสฺส อุปปตฺติจิตฺตสฺส.\csromanspacer{} สุขสหคตํ ภวงฺคญฺจ สุขญฺจ อาวชฺชนาย.\csromanspacer{} สุขสหคตํ กายวิญฺญาณญฺจ สุขญฺจ วิปากมโนธาตุยา.\csromanspacer{} สุขสหคตา วิปากมโนวิญฺญาณธาตุ จ สุขญฺจ กิริยมโนวิญฺญาณธาตุยา.\csromanspacer{} สุขสหคตํ ภวงฺคญฺจ สุขญฺจ นสุขสหคตสฺส ภวงฺคสฺส.\csromanspacer{} สุขสหคตํ กุสลากุสลญฺจ สุขญฺจ นสุขสหคตสฺส วุฏฺฐานสฺส.\csromanspacer{} กิริยํ วุฏฺฐานสฺส.\csromanspacer{} ผลํ วุฏฺฐานสฺส อนนฺตรปจฺจเยน ปจฺจโย. (มูลํ กาตพฺพํ.) ปุริมา ปุริมา สุขสหคตา ขนฺธา จ สุขญฺจ ปจฺฉิมานํ ปจฺฉิมานํ สุขสหคตานํ ขนฺธานํ สุขสฺส จ อนนฺตรปจฺจเยน ปจฺจโย. (3)}

\prose{สมนนฺตรปจฺจเยน ปจฺจโย.\csromanspacer{} สหชาตปจฺจเยน ปจฺจโย.\csromanspacer{} อญฺญมญฺญปจฺจเยน ปจฺจโย.\csromanspacer{} นิสฺสยปจฺจเยน ปจฺจโย.}

\csromanheader{2}{\textbf{อุปนิสฺสย}}
\proseitem{173}{สุขสหคโต ธมฺโม สุขสหคตสฺส ธมฺมสฺส อุปนิสฺสยปจฺจเยน ปจฺจโย.\csromanspacer{} ตีณิ.}

\csromanheader{2}{91. สุขสหคตทุก}
\prose{\csromanfolio{115}นสุขสหคโต ธมฺโม นสุขสหคตสฺส ธมฺมสฺส อุปนิสฺสยปจฺจเยน ปจฺจโย.\csromanspacer{} อารมฺมณูปนิสฺสโย, อนนฺตรูปนิสฺสโย, ปกตูปนิสฺสโย ฯเปฯ.\csromanspacer{} ปกตูปนิสฺสโย\csromandash{}นสุขสหคตํ สทฺธํ อุปนิสฺสาย นสุขสหคเตน จิตฺเตน ทานํ เทติ.\csromanspacer{} สีลํ ฯเปฯ สมาปตฺติํ อุปฺปาเทติ, มานํ ชปฺเปติ, ทิฏฺฐิํ คณฺหาติ.\csromanspacer{} นสุขสหคตํ สีลํ ฯเปฯ ปญฺญํ.\csromanspacer{} ราคํ ฯเปฯ.\csromanspacer{} ปตฺถนํ.\csromanspacer{} กายิกํ สุขํ.\csromanspacer{} กายิกํ ทุกฺขํ.\csromanspacer{} อุตุํ.\csromanspacer{} โภชนํ.\csromanspacer{} เสนาสนํ สุขํ อุปนิสฺสาย นสุขสหคเตน จิตฺเตน ทานํ เทติ ฯเปฯ สมาปตฺติํ อุปฺปาเทติ.\csromanspacer{} ปาณํ หนติ ฯเปฯ สํฆํ ภินฺทติ.\csromanspacer{} นสุขสหคตา สทฺธา ฯเปฯ.\csromanspacer{} เสนาสนํ สุขญฺจ นสุขสหคตาย สทฺธาย ฯเปฯ ปญฺญาย.\csromanspacer{} ราคสฺส.\csromanspacer{} โทสสฺส ฯเปฯ ปตฺถนาย.\csromanspacer{} กายิกสฺส สุขสฺส, กายิกสฺส ทุกฺขสฺส, มคฺคสฺส, ผลสมาปตฺติยา สุขสฺส จ อุปนิสฺสยปจฺจเยน ปจฺจโย. (1)}

\prose{นสุขสหคโต ธมฺโม สุขสหคตสฺส ธมฺมสฺส อุปนิสฺสยปจฺจเยน ปจฺจโย.\csromanspacer{} อารมฺมณูปนิสฺสโย, อนนฺตรูปนิสฺสโย, ปกตูปนิสฺสโย ฯเปฯ.\csromanspacer{} ปกตูปนิสฺสโย\csromandash{}นสุขสหคตํ สทฺธํ อุปนิสฺสาย สุขสหคเตน จิตฺเตน ทานํ เทติ ฯเปฯ สมาปตฺติํ อุปฺปาเทติ.\csromanspacer{} มานํ ชปฺเปติ, ทิฏฺฐิํ คณฺหาติ.\csromanspacer{} นสุขสหคตํ สีลํ ฯเปฯ ปญฺญํ.\csromanspacer{} ราคํ ฯเปฯ.\csromanspacer{} ปตฺถนํ.\csromanspacer{} กายิกํ สุขํ.\csromanspacer{} กายิกํ ทุกฺขํ.\csromanspacer{} อุตุํ.\csromanspacer{} โภชนํ.\csromanspacer{} เสนาสนํ สุขํ อุปนิสฺสาย สุขสหคเตน จิตฺเตน ทานํ เทติ ฯเปฯ.\csromanspacer{} สมาปตฺติํ อุปฺปาเทติ.\csromanspacer{} สุขสหคเตน จิตฺเตน อทินฺนํ อาทิยติ.\csromanspacer{} มุสา ฯเปฯ ปิสุณํ ฯเปฯ สมฺผํ ฯเปฯ สนฺธิํ ฯเปฯ นิลฺโลปํ ฯเปฯ เอกาคาริกํ ฯเปฯ ปริปนฺเถ ฯเปฯ ปรทารํ ฯเปฯ คามฆาตํ ฯเปฯ นิคมฆาตํ กโรติ.\csromanspacer{} นสุขสหคตา สทฺธา ฯเปฯ.\csromanspacer{} เสนาสนํ สุขญฺจ สุขสหคตาย สทฺธาย ฯเปฯ ปญฺญาย.\csromanspacer{} ราคสฺส.\csromanspacer{} มานสฺส.\csromanspacer{} ทิฏฺฐิยา.\csromanspacer{} ปตฺถนาย.\csromanspacer{} กายิกสฺส สุขสฺส.\csromanspacer{} มคฺคสฺส, ผลสมาปตฺติยา อุปนิสฺสยปจฺจเยน ปจฺจโย. (2)}

\prose{นสุขสหคโต ธมฺโม สุขสหคตสฺส จ นสุขสหคตสฺส จ ธมฺมสฺส อุปนิสฺสยปจฺจเยน ปจฺจโย.\csromanspacer{} อารมฺมณูปนิสฺสโย, อนนฺตรูปนิสฺสโย, ปกตูปนิสฺสโย ฯเปฯ.\csromanspacer{} ปกตูปนิสฺสโย\csromandash{} นสุขสหคตํ สทฺธํ อุปนิสฺสาย สุขสหคเตน จิตฺเตน ทานํ เทติ ฯเปฯ. (ทุติยคมนสทิสํ.) มานํ ชปฺเปติ, ทิฏฺฐิํ คณฺหาติ.\csromanspacer{} นสุขสหคตํ สีลํ ฯเปฯ.\csromanspacer{} เสนาสนํ สุขํ อุปนิสฺสาย ทานํ เทติ ฯเปฯ สมาปตฺติํ อุปฺปาเทติ.\csromanspacer{} สุขสหคเตน จิตฺเตน อทินฺนํ อาทิยติ ฯเปฯ.\csromanspacer{} นสุขสหคตา สทฺธา ฯเปฯ.\csromanspacer{} เสนาสนํ สุขญฺจ}

\prose{\csromanfolio{116}สุขสหคตาย สทฺธาย ฯเปฯ ปตฺถนาย.\csromanspacer{} กายิกสฺส สุขสฺส, มคฺคสฺส, ผลสมาปตฺติยา สุขสฺส จ อุปนิสฺสยปจฺจเยน ปจฺจโย. (3)}

\prose{สุขสหคโต จ นสุขสหคโต จ ธมฺมา สุขสหคตสฺส ธมฺมสฺส อุปนิสฺสยปจฺจเยน ปจฺจโย.\csromanspacer{} ตีณิ.}

\csromanheader{2}{\textbf{ปุเรชาตาทิ}}
\proseitem{174}{นสุขสหคโต ธมฺโม นสุขสหคตสฺส ธมฺมสฺส ปุเรชาตปจฺจเยน ปจฺจโย.\csromanspacer{} อารมฺมณปุเรชาตํ, วตฺถุปุเรชาตํ.\csromanspacer{}\csromanspacer{}\textbf{อารมฺมณปุเรชาตํ\csromandash{}}จกฺขุํ ฯเปฯ.\csromanspacer{} วตฺถุํ นสุขสหคเตน จิตฺเตน อนิจฺจโต ฯเปฯ วิปสฺสติ, อสฺสาเทติ อภินนฺทติ, ตํ อารพฺภ นสุขสหคโต ราโค อุปฺปชฺชติ ฯเปฯ โทมนสฺสํ อุปฺปชฺชติ, สุขํ อุปฺปชฺชติ.\csromanspacer{} ทิพฺเพน จกฺขุนา รูปํ ปสฺสติ ฯเปฯ.\csromanspacer{} โผฏฺฐพฺพายตนํ กายวิญฺญาณสฺส ฯเปฯ.\csromanspacer{} \textbf{วตฺถุปุเรชาตํ\csromandash{}}จกฺขายตนํ จกฺขุวิญฺญาณสฺส ฯเปฯ.\csromanspacer{} กายายตนํ กายวิญฺญาณสฺส ฯเปฯ.\csromanspacer{} วตฺถุ นสุขสหคตานํ ขนฺธานํ สุขสฺส จ ปุเรชาตปจฺจเยน ปจฺจโย. (1)}

\prose{นสุขสหคโต ธมฺโม สุขสหคตสฺส ธมฺมสฺส ปุเรชาตปจฺจเยน ปจฺจโย.\csromanspacer{} อารมฺมณปุเรชาตํ, วตฺถุปุเรชาตํ.\csromanspacer{}\csromanspacer{}\textbf{อารมฺมณปุเรชาตํ\csromandash{}}จกฺขุํ ฯเปฯ.\csromanspacer{} วตฺถุํ สุขสหคเตน จิตฺเตน อนิจฺจโต ฯเปฯ วิปสฺสติ, อสฺสาเทติ อภินนฺทติ, ตํ อารพฺภ สุขสหคโต ราโค อุปฺปชฺชติ.\csromanspacer{} ทิฏฺฐิ อุปฺปชฺชติ.\csromanspacer{} \textbf{วตฺถุปุเรชาตํ\csromandash{}}วตฺถุ สุขสหคตานํ ขนฺธานํ ปุเรชาตปจฺจเยน ปจฺจโย. (2)}

\prose{นสุขสหคโต ธมฺโม สุขสหคตสฺส จ นสุขสหคตสฺส จ ธมฺมสฺส ปุเรชาตปจฺจเยน ปจฺจโย.\csromanspacer{} อารมฺมณปุเรชาตํ, วตฺถุปุเรชาตํ.\csromanspacer{}\csromanspacer{}\textbf{อารมฺมณปุเรชาตํ\csromandash{}}จกฺขุํ ฯเปฯ.\csromanspacer{} วตฺถุํ สุขสหคเตน จิตฺเตน อนิจฺจโต ฯเปฯ วิปสฺสติ, อสฺสาเทติ อภินนฺทติ, ตํ อารพฺภ สุขญฺจ สมฺปยุตฺตกา ขนฺธา จ อุปฺปชฺชนฺติ.\csromanspacer{} \textbf{วตฺถุปุเรชาตํ\csromandash{}} วตฺถุ สุขสหคตานํ ขนฺธานํ สุขสฺส จ ปุเรชาตปจฺจเยน ปจฺจโย. (3)}

\prose{ปจฺฉาชาตปจฺจเยน ปจฺจโย.\csromanspacer{} ตีณิ.\csromanspacer{} อาเสวนปจฺจเยน ปจฺจโย.\csromanspacer{} นว.}

\csromanheader{2}{92. อุเปกฺขาสหคตทุก \textbf{กมฺมาทิ}}
\proseitem{175}{\csromanfolio{117}สุขสหคโต ธมฺโม สุขสหคตสฺส ธมฺมสฺส กมฺมปจฺจเยน ปจฺจโย.\csromanspacer{} สหชาตา, นานากฺขณิกา.\csromanspacer{}\csromanspacer{}(ฉ กมฺมานิ จตฺตาริ สหชาตาปิ นานากฺขณิกาปิ กาตพฺพา.\csromanspacer{} เทฺว นานากฺขณิกา.) วิปากปจฺจเยน ปจฺจโย.\csromanspacer{} นว.\csromanspacer{} อาหารปจฺจเยน ปจฺจโย.\csromanspacer{} จตฺตาริ.\csromanspacer{} อินฺทฺริยปจฺจเยน ปจฺจโย.\csromanspacer{} นว.\csromanspacer{} ฌานปจฺจเยน ปจฺจโย.\csromanspacer{} นว.\csromanspacer{} มคฺคปจฺจเยน ปจฺจโย.\csromanspacer{} จตฺตาริ.\csromanspacer{} สมฺปยุตฺตปจฺจเยน ปจฺจโย.\csromanspacer{} ฉ.\csromanspacer{} วิปฺปยุตฺตปจฺจเยน ปจฺจโย.\csromanspacer{} ปญฺจ.\csromanspacer{} อตฺถิปจฺจเยน ปจฺจโย.\csromanspacer{} นว.\csromanspacer{} นตฺถิปจฺจเยน ปจฺจโย.\csromanspacer{} นว.\csromanspacer{} วิคตปจฺจเยน ปจฺจโย.\csromanspacer{} นว.\csromanspacer{} อวิคตปจฺจเยน ปจฺจโย.\csromanspacer{} นว.}

\csromanheader{2}{1. ปจฺจยานุโลม 2. สงฺขฺยาวาร}
\csromanheader{2}{\textbf{สุทฺธ}}
\proseitem{176}{\textbf{เหตุ}ยา จตฺตาริ, อารมฺมเณ นว, อธิปติยา นว, อนนฺตเร นว, สมนนฺตเร นว, สหชาเต นว, อญฺญมญฺเญ นว, นิสฺสเย นว, อุปนิสฺสเย นว, ปุเรชาเต ตีณิ, ปจฺฉาชาเต ตีณิ, อาเสวเน นว, กมฺเม ฉ, วิปาเก นว, อาหาเร จตฺตาริ, อินฺทฺริเย นว, ฌาเน นว, มคฺเค จตฺตาริ, สมฺปยุตฺเต ฉ, วิปฺปยุตฺเต ปญฺจ, อตฺถิยา นว, นตฺถิยา นว, วิคเต นว, อวิคเต นว.}

\prose{(เอวํ ปจฺจนียวิภงฺโคปิ คณนาปิ สปฺปีติกทุกสทิสํ กาตพฺพํ, ยทิปิ วิมติ อตฺถิ, อนุโลมํ ปสฺสิตฺวา คเณตพฺพํ.)}

\nitthitam{สุขสหคตทุกํ นิฏฺฐิตํ.}
\csromanmatikaanchor{mtk.21}
\csromantocmarkref{section}{92. อุเปกฺขาสหคตทุก}{mtk.21}
\bookmark[dest={mtk.21},level=1]{92. อุเปกฺขาสหคตทุก}
\csromanheader{1}{92. อุเปกฺขาสหคตทุก 1. ปฏิจฺจวาร}
\proseitem{177}{อุเปกฺขาสหคตํ ธมฺมํ ปฏิจฺจ อุเปกฺขาสหคโต ธมฺโม อุปฺปชฺชติ เหตุปจฺจยา.\csromanspacer{} อุเปกฺขาสหคตํ เอกํ ขนฺธํ ปฏิจฺจ เทฺว ขนฺธา.\csromanspacer{} เทฺว \csromanfolio{118}ขนฺเธ ฯเปฯ.\csromanspacer{} ปฏิสนฺธิกฺขเณ ฯเปฯ.\csromanspacer{} อุเปกฺขาสหคตํ ธมฺมํ ปฏิจฺจ น-อุเปกฺขาสหคโต ธมฺโม อุปฺปชฺชติ เหตุปจฺจยา.\csromanspacer{} อุเปกฺขาสหคเต ขนฺเธ ปฏิจฺจ อุเปกฺขา จิตฺตสมุฏฺฐานญฺจ รูปํ.\csromanspacer{} ปฏิสนฺธิกฺขเณ ฯเปฯ.\csromanspacer{} อุเปกฺขาสหคตํ ธมฺมํ ปฏิจฺจ อุเปกฺขาสหคโต จ น-อุเปกฺขาสหคโต จ ธมฺมา อุปฺปชฺชนฺติ เหตุปจฺจยา.\csromanspacer{} อุเปกฺขาสหคตํ เอกํ ขนฺธํ ปฏิจฺจ เทฺว ขนฺธา อุเปกฺขา จ จิตฺตสมุฏฺฐานญฺจ รูปํ, เทฺว ขนฺเธ ฯเปฯ.\csromanspacer{} ปฏิสนฺธิกฺขเณ ฯเปฯ. (3)}

\prose{น-อุเปกฺขาสหคตํ ธมฺมํ ปฏิจฺจ น-อุเปกฺขาสหคโต ธมฺโม อุปฺปชฺชติ เหตุปจฺจยา.\csromanspacer{} น-อุเปกฺขาสหคตํ เอกํ ขนฺธํ ปฏิจฺจ ตโย ขนฺธา จิตฺตสมุฏฺฐานญฺจ รูปํ ฯเปฯ เทฺว ขนฺเธ ฯเปฯ.\csromanspacer{} อุเปกฺขํ ปฏิจฺจ จิตฺตสมุฏฺฐานํ รูปํ.\csromanspacer{} ปฏิสนฺธิกฺขเณ น-อุเปกฺขาสหคตํ เอกํ ขนฺธํ ปฏิจฺจ ตโย ขนฺธา กฏตฺตา จ รูปํ ฯเปฯ เทฺว ขนฺเธ ฯเปฯ.\csromanspacer{} ปฏิสนฺธิกฺขเณ อุเปกฺขํ ปฏิจฺจ กฏตฺตารูปํ.\csromanspacer{} ขนฺเธ ปฏิจฺจ วตฺถุ, วตฺถุํ ปฏิจฺจ ขนฺธา.\csromanspacer{} อุเปกฺขํ ปฏิจฺจ วตฺถุ, วตฺถุํ ปฏิจฺจ อุเปกฺขา, เอกํ มหาภูตํ ฯเปฯ. (สปฺปีติกทุกสทิสํ.\csromanspacer{} อนุโลเม นวปิ ปญฺหา.)}

\csromanheader{2}{1. ปจฺจยานุโลม 2. สงฺขฺยาวาร}
\proseitem{178}{เหตุยา นว, อารมฺมเณ นว, อธิปติยา นว ฯเปฯ ปุเรชาเต ฉ, อาเสวเน ฉ, กมฺเม นว, (สพฺพตฺถ นว.) อวิคเต นว.}

\csromanheader{2}{2. ปจฺจยปจฺจนีย 1. วิภงฺควาร}
\csromanheader{2}{\textbf{นเหตุ}}
\proseitem{179}{อุเปกฺขาสหคตํ ธมฺมํ ปฏิจฺจ อุเปกฺขาสหคโต ธมฺโม อุปฺปชฺชติ นเหตุปจฺจยา.\csromanspacer{} อเหตุกํ อุเปกฺขาสหคตํ เอกํ ขนฺธํ ปฏิจฺจ เทฺว ขนฺธา.\csromanspacer{} เทฺว ขนฺเธ ฯเปฯ.\csromanspacer{} อเหตุกปฏิสนฺธิขเณ ฯเปฯ.\csromanspacer{} วิจิกิจฺฉาสหคเต อุทฺธจฺจสหคเต ขนฺเธ ปฏิจฺจ วิจิกิจฺฉาสหคโต อุทฺธจฺจสหคโต โมโห.\csromanspacer{}\csromanspacer{}อุเปกฺขาสหคตํ ธมฺมํ ปฏิจฺจ น-อุเปกฺขาสหคโต ธมฺโม อุปฺปชฺชติ นเหตุปจฺจยา.\csromanspacer{} อเหตุเก อุเปกฺขาสหคเต ขนฺเธ ปฏิจฺจ อุเปกฺขา จิตฺตสมุฏฺฐานญฺจ รูปํ.\csromanspacer{} อเหตุกปฏิสนฺธิกฺขเณ ฯเปฯ.\csromanspacer{} อุเปกฺขาสหคตํ}

\vspace{\tipitakaparskip}
\csromancenter{อุเปกฺขาสหคตทุก ธมฺมํ ปฏิจฺจ อุเปกฺขาสหคโต จ น-อุเปกฺขาสหคโต จ ธมฺมา อุปฺปชฺชนฺติ นเหตุปจฺจยา.\csromanspacer{} อเหตุกํ อุเปกฺขาสหคตํ เอกํ ขนฺธํ ปฏิจฺจ เทฺว ขนฺธา อุเปกฺขา จ จิตฺตสมุฏฺฐานญฺจ รูปํ.\csromanspacer{} เทฺว ขนฺเธ ฯเปฯ.\csromanspacer{} อเหตุกปฏิสนฺธิกฺขเณ ฯเปฯ. (3)}
\proseitem{180}{\csromanfolio{119}น-อุเปกฺขาสหคตํ ธมฺมํ ปฏิจฺจ น-อุเปกฺขาสหคโต ธมฺโม อุปฺปชฺชติ นเหตุปจฺจยา.\csromanspacer{} อเหตุกํ น-อุเปกฺขาสหคตํ เอกํ ขนฺธํ ปฏิจฺจ ตโย ขนฺธา จิตฺตสมุฏฺฐานญฺจ รูปํ ฯเปฯ เทฺว ขนฺเธ ฯเปฯ.\csromanspacer{} อเหตุกํ อุเปกฺขํ ปฏิจฺจ จิตฺตสมุฏฺฐานํ รูปํ.\csromanspacer{} อเหตุกปฏิสนฺธิกฺขเณ ฯเปฯ.\csromanspacer{} อุเปกฺขํ ปฏิจฺจ วตฺถุ, วตฺถุํ ปฏิจฺจ อุเปกฺขา.\csromanspacer{} เอกํ มหาภูตํ ฯเปฯ. (ยาว อสญฺญสตฺตา.) (1)}

\prose{น-อุเปกฺขาสหคตํ ธมฺมํ ปฏิจฺจ อุเปกฺขาสหคโต ธมฺโม อุปฺปชฺชติ นเหตุปจฺจยา.\csromanspacer{} อเหตุกํ อุเปกฺขํ ปฏิจฺจ สมฺปยุตฺตกา ขนฺธา.\csromanspacer{} อเหตุกปฏิสนฺธิกฺขเณ อุเปกฺขํ ปฏิจฺจ สมฺปยุตฺตกา ขนฺธา.\csromanspacer{} อเหตุกปฏิสนฺธิกฺขเณ วตฺถุํ ปฏิจฺจ อุเปกฺขาสหคตา ขนฺธา.\csromanspacer{} วิจิกิจฺฉาสหคตํ อุทฺธจฺจสหคตํ อุเปกฺขํ ปฏิจฺจ วิจิกิจฺฉาสหคโต อุทฺธจฺจสหคโต โมโห. (2)}

\prose{น-อุเปกฺขาสหคตํ ธมฺมํ ปฏิจฺจ อุเปกฺขาสหคโต จ น-อุเปกฺขาสหคโต จ ธมฺมา อุปฺปชฺชนฺติ นเหตุปจฺจยา.\csromanspacer{} อเหตุกํ อุเปกฺขํ ปฏิจฺจ สมฺปยุตฺตกา ขนฺธา จิตฺตสมุฏฺฐานญฺจ รูปํ.\csromanspacer{} อเหตุกํ อุเปกฺขํ ปฏิจฺจ สมฺปยุตฺตกา ขนฺธา.\csromanspacer{} มหาภูเต ปฏิจฺจ จิตฺตสมุฏฺฐานํ รูปํ.\csromanspacer{} อเหตุกปฏิสนฺธิกฺขเณ อุเปกฺขํ ปฏิจฺจ สมฺปยุตฺตกา ขนฺธา กฏตฺตา จ รูปํ.\csromanspacer{} อเหตุกปฏิสนฺธิกฺขเณ อุเปกฺขํ ปฏิจฺจ สมฺปยุตฺตกา ขนฺธา.\csromanspacer{} มหาภูเต ปฏิจฺจ กฏตฺตารูปํ.\csromanspacer{} อเหตุกปฏิสนฺธิกฺขเณ วตฺถุํ ปฏิจฺจ อุเปกฺขาสหคตา ขนฺธา.\csromanspacer{} มหาภูเต ปฏิจฺจ กฏตฺตารูปํ.\csromanspacer{} อเหตุกปฏิสนฺธิกฺขเณ วตฺถุํ ปฏิจฺจ อุเปกฺขา จ สมฺปยุตฺตกา จ ขนฺธา. (3)}

\proseitem{181}{อุเปกฺขาสหคตญฺจ น-อุเปกฺขาสหคตญฺจ ธมฺมํ ปฏิจฺจ อุเปกฺขาสหคโต ธมฺโม อุปฺปชฺชติ นเหตุปจฺจยา.\csromanspacer{} อเหตุกํ อุเปกฺขาสหคตํ เอกํ ขนฺธญฺจ อุเปกฺขญฺจ ปฏิจฺจ เทฺว ขนฺธา.\csromanspacer{} เทฺว ขนฺเธ ฯเปฯ.\csromanspacer{} อเหตุกปฏิสนฺธิกฺขเณ อุเปกฺขาสหคตํ เอกํ ขนฺธญฺจ อุเปกฺขญฺจ ปฏิจฺจ เทฺว ขนฺธา.\csromanspacer{} เทฺว ขนฺเธ ฯเปฯ.\csromanspacer{} อเหตุกปฏิสนฺธิกฺขเณ อุเปกฺขาสหคตํ เอกํ ขนฺธญฺจ วตฺถุญฺจ ปฏิจฺจ เทฺว ขนฺธา.\csromanspacer{} เทฺว ขนฺเธ ฯเปฯ.\csromanspacer{} วิจิกิจฺฉาสหคเต อุทฺธจฺจสหคเต ขนฺเธ จ อุเปกฺขญฺจ ปฏิจฺจ วิจิกิจฺฉาสหคโต อุทฺธจฺจสหคโต โมโห. (1)}

\prose{\csromanfolio{120}อุเปกฺขาสหคตญฺจ น-อุเปกฺขาสหคตญฺจ ธมฺมํ ปฏิจฺจ น-อุเปกฺขาสหคโต ธมฺโม อุปฺปชฺชติ นเหตุปจฺจยา.\csromanspacer{} อเหตุเก อุเปกฺขาสหคเต ขนฺเธ จ อุเปกฺขญฺจ ปฏิจฺจ จิตฺตสมุฏฺฐานํ รูปํ.\csromanspacer{} อเหตุเก อุเปกฺขาสหคเต ขนฺเธ จ มหาภูเต จ ปฏิจฺจ จิตฺตสมุฏฺฐานํ รูปํ.\csromanspacer{} อเหตุกปฏิสนฺธิกฺขเณ อุเปกฺขาสหคเต ขนฺเธ จ อุเปกฺขญฺจ ปฏิจฺจ กฏตฺตารูปํ.\csromanspacer{} อเหตุกปฏิสนฺธิกฺขเณ อุเปกฺขาสหคเต ขนฺเธ จ มหาภูเต จ ปฏิจฺจ กฏตฺตารูปํ.\csromanspacer{} อเหตุกปฏิสนฺธิกฺขเณ อุเปกฺขาสหคเต ขนฺเธ จ วตฺถุญฺจ ปฏิจฺจ อุเปกฺขา. (2)}

\prose{อุเปกฺขาสหคตญฺจ น-อุเปกฺขาสหคตญฺจ ธมฺมํ ปฏิจฺจ อุเปกฺขาสหคโต จ น-อุเปกฺขาสหคโต จ ธมฺมา อุปฺปชฺชนฺติ นเหตุปจฺจยา.\csromanspacer{} อเหตุกํ อุเปกฺขาสหคตํ เอกํ ขนฺธญฺจ อุเปกฺขญฺจ ปฏิจฺจ เทฺว ขนฺธา จิตฺตสมุฏฺฐานญฺจ รูปํ.\csromanspacer{} เทฺว ขนฺเธ ฯเปฯ.\csromanspacer{} อเหตุกํ อุเปกฺขาสหคตํ เอกํ ขนฺธญฺจ อุเปกฺขญฺจ ปฏิจฺจ เทฺว ขนฺธา.\csromanspacer{} เทฺว ขนฺเธ ฯเปฯ.\csromanspacer{} อเหตุเก อุเปกฺขาสหคเต ขนฺเธ จ อุเปกฺขญฺจ มหาภูเต จ ปฏิจฺจ จิตฺตสมุฏฺฐานํ รูปํ.\csromanspacer{} อเหตุกปฏิสนฺธิกฺขเณ อุเปกฺขาสหคตํ เอกํ ขนฺธญฺจ อุเปกฺขญฺจ ปฏิจฺจ เทฺว ขนฺธา กฏตฺตา จ รูปํ.\csromanspacer{} เทฺว ขนฺเธ ฯเปฯ.\csromanspacer{} อเหตุกปฏิสนฺธิกฺขเณ อุเปกฺขาสหคตํ เอกํ ขนฺธญฺจ อุเปกฺขญฺจ ปฏิจฺจ เทฺว ขนฺธา.\csromanspacer{} เทฺว ขนฺเธ ฯเปฯ.\csromanspacer{} อเหตุเก อุเปกฺขาสหคเต ขนฺเธ จ อุเปกฺขญฺจ มหาภูเต จ ปฏิจฺจ กฏตฺตารูปํ.\csromanspacer{} อเหตุกปฏิสนฺธิกฺขเณ อุเปกฺขาสหคตํ เอกํ ขนฺธญฺจ วตฺถุญฺจ ปฏิจฺจ เทฺว ขนฺธา อุเปกฺขา จ.\csromanspacer{} เทฺว ขนฺเธ ฯเปฯ. (สํขิตฺตํ.) (3)}

\csromanheader{2}{2. ปจฺจยปจฺจนีย 2. สงฺขฺยาวาร}
\csromanheader{2}{\textbf{สุทฺธ}}
\proseitem{182}{นเหตุยา นว, น-อารมฺมเณ ตีณิ, น-อธิปติยา นว, น-อนนฺตเร ตีณิ, น-อญฺญมญฺเญ ตีณิ, น-อุปนิสฺสเย ตีณิ, นปุเรชาเต นว, นปจฺฉาชาเต นว, น-อาเสวเน นว, นกมฺเม จตฺตาริ, นวิปาเก นว, น-อาหาเร เอกํ, น-อินฺทฺริเย เอกํ, นฌาเน นว, นมคฺเค นว, นสมฺปยุตฺเต ตีณิ, นวิปฺปยุตฺเต ฉ, โนนตฺถิยา ตีณิ, โนวิคเต ตีณิ.}

\vspace{\tipitakaparskip}
\csromancenter{(เอวํ อิตเร เทฺว คณนาปิ สหชาตวาโรปิ กาตพฺโพ.)}
\csromanheader{2}{92. อุเปกฺขาสหคตทุก \textbf{92. อุเปกฺขาสหคตทุก 3. ปจฺจยวาร}}
\csromanheader{2}{1. ปจฺจยานุโลมาทิ}
\proseitem{183}{\csromanfolio{121}อุเปกฺขาสหคตํ ธมฺมํ ปจฺจยา อุเปกฺขาสหคโต ธมฺโม อุปฺปชฺชติ เหตุปจฺจยา.\csromanspacer{} อุเปกฺขาสหคตํ เอกํ ขนฺธํ ปจฺจยา เทฺว ขนฺธา.\csromanspacer{} เทฺว ขนฺเธ ฯเปฯ.\csromanspacer{} ปฏิสนฺธิกฺขเณ ฯเปฯ. (ยถา สวิตกฺกทุกสทิสํ, ปจฺจยวาเร นานากรณํ, “อุเปกฺขนฺ”ติ นวปิ ปญฺหา กาตพฺพา.\csromanspacer{} ปฏิสนฺธิปวตฺติปิ วตฺถุปิ.)}

\vspace{\tipitakaparskip}
\csromancenter{\textbf{เหตุ}ยา นว, อารมฺมเณ นว ฯเปฯ ปุเรชาเต นว, อาเสวเน นว, (สพฺพตฺถ นว.) อวิคเต นว.\csromanspacer{} อนุโลมํ.}
\proseitem{184}{อุเปกฺขาสหคตํ ธมฺมํ ปจฺจยา อุเปกฺขาสหคโต ธมฺโม อุปฺปชฺชติ นเหตุปจฺจยา.\csromanspacer{} อเหตุกํ อุเปกฺขาสหคตํ เอกํ ขนฺธํ ปจฺจยา เทฺว ขนฺธา.\csromanspacer{} เทฺว ขนฺเธ ฯเปฯ.\csromanspacer{} อเหตุกปฏิสนฺธิกฺขเณ ฯเปฯ.\csromanspacer{} วิจิกิจฺฉาสหคเต อุทฺธจฺจสหคเต ขนฺเธ ปจฺจยา วิจิกิจฺฉาสหคโต อุทฺธจฺจสหคโต โมโห. (เอวํ นวปิ ปญฺหา ปวตฺติปฏิสนฺธิโย ยถา สวิตกฺกทุกสฺส เอวํ กาตพฺพา.\csromanspacer{} ตีณิเยว โมโห.\csromanspacer{} ปวตฺเต วตฺถุปิ กาตพฺพา.)}

\prose{นเหตุยา นว, น-อารมฺมเณ ตีณิ, น-อธิปติยา นว, น-อนนฺตเร ตีณิ, นสมนนฺตเร ตีณิ, น-อุปนิสฺสเย ตีณิ, นปุเรชาเต นว, นปจฺฉาชาเต นว, น-อาเสวเน นว, นกมฺเม จตฺตาริ, นวิปาเก นว, น-อาหาเร เอกํ, น-อินฺทฺริเย เอกํ, นฌาเน เอกํ, นมคฺเค นว, นสมฺปยุตฺเต ตีณิ, นวิปฺปยุตฺเต ฉ, โนนตฺถิยา ตีณิ, โนวิคเต ตีณิ.}

\vspace{\tipitakaparskip}
\csromancenter{ปจฺจนียํ.}
\csromancenter{(เอวํ อิตเร เทฺว คณนาปิ นิสฺสยวาโรปิ กาตพฺโพ.)}
\csromanheader{2}{92. อุเปกฺขาสหคตทุก 5. สํสฏฺฐวาร}
\csromanheader{2}{1. ปจฺจยานุโลมาทิ}
\proseitem{185}{อุเปกฺขาสหคตํ ธมฺมํ สํสฏฺโฐ อุเปกฺขาสหคโต ธมฺโม อุปฺปชฺชติ เหตุปจฺจยา.\csromanspacer{} อุเปกฺขาสหคตํ เอกํ ขนฺธํ สํสฏฺฐา เทฺว ขนฺธา.\csromanspacer{} เทฺว \csromanfolio{122}ขนฺเธ ฯเปฯ.\csromanspacer{} ปฏิสนฺธิกฺขเณ ฯเปฯ. (ยถา สวิตกฺกทุกํ สมฺปยุตฺตวาโร, เอวํ กาตพฺโพ.)}

\prose{เหตุยา ฉ, อารมฺมเณ ฉ, (สพฺพตฺถ ฉ.) อวิคเต ฉ.\csromanspacer{} อนุโลมํ.}

\prose{อุเปกฺขาสหคตํ ธมฺมํ สํสฏฺโฐ อุเปกฺขาสหคโต ธมฺโม อุปฺปชฺชติ นเหตุปจฺจยา.\csromanspacer{} อเหตุกํ อุเปกฺขาสหคตํ เอกํ ขนฺธํ สํสฏฺฐา เทฺว ขนฺธา.\csromanspacer{} เทฺว ขนฺเธ ฯเปฯ.\csromanspacer{} อเหตุกปฏิสนฺธิกฺขเณ ฯเปฯ.\csromanspacer{} วิจิกิจฺฉาสหคเต อุทฺธจฺจสหคเต ขนฺเธ สํสฏฺโฐ วิจิกิจฺฉาสหคโต อุทฺธจฺจสหคโต โมโห. (เอวํ ปญฺจ ปญฺหา ยถา สวิตกฺกทุเก, เอวํ กาตพฺพา.)}

\prose{นเหตุยา ฉ, น-อธิปติยา ฉ, นปุเรชาเต ฉ, นปจฺฉาชาเต ฉ, น-อาเสวเน ฉ, นกมฺเม จตฺตาริ, นวิปาเก ฉ, นวิปฺปยุตฺเต ฉ.}

\vspace{\tipitakaparskip}
\csromancenter{ปจฺจนียํ.}
\csromancenter{(เอวํ อิตเร เทฺว คณนาปิ สมฺปยุตฺตวาโรปิ กาตพฺโพ.)}
\csromanheader{2}{92. อุเปกฺขาสหคตทุก 7. ปญฺหาวาร}
\csromanheader{2}{\textbf{เหตุ-อาทิ}}
\proseitem{186}{อุเปกฺขาสหคโต ธมฺโม อุเปกฺขาสหคตสฺส ธมฺมสฺส เหตุปจฺจเยน ปจฺจโย.\csromanspacer{} อุเปกฺขาสหคตา เหตู สมฺปยุตฺตกานํ ขนฺธานํ เหตุปจฺจเยน ปจฺจโย.\csromanspacer{} ปฏิสนฺธิกฺขเณ ฯเปฯ. (เอวํ จตฺตาริ ปญฺหา ยถา สวิตกฺกทุกสฺส.)}

\prose{อุเปกฺขาสหคโต ธมฺโม อุเปกฺขาสหคตสฺส ธมฺมสฺส อารมฺมณปจฺจเยน ปจฺจโย.\csromanspacer{} อธิปติปจฺจเยน ปจฺจโย. (ยถา สปฺปีติกทุกํ, เอวํ อารมฺมณมฺปิ อธิปติปิ วิตฺถาเรตพฺพา, “อุเปกฺขา”ติ นานํ.)}

\csromanheader{2}{92. อุเปกฺขาสหคตทุก \textbf{อนนฺตราทิ}}
\proseitem{187}{\csromanfolio{123}อุเปกฺขาสหคโต ธมฺโม อุเปกฺขาสหคตสฺส ธมฺมสฺส อนนฺตรปจฺจเยน ปจฺจโย.\csromanspacer{} ปุริมา ปุริมา อุเปกฺขาสหคตา ขนฺธา ปจฺฉิมานํ ปจฺฉิมานํ อุเปกฺขาสหคตานํ ขนฺธานํ อนนฺตรปจฺจเยน ปจฺจโย. (1)}

\prose{อุเปกฺขาสหคโต ธมฺโม น-อุเปกฺขาสหคตสฺส ธมฺมสฺส อนนฺตรปจฺจเยน ปจฺจโย.\csromanspacer{} ปุริมา ปุริมา อุเปกฺขาสหคตา ขนฺธา ปจฺฉิมาย ปจฺฉิมาย อุเปกฺขาย อนนฺตรปจฺจเยน ปจฺจโย.\csromanspacer{} อุเปกฺขาสหคตํ จุติจิตฺตํ น-อุเปกฺขาสหคตสฺส อุปปตฺติจิตฺตสฺส.\csromanspacer{} อาวชฺชนา น-อุเปกฺขาสหคตานํ ขนฺธานํ.\csromanspacer{} วิปากมโนธาตุ น-อุเปกฺขาสหคตาย วิปากมโนวิญฺญาณธาตุยา.\csromanspacer{} อุเปกฺขาสหคตํ ภวงฺคํ น-อุเปกฺขาสหคตสฺส ภวงฺคสฺส.\csromanspacer{} อุเปกฺขาสหคตํ กุสลากุสลํ น-อุเปกฺขาสหคตสฺส วุฏฺฐานสฺส.\csromanspacer{} กิริยํ วุฏฺฐานสฺส.\csromanspacer{} ผลํ วุฏฺฐานสฺส.\csromanspacer{} นิโรธา วุฏฺฐหนฺตสฺส เนวสญฺญานาสญฺญายตนํ น-อุเปกฺขาสหคตาย ผลสมาปตฺติยา อนนฺตรปจฺจเยน ปจฺจโย. (2)}

\prose{อุเปกฺขาสหคโต ธมฺโม อุเปกฺขาสหคตสฺส จ น-อุเปกฺขาสหคตสฺส จ ธมฺมสฺส อนนฺตรปจฺจเยน ปจฺจโย.\csromanspacer{} ปุริมา ปุริมา อุเปกฺขาสหคตา ขนฺธา ปจฺฉิมานํ ปจฺฉิมานํ อุเปกฺขาสหคตานํ ขนฺธานํ อุเปกฺขาย จ อนนฺตรปจฺจเยน.\csromanspacer{} ปจฺจโย. (3)}

\proseitem{188}{น-อุเปกฺขาสหคโต ธมฺโม น-อุเปกฺขาสหคตสฺส ธมฺมสฺส อนนฺตรปจฺจเยน ปจฺจโย.\csromanspacer{} ปุริมา ปุริมา อุเปกฺขา ปจฺฉิมาย ปจฺฉิมาย อุเปกฺขาย อนนฺตรปจฺจเยน ปจฺจโย.\csromanspacer{} ปุริมา ปุริมา น-อุเปกฺขาสหคตา ขนฺธา ปจฺฉิมานํ ปจฺฉิมานํ น-อุเปกฺขาสหคตานํ ขนฺธานํ อนนฺตรปจฺจเยน ปจฺจโย.\csromanspacer{} อนุโลมํ โคตฺรภุสฺส ฯเปฯ.\csromanspacer{} อนุโลมํ ผลสมาปตฺติยา อนนฺตรปจฺจเยน ปจฺจโย. (มูลํ กาตพฺพํ.) ปุริมา ปุริมา อุเปกฺขา ปจฺฉิมานํ ปจฺฉิมานํ อุเปกฺขาสหคตานํ ขนฺธานํ อนนฺตรปจฺจเยน ปจฺจโย.\csromanspacer{} น-อุเปกฺขาสหคตํ จุติจิตฺตํ อุเปกฺขาสหคตสฺส อุปปตฺติจิตฺตสฺส.\csromanspacer{} น-อุปกฺขาสหคตํ ภวงฺคํ อาวชฺชนาย.\csromanspacer{} กายวิญฺญาณธาตุ วิปากมโนธาตุยา.\csromanspacer{} น-อุเปกฺขาสหคตา วิปากมโนวิญฺญาณธาตุ กิริยมโนวิญฺญาณธาตุยา. \csromanfolio{124}น-อุเปกฺขาสหคตํ ภวงฺคํ อุเปกฺขาสหคตสฺส ภวงฺคสฺส.\csromanspacer{} น-อุเปกฺขาสหคตํ กุสลากุสลํ อุเปกฺขาสหคตสฺส วุฏฺฐานสฺส.\csromanspacer{} กิริยํ วุฏฺฐานสฺส.\csromanspacer{} ผลํ วุฏฺฐานสฺส อนนฺตรปจฺจเยน ปจฺจโย. (2)}

\prose{น-อุเปกฺขาสหคโต ธมฺโม อุเปกฺขาสหคตสฺส จ น-อุเปกฺขาสหคตสฺส จ ธมฺมสฺส อนนฺตรปจฺจเยน ปจฺจโย.\csromanspacer{} ปุริมา ปุริมา อุเปกฺขา ปจฺฉิมานํ ปจฺฉิมานํ อุเปกฺขาสหคตานํ ขนฺธานํ อุเปกฺขาย จ อนนฺตรปจฺจเยน ปจฺจโย. (3)}

\proseitem{189}{อุเปกฺขาสหคโต จ น-อุเปกฺขาสหคโต จ ธมฺมา อุเปกฺขาสหคตสฺส ธมฺมสฺส อนนฺตรปจฺจเยน ปจฺจโย.\csromanspacer{} ปุริมา ปุริมา อุเปกฺขาสหคตา ขนฺธา จ อุเปกฺขา จ ปจฺฉิมานํ ปจฺฉิมานํ อุเปกฺขาสหคตานํ ขนฺธานํ อนนฺตรปจฺจเยน ปจฺจโย. (มูลํ กาตพฺพํ.) ปุริมา ปุริมา อุเปกฺขาสหคตา ขนฺธา จ อุเปกฺขา จ ปจฺฉิมาย ปจฺฉิมาย อุเปกฺขาย อนนฺตรปจฺจเยน ปจฺจโย.\csromanspacer{} อุเปกฺขาสหคตํ จุติจิตฺตญฺจ อุเปกฺขา จ น-อุเปกฺขาสหคตสฺส อุปปตฺติจิตฺตสฺส.\csromanspacer{} อาวชฺชนา จ อุเปกฺขา จ น-อุเปกฺขาสหคตานํ ขนฺธานํ.\csromanspacer{} วิปากมโนธาตุ จ อุเปกฺขา จ น-อุเปกฺขาสหคตาย วิปากมโนวิญฺญาณธาตุยา.\csromanspacer{} อุเปกฺขาสหคตํ ภวงฺคญฺจ อุเปกฺขา จ น-อุเปกฺขาสหคตสฺส ภวงฺคสฺส.\csromanspacer{} อุเปกฺขาสหคตํ กุสลากุสลญฺจ อุเปกฺขา จ น-อุเปกฺขาสหคตสฺส วุฏฺฐานสฺส.\csromanspacer{} กิริยญฺจ อุเปกฺขา จ วุฏฺฐานสฺส.\csromanspacer{} ผลญฺจ อุเปกฺขา จ วุฏฺฐานสฺส.\csromanspacer{} นิโรธา วุฏฺฐหนฺตสฺส เนวสญฺญานาสญฺญายตนญฺจ อุเปกฺขา จ น-อุเปกฺขาสหคตาย ผลสมาปตฺติยา อนนฺตรปจฺจเยน ปจฺจโย. (มูลํ กาตพฺพํ.) ปุริมา ปุริมา อุเปกฺขาสหคตา ขนฺธา จ อุเปกฺขา จ ปจฺฉิมานํ ปจฺฉิมานํ อุเปกฺขาสหคตานํ ขนฺธานํ อุเปกฺขาย จ อนนฺตรปจฺจเยน ปจฺจโย. (3)}

\prose{สมนนฺตรปจฺจเยน ปจฺจโย.\csromanspacer{} สหชาตปจฺจเยน ปจฺจโย.\csromanspacer{} นว.\csromanspacer{} อญฺญมญฺญปจฺจเยน ปจฺจโย.\csromanspacer{} นว.\csromanspacer{} นิสฺสยปจฺจเยน ปจฺจโย.\csromanspacer{} นว.}

\csromanheader{2}{\textbf{อุปนิสฺสย}}
\proseitem{190}{อุเปกฺขาสหคโต ธมฺโม อุเปกฺขาสหคตสฺส ธมฺมสฺส อุปนิสฺสยปจฺจเยน ปจฺจโย.\csromanspacer{} ตีณิ.}

\csromanheader{2}{92. อุเปกฺขาสหคตทุก}
\prose{\csromanfolio{125}น-อุเปกฺขาสหคโต ธมฺโม น-อุเปกฺขาสหคตสฺส ธมฺมสฺส อุปนิสฺสยปจฺจเยน ปจฺจโย.\csromanspacer{} อารมฺมณูปนิสฺสโย, อนนฺตรูปนิสฺสโย, ปกตูปนิสฺสโย ฯเปฯ.\csromanspacer{} ปกตูปนิสฺสโย\csromandash{}น-อุเปกฺขาสหคตํ สทฺธํ อุปนิสฺสาย น-อุเปกฺขาสหคเตน จิตฺเตน ทานํ เทติ.\csromanspacer{} สีลํ ฯเปฯ อุโปสถกมฺมํ ฯเปฯ.\csromanspacer{} น-อุเปกฺขาสหคตํ ฌานํ ฯเปฯ วิปสฺสนํ ฯเปฯ มคฺคํ ฯเปฯ สมาปตฺติํ อุปฺปาเทติ.\csromanspacer{} มานํ ชปฺเปติ.\csromanspacer{} ทิฏฺฐิํ คณฺหาติ.\csromanspacer{} น-อุเปกฺขาสหคตํ สีลํ ฯเปฯ ปญฺญํ.\csromanspacer{} ราคํ.\csromanspacer{} โทสํ.\csromanspacer{} โมหํ.\csromanspacer{} มานํ.\csromanspacer{} ทิฏฺฐิํ.\csromanspacer{} ปตฺถนํ.\csromanspacer{} กายิกํ สุขํ.\csromanspacer{} กายิกํ ทุกฺขํ.\csromanspacer{} อุตุํ.\csromanspacer{} โภชนํ.\csromanspacer{} เสนาสนํ อุเปกฺขํ อุปนิสฺสาย น-อุเปกฺขาสหคเตน จิตฺเตน ทานํ เทติ ฯเปฯ สมาปตฺติํ อุปฺปาเทติ.\csromanspacer{} ปาณํ หนติ ฯเปฯ สํฆํ ภินฺทติ.\csromanspacer{} น-อุเปกฺขาสหคตา สทฺธา ฯเปฯ.\csromanspacer{} เสนาสนํ อุเปกฺขา จ น-อุเปกฺขาสหคตาย สทฺธาย ฯเปฯ ปญฺญาย.\csromanspacer{} ราคสฺส ฯเปฯ ปตฺถนาย.\csromanspacer{} กายิกสฺส สุขสฺส.\csromanspacer{} กายิกสฺส ทุกฺขสฺส.\csromanspacer{} มคฺคสฺส, ผลสมาปตฺติยา อุเปกฺขาย จ อุปนิสฺสยปจฺจเยน ปจฺจโย. (1)}

\prose{น-อุเปกฺขาสหคโต ธมฺโม อุเปกฺขาสหคตสฺส ธมฺมสฺส อุปนิสฺสยปจฺจเยน ปจฺจโย. (ตีณิปิ อุปนิสฺสยา.) น-อุเปกฺขาสหคตํ สทฺธํ อุปนิสฺสาย อุเปกฺขาสหคเตน จิตฺเตน ทานํ เทติ ฯเปฯ สมาปตฺติํ อุปฺปาเทติ.\csromanspacer{} มานํ ชปฺเปติ.\csromanspacer{} ทิฏฺฐิํ คณฺหาติ.\csromanspacer{} น-อุเปกฺขาสหคตํ สีลํ ฯเปฯ.\csromanspacer{} เสนาสนํ อุเปกฺขํ อุปนิสฺสาย อุเปกฺขาสหคเตน จิตฺเตน ทานํ เทติ ฯเปฯ สมาปตฺติํ อุปฺปาเทติ.\csromanspacer{} อุเปกฺขาสหคเตน จิตฺเตน อทินฺนํ อาทิยติ.\csromanspacer{} มุสา ภณติ.\csromanspacer{} ปิสุณํ ฯเปฯ สมฺผํ ฯเปฯ สนฺธิํ ฯเปฯ นิลฺโลปํ ฯเปฯ เอกาคาริกํ ฯเปฯ ปริปนฺเถ ฯเปฯ ปรทารํ ฯเปฯ คามฆาตํ ฯเปฯ นิคมฆาตํ กโรติ.\csromanspacer{} น-อุเปกฺขาสหคตา สทฺธา ฯเปฯ.\csromanspacer{} เสนาสนํ อุเปกฺขา จ อุเปกฺขาสหคตาย สทฺธาย ฯเปฯ ปตฺถนาย.\csromanspacer{} มคฺคสฺส, ผลสมาปตฺติยา อุปนิสฺสยปจฺจเยน ปจฺจโย.}

\prose{น-อุเปกฺขาสหคโต ธมฺโม อุเปกฺขาสหคตสฺส จ น-อุเปกฺขาสหคตสฺส จ ธมฺมสฺส อุเปนิสฺสยปจฺจเยน ปจฺจโย. (ตีณิ อุปนิสฺสยา, ทุติยคมนสทิสา.) (3)}

\prose{อุเปกฺขาสหคโต จ น-อุเปกฺขาสหคโต จ ธมฺมา อุเปกฺขาสหคตสฺส ธมฺมสฺส อุปนิสฺสยปจฺจเยน ปจฺจโย.\csromanspacer{} ตีณิ.}

\csromanheader{2}{\textbf{ปุเรชาต}}
\proseitem{191}{\csromanfolio{126}น-อุเปกฺขาสหคโต ธมฺโม น-อุเปกฺขาสหคตสฺส ธมฺมสฺส ปุเรชาตปจฺจเยน ปจฺจโย.\csromanspacer{} ตีณิ. (สปฺปีติกทุกสทิสา.)}

\csromanheader{2}{\textbf{ปจฺฉาชาตาทิ}}
\proseitem{192}{อุเปกฺขาสหคโต ธมฺโม น-อุเปกฺขาสหคตสฺส ธมฺมสฺส ปจฺฉาชาตปจฺจเยน ปจฺจโย.\csromanspacer{} ตีณิ.\csromanspacer{} อาเสวนปจฺจเยน ปจฺจโย.\csromanspacer{} นว.\csromanspacer{} กมฺมปจฺจเยน ปจฺจโย.\csromanspacer{} ฉ. (จตฺตาริ สหชาตา นานากฺขณิกา กาตพฺพา.\csromanspacer{} เทฺว นานากฺขณิกา จ.) วิปากปจฺจเยน ปจฺจโย.\csromanspacer{} นว.\csromanspacer{} อาหารปจฺจเยน ปจฺจโย.\csromanspacer{} จตฺตาริ.\csromanspacer{} อินฺทฺริยปจฺจเยน ปจฺจโย.\csromanspacer{} นว.\csromanspacer{} ฌานปจฺจเยน ปจฺจโย.\csromanspacer{} นว.\csromanspacer{} มคฺคปจฺจเยน ปจฺจโย.\csromanspacer{} จตฺตาริ.\csromanspacer{} สมฺปยุตฺตปจฺจเยน ปจฺจโย.\csromanspacer{} ฉ.\csromanspacer{} วิปฺปยุตฺตปจฺจเยน ปจฺจโย.\csromanspacer{} ปญฺจ.\csromanspacer{} อตฺถิปจฺจเยน ปจฺจโย.\csromanspacer{} นว.\csromanspacer{} นตฺถิปจฺจเยน ปจฺจโย.\csromanspacer{} วิคตปจฺจเยน ปจฺจโย.\csromanspacer{} อวิคตปจฺจเยน ปจฺจโย. (อิมานิ ปจฺจยานิ สปฺปีติกกรเณน วิภชิตพฺพานิ.)}

\csromanheader{2}{1. ปจฺจยานุโลม 2. สงฺขฺยาวาร}
\csromanheader{2}{\textbf{สุทฺธ}}
\proseitem{193}{เหตุยา จตฺตาริ, อารมฺมเณ นว, อธิปติยา นว, อนนฺตเร นว, สมนนฺตเร นว, สหชาเต นว, อญฺญมญฺเญ นว, นิสฺสเย นว, อุปนิสฺสเย นว, ปุเรชาเต ตีณิ, ปจฺฉาชาเต ตีณิ, อาเสวเน นว, กมฺเม ฉ, วิปาเก นว, อาหาเร จตฺตาริ, อินฺทฺริเย นว, ฌาเน นว, มคฺเค จตฺตาริ, สมฺปยุตฺเต ฉ, วิปฺปยุตฺเต ปญฺจ, อตฺถิยา นว, นตฺถิยา นว, วิคเต นว, อวิคเต นว.}

\prose{(เอวํ ปจฺจนียวิภงฺโคปิ อิตเร ตีณิ คณนาปิ สปฺปีติกทุกสทิสา กาตพฺพา.)}

\nitthitam{อุเปกฺขาสหคตทุกํ นิฏฺฐิตํ.}
\csromanmatikaanchor{mtk.22}
\csromantocmarkref{section}{93. กามาวจรทุก}{mtk.22}
\bookmark[dest={mtk.22},level=1]{93. กามาวจรทุก}
\csromanheader{1}{93. กามาวจรทุก \textbf{93. กามาวจรทุก 1. ปฏิจฺจวาร}}
\proseitem{194}{\csromanfolio{127}กามาวจรํ ธมฺมํ ปฏิจฺจ กามาวจโร ธมฺโม อุปฺปชฺชติ เหตุปจฺจยา.\csromanspacer{} กามาวจรํ เอกํ ขนฺธํ ปฏิจฺจ ตโย ขนฺธา จิตฺตสมุฏฺฐานญฺจ รูปํ ฯเปฯ.\csromanspacer{} เทฺว ขนฺเธ ฯเปฯ.\csromanspacer{} ปฏิสนฺธิกฺขเณ ฯเปฯ.\csromanspacer{} เอกํ มหาภูตํ ฯเปฯ. (สํขิตฺตํ.).\csromanspacer{} กามาวจรํ ธมฺมํ ปฏิจฺจ นกามาวจโร ธมฺโม อุปฺปชฺชติ เหตุปจฺจยา.\csromanspacer{} ปฏิสนฺธิกฺขเณ วตฺถุํ ปฏิจฺจ นกามาวจรา ขนฺธา.\csromanspacer{}\csromanspacer{}กามาวจรํ ธมฺมํ ปฏิจฺจ กามาวจโร จ นกามาวจโร จ ธมฺมา อุปฺปชฺชนฺติ เหตุปจฺจยา.\csromanspacer{} ปฏิสนฺธิกฺขเณ วตฺถุํ ปฏิจฺจ นกามาวจรา ขนฺธา.\csromanspacer{} มหาภูเต ปฏิจฺจ กฏตฺตารูปํ. (3)}

\proseitem{195}{นกามาวจรํ ธมฺมํ ปฏิจฺจ นกามาวจโร ธมฺโม อุปฺปชฺชติ เหตุปจฺจยา.\csromanspacer{} นกามาวจรํ เอกํ ขนฺธํ ปฏิจฺจ ตโย ขนฺธา ฯเปฯ เทฺว ขนฺเธ ฯเปฯ.\csromanspacer{} ปฏิสนฺธิกฺขเณ ฯเปฯ.\csromanspacer{} นกามาวจรํ ธมฺมํ ปฏิจฺจ กามาวจโร ธมฺโม อุปฺปชฺชติ เหตุปจฺจยา.\csromanspacer{} นกามาวจเร ขนฺเธ ปฏิจฺจ จิตฺตสมุฏฺฐานํ รูปํ.\csromanspacer{} ปฏิสนฺธิกฺขเณ นกามาวจเร ขนฺเธ ปฏิจฺจ กฏตฺตารูปํ.\csromanspacer{}\csromanspacer{}นกามาวจรํ ธมฺมํ ปฏิจฺจ กามาวจโร จ นกามาวจโร จ ธมฺมา อุปฺปชฺชนฺติ เหตุปจฺจยา.\csromanspacer{} นกามาวจรํ เอกํ ขนฺธํ ปฏิจฺจ ตโย ขนฺธา จิตฺตสมุฏฺฐานญฺจ รูปํ ฯเปฯ เทฺว ขนฺเธ ฯเปฯ.\csromanspacer{} ปฏิสนฺธิกฺขเณ ฯเปฯ. (3)}

\proseitem{196}{กามาวจรญฺจ นกามาวจรญฺจ ธมฺมํ ปฏิจฺจ กามาวจโร ธมฺโม อุปฺปชฺชติ เหตุปจฺจยา.\csromanspacer{} นกามาวจเร ขนฺเธ จ มหาภูเต จ ปฏิจฺจ จิตฺตสมุฏฺฐานํ รูปํ.\csromanspacer{} ปฏิสนฺธิกฺขเณ นกามาวจเร ขนฺเธ จ มหาภูเต จ ปฏิจฺจ กฏตฺตารูปํ.\csromanspacer{}\csromanspacer{}กามาวจรญฺจ นกามาวจรญฺจ ธมฺมํ ปฏิจฺจ นกามาวจโร ธมฺโม อุปฺปชฺชติ เหตุปจฺจยา.\csromanspacer{} ปฏิสนฺธิกฺขเณ นกามาวจรํ เอกํ ขนฺธญฺจ วตฺถุญฺจ ปฏิจฺจ ตโย ขนฺธา ฯเปฯ เทฺว ขนฺเธ จ ฯเปฯ.\csromanspacer{} กามาวจรญฺจ นกามาวจรญฺจ ธมฺมํ ปฏิจฺจ กามาวจโร จ นกามาวจโร จ ธมฺมา อุปฺปชฺชนฺติ เหตุปจฺจยา.\csromanspacer{} ปฏิสนฺธิกฺขเณ นกามาวจรํ เอกํ ขนฺธญฺจ วตฺถุญฺจ ปฏิจฺจ ตโย ขนฺธา ฯเปฯ เทฺว ขนฺเธ จ ฯเปฯ.\csromanspacer{} นกามาวจเร ขนฺเธ จ มหาภูเต จ ปฏิจฺจ กฏตฺตารูปํ. (3) (สํขิตฺตํ.)}

\csromanheader{2}{1. ปจฺจยานุโลม 2. สงฺขฺยาวาร}
\csromanheader{2}{\textbf{สุทฺธ}}
\proseitem{197}{\csromanfolio{128}เหตุยา นว, อารมฺมเณ จตฺตาริ, อธิปติยา ปญฺจ, อนนฺตเร จตฺตาริ, สมนนฺตเร จตฺตาริ, สหชาเต นว, อญฺญมญฺเญ ฉ, นิสฺสเย นว, อุปนิสฺสเย จตฺตาริ, ปุเรชาเต เทฺว, อาเสวเน เทฺว, กมฺเม นว, วิปาเก นว, อาหาเร นว, อินฺทฺริเย นว, ฌาเน นว, มคฺเค นว, สมฺปยุตฺเต จตฺตาริ, วิปฺปยุตฺเต นว, อตฺถิยา นว, นตฺถิยา จตฺตาริ, วิคเต จตฺตาริ, อวิคเต นว.}

\csromanheader{2}{2. ปจฺจยปจฺจนีย 1. วิภงฺควาร}
\csromanheader{2}{\textbf{นเหตุ}}
\proseitem{198}{กามาวจรํ ธมฺมํ ปฏิจฺจ กามาวจโร ธมฺโม อุปฺปชฺชติ นเหตุปจฺจยา.\csromanspacer{} อเหตุกํ กามาวจรํ เอกํ ขนฺธํ ปฏิจฺจ ตโย ขนฺธา จิตฺตสมุฏฺฐานญฺจ รูปํ ฯเปฯ เทฺว ขนฺเธ ฯเปฯ.\csromanspacer{} อเหตุกปฏิสนฺธิกฺขเณ ฯเปฯ. (ยาว อสญฺญสตฺตา.) วิจิกิจฺฉาสหคเต อุทฺธจฺจสหคเต ขนฺเธ ปฏิจฺจ วิจิกิจฺฉาสหคโต อุทฺธจฺจสหคโต โมโห. (1).\csromanspacer{} น-อารมฺมณปจฺจยา ตีณิ.}

\csromanheader{2}{\textbf{น-อธิปตฺยาทิ}}
\proseitem{199}{กามาวจรํ ธมฺมํ ปฏิจฺจ กามาวจโร ธมฺโม อุปฺปชฺชติ น-อธิปติปจฺจยา.\csromanspacer{} ตีณิ.}

\prose{นกามาวจรํ ธมฺมํ ปฏิจฺจ นกามาวจโร ธมฺโม อุปฺปชฺชติ น-อธิปติปจฺจยา.\csromanspacer{} นกามาวจเร ขนฺเธ ปฏิจฺจ นกามาวจราธิปติ.\csromanspacer{} วิปากํ นกามาวจรํ เอกํ ขนฺธํ ปฏิจฺจ ตโย ขนฺธา ฯเปฯ เทฺว ขนฺเธ ฯเปฯ.\csromanspacer{} ปฏิสนฺธิกฺขเณ ฯเปฯ.\csromanspacer{} นกามาวจรํ ธมฺมํ ปฏิจฺจ กามาวจโร ธมฺโม อุปฺปชฺชติ น-อธิปติปจฺจยา.\csromanspacer{} วิปาเก นกามาวจเร ขนฺเธ ปฏิจฺจ จิตฺตสมุฏฺฐานํ รูปํ.\csromanspacer{} ปฏิสนฺธิกฺขเณ ฯเปฯ.\csromanspacer{} นกามาวจรํ ธมฺมํ ปฏิจฺจ กามาวจโร จ นกามาวจโร จ ธมฺมา อุปฺปชฺชนฺติ น-อธิปติปจฺจยา.\csromanspacer{} วิปากํ นกามาวจรํ เอกํ ขนฺธํ ปฏิจฺจ ตโย ขนฺธา จิตฺตสมุฏฺฐานญฺจ รูปํ ฯเปฯ เทฺว ขนฺเธ ฯเปฯ.\csromanspacer{} ปฏิสนฺธิกฺขเณ ฯเปฯ. (3)}

\csromanheader{2}{93. กามาวจรทุก}
\prose{\csromanfolio{129}กามาวจรญฺจ นกามาวจรญฺจ ธมฺมํ ปฏิจฺจ กามาวจโร ธมฺโม อุปฺปชฺชติ น-อธิปติปจฺจยา.\csromanspacer{} วิปาเก นกามาวจเร ขนฺเธ จ มหาภูเต จ ปฏิจฺจ จิตฺตสมุฏฺฐานํ รูปํ.\csromanspacer{} ปฏิสนฺธิกฺขเณ ฯเปฯ. (อิตเร เทฺว ปากติกา.) น-อนนฺตรปจฺจยา ฯเปฯ.\csromanspacer{} นปุเรชาตปจฺจยา.\csromanspacer{} นปจฺฉาชาตปจฺจยา.}

\csromanheader{2}{\textbf{น-อาเสวน}}
\proseitem{200}{กามาวจรํ ธมฺมํ ปฏิจฺจ กามาวจโร ธมฺโม อุปฺปชฺชติ น-อาเสวนปจฺจยา.\csromanspacer{} ตีณิ.}

\prose{นกามาวจรํ ธมฺมํ ปฏิจฺจ นกามาวจโร ธมฺโม อุปฺปชฺชติ น-อาเสวนปจฺจยา.\csromanspacer{} วิปากํ นกามาวจรํ เอกํ ขนฺธํ ปฏิจฺจ ตโย ขนฺธา ฯเปฯ เทฺว ขนฺเธ ฯเปฯ.\csromanspacer{} ปฏิสนฺธิกฺขเณ ฯเปฯ.\csromanspacer{} นกามาวจรํ ธมฺมํ ปฏิจฺจ กามาวจโร ธมฺโม อุปฺปชฺชติ น-อาเสวนปจฺจยา.\csromanspacer{} นกามาวจเร ขนฺเธ ปฏิจฺจ จิตฺตสมุฏฺฐานํ รูปํ.\csromanspacer{} ปฏิสนฺธิกฺขเณ ฯเปฯ. (มูลํ กาตพฺพํ.) วิปากํ นกามาวจรํ เอกํ ขนฺธํ ปฏิจฺจ ตโย ขนฺธา จิตฺตสมุฏฺฐานญฺจ รูปํ ฯเปฯ เทฺว ขนฺเธ ฯเปฯ.\csromanspacer{} ปฏิสนฺธิกฺขเณ ฯเปฯ. (3) (อวเสสา ตีณิ ปากติกา.\csromanspacer{} สํขิตฺตํ.)}

\csromanheader{2}{2. ปจฺจยปจฺจนีย 2. สงฺขฺยาวาร}
\proseitem{201}{นเหตุยา เอกํ, น-อารมฺมเณ ตีณิ, น-อธิปติยา นว, น-อนนฺตเร ตีณิ, นสมนนฺตเร ตีณิ, น-อุปนิสฺสเย ตีณิ, นปุเรชาเต นว, นปจฺฉาชาเต นว, น-อาเสวเน นว, นกมฺเม เทฺว, นวิปาเก ปญฺจ, น-อาหาเร เอกํ, น-อินฺทฺริเย เอกํ, นฌาเน เอกํ, นมคฺเค เอกํ, นสมฺปยุตฺเต ตีณิ, นวิปฺปยุตฺเต เทฺว, โนนตฺถิยา ตีณิ, โนวิคเต ตีณิ. (อวเสสา คณนาปิ สหชาตวาโรปิ กาตพฺโพ.)}

\csromanheader{2}{93. กามาวจรทุก 3. ปจฺจยวาร}
\csromanheader{2}{1. ปจฺจยานุโลมาทิ}
\proseitem{202}{กามาวจรํ ธมฺมํ ปจฺจยา กามาวจโร ธมฺโม อุปฺปชฺชติ เหตุปจฺจยา.\csromanspacer{} กามาวจรํ เอกํ ขนฺธํ ปจฺจยา ตโย ขนฺธา จิตฺตสมุฏฺฐานญฺจ \csromanfolio{130}รูปํ ฯเปฯ เทฺว ขนฺเธ ฯเปฯ.\csromanspacer{} ปฏิสนฺธิกฺขเณ ฯเปฯ. (ยาว อชฺฌตฺติกา มหาภูตา.) วตฺถุํ ปจฺจยา กามาวจรา ขนฺธา.\csromanspacer{}\csromanspacer{}กามาวจรํ ธมฺมํ ปจฺจยา นกามาวจโร ธมฺโม อุปฺปชฺชติ เหตุปจฺจยา.\csromanspacer{} วตฺถุํ ปจฺจยา นกามาวจรา ขนฺธา.\csromanspacer{} ปฏิสนฺธิกฺขเณ ฯเปฯ.\csromanspacer{} กามาวจรํ ธมฺมํ ปจฺจยา กามาวจโร จ นกามาวจโร จ ธมฺมา อุปฺปชฺชนฺติ เหตุปจฺจยา.\csromanspacer{} วตฺถุํ ปจฺจยา นกามาวจรา ขนฺธา.\csromanspacer{} มหาภูเต ปจฺจยา จิตฺตสมุฏฺฐานํ รูปํ.\csromanspacer{} ปฏิสนฺธิกฺขเณ ฯเปฯ. (3)}

\prose{นกามาวจรํ ธมฺมํ ปจฺจยา นกามาวจโร ธมฺโม อุปฺปชฺชติ เหตุปจฺจยา.\csromanspacer{} ตีณิ. (ปฏิจฺจสทิสา.)}

\prose{กามาวจรญฺจ นกามาวจรญฺจ ธมฺมํ ปจฺจยา กามาวจโร ธมฺโม อุปฺปชฺชติ เหตุปจฺจยา. (ปฏิจฺจสทิสา.).\csromanspacer{} กามาวจรญฺจ นกามาวจรญฺจ ธมฺมํ ปจฺจยา นกามาวจโร ธมฺโม อุปฺปชฺชติ เหตุปจฺจยา.\csromanspacer{} นกามาวจรํ เอกํ ขนฺธญฺจ วตฺถุญฺจ ปจฺจยา ตโย ขนฺธา ฯเปฯ เทฺว ขนฺเธ จ ฯเปฯ.\csromanspacer{} ปฏิสนฺธิกฺขเณ ฯเปฯ.\csromanspacer{} กามาวจรญฺจ นกามาวจรญฺจ ธมฺมํ ปจฺจยา กามาวจโร จ นกามาวจโร จ ธมฺมา อุปฺปชฺชนฺติ เหตุปจฺจยา.\csromanspacer{} นกามาวจรํ เอกํ ขนฺธญฺจ วตฺถุญฺจ ปจฺจยา ตโย ขนฺธา ฯเปฯ เทฺว ขนฺเธ จ ฯเปฯ.\csromanspacer{} นกามาวจเร ขนฺเธ จ มหาภูเต จ ปจฺจยา จิตฺตสมุฏฺฐานํ รูปํ.\csromanspacer{} ปฏิสนฺธิกฺขเณ ฯเปฯ. (3) (สํขิตฺตํ.)}

\prose{เหตุยา นว, อารมฺมเณ จตฺตาริ, อธิปติยา นว, อนนฺตเร จตฺตาริ, สมนนฺตเร จตฺตาริ, สหชาเต นว, อญฺญมญฺเญ ฉ, นิสฺสเย นว, อุปนิสฺสเย จตฺตาริ, ปุเรชาเต จตฺตาริ, อาเสวเน จตฺตาริ, กมฺเม นว ฯเปฯ อวิคเต นว.}

\vspace{\tipitakaparskip}
\csromancenter{อนุโลมํ.}
\csromanheader{2}{\textbf{นเหตุ}}
\proseitem{203}{กามาวจรํ ธมฺมํ ปจฺจยา กามาวจโร ธมฺโม อุปฺปชฺชติ นเหตุปจฺจยา.\csromanspacer{} อเหตุกํ กามาวจรํ เอกํ ขนฺธํ ปจฺจยา ฯเปฯ.\csromanspacer{} อเหตุกปฏิสนฺธิกฺขเณ ฯเปฯ. (ยาว อสญฺญสตฺตา.) จกฺขายตนํ ปจฺจยา จกฺขุวิญฺญาณํ ฯเปฯ.\csromanspacer{} กายายตนํ ปจฺจยา กายวิญฺญาณํ.\csromanspacer{} วตฺถุํ ปจฺจยา อเหตุกา กามาวจรา ขนฺธา.\csromanspacer{} วิจิกิจฺฉาสหคเต อุทฺธจฺจสหคเต ขนฺเธ จ วตฺถุญฺจ ปจฺจยา วิจิกิจฺฉาสหคโต อุทฺธจฺจสหคโต โมโห. (1) น-อารมฺมณปจฺจยา.\csromanspacer{} ตีณิ.}

\csromanheader{2}{93. กามาวจรทุก \textbf{น-อธิปตฺยาทิ}}
\proseitem{204}{\csromanfolio{131}กามาวจรํ ธมฺมํ ปจฺจยา กามาวจโร ธมฺโม อุปฺปชฺชติ น-อธิปติปจฺจยา.\csromanspacer{} เอกํ. (ยาว อสญฺญสตฺตา.).\csromanspacer{} กามาวจรํ ธมฺมํ ปจฺจยา นกามาวจโร ธมฺโม อุปฺปชฺชติ น-อธิปติปจฺจยา.\csromanspacer{} วตฺถุํ ปจฺจยา นกามาวจราธิปติ.\csromanspacer{} วตฺถุํ ปจฺจยา วิปากา นกามาวจรา ขนฺธา.\csromanspacer{} ปฏิสนฺธิกฺขเณ ฯเปฯ.\csromanspacer{} กามาวจรํ ธมฺมํ ปจฺจยา กามาวจโร จ นกามาวจโร จ ธมฺมา อุปฺปชฺชนฺติ น-อธิปติปจฺจยา.\csromanspacer{} วตฺถุํ ปจฺจยา วิปากา นกามาวจรา ขนฺธา.\csromanspacer{} มหาภูเต ปจฺจยา จิตฺตสมุฏฺฐานํ รูปํ.\csromanspacer{} ปฏิสนฺธิกฺขเณ ฯเปฯ. (3)}

\prose{นกามาวจรํ ธมฺมํ ปจฺจยา นกามาวจโร ธมฺโม อุปฺปชฺชติ น-อธิปติปจฺจยา.\csromanspacer{} ตีณิ. (ปฏิจฺจสทิสา.)}

\prose{กามาวจรญฺจ นกามาวจรญฺจ ธมฺมํ ปจฺจยา กามาวจโร ธมฺโม อุปฺปชฺชติ น-อธิปติปจฺจยา. (ปฏิจฺจสทิสา.\csromanspacer{} มูลํ กาตพฺพํ.) นกามาวจเร ขนฺเธ จ วตฺถุญฺจ ปจฺจยา นกามาวจราธิปติ.\csromanspacer{} วิปากํ นกามาวจรํ เอกํ ขนฺธญฺจ วตฺถุญฺจ ปจฺจยา ตโย ขนฺธา ฯเปฯ เทฺว ขนฺเธ จ ฯเปฯ.\csromanspacer{} ปฏิสนฺธิกฺขเณ ฯเปฯ. (มูลํ กาตพฺพํ.) วิปากํ นกามาวจรํ เอกํ ขนฺธญฺจ วตฺถุญฺจ ปจฺจยา ตโย ขนฺธา ฯเปฯ เทฺว ขนฺเธ จ ฯเปฯ.\csromanspacer{} วิปาเก นกามาวจเร ขนฺเธ จ มหาภูเต จ ปจฺจยา จิตฺตสมุฏฺฐานํ รูปํ.\csromanspacer{} ปฏิสนฺธิกฺขเณ ฯเปฯ. (3) น-อุปนิสฺสยปจฺจยา.\csromanspacer{} ตีณิ.\csromanspacer{} น-อาเสวนปจฺจยา. (\textbf{สุทฺธ}เก อรูปมิสฺสเก จ “วิปากนฺ”ติ นิยาเมตพฺพํ.\csromanspacer{} รูปมิสฺสเก นตฺถิ.\csromanspacer{} สํขิตฺตํ.)}

\csromanheader{2}{\textbf{สุทฺธ}}
\prose{นเหตุยา เอกํ, น-อารมฺมเณ ตีณิ, น-อธิปติยา นว, น-อนนฺตเร ตีณิ, นสมนนฺตเร ตีณิ, น-อุปนิสฺสเย ตีณิ, นปุเรชาเต นว, นปจฺฉาชาเต นว, น-อาเสวเน นว, นกมฺเม จตฺตาริ, นวิปาเก นว, น-อาหาเร เอกํ, น-อินฺทฺริเย เอกํ, นฌาเน เอกํ, นมคฺเค เอกํ, นสมฺปยุตฺเต ตีณิ, นวิปฺปยุตฺเต เทฺว, โนนตฺถิยา ตีณิ, โนวิคเต ตีณิ.}

\vspace{\tipitakaparskip}
\csromancenter{ปจฺจนียํ.}
\csromancenter{(เอวํ อิตเร เทฺว คณนาปิ นิสฺสยวาโรปิ กาตพฺโพ.)}
\csromanheader{2}{93. กามาวจรทุก 5. สํสฏฺฐวาร}
\csromanheader{2}{1. ปจฺจยานุโลมาทิ}
\proseitem{205}{\csromanfolio{132}กามาวจรํ ธมฺมํ สํสฏฺโฐ กามาวจโร ธมฺโม อุปฺปชฺชติ เหตุปจฺจยา.\csromanspacer{} กามาวจรํ เอกํ ขนฺธํ สํสฏฺฐา ตโย ขนฺธา ฯเปฯ เทฺว ขนฺเธ ฯเปฯ.\csromanspacer{} ปฏิสนฺธิกฺขเณ ฯเปฯ. (1)}

\prose{นกามาวจรํ ธมฺมํ สํสฏฺโฐ นกามาวจโร ธมฺโม อุปฺปชฺชติ เหตุปจฺจยา.\csromanspacer{} นกามาวจรํ เอกํ ขนฺธํ สํสฏฺฐา ตโย ขนฺธา ฯเปฯ เทฺว ขนฺเธ ฯเปฯ.\csromanspacer{} ปฏิสนฺธิกฺขเณ ฯเปฯ. (1)}

\vspace{\tipitakaparskip}
\csromancenter{\textbf{เหตุ}ยา เทฺว, อารมฺมเณ เทฺว, อธิปติยา เทฺว, (สพฺพตฺถ เทฺว,) อวิคเต เทฺว.\csromanspacer{} อนุโลมํ.}
\prose{กามาวจรํ ธมฺมํ สํสฏฺโฐ กามาวจโร ธมฺโม อุปฺปชฺชติ นเหตุปจฺจยา.\csromanspacer{} อเหตุกํ กามาวจรํ เอกํ ขนฺธํ สํสฏฺฐา ตโย ขนฺธา ฯเปฯ เทฺว ขนฺเธ ฯเปฯ.\csromanspacer{} อเหตุกปฏิสนฺธิกฺขเณ ฯเปฯ.\csromanspacer{} วิจิกิจฺฉาสหคเต อุทฺธจฺจสหคเต ขนฺเธ สํสฏฺโฐ วิจิกิจฺฉาสหคโต อุทฺธจฺจสหคโต โมโห.}

\prose{นเหตุยา เอกํ, น-อธิปติยา เทฺว, นปุเรชาเต เทฺว, นปจฺฉาชาเต เทฺว, น-อาเสวเน เทฺว, นกมฺเม เทฺว, นวิปาเก เทฺว, นฌาเน เอกํ, นมคฺเค เอกํ, นวิปฺปยุตฺเต เทฺว.}

\vspace{\tipitakaparskip}
\csromancenter{ปจฺจนียํ.}
\csromancenter{(เอวํ อิตเร เทฺว คณนาปิ สมฺปยุตฺตวาโรปิ กาตพฺโพ.)}
\csromanheader{2}{93. กามาวจรทุก 7. ปญฺหาวาร}
\proseitem{206}{กามาวจโร ธมฺโม กามาวจรสฺส ธมฺมสฺส เหตุปจฺจเยน ปจฺจโย.\csromanspacer{} กามาวจรา เหตู สมฺปยุตฺตกานํ ขนฺธานํ จิตฺตสมุฏฺฐานานญฺจ รูปานํ เหตุปจฺจเยน ปจฺจโย.\csromanspacer{} ปฏิสนฺธิกฺขเณ ฯเปฯ. (1)}

\csromanheader{2}{93. กามาวจรทุก}
\prose{\csromanfolio{133}นกามาวจโร ธมฺโม นกามาวจรสฺส ธมฺมสฺส เหตุปจฺจเยน ปจฺจโย.\csromanspacer{} นกามาวจรา เหตู สมฺปยุตฺตกานํ ขนฺธานํ เหตุปจฺจเยน ปจฺจโย.\csromanspacer{} ปฏิสนฺธิกฺขเณ ฯเปฯ.\csromanspacer{} นกามาวจโร ธมฺโม กามาวจรสฺส ธมฺมสฺส เหตุปจฺจเยน ปจฺจโย.\csromanspacer{} นกามาวจรา เหตู จิตฺตสมุฏฺฐานานํ รูปานํ เหตุปจฺจเยน ปจฺจโย.\csromanspacer{} ปฏิสนฺธิกฺขเณ ฯเปฯ. (มูลํ กาตพฺพํ.) นกามาวจรา เหตู สมฺปยุตฺตกานํ ขนฺธานํ จิตฺตสมุฏฺฐานานญฺจ รูปานํ เหตุปจฺจเยน ปจฺจโย.\csromanspacer{} ปฏิสนฺธิกฺขเณ ฯเปฯ. (3)}

\csromanheader{2}{\textbf{อารมฺมณ}}
\proseitem{207}{กามาวจโร ธมฺโม กามาวจรสฺส ธมฺมสฺส อารมฺมณปจฺจเยน ปจฺจโย.\csromanspacer{} ทานํ ทตฺวา, สีลํ ฯเปฯ อุโปสถกมฺมํ ฯเปฯ.\csromanspacer{} ปุพฺเพ สุจิณฺณานิ ปจฺจเวกฺขติ, อสฺสาเทติ อภินนฺทติ, ตํ อารพฺภ ราโค ฯเปฯ โทมนสฺสํ อุปฺปชฺชติ.\csromanspacer{} อาริยา โคตฺรภุํ ปจฺจเวกฺขนฺติ, โวทานํ ปจฺจเวกฺขนฺติ.\csromanspacer{} ปหีเน กิเลเส ฯเปฯ.\csromanspacer{} วิกฺขมฺภิเต กิเลเส ฯเปฯ.\csromanspacer{} ปุพฺเพ สุจิณฺณานิ ฯเปฯ.\csromanspacer{} จกฺขุํ ฯเปฯ.\csromanspacer{} วตฺถุํ กามาวจเร ขนฺเธ อนิจฺจโต ฯเปฯ โทมนสฺสํ อุปฺปชฺชติ.\csromanspacer{} รูปายตนํ จกฺขุวิญฺญาณสฺส ฯเปฯ.\csromanspacer{} โผฏฺฐพฺพายตนํ กายวิญฺญาณสฺส อารมฺมณปจฺจเยน ปจฺจโย. (1)}

\prose{กามาวจโร ธมฺโม นกามาวจรสฺส ธมฺมสฺส อารมฺมณปจฺจเยน ปจฺจโย.\csromanspacer{} ทิพฺเพน จกฺขุนา รูปํ ปสฺสติ.\csromanspacer{} ทิพฺพาย โสตธาตุยา สทฺทํ สุณาติ.\csromanspacer{} เจโตปริยญาเณน กามาวจรจิตฺตสมงฺคิสฺส จิตฺตํ ชานาติ.\csromanspacer{} กามาวจรา ขนฺธา อิทฺธิวิธญาณสฺส, เจโตปริยญาณสฺส, ปุพฺเพนิวาสานุสฺสติญาณสฺส, ยถากมฺมูปคญาณสฺส, อนาคตํสญาณสฺส อารมฺมณปจฺจเยน ปจฺจโย. (2)}

\proseitem{208}{นกามาวจโร ธมฺโม นกามาวจรสฺส ธมฺมสฺส อารมฺมณปจฺจเยน ปจฺจโย.\csromanspacer{} นิพฺพานํ มคฺคสฺส, ผลสฺส อารมฺมณปจฺจเยน ปจฺจโย.\csromanspacer{} เจโตปริยญาเณน นกามาวจรจิตฺตสมงฺคิสฺส จิตฺตํ ชานาติ.\csromanspacer{} อากาสานญฺจายตนํ วิญฺญาณญฺจายตนสฺส ฯเปฯ.\csromanspacer{} อากิญฺจญฺญายตนํ เนวสญฺญานาสญฺญายตนสฺส ฯเปฯ.\csromanspacer{} นกามาวจรา ขนฺธา อิทฺธิวิธญาณสฺส, เจโตปริยญาณสฺส, ปุพฺเพนิวาสานุสฺสติญาณสฺส, ยถากมฺมูปคญาณสฺส, อนาคตํสญาณสฺส อารมฺมณปจฺจเยน ปจฺจโย. (1)}

\prose{\csromanfolio{134}นกามาวจโร ธมฺโม กามาวจรสฺส ธมฺมสฺส อารมฺมณปจฺจเยน ปจฺจโย.\csromanspacer{} ฌานา วุฏฺฐหิตฺวา ฌานํ ปจฺจเวกฺขติ, อสฺสาเทติ อภินนฺทติ, ตํ อารพฺภ ราโค อุปฺปชฺชติ.\csromanspacer{} ทิฏฺฐิ อุปฺปชฺชติ ฯเปฯ.\csromanspacer{} อริยา มคฺคา วุฏฺฐหิตฺวา มคฺคํ ปจฺจเวกฺขนฺติ.\csromanspacer{} ผลํ ฯเปฯ.\csromanspacer{} นิพฺพานํ ปจฺจเวกฺขนฺติ.\csromanspacer{} นิพฺพานํ โคตฺรภุสฺส, โวทานสฺส, อาวชฺชนาย อารมฺมณปจฺจเยน ปจฺจโย.\csromanspacer{} อากาสานญฺจายตนํ ปจฺจเวกฺขติ.\csromanspacer{} วิญฺญาณญฺจายตนํ ปจฺจเวกฺขติ.\csromanspacer{} อากิญฺจญฺญายตนํ ปจฺจเวกฺขติ.\csromanspacer{} เนวสญฺญานาสญฺญายตนํ ปจฺจเวกฺขติ.\csromanspacer{} ทิพฺพํ จกฺขุํ ปจฺจเวกฺขติ.\csromanspacer{} ทิพฺพํ โสตธาตุํ ปจฺจเวกฺขติ.\csromanspacer{} อิทฺธิวิธญาณํ ปจฺจเวกฺขติ.\csromanspacer{} เจโตปริยญาณํ ฯเปฯ.\csromanspacer{} ปุพฺเพ นิวาสานุสฺสติญาณํ ฯเปฯ.\csromanspacer{} ยถากมฺมูปคญาณํ ฯเปฯ.\csromanspacer{} อนาคตํสญาณํ ปจฺจเวกฺขติ.\csromanspacer{} นกามาวจเร ขนฺเธ อนิจฺจโต ฯเปฯ โทมนสฺสํ อุปฺปชฺชติ. (2)}

\csromanheader{2}{\textbf{อธิปติ}}
\proseitem{209}{กามาวจโร ธมฺโม กามาวจรสฺส ธมฺมสฺส อธิปติปจฺจเยน ปจฺจโย.\csromanspacer{} อารมฺมณาธิปติ, สหชาตาธิปติ.\csromanspacer{}\csromanspacer{}\textbf{อารมฺมณาธิปติ\csromandash{}}ทานํ ฯเปฯ สีลํ ฯเปฯ อุโปสถกมฺมํ ฯเปฯ.\csromanspacer{} ปุพฺเพ สุจิณฺณานิ ครุํ กตฺวา ปจฺจเวกฺขติ, อสฺสาเทติ อภินนฺทติ, ตํ ครุํ กตฺวา ราโค อุปฺปชฺชติ.\csromanspacer{} ทิฏฺฐิ อุปฺปชฺชติ.\csromanspacer{} เสกฺขา โคตฺรภุํ ครุํ กตฺวา ปจฺจเวกฺขนฺติ.\csromanspacer{} โวทานํ ครุํ กตฺวา ปจฺจเวกฺขนฺติ.\csromanspacer{} จกฺขุํ ฯเปฯ.\csromanspacer{} วตฺถุํ กามาวจเร ขนฺเธ ครุํ กตฺวา อสฺสาเทติ อภินนฺทติ, ตํ ครุํ กตฺวา ราโค อุปฺปชฺชติ.\csromanspacer{} ทิฏฺฐิ อุปฺปชฺชติ.\csromanspacer{} \textbf{สหชาตาธิปติ\csromandash{}}กามาวจราธิปติ สมฺปยุตฺตกานํ ขนฺธานํ จิตฺตสมุฏฺฐานานญฺจ รูปานํ อธิปติปจฺจเยน ปจฺจโย. (1)}

\proseitem{210}{นกามาวจโร ธมฺโม นกามาวจรสฺส ธมฺมสฺส อธิปติปจฺจเยน ปจฺจโย.\csromanspacer{} อารมฺมณาธิปติ, สหชาตาธิปติ.\csromanspacer{}\csromanspacer{}\textbf{อารมฺมณาธิปติ\csromandash{}}นิพฺพานํ มคฺคสฺส, ผลสฺส อธิปติปจฺจเยน ปจฺจโย.\csromanspacer{} \textbf{สหชาตาธิปติ\csromandash{}}นกามาวจราธิปติ สมฺปยุตฺตกานํ ขนฺธานํ อธิปติปจฺจเยน ปจฺจโย.\csromanspacer{}\csromanspacer{}นกามาวจโร ธมฺโม กามาวจรสฺส ธมฺมสฺส อธิปติปจฺจเยน ปจฺจโย.\csromanspacer{} อารมฺมณาธิปติ, สหชาตาธิปติ.\csromanspacer{}\csromanspacer{}\textbf{อารมฺมณาธิปติ\csromandash{}}ฌานา วุฏฺฐหิตฺวา ฌานํ ครุํ กตฺวา ปจฺจเวกฺขติ, อสฺสาเทติ อภินนฺทติ, ตํ ครุํ กตฺวา ราโค อุปฺปชฺชติ.\csromanspacer{} ทิฏฺฐิ อุปฺปชฺชติ.\csromanspacer{} อริยา มคฺคา วุฏฺฐหิตฺวา มคฺคํ ครุํ กตฺวา ปจฺจเวกฺขนฺติ.\csromanspacer{} ผลํ ฯเปฯ.\csromanspacer{} นิพฺพานํ ครุํ กตฺวา ปจฺจเวกฺขนฺติ.\csromanspacer{} นิพฺพานํ โคตฺรภุสฺส, โวทานสฺส อธิปติปจฺจเยน ปจฺจโย.\csromanspacer{} อากาสานญฺจายตนํ ครุํ กตฺวา}

\vspace{\tipitakaparskip}
\csromancenter{กามาวจรทุก ปจฺจเวกฺขติ.\csromanspacer{} วิญฺญาณญฺจายตนํ ฯเปฯ.\csromanspacer{} อากิญฺจญฺญายตนํ ฯเปฯ.\csromanspacer{} เนวสญฺญานาสญฺญายตนํ ครุํ กตฺวา ปจฺจเวกฺขติ.\csromanspacer{} ทิพฺพํ จกฺขุํ ฯเปฯ.\csromanspacer{} ทิพฺพํ โสตธาตุํ ฯเปฯ.\csromanspacer{} อิทฺธิวิธญาณํ ฯเปฯ.\csromanspacer{} อนาคตํสญาณํ ครุํ กตฺวา ปจฺจเวกฺขติ.\csromanspacer{} นกามาวจเร ขนฺเธ ครุํ กตฺวา อสฺสาเทติ อภินนฺทติ, ตํ ครุํ กตฺวา ราโค อุปฺปชฺชติ.\csromanspacer{} ทิฏฺฐิ อุปฺปชฺชติ.\csromanspacer{} \textbf{สหชาตาธิปติ\csromandash{}}นกามาวจราธิปติ จิตฺตสมุฏฺฐานานํ รูปานํ อธิปติปจฺจเยน ปจฺจโย.\csromanspacer{}\csromanspacer{}นกามาวจโร ธมฺโม กามาวจรสฺส จ นกามาวจรสฺส จ ธมฺมสฺส อธิปติปจฺจเยน ปจฺจโย.\csromanspacer{} \textbf{สหชาตาธิปติ\csromandash{}}นกามาวจราธิปติ สมฺปยุตฺตกานํ ขนฺธานํ จิตฺตสมุฏฺฐานานญฺจ รูปานํ อธิปติปจฺจเยน ปจฺจโย. (3)}
\csromanheader{2}{\textbf{อนนฺตราทิ}}
\proseitem{211}{\csromanfolio{135}กามาวจโร ธมฺโม กามาวจรสฺส ธมฺมสฺส อนนฺตรปจฺจเยน ปจฺจโย.\csromanspacer{} ปุริมา ปุริมา กามาวจรา ขนฺธา ปจฺฉิมานํ ปจฺฉิมานํ กามาวจรานํ ขนฺธานํ อนนฺตรปจฺจเยน ปจฺจโย.\csromanspacer{} อนุโลมํ โคตฺรภุสฺส.\csromanspacer{} อนุโลมํ โวทานสฺส.\csromanspacer{} อาวชฺชนา กามาวจรานํ ขนฺธานํ อนนฺตรปจฺจเยน ปจฺจโย. (1)}

\prose{กามาวจโร ธมฺโม นกามาวจรสฺส ธมฺมสฺส อนนฺตรปจฺจเยน ปจฺจโย.\csromanspacer{} กามาวจรํ จุติจิตฺตํ นกามาวจรสฺส อุปปตฺติจิตฺตสฺส.\csromanspacer{} อาวชฺชนา นกามาวจรานํ ขนฺธานํ อนนฺตรปจฺจเยน ปจฺจโย.\csromanspacer{} กามาวจรา ขนฺธา นกามาวจรสฺส วุฏฺฐานสฺส อนนฺตรปจฺจเยน ปจฺจโย.\csromanspacer{} ปฐมสฺส ฌานสฺส ปริกมฺมํ ปฐมสฺส ฌานสฺส อนนฺตรปจฺจเยน ปจฺจโย ฯเปฯ.\csromanspacer{} จตุตฺถสฺส ฌานสฺส ฯเปฯ.\csromanspacer{} เนวสญฺญานาสญฺญายตนสฺส ปริกมฺมํ ฯเปฯ.\csromanspacer{} ทิพฺพสฺส จกฺขุสฺส ฯเปฯ.\csromanspacer{} ทิพฺพาย โสตธาตุยา ฯเปฯ.\csromanspacer{} อิธิวิธญาณสฺส ฯเปฯ.\csromanspacer{} เจโตปริยญาณสฺส ฯเปฯ.\csromanspacer{} ปุพฺเพนิวาสานุสฺสติญาณสฺส ฯเปฯ.\csromanspacer{} ยถากมฺมูปคญาณสฺส ฯเปฯ.\csromanspacer{} อนาคตํสญาณสฺส ปริกมฺมํ อนาคตํสญาณสฺส อนนฺตรปจฺจเยน ปจฺจโย.\csromanspacer{} โคตฺรภุ มคฺคสฺส.\csromanspacer{} โวทานํ มคฺคสฺส.\csromanspacer{} อนุโลมํ ผลสมาปตฺติยา อนนฺตรปจฺจเยน ปจฺจโย. (2)}

\proseitem{212}{นกามาวจโร ธมฺโม นกามาวจรสฺส ธมฺมสฺส อนนฺตรปจฺจเยน ปจฺจโย.\csromanspacer{} ปุริมา ปุริมา นกามาวจรา ขนฺธา ปจฺฉิมานํ ปจฺฉิมานํ นกามาวจรานํ ขนฺธานํ อนนฺตรปจฺจเยน ปจฺจโย.\csromanspacer{} มคฺโค ผลสฺส.\csromanspacer{} ผลํ ผลสฺส.\csromanspacer{} นิโรธา วุฏฺฐหนฺตสฺส เนวสญฺญานาสญฺญายตนํ ผลสมาปตฺติยา อนนฺตรปจฺจเยน ปจฺจโย.\csromanspacer{}\csromanspacer{}นกามาวจโร ธมฺโม กามาวจรสฺส ธมฺมสฺส อนนฺตรปจฺจเยน \csromanfolio{136}ปจฺจโย.\csromanspacer{} นกามาวจรํ จุติจิตฺตํ กามาวจรสฺส อุปปตฺติจิตฺตสฺส.\csromanspacer{} นกามาวจรํ ภวงฺคํ อาวชฺชนาย.\csromanspacer{} นกามาวจรา ขนฺธา กามาวจรสฺส วุฏฺฐานสฺส อนนฺตรปจฺจเยน ปจฺจโย. (2)}

\prose{สมนนฺตรปจฺจเยน ปจฺจโย.\csromanspacer{} สหชาตปจฺจเยน ปจฺจโย.\csromanspacer{} สตฺต.\csromanspacer{} อญฺญมญฺญปจฺจเยน ปจฺจโย.\csromanspacer{} ฉ.\csromanspacer{} นิสฺสยปจฺจเยน ปจฺจโย.\csromanspacer{} สตฺต.}

\csromanheader{2}{\textbf{อุปนิสฺสย}}
\proseitem{213}{กามาวจโร ธมฺโม กามาวจรสฺส ธมฺมสฺส อุปนิสฺสยปจฺจเยน ปจฺจโย.\csromanspacer{} \textbf{อารมฺมณูปนิสฺสโย, อนนฺตรูปนิสฺสโย, ปกตูปนิสฺสโย ฯเปฯ.}\csromanspacer{} ปกตูปนิสฺสโย\csromandash{}กามาวจรํ สทฺธํ อุปนิสฺสาย ทานํ เทติ ฯเปฯ สีลํ ฯเปฯ อุโปสถกมฺมํ ฯเปฯ.\csromanspacer{} วิปสฺสนํ อุปฺปาเทติ.\csromanspacer{} มานํ ชปฺเปติ.\csromanspacer{} ทิฏฺฐิํ คณฺหาติ.\csromanspacer{} กามาวจรํ สีลํ ฯเปฯ ปญฺญํ.\csromanspacer{} ราคํ ฯเปฯ.\csromanspacer{} ปตฺถนํ.\csromanspacer{} กายิกํ สุขํ.\csromanspacer{} กายิกํ ทุกฺขํ.\csromanspacer{} อุตุํ.\csromanspacer{} โภชนํ.\csromanspacer{} เสนาสนํ อุปนิสฺสาย ทานํ เทติ ฯเปฯ.\csromanspacer{} วิปสฺสนํ อุปฺปาเทติ.\csromanspacer{} ปาณํ หนติ ฯเปฯ สํฆํ ภินฺทติ.\csromanspacer{} กามาวจรา สทฺธา ฯเปฯ.\csromanspacer{} เสนาสนํ กามาวจราย สทฺธาย ฯเปฯ ปตฺถนาย.\csromanspacer{} กายิกสฺส สุขสฺส.\csromanspacer{} กายิกสฺส ทุกฺขสฺส อุปนิสฺสยปจฺจเยน ปจฺจโย. (1)}

\prose{กามาวจโร ธมฺโม นกามาวจรสฺส ธมฺมสฺส อุปนิสฺสยปจฺจเยน ปจฺจโย.\csromanspacer{} \textbf{อนนฺตรูปนิสฺสโย, ปกตูปนิสฺสโย ฯเปฯ.}\csromanspacer{} ปกตูปนิสฺสโย\csromandash{}กามาวจรํ สทฺธํ อุปนิสฺสาย นกามาวจรํ ฌานํ อุปฺปาเทติ.\csromanspacer{} มคฺคํ ฯเปฯ.\csromanspacer{} อภิญฺญํ ฯเปฯ.\csromanspacer{} สมาปตฺติํ อุปฺปาเทติ.\csromanspacer{} กามาวจรํ สีลํ ฯเปฯ.\csromanspacer{} เสนาสนํ อุปนิสฺสาย นกามาวจรํ ฌานํ อุปฺปาเทติ.\csromanspacer{} มคฺคํ ฯเปฯ.\csromanspacer{} อภิญฺญํ ฯเปฯ.\csromanspacer{} สมาปตฺติํ อุปฺปาเทติ.\csromanspacer{} กามาวจรา สทฺธา ฯเปฯ.\csromanspacer{} เสนาสนํ นกามาวจราย สทฺธาย ฯเปฯ ปญฺญาย.\csromanspacer{} มคฺคสฺส, ผลสมาปตฺติยา อุปนิสฺสยปจฺจเยน ปจฺจโย.\csromanspacer{} ปฐมสฺส ฌานสฺส ปริกมฺมํ ปฐมสฺส ฌานสฺส ฯเปฯ.\csromanspacer{} จตุตฺถสฺส ฌานสฺส ฯเปฯ.\csromanspacer{} อากาสานญฺจายตนสฺส ฯเปฯ.\csromanspacer{} ปฐมสฺส มคฺคสฺส ฯเปฯ.\csromanspacer{} จตุตฺถสฺส มคฺคสฺส ปริกมฺมํ จตุตฺถสฺส มคฺคสฺส อุปนิสฺสยปจฺจเยน ปจฺจโย. (2)}

\proseitem{214}{นกามาวจโร ธมฺโม นกามาวจรสฺส ธมฺมสฺส อุปนิสฺสยปจฺจเยน ปจฺจโย.\csromanspacer{} \textbf{อารมฺมณูปนิสฺสโย, อนนฺตรูปนิสฺสโย, ปกตูปนิสฺสโย ฯเปฯ.}\csromanspacer{} ปกตูปนิสฺสโย\csromandash{}นกามาวจรํ สทฺธํ อุปนิสฺสาย ฌานํ อุปฺปาเทติ.\csromanspacer{} มคฺคํ ฯเปฯ.\csromanspacer{} อภิญฺญํ ฯเปฯ.\csromanspacer{} สมาปตฺติํ อุปฺปาเทติ.\csromanspacer{} นกามาวจรํ}

\vspace{\tipitakaparskip}
\csromancenter{กามาวจรทุก สีลํ ฯเปฯ ปญฺญํ อุปนิสฺสาย นกามาวจรํ ฌานํ อุปฺปาเทติ ฯเปฯ.\csromanspacer{} มคฺคํ ฯเปฯ.\csromanspacer{} อภิญฺญํ ฯเปฯ.\csromanspacer{} สมาปตฺติํ อุปฺปาเทติ.\csromanspacer{} นกามาวจรา สทฺธา ฯเปฯ ปญฺญา นกามาวจราย สทฺธาย ฯเปฯ ปญฺญาย.\csromanspacer{} มคฺคสฺส, ผลสมาปตฺติยา อุปนิสฺสยปจฺจเยน ปจฺจโย.\csromanspacer{} ปฐมํ ฌานํ ทุติยสฺส ฌานสฺส อุปนิสฺสยปจฺจเยน ปจฺจโย ฯเปฯ.\csromanspacer{} ตติยํ ฌานํ ฯเปฯ.\csromanspacer{} จตุตฺถํ ฌานํ ฯเปฯ.\csromanspacer{} อากาสานญฺจายตนํ วิญฺญาณญฺจายตนสฺส ฯเปฯ.\csromanspacer{} อากิญฺจญฺญายตนํ เนวสญฺญานาสญฺญายตนสฺส อุปนิสฺสยปจฺจเยน ปจฺจโย.\csromanspacer{} ปฐโม มคฺโค ทุติยสฺส มคฺคสฺส ฯเปฯ.\csromanspacer{} ทุติโย มคฺโค ตติยสฺส มคฺคสฺส ฯเปฯ.\csromanspacer{} ตติโย มคฺโค จตุตฺถสฺส มคฺคสฺส อุปนิสฺสยปจฺจเยน ปจฺจโย.\csromanspacer{} มคฺโค ผลสมาปตฺติยา อุปนิสฺสยปจฺจเยน ปจฺจโย. (1)}
\prose{\csromanfolio{137}นกามาวจรา ธมฺโม กามาวจรสฺส ธมฺมสฺส อุปนิสฺสยปจฺจเยน ปจฺจโย.\csromanspacer{} \textbf{อารมฺมณูปนิสฺสโย, อนนฺตรูปนิสฺสโย,} \textbf{ปกตูปนิสฺสโย ฯเปฯ.}\csromanspacer{} ปกตูปนิสฺสโย\csromandash{}นกามาวจรํ สทฺธํ อุปนิสฺสาย ทานํ เทติ.\csromanspacer{} สีลํ ฯเปฯ อุโปสถกมฺมํ ฯเปฯ.\csromanspacer{} วิปสฺสนํ อุปฺปาเทติ.\csromanspacer{} มานํ ชปฺเปติ.\csromanspacer{} ทิฏฺฐิํ คณฺหาติ.\csromanspacer{} นกามาวจรํ สีลํ ฯเปฯ ปญฺญํ อุปนิสฺสาย ทานํ เทติ.\csromanspacer{} สีลํ ฯเปฯ.\csromanspacer{} อุโปสถกมฺมํ ฯเปฯ.\csromanspacer{} วิปสฺสนํ อุปฺปาเทติ.\csromanspacer{} มานํ ชปฺเปติ.\csromanspacer{} ทิฏฺฐิํ คณฺหาติ.\csromanspacer{} นกามาวจรา สทฺธา ฯเปฯ ปญฺญา กามาวจราย สทฺธาย ฯเปฯ ปญฺญาย.\csromanspacer{} ราคสฺส ฯเปฯ.\csromanspacer{} ปตฺถนาย.\csromanspacer{} กายิกสฺส สุขสฺส.\csromanspacer{} กายิกสฺส ทุกฺขสฺส อุปนิสฺสยปจฺจเยน ปจฺจโย.\csromanspacer{} อริยา มคฺคํ อุปนิสฺสาย สงฺขาเร อนิจฺจโต ฯเปฯ วิปสฺสนฺติ.\csromanspacer{} มคฺโค อริยานํ อตฺถปฏิสมฺภิทาย.\csromanspacer{} ธมฺมปฏิสมฺภิทาย.\csromanspacer{} นิรุตฺติปฏิสมฺภิทาย.\csromanspacer{} ปฏิภานปฏิสมฺภิทาย.\csromanspacer{} ฐานาฐานโกสลฺลสฺส อุปนิสฺสยปจฺจเยน ปจฺจโย.\csromanspacer{} ผลสมาปตฺติ กายิกสฺส สุขสฺส อุปนิสฺสยปจฺจเยน ปจฺจโย. (2)}

\csromanheader{2}{\textbf{ปุเรชาตาทิ}}
\proseitem{215}{กามาวจโร ธมฺโม กามาวจรสฺส ธมฺมสฺส ปุเรชาตปจฺจเยน ปจฺจโย.\csromanspacer{} อารมฺมณปุเรชาตํ, วตฺถุปุเรชาตํ.\csromanspacer{}\csromanspacer{}\textbf{อารมฺมณปุเรชาตํ\csromandash{}}จกฺขุํ ฯเปฯ.\csromanspacer{} วตฺถุํ อนิจฺจโต ฯเปฯ โทมนสฺสํ อุปฺปชฺชติ.\csromanspacer{} รูปายตนํ จกฺขุวิญฺญาณสฺส ฯเปฯ.\csromanspacer{} โผฏฺฐพฺพายตนํ กายวิญฺญาณสฺส ฯเปฯ.\csromanspacer{} \textbf{วตฺถุปุเรชาตํ\csromandash{}}จกฺขายตนํ จกฺขุวิญฺญาณสฺส ฯเปฯ.\csromanspacer{} กายายตนํ กายวิญฺญาณสฺส ฯเปฯ.\csromanspacer{} วตฺถุ กามาวจรานํ ขนฺธานํ ปุเรชาตปจฺจเยน ปจฺจโย.\csromanspacer{}\csromanspacer{}กามาวจโร \csromanfolio{138}ธมฺโม นกามาวจรสฺส ธมฺมสฺส ปุเรชาตปจฺจเยน ปจฺจโย.\csromanspacer{} อารมฺมณปุเรชาตํ, วตฺถุปุเรชาตํ.\csromanspacer{}\csromanspacer{}\textbf{อารมฺมณปุเรชาตํ\csromandash{}}ทิพฺเพน จกฺขุนา รูปํ ปสฺสติ.\csromanspacer{} ทิพฺพาย โสตธาตุยา สทฺทํ สุณาติ.\csromanspacer{} \textbf{วตฺถุปุเรชาตํ\csromandash{}}วตฺถุ นกามาวจรานํ ขนฺธานํ ปุเรชาตปจฺจเยน ปจฺจโย. (2)}

\prose{ปจฺฉาชาตปจฺจเยน ปจฺจโย.\csromanspacer{} เทฺว.\csromanspacer{} อาเสวนปจฺจเยน ปจฺจโย.\csromanspacer{} ตีณิ.}

\csromanheader{2}{\textbf{กมฺมาทิ}}
\proseitem{216}{กามาวจโร ธมฺโม กามาวจรสฺส ธมฺมสฺส กมฺมปจฺจเยน ปจฺจโย.\csromanspacer{} \textbf{สหชาตา}, นานากฺขณิกา. (สํขิตฺตํ.)}

\prose{นกามาวจโร ธมฺโม นกามาวจรสฺส ธมฺมสฺส กมฺมปจฺจเยน ปจฺจโย.\csromanspacer{} \textbf{สหชาตา}, นานากฺขณิกา. (สํขิตฺตํ.).\csromanspacer{} นกามาวจโร ธมฺโม กามาวจรสฺส ธมฺมสฺส กมฺมปจฺจเยน ปจฺจโย.\csromanspacer{} \textbf{สหชาตา}, นานากฺขณิกา. (สํขิตฺตํ.).\csromanspacer{} นกามาวจโร ธมฺโม กามาวจรสฺส จ นกามาวจรสฺส จ ธมฺมสฺส กมฺมปจฺจเยน ปจฺจโย.\csromanspacer{} \textbf{สหชาตา}, นานากฺขณิกา. (สํขิตฺตํ.) (3)}

\prose{วิปากปจฺจเยน ปจฺจโย.\csromanspacer{} จตฺตาริ.\csromanspacer{} อาหารปจฺจเยน ปจฺจโย.\csromanspacer{} จตฺตาริ.\csromanspacer{} อินฺทฺริยปจฺจเยน ปจฺจโย.\csromanspacer{} จตฺตาริ.\csromanspacer{} ฌานปจฺจเยน ปจฺจโย.\csromanspacer{} จตฺตาริ.\csromanspacer{} มคฺคปจฺจเยน ปจฺจโย.\csromanspacer{} จตฺตาริ.\csromanspacer{} สมฺปยุตฺตปจฺจเยน ปจฺจโย.\csromanspacer{} เทฺว.}

\csromanheader{2}{\textbf{วิปฺปยุตฺต}}
\proseitem{217}{กามาวจโร ธมฺโม กามาวจรสฺส ธมฺมสฺส วิปฺปยุตฺตปจฺจเยน ปจฺจโย.\csromanspacer{} สหชาตํ, ปุเรชาตํ, ปจฺฉาชาตํ. (สํขิตฺตํ.).\csromanspacer{} กามาวจโร ธมฺโม นกามาวจรสฺส ธมฺมสฺส วิปฺปยุตฺตปจฺจเยน ปจฺจโย.\csromanspacer{} สหชาตํ, ปุเรชาตํ.\csromanspacer{}\csromanspacer{}สหชาตํ\csromandash{} (สํขิตฺตํ.) ปฏิสนฺธิกฺขเณ วตฺถุ นกามาวจรานํ ขนฺธานํ วิปฺปยุตฺตปจฺจเยน ปจฺจโย.\csromanspacer{} \textbf{ปุเรชาตํ} วตฺถุ นกามาวจรานํ ขนฺธานํ วิปฺปยุตฺตปจฺจเยน ปจฺจโย. (2)}

\prose{นกามาวจโร ธมฺโม กามาวจรสฺส ธมฺมสฺส วิปฺปยุตฺตปจฺจเยน ปจฺจโย.\csromanspacer{} สหชาตํ, ปจฺฉาชาตํ.\csromanspacer{}\csromanspacer{}\textbf{สหชาตา} นกามาวจรา ขนฺธา จิตฺตสมุฏฺฐานานํ รูปานํ วิปฺปยุตฺตปจฺจเยน ปจฺจโย.\csromanspacer{} ปฏิสนฺธิกฺขเณ ฯเปฯ.\csromanspacer{} \textbf{ปจฺฉาชาตา} นกามาวจรา ขนฺธา ปุเรชาตสฺส อิมสฺส กายสฺส วิปฺปยุตฺตปจฺจเยน ปจฺจโย. (1)}

\csromanheader{2}{93. กามาวจรทุก \textbf{อตฺถิ-อาทิ}}
\proseitem{218}{\csromanfolio{139}กามาวจโร ธมฺโม กามาวจรสฺส ธมฺมสฺส อตฺถิปจฺจเยน ปจฺจโย.\csromanspacer{} \textbf{สหชาตํ}, ปุเรชาตํ, ปจฺฉาชาตํ, อาหารํ, อินฺทฺริยํ. (สํขิตฺตํ.).\csromanspacer{} กามาวจโร ธมฺโม นกามาวจรสฺส ธมฺมสฺส อตฺถิปจฺจเยน ปจฺจโย.\csromanspacer{} \textbf{สหชาตํ}, ปุเรชาตํ.\csromanspacer{}\csromanspacer{}\textbf{สหชาตํ} ปฏิสนฺธิกฺขเณ วตฺถุ นกามาวจรานํ ขนฺธานํ อตฺถิปจฺจเยน ปจฺจโย.\csromanspacer{} \textbf{ปุเรชาตํ} ทิพฺเพน จกฺขุนา รูปํ ปสฺสติ ฯเปฯ. (ปุเรชาตสทิสํ.) (2)}

\prose{นกามาวจโร ธมฺโม นกามาวจรสฺส ธมฺมสฺส อตฺถิปจฺจเยน ปจฺจโย.\csromanspacer{} \textbf{สหชาตํ}. (สํขิตฺตํ.).\csromanspacer{} นกามาวจโร ธมฺโม กามาวจรสฺส ธมฺมสฺส อตฺถิปจฺจเยน ปจฺจโย.\csromanspacer{} \textbf{สหชาตํ}, ปจฺฉาชาตํ. (วิปฺปยุตฺตสทิสํ.).\csromanspacer{} นกามาวจโร ธมฺโม กามาวจรสฺส จ นกามาวจรสฺส จ ธมฺมสฺส อตฺถิปจฺจเยน ปจฺจโย. (ปฏิจฺจสทิสํ.) (3)}

\proseitem{219}{กามาวจโร จ นกามาวจโร จ ธมฺมา กามาวจรสฺส ธมฺมสฺส อตฺถิปจฺจเยน ปจฺจโย.\csromanspacer{} \textbf{สหชาตํ}, ปจฺฉาชาตํ, อาหารํ, อินฺทฺริยํ.\csromanspacer{}\csromanspacer{}\textbf{สหชาตา} นกามาวจรา ขนฺธา จ มหาภูตา จ จิตฺตสมุฏฺฐานานํ รูปานํ อตฺถิปจฺจเยน ปจฺจโย.\csromanspacer{} ปฏิสนฺธิกฺขเณ ฯเปฯ.\csromanspacer{} \textbf{ปจฺฉาชาตา} นกามาวจรา ขนฺธา จ กพฬีกาโร อาหาโร จ อิมสฺส กายสฺส อตฺถิปจฺจเยน ปจฺจโย.\csromanspacer{} \textbf{ปจฺฉาชาตา} นกามาวจรา ขนฺธา จ รูปชีวิตินฺทฺริยญฺจ กฏตฺตารูปานํ อตฺถิปจฺจเยน ปจฺจโย.\csromanspacer{}\csromanspacer{}กามาวจโร จ นกามาวจโร จ ธมฺมา นกามาวจรสฺส ธมฺมสฺส อตฺถิปจฺจเยน ปจฺจโย.\csromanspacer{} \textbf{สหชาตํ}, ปุเรชาตํ.\csromanspacer{}\csromanspacer{}สหชาโต นกามาวจโร เอโก ขนฺโธ จ วตฺถุ จ ติณฺณนฺนํ ขนฺธานํ อตฺถิปจฺจเยน ปจฺจโย ฯเปฯ เทฺว ขนฺธา จ ฯเปฯ.\csromanspacer{} ปฏิสนฺธิกฺขเณ ฯเปฯ. (2)}

\prose{นตฺถิปจฺจเยน ปจฺจโย.\csromanspacer{} วิคตปจฺจเยน ปจฺจโย.\csromanspacer{} อวิคตปจฺจเยน ปจฺจโย.}

\csromanheader{2}{1. ปจฺจยานุโลม 2. สงฺขฺยาวาร}
\csromanheader{2}{\textbf{สุทฺธ}}
\proseitem{220}{เหตุยา จตฺตาริ, อารมฺมเณ จตฺตาริ, อธิปติยา จตฺตาริ, อนนฺตเร จตฺตาริ, สมนนฺตเร จตฺตาริ, สหชาเต สตฺต, อญฺญมญฺเญ ฉ, นิสฺสเย สตฺต, อุปนิสฺสเย จตฺตาริ, ปุเรชาเต เทฺว, ปจฺฉาชาเต เทฺว, \csromanfolio{140}อาเสวเน ตีณิ, กมฺเม จตฺตาริ, วิปาเก จตฺตาริ, อาหาเร จตฺตาริ, อินฺทฺริเย จตฺตาริ, ฌาเน จตฺตาริ, มคฺเค จตฺตาริ, สมฺปยุตฺเต เทฺว, วิปฺปยุตฺเต ตีณิ, อตฺถิยา สตฺต, นตฺถิยา จตฺตาริ, วิคเต จตฺตาริ, อวิคเต สตฺต.}

\csromanheader{2}{2. ปจฺจนียุทฺธาร}
\proseitem{221}{กามาวจโร ธมฺโม กามาวจรสฺส ธมฺมสฺส อารมฺมณปจฺจเยน ปจฺจโย.\csromanspacer{} สหชาตปจฺจเยน ปจฺจโย.\csromanspacer{} อุปนิสฺสยปจฺจเยน ปจฺจโย.\csromanspacer{} ปุเรชาตปจฺจเยน ปจฺจโย.\csromanspacer{} ปจฺฉาชาตปจฺจเยน ปจฺจโย.\csromanspacer{} กมฺมปจฺจเยน ปจฺจโย.\csromanspacer{} อาหารปจฺจเยน ปจฺจโย.\csromanspacer{} อินฺทฺริยปจฺจเยน ปจฺจโย.\csromanspacer{}\csromanspacer{}กามาวจโร ธมฺโม นกามาวจรสฺส ธมฺมสฺส อารมฺมณปจฺจเยน ปจฺจโย.\csromanspacer{} สหชาตปจฺจเยน ปจฺจโย.\csromanspacer{} อุปนิสฺสยปจฺจเยน ปจฺจโย.\csromanspacer{} ปุเรชาตปจฺจเยน ปจฺจโย. (2)}

\proseitem{222}{นกามาวจโร ธมฺโม นกามาวจรสฺส ธมฺมสฺส อารมฺมณปจฺจเยน ปจฺจโย.\csromanspacer{} สหชาตปจฺจเยน ปจฺจโย.\csromanspacer{} อุปนิสฺสยปจฺจเยน ปจฺจโย.\csromanspacer{}\csromanspacer{}นกามาวจโร ธมฺโม กามาวจรสฺส ธมฺมสฺส อารมฺมณปจฺจเยน ปจฺจโย.\csromanspacer{} สหชาตปจฺจเยน ปจฺจโย.\csromanspacer{} อุปนิสฺสยปจฺจเยน ปจฺจโย.\csromanspacer{} ปจฺฉาชาตปจฺจเยน ปจฺจโย.\csromanspacer{} กมฺมปจฺจเยน ปจฺจโย.\csromanspacer{}\csromanspacer{}นกามาวจโร ธมฺโม กามาวจรสฺส จ นกามาวจรสฺส จ ธมฺมสฺส สหชาตปจฺจเยน ปจฺจโย.\csromanspacer{} กมฺมปจฺจเยน ปจฺจโย. (3)}

\prose{กามาวจโร จ นกามาวจโร จ ธมฺมา กามาวจรสฺส ธมฺมสฺส สหชาตํ, ปจฺฉาชาตํ, อาหารํ, อินฺทฺริยํ.\csromanspacer{}\csromanspacer{}กามาวจโร จ นกามาวจโร จ ธมฺมา นกามาวจรสฺส ธมฺมสฺส สหชาตํ, ปุเรชาตํ. (2)}

\csromanheader{2}{2. ปจฺจยปจฺจนีย 2. สงฺขฺยาวาร}
\csromanheader{2}{\textbf{สุทฺธ}}
\proseitem{223}{นเหตุยา สตฺต, น-อารมฺมเณ สตฺต, น-อธิปติยา สตฺต, น-อนนฺตเร สตฺต, นสมนนฺตเร สตฺต, น สหชาเต ฉ, น-อญฺญมญฺเญ ฉ, นนิสฺสเย ฉ, น-อุปนิสฺสเย สตฺต, นปุเรชาเต ฉ ฯเปฯ นสมฺปยุตฺเต ฉ,}

\vspace{\tipitakaparskip}
\csromancenter{กามาวจรทุก นวิปฺปยุตฺเต ปญฺจ, โน-อตฺถิยา ปญฺจ, โนนตฺถิยา สตฺต, โนวิคเต สตฺต, โน-อวิคเต ปญฺจ.}
\csromanheader{2}{3. ปจฺจยานุโลมปจฺจนีย}
\proseitem{224}{\csromanfolio{141}\textbf{เหตุ}ปจฺจยา น-อารมฺมเณ จตฺตาริ, น-อธิปติยา จตฺตาริ, น-อนนฺตเร จตฺตาริ, นสมนนฺตเร จตฺตาริ, น-อญฺญมญฺเญ เทฺว, น-อุปนิสฺสเย จตฺตาริ ฯเปฯ นสมฺปยุตฺเต เทฺว, นวิปฺปยุตฺเต เทฺว, โนนตฺถิยา จตฺตาริ, โนวิคเต จตฺตาริ.}

\csromanheader{2}{4. ปจฺจยปจฺจนียานุโลม}
\proseitem{225}{นเหตุปจฺจยา อารมฺมเณ จตฺตาริ, อธิปติยา จตฺตาริ, (อนุโลมมาติกา กาตพฺพา.) ฯเปฯ อวิคเต สตฺต.}

\nitthitam{กามาวจรทุกํ นิฏฺฐิตํ.}
\csromanmatikaanchor{mtk.23}
\csromantocmarkref{section}{94. รูปาวจรทุก}{mtk.23}
\bookmark[dest={mtk.23},level=1]{94. รูปาวจรทุก}
\csromanheader{1}{94. รูปาวจรทุก 1. ปฏิจฺจวาร}
\proseitem{226}{รูปาวจรํ ธมฺมํ ปฏิจฺจ รูปาวจโร ธมฺโม อุปฺปชฺชติ เหตุปจฺจยา.\csromanspacer{} รูปาวจรํ เอกํ ขนฺธํ ปฏิจฺจ ตโย ขนฺธา ฯเปฯ เทฺว ขนฺเธ ฯเปฯ.\csromanspacer{} ปฏิสนฺธิกฺขเณ ฯเปฯ.\csromanspacer{} รูปาวจรํ ธมฺมํ ปฏิจฺจ นรูปาวจโร ธมฺโม อุปฺปชฺชติ เหตุปจฺจยา.\csromanspacer{} รูปาวจเร ขนฺเธ ปฏิจฺจ จิตฺตสมุฏฺฐานํ รูปํ.\csromanspacer{} ปฏิสนฺธิกฺขเณ ฯเปฯ.\csromanspacer{} รูปาวจรํ ธมฺมํ ปฏิจฺจ รูปาวจโร จ นรูปาวจโร จ ธมฺมา อุปฺปชฺชนฺติ เหตุปจฺจยา.\csromanspacer{} รูปาวจรํ เอกํ ขนฺธํ ปฏิจฺจ ตโย ขนฺธา จิตฺตสมุฏฺฐานญฺจ รูปํ ฯเปฯ เทฺว ขนฺเธ ฯเปฯ.\csromanspacer{} ปฏิสนฺธิกฺขเณ ฯเปฯ. (3)}

\proseitem{227}{นรูปาวจรํ ธมฺมํ ปฏิจฺจ นรูปาวจโร ธมฺโม อุปฺปชฺชติ เหตุปจฺจยา.\csromanspacer{} นรูปาวจรํ เอกํ ขนฺธํ ปฏิจฺจ ตโย ขนฺธา จิตฺตสมุฏฺฐานญฺจ รูปํ ฯเปฯ เทฺว ขนฺเธ ฯเปฯ. \csromanfolio{142}ปฏิสนฺธิกฺขเณ ฯเปฯ.\csromanspacer{} เอกํ มหาภูตํ ฯเปฯ.\csromanspacer{} นรูปาวจรํ ธมฺมํ ปฏิจฺจ รูปาวจโร ธมฺโม อุปฺปชฺชติ เหตุปจฺจยา.\csromanspacer{} ปฏิสนฺธิกฺขเณ วตฺถุํ ปฏิจฺจ รูปาวจรา ขนฺธา.\csromanspacer{}\csromanspacer{}นรูปาวจรํ ธมฺมํ ปฏิจฺจ รูปาวจโร จ นรูปาวจโร จ ธมฺมา อุปฺปชฺชนฺติ เหตุปจฺจยา.\csromanspacer{} ปฏิสนฺธิกฺขเณ วตฺถุํ ปฏิจฺจ รูปาวจรา ขนฺธา.\csromanspacer{} มหาภูเต ปฏิจฺจ กฏตฺตารูปํ. (3)}

\proseitem{228}{รูปาวจรญฺจ นรูปาวจรญฺจ ธมฺมํ ปฏิจฺจ รูปาวจโร ธมฺโม อุปฺปชฺชติ เหตุปจฺจยา.\csromanspacer{} ปฏิสนฺธิกฺขเณ รูปาวจรํ เอกํ ขนฺธญฺจ วตฺถุญฺจ ปฏิจฺจ ตโย ขนฺธา ฯเปฯ เทฺว ขนฺเธ จ ฯเปฯ.\csromanspacer{} รูปาวจรญฺจ นรูปาวจรญฺจ ธมฺมํ ปฏิจฺจ นรูปาวจโร ธมฺโม อุปฺปชฺชติ เหตุปจฺจยา.\csromanspacer{} รูปาวจเร ขนฺเธ จ มหาภูเต จ ปฏิจฺจ จิตฺตสมุฏฺฐานํ รูปํ.\csromanspacer{} ปฏิสนฺธิกฺขเณ รูปาวจเร ขนฺเธ จ มหาภูเต จ ปฏิจฺจ กฏตฺตารูปํ.\csromanspacer{}\csromanspacer{}รูปาวจรญฺจ นรูปาวจรญฺจ ธมฺมํ ปฏิจฺจ รูปาวจโร จ นรูปาวจโร จ ธมฺมา อุปฺปชฺชนฺติ เหตุปจฺจยา.\csromanspacer{} ปฏิสนฺธิกฺขเณ รูปาวจรํ เอกํ ขนฺธญฺจ วตฺถุญฺจ ปฏิจฺจ ตโย ขนฺธา ฯเปฯ เทฺว ขนฺเธ จ ฯเปฯ.\csromanspacer{} รูปาวจเร ขนฺเธ จ มหาภูเต จ ปฏิจฺจ กฏตฺตารูปํ. (สํขิตฺตํ.) (3)}

\csromanheader{2}{1. ปจฺจยานุโลม 2. สงฺขฺยาวาร}
\csromanheader{2}{\textbf{สุทฺธ}}
\proseitem{229}{เหตุยา นว, อารมฺมเณ จตฺตาริ, อธิปติยา ปญฺจ, อนนฺตเร จตฺตาริ, สมนนฺตเร จตฺตาริ, สหชาเต นว, อญฺญมญฺเญ ฉ, นิสฺสเย นว, อุปนิสฺสเย จตฺตาริ, ปุเรชาเต เทฺว, อาเสวเน เทฺว, กมฺเม นว, วิปาเก นว, ฌาเน นว, มคฺเค นว, สมฺปยุตฺเต จตฺตาริ, วิปฺปยุตฺเต นว, อตฺถิยา นว, นตฺถิยา จตฺตาริ, วิคเต จตฺตาริ, อวิคเต นว.}

\csromanheader{2}{2. ปจฺจยปจฺจนีย 1. วิภงฺควาร}
\csromanheader{2}{\textbf{นเหตุ น-อารมฺมณ}}
\proseitem{230}{นรูปาวจรํ ธมฺมํ ปฏิจฺจ นรูปาวจโร ธมฺโม อุปฺปชฺชติ นเหตุปจฺจยา.\csromanspacer{} อเหตุกํ นรูปาวจรํ เอกํ ขนฺธํ ปฏิจฺจ ตโย ขนฺธา จิตฺตสมุฏฺฐานญฺจ รูปํ ฯเปฯ เทฺว ขนฺเธ ฯเปฯ.\csromanspacer{} อเหตุกปฏิสนฺธิกฺขเณ ฯเปฯ.\csromanspacer{} ขนฺเธ ปฏิจฺจ วตฺถุ, วตฺถุํ ปฏิจฺจ ขนฺธา.\csromanspacer{} เอกํ มหาภูตํ ฯเปฯ. (ยาว อสญฺญสตฺตา.) วิจิกิจฺฉาสหคเต}

\vspace{\tipitakaparskip}
\csromancenter{รูปาวจรทุก อุทฺธจฺจสหคเต ขนฺเธ ปฏิจฺจ วิจิกิจฺฉาสหคโต อุทฺธจฺจสหคโต โมโห. (1).\csromanspacer{}\csromanspacer{}น-อารมฺมณปจฺจยา ตีณิ.}
\csromanheader{2}{\textbf{น-อธิปตฺยาทิ}}
\proseitem{231}{\csromanfolio{143}รูปาวจรํ ธมฺมํ ปฏิจฺจ รูปาวจโร ธมฺโม อุปฺปชฺชติ น-อธิปติปจฺจยา.\csromanspacer{} รูปาวจเร ขนฺเธ ปฏิจฺจ รูปาวจราธิปติ.\csromanspacer{} วิปากํ รูปาวจรํ เอกํ ขนฺธํ ปฏิจฺจ ตโย ขนฺธา ฯเปฯ เทฺว ขนฺเธ ฯเปฯ.\csromanspacer{} ปฏิสนฺธิกฺขเณ ฯเปฯ.\csromanspacer{} รูปาวจรํ ธมฺมํ ปฏิจฺจ นรูปาวจโร ธมฺโม อุปฺปชฺชติ น-อธิปติปจฺจยา.\csromanspacer{} วิปาเก รูปาวจเร ขนฺเธ ปฏิจฺจ จิตฺตสมุฏฺฐานํ รูปํ.\csromanspacer{} ปฏิสนฺธิกฺขเณ ฯเปฯ.\csromanspacer{} รูปาวจรํ ธมฺมํ ปฏิจฺจ รูปาวจโร จ นรูปาวจโร จ ธมฺมา อุปฺปชฺชนฺติ น-อธิปติปจฺจยา.\csromanspacer{} วิปากํ รูปาวจรํ เอกํ ขนฺธํ ปฏิจฺจ ตโย ขนฺธา จิตฺตสมุฏฺฐานญฺจ รูปํ ฯเปฯ เทฺว ขนฺเธ ฯเปฯ.\csromanspacer{} ปฏิสนฺธิกฺขเณ ฯเปฯ. (3)}

\prose{นรูปาวจรํ ธมฺมํ ปฏิจฺจ นรูปาวจโร ธมฺโม อุปฺปชฺชติ น-อธิปติปจฺจยา.\csromanspacer{} นรูปาวจรํ เอกํ ขนฺธํ ปฏิจฺจ ตโย ขนฺธา จิตฺตสมุฏฺฐานญฺจ รูปํ ฯเปฯ เทฺว ขนฺเธ ฯเปฯ.\csromanspacer{} ตีณิ. (นกามาวจรํ ปฏิจฺจวารสทิสํ นินฺนานํ, อิธ สพฺเพ มหาภูตา กาตพฺพา.) (3)}

\proseitem{232}{รูปาวจรญฺจ นรูปาวจรญฺจ ธมฺมํ ปฏิจฺจ รูปาวจโร ธมฺโม อุปฺปชฺชติ น-อธิปติปจฺจยา.\csromanspacer{} ปฏิสนฺธิกฺขเณ รูปาวจรํ เอกํ ขนฺธญฺจ วตฺถุญฺจ ปฏิจฺจ ตโย ขนฺธา ฯเปฯ เทฺว ขนฺเธ จ ฯเปฯ.\csromanspacer{} รูปาวจรญฺจ นรูปาวจรญฺจ ธมฺมํ ปฏิจฺจ นรูปาวจโร ธมฺโม อุปฺปชฺชติ น-อธิปติปจฺจยา.\csromanspacer{} วิปาเก รูปาวจเร ขนฺเธ จ มหาภูเต จ ปฏิจฺจ จิตฺตสมุฏฺฐานํ รูปํ.\csromanspacer{} ปฏิสนฺธิกฺขเณ ฯเปฯ.\csromanspacer{} รูปาวจรญฺจ นรูปาวจรญฺจ ธมฺมํ ปฏิจฺจ รูปาวจโร จ นรูปาวจโร จ ธมฺมา อุปฺปชฺชนฺติ น-อธิปติปจฺจยา.\csromanspacer{} ปฏิสนฺธิกฺขเณ รูปาวจรํ เอกํ ขนฺธญฺจ วตฺถุญฺจ ปฏิจฺจ ตโย ขนฺธา ฯเปฯ เทฺว ขนฺเธ จ ฯเปฯ.\csromanspacer{} รูปาวจเร ขนฺเธ จ มหาภูเต จ ปฏิจฺจ กฏตฺตารูปํ. (3).\csromanspacer{}\csromanspacer{}น-อนนฺตรปจฺจยา ฯเปฯ.\csromanspacer{} น-อุปนิสฺสยปจฺจยา.}

\csromanheader{2}{\textbf{นปุเรชาตาทิ}}
\proseitem{233}{รูปาวจรํ ธมฺมํ ปฏิจฺจ รูปาวจโร ธมฺโม อุปฺปชฺชติ นปุเรชาตปจฺจยา.\csromanspacer{} ปฏิสนฺธิกฺขเณ รูปาวจรํ เอกํ ขนฺธํ ปฏิจฺจ ตโย ขนฺธา ฯเปฯ เทฺว ขนฺเธ ฯเปฯ.\csromanspacer{} รูปาวจรํ ธมฺมํ ปฏิจฺจ นรูปาวจโร ธมฺโม อุปฺปชฺชติ นปุเรชาตปจฺจยา.\csromanspacer{} รูปาวจเร ขนฺเธ ปฏิจฺจ จิตฺตสมุฏฺฐานํ รูปํ.\csromanspacer{} ปฏิสนฺธิกฺขเณ ฯเปฯ. \csromanfolio{144}รูปาวจรํ ธมฺมํ ปฏิจฺจ รูปาวจโร จ นรูปาวจโร จ ธมฺมา อุปฺปชฺชนฺติ นปุเรชาตปจฺจยา.\csromanspacer{} ปฏิสนฺธิกฺขเณ รูปาวจรํ เอกํ ขนฺธํ ปฏิจฺจ ตโย ขนฺธา กฏตฺตา จ รูปํ ฯเปฯ เทฺว ขนฺเธ ฯเปฯ. (3)}

\prose{นรูปาวจรํ ธมฺมํ ปฏิจฺจ นรูปาวจโร ธมฺโม อุปฺปชฺชติ นปุเรชาตปจฺจยา.\csromanspacer{} อรูเป นรูปาวจรํ เอกํ ขนฺธํ ปฏิจฺจ ตโย ขนฺธา ฯเปฯ เทฺว ขนฺเธ ฯเปฯ.\csromanspacer{} นรูปาวจเร ขนฺเธ ปฏิจฺจ จิตฺตสมุฏฺฐานํ รูปํ.\csromanspacer{} ปฏิสนฺธิกฺขเณ ฯเปฯ. (ยาว อสญฺญสตฺตา.\csromanspacer{} อิตเร ปญฺจปิ ปญฺหา.\csromanspacer{} อนุโลมํ กาตพฺพํ.) นปจฺฉาชาตปจฺจยา.\csromanspacer{} นว.}

\csromanheader{2}{\textbf{น-อาเสวนาทิ}}
\proseitem{234}{รูปาวจรํ ธมฺมํ ปฏิจฺจ รูปาวจโร ธมฺโม อุปฺปชฺชติ น-อาเสวนปจฺจยา.\csromanspacer{} วิปากํ รูปาวจรํ เอกํ ขนฺธํ ปฏิจฺจ ตโย ขนฺธา ฯเปฯ เทฺว ขนฺเธ ฯเปฯ.\csromanspacer{} ปฏิสนฺธิกฺขเณ ฯเปฯ.\csromanspacer{} รูปาวจรํ ธมฺมํ ปฏิจฺจ นรูปาวจโร ธมฺโม อุปฺปชฺชติ น-อาเสวนปจฺจยา.\csromanspacer{} รูปาวจเร ขนฺเธ ปฏิจฺจ จิตฺตสมุฏฺฐานํ รูปํ.\csromanspacer{} ปฏิสนฺธิกฺขเณ ฯเปฯ.\csromanspacer{} รูปาวจรํ ธมฺมํ ปฏิจฺจ รูปาวจโร จ นรูปาวจโร จ ธมฺมา อุปฺปชฺชนฺติ น-อาเสวนปจฺจยา.\csromanspacer{} วิปากํ รูปาวจรํ เอกํ ขนฺธํ ปฏิจฺจ ตโย ขนฺธา จิตฺตสมุฏฺฐานญฺจ รูปํ ฯเปฯ เทฺว ขนฺเธ ฯเปฯ.\csromanspacer{} ปฏิสนฺธิกฺขเณ ฯเปฯ. (3)}

\prose{นรูปาวจรํ ธมฺมํ ปฏิจฺจ นรูปาวจโร ธมฺโม อุปฺปชฺชติ น-อาเสวนปจฺจยา.\csromanspacer{} ตีณิ.}

\prose{รูปาวจรญฺจ นรูปาวจรญฺจ ธมฺมํ ปฏิจฺจ รูปาวจโร ธมฺโม อุปฺปชฺชติ น-อาเสวนปจฺจยา. (สํขิตฺตํ.\csromanspacer{} มูลํ.\csromanspacer{} อิตเรปิ ปญฺหา กาตพฺพา.) นกมฺมปจฺจยา.\csromanspacer{} เทฺว ฯเปฯ.\csromanspacer{} นสมฺปยุตฺตปจฺจยา.}

\proseitem{235}{นรูปาวจรํ ธมฺมํ ปฏิจฺจ นรูปาวจโร ธมฺโม อุปฺปชฺชติ นวิปฺปยุตฺตปจฺจยา.\csromanspacer{} อรูเป นรูปาวจรํ เอกํ ขนฺธํ ปฏิจฺจ ตโย ขนฺธา ฯเปฯ เทฺว ขนฺเธ ฯเปฯ.\csromanspacer{} พาหิรํ.\csromanspacer{} อาหารสมุฏฺฐานํ.\csromanspacer{} อุตุสมุฏฺฐานํ.\csromanspacer{} อสญฺญสตฺตานํ ฯเปฯ. (สํขิตฺตํ.)}

\csromanheader{2}{2. ปจฺจยปจฺจนีย 2. สงฺขฺยาวาร}
\csromanheader{2}{\textbf{สุทฺธ}}
\proseitem{236}{นเหตุยา เอกํ, น-อารมฺมเณ ตีณิ, น-อธิปติยา นว, น-อนนฺตเร ตีณิ, นสมนนฺตเร ตีณิ, น-อุปนิสฺสเย ตีณิ, นปุเรชาเต นว, นปจฺฉาชาเต นว, น-อาเสวเน นว, นกมฺเม เทฺว, นวิปาเก}

\vspace{\tipitakaparskip}
\csromancenter{รูปาวจรทุก ปญฺจ, น-อาหาเร เอกํ, น-อินฺทฺริเย เอกํ, นฌาเน เอกํ, นมคฺเค เอกํ, นสมฺปยุตฺเต ตีณิ, นวิปฺปยุตฺเต เอกํ, โนนตฺถิยา ตีณิ, โนวิคเต ตีณิ.}
\csromancenter{(เอวํ อิตเร เทฺว คณนาปิ สหชาตวาโรปิ กาตพฺโพ.)}
\csromanheader{2}{94. รูปาวจรทุก 3. ปจฺจยวาร}
\proseitem{237}{\csromanfolio{145}รูปาวจรํ ธมฺมํ ปจฺจยา รูปาวจโร ธมฺโม อุปฺปชฺชติ เหตุปจฺจยา.\csromanspacer{} ตีณิ. (ปฏิจฺจสทิสํ.)}

\prose{นรูปาวจรํ ธมฺมํ ปจฺจยา นรูปาวจโร ธมฺโม อุปฺปชฺชติ เหตุปจฺจยา.\csromanspacer{} นรูปาวจรํ เอกํ ขนฺธํ ปจฺจยา ตโย ขนฺธา จิตฺตสมุฏฺฐานญฺจ รูปํ ฯเปฯ เทฺว ขนฺเธ ฯเปฯ. (ยาว อชฺฌตฺติกํ มหาภูตํ.) วตฺถุํ ปจฺจยา นรูปาวจรา ขนฺธา.\csromanspacer{}\csromanspacer{}นรูปาวจรํ ธมฺมํ ปจฺจยา รูปาวจโร ธมฺโม อุปฺปชฺชติ เหตุปจฺจยา.\csromanspacer{} วตฺถุํ ปจฺจยา รูปาวจรา ขนฺธา.\csromanspacer{} ปฏิสนฺธิกฺขเณ ฯเปฯ.\csromanspacer{} นรูปาวจรํ ธมฺมํ ปจฺจยา รูปาวจโร จ นรูปาวจโร จ ธมฺมา อุปฺปชฺชนฺติ เหตุปจฺจยา.\csromanspacer{} วตฺถุํ ปจฺจยา รูปาวจรา ขนฺธา.\csromanspacer{} มหาภูเต ปจฺจยา จิตฺตสมุฏฺฐานํ รูปํ.\csromanspacer{} ปฏิสนฺธิกฺขเณ ฯเปฯ. (3)}

\proseitem{238}{รูปาวจรญฺจ นรูปาวจรญฺจ ธมฺมํ ปจฺจยา รูปาวจโร ธมฺโม อุปฺปชฺชติ เหตุปจฺจยา.\csromanspacer{} รูปาวจรํ เอกํ ขนฺธญฺจ วตฺถุญฺจ ปจฺจยา ตโย ขนฺธา ฯเปฯ เทฺว ขนฺเธ จ ฯเปฯ.\csromanspacer{} ปฏิสนฺธิกฺขเณ ฯเปฯ.\csromanspacer{} รูปาวจรญฺจ นรูปาวจรญฺจ ธมฺมํ ปจฺจยา นรูปาวจโร ธมฺโม อุปฺปชฺชติ เหตุปจฺจยา.\csromanspacer{} รูปาวจเร ขนฺเธ จ มหาภูเต จ ปจฺจยา จิตฺตสมุฏฺฐานํ รูปํ.\csromanspacer{} ปฏิสนฺธิกฺขเณ ฯเปฯ.\csromanspacer{} รูปาวจรญฺจ นรูปาวจรญฺจ ธมฺมํ ปจฺจยา รูปาวจโร จ นรูปาวจโร จ ธมฺมา อุปฺปชฺชนฺติ เหตุปจฺจยา.\csromanspacer{} รูปาวจรํ เอกํ ขนฺธญฺจ วตฺถุญฺจ ปจฺจยา ตโย ขนฺธา ฯเปฯ เทฺว ขนฺเธ จ ฯเปฯ.\csromanspacer{} รูปาวจเร ขนฺเธ จ มหาภูเต จ ปจฺจยา จิตฺตสมุฏฺฐานํ รูปํ.\csromanspacer{} ปฏิสนฺธิกฺขเณ ฯเปฯ. (3) (สํขิตฺตํ.)}

\csromanheader{2}{1. ปจฺจยานุโลม 2. สงฺขฺยาวาร}
\proseitem{239}{\csromanfolio{146}เหตุยา นว, อารมฺมเณ จตฺตาริ, อธิปติยา นว, อนนฺตเร จตฺตาริ, สมนนฺตเร จตฺตาริ, สหชาเต นว, อญฺญมญฺเญ ฉ, นิสฺสเย นว, อุปนิสฺสเย จตฺตาริ, ปุเรชาเต จตฺตาริ, อาเสวเน จตฺตาริ, กมฺเม นว ฯเปฯ อวิคเต นว.}

\csromanheader{2}{2. ปจฺจยปจฺจนีย 1. วิภงฺควาร}
\csromanheader{2}{\textbf{นเหตุ}}
\proseitem{240}{นรูปาวจรํ ธมฺมํ ปจฺจยา นรูปาวจโร ธมฺโม อุปฺปชฺชติ นเหตุปจฺจยา.\csromanspacer{} อเหตุกํ นรูปาวจรํ เอกํ ขนฺธํ ปจฺจยา ตโย ขนฺธา จิตฺตสมุฏฺฐานญฺจ รูปํ ฯเปฯ เทฺว ขนฺเธ ฯเปฯ.\csromanspacer{} อเหตุกปฏิสนฺธิกฺขเณ ฯเปฯ. (ยาว อสญฺญสตฺตา.) จกฺขายตนํ ปจฺจยา จกฺขุวิญฺญาณํ ฯเปฯ.\csromanspacer{} กายายตนํ ปจฺจยา กายวิญฺญาณํ วตฺถุํ ปจฺจยา อเหตุกา นรูปาวจรา ขนฺธา.\csromanspacer{} วิจิกิจฺฉาสหคเต อุทฺธจฺจสหคเต ขนฺเธ จ วตฺถุญฺจ ปจฺจยา วิจิกิจฺฉาสหคโต อุทฺธจฺจสหคโต โมโห. (1) (สํขิตฺตํ.)}

\csromanheader{2}{\textbf{น-อธิปติ}}
\proseitem{241}{รูปาวจรํ ธมฺมํ ปจฺจยา รูปาวจโร ธมฺโม อุปฺปชฺชติ น-อธิปติปจฺจยา.\csromanspacer{} ตีณิ. (ปฏิจฺจวารสทิสํ.)}

\prose{นรูปาวจรํ ธมฺมํ ปจฺจยา นรูปาวจโร ธมฺโม อุปฺปชฺชติ น-อธิปติปจฺจยา. (ปฏิจฺจวารสทิสํ.).\csromanspacer{} นรูปาวจรํ ธมฺมํ ปจฺจยา รูปาวจโร ธมฺโม อุปฺปชฺชติ น-อธิปติปจฺจยา.\csromanspacer{} วตฺถุํ ปจฺจยา รูปาวจราธิปติ.\csromanspacer{} วตฺถุํ ปจฺจยา วิปากา รูปาวจรา ขนฺธา.\csromanspacer{} ปฏิสนฺธิกฺขเณ ฯเปฯ.\csromanspacer{} นรูปาวจรํ ธมฺมํ ปจฺจยา รูปาวจโร จ นรูปาวจโร จ ธมฺมา อุปฺปชฺชนฺติ น-อธิปติปจฺจยา.\csromanspacer{} วตฺถุํ ปจฺจยา วิปากา รูปาวจรา ขนฺธา.\csromanspacer{} มหาภูเต ปจฺจยา จิตฺตสมุฏฺฐานํ รูปํ.\csromanspacer{} ปฏิสนฺธิกฺขเณ ฯเปฯ. (3)}

\proseitem{242}{รูปาวจรญฺจ นรูปาวจรญฺจ ธมฺมํ ปจฺจยา รูปาวจโร ธมฺโม อุปฺปชฺชติ น-อธิปติปจฺจยา.\csromanspacer{} รูปาวจเร ขนฺเธ จ วตฺถุญฺจ ปจฺจยา รูปาวจราธิปติ.}

\vspace{\tipitakaparskip}
\csromancenter{รูปาวจรทุก วิปากํ รูปาวจรํ เอกํ ขนฺธญฺจ วตฺถุญฺจ ปจฺจยา ตโย ขนฺธา ฯเปฯ เทฺว ขนฺเธ จ ฯเปฯ.\csromanspacer{} ปฏิสนฺธิกฺขเณ รูปาวจรํ เอกํ ขนฺธญฺจ วตฺถุญฺจ ปจฺจยา ตโย ขนฺธา ฯเปฯ เทฺว ขนฺเธ จ ฯเปฯ.\csromanspacer{} รูปาวจรญฺจ นรูปาวจรญฺจ ธมฺมํ ปจฺจยา นรูปาวจโร ธมฺโม อุปฺปชฺชติ น-อธิปติปจฺจยา.\csromanspacer{} วิปาเก รูปาวจเร ขนฺเธ จ มหาภูเต จ ปจฺจยา จิตฺตสมุฏฺฐานํ รูปํ.\csromanspacer{} ปฏิสนฺธิกฺขเณ ฯเปฯ.\csromanspacer{} รูปาวจรญฺจ นรูปาวจรญฺจ ธมฺมํ ปจฺจยา รูปาวจโร จ นรูปาวจโร จ ธมฺมา อุปฺปชฺชนฺติ น-อธิปติปจฺจยา.\csromanspacer{} วิปากํ รูปาวจรํ เอกํ ขนฺธญฺจ วตฺถุญฺจ ปจฺจยา ตโย ขนฺธา ฯเปฯ เทฺว ขเธ จ ฯเปฯ.\csromanspacer{} วิปาเก รูปาวจเร ขนฺเธ จ มหาภูเต จ ปจฺจยา จิตฺตสมุฏฺฐานํ รูปํ.\csromanspacer{} ปฏิสนฺธิกฺขเณ ฯเปฯ. (3) (สํขิตฺตํ.)}
\csromanheader{2}{2. ปจฺจยปจฺจนีย 2. สงฺขฺยาวาร}
\csromanheader{2}{\textbf{สุทฺธ}}
\proseitem{243}{\csromanfolio{147}นเหตุยา เอกํ, น-อารมฺมเณ ตีณิ, น-อธิปติยา นว, น-อนนฺตเร ตีณิ ฯเปฯ น-อุปนิสฺสเย ตีณิ, นปุเรชาเต นว, นปจฺฉาชาเต นว, น-อาเสวเน นว, (สุทฺธิเก อรูเป จ มิสฺสเก จ “วิปากนฺ”ติ นิยาเมตพฺพํ.) นกมฺเม จตฺตาริ, นวิปาเก นว, น-อาหาเร เอกํ, น-อินฺทฺริเย เอกํ, นฌาเน เอกํ, นมคฺเค เอกํ, นสมฺปยุตฺเต ตีณิ, นวิปฺปยุตฺเต เอกํ, โนนตฺถิยา ตีณิ, โนวิคเต ตีณิ.}

\vspace{\tipitakaparskip}
\csromancenter{(เอวํ อิตเร เทฺว คณนาปิ นิสฺสยวาโรปิ กาตพฺโพ.)}
\csromanheader{2}{94. รูปาวจรทุก 5. สํสฏฺฐวาร}
\csromanheader{2}{1. ปจฺจยานุโลมาทิ}
\proseitem{244}{รูปาวจรํ ธมฺมํ สํสฏฺโฐ รูปาวจโร ธมฺโม อุปฺปชฺชติ เหตุปจฺจยา.\csromanspacer{} รูปาวจรํ เอกํ ขนฺธํ สํสฏฺฐา ตโย ขนฺธา ฯเปฯ เทฺว ขนฺเธ ฯเปฯ.\csromanspacer{} ปฏิสนฺธิกฺขเณ ฯเปฯ. (1)}

\prose{\csromanfolio{148}นรูปาวจรํ ธมฺมํ สํสฏฺโฐ นรูปาวจโร ธมฺโม อุปฺปชฺชติ เหตุปจฺจยา.\csromanspacer{} นรูปาวจรํ เอกํ ขนฺธํ สํสฏฺฐา ตโย ขนฺธา ฯเปฯ เทฺว ขนฺเธ ฯเปฯ.\csromanspacer{} ปฏิสนฺธิกฺขเณ ฯเปฯ. (1)}

\prose{\textbf{เหตุ}ยา เทฺว, อารมฺมเณ เทฺว, (สพฺพตฺถ เทฺว.) อวิคเต เทฺว.\csromanspacer{} อนุโลมํ.}

\prose{นรูปาวจรํ ธมฺมํ สํสฏฺโฐ นรูปาวจโร ธมฺโม อุปฺปชฺชติ นเหตุปจฺจยา. (สํขิตฺตํ.)}

\prose{นเหตุยา เอกํ, น-อธิปติยา เทฺว, นปุเรชาเต เทฺว, นปจฺฉาชาเต เทฺว, น-อาเสวเน เทฺว, นกมฺเม เทฺว, นวิปาเก เทฺว, นฌาเน เอกํ, นมคฺเค เอกํ, นวิปฺปยุตฺเต เอกํ.}

\vspace{\tipitakaparskip}
\csromancenter{ปจฺจนียํ.}
\csromancenter{(เอวํ อิตเร เทฺว คณนาปิ สมฺปยุตฺตวาโรปิ กาตพฺโพ.)}
\csromanheader{2}{94. รูปาวจรทุก 7. ปญฺหาวาร}
\proseitem{245}{รูปาวจโร ธมฺโม รูปาวจรสฺส ธมฺมสฺส เหตุปจฺจเยน ปจฺจโย.\csromanspacer{} รูปาวจรา เหตู สมฺปยุตฺตกานํ ขนฺธานํ เหตุปจฺจเยน ปจฺจโย.\csromanspacer{} ปฏิสนฺธิกฺขเณ ฯเปฯ. (มูลํ กาตพฺพํ.) รูปาวจรา เหตู จิตฺตสมุฏฺฐานานํ รูปานํ เหตุปจฺจเยน ปจฺจโย.\csromanspacer{} ปฏิสนฺธิกฺขเณ ฯเปฯ. (มูลํ กาตพฺพํ.) รูปาวจรา เหตู สมฺปยุตฺตกานํ ขนฺธานํ จิตฺตสมุฏฺฐานานญฺจ รูปานํ เหตุปจฺจเยน ปจฺจโย.\csromanspacer{} ปฏิสนฺธิกฺขเณ ฯเปฯ. (3)}

\prose{นรูปาวจโร ธมฺโม นรูปาวจรสฺส ธมฺมสฺส เหตุปจฺจเยน ปจฺจโย.\csromanspacer{} นรูปาวจรา เหตู สมฺปยุตฺตกานํ ขนฺธานํ จิตฺตสมุฏฺฐานานญฺจ รูปานํ เหตุปจฺจเยน ปจฺจโย.\csromanspacer{} ปฏิสนฺธิกฺขเณ ฯเปฯ. (1)}

\csromanheader{2}{94. รูปาวจรทุก \textbf{อารมฺมณ}}
\proseitem{246}{\csromanfolio{149}รูปาวจโร ธมฺโม รูปาวจรสฺส ธมฺมสฺส อารมฺมณปจฺจเยน ปจฺจโย.\csromanspacer{} เจโตปริยญาเณน รูปาวจรจิตฺตสมงฺคิสฺส จิตฺตํ ชานาติ.\csromanspacer{} รูปาวจรา ขนฺธา อิทฺธิวิธญาณสฺส, เจโตปริยญาณสฺส, ปุพฺเพนิวาสานุสฺสติญาณสฺส, ยถากมฺมูปคญาณสฺส, อนาคตํสญาณสฺส อารมฺมณปจฺจเยน ปจฺจโย.\csromanspacer{}\csromanspacer{}รูปาวจโร ธมฺโม นรูปาวจรสฺส ธมฺมสฺส อารมฺมณปจฺจเยน ปจฺจโย.\csromanspacer{} ปฐมํ ฌานํ ปจฺจเวกฺขติ ฯเปฯ.\csromanspacer{} จตุตฺถํ ฌานํ ปจฺจเวกฺขติ.\csromanspacer{} ทิพฺพํ จกฺขุํ ปจฺจเวกฺขติ.\csromanspacer{} ทิพฺพํ โสตธาตุํ ฯเปฯ.\csromanspacer{} อิทฺธิวิธญาณํ ฯเปฯ.\csromanspacer{} เจโตปริยญาณํ ฯเปฯ.\csromanspacer{} ปุพฺเพนิวาสานุสฺสติญาณํ ฯเปฯ.\csromanspacer{} ยถากมฺมูปคญาณํ ฯเปฯ.\csromanspacer{} อนาคตํสญาณํ ปจฺจเวกฺขติ.\csromanspacer{} รูปาวจเร ขนฺเธ อนิจฺจโต ฯเปฯ โทมนสฺสํ อุปฺปชฺชติ. (2)}

\proseitem{247}{นรูปาวจโร ธมฺโม นรูปาวจรสฺส ธมฺมสฺส อารมฺมณปจฺจเยน ปจฺจโย.\csromanspacer{} ทานํ ทตฺวา, สีลํ ฯเปฯ อุโปสถกมฺมํ ฯเปฯ.\csromanspacer{} ปุพฺเพ สุจิณฺณานิ ปจฺจเวกฺขติ, อสฺสาเทติ อภินนฺทติ, ตํ อารพฺภ ราโค ฯเปฯ โทมนสฺสํ อุปฺปชฺชติ.\csromanspacer{} อริยา มคฺคา วุฏฺฐหิตฺวา มคฺคํ ปจฺจเวกฺขนฺติ, ผลํ ปจฺจเวกฺขนฺติ, นิพฺพานํ ปจฺจเวกฺขนฺติ.\csromanspacer{} นิพฺพานํ โคตฺรภุสฺส, โวทานสฺส, มคฺคสฺส, ผลสฺส, อาวชฺชนาย อารมฺมณปจฺจเยน ปจฺจโย.\csromanspacer{} อริยา ปหีเน กิเลเส ฯเปฯ.\csromanspacer{} วิกฺขมฺภิเต กิเลเส ฯเปฯ.\csromanspacer{} ปุพฺเพ ฯเปฯ.\csromanspacer{} จกฺขุํ ฯเปฯ.\csromanspacer{} วตฺถุํ นรูปาวจเร ขนฺเธ อนิจฺจโต ฯเปฯ โทมนสฺสํ อุปฺปชฺชติ.\csromanspacer{} อากาสานญฺจายตนํ วิญฺญาณญฺจายตนสฺส ฯเปฯ.\csromanspacer{} อากิญฺจญฺญายตนํ เนวสญฺญานาสญฺญายตนสฺส ฯเปฯ.\csromanspacer{} รูปายตนํ จกฺขุวิญฺญาณสฺส ฯเปฯ.\csromanspacer{} โผฏฺฐพฺพายตนํ กายวิญฺญาณสฺส อารมฺมณปจฺจเยน ปจฺจโย. (1)}

\prose{นรูปาวจโร ธมฺโม รูปาวจรสฺส ธมฺมสฺส อารมฺมณปจฺจเยน ปจฺจโย.\csromanspacer{} ทิพฺเพน จกฺขุนา รูปํ ปสฺสติ, ทิพฺพาย โสตธาตุยา สทฺทํ สุณาติ.\csromanspacer{} เจโตปริยญาเณน นรูปาวจรจิตฺตสมงฺคิสฺส จิตฺตํ ชานาติ.\csromanspacer{} นรูปาวจรา ขนฺธา อิทฺธิวิธญาณสฺส, เจโตปริยญาณสฺส, ปุพฺเพนิวาสานุสฺสติญาณสฺส, ยถากมฺมูปคญาณสฺส, อนาคตํสญาณสฺส อารมฺมณปจฺจเยน ปจฺจโย. (2)}

\csromanheader{2}{\textbf{อธิปติ}}
\proseitem{248}{รูปาวจโร ธมฺโม รูปาวจรสฺส ธมฺมสฺส อธิปติปจฺจเยน ปจฺจโย.\csromanspacer{} \textbf{สหชาตาธิปติ\csromandash{}}รูปาวจราธิปติ สมฺปยุตฺตกานํ ขนฺธานํ อธิปติปจฺจเยน \csromanfolio{150}ปจฺจโย.\csromanspacer{}\csromanspacer{}รูปาวจโร ธมฺโม นรูปาวจรสฺส ธมฺมสฺส อธิปติปจฺจเยน ปจฺจโย.\csromanspacer{} อารมฺมณาธิปติ, สหชาตาธิปติ.\csromanspacer{}\csromanspacer{}\textbf{อารมฺมณาธิปติ\csromandash{}}ปฐมํ ฌานํ ครุํ กตฺวา ปจฺจเวกฺขติ, อสฺสาเทติ อภินนฺทติ, ตํ ครุํ กตฺวา ราโค อุปฺปชฺชติ, ทิฏฺฐิ อุปฺปชฺชติ ฯเปฯ.\csromanspacer{} จตุตฺถํ ฌานํ ฯเปฯ.\csromanspacer{} ทิพฺพํ จกฺขุํ ฯเปฯ.\csromanspacer{} ทิพฺพํ โสตธาตุํ ฯเปฯ.\csromanspacer{} อิทฺธิวิธญาณํ ฯเปฯ.\csromanspacer{} เจโตปริยญาณํ ฯเปฯ.\csromanspacer{} ปุพฺเพนิวาสานุสฺสติญาณํ ฯเปฯ.\csromanspacer{} ยถากมฺมูปคญาณํ ฯเปฯ.\csromanspacer{} อนาคตํสญาณํ ครุํ กตฺวา ปจฺจเวกฺขติ, อสฺสาเทติ อภินนฺทติ, ตํ ครุํ กตฺวา ราโค อุปฺปชฺชติ, ทิฏฺฐิ อุปฺปชฺชติ.\csromanspacer{} รูปาวจเร ขนฺเธ ครุํ กตฺวา อสฺสาเทติ อภินนฺทติ, ตํ ครุํ กตฺวา ราโค อุปฺปชฺชติ.\csromanspacer{} ทิฏฺฐิ อุปฺปชฺชติ.\csromanspacer{} \textbf{สหชาตาธิปติ\csromandash{}}รูปาวจราธิปติ จิตฺตสมุฏฺฐานานํ รูปานํ อธิปติปจฺจเยน ปจฺจโย.\csromanspacer{}\csromanspacer{}รูปาวจโร ธมฺโม รูปาวจรสฺส จ นรูปาวจรสฺส จ ธมฺมสฺส อธิปติปจฺจเยน ปจฺจโย.\csromanspacer{} \textbf{สหชาตาธิปติ\csromandash{}}รูปาวจราธิปติ สมฺปยุตฺตกานํ ขนฺธานํ จิตฺตสมุฏฺฐานานญฺจ รูปานํ อธิปติปจฺจเยน ปจฺจโย. (3)}

\proseitem{249}{นรูปาวจโร ธมฺโม นรูปาวจรสฺส ธมฺมสฺส อธิปติปจฺจเยน ปจฺจโย.\csromanspacer{} อารมฺมณาธิปติ, สหชาตาธิปติ.\csromanspacer{}\csromanspacer{}\textbf{อารมฺมณาธิปติ\csromandash{}}ทานํ ฯเปฯ สีลํ ฯเปฯ อุโปสถกมฺมํ ฯเปฯ.\csromanspacer{} ปุพฺเพ สุจิณฺณานิ ครุํ กตฺวา ปจฺจเวกฺขติ, อสฺสาเทติ อภินนฺทติ, ตํ ครุํ กตฺวา ราโค อุปฺปชฺชติ.\csromanspacer{} ทิฏฺฐิ อุปฺปชฺชติ.\csromanspacer{} อริยา มคฺคา วุฏฺฐหิตฺวา มคฺคํ ครุํ กตฺวา ฯเปฯ ผลํ ฯเปฯ นิพฺพานํ ครุํ กตฺวา ปจฺจเวกฺขนฺติ.\csromanspacer{} นิพฺพานํ โคตฺรภุสฺส, โวทานสฺส, มคฺคสฺส, ผลสฺส อธิปติปจฺจเยน ปจฺจโย.\csromanspacer{} จกฺขุํ ฯเปฯ.\csromanspacer{} วตฺถุํ นรูปาวจเร ขนฺเธ ครุํ กตฺวา อสฺสาเทติ อภินนฺทติ, ตํ ครุํ กตฺวา ราโค อุปฺปชฺชติ.\csromanspacer{} ทิฏฺฐิ อุปฺปชฺชติ.\csromanspacer{} อากาสานญฺจายตนํ ครุํ กตฺวา ปจฺจเวกฺขติ ฯเปฯ.\csromanspacer{} เนวสญฺญานาสญฺญายตนํ ครุํ กตฺวา ปจฺจเวกฺขติ.\csromanspacer{} \textbf{สหชาตาธิปติ\csromandash{}} นรูปาวจราธิปติ สมฺปยุตฺตกานํ ขนฺธานํ จิตฺตสมุฏฺฐานานญฺจ รูปานํ อธิปติปจฺจเยน ปจฺจโย. (1)}

\csromanheader{2}{\textbf{อนนฺตราทิ}}
\proseitem{250}{รูปาวจโร ธมฺโม รูปาวจรสฺส ธมฺมสฺส อนนฺตรปจฺจเยน ปจฺจโย.\csromanspacer{} ปุริมา ปุริมา รูปาวจรา ขนฺธา ปจฺฉิมานํ ปจฺฉิมานํ รูปาวจรานํ ขนฺธานํ อนนฺตรปจฺจเยน ปจฺจโย.\csromanspacer{}\csromanspacer{}รูปาวจโร ธมฺโม นรูปาวจรสฺส ธมฺมสฺส อนนฺตรปจฺจเยน ปจฺจโย.\csromanspacer{} รูปาวจรํ จุติจิตฺตํ นรูปาวจรสฺส}

\vspace{\tipitakaparskip}
\csromancenter{รูปาวจรทุก อุปปตฺติจิตฺตสฺส อนนฺตรปจฺจเยน ปจฺจโย.\csromanspacer{} รูปาวจรํ ภวงฺคํ อาวชฺชนาย.\csromanspacer{} รูปาวจรา ขนฺธา นรูปาวจรสฺส วุฏฺฐานสฺส อนนฺตรปจฺจเยน ปจฺจโย. (2)}
\proseitem{251}{\csromanfolio{151}นรูปาวจโร ธมฺโม นรูปาวจรสฺส ธมฺมสฺส อนนฺตรปจฺจเยน ปจฺจโย.\csromanspacer{} ปุริมา ปุริมา นรูปาวจรา ขนฺธา ปจฺฉิมานํ ปจฺฉิมานํ นรูปาวจรานํ ขนฺธานํ อนนฺตรปจฺจเยน ปจฺจโย.\csromanspacer{} อนุโลมํ โคตฺรภุสฺส ฯเปฯ.\csromanspacer{} เนวสญฺญานาสญฺญายตนํ ผลสมาปตฺติยา อนนฺตรปจฺจเยน ปจฺจโย.\csromanspacer{}\csromanspacer{}นรูปาวจโร ธมฺโม รูปาวจรสฺส ธมฺมสฺส อนนฺตรปจฺจเยน ปจฺจโย.\csromanspacer{} นรูปาวจรํ จุติจิตฺตํ รูปาวจรสฺส อุปปตฺติจิตฺตสฺส อนนฺตรปจฺจเยน ปจฺจโย.\csromanspacer{} นรูปาวจรา ขนฺธา รูปาวจรสฺส วุฏฺฐานสฺส อนนฺตรปจฺจเยน ปจฺจโย.\csromanspacer{} ปฐมสฺส ฌานสฺส ปริกมฺมํ ปฐมสฺส ฌานสฺส อนนฺตรปจฺจเยน ปจฺจโย ฯเปฯ.\csromanspacer{} จตุตฺถสฺส ฌานสฺส ฯเปฯ.\csromanspacer{} ทิพฺพสฺส จกฺขุสฺส ฯเปฯ.\csromanspacer{} ทิพฺพาย โสตธาตุยา ฯเปฯ.\csromanspacer{} อิทฺธิวิธญาณสฺส ฯเปฯ.\csromanspacer{} เจโตปริยญาณสฺส ฯเปฯ.\csromanspacer{} ปุพฺเพนิวาสานุสฺสติญาณสฺส ฯเปฯ.\csromanspacer{} ยถากมฺมูปคญาณสฺส ฯเปฯ.\csromanspacer{} อนาคตํสญาณสฺส ปริกมฺมํ อนาคตํสญาณสฺส อนนฺตรปจฺจเยน ปจฺจโย. (2)}

\prose{สมนนฺตรปจฺจเยน ปจฺจโย.\csromanspacer{} สหชาตปจฺจเยน ปจฺจโย.\csromanspacer{} สตฺต.\csromanspacer{} อญฺญมญฺญปจฺจเยน ปจฺจโย.\csromanspacer{} ฉ.\csromanspacer{} นิสฺสยปจฺจเยน ปจฺจโย.\csromanspacer{} สตฺต.}

\csromanheader{2}{\textbf{อุปนิสฺสย}}
\proseitem{252}{รูปาวจโร ธมฺโม รูปาวจรสฺส ธมฺมสฺส อุปนิสฺสยปจฺจเยน ปจฺจโย.\csromanspacer{} \textbf{อนนฺตรูปนิสฺสโย, ปกตูปนิสฺสโย ฯเปฯ.}\csromanspacer{} ปกตูปนิสฺสโย\csromandash{}รูปาวจรํ สทฺธํ อุปนิสฺสาย รูปาวจรํ ฌานํ อุปฺปาเทติ, อภิญฺญํ ฯเปฯ สมาปตฺติํ อุปฺปาเทติ.\csromanspacer{} รูปาวจรํ สีลํ ฯเปฯ ปญฺญํ อุปนิสฺสาย รูปาวจรํ ฌานํ ฯเปฯ อภิญฺญํ ฯเปฯ สมาปตฺติํ อุปฺปาเทติ.\csromanspacer{} รูปาวจรา สทฺธา ฯเปฯ ปญฺญา รูปาวจราย สทฺธาย ฯเปฯ ปญฺญาย อุปนิสฺสยปจฺจเยน ปจฺจโย.\csromanspacer{} ปฐมํ ฌานํ ทุติยสฺส ฌานสฺส อุปนิสฺสยปจฺจเยน ปจฺจโย.\csromanspacer{} ทุติยํ ฌานํ ตติยสฺส ฌานสฺส ฯเปฯ.\csromanspacer{} ตติยํ ฌานํ จตุตฺถสฺส ฌานสฺส อุปนิสฺสยปจฺจเยน ปจฺจโย. (1)}

\prose{รูปาวจโร ธมฺโม นรูปาวจรสฺส ธมฺมสฺส อุปนิสฺสยปจฺจเยน ปจฺจโย.\csromanspacer{} \textbf{อารมฺมณูปนิสฺสโย, อนนฺตรูปนิสฺสโย, ปกตูปนิสฺสโย ฯเปฯ.}\csromanspacer{} ปกตูปนิสฺสโย\csromandash{}รูปาวจรํ สทฺธํ อุปนิสฺสาย ทานํ เทติ.\csromanspacer{} สีลํ ฯเปฯ \csromanfolio{152}อุโปสถกมฺมํ กโรติ.\csromanspacer{} นรูปาวจรํ ฌานํ ฯเปฯ วิปสฺสนํ ฯเปฯ มคฺคํ ฯเปฯ สมาปตฺติํ อุปฺปาเทติ.\csromanspacer{} มานํ ชปฺเปติ, ทิฏฺฐิํ คณฺหาติ.\csromanspacer{} รูปาวจรํ สีลํ ฯเปฯ ปญฺญํ อุปนิสฺสาย ทานํ เทติ ฯเปฯ มานํ ชปฺเปติ, ทิฏฺฐิํ คณฺหาติ.\csromanspacer{} รูปาวจรา สทฺธา ฯเปฯ ปญฺญา นรูปาวจราย สทฺธาย ฯเปฯ ปญฺญาย.\csromanspacer{} ราคสฺส ฯเปฯ ปตฺถนาย.\csromanspacer{} กายิกสฺส สุขสฺส.\csromanspacer{} กายิกสฺส ทุกฺขสฺส.\csromanspacer{} มคฺคสฺส, ผลสมาปตฺติยา อุปนิสฺสยปจฺจเยน ปจฺจโย. (2)}

\proseitem{253}{นรูปาวจโร ธมฺโม นรูปาวจรสฺส ธมฺมสฺส อุปนิสฺสยปจฺจเยน ปจฺจโย.\csromanspacer{} \textbf{อารมฺมณูปนิสฺสโย, อนนฺตรูปนิสฺสโย, ปกตูปนิสฺสโย ฯเปฯ.}\csromanspacer{} ปกตูปนิสฺสโย\csromandash{}นรูปาวจรํ สทฺธํ อุปนิสฺสาย ทานํ ฯเปฯ สีลํ ฯเปฯ อุโปสถกมฺมํ กโรติ.\csromanspacer{} นรูปาวจรํ ฌานํ ฯเปฯ วิปสฺสนํ ฯเปฯ มคฺคํ ฯเปฯ สมาปตฺติํ อุปฺปาเทติ.\csromanspacer{} มานํ ชปฺเปติ, ทิฏฺฐิํ คณฺหาติ.\csromanspacer{} นรูปาวจรํ สีลํ ฯเปฯ ปญฺญํ.\csromanspacer{} ราคํ ฯเปฯ.\csromanspacer{} โภชนํ.\csromanspacer{} เสนาสนํ อุปนิสฺสาย ทานํ เทติ ฯเปฯ ปาณํ หนติ ฯเปฯ สํฆํ ภินฺทติ.\csromanspacer{} นรูปาวจรา สทฺธา ฯเปฯ ปญฺญา.\csromanspacer{} ราโค ฯเปฯ.\csromanspacer{} เสนาสนํ นรูปาวจราย สทฺธาย ฯเปฯ ปญฺญาย.\csromanspacer{} ราคสฺส ฯเปฯ ปตฺถนาย.\csromanspacer{} กายิกสฺส สุขสฺส.\csromanspacer{} กายิกสฺส ทุกฺขสฺส.\csromanspacer{} มคฺคสฺส, ผลสมาปตฺติยา อุปนิสฺสยปจฺจเยน ปจฺจโย. (1)}

\prose{นรูปาวจโร ธมฺโม รูปาวจรสฺส ธมฺมสฺส อุปนิสฺสยปจฺจเยน ปจฺจโย.\csromanspacer{} \textbf{อนนฺตรูปนิสฺสโย, ปกตูปนิสฺสโย ฯเปฯ.}\csromanspacer{} ปกตูปนิสฺสโย\csromandash{} นรูปาวจรํ สทฺธํ อุปนิสฺสาย รูปาวจรํ ฌานํ ฯเปฯ อภิญฺญํ ฯเปฯ สมาปตฺติํ อุปฺปาเทติ.\csromanspacer{} นรูปาวจรํ สีลํ ฯเปฯ.\csromanspacer{} เสนาสนํ อุปนิสฺสาย รูปาวจรํ ฌานํ ฯเปฯ อภิญฺญํ ฯเปฯ สมาปตฺติํ อุปฺปาเทติ.\csromanspacer{} นรูปาวจรา สทฺธา ฯเปฯ.\csromanspacer{} เสนาสนํ รูปาวจราย สทฺธาย ฯเปฯ ปญฺญาย อุปนิสฺสยปจฺจเยน ปจฺจโย.\csromanspacer{} ปฐมสฺส ฌานสฺส ปริกมฺมํ ปฐมสฺส ฌานสฺส อุปนิสฺสยปจฺจเยน ปจฺจโย ฯเปฯ.\csromanspacer{} จตุตฺถสฺส ฌานสฺส ฯเปฯ.\csromanspacer{} ทิพฺพสฺส จกฺขุสฺส ปริกมฺมํ ฯเปฯ.\csromanspacer{} ทิพฺพาย โสตธาตุยา ฯเปฯ.\csromanspacer{} อิทฺธิวิธญาณสฺส ฯเปฯ.\csromanspacer{} เจโตปริยญาณสฺส ฯเปฯ.\csromanspacer{} ปุพฺเพนิวาสานุสฺสติญาณสฺส ฯเปฯ.\csromanspacer{} ยถากมฺมูปคญาณสฺส ฯเปฯ.\csromanspacer{} อนาคตํสญาณสฺส ปริกมฺมํ อนาคตํสญาณสฺส อุปนิสฺสยปจฺจเยน ปจฺจโย. (2)}

\csromanheader{2}{\textbf{ปุเรชาตาทิ}}
\proseitem{254}{นรูปาวจโร ธมฺโม นรูปาวจรสฺส ธมฺมสฺส ปุเรชาตปจฺจเยน ปจฺจโย.\csromanspacer{} อารมฺมณปุเรชาตํ, วตฺถุปุเรชาตํ.\csromanspacer{}\csromanspacer{}\textbf{อารมฺมณปุเรชาตํ\csromandash{}}}

\vspace{\tipitakaparskip}
\csromancenter{รูปาวจรทุก จกฺขุํ ฯเปฯ.\csromanspacer{} วตฺถุํ อนิจฺจโต ฯเปฯ โทมนสฺสํ อุปฺปชฺชติ.\csromanspacer{} รูปายตนํ จกฺขุวิญฺญาณสฺส ฯเปฯ.\csromanspacer{} โผฏฺฐพฺพายตนํ กายวิญฺญาณสฺส ปุเรชาตปจฺจเยน ปจฺจโย.\csromanspacer{} \textbf{วตฺถุปุเรชาตํ\csromandash{}}จกฺขายตนํ จกฺขุวิญฺญาณสฺส ฯเปฯ.\csromanspacer{} กายายตนํ กายวิญฺญาณสฺส ปุเรชาตปจฺจเยน ปจฺจโย.\csromanspacer{} วตฺถุ นรูปาวจรานํ ขนฺธานํ ปุเรชาตปจฺจเยน ปจฺจโย.\csromanspacer{}\csromanspacer{}นรูปาวจโร ธมฺโม รูปาวจรสฺส ธมฺมสฺส ปุเรชาตปจฺจเยน ปจฺจโย.\csromanspacer{} อารมฺมณปุเรชาตํ, วตฺถุปุเรชาตํ.\csromanspacer{}\csromanspacer{}\textbf{อารมฺมณปุเรชาตํ\csromandash{}} ทิพฺเพน จกฺขุนา รูปํ ปสฺสติ.\csromanspacer{} ทิพฺพาย โสตธาตุยา สทฺทํ สุณาติ.\csromanspacer{} \textbf{วตฺถุปุเรชาตํ\csromandash{}}วตฺถุ รูปาวจรานํ ขนฺธานํ ปุเรชาตปจฺจเยน ปจฺจโย. (2)}
\prose{\csromanfolio{153}ปจฺฉาชาตปจฺจเยน ปจฺจโย.\csromanspacer{} เทฺว.\csromanspacer{} อาเสวนปจฺจเยน ปจฺจโย.\csromanspacer{} ตีณิ.}

\csromanheader{2}{\textbf{กมฺมาทิ}}
\proseitem{255}{รูปาวจโร ธมฺโม รูปาวจรสฺส ธมฺมสฺส กมฺมปจฺจเยน ปจฺจโย.\csromanspacer{} \textbf{สหชาตา}, นานากฺขณิกา.\csromanspacer{}\csromanspacer{}\textbf{สหชาตา} รูปาวจรา เจตนา สมฺปยุตฺตกานํ ขนฺธานํ กมฺมปจฺจเยน ปจฺจโย.\csromanspacer{} ปฏิสนฺธิกฺขเณ ฯเปฯ.\csromanspacer{} \textbf{นานากฺขณิกา} รูปาวจรา เจตนา วิปากานํ รูปาวจรานํ ขนฺธานํ กมฺมปจฺจเยน ปจฺจโย.\csromanspacer{}\csromanspacer{}รูปาวจโร ธมฺโม นรูปาวจรสฺส ธมฺมสฺส กมฺมปจฺจเยน ปจฺจโย.\csromanspacer{} \textbf{สหชาตา}, นานากฺขณิกา.\csromanspacer{}\csromanspacer{}\textbf{สหชาตา} รูปาวจรา เจตนา จิตฺตสมุฏฺฐานานํ รูปานํ กมฺมปจฺจเยน ปจฺจโย.\csromanspacer{} ปฏิสนฺธิกฺขเณ ฯเปฯ.\csromanspacer{} \textbf{นานากฺขณิกา} รูปาวจรา เจตนา กฏตฺตารูปานํ กมฺมปจฺจเยน ปจฺจโย.\csromanspacer{}\csromanspacer{}รูปาวจโร ธมฺโม รูปาวจรสฺส จ นรูปาวจรสฺส จ ธมฺมสฺส กมฺมปจฺจเยน ปจฺจโย.\csromanspacer{} \textbf{สหชาตา}, นานากฺขณิกา.\csromanspacer{}\csromanspacer{}\textbf{สหชาตา} รูปาวจรา เจตนา สมฺปยุตฺตกานํ ขนฺธานํ จิตฺตสมุฏฺฐานานญฺจ รูปานํ กมฺมปจฺจเยน ปจฺจโย.\csromanspacer{} ปฏิสนฺธิกฺขเณ ฯเปฯ.\csromanspacer{} \textbf{นานากฺขณิกา} รูปาวจรา เจตนา วิปากานํ รูปาวจรานํ ขนฺธานํ กฏตฺตา จ รูปานํ กมฺมปจฺจเยน ปจฺจโย. (3)}

\prose{นรูปาวจโร ธมฺโม นรูปาวจรสฺส ธมฺมสฺส กมฺมปจฺจเยน ปจฺจโย.\csromanspacer{} \textbf{สหชาตา}, นานากฺขณิกา.\csromanspacer{}\csromanspacer{}\textbf{สหชาตา} นรูปาวจรา เจตนา สมฺปยุตฺตกานํ ขนฺธานํ จิตฺตสมุฏฺฐานานญฺจ รูปานํ กมฺมปจฺจเยน ปจฺจโย.\csromanspacer{} ปฏิสนฺธิกฺขเณ ฯเปฯ.\csromanspacer{} \textbf{นานากฺขณิกา} นรูปาวจรา เจตนา วิปากานํ นรูปาวจรานํ ขนฺธานํ กฏตฺตา จ รูปานํ กมฺมปจฺจเยน ปจฺจโย. (1)}

\prose{\csromanfolio{154}วิปากปจฺจเยน ปจฺจโย.\csromanspacer{} จตฺตาริ.\csromanspacer{} อาหารปจฺจเยน ปจฺจโย.\csromanspacer{} จตฺตาริ.\csromanspacer{} อินฺทฺริยปจฺจเยน ปจฺจโย.\csromanspacer{} จตฺตาริ.\csromanspacer{} ฌานปจฺจเยน ปจฺจโย.\csromanspacer{} จตฺตาริ.\csromanspacer{} มคฺคปจฺจเยน ปจฺจโย.\csromanspacer{} จตฺตาริ.\csromanspacer{} สมฺปยุตฺตปจฺจเยน ปจฺจโย.\csromanspacer{} เทฺว.}

\csromanheader{2}{\textbf{วิปฺปยุตฺต}}
\proseitem{256}{รูปาวจโร ธมฺโม นรูปาวจรสฺส ธมฺมสฺส วิปฺปยุตฺตปจฺจเยน ปจฺจโย.\csromanspacer{} \textbf{สหชาตํ}, ปจฺฉาชาตํ. (สํขิตฺตํ.) (1)}

\prose{นรูปาวจโร ธมฺโม นรูปาวจรสฺส ธมฺมสฺส วิปฺปยุตฺตปจฺจเยน ปจฺจโย.\csromanspacer{} \textbf{สหชาตํ}, ปุเรชาตํ, ปจฺฉาชาตํ. (สํขิตฺตํ.).\csromanspacer{} นรูปาวจโร ธมฺโม รูปาวจรสฺส ธมฺมสฺส วิปฺปยุตฺตปจฺจเยน ปจฺจโย.\csromanspacer{} \textbf{สหชาตํ}, ปุเรชาตํ.\csromanspacer{}\csromanspacer{}\textbf{สหชาตํ} ปฏิสนฺธิกฺขเณ วตฺถุ รูปาวจรานํ ขนฺธานํ วิปฺปยุตฺตปจฺจเยน ปจฺจโย.\csromanspacer{} \textbf{ปุเรชาตํ} วตฺถุ รูปาวจรานํ ขนฺธานํ วิปฺปยุตฺตปจฺจเยน ปจฺจโย. (2)}

\csromanheader{2}{\textbf{อตฺถิ-อาทิ}}
\proseitem{257}{รูปาวจโร ธมฺโม รูปาวจรสฺส ธมฺมสฺส อตฺถิปจฺจเยน ปจฺจโย. (ปฏิจฺจสทิสํ.).\csromanspacer{} รูปาวจโร ธมฺโม นรูปาวจรสฺส ธมฺมสฺส อตฺถิปจฺจเยน ปจฺจโย.\csromanspacer{} \textbf{สหชาตํ}, ปุเรชาตํ, ปจฺฉาชาตํ.\csromanspacer{}\csromanspacer{}รูปาวจโร ธมฺโม รูปาวจรสฺส จ นรูปาวจรสฺส จ ธมฺมสฺส อตฺถิปจฺจเยน ปจฺจโย. (ปฏิจฺจสทิสํ.) (3)}

\prose{นรูปาวจโร ธมฺโม นรูปาวจรสฺส ธมฺมสฺส อตฺถิปจฺจเยน ปจฺจโย.\csromanspacer{} \textbf{สหชาตํ}, ปุเรชาตํ, ปจฺฉาชาตํ, อาหารํ, อินฺทฺริยํ. (สํขิตฺตํ.).\csromanspacer{} นรูปาวจโร ธมฺโม รูปาวจรสฺส ธมฺมสฺส อตฺถิปจฺจเยน ปจฺจโย.\csromanspacer{} \textbf{สหชาตํ}, ปุเรชาตํ.\csromanspacer{}\csromanspacer{}\textbf{สหชาตํ} ปฏิสนฺธิกฺขเณ วตฺถุ รูปาวจรานํ ขนฺธานํ อตฺถิปจฺจเยน ปจฺจโย.\csromanspacer{} \textbf{ปุเรชาตํ} ทิพฺเพน จกฺขุนา รูปํ ปสฺสติ.\csromanspacer{} ทิพฺพาย โสตธาตุยา สทฺทํ สุณาติ.\csromanspacer{} วตฺถุ รูปาวจรานํ ขนฺธานํ อตฺถิปจฺจเยน ปจฺจโย. (2)}

\proseitem{258}{รูปาวจโร จ นรูปาวจโร จ ธมฺมา รูปาวจรสฺส ธมฺมสฺส อตฺถิปจฺจเยน ปจฺจโย.\csromanspacer{} \textbf{สหชาตํ}, ปุเรชาตํ.\csromanspacer{}\csromanspacer{}สหชาโต รูปาวจโร เอโก ขนฺโธ จ วตฺถุ จ ติณฺณนฺนํ ขนฺธานํ อตฺถิปจฺจเยน ปจฺจโย ฯเปฯ เทฺว ขนฺธา จ ฯเปฯ.\csromanspacer{} ปฏิสนฺธิกฺขเณ ฯเปฯ.\csromanspacer{} รูปาวจโร จ นรูปาวจโร จ ธมฺมา}

\vspace{\tipitakaparskip}
\csromancenter{รูปาวจรทุก นรูปาวจรสฺส ธมฺมสฺส อตฺถิปจฺจเยน ปจฺจโย.\csromanspacer{} สหชาตํ, ปจฺฉาชาตํ, อาหารํ, อินฺทฺริยํ.\csromanspacer{}\csromanspacer{}\textbf{สหชาตา} รูปาวจรา ขนฺธา จ มหาภูตา จ จิตฺตสมุฏฺฐานานํ รูปานํ อตฺถิปจฺจเยน ปจฺจโย.\csromanspacer{} ปฏิสนฺธิกฺขเณ ฯเปฯ.\csromanspacer{} \textbf{ปจฺฉาชาตา} รูปาวจรา ขนฺธา จ กพฬีกาโร อาหาโร จ ปุเรชาตสฺส อิมสฺส กายสฺส อตฺถิปจฺจเยน ปจฺจโย.\csromanspacer{} \textbf{ปจฺฉาชาตา} รูปาวจรา ขนฺธา จ รูปาชีวิตินฺทฺริยญฺจ กฏตฺตารูปานํ อตฺถิปจฺจเยน ปจฺจโย. (2)}
\prose{\csromanfolio{155}นตฺถิปจฺจเยน ปจฺจโย.\csromanspacer{} วิคตปจฺจเยน ปจฺจโย.\csromanspacer{} อวิคตปจฺจเยน ปจฺจโย.}

\csromanheader{2}{1. ปจฺจยานุโลม 2. สงฺขฺยาวาร}
\csromanheader{2}{\textbf{สุทฺธ}}
\proseitem{259}{เหตุยา จตฺตาริ, อารมฺมเณ จตฺตาริ, อธิปติยา จตฺตาริ, อนนฺตเร จตฺตาริ, สมนนฺตเร จตฺตาริ, สหชาเต สตฺต, อญฺญมญฺเญ ฉ, นิสฺสเย สตฺต, อุปนิสฺสเย จตฺตาริ, ปุเรชาเต เทฺว, ปจฺฉาชาเต เทฺว, อาเสวเน ตีณิ, กมฺเม จตฺตาริ, วิปาเก จตฺตาริ, อาหาเร จตฺตาริ, อินฺทฺริเย จตฺตาริ, ฌาเน จตฺตาริ, มคฺเค จตฺตาริ, สมฺปยุตฺเต เทฺว, วิปฺปยุตฺเต ตีณิ, อตฺถิยา สตฺต, นตฺถิยา จตฺตาริ, วิคเต จตฺตาริ, อวิคเต สตฺต.}

\csromanheader{2}{2. ปจฺจนียุทฺธาร}
\proseitem{260}{รูปาวจโร ธมฺโม รูปาวจรสฺส ธมฺมสฺส อารมฺมณปจฺจเยน ปจฺจโย.\csromanspacer{} สหชาตปจฺจเยน ปจฺจโย.\csromanspacer{} อุปนิสฺสยปจฺจเยน ปจฺจโย.\csromanspacer{} รูปาวจโร ธมฺโม นรูปาวจรสฺส ธมฺมสฺส อารมฺมณปจฺจเยน ปจฺจโย.\csromanspacer{} สหชาตปจฺจเยน ปจฺจโย.\csromanspacer{} อุปนิสฺสยปจฺจเยน ปจฺจโย.\csromanspacer{} ปจฺฉาชาตปจฺจเยน ปจฺจโย.\csromanspacer{} กมฺมปจฺจเยน ปจฺจโย.\csromanspacer{}\csromanspacer{}รูปาวจโร ธมฺโม รูปาวจรสฺส จ นรูปาวจรสฺส จ ธมฺมสฺส สหชาตปจฺจเยน ปจฺจโย.\csromanspacer{} กมฺมปจฺจเยน ปจฺจโย. (3)}

\proseitem{261}{นรูปาวจโร ธมฺโม นรูปาวจรสฺส ธมฺมสฺส อารมฺมณปจฺจเยน ปจฺจโย.\csromanspacer{} สหชาตปจฺจเยน ปจฺจโย.\csromanspacer{} อุปนิสฺสยปจฺจเยน ปจฺจโย. \csromanfolio{156}ปุเรชาตปจฺจเยน ปจฺจโย.\csromanspacer{} ปจฺฉาชาตปจฺจเยน ปจฺจโย.\csromanspacer{} กมฺมปจฺจเยน ปจฺจโย.\csromanspacer{} อาหารปจฺจเยน ปจฺจโย.\csromanspacer{} อินฺทฺริยปจฺจเยน ปจฺจโย.\csromanspacer{}\csromanspacer{}นรูปาวจโร ธมฺโม รูปาวจรสฺส ธมฺมสฺส อารมฺมณปจฺจเยน ปจฺจโย.\csromanspacer{} สหชาตปจฺจเยน ปจฺจโย.\csromanspacer{} อุปนิสฺสยปจฺจเยน ปจฺจโย.\csromanspacer{} ปุเรชาตปจฺจเยน ปจฺจโย. (2)}

\prose{รูปาวจโร จ นรูปาวจโร จ ธมฺมา รูปาวจรสฺส ธมฺมสฺส สหชาตํ, ปุเรชาตํ.\csromanspacer{}\csromanspacer{}รูปาวจโร จ นรูปาวจโร จ ธมฺมา นรูปาวจรสฺส ธมฺมสฺส สหชาตํ, ปจฺฉาชาตํ, อาหารํ, อินฺทฺริยํ. (2)}

\csromanheader{2}{2. ปจฺจยปจฺจนีย 2. สงฺขฺยาวาร}
\proseitem{262}{นเหตุยา สตฺต, น-อารมฺมเณ สตฺต, น-อธิปติยา สตฺต, น-อนนฺตเร สตฺต, นสมนนฺตเร สตฺต, นสหชาเต ฉ, น-อญฺญมญฺเญ ฉ, นนิสฺสเย ฉ, น-อุปนิสฺสเย สตฺต, นปุเรชาเต สตฺต, นปจฺฉาชาเต สตฺต ฯเปฯ นสมฺปยุตฺเต ฉ, นวิปฺปยุตฺเต ปญฺจ, โน-อตฺถิยา ปญฺจ, โนนตฺถิยา สตฺต, โนวิคเต สตฺต, โน-อวิคเต ปญฺจ.}

\csromanheader{2}{3. ปจฺจยานุโลมปจฺจนีย}
\proseitem{263}{เหตุปจฺจยา น-อารมฺมเณ จตฺตาริ, น-อธิปติยา จตฺตาริ, น-อนนฺตเร นสมนนฺตเร จตฺตาริ, น-อญฺญมญฺเญ เทฺว, น-อุปนิสฺสเย จตฺตาริ ฯเปฯ นสมฺปยุตฺเต เทฺว, นวิปฺปยุตฺเต เทฺว, โนนตฺถิยา จตฺตาริ, โนวิคเต จตฺตาริ.}

\csromanheader{2}{4. ปจฺจยปจฺจนียานุโลม}
\proseitem{264}{นเหตุปจฺจยา อารมฺมเณ จตฺตาริ, อธิปติยา จตฺตาริ, (อนุโลมมาติกา กาตพฺพา.) ฯเปฯ อวิคเต สตฺต.}

\nitthitam{รูปาวจรทุกํ นิฏฺฐิตํ.}
\csromanmatikaanchor{mtk.24}
\csromantocmarkref{section}{95. อรูปาวจรทุก}{mtk.24}
\bookmark[dest={mtk.24},level=1]{95. อรูปาวจรทุก}
\csromanheader{1}{95. อรูปาวจรทุก \textbf{95. อรูปาวจรทุก 1. ปฏิจฺจวาร}}
\proseitem{265}{\csromanfolio{157}อรูปาวจรํ ธมฺมํ ปฏิจฺจ อรูปาวจโร ธมฺโม อุปฺปชฺชติ เหตุปจฺจยา.\csromanspacer{} อรูปาวจรํ เอกํ ขนฺธํ ปฏิจฺจ ตโย ขนฺธา ฯเปฯ เทฺว ขนฺเธ ฯเปฯ.\csromanspacer{} ปฏิสนฺธิกฺขเณ ฯเปฯ.\csromanspacer{} อรูปาวจรํ ธมฺมํ ปฏิจฺจ น-อรูปาวจโร ธมฺโม อุปฺปชฺชติ เหตุปจฺจยา.\csromanspacer{} อรูปาวจเร ขนฺเธ ปฏิจฺจ จิตฺตสมุฏฺฐานํ รูปํ.\csromanspacer{}\csromanspacer{}อรูปาวจรํ ธมฺมํ ปฏิจฺจ อรูปาวจโร จ น-อรูปาวจโร จ ธมฺมา อุปฺปชฺชนฺติ เหตุปจฺจยา.\csromanspacer{} อรูปาวจรํ เอกํ ขนฺธํ ปฏิจฺจ ตโย ขนฺธา จิตฺตสมุฏฺฐานญฺจ รูปํ ฯเปฯ เทฺว ขนฺเธ ฯเปฯ. (3)}

\proseitem{266}{น-อรูปาวจรํ ธมฺมํ ปฏิจฺจ น-อรูปาวจโร ธมฺโม อุปฺปชฺชติ เหตุปจฺจยา.\csromanspacer{} น-อรูปาวจรํ เอกํ ขนฺธํ ปฏิจฺจ ตโย ขนฺธา จิตฺตสมุฏฺฐานญฺจ รูปํ ฯเปฯ เทฺว ขนฺเธ ฯเปฯ.\csromanspacer{} ปฏิสนฺธิกฺขเณ น-อรูปาวจรํ เอกํ ขนฺธํ ปฏิจฺจ ตโย ขนฺธา กฏตฺตา จ รูปํ ฯเปฯ เทฺว ขนฺเธ ฯเปฯ.\csromanspacer{} ขนฺเธ ปฏิจฺจ วตฺถุ, วตฺถุํ ปฏิจฺจ ขนฺธา.\csromanspacer{} เอกํ มหาภูตํ ฯเปฯ. (1)}

\prose{อรูปาวจรญฺจ น-อรูปาวจรญฺจ ธมฺมํ ปฏิจฺจ น-อรูปาวจโร ธมฺโม อุปฺปชฺชติ เหตุปจฺจยา.\csromanspacer{} อรูปาวจเร ขนฺเธ จ มหาภูเต จ ปฏิจฺจ จิตฺตสมุฏฺฐานํ รูปํ. (1) (สํขิตฺตํ.)}

\csromanheader{2}{1. ปจฺจยานุโลม 2. สงฺขฺยาวาร}
\vspace{\tipitakaparskip}
\csromancenter{\textbf{เหตุ}ยา ปญฺจ, อารมฺมเณ เทฺว, อธิปติยา ปญฺจ, อนนฺตเร เทฺว, สมนนฺตเร เทฺว, สหชาเต ปญฺจ, อญฺญมญฺเญ เทฺว, นิสฺสเย ปญฺจ, อุปนิสฺสเย เทฺว, ปุเรชาเต เทฺว, อาเสวเน เทฺว, กมฺเม ปญฺจ, วิปาเก เทฺว, อาหาเร ปญฺจ ฯเปฯ อวิคเต ปญฺจ.}
\csromanheader{2}{2. ปจฺจยปจฺจนีย 1. วิภงฺควาร}
\csromanheader{2}{\textbf{นเหตุ น-อารมฺมณ}}
\proseitem{268}{น-อรูปาวจรํ ธมฺมํ ปฏิจฺจ น-อรูปาวจโร ธมฺโม อุปฺปชฺชติ นเหตุปจฺจยา.\csromanspacer{} อเหตุกํ น-อรูปาวจรํ เอกํ ขนฺธํ ปฏิจฺจ ตโย ขนฺธา จิตฺตสมุฏฺฐานญฺจ \csromanfolio{158}รูปํ ฯเปฯ เทฺว ขนฺเธ ฯเปฯ.\csromanspacer{} อเหตุกปฏิสนฺธิกฺขเณ ฯเปฯ. (ยาว อสญฺญสตฺตา.) วิจิกิจฺฉาสหคเต อุทฺธจฺจสหคเต ขนฺเธ ปฏิจฺจ วิจิกิจฺฉาสหคโต อุทฺธจฺจสหคโต โมโห. (1).\csromanspacer{} น-อารมฺมณปจฺจยา.\csromanspacer{} ตีณิ.}

\csromanheader{2}{\textbf{น-อธิปตฺยาทิ}}
\proseitem{269}{อรูปาวจรํ ธมฺมํ ปฏิจฺจ อรูปาวจโร ธมฺโม อุปฺปชฺชติ น-อธิปติปจฺจยา.\csromanspacer{} อรูปาวจเร ขนฺเธ ปฏิจฺจ อรูปาวจราธิปติ.\csromanspacer{} วิปากํ อรูปาวจรํ เอกํ ขนฺธํ ปฏิจฺจ ตโย ขนฺธา ฯเปฯ เทฺว ขนฺเธ ฯเปฯ.\csromanspacer{} ปฏิสนฺธิกฺขเณ ฯเปฯ. (1)}

\prose{น-อรูปาวจรํ ธมฺมํ ปฏิจฺจ น-อรูปาวจโร ธมฺโม อุปฺปชฺชติ น-อธิปติปจฺจยา.\csromanspacer{} น-อรูปาวจรํ เอกํ ขนฺธํ ปฏิจฺจ ตโย ขนฺธา จิตฺตสมุฏฺฐานญฺจ รูปํ ฯเปฯ เทฺว ขนฺเธ ฯเปฯ.\csromanspacer{} ปฏิสนฺธิกฺขเณ ฯเปฯ. (ยาว อสญฺญสตฺตา.) (1).\csromanspacer{} น-อนนฺตรปจฺจยา ฯเปฯ.\csromanspacer{} น-อุปนิสฺสยปจฺจยา.}

\csromanheader{2}{\textbf{นปุเรชาตาทิ}}
\proseitem{270}{อรูปาวจรํ ธมฺมํ ปฏิจฺจ อรูปาวจโร ธมฺโม อุปฺปชฺชติ นปุเรชาตปจฺจยา.\csromanspacer{} อรูเป อรูปาวจรํ เอกํ ขนฺธํ ปฏิจฺจ ตโย ขนฺธา ฯเปฯ เทฺว ขนฺเธ ฯเปฯ. (มูลํ กาตพฺพํ.) อรูปาวจเร ขนฺเธ ปฏิจฺจ จิตฺตสมุฏฺฐานํ รูปํ. (2)}

\prose{น-อรูปาวจรํ ธมฺมํ ปฏิจฺจ น-อรูปาวจโร ธมฺโม อุปฺปชฺชติ นปุเรชาตปจฺจยา.\csromanspacer{} อรูเป น-อรูปาวจรํ เอกํ ขนฺธํ ปฏิจฺจ ตโย ขนฺธา ฯเปฯ เทฺว ขนฺเธ ฯเปฯ.\csromanspacer{} ปฏิสนฺธิกฺขเณ ฯเปฯ. (ยาว อสญฺญสตฺตา.) (1)}

\prose{อรูปาวจรญฺจ น-อรูปาวจรญฺจ ธมฺมํ ปฏิจฺจ น-อรูปาวจโร ธมฺโม อุปฺปชฺชติ นปุเรชาตปจฺจยา.\csromanspacer{} อรูปาวจเร ขนฺเธ จ มหาภูเต จ ปฏิจฺจ จิตฺตสมุฏฺฐานํ รูปํ. (1).\csromanspacer{} นปจฺฉาชาตปจฺจยา.}

\csromanheader{2}{\textbf{น-อาเสวน}}
\proseitem{271}{อรูปาวจรํ ธมฺมํ ปฏิจฺจ อรูปาวจโร ธมฺโม อุปฺปชฺชติ น-อาเสวนปจฺจยา.\csromanspacer{} วิปากํ อรูปาวจรํ เอกํ ขนฺธํ ปฏิจฺจ ตโย ขนฺธา ฯเปฯ เทฺว ขนฺเธ ฯเปฯ.\csromanspacer{} ปฏิสนฺธิกฺขเณ ฯเปฯ.\csromanspacer{} อรูปาวจรํ ธมฺมํ ปฏิจฺจ น-อรูปาวจโร ธมฺโม อุปฺปชฺชติ น-อาเสวนปจฺจยา.\csromanspacer{} อรูปาวจเร ขนฺเธ ปฏิจฺจ จิตฺตสมุฏฺฐานํ รูปํ. (2)}

\csromanheader{2}{95. อรูปาวจรทุก}
\prose{\csromanfolio{159}น-อรูปาวจรํ ธมฺมํ ปฏิจฺจ น-อรูปาวจโร ธมฺโม อุปฺปชฺชติ น-อาเสวนปจฺจยา.\csromanspacer{} น-อรูปาวจรํ เอกํ ขนฺธํ ปฏิจฺจ ตโย ขนฺธา จิตฺตสมุฏฺฐานญฺจ รูปํ ฯเปฯ เทฺว ขนฺเธ ฯเปฯ. (ยาว อสญฺญสตฺตา.) (1)}

\prose{อรูปาวจรญฺจ น-อรูปาวจรญฺจ ธมฺมํ ปฏิจฺจ น-อรูปาวจโร ธมฺโม อุปฺปชฺชติ น-อาเสวนปจฺจยา.\csromanspacer{} อรูปาวจเร ขนฺเธ จ มหาภูเต จ ปฏิจฺจ จิตฺตสมุฏฺฐานํ รูปํ. (1) (สํขิตฺตํ.)}

\csromanheader{2}{2. ปจฺจยปจฺจนีย 2. สงฺขฺยาวาร}
\proseitem{272}{นเหตุยา เอกํ, น-อารมฺมเณ ตีณิ, น-อธิปติยา เทฺว, น-อนนฺตเร ตีณิ, นสมนนฺตเร ตีณิ, น-อญฺญมญฺเญ ตีณิ, น-อุปนิสฺสเย ตีณิ, นปุเรชาเต จตฺตาริ, นปจฺฉาชาเต ปญฺจ, น-อาเสวเน จตฺตาริ, นกมฺเม เทฺว, นวิปาเก ปญฺจ, น-อาหาเร เอกํ, น-อินฺทฺริเย เอกํ, นฌาเน เอกํ, นมคฺเค เอกํ, นสมฺปยุตฺเต ตีณิ, นวิปฺปยุตฺเต เทฺว, โนนตฺถิยา ตีณิ, โนวิคเต ตีณิ.}

\vspace{\tipitakaparskip}
\csromancenter{(เอวํ อิตเร เทฺว คณนาปิ สหชาตวาโรปิ กาตพฺโพ.)}
\csromanheader{2}{95. อรูปาวจรทุก 3. ปจฺจยวาร}
\proseitem{273}{อรูปาวจรํ ธมฺมํ ปจฺจยา อรูปาวจโร ธมฺโม อุปฺปชฺชติ เหตุปจฺจยา.\csromanspacer{} ตีณิ. (ปฏิจฺจสทิสา.)}

\prose{น-อรูปาวจรํ ธมฺมํ ปจฺจยา น-อรูปาวจโร ธมฺโม อุปฺปชฺชติ เหตุปจฺจยา.\csromanspacer{} น-อรูปาวจรํ เอกํ ขนฺธํ ฯเปฯ. (ยาว อชฺฌตฺติกา มหาภูตา.) วตฺถุํ ปจฺจยา น-อรูปาวจรา ขนฺธา.\csromanspacer{}\csromanspacer{}น-อรูปาวจรํ ธมฺมํ ปจฺจยา อรูปาวจโร ธมฺโม อุปฺปชฺชติ เหตุปจฺจยา.\csromanspacer{} วตฺถุํ ปจฺจยา อรูปาวจรา ขนฺธา.\csromanspacer{}\csromanspacer{}น-อรูปาวจรํ ธมฺมํ ปจฺจยา อรูปาวจโร จ น-อรูปาวจโร จ ธมฺมา อุปฺปชฺชนฺติ เหตุปจฺจยา.\csromanspacer{} วตฺถุํ ปจฺจยา อรูปาวจรา ขนฺธา.\csromanspacer{} มหาภูเต ปจฺจยา จิตฺตสมุฏฺฐานํ รูปํ. (3)}

\proseitem{274}{\csromanfolio{160}อรูปาวจรญฺจ น-อรูปาวจรญฺจ ธมฺมํ ปจฺจยา อรูปาวจโร ธมฺโม อุปฺปชฺชติ เหตุปจฺจยา.\csromanspacer{} อรูปาวจรํ เอกํ ขนฺธญฺจ วตฺถุญฺจ ปจฺจยา ตโย ขนฺธา ฯเปฯ เทฺว ขนฺเธ จ ฯเปฯ.\csromanspacer{} อรูปาวจรญฺจ น-อรูปาวจรญฺจ ธมฺมํ ปจฺจยา น-อรูปาวจโร ธมฺโม อุปฺปชฺชติ เหตุปจฺจยา.\csromanspacer{} อรูปาวจเร ขนฺเธ จ มหาภูเต จ ปจฺจยา จิตฺตสมุฏฺฐานํ รูปํ.\csromanspacer{}\csromanspacer{}อรูปาวจรญฺจ น-อรูปาวจรญฺจ ธมฺมํ ปจฺจยา อรูปาวจโร จ น-อรูปาวจโร จ ธมฺมา อุปฺปชฺชนฺติ เหตุปจฺจยา.\csromanspacer{} อรูปาวจรํ เอกํ ขนฺธญฺจ วตฺถุญฺจ ปจฺจยา ตโย ขนฺธา ฯเปฯ เทฺว ขนฺเธ จ ฯเปฯ.\csromanspacer{} อรูปาวจเร ขนฺเธ จ มหาภูเต จ ปจฺจยา จิตฺตสมุฏฺฐานํ รูปํ. (3)}

\csromanheader{2}{1. ปจฺจยานุโลม 2. สงฺขฺยาวาร}
\proseitem{275}{เหตุยา นว, อารมฺมเณ จตฺตาริ, อธิปติยา นว, อนนฺตเร จตฺตาริ, สมนนฺตเร จตฺตาริ, สหชาเต นว, อญฺญมญฺเญ จตฺตาริ, นิสฺสเย นว, อุปนิสฺสเย จตฺตาริ, ปุเรชาเต จตฺตาริ, อาเสวเน จตฺตาริ, กมฺเม นว, วิปาเก เทฺว ฯเปฯ อวิคเต นว.}

\csromanheader{2}{2. ปจฺจยปจฺจนีย 1. วิภงฺควาร}
\csromanheader{2}{\textbf{นเหตฺวาทิ}}
\proseitem{276}{น-อรูปาวจรํ ธมฺมํ ปจฺจยา น-อรูปาวจโร ธมฺโม อุปฺปชฺชติ นเหตุปจฺจยา.\csromanspacer{} อเหตุกํ น-อรูปาวจรํ เอกํ ขนฺธํ ฯเปฯ. (ยาว อสญฺญสตฺตา.) จกฺขายตนํ ปจฺจยา จกฺขุวิญฺญาณํ ฯเปฯ.\csromanspacer{} กายายตนํ ปจฺจยา กายวิญฺญาณํ.\csromanspacer{} วตฺถุํ ปจฺจยา อเหตุกา น-อรูปาวจรา ขนฺธา.\csromanspacer{} วิจิกิจฺฉาสหคเต อุทฺธจฺจสหคเต ขนฺเธ จ วตฺถุญฺจ ปจฺจยา วิจิกิจฺฉาสหคโต อุทฺธจฺจสหคโต โมโห. (1).\csromanspacer{} น-อารมฺมณปจฺจยา.\csromanspacer{} ตีณิ.}

\prose{อรูปาวจรํ ธมฺมํ ปจฺจยา อรูปาวจโร ธมฺโม อุปฺปชฺชติ น-อธิปติปจฺจยา.\csromanspacer{} อรูปาวจเร ขนฺเธ ปจฺจยา อรูปาวจราธิปติ.\csromanspacer{} วิปากํ อรูปาวจรํ เอกํ ขนฺธํ ฯเปฯ.\csromanspacer{} ปฏิสนฺธิกฺขเณ ฯเปฯ. (1)}

\prose{น-อรูปาวจรํ ธมฺมํ ปจฺจยา น-อรูปาวจโร ธมฺโม อุปฺปชฺชติ น-อธิปติปจฺจยา.\csromanspacer{} น-อรูปาวจรํ เอกํ ขนฺธํ ฯเปฯ.\csromanspacer{} ปฏิสนฺธิกฺขเณ ฯเปฯ. (ยาว อสญฺญสตฺตา.) จกฺขายตนํ ปจฺจยา จกฺขุวิญฺญาณํ ฯเปฯ.\csromanspacer{} กายายตนํ ปจฺจยา กายวิญฺญาณํ.\csromanspacer{} วตฺถุํ ปจฺจยา น-อรูปาวจรา ขนฺธา. (1)}

\csromanheader{2}{95. อรูปาวจรทุก \textbf{2. ปจฺจยปจฺจนีย 2. สงฺขฺยาวาร}}
\proseitem{277}{\csromanfolio{161}นเหตุยา เอกํ, น-อารมฺมเณ ตีณิ, น-อธิปติยา จตฺตาริ, น-อนนฺตเร ตีณิ, นสมนนฺตเร น-อญฺญมญฺเญ น-อุปนิสฺสเย ตีณิ, นปุเรชาเต จตฺตาริ, นปจฺฉาชาเต นว, น-อาเสวเน จตฺตาริ, นกมฺเม จตฺตาริ, นวิปาเก นว, น-อาหาเร เอกํ, น-อินฺทฺริเย เอกํ, นฌาเน เอกํ, นมคฺเค เอกํ, นสมฺปยุตฺเต ตีณิ, นวิปฺปยุตฺเต เทฺว, โนนตฺถิยา ตีณิ, โนวิคเต ตีณิ.}

\vspace{\tipitakaparskip}
\csromancenter{(เอวํ อิตเร เทฺว คณนาปิ นิสฺสยวาโรปิ กาตพฺโพ.)}
\csromanheader{2}{95. อรูปาวจรทุก 5. สํสฏฺฐวาร}
\csromanheader{2}{1. ปจฺจยานุโลมาทิ}
\proseitem{278}{อรูปาวจรํ ธมฺมํ สํสฏฺโฐ อรูปาวจโร ธมฺโม อุปฺปชฺชติ เหตุปจฺจยา.\csromanspacer{} อรูปาวจรํ เอกํ ขนฺธํ สํสฏฺฐา ตโย ขนฺธา ฯเปฯ เทฺว ขนฺเธ ฯเปฯ.\csromanspacer{} ปฏิสนฺธิกฺขเณ ฯเปฯ. (1)}

\prose{น-อรูปาวจรํ ธมฺมํ สํสฏฺโฐ น-อรูปาวจโร ธมฺโม อุปฺปชฺชติ เหตุปจฺจยา.\csromanspacer{} น-อรูปาวจรํ เอกํ ขนฺธํ สํสฏฺฐา ตโย ขนฺธา ฯเปฯ เทฺว ขนฺเธ ฯเปฯ.\csromanspacer{} ปฏิสนฺธิกฺขเณ ฯเปฯ. (1)}

\csromanheader{2}{\textbf{เหตุ}ยา เทฺว ฯเปฯ อวิคเต เทฺว. (อนุโลมํ.)}
\prose{นเหตุยา เอกํ, น-อธิปติยา เทฺว, นปุเรชาเต เทฺว, นปจฺฉาชาเต เทฺว, น-อาเสวเน เทฺว, นกมฺเม เทฺว, นวิปาเก เทฺว, นฌาเน เอกํ, นมคฺเค เอกํ, นวิปฺปยุตฺเต เทฺว. (ปจฺจนียํ.)}

\vspace{\tipitakaparskip}
\csromancenter{(เอวํ อิตเร เทฺว คณนาปิ สมฺปยุตฺตวาโรปิ กาตพฺโพ.)}
\csromanheader{2}{95. อรูปาวจรทุก 7. ปญฺหาวาร}
\proseitem{279}{\csromanfolio{162}อรูปาวจโร ธมฺโม อรูปาวจรสฺส ธมฺมสฺส เหตุปจฺจเยน ปจฺจโย.\csromanspacer{} อรูปาวจรา เหตู สมฺปยุตฺตกานํ ขนฺธานํ เหตุปจฺจเยน ปจฺจโย.\csromanspacer{} ปฏิสนฺธิกฺขเณ ฯเปฯ. (มูลํ กาตพฺพํ.) อรูปาวจรา เหตู จิตฺตสมุฏฺฐานานํ รูปานํ เหตุปจฺจเยน ปจฺจโย. (มูลํ กาตพฺพํ.) อรูปาวจรา เหตู สมฺปยุตฺตกานํ ขนฺธานํ จิตฺตสมุฏฺฐานานญฺจ รูปานํ เหตุปจฺจเยน ปจฺจโย. (3)}

\prose{น-อรูปาวจโร ธมฺโม น-อรูปาวจรสฺส ธมฺมสฺส เหตุปจฺจเยน ปจฺจโย.\csromanspacer{} น-อรูปาวจรา เหตู สมฺปยุตฺตกานํ ขนฺธานํ จิตฺตสมุฏฺฐานานญฺจ รูปานํ เหตุปจฺจเยน ปจฺจโย.\csromanspacer{} ปฏิสนฺธิกฺขเณ ฯเปฯ. (1)}

\csromanheader{2}{\textbf{อารมฺมณ}}
\proseitem{280}{อรูปาวจโร ธมฺโม อรูปาวจรสฺส ธมฺมสฺส อารมฺมณปจฺจเยน ปจฺจโย.\csromanspacer{} อากาสานญฺจายตนํ วิญฺญาณญฺจายตนสฺส อารมฺมณปจฺจเยน ปจฺจโย.\csromanspacer{} อากิญฺจญฺญายตนํ เนวสญฺญานาสญฺญายตนสฺส อารมฺมณปจฺจเยน ปจฺจโย.\csromanspacer{}\csromanspacer{}อรูปาวจโร ธมฺโม น-อรูปาวจรสฺส ธมฺมสฺส อารมฺมณปจฺจเยน ปจฺจโย.\csromanspacer{} อากาสานญฺจายตนํ ปจฺจเวกฺขติ.\csromanspacer{} วิญฺญาณญฺจายตนํ ฯเปฯ.\csromanspacer{} อากิญฺจญฺญายตนํ ฯเปฯ.\csromanspacer{} เนวสญฺญานาสญฺญายตนํ ปจฺจเวกฺขติ.\csromanspacer{} อรูปาวจเร ขนฺเธ อนิจฺจโต ฯเปฯ โทมนสฺสํ อุปฺปชฺชติ.\csromanspacer{} เจโตปริยญาเณน อรูปาวจรจิตฺตสมงฺคิสฺส จิตฺตํ ชานาติ.\csromanspacer{} อรูปาวจรา ขนฺธา เจโตปริยญาณสฺส, ปุพฺเพนิวาสานุสฺสติญาณสฺส, ยถากมฺมูปคญาณสฺส, อนาคตํสญาณสฺส, อาวชฺชนาย อารมฺมณปจฺจเยน ปจฺจโย. (2)}

\prose{น-อรูปาวจโร ธมฺโม น-อรูปาวจรสฺส ธมฺมสฺส อารมฺมณปจฺจเยน ปจฺจโย.\csromanspacer{} ทานํ ฯเปฯ สีลํ ฯเปฯ อุโปสถกมฺมํ กตฺวา ตํ ปจฺจเวกฺขติ, อสฺสาเทติ อภินนฺทติ, ตํ อารพฺภ ราโค ฯเปฯ โทมนสฺสํ อุปฺปชฺชติ.\csromanspacer{} ปุพฺเพ สุจิณฺณานิ ฯเปฯ.\csromanspacer{} ฌานา ฯเปฯ.\csromanspacer{} อริยา มคฺคา วุฏฺฐหิตฺวา มคฺคํ ปจฺจเวกฺขนฺติ, ผลํ}

\vspace{\tipitakaparskip}
\csromancenter{อรูปาวจรทุก ปจฺจเวกฺขนฺติ, นิพฺพานํ ปจฺจเวกฺขนฺติ.\csromanspacer{} นิพฺพานํ โคตฺรภุสฺส, โวทานสฺส, มคฺคสฺส, ผลสฺส, อาวชฺชนาย อารมฺมณปจฺจเยน ปจฺจโย.\csromanspacer{} อริยา ปหีเน กิเลเส ปจฺจเวกฺขนฺติ, วิกฺขมฺภิเต กิเลเส ฯเปฯ.\csromanspacer{} ปุพฺเพ ฯเปฯ.\csromanspacer{} จกฺขุํ ฯเปฯ.\csromanspacer{} วตฺถุํ น-อรูปาวจเร ขนฺเธ อนิจฺจโต ฯเปฯ โทมนสฺสํ อุปฺปชฺชติ.\csromanspacer{} ทิพฺเพน จกฺขุนา รูปํ ปสฺสติ.\csromanspacer{} ทิพฺพาย โสตธาตุยา สทฺทํ สุณาติ.\csromanspacer{} เจโตปริยญาเณน น-อรูปาวจรจิตฺตสมงฺคิสฺส จิตฺตํ ชานาติ.\csromanspacer{} รูปายตนํ จกฺขุวิญฺญาณสฺส ฯเปฯ.\csromanspacer{} โผฏฺฐพฺพายตนํ กายวิญฺญาณสฺส ฯเปฯ.\csromanspacer{} น-อรูปาวจรา ขนฺธา อิทฺธิวิธญาณสฺส, เจโตปริยญาณสฺส, ปุพฺเพนิวาสานุสฺสติญาณสฺส, ยถากมฺมูปคญาณสฺส, อนาคตํสญาณสฺส, อาวชฺชนาย อารมฺมณปจฺจเยน ปจฺจโย. (1)}
\csromanheader{2}{\textbf{อธิปติ}}
\proseitem{281}{\csromanfolio{163}อรูปาวจโร ธมฺโม อรูปาวจรสฺส ธมฺมสฺส อธิปติปจฺจเยน ปจฺจโย.\csromanspacer{} \textbf{สหชาตาธิปติ\csromandash{}}อรูปาวจราธิปติ สมฺปยุตฺตกานํ ขนฺธานํ อธิปติปจฺจเยน ปจฺจโย.\csromanspacer{}\csromanspacer{}อรูปาวจโร ธมฺโม น-อรูปาวจรสฺส ธมฺมสฺส อธิปติปจฺจเยน ปจฺจโย.\csromanspacer{} อารมฺมณาธิปติ, สหชาตาธิปติ.\csromanspacer{}\csromanspacer{}\textbf{อารมฺมณาธิปติ\csromandash{}} อากาสานญฺจายตนํ ครุํ กตฺวา ปจฺจเวกฺขติ ฯเปฯ.\csromanspacer{} เนวสญฺญานาสญฺญายตนํ ครุํ กตฺวา ปจฺจเวกฺขติ.\csromanspacer{} อรูปาวจเร ขนฺเธ ครุํ กตฺวา อสฺสาเทติ อภินนฺทติ, ตํ ครุํ กตฺวา ราโค อุปฺปชฺชติ.\csromanspacer{} ทิฏฺฐิ อุปฺปชฺชติ.\csromanspacer{} \textbf{สหชาตาธิปติ\csromandash{}}อรูปาวจราธิปติ จิตฺตสมุฏฺฐานานํ รูปานํ อธิปติปจฺจเยน ปจฺจโย. (มูลํ กาตพฺพํ.) อรูปาวจราธิปติ สมฺปยุตฺตกานํ ขนฺธานํ จิตฺตสมุฏฺฐานานญฺจ รูปานํ อธิปติปจฺจเยน ปจฺจโย. (3)}

\proseitem{282}{น-อรูปาวจโร ธมฺโม น-อรูปาวจรสฺส ธมฺมสฺส อธิปติปจฺจเยน ปจฺจโย.\csromanspacer{} อารมฺมณาธิปติ, สหชาตาธิปติ.\csromanspacer{}\csromanspacer{}\textbf{อารมฺมณาธิปติ\csromandash{}}ทานํ ฯเปฯ สีลํ ฯเปฯ อุโปสถกมฺมํ กตฺวา ตํ ครุํ กตฺวา ปจฺจเวกฺขติ, อสฺสาเทติ อภินนฺทติ, ตํ ครุํ กตฺวา ราโค อุปฺปชฺชติ.\csromanspacer{} ทิฏฺฐิ อุปฺปชฺชติ.\csromanspacer{} ปุพฺเพ สุจิณฺณานิ ฯเปฯ.\csromanspacer{} ฌานา ฯเปฯ.\csromanspacer{} อริยา มคฺคา วุฏฺฐหิตฺวา มคฺคํ ครุํ กตฺวา ปจฺจเวกฺขนฺติ, ผลํ ฯเปฯ.\csromanspacer{} นิพฺพานํ ครุํ กตฺวา ปจฺจเวกฺขนฺติ.\csromanspacer{} นิพฺพานํ โคตฺรภุสฺส, โวทานสฺส, มคฺคสฺส, ผลสฺส อธิปติปจฺจเยน ปจฺจโย.\csromanspacer{} จกฺขุํ ฯเปฯ.\csromanspacer{} วตฺถุํ น-อรูปาวจเร ขนฺเธ ครุํ กตฺวา อสฺสาเทติ อภินนฺทติ, ตํ ครุํ กตฺวา ราโค อุปฺปชฺชติ.\csromanspacer{} ทิฏฺฐิ อุปฺปชฺชติ.\csromanspacer{} \textbf{สหชาตาธิปติ\csromandash{}}น-อรูปาวจราธิปติ สมฺปยุตฺตกานํ ขนฺธานํ จิตฺตสมุฏฺฐานานญฺจ รูปานํ อธิปติปจฺจเยน ปจฺจโย. (1)}

\csromanheader{2}{\textbf{อนนฺตราทิ}}
\proseitem{283}{\csromanfolio{164}อรูปาวจโร ธมฺโม อรูปาวจรสฺส ธมฺมสฺส อนนฺตรปจฺจเยน ปจฺจโย.\csromanspacer{} ปุริมา ปุริมา อรูปาวจรา ขนฺธา ปจฺฉิมานํ ปจฺฉิมานํ อรูปาวจรานํ ขนฺธานํ อนนฺตรปจฺจเยน ปจฺจโย.\csromanspacer{}\csromanspacer{}อรูปาวจโร ธมฺโม น-อรูปาวจรสฺส ธมฺมสฺส อนนฺตรปจฺจเยน ปจฺจโย.\csromanspacer{} อรูปาวจรํ จุติจิตฺตํ น-อรูปาวจรสฺส อุปปตฺติจิตฺตสฺส.\csromanspacer{} อรูปาวจรํ ภวงฺคํ อาวชฺชนาย.\csromanspacer{} อรูปาวจรา ขนฺธา น-อรูปาวจรสฺส วุฏฺฐานสฺส.\csromanspacer{} นิโรธา วุฏฺฐหนฺตสฺส เนวสญฺญานาสญฺญายตนํ ผลสมาปตฺติยา อนนฺตรปจฺจเยน ปจฺจโย. (2)}

\proseitem{284}{น-อรูปาวจโร ธมฺโม น-อรูปาวจรสฺส ธมฺมสฺส อนนฺตรปจฺจเยน ปจฺจโย.\csromanspacer{} ปุริมา ปุริมา น-อรูปาวจรา ขนฺธา ปจฺฉิมานํ ปจฺฉิมานํ น-อรูปาวจรานํ ขนฺธานํ อนนฺตรปจฺจเยน ปจฺจโย.\csromanspacer{} อนุโลมํ โคตฺรภุสฺส ฯเปฯ.\csromanspacer{} อนุโลมํ ผลสมาปตฺติยา อนนฺตรปจฺจเยน ปจฺจโย.\csromanspacer{}\csromanspacer{}น-อรูปาวจโร ธมฺโม อรูปาวจรสฺส ธมฺมสฺส อนนฺตรปจฺจเยน ปจฺจโย.\csromanspacer{} น-อรูปาวจรํ จุติจิตฺตํ อรูปาวจรสฺส อุปปตฺติจิตฺตสฺส อนนฺตรปจฺจเยน ปจฺจโย.\csromanspacer{} น-อรูปาวจรา ขนฺธา อรูปาวจรสฺส วุฏฺฐานสฺส อนนฺตรปจฺจเยน ปจฺจโย.\csromanspacer{} อากาสานญฺจายตนสฺส ปริกมฺมํ อากาสานญฺจายตนสฺส อนนฺตรปจฺจเยน ปจฺจโย.\csromanspacer{} วิญฺญาณญฺจายตนสฺส ฯเปฯ.\csromanspacer{} อากิญฺจญฺญายตนสฺส ฯเปฯ.\csromanspacer{} เนวสญฺญานาสญฺญายตนสฺส ปริกมฺมํ เนวสญฺญานาสญฺญายตนสฺส อนนฺตรปจฺจเยน ปจฺจโย. (2)}

\prose{สมนนฺตรปจฺจเยน ปจฺจโย.\csromanspacer{} สหชาตปจฺจเยน ปจฺจโย.\csromanspacer{} ปญฺจ.\csromanspacer{} อญฺญมญฺญปจฺจเยน ปจฺจโย.\csromanspacer{} เทฺว.\csromanspacer{} นิสฺสยปจฺจเยน ปจฺจโย.\csromanspacer{} สตฺต.}

\csromanheader{2}{\textbf{อุปนิสฺสย}}
\proseitem{285}{อรูปาวจโร ธมฺโม อรูปาวจรสฺส ธมฺมสฺส อุปนิสฺสยปจฺจเยน ปจฺจโย.\csromanspacer{} \textbf{อนนฺตรูปนิสฺสโย, ปกตูปนิสฺสโย ฯเปฯ.}\csromanspacer{} ปกตูปนิสฺสโย\csromandash{}อากาสานญฺจายตนํ วิญฺญาณญฺจายตนสฺส อุปนิสฺสยปจฺจเยน ปจฺจโย.\csromanspacer{} วิญฺญาณญฺจายตนํ อากิญฺจญฺญายตนสฺส ฯเปฯ.\csromanspacer{} อากิญฺจญฺญายตนํ เนวสญฺญานาสญฺญายตนสฺส อุปนิสฺสยปจฺจเยน ปจฺจโย.\csromanspacer{}\csromanspacer{}อรูปาวจโร ธมฺโม น-อรูปาวจรสฺส ธมฺมสฺส อุปนิสฺสยปจฺจเยน ปจฺจโย.\csromanspacer{} \textbf{อารมฺมณูปนิสฺสโย,} \textbf{อนนฺตรูปนิสฺสโย, ปกตูปนิสฺสโย ฯเปฯ.}}

\vspace{\tipitakaparskip}
\csromancenter{อรูปาวจรทุก ปกตูปนิสฺสโย\csromandash{}อรูปาวจรํ สทฺธํ อุปนิสฺสาย ทานํ เทติ.\csromanspacer{} สีลํ ฯเปฯ อุโปสถกมฺมํ กโรติ.\csromanspacer{} น-อรูปาวจรํ ฌานํ อุปฺปาเทติ.\csromanspacer{} วิปสฺสนํ ฯเปฯ มคฺคํ ฯเปฯ อภิญฺญํ ฯเปฯ สมาปตฺติํ อุปฺปาเทติ.\csromanspacer{} มานํ ชปฺเปติ.\csromanspacer{} ทิฏฺฐิํ คณฺหาติ.\csromanspacer{} อรูปาวจรํ สีลํ ฯเปฯ.\csromanspacer{} ปญฺญํ อุปนิสฺสาย ทานํ เทติ ฯเปฯ.\csromanspacer{} สมาปตฺติํ อุปฺปาเทติ.\csromanspacer{} มานํ ชปฺเปติ.\csromanspacer{} ทิฏฺฐิํ คณฺหาติ.\csromanspacer{} อรูปาวจรา สทฺธา ฯเปฯ ปญฺญา น-อรูปาวจราย สทฺธาย ฯเปฯ ปญฺญาย.\csromanspacer{} ราคสฺส ฯเปฯ ปตฺถนาย.\csromanspacer{} กายิกสฺส สุขสฺส.\csromanspacer{} กายิกสฺส ทุกฺขสฺส, มคฺคสฺส, ผลสมาปตฺติยา อุปนิสฺสยปจฺจเยน ปจฺจโย. (2)}
\proseitem{286}{\csromanfolio{165}น-อรูปาวจโร ธมฺโม น-อรูปาวจรสฺส ธมฺมสฺส อุปนิสฺสยปจฺจเยน ปจฺจโย.\csromanspacer{} \textbf{อารมฺมณูปนิสฺสโย, อนนฺตรูปนิสฺสโย,} \textbf{ปกตูปนิสฺสโย ฯเปฯ.}\csromanspacer{} ปกตูปนิสฺสโย\csromandash{}น-อรูปาวจรํ สทฺธํ อุปนิสฺสาย ทานํ เทติ.\csromanspacer{} สีลํ ฯเปฯ.\csromanspacer{} อุโปสถกมฺมํ กโรติ.\csromanspacer{} น-อรูปาวจรํ ฌานํ อุปฺปาเทติ.\csromanspacer{} วิปสฺสนํ ฯเปฯ.\csromanspacer{} มคฺคํ ฯเปฯ.\csromanspacer{} อภิญฺญํ ฯเปฯ.\csromanspacer{} สมาปตฺติํ อุปฺปาเทติ.\csromanspacer{} มานํ ชปฺเปติ.\csromanspacer{} ทิฏฺฐิํ คณฺหาติ.\csromanspacer{} น-อรูปาวจรํ สีลํ ฯเปฯ ปญฺญํ ฯเปฯ ราคํ ฯเปฯ ปตฺถนํ ฯเปฯ กายิกํ สุขํ.\csromanspacer{} กายิกํ ทุกฺขํ.\csromanspacer{} อุตุํ.\csromanspacer{} โภชนํ.\csromanspacer{} เสนาสนํ อุปนิสฺสาย ทานํ เทติ ฯเปฯ.\csromanspacer{} สมาปตฺติํ อุปฺปาเทติ.\csromanspacer{} ปาณํ หนติ ฯเปฯ สํฆํ ภินฺทติ.\csromanspacer{} น-อรูปาวจรา สทฺธา ฯเปฯ.\csromanspacer{} เสนาสนํ น-อรูปาวจราย สทฺธาย ฯเปฯ ปตฺถนาย.\csromanspacer{} กายิกสฺส สุขสฺส.\csromanspacer{} กายิกสฺส ทุกฺขสฺส, มคฺคสฺส, ผลสมาปตฺติยา อุปนิสฺสยปจฺจเยน ปจฺจโย. (1)}

\prose{น-อรูปาวจโร ธมฺโม อรูปาวจรสฺส ธมฺมสฺส อุปนิสฺสยปจฺจเยน ปจฺจโย.\csromanspacer{} \textbf{อนนฺตรูปนิสฺสโย, ปกตูปนิสฺสโย ฯเปฯ.}\csromanspacer{} ปกตูปนิสฺสโย\csromandash{}อากาสานญฺจายตนสฺส ปริกมฺมํ อากาสานญฺจายตนสฺส อุปนิสฺสยปจฺจเยน ปจฺจโย ฯเปฯ.\csromanspacer{} เนวสญฺญานาสญฺญายตนสฺส ปริกมฺมํ เนวสญฺญานาสญฺญายตนสฺส อุปนิสฺสยปจฺจเยน ปจฺจโย. (2)}

\csromanheader{2}{\textbf{ปุเรชาต}}
\proseitem{287}{น-อรูปาวจโร ธมฺโม น-อรูปาวจรสฺส ธมฺมสฺส ปุเรชาตปจฺจเยน ปจฺจโย.\csromanspacer{} อารมฺมณปุเรชาตํ, วตฺถุปุเรชาตํ.\csromanspacer{}\csromanspacer{}\textbf{อารมฺมณปุเรชาตํ\csromandash{}}จกฺขุํ ฯเปฯ.\csromanspacer{} วตฺถุํ อนิจฺจโต ฯเปฯ โทมนสฺสํ อุปฺปชฺชติ.\csromanspacer{} ทิพฺเพน จกฺขุนา รูปํ ปสฺสติ.\csromanspacer{} ทิพฺพาย โสตธาตุยา สทฺทํ สุณาติ.\csromanspacer{} รูปายตนํ จกฺขุวิญฺญาณสฺส ฯเปฯ.\csromanspacer{} โผฏฺฐพฺพายตนํ กายวิญฺญาณสฺส ฯเปฯ.\csromanspacer{} \textbf{วตฺถุปุเรชาตํ\csromandash{}} \csromanfolio{166}จกฺขายตนํ จกฺขุวิญฺญาณสฺส ฯเปฯ.\csromanspacer{} กายายตนํ กายวิญฺญาณสฺส ฯเปฯ.\csromanspacer{} วตฺถุ น-อรูปาวจรานํ ขนฺธานํ ปุเรชาตปจฺจเยน ปจฺจโย.\csromanspacer{}\csromanspacer{}น-อรูปาวจโร ธมฺโม อรูปาวจรสฺส ธมฺมสฺส ปุเรชาตปจฺจเยน ปจฺจโย.\csromanspacer{} \textbf{วตฺถุปุเรชาตํ\csromandash{}}วตฺถุ อรูปาวจรานํ ขนฺธานํ ปุเรชาตปจฺจเยน ปจฺจโย. (2)}

\prose{ปจฺฉาชาตปจฺจเยน ปจฺจโย.\csromanspacer{} เทฺว.\csromanspacer{} อาเสวนปจฺจเยน ปจฺจโย.\csromanspacer{} ตีณิ.}

\csromanheader{2}{\textbf{กมฺม}}
\proseitem{288}{อรูปาวจโร ธมฺโม อรูปาวจรสฺส ธมฺมสฺส กมฺมปจฺจเยน ปจฺจโย.\csromanspacer{} \textbf{สหชาตา}, นานากฺขณิกา ฯเปฯ.\csromanspacer{} อรูปาวจโร ธมฺโม น-อรูปาวจรสฺส ธมฺมสฺส กมฺมปจฺจเยน ปจฺจโย.\csromanspacer{} อรูปาวจรา เจตนา จิตฺตสมุฏฺฐานานํ รูปานํ กมฺมปจฺจเยน ปจฺจโย. (มูลํ กาตพฺพํ.) อรูปาวจรา เจตนา สมฺปยุตฺตกานํ ขนฺธานํ จิตฺตสมุฏฺฐานานญฺจ รูปานํ กมฺมปจฺจเยน ปจฺจโย. (3)}

\prose{น-อรูปาวจโร ธมฺโม น-อรูปาวจรสฺส ธมฺมสฺส กมฺมปจฺจเยน ปจฺจโย.\csromanspacer{} \textbf{สหชาตา}, นานากฺขณิกา.\csromanspacer{}\csromanspacer{}\textbf{สหชาตา} น-อรูปาวจรา เจตนา สมฺปยุตฺตกานํ ขนฺธานํ จิตฺตสมุฏฺฐานานญฺจ รูปานํ กมฺมปจฺจเยน ปจฺจโย. (สํขิตฺตํ.) (1)}

\csromanheader{2}{\textbf{วิปากาทิ}}
\proseitem{289}{อรูปาวจโร ธมฺโม อรูปาวจรสฺส ธมฺมสฺส วิปากปจฺจเยน ปจฺจโย. (สํขิตฺตํ.)}

\prose{น-อรูปาวจโร ธมฺโม น-อรูปาวจรสฺส ธมฺมสฺส วิปากปจฺจเยน ปจฺจโย. (สํขิตฺตํ.) อาหารปจฺจเยน ปจฺจโย.\csromanspacer{} จตฺตาริ.\csromanspacer{} อินฺทฺริยปจฺจเยน ปจฺจโย.\csromanspacer{} จตฺตาริ.\csromanspacer{} ฌานปจฺจเยน ปจฺจโย.\csromanspacer{} จตฺตาริ.\csromanspacer{} มคฺคปจฺจเยน ปจฺจโย.\csromanspacer{} จตฺตาริ.\csromanspacer{} สมฺปยุตฺตปจฺจเยน ปจฺจโย.\csromanspacer{} เทฺว.}

\csromanheader{2}{\textbf{วิปฺปยุตฺต}}
\proseitem{290}{อรูปาวจโร ธมฺโม น-อรูปาวจรสฺส ธมฺมสฺส วิปฺปยุตฺตปจฺจเยน ปจฺจโย.\csromanspacer{} สหชาตํ, ปจฺฉาชาตํ. (สํขิตฺตํ.) (1)}

\prose{น-อรูปาวจโร ธมฺโม น-อรูปาวจรสฺส ธมฺมสฺส วิปฺปยุตฺตปจฺจเยน ปจฺจโย.\csromanspacer{} สหชาตํ, ปุเรชาตํ, ปจฺฉาชาตํ. (สํขิตฺตํ.) น-อรูปาวจโร}

\vspace{\tipitakaparskip}
\csromancenter{อรูปาวจรทุก ธมฺโม อรูปาวจรสฺส ธมฺมสฺส วิปฺปยุตฺตปจฺจเยน ปจฺจโย.\csromanspacer{} \textbf{ปุเรชาตํ} วตฺถุ อรูปาวจรานํ ขนฺธานํ วิปฺปยุตฺตปจฺจเยน ปจฺจโย. (2)}
\csromanheader{2}{\textbf{อตฺถิ-อาทิ}}
\proseitem{291}{\csromanfolio{167}อรูปาวจโร ธมฺโม อรูปาวจรสฺส ธมฺมสฺส อตฺถิปจฺจเยน ปจฺจโย.\csromanspacer{} สหชาตํ ฯเปฯ.\csromanspacer{} อรูปาวจโร ธมฺโม น-อรูปาวจรสฺส ธมฺมสฺส อตฺถิปจฺจเยน ปจฺจโย.\csromanspacer{} สหชาตํ, ปจฺฉาชาตํ ฯเปฯ.\csromanspacer{} อรูปาวจโร ธมฺโม อรูปาวจรสฺส จ น-อรูปาวจรสฺส จ ธมฺมสฺส อตฺถิปจฺจเยน ปจฺจโย.\csromanspacer{} สหชาตํ. (สํขิตฺตํ.) (3)}

\prose{น-อรูปาวจโร ธมฺโม น-อรูปาวจรสฺส ธมฺมสฺส อตฺถิปจฺจเยน ปจฺจโย.\csromanspacer{} สหชาตํ, ปุเรชาตํ, ปจฺฉาชาตํ, อาหารํ, อินฺทฺริยํ. (สํขิตฺตํ).\csromanspacer{}\csromanspacer{}น-อรูปาวจโร ธมฺโม อรูปาวจรสฺส ธมฺมสฺส อตฺถิปจฺจเยน ปจฺจโย.\csromanspacer{} \textbf{ปุเรชาตํ}\csromandash{}วตฺถุ อรูปาวจรานํ ขนฺธานํ อตฺถิปจฺจเยน ปจฺจโย. (2)}

\proseitem{292}{อรูปาวจโร จ น-อรูปาวจโร จ ธมฺมา อรูปาวจรสฺส ธมฺมสฺส อตฺถิปจฺจเยน ปจฺจโย.\csromanspacer{} สหชาตํ, ปุเรชาตํ.\csromanspacer{}\csromanspacer{}สหชาโต\csromandash{} อรูปาวจโร เอโก ขนฺโธ จ วตฺถุ จ ติณฺณนฺนํ ขนฺธานํ อตฺถิปจฺจเยน ปจฺจโย ฯเปฯ เทฺว ขนฺธา จ ฯเปฯ.\csromanspacer{}\csromanspacer{}อรูปาวจโร จ น-อรูปาวจโร จ ธมฺมา น-อรูปาวจรสฺส ธมฺมสฺส อตฺถิปจฺจเยน ปจฺจโย.\csromanspacer{} สหชาตํ, ปจฺฉาชาตํ, อาหารํ, อินฺทฺริยํ.\csromanspacer{}\csromanspacer{}สหชาตา\csromandash{} อรูปาวจรา ขนฺธา จ มหาภูตา จ จิตฺตสมุฏฺฐานานํ รูปานํ อตฺถิปจฺจเยน ปจฺจโย.\csromanspacer{} ปจฺฉาชาตา\csromandash{}อรูปาวจรา ขนฺธา จ กพฬีกาโร อาหาโร จ อิมสฺส กายสฺส อตฺถิปจฺจเยน ปจฺจโย.\csromanspacer{} ปจฺฉาชาตา\csromandash{}อรูปาวจรา ขนฺธา จ รูปชีวิตินฺทฺริยญฺจ กฏตฺตารูปานํ อตฺถิปจฺจเยน ปจฺจโย. (2)}

\prose{นตฺถิปจฺจเยน ปจฺจโย.\csromanspacer{} วิคตปจฺจเยน ปจฺจโย.\csromanspacer{} อวิคตปจฺจเยน ปจฺจโย.}

\csromanheader{2}{1. ปจฺจยานุโลม 2. สงฺขฺยาวาร}
\proseitem{293}{เหตุยา จตฺตาริ, อารมฺมเณ ตีณิ, อธิปติยา จตฺตาริ, อนนฺตเร จตฺตาริ, สมนนฺตเร จตฺตาริ, สหชาเต ปญฺจ, อญฺญมญฺเญ เทฺว, นิสฺสเย สตฺต, อุปนิสฺสเย จตฺตาริ, ปุเรชาเต เทฺว, ปจฺฉาชาเต เทฺว, อาเสวเน ตีณิ, กมฺเม จตฺตาริ, วิปาเก เทฺว, อาหาเร จตฺตาริ, \csromanfolio{168}อินฺทฺริเย จตฺตาริ, ฌาเน จตฺตาริ, มคฺเค จตฺตาริ, สมฺปยุตฺเต เทฺว, วิปฺปยุตฺเต ตีณิ, อตฺถิยา สตฺต, นตฺถิยา จตฺตาริ, วิคเต จตฺตาริ, อวิคเต สตฺต.}

\csromanheader{2}{2. ปจฺจนียุทฺธาร}
\proseitem{294}{อรูปาวจโร ธมฺโม อรูปาวจรสฺส ธมฺมสฺส อารมฺมณปจฺจเยน ปจฺจโย.\csromanspacer{} สหชาตปจฺจเยน ปจฺจโย.\csromanspacer{} อุปนิสฺสยปจฺจเยน ปจฺจโย.\csromanspacer{}\csromanspacer{}อรูปาวจโร ธมฺโม น-อรูปาวจรสฺส ธมฺมสฺส อารมฺมณปจฺจเยน ปจฺจโย.\csromanspacer{} สหชาตปจฺจเยน ปจฺจโย.\csromanspacer{} อุปนิสยปจฺจเยน ปจฺจโย.\csromanspacer{} ปจฺฉาชาตปจฺจเยน ปจฺจโย.\csromanspacer{}\csromanspacer{}อรูปาวจโร ธมฺโม อรูปาวจรสฺส จ น-อรูปาวจรสฺส จ ธมฺมสฺส สหชาตปจฺจเยน ปจฺจโย. (3)}

\proseitem{295}{น-อรูปาวจโร ธมฺโม น-อรูปาวจรสฺส ธมฺมสฺส อารมฺมณปจฺจเยน ปจฺจโย.\csromanspacer{} สหชาตปจฺจเยน ปจฺจโย.\csromanspacer{} อุปนิสฺสยปจฺจเยน ปจฺจโย.\csromanspacer{} ปุเรชาตปจฺจเยน ปจฺจโย.\csromanspacer{} ปจฺฉาชาตปจฺจเยน ปจฺจโย.\csromanspacer{} กมฺมปจฺจเยน ปจฺจโย.\csromanspacer{} อาหารปจฺจเยน ปจฺจโย.\csromanspacer{} อินฺทฺริยปจฺจเยน ปจฺจโย.\csromanspacer{}\csromanspacer{}น-อรูปาวจโร ธมฺโม อรูปาวจรสฺส ธมฺมสฺส อุปนิสฺสยปจฺจเยน ปจฺจโย.\csromanspacer{} ปุเรชาตปจฺจเยน ปจฺจโย. (2)}

\prose{อรูปาวจโร จ น-อรูปาวจโร จ ธมฺมา อรูปาวจรสฺส ธมฺมสฺส สหชาตํ.\csromanspacer{} ปุเรชาตํ.\csromanspacer{}\csromanspacer{}อรูปาวจโร จ น-อรูปาวจโร จ ธมฺมา น-อรูปาวจรสฺส ธมฺมสฺส สหชาตํ, ปจฺฉาชาตํ, อาหารํ, อินฺทฺริยํ. (2)}

\csromanheader{2}{2. ปจฺจยปจฺจนีย 2. สงฺขฺยาวาร}
\proseitem{296}{นเหตุยา สตฺต, น-อารมฺมเณ สตฺต, น-อธิปติยา สตฺต, น-อนนฺตเร สตฺต, นสมนนฺตเร สตฺต, นสหชาเต ปญฺจ, น-อญฺญมญฺเญ ปญฺจ, นนิสฺสเย ปญฺจ, น-อุปนิสฺสเย สตฺต, นปุเรชาเต ฉ, นปจฺฉาชาเต สตฺต ฯเปฯ นสมฺปยุตฺเต ปญฺจ, นวิปฺปยุตฺเต จตฺตาริ, โน-อตฺถิยา จตฺตาริ, โนนตฺถิยา สตฺต, โนวิคเต สตฺต, โน-อวิคเต จตฺตาริ.}

\csromanmatikaanchor{mtk.25}
\csromanmatikaanchor{mtk.26}
\csromantocmarkref{section}{96. ปริยาปนฺนทุก}{mtk.25}
\bookmark[dest={mtk.25},level=1]{96. ปริยาปนฺนทุก}
\csromantocmarkref{section}{97. นิยฺยานิกทุก}{mtk.26}
\bookmark[dest={mtk.26},level=1]{97. นิยฺยานิกทุก}
\csromanheader{1}{97. นิยฺยานิกทุก \textbf{3. ปจฺจยานุโลมปจฺจนีย}}
\proseitem{297}{\csromanfolio{169}เหตุปจฺจยา น-อารมฺมเณ จตฺตาริ, น-อธิปติยา จตฺตาริ, น-อนนฺตเร นสมนนฺตเร จตฺตาริ, น-อญฺญมญฺเญ เทฺว, น-อุปนิสฺสเย จตฺตาริ ฯเปฯ นสมฺปยุตฺเต เทฺว, นวิปฺปยุตฺเต เทฺว, โนนตฺถิยา จตฺตาริ, โนวิคเต จตฺตาริ.}

\csromanheader{2}{4. ปจฺจยปจฺจนียานุโลม}
\proseitem{298}{นเหตุปจฺจยา อารมฺมเณ ตีณิ, อธิปติยา จตฺตาริ, (อนุโลมมาติกา กาตพฺพา) ฯเปฯ อวิคเต สตฺต.}

\nitthitam{อรูปาวจรทุกํ นิฏฺฐิตํ.}
\csromanheader{1}{96. ปริยาปนฺนทุก 1. ปฏิจฺจวาร}
\proseitem{299}{ปริยาปนฺนํ ธมฺมํ ปฏิจฺจ ปริยาปนฺโน ธมฺโม อุปฺปชฺชติ เหตุปจฺจยา.\csromanspacer{} ปริยาปนฺนํ เอกํ ขนฺธํ ปฏิจฺจ ตโย ขนฺธา จิตฺตสมุฏฺฐานญฺจ รูปํ ฯเปฯ เทฺว ขนฺเธ ฯเปฯ.\csromanspacer{} ปฏิสนฺธิกฺขเณ ฯเปฯ. (ยถา จูฬนฺตรทุเก โลกิยทุกํ.\csromanspacer{} เอวํ อิมมฺปิ ทุกํ กาตพฺพํ, นินฺนานากรณํ.)}

\nitthitam{ปริยาปนฺนทุกํ นิฏฺฐิตํ.}
\csromanheader{2}{97. นิยฺยานิกทุก 1. ปฏิจฺจวาร}
\proseitem{300}{นิยฺยานิกํ ธมฺมํ ปฏิจฺจ นิยฺยานิโก ธมฺโม อุปฺปชฺชติ เหตุปจฺจยา.\csromanspacer{} นิยฺยานิกํ เอกํ ขนฺธํ ปฏิจฺจ ตโย ขนฺธา ฯเปฯ เทฺว ขนฺเธ ฯเปฯ.\csromanspacer{} นิยฺยานิกํ ธมฺมํ ปฏิจฺจ อนิยฺยานิโก ธมฺโม อุปฺปชฺชติ เหตุปจฺจยา.\csromanspacer{} นิยฺยานิเก ขนฺเธ ปฏิจฺจ จิตฺตสมุฏฺฐานํ รูปํ.\csromanspacer{}\csromanspacer{}นิยฺยานิกํ ธมฺมํ ปฏิจฺจ นิยฺยานิโก จ อนิยฺยานิโก จ ธมฺมา อุปฺปชฺชนฺติ เหตุปจฺจยา.\csromanspacer{} นิยฺยานิกํ เอกํ ขนฺธํ ปฏิจฺจ ตโย ขนฺธา จิตฺตสมุฏฺฐานญฺจ รูปํ ฯเปฯ เทฺว ขนฺเธ ฯเปฯ. (3)}

\prose{\csromanfolio{170}อนิยฺยานิกํ ธมฺมํ ปฏิจฺจ อนิยฺยานิโก ธมฺโม อุปฺปชฺชติ เหตุปจฺจยา.\csromanspacer{} อนิยฺยานิกํ เอกํ ขนฺธํ ปฏิจฺจ ตโย ขนฺธา จิตฺตสมุฏฺฐานญฺจ รูปํ ฯเปฯ เทฺว ขนฺเธ ฯเปฯ.\csromanspacer{} ปฏิสนฺธิกฺขเณ ฯเปฯ.\csromanspacer{} ขนฺเธ ปฏิจฺจ วตฺถุ, วตฺถุํ ปฏิจฺจ ขนฺธา.\csromanspacer{} เอกํ มหาภูตํ ฯเปฯ. (1)}

\prose{นิยฺยานิกญฺจ อนิยฺยานิกญฺจ ธมฺมํ ปฏิจฺจ อนิยฺยานิโก ธมฺโม อุปฺปชฺชติ เหตุปจฺจยา.\csromanspacer{} นิยฺยานิเก ขนฺเธ จ มหาภูเต จ ปฏิจฺจ จิตฺตสมุฏฺฐานํ รูปํ. (1)}

\csromanheader{2}{1. ปจฺจยานุโลม 2. สงฺขฺยาวาร}
\proseitem{301}{เหตุยา ปญฺจ, อารมฺมเณ เทฺว, อธิปติยา ปญฺจ, อนนฺตเร สมนนฺตเร เทฺว, สหชาเต ปญฺจ, อญฺญมญฺเญ เทฺว, นิสฺสเย ปญฺจ, อุปนิสฺสเย เทฺว, ปุเรชาเต เทฺว, อาเสวเน เทฺว, กมฺเม ปญฺจ, วิปาเก เอกํ, อาหาเร ปญฺจ ฯเปฯ อวิคเต ปญฺจ.}

\csromanheader{2}{2. ปจฺจยปจฺจนีย 1. วิภงฺควาร}
\csromanheader{2}{\textbf{นเหตุ น-อารมฺมณ}}
\proseitem{302}{อนิยฺยานิกํ ธมฺมํ ปฏิจฺจ อนิยฺยานิโก ธมฺโม อุปฺปชฺชติ นเหตุปจฺจยา.\csromanspacer{} อเหตุกํ อนิยฺยานิกํ เอกํ ขนฺธํ ปฏิจฺจ ตโย ขนฺธา จิตฺตสมุฏฺฐานญฺจ รูปํ ฯเปฯ เทฺว ขนฺเธ ฯเปฯ.\csromanspacer{} อเหตุกปฏิสนฺธิกฺขเณ ฯเปฯ.\csromanspacer{} ขนฺเธ ปฏิจฺจ วตฺถุ, วตฺถุํ ปฏิจฺจ ขนฺธา.\csromanspacer{} เอกํ มหาภูตํ ฯเปฯ. (ยาว อสญฺญสตฺตา.) วิจิกิจฺฉาสหคเต อุทฺธจฺจสหคเต ขนฺเธ ปฏิจฺจ วิจิกิจฺฉาสหคโต อุทฺธจฺจสหคโต โมโห. (1).\csromanspacer{} น-อารมฺมณปจฺจยา.\csromanspacer{} ตีณิ.}

\csromanheader{2}{\textbf{น-อธิปตฺยาทิ}}
\proseitem{303}{นิยฺยานิกํ ธมฺมํ ปฏิจฺจ นิยฺยานิโก ธมฺโม อุปฺปชฺชติ น-อธิปติปจฺจยา.\csromanspacer{} นิยฺยานิเก ขนฺเธ ปฏิจฺจ นิยฺยานิกาธิปติ. (1)}

\prose{อนิยฺยานิกํ ธมฺมํ ปฏิจฺจ อนิยฺยานิโก ธมฺโม อุปฺปชฺชติ น-อธิปติปจฺจยา.\csromanspacer{} อนิยฺยานิกํ เอกํ ขนฺธํ ปฏิจฺจ ตโย ขนฺธา จิตฺตสมุฏฺฐานญฺจ รูปํ ฯเปฯ เทฺว ขนฺเธ ฯเปฯ.\csromanspacer{} ปฏิสนฺธิกฺขเณ ฯเปฯ. (ยาว อสญฺญสตฺตา กาตพฺพา.) (1)}

\prose{น-อนนฺตรปจฺจยา.\csromanspacer{} นสมนนฺตรปจฺจยา.\csromanspacer{} น-อญฺญมญฺญปจฺจยา ฯเปฯ.}

\csromanheader{2}{97. นิยฺยานิกทุก \textbf{นปุเรชาต}}
\proseitem{304}{\csromanfolio{171}นิยฺยานิกํ ธมฺมํ ปฏิจฺจ นิยฺยานิโก ธมฺโม อุปฺปชฺชติ นปุเรชาตปจฺจยา.\csromanspacer{} อรูเป นิยฺยานิกํ เอกํ ขนฺธํ ฯเปฯ.\csromanspacer{} นิยฺยานิกํ ธมฺมํ ปฏิจฺจ อนิยฺยานิโก ธมฺโม อุปฺปชฺชติ นปุเรชาตปจฺจยา.\csromanspacer{} นิยฺยานิเก ขนฺเธ ปฏิจฺจ จิตฺตสมุฏฺฐานํ รูปํ. (2)}

\prose{อนิยฺยานิกํ ธมฺมํ ปฏิจฺจ อนิยฺยานิโก ธมฺโม อุปฺปชฺชติ นปุเรชาตปจฺจยา.\csromanspacer{} อรูเป อนิยฺยานิกํ เอกํ ขนฺธํ ปฏิจฺจ ตโย ขนฺธา ฯเปฯ เทฺว ขนฺเธ ฯเปฯ.\csromanspacer{} อนิยฺยานิเก ขนฺเธ ปฏิจฺจ จิตฺตสมุฏฺฐานํ รูปํ.\csromanspacer{} ปฏิสนฺธิกฺขเณ ฯเปฯ. (ยาว อสญฺญสตฺตา.) (1)}

\prose{นิยฺยานิกญฺจ อนิยฺยานิกญฺจ ธมฺมํ ปฏิจฺจ อนิยฺยานิโก ธมฺโม อุปฺปชฺชติ นปุเรชาตปจฺจยา.\csromanspacer{} นิยฺยานิเก ขนฺเธ จ มหาภูเต จ ปฏิจฺจ จิตฺตสมุฏฺฐานํ รูปํ. (1)}

\csromanheader{2}{2. ปจฺจยปจฺจนีย 2. สงฺขฺยาวาร}
\proseitem{305}{นเหตุยา เอกํ, น-อารมฺมเณ ตีณิ, น-อธิปติยา เทฺว, น-อนนฺตเร ตีณิ, น-อุปนิสฺสเย ตีณิ, นปุเรชาเต จตฺตาริ, นปจฺฉาชาเต ปญฺจ, น-อาเสวเน เอกํ, นกมฺเม เทฺว, นวิปาเก ปญฺจ, น-อาหาเร เอกํ, น-อินฺทฺริเย เอกํ, นฌาเน เอกํ, นมคฺเค เอกํ, นสมฺปยุตฺเต ตีณิ, นวิปฺปยุตฺเต เทฺว, โนนตฺถิยา ตีณิ, โนวิคเต ตีณิ.}

\vspace{\tipitakaparskip}
\csromancenter{(เอวํ อิตเร เทฺว คณนาปิ สหชาตวาโรปิ กาตพฺโพ.)}
\csromanheader{2}{97. นิยฺยานิกทุก 3. ปจฺจยวาร}
\proseitem{306}{นิยฺยานิกํ ธมฺมํ ปจฺจยา นิยฺยานิโก ธมฺโม อุปฺปชฺชติ เหตุปจฺจยา.\csromanspacer{} ตีณิ. (ปฏิจฺจสทิสา.)}

\prose{\csromanfolio{172}อนิยฺยานิกํ ธมฺมํ ปจฺจยา อนิยฺยานิโก ธมฺโม อุปฺปชฺชติ เหตุปจฺจยา.\csromanspacer{} อนิยฺยานิกํ เอกํ ขนฺธํ ปจฺจยา ตโย ขนฺธา จิตฺตสมุฏฺฐานญฺจ รูปํ ฯเปฯ เทฺว ขนฺเธ ฯเปฯ. (ยาว อชฺฌตฺติกา มหาภูตา.) วตฺถุํ ปจฺจยา อนิยฺยานิกา ขนฺธา.\csromanspacer{}\csromanspacer{}อนิยฺยานิกํ ธมฺมํ ปจฺจยา นิยฺยานิโก ธมฺโม อุปฺปชฺชติ เหตุปจฺจยา.\csromanspacer{} วตฺถุํ ปจฺจยา นิยฺยานิกา ขนฺธา.\csromanspacer{}\csromanspacer{}อนิยฺยานิกํ ธมฺมํ ปจฺจยา นิยฺยานิโก จ อนิยฺยานิโก จ ธมฺมา อุปฺปชฺชนฺติ เหตุปจฺจยา.\csromanspacer{} วตฺถุํ ปจฺจยา นิยฺยานิกา ขนฺธา.\csromanspacer{} มหาภูเต ปจฺจยา จิตฺตสมุฏฺฐานํ รูปํ. (3)}

\proseitem{307}{นิยฺยานิกญฺจ อนิยฺยานิกญฺจ ธมฺมํ ปจฺจยา นิยฺยานิโก ธมฺโม อุปฺปชฺชติ เหตุปจฺจยา.\csromanspacer{} นิยฺยานิกํ เอกํ ขนฺธญฺจ วตฺถุญฺจ ปจฺจยา ตโย ขนฺธา ฯเปฯ เทฺว ขนฺเธ จ ฯเปฯ.\csromanspacer{} นิยฺยานิกญฺจ อนิยฺยานิกญฺจ ธมฺมํ ปจฺจยา อนิยฺยานิโก ธมฺโม อุปฺปชฺชติ เหตุปจฺจยา.\csromanspacer{} นิยฺยานิเก ขนฺเธ จ มหาภูเต จ ปจฺจยา จิตฺตสมุฏฺฐานํ รูปํ.\csromanspacer{}\csromanspacer{}นิยฺยานิกญฺจ อนิยฺยานิกญฺจ ธมฺมํ ปจฺจยา นิยฺยานิโก จ อนิยฺยานิโก จ ธมฺมา อุปฺปชฺชนฺติ เหตุปจฺจยา.\csromanspacer{} นิยฺยานิกํ เอกํ ขนฺธญฺจ วตฺถุญฺจ ปจฺจยา ตโย ขนฺธา ฯเปฯ เทฺว ขนฺเธ จ ฯเปฯ.\csromanspacer{} นิยฺยานิเก ขนฺเธ จ มหาภูเต จ ปจฺจยา จิตฺตสมุฏฺฐานํ รูปํ. (3)}

\csromanheader{2}{1. ปจฺจยานุโลม 2. สงฺขฺยาวาร}
\proseitem{308}{เหตุยา นว, อารมฺมเณ จตฺตาริ, อธิปติยา นว, อนนฺตเร จตฺตาริ, สมนนฺตเร จตฺตาริ, สหชาเต นว, อญฺญมญฺเญ จตฺตาริ, นิสฺสเย นว, อุปนิสฺสเย จตฺตาริ, ปุเรชาเต จตฺตาริ, อาเสวเน จตฺตาริ, กมฺเม นว, วิปาเก เอกํ ฯเปฯ อวิคเต นว.}

\csromanheader{2}{2. ปจฺจยปจฺจนีย 1. วิภงฺควาร}
\csromanheader{2}{\textbf{นเหตุ}}
\proseitem{309}{อนิยฺยานิกํ ธมฺมํ ปจฺจยา อนิยฺยานิโก ธมฺโม อุปฺปชฺชติ นเหตุปจฺจยา.\csromanspacer{} อเหตุกํ อนิยฺยานิกํ เอกํ ขนฺธํ ปจฺจยา ตโย ขนฺธา จิตฺตสมุฏฺฐานญฺจ รูปํ ฯเปฯ เทฺว ขนฺเธ ฯเปฯ. (ยาว อสญฺญสตฺตา.) จกฺขายตนํ ปจฺจยา จกฺขุวิญฺญาณํ ฯเปฯ.\csromanspacer{} กายายตนํ ปจฺจยา กายวิญฺญาณํ.\csromanspacer{} วตฺถุํ ปจฺจยา อเหตุกา อนิยฺยานิกา ขนฺธา.\csromanspacer{} วิจิกิจฺฉาสหคเต อุทฺธจฺจสหคเต ขนฺเธ จ วตฺถุญฺจ ปจฺจยา วิจิกิจฺฉาสหคโต อุทฺธจฺจสหคโต โมโห. (1)}

\csromanheader{2}{97. นิยฺยานิกทุก \textbf{น-อารมฺมณาทิ}}
\proseitem{310}{\csromanfolio{173}นิยฺยานิกํ ธมฺมํ ปจฺจยา อนิยฺยานิโก ธมฺโม อุปฺปชฺชติ น-อารมฺมณปจฺจยา.\csromanspacer{} ตีณิ.}

\prose{นิยฺยานิกํ ธมฺมํ ปจฺจยา นิยฺยานิโก ธมฺโม อุปฺปชฺชติ น-อธิปติปจฺจยา.\csromanspacer{} นิยฺยานิเก ขนฺเธ ปจฺจยา นิยฺยานิกาธิปติ. (1)}

\prose{อนิยฺยานิกํ ธมฺมํ ปจฺจยา อนิยฺยานิโก ธมฺโม อุปฺปชฺชติ น-อธิปติปจฺจยา.\csromanspacer{} อนิยฺยานิกํ เอกํ ขนฺธํ ปจฺจยา ฯเปฯ. (ยาว อสญฺญสตฺตา.) จกฺขายตนํ ปจฺจยา จกฺขุวิญฺญาณํ ฯเปฯ.\csromanspacer{} กายายตนํ ปจฺจยา กายวิญฺญาณํ.\csromanspacer{} วตฺถุํ ปจฺจยา อนิยฺยานิกา ขนฺธา.\csromanspacer{}\csromanspacer{}อนิยฺยานิกํ ธมฺมํ ปจฺจยา นิยฺยานิโก ธมฺโม อุปฺปชฺชติ น-อธิปติปจฺจยา.\csromanspacer{} วตฺถุํ ปจฺจยา นิยฺยานิกาธิปติ. (2)}

\prose{นิยฺยานิกญฺจ อนิยฺยานิกญฺจ ธมฺมํ ปจฺจยา นิยฺยานิโก ธมฺโม อุปฺปชฺชติ น-อธิปติปจฺจยา.\csromanspacer{} นิยฺยานิเก ขนฺเธ จ วตฺถุญฺจ ปจฺจยา นิยฺยานิกาธิปติ. (1)}

\csromanheader{2}{2. ปจฺจยปจฺจนีย 2. สงฺขฺยาวาร}
\proseitem{311}{นเหตุยา เอกํ, น-อารมฺมเณ ตีณิ, น-อธิปติยา จตฺตาริ, น-อนนฺตเร ตีณิ ฯเปฯ น-อุปนิสฺสเย ตีณิ, นปุเรชาเต จตฺตาริ, นปจฺฉาชาเต นว, น-อาเสวเน เอกํ, นกมฺเม จตฺตาริ, นวิปาเก นว, น-อาหาเร เอกํ, น-อินฺทฺริเย เอกํ, นฌาเน เอกํ, นมคฺเค เอกํ, นสมฺปยุตฺเต ตีณิ, นวิปฺปยุตฺเต เทฺว, โนนตฺถิยา ตีณิ, โนวิคเต ตีณิ.}

\vspace{\tipitakaparskip}
\csromancenter{(เอวํ อิตเร เทฺว คณนาปิ นิสฺสยวาโรปิ กาตพฺโพ.)}
\csromanheader{2}{97. นิยฺยานิกทุก 5. สํสฏฺฐวาร}
\csromanheader{2}{1. ปจฺจยานุโลมาทิ}
\proseitem{312}{นิยฺยานิกํ ธมฺมํ สํสฏฺโฐ นิยฺยานิโก ธมฺโม อุปฺปชฺชติ เหตุปจฺจยา.\csromanspacer{} นิยฺยานิกํ เอกํ ขนฺธํ สํสฏฺฐา ตโย ขนฺธา ฯเปฯ เทฺว ขนฺเธ ฯเปฯ. (1)}

\prose{\csromanfolio{174}อนิยฺยานิกํ ธมฺมํ สํสฏฺโฐ อนิยฺยานิโก ธมฺโม อุปฺปชฺชติ เหตุปจฺจยา.\csromanspacer{} อนิยฺยานิกํ เอกํ ขนฺธํ สํสฏฺฐา ตโย ขนฺธา ฯเปฯ เทฺว ขนฺเธ ฯเปฯ.\csromanspacer{} ปฏิสนฺธิกฺขเณ ฯเปฯ. (1)}

\prose{\textbf{เหตุ}ยา เทฺว, อารมฺมเณ เทฺว, (สพฺพตฺถ เทฺว,) วิปาเก เอกํ ฯเปฯ อวิคเต เทฺว. (อนุโลมํ.)}

\prose{นเหตุยา เอกํ, น-อธิปติยา เทฺว, นปุเรชาเต เทฺว, นปจฺฉาชาเต เทฺว, น-อาเสวเน เอกํ, นกมฺเม เทฺว, นวิปาเก เทฺว, นฌาเน เอกํ, นมคฺเค เอกํ, นวิปฺปยุตฺเต เทฺว. (ปจฺจนียํ.)}

\vspace{\tipitakaparskip}
\csromancenter{(เอวํ อิตเร เทฺว คณนาปิ สมฺปยุตฺตวาโรปิ กาตพฺโพ.)}
\csromanheader{2}{97. นิยฺยานิกทุก 7. ปญฺหาวาร}
\proseitem{313}{นิยฺยานิโก ธมฺโม นิยฺยานิกสฺส ธมฺมสฺส เหตุปจฺจเยน ปจฺจโย.\csromanspacer{} นิยฺยานิกา เหตู สมฺปยุตฺตกานํ ขนฺธานํ เหตุปจฺจเยน ปจฺจโย. (มูลํ กาตพฺพํ.) นิยฺยานิกา เหตู จิตฺตสมุฏฺฐานานํ รูปานํ เหตุปจฺจเยน ปจฺจโย. (มูลํ กาตพฺพํ.) นิยฺยานิกา เหตู สมฺปยุตฺตกานํ ขนฺธานํ จิตฺตสมุฏฺฐานานญฺจ รูปานํ เหตุปจฺจเยน ปจฺจโย. (3)}

\prose{อนิยฺยานิโก ธมฺโม อนิยฺยานิกสฺส ธมฺมสฺส เหตุปจฺจเยน ปจฺจโย.\csromanspacer{} อนิยฺยานิกา เหตู สมฺปยุตฺตกานํ ขนฺธานํ จิตฺตสมุฏฺฐานานญฺจ รูปานํ เหตุปจฺจเยน ปจฺจโย.\csromanspacer{} ปฏิสนฺธิกฺขเณ ฯเปฯ. (1)}

\csromanheader{2}{\textbf{อารมฺมณ}}
\proseitem{314}{นิยฺยานิโก ธมฺโม อนิยฺยานิกสฺส ธมฺมสฺส อารมฺมณปจฺจเยน ปจฺจโย.\csromanspacer{} อริยา มคฺคา วุฏฺฐหิตฺวา มคฺคํ ปจฺจเวกฺขนฺติ.\csromanspacer{} เจโตปริยญาเณน นิยฺยานิกจิตฺตสมงฺคิสฺส จิตฺตํ ชานาติ.\csromanspacer{} นิยฺยานิกา ขนฺธา เจโตปริยญาณสฺส, ปุพฺเพนิวาสานุสฺสติญาณสฺส, อนาคตํสญาณสฺส, อาวชฺชนาย อารมฺมณปจฺจเยน ปจฺจโย. (1)}

\csromanheader{2}{97. นิยฺยานิกทุก}
\prose{\csromanfolio{175}อนิยฺยานิโก ธมฺโม อนิยฺยานิกสฺส ธมฺมสฺส อารมฺมณปจฺจเยน ปจฺจโย.\csromanspacer{} ทานํ ฯเปฯ.\csromanspacer{} สีลํ ฯเปฯ อุโปสถกมฺมํ กตฺวา ตํ ปจฺจเวกฺขติ, อสฺสาเทติ อภินนฺทติ, ตํ อารพฺภ ราโค อุปฺปชฺชติ ฯเปฯ โทมนสฺสํ อุปฺปชฺชติ.\csromanspacer{} ปุพฺเพ สุจิณฺณานิ ฯเปฯ.\csromanspacer{} ฌานา ฯเปฯ.\csromanspacer{} อริยา ผลํ ปจฺจเวกฺขนฺติ, นิพฺพานํ ปจฺจเวกฺขนฺติ.\csromanspacer{} นิพฺพานํ โคตฺรภุสฺส, โวทานสฺส, ผลสฺส, อาวชฺชนาย อารมฺมณปจฺจเยน ปจฺจโย.\csromanspacer{} อริยา ปหีเน กิเลเส ฯเปฯ.\csromanspacer{} วิกฺขมฺภิเต ฯเปฯ.\csromanspacer{} ปุพฺเพ สมุทาจิณฺเณ ฯเปฯ.\csromanspacer{} จกฺขุํ ฯเปฯ.\csromanspacer{} วตฺถุํ อนิยฺยานิเก ขนฺเธ อนิจฺจโต ฯเปฯ โทมนสฺสํ อุปฺปชฺชติ.\csromanspacer{} ทิพฺเพน จกฺขุนา รูปํ ปสฺสติ.\csromanspacer{} ทิพฺพาย โสตธาตุยา สทฺทํ สุณาติ.\csromanspacer{} เจโตปริยญาเณน อนิยฺยานิกจิตฺตสมงฺคิสฺส จิตฺตํ ชานาติ.\csromanspacer{} อากาสานญฺจายตนํ วิญฺญาณญฺจายตนสฺส ฯเปฯ.\csromanspacer{} อากิญฺจญฺญายตนํ เนวสญฺญานาสญฺญายตนสฺส ฯเปฯ.\csromanspacer{} รูปายตนํ จกฺขุวิญฺญาณสฺส ฯเปฯ.\csromanspacer{} โผฏฺฐพฺพายตนํ กายวิญฺญาณสฺส ฯเปฯ.\csromanspacer{} อนิยฺยานิกา ขนฺธา อิทฺธิวิธญาณสฺส, เจโตปริยญาณสฺส, ปุพฺเพนิวาสานุสฺสติญาณสฺส, ยถากมฺมูปคญาณสฺส, อนาคตํสญาณสฺส, อาวชฺชนาย อารมฺมณปจฺจเยน ปจฺจโย. (1)}

\prose{อนิยฺยานิโก ธมฺโม นิยฺยานิกสฺส ธมฺมสฺส อารมฺมณปจฺจเยน ปจฺจโย.\csromanspacer{} นิพฺพานํ มคฺคสฺส อารมฺมณปจฺจเยน ปจฺจโย. (2)}

\csromanheader{2}{\textbf{อธิปติ}}
\proseitem{315}{นิยฺยานิโก ธมฺโม นิยฺยานิกสฺส ธมฺมสฺส อธิปติปจฺจเยน ปจฺจโย.\csromanspacer{} \textbf{สหชาตาธิปติ\csromandash{}}นิยฺยานิกาธิปติ สมฺปยุตฺตกานํ ขนฺธานํ อธิปติปจฺจเยน ปจฺจโย.\csromanspacer{}\csromanspacer{}นิยฺยานิโก ธมฺโม อนิยฺยานิกสฺส ธมฺมสฺส อธิปติปจฺจเยน ปจฺจโย.\csromanspacer{} อารมฺมณาธิปติ, สหชาตาธิปติ.\csromanspacer{}\csromanspacer{}\textbf{อารมฺมณาธิปติ\csromandash{}}อริยา มคฺคา วุฏฺฐหิตฺวา มคฺคํ ครุํ กตฺวา ปจฺจเวกฺขนฺติ.\csromanspacer{} \textbf{สหชาตาธิปติ\csromandash{}}นิยฺยานิกาธิปติ จิตฺตสมุฏฺฐานานํ รูปานํ อธิปติปจฺจเยน ปจฺจโย. (มูลํ กาตพฺพํ.) นิยฺยานิกาธิปติ สมฺปยุตฺตกานํ ขนฺธานํ จิตฺตสมุฏฺฐานานญฺจ รูปานํ อธิปติปจฺจเยน ปจฺจโย. (3)}

\proseitem{316}{อนิยฺยานิโก ธมฺโม อนิยฺยานิกสฺส ธมฺมสฺส อธิปติปจฺจเยน ปจฺจโย.\csromanspacer{} อารมฺมณาธิปติ, สหชาตาธิปติ.\csromanspacer{}\csromanspacer{}\textbf{อารมฺมณาธิปติ\csromandash{}}ทานํ ฯเปฯ สีลํ ฯเปฯ อุโปสถกมฺมํ กตฺวา ตํ ครุํ กตฺวา ปจฺจเวกฺขติ, อสฺสาเทติ อภินนฺทติ, ตํ ครุํ กตฺวา ราโค อุปฺปชฺชติ.\csromanspacer{} ทิฏฺฐิ อุปฺปชฺชติ.\csromanspacer{} ปุพฺเพ สุจิณฺณานิ ฯเปฯ.\csromanspacer{} ฌานา วุฏฺฐหิตฺวา ฯเปฯ.\csromanspacer{} อริยา ผลํ ครุํ กตฺวา \csromanfolio{176}ปจฺจเวกฺขนฺติ, นิพฺพานํ ครุํ กตฺวา ปจฺจเวกฺขนฺติ.\csromanspacer{} นิพฺพานํ โคตฺรภุสฺส, โวทานสฺส, ผลสฺส อธิปติปจฺจเยน ปจฺจโย.\csromanspacer{} จกฺขุํ ฯเปฯ.\csromanspacer{} วตฺถุํ อนิยฺยานิเก ขนฺเธ ครุํ กตฺวา อสฺสาเทติ อภินนฺทติ, ตํ ครุํ กตฺวา ราโค อุปฺปชฺชติ.\csromanspacer{} ทิฏฺฐิ อุปฺปชฺชติ.\csromanspacer{} \textbf{สหชาตาธิปติ\csromandash{}} อนิยฺยานิกาธิปติ สมฺปยุตฺตกานํ ขนฺธานํ จิตฺตสมุฏฺฐานานญฺจ รูปานํ อธิปติปจฺจเยน ปจฺจโย.\csromanspacer{}\csromanspacer{}อนิยฺยานิโก ธมฺโม นิยฺยานิกสฺส ธมฺมสฺส อธิปติปจฺจเยน ปจฺจโย.\csromanspacer{} \textbf{อารมฺมณาธิปติ\csromandash{}}นิพฺพานํ มคฺคสฺส อธิปติปจฺจเยน ปจฺจโย. (2)}

\csromanheader{2}{\textbf{อนนฺตราทิ}}
\proseitem{317}{นิยฺยานิโก ธมฺโม อนิยฺยานิกสฺส ธมฺมสฺส อนนฺตรปจฺจเยน ปจฺจโย.\csromanspacer{} มคฺโค ผลสฺส อนนฺตรปจฺจเยน ปจฺจโย. (1)}

\prose{อนิยฺยานิโก ธมฺโม อนิยฺยานิกสฺส ธมฺมสฺส อนนฺตรปจฺจเยน ปจฺจโย.\csromanspacer{} ปุริมา ปุริมา อนิยฺยานิกา ขนฺธา ปจฺฉิมานํ ปจฺฉิมานํ อนิยฺยานิกานํ ขนฺธานํ อนนฺตรปจฺจเยน ปจฺจโย.\csromanspacer{} อนุโลมํ โคตฺรภุสฺส.\csromanspacer{} อนุโลมํ โวทานสฺส.\csromanspacer{} ผลํ ผลสฺส.\csromanspacer{} อนุโลมํ ผลสมาปตฺติยา.\csromanspacer{} นิโรธา วุฏฺฐหนฺตสฺส เนวสญฺญานาสญฺญายตนํ ผลสมาปตฺติยา อนนฺตรปจฺจเยน ปจฺจยา.\csromanspacer{}\csromanspacer{}อนิยฺยานิโก ธมฺโม นิยฺยานิกสฺส ธมฺมสฺส อนนฺตรปจฺจเยน ปจฺจโย.\csromanspacer{} โคตฺรภุ มคฺคสฺส.\csromanspacer{} โวทานํ มคฺคสฺส อนนฺตรปจฺจเยน ปจฺจโย. (2)}

\prose{สมนนฺตรปจฺจเยน ปจฺจโย.\csromanspacer{} สหชาตปจฺจเยน ปจฺจโย.\csromanspacer{} ปญฺจ.\csromanspacer{} อญฺญมญฺญปจฺจเยน ปจฺจโย.\csromanspacer{} เทฺว.\csromanspacer{} นิสฺสยปจฺจเยน ปจฺจโย.\csromanspacer{} สตฺต.}

\csromanheader{2}{\textbf{อุปนิสฺสย}}
\proseitem{318}{นิยฺยานิโก ธมฺโม นิยฺยานิกสฺส ธมฺมสฺส อุปนิสฺสยปจฺจเยน ปจฺจโย.\csromanspacer{} \textbf{ปกตูปนิสฺสโย\csromandash{}}ปฐโม มคฺโค ทุติยสฺส มคฺคสฺส อุปนิสฺสยปจฺจเยน ปจฺจโย ฯเปฯ.\csromanspacer{} ตติโย มคฺโค จตุตฺถสฺส มคฺคสฺส อุปนิสฺสยปจฺจเยน ปจฺจโย.\csromanspacer{}\csromanspacer{}นิยฺยานิโก ธมฺโม อนิยฺยานิกสฺส ธมฺมสฺส อุปนิสฺสยปจฺจเยน ปจฺจโย.\csromanspacer{} \textbf{อารมฺมณูปนิสฺสโย, อนนฺตรูปนิสฺสโย,} \textbf{ปกตูปนิสฺสโย ฯเปฯ.}\csromanspacer{} ปกตูปนิสฺสโย\csromandash{}อริยา มคฺคํ อุปนิสฺสาย อนุปฺปนฺนํ สมาปตฺติํ อุปฺปาเทนฺติ, อุปฺปนฺนํ สมาปชฺชนฺติ, สงฺขาเร อนิจฺจโต ฯเปฯ วิปสฺสนฺติ.\csromanspacer{} มคฺโค อริยานํ อตฺถปฏิสมฺภิทาย ฯเปฯ.\csromanspacer{} ปฏิภานปฏิสมฺภิทาย.}

\vspace{\tipitakaparskip}
\csromancenter{นิยฺยานิกทุก ฐานาฐานโกสลฺลสฺส อุปนิสฺสยปจฺจเยน ปจฺจโย.\csromanspacer{} มคฺโค ผลสมาปตฺติยา อุปนิสฺสยปจฺจเยน ปจฺจโย. (2)}
\proseitem{319}{\csromanfolio{177}อนิยฺยานิโก ธมฺโม อนิยฺยานิกสฺส ธมฺมสฺส อุปนิสฺสยปจฺจเยน ปจฺจโย.\csromanspacer{} \textbf{อารมฺมณูปนิสฺสโย, อนนฺตรูปนิสฺสโย, ปกตูปนิสฺสโย ฯเปฯ.}\csromanspacer{} ปกตูปนิสฺสโย\csromandash{}อนิยฺยานิกํ สทฺธํ อุปนิสฺสาย ทานํ เทติ.\csromanspacer{} สีลํ ฯเปฯ.\csromanspacer{} อุโปสถกมฺมํ กโรติ.\csromanspacer{} ฌานํ อุปฺปาเทติ.\csromanspacer{} วิปสฺสนํ.\csromanspacer{} อภิญฺญํ.\csromanspacer{} สมาปตฺติํ อุปฺปาเทติ.\csromanspacer{} มานํ ชปฺเปติ.\csromanspacer{} ทิฏฺฐิํ คณฺหาติ.\csromanspacer{} อนิยฺยานิกํ สีลํ ฯเปฯ ปญฺญํ.\csromanspacer{} ราคํ ฯเปฯ ปตฺถนํ.\csromanspacer{} อุตุํ.\csromanspacer{} โภชนํ.\csromanspacer{} เสนาสนํ อุปนิสฺสาย ทานํ เทติ ฯเปฯ.\csromanspacer{} สมาปตฺติํ อุปฺปาเทติ.\csromanspacer{} ปาณํ หนติ ฯเปฯ สํฆํ ภินฺทติ.\csromanspacer{} อนิยฺยานิกา สทฺธา ฯเปฯ.\csromanspacer{} เสนาสนํ อนิยฺยานิกาย สทฺธาย ฯเปฯ ปตฺถนาย.\csromanspacer{} กายิกสฺส สุขสฺส.\csromanspacer{} กายิกสฺส ทุกฺขสฺส.\csromanspacer{} ผลสมาปตฺติยา อุปนิสฺสยปจฺจเยน ปจฺจโย.\csromanspacer{}\csromanspacer{}อนิยฺยานิโก ธมฺโม นิยฺยานิกสฺส ธมฺมสฺส อุปนิสฺสยปจฺจเยน ปจฺจโย.\csromanspacer{} \textbf{อารมฺมณูปนิสฺสโย, อนนฺตรูปนิสฺสโย, ปกตูปนิสฺสโย ฯเปฯ.}\csromanspacer{} ปกตูปนิสฺสโย\csromandash{}ปฐมสฺส มคฺคสฺส ปริกมฺมํ ปฐมสฺส มคฺคสฺส อุปนิสฺสยปจฺจเยน ปจฺจโย ฯเปฯ.\csromanspacer{} จตุตฺถสฺส มคฺคสฺส ปริกมฺมํ จตุตฺถสฺส มคฺคสฺส อุปนิสฺสยปจฺจเยน ปจฺจโย. (2)}

\csromanheader{2}{\textbf{ปุเรชาตาทิ}}
\proseitem{320}{อนิยฺยานิโก ธมฺโม อนิยฺยานิกสฺส ธมฺมสฺส ปุเรชาตปจฺจเยน ปจฺจโย.\csromanspacer{} เทฺว. (อรูปทุกสทิสา กาตพฺพา.) ปจฺฉาชาตปจฺจเยน ปจฺจโย.\csromanspacer{} เทฺว.\csromanspacer{} อาเสวนปจฺจเยน ปจฺจโย.\csromanspacer{} เทฺว.}

\csromanheader{2}{\textbf{กมฺม}}
\proseitem{321}{นิยฺยานิโก ธมฺโม นิยฺยานิกสฺส ธมฺมสฺส กมฺมปจฺจเยน ปจฺจโย.\csromanspacer{} นิยฺยานิกา เจตนา สมฺปยุตฺตกานํ ขนฺธานํ กมฺมปจฺจเยน ปจฺจโย.\csromanspacer{}\csromanspacer{}นิยฺยานิโก ธมฺโม อนิยฺยานิกสฺส ธมฺมสฺส กมฺมปจฺจเยน ปจฺจโย.\csromanspacer{} \textbf{สหชาตา}, นานากฺขณิกา.\csromanspacer{}\csromanspacer{}\textbf{สหชาตา} นิยฺยานิกา เจตนา จิตฺตสมุฏฺฐานานํ รูปานํ กมฺมปจฺจเยน ปจฺจโย.\csromanspacer{} \textbf{นานากฺขณิกา} นิยฺยานิกา เจตนา ผลสฺส กมฺมปจฺจเยน ปจฺจโย. (มูลํ กาตพฺพํ.) นิยฺยานิกา เจตนา สมฺปยุตฺตกานํ ขนฺธานํ จิตฺตสมุฏฺฐานานญฺจ รูปานํ กมฺมปจฺจเยน ปจฺจโย. (3)}

\proseitem{322}{\csromanfolio{178}อนิยฺยานิโก ธมฺโม อนิยฺยานิกสฺส ธมฺมสฺส กมฺมปจฺจเยน ปจฺจโย.\csromanspacer{} \textbf{สหชาตา}, นานากฺขณิกา.\csromanspacer{}\csromanspacer{}\textbf{สหชาตา} (สํขิตฺตํ.) \textbf{นานากฺขณิกา} อนิยฺยานิกา เจตนา วิปากานํ ขนฺธานํ กฏตฺตา จ รูปานํ กมฺมปจฺจเยน ปจฺจโย.\csromanspacer{} วิปากปจฺจเยน ปจฺจโย.\csromanspacer{} เอกํเยว.\csromanspacer{} อาหารปจฺจเยน ปจฺจโย.\csromanspacer{} จตฺตาริ.\csromanspacer{} อินฺทฺริยปจฺจเยน ปจฺจโย.\csromanspacer{} จตฺตาริ.\csromanspacer{} ฌานปจฺจเยน ปจฺจโย.\csromanspacer{} จตฺตาริ.\csromanspacer{} มคฺคปจฺจเยน ปจฺจโย.\csromanspacer{} จตฺตาริ.\csromanspacer{} สมฺปยุตฺตปจฺจเยน ปจฺจโย.\csromanspacer{} เทฺว.\csromanspacer{} วิปฺปยุตฺตปจฺจเยน ปจฺจโย.\csromanspacer{} ตีณิ. (อรูปทุกสทิสา กาตพฺพา.) อตฺถิปจฺจเยน ปจฺจโย.\csromanspacer{} สตฺต. (อรูปทุกสทิสา กาตพฺพา.\csromanspacer{} อามสนา นานาปทาเยว.) นตฺถิปจฺจเยน ปจฺจโย.\csromanspacer{} ตีณิ.\csromanspacer{} วิคตปจฺจเยน ปจฺจโย.\csromanspacer{} ตีณิ.\csromanspacer{} อวิคตปจฺจเยน ปจฺจโย.\csromanspacer{} สตฺต.}

\csromanheader{2}{1. ปจฺจยานุโลม 2. สงฺขฺยาวาร}
\proseitem{323}{เหตุยา จตฺตาริ, อารมฺมเณ ตีณิ, อธิปติยา ปญฺจ, อนนฺตเร ตีณิ, สมนนฺตเร ตีณิ, สหชาเต ปญฺจ, อญฺญมญฺเญ เทฺว, นิสฺสเย สตฺต, อุปนิสฺสเย จตฺตาริ, ปุเรชาเต เทฺว, ปจฺฉาชาเต เทฺว, อาเสวเน เทฺว, กมฺเม จตฺตาริ, วิปาเก เอกํ, อาหาเร จตฺตาริ, อินฺทฺริเย จตฺตาริ, ฌาเน จตฺตาริ, มคฺเค จตฺตาริ, สมฺปยุตฺเต เทฺว, วิปฺปยุตฺเต ตีณิ, อตฺถิยา สตฺต, นตฺถิยา ตีณิ, วิคเต ตีณิ, อวิคเต สตฺต.}

\vspace{\tipitakaparskip}
\csromancenter{อนุโลมํ.}
\csromanheader{2}{2. ปจฺจนียุทฺธาร}
\proseitem{324}{นิยฺยานิโก ธมฺโม นิยฺยานิกสฺส ธมฺมสฺส สหชาตปจฺจเยน ปจฺจโย.\csromanspacer{} อุปนิสฺสยปจฺจเยน ปจฺจโย.\csromanspacer{}\csromanspacer{}นิยฺยานิโก ธมฺโม อนิยฺยานิกสฺส ธมฺมสฺส อารมฺมณปจฺจเยน ปจฺจโย.\csromanspacer{} สหชาตปจฺจเยน ปจฺจโย.\csromanspacer{} อุปนิสฺสยปจฺจเยน ปจฺจโย.\csromanspacer{} ปจฺฉาชาตปจฺจเยน ปจฺจโย.\csromanspacer{}\csromanspacer{}นิยฺยานิโก ธมฺโม นิยฺยานิกสฺส จ อนิยฺยานิกสฺส จ ธมฺมสฺส สหชาตปจฺจเยน ปจฺจโย. (3)}

\proseitem{325}{อนิยฺยานิโก ธมฺโม อนิยฺยานิกสฺส ธมฺมสฺส อารมฺมณปจฺจเยน ปจฺจโย.\csromanspacer{} สหชาตปจฺจเยน ปจฺจโย.\csromanspacer{} อุปนิสฺสยปจฺจเยน ปจฺจโย.}

\vspace{\tipitakaparskip}
\csromancenter{นิยฺยานิกทุก ปุเรชาตปจฺจเยน ปจฺจโย.\csromanspacer{} ปจฺฉาชาตปจฺจเยน ปจฺจโย.\csromanspacer{} กมฺมปจฺจเยน ปจฺจโย.\csromanspacer{} อาหารปจฺจเยน ปจฺจโย.\csromanspacer{} อินฺทฺริยปจฺจเยน ปจฺจโย.\csromanspacer{}\csromanspacer{}อนิยฺยานิโก ธมฺโม นิยฺยานิกสฺส ธมฺมสฺส อุปนิสฺสยปจฺจเยน ปจฺจโย.\csromanspacer{} ปุเรชาตปจฺจเยน ปจฺจโย. (2)}
\prose{\csromanfolio{179}นิยฺยานิโก จ อนิยฺยานิโก จ ธมฺมา นิยฺยานิกสฺส ธมฺมสฺส สหชาตํ, ปุเรชาตํ.\csromanspacer{}\csromanspacer{}นิยฺยานิโก จ อนิยฺยานิโก จ ธมฺมา อนิยฺยานิกสฺส ธมฺมสฺส สหชาตํ, ปจฺฉาชาตํ, อาหารํ, อินฺทฺริยํ. (2)}

\csromanheader{2}{2. ปจฺจยปจฺจนีย 2. สงฺขฺยาวาร}
\proseitem{326}{นเหตุยา สตฺต, น-อารมฺมเณ สตฺต, น-อธิปติยา สตฺต, น-อนนฺตเร นสมนนฺตเร สตฺต, นสหชาเต ปญฺจ, น-อญฺญมญฺเญ ปญฺจ, นนิสฺสเย ปญฺจ, น-อุปนิสฺสเย สตฺต, นปุเรชาเต ฉ, นปจฺฉาชาเต สตฺต ฯเปฯ นสมฺปยุตฺเต ปญฺจ, นวิปฺปยุตฺเต จตฺตาริ, โน-อตฺถิยา จตฺตาริ, โนนตฺถิยา สตฺต, โนวิคเต สตฺต, โน-อวิคเต จตฺตาริ.}

\csromanheader{2}{3. ปจฺจยานุโลมปจฺจนีย}
\proseitem{327}{เหตุปจฺจยา น-อารมฺมเณ จตฺตาริ, น-อธิปติยา จตฺตาริ, น-อนนฺตเร นสมนนฺตเร จตฺตาริ, น-อญฺญมญฺเญ เทฺว, น-อุปนิสฺสเย จตฺตาริ ฯเปฯ นสมฺปยุตฺเต เทฺว, นวิปฺปยุตฺเต เทฺว, โนนตฺถิยา จตฺตาริ, โนวิคเต จตฺตาริ.}

\csromanheader{2}{4. ปจฺจยปจฺจนียานุโลม}
\proseitem{328}{นเหตุปจฺจยา อารมฺมเณ ตีณิ, อธิปติยา ปญฺจ, (อนุโลมมาติกา วิตฺถาเรตพฺพา.) ฯเปฯ อวิคเต สตฺต.}

\nitthitam{นิยฺยานิกทุกํ นิฏฺฐิตํ.}
\csromanheader{2}{98. นิยตทุก 1-6. ปฏิจฺจาทิวาร}
\proseitem{329}{\csromanfolio{180}นิยตํ ธมฺมํ ปฏิจฺจ นิยโต ธมฺโม อุปฺปชฺชติ เหตุปจฺจยา.\csromanspacer{} นิยตํ เอกํ ขนฺธํ ปฏิจฺจ ตโย ขนฺธา ฯเปฯ เทฺว ขนฺเธ ฯเปฯ.\csromanspacer{} นิยตํ ธมฺมํ ปฏิจฺจ อนิยโต ธมฺโม อุปฺปชฺชติ เหตุปจฺจยา.\csromanspacer{} นิยเต ขนฺเธ ปฏิจฺจ จิตฺตสมุฏฺฐานํ รูปํ. (สํขิตฺตํ.\csromanspacer{} ปญฺจปิ ปญฺหา กาตพฺพา, ยถา นิยฺยานิกทุกํ.\csromanspacer{} เอวํ ปฏิจฺจวาโรปิ สหชาตวาโรปิ ปจฺจยวาโรปิ นิสฺสยวาโรปิ สํสฏฺฐวาโรปิ สมฺปยุตฺตวาโรปิ กาตพฺพา.\csromanspacer{}\csromanspacer{}นินฺนานากรณํ อามสนํ นานํ.)}

\csromanmatikaanchor{mtk.27}
\csromantocmarkref{section}{98. นิยตทุก}{mtk.27}
\bookmark[dest={mtk.27},level=1]{98. นิยตทุก}
\csromanheader{1}{98. นิยตทุก 7. ปญฺหาวาร}
\proseitem{330}{นิยโต ธมฺโม นิยตสฺส ธมฺมสฺส เหตุปจฺจเยน ปจฺจโย.\csromanspacer{} จตฺตาริ. (นิยฺยานิกทุกสทิสา นินฺนานากรณา.)}

\csromanheader{2}{\textbf{อารมฺมณ}}
\proseitem{331}{นิยโต ธมฺโม อนิยตสฺส ธมฺมสฺส อารมฺมณปจฺจเยน ปจฺจโย.\csromanspacer{} อริยา มคฺคา วุฏฺฐหิตฺวา มคฺคํ ปจฺจเวกฺขนฺติ.\csromanspacer{} นิยเต ปหีเน กิเลเส ปจฺจเวกฺขนฺติ.\csromanspacer{} ปุพฺเพ สมุทาจิณฺเณ ฯเปฯ.\csromanspacer{} นิยเต ขนฺเธ อนิจฺจโต ฯเปฯ วิปสฺสติ.\csromanspacer{} เจโตปริยญาเณน นิยตจิตฺตสมงฺคิสฺส จิตฺตํ ชานาติ.\csromanspacer{} นิยตา ขนฺธา เจโตปริยญาณสฺส, ปุพฺเพนิวาสานุสฺสติญาณสฺส, ยถากมฺมูปคญาณสฺส, อนาคตํสญาณสฺส, อาวชฺชนาย อารมฺมณปจฺจเยน ปจฺจโย. (1)}

\prose{อนิยโต ธมฺโม อนิยตสฺส ธมฺมสฺส อารมฺมณปจฺจเยน ปจฺจโย.\csromanspacer{} ทานํ ทตฺวา, สีลํ ฯเปฯ อุโปสถกมฺมํ ฯเปฯ.\csromanspacer{} ปุพฺเพ สุจิณฺณานิ ฯเปฯ.\csromanspacer{} ฌานา วุฏฺฐหิตฺวา ฌานํ ปจฺจเวกฺขติ, อสฺสาเทติ อภินนฺทติ, ตํ อารพฺภ อนิยโต ราโค อุปฺปชฺชติ.\csromanspacer{} ทิฏฺฐิ ฯเปฯ วิจิกิจฺฉา ฯเปฯ อุทฺธจฺจํ ฯเปฯ.\csromanspacer{} อนิยตํ โทมนสฺสํ อุปฺปชฺชติ.\csromanspacer{} อริยา ผลํ ปจฺจเวกฺขนฺติ, นิพฺพานํ ปจฺจเวกฺขนฺติ.\csromanspacer{} นิพฺพานํ}

\vspace{\tipitakaparskip}
\csromancenter{นิยตทุก โคตฺรภุสฺส, โวทานสฺส, ผลสฺส, อาวชฺชนาย อารมฺมณปจฺจเยน ปจฺจโย.\csromanspacer{} อริยา อนิยเต ปหีเน กิเลเส ปจฺจเวกฺขนฺติ.\csromanspacer{} วิกฺขมฺภิเต กิเลเส ฯเปฯ.\csromanspacer{} ปุพฺเพ สมุทาจิณฺเณ ฯเปฯ.\csromanspacer{} จกฺขุํ ฯเปฯ.\csromanspacer{} วตฺถุํ อนิยเต ขนฺเธ อนิจฺจโต ฯเปฯ วิปสฺสติ, อสฺสาเทติ อภินนฺทติ, ตํ อารพฺภ อนิยโต ราโค ฯเปฯ โทมนสฺสํ อุปฺปชฺชติ.\csromanspacer{} ทิพฺเพน จกฺขุนา รูปํ ปสฺสติ ฯเปฯ.\csromanspacer{} อาวชฺชนาย อารมฺมณปจฺจเยน ปจฺจโย.\csromanspacer{}\csromanspacer{}อนิยโต ธมฺโม นิยตสฺส ธมฺมสฺส อารมฺมณปจฺจเยน ปจฺจโย.\csromanspacer{} นิพฺพานํ มคฺคสฺส อารมฺมณปจฺจเยน ปจฺจโย.\csromanspacer{} รูปชีวิตินฺทฺริยํ มาตุฆาติกมฺมสฺส.\csromanspacer{} ปิตุฆาติกมฺมสฺส.\csromanspacer{} อรหนฺตฆาติกมฺมสฺส.\csromanspacer{} รุหิรุปฺปาทกมฺมสฺส อารมฺมณปจฺจเยน ปจฺจโย.\csromanspacer{} ยํ วตฺถุํ อามสนฺตสฺส มิจฺฉตฺตนิยตา ขนฺธา อุปฺปชฺชนฺติ, ตํ วตฺถุ มิจฺฉตฺตนิยตานํ ขนฺธานํ อารมฺมณปจฺจเยน ปจฺจโย. (2)}
\csromanheader{2}{\textbf{อธิปติ}}
\proseitem{332}{\csromanfolio{181}นิยโต ธมฺโม นิยตสฺส ธมฺมสฺส อธิปติปจฺจเยน ปจฺจโย.\csromanspacer{} \textbf{สหชาตาธิปติ\csromandash{}}นิยตาธิปติ สมฺปยุตฺตกานํ ขนฺธานํ อธิปติปจฺจเยน ปจฺจโย.\csromanspacer{}\csromanspacer{}นิยโต ธมฺโม อนิยตสฺส ธมฺมสฺส อธิปติปจฺจเยน ปจฺจโย.\csromanspacer{} อารมฺมณาธิปติ, สหชาตาธิปติ.\csromanspacer{}\csromanspacer{}\textbf{อารมฺมณาธิปติ\csromandash{}}อริยา มคฺคา วุฏฺฐหิตฺวา มคฺคํ ครุํ กตฺวา ปจฺจเวกฺขนฺติ.\csromanspacer{} \textbf{สหชาตาธิปติ\csromandash{}}นิยตาธิปติ จิตฺตสมุฏฺฐานานํ รูปานํ อธิปติปจฺจเยน ปจฺจโย.\csromanspacer{}\csromanspacer{}นิยโต ธมฺโม นิยตสฺส จ อนิยตสฺส จ ธมฺมสฺส อธิปติปจฺจเยน ปจฺจโย.\csromanspacer{} \textbf{สหชาตาธิปติ\csromandash{}} นิยตาธิปติ สมฺปยุตฺตกานํ ขนฺธานํ จิตฺตสมุฏฺฐานานญฺจ รูปานํ อธิปติปจฺจเยน ปจฺจโย. (3)}

\proseitem{333}{อนิยโต ธมฺโม อนิยตสฺส ธมฺมสฺส อธิปติปจฺจเยน ปจฺจโย.\csromanspacer{} อารมฺมณาธิปติ, สหชาตาธิปติ.\csromanspacer{}\csromanspacer{}\textbf{อารมฺมณาธิปติ\csromandash{}}ทานํ ทตฺวา, สีลํ ฯเปฯ อุโปสถกมฺมํ กตฺวา ตํ ครุํ กตฺวา ปจฺจเวกฺขติ, อสฺสาเทติ อภินนฺทติ, ตํ ครุํ กตฺวา อนิยโต ราโค อุปฺปชฺชติ.\csromanspacer{} ทิฏฺฐิ อุปฺปชฺชติ.\csromanspacer{} ปุพฺเพ ฯเปฯ.\csromanspacer{} ฌานา ฯเปฯ.\csromanspacer{} อริยา ผลํ ครุํ กตฺวา ปจฺจเวกฺขนฺติ, นิพฺพานํ ครุํ กตฺวา ปจฺจเวกฺขนฺติ.\csromanspacer{} นิพฺพานํ โคตฺรภุสฺส, โวทานสฺส, ผลสฺส อธิปติปจฺจเยน ปจฺจโย.\csromanspacer{} จกฺขุํ ฯเปฯ.\csromanspacer{} วตฺถุํ อนิยเต ขนฺเธ ครุํ กตฺวา อสฺสาเทติ อภินนฺทติ, ตํ ครุํ กตฺวา อนิยโต ราโค อุปฺปชฺชติ.\csromanspacer{} ทิฏฺฐิ อุปฺปชฺชติ.\csromanspacer{} \textbf{สหชาตาธิปติ\csromandash{}} อนิยตาธิปติ สมฺปยุตฺตกานํ ขนฺธานํ \csromanfolio{182}จิตฺตสมุฏฺฐานานญฺจ รูปานํ อธิปติปจฺจเยน ปจฺจยา.\csromanspacer{}\csromanspacer{}อนิยโต ธมฺโม นิยตสฺส ธมฺมสฺส อธิปติปจฺจเยน ปจฺจโย.\csromanspacer{} \textbf{อารมฺมณาธิปติ\csromandash{}}นิพฺพานํ มคฺคสฺส อธิปติปจฺจเยน ปจฺจโย. (2)}

\csromanheader{2}{\textbf{อนนฺตราทิ}}
\proseitem{334}{นิยโต ธมฺโม อนิยตสฺส ธมฺมสฺส อนนฺตรปจฺจเยน ปจฺจโย.\csromanspacer{} มคฺโค ผลสฺส อนนฺตรปจฺจเยน ปจฺจโย.\csromanspacer{} นิยตา ขนฺธา วุฏฺฐานสฺส อนนฺตรปจฺจเยน ปจฺจโย. (1)}

\prose{อนิยโต ธมฺโม อนิยตสฺส ธมฺมสฺส อนนฺตรปจฺจเยน ปจฺจโย.\csromanspacer{} ปุริมา ปุริมา อนิยตา ขนฺธา ปจฺฉิมานํ ปจฺฉิมานํ อนิยตานํ ขนฺธานํ อนนฺตรปจฺจเยน ปจฺจโย.\csromanspacer{} อนุโลมํ โคตฺรภุสฺส.\csromanspacer{} อนุโลมํ โวทานสฺส.\csromanspacer{} ผลํ ผลสฺส.\csromanspacer{} อนุโลมํ ผลสมาปตฺติยา.\csromanspacer{} นิโรธา วุฏฺฐหนฺตสฺส เนวสญฺญานาสญฺญายตนํ ผลสมาปตฺติยา อนนฺตรปจฺจเยน ปจฺจโย.\csromanspacer{}\csromanspacer{}อนิยโต ธมฺโม นิยตสฺส ธมฺมสฺส อนนฺตรปจฺจเยน ปจฺจโย.\csromanspacer{} อนิยตํ โทมนสฺสํ นิยตสฺส โทมนสฺสสฺส.\csromanspacer{} อนิยตา มิจฺฉาทิฏฺฐิ นิยตมิจฺฉาทิฏฺฐิยา อนนฺตรปจฺจเยน ปจฺจโย.\csromanspacer{} โคตฺรภุ มคฺคสฺส.\csromanspacer{} โวทานํ มคฺคสฺส อนนฺตรปจฺจเยน ปจฺจโย. (2)}

\prose{สมนนฺตรปจฺจเยน ปจฺจโย.\csromanspacer{} สหชาตปจฺจเยน ปจฺจโย.\csromanspacer{} ปญฺจ.\csromanspacer{} อญฺญมญฺญปจฺจเยน ปจฺจโย.\csromanspacer{} เทฺว.\csromanspacer{} นิสฺสยปจฺจเยน ปจฺจโย.\csromanspacer{} สตฺต.}

\csromanheader{2}{\textbf{อุปนิสฺสย}}
\proseitem{335}{นิยโต ธมฺโม นิยตสฺส ธมฺมสฺส อุปนิสฺสยปจฺจเยน ปจฺจโย.\csromanspacer{} ปฺ\textbf{อกตูปนิสฺสโย\csromandash{}}มาตุฆาติกมฺมํ มาตุฆาติกมฺมสฺส.\csromanspacer{} ปิตุฆาติกมฺมสฺส.\csromanspacer{} อรหนฺตฆาติกมฺมสฺส.\csromanspacer{} รุหิรุปฺปาทกมฺมสฺส.\csromanspacer{} สํฆเภทกมฺมสฺส.\csromanspacer{} นิยตมิจฺฉาทิฏฺฐิยา อุปนิสฺสยปจฺจเยน ปจฺจโย. (จกฺกํ.) ปฐโม มคฺโค ทุติยสฺส มคฺคสฺส ฯเปฯ.\csromanspacer{} ตติโย มคฺโค จตุตฺถสฺส มคฺคสฺส อุปนิสฺสยปจฺจเยน ปจฺจโย.\csromanspacer{}\csromanspacer{}นิยโต ธมฺโม อนิยตสฺส ธมฺมสฺส อุปนิสฺสยปจฺจเยน ปจฺจโย.\csromanspacer{} \textbf{อารมฺมณูปนิสฺสโย, อนนฺตรูปนิสฺสโย, ปกตูปนิสฺสโย ฯเปฯ.}\csromanspacer{} ปฺ\textbf{อกตูปนิสฺสโย\csromandash{}}มาตรํ ชีวิตา โวโรเปตฺวา ฯเปฯ สํฆํ ภินฺทิตฺวา, ตสฺส ปฏิฆาตตฺถาย ทานํ เทติ.\csromanspacer{} สีลํ สมาทิยติ.}

\vspace{\tipitakaparskip}
\csromancenter{นิยตทุก อุโปสถกมฺมํ กโรติ.\csromanspacer{} อริยา มคฺคํ อุปนิสฺสาย อนุปฺปนฺนํ สมาปตฺติํ อุปฺปาเทนฺติ.\csromanspacer{} อุปฺปนฺนํ สมาปชฺชนฺติ ฯเปฯ.\csromanspacer{} ฐานาฐานโกสลฺลสฺส อุปนิสฺสยปจฺจเยน ปจฺจโย.\csromanspacer{} มคฺโค ผลสมาปตฺติยา อุปนิสฺสยปจฺจเยน ปจฺจโย. (2)}
\proseitem{336}{\csromanfolio{183}อนิยโต ธมฺโม อนิยตสฺส ธมฺมสฺส อุปนิสฺสยปจฺจเยน ปจฺจโย.\csromanspacer{} \textbf{อารมฺมณูปนิสฺสโย, อนนฺตรูปนิสฺสโย, ปกตูปนิสฺสโย ฯเปฯ.}\csromanspacer{} ปกตูปนิสฺสโย\csromandash{}อนิยตํ สทฺธํ อุปนิสฺสาย ทานํ เทติ.\csromanspacer{} สีลํ ฯเปฯ.\csromanspacer{} อุโปสถกมฺมํ กโรติ.\csromanspacer{} ฌานํ อุปฺปาเทติ.\csromanspacer{} วิปสฺสนํ ฯเปฯ.\csromanspacer{} อภิญฺญํ ฯเปฯ.\csromanspacer{} สมาปตฺติํ อุปฺปาเทติ.\csromanspacer{} มานํ ชปฺเปติ.\csromanspacer{} ทิฏฺฐิํ คณฺหาติ.\csromanspacer{} อนิยตํ สีลํ ฯเปฯ ปญฺญํ.\csromanspacer{} ราคํ ฯเปฯ.\csromanspacer{} ปตฺถนํ.\csromanspacer{} กายิกํ สุขํ.\csromanspacer{} กายิกํ ทุกฺขํ.\csromanspacer{} อุตุํ.\csromanspacer{} โภชนํ.\csromanspacer{} เสนาสนํ อุปนิสฺสาย ทานํ เทติ ฯเปฯ.\csromanspacer{} สมาปตฺติํ อุปฺปาเทติ.\csromanspacer{} ปาณํ หนติ ฯเปฯ นิคมฆาตํ กโรติ.\csromanspacer{} อนิยตา สทฺธา ฯเปฯ.\csromanspacer{} เสนาสนํ อนิยตาย สทฺธาย ฯเปฯ ปตฺถนาย.\csromanspacer{} กายิกสฺส สุขสฺส.\csromanspacer{} กายิกสฺส ทุกฺขสฺส, ผลสมาปตฺติยา อุปนิสฺสยปจฺจเยน ปจฺจโย. (1)}

\prose{อนิยโต ธมฺโม นิยตสฺส ธมฺมสฺส อุปนิสฺสยปจฺจเยน ปจฺจโย.\csromanspacer{} \textbf{อารมฺมณูปนิสฺสโย, อนนฺตรูปนิสฺสโย, ปกตูปนิสฺสโย ฯเปฯ.}\csromanspacer{} ปกตูปนิสฺสโย\csromandash{}อนิยตํ ราคํ อุปนิสฺสาย มาตรํ ชีวิตา โวโรเปติ ฯเปฯ สํฆํ ภินฺทติ.\csromanspacer{} อนิยตํ โทสํ ฯเปฯ.\csromanspacer{} เสนาสนํ อุปนิสฺสาย มาตรํ ชีวิตา โวโรเปติ ฯเปฯ สํฆํ ภินฺทติ.\csromanspacer{} อนิยโต ราโค.\csromanspacer{} โทสํ ฯเปฯ.\csromanspacer{} เสนาสนํ มาตุฆาติกมฺมสฺส ฯเปฯ สํฆเภทกมฺมสฺส อุปนิสฺสยปจฺจเยน ปจฺจโย.\csromanspacer{} ปฐมสฺส มคฺคสฺส ปริกมฺมํ ปฐมสฺส มคฺคสฺส ฯเปฯ.\csromanspacer{} จตุตฺถสฺส มคฺคสฺส ปริกมฺมํ จตุตฺถสฺส มคฺคสฺส อุปนิสฺสยปจฺจเยน ปจฺจโย. (2)}

\csromanheader{2}{\textbf{ปุเรชาตาทิ}}
\proseitem{337}{อนิยโต ธมฺโม อนิยตสฺส ธมฺมสฺส ปุเรชาตปจฺจเยน ปจฺจโย.\csromanspacer{} อารมฺมณปุเรชาตํ, วตฺถุปุเรชาตํ. (สํขิตฺตํ.) อนิยโต ธมฺโม นิยตสฺส ธมฺมสฺส ปุเรชาตปจฺจเยน ปจฺจโย.\csromanspacer{} อารมฺมณปุเรชาตํ, วตฺถุปุเรชาตํ.\csromanspacer{}\csromanspacer{}\textbf{อารมฺมณปุเรชาตํ\csromandash{}} รูปชีวิตินฺทฺริยํ มาตุฆาติกมฺมสฺส.\csromanspacer{} ปิตุฆาติกมฺมสฺส.\csromanspacer{} อรหนฺตฆาติกมฺมสฺส.\csromanspacer{} รุหิรุปฺปาทกมฺมสฺส ปุเรชาตปจฺจเยน ปจฺจโย.\csromanspacer{} \textbf{วตฺถุปุเรชาตํ\csromandash{}}วตฺถุ นิยตานํ ขนฺธานํ ปุเรชาตปจฺจเยน ปจฺจโย. (2)}

\prose{ปจฺฉาชาตปจฺจเยน ปจฺจโย.\csromanspacer{} เทฺว.\csromanspacer{} อาเสวนปจฺจเยน ปจฺจโย.\csromanspacer{} เทฺว.}

\csromanheader{2}{\textbf{กมฺม}}
\proseitem{338}{\csromanfolio{184}นิยโต ธมฺโม นิยตสฺส ธมฺมสฺส กมฺมปจฺจเยน ปจฺจโย.\csromanspacer{} นิยตา เจตนา สมฺปยุตฺตกานํ ขนฺธานํ กมฺมปจฺจเยน ปจฺจโย.\csromanspacer{}\csromanspacer{}นิยโต ธมฺโม อนิยตสฺส ธมฺมสฺส กมฺมปจฺจเยน ปจฺจโย.\csromanspacer{} \textbf{สหชาตา}, นานากฺขณิกา.\csromanspacer{}\csromanspacer{}\textbf{สหชาตา} นิยตา เจตนา จิตฺตสมุฏฺฐานานํ รูปานํ กมฺมปจฺจเยน ปจฺจโย.\csromanspacer{} \textbf{นานากฺขณิกา} นิยตา เจตนา วิปากานํ ขนฺธานํ กฏตฺตา จ รูปานํ กมฺมปจฺจเยน ปจฺจโย. (มูลํ กาตพฺพํ.) นิยตา เจตนา สมฺปยุตฺตกานํ ขนฺธานํ จิตฺตสมุฏฺฐานานญฺจ รูปานํ กมฺมปจฺจเยน ปจฺจโย. (3)}

\prose{อนิยโต ธมฺโม อนิยตสฺส ธมฺมสฺส กมฺมปจฺจเยน ปจฺจโย.\csromanspacer{} \textbf{สหชาตา}, นานากฺขณิกา. (สํขิตฺตํ.)}

\prose{วิปากปจฺจเยน ปจฺจโย.\csromanspacer{} เอกํ.\csromanspacer{} อาหารปจฺจเยน ปจฺจโย.\csromanspacer{} จตฺตาริ.\csromanspacer{} อินฺทฺริยปจฺจเยน ปจฺจโย.\csromanspacer{} จตฺตาริ.\csromanspacer{} ฌานปจฺจเยน ปจฺจโย.\csromanspacer{} จตฺตาริ.\csromanspacer{} มคฺคปจฺจเยน ปจฺจโย.\csromanspacer{} จตฺตาริ.\csromanspacer{} สมฺปยุตฺตปจฺจเยน ปจฺจโย.\csromanspacer{} เทฺว.\csromanspacer{} วิปฺปยุตฺตปจฺจเยน ปจฺจโย.\csromanspacer{} ตีณิ. (อรูปทุกสทิสํ.) อตฺถิปจฺจเยน ปจฺจโย.\csromanspacer{} สตฺต. (อรูปาวจรทุกสทิสํ.) นตฺถิปจฺจเยน ปจฺจโย.\csromanspacer{} วิคตปจฺจเยน ปจฺจโย.\csromanspacer{} อวิคตปจฺจเยน ปจฺจโย.\csromanspacer{} สตฺต.}

\csromanheader{2}{1. ปจฺจยานุโลม 2. สงฺขฺยาวาร}
\proseitem{339}{เหตุยา จตฺตาริ, อารมฺมเณ ตีณิ, อธิปติยา ปญฺจ, อนนฺตเร ตีณิ, สมนนฺตเร ตีณิ, สหชาเต ปญฺจ, อญฺญมญฺเญ เทฺว, นิสฺสเย สตฺต, อุปนิสฺสเย จตฺตาริ, ปุเรชาเต เทฺว, ปจฺฉาชาเต เทฺว, อาเสวเน เทฺว, กมฺเม จตฺตาริ, วิปาเก เอกํ, อาหาเร จตฺตาริ, อินฺทฺริเย ฌาเน จตฺตาริ, มคฺเค จตฺตาริ, สมฺปยุตฺเต เทฺว, วิปฺปยุตฺเต ตีณิ, อตฺถิยา สตฺต, นตฺถิยา ตีณิ, วิคเต ตีณิ, อวิคเต สตฺต.}

\csromanheader{2}{2. ปจฺจนียุทฺธาร}
\proseitem{340}{นิยโต ธมฺโม นิยตสฺส ธมฺมสฺส สหชาตปจฺจเยน ปจฺจโย.\csromanspacer{} อุปนิสฺสยปจฺจเยน ปจฺจโย.\csromanspacer{}\csromanspacer{}นิยโต ธมฺโม อนิยตสฺส ธมฺมสฺส}

\vspace{\tipitakaparskip}
\csromancenter{นิยตทุก อารมฺมณปจฺจเยน ปจฺจโย.\csromanspacer{} สหชาตปจฺจเยน ปจฺจโย.\csromanspacer{} อุปนิสฺสยปจฺจเยน ปจฺจโย.\csromanspacer{} ปจฺฉาชาตปจฺจเยน ปจฺจโย.\csromanspacer{} กมฺมปจฺจเยน ปจฺจโย.\csromanspacer{}\csromanspacer{}นิยโต ธมฺโม นิยตสฺส จ อนิยตสฺส จ ธมฺมสฺส สหชาตปจฺจเยน ปจฺจโย. (3)}
\proseitem{341}{\csromanfolio{185}อนิยโต ธมฺโม อนิยตสฺส ธมฺมสฺส อารมฺมณปจฺจเยน ปจฺจโย.\csromanspacer{} สหชาตปจฺจเยน ปจฺจโย.\csromanspacer{} อุปนิสฺสยปจฺจเยน ปจฺจโย.\csromanspacer{} ปุเรชาตปจฺจเยน ปจฺจโย.\csromanspacer{} ปจฺฉาชาตปจฺจเยน ปจฺจโย.\csromanspacer{} กมฺมปจฺจเยน ปจฺจโย.\csromanspacer{} อาหารปจฺจเยน ปจฺจโย.\csromanspacer{} อินฺทฺริยปจฺจเยน ปจฺจโย.\csromanspacer{}\csromanspacer{}อนิยโต ธมฺโม นิยตสฺส ธมฺมสฺส อารมฺมณปจฺจเยน ปจฺจโย.\csromanspacer{} อุปนิสฺสยปจฺจเยน ปจฺจโย.\csromanspacer{} ปุเรชาตปจฺจเยน ปจฺจโย. (2)}

\prose{นิยโต จ อนิยโต จ ธมฺมา นิยตสฺส ธมฺมสฺส สหชาตํ, ปุเรชาตํ.\csromanspacer{}\csromanspacer{}นิยโต จ อนิยโต จ ธมฺมา อนิยตสฺส ธมฺมสฺส สหชาตํ, ปจฺฉาชาตํ, อาหารํ, อินฺทฺริยํ. (2)}

\csromanheader{2}{2. ปจฺจยปจฺจนีย 2. สงฺขฺยาวาร}
\proseitem{342}{นเหตุยา สตฺต, น-อารมฺมเณ สตฺต, น-อธิปติยา สตฺต, น-อนนฺตเร นสมนนฺตเร สตฺต, นสหชาเต ปญฺจ, น-อญฺญมญฺเญ ปญฺจ, นนิสฺสเย ปญฺจ, น-อุปนิสฺสเย สตฺต, นปุเรชาเต ฉ, นปจฺฉาชาเต สตฺต ฯเปฯ นสมฺปยุตฺเต ปญฺจ, นวิปฺปยุตฺเต จตฺตาริ, โน-อตฺถิยา จตฺตาริ, โนนตฺถิยา สตฺต, โนวิคเต สตฺต, โน-อวิคเต จตฺตาริ.}

\csromanheader{2}{3. ปจฺจยานุโลมปจฺจนีย}
\proseitem{343}{เหตุปจฺจยา น-อารมฺมเณ จตฺตาริ, น-อธิปติยา จตฺตาริ, น-อนนฺตเร นสมนนฺตเร จตฺตาริ, น-อญฺญมญฺเญ เทฺว, น-อุปนิสฺสเย จตฺตาริ ฯเปฯ นสมฺปยุตฺเต เทฺว, นวิปฺปยุตฺเต เทฺว, โนนตฺถิยา จตฺตาริ, โนวิคเต จตฺตาริ.}

\csromanheader{2}{4. ปจฺจยปจฺจนียานุโลม}
\proseitem{344}{นเหตุปจฺจยา อารมฺมเณ ตีณิ, อธิปติยา ปญฺจ, (อนุโลมมาติกา วิตฺถาเรตพฺพา.) ฯเปฯ อวิคเต สตฺต.}

\nitthitam{นิยตทุกํ นิฏฺฐิตํ.}
\csromanmatikaanchor{mtk.28}
\csromantocmarkref{section}{99. ส-อุตฺตรทุก}{mtk.28}
\bookmark[dest={mtk.28},level=1]{99. ส-อุตฺตรทุก}
\csromanheader{1}{99. ส-อุตฺตรทุก 1. ปฏิจฺจวาร}
\proseitem{345}{\csromanfolio{186}ส-อุตฺตรํ ธมฺมํ ปฏิจฺจ ส-อุตฺตโร ธมฺโม อุปฺปชฺชติ เหตุปจฺจยา.\csromanspacer{} ส-อุตฺตรํ เอกํ ขนฺธํ ปฏิจฺจ ตโย ขนฺธา จิตฺตสมุฏฺฐานญฺจ รูปํ ฯเปฯ เทฺว ขนฺเธ ฯเปฯ.\csromanspacer{} ปฏิสนฺธิกฺขเณ ฯเปฯ. (ยาว อชฺฌตฺติกา มหาภูตา.\csromanspacer{} ยถา จูฬนฺตรทุเก โลกิยทุกสทิสํ, นินฺนานากรณํ.)}

\nitthitam{ส-อุตฺตรทุกํ นิฏฺฐิตํ.}
\csromanmatikaanchor{mtk.29}
\csromantocmarkref{subsection}{100. สรณทุก}{mtk.29}
\bookmark[dest={mtk.29},level=2]{100. สรณทุก}
\csromanheader{2}{100. สรณทุก 1. ปฏิจฺจวาร}
\csromanheader{2}{1. ปจฺจยานุโลมาทิ}
\proseitem{346}{สรณํ ธมฺมํ ปฏิจฺจ สรโณ ธมฺโม อุปฺปชฺชติ เหตุปจฺจยา.\csromanspacer{} สรณํ เอกํ ขนฺธํ ปฏิจฺจ ตโย ขนฺธา ฯเปฯ เทฺว ขนฺเธ ฯเปฯ.}

\prose{(ปญฺจ ปญฺหา อรูปาวจรทุกสทิสา.\csromanspacer{}\csromanspacer{}อนุโลมปฏิจฺจสทิสา.)}

\vspace{\tipitakaparskip}
\csromancenter{\textbf{เหตุ}ยา ปญฺจ, อารมฺมเณ เทฺว, อธิปติยา ปญฺจ ฯเปฯ อวิคเต ปญฺจ.\csromanspacer{} อนุโลมํ.}
\prose{สรณํ ธมฺมํ ปฏิจฺจ สรโณ ธมฺโม อุปฺปชฺชติ นเหตุปจฺจยา.\csromanspacer{} วิจิกิจฺฉาสหคเต อุทฺธจฺจสหคเต ขนฺเธ ปฏิจฺจ วิจิกิจฺฉาสหคโต อุทฺธจฺจสหคโต โมโห. (1)}

\prose{อรณํ ธมฺมํ ปฏิจฺจ อรโณ ธมฺโม อุปฺปชฺชติ นเหตุปจฺจยา.\csromanspacer{} อเหตุกํ อรณํ เอกํ ขนฺธํ ปฏิจฺจ ตโย ขนฺธา จิตฺตสมุฏฺฐานญฺจ รูปํ ฯเปฯ เทฺว ขนฺเธ ฯเปฯ.\csromanspacer{} อเหตุกปฏิสนฺธิกฺขเณ ฯเปฯ. (ยาว อสญฺญสตฺตา.) (1)}

\prose{นเหตุยา เทฺว, น-อารมฺมเณ ตีณิ, น-อธิปติยา ปญฺจ, น-อนนฺตเร ตีณิ, น-อุปนิสฺสเย ตีณิ, นปุเรชาเต จตฺตาริ, นปจฺฉาชาเต ปญฺจ, น-อาเสวเน ปญฺจ, นกมฺเม เทฺว, นวิปาเก ปญฺจ, น-อาหาเร เอกํ,}

\vspace{\tipitakaparskip}
\csromancenter{สรณทุก น-อินฺทฺริเย เอกํ, นฌาเน เอกํ, นมคฺเค เอกํ, นสมฺปยุตฺเต ตีณิ, นวิปฺปยุตฺเต เทฺว, โนนตฺถิยา ตีณิ, โนวิคเต ตีณิ.}
\csromancenter{ปจฺจนียํ.}
\csromancenter{(เอวํ อิตเร เทฺว คณนาปิ สหชาตวาโรปิ กาตพฺโพ.)}
\csromanheader{2}{100. สรณทุก 3. ปจฺจยวาร}
\csromanheader{2}{\textbf{ปจฺจยจตุกฺก}}
\proseitem{347}{\csromanfolio{187}สรณํ ธมฺมํ ปจฺจยา สรโณ ธมฺโม อุปฺปชฺชติ เหตุปจฺจยา.\csromanspacer{} สรณํ เอกํ ขนฺธํ ปจฺจยา ตโย ขนฺธา ฯเปฯ เทฺว ขนฺเธ ฯเปฯ. (ยถา อรูปาวจรทุกสฺส ปจฺจยวาโรปิ เอวํ กาตพฺโพ.)}

\prose{เหตุยา นว, อารมฺมเณ จตฺตาริ, อธิปติยา นว ฯเปฯ อวิคเต นว.\csromanspacer{} อนุโลมํ.}

\proseitem{348}{สรณํ ธมฺมํ ปจฺจยา สรโณ ธมฺโม อุปฺปชฺชติ นเหตุปจฺจยา.\csromanspacer{} วิจิกิจฺฉาสหคเต อุทฺธจฺจสหคเต ขนฺเธ ปจฺจยา วิจิกิจฺฉาสหคโต อุทฺธจฺจสหคโต โมโห. (1)}

\prose{อรณํ ธมฺมํ ปจฺจยา อรโณ ธมฺโม อุปฺปชฺชติ นเหตุปจฺจยา.\csromanspacer{} อเหตุกํ อรณํ เอกํ ขนฺธํ ปจฺจยา ตโย ขนฺธา จิตฺตสมุฏฺฐานญฺจ รูปํ ฯเปฯ เทฺว ขนฺเธ ฯเปฯ.\csromanspacer{} อเหตุกปฏิสนฺธิกฺขเณ ฯเปฯ. (ยาว อสญฺญสตฺตา.) จกฺขายตนํ ปจฺจยา จกฺขุวิญฺญาณํ ฯเปฯ.\csromanspacer{} กายายตนํ ปจฺจยา กายวิญฺญาณํ.\csromanspacer{} วตฺถุํ ปจฺจยา อเหตุกา ขนฺธา.\csromanspacer{}\csromanspacer{}อรณํ ธมฺมํ ปจฺจยา สรโณ ธมฺโม อุปฺปชฺชติ นเหตุปจฺจยา.\csromanspacer{} วตฺถุํ ปจฺจยา วิจิกิจฺฉาสหคโต อุทฺธจฺจสหคโต โมโห. (2)}

\prose{สรณญฺจ อรณญฺจ ธมฺมํ ปจฺจยา สรโณ ธมฺโม อุปฺปชฺชติ นเหตุปจฺจยา.\csromanspacer{} วิจิกิจฺฉาสหคเต อุทฺธจฺจสหคเต ขนฺเธ จ วตฺถุญฺจ ปจฺจยา วิจิกิจฺฉาสหคโต อุทฺธจฺจสหคโต โมโห. (1)}

\prose{\csromanfolio{188}นเหตุยา จตฺตาริ, น-อารมฺมเณ ตีณิ, น-อธิปติยา นว, น-อนนฺตเร ตีณิ, นสมนนฺตเร ตีณิ, น-อญฺญมญฺเญ ตีณิ, น-อุปนิสฺสเย ตีณิ, นปุเรชาเต จตฺตาริ, นปจฺฉาชาเต นว, น-อาเสวเน นว, นกมฺเม จตฺตาริ, นวิปาเก นว, น-อาหาเร เอกํ, น-อินฺทฺริเย เอกํ, นฌาเน เอกํ, นมคฺเค เอกํ, นสมฺปยุตฺเต ตีณิ, นวิปฺปยุตฺเต เทฺว, โนนตฺถิยา ตีณิ, โนวิคเต ตีณิ.\csromanspacer{} ปจฺจนียํ.}

\prose{(เอวํ อิตเร เทฺว คณนาปิ นิสฺสยวาโรปิ กาตพฺโพ, สํสฏฺฐวาเรปิ เทฺว ปญฺหา กาตพฺพา.\csromanspacer{} สพฺพตฺถ.)}

\prose{\textbf{เหตุ}ยา เทฺว, อารมฺมเณ เทฺว, (สพฺพตฺถ เทฺว.) วิปาเก เอกํ, อวิคเต เทฺว. (อนุโลมํ.)}

\prose{นเหตุยา เทฺว, น-อธิปติยา เทฺว, นปุเรชาเต เทฺว, นปจฺฉาชาเต เทฺว, น-อาเสวเน เทฺว, นกมฺเม เทฺว, นวิปาเก เทฺว, นฌาเน เอกํ, นมคฺเค เอกํ, นวิปฺปยุตฺเต เทฺว. (ปจฺจนียํ.)}

\vspace{\tipitakaparskip}
\csromancenter{(เอวํ อิตเร เทฺว คณนาปิ สมฺปยุตฺตวาโรปิ กาตพฺโพ.)}
\csromanheader{2}{100. สรณทุก 7. ปญฺหาวาร}
\proseitem{349}{สรโณ ธมฺโม สรณสฺส ธมฺมสฺส เหตุปจฺจเยน ปจฺจโย. (อรูปทุกสทิสํ, จตฺตาริ.)}

\csromanheader{2}{\textbf{อารมฺมณ}}
\proseitem{350}{สรโณ ธมฺโม สรณสฺส ธมฺมสฺส อารมฺมณปจฺจเยน ปจฺจโย.\csromanspacer{} ราคํ อสฺสาเทติ อภินนฺทติ, ตํ อรพฺภ ราโค อุปฺปชฺชติ.\csromanspacer{} ทิฏฺฐิ ฯเปฯ วิจิกิจฺฉา ฯเปฯ อุทฺธจฺจํ ฯเปฯ โทมนสฺสํ อุปฺปชฺชติ.\csromanspacer{} ทิฏฺฐิํ อสฺสาเทติ อภินนฺทติ, ตํ อารพฺภ ราโค อุปฺปชฺชติ ฯเปฯ โทมนสฺสํ อุปฺปชฺชติ.\csromanspacer{} วิจิกิจฺฉํ}

\vspace{\tipitakaparskip}
\csromancenter{สรณทุก อารพฺภ ฯเปฯ อุทฺธจฺจํ อารพฺภ ฯเปฯ.\csromanspacer{} โทมนสฺสํ อารพฺภ โทมนสฺสํ อุปฺปชฺชติ.\csromanspacer{} ทิฏฺฐิ ฯเปฯ วิจิกิจฺฉา ฯเปฯ อุทฺธจฺจํ อุปฺปชฺชติ.\csromanspacer{}\csromanspacer{}สรโณ ธมฺโม อรณสฺส ธมฺมสฺส อรมฺมณปจฺจเยน ปจฺจโย.\csromanspacer{} อริยา ปหีเน กิเลเส ฯเปฯ.\csromanspacer{} วิกฺขมฺภิเต กิเลเส ฯเปฯ.\csromanspacer{} ปุพฺเพ สมุทาจิณฺเณ ฯเปฯ.\csromanspacer{} สรเณ ขนฺเธ อนิจฺจโต ฯเปฯ วิปสฺสติ.\csromanspacer{} เจโตปริยญาเณน สรณจิตฺตสมงฺคิสฺส จิตฺตํ ชานาติ.\csromanspacer{} สรณา ขนฺธา เจโตปริยญาณสฺส, ปุพฺเพนิวาสานุสฺสติญาณสฺส, ยถากมฺมูปคญาณสฺส, อนาคตํสญาณสฺส อาวชฺชนาย อารมฺมณปจฺจเยน ปจฺจโย. (2)}
\proseitem{351}{\csromanfolio{189}อรโณ ธมฺโม อรณสฺส ธมฺมสฺส อารมฺมณปจฺจเยน ปจฺจโย.\csromanspacer{} ทานํ ฯเปฯ.\csromanspacer{} สีลํ ฯเปฯ.\csromanspacer{} อุโปสถกมฺมํ กตฺวา ตํ ปจฺจเวกฺขติ.\csromanspacer{} ปุพฺเพ ฯเปฯ.\csromanspacer{} ฌานา ฯเปฯ.\csromanspacer{} อริยา มคฺคา วุฏฺฐหิตฺวา มคฺคํ ปจฺจเวกฺขนฺติ.\csromanspacer{} ผลํ ฯเปฯ.\csromanspacer{} นิพฺพานํ ปจฺจเวกฺขนฺติ.\csromanspacer{} นิพฺพานํ โคตฺรภุสฺส, โวทานสฺส, มคฺคสฺส, ผลสฺส, อาวชฺชนาย อารมฺมณปจฺจเยน ปจฺจโย.\csromanspacer{} จกฺขุํ ฯเปฯ.\csromanspacer{} วตฺถุํ อรเณ ขนฺเธ อนิจฺจโต ฯเปฯ วิปสฺสติ.\csromanspacer{} ทิพฺเพน จกฺขุนา รูปํ ปสฺสติ ฯเปฯ.\csromanspacer{} อนาคตํสญาณสฺส, อาวชฺชนาย อารมฺมณปจฺจเยน ปจฺจโย.\csromanspacer{}\csromanspacer{}อรโณ ธมฺโม สรณสฺส ธมฺมสฺส อารมฺมณปจฺจเยน ปจฺจโย.\csromanspacer{} ทานํ ฯเปฯ.\csromanspacer{} สีลํ ฯเปฯ.\csromanspacer{} อุโปสถกมฺมํ กตฺวา ตํ อสฺสาเทติ อภินนฺทติ, ตํ อารพฺภ ราโค อุปฺปชฺชติ ฯเปฯ.\csromanspacer{} ปุพฺเพ ฯเปฯ.\csromanspacer{} ฌานา ฯเปฯ.\csromanspacer{} จกฺขุํ ฯเปฯ.\csromanspacer{} วตฺถุํ.\csromanspacer{} อรเณ ขนฺเธ อสฺสาเทติ อภินนฺทติ, ตํ อารพฺภ ราโค อุปฺปชฺชติ ฯเปฯ โทมนสฺสํ อุปฺปชฺชติ. (2)}

\csromanheader{2}{\textbf{อธิปติ}}
\proseitem{352}{สรโณ ธมฺโม สรณสฺส ธมฺมสฺส อธิปติปจฺจเยน ปจฺจโย.\csromanspacer{} อารมฺมณาธิปติ, สหชาตาธิปติ.\csromanspacer{}\csromanspacer{}\textbf{อารมฺมณาธิปติ\csromandash{}}ราคํ ครุํ กตฺวา อสฺสาเทติ อภินนฺทติ, ตํ ครุํ กตฺวา ราโค อุปฺปชฺชติ.\csromanspacer{} ทิฏฺฐิ อุปฺปชฺชติ.\csromanspacer{} ทิฏฺฐิํ ครุํ กตฺวา อสฺสาเทติ อภินนฺทติ.\csromanspacer{} \textbf{สหชาตาธิปติ\csromandash{}}สรณาธิปติ สมฺปยุตฺตกานํ ขนฺธานํ อธิปติปจฺจเยน ปจฺจโย.\csromanspacer{}\csromanspacer{}สรโณ ธมฺโม อรณสฺส ธมฺมสฺส อธิปติปจฺจเยน ปจฺจโย.\csromanspacer{} \textbf{สหชาตาธิปติ\csromandash{}}สรณาธิปติ จิตฺตสมุฏฺฐานานํ รูปานํ อธิปติปจฺจเยน ปจฺจโย.\csromanspacer{}\csromanspacer{}สรโณ ธมฺโม สรณสฺส จ อรณสฺส จ ธมฺมสฺส อธิปติปจฺจเยน ปจฺจโย.\csromanspacer{} \textbf{สหชาตาธิปติ\csromandash{}}สรณาธิปติ สมฺปยุตฺตกานํ ขนฺธานํ จิตฺตสมุฏฺฐานานญฺจ รูปานํ อธิปติปจฺจเยน ปจฺจโย. (3)}

\proseitem{353}{\csromanfolio{190}อรโณ ธมฺโม อรณสฺส ธมฺมสฺส อธิปติปจฺจเยน ปจฺจโย.\csromanspacer{} อารมฺมณธิปติ, สหชาตาธิปติ.\csromanspacer{}\csromanspacer{}\textbf{อารมฺมณาธิปติ\csromandash{}}ทานํ ฯเปฯ สีลํ ฯเปฯ อุโปสถกมฺมํ กตฺวา ตํ ครุํ กตฺวา ปจฺจเวกฺขติ.\csromanspacer{} ปุพฺเพ ฯเปฯ.\csromanspacer{} ฌานา ฯเปฯ.\csromanspacer{} อริยา มคฺคา วุฏฺฐหิตฺวา มคฺคํ ครุํ กตฺวา ปจฺจเวกฺขนฺติ.\csromanspacer{} ผลํ ฯเปฯ.\csromanspacer{} นิพฺพานํ ครุํ กตฺวา ปจฺจเวกฺขนฺติ.\csromanspacer{} นิพฺพานํ โคตฺรภุสฺส, โวทานสฺส, มคฺคสฺส, ผลสฺส อธิปติปจฺจเยน ปจฺจโย.\csromanspacer{} \textbf{สหชาตาธิปติ\csromandash{}}อรณาธิปติ สมฺปยุตฺตกานํ ขนฺธานํ จิตฺตสมุฏฺฐานานญฺจ รูปานํ อธิปติปจฺจเยน ปจฺจโย.\csromanspacer{}\csromanspacer{}อรโณ ธมฺโม สรณสฺส ธมฺมสฺส อธิปติปจฺจเยน ปจฺจโย.\csromanspacer{} \textbf{อารมฺมณาธิปติ\csromandash{}}ทานํ ฯเปฯ สีลํ ฯเปฯ อุโปสถกมฺมํ กตฺวา ตํ ครุํ กตฺวา อสฺสาเทติ อภินนฺทติ, ตํ ครุํ กตฺวา ราโค อุปฺปชฺชติ.\csromanspacer{} ทิฏฺฐิ อุปฺปชฺชติ.\csromanspacer{} ปุพฺเพ สุจิณฺณานิ ฯเปฯ.\csromanspacer{} ฌานา ฯเปฯ.\csromanspacer{} จกฺขุํ ฯเปฯ.\csromanspacer{} วตฺถุํ อรเณ ขนฺเธ ครุํ กตฺวา อสฺสาเทติ อภินนฺทติ, ตํ ครุํ กตฺวา ราโค อุปฺปชฺชติ.\csromanspacer{} ทิฏฺฐิ อุปฺปชฺชติ. (2)}

\csromanheader{2}{\textbf{อนนฺตราทิ}}
\proseitem{354}{สรโณ ธมฺโม สรณสฺส ธมฺมสฺส อนนฺตรปจฺจเยน ปจฺจโย.\csromanspacer{} ปุริมา ปุริมา สรณา ขนฺธา ปจฺฉิมานํ ปจฺฉิมานํ สรณานํ ขนฺธานํ อนนฺตรปจฺจเยน ปจฺจโย.\csromanspacer{}\csromanspacer{}สรโณ ธมฺโม อรณสฺส ธมฺมสฺส อนนฺตรปจฺจเยน ปจฺจโย.\csromanspacer{} สรณา ขนฺธา วุฏฺฐานสฺส อนนฺตรปจฺจเยน ปจฺจโย. (2)}

\proseitem{355}{อรโณ ธมฺโม อรณสฺส ธมฺมสฺส อนนฺตรปจฺจเยน ปจฺจโย.\csromanspacer{} ปุริมา ปุริมา อรณา ขนฺธา ปจฺฉิมานํ ปจฺฉิมานํ อรณานํ ขนฺธานํ อนนฺตรปจฺจเยน ปจฺจโย.\csromanspacer{} อนุโลมํ โคตฺรภุสฺส ฯเปฯ ผลสมาปตฺติยา อนนฺตรปจฺจเยน ปจฺจโย.\csromanspacer{}\csromanspacer{}อรโณ ธมฺโม สรณสฺส ธมฺมสฺส อนนฺตรปจฺจเยน ปจฺจโย.\csromanspacer{} อาวชฺชนา สรณานํ ขนฺธานํ อนนฺตรปจฺจเยน ปจฺจโย. (2)}

\prose{สมนนฺตรปจฺจเยน ปจฺจโย.\csromanspacer{} สหชาตปจฺจเยน ปจฺจโย.\csromanspacer{} ปญฺจ.\csromanspacer{} อญฺญมญฺญปจฺจเยน ปจฺจโย.\csromanspacer{} เทฺว.\csromanspacer{} นิสฺสยปจฺจเยน ปจฺจโย.\csromanspacer{} สตฺต.}

\csromanheader{2}{\textbf{อุปนิสฺสย}}
\proseitem{356}{สรโณ ธมฺโม สรณสฺส ธมฺมสฺส อุปนิสฺสยปจฺจเยน ปจฺจโย.\csromanspacer{} \textbf{อารมฺมณูปนิสฺสโย, อนนฺตรูปนิสฺสโย, ปกตูปนิสฺสโย ฯเปฯ.}}

\vspace{\tipitakaparskip}
\csromancenter{สรณทุก ปกตูปนิสฺสโย ราคํ อุปนิสฺสาย ปาณํ หนติ ฯเปฯ สํฆํ ภินฺทติ.\csromanspacer{} โทสํ ฯเปฯ.\csromanspacer{} ปตฺถนํ อุปนิสฺสาย ปาณํ หนติ ฯเปฯ สํฆํ ภินฺทติ.\csromanspacer{} ราโค ฯเปฯ.\csromanspacer{} ปตฺถนา ราคสฺส ฯเปฯ ปตฺถนาย อุปนิสฺสยปจฺจเยน ปจฺจโย.\csromanspacer{}\csromanspacer{}สรโณ ธมฺโม อรณสฺส ธมฺมสฺส อุปนิสฺสยปจฺจเยน ปจฺจโย.\csromanspacer{} \textbf{อนนฺตรูปนิสฺสโย, ปกตูปนิสฺสโย ฯเปฯ.}\csromanspacer{} อกตูปนิสฺสโย\csromandash{}ราคํ อุปนิสฺสาย ทานํ เทติ.\csromanspacer{} สีลํ ฯเปฯ อุโปสถกมฺมํ กโรติ.\csromanspacer{} ฌานํ อุปฺปาเทติ.\csromanspacer{} วิปสฺสนํ ฯเปฯ.\csromanspacer{} มคฺคํ ฯเปฯ.\csromanspacer{} อภิญฺญํ ฯเปฯ.\csromanspacer{} สมาปตฺติํ อุปฺปาเทติ.\csromanspacer{} โทสํ ฯเปฯ.\csromanspacer{} ปตฺถนํ อุปนิสฺสาย ทานํ เทติ ฯเปฯ.\csromanspacer{} สมาปตฺติํ อุปฺปาเทติ.\csromanspacer{} ราโค ฯเปฯ.\csromanspacer{} ปตฺถนา สทฺธาย ฯเปฯ ปญฺญาย.\csromanspacer{} กายิกสฺส สุขสฺส.\csromanspacer{} กายิกสฺส ทุกฺขสฺส.\csromanspacer{} มคฺคสฺส.\csromanspacer{} ผลสมาปตฺติยา อุปนิสฺสยปจฺจเยน ปจฺจโย. (2)}
\proseitem{357}{\csromanfolio{191}อรโณ ธมฺโม อรณสฺส ธมฺมสฺส อุปนิสฺสยปจฺจเยน ปจฺจโย.\csromanspacer{} \textbf{อารมฺมณูปนิสฺสโย, อนนฺตรูปนิสฺสโย, ปกตูปนิสฺสโย ฯเปฯ.}\csromanspacer{} ปกตูปนิสฺสโย\csromandash{}สทฺธํ อุปนิสฺสาย ทานํ เทติ ฯเปฯ.\csromanspacer{} สมาปตฺติํ อุปฺปาเทติ.\csromanspacer{} สีลํ ฯเปฯ.\csromanspacer{} ปญฺญํ.\csromanspacer{} กายิกํ สุขํ.\csromanspacer{} กายิกํ ทุกฺขํ.\csromanspacer{} อุตุํ.\csromanspacer{} โภชนํ.\csromanspacer{} เสนาสนํ อุปนิสฺสาย ทานํ เทติ ฯเปฯ.\csromanspacer{} สมาปตฺติํ อุปฺปาเทติ.\csromanspacer{} สทฺธา ฯเปฯ.\csromanspacer{} เสนาสนํ สทฺธาย ฯเปฯ ปญฺญาย.\csromanspacer{} กายิกสฺส สุขสฺส.\csromanspacer{} กายิกสฺส ทุกฺขสฺส.\csromanspacer{} มคฺคสฺส.\csromanspacer{} ผลสมาปตฺติยา อุปนิสฺสยปจฺจเยน ปจฺจโย. (1)}

\prose{อรโณ ธมฺโม สรณสฺส ธมฺมสฺส อุปนิสฺสยปจฺจเยน ปจฺจโย.\csromanspacer{} \textbf{อารมฺมณูปนิสฺสโย, อนนฺตรูปนิสฺสโย, ปกตูปนิสฺสโย ฯเปฯ.}\csromanspacer{} ปกตูปนิสฺสโย\csromandash{}สทฺธํ อุปนิสฺสาย มานํ ชปฺเปติ.\csromanspacer{} ทิฏฺฐิํ คณฺหาติ.\csromanspacer{} สีลํ ฯเปฯ.\csromanspacer{} เสนาสนํ อุปนิสฺสาย ปาณํ หนติ ฯเปฯ สํฆํ ภินฺทติ.\csromanspacer{} สทฺธา ฯเปฯ.\csromanspacer{} เสนาสนํ ราคสฺส ฯเปฯ ปตฺถนาย อุปนิสฺสยปจฺจเยน ปจฺจโย. (2)}

\csromanheader{2}{\textbf{ปุเรชาตาทิ}}
\proseitem{358}{อรโณ ธมฺโม อรณสฺส ธมฺมสฺส ปุเรชาตปจฺจเยน ปจฺจโย.\csromanspacer{} อารมฺมณปุเรชาตํ, วตฺถุปุเรชาตํ.\csromanspacer{}\csromanspacer{}\textbf{อารมฺมณปุเรชาตํ\csromandash{}}จกฺขุํ ฯเปฯ.\csromanspacer{} วตฺถุํ อนิจฺจโต ฯเปฯ วิปสฺสติ ฯเปฯ.\csromanspacer{} ทิพฺเพน จกฺขุนา รูปํ ปสฺสติ.\csromanspacer{} ทิพฺพาย โสตธาตุยา สทฺทํ สุณาติ.\csromanspacer{} รูปายตนํ จกฺขุวิญฺญาณสฺส ฯเปฯ.\csromanspacer{} โผฏฺฐพฺพายตนํ กายวิญฺญาณสฺส ฯเปฯ.\csromanspacer{} \textbf{วตฺถุปุเรชาตํ\csromandash{}}จกฺขายตนํ \csromanfolio{192}จกฺขุวิญฺญาณสฺส ฯเปฯ.\csromanspacer{} กายายตนํ กายวิญฺญาณสฺส ฯเปฯ.\csromanspacer{} วตฺถุ อรณานํ ขนฺธานํ ปุเรชาตปจฺจเยน ปจฺจโย.\csromanspacer{}\csromanspacer{}อรโณ ธมฺโม สรณสฺส ธมฺมสฺส ปุเรชาตปจฺจเยน ปจฺจโย.\csromanspacer{} อารมฺมณปุเรชาตํ, วตฺถุปุเรชาตํ.\csromanspacer{}\csromanspacer{}\textbf{อารมฺมณปุเรชาตํ\csromandash{}}จกฺขุํ ฯเปฯ.\csromanspacer{} วตฺถุํ อสฺสาเทติ อภินนฺทติ, ตํ อารพฺภ ราโค อุปฺปชฺชติ ฯเปฯ โทมนสฺสํ อุปฺปชฺชติ.\csromanspacer{} \textbf{วตฺถุปุเรชาตํ\csromandash{}}วตฺถุ สรณานํ ขนฺธานํ ปุเรชาตปจฺจเยน ปจฺจโย. (2)}

\prose{ปจฺฉาชาตปจฺจเยน ปจฺจโย.\csromanspacer{} เทฺว.\csromanspacer{} อาเสวนปจฺจเยน ปจฺจโย.\csromanspacer{} เทฺว.}

\csromanheader{2}{\textbf{กมฺม}}
\proseitem{359}{สรโณ ธมฺโม สรณสฺส ธมฺมสฺส กมฺมปจฺจเยน ปจฺจโย.\csromanspacer{} สรณา เจตนา สมฺปยุตฺตกานํ ขนฺธานํ กมฺมปจฺจเยน ปจฺจโย. (มูลํ กาตพฺพํ.) สหชาตา, นานากฺขณิกา.\csromanspacer{}\csromanspacer{}\textbf{สหชาตา} สรณา เจตนา จิตฺตสมุฏฺฐานานํ รูปานํ กมฺมปจฺจเยน ปจฺจโย.\csromanspacer{} \textbf{นานากฺขณิกา} สรณา เจตนา วิปากานํ ขนฺธานํ กฏตฺตา จ รูปานํ กมฺมปจฺจเยน ปจฺจโย. (มูลํ กาตพฺพํ.) สรณา เจตนา สมฺปยุตฺตกานํ ขนฺธานํ จิตฺตสมุฏฺฐานานญฺจ รูปานํ กมฺมปจฺจเยน ปจฺจโย. (3)}

\prose{อรโณ ธมฺโม อรณสฺส ธมฺมสฺส กมฺมปจฺจเยน ปจฺจโย.\csromanspacer{} \textbf{สหชาตา}, นานากฺขณิกา.\csromanspacer{}\csromanspacer{}\textbf{สหชาตา}\csromandash{}(สํขิตฺตํ.) (1)}

\prose{วิปากปจฺจเยน ปจฺจโย.\csromanspacer{} เอกํ.\csromanspacer{} อาหารปจฺจเยน ปจฺจโย.\csromanspacer{} จตฺตาริ.\csromanspacer{} อินฺทฺริยปจฺจเยน ปจฺจโย.\csromanspacer{} จตฺตาริ.\csromanspacer{} ฌานปจฺจเยน ปจฺจโย.\csromanspacer{} จตฺตาริ.\csromanspacer{} มคฺคปจฺจเยน ปจฺจโย.\csromanspacer{} จตฺตาริ.\csromanspacer{} สมฺปยุตฺตปจฺจเยน ปจฺจโย.\csromanspacer{} เทฺว.\csromanspacer{} วิปฺปยุตฺตปจฺจเยน ปจฺจโย.\csromanspacer{} ตีณิ. (อรูปทุกสทิสา.)}

\csromanheader{2}{\textbf{อตฺถิ}}
\proseitem{360}{สรโณ ธมฺโม สรณสฺส ธมฺมสฺส อตฺถิปจฺจเยน ปจฺจโย. (สํขิตฺตํ.).\csromanspacer{} สรโณ ธมฺโม อรณสฺส ธมฺมสฺส อตฺถิปจฺจเยน ปจฺจโย.\csromanspacer{} สหชาตํ, ปจฺฉาชาตํ. (สํขิตฺตํ.).\csromanspacer{} สรโณ ธมฺโม สรณสฺส จ อรณสฺส จ ธมฺมสฺส อตฺถิปจฺจเยน ปจฺจโย. (สํขิตฺตํ.) (3)}

\proseitem{361}{อรโณ ธมฺโม อรณสฺส ธมฺมสฺส อตฺถิปจฺจเยน ปจฺจโย.\csromanspacer{} สหชาตํ, ปุเรชาตํ, ปจฺฉาชาตํ, อาหารํ, อินฺทฺริยํ. (สํขิตฺตํ.).\csromanspacer{} อรโณ ธมฺโม สรณสฺส ธมฺมสฺส อตฺถิปจฺจเยน ปจฺจโย.\csromanspacer{} \textbf{ปุเรชาตํ}}

\vspace{\tipitakaparskip}
\csromancenter{สรณทุก จกฺขุํ ฯเปฯ.\csromanspacer{} วตฺถุํ อสฺสาเทติ อภินนฺทติ, ตํ อารพฺภ ราโค อุปฺปชฺชติ ฯเปฯ โทมนสฺสํ อุปฺปชฺชติ.\csromanspacer{} วตฺถุ สรณานํ ขนฺธานํ อตฺถิปจฺจเยน ปจฺจโย. (2)}
\prose{\csromanfolio{193}สรโณ จ อรโณ จ ธมฺมา สรณสฺส ธมฺมสฺส อตฺถิปจฺจเยน ปจฺจโย.\csromanspacer{} สหชาตํ, ปุเรชาตํ. (สํขิตฺตํ.).\csromanspacer{} สรโณ จ อรโณ จ ธมฺมา อรณสฺส ธมฺมสฺส อตฺถิปจฺจเยน ปจฺจโย.\csromanspacer{} สหชาตํ, ปจฺฉาชาตํ, อาหารํ, อินฺทฺริยํ.\csromanspacer{}\csromanspacer{}\textbf{สหชาตา} สรณา ขนฺธา จ มหาภูตา จ จิตฺตสมุฏฺฐานานํ รูปานํ อตฺถิปจฺจเยน ปจฺจโย.\csromanspacer{} \textbf{ปจฺฉาชาตา} สรณา ขนฺธา จ กพฬีกาโร อาหาโร จ อิมสฺส กายสฺส อตฺถิปจฺจเยน ปจฺจโย.\csromanspacer{} \textbf{ปจฺฉาชาตา} สรณา ขนฺธา จ รูปชีวิตินฺทฺริยญฺจ กฏตฺตารูปานํ อตฺถิปจฺจเยน ปจฺจโย. (2)}

\csromanheader{2}{1. ปจฺจยานุโลม 2. สงฺขฺยาวาร}
\csromanheader{2}{\textbf{สุทฺธ}}
\proseitem{362}{เหตุยา จตฺตาริ, อารมฺมเณ จตฺตาริ, อธิปติยา ปญฺจ, อนนฺตเร จตฺตาริ, สมนนฺตเร จตฺตาริ, สหชาเต ปญฺจ, อญฺญมญฺเญ เทฺว, นิสฺสเย สตฺต, อุปนิสฺสเย จตฺตาริ, ปุเรชาเต เทฺว, ปจฺฉาชาเต เทฺว, อาเสวเน เทฺว, กมฺเม จตฺตาริ, วิปาเก เอกํ, อาหาเร จาตฺตาริ, อินฺทฺริเย จตฺตาริ, ฌาเน จตฺตาริ, มคฺเค จตฺตาริ, สมฺปยุตฺเต เทฺว, วิปฺปยุตฺเต ตีณิ, อตฺถิยา สตฺต, นตฺถิยา จตฺตาริ, วิคเต จตฺตาริ, อวิคเต สตฺต.}

\csromanheader{2}{2. ปจฺจนียุทฺธาร}
\proseitem{363}{สรโณ ธมฺโม สรณสฺส ธมฺมสฺส อารมฺมณปจฺจเยน ปจฺจโย.\csromanspacer{} สหชาตปจฺจเยน ปจฺจโย.\csromanspacer{} อุปนิสฺสยปจฺจเยน ปจฺจโย.\csromanspacer{}\csromanspacer{}สรโณ ธมฺโม อรณสฺส ธมฺมสฺส อารมฺมณปจฺจเยน ปจฺจโย.\csromanspacer{} สหชาตปจฺจเยน ปจฺจโย.\csromanspacer{} อุปนิสฺสยปจฺจเยน ปจฺจโย.\csromanspacer{} ปจฺฉาชาตปจฺจเยน ปจฺจโย.\csromanspacer{} กมฺมปจฺจเยน ปจฺจโย.\csromanspacer{}\csromanspacer{}สรโณ ธมฺโม สรณสฺส จ อรณสฺส จ ธมฺมสฺส สหชาตปจฺจเยน ปจฺจโย. (3)}

\proseitem{364}{อรโณ ธมฺโม อรณสฺส ธมฺมสฺส อารมฺมณปจฺจเยน ปจฺจโย.\csromanspacer{} สหชาตปจฺจเยน ปจฺจโย.\csromanspacer{} อุปนิสฺสยปจฺจเยน ปจฺจโย.\csromanspacer{} ปุเรชาตปจฺจเยน ปจฺจโย.\csromanspacer{} ปจฺฉาชาตปจฺจเยน ปจฺจโย.\csromanspacer{} กมฺมปจฺจเยน ปจฺจโย. \csromanfolio{194}อาหารปจฺจเยน ปจฺจโย.\csromanspacer{} อินฺทฺริยปจฺจเยน ปจฺจโย.\csromanspacer{}\csromanspacer{}อรโณ ธมฺโม สรณสฺส ธมฺมสฺส อารมฺมณปจฺจเยน ปจฺจโย.\csromanspacer{} อุปนิสฺสยปจฺจเยน ปจฺจโย.\csromanspacer{} ปุเรชาตปจฺจเยน ปจฺจโย. (2)}

\prose{สรโณ จ อรโณ จ ธมฺมา สรณสฺส ธมฺมสฺส สหชาตํ, ปุเรชาตํ.\csromanspacer{}\csromanspacer{}สรโณ จ อรโณ จ ธมฺมา อรณสฺส ธมฺมสฺส สหชาตํ, ปจฺฉาชาตํ, อาหารํ, อินฺทฺริยํ. (2)}

\csromanheader{2}{2. ปจฺจยปจฺจนีย 2. สงฺขฺยาวาร}
\csromanheader{2}{\textbf{สุทฺธ}}
\proseitem{365}{นเหตุยา สตฺต, นสมนนฺตเร สตฺต, นสหชาเต ปญฺจ, น-อญฺญมญฺเญ ปญฺจ, นนิสฺสเย ปญฺจ, น-อุปนิสฺสเย สตฺต, นปุเรชาเต ฉ, นปจฺฉาชาเต สตฺต ฯเปฯ นสมฺปยุตฺเต ปญฺจ, นวิปฺปยุตฺเต จตฺตาริ, โน-อตฺถิยา จตฺตาริ, โนนตฺถิยา สตฺต, โนวิคเต สตฺต, โน-อวิคเต จตฺตาริ.}

\csromanheader{2}{3. ปจฺจยานุโลมปจฺจนีย}
\proseitem{366}{เหตุปจฺจยา น-อารมฺมเณ จตฺตาริ, น-อธิปติยา จตฺตาริ ฯเปฯ นสมนนฺตเร จตฺตาริ, น-อญฺญมญฺเญ เทฺว, น-อุปนิสฺสเย จตฺตาริ ฯเปฯ นสมฺปยุตฺเต เทฺว, นวิปฺปยุตฺเต เทฺว, โนนตฺถิยา จตฺตาริ, โนวิคเต จตฺตาริ.}

\csromanheader{2}{4. ปจฺจยปจฺจนียานุโลม}
\proseitem{367}{นเหตุปจฺจยา อารมฺมเณ จตฺตาริ, อธิปติยา ปญฺจ, อนนฺตเร จตฺตาริ, (อนุโลมมาติกา คเหตพฺพา.) ฯเปฯ อวิคเต สตฺต.}

\nitthitam{สรณทุกํ นิฏฺฐิตํ.}
\nitthitam{ปิฏฺฐิทุกํ นิฏฺฐิตํ.}
\nitthitam{\textbf{ธมฺมานุโลเม ทุกปฏฺฐานํ นิฏฺฐิตํ.}}
\pitaka{\textbf{อภิธมฺมปิฏก}}
\csromanheader{2}{\textbf{ธมฺมานุโลม}}
\csromanmatikaanchor{mtk.30}
\csromantocmarknopage{chapter}{ทุกติกปฏฺฐานปาฬิ}{mtk.30}
\bookmark[dest={mtk.30},level=0]{ทุกติกปฏฺฐานปาฬิ}
\gambhira{\textbf{ทุกติกปฏฺฐานปาฬิ}}
\namakkaram{นโม ตสฺส ภควโต อรหโต สมฺมาสมฺพุทฺธสฺส.}
\csromanmatikaanchor{mtk.31}
\csromanmatikaanchor{mtk.32}
\csromanmatikaanchor{mtk.92}
\csromantocmarknopage{chapter}{1. เหตุทุก}{mtk.31}
\bookmark[dest={mtk.31},level=0]{1. เหตุทุก}
\csromantocmarkref{section}{1. กุสลตฺติก}{mtk.32}
\bookmark[dest={mtk.32},level=1]{1. กุสลตฺติก}
\csromanheader{1}{1. เหตุทุก 1. กุสลตฺติก}
\csromanheader{2}{1. กุสลปท 1. ปฏิจฺจวาร}
\csromanheader{2}{\textbf{ปจฺจยจตุกฺก-เหตุ}}
\proseitem{1}{\csromanfolio{195}เหตุํ กุสลํ ธมฺมํ ปฏิจฺจ เหตุ กุสโล ธมฺโม อุปฺปชฺชติ เหตุปจฺจยา.\csromanspacer{}\csromanspacer{}เหตุํ กุสลํ ธมฺมํ ปฏิจฺจ นเหตุ กุสโล ธมฺโม อุปฺปชฺชติ เหตุปจฺจยา.\csromanspacer{}\csromanspacer{}เหตุํ กุสลํ ธมฺมํ ปฏิจฺจ เหตุ กุสโล จ นเหตุ กุสโล จ ธมฺมา อุปฺปชฺชนฺติ เหตุปจฺจยา. (3)}

\proseitem{2}{นเหตุํ กุสลํ ธมฺมํ ปฏิจฺจ นเหตุ กุสโล ธมฺโม อุปฺปชฺชติ เหตุปจฺจยา.\csromanspacer{}\csromanspacer{}นเหตุํ กุสลํ ธมฺมํ ปฏิจฺจ เหตุ กุสโล ธมฺโม อุปฺปชฺชติ เหตุปจฺจยา.\csromanspacer{}\csromanspacer{}นเหตุํ กุสลํ ธมฺมํ ปฏิจฺจ เหตุ กุสโล จ นเหตุ กุสโล จ ธมฺมา อุปฺปชฺชนฺติ เหตุปจฺจยา. (3)}

\proseitem{3}{เหตุํ กุสลญฺจ นเหตุํ กุสลญฺจ ธมฺมํ ปฏิจฺจ เหตุ กุสโล ธมฺโม อุปฺปชฺชติ เหตุปจฺจยา.\csromanspacer{}\csromanspacer{}เหตุํ กุสลญฺจ นเหตุํ กุสลญฺจ ธมฺมํ ปฏิจฺจ นเหตุ กุสโล ธมฺโม อุปฺปชฺชติ เหตุปจฺจยา.\csromanspacer{}\csromanspacer{}เหตุํ กุสลญฺจ นเหตุํ กุสลญฺจ ธมฺมํ ปฏิจฺจ เหตุ กุสโล จ นเหตุ กุสโล จ ธมฺมา อุปฺปชฺชนฺติ เหตุปจฺจยา. (3)}

\csromanheader{2}{ธมฺมานุโลม ทุกติกปฏฺฐานปาฬิ \textbf{อารมฺมณ}}
\proseitem{4}{\csromanfolio{196}เหตุํ กุสลํ ธมฺมํ ปฏิจฺจ เหตุ กุสโล ธมฺโม อุปฺปชฺชติ อารมฺมณปจฺจยา. (สํขิตฺตํ.)}

\proseitem{5}{เหตุยา นว, อารมฺมเณ นว, อธิปติยา นว, อนนฺตเร นว, สมนนฺตเร นว, สหชาเต นว, อญฺญมญฺเญ นว, นิสฺสเย นว, อุปนิสฺสเย นว, ปุเรชาเต นว, อาเสวเน นว, กมฺเม นว, อาหาเร นว, อินฺทฺริเย นว, ฌาเน นว, มคฺเค นว, สมฺปยุตฺเต นว, วิปฺปยุตฺเต นว, อตฺถิยา นว, นตฺถิยา นว, วิคเต นว, อวิคเต นว. (อนุโลมํ.)}

\csromanheader{2}{\textbf{น-อธิปติ}}
\proseitem{6}{เหตุํ กุสลํ ธมฺมํ ปฏิจฺจ เหตุ กุสโล ธมฺโม อุปฺปชฺชติ น-อธิปติปจฺจยา.\csromanspacer{}\csromanspacer{}เหตุํ กุสลํ ธมฺมํ ปฏิจฺจ นเหตุ กุสโล ธมฺโม อุปฺปชฺชติ น-อธิปติปจฺจยา.\csromanspacer{}\csromanspacer{}เหตุํ กุสลํ ธมฺมํ ปฏิจฺจ เหตุ กุสโล จ นเหตุ กุสโล จ ธมฺมา อุปฺปชฺชนฺติ น-อธิปติปจฺจยา. (3)}

\proseitem{7}{นเหตุํ กุสลํ ธมฺมํ ปฏิจฺจ นเหตุ กุสโล ธมฺโม อุปฺปชฺชติ น-อธิปติปจฺจยา.\csromanspacer{} ตีณิ.}

\prose{เหตุํ กุสลญฺจ นเหตุํ กุสลญฺจ ธมฺมํ ปฏิจฺจ เหตุ กุสโล ธมฺโม อุปฺปชฺชติ น-อธิปติปจฺจยา.\csromanspacer{} ตีณิ.}

\csromanheader{2}{\textbf{นปุเรชาตาทิ}}
\proseitem{8}{เหตุํ กุสลํ ธมฺมํ ปฏิจฺจ เหตุ กุสโล ธมฺโม อุปฺปชฺชติ นปุเรชาตปจฺจยา.\csromanspacer{} นว.\csromanspacer{} นปจฺฉาชาตปจฺจยา.\csromanspacer{} นว.\csromanspacer{} น-อาเสวนปจฺจยา.\csromanspacer{} นว.}

\csromanheader{2}{\textbf{นกมฺม}}
\proseitem{9}{เหตุํ กุสลํ ธมฺมํ ปฏิจฺจ นเหตุ กุสโล ธมฺโม อุปฺปชฺชติ นกมฺมปจฺจยา. (1)}

\prose{นเหตุํ กุสลํ ธมฺมํ ปฏิจฺจ นเหตุ กุสโล ธมฺโม อุปฺปชฺชติ นกมฺมปจฺจยา. (1)}

\csromanheader{2}{1. เหตุทุก 1. กุสลตฺติก}
\prose{\csromanfolio{197}เหตุํ กุสลญฺจ นเหตุํ กุสลญฺจ ธมฺมํ ปฏิจฺจ นเหตุ กุสโล ธมฺโม อุปฺปชฺชติ นกมฺมปจฺจยา. (1)}

\csromanheader{2}{\textbf{นวิปากาทิ}}
\proseitem{10}{เหตุํ กุสลํ ธมฺมํ ปฏิจฺจ เหตุ กุสโล ธมฺโม อุปฺปชฺชติ นวิปากปจฺจยา.\csromanspacer{} นว ฯเปฯ.\csromanspacer{} นวิปฺปยุตฺตปจฺจยา.\csromanspacer{} นว.}

\proseitem{11}{น-อธิปติยา นว, นปุเรชาเต นว, นปจฺฉาชาเต นว, น-อาเสวเน นว, นกมฺเม ตีณิ, นวิปาเก นว, นวิปฺปยุตฺเต นว. (ปจฺจนียํ.)}

\prose{เหตุปจฺจยา น-อธิปติยา นว. (สํขิตฺตํ.\csromanspacer{} อนุโลม ปจฺจนียํ.)}

\prose{น-อธิปติปจฺจยา เหตุยา นว.\csromanspacer{} อารมฺมเณ นว. (สํขิตฺตํ.\csromanspacer{} ปจฺจนียานุโลมํ.)}

\prose{(สหชาตวารมฺปิ ปจฺจยวารมฺปิ นิสฺสยวารมฺปิ สํสฏฺฐวารมฺปิ สมฺปยุตฺตวารมฺปิ ปฏิจฺจวารสทิสา วิตฺถาเรตพฺพา.)}

\csromanheader{2}{1. กุสลปท 7. ปญฺหาวาร}
\csromanheader{2}{\textbf{ปจฺจยจตุกฺก-เหตุ}}
\proseitem{12}{เหตุ กุสโล ธมฺโม เหตุสฺส กุสลสฺส ธมฺมสฺส เหตุปจฺจเยน ปจฺจโย.\csromanspacer{}\csromanspacer{}เหตุ กุสโล ธมฺโม นเหตุสฺส กุสลสฺส ธมฺมสฺส เหตุปจฺจเยน ปจฺจโย.\csromanspacer{}\csromanspacer{}เหตุ กุสโล ธมฺโม เหตุสฺส กุสลสฺส จ นเหตุสฺส กุสลสฺส จ ธมฺมสฺส เหตุปจฺจเยน ปจฺจโย. (3)}

\csromanheader{2}{\textbf{อารมฺมณาทิ}}
\proseitem{13}{เหตุ กุสโล ธมฺโม เหตุสฺส กุสลสฺส ธมฺมสฺส อารมฺมณปจฺจเยน ปจฺจโย.\csromanspacer{} ตีณิ.}

\prose{นเหตุ กุสโล ธมฺโม นเหตุสฺส กุสลสฺส ธมฺมสฺส อารมฺมณปจฺจเยน ปจฺจโย.\csromanspacer{} ตีณิ.}

\gambhira{ธมฺมานุโลม ทุกติกปฏฺฐานปาฬิ}
\prose{\csromanfolio{198}เหตุ กุสโล จ นเหตุ กุสโล จ ธมฺมา เหตุสฺส กุสลสฺส ธมฺมสฺส อารมฺมณปจฺจเยน ปจฺจโย.\csromanspacer{} ตีณิ.}

\proseitem{14}{เหตุ กุสโล ธมฺโม เหตุสฺส กุสลสฺส ธมฺมสฺส อธิปติปจฺจเยน ปจฺจโย.\csromanspacer{} อารมฺมณาธิปติ, สหชาตาธิปติ.\csromanspacer{} ตีณิ.}

\prose{นเหตุ กุสโล ธมฺโม นเหตุสฺส กุสลสฺส ธมฺมสฺส อธิปติปจฺจเยน ปจฺจโย.\csromanspacer{} อารมฺมณาธิปติ, สหชาตาธิปติ.\csromanspacer{} ตีณิ.}

\prose{เหตุ กุสโล จ นเหตุ กุสโล จ ธมฺมา เหตุสฺส กุสลสฺส ธมฺมสฺส อธิปติปจฺจเยน ปจฺจโย.\csromanspacer{} อารมฺมณาธิปติ.\csromanspacer{} ตีณิ.}

\csromanheader{2}{\textbf{อนนฺตราทิ}}
\proseitem{15}{เหตุ กุสโล ธมฺโม เหตุสฺส กุสลสฺส ธมฺมสฺส อนนฺตรปจฺจเยน ปจฺจโย.\csromanspacer{} สมนนฺตรปจฺจเยน ปจฺจโย.\csromanspacer{} สหชาตปจฺจเยน ปจฺจโย.\csromanspacer{} อญฺญมญฺญปจฺจเยน ปจฺจโย.\csromanspacer{} นิสฺสยปจฺจเยน ปจฺจโย.}

\proseitem{16}{เหตุ กุสโล ธมฺโม เหตุสฺส กุสลสฺส ธมฺมสฺส อุปนิสฺสยปจฺจเยน ปจฺจโย.\csromanspacer{} อารมฺมณูปนิสฺสโย, อนนฺตรูปนิสฺสโย, ปกตูปนิสฺสโย.\csromanspacer{} ตีณิ.}

\prose{นเหตุ กุสโล ธมฺโม นเหตุสฺส กุสลสฺส ธมฺมสฺส อุปนิสฺสยปจฺจเยน ปจฺจโย.\csromanspacer{} อารมฺมณูปนิสฺสโย, อนนฺตรูปนิสฺสโย, ปกตูปนิสฺสโย.\csromanspacer{} ตีณิ.}

\prose{เหตุ กุสโล จ นเหตุ กุสโล จ ธมฺมา เหตุสฺส กุสลสฺส ธมฺมสฺส อุปนิสฺสยปจฺจเยน ปจฺจโย.\csromanspacer{} อารมฺมณูปนิสฺสโย, อนนฺตรูปนิสฺสโย, ปกตูปนิสฺสโย.\csromanspacer{} ตีณิ.\csromanspacer{} อาเสวนปจฺจเยน ปจฺจโย.\csromanspacer{} นว.}

\proseitem{17}{นเหตุ กุสโล ธมฺโม นเหตุสฺส กุสลสฺส ธมฺมสฺส กมฺมปจฺจเยน ปจฺจโย.\csromanspacer{}\csromanspacer{}นเหตุ กุสโล ธมฺโม เหตุสฺส กุสลสฺส ธมฺมสฺส กมฺมปจฺจเยน ปจฺจโย.\csromanspacer{}\csromanspacer{}นเหตุ กุสโล ธมฺโม เหตุสฺส กุสลสฺส จ นเหตุสฺส กุสลสฺส จ ธมฺมสฺส กมฺมปจฺจเยน ปจฺจโย. (3)}

\csromanheader{2}{1. เหตุทุก 1. กุสลตฺติก \textbf{อาหาราทิ}}
\proseitem{18}{\csromanfolio{199}นเหตุ กุสโล ธมฺโม นเหตุสฺส กุสลสฺส ธมฺมสฺส อาหารปจฺจเยน ปจฺจโย.\csromanspacer{} ตีณิ.}

\proseitem{19}{เหตุ กุสโล ธมฺโม เหตุสฺส กุสลสฺส ธมฺมสฺส อินฺทฺริยปจฺจเยน ปจฺจโย.\csromanspacer{} นว.}

\proseitem{20}{นเหตุ กุสโล ธมฺโม นเหตุสฺส กุสลสฺส ธมฺมสฺส ฌานปจฺจเยน ปจฺจโย.\csromanspacer{} ตีณิ.}

\proseitem{21}{เหตุ กุสโล ธมฺโม นเหตุสฺส กุสลสฺส ธมฺมสฺส มคฺคปจฺจเยน ปจฺจโย.\csromanspacer{} นว.}

\proseitem{22}{เหตุยา ตีณิ, อารมฺมเณ นว, อธิปติยา นว, อนนฺตเร นว, สมนนฺตเร นว, สหชาเต นว, อญฺญมญฺเญ นว, นิสฺสเย นว, อุปนิสฺสเย นว, อาเสวเน นว, กมฺเม ตีณิ, อาหาเร ตีณิ, อินฺทฺริเย นว, ฌาเน ตีณิ, มคฺเค นว, สมฺปยุตฺเต นว, อตฺถิยา นว, นตฺถิยา นว, วิคเต นว, อวิคเต นว.}

\csromanheader{2}{\textbf{ปจฺจนียุทฺธาร}}
\proseitem{23}{เหตุ กุสโล ธมฺโม เหตุสฺส กุสลสฺส ธมฺมสฺส อารมฺมณปจฺจเยน ปจฺจโย.\csromanspacer{} สหชาตปจฺจเยน ปจฺจโย.\csromanspacer{} อุปนิสฺสยปจฺจเยน ปจฺจโย.\csromanspacer{} ตีณิ. (3)}

\proseitem{24}{นเหตุ กุสโล ธมฺโม นเหตุสฺส กุสลสฺส ธมฺมสฺส อารมฺมณปจฺจเยน ปจฺจโย.\csromanspacer{} สหชาตปจฺจเยน ปจฺจโย.\csromanspacer{} อุปนิสฺสยปจฺจเยน ปจฺจโย.\csromanspacer{}\csromanspacer{}นเหตุ กุสโล ธมฺโม เหตุสฺส กุสลสฺส ธมฺมสฺส อารมฺมณปจฺจเยน ปจฺจโย.\csromanspacer{} สหชาตปจฺจเยน ปจฺจโย.\csromanspacer{} อุปนิสฺสยปจฺจเยน ปจฺจโย.\csromanspacer{}\csromanspacer{}นเหตุ กุสโล ธมฺโม เหตุสฺส กุสลสฺส จ นเหตุสฺส กุสลสฺส จ ธมฺมสฺส อารมฺมณปจฺจเยน ปจฺจโย.\csromanspacer{} สหชาตปจฺจเยน ปจฺจโย.\csromanspacer{} อุปนิสฺสยปจฺจเยน ปจฺจโย. (3)}

\gambhira{ธมฺมานุโลม ทุกติกปฏฺฐานปาฬิ}
\proseitem{25}{\csromanfolio{200}เหตุ กุสโล จ นเหตุ กุสโล จ ธมฺมา เหตุสฺส กุสลสฺส ธมฺมสฺส อารมฺมณปจฺจเยน ปจฺจโย.\csromanspacer{} สหชาตปจฺจเยน ปจฺจโย.\csromanspacer{} อุปนิสฺสยปจฺจเยน ปจฺจโย.\csromanspacer{} ตีณิ. (3)}

\proseitem{26}{นเหตุยา นว, น-อารมฺมเณ นว, น-อธิปติยา นว ฯเปฯ โน-อวิคเต นว. (สํขิตฺตํ.\csromanspacer{} ปจฺจนียํ.)}

\prose{เหตุปจฺจยา น-อารมฺมเณ ตีณิ ฯเปฯ โนนตฺถิยา ตีณิ, โนวิคเต ตีณิ. (สํขิตฺตํ.\csromanspacer{} อนุโลมปจฺจนียํ.)}

\prose{นเหตุปจฺจยา อารมฺมเณ นว ฯเปฯ อวิคเต นว. (สํขิตฺตํ.\csromanspacer{} ปจฺจนียานุโลมํ.)}

\prose{ยถา กุสลตฺติเก ปญฺหาวารสฺส อนุโลมมฺปิ ปจฺจนียมฺปิ อนุโลมปจฺจนียมฺปิ ปจฺจนียานุโลมมฺปิ คณิตํ, เอวํ คเณตพฺพํ.)}

\csromanheader{2}{2. อกุสลปท 1. ปฏิจฺจวาร}
\csromanheader{2}{\textbf{ปจฺจยจตุกฺก-เหตุ}}
\proseitem{27}{เหตุํ อกุสลํ ธมฺมํ ปฏิจฺจ เหตุ อกุสโล ธมฺโม อุปฺปชฺชติ เหตุปจฺจยา.\csromanspacer{}\csromanspacer{}เหตุํ อกุสลํ ธมฺมํ ปฏิจฺจ นเหตุ อกุสโล ธมฺโม อุปฺปชฺชติ เหตุปจฺจยา.\csromanspacer{}\csromanspacer{}เหตุํ อกุสลํ ธมฺมํ ปฏิจฺจ เหตุ อกุสโล จ นเหตุ อกุสโล จ ธมฺมา อุปฺปชฺชนฺติ เหตุปจฺจยา. (3)}

\proseitem{28}{นเหตุํ อกุสลํ ธมฺมํ ปฏิจฺจ นเหตุ อกุสโล ธมฺโม อุปฺปชฺชติ เหตุปจฺจยา.\csromanspacer{}\csromanspacer{}นเหตุํ อกุสลํ ธมฺมํ ปฏิจฺจ เหตุ อกุสโล ธมฺโม อุปฺปชฺชติ เหตุปจฺจยา.\csromanspacer{}\csromanspacer{}นเหตุํ อกุสลํ ธมฺมํ ปฏิจฺจ เหตุ อกุสโล จ นเหตุ อกุสโล จ ธมฺมา อุปฺปชฺชนฺติ เหตุปจฺจยา. (3)}

\proseitem{29}{เหตุํ อกุสลญฺจ นเหตุํ อกุสลญฺจ ธมฺมํ ปฏิจฺจ เหตุ อกุสโล ธมฺโม อุปฺปชฺชติ เหตุปจฺจยา.\csromanspacer{}\csromanspacer{}เหตุํ อกุสลญฺจ นเหตุํ อกุสลญฺจ ธมฺมํ ปฏิจฺจ นเหตุ อกุสโล ธมฺโม อุปฺปชฺชติ เหตุปจฺจยา.\csromanspacer{}\csromanspacer{}เหตุํ อกุสลญฺจ นเหตุํ อกุสลญฺจ ธมฺมํ ปฏิจฺจ เหตุ อกุสโล จ นเหตุ อกุสโล จ ธมฺมา อุปฺปชฺชนฺติ เหตุปจฺจยา. (3)}

\csromanheader{2}{1. เหตุทุก 1. กุสลตฺติก}
\proseitem{30}{\csromanfolio{201}เหตุยา นว, อารมฺมเณ นว, อธิปติยา นว ฯเปฯ กมฺเม นว, อาหาเร นว ฯเปฯ อวิคเต นว. (สํขิตฺตํ.\csromanspacer{} อนุโลมํ.)}

\csromanheader{2}{\textbf{นเหตุ}}
\vspace{\tipitakaparskip}
\csromancenter{\textbf{นเหตุ}ํ อกุสลํ ธมฺมํ ปฏิจฺจ เหตุ อกุสโล ธมฺโม อุปฺปชฺชติ นเหตุปจฺจยา. (1)}
\csromanheader{2}{\textbf{น-อธิปตฺยาทิ}}
\proseitem{32}{เหตุํ อกุสลํ ธมฺมํ ปฏิจฺจ เหตุ อกุสโล ธมฺโม อุปฺปชฺชติ น-อธิปติปจฺจยา.\csromanspacer{} นว ฯเปฯ.}

\prose{เหตุํ อกุสลํ ธมฺมํ ปฏิจฺจ นเหตุ อกุสโล ธมฺโม อุปฺปชฺชติ นกมฺมปจฺจยา. (1)}

\vspace{\tipitakaparskip}
\csromancenter{\textbf{นเหตุ}ํ อกุสลํ ธมฺมํ ปฏิจฺจ นเหตุ อกุสโล ธมฺโม อุปฺปชฺชติ นกมฺมปจฺจยา. (1)}
\prose{เหตุํ อกุสลญฺจ นเหตุํ อกุสลญฺจ ธมฺมํ ปฏิจฺจ นเหตุ อกุสโล ธมฺโม อุปฺปชฺชติ นกมฺมปจฺจยา. (1) (สํขิตฺตํ.)}

\vspace{\tipitakaparskip}
\csromancenter{\textbf{นเหตุ}ยา เอกํ, น-อธิปติยา นว, นปุเรชาเต นว, นปจฺฉาชาเต นว, น-อาเสวเน นว, นกมฺเม ตีณิ, นวิปาเก นว, นวิปฺปยุตฺเต นว. (ปจฺจนียํ.)}
\prose{เหตุปจฺจยา น-อธิปติยา นว, (สํขิตฺตํ.\csromanspacer{} อนุโลมปจฺจนียํ.)}

\csromanheader{2}{\textbf{นเหตุ}ปจฺจยา อารมฺมเณ เอกํ. (สํขิตฺตํ. ปจฺจนียานุโลมํ.)}
\prose{(สหชาตวารมฺปิ ปจฺจยวารมฺปิ นิสฺสยวารมฺปิ สํสฏฺฐวารมฺปิ สมฺปยุตฺตวารมฺปิ ปฏิจฺจวารสทิสา วิตฺถาเรตพฺพา.)}

\csromanheader{2}{2. อกุสลปท 7. ปญฺหาวาร}
\csromanheader{2}{\textbf{ปจฺจยจตุกฺก-เหตุ}}
\proseitem{34}{เหตุ อกุสโล ธมฺโม เหตุสฺส อกุสลสฺส ธมฺมสฺส เหตุปจฺจเยน ปจฺจโย.\csromanspacer{}\csromanspacer{}เหตุ อกุสโล ธมฺโม นเหตุสฺส}

\vspace{\tipitakaparskip}
\csromancenter{ธมฺมานุโลม ทุกติกปฏฺฐานปาฬิ อกุสลสฺส ธมฺมสฺส เหตุปจฺจเยน ปจฺจโย.\csromanspacer{}\csromanspacer{}เหตุ อกุสโล ธมฺโม เหตุสฺส อกุสลสฺส จ นเหตุสฺส อกุสลสฺส จ ธมฺมสฺส เหตุปจฺจเยน ปจฺจโย. (3)}
\csromanheader{2}{\textbf{อารมฺมณาทิ}}
\proseitem{35}{\csromanfolio{202}เหตุ อกุสโล ธมฺโม เหตุสฺส อกุสลสฺส ธมฺมสฺส อารมฺมณปจฺจเยน ปจฺจโย.\csromanspacer{} ตีณิ.}

\prose{นเหตุ อกุสโล ธมฺโม นเหตุสฺส อกุสลสฺส ธมฺมสฺส อารมฺมณปจฺจเยน ปจฺจโย.\csromanspacer{} ตีณิ.}

\prose{เหตุ อกุสโล จ นเหตุ อกุสโล จ ธมฺมา เหตุสฺส อกุสลสฺส ธมฺมสฺส อารมฺมณปจฺจเยน ปจฺจโย.\csromanspacer{} ตีณิ.}

\proseitem{36}{เหตุ อกุสโล ธมฺโม เหตุสฺส อกุสลสฺส ธมฺมสฺส อธิปติปจฺจเยน ปจฺจโย.\csromanspacer{} อารมฺมณาธิปติ.\csromanspacer{} ตีณิ.}

\prose{นเหตุ อกุสโล ธมฺโม นเหตุสฺส อกุสลสฺส ธมฺมสฺส อธิปติปจฺจเยน ปจฺจโย.\csromanspacer{} อารมฺมณาธิปติ, สหชาตาธิปติ.\csromanspacer{} ตีณิ.}

\prose{เหตุ อกุสโล จ นเหตุ อกุสโล จ ธมฺมา เหตุสฺส อกุสลสฺส ธมฺมสฺส อธิปติปจฺจเยน ปจฺจโย.\csromanspacer{} อารมฺมณาธิปติ.\csromanspacer{} ตีณิ ฯเปฯ.}

\proseitem{37}{นเหตุ อกุสโล ธมฺโม นเหตุสฺส อกุสลสฺส ธมฺมสฺส กมฺมปจฺจเยน ปจฺจโย.\csromanspacer{}\csromanspacer{}นเหตุ อกุสโล ธมฺโม เหตุสฺส อกุสลสฺส ธมฺมสฺส กมฺมปจฺจเยน ปจฺจโย.\csromanspacer{}\csromanspacer{}นเหตุ อกุสโล ธมฺโม เหตุสฺส อกุสลสฺส จ นเหตุสฺส อกุสลสฺส จ ธมฺมสฺส กมฺมปจฺจเยน ปจฺจโย. (3)}

\csromanheader{2}{\textbf{อาหาราทิ}}
\proseitem{38}{นเหตุ อกุสโล ธมฺโม นเหตุสฺส อกุสลสฺส ธมฺมสฺส อาหารปจฺจเยน ปจฺจโย.\csromanspacer{} ตีณิ.}

\proseitem{39}{นเหตุ อกุสโล ธมฺโม นเหตุสฺส อกุสลสฺส ธมฺมสฺส อินฺทฺริยปจฺจเยน ปจฺจโย.\csromanspacer{} ตีณิ.}

\csromanheader{2}{1. เหตุทุก 1. กุสลตฺติก}
\proseitem{40}{\csromanfolio{203}นเหตุ อกุสโล ธมฺโม นเหตุสฺส อกุสลสฺส ธมฺมสฺส ฌานปจฺจเยน ปจฺจโย.\csromanspacer{} ตีณิ. (สํขิตฺตํ.)}

\proseitem{41}{เหตุยา ตีณิ, อารมฺมเณ นว, อธิปติยา นว, อนนฺตเร นว, สมนนฺตเร นว, สหชาเต นว, อญฺญมญฺเญ นว, นิสฺสเย นว, อุปนิสฺสเย นว, อาเสวเน นว, กมฺเม ตีณิ, อาหาเร ตีณิ, อินฺทฺริเย ตีณิ, ฌาเน ตีณิ, มคฺเค ตีณิ, สมฺปยุตฺเต นว, อตฺถิยา นว, นตฺถิยา นว, วิคเต นว, อวิคเต นว.}

\csromanheader{2}{\textbf{ปจฺจนียุทฺธาร}}
\proseitem{42}{เหตุ อกุสโล ธมฺโม เหตุสฺส อกุสลสฺส ธมฺมสฺส อารมฺมณปจฺจเยน ปจฺจโย.\csromanspacer{} สหชาตปจฺจเยน ปจฺจโย.\csromanspacer{} อุปนิสฺสยปจฺจเยน ปจฺจโย. (สํขิตฺตํ.)}

\proseitem{43}{นเหตุยา นว, น-อารมฺมเณ นว, น-อธิปติยา นว, น-อนนฺตเร นว, นสมนนฺตเร นว, นสหชาเต นว, น-อญฺญมญฺเญ นว ฯเปฯ โน-อวิคเต นว. (ปจฺจนียํ.)}

\prose{เหตุปจฺจยา น-อารมฺมเณ ตีณิ. (สํขิตฺตํ.\csromanspacer{} อนุโลมปจฺจนียํ.)}

\csromanheader{2}{นเหตุปจฺจยา อารมฺมเณ นว. (สํขิตฺตํ. ปจฺจนียานุโลมํ.)}
\prose{(ยถา กุสลตฺติเก ปญฺหาวารสฺส อนุโลมมฺปิ ปจฺจนียมฺปิ อนุโลมปจฺจนียมฺปิ ปจฺจนียานุโลมมฺปิ คณิตํ, เอวํ คเณตพฺพํ.)}

\csromanheader{2}{3. อพฺยากตปท 1. ปฏิจฺจวาร}
\csromanheader{2}{\textbf{ปจฺจยจตุกฺก-เหตฺวาทิ}}
\proseitem{44}{เหตุํ อพฺยากตํ ธมฺมํ ปฏิจฺจ เหตุ อพฺยากโต ธมฺโม อุปฺปชฺชติ เหตุปจฺจยา.\csromanspacer{}\csromanspacer{}เหตุํ อพฺยากตํ ธมฺมํ ปฏิจฺจ นเหตุ อพฺยากโต ธมฺโม อุปฺปชฺชติ เหตุปจฺจยา.\csromanspacer{}\csromanspacer{}เหตุํ อพฺยากตํ ธมฺมํ ปฏิจฺจ เหตุ อพฺยากโต จ นเหตุ อพฺยากโต จ ธมฺมา อุปฺปชฺชนฺติ เหตุปจฺจยา. (3)}

\gambhira{ธมฺมานุโลม ทุกติกปฏฺฐานปาฬิ}
\proseitem{45}{\csromanfolio{204}นเหตุํ อพฺยากตํ ธมฺมํ ปฏิจฺจ นเหตุ อพฺยากโต ธมฺโม อุปฺปชฺชติ เหตุปจฺจยา.\csromanspacer{} ตีณิ.}

\prose{เหตุํ อพฺยากตญฺจ นเหตุํ อพฺยากตญฺจ ธมฺมํ ปฏิจฺจ เหตุ อพฺยากโต ธมฺโม อุปฺปชฺชติ เหตุปจฺจยา.\csromanspacer{} ตีณิ.}

\proseitem{46}{เหตุํ อพฺยากตํ ธมฺมํ ปฏิจฺจ เหตุ อพฺยากโต ธมฺโม อุปฺปชฺชติ อารมฺมณปจฺจยา. (สํขิตฺตํ.)}

\proseitem{47}{เหตุยา นว, อารมฺมเณ นว, อธิปติยา นว, อนนฺตเร นว, สมนนฺตเร นว, สหชาเต นว, อญฺญมญฺเญ นว, นิสฺสเย นว, อุปนิสฺสเย นว, ปุเรชาเต นว, อาเสวเน นว, กมฺเม นว, วิปาเก นว, อาหาเร นว, อินฺทฺริเย นว, ฌาเน นว, มคฺเค นว, สมฺปยุตฺเต นว, วิปฺปยุตฺเต นว, อตฺถิยา นว, นตฺถิยา นว, วิคเต นว, อวิคเต นว. (อนุโลมํ.)}

\csromanheader{2}{\textbf{นเหตุ น-อารมฺมณ}}
\proseitem{48}{นเหตุํ อพฺยากตํ ธมฺมํ ปฏิจฺจ นเหตุ อพฺยากโต ธมฺโม อุปฺปชฺชติ นเหตุปจฺจยา. (1)}

\proseitem{49}{เหตุํ อพฺยากตํ ธมฺมํ ปฏิจฺจ นเหตุ อพฺยากโต ธมฺโม อุปฺปชฺชติ น-อารมฺมณปจฺจยา. (1)}

\prose{นเหตุํ อพฺยากตํ ธมฺมํ ปฏิจฺจ นเหตุ อพฺยากโต ธมฺโม อุปฺปชฺชติ น-อารมฺมณปจฺจยา. (1)}

\prose{เหตุํ อพฺยากตญฺจ นเหตุํ อพฺยากตญฺจ ธมฺมํ ปฏิจฺจ นเหตุ อพฺยากโต ธมฺโม อุปฺปชฺชติ น-อารมฺมณปจฺจยา. (1)}

\csromanheader{2}{\textbf{น-อธิปตฺยาทิ}}
\proseitem{50}{เหตุํ อพฺยากตํ ธมฺมํ ปฏิจฺจ เหตุ อพฺยากโต ธมฺโม อุปฺปชฺชติ น-อธิปติปจฺจยา.\csromanspacer{} นว ฯเปฯ.}

\proseitem{51}{เหตุํ อพฺยากตํ ธมฺมํ ปฏิจฺจ นเหตุ อพฺยากโต ธมฺโม อุปฺปชฺชติ นกมฺมปจฺจยา. (1)}

\csromanheader{2}{1. เหตุทุก 1. กุสลตฺติก}
\prose{\csromanfolio{205}นเหตุํ อพฺยากตํ ธมฺมํ ปฏิจฺจ นเหตุ อพฺยากโต ธมฺโม อุปฺปชฺชติ นกมฺมปจฺจยา. (1)}

\prose{เหตุํ อพฺยากตญฺจ นเหตุํ อพฺยากตญฺจ ธมฺมํ ปฏิจฺจ นเหตุ อพฺยากโต ธมฺโม อุปฺปชฺชติ นกมฺมปจฺจยา. (1)}

\proseitem{52}{นเหตุํ อพฺยากตํ ธมฺมํ ปฏิจฺจ นเหตุ อพฺยากโต ธมฺโม อุปฺปชฺชติ น-อาหารปจฺจยา.\csromanspacer{} น-อินฺทฺริยปจฺจยา.\csromanspacer{} นฌานปจฺจยา. (สํขิตฺตํ.)}

\proseitem{53}{นหุตุยา เอกํ, น-อารมฺมเณ ตีณิ, น-อธิปติยา นว, น-อนนฺตเร ตีณิ, นสมนนฺตเร ตีณิ, น-อญฺญมญฺเญ ตีณิ.\csromanspacer{} น-อุปนิสฺสเย ตีณิ, นปุเรชาเต นว, นปจฺฉาชาเต นว, น-อาเสวเน นว, นกมฺเม ตีณิ, นวิปาเก นว, น-อาหาเร เอกํ, น-อินฺทฺริเย เอกํ, นฌาเน เอกํ, นมคฺเค เอกํ, นสมฺปยุตฺเต ตีณิ, นวิปฺปยุตฺเต นว, โนนตฺถิยา ตีณิ, โนวิคเต ตีณิ. (ปจฺจนียํ.)}

\prose{เหตุปจฺจยา น-อารมฺมเณ ตีณิ, น-อธิปติยา นว. (สํขิตฺตํ.\csromanspacer{} อนุโลมปจฺจนียํ.)}

\csromanheader{2}{นเหตุปจฺจยา อารมฺมเณ เอกํ. (สํขิตฺตํ. ปจฺจนียานุโลมํ.)}
\prose{(สหชาตวาโรปิ ปจฺจยวาโรปิ นิสฺสยวาโรปิ สํสฏฺฐวาโรปิ สมฺปยุตฺตวาโรปิ ปฏิจฺจวารสทิสา วิตฺถาเรตพฺพา.)}

\csromanheader{2}{3. อพฺยากตปท 7. ปญฺหาวาร}
\csromanheader{2}{\textbf{ปจฺจยจตุกฺก-เหตุ}}
\proseitem{54}{เหตุ อพฺยากโต ธมฺโม เหตุสฺส อพฺยากตสฺส ธมฺมสฺส เหตุปจฺจเยน ปจฺจโย.\csromanspacer{}\csromanspacer{}เหตุ อพฺยากโต ธมฺโม นเหตุสฺส อพฺยากตสฺส ธมฺมสฺส เหตุปจฺจเยน ปจฺจโย.\csromanspacer{}\csromanspacer{}เหตุ อพฺยากโต ธมฺโม เหตุสฺส อพฺยากตสฺส จ นเหตุสฺส อพฺยากตสฺส จ ธมฺมสฺส เหตุปจฺจเยน ปจฺจโย. (3)}

\csromanheader{2}{ธมฺมานุโลม ทุกติกปฏฺฐานปาฬิ \textbf{อารมฺมณาทิ}}
\proseitem{55}{\csromanfolio{206}เหตุ อพฺยากโต ธมฺโม เหตุสฺส อพฺยากตสฺส ธมฺมสฺส อารมฺมณปจฺจเยน ปจฺจโย.\csromanspacer{} นว.}

\proseitem{56}{เหตุ อพฺยากโต ธมฺโม เหตุสฺส อพฺยากตสฺส ธมฺมสฺส อธิปติปจฺจเยน ปจฺจโย.\csromanspacer{} อารมฺมณาธิปติ, สหชาตาธิปติ.\csromanspacer{} ตีณิ.}

\prose{นเหตุ อพฺยากโต ธมฺโม นเหตุสฺส อพฺยากตสฺส ธมฺมสฺส อธิปติปจฺจเยน ปจฺจโย.\csromanspacer{} อารมฺมณาธิปติ, สหชาตาธิปติ.\csromanspacer{} ตีณิ.}

\prose{เหตุ อพฺยากโต จ นเหตุ อพฺยากโต จ ธมฺมา เหตุสฺส อพฺยากตสฺส ธมฺมสฺส อธิปติปจฺจเยน ปจฺจโย.\csromanspacer{} อารมฺมณาธิปติ.\csromanspacer{} ตีณิ ฯเปฯ.}

\csromanheader{2}{\textbf{ปุเรชาตาทิ}}
\proseitem{57}{นเหตุ อพฺยากโต ธมฺโม นเหตุสฺส อพฺยากตสฺส ธมฺมสฺส ปุเรชาตปจฺจเยน ปจฺจโย.\csromanspacer{} ตีณิ.}

\proseitem{58}{เหตุ อพฺยากโต ธมฺโม นเหตุสฺส อพฺยากตสฺส ธมฺมสฺส ปจฺฉาชาตปจฺจเยน ปจฺจโย. (1)}

\prose{นเหตุ อพฺยากโต ธมฺโม นเหตุสฺส อพฺยากตสฺส ธมฺมสฺส ปจฺฉาชาตปจฺจเยน ปจฺจโย. (1)}

\prose{เหตุ อพฺยากโต จ นเหตุ อพฺยากโต จ ธมฺมา นเหตุสฺส อพฺยากตสฺส ธมฺมสฺส ปจฺฉาชาตปจฺจเยน ปจฺจโย. (1)}

\proseitem{59}{นเหตุ อพฺยากโต ธมฺโม นเหตุสฺส อพฺยากตสฺส ธมฺมสฺส กมฺมปจฺจเยน ปจฺจโย.\csromanspacer{} ตีณิ.}

\proseitem{60}{เหตุ อพฺยากโต ธมฺโม เหตุสฺส อพฺยากตสฺส ธมฺมสฺส วิปากปจฺจเยน ปจฺจโย.\csromanspacer{} นว.}

\proseitem{61}{นเหตุ อพฺยากโต ธมฺโม นเหตุสฺส อพฺยากตสฺส ธมฺมสฺส อาหารปจฺจเยน ปจฺจโย.\csromanspacer{} อินฺทฺริยปจฺจเยน ปจฺจโย.\csromanspacer{} ฌานปจฺจเยน ปจฺจโย.\csromanspacer{} มคฺคปจฺจเยน ปจฺจโย.\csromanspacer{} สมฺปยุตฺตปจฺจเยน ปจฺจโย.}

\csromanheader{2}{1. เหตุทุก 1. กุสลตฺติก \textbf{วิปฺปยุตฺต}}
\proseitem{62}{\csromanfolio{207}เหตุ อพฺยากโต ธมฺโม นเหตุสฺส อพฺยากตสฺส ธมฺมสฺส วิปฺปยุตฺตปจฺจเยน ปจฺจโย. (1)}

\prose{นเหตุ อพฺยากโต ธมฺโม นเหตุสฺส อพฺยากตสฺส ธมฺมสฺส วิปฺปยุตฺตปจฺจเยน ปจฺจโย.\csromanspacer{} ตีณิ.}

\prose{เหตุ อพฺยากโต จ นเหตุ อพฺยากโต จ ธมฺมา นเหตุสฺส อพฺยากตสฺส ธมฺมสฺส วิปฺปยุตฺตปจฺจเยน ปจฺจโย. (1) (สํขิตฺตํ.)}

\proseitem{63}{เหตุยา ตีณิ, อารมฺมเณ นว, อธิปติยา นว, อนนฺตเร นว, สมนนฺตเร นว, สหชาเต นว, อญฺญมญฺเญ นว, นิสฺสเย นว, อุปนิสฺสเย นว, ปุเรชาเต ตีณิ, ปจฺฉาชาเต ตีณิ, อาเสวเน นว, กมฺเม ตีณิ, วิปาเก นว, อาหาเร ตีณิ, อินฺทฺริเย นว, ฌาเน ตีณิ, มคฺเค นว, สมฺปยุตฺเต นว, วิปฺปยุตฺเต ปญฺจ, อตฺถิยา นว, นตฺถิยา นว, วิคเต นว, อวิคเต นว.}

\csromanheader{2}{\textbf{ปจฺจนียุทฺธาร}}
\proseitem{64}{เหตุ อพฺยากโต ธมฺโม เหตุสฺส อพฺยากตสฺส ธมฺมสฺส อารมฺมณปจฺจเยน ปจฺจโย.\csromanspacer{} สหชาตปจฺจเยน ปจฺจโย.\csromanspacer{} อุปนิสฺสยปจฺจเยน ปจฺจโย.\csromanspacer{}\csromanspacer{}เหตุ อพฺยากโต ธมฺโม นเหตุสฺส อพฺยากตสฺส ธมฺมสฺส อารมฺมณปจฺจเยน ปจฺจโย.\csromanspacer{} สหชาตปจฺจเยน ปจฺจโย.\csromanspacer{} อุปนิสฺสยปจฺจเยน ปจฺจโย.\csromanspacer{} ปจฺฉาชาตปจฺจเยน ปจฺจโย.\csromanspacer{}\csromanspacer{}เหตุ อพฺยากโต ธมฺโม เหตุสฺส อพฺยากตสฺส จ นเหตุสฺส อพฺยากตสฺส จ ธมฺมสฺส อารมฺมณปจฺจเยน ปจฺจโย.\csromanspacer{} สหชาตปจฺจเยน ปจฺจโย.\csromanspacer{} อุปนิสฺสยปจฺจเยน ปจฺจโย. (3)}

\proseitem{65}{นเหตุ อพฺยากโต ธมฺโม นเหตุสฺส อพฺยากตสฺส ธมฺมสฺส อารมฺมณปจฺจเยน ปจฺจโย.\csromanspacer{} สหชาตปจฺจเยน ปจฺจโย.\csromanspacer{} อุปนิสฺสยปจฺจเยน ปจฺจโย.\csromanspacer{} ปุเรชาตปจฺจเยน ปจฺจโย.\csromanspacer{} ปจฺฉาชาตปจฺจเยน ปจฺจโย.\csromanspacer{} อาหารปจฺจเยน ปจฺจโย.\csromanspacer{} อินฺทฺริยปจฺจเยน ปจฺจโย.}

\proseitem{66}{นเหตุยา นว, น-อารมฺมเณ นว, น-อธิปติยา นว. (สํขิตฺตํ.\csromanspacer{} ปจฺจนียํ.)}

\gambhira{ธมฺมานุโลม ทุกติกปฏฺฐานปาฬิ}
\prose{\csromanfolio{208}เหตุปจฺจยา น-อารมฺมเณ ตีณิ. (สํขิตฺตํ.\csromanspacer{} อนุโลมปจฺจนียํ.)}

\csromanheader{2}{นเหตุปจฺจยา อารมฺมเณ นว. (สํขิตฺตํ. ปจฺจนียานุโลมํ.)}
\prose{(ยถา กุสลตฺติเก ปญฺหาวารสฺส อนุโลมมฺปิ ปจฺจนียมฺปิ อนุโลมปจฺจนียมฺปิ ปจฺจนียานุโลมมฺปิ คณิตํ, เอวํ คเณตพฺพํ.)}

\nitthitam{เหตุทุกสลตฺติกํ นิฏฺฐิตํ.}
\csromanmatikaanchor{mtk.33}
\csromantocmarkref{section}{2. เวทนาตฺติก}{mtk.33}
\bookmark[dest={mtk.33},level=1]{2. เวทนาตฺติก}
\csromanheader{1}{1. เหตุทุก 2. เวทนาตฺติก}
\csromanheader{2}{1. สุขายเวทนายสมฺปยุตฺตปท 1. ปฏิจฺจวาร}
\csromanheader{2}{\textbf{ปจฺจยจตุกฺก-เหตุ}}
\proseitem{67}{เหตุํ สุขาย เวทนาย สมฺปยุตฺตํ ธมฺมํ ปฏิจฺจ เหตุ สุขาย เวทนาย สมฺปยุตฺโต ธมฺโม อุปฺปชฺชติ เหตุปจฺจยา.\csromanspacer{}\csromanspacer{}เหตุํ สุขาย เวทนาย สมฺปยุตฺตํ ธมฺมํ ปฏิจฺจ นเหตุ สุขาย เวทนาย สมฺปยุตฺโต ธมฺโม อุปฺปชฺชติ เหตุปจฺจยา.\csromanspacer{}\csromanspacer{}เหตุํ สุขาย เวทนาย สมฺปยุตฺตํ ธมฺมํ ปฏิจฺจ เหตุ สุขาย เวทนาย สมฺปยุตฺโต จ นเหตุ สุขาย เวทนาย สมฺปยุตฺโต จ ธมฺมา อุปฺปชฺชติ เหตุปจฺจยา. (3)}

\proseitem{68}{นเหตุํ สุขาย เวทนาย สมฺปยุตฺตํ ธมฺมํ ปฏิจฺจ นเหตุ สุขาย เวทนาย สมฺปยุตฺโต ธมฺโม อุปฺปชฺชติ เหตุปจฺจยา.\csromanspacer{}\csromanspacer{}นเหตุํ สุขาย เวทนาย สมฺปยุตฺตํ ธมฺมํ ปฏิจฺจ เหตุ สุขาย เวทนาย สมฺปยุตฺโต ธมฺโม อุปฺปชฺชติ เหตุปจฺจยา.\csromanspacer{}\csromanspacer{}นเหตุํ สุขาย เวทนาย สมฺปยุตฺตํ ธมฺมํ ปฏิจฺจ เหตุ สุขาย เวทนาย สมฺปยุตฺโต จ นเหตุ สุขาย เวทนาย สมฺปยุตฺโต จ ธมฺโม อุปฺปชฺชนฺติ เหตุปจฺจยา. (3)}

\proseitem{69}{เหตุํ สุขาย เวทนาย สมฺปยุตฺตญฺจ นเหตุํ สุขาย เวทนาย สมฺปยุตฺตญฺจ ธมฺมํ ปฏิจฺจ เหตุ สุขาย เวทนาย สมฺปยุตฺโต ธมฺโม อุปฺปชฺชติ เหตุปจฺจยา.\csromanspacer{}\csromanspacer{}เหตุํ สุขาย เวทนาย สมฺปยุตฺตญฺจ นเหตุํ สุขาย เวทนาย สมฺปยุตฺตญฺจ ธมฺมํ ปฏิจฺจ นเหตุ สุขาย เวทนาย สมฺปยุตฺโต ธมฺโม อุปฺปชฺชติ เหตุปจฺจยา.\csromanspacer{}\csromanspacer{}เหตุํ สุขาย เวทนาย}

\vspace{\tipitakaparskip}
\csromancenter{เหตุทุก 2.\csromanspacer{} เวทนาตฺติก สมฺปยุตฺตญฺจ นเหตุํ สุขาย เวทนาย สมฺปยุตฺตญฺจ ธมฺมํ ปฏิจฺจ เหตุ สุขาย เวทนาย สมฺปยุตฺโต จ นเหตุ สุขาย เวทนาย สมฺปยุตฺโต จ ธมฺมา อุปฺปชฺชนฺติ เหตุปจฺจยา. (3) (สํขิตฺตํ.)}
\proseitem{70}{\csromanfolio{209}เหตุยา นว, อารมฺมเณ นว, อธิปติยา นว, อนนฺตเร นว, สมนนฺตเร นว, สหชาเต นว, อญฺญมญฺเญ นว, นิสฺสเย นว, อุปนิสฺสเย นว, ปุเรชาเต นว, อาเสวเน นว, กมฺเม นว, วิปาเก นว, อาหาเร นว ฯเปฯ อวิคเต นว. (อนุโลมํ.)}

\csromanheader{2}{\textbf{นเหตุ น-อธิปติ}}
\proseitem{71}{นเหตุํ สุขาย เวทนาย สมฺปยุตฺตํ ธมฺมํ ปฏิจฺจ นเหตุ สุขาย เวทนาย สมฺปยุตฺโต ธมฺโม อุปฺปชฺชติ นเหตุปจฺจยา. (1)}

\proseitem{72}{เหตุํ สุขาย เวทนาย สมฺปยุตฺตํ ธมฺมํ ปฏิจฺจ เหตุ สุขาย เวทนาย สมฺปยุตฺโต ธมฺโม อุปฺปชฺชติ น-อธิปติปจฺจยา.\csromanspacer{}\csromanspacer{}เหตุํ สุขาย เวทนาย สมฺปยุตฺตํ ธมฺมํ ปฏิจฺจ นเหตุ สุขาย เวทนาย สมฺปยุตฺโต ธมฺโม อุปฺปชฺชติ น-อธิปติปจฺจยา. (สํขิตฺตํ.)}

\proseitem{73}{นเหตุยา เอกํ, น-อธิปติยา นว, นปุเรชาเต นว, นปจฺฉาชาเต นว, น-อาเสวเน นว, นกมฺเม ตีณิ, นวิปาเก นว, นฌาเน เอกํ, นมคฺเค เอกํ, นวิปฺปยุตฺเต นว. (ปจฺจนียํ.)}

\prose{เหตุปจฺจยา น-อธิปติยา นว. (สํขิตฺตํ.\csromanspacer{} อนุโลมปจฺจนียํ.)}

\csromanheader{2}{นเหตุปจฺจยา อารมฺมเณ เอกํ. (สํขิตฺตํ. ปจฺจนียานุโลมํ.)}
\prose{(สหชาตวาโรปิ ปจฺจยวาโรปิ นิสฺสยวาโรปิ สํสฏฺฐวาโรปิ สมฺปยุตฺตวาโรปิ ปฏิจฺจวารสทิสา วิตฺถาเรตพฺพา.)}

\csromanheader{2}{1. สุขายเวทนายสมฺปยุตฺตปท 7. ปญฺหาวาร}
\csromanheader{2}{\textbf{ปจฺจยจตุกฺก-เหตุ}}
\proseitem{74}{เหตุ สุขาย เวทนาย สมฺปยุตฺโต ธมฺโม เหตุสฺส สุขาย เวทนาย สมฺปยุตฺตสฺส ธมฺมสฺส เหตุปจฺจเยน ปจฺจโย.\csromanspacer{} ตีณิ.}

\csromanheader{2}{ธมฺมานุโลม ทุกติกปฏฺฐานปาฬิ \textbf{อารมฺมณาทิ}}
\proseitem{75}{\csromanfolio{210}เหตุ สุขาย เวทนาย สมฺปยุตฺโต ธมฺโม เหตุสฺส สุขาย เวทนาย สมฺปยุตฺตสฺส ธมฺมสฺส อารมฺมณปจฺจเยน ปจฺจโย.\csromanspacer{} ตีณิ.}

\prose{นเหตุ สุขาย เวทนาย สมฺปยุตฺโต ธมฺโม นเหตุสฺส สุขาย เวทนาย สมฺปยุตฺตสฺส ธมฺมสฺส อารมฺมณปจฺจเยน ปจฺจโย.\csromanspacer{} ตีณิ.}

\prose{เหตุ สุขาย เวทนาย สมฺปยุตฺโต จ นเหตุ สุขาย เวทนาย สมฺปยุตฺโต จ ธมฺมา เหตุสฺส สุขาย เวทนาย สมฺปยุตฺตสฺส ธมฺมสฺส อารมฺมณปจฺจเยน ปจฺจโย.\csromanspacer{} ตีณิ.}

\proseitem{76}{เหตุ สุขาย เวทนาย สมฺปยุตฺโต ธมฺโม เหตุสฺส สุขาย เวทนาย สมฺปยุตฺตสฺส ธมฺมสฺส อธิปติปจฺจเยน ปจฺจโย.\csromanspacer{} อารมฺมณาธิปติ, สหชาตาธิปติ.\csromanspacer{} ตีณิ.}

\prose{นเหตุ สุขาย เวทนาย สมฺปยุตฺโต ธมฺโม นเหตุสฺส สุขาย เวทนาย สมฺปยุตฺตสฺส ธมฺมสฺส อธิปติปจฺจเยน ปจฺจโย.\csromanspacer{} อารมฺมณาธิปติ, สหชาตาธิปติ.\csromanspacer{} ตีณิ.}

\prose{เหตุ สุขาย เวทนาย สมฺปยุตฺโต จ นเหตุ สุขาย เวทนาย สมฺปยุตฺโต จ ธมฺมา เหตุสฺส สุขาย เวทนาย สมฺปยุตฺตสฺส ธมฺมสฺส อธิปติปจฺจเยน ปจฺจโย.\csromanspacer{} อารมฺมณาธิปติ.\csromanspacer{} ตีณิ ฯเปฯ.}

\csromanheader{2}{\textbf{อุปนิสฺสยาทิ}}
\proseitem{77}{เหตุ สุขาย เวทนาย สมฺปยุตฺโต ธมฺโม เหตุสฺส สุขาย เวทนาย สมฺปยุตฺตสฺส ธมฺมสฺส อุปนิสฺสยปจฺจเยน ปจฺจโย.\csromanspacer{} อารมฺมณูปนิสฺสโย, อนนฺตรูปนิสฺสโย, ปกตูปนิสฺสโย.\csromanspacer{} นว. (สํขิตฺตํ.)}

\proseitem{78}{นเหตุ สุขาย เวทนาย สมฺปยุตฺโต ธมฺโม นเหตุสฺส สุขาย เวทนาย สมฺปยุตฺตสฺส ธมฺมสฺส กมฺมปจฺจเยน ปจฺจโย.\csromanspacer{}\csromanspacer{}นเหตุ สุขาย เวทนาย สมฺปยุตฺโต ธมฺโม เหตุสฺส สุขาย เวทนาย สมฺปยุตฺตสฺส ธมฺมสฺส กมฺมปจฺจเยน ปจฺจโย.\csromanspacer{}\csromanspacer{}นเหตุ สุขาย เวทนาย สมฺปยุตฺโต ธมฺโม เหตุสฺส สุขาย เวทนาย สมฺปยุตฺตสฺส จ}

\vspace{\tipitakaparskip}
\csromancenter{เหตุทุก 2.\csromanspacer{} เวทนาตฺติก นเหตุสฺส สุขาย เวทนาย สมฺปยุตฺตสฺส จ ธมฺมสฺส กมฺมปจฺจเยน ปจฺจโย. (3).\csromanspacer{} วิปากปจฺจเยน ปจฺจโย.}
\proseitem{79}{\csromanfolio{211}นเหตุ สุขาย เวทนาย สมฺปยุตฺโต ธมฺโม นเหตุสฺส สุขาย เวทนาย สมฺปยุตฺตสฺส ธมฺมสฺส อาหารปจฺจเยน ปจฺจโย.\csromanspacer{} ตีณิ ฯเปฯ.\csromanspacer{} เหตุ สุขาย เวทนาย สมฺปยุตฺโต ธมฺโม เหตุสฺส สุขาย เวทนาย สมฺปยุตฺตสฺส ธมฺมสฺส อวิคตปจฺจเยน ปจฺจโย.}

\proseitem{80}{เหตุยา ตีณิ, อารมฺมเณ นว, อธิปติยา นว, อนนฺตเร นว, สมนนฺตเร นว, สหชาเต นว, อญฺญมญฺเญ นว, นิสฺสเย นว, อุปนิสฺสเย นว, อาเสวเน นว, กมฺเม ตีณิ, วิปาเก นว, อาหาเร ตีณิ, อินฺทฺริเย นว, ฌาเน ตีณิ, มคฺเค นว, สมฺปยุตฺเต นว, อตฺถิยา นว, นตฺถิยา นว, วิคเต นว, อวิคเต นว.}

\csromanheader{2}{\textbf{ปจฺจนียุทฺธาร}}
\prosecont{เหตุ สุขาย เวทนาย สมฺปยุตฺโต ธมฺโม เหตุสฺส สุขาย เวทนาย สมฺปยุตฺตสฺส ธมฺมสฺส อารมฺมณปจฺจเยน ปจฺจโย.\csromanspacer{} สหชาตปจฺจเยน ปจฺจโย.\csromanspacer{} อุปนิสฺสยปจฺจเยน ปจฺจโย. (สํขิตฺตํ.)}

\proseitem{82}{นเหตุยา นว, น-อารมฺมเณ นว ฯเปฯ โน-อวิคเต นว. (ปจฺจนียํ.)}

\prose{เหตุปจฺจยา น-อารมฺมเณ ตีณิ. (สํขิตฺตํ.\csromanspacer{} อนุโลมปจฺจนียํ.)}

\csromanheader{2}{นเหตุปจฺจยา อารมฺมเณ นว. (สํขิตฺตํ. ปจฺจนียานุโลมํ.)}
\prose{(ยถา กุสลตฺติเก ปญฺหาวารสฺส อนุโลมมฺปิ ปจฺจนียมฺปิ อนุโลมปจฺจนียมฺปิ ปจฺจนียานุโลมมฺปิ คณิตํ, เอวํ คเณตพฺพํ.)}

\csromanheader{2}{2. ทุกฺขายเวทนายสมฺปยุตฺตปท 1. ปฏิจฺจวาร}
\csromanheader{2}{\textbf{ปจฺจยจตุกฺก-เหตุ}}
\proseitem{83}{เหตุํ ทุกฺขาย เวทนาย สมฺปยุตฺตํ ธมฺมํ ปฏิจฺจ เหตุ ทุกฺขาย เวทนาย สมฺปยุตฺโต ธมฺโม อุปฺปชฺชติ เหตุปจฺจยา.\csromanspacer{}\csromanspacer{}เหตุํ ทุกฺขาย}

\vspace{\tipitakaparskip}
\csromancenter{ธมฺมานุโลม ทุกติกปฏฺฐานปาฬิ เวทนาย สมฺปยุตฺตํ ธมฺมํ ปฏิจฺจ นเหตุ ทุกฺขาย เวทนาย สมฺปยุตฺโต ธมฺโม อุปฺปชฺชติ เหตุปจฺจยา.\csromanspacer{}\csromanspacer{}เหตุํ ทุกฺขาย เวทนาย สมฺปยุตฺตํ ธมฺมํ ปฏิจฺจ เหตุ ทุกฺขาย เวทนาย สมฺปยุตฺโต จ นเหตุ ทุกฺขาย เวทนาย สมฺปยุตฺโต จ ธมฺมา อุปฺปชฺชนฺติ เหตุปจฺจยา. (3)}
\prose{\csromanfolio{212}นเหตุํ ทุกฺขาย เวทนาย สมฺปยุตฺตํ ธมฺมํ ปฏิจฺจ นเหตุ ทุกฺขาย เวทนาย สมฺปยุตฺโต ธมฺโม อุปฺปชฺชติ เหตุปจฺจยา.\csromanspacer{} ตีณิ.}

\prose{เหตุํ ทุกฺขาย เวทนาย สมฺปยุตฺตญฺจ นเหตุํ ทุกฺขาย เวทนาย สมฺปยุตฺตญฺจ ธมฺมํ ปฏิจฺจ เหตุ ทุกฺขาย เวทนาย สมฺปยุตฺโต ธมฺโม อุปฺปชฺชติ เหตุปจฺจยา.\csromanspacer{} ตีณิ. (สํขิตฺตํ.)}

\proseitem{84}{เหตุยา นว, อารมฺมเณ นว, อธิปติยา นว, อนนฺตเร นว, สมนนฺตเร นว, สหชาเต นว, อญฺญมญฺเญ นว, นิสฺสเย นว, อุปนิสฺสเย นว, ปุเรชาเต นว, อาเสวเน นว, กมฺเม นว, วิปาเก เอกํ, อาหาเร นว, อินฺทฺริเย นว, ฌาเน นว, มคฺเค นว, สมฺปยุตฺเต นว, วิปฺปยุตฺเต นว, อตฺถิยา นว, นตฺถิยา นว, วิคเต นว, อวิคเต นว. (อนุโลมํ.)}

\csromanheader{2}{\textbf{นเหตุ น-อธิปติ}}
\proseitem{85}{นเหตุํ ทุกฺขาย เวทนาย สมฺปยุตฺตํ ธมฺมํ ปฏิจฺจ นเหตุ ทุกฺขาย เวทนาย สมฺปยุตฺโต ธมฺโม อุปฺปชฺชติ นเหตุปจฺจยา. (1)}

\proseitem{86}{เหตุํ ทุกฺขาย เวทนาย สมฺปยุตฺตํ ธมฺมํ ปฏิจฺจ เหตุ ทุกฺขาย เวทนาย สมฺปยุตฺโต ธมฺโม อุปฺปชฺชติ น-อธิปติปจฺจยา. (สํขิตฺตํ.)}

\proseitem{87}{นเหตุยา เอกํ, น-อธิปติยา นว, นปจฺฉาชาเต นว, น-อาเสวเน นว, นกมฺเม ตีณิ, นวิปาเก นว, นฌาเน เอกํ, นมคฺเค เอกํ. (ปจฺจนียํ.)}

\prose{เหตุปจฺจยา น-อธิปติยา นว. (สํขิตฺตํ.\csromanspacer{} อนุโลมปจฺจนียํ.)}

\csromanheader{2}{นเหตุปจฺจยา อารมฺมเณ เอกํ. (สํขิตฺตํ. ปจฺจนียานุโลมํ)}
\prose{(สหชาตวาโรปิ ปจฺจยวาโรปิ นิสฺสยวาโรปิ สํสฏฺฐวาโรปิ สมฺปยุตฺตวาโรปิ ปฏิจฺจวารสทิสา วิตฺถาเรตพฺพา.)}

\vspace{\tipitakaparskip}
\csromancenter{เหตุทุก 2.\csromanspacer{} เวทนาตฺติก \textbf{2.}\csromanspacer{}\textbf{ ทุกฺขายเวทนายสมฺปยุตฺตปท 7.}\csromanspacer{}\textbf{ ปญฺหาวาร}}
\csromanheader{2}{\textbf{ปจฺจยจตุกฺก-เหตุ}}
\proseitem{88}{\csromanfolio{213}เหตุ ทุกฺขาย เวทนาย สมฺปยุตฺโต ธมฺโม เหตุสฺส ทุกฺขาย เวทนาย สมฺปยุตฺตสฺส ธมฺมสฺส เหตุปจฺจเยน ปจฺจโย.\csromanspacer{}\csromanspacer{}เหตุ ทุกฺขาย เวทนาย สมฺปยุตฺโต ธมฺโม นเหตุสฺส ทุกฺขาย เวทนาย สมฺปยุตฺตสฺส ธมฺมสฺส เหตุปจฺจเยน ปจฺจโย.\csromanspacer{}\csromanspacer{}เหตุ ทุกฺขาย เวทนาย สมฺปยุตฺโต ธมฺโม เหตุสฺส ทุกฺขาย เวทนาย สมฺปยุตฺตสฺส จ นเหตุสฺส ทุกฺขาย เวทนาย สมฺปยุตฺตสฺส จ ธมฺมสฺส เหตุปจฺจเยน ปจฺจโย. (3)}

\csromanheader{2}{\textbf{อารมฺมณาทิ}}
\proseitem{89}{เหตุ ทุกฺขาย เวทนาย สมฺปยุตฺโต ธมฺโม เหตุสฺส ทุกฺขาย เวทนาย สมฺปยุตฺตสฺส ธมฺมสฺส อารมฺมณปจฺจเยน ปจฺจโย.\csromanspacer{} ตีณิ.}

\prose{นเหตุ ทุกฺขาย เวทนาย สมฺปยุตฺโต ธมฺโม นเหตุสฺส ทุกฺขาย เวทนาย สมฺปยุตฺตสฺส ธมฺมสฺส อารมฺมณปจฺจเยน ปจฺจโย.\csromanspacer{} ตีณิ.}

\prose{เหตุ ทุกฺขาย เวทนาย สมฺปยุตฺโต จ นเหตุ ทุกฺขาย เวทนาย สมฺปยุตฺโต จ ธมฺมา เหตุสฺส ทุกฺขาย เวทนาย สมฺปยุตฺตสฺส ธมฺมสฺส อารมฺมณปจฺจเยน ปจฺจโย.\csromanspacer{} ตีณิ. (สํขิตฺตํ.)}

\proseitem{90}{นเหตุ ทุกฺขาย เวทนาย สมฺปยุตฺโต ธมฺโม นเหตุสฺส ทุกฺขาย เวทนาย สมฺปยุตฺตสฺส ธมฺมสฺส กมฺมปจฺจเยน ปจฺจโย.\csromanspacer{} ตีณิ. (สํขิตฺตํ.)}

\proseitem{91}{เหตุยา ตีณิ, อารมฺมเณ นว, อธิปติยา ตีณิ, อนนฺตเร นว, สมนนฺตเร นว, สหชาเต นว, อญฺญมญฺเญ นว, นิสฺสเย นว, อุปนิสฺสเย นว, อาเสวเน นว, กมฺเม ตีณิ, วิปาเก เอกํ, อาหาเร ตีณิ, อินฺทฺริเย ตีณิ, ฌาเน ตีณิ, มคฺเค ตีณิ, สมฺปยุตฺเต นว, อตฺถิยา นว, นตฺถิยา นว, วิคเต นว, อวิคเต นว. (อนุโลมํ.)}

\csromanheader{2}{ธมฺมานุโลม ทุกติกปฏฺฐานปาฬิ \textbf{ปจฺจนียุทฺธาร}}
\proseitem{92}{\csromanfolio{214}เหตุ ทุกฺขาย เวทนาย สมฺปยุตฺโต ธมฺโม เหตุสฺส ทุกฺขาย เวทนาย สมฺปยุตฺตสฺส ธมฺมสฺส อารมฺมณปจฺจเยน ปจฺจโย.\csromanspacer{} สหชาตปจฺจเยน ปจฺจโย.\csromanspacer{} อุปนิสฺสยปจฺจเยน ปจฺจโย. (สํขิตฺตํ.)}

\proseitem{93}{นเหตุยา นว, น-อารมฺมเณ นว, น-อธิปติยา นว. (สํขิตฺตํ.\csromanspacer{} ปจฺจนียํ.)}

\prose{เหตุปจฺจยา น-อารมฺมเณ ตีณิ. (สํขิตฺตํ.\csromanspacer{} อนุโลมปจฺจนียํ.)}

\csromanheader{2}{นเหตุปจฺจยา อารมฺมเณ นว. (สํขิตฺตํ. ปจฺจนียานุโลมํ.)}
\prose{(ยถา กุสลตฺติเก ปญฺหาวารสฺส อนุโลมมฺปิ ปจฺจนียมฺปิ อนุโลมปจฺจนียมฺปิ ปจฺจนียานุโลมมฺปิ คณิตํ, เอวํ คเณตพฺพํ.)}

\csromanheader{2}{3. อทุกฺขมสุขายเวทนายสมฺปยุตฺตปท 1. ปฏิจฺจวาร}
\csromanheader{2}{\textbf{ปจฺจยจตุกฺก-เหตุ}}
\proseitem{94}{เหตุํ อทุกฺขมสุขาย เวทนาย สมฺปยุตฺตํ ธมฺมํ ปฏิจฺจ เหตุ อทุกฺขมสุขาย เวทนาย สมฺปยุตฺโต ธมฺโม อุปฺปชฺชติ เหตุปจฺจยา.\csromanspacer{} ตีณิ.}

\prose{นเหตุํ อทุกฺขมสุขาย เวทนาย สมฺปยุตฺตํ ธมฺมํ ปฏิจฺจ นเหตุ อทุกฺขมสุขาย เวทนาย สมฺปยุตฺโต ธมฺโม อุปฺปชฺชติ เหตุปจฺจยา.\csromanspacer{} ตีณิ.}

\prose{เหตุํ อทุกฺขมสุขาย เวทนาย สมฺปยุตฺตญฺจ นเหตุํ อทุกฺขมสุขาย เวทนาย สมฺปยุตฺตญฺจ ธมฺมํ ปฏิจฺจ เหตุ อทุกฺขมสุขาย เวทนาย สมฺปยุตฺโต ธมฺโม อุปฺปชฺชติ เหตุปจฺจยา.\csromanspacer{} ตีณิ. (สํขิตฺตํ.)}

\proseitem{95}{เหตุยา นว, อารมฺมเณ นว, อธิปติยา นว, อนนฺตเร นว, สมนนฺตเร นว, สหชาเต นว, อญฺญมญฺเญ นว, นิสฺสเย นว, อุปนิสฺสเย นว, ปุเรชาเต นว, อาเสวเน นว, กมฺเม นว, วิปาเก นว, อาหาเร}

\vspace{\tipitakaparskip}
\csromancenter{เหตุทุก 2.\csromanspacer{} เวทนาตฺติก นว, อินฺทฺริเย นว, ฌาเน นว, มคฺเค นว, สมฺปยุตฺเต นว, วิปฺปยุตฺเต นว, อตฺถิยา นว, นตฺถิยา นว, วิคเต นว, อวิคเต นว. (อนุโลมํ.)}
\csromanheader{2}{\textbf{นเหตุ น-อธิปติ}}
\proseitem{96}{\csromanfolio{215}นเหตุํ อทุกฺขมสุขาย เวทนาย สมฺปยุตฺตํ ธมฺมํ ปฏิจฺจ นเหตุ อทุกฺขมสุขาย เวทนาย สมฺปยุตฺโต ธมฺโม อุปฺปชฺชติ นเหตุปจฺจยา.\csromanspacer{} เทฺว.}

\proseitem{97}{เหตุํ อทุกฺขมสุขาย เวทนาย สมฺปยุตฺตํ ธมฺมํ ปฏิจฺจ เหตุ อทุกฺขมสุขาย เวทนาย สมฺปยุตฺโต ธมฺโม อุปฺปชฺชติ น-อธิปติปจฺจยา. (สํขิตฺตํ.)}

\proseitem{98}{นเหตุยา เทฺว, น-อธิปติยา นว, นปุเรชาเต นว, นปจฺฉาชาเต นว, น-อาเสวเน นว, นกมฺเม ตีณิ, นวิปาเก นว, นฌาเน เอกํ, นมคฺเค เอกํ, นวิปฺปยุตฺเต นว. (สํขิตฺตํ.\csromanspacer{} ปจฺจนียํ.)}

\prose{เหตุปจฺจยา น-อธิปติยา นว. (สํขิตฺตํ.\csromanspacer{} อนุโลมปจฺจนียํ.)}

\csromanheader{2}{นเหตุปจฺจยา อารมฺมเณ เทฺว. (สํขิตฺตํ. ปจฺจนียานุโลมํ.)}
\prose{(สหชาตวาโรปิ ปจฺจยวาโรปิ นิสฺสยวาโรปิ สํสฏฺฐวาโรปิ สมฺปยุตฺตวาโรปิ ปฏิจฺจวารสทิสา วิตฺถาเรตพฺพา.)}

\csromanheader{2}{3. อทุกฺขมสุขายเวทนายสมฺปยุตฺตปท 7. ปญฺหาวาร}
\csromanheader{2}{\textbf{ปจฺจยจตุกฺก-เหตฺวาทิ}}
\proseitem{99}{เหตุ อทุกฺขมสุขาย เวทนาย สมฺปยุตฺโต ธมฺโม เหตุสฺส อทุกฺขมสุขาย เวทนาย สมฺปยุตฺตสฺส ธมฺมสฺส เหตุปจฺจเยน ปจฺจโย.\csromanspacer{} ตีณิ.}

\proseitem{100}{เหตุ อทุกฺขมสุขาย เวทนาย สมฺปยุตฺโต ธมฺโม เหตุสฺส อทุกฺขมสุขาย เวทนาย สมฺปยุตฺตสฺส ธมฺมสฺส อารมฺมณปจฺจเยน ปจฺจโย.\csromanspacer{} นว.}

\gambhira{ธมฺมานุโลม ทุกติกปฏฺฐานปาฬิ}
\proseitem{101}{\csromanfolio{216}เหตุ อทุกฺขมสุขาย เวทนาย สมฺปยุตฺโต ธมฺโม เหตุสฺส อทุกฺขมสุขาย เวทนาย สมฺปยุตฺตสฺส ธมฺมสฺส อธิปติปจฺจเยน ปจฺจโย.\csromanspacer{} อารมฺมณาธิปติ, สหชาตาธิปติ.\csromanspacer{} ตีณิ.}

\prose{นเหตุ อทุกฺขมสุขาย เวทนาย สมฺปยุตฺโต ธมฺโม นเหตุสฺส อทุกฺขมสุขาย เวทนาย สมฺปยุตฺตสฺส ธมฺมสฺส อธิปติปจฺจเยน ปจฺจโย.\csromanspacer{} อารมฺมณาธิปติ, สหชาตาธิปติ.\csromanspacer{} ตีณิ.}

\prose{เหตุ อทุกฺขมสุขาย เวทนาย สมฺปยุตฺโต จ นเหตุ อทุกฺขมสุขาย เวทนาย สมฺปยุตฺโต จ ธมฺมา เหตุสฺส อทุกฺขมสุขาย เวทนาย สมฺปยุตฺตสฺส ธมฺมสฺส อธิปติปจฺจเยน ปจฺจโย.\csromanspacer{} อารมฺมณาธิปติ.\csromanspacer{} ตีณิ. (อารมฺมณาธิปติเยว.) ฯเปฯ.}

\csromanheader{2}{\textbf{อุปนิสฺสยาทิ}}
\proseitem{102}{เหตุ อทุกฺขมสุขาย เวทนาย สมฺปยุตฺโต ธมฺโม เหตุสฺส อทุกฺขมสุขาย เวทนาย สมฺปยุตฺตสฺส ธมฺมสฺส อุปนิสฺสยปจฺจเยน ปจฺจโย.\csromanspacer{} อารมฺมณูปนิสฺสโย, อนนฺตรูปนิสฺสโย, ปกตูปนิสฺสโย.\csromanspacer{} นว.\csromanspacer{} อาเสวนปจฺจเยน ปจฺจโย.\csromanspacer{} นว.}

\proseitem{103}{นเหตุ อทุกฺขมสุขาย เวทนาย สมฺปยุตฺโต ธมฺโม นเหตุสฺส อทุกฺขมสุขาย เวทนาย สมฺปยุตฺตสฺส ธมฺมสฺส กมฺมปจฺจเยน ปจฺจโย.\csromanspacer{} ตีณิ.}

\proseitem{104}{เหตุ อทุกฺขมสุขาย เวทนาย สมฺปยุตฺโต ธมฺโม เหตุสฺส อทุกฺขมสุขาย เวทนาย สมฺปยุตฺตสฺส ธมฺมสฺส วิปากปจฺจเยน ปจฺจโย.\csromanspacer{} นว ฯเปฯ.\csromanspacer{} อวิคตปจฺจเยน ปจฺจโย.\csromanspacer{} นว.}

\proseitem{105}{เหตุยา ตีณิ, อารมฺมเณ นว, อธิปติยา นว, อนนฺตเร นว, สมนนฺตเร นว, สหชาเต นว, อญฺญมญฺเญ นว, นิสฺสเย นว, อุปนิสฺสเย นว, อาเสวเน นว, กมฺเม ตีณิ, วิปาเก นว, อาหาเร ตีณิ, อินฺทฺริเย นว, ฌาเน ตีณิ, มคฺเค นว, สมฺปยุตฺเต นว, อตฺถิยา นว, นตฺถิยา นว, วิคเต นว, อวิคเต นว. (อนุโลมํ.)}

\csromanheader{2}{1. เหตุทุก 3. วิปากตฺติก \textbf{ปจฺจนียุทฺธาร}}
\proseitem{106}{\csromanfolio{217}เหตุ อทุกฺขมสุขาย เวทนาย สมฺปยุตฺโต ธมฺโม เหตุสฺส อทุกฺขมสุขาย เวทนาย สมฺปยุตฺตสฺส ธมฺมสฺส อารมฺมณปจฺจเยน ปจฺจโย.\csromanspacer{} สหชาตปจฺจเยน ปจฺจโย.\csromanspacer{} อุปนิสฺสยปจฺจเยน ปจฺจโย. (สํขิตฺตํ.)}

\proseitem{107}{นเหตุยา นว, น-อารมฺมเณ นว. (สํขิตฺตํ.\csromanspacer{} ปจฺจนียํ.)}

\prose{เหตุปจฺจยา น-อารมฺมเณ ตีณิ. (สํขิตฺตํ.\csromanspacer{} อนุโลมปจฺจนียํ.)}

\csromanheader{2}{นเหตุปจฺจยา อารมฺมเณ นว. (สํขิตฺตํ. ปจฺจนียานุโลมํ.)}
\prose{(ยถา กุสลตฺติเก ปญฺหาวารสฺส อนุโลมมฺปิ ปจฺจนียมฺปิ อนุโลมปจฺจนียมฺปิ ปจฺจนียานุโลมมฺปิ คณิตํ, เอวํ คเณตพฺพํ.)}

\nitthitam{เหตุทุกเวทนาตฺติกํ นิฏฺฐิตํ.}
\csromanmatikaanchor{mtk.34}
\csromantocmarkref{section}{3. วิปากตฺติก}{mtk.34}
\bookmark[dest={mtk.34},level=1]{3. วิปากตฺติก}
\csromanheader{1}{1. เหตุทุก 3. วิปากตฺติก}
\csromanheader{2}{1. วิปากปท 1. ปฏิจฺจวาร}
\csromanheader{2}{\textbf{ปจฺจยจตุกฺก-เหตุ}}
\proseitem{108}{เหตุํ วิปากํ ธมฺมํ ปฏิจฺจ เหตุ วิปาโก ธมฺโม อุปฺปชฺชติ เหตุปจฺจยา.\csromanspacer{}\csromanspacer{}เหตุํ วิปากํ ธมฺมํ ปฏิจฺจ นเหตุ วิปาโก ธมฺโม อุปฺปชฺชติ เหตุปจฺจยา.\csromanspacer{}\csromanspacer{}เหตุํ วิปากํ ธมฺมํ ปฏิจฺจ เหตุ วิปาโก จ นเหตุ วิปาโก จ ธมฺมา อุปฺปชฺชนฺติ เหตุปจฺจยา. (3)}

\prose{นเหตุํ วิปากํ ธมฺมํ ปฏิจฺจ นเหตุ วิปาโก ธมฺโม อุปฺปชฺชติ เหตุปจฺจยา.\csromanspacer{} ตีณิ.}

\prose{เหตุํ วิปากญฺจ นเหตุํ วิปากญฺจ ธมฺมํ ปฏิจฺจ เหตุ วิปาโก ธมฺโม อุปฺปชฺชติ เหตุปจฺจยา.\csromanspacer{} ตีณิ. (สํขิตฺตํ.)}

\proseitem{109}{เหตุยา นว, อารมฺมเณ นว, อธิปติยา นว, อนนฺตเร นว, สมนนฺตเร นว, สหชาเต นว, อญฺญมญฺเญ นว, นิสฺสเย นว, อุปนิสฺสเย}

\vspace{\tipitakaparskip}
\csromancenter{ธมฺมานุโลม ทุกติกปฏฺฐานปาฬิ นว, ปุเรชาเต นว, กมฺเม นว, วิปาเก นว, อาหาเร นว, อินฺทฺริเย นว, ฌาเน นว, มคฺเค นว, สมฺปยุตฺเต นว, วิปฺปยุตฺเต นว, อตฺถิยา นว, นตฺถิยา นว, วิคเต นว, อวิคเต นว. (อนุโลมํ.)}
\csromanheader{2}{\textbf{นเหตุ น-อธิปติ}}
\proseitem{110}{\csromanfolio{218}นเหตุํ วิปากํ ธมฺมํ ปฏิจฺจ นเหตุ วิปาโก ธมฺโม อุปฺปชฺชติ นเหตุปจฺจยา. (1)}

\proseitem{111}{เหตุํ วิปากํ ธมฺมํ ปฏิจฺจ เหตุ วิปาโก ธมฺโม อุปฺปชฺชติ น-อธิปติปจฺจยา. (สํขิตฺตํ.)}

\proseitem{112}{นเหตุยา เอกํ, น-อธิปติยา นว, นปุเรชาเต นว, นปจฺฉาชาเต นว, น-อาเสวเน นว, นฌาเน เอกํ, นมคฺเค เอกํ, นวิปฺปยุตฺเต นว. (สํขิตฺตํ.\csromanspacer{} ปจฺจนียํ.)}

\prose{เหตุปจฺจยา น-อธิปติยา นว. (สํขิตฺตํ.\csromanspacer{} อนุโลมฺมปจฺจนียํ.)}

\csromanheader{2}{นเหตุปจฺจยา อารมฺมเณ เอกํ. (สํขิตฺตํ. ปจฺจนียานุโลมํ.)}
\prose{(สหชาตวาโรปิ ปจฺจยวาโรปิ นิสฺสยวาโรปิ สํสฏฺฐวาโรปิ สมฺปยุตฺตวาโรปิ ปฏิจฺจวารสทิสา วิตฺถาเรตพฺพา.)}

\csromanheader{2}{1. วิปากปท 7. ปญฺหาวาร}
\csromanheader{2}{\textbf{ปจฺจยจตุกฺก-เหตฺวาทิ}}
\proseitem{113}{เหตุ วิปาโก ธมฺโม เหตุสฺส วิปากสฺส ธมฺมสฺส เหตุปจฺจเยน ปจฺจโย.\csromanspacer{} ตีณิ.}

\proseitem{114}{เหตุ วิปาโก ธมฺโม เหตุสฺส วิปากสฺส ธมฺมสฺส อารมฺมณปจฺจเยน ปจฺจโย.\csromanspacer{} ตีณิ.}

\prose{นเหตุ วิปาโก ธมฺโม นเหตุสฺส วิปากสฺส ธมฺมสฺส อารมฺมณปจฺจเยน ปจฺจโย.\csromanspacer{} ตีณิ.}

\csromanheader{2}{1. เหตุทุก 3. วิปากตฺติก}
\prose{\csromanfolio{219}เหตุ วิปาโก จ นเหตุ วิปาโก จ ธมฺมา เหตุสฺส วิปากสฺส ธมฺมสฺส อารมฺมณปจฺจเยน ปจฺจโย.\csromanspacer{} ตีณิ. (ตทารมฺมณาเยว ลพฺภนฺติ.)}

\proseitem{115}{เหตุ วิปาโก ธมฺโม เหตุสฺส วิปากสฺส ธมฺมสฺส อธิปติปจฺจเยน ปจฺจโย.\csromanspacer{} ตีณิ. (สหชาตาธิปติเยว ลพฺภติ.\csromanspacer{} อารมฺมณาธิปติ นตฺถิ.) ฯเปฯ.}

\csromanheader{2}{\textbf{อุปนิสฺสยาทิ}}
\proseitem{116}{เหตุ วิปาโก ธมฺโม เหตุสฺส วิปากสฺส ธมฺมสฺส อุปนิสฺสยปจฺจเยน ปจฺจโย.\csromanspacer{} อนนฺตรูปนิสฺสโย, ปกตูปนิสฺสโย.\csromanspacer{}\csromanspacer{}เหตุ วิปาโก ธมฺโม นเหตุสฺส วิปากสฺส ธมฺมสฺส อุปนิสฺสยปจฺจเยน ปจฺจโย.\csromanspacer{} อนนฺตรูปนิสฺสโย, ปกตูปนิสฺสโย.\csromanspacer{}\csromanspacer{}เหตุ วิปาโก ธมฺโม เหตุสฺส วิปากสฺส จ นเหตุสฺส วิปากสฺส จ ธมฺมสฺส อุปนิสฺสยปจฺจเยน ปจฺจโย.\csromanspacer{} อนนฺตรูปนิสฺสโย, ปกตูปนิสฺสโย. (3)}

\prose{นเหตุ วิปาโก ธมฺโม นเหตุสฺส วิปากสฺส ธมฺมสฺส อุปนิสฺสยปจฺจเยน ปจฺจโย.\csromanspacer{} อนนฺตรูปนิสฺสโย, ปกตูปนิสฺสโย ฯเปฯ. (อิตเร เทฺว อนนฺตรูปนิสฺสโย ปกตูปนิสฺสโยเยว.)}

\proseitem{117}{นเหตุ วิปาโก ธมฺโม นเหตุสฺส วิปากสฺส ธมฺมสฺส กมฺมปจฺจเยน ปจฺจโย.\csromanspacer{} ตีณิ. (สหชาตกมฺมเมว.\csromanspacer{} สํขิตฺตํ.)}

\prose{เหตุ วิปาโก ธมฺโม เหตุสฺส วิปากสฺส ธมฺมสฺส วิปากปจฺจเยน ปจฺจโย.\csromanspacer{} นว.}

\prose{นเหตุ วิปาโก ธมฺโม นเหตุสฺส วิปากสฺส ธมฺมสฺส อาหารปจฺจเยน ปจฺจโย. (สํขิตฺตํ.)}

\proseitem{118}{เหตุยา ตีณิ, อารมฺมเณ นว, อธิปติยา ฉ, อนนฺตเร นว, สมนนฺตเร นว, สหชาเต นว, อญฺญมญฺเญ นว, นิสฺสเย นว, อุปนิสฺสเย นว, กมฺเม ตีณิ, วิปาเก นว, อาหาเร ตีณิ, อินฺทฺริเย นว, ฌาเน ตีณิ, มคฺเค นว, สมฺปยุตฺเต นว, อตฺถิยา นว, นตฺถิยา นว, วิคเต นว, อวิคเต นว. (อนุโลมํ.)}

\csromanheader{2}{ธมฺมานุโลม ทุกติกปฏฺฐานปาฬิ \textbf{ปจฺจนียุทฺธาร}}
\proseitem{119}{\csromanfolio{220}เหตุ วิปาโก ธมฺโม เหตุสฺส วิปากสฺส ธมฺมสฺส อารมฺมณปจฺจเยน ปจฺจโย.\csromanspacer{} สหชาตปจฺจเยน ปจฺจโย.\csromanspacer{} อุปนิสฺสยปจฺจเยน ปจฺจโย. (สํขิตฺตํ.)}

\proseitem{120}{นเหตุยา นว, น-อารมฺมเณ นว. (สํขิตฺตํ.\csromanspacer{} ปจฺจนียํ.)}

\prose{เหตุปจฺจยา น-อารมฺมเณ ตีณิ. (สํขิตฺตํ.\csromanspacer{} อนุโลมปจฺจนียํ.)}

\csromanheader{2}{นเหตุปจฺจยา อารมฺมเณ นว. (สํขิตฺตํ. ปจฺจนียานุโลมํ.)}
\prose{(ยถา กุสลตฺติเก ปญฺหาวารสฺส อนุโลมมฺปิ ปจฺจนียมฺปิ อนุโลมปจฺจนียมฺปิ ปจฺจนียานุโลมมฺปิ คณิตํ, เอวํ คเณตพฺพํ.)}

\csromanheader{2}{2. วิปากธมฺมปท 1. ปฏิจฺจวาร}
\csromanheader{2}{\textbf{ปจฺจยจตุกฺก-เหตุ}}
\proseitem{121}{เหตุํ วิปากธมฺมธมฺมํ ปฏิจฺจ เหตุ วิปากธมฺมธมฺโม อุปฺปชฺชติ เหตุปจฺจยา.\csromanspacer{}\csromanspacer{}เหตุํ วิปากธมฺมธมฺมํ ปฏิจฺจ นเหตุ วิปากธมฺมธมฺโม อุปฺปชฺชติ เหตุปจฺจยา.\csromanspacer{}\csromanspacer{}เหตุํ วิปากธมฺมธมฺมํ ปฏิจฺจ เหตุ วิปากธมฺมธมฺโม จ นเหตุ วิปากธมฺมธมฺโม จ ธมฺมา อุปฺปชฺชนฺติ เหตุปจฺจยา. (3)}

\prose{นเหตุํ วิปากธมฺมธมฺมํ ปฏิจฺจ นเหตุ วิปากธมฺมธมฺโม อุปฺปชฺชติ เหตุปจฺจยา.\csromanspacer{} ตีณิ.}

\prose{เหตุํ วิปากธมฺมธมฺมญฺจ นเหตุํ วิปากธมฺมธมฺมญฺจ ปฏิจฺจ เหตุ วิปากธมฺมธมฺโม อุปฺปชฺชติ เหตุปจฺจยา.\csromanspacer{} ตีณิ. (สํขิตฺตํ.)}

\proseitem{122}{เหตุยา นว, อารมฺมเณ นว, อธิปติยา นว ฯเปฯ กมฺเม นว, อาหาเร นว, อินฺทฺริเย นว, ฌาเน นว, มคฺเค นว, สมฺปยุตฺเต นว, วิปฺปยุตฺเต นว, อตฺถิยา นว, นตฺถิยา นว, วิคเต นว, อวิคเต นว.}

\csromanheader{2}{\textbf{นเหตุ น-อธิปติ}}
\proseitem{123}{นเหตุํ วิปากธมฺมธมฺมํ ปฏิจฺจ เหตุ วิปากธมฺมธมฺโม อุปฺปชฺชติ นเหตุปจฺจยา. (1)}

\csromanheader{2}{1. เหตุทุก 3. วิปากตฺติก}
\proseitem{124}{\csromanfolio{221}เหตุํ วิปากธมฺมธมฺมํ ปฏิจฺจ เหตุ วิปากธมฺมธมฺโม อุปฺปชฺชติ น-อธิปติปจฺจยา.\csromanspacer{} นว.}

\csromanheader{2}{\textbf{นปุเรชาตาทิ}}
\proseitem{125}{เหตุํ วิปากธมฺมธมฺมํ ปฏิจฺจ เหตุ วิปากธมฺมธมฺโม อุปฺปชฺชติ นปุเรชาตปจฺจยา.\csromanspacer{} นว.\csromanspacer{} นปจฺฉาชาตปจฺจยา.\csromanspacer{} นว.\csromanspacer{} น-อาเสวนปจฺจยา.\csromanspacer{} นว.}

\proseitem{126}{เหตุํ วิปากธมฺมธมฺมํ ปฏิจฺจ นเหตุ วิปากธมฺมธมฺโม อุปฺปชฺชติ นกมฺมปจฺจยา. (1)}

\prose{นเหตุํ วิปากธมฺมธมฺมํ ปฏิจฺจ นเหตุ วิปากธมฺมธมฺโม อุปฺปชฺชติ นกมฺมปจฺจยา. (1)}

\prose{เหตุํ วิปากธมฺมธมฺมญฺจ นเหตุํ วิปากธมฺมธมฺมญฺจ ปฏิจฺจ นเหตุ วิปากธมฺมธมฺโม อุปฺปชฺชติ นกมฺมปจฺจยา. (1) (สํขิตฺตํ.)}

\proseitem{127}{นเหตุยา เอกํ, น-อธิปติยา นว, นปุเรชาเต นว, นปจฺฉาชาเต นว, น-อาเสวเน นว, นกมฺเม ตีณิ, นวิปาเก นว, นวิปฺปยุตฺเต นว. (สํขิตฺตํ.\csromanspacer{} ปจฺจนียํ.)}

\prose{เหตุปจฺจยา น-อธิปติยา นว. (สํขิตฺตํ.\csromanspacer{} อนุโลมปจฺจนียํ.)}

\csromanheader{2}{นเหตุปจฺจยา อารมฺมเณ เอกํ. (สํขิตฺตํ. ปจฺจนียานุโลมํ.)}
\prose{(สหชาตวาโรปิ ปจฺจยวาโรปิ นิสฺสยวาโรปิ สํสฏฺฐวาโรปิ สมฺปยุตฺตวาโรปิ ปฏิจฺจวารสทิสา วิตฺถาเรตพฺพา.)}

\csromanheader{2}{2. วิปากธมฺมปท 7. ปญฺหาวาร}
\csromanheader{2}{\textbf{ปจฺจยจตุกฺก-เหตฺวาทิ}}
\proseitem{128}{เหตุ วิปากธมฺมธมฺโม เหตุสฺส วิปากธมฺมธมฺมสฺส เหตุปจฺจเยน ปจฺจโย.\csromanspacer{} ตีณิ.}

\proseitem{129}{เหตุ วิปากธมฺมธมฺโม เหตุสฺส วิปากธมฺมธมฺมสฺส อารมฺมณปจฺจเยน ปจฺจโย.\csromanspacer{} นว.}

\gambhira{ธมฺมานุโลม ทุกติกปฏฺฐานปาฬิ}
\proseitem{130}{\csromanfolio{222}เหตุ วิปากธมฺมธมฺโม เหตุสฺส วิปากธมฺมธมฺมสฺส อธิปติปจฺจเยน ปจฺจโย.\csromanspacer{} อารมฺมณาธิปติ, สหชาตาธิปติ.\csromanspacer{} นว.}

\proseitem{131}{เหตุ วิปากธมฺมธมฺโม เหตุสฺส วิปากธมฺมธมฺมสฺส อนนฺตรปจฺจเยน ปจฺจโย ฯเปฯ อุปนิสฺสยปจฺจเยน ปจฺจโย.\csromanspacer{} อารมฺมณูปนิสฺสโย, อนนฺตรูปนิสฺสโย, ปกตูปนิสฺสโย.\csromanspacer{} นว.\csromanspacer{} อาเสวนปจฺจเยน ปจฺจโย.\csromanspacer{} นว.}

\proseitem{132}{นเหตุ วิปากธมฺมธมฺโม นเหตุสฺส วิปากธมฺมธมฺมสฺส กมฺมปจฺจเยน ปจฺจโย.\csromanspacer{} ตีณิ.}

\prose{นเหตุ วิปากธมฺมธมฺโม นเหตุสฺส วิปากธมฺมธมฺมสฺส อาหารปจฺจเยน ปจฺจโย. (สํขิตฺตํ.)}

\proseitem{133}{เหตุยา ตีณิ, อารมฺมเณ นว, อธิปติยา นว, อนนฺตเร นว, สมนนฺตเร นว, สหชาเต นว, อญฺญมญฺเญ นว, นิสฺสเย นว, อุปนิสฺสเย นว, อาเสวเน นว, กมฺเม ตีณิ, อาหาเร ตีณิ, อินฺทฺริเย นว, ฌาเน ตีณิ, มคฺเค นว, สมฺปยุตฺเต นว, อตฺถิยา นว, นตฺถิยา นว, วิคเต นว, อวิคเต นว. (สํขิตฺตํ.\csromanspacer{} อนุโลมํ.)}

\csromanheader{2}{\textbf{ปจฺจนียุทฺธาร}}
\proseitem{134}{เหตุ วิปากธมฺมธมฺโม เหตุสฺส วิปากธมฺมธมฺมสฺส อารมฺมณปจฺจเยน ปจฺจโย.\csromanspacer{} สหชาตปจฺจเยน ปจฺจโย.\csromanspacer{} อุปนิสฺสยปจฺจเยน ปจฺจโย. (สํขิตฺตํ.)}

\proseitem{135}{นเหตุยา นว, น-อารมฺมเณ นว. (สํขิตฺตํ.\csromanspacer{} ปจฺจนียํ.)}

\prose{เหตุปจฺจยา น-อารมฺมเณ ตีณิ. (สํขิตฺตํ.\csromanspacer{} อนุโลมปจฺจนียํ.)}

\csromanheader{2}{นเหตุปจฺจยา อารมฺมเณ นว. (สํขิตฺตํ. ปจฺจนียานุโลมํ.)}
\prose{(ยถา กุสลตฺติเก ปญฺหาวารสฺส อนุโลมมฺปิ ปจฺจนียมฺปิ อนุโลมปจฺจนียมฺปิ ปจฺจนียานุโลมมฺปิ คณิตํ, เอวํ คเณตพฺพํ.)}

\vspace{\tipitakaparskip}
\csromancenter{เหตุทุก 3.\csromanspacer{} วิปากตฺติก \textbf{3.}\csromanspacer{}\textbf{ เนววิปากนวิปากธมฺมปท 1.}\csromanspacer{}\textbf{ ปฏิจฺจวาร}}
\csromanheader{2}{\textbf{ปจฺจยจตุกฺก-เหตุ}}
\proseitem{136}{\csromanfolio{223}เหตุํ เนววิปากนวิปากธมฺมธมฺมํ ปฏิจฺจ เหตุ เนววิปากนวิปากธมฺมธมฺโม อุปฺปชฺชติ เหตุปจฺจยา.\csromanspacer{} ตีณิ.}

\prose{นเหตุํ เนววิปากนวิปากธมฺมธมฺมํ ปฏิจฺจ นเหตุ เนววิปากนวิปากธมฺมธมฺโม อุปฺปชฺชติ เหตุปจฺจยา.\csromanspacer{} ตีณิ.}

\prose{เหตุํ เนววิปากนวิปากธมฺมธมฺมญฺจ นเหตุํ เนววิปากนวิปากธมฺมธมฺมญฺจ ปฏิจฺจ เหตุ เนววิปากนวิปากธมฺมธมฺโม อุปฺปชฺชติ เหตุปจฺจยา.\csromanspacer{} ตีณิ. (สํขิตฺตํ.)}

\proseitem{137}{เหตุยา นว, อารมฺมเณ นว, อธิปติยา นว, อนนฺตเร นว, สมนนฺตเร นว, สหชาเต นว, อญฺญมญฺเญ นว, นิสฺสเย นว, อุปนิสฺสเย นว, ปุเรชาเต นว, อาเสวเน นว, กมฺเม นว, วิปาเก เอกํ, อาหาเร นว, อินฺทฺริเย นว, ฌาเน นว, มคฺเค นว, สมฺปยุตฺเต นว, วิปฺปยุตฺเต นว, อตฺถิยา นว, นตฺถิยา นว, วิคเต นว, อวิคเต นว.}

\csromanheader{2}{\textbf{นเหตฺวาทิ}}
\proseitem{138}{นเหตุํ เนววิปากนวิปากธมฺมธมฺมํ ปฏิจฺจ นเหตุ เนววิปากนวิปากธมฺมธมฺโม อุปฺปชฺชติ นเหตุปจฺจยา. (1)}

\proseitem{139}{เหตุํ เนววิปากนวิปากธมฺมธมฺมํ ปฏิจฺจ นเหตุ เนววิปากนวิปากธมฺมธมฺโม อุปฺปชฺชติ น-อารมฺมณปจฺจยา. (1)}

\prose{นเหตุํ เนววิปากนวิปากธมฺมธมฺมํ ปฏิจฺจ นเหตุ เนววิปากนวิปากธมฺมธมฺโม อุปฺปชฺชติ น-อารมฺมณปจฺจยา. (1)}

\prose{เหตุํ เนววิปากนวิปากธมฺมธมฺมญฺจ นเหตุํ เนววิปากนวิปากธมฺมธมฺมญฺจ ปฏิจฺจ นเหตุ เนววิปากนวิปากธมฺมธมฺโม อุปฺปชฺชติ น-อารมฺมณปจฺจยา. (1)}

\proseitem{140}{เหตุํ เนววิปากนวิปากธมฺมธมฺมํ ปฏิจฺจ เหตุ เนววิปากนวิปากธมฺมธมฺโม อุปฺปชฺชติ น-อธิปติปจฺจยา.\csromanspacer{} นว ฯเปฯ.}

\gambhira{ธมฺมานุโลม ทุกติกปฏฺฐานปาฬิ}
\proseitem{141}{\csromanfolio{224}เหตุํ เนววิปากนวิปากธมฺมธมฺมํ ปฏิจฺจ นเหตุ เนววิปากนวิปากธมฺมธมฺโม อุปฺปชฺชติ นกมฺมปจฺจยา.\csromanspacer{} ตีณิ ฯเปฯ.}

\proseitem{142}{นเหตุํ เนววิปากนวิปากธมฺมธมฺมํ ปฏิจฺจ นเหตุ เนววิปากนวิปากธมฺมธมฺโม อุปฺปชฺชติ น-อาหารปจฺจยา. (1) (สํขิตฺตํ.)}

\proseitem{143}{นเหตุยา เอกํ, น-อารมฺมเณ ตีณิ, น-อธิปติยา นว, น-อนนฺตเร ตีณิ, นสมนนฺตเร ตีณิ, น-อญฺญมญฺเญ ตีณิ, น-อุปนิสฺสเย ตีณิ, นปุเรชาเต นว, นปจฺฉาชาเต นว, น-อาเสวเน นว, นกมฺเม ตีณิ, นวิปาเก นว, น-อาหาเร เอกํ, น-อินฺทฺริเย เอกํ, นฌาเน เอกํ, นมคฺเค เอกํ, นสมฺปยุตฺเต ตีณิ, นวิปฺปยุตฺเต นว, โนนตฺถิยา ตีณิ, โนวิคเต ตีณิ. (สํขิตฺตํ.\csromanspacer{} ปจฺจนียํ.)}

\prose{เหตุปจฺจยา น-อารมฺมเณ ตีณิ. (สํขิตฺตํ.\csromanspacer{} อนุโลมปจฺจนียํ.)}

\csromanheader{2}{นเหตุปจฺจยา อารมฺมเณ เอกํ. (สํขิตฺตํ. ปจฺจนียานุโลมํ.)}
\prose{(สหชาตวาโรปิ ปจฺจยวาโรปิ นิสฺสยวาโรปิ สํสฏฺฐวาโรปิ สมฺปยุตฺตวาโรปิ ปฏิจฺจวารสทิสา วิตฺถาเรตพฺพา.)}

\csromanheader{2}{3. เนววิปากนวิปากธมฺมปท 7. ปญฺหาวาร}
\csromanheader{2}{\textbf{ปจฺจยจตุกฺก-เหตฺวาทิ}}
\proseitem{144}{เหตุ เนววิปากนวิปากธมฺมธมฺโม เหตุสฺส เนววิปากนวิปากธมฺมธมฺมสฺส เหตุปจฺจเยน ปจฺจโย.\csromanspacer{} ตีณิ ฯเปฯ.}

\proseitem{145}{เหตุ เนววิปากนวิปากธมฺมธมฺโม เหตุสฺส เนววิปากนวิปากธมฺมธมฺมสฺส อธิปติปจฺจเยน ปจฺจโย.\csromanspacer{} สหชาตาธิปติ.\csromanspacer{} ตีณิ.}

\prose{นเหตุ เนววิปากนวิปากธมฺมธมฺโม นเหตุสฺส เนววิปากนวิปากธมฺมธมฺมสฺส อธิปติปจฺจเยน ปจฺจโย.\csromanspacer{} อารมฺมณาธิปติ, สหชาตาธิปติ.\csromanspacer{} ตีณิ ฯเปฯ.}

\csromanheader{2}{1. เหตุทุก 3. วิปากตฺติก \textbf{ปุเรชาตาทิ}}
\proseitem{146}{\csromanfolio{225}นเหตุ เนววิปากนวิปากธมฺมธมฺโม นเหตุสฺส เนววิปากนวิปากธมฺมธมฺมสฺส ปุเรชาตปจฺจเยน ปจฺจโย.\csromanspacer{} อารมฺมณปุเรชาตํ, วตฺถุปุเรชาตํ.\csromanspacer{} ตีณิ.}

\proseitem{147}{เหตุ เนววิปากนวิปากธมฺมธมฺโม นเหตุสฺส เนววิปากนวิปากธมฺมธมฺมสฺส ปจฺฉาชาตปจฺจเยน ปจฺจโย. (1)}

\prose{นเหตุ เนววิปากนวิปากธมฺมธมฺโม นเหตุสฺส เนววิปากนวิปากธมฺมธมฺมสฺส ปจฺฉาชาตปจฺจเยน ปจฺจโย. (1)}

\prose{เหตุ เนววิปากนวิปากธมฺมธมฺโม จ นเหตุ เนววิปากนวิปากธมฺมธมฺโม จ นเหตุสฺส เนววิปากนวิปากธมฺมธมฺมสฺส ปจฺฉาชาตปจฺจเยน ปจฺจโย. (1) อาเสวนปจฺจเยน ปจฺจโย.\csromanspacer{} นว.}

\proseitem{148}{นเหตุ เนววิปากนวิปากธมฺมธมฺโม นเหตุสฺส เนววิปากนวิปากธมฺมธมฺมสฺส กมฺมปจฺจเยน ปจฺจโย.\csromanspacer{} ตีณิ. (สํขิตฺตํ.)}

\proseitem{149}{เหตุยา ตีณิ, อารมฺมเณ นว, อธิปติยา ฉ, อนนฺตเร นว, สมนนฺตเร นว, สหชาเต นว, อญฺญมญฺเญ นว, นิสฺสเย นว, อุปนิสฺสเย นว, ปุเรชาเต ตีณิ, ปจฺฉาชาเต ตีณิ, อาเสวเน นว, กมฺเม ตีณิ, อาหาเร ตีณิ, อินฺทฺริเย นว, ฌาเน ตีณิ, มคฺเค นว, สมฺปยุตฺเต นว, วิปฺปยุตฺเต ปญฺจ, อตฺถิยา นว, นตฺถิยา นว, วิคเต นว, อวิคเต นว.}

\csromanheader{2}{\textbf{ปจฺจนียุทฺธาร}}
\proseitem{150}{เหตุ เนววิปากนวิปากธมฺมธมฺโม เหตุสฺส เนววิปากนวิปากธมฺมธมฺมสฺส อารมฺมณปจฺจเยน ปจฺจโย.\csromanspacer{} สหชาตปจฺจเยน ปจฺจโย.\csromanspacer{} อุปนิสฺสยปจฺจเยน ปจฺจโย. (สํขิตฺตํ.)}

\proseitem{151}{นเหตุยา นว, น-อารมฺมเณ นว, น-อธิปติยา นว. (สํขิตฺตํ.)}

\csromanheader{2}{เหตุปจฺจยา น-อารมฺมเณ ตีณิ. (สํขิตฺตํ.)}
\csromanheader{2}{นเหตุปจฺจยา อารมฺมเณ นว. (สํขิตฺตํ.)}
\gambhira{ธมฺมานุโลม ทุกติกปฏฺฐานปาฬิ}
\prose{\csromanfolio{226}(ยถา กุสลตฺติเก ปญฺหาวารสฺส อนุโลมมฺปิ ปจฺจนียมฺปิ อนุโลมปจฺจนียมฺปิ ปจฺจนียานุโลมมฺปิ คณิตํ, เอวํ คเณตพฺพํ.)}

\nitthitam{เหตุทุกวิปากตฺติกํ นิฏฺฐิตํ.}
\csromanmatikaanchor{mtk.35}
\csromantocmarkref{section}{4. อุปาทินฺนตฺติก}{mtk.35}
\bookmark[dest={mtk.35},level=1]{4. อุปาทินฺนตฺติก}
\csromanheader{1}{1. เหตุทุก 4. อุปาทินฺนตฺติก}
\csromanheader{2}{1. อุปาทินฺนุปาทานิยปท 1. ปฏิจฺจวาร}
\csromanheader{2}{\textbf{ปจฺจยจตุกฺก-เหตุ}}
\proseitem{152}{เหตุํ อุปาทินฺนุปาทานิยํ ธมฺมํ ปฏิจฺจ เหตุ อุปาทินฺนุปาทานิโย ธมฺโม อุปฺปชฺชติ เหตุปจฺจยา.\csromanspacer{} ตีณิ.}

\prose{นเหตุํ อุปาทินฺนุปาทานิยํ ธมฺมํ ปฏิจฺจ นเหตุ อุปาทินฺนุปาทานิโย ธมฺโม อุปฺปชฺชติ เหตุปจฺจยา.\csromanspacer{} ตีณิ.}

\prose{เหตุํ อุปาทินฺนุปาทานิยญฺจ นเหตุํ อุปาทินฺนุปาทานิยญฺจ ธมฺมํ ปฏิจฺจ เหตุ อุปาทินฺนุปาทานิโย ธมฺโม อุปฺปชฺชติ เหตุปจฺจยา.\csromanspacer{} ตีณิ. (สํขิตฺตํ.)}

\proseitem{153}{เหตุยา นว, อารมฺมเณ นว, อนนฺตเร นว, สมนนฺตเร นว ฯเปฯ กมฺเม นว, วิปาเก นว ฯเปฯ อวิคเต นว. (สํขิตฺตํ.)}

\csromanheader{2}{\textbf{นเหตฺวาทิ}}
\proseitem{154}{นเหตุํ อุปาทินฺนุปาทานิยํ ธมฺมํ ปฏิจฺจ นเหตุ อุปาทินฺนุปาทานิโย ธมฺโม อุปฺปชฺชติ นเหตุปจฺจยา. (1)}

\proseitem{155}{เหตุํ อุปาทินฺนุปาทานิยํ ธมฺมํ ปฏิจฺจ นเหตุ อุปาทินฺนุปาทานิโย ธมฺโม อุปฺปชฺชติ น-อารมฺมณปจฺจยา. (1)}

\prose{นเหตุํ อุปาทินฺนุปาทานิยํ ธมฺมํ ปฏิจฺจ นเหตุ อุปาทินฺนุปาทานิโย ธมฺโม อุปฺปชฺชติ น-อารมฺมณปจฺจยา. (1)}

\prose{เหตุํ อุปาทินฺนุปาทานิยญฺจ นเหตุํ อุปาทินฺนุปาทานิยญฺจ ธมฺมํ ปฏิจฺจ นเหตุ อุปาทินฺนุปาทานิโย ธมฺโม อุปฺปชฺชติ น-อารมฺมณปจฺจยา. (1)}

\csromanheader{2}{1. เหตุทุก 4. อุปาทินฺนตฺติก}
\proseitem{156}{\csromanfolio{227}เหตุํ อุปาทินฺนุปาทานิยํ ธมฺมํ ปฏิจฺจ เหตุ อุปาทินฺนุปาทานิโย ธมฺโม อุปฺปชฺชติ น-อธิปติปจฺจยา.\csromanspacer{} นว ฯเปฯ.}

\proseitem{157}{นเหตุํ อุปาทินฺนุปาทานิยํ ธมฺมํ ปฏิจฺจ นเหตุ อุปาทินฺนุปาทานิโย ธมฺโม อุปฺปชฺชติ นวิปากปจฺจยา.\csromanspacer{} น-อาหารปจฺจยา. (สํขิตฺตํ.)}

\proseitem{158}{นเหตุยา เอกํ, น-อารมฺมเณ ตีณิ, น-อธิปติยา นว, น-อนนฺตเร ตีณิ, นสมนนฺตเร ตีณิ, น-อญฺญมญฺเญ ตีณิ, น-อุปนิสฺสเย ตีณิ, นปุเรชาเต นว, นปจฺฉาชาเต นว, น-อาเสวเน นว, นวิปาเก เอกํ, น-อาหาเร เอกํ, น-อินฺทฺริเย เอกํ, นฌาเน เอกํ, นมคฺเค เอกํ, นสมฺปยุตฺเต ตีณิ, นวิปฺปยุตฺเต นว, โนนตฺถิยา ตีณิ, โนวิคเต ตีณิ.}

\csromanheader{2}{เหตุปจฺจยา น-อารมฺมเณ ตีณิ. (สํขิตฺตํ.)}
\csromanheader{2}{นเหตุปจฺจยา อารมฺมเณ เอกํ. (สํขิตฺตํ.)}
\prose{(สหชาตวาโรปิ ปจฺจยวาโรปิ นิสฺสยวาโรปิ สํสฏฺฐวาโรปิ สมฺปยุตฺตวาโรปิ ปฏิจฺจวารสทิสา วิตฺถาเรตพฺพา.)}

\csromanheader{2}{1. อุปาทินฺนุปาทานิยปท 7. ปญฺหาวาร}
\csromanheader{2}{\textbf{ปจฺจยจตุกฺก-เหตุ อารมฺมณ}}
\proseitem{159}{เหตุ อุปาทินฺนุปาทานิโย ธมฺโม เหตุสฺส อุปาทินฺนุปาทานิยสฺส ธมฺมสฺส เหตุปจฺจเยน ปจฺจโย.\csromanspacer{} ตีณิ.}

\proseitem{160}{เหตุ อุปาทินฺนุปาทานิโย ธมฺโม เหตุสฺส อุปาทินฺนุปาทานิยสฺส ธมฺมสฺส อารมฺมณปจฺจเยน ปจฺจโย.\csromanspacer{} ตีณิ.}

\prose{นเหตุ อุปาทินฺนุปาทานิโย ธมฺโม นเหตุสฺส อุปาทินฺนุปาทานิยสฺส ธมฺมสฺส อารมฺมณปจฺจเยน ปจฺจโย.\csromanspacer{} ตีณิ.}

\prose{เหตุ อุปาทินฺนุปาทานิโย จ นเหตุ อุปาทินฺนุปาทานิโย จ ธมฺมา เหตุสฺส อุปาทินฺนุปาทานิยสฺส ธมฺมสฺส อารมฺมณปจฺจเยน ปจฺจโย.\csromanspacer{} ตีณิ. (สํขิตฺตํ.)}

\gambhira{ธมฺมานุโลม ทุกติกปฏฺฐานปาฬิ}
\proseitem{161}{\csromanfolio{228}เหตุยา ตีณิ, อารมฺมเณ นว, อนนฺตเร นว, สมนนฺตเร นว, สหชาเต นว, อญฺญมญฺเญ นว, นิสฺสเย นว, อุปนิสฺสเย นว, ปุเรชาเต ตีณิ, ปจฺฉาชาเต ตีณิ, กมฺเม ตีณิ, วิปาเก นว, อาหาเร ตีณิ, อินฺทฺริเย นว, ฌาเน ตีณิ, มคฺเค นว, สมฺปยุตฺเต นว, วิปฺปยุตฺเต ปญฺจ, อตฺถิยา นว, นตฺถิยา นว, วิคเต นว, อวิคเต นว.}

\csromanheader{2}{\textbf{ปจฺจนียุทฺธาร}}
\proseitem{162}{เหตุ อุปาทินฺนุปาทานิโย ธมฺโม เหตุสฺส อุปาทินฺนุปาทานิยสฺส ธมฺมสฺส อารมฺมณปจฺจเยน ปจฺจโย.\csromanspacer{} สหชาตปจฺจเยน ปจฺจโย.\csromanspacer{} อุปนิสฺสยปจฺจเยน ปจฺจโย.\csromanspacer{} ตีณิ.}

\prose{นเหตุ อุปาทินฺนุปาทานิโย ธมฺโม นเหตุสฺส อุปาทินฺนุปาทานิยสฺส ธมฺมสฺส อารมฺมณปจฺจเยน ปจฺจโย.\csromanspacer{} สหชาตปจฺจเยน ปจฺจโย.\csromanspacer{} อุปนิสฺสยปจฺจเยน ปจฺจโย.\csromanspacer{} ปุเรชาตปจฺจเยน ปจฺจโย.\csromanspacer{} ปจฺฉาชาตปจฺจเยน ปจฺจโย.\csromanspacer{} อาหารปจฺจเยน ปจฺจโย.\csromanspacer{} อินฺทฺริยปจฺจเยน ปจฺจโย. (สํขิตฺตํ.)}

\proseitem{163}{นเหตุยา นว, น-อารมฺมเณ นว, น-อธิปติยา นว. (สํขิตฺตํ.)}

\csromanheader{2}{เหตุปจฺจยา น-อารมฺมเณ ตีณิ. (สํขิตฺตํ.)}
\csromanheader{2}{นเหตุปจฺจยา อารมฺมเณ นว. (สํขิตฺตํ.)}
\prose{(ยถา กุสลตฺติเก ปญฺหาวารสฺส อนุโลมมฺปิ ปจฺจนียมฺปิ อนุโลมปจฺจนียมฺปิ ปจฺจนียานุโลมมฺปิ คณิตํ, เอวํ คเณตพฺพํ.)}

\csromanheader{2}{2. อนุปาทินฺนุปาทานิยปท 1. ปฏิจฺจวาร}
\csromanheader{2}{\textbf{ปจฺจยจตุกฺก-เหตุ}}
\proseitem{164}{เหตุํ อนุปาทินฺนุปาทานิยํ ธมฺมํ ปฏิจฺจ เหตุ อนุปาทินฺนุปาทานิโย ธมฺโม อุปฺปชฺชติ เหตุปจฺจยา.\csromanspacer{} ตีณิ.}

\prose{นเหตุํ อนุปาทินฺนุปาทานิยํ ธมฺมํ ปฏิจฺจ นเหตุ อนุปาทินฺนุปาทานิโย ธมฺโม อุปฺปชฺชติ เหตุปจฺจยา.\csromanspacer{} ตีณิ.}

\csromanheader{2}{1. เหตุทุก 4. อุปาทินฺนตฺติก}
\prose{\csromanfolio{229}เหตุํ อนุปาทินฺนุปาทานิยญฺจ นเหตุํ อนุปาทินฺนุปาทานิยญฺจ ธมฺมํ ปฏิจฺจ เหตุ อนุปาทินฺนุปาทานิโย ธมฺโม อุปฺปชฺชติ เหตุปจฺจยา.\csromanspacer{} ตีณิ. (สํขิตฺตํ.)}

\proseitem{165}{เหตุยา นว, อารมฺมเณ นว, อธิปติยา นว, อนนฺตเร นว, สมนนฺตเร นว, สหชาเต นว, อญฺญมญฺเญ นว, นิสฺสเย นว, อุปนิสฺสเย นว, ปุเรชาเต นว, อาเสวเน นว, กมฺเม นว, วิปาเก เอกํ, อาหาเร นว, อินฺทฺริเย นว, สมฺปยุตฺเต นว, วิปฺปยุตฺเต นว, อตฺถิยา นว, นตฺถิยา นว, วิคเต นว, อวิคเต นว.}

\csromanheader{2}{\textbf{นเหตุ น-อารมฺมณ}}
\proseitem{166}{นเหตุํ อนุปาทินฺนุปาทานิยํ ธมฺมํ ปฏิจฺจ นเหตุ อนุปาทินฺนุปาทานิโย ธมฺโม อุปฺปชฺชติ นเหตุปจฺจยา.\csromanspacer{} เทฺว. (2)}

\proseitem{167}{เหตุํ อนุปาทินฺนุปาทานิยํ ธมฺมํ ปฏิจฺจ นเหตุ อนุปาทินฺนุปาทานิโย ธมฺโม อุปฺปชฺชติ น-อารมฺมณปจฺจยา. (1)}

\prose{นเหตุํ อนุปาทินฺนุปาทานิยํ ธมฺมํ ปฏิจฺจ นเหตุ อนุปาทินฺนุปาทานิโย ธมฺโม อุปฺปชฺชติ น-อารมฺมณปจฺจยา. (1)}

\prose{เหตุํ อนุปาทินฺนุปาทานิยญฺจ นเหตุํ อนุปาทินฺนุปาทานิยญฺจ ธมฺมํ ปฏิจฺจ นเหตุ อนุปาทินฺนุปาทานิโย ธมฺโม อุปฺปชฺชติ น-อารมฺมณปจฺจยา. (1) (สํขิตฺตํ.)}

\proseitem{168}{นเหตุยา เทฺว, น-อารมฺมเณ ตีณิ, น-อธิปติยา นว, น-อนนฺตเร ตีณิ, นสมนนฺตเร ตีณิ, น-อญฺญมญฺเญ ตีณิ, น-อุปนิสฺสเย ตีณิ, นปุเรชาเต นว, นปจฺฉาชาเต นว, น-อาเสวเน นว, นกมฺเม ตีณิ, นวิปาเก นว, น-อาหาเร เอกํ, น-อินฺทฺริเย เอกํ, นฌาเน เอกํ, นมคฺเค เอกํ, นสมฺปยุตฺเต ตีณิ, นวิปฺปยุตฺเต นว, โนนตฺถิยา ตีณิ, โนวิคเต ตีณิ. (สํขิตฺตํ.)}

\csromanheader{2}{เหตุปจฺจยา น-อารมฺมเณ ตีณิ. (สํขิตฺตํ.)}
\csromanheader{2}{นเหตุปจฺจยา อารมฺมเณ เทฺว. (สํขิตฺตํ.)}
\prose{(สหชาตวาโรปิ ปจฺจยวาโรปิ นิสฺสยวาโรปิ สํสฏฺฐวาโรปิ สมฺปยุตฺตวาโรปิ ปฏิจฺจวารสทิสา วิตฺถาเรตพฺพา.)}

\csromanheader{2}{ธมฺมานุโลม ทุกติกปฏฺฐานปาฬิ \textbf{2. อนุปาทินฺนุปาทานิยปท 7. ปญฺหาวาร}}
\csromanheader{2}{\textbf{ปจฺจยจตุกฺก-เหตฺวาทิ}}
\proseitem{169}{\csromanfolio{230}เหตุ อนุปาทินฺนุปาทานิโย ธมฺโม เหตุสฺส อนุปาทินฺนุปาทานิยสฺส ธมฺมสฺส เหตุปจฺจเยน ปจฺจโย.\csromanspacer{} ตีณิ.}

\proseitem{170}{เหตุ อนุปาทินฺนุปาทานิโย ธมฺโม เหตุสฺส อนุปาทินฺนุปาทานิยสฺส ธมฺมสฺส อารมฺมณปจฺจเยน ปจฺจโย.\csromanspacer{} ตีณิ.}

\prose{นเหตุ อนุปาทินฺนุปาทานิโย ธมฺโม นเหตุสฺส อนุปาทินฺนุปาทานิยสฺส ธมฺมสฺส อารมฺมณปจฺจเยน ปจฺจโย.\csromanspacer{} นว.}

\proseitem{171}{เหตุ อนุปาทินฺนุปาทานิโย ธมฺโม เหตุสฺส อนุปาทินฺนุปาทานิยสฺส ธมฺมสฺส อธิปติปจฺจเยน ปจฺจโย.\csromanspacer{} อารมฺมณาธิปติ, สหชาตาธิปติ.\csromanspacer{} ตีณิ.}

\prose{นเหตุ อนุปาทินฺนุปาทานิโย ธมฺโม นเหตุสฺส อนุปาทินฺนุปาทานิยสฺส ธมฺมสฺส อธิปติปจฺจเยน ปจฺจโย.\csromanspacer{} อารมฺมณาธิปติ, สหชาตาธิปติ.\csromanspacer{} นว ฯเปฯ.}

\proseitem{172}{เหตุ อนุปาทินฺนุปาทานิโย ธมฺโม เหตุสฺส อนุปาทินฺนุปาทานิยสฺส ธมฺมสฺส อุปนิสฺสยปจฺจเยน ปจฺจโย.\csromanspacer{} อารมฺมณูปนิสฺสโย, อนนฺตรูปนิสฺสโย, ปกตูปนิสฺสโย.\csromanspacer{} นว.}

\proseitem{173}{นเหตุ อนุปาทินฺนุปาทานิโย ธมฺโม นเหตุสฺส อนุปาทินฺนุปาทานิยสฺส ธมฺมสฺส ปุเรชาตปจฺจเยน ปจฺจโย.\csromanspacer{} อารมฺมณปุเรชาตํ.\csromanspacer{} ตีณิ. (สํขิตฺตํ.)}

\proseitem{174}{เหตุยา ตีณิ, อารมฺมเณ นว, อธิปติยา นว, อนนฺตเร นว, สมนนฺตเร นว, สหชาเต นว, อญฺญมญฺเญ นว, นิสฺสเย นว, อุปนิสฺสเย นว, ปุเรชาเต ตีณิ, ปจฺฉาชาเต ตีณิ, อาเสวเน นว, กมฺเม ตีณิ, อาหาเร ตีณิ, อินฺทฺริเย นว, ฌาเน ตีณิ, มคฺเค นว, สมฺปยุตฺเต นว, วิปฺปยุตฺเต ตีณิ, อตฺถิยา นว, นตฺถิยา นว, วิคเต นว, อวิคเต นว.}

\csromanheader{2}{1. เหตุทุก 4. อุปาทินฺนตฺติก \textbf{ปจฺจนียุทฺธาร}}
\proseitem{175}{\csromanfolio{231}เหตุ อนุปาทินฺนุปาทานิโย ธมฺโม เหตุสฺส อนุปาทินฺนุปาทานิยสฺส ธมฺมสฺส อารมฺมณปจฺจเยน ปจฺจโย.\csromanspacer{} สหชาตปจฺจเยน ปจฺจโย.\csromanspacer{} อุปนิสฺสยปจฺจเยน ปจฺจโย. (สํขิตฺตํ.)}

\proseitem{176}{นเหตุยา นว, น-อารมฺมเณ นว. (สํขิตฺตํ.)}

\csromanheader{2}{เหตุปจฺจยา น-อารมฺมเณ ตีณิ. (สํขิตฺตํ.)}
\csromanheader{2}{นเหตุปจฺจยา อารมฺมเณ นว. (สํขิตฺตํ.)}
\prose{(ยถา กุสลตฺติเก ปญฺหาวารสฺส อนุโลมมฺปิ ปจฺจนียมฺปิ อนุโลมปจฺจนียมฺปิ ปจฺจนียานุโลมมฺปิ คณิตํ, เอวํ คเณตพฺพํ.)}

\csromanheader{2}{3. อนุปาทินฺน-อนุปาทานิยปท 1. ปฏิจฺจวาร}
\csromanheader{2}{\textbf{ปจฺจยจตุกฺก-เหตุ}}
\proseitem{177}{เหตุํ อนุปาทินฺน-อนุปาทานิยํ ธมฺมํ ปฏิจฺจ เหตุ อนุปาทินฺน-อนุปาทานิโย ธมฺโม อุปฺปชฺชติ เหตุปจฺจยา.\csromanspacer{} ตีณิ.}

\prose{นเหตุํ อนุปาทินฺน-อนุปาทานิยํ ธมฺมํ ปฏิจฺจ นเหตุ อนุปาทินฺน-อนุปาทานิโย ธมฺโม อุปฺปชฺชติ เหตุปจฺจยา.\csromanspacer{} ตีณิ.}

\prose{เหตุํ อนุปาทินฺน-อนุปาทานิยญฺจ นเหตุํ อนุปาทินฺน-อนุปาทานิยญฺจ ธมฺมํ ปฏิจฺจ เหตุ อนุปาทินฺน-อนุปาทานิโย ธมฺโม อุปฺปชฺชติ เหตุปจฺจยา.\csromanspacer{} ตีณิ. (สํขิตฺตํ.)}

\proseitem{178}{เหตุยา นว, อารมฺมเณ นว, อธิปติยา นว, อนนฺตเร นว, สมนนฺตเร นว, สหชาเต นว, อญฺญมญฺเญ นว, นิสฺสเย นว, อุปนิสฺสเย นว, ปุเรชาเต นว, อาเสวเน นว, กมฺเม นว, วิปาเก นว, อาหาเร นว, อินฺทฺริเย นว, ฌาเน นว, มคฺเค นว, สมฺปยุตฺเต นว, วิปฺปยุตฺเต นว, อตฺถิยา นว, นตฺถิยา นว, วิคเต นว, อวิคเต นว.}

\csromanheader{2}{ธมฺมานุโลม ทุกติกปฏฺฐานปาฬิ \textbf{น-อธิปติ}}
\proseitem{179}{\csromanfolio{232}เหตุํ อนุปาทินฺน-อนุปาทานิยํ ธมฺมํ ปฏิจฺจ เหตุ อนุปาทินฺน-อนุปาทานิโย ธมฺโม อุปฺปชฺชติ น-อธิปติปจฺจยา.\csromanspacer{}\csromanspacer{}เหตุํ อนุปาทินฺน-อนุปาทานิยํ ธมฺมํ ปฏิจฺจ นเหตุ อนุปาทินฺน-อนุปาทานิโย ธมฺโม อุปฺปชฺชติ น-อธิปติปจฺจยา. (2)}

\prose{นเหตุํ อนุปาทินฺน-อนุปาทานิยํ ธมฺมํ ปฏิจฺจ นเหตุ อนุปาทินฺน-อนุปาทานิโย ธมฺโม อุปฺปชฺชติ น-อธิปติปจฺจยา.\csromanspacer{}\csromanspacer{}นเหตุํ อนุปาทินฺน-อนุปาทานิย ธมฺมํ ปฏิจฺจ เหตุ อนุปาทินฺน-อนุปาทานิโย ธมฺโม อุปฺปชฺชติ น-อธิปติปจฺจยา. (2)}

\prose{เหตุํ อนุปาทินฺน-อนุปาทานิยญฺจ นเหตุํ อนุปาทินฺน-อนุปาทานิยญฺจ ธมฺมํ ปฏิจฺจ เหตุ อนุปาทินฺน-อนุปาทานิโย ธมฺโม อุปฺปชฺชติ น-อธิปติปจฺจยา.\csromanspacer{}\csromanspacer{}เหตุํ อนุปาทินฺน-อนุปาทานิยญฺจ นเหตุํ อนุปาทินฺน-อนุปาทานิยญฺจ ธมฺมํ ปฏิจฺจ นเหตุ อนุปาทินฺน-อนุปาทานิโย ธมฺโม อุปฺปชฺชติ น-อธิปติปจฺจยา. (2) (สํขิตฺตํ.)}

\vspace{\tipitakaparskip}
\csromancenter{\textbf{น-อธิปติ}ยา ฉ, นปุเรชาเต นว, นปจฺฉาชาเต นว, น-อาเสวเน นว, นกมฺเม ตีณิ, นวิปาเก นว, นวิปฺปยุตฺเต นว. (สํขิตฺตํ.)}
\csromanheader{2}{เหตุปจฺจยา น-อธิปติยา ฉ. (สํขิตฺตํ.)}
\csromanheader{2}{\textbf{น-อธิปติ}ปจฺจยา เหตุยา ฉ. (สํขิตฺตํ.)}
\prose{(สหชาตวาโรปิ ปจฺจยวาโรปิ นิสฺสยวาโรปิ สํสฏฺฐวาโรปิ สมฺปยุตฺตวาโรปิ ปฏิจฺจวารสทิสา วิตฺถาเรตพฺพา.)}

\csromanheader{2}{3. อนุปาทินฺน-อนุปาทานิยปท 7. ปญฺหาวาร}
\csromanheader{2}{\textbf{ปจฺจยจตุกฺก-เหตุ อารมฺมณ}}
\proseitem{181}{เหตุ อนุปาทินฺน-อนุปาทานิโย ธมฺโม เหตุสฺส อนุปาทินฺน-อนุปาทานิยสฺส ธมฺมสฺส เหตุปจฺจเยน ปจฺจโย.\csromanspacer{} ตีณิ.}

\csromanmatikaanchor{mtk.36}
\csromantocmarkref{section}{5. สํกิลิฏฺฐตฺติก}{mtk.36}
\bookmark[dest={mtk.36},level=1]{5. สํกิลิฏฺฐตฺติก}
\csromanheader{1}{1. เหตุทุก 5. สํกิลิฏฺฐตฺติก}
\proseitem{182}{\csromanfolio{233}นเหตุ อนุปาทินฺน-อนุปาทานิโย ธมฺโม นเหตุสฺส อนุปาทินฺน-อนุปาทานิยสฺส ธมฺมสฺส อารมฺมณปจฺจเยน ปจฺจโย.\csromanspacer{} ตีณิ. (สํขิตฺตํ.)}

\proseitem{183}{เหตุยา ตีณิ, อารมฺมเณ ตีณิ, อธิปติยา ฉ, อนนฺตเร นว, สมนนฺตเร นว, สหชาเต นว, อญฺญมญฺเญ นว, นิสฺสเย นว, อุปนิสฺสเย นว, กมฺเม ตีณิ, วิปาเก นว, อาหาเร ตีณิ, อินฺทฺริเย นว, ฌาเน ตีณิ, มคฺเค นว, สมฺปยุตฺเต นว, อตฺถิยา นว, นตฺถิยา นว, วิคเต นว, อวิคเต นว.}

\csromanheader{2}{\textbf{ปจฺจนียุทฺธาร}}
\proseitem{184}{เหตุ อนุปาทินฺน-อนุปาทานิโย ธมฺโม เหตุสฺส อนุปาทินฺน-อนุปาทานิยสฺส ธมฺมสฺส สหชาตปจฺจเยน ปจฺจโย.\csromanspacer{} อุปนิสฺสยปจฺจเยน ปจฺจโย. (สํขิตฺตํ.)}

\proseitem{185}{นเหตุยา นว, น-อารมฺมเณ นว. (สํขิตฺตํ.)}

\csromanheader{2}{เหตุปจฺจยา น-อารมฺมเณ ตีณิ. (สํขิตฺตํ.)}
\csromanheader{2}{นเหตุปจฺจยา อารมฺมเณ ตีณิ. (สํขิตฺตํ.)}
\prose{(ยถา กุสลตฺติเก ปญฺหาวารสฺส อนุโลมมฺปิ ปจฺจนียมฺปิ อนุโลมปจฺจนียมฺปิ ปจฺจนียานุโลมมฺปิ คณิตํ, เอวํ คเณตพฺพํ.)}

\nitthitam{เหตุทุก-อุปาทินฺนตฺติกํ นิฏฺฐิตํ.}
\csromanheader{2}{1. เหตุทุก 5. สํกิลิฏฺฐตฺติก}
\csromanheader{2}{1. สํกิลิฏฺฐสํกิเลสิกปท 1. ปฏิจฺจวาร}
\csromanheader{2}{\textbf{ปจฺจยจตุกฺก-เหตุ}}
\proseitem{186}{เหตุํ สํกิลิฏฺฐสํกิเลสิกํ ธมฺมํ ปฏิจฺจ เหตุ สํกิลิฏฺฐสํกิเลสิโก ธมฺโม อุปฺปชฺชติ เหตุปจฺจยา.\csromanspacer{} ตีณิ.}

\prose{นเหตุํ สํกิลิฏฺฐสํกิเลสิกํ ธมฺมํ ปฏิจฺจ นเหตุ สํกิลิฏฺฐสํกิเลสิโก ธมฺโม อุปฺปชฺชติ เหตุปจฺจยา.\csromanspacer{} ตีณิ.}

\gambhira{ธมฺมานุโลม ทุกติกปฏฺฐานปาฬิ}
\prose{\csromanfolio{234}เหตุํ สํกิลิฏฺฐสํกิเลสิกญฺจ นเหตุํ สํกิลิฏฺฐสํกิเลสิกญฺจ ธมฺมํ ปฏิจฺจ เหตุ สํกิลิฏฺฐสํกิเลสิโก ธมฺโม อุปฺปชฺชติ เหตุปจฺจยา.\csromanspacer{} ตีณิ. (สํขิตฺตํ.)}

\proseitem{187}{เหตุยา นว, อารมฺมเณ นว, อธิปติยา นว ฯเปฯ กมฺเม นว, อาหาเร นว, อวิคเต นว. (สํขิตฺตํ.)}

\csromanheader{2}{\textbf{นเหตุ น-อธิปติ}}
\proseitem{188}{นเหตุํ สํกิลิฏฺฐสํกิเลสิกํ ธมฺมํ ปฏิจฺจ เหตุ สํกิลิฏฺฐสํกิเลสิโก ธมฺโม อุปฺปชฺชติ นเหตุปจฺจยา.}

\prose{เหตุํ สํกิลิฏฺฐสํกิเลสิกํ ธมฺมํ ปฏิจฺจ เหตุ สํกิลิฏฺฐสํกิเลสิโก ธมฺโม อุปฺปชฺชติ น-อธิปติปจฺจยา. (สํขิตฺตํ.)}

\proseitem{189}{นเหตุยา เอกํ, น-อธิปติยา นว, นปุเรชาเต นว, นปจฺฉาชาเต นว, น-อาเสวเน นว, นกมฺเม ตีณิ, นวิปาเก นว, นวิปฺปยุตฺเต นว. (สํขิตฺตํ.)}

\csromanheader{2}{เหตุปจฺจยา น-อธิปติยา นว. (สํขิตฺตํ.)}
\csromanheader{2}{นเหตุปจฺจยา อารมฺมเณ เอกํ. (สํขิตฺตํ.)}
\prose{(สหชาตวาโรปิ ปจฺจยวาโรปิ นิสฺสยวาโรปิ สํสฏฺฐวาโรปิ สมฺปยุตฺตวาโรปิ ปฏิจฺจวารสทิสา วิตฺถาเรตพฺพา.)}

\csromanheader{2}{1. สํกิลิฏฺฐสํกิเลสิกปท 7. ปญฺหาวาร}
\csromanheader{2}{\textbf{ปจฺจยจตุกฺก-เหตฺวาทิ}}
\proseitem{190}{เหตุ สํกิลิฏฺฐสํกิเลสิโก ธมฺโม เหตุสฺส สํกิลิฏฺฐสํกิเลสิกสฺส ธมฺมสฺส เหตุปจฺจเยน ปจฺจโย.\csromanspacer{} ตีณิ.}

\prose{เหตุ สํกิลิฏฺฐสํกิเลสิโก ธมฺโม เหตุสฺส สํกิลิฏฺฐสํกิเลสิกสฺส ธมฺมสฺส อารมฺมณปจฺจเยน ปจฺจโย.\csromanspacer{} นว.}

\prose{เหตุ สํกิลิฏฺฐสํกิเลสิโก ธมฺโม เหตุสฺส สํกิลิฏฺฐสํกิเลสิกสฺส ธมฺมสฺส อธิปติปจฺจเยน ปจฺจโย.\csromanspacer{} อารมฺมณาธิปติ.\csromanspacer{} นว. (สํขิตฺตํ.)}

\csromanheader{2}{1. เหตุทุก 5. สํกิลิฏฺฐตฺติก}
\proseitem{191}{\csromanfolio{235}เหตุยา ตีณิ, อารมฺมเณ นว, อธิปติยา นว, อนนฺตเร นว, สมนนฺตเร นว, สหชาเต นว, อญฺญมญฺเญ นว, นิสฺสเย นว, อุปนิสฺสเย นว, อาเสวเน นว, กมฺเม ตีณิ, อาหาเร ตีณิ, อินฺทฺริเย ตีณิ, ฌาเน ตีณิ, มคฺเค ตีณิ, สมฺปยุตฺเต นว, อตฺถิยา นว, นตฺถิยา นว, วิคเต นว, อวิคเต นว. (สํขิตฺตํ.)}

\csromanheader{2}{\textbf{ปจฺจนียุทฺธาร}}
\proseitem{192}{เหตุ สํกิลิฏฺฐสํกิเลสิโก ธมฺโม เหตุสฺส สํกิลิฏฺฐสํกิเลสิกสฺส ธมฺมสฺส อารมฺมณปจฺจเยน ปจฺจโย.\csromanspacer{} สหชาตปจฺจเยน ปจฺจโย.\csromanspacer{} อุปนิสฺสยปจฺจเยน ปจฺจโย. (สํขิตฺตํ.)}

\proseitem{193}{นเหตุยา นว, น-อารมฺมเณ นว. (สํขิตฺตํ.)}

\csromanheader{2}{เหตุปจฺจยา น-อารมฺมเณ ตีณิ. (สํขิตฺตํ.)}
\csromanheader{2}{นเหตุปจฺจยา อารมฺมเณ นว. (สํขิตฺตํ.)}
\prose{(ยถา กุสลตฺติเก ปญฺหาวารสฺส อนุโลมมฺปิ ปจฺจนียมฺปิ อนุโลมปจฺจนียมฺปิ ปจฺจนียานุโลมมฺปิ คณิตํ, เอวํ คเณตพฺพํ.)}

\csromanheader{2}{2. อสํกิลิฏฺฐสํกิเลสิกปท 1. ปฏิจฺจวาร}
\csromanheader{2}{\textbf{ปจฺจยจตุกฺก-เหตุ}}
\proseitem{194}{เหตุํ อสํกิลิฏฺฐสํกิเลสิกํ ธมฺมํ ปฏิจฺจ เหตุ อสํกิลิฏฺฐสํกิเลสิโก ธมฺโม อุปฺปชฺชติ เหตุปจฺจยา.\csromanspacer{} ตีณิ.}

\prose{นเหตุํ อสํกิลิฏฺฐสํกิเลสิกํ ธมฺมํ ปฏิจฺจ นเหตุ อสํกิลิฏฺฐสํกิเลสิโก ธมฺโม อุปฺปชฺชติ เหตุปจฺจยา.\csromanspacer{} ตีณิ.}

\prose{เหตุํ อสํกิลิฏฺฐสํกิเลสิกญฺจ นเหตุํ อสํกิลิฏฺฐสํกิเลสิกญฺจ ธมฺมํ ปฏิจฺจ เหตุ อสํกิลิฏฺฐสํกิเลสิโก ธมฺโม อุปฺปชฺชติ เหตุปจฺจยา.\csromanspacer{} ตีณิ. (สํขิตฺตํ.)}

\proseitem{195}{เหตุยา นว, อารมฺมเณ นว, อธิปติยา นว ฯเปฯ อวิคเต นว. (สํขิตฺตํ.)}

\gambhira{ธมฺมานุโลม ทุกติกปฏฺฐานปาฬิ \textbf{นเหตุ}}
\vspace{\tipitakaparskip}
\csromancenter{\textbf{นเหตุ}ํ อสํกิลิฏฺฐสํกิเลสิกํ ธมฺมํ ปฏิจฺจ นเหตุ อสํกิลิฏฺฐสํกิเลสิโก ธมฺโม อุปฺปชฺชติ นเหตุปจฺจยา. (สํขิตฺตํ.)}
\csromancenter{\textbf{นเหตุ}ยา เอกํ, น-อารมฺมเณ ตีณิ, น-อธิปติยา นว, น-อนนฺตเร ตีณิ, นสมนนฺตเร ตีณิ, น-อญฺญมญฺเญ ตีณิ, น-อุปนิสฺสเย ตีณิ, นปุเรชาเต นว, นปจฺฉาชาเต นว, น-อาเสวเน นว, นกมฺเม ตีณิ, นวิปาเก นว, น-อาหาเร เอกํ, น-อินฺทฺริเย เอกํ, นฌาเน เอกํ, นมคฺเค เอกํ, นสมฺปยุตฺเต ตีณิ, นวิปฺปยุตฺเต นว, โนนตฺถิยา ตีณิ, โนวิคเต ตีณิ. (สํขิตฺตํ.)}
\csromanheader{2}{เหตุปจฺจยา น-อารมฺมเณ ตีณิ. (สํขิตฺตํ.)}
\csromanheader{2}{\textbf{นเหตุ}ปจฺจยา อารมฺมเณ เอกํ. (สํขิตฺตํ.)}
\prose{\csromanfolio{236}(สหชาตวาโรปิ ปจฺจยวาโรปิ นิสฺสยวาโรปิ สํสฏฺฐวาโรปิ สมฺปยุตฺตวาโรปิ ปฏิจฺจวารสทิสา วิตฺถาเรตพฺพา.)}

\csromanheader{2}{2. อสํกิลิฏฺฐสํกิเลสิกปท 7. ปญฺหาวาร}
\csromanheader{2}{\textbf{ปจฺจยจตุกฺก-เหตฺวาทิ}}
\proseitem{198}{เหตุ อสํกิลิฏฺฐสํกิเลสิโก ธมฺโม เหตุสฺส อสํกิลิฏฺฐสํกิเลสิกสฺส ธมฺมสฺส เหตุปจฺจเยน ปจฺจโย.\csromanspacer{} ตีณิ.}

\prose{เหตุ อสํกิลิฏฺฐสํกิเลสิโก ธมฺโม เหตุสฺส อสํกิลิฏฺฐสํกิเลสิกสฺส ธมฺมสฺส อารมฺมณปจฺจเยน ปจฺจโย.\csromanspacer{} นว.}

\prose{เหตุ อสํกิลิฏฺฐสํกิเลสิโก ธมฺโม เหตุสฺส อสํกิลิฏฺฐสํกิเลสิกสฺส ธมฺมสฺส อธิปติปจฺจเยน ปจฺจโย.\csromanspacer{} อารมฺมณาธิปติ, สหชาตาธิปติ.\csromanspacer{} ตีณิ.}

\vspace{\tipitakaparskip}
\csromancenter{\textbf{นเหตุ} อสํกิลิฏฺฐสํกิเลสิโก ธมฺโม นเหตุสฺส อสํกิลิฏฺฐสํกิเลสิกสฺส ธมฺมสฺส อธิปติปจฺจเยน ปจฺจโย.\csromanspacer{} อารมฺมณาธิปติ, สหชาตาธิปติ.\csromanspacer{} ตีณิ.}
\csromanheader{2}{1. เหตุทุก 5. สํกิลิฏฺฐตฺติก}
\prose{\csromanfolio{237}เหตุ อสํกิลิฏฺฐสํกิเลสิโก จ นเหตุ อสํกิลิฏฺฐสํกิเลสิโก จ ธมฺมา เหตุสฺส อสํกิลิฏฺฐสํกิเลสิกสฺส ธมฺมสฺส อธิปติปจฺจเยน ปจฺจโย.\csromanspacer{} อารมฺมณาธิปติ.\csromanspacer{} ตีณิ. (สํขิตฺตํ.)}

\proseitem{199}{เหตุยา ตีณิ, อารมฺมเณ นว, อธิปติยา นว, อนนฺตเร นว, สมนนฺตเร นว, สหชาเต นว, อญฺญมญฺเญ นว, นิสฺสเย นว, อุปนิสฺสเย นว, ปุเรชาเต ตีณิ, ปจฺฉาชาเต ตีณิ, อาเสวเน นว, กมฺเม ตีณิ, วิปาเก นว, อาหาเร ตีณิ, อินฺทฺริเย นว, ฌาเน ตีณิ, มคฺเค นว, สมฺปยุตฺเต นว, วิปฺปยุตฺเต ปญฺจ, อตฺถิยา นว, นตฺถิยา นว, วิคเต นว, อวิคเต นว. (สํขิตฺตํ.)}

\csromanheader{2}{\textbf{ปจฺจนียุทฺธาร}}
\proseitem{200}{เหตุ อสํกิลิฏฺฐสํกิเลสิโก ธมฺโม เหตุสฺส อสํกิลิฏฺฐสํกิเลสิกสฺส ธมฺมสฺส อารมฺมณปจฺจเยน ปจฺจโย.\csromanspacer{} สหชาตปจฺจเยน ปจฺจโย.\csromanspacer{} อุปนิสฺสยปจฺจเยน ปจฺจโย. (สํขิตฺตํ.)}

\proseitem{201}{นเหตุยา นว, น-อารมฺมเณ นว. (สํขิตฺตํ.)}

\csromanheader{2}{เหตุปจฺจยา น-อารมฺมเณ ตีณิ. (สํขิตฺตํ.)}
\csromanheader{2}{นเหตุปจฺจยา อารมฺมเณ นว. (สํขิตฺตํ.)}
\prose{(ยถา กุสลตฺติเก ปญฺหาวารสฺส อนุโลมมฺปิ ปจฺจนียมฺปิ อนุโลมปจฺจนียมฺปิ ปจฺจนียานุโลมมฺปิ คณิตํ, เอวํ คเณตพฺพํ.)}

\csromanheader{2}{3. อสํกิลิฏฺฐ-อสํกิเลสิกปท 1. ปฏิจฺจวาร}
\csromanheader{2}{\textbf{ปจฺจยจตุกฺก-เหตุ}}
\proseitem{202}{เหตุํ อสํกิลิฏฺฐ-อสํกิเลสิกํ ธมฺมํ ปฏิจฺจ เหตุ อสํกิลิฏฺฐ-อสํกิเลสิโก ธมฺโม อุปฺปชฺชติ เหตุปจฺจยา. (สํขิตฺตํ.)}

\proseitem{203}{เหตุยา นว, อารมฺมเณ นว, อธิปติยา นว, อนนฺตเร นว, สมนนฺตเร นว, สหชาเต นว, อญฺญมญฺเญ นว, นิสฺสเย นว,}

\vspace{\tipitakaparskip}
\csromancenter{ธมฺมานุโลม ทุกติกปฏฺฐานปาฬิ อุปนิสฺสเย นว, ปุเรชาเต นว, อาเสวเน นว, กมฺเม นว, วิปาเก นว, อาหาเร นว, อินฺทฺริเย นว, ฌาเน นว, มคฺเค นว, สมฺปยุตฺเต นว, วิปฺปยุตฺเต นว, อตฺถิยา นว, นตฺถิยา นว, วิคเต นว, อวิคเต นว.}
\csromanheader{2}{\textbf{น-อธิปติ}}
\proseitem{204}{\csromanfolio{238}เหตุํ อสํกิลิฏฺฐ-อสํกิเลสิกํ ธมฺมํ ปฏิจฺจ เหตุ อสํกิลิฏฺฐ-อสํกิเลสิโก ธมฺโม อุปฺปชฺชติ น-อธิปติปจฺจยา.\csromanspacer{}\csromanspacer{}เหตุํ อสํกิลิฏฺฐ-อสํกิเลสิกํ ธมฺมํ ปฏิจฺจ นเหตุ อสํกิลิฏฺฐ-อสํกิเลสิโก ธมฺโม อุปฺปชฺชติ น-อธิปติปจฺจยา. (2)}

\prose{นเหตุํ อสํกิลิฏฺฐ-อสํกิเลสิกํ ธมฺมํ ปฏิจฺจ นเหตุ อสํกิลิฏฺฐ-อสํกิเลสิโก ธมฺโม อุปฺปชฺชติ น-อธิปติปจฺจยา.\csromanspacer{}\csromanspacer{}นเหตุํ อสํกิลิฏฺฐ-อสํกิเลสิกํ ธมฺมํ ปฏิจฺจ เหตุ อสํกิลิฏฺฐ-อสํกิเลสิโก ธมฺโม อุปฺปชฺชติ น-อธิปติปจฺจยา. (2) (สํขิตฺตํ.)}

\vspace{\tipitakaparskip}
\csromancenter{\textbf{น-อธิปติ}ยา ฉ, นปุเรชาเต นว, นปจฺฉาชาเต นว, น-อาเสวเน นว, นกมฺเม ตีณิ, นวิปาเก นว, นวิปฺปยุตฺเต นว. (สํขิตฺตํ.)}
\csromanheader{2}{เหตุปจฺจยา น-อธิปติยา ฉ. (สํขิตฺตํ.)}
\csromanheader{2}{\textbf{น-อธิปติ}ปจฺจยา เหตุยา ฉ. (สํขิตฺตํ.)}
\prose{(สหชาตวาโรปิ ปจฺจยวาโรปิ นิสฺสยวาโรปิ สํสฏฺฐวาโรปิ สมฺปยุตฺตวาโรปิ ปฏิจฺจวารสทิสา วิตฺถาเรตพฺพา.)}

\csromanheader{2}{3. อสํกิลิฏฺฐ-อสํกิเลสิกปท 7. ปญฺหาวาร}
\csromanheader{2}{\textbf{ปจฺจยจตุกฺก-เหตุ อารมฺมณ}}
\proseitem{206}{เหตุ อสํกิลิฏฺฐ-อสํกิเลสิโก ธมฺโม เหตุสฺส อสํกิลิฏฺฐ-อสํกิเลสิกสฺส ธมฺมสฺส เหตุปจฺจเยน ปจฺจโย.\csromanspacer{} ตีณิ.}

\prose{นเหตุ อสํกิลิฏฺฐ-อสํกิเลสิโก ธมฺโม นเหตุสฺส อสํกิลิฏฺฐ-อสํกิเลสิกสฺส ธมฺมสฺส อารมฺมณปจฺจเยน ปจฺจโย.\csromanspacer{} ตีณิ. (สํขิตฺตํ.)}

\csromanmatikaanchor{mtk.37}
\csromantocmarkref{section}{6. วิตกฺกตฺติก}{mtk.37}
\bookmark[dest={mtk.37},level=1]{6. วิตกฺกตฺติก}
\csromanheader{1}{1. เหตุทุก 6. วิตกฺกตฺติก}
\proseitem{207}{\csromanfolio{239}เหตุยา ตีณิ, อารมฺมเณ ตีณิ, อธิปติยา ฉ, อนนฺตเร นว, สมนฺตเร นว, สหชาเต นว, อญฺญมญฺเญ นว, นิสฺสเย นว, อุปนิสฺสเย นว, กมฺเม ตีณิ, วิปาเก นว, อาหาเร ตีณิ, อินฺทฺริเย นว, ฌาเน ตีณิ, มคฺเค นว, สมฺปยุตฺเต นว, อตฺถิยา นว, นตฺถิยา นว, วิคเต นว, อวิคเต นว. (สํขิตฺตํ.)}

\csromanheader{2}{\textbf{ปจฺจนียุทฺธาร}}
\proseitem{208}{เหตุ อสํกิลิฏฺฐ-อสํกิเลสิโก ธมฺโม เหตุสฺส อสํกิลิฏฺฐ-อสํกิเลสิกสฺส ธมฺมสฺส สหชาตปจฺจเยน ปจฺจโย.\csromanspacer{} อุปนิสฺสยปจฺจเยน ปจฺจโย. (สํขิตฺตํ.)}

\proseitem{209}{นเหตุยา นว, น-อารมฺมเณ นว. (สํขิตฺตํ.)}

\csromanheader{2}{เหตุปจฺจยา น-อารมฺมเณ ตีณิ. (สํขิตฺตํ.)}
\csromanheader{2}{นเหตุปจฺจยา อารมฺมเณ ตีณิ. (สํขิตฺตํ.)}
\prose{(ยถา กุสลตฺติเก ปญฺหาวารสฺส อนุโลมมฺปิ ปจฺจนียมฺปิ อนุโลมปจฺจนียมฺปิ ปจฺจนียานุโลมมฺปิ คณิตํ, เอวํ คเณตพฺพํ.)}

\nitthitam{เหตุทุกสํกิลิฏฺฐตฺติกํ นิฏฺฐิตํ.}
\csromanheader{2}{1. เหตุทุก 6. วิตกฺกตฺติก}
\csromanheader{2}{1. สวิตกฺกสวิจารปท 1. ปฏิจฺจวาร}
\csromanheader{2}{\textbf{ปจฺจยจตุกฺก-เหตุ}}
\proseitem{210}{เหตุํ สวิตกฺกสวิจารํ ธมฺมํ ปฏิจฺจ เหตุ สวิตกฺกสวิจาโร ธมฺโม อุปฺปชฺชติ เหตุปจฺจยา.\csromanspacer{}\csromanspacer{}เหตุํ สวิตกฺกสวิจารํ ธมฺมํ ปฏิจฺจ นเหตุ สวิตกฺกสวิจาโร ธมฺโม อุปฺปชฺชติ เหตุปจฺจยา.\csromanspacer{}\csromanspacer{}เหตุํ สวิตกฺกสวิจารํ ธมฺมํ ปฏิจฺจ เหตุ สวิตกฺกสวิจาโร จ นเหตุ สวิตกฺกสวิจาโร จ ธมฺมา อุปฺปชฺชนฺติ เหตุปจฺจยา. (3)}

\prose{นเหตุํ สวิตกฺกสวิจารํ ธมฺมํ ปฏิจฺจ นเหตุ สวิตกฺกสวิจาโร ธมฺโม อุปฺปชฺชติ เหตุปจฺจยา.\csromanspacer{} ตีณิ.}

\gambhira{ธมฺมานุโลม ทุกติกปฏฺฐานปาฬิ}
\prose{\csromanfolio{240}เหตุํ สวิตกฺกสวิจารญฺจ นเหตุํ สวิตกฺกสวิจารญฺจ ธมฺมํ ปฏิจฺจ เหตุ สวิตกฺกสวิจาโร ธมฺโม อุปฺปชฺชติ เหตุปจฺจยา.\csromanspacer{} ตีณิ. (สํขิตฺตํ.)}

\proseitem{211}{เหตุยา นว, อารมฺมเณ นว, อธิปติยา นว, อนนฺตเร นว, สมนนฺตเร นว ฯเปฯ กมฺเม นว, วิปาเก นว ฯเปฯ อวิคเต นว. (สํขิตฺตํ.)}

\csromanheader{2}{\textbf{นเหตุน-อธิปติ}}
\proseitem{212}{นเหตุํ สวิตกฺกสวิจารํ ธมฺมํ ปฏิจฺจ นเหตุ สวิตกฺกสวิจาโร ธมฺโม อุปฺปชฺชติ นเหตุปจฺจยา.\csromanspacer{} เทฺว.}

\prose{เหตุํ สวิตกฺกสวิจารํ ธมฺมํ ปฏิจฺจ เหตุ สวิตกฺกสวิจาโร ธมฺโม อุปฺปชฺชติ น-อธิปติปจฺจยา. (สํขิตฺตํ.)}

\proseitem{213}{นเหตุยา เทฺว, น-อธิปติยา นว, นปุเรชาเต นว, นปจฺฉาชาเต นว, น-อาเสวเน นว, นกมฺเม ตีณิ, นวิปาเก นว, นมคฺเค เอกํ, นวิปฺปยุตฺเต นว. (สํขิตฺตํ.)}

\csromanheader{2}{เหตุปจฺจยา น-อธิปติยา นว. (สํขิตฺตํ.)}
\csromanheader{2}{นเหตุปจฺจยา อารมฺมเณ เทฺว. (สํขิตฺตํ.)}
\prose{(สหชาตวาโรปิ ปจฺจยวาโรปิ นิสฺสยวาโรปิ สํสฏฺฐวาโรปิ สมฺปยุตฺตวาโรปิ ปฏิจฺจวารสทิสา วิตฺถาเรตพฺพา.)}

\csromanheader{2}{1. สวิตกฺกสวิจารปท 7. ปญฺหาวาร}
\csromanheader{2}{\textbf{ปจฺจยจตุกฺก-เหตฺวาทิ}}
\proseitem{214}{เหตุ สวิตกฺกสวิจาโร ธมฺโม เหตุสฺส สวิตกฺกสวิจารสฺส ธมฺมสฺส เหตุปจฺจเยน ปจฺจโย.\csromanspacer{} ตีณิ.}

\prose{เหตุ สวิตกฺกสวิจาโร ธมฺโม เหตุสฺส สวิตกฺกสวิจารสฺส ธมฺมสฺส อารมฺมณปจฺจเยน ปจฺจโย.\csromanspacer{} นว.}

\prose{เหตุ สวิตกฺกสวิจาโร ธมฺโม เหตุสฺส สวิตกฺกสวิจารสฺส ธมฺมสฺส อธิปติปจฺจเยน ปจฺจโย.\csromanspacer{} อารมฺมณาธิปติ, สหชาตาธิปติ.\csromanspacer{} ตีณิ.}

\csromanheader{2}{1. เหตุทุก 6. วิตกฺกตฺติก}
\prose{\csromanfolio{241}นเหตุ สวิตกฺกสวิจาโร ธมฺโม นเหตุสฺส สวิตกฺกสวิจารสฺส ธมฺมสฺส อธิปติปจฺจเยน ปจฺจโย.\csromanspacer{} อารมฺมณาธิปติ, สหชาตาธิปติ.\csromanspacer{} ตีณิ.}

\prose{เหตุ สวิตกฺกสวิจาโร จ นเหตุ สวิตกฺกสวิจาโร จ ธมฺมา เหตุสฺส สวิตกฺกสวิจารสฺส ธมฺมสฺส อธิปติปจฺจเยน ปจฺจโย.\csromanspacer{} อารมฺมณาธิปติ.\csromanspacer{} ตีณิ. (สํขิตฺตํ.)}

\proseitem{215}{เหตุยา ตีณิ, อารมฺมเณ นว, อธิปติยา นว, อนนฺตเร นว, สมนนฺตเร นว, สหชาเต นว, อญฺญมญฺเญ นว, นิสฺสเย นว, อุปนิสฺสเย นว, อาเสวเน นว, กมฺเม ตีณิ, วิปาเก นว, อาหาเร ตีณิ, อินฺทฺริเย นว, ฌาเน ตีณิ, มคฺเค นว, สมฺปยุตฺเต นว, อตฺถิยา นว, นตฺถิยา นว, วิคเต นว, อวิคเต นว. (สํขิตฺตํ.)}

\csromanheader{2}{\textbf{ปจฺจนียุทฺธาร}}
\proseitem{216}{เหตุ สวิตกฺกสวิจาโร ธมฺโม เหตุสฺส สวิตกฺกสวิจารสฺส ธมฺมสฺส อารมฺมณปจฺจเยน ปจฺจโย.\csromanspacer{} สหชาตปจฺจเยน ปจฺจโย.\csromanspacer{} อุปนิสฺสยปจฺจเยน ปจฺจโย. (สํขิตฺตํ.)}

\proseitem{217}{นเหตุยา นว, น-อารมฺมเณ นว. (สํขิตฺตํ.)}

\csromanheader{2}{เหตุปจฺจยา น-อารมฺมเณ ตีณิ. (สํขิตฺตํ.)}
\csromanheader{2}{นเหตุปจฺจยา อารมฺมเณ นว. (สํขิตฺตํ.)}
\prose{(ยถา กุสลตฺติเก ปญฺหาวารสฺส อนุโลมมฺปิ ปจฺจนียมฺปิ อนุโลมปจฺจนียมฺปิ ปจฺจนียานุโลมมฺปิ คณิตํ, เอวํ คเณตพฺพํ.)}

\csromanheader{2}{2. อวิตกฺกวิจารมตฺตปท 1. ปฏิจฺจวาร}
\csromanheader{2}{\textbf{ปจฺจยจตุกฺก-เหตุ}}
\proseitem{218}{เหตุํ อวิตกฺกวิจารมตฺตํ ธมฺมํ ปฏิจฺจ เหตุ อวิตกฺกวิจารมตฺโต ธมฺโม อุปฺปชฺชติ เหตุปจฺจยา.\csromanspacer{} ตีณิ.}

\gambhira{ธมฺมานุโลม ทุกติกปฏฺฐานปาฬิ}
\prose{\csromanfolio{242}นเหตุํ อวิตกฺกวิจารมตฺตํ ธมฺมํ ปฏิจฺจ นเหตุ อวิตกฺกวิจารมตฺโต ธมฺโม อุปฺปชฺชติ เหตุปจฺจยา.\csromanspacer{} ตีณิ.}

\prose{เหตุํ อวิตกฺกวิจารมตฺตญฺจ นเหตุํ อวิตกฺกวิจารมตฺตญฺจ ธมฺมํ ปฏิจฺจ เหตุ อวิตกฺกวิจารมตฺโต ธมฺโม อุปฺปชฺชติ เหตุปจฺจยา.\csromanspacer{} ตีณิ. (สํขิตฺตํ.)}

\proseitem{219}{เหตุยา นว, อารมฺมเณ นว, อธิปติยา นว ฯเปฯ กมฺเม นว, วิปาเก นว ฯเปฯ อวิคเต นว. (สํขิตฺตํ.)}

\csromanheader{2}{\textbf{น-อธิปติ}}
\proseitem{220}{เหตุํ อวิตกฺกวิจารมตฺตํ ธมฺมํ ปฏิจฺจ เหตุ อวิตกฺกวิจารมตฺโต ธมฺโม อุปฺปชฺชติ น-อธิปติปจฺจยา. (สํขิตฺตํ.)}

\vspace{\tipitakaparskip}
\csromancenter{\textbf{น-อธิปติ}ยา นว, นปุเรชาเต นว, นปจฺฉาชาเต นว, น-อาเสวเน นว, นกมฺเม ตีณิ, นวิปาเก นว, นวิปฺปยุตฺเต นว. (สํขิตฺตํ.)}
\csromanheader{2}{เหตุปจฺจยา น-อธิปติยา นว. (สํขิตฺตํ.)}
\csromanheader{2}{\textbf{น-อธิปติ}ปจฺจยา เหตุยา นว. (สํขิตฺตํ.)}
\prose{(สหชาตวาโรปิ ปจฺจยวาโรปิ นิสฺสยวาโรปิ สํสฏฺฐวาโรปิ สมฺปยุตฺตวาโรปิ ปฏิจฺจวารสทิสา วิตฺถาเรตพฺพา.)}

\csromanheader{2}{2. อวิตกฺกวิจารมตฺตปท 7. ปญฺหาวาร}
\csromanheader{2}{\textbf{ปจฺจยจตุกฺก-เหตฺวาทิ}}
\proseitem{222}{เหตุ อวิตกฺกวิจารมตฺโต ธมฺโม เหตุสฺส อวิตกฺกวิจารมตฺตสฺส ธมฺมสฺส เหตุปจฺจเยน ปจฺจโย.\csromanspacer{} ตีณิ.}

\prose{เหตุ อวิตกฺกวิจารมตฺโต ธมฺโม นเหตุสฺส อวิตกฺกวิจารมตฺตสฺส ธมฺมสฺส อารมฺมณปจฺจเยน ปจฺจโย.\csromanspacer{} ตีณิ.}

\prose{เหตุ อวิตกฺกวิจารมตฺโต ธมฺโม เหตุสฺส อวิตกฺกวิจารมตฺตสฺส ธมฺมสฺส อธิปติปจฺจเยน ปจฺจโย.\csromanspacer{} สหชาตาธิปติ.\csromanspacer{} ตีณิ.}

\csromanheader{2}{1. เหตุทุก 6. วิตกฺกตฺติก}
\prose{\csromanfolio{243}นเหตุ อวิตกฺกวิจารมตฺโต ธมฺโม นเหตุสฺส อวิตกฺกวิจารมตฺตสฺส ธมฺมสฺส อธิปติปจฺจเยน ปจฺจโย.\csromanspacer{} อารมฺมณาธิปติ, สหชาตาธิปติ.\csromanspacer{} ตีณิ.}

\prose{เหตุ อวิตกฺกวิจารมตฺโต จ นเหตุ อวิตกฺกวิจารมตฺโต จ ธมฺมา นเหตุสฺส อวิตกฺกวิจารมตฺตสฺส ธมฺมสฺส อธิปติปจฺจเยน ปจฺจโย.\csromanspacer{} อารมฺมณาธิปติ. (สํขิตฺตํ.)}

\proseitem{223}{เหตุยา ตีณิ, อารมฺมเณ ตีณิ, อธิปติยา สตฺต, อนนฺตเร นว, สมนนฺตเร นว, สหชาเต นว, อญฺญมญฺเญ นว, นิสฺสเย นว, อุปนิสฺสเย นว, อาเสวเน นว, กมฺเม ตีณิ, วิปาเก นว, อาหาเร ตีณิ, อินฺทฺริเย นว, ฌาเน ตีณิ, มคฺเค นว, สมฺปยุตฺเต นว, อตฺถิยา นว, นตฺถิยา นว, วิคเต นว, อวิคเต นว. (สํขิตฺตํ.)}

\csromanheader{2}{\textbf{ปจฺจนียุทฺธาร}}
\proseitem{224}{เหตุ อวิตกฺกวิจารมตฺโต ธมฺโม เหตุสฺส อวิตกฺกวิจารมตฺตสฺส ธมฺมสฺส สหชาตปจฺจเยน ปจฺจโย.\csromanspacer{} อุปนิสฺสยปจฺจเยน ปจฺจโย. (สํขิตฺตํ.)}

\proseitem{225}{นเหตุยา นว, น-อารมฺมเณ นว. (สํขิตฺตํ.)}

\csromanheader{2}{เหตุปจฺจยา น-อารมฺมเณ ตีณิ. (สํขิตฺตํ.)}
\csromanheader{2}{นเหตุปจฺจยา อารมฺมเณ ตีณิ. (สํขิตฺตํ.)}
\prose{(ยถา กุสลตฺติเก ปญฺหาวารสฺส อนุโลมมฺปิ ปจฺจนียมฺปิ อนุโลมปจฺจนียมฺปิ ปจฺจนียานุโลมมฺปิ คณิตํ, เอวํ คเณตพฺพํ.)}

\csromanheader{2}{3. อวิตกฺก-อวิจารปท 1. ปฏิจฺจวาร}
\csromanheader{2}{\textbf{ปจฺจยจตุกฺก-เหตุ}}
\proseitem{226}{เหตุํ อวิตกฺก-อวิจารํ ธมฺมํ ปฏิจฺจ เหตุ อวิตกฺก-อวิจาโร ธมฺโม อุปฺปชฺชติ เหตุปจฺจยา. (สํขิตฺตํ.)}

\gambhira{ธมฺมานุโลม ทุกติกปฏฺฐานปาฬิ}
\proseitem{227}{\csromanfolio{244}เหตุยา นว, อารมฺมเณ นว, อธิปติยา นว ฯเปฯ กมฺเม นว, วิปาเก นว ฯเปฯ วิคเต นว, อวิคเต นว. (สํขิตฺตํ.)}

\csromanheader{2}{\textbf{นเหตฺวาทิ}}
\proseitem{228}{นเหตุํ อวิตกฺก-อวิจารํ ธมฺมํ ปฏิจฺจ นเหตุ อวิตกฺก-อวิจาโร ธมฺโม อุปฺปชฺชติ นเหตุปจฺจยา. (1)}

\prose{เหตุํ อวิตกฺก-อวิจารํ ธมฺมํ ปฏิจฺจ นเหตุ อวิตกฺก-อวิจาโร ธมฺโม อุปฺปชฺชติ น-อารมฺมณปจฺจยา. (1)}

\prose{นเหตุํ อวิตกฺก-อวิจารํ ธมฺมํ ปฏิจฺจ นเหตุ อวิตกฺก-อวิจาโร ธมฺโม อุปฺปชฺชติ น-อารมฺมณปจฺจยา. (1)}

\prose{เหตุํ อวิตกฺก-อวิจารญฺจ นเหตุํ อวิตกฺก-อวิจารญฺจ ธมฺมํ ปฏิจฺจ นเหตุ อวิตกฺก-อวิจาโร ธมฺโม อุปฺปชฺชติ น-อารมฺมณปจฺจยา. (1)}

\prose{เหตุํ อวิตกฺก-อวิจารํ ธมฺมํ ปฏิจฺจ เหตุ อวิตกฺก-อวิจาโร ธมฺโม อุปฺปชฺชติ น-อธิปติปจฺจยา.\csromanspacer{} นว. (สํขิตฺตํ.)}

\proseitem{229}{นเหตุ เอกํ, น-อารมฺมเณ ตีณิ, น-อธิปติยา นว, น-อนนฺตเร ตีณิ, นสมนนฺตเร ตีณิ, น-อญฺญมญฺเญ ตีณิ, น-อุปนิสฺสเย ตีณิ, นปุเรชาเต นว, นปจฺฉาชาเต นว, น-อาเสวเน นว, นกมฺเม ตีณิ, นวิปาเก นว, น-อาหาเร เอกํ, น-อินฺทฺริเย เอกํ, นฌาเน เอกํ, นมคฺเค เอกํ, นสมฺปยุตฺเต ตีณิ, นวิปฺปยุตฺเต นว, โนนตฺถิยา ตีณิ, โนวิคเต ตีณิ. (สํขิตฺตํ.)}

\csromanheader{2}{เหตุปจฺจยา น-อารมฺมเณ ตีณิ. (สํขิตฺตํ.)}
\csromanheader{2}{นเหตุปจฺจยา อารมฺมเณ เอกํ. (สํขิตฺตํ.)}
\prose{(สหชาตวาโรปิ ปจฺจยวาโรปิ นิสฺสยวาโรปิ สํสฏฺฐวาโรปิ สมฺปยุตฺตวาโรปิ ปฏิจฺจวารสทิสา วิตฺถาเรตพฺพา.)}

\vspace{\tipitakaparskip}
\csromancenter{เหตุทุก 6.\csromanspacer{} วิตกฺกตฺติก \textbf{3.}\csromanspacer{}\textbf{ อวิตกฺก-อวิจารปท 7.}\csromanspacer{}\textbf{ ปญฺหาวาร}}
\csromanheader{2}{\textbf{ปจฺจยจตุกฺก-เหตุ อารมฺมณ}}
\proseitem{230}{\csromanfolio{245}เหตุ อวิตกฺก-อวิจาโร ธมฺโม เหตุสฺส อวิตกฺก-อวิจารสฺส ธมฺมสฺส เหตุปจฺจเยน ปจฺจโย.\csromanspacer{} ตีณิ.}

\prose{เหตุ อวิตกฺก-อวิจาโร ธมฺโม เหตุสฺส อวิตกฺก-อวิจารสฺส ธมฺมสฺส อารมฺมณปจฺจเยน ปจฺจโย.\csromanspacer{} นว. (สํขิตฺตํ.)}

\proseitem{231}{เหตุยา ตีณิ, อารมฺมเณ นว, อธิปติยา ฉ, อนนฺตเร นว, สมนนฺตเร นว, สหชาเต นว, อญฺญมญฺเญ นว, นิสฺสเย นว, อุปนิสฺสเย นว, ปุเรชาเต ตีณิ, ปจฺฉาชาเต ตีณิ, อาเสวเน นว, กมฺเม ตีณิ, วิปาเก นว, อาหาเร ตีณิ, อินฺทฺริเย นว, ฌาเน ตีณิ, มคฺเค นว, สมฺปยุตฺเต นว, วิปฺปยุตฺเต ปญฺจ, อตฺถิยา นว, นตฺถิยา นว, วิคเต นว, อวิคเต นว. (สํขิตฺตํ.)}

\csromanheader{2}{\textbf{ปจฺจนียุทฺธาร}}
\proseitem{232}{เหตุ อวิตกฺก-อวิจาโร ธมฺโม เหตุสฺส อวิตกฺก-อวิจารสฺส ธมฺมสฺส อารมฺมณปจฺจเยน ปจฺจโย.\csromanspacer{} สหชาตปจฺจเยน ปจฺจโย.\csromanspacer{} อุปนิสฺสยปจฺจเยน ปจฺจโย. (สํขิตฺตํ.)}

\proseitem{233}{นเหตุยา นว, น-อารมฺมเณ นว, น-อธิปติยา นว. (สํขิตฺตํ.)}

\csromanheader{2}{เหตุปจฺจยา น-อารมฺมเณ ตีณิ. (สํขิตฺตํ.)}
\csromanheader{2}{นเหตุปจฺจยา อารมฺมเณ นว. (สํขิตฺตํ.)}
\prose{(ยถา กุสลตฺติเก ปญฺหาวารสฺส อนิโลมมฺปิ ปจฺจนียมฺปิ อนุโลมปจฺจนียมฺปิ ปจฺจนียานุโลมมฺปิ คณิตํ, เอวํ คเณตพฺพํ.)}

\nitthitam{เหตุทุกวิตกฺกตฺติกํ นิฏฺฐิตํ.}
\csromanmatikaanchor{mtk.38}
\csromantocmarkref{section}{7. ปีติตฺติก}{mtk.38}
\bookmark[dest={mtk.38},level=1]{7. ปีติตฺติก}
\csromanheader{1}{ธมฺมานุโลม ทุกติกปฏฺฐานปาฬิ \textbf{1. เหตุทุก 7. ปีติตฺติก}}
\csromanheader{2}{1. ปีติสหคตปท 1. ปฏิจฺจวาร}
\csromanheader{2}{\textbf{ปจฺจยจตุกฺก-เหตุ}}
\proseitem{234}{\csromanfolio{246}เหตุํ ปีติสหคตํ ธมฺมํ ปฏิจฺจ เหตุ ปีติสหคโต ธมฺโม อุปฺปชฺชติ เหตุปจฺจยา.\csromanspacer{} ตีณิ.}

\prose{นเหตุํ ปีติสหคตํ ธมฺมํ ปฏิจฺจ นเหตุ ปีติสหคโต ธมฺโม อุปฺปชฺชติ เหตุปจฺจยา. (สํขิตฺตํ.)}

\proseitem{235}{เหตุยา นว, อารมฺมเณ นว, อธิปติยา นว ฯเปฯ กมฺเม นว, วิปาเก นว, อาหาเร นว ฯเปฯ อวิคเต นว. (สํขิตฺตํ.)}

\csromanheader{2}{\textbf{นเหตุ น-อธิปติ}}
\proseitem{236}{นเหตุํ ปีติสหคตํ ธมฺมํ ปฏิจฺจ นเหตุ ปีติสหคโต ธมฺโม อุปฺปชฺชติ นเหตุปจฺจยา.}

\prose{เหตุํ ปีติสหคตํ ธมฺมํ ปฏิจฺจ เหตุ ปีติสหคโต ธมฺโม อุปฺปชฺชติ น-อธิปติปจฺจยา. (สํขิตฺตํ.)}

\proseitem{237}{นเหตุยา เอกํ, น-อธิปติยา นว, นปุเรชาเต นว, นปจฺฉาชาเต นว, น-อาเสวเน นว, นกมฺเม ตีณิ, นวิปาเก นว, นมคฺเค เอกํ, นวิปฺปยุตฺเต นว. (สํขิตฺตํ.)}

\csromanheader{2}{เหตุปจฺจยา น-อธิปติยา นว. (สํขิตฺตํ.)}
\csromanheader{2}{นเหตุปจฺจยา อารมฺมเณ เอกํ. (สํขิตฺตํ.)}
\prose{(สหชาตวาโรปิ ปจฺจยวาโรปิ นิสฺสยวาโรปิ สํสฏฺฐวาโรปิ สมฺปยุตฺตวาโรปิ ปฏิจฺจวารสทิสา วิตฺถาเรตพฺพา.)}

\csromanheader{2}{1. เหตุทุก 7. ปีติตฺติก \textbf{1. ปีติสหคตปท 7. ปญฺหาวาร}}
\csromanheader{2}{\textbf{ปจฺจยจตุกฺก-เหตฺวาทิ}}
\proseitem{238}{\csromanfolio{247}เหตุ ปีติสหคโต ธมฺโม เหตุสฺส ปีติสหคตสฺส ธมฺมสฺส เหตุปจฺจเยน ปจฺจโย.\csromanspacer{} ตีณิ.}

\prose{เหตุ ปีติสหคโต ธมฺโม เหตุสฺส ปีติสหคตสฺส ธมฺมสฺส อารมฺมณปจฺจเยน ปจฺจโย.\csromanspacer{} ตีณิ.}

\prose{นเหตุ ปีติสหคโต ธมฺโม นเหตุสฺส ปีติสหคตสฺส ธมฺมสฺส อารมฺมณปจฺจเยน ปจฺจโย.\csromanspacer{} ตีณิ.}

\prose{เหตุ ปีติสหคโต จ นเหตุ ปีติสหคโต จ ธมฺมา เหตุสฺส ปีติสหคตสฺส ธมฺมสฺส อารมฺมณปจฺจเยน ปจฺจโย.\csromanspacer{} ตีณิ.}

\prose{เหตุ ปีติสหคโต ธมฺโม เหตุสฺส ปีติสหคตสฺส ธมฺมสฺส อธิปติปจฺจเยน ปจฺจโย.\csromanspacer{} อารมฺมณาธิปติ, สหชาตาธิปติ.\csromanspacer{} ตีณิ.}

\prose{นเหตุ ปีติสหคโต ธมฺโม นเหตุสฺส ปีติสหคตสฺส ธมฺมสฺส อธิปติปจฺจเยน ปจฺจโย.\csromanspacer{} อารมฺมณาธิปติ, สหชาตาธิปติ.\csromanspacer{} ตีณิ.}

\prose{เหตุ ปีติสหคโต จ นเหตุ ปีติสหคโต จ ธมฺมา เหตุสฺส ปีติสหคตสฺส ธมฺมสฺส อธิปติปจฺจเยน ปจฺจโย.\csromanspacer{} อารมฺมณาธิปติ.\csromanspacer{} ตีณิ. (สํขิตฺตํ.)}

\proseitem{239}{เหตุยา ตีณิ, อารมฺมเณ นว, อธิปติยา นว, อนนฺตเร นว, สมนนฺตเร นว, สหชาเต นว, อญฺญมญฺเญ นว, นิสฺสเย นว, อุปนิสฺสเย นว, อาเสวเน นว, กมฺเม ตีณิ, วิปาเก นว, อาหาเร ตีณิ, อินฺทฺริเย นว, ฌาเน ตีณิ, มคฺเค นว, สมฺปยุตฺเต นว, อตฺถิยา นว, นตฺถิยา นว, วิคเต นว, อวิคเต นว. (สํขิตฺตํ.)}

\csromanheader{2}{\textbf{ปจฺจนียุทฺธาร}}
\proseitem{240}{เหตุ ปีติสหคโต ธมฺโม เหตุสฺส ปีติสหคตสฺส ธมฺมสฺส อารมฺมณปจฺจเยน ปจฺจโย.\csromanspacer{} สหชาตปจฺจเยน ปจฺจโย.\csromanspacer{} อุปนิสฺสยปจฺจเยน ปจฺจโย. (สํขิตฺตํ.)}

\gambhira{ธมฺมานุโลม ทุกติกปฏฺฐานปาฬิ}
\proseitem{241}{\csromanfolio{248}นเหตุยา นว, น-อารมฺมเณ นว. (สํขิตฺตํ.)}

\csromanheader{2}{เหตุปจฺจยา น-อารมฺมเณ ตีณิ. (สํขิตฺตํ.)}
\csromanheader{2}{นเหตุปจฺจยา อารมฺมเณ นว. (สํขิตฺตํ.)}
\prose{(ยถา กุสลตฺติเก ปญฺหาวารสฺส อนุโลมมฺปิ ปจฺจนียมฺปิ อนุโลมปจฺจนียมฺปิ ปจฺจนียานุโลมมฺปิ คณิตํ, เอวํ คเณตพฺพํ.)}

\csromanheader{2}{2. สุขสหคตปท 1. ปฏิจฺจวาร}
\csromanheader{2}{\textbf{ปจฺจยจตุกฺก-เหตุ}}
\proseitem{242}{เหตุํ สุขสหคตํ ธมฺมํ ปฏิจฺจ เหตุ สุขสหคโต ธมฺโม อุปฺปชฺชติ เหตุปจฺจยา.\csromanspacer{} ตีณิ.}

\prose{นเหตุํ สุขสหคตํ ธมฺมํ ปฏิจฺจ นเหตุ สุขสหคโต ธมฺโม อุปฺปชฺชติ เหตุปจฺจยา.\csromanspacer{} ตีณิ.}

\prose{เหตุํ สุขสหคตญฺจ นเหตุํ สุขสหคตญฺจ ธมฺมํ ปฏิจฺจ เหตุ สุขสหคโต ธมฺโม อุปฺปชฺชติ เหตุปจฺจยา.\csromanspacer{} ตีณิ. (สํขิตฺตํ.)}

\proseitem{243}{เหตุยา นว, อารมฺมเณ นว, อธิปติยา นว ฯเปฯ อวิคเต นว. (สํขิตฺตํ.)}

\csromanheader{2}{\textbf{นเหตุ น-อธิปติ}}
\proseitem{244}{นเหตุํ สุขสหคตํ ธมฺมํ ปฏิจฺจ นเหตุ สุขสหคโต ธมฺโม อุปฺปชฺชติ นเหตุปจฺจยา.}

\prose{เหตุํ สุขสหคตํ ธมฺมํ ปฏิจฺจ เหตุ สุขสหคโต ธมฺโม อุปฺปชฺชติ น-อธิปติปจฺจยา.\csromanspacer{} นว. (สํขิตฺตํ.)}

\proseitem{245}{นเหตุยา เอกํ, น-อธิปติยา นว, นปุเรชาเต นว, นปจฺฉาชาเต นว, น-อาเสวเน นว, นกมฺเม ตีณิ, นวิปาเก นว, นฌาเน เอกํ, นมคฺเค เอกํ, นวิปฺปยุตฺเต นว. (สํขิตฺตํ.)}

\csromanheader{2}{1. เหตุทุก 7. ปีติตฺติก}
\csromanheader{2}{เหตุปจฺจยา น-อธิปติยา นว. (สํขิตฺตํ.)}
\csromanheader{2}{นเหตุปจฺจยา อารมฺมเณ เอกํ. (สํขิตฺตํ.)}
\prose{\csromanfolio{249}(สหชาตวาโรปิ ปจฺจยวาโรปิ นิสฺสยวาโรปิ สํสฏฺฐวาโรปิ สมฺปยุตฺตวาโรปิ ปฏิจฺจวารสทิสา วิตฺถาเรตพฺพา.)}

\csromanheader{2}{2. สุขสหคตปท 7. ปญฺหาวาร}
\csromanheader{2}{\textbf{ปจฺจยจตุกฺก-เหตฺวาทิ}}
\proseitem{246}{เหตุ สุขสหคโต ธมฺโม เหตุสฺส สุขสหคตสฺส ธมฺมสฺส เหตุปจฺจเยน ปจฺจโย.\csromanspacer{} ตีณิ.}

\prose{เหตุ สุขสหคโต ธมฺโม เหตุสฺส สุขสหคตสฺส ธมฺมสฺส อารมฺมณปจฺจเยน ปจฺจโย.\csromanspacer{} นว.}

\prose{เหตุ สุขสหคโต ธมฺโม เหตุสฺส สุขสหคตสฺส ธมฺมสฺส อธิปติปจฺจเยน ปจฺจโย.\csromanspacer{} นว. (สํขิตฺตํ.)}

\proseitem{247}{เหตุยา ตีณิ, อารมฺมเณ นว, อธิปติยา นว, อนนฺตเร นว, สมนนฺตเร นว, สหชาเต นว, อญฺญมญฺเญ นว, นิสฺสเย นว, อุปนิสฺสเย นว, อาเสวเน นว, กมฺเม ตีณิ, วิปาเก นว, อาหาเร ตีณิ, อินฺทฺริเย นว, ฌาเน ตีณิ, มคฺเค นว, สมฺปยุตฺเต นว, อตฺถิยา นว, นตฺถิยา นว, วิคเต นว, อวิคเต นว. (สํขิตฺตํ.)}

\csromanheader{2}{\textbf{ปจฺจนียุทฺธาร}}
\proseitem{248}{เหตุ สุขสหคโต ธมฺโม เหตุสฺส สุขสหคตสฺส ธมฺมสฺส อารมฺมณปจฺจเยน ปจฺจโย.\csromanspacer{} สหชาตปจฺจเยน ปจฺจโย.\csromanspacer{} อุปนิสฺสยปจฺจเยน ปจฺจโย. (สํขิตฺตํ.)}

\proseitem{249}{นเหตุยา นว, น-อารมฺมเณ นว. (สํขิตฺตํ.)}

\csromanheader{2}{เหตุปจฺจยา น-อารมฺมเณ ตีณิ. (สํขิตฺตํ.)}
\csromanheader{2}{นเหตุปจฺจยา อารมฺมเณ นว. (สํขิตฺตํ.)}
\gambhira{ธมฺมานุโลม ทุกติกปฏฺฐานปาฬิ}
\prose{\csromanfolio{250}(ยถา กุสลตฺติเก ปญฺหาวารสฺส อนุโลมมฺปิ ปจฺจนียมฺปิ อนุโลมปจฺจนียมฺปิ ปจฺจนียานุโลมมฺปิ คณิตํ, เอวํ คเณตพฺพํ.)}

\csromanheader{2}{3. อุเปกฺขาสหคตปท 1. ปฏิจฺจวาร}
\csromanheader{2}{\textbf{ปจฺจยจตุกฺก-เหตุ}}
\proseitem{250}{เหตุํ อุเปกฺขาสหคตํ ธมฺมํ ปฏิจฺจ เหตุ อุเปกฺขาสหคโต ธมฺโม อุปฺปชฺชติ เหตุปจฺจยา.\csromanspacer{} ตีณิ.}

\prose{นเหตุํ อุเปกฺขาสหคตํ ธมฺมํ ปฏิจฺจ นเหตุ อุเปกฺขาสหคโต ธมฺโม อุปฺปชฺชติ เหตุปจฺจยา.\csromanspacer{} ตีณิ.}

\prose{เหตุํ อุเปกฺขาสหคตญฺจ นเหตุํ อุเปกฺขาสหคตญฺจ ธมฺมํ ปฏิจฺจ เหตุ อุเปกฺขาสหคโต ธมฺโม อุปฺปชฺชติ เหตุปจฺจยา.\csromanspacer{} ตีณิ. (สํขิตฺตํ.)}

\proseitem{251}{เหตุยา นว, อารมฺมเณ นว, อธิปติยา นว ฯเปฯ กมฺเม นว, วิปาเก นว ฯเปฯ อวิคเต นว. (สํขิตฺตํ.)}

\csromanheader{2}{\textbf{นเหตฺวาทิ}}
\proseitem{252}{นเหตุํ อุเปกฺขาสหคตํ ธมฺมํ ปฏิจฺจ นเหตุ อุเปกฺขาสหคโต ธมฺโม อุปฺปชฺชติ นเหตุปจฺจยา.\csromanspacer{} เทฺว.}

\prose{เหตุํ อุเปกฺขาสหคตํ ธมฺมํ ปฏิจฺจ เหตุ อุเปกฺขาสหคโต ธมฺโม อุปฺปชฺชติ น-อธิปติปจฺจยา.\csromanspacer{} นว.}

\prose{เหตุํ อุเปกฺขาสหคตํ ธมฺมํ ปฏิจฺจ เหตุ อุเปกฺขาสหคโต ธมฺโม อุปฺปชฺชติ นปุเรชาตปจฺจยา.\csromanspacer{} นว.}

\prose{เหตุํ อุเปกฺขาสหคตํ ธมฺมํ ปฏิจฺจ เหตุ อุเปกฺขาสหคโต ธมฺโม อุปฺปชฺชติ นปจฺฉาชาตปจฺจยา.\csromanspacer{} นว. (สํขิตฺตํ.)}

\proseitem{253}{นเหตุยา เทฺว, น-อธิปติยา นว, นปุเรชาเต นว, นปจฺฉาชาเต นว, น-อาเสวเน นว, นกมฺเม ตีณิ, นวิปาเก นว, นฌาเน เอกํ, นมคฺเค เอกํ, นวิปฺปยุตฺเต นว. (สํขิตฺตํ.)}

\csromanheader{2}{1. เหตุทุก 7. ปีติตฺติก}
\csromanheader{2}{เหตุปจฺจยา น-อธิปติยา นว. (สํขิตฺตํ.)}
\csromanheader{2}{นเหตุปจฺจยา อารมฺมเณ เทฺว. (สํขิตฺตํ.)}
\prose{\csromanfolio{251}(สหชาตวาโรปิ ปจฺจยวาโรปิ นิสฺสยวาโรปิ สํสฏฺฐวาโรปิ สมฺปยุตฺตวาโรปิ ปฏิจฺจวารสทิสา วิตฺถาเรตพฺพา.)}

\csromanheader{2}{3. อุเปกฺขาสหคตปท 7. ปญฺหาวาร}
\csromanheader{2}{\textbf{ปจฺจยจตุกฺก-เหตฺวาทิ}}
\proseitem{254}{เหตุ อุเปกฺขาสหคโต ธมฺโม เหตุสฺส อุเปกฺขาสหคตสฺส ธมฺมสฺส เหตุปจฺจเยน ปจฺจโย.\csromanspacer{} ตีณิ.}

\prose{เหตุ อุเปกฺขาสหคโต ธมฺโม เหตุสฺส อุเปกฺขาสหคตสฺส ธมฺมสฺส อารมฺมณปจฺจเยน ปจฺจโย.\csromanspacer{} นว.}

\prose{เหตุ อุเปกฺขาสหคโต ธมฺโม เหตุสฺส อุเปกฺขาสหคตสฺส ธมฺมสฺส อธิปติปจฺจเยน ปจฺจโย.\csromanspacer{} นว. (สํขิตฺตํ.)}

\proseitem{255}{เหตุยา ตีณิ, อารมฺมเณ นว, อธิปติยา นว, อนนฺตเร นว, สมนนฺตเร นว, สหชาเต นว, อญฺญมญฺเญ นว, นิสฺสเย นว, อุปนิสฺสเย นว, อาเสวเน นว, กมฺเม ตีณิ, วิปาเก นว, อาหาเร ตีณิ, อินฺทฺริเร นว, ฌาเน ตีณิ, มคฺเค นว, สมฺปยุตฺเต นว, อตฺถิยา นว, นตฺถิยา นว, วิคเต นว.\csromanspacer{} อวิคเต นว. (สํขิตฺตํ.)}

\csromanheader{2}{\textbf{ปจฺจนียุทฺธาร}}
\proseitem{256}{เหตุ อุเปกฺขาสหคโต ธมฺโม เหตุสฺส อุเปกฺขาสหคตสฺส ธมฺมสฺส อารมฺมณปจฺจเยน ปจฺจโย.\csromanspacer{} สหชาตปจฺจเยน ปจฺจโย.\csromanspacer{} อุปนิสฺสยปจฺจเยน ปจฺจโย. (สํขิตฺตํ.)}

\proseitem{257}{นเหตุยา นว, น-อารมฺมเณ นว. (สํขิตฺตํ.)}

\csromanheader{2}{เหตุปจฺจยา น-อารมฺมเณ ตีณิ. (สํขิตฺตํ.)}
\csromanheader{2}{นเหตุปจฺจยา อารมฺมเณ นว. (สํขิตฺตํ.)}
\gambhira{ธมฺมานุโลม ทุกติกปฏฺฐานปาฬิ}
\prose{\csromanfolio{252}(ยถา กุสลตฺติเก ปญฺหาวารสฺส อนุโลมมฺปิ ปจฺจนียมฺปิ อนุโลมปจฺจนียมฺปิ ปจฺจนียานุโลมมฺปิ คณิตํ, เอวํ คเณตพฺพํ.)}

\nitthitam{เหตุทุกปีติตฺติกํ นิฏฺฐิตํ.}
\csromanmatikaanchor{mtk.39}
\csromantocmarkref{section}{8. ทสฺสเนนปหาตพฺพตฺติก}{mtk.39}
\bookmark[dest={mtk.39},level=1]{8. ทสฺสเนนปหาตพฺพตฺติก}
\csromanheader{1}{1. เหตุทุก 8. ทสฺสเนนปหาตพฺพตฺติก}
\csromanheader{2}{1. ทสฺสเนนปหาตพฺพปท 1. ปฏิจฺจวาร}
\csromanheader{2}{\textbf{ปจฺจยจตุกฺก-เหตุ}}
\proseitem{258}{เหตุํ ทสฺสเนน ปหาตพฺพํ ธมฺมํ ปฏิจฺจ เหตุ ทสฺสเนน ปหาตพฺโพ ธมฺโม อุปฺปชฺชติ เหตุปจฺจยา.\csromanspacer{} ตีณิ.}

\prose{นเหตุํ ทสฺสเนน ปหาตพฺพํ ธมฺมํ ปฏิจฺจ นเหตุ ทสฺสเนน ปหาตพฺโพ ธมฺโม อุปฺปชฺชติ เหตุปจฺจยา.\csromanspacer{} ตีณิ.}

\prose{เหตุํ ทสฺสเนน ปหาตพฺพญฺจ นเหตุํ ทสฺสเนน ปหาตพฺพญฺจ ธมฺมํ ปฏิจฺจ เหตุ ทสฺสเนน ปหาตพฺโพ ธมฺโม อุปฺปชฺชติ เหตุปจฺจยา.\csromanspacer{} ตีณิ. (สํขิตฺตํ.)}

\proseitem{259}{เหตุยา นว, อารมฺมเณ นว, อธิปติยา นว ฯเปฯ กมฺเม นว, อาหาเร นว ฯเปฯ อวิคเต นว. (สํขิตฺตํ.)}

\csromanheader{2}{\textbf{นเหตุ น-อธิปติ}}
\proseitem{260}{นเหตุํ ทสฺสเนน ปหาตพฺพํ ธมฺมํ ปฏิจฺจ เหตุ ทสฺสเนน ปหาตพฺโพ ธมฺโม อุปฺปชฺชติ นเหตุปจฺจยา. (1)}

\prose{เหตุํ ทสฺสเนน ปหาตพฺพํ ธมฺมํ ปฏิจฺจ เหตุ ทสฺสเนน ปหาตพฺโพ ธมฺโม อุปฺปชฺชติ น-อธิปติปจฺจยา.\csromanspacer{} นว.}

\proseitem{261}{นเหตุยา เอกํ, น-อธิปติยา นว, นปุเรชาเต นว, นปจฺฉาชาเต นว, น-อาเสวเน นว, นกมฺเม ตีณิ, นวิปาเก นว, นวิปฺปยุตฺเต นว. (สํขิตฺตํ.)}

\csromanheader{2}{เหตุปจฺจยา น-อธิปติยา นว. (สํขิตฺตํ.)}
\csromanheader{2}{1. เหตุทุก 8. ทสฺสเนนปหาตพฺพตฺติก}
\csromanheader{2}{นเหตุปจฺจยา อารมฺมเณ เอกํ. (สํขิตฺตํ.)}
\prose{\csromanfolio{253}(สหชาตวาโรปิ ปจฺจยวาโรปิ นิสฺสยวาโรปิ สํสฏฺฐวาโรปิ สมฺปยุตฺตวาโรปิ ปฏิจฺจวารสทิสา วิตฺถาเรตพฺพา.)}

\csromanheader{2}{1. ทสฺสเนนปหาตพฺพปท 7. ปญฺหาวาร}
\csromanheader{2}{\textbf{ปจฺจยจตุกฺก-เหตุ อารมฺมณ}}
\proseitem{262}{เหตุ ทสฺสเนน ปหาตพฺโพ ธมฺโม เหตุสฺส ทสฺสเนน ปหาตพฺพสฺส ธมฺมสฺส เหตุปจฺจเยน ปจฺจโย.\csromanspacer{} ตีณิ.}

\prose{เหตุ ทสฺสเนน ปหาตพฺโพ ธมฺโม เหตุสฺส ทสฺสเนน ปหาตพฺพสฺส ธมฺมสฺส อารมฺมณปจฺจเยน ปจฺจโย.\csromanspacer{} นว. (สํขิตฺตํ.)}

\proseitem{263}{เหตุยา ตีณิ, อารมฺมเณ นว, อธิปติยา นว, อนนฺตเร นว, สมนนฺตเร นว, สหชาเต นว, อญฺญมญฺเญ นว, นิสฺสเย นว, อุปนิสฺสเย นว, อาเสวเน นว, กมฺเม ตีณิ, อาหาเร ตีณิ, อินฺทฺริเย ตีณิ, ฌาเน ตีณิ, มคฺเค ตีณิ, สมฺปยุตฺเต นว ฯเปฯ อวิคเต นว. (สํขิตฺตํ.)}

\csromanheader{2}{\textbf{ปจฺจนียุทฺธาร}}
\proseitem{264}{เหตุ ทสฺสเนน ปหาตพฺโพ ธมฺโม เหตุสฺส ทสฺสเนน ปหาตพฺพสฺส ธมฺมสฺส อารมฺมณปจฺจเยน ปจฺจโย.\csromanspacer{} สหชาตปจฺจเยน ปจฺจโย.\csromanspacer{} อุปนิสฺสยปจฺจเยน ปจฺจโย. (สํขิตฺตํ.)}

\proseitem{265}{นเหตุยา นว, น-อารมฺมเณ นว. (สํขิตฺตํ.)}

\csromanheader{2}{เหตุปจฺจยา น-อารมฺมเณ ตีณิ. (สํขิตฺตํ.)}
\csromanheader{2}{นเหตุปจฺจยา อารมฺมเณ นว. (สํขิตฺตํ.)}
\prose{(ยถา กุสลตฺติเก ปญฺหาวารสฺส อนุโลมมฺปิ ปจฺจนียมฺปิ อนุโลมปจฺจนียมฺปิ ปจฺจนียานุโลมมฺปิ คณิตํ, เอวํ คเณตพฺพํ.)}

\csromanheader{2}{ธมฺมานุโลม ทุกติกปฏฺฐานปาฬิ \textbf{2. ภาวนายปหาตพฺพปท 1. ปฏิจฺจวาร}}
\csromanheader{2}{\textbf{ปจฺจยจตุกฺก-เหตุ}}
\proseitem{266}{\csromanfolio{254}เหตุํ ภาวนาย ปหาตพฺพํ ธมฺมํ ปฏิจฺจ เหตุ ภาวนาย ปหาตพฺโพ ธมฺโม อุปฺปชฺชติ เหตุปจฺจยา.\csromanspacer{} ตีณิ.}

\prose{\textbf{นเหตุ}ํ ภาวนาย ปหาตพฺพํ ธมฺมํ ปฏิจฺจ นเหตุ ภาวนาย ปหาตพฺโพ ธมฺโม อุปฺปชฺชติ เหตุปจฺจยา.\csromanspacer{} ตีณิ.}

\prose{เหตุํ ภาวนาย ปหาตพฺพญฺจ นเหตุํ ภาวนาย ปหาตพฺพญฺจ ธมฺมํ ปฏิจฺจ เหตุ ภาวนาย ปหาตพฺโพ ธมฺโม อุปฺปชฺชติ เหตุปจฺจยา.\csromanspacer{} ตีณิ. (สํขิตฺตํ.)}

\proseitem{267}{เหตุยา นว, อารมฺมเณ นว ฯเปฯ อวิคเต นว. (สํขิตฺตํ.)}

\csromanheader{2}{\textbf{นเหตุ}}
\proseitem{268}{\textbf{นเหตุ}ํ ภาวนาย ปหาตพฺพํ ธมฺมํ ปฏิจฺจ เหตุ ภาวนาย ปหาตพฺโพ ธมฺโม อุปฺปชฺชติ นเหตุปจฺจยา. (1) (สํขิตฺตํ.)}

\vspace{\tipitakaparskip}
\csromancenter{\textbf{นเหตุ}ยา เอกํ, น-อธิปติยา นว, นปุเรชาเต นว, นปจฺฉาชาเต นว, น-อาเสวเน นว, นกมฺเม ตีณิ, นวิปาเก นว, นวิปฺปยุตฺเต นว. (สํขิตฺตํ.)}
\csromanheader{2}{เหตุปจฺจยา น-อธิปติยา นว. (สํขิตฺตํ.)}
\csromanheader{2}{\textbf{นเหตุ}ปจฺจยา อารมฺมเณ เอกํ. (สํขิตฺตํ.)}
\prose{(สหชาตวาโรปิ ปจฺจยวาโรปิ นิสฺสยวาโรปิ สํสฏฺฐวาโรปิ สมฺปยุตฺตวาโรปิ ปฏิจฺจวารสทิสา วิตฺถาเรตพฺพา.)}

\csromanheader{2}{2. ภาวนายปหาตพฺพปท 7. ปญฺหาวาร}
\csromanheader{2}{\textbf{ปจฺจยจตุกฺก-เหตุ อารมฺมณ}}
\proseitem{270}{เหตุ ภาวนาย ปหาตพฺโพ ธมฺโม เหตุสฺส ภาวนาย ปหาตพฺพสฺส ธมฺมสฺส เหตุปจฺจเยน ปจฺจโย.\csromanspacer{} ตีณิ.}

\csromanheader{2}{1. เหตุทุก 8. ทสฺสเนนปหาตพฺพตฺติก}
\prose{\csromanfolio{255}เหตุ ภาวนาย ปหาตพฺโพ ธมฺโม เหตุสฺส ภาวนาย ปหาตพฺพสฺส ธมฺมสฺส อารมฺมณปจฺจเยน ปจฺจโย.\csromanspacer{} นว. (สํขิตฺตํ.)}

\proseitem{271}{เหตุยา ตีณิ, อารมฺมเณ นว, อธิปติยา นว, อนนฺตเร นว, สมนนฺตเร นว, สหชาเต นว, อญฺญมญฺเญ นว, นิสฺสเย นว, อุปนิสฺสเย นว, อาเสวเน นว, กมฺเม ตีณิ, อาหาเร ตีณิ, อินฺทฺริเย ตีณิ ฯเปฯ อวิคเต นว. (สํขิตฺตํ.)}

\csromanheader{2}{\textbf{ปจฺจนียุทฺธาร}}
\proseitem{272}{เหตุ ภาวนาย ปหาตพฺโพ ธมฺโม เหตุสฺส ภาวนาย ปหาตพฺพสฺส ธมฺมสฺส อารมฺมณปจฺจเยน ปจฺจโย.\csromanspacer{} สหชาตปจฺจเยน ปจฺจโย.\csromanspacer{} อุปนิสฺสยปจฺจเยน ปจฺจโย. (สํขิตฺตํ.)}

\proseitem{273}{นเหตุยา นว, น-อารมฺมเณ นว. (สํขิตฺตํ.)}

\csromanheader{2}{เหตุปจฺจยา น-อารมฺมเณ ตีณิ. (สํขิตฺตํ.)}
\csromanheader{2}{นเหตุปจฺจยา อารมฺมเณ นว. (สํขิตฺตํ.)}
\prose{(ยถา กุสลตฺติเก ปญฺหาวารสฺส อนุโลมมฺปิ ปจฺจนียมฺปิ อนุโลมปจฺจนียมฺปิ ปจฺจนียานุโลมมฺปิ คณิตํ, เอวํ คเณตพฺพํ.)}

\csromanheader{2}{3. เนวทสฺสเนนภาวนายปหาตพฺพปท 1. ปฏิจฺจวาร}
\csromanheader{2}{\textbf{ปจฺจยจตุกฺก-เหตุ}}
\proseitem{274}{เหตุํ เนวทสฺสเนน นภาวนาย ปหาตพฺพํ ธมฺมํ ปฏิจฺจ เหตุ เนวทสฺสเนน นภาวนาย ปหาตพฺโพ ธมฺโม อุปฺปชฺชติ เหตุปจฺจยา.\csromanspacer{} ตีณิ.}

\prose{นเหตุํ เนวทสฺสเนน นภาวนาย ปหาตพฺพํ ธมฺมํ ปฏิจฺจ นเหตุ เนวทสฺสเนน นภาวนาย ปหาตพฺโพ ธมฺโม อุปฺปชฺชติ เหตุปจฺจยา. (สํขิตฺตํ.)}

\gambhira{ธมฺมานุโลม ทุกติกปฏฺฐานปาฬิ}
\proseitem{275}{\csromanfolio{256}เหตุยา นว, อารมฺมเณ นว, อธิปติยา นว ฯเปฯ กมฺเม นว, วิปาเก นว ฯเปฯ อวิคเต นว. (สํขิตฺตํ.)}

\csromanheader{2}{\textbf{นเหตฺวาทิ}}
\proseitem{276}{นเหตุํ เนวทสฺสเนน นภาวนาย ปหาตพฺพํ ธมฺมํ ปฏิจฺจ นเหตุ เนวทสฺสเนน นภาวนาย ปหาตพฺโพ ธมฺโม อุปฺปชฺชติ นเหตุปจฺจยา.}

\prose{เหตุํ เนวทสฺสเนน นภาวนาย ปหาตพฺพํ ธมฺมํ ปฏิจฺจ นเหตุ เนวทสฺสเนน นภาวนาย ปหาตพฺโพ ธมฺโม อุปฺปชฺชติ น-อารมฺมณปจฺจยา. (1)}

\prose{นเหตุํ เนวทสฺสเนน นภาวนาย ปหาตพฺพํ ธมฺมํ ปฏิจฺจ นเหตุ เนวทสฺสเนน นภาวนาย ปหาตพฺโพ ธมฺโม อุปฺปชฺชติ น-อารมฺมณปจฺจยา. (1)}

\prose{เหตุํ เนวทสฺสเนน นภาวนาย ปหาตพฺพญฺจ นเหตุํ เนวทสฺสเนน นภาวนาย ปหาตพฺพญฺจ ธมฺมํ ปฏิจฺจ นเหตุ เนวทสฺสเนน นภาวนาย ปหาตพฺโพ ธมฺโม อุปฺปชฺชติ น-อารมฺมณปจฺจยา. (1)}

\prose{เหตุํ เนวทสฺสเนน นภาวนาย ปหาตพฺพํ ธมฺมํ ปฏิจฺจ เหตุ เนวทสฺสเนน นภาวนาย ปหาตพฺโพ ธมฺโม อุปฺปชฺชติ น-อธิปติปจฺจยา. (สํขิตฺตํ.)}

\proseitem{277}{นเหตุยา เอกํ, น-อารมฺมเณ ตีณิ, น-อธิปติยา นว, น-อนนฺตเร ตีณิ, นสมนนฺตเร ตีณิ, น-อญฺญมญฺเญ ตีณิ, น-อุปนิสฺสเย ตีณิ, นปุเรชาเต นว, นปจฺฉาชาเต นว, น-อาเสวเน นว, นกมฺเม ตีณิ, นวิปาเก นว, น-อาหาเร เอกํ, น-อินฺทฺริเย เอกํ, นฌาเน เอกํ, นมคฺเค เอกํ, นสมฺปยุตฺเต ตีณิ, นวิปฺปยุตฺเต นว, โนนตฺถิยา ตีณิ, โนวิคเต ตีณิ. (สํขิตฺตํ.)}

\csromanheader{2}{เหตุปจฺจยา น-อารมฺมเณ ตีณิ. (สํขิตฺตํ.)}
\csromanheader{2}{นเหตุปจฺจยา อารมฺมเณ เอกํ. (สํขิตฺตํ.)}
\prose{(สหชาตวาโรปิ ปจฺจยวาโรปิ นิสฺสยวาโรปิ สํสฏฺฐวาโรปิ สมฺปยุตฺตวาโรปิ ปฏิจฺจวารสทิสา วิตฺถาเรตพฺพา.)}

\vspace{\tipitakaparskip}
\csromancenter{เหตุทุก 8.\csromanspacer{} ทสฺสเนนปหาตพฺพตฺติก \textbf{3.}\csromanspacer{}\textbf{ เนวทสฺสเนนนภาวนายปหาตพฺพปท 7.}\csromanspacer{}\textbf{ ปญฺหาวาร}}
\csromanheader{2}{\textbf{ปจฺจยจตุกฺก-เหตุ อารมฺมณ}}
\proseitem{278}{\csromanfolio{257}เหตุ เนวทสฺสเนน นภาวนาย ปหาตพฺโพ ธมฺโม เหตุสฺส เนวทสฺสเนน นภาวนาย ปหาตพฺพสฺส ธมฺมสฺส เหตุปจฺจเยน ปจฺจโย.\csromanspacer{} ตีณิ.}

\prose{เหตุ เนวทสฺสเนน นภาวนาย ปหาตพฺโพ ธมฺโม เหตุสฺส เนวทสฺสเนน นภาวนาย ปหาตพฺพสฺส ธมฺมสฺส อารมฺมณปจฺจเยน ปจฺจโย.\csromanspacer{} นว. (สํขิตฺตํ.)}

\proseitem{279}{เหตุยา ตีณิ, อารมฺมเณ นว, อธิปติยา นว, อนนฺตเร นว, สมนนฺตเร นว, สหชาเต นว, อญฺญมญฺเญ นว, นิสฺสเย นว, อุปนิสฺสเย นว, ปุเรชาเต ตีณิ, ปจฺฉาชาเต ตีณิ, อาเสวเน นว, กมฺเม ตีณิ, วิปาเก นว, อาหาเร ตีณิ, อินฺทฺริเย นว, ฌาเน ตีณิ, มคฺเค นว, สมฺปยุตฺเต นว, วิปฺปยุตฺเต ปญฺจ, อตฺถิยา นว, นตฺถิยา นว, วิคเต นว, อวิคเต นว.}

\csromanheader{2}{\textbf{ปจฺจนียุทฺธาร}}
\proseitem{280}{เหตุ เนวทสฺสเนน นภาวนาย ปหาตพฺโพ ธมฺโม เหตุสฺส เนวทสฺสเนน นภาวนาย ปหาตพฺพสฺส ธมฺมสฺส อารมฺมณปจฺจเยน ปจฺจโย.\csromanspacer{} สหชาตปจฺจเยน ปจฺจโย.\csromanspacer{} อุปนิสฺสยปจฺจเยน ปจฺจโย. (สํขิตฺตํ.)}

\proseitem{281}{นเหตุยา นว, น-อารมฺมเณ นว. (สํขิตฺตํ.)}

\csromanheader{2}{เหตุปจฺจยา น-อารมฺมเณ ตีณิ. (สํขิตฺตํ.)}
\csromanheader{2}{นเหตุปจฺจยา อารมฺมเณ นว. (สํขิตฺตํ.)}
\prose{(ยถา กุสลตฺติเก ปญฺหาวารสฺส อนุโลมมฺปิ ปจฺจนียมฺปิ อนุโลมปจฺจนียมฺปิ ปจฺจนียานุโลมมฺปิ คณิตํ, เอวํ คเณตพฺพํ.)}

\nitthitam{เหตุทุก ทสฺสเนน ปหาตพฺพตฺติกํ นิฏฺฐิตํ.}
\csromanmatikaanchor{mtk.40}
\csromantocmarkref{section}{9. ทสฺสเนนปหาตพฺพเหตุกตฺติก}{mtk.40}
\bookmark[dest={mtk.40},level=1]{9. ทสฺสเนนปหาตพฺพเหตุกตฺติก}
\csromanheader{1}{ธมฺมานุโลม ทุกติกปฏฺฐานปาฬิ \textbf{1. เหตุทุก 9. ทสฺสเนนปหาตพฺพเหตุกตฺติก}}
\csromanheader{2}{1. ทสฺสเนนปหาตพฺพเหตุกปท 1. ปฏิจฺจวาร}
\csromanheader{2}{\textbf{ปจฺจยจตุกฺก-เหตุ}}
\proseitem{282}{\csromanfolio{258}เหตุํ ทสฺสเนน ปหาตพฺพเหตุกํ ธมฺมํ ปฏิจฺจ เหตุ ทสฺสเนน ปหาตพฺพเหตุโก ธมฺโม อุปฺปชฺชติ เหตุปจฺจยา.\csromanspacer{} ตีณิ.}

\prose{นเหตุํ ทสฺสเนน ปหาตพฺพเหตุกํ ธมฺมํ ปฏิจฺจ นเหตุ ทสฺสเนน ปหาตพฺพเหตุโก ธมฺโม อุปฺปชฺชติ เหตุปจฺจยา.\csromanspacer{} ตีณิ.}

\prose{เหตุํ ทสฺสเนน ปหาตพฺพเหตุกญฺจ นเหตุํ ทสฺสเนน ปหาตพฺพเหตุกญฺจ ธมฺมํ ปฏิจฺจ เหตุ ทสฺสเนน ปหาตพฺพเหตุโก ธมฺโม อุปฺปชฺชติ เหตุปจฺจยา.\csromanspacer{} ตีณิ. (สํขิตฺตํ.)}

\proseitem{283}{เหตุยา นว, อารมฺมเณ นว ฯเปฯ อวิคเต นว. (สํขิตฺตํ.)}

\csromanheader{2}{\textbf{น-อธิปติ}}
\proseitem{284}{เหตุํ ทสฺสเนน ปหาตพฺพเหตุกํ ธมฺมํ ปฏิจฺจ เหตุ ทสฺสเนน ปหาตพฺพเหตุโก ธมฺโม อุปฺปชฺชติ น-อธิปติปจฺจยา. (สํขิตฺตํ.)}

\vspace{\tipitakaparskip}
\csromancenter{\textbf{น-อธิปติ}ยา นว, นปุเรชาเต นว, นปจฺฉาชาเต นว, น-อาเสวเน นว, นกมฺเม ตีณิ, นวิปาเก นว, นวิปฺปยุตฺเต นว. (สํขิตฺตํ.)}
\csromanheader{2}{เหตุปจฺจยา น-อธิปติยา นว. (สํขิตฺตํ.)}
\csromanheader{2}{\textbf{น-อธิปติ}ปจฺจยา เหตุยา นว. (สํขิตฺตํ.)}
\prose{(สหชาตวาโรปิ ปจฺจยวาโรปิ นิสฺสยวาโรปิ สํสฏฺฐวาโรปิ สมฺปยุตฺตวาโรปิ ปฏิจฺจวารสทิสา วิตฺถาเรตพฺพา.)}

\csromanheader{2}{1. ทสฺสเนนปหาตพฺพเหตุกปท 7. ปญฺหาวาร}
\csromanheader{2}{\textbf{ปจฺจยจตุกฺก-เหตุ อารมฺมณ}}
\proseitem{286}{เหตุ ทสฺสเนน ปหาตพฺพเหตุโก ธมฺโม เหตุสฺส ทสฺสเนน ปหาตพฺพเหตุกสฺส ธมฺมสฺส เหตุปจฺจเยน ปจฺจโย.\csromanspacer{} ตีณิ.}

\csromanheader{2}{1. เหตุทุก 9. ทสฺสเนนปหาตพฺพเหตุกตฺติก}
\prose{\csromanfolio{259}เหตุ ทสฺสเนน ปหาตพฺพเหตุโก ธมฺโม เหตุสฺส ทสฺสเนน ปหาตพฺพเหตุกสฺส ธมฺมสฺส อารมฺมณปจฺจเยน ปจฺจโย.\csromanspacer{} นว. (สํขิตฺตํ.)}

\proseitem{287}{เหตุยา ตีณิ, อารมฺมเณ นว, อธิปติยา นว ฯเปฯ อุปนิสฺสเย นว, อาเสวเน นว, กมฺเม ตีณิ, อาหาเร ตีณิ ฯเปฯ มคฺเค ตีณิ, สมฺปยุตฺเต นว ฯเปฯ อวิคเต นว. (สํขิตฺตํ.)}

\csromanheader{2}{\textbf{ปจฺจนียุทฺธาร}}
\proseitem{288}{เหตุ ทสฺสเนน ปหาตพฺพเหตุโก ธมฺโม เหตุสฺส ทสฺสเนน ปหาตพฺพเหตุกสฺส ธมฺมสฺส อารมฺมณปจฺจเยน ปจฺจโย.\csromanspacer{} สหชาตปจฺจเยน ปจฺจโย.\csromanspacer{} อุปนิสฺสยปจฺจเยน ปจฺจโย. (สํขิตฺตํ.)}

\proseitem{289}{นเหตุยา นว, น-อารมฺมเณ นว. (สํขิตฺตํ.)}

\csromanheader{2}{เหตุปจฺจยา น-อารมฺมเณ ตีณิ. (สํขิตฺตํ.)}
\csromanheader{2}{นเหตุปจฺจยา อารมฺมเณ นว. (สํขิตฺตํ.)}
\prose{(ยถา กุสลตฺติเก ปญฺหาวารสฺส อนุโลมมฺปิ ปจฺจนียมฺปิ อนุโลมปจฺจนียมฺปิ ปจฺจนียานุโลมมฺปิ คณิตํ, เอวํ คเณตพฺพํ.)}

\csromanheader{2}{2. ภาวนายปหาตพฺพเหตุกปท 1. ปฏิจฺจวาร}
\csromanheader{2}{\textbf{ปจฺจยจตุกฺก-เหตุ}}
\proseitem{290}{เหตุํ ภาวนาย ปหาตพฺพเหตุกํ ธมฺมํ ปฏิจฺจ เหตุ ภาวนาย ปหาตพฺพเหตุโก ธมฺโม อุปฺปชฺชติ เหตุปจฺจยา.\csromanspacer{} ตีณิ.}

\prose{นเหตุํ ภาวนาย ปหาตพฺพเหตุกํ ธมฺมํ ปฏิจฺจ นเหตุ ภาวนาย ปหาตพฺพเหตุโก ธมฺโม อุปฺปชฺชติ เหตุปจฺจยา. (สํขิตฺตํ.)}

\proseitem{291}{เหตุยา นว, อารมฺมเณ นว, อธิปติยา นว ฯเปฯ อวิคเต นว. (สํขิตฺตํ.)}

\csromanheader{2}{ธมฺมานุโลม ทุกติกปฏฺฐานปาฬิ \textbf{น-อธิปติ}}
\proseitem{292}{\csromanfolio{260}เหตุํ ภาวนาย ปหาตพฺพเหตุกํ ธมฺมํ ปฏิจฺจ เหตุ ภาวนาย ปหาตพฺพเหตุโก ธมฺโม อุปฺปชฺชติ น-อธิปติปจฺจยา.\csromanspacer{} นว. (สํขิตฺตํ.)}

\vspace{\tipitakaparskip}
\csromancenter{\textbf{น-อธิปติ}ยา นว, นปุเรชาเต นว, นปจฺฉาชาเต นว, น-อาเสวเน นว, นกมฺเม ตีณิ, นวิปาเก นว, นวิปฺปยุตฺเต นว. (สํขิตฺตํ.)}
\csromanheader{2}{เหตุปจฺจยา น-อธิปติยา นว. (สํขิตฺตํ.)}
\csromanheader{2}{\textbf{น-อธิปติ}ปจฺจยา เหตุยา นว. (สํขิตฺตํ.)}
\prose{(สหชาตวาโรปิ ปจฺจยวาโรปิ นิสฺสยวาโรปิ สํสฏฺฐวาโรปิ สมฺปยุตฺตวาโรปิ ปฏิจฺจวารสทิสา วิตฺถาเรตพฺพา.)}

\csromanheader{2}{2. ภาวนายปหาตพฺพเหตุกปท 7. ปญฺหาวาร}
\csromanheader{2}{\textbf{ปจฺจยจตุกฺก-เหตุ อารมฺมณ}}
\proseitem{294}{เหตุ ภาวนาย ปหาตพฺพเหตุโก ธมฺโม เหตุสฺส ภาวนาย ปหาตพฺพเหตุกสฺส ธมฺมสฺส เหตุปจฺจเยน ปจฺจโย.\csromanspacer{} ตีณิ.}

\prose{เหตุ ภาวนาย ปหาตพฺพเหตุโก ธมฺโม เหตุสฺส ภาวนาย ปหาตพฺพเหตุกสฺส ธมฺมสฺส อารมฺมณปจฺจเยน ปจฺจโย.\csromanspacer{} นว. (สํขิตฺตํ.)}

\proseitem{295}{เหตุยา ตีณิ, อารมฺมเณ นว, อธิปติยา นว, อนนฺตเร นว, สมนนฺตเร นว, สหชาเต นว, อญฺญมญฺเญ นว, นิสฺสเย นว, อุปนิสฺสเย นว, อาเสวเน นว, กมฺเม ตีณิ, อาหาเร ตีณิ ฯเปฯ มคฺเค ตีณิ, สมฺปยุตฺเต นว ฯเปฯ อวิคเต นว. (สํขิตฺตํ.)}

\csromanheader{2}{\textbf{ปจฺจนียุทฺธาร}}
\proseitem{296}{เหตุ ภาวนาย ปหาตพฺพเหตุโก ธมฺโม เหตุสฺส ภาวนาย ปหาตพฺพเหตุกสฺส ธมฺมสฺส อารมฺมณปจฺจเยน ปจฺจโย.\csromanspacer{} สหชาตปจฺจเยน ปจฺจโย.\csromanspacer{} อุปนิสฺสยปจฺจเยน ปจฺจโย. (สํขิตฺตํ.)}

\csromanheader{2}{1. เหตุทุก 9. ทสฺสเนนปหาตพฺพเหตุกตฺติก}
\proseitem{297}{\csromanfolio{261}\textbf{นเหตุ}ยา นว, น-อารมฺมเณ นว. (สํขิตฺตํ.)}

\csromanheader{2}{เหตุปจฺจยา น-อารมฺมเณ ตีณิ. (สํขิตฺตํ.)}
\csromanheader{2}{\textbf{นเหตุ}ปจฺจยา อารมฺมเณ นว. (สํขิตฺตํ.)}
\prose{(ยถา กุสลตฺติเก ปญฺหาวารสฺส อนุโลมมฺปิ ปจฺจนียมฺปิ อนุโลมปจฺจนียมฺปิ ปจฺจนียานุโลมมฺปิ คณิตํ, เอวํ คเณตพฺพํ.)}

\csromanheader{2}{3. เนวทสฺสเนนนภาวนายปหาตพฺพเหตุกปท 1. ปฏิจฺจวาร}
\csromanheader{2}{\textbf{ปจฺจยจตุกฺก-เหตุ}}
\proseitem{298}{เหตุํ เนวทสฺสเนน นภาวนาย ปหาตพฺพเหตุกํ ธมฺมํ ปฏิจฺจ เหตุ เนวทสฺสเนน นภาวนาย ปหาตพฺพเหตุโก ธมฺโม อุปฺปชฺชติ เหตุปจฺจยา. (สํขิตฺตํ.)}

\proseitem{299}{เหตุยา นว, อารมฺมเณ นว ฯเปฯ อวิคเต นว. (สํขิตฺตํ.)}

\csromanheader{2}{\textbf{นเหตุ}}
\proseitem{300}{เหตุํ เนวทสฺสเนน นภาวนาย ปหาตพฺพเหตุกํ ธมฺมํ ปฏิจฺจ นเหตุ เนวทสฺสเนน นภาวนาย ปหาตพฺพเหตุโก ธมฺโม อุปฺปชฺชติ นเหตุปจฺจยา. (สํขิตฺตํ.)}

\vspace{\tipitakaparskip}
\csromancenter{\textbf{นเหตุ}ยา เอกํ, น-อารมฺมเณ ตีณิ, น-อธิปติยา นว, น-อนนฺตเร ตีณิ, นสมนนฺตเร ตีณิ, น-อญฺญมญฺเญ ตีณิ, น-อุปนิสฺสเย ตีณิ, นปุเรชาเต นว, นปจฺฉาชาเต นว, น-อาเสวเน นว, นกมฺเม ตีณิ, นวิปาเก นว, น-อาหาเร เอกํ, น-อินฺทฺริเย เอกํ, นฌาเน เอกํ, นมคฺเค เอกํ, นสมฺปยุตฺเต ตีณิ, นวิปฺปยุตฺเต นว, โนนตฺถิยา ตีณิ, โนวิคเต ตีณิ. (สํขิตฺตํ.)}
\csromanheader{2}{เหตุปจฺจยา น-อารมฺมเณ ตีณิ. (สํขิตฺตํ.)}
\csromanheader{2}{\textbf{นเหตุ}ปจฺจยา อารมฺมเณ เอกํ. (สํขิตฺตํ.)}
\prose{(สหชาตวาโรปิ ปจฺจยวาโรปิ นิสฺสยวาโรปิ สํสฏฺฐวาโรปิ สมฺปยุตฺตวาโรปิ ปฏิจฺจวารสทิสา วิตฺถาเรตพฺพา.)}

\vspace{\tipitakaparskip}
\csromancenter{ธมฺมานุโลม ทุกติกปฏฺฐานปาฬิ \textbf{3.}\csromanspacer{}\textbf{ เนวทสฺสเนนนภาวนายปหาตพฺพเหตุกปท 7.}\csromanspacer{}\textbf{ ปญฺหาวาร}}
\csromanheader{2}{\textbf{ปจฺจยจตุกฺก-เหตุ}}
\proseitem{302}{\csromanfolio{262}เหตุ เนวทสฺสเนน นภาวนาย ปหาตพฺพเหตุโก ธมฺโม เหตุสฺส เนวทสฺสเนน นภาวนาย ปหาตพฺพเหตุกสฺส ธมฺมสฺส เหตุปจฺจเยน ปจฺจโย.\csromanspacer{} ตีณิ. (สํขิตฺตํ.)}

\proseitem{303}{เหตุยา ตีณิ, อารมฺมเณ นว, อธิปติยา นว, อนนฺตเร นว, สมนนฺตเร นว, สหชาเต นว, อญฺญมญฺเญ นว, นิสฺสเย นว, อุปนิสฺสเย นว, ปุเรชาเต ตีณิ, ปจฺฉาชาเต ตีณิ, อาเสวเน นว, กมฺเม ตีณิ, วิปาเก นว, อาหาเร ตีณิ, อินฺทฺริเย นว, ฌาเน ตีณิ, มคฺเค นว, สมฺปยุตฺเต นว, วิปฺปยุตฺเต ปญฺจ, อตฺถิยา นว, นตฺถิยา นว, วิคเต นว, อวิคเต นว. (สํขิตฺตํ.)}

\csromanheader{2}{\textbf{ปจฺจนียุทฺธาร}}
\proseitem{304}{เหตุ เนวทสฺสเนน นภาวนาย ปหาตพฺพเหตุโก ธมฺโม เหตุสฺส เนวทสฺสเนน นภาวนาย ปหาตพฺพเหตุกสฺส ธมฺมสฺส อารมฺมณปจฺจเยน ปจฺจโย.\csromanspacer{} สหชาตปจฺจเยน ปจฺจโย.\csromanspacer{} อุปนิสฺสยปจฺจเยน ปจฺจโย. (สํขิตฺตํ.)}

\proseitem{305}{นเหตุยา นว, น-อารมฺมเณ นว. (สํขิตฺตํ.)}

\csromanheader{2}{เหตุปจฺจยา น-อารมฺมเณ ตีณิ. (สํขิตฺตํ.)}
\csromanheader{2}{นเหตุปจฺจยา อารมฺมเณ นว. (สํขิตฺตํ.)}
\prose{(ยถา กุสลตฺติเก ปญฺหาวารสฺส อนุโลมมฺปิ ปจฺจนียมฺปิ อนุโลมปจฺจนียมฺปิ ปจฺจนียานุโลมมฺปิ คณิตํ, เอวํ คเณตพฺพํ.)}

\nitthitam{เหตุทุก ทสฺสเนน ปหาตพฺพเหตุกตฺติกํ นิฏฺฐิตํ.}
\csromanmatikaanchor{mtk.41}
\csromantocmarkref{section}{10. อาจยคามิตฺติก}{mtk.41}
\bookmark[dest={mtk.41},level=1]{10. อาจยคามิตฺติก}
\csromanheader{1}{1. เหตุทุก 10. อาจยคามิตฺติก \textbf{1. เหตุทุก 10. อาจยคามิตฺติก}}
\csromanheader{2}{1. อาจยคามิปท 1. ปฏิจฺจวาร}
\csromanheader{2}{\textbf{ปจฺจยจตุกฺก-เหตุ}}
\proseitem{306}{\csromanfolio{263}เหตุํ อาจยคามิํ ธมฺมํ ปฏิจฺจ เหตุ อาจยคามี ธมฺโม อุปฺปชฺชติ เหตุปจฺจยา.\csromanspacer{} ตีณิ.}

\prose{นเหตุํ อาจยคามิํ ธมฺมํ ปฏิจฺจ นเหตุ อาจยคามี ธมฺโม อุปฺปชฺชติ เหตุปจฺจยา.\csromanspacer{} ตีณิ.}

\prose{เหตุํ อาจยคามิญฺจ นเหตุํ อาจยคามิญฺจ ธมฺมํ ปฏิจฺจ เหตุ อาจยคามี ธมฺโม อุปฺปชฺชติ เหตุปจฺจยา.\csromanspacer{} ตีณิ. (สํขิตฺตํ.)}

\proseitem{307}{เหตุยา นว, อารมฺมเณ นว, อธิปติยา นว, อนนฺตเร นว, สมนนฺตเร นว ฯเปฯ อวิคเต นว. (สํขิตฺตํ.)}

\csromanheader{2}{\textbf{นเหตุ น-อธิปติ}}
\proseitem{308}{นเหตุํ อาจยคามิํ ธมฺมํ ปฏิจฺจ เหตุ อาจยคามี ธมฺโม อุปฺปชฺชติ นเหตุปจฺจยา. (1)}

\prose{เหตุํ อาจยคามิํ ธมฺมํ ปฏิจฺจ เหตุ อาจยคามี ธมฺโม อุปฺปชฺชติ น-อธิปติปจฺจยา. (สํขิตฺตํ.)}

\proseitem{309}{นเหตุยา เอกํ, น-อธิปติยา นว, นปุเรชาเต นว, นปจฺฉาชาเต นว, น-อาเสวเน นว, นกมฺเม ตีณิ, นวิปาเก นว, นวิปฺปยุตฺเต นว. (สํขิตฺตํ.)}

\csromanheader{2}{เหตุปจฺจยา น-อธิปติยา นว. (สํขิตฺตํ.)}
\csromanheader{2}{นเหตุปจฺจยา อารมฺมเณ เอกํ. (สํขิตฺตํ.)}
\prose{(สหชาตวาโรปิ ปจฺจยวาโรปิ นิสฺสยวาโรปิ สํสฏฺฐวาโรปิ สมฺปยุตฺตวาโรปิ ปฏิจฺจวารสทิสา วิตฺถาเรตพฺพา.)}

\csromanheader{2}{ธมฺมานุโลม ทุกติกปฏฺฐานปาฬิ \textbf{1. อาจยคามิปท 7. ปญฺหาวาร}}
\csromanheader{2}{\textbf{ปจฺจยจตุกฺก-เหตฺวาทิ}}
\proseitem{310}{\csromanfolio{264}เหตุ อาจยคามี ธมฺโม เหตุสฺส อาจยคามิสฺส ธมฺมสฺส เหตุปจฺจเยน ปจฺจโย.\csromanspacer{} ตีณิ.}

\prose{เหตุ อาจยคามี ธมฺโม เหตุสฺส อาจยคามิสฺส ธมฺมสฺส อารมฺมณปจฺจเยน ปจฺจโย.\csromanspacer{} นว.}

\prose{เหตุ อาจยคามี ธมฺโม เหตุสฺส อาจยคามิสฺส ธมฺมสฺส อธิปติปจฺจเยน ปจฺจโย.\csromanspacer{} อารมฺมณาธิปติ, สหชาตาธิปติ.\csromanspacer{} ตีณิ.}

\prose{นเหตุ อาจยคามี ธมฺโม นเหตุสฺส อาจยคามิสฺส ธมฺมสฺส อธิปติปจฺจเยน ปจฺจโย.\csromanspacer{} อารมฺมณาธิปติ, สหชาตาธิปติ.\csromanspacer{} ตีณิ.}

\prose{เหตุ อาจยคามี จ นเหตุ อาจยคามี จ ธมฺมา เหตุสฺส อาจยคามิสฺส ธมฺมสฺส อธิปติปจฺจเยน ปจฺจโย.\csromanspacer{} อารมฺมณาธิปติ.\csromanspacer{} ตีณิ. (สํขิตฺตํ.)}

\proseitem{311}{เหตุยา ตีณิ, อารมฺมเณ นว, อธิปติยา นว ฯเปฯ อุปนิสฺสเย นว, อาเสวเน นว, กมฺเม ตีณิ, อาหาเร ตีณิ ฯเปฯ อวิคเต นว. (สํขิตฺตํ.)}

\csromanheader{2}{\textbf{ปจฺจนียุทฺธาร}}
\proseitem{312}{เหตุ อาจยคามี ธมฺโม เหตุสฺส อาจยคามิสฺส ธมฺมสฺส อารมฺมณปจฺจเยน ปจฺจโย.\csromanspacer{} สหชาตปจฺจเยน ปจฺจโย.\csromanspacer{} อุปนิสฺสยปจฺจเยน ปจฺจโย. (สํขิตฺตํ.)}

\proseitem{313}{นเหตุยา นว, น-อารมฺมเณ นว. (สํขิตฺตํ.)}

\csromanheader{2}{เหตุปจฺจยา น-อารมฺมเณ ตีณิ. (สํขิตฺตํ.)}
\csromanheader{2}{นเหตุปจฺจยา อารมฺมเณ นว. (สํขิตฺตํ.)}
\prose{(ยถา กุสลตฺติเก ปญฺหาวารสฺส อนุโลมมฺปิ ปจฺจนียมฺปิ อนุโลมปจฺจนียมฺปิ ปจฺจนียานุโลมมฺปิ คณิตํ, เอวํ คเณตพฺพํ.)}

\vspace{\tipitakaparskip}
\csromancenter{เหตุทุก 10.\csromanspacer{} อาจยคามิตฺติก \textbf{2.}\csromanspacer{}\textbf{ อปจยคามิปท 1.}\csromanspacer{}\textbf{ ปฏิจฺจวาร}}
\csromanheader{2}{\textbf{ปจฺจยจตุกฺก-เหตุ}}
\proseitem{314}{\csromanfolio{265}เหตุํ อปจยคามิํ ธมฺมํ ปฏิจฺจ เหตุ อปจยคามี ธมฺโม อุปฺปชฺชติ เหตุปจฺจยา.\csromanspacer{} ตีณิ.}

\prose{นเหตุํ อปจยคามิํ ธมฺมํ ปฏิจฺจ นเหตุ อปจยคามี ธมฺโม อุปฺปชฺชติ เหตุปจฺจยา.\csromanspacer{} ตีณิ.}

\prose{เหตุํ อปจยคามิญฺจ นเหตุํ อปจยคามิญฺจ ธมฺมํ ปฏิจฺจ เหตุ อปจยคามี ธมฺโม อุปฺปชฺชติ เหตุปจฺจยา.\csromanspacer{} ตีณิ. (สํขิตฺตํ.)}

\proseitem{315}{เหตุยา นว, อารมฺมเณ นว ฯเปฯ กมฺเม นว, อาหาเร นว ฯเปฯ อวิคเต นว. (สํขิตฺตํ.)}

\csromanheader{2}{\textbf{น-อธิปติ}}
\proseitem{316}{เหตุํ อปจยคามิํ ธมฺมํ ปฏิจฺจ เหตุ อปจยคามี ธมฺโม อุปฺปชฺชติ น-อธิปติปจฺจยา. (สํขิตฺตํ.)}

\vspace{\tipitakaparskip}
\csromancenter{\textbf{น-อธิปติ}ยา ฉ, นปุเรชาเต นว, นปจฺฉาชาเต นว, นกมฺเม ตีณิ, นวิปาเก นว, นวิปฺปยุตฺเต นว. (สํขิตฺตํ.)}
\csromanheader{2}{เหตุปจฺจยา น-อธิปติยา ฉ. (สํขิตฺตํ.)}
\csromanheader{2}{\textbf{น-อธิปติ}ปจฺจยา เหตุยา ฉ. (สํขิตฺตํ.)}
\prose{(สหชาตวาโรปิ ปจฺจยวาโรปิ นิสฺสยวาโรปิ สํสฏฺฐวาโรปิ สมฺปยุตฺตวาโรปิ ปฏิจฺจวารสทิสา วิตฺถาเรตพฺพา.)}

\csromanheader{2}{2. อปจยคามิปท 7. ปญฺหาวาร}
\csromanheader{2}{\textbf{ปจฺจยจตุกฺก-เหตุ อธิปติ}}
\proseitem{318}{เหตุ อปจยคามี ธมฺโม เหตุสฺส อปจยคามิสฺส ธมฺมสฺส เหตุปจฺจเยน ปจฺจโย.\csromanspacer{} ตีณิ.}

\gambhira{ธมฺมานุโลม ทุกติกปฏฺฐานปาฬิ}
\prose{\csromanfolio{266}เหตุ อปจยคามี ธมฺโม เหตุสฺส อปจยคามิสฺส ธมฺมสฺส อธิปติปจฺจเยน ปจฺจโย.\csromanspacer{} สหชาตาธิปติ.\csromanspacer{} ตีณิ.}

\prose{นเหตุ อปจยคามี ธมฺโม นเหตุสฺส อปจยคามิสฺส ธมฺมสฺส อธิปติปจฺจเยน ปจฺจโย.\csromanspacer{} สหชาตาธิปติ.\csromanspacer{} ตีณิ. (สํขิตฺตํ.)}

\proseitem{319}{เหตุยา ตีณิ, อธิปติยา ฉ, สหชาเต นว ฯเปฯ อุปนิสฺสเย นว, กมฺเม ตีณิ, อาหาเร ตีณิ, อินฺทฺริเย นว, ฌาเน ตีณิ, มคฺเค นว, สมฺปยุตฺเต นว, อตฺถิยา นว, อวิคเต นว. (สํขิตฺตํ.)}

\csromanheader{2}{\textbf{ปจฺจนียุทฺธาร}}
\proseitem{320}{เหตุ อปจยคามี ธมฺโม เหตุสฺส อปจยคามิสฺส ธมฺมสฺส สหชาตปจฺจเยน ปจฺจโย.\csromanspacer{} อุปนิสฺสยปจฺจเยน ปจฺจโย. (สํขิตฺตํ.)}

\proseitem{321}{นเหตุยา นว, น-อารมฺมเณ นว. (สํขิตฺตํ.)}

\csromanheader{2}{เหตุปจฺจยา น-อารมฺมเณ ตีณิ. (สํขิตฺตํ.)}
\csromanheader{2}{นเหตุปจฺจยา อธิปติยา ตีณิ. (สํขิตฺตํ.)}
\prose{(ยถา กุสลตฺติเก ปญฺหาวารสฺส อนุโลมมฺปิ ปจฺจนียมฺปิ อนุโลมปจฺจนียมฺปิ ปจฺจนียานุโลมมฺปิ คณิตํ, เอวํ คเณตพฺพํ.)}

\csromanheader{2}{3. เนวาจยคามินาปจยคามิปท 1. ปฏิจฺจวาร}
\csromanheader{2}{\textbf{ปจฺจยจตุกฺก-เหตุ}}
\proseitem{322}{เหตุํ เนวาจยคามินาปจยคามิํ ธมฺมํ ปฏิจฺจ เหตุ เนวาจยคามินาปจยคามี ธมฺโม อุปฺปชฺชติ เหตุปจฺจยา.\csromanspacer{} ตีณิ.}

\prose{นเหตุํ เนวาจยคามินาปจยคามิํ ธมฺมํ ปฏิจฺจ นเหตุ เนวาจยคามินาปจยคามี ธมฺโม อุปฺปชฺชติ เหตุปจฺจยา.\csromanspacer{} ตีณิ.}

\prose{เหตุํ เนวาจยคามินาปจยคามิญฺจ นเหตุํ เนวาจยคามินาปจยคามิญฺจ ธมฺมํ ปฏิจฺจ เหตุ เนวาจยคามินาปจยคามี ธมฺโม อุปฺปชฺชติ เหตุปจฺจยา.\csromanspacer{} ตีณิ. (สํขิตฺตํ.)}

\csromanheader{2}{1. เหตุทุก 10. อาจยคามิตฺติก}
\proseitem{323}{\csromanfolio{267}เหตุยา นว, อารมฺมเณ นว, อธิปติยา นว ฯเปฯ อุปนิสฺสเย นว, ปุเรชาเต นว, อาเสวเน นว, กมฺเม นว, วิปาเก นว, อาหาเร นว, อินฺทฺริเย นว ฯเปฯ อวิคเต นว. (สํขิตฺตํ.)}

\csromanheader{2}{\textbf{นเหตฺวาทิ}}
\proseitem{324}{นเหตุํ เนวาจยคามินาปจยคามิํ ธมฺมํ ปฏิจฺจ นเหตุ เนวาจยคามินาปจยคามี ธมฺโม อุปฺปชฺชติ นเหตุปจฺจยา. (1)}

\prose{เหตุํ เนวาจยคามินาปจยคามิํ ธมฺมํ ปฏิจฺจ นเหตุ เนวาจยคามินาปจยคามี ธมฺโม อุปฺปชฺชติ น-อารมฺมณปจฺจยา. (1)}

\prose{นเหตุํ เนวาจยคามินาปจยคามิํ ธมฺมํ ปฏิจฺจ นเหตุ เนวาจยคามินาปจยคามี ธมฺโม อุปฺปชฺชติ น-อารมฺมณปจฺจยา. (1)}

\prose{เหตุํ เนวาจยคามินาปจยคามิญฺจ นเหตุํ เนวาจยคามินาปจยคามิญฺจ ธมฺมํ ปฏิจฺจ นเหตุ เนวาจยคามินาปจยคามี ธมฺโม อุปฺปชฺชติ น-อารมฺมณปจฺจยา. (1)}

\prose{เหตุํ เนวาจยคามินาปจยคามิํ ธมฺมํ ปฏิจฺจ เหตุ เนวาจยคามินาปจยคามี ธมฺโม อุปฺปชฺชติ น-อธิปติปจฺจยา ฯเปฯ.}

\prose{เหตุํ เนวาจยคามินาปจยคามิํ ธมฺมํ ปฏิจฺจ เหตุ เนวาจยคามินาปจยคามี ธมฺโม อุปฺปชฺชติ นปุเรชาตปจฺจยา.\csromanspacer{} ตีณิ.}

\prose{เหตุํ เนวาจยคามินาปจยคามิํ ธมฺมํ ปฏิจฺจ เหตุ เนวาจยคามินาปจยคามี ธมฺโม อุปฺปชฺชติ นปจฺฉาชาตปจฺจยา ฯเปฯ.}

\prose{เหตุํ เนวาจยคามินาปจยคามิํ ธมฺมํ ปฏิจฺจ นเหตุ เนวาจยคามินาปจยคามี ธมฺโม อุปฺปชฺชติ นกมฺมปจฺจยา. (สํขิตฺตํ.)}

\proseitem{325}{นเหตุยา เอกํ, น-อารมฺมเณ ตีณิ, น-อธิปติยา นว, น-อนนฺตเร ตีณิ, นสมนนฺตเร ตีณิ, น-อญฺญมญฺเญ ตีณิ, น-อุปนิสฺสเย ตีณิ, นปุเรชาเต นว, นปจฺฉาชาเต นว, น-อาเสวเน นว, นกมฺเม ตีณิ, นวิปาเก นว, น-อาหาเร เอกํ, น-อินฺทฺริเย เอกํ, นฌาเน เอกํ, นมคฺเค เอกํ, นสมฺปยุตฺเต ตีณิ, นวิปฺปยุตฺเต นว, โนนตฺถิยา ตีณิ, โนวิคเต ตีณิ. (สํขิตฺตํ.)}

\gambhira{ธมฺมานุโลม ทุกติกปฏฺฐานปาฬิ}
\csromanheader{2}{เหตุปจฺจยา น-อารมฺมเณ ตีณิ. (สํขิตฺตํ.)}
\csromanheader{2}{นเหตุปจฺจยา อารมฺมเณ เอกํ. (สํขิตฺตํ.)}
\prose{\csromanfolio{268}(สหชาตวาโรปิ ปจฺจยวาโรปิ นิสฺสยวาโรปิ สํสฏฺฐวาโรปิ สมฺปยุตฺตวาโรปิ ปฏิจฺจวารสทิสา วิตฺถาเรตพฺพา.)}

\csromanheader{2}{3. เนวาจยคามินาปจยคามิปท 7. ปญฺหาวาร}
\csromanheader{2}{\textbf{ปจฺจยจตุกฺก-เหตุ อารมฺมณ}}
\proseitem{326}{เหตุ เนวาจยคามินาปจยคามี ธมฺโม เหตุสฺส เนวาจยคามินาปจยคามิสฺส ธมฺมสฺส เหตุปจฺจเยน ปจฺจโย.\csromanspacer{} ตีณิ.}

\prose{เหตุ เนวาจยคามินาปจยคามี ธมฺโม เหตุสฺส เนวาจยคามินาปจยคามิสฺส ธมฺมสฺส อารมฺมณปจฺจเยน ปจฺจโย. (สํขิตฺตํ.)}

\proseitem{327}{เหตุยา ตีณิ, อารมฺมเณ นว, อธิปติยา นว, อนนฺตเร นว, สมนนฺตเร นว, สหชาเต นว, อญฺญมญฺเญ นว, นิสฺสเย นว, อุปนิสฺสเย นว, ปุเรชาเต ตีณิ, ปจฺฉาชาเต ตีณิ, อาเสวเน นว, กมฺเม ตีณิ, วิปาเก นว, อาหาเร ตีณิ, อินฺทฺริเย นว, ฌาเน ตีณิ, มคฺเค นว, สมฺปยุตฺเต นว, วิปฺปยุตฺเต ปญฺจ, อตฺถิยา นว, นตฺถิยา นว, วิคเต นว, อวิคเต นว. (สํขิตฺตํ.)}

\csromanheader{2}{\textbf{ปจฺจนียุทฺธาร}}
\proseitem{328}{เหตุ เนวาจยคามินาปจยคามี ธมฺโม เหตุสฺส เนวาจยคามินาปจยคามิสฺส ธมฺมสฺส อารมฺมณปจฺจเยน ปจฺจโย.\csromanspacer{} สหชาตปจฺจเยน ปจฺจโย.\csromanspacer{} อุปนิสฺสยปจฺจเยน ปจฺจโย.\csromanspacer{} ตีณิ.}

\prose{นเหตุ เนวาจยคามินาปจยคามี ธมฺโม นเหตุสฺส เนวาจยคามินาปจยคามิสฺส ธมฺมสฺส อารมฺมณปจฺจเยน ปจฺจโย.\csromanspacer{} สหชาตปจฺจเยน ปจฺจโย.\csromanspacer{} อุปนิสฺสยปจฺจเยน ปจฺจโย.\csromanspacer{} ปุเรชาตปจฺจเยน ปจฺจโย.\csromanspacer{} ปจฺฉาชาตปจฺจเยน ปจฺจโย.\csromanspacer{} อาหารปจฺจเยน ปจฺจโย.\csromanspacer{} อินฺทฺริยปจฺจเยน ปจฺจโย. (สํขิตฺตํ.)}

\csromanmatikaanchor{mtk.42}
\csromantocmarkref{section}{11. เสกฺขตฺติก}{mtk.42}
\bookmark[dest={mtk.42},level=1]{11. เสกฺขตฺติก}
\csromanheader{1}{1. เหตุทุก 11. เสกฺขตฺติก}
\proseitem{329}{\csromanfolio{269}นเหตุยา นว, น-อารมฺมเณ นว. (สํขิตฺตํ.)}

\csromanheader{2}{เหตุปจฺจยา น-อารมฺมเณ ตีณิ. (สํขิตฺตํ.)}
\csromanheader{2}{นเหตุปจฺจยา อารมฺมเณ นว. (สํขิตฺตํ.)}
\prose{(ยถา กุสลตฺติเก ปญฺหาวารสฺส อนุโลมมฺปิ ปจฺจนียมฺปิ อนุโลมปจฺจนียมฺปิ ปจฺจนียานุโลมมฺปิ คณิตํ, เอวํ คเณตพฺพํ.)}

\nitthitam{เหตุทุก-อาจยคามิตฺติกํ นิฏฺฐิตํ.}
\csromanheader{2}{1. เหตุทุก 11. เสกฺขตฺติก}
\csromanheader{2}{1. เสกฺขปท 1. ปฏิจฺจวาร}
\csromanheader{2}{\textbf{ปจฺจยจตุกฺก-เหตุ}}
\proseitem{330}{เหตุํ เสกฺขํ ธมฺมํ ปฏิจฺจ เหตุ เสกฺโข ธมฺโม อุปฺปชฺชติ เหตุปจฺจยา.\csromanspacer{} ตีณิ.}

\prose{นเหตุํ เสกฺขํ ธมฺมํ ปฏิจฺจ นเหตุ เสกฺโข ธมฺโม อุปฺปชฺชติ เหตุปจฺจยา.\csromanspacer{} ตีณิ.}

\prose{เหตุํ เสกฺขญฺจ นเหตุํ เสกฺขญฺจ ธมฺมํ ปฏิจฺจ เหตุ เสกฺโข ธมฺโม อุปฺปชฺชติ เหตุปจฺจยา.\csromanspacer{} ตีณิ. (สํขิตฺตํ.)}

\proseitem{331}{เหตุยา นว, อารมฺมเณ นว, อธิปติยา นว, อนนฺตเร นว, สมนนฺตเร นว, สหชาเต นว, อญฺญมญฺเญ นว, นิสฺสเย นว, อุปนิสฺสเย นว, ปุเรชาเต นว, อาเสวเน นว, กมฺเม นว, วิปาเก นว, อาหาเร นว ฯเปฯ อวิคเต นว. (สํขิตฺตํ.)}

\csromanheader{2}{\textbf{น-อธิปติ}}
\proseitem{332}{เหตุํ เสกฺขํ ธมฺมํ ปฏิจฺจ เหตุ เสกฺโข ธมฺโม อุปฺปชฺชติ น-อธิปติปจฺจยา. (สํขิตฺตํ.)}

\gambhira{ธมฺมานุโลม ทุกติกปฏฺฐานปาฬิ}
\proseitem{333}{\csromanfolio{270}น-อธิปติยา ฉ, นปุเรชาเต นว, นปจฺฉาชาเต นว, น-อาเสวเน นว, นกมฺเม ตีณิ, นวิปาเก นว, นวิปฺปยุตฺเต นว.}

\csromanheader{2}{เหตุปจฺจยา น-อธิปติยา ฉ. (สํขิตฺตํ.)}
\csromanheader{2}{น-อธิปติปจฺจยา เหตุยา ฉ. (สํขิตฺตํ.)}
\prose{(สหชาตวาโรปิ ปจฺจยวาโรปิ นิสฺสยวาโรปิ สํสฏฺฐวาโรปิ สมฺปยุตฺตวาโรปิ ปฏิจฺจวารสทิสา วิตฺถาเรตพฺพา.)}

\csromanheader{2}{1. เสกฺขปท 7. ปญฺหาวาร}
\csromanheader{2}{\textbf{ปจฺจยจตุกฺก-เหตฺวาทิ}}
\proseitem{334}{เหตุ เสกฺโข ธมฺโม เหตุสฺส เสกฺขสฺส ธมฺมสฺส เหตุปจฺจเยน ปจฺจโย.\csromanspacer{} ตีณิ.}

\prose{เหตุ เสกฺโข ธมฺโม เหตุสฺส เสกฺขสฺส ธมฺมสฺส อธิปติปจฺจเยน ปจฺจโย.\csromanspacer{} สหชาตาธิปติ.\csromanspacer{} ตีณิ.}

\prose{นเหตุ เสกฺโข ธมฺโม นเหตุสฺส เสกฺขสฺส ธมฺมสฺส อธิปติปจฺจเยน ปจฺจโย.\csromanspacer{} สหชาตาธิปติ.\csromanspacer{} ตีณิ.}

\prose{เหตุ เสกฺโข ธมฺโม เหตุสฺส เสกฺขสฺส ธมฺมสฺส อนนฺตรปจฺจเยน ปจฺจโย. (สํขิตฺตํ.)}

\proseitem{335}{เหตุยา ตีณิ, อธิปติยา ฉ, อนนฺตเร นว, สมนนฺตเร นว, สหชาเต นว, อญฺญมญฺเญ นว, นิสฺสเย นว, อุปนิสฺสเย นว, กมฺเม ตีณิ, วิปาเก นว, อาหาเร ตีณิ, อินฺทฺริเย นว, ฌาเน ตีณิ, มคฺเค นว, สมฺปยุตฺเต นว, อตฺถิยา นว, นตฺถิยา นว, วิคเต นว, อวิคเต นว. (สํขิตฺตํ.)}

\csromanheader{2}{\textbf{ปจฺจนียุทฺธาร}}
\proseitem{336}{เหตุ เสกฺโข ธมฺโม เหตุสฺส เสกฺขสฺส ธมฺมสฺส สหชาตปจฺจเยน ปจฺจโย.\csromanspacer{} อุปนิสฺสยปจฺจเยน ปจฺจโย. (สํขิตฺตํ.)}

\csromanheader{2}{1. เหตุทุก 11. เสกฺขตฺติก}
\proseitem{337}{\csromanfolio{271}นเหตุยา นว, น-อารมฺมเณ นว. (สํขิตฺตํ.)}

\csromanheader{2}{เหตุปจฺจยา น-อารมฺมเณ ตีณิ. (สํขิตฺตํ.)}
\csromanheader{2}{นเหตุปจฺจยา อธิปติยา ตีณิ. (สํขิตฺตํ.)}
\prose{(ยถา กุสลตฺติเก ปญฺหาวารสฺส อนุโลมมฺปิ ปจฺจนียมฺปิ อนุโลมปจฺจนียมฺปิ ปจฺจนียานุโลมมฺปิ คณิตํ, เอวํ คเณตพฺพํ.)}

\csromanheader{2}{2. อเสกฺขปท 1. ปฏิจฺจวาร}
\csromanheader{2}{\textbf{ปจฺจยจตุกฺก-เหตุ}}
\proseitem{338}{เหตุํ อเสกฺขํ ธมฺมํ ปฏิจฺจ เหตุ อเสกฺโข ธมฺโม อุปฺปชฺชติ เหตุปจฺจยา.\csromanspacer{} ตีณิ.}

\prose{นเหตุํ อเสกฺขํ ธมฺมํ ปฏิจฺจ นเหตุ อเสกฺโข ธมฺโม อุปฺปชฺชติ เหตุปจฺจยา.\csromanspacer{} ตีณิ.}

\prose{เหตุํ อเสกฺขญฺจ นเหตุํ อเสกฺขญฺจ ธมฺมํ ปฏิจฺจ เหตุ อเสกฺโข ธมฺโม อุปฺปชฺชติ เหตุปจฺจยา.\csromanspacer{} ตีณิ. (สํขิตฺตํ.)}

\proseitem{339}{เหตุยา นว, อารมฺมเณ นว, อธิปติยา นว, อนนฺตเร นว, สมนนฺตเร นว, สหชาเต นว, อญฺญมญฺเญ นว, นิสฺสเย นว, อุปนิสฺสเย นว, ปุเรชาเต นว, กมฺเม นว, วิปาเก นว ฯเปฯ อวิคเต นว. (สํขิตฺตํ.)}

\csromanheader{2}{\textbf{น-อธิปติ}}
\proseitem{340}{เหตุํ อเสกฺขํ ธมฺมํ ปฏิจฺจ เหตุ อเสกฺโข ธมฺโม อุปฺปชฺชติ น-อธิปติปจฺจยา. (สํขิตฺตํ.)}

\vspace{\tipitakaparskip}
\csromancenter{\textbf{น-อธิปติ}ยา ฉ, นปุเรชาเต นว, นปจฺฉาชาเต นว, น-อาเสวเน นว, นวิปฺปยุตฺเต นว. (สํขิตฺตํ.)}
\csromanheader{2}{เหตุปจฺจยา น-อธิปติยา ฉ. (สํขิตฺตํ.)}
\csromanheader{2}{\textbf{น-อธิปติ}ปจฺจยา เหตุยา ฉ. (สํขิตฺตํ.)}
\gambhira{ธมฺมานุโลม ทุกติกปฏฺฐานปาฬิ}
\prose{\csromanfolio{272}(สหชาตวาโรปิ ปจฺจยวาโรปิ นิสฺสยวาโรปิ สํสฏฺฐวาโรปิ สมฺปยุตฺตวาโรปิ ปฏิจฺจวารสทิสา วิตฺถาเรตพฺพา.)}

\csromanheader{2}{2. อเสกฺขปท 7. ปญฺหาวาร}
\csromanheader{2}{\textbf{ปจฺจยจตุกฺก-เหตฺวาทิ}}
\proseitem{342}{เหตุ อเสกฺโข ธมฺโม เหตุสฺส อเสกฺขสฺส ธมฺมสฺส เหตุปจฺจเยน ปจฺจโย.\csromanspacer{} ตีณิ.}

\prose{เหตุ อเสกฺโข ธมฺโม เหตุสฺส อเสกฺขสฺส ธมฺมสฺส อธิปติปจฺจเยน ปจฺจโย.\csromanspacer{} สหชาตาธิปติ.\csromanspacer{} ตีณิ.}

\prose{นเหตุ อเสกฺโข ธมฺโม นเหตุสฺส อเสกฺขสฺส ธมฺมสฺส อธิปติปจฺจเยน ปจฺจโย.\csromanspacer{} สหชาตาธิปติ.\csromanspacer{} ตีณิ ฯเปฯ.}

\prose{เหตุ อเสกฺโข ธมฺโม เหตุสฺส อเสกฺขสฺส ธมฺมสฺส อุปนิสฺสยปจฺจเยน ปจฺจโย.\csromanspacer{} อนนฺตรูปนิสฺสโย. (สํขิตฺตํ.)}

\proseitem{343}{เหตุยา ตีณิ, อธิปติยา ฉ, อนนฺตเร นว, สมนนฺตเร นว, สหชาเต นว, อญฺญมญฺเญ นว, นิสฺสเย นว, อุปนิสฺสเย นว, กมฺเม ตีณิ, วิปาเก นว, อาหาเร ตีณิ, อินฺทฺริเย นว, ฌาเน ตีณิ, มคฺเค นว, สมฺปยุตฺเต นว, อตฺถิยา นว, นตฺถิยา นว, วิคเต นว, อวิคเต นว.}

\csromanheader{2}{\textbf{ปจฺจนียุทฺธาร}}
\proseitem{344}{เหตุ อเสกฺโข ธมฺโม เหตุสฺส อเสกฺขสฺส ธมฺมสฺส สหชาตปจฺจเยน ปจฺจโย.\csromanspacer{} อุปนิสฺสยปจฺจเยน ปจฺจโย. (สํขิตฺตํ.)}

\proseitem{345}{นเหตุยา นว, น-อารมฺมเณ นว. (สํขิตฺตํ.)}

\csromanheader{2}{เหตุปจฺจยา น-อารมฺมเณ ตีณิ. (สํขิตฺตํ.)}
\csromanheader{2}{นเหตุปจฺจยา อธิปติยา ตีณิ. (สํขิตฺตํ.)}
\prose{(ยถา กุสลตฺติเก ปญฺหาวารสฺส อนุโลมมฺปิ ปจฺจนียมฺปิ อนุโลมปจฺจนียมฺปิ ปจฺจนียานุโลมมฺปิ คณิตํ, เอวํ คเณตพฺพํ.)}

\vspace{\tipitakaparskip}
\csromancenter{เหตุทุก 11.\csromanspacer{} เสกฺขตฺติก \textbf{3.}\csromanspacer{}\textbf{ เนวเสกฺขนาเสกฺขปท 1.}\csromanspacer{}\textbf{ ปฏิจฺจวาร}}
\csromanheader{2}{\textbf{ปจฺจยจตุกฺก-เหตุ}}
\proseitem{346}{\csromanfolio{273}เหตุํ เนวเสกฺขนาเสกฺขํ ธมฺมํ ปฏิจฺจ เหตุ เนวเสกฺขนาเสกฺโข ธมฺโม อุปฺปชฺชติ เหตุปจฺจยา.\csromanspacer{} ตีณิ.}

\prose{นเหตุํ เนวเสกฺขนาเสกฺขํ ธมฺมํ ปฏิจฺจ นเหตุ เนวเสกฺขนาเสกฺโข ธมฺโม อุปฺปชฺชติ เหตุปจฺจยา.\csromanspacer{} ตีณิ.}

\prose{เหตุํ เนวเสกฺขนาเสกฺขญฺจ นเหตุํ เนวเสกฺขนาเสกฺขญฺจ ธมฺมํ ปฏิจฺจ เหตุ เนวเสกฺขนาเสกฺโข ธมฺโม อุปฺปชฺชติ เหตุปจฺจยา.\csromanspacer{} ตีณิ. (สํขิตฺตํ.)}

\proseitem{347}{เหตุยา นว, อารมฺมเณ นว, อธิปติยา นว, อนนฺตเร นว, สมนนฺตเร นว, สหชาเต นว, อญฺญมญฺเญ นว, นิสฺสเย นว, อุปนิสฺสเย นว, ปุเรชาเต นว, อาเสวเน นว, กมฺเม นว, วิปาเก นว ฯเปฯ อวิคเต นว. (สํขิตฺตํ.)}

\csromanheader{2}{\textbf{นเหตุ น-อารมฺมณาทิ}}
\proseitem{348}{นเหตุํ เนวเสกฺขนาเสกฺขํ ธมฺมํ ปฏิจฺจ นเหตุ เนวเสกฺขนาเสกฺโข ธมฺโม อุปฺปชฺชติ นเหตุปจฺจยา.\csromanspacer{}\csromanspacer{}นเหตุํ เนวเสกฺขนาเสกฺขํ ธมฺมํ ปฏิจฺจ เหตุ เนวเสกฺขนาเสกฺโข ธมฺโม อุปฺปชฺชติ นเหตุปจฺจยา. (2)}

\prose{เหตุํ เนวเสกฺขนาเสกฺขํ ธมฺมํ ปฏิจฺจ นเหตุ เนวเสกฺขนาเสกฺโข ธมฺโม อุปฺปชฺชติ น-อารมฺมณปจฺจยา. (1)}

\prose{นเหตุํ เนวเสกฺขนาเสกฺขํ ธมฺมํ ปฏิจฺจ นเหตุ เนวเสกฺขนาเสกฺโข ธมฺโม อุปฺปชฺชติ น-อารมฺมณปจฺจยา. (1)}

\prose{เหตุํ เนวเสกฺขนาเสกฺขญฺจ นเหตุํ เนวเสกฺขนาเสกฺขญฺจ ธมฺมํ ปฏิจฺจ นเหตุ เนวเสกฺขนาเสกฺโข ธมฺโม อุปฺปชฺชติ น-อารมฺมณปจฺจยา. (1)}

\proseitem{349}{เหตุํ เนวเสกฺขนาเสกฺขํ ธมฺมํ ปฏิจฺจ เหตุ เนวเสกฺขนาเสกฺโข ธมฺโม อุปฺปชฺชติ น-อธิปติปจฺจยา ฯเปฯ.}

\gambhira{ธมฺมานุโลม ทุกติกปฏฺฐานปาฬิ}
\prose{\csromanfolio{274}เหตุํ เนวเสกฺขนาเสกฺขํ ธมฺมํ ปฏิจฺจ เหตุ เนวเสกฺขนาเสกฺโข ธมฺโม อุปฺปชฺชติ นปุเรชาตปจฺจยา.\csromanspacer{} ตีณิ. (สํขิตฺตํ.)}

\proseitem{350}{นเหตุยา เทฺว, น-อารมฺมเณ ตีณิ, น-อธิปติยา นว, น-อนนฺตเร ตีณิ, นสมนนฺตเร ตีณิ, น-อญฺญมญฺเญ ตีณิ, น-อุปนิสฺสเย ตีณิ, นปุเรชาเต นว, นปจฺฉาชาเต นว, น-อาเสวเน นว, นกมฺเม ตีณิ, นวิปาเก นว, น-อาหาเร เอกํ, น-อินฺทฺริเย เอกํ, นฌาเน เอกํ, นมคฺเค เอกํ, นสมฺปยุตฺเต ตีณิ, นวิปฺปยุตฺเต นว, โนนตฺถิยา ตีณิ, โนวิคเต ตีณิ. (สํขิตฺตํ.)}

\csromanheader{2}{เหตุปจฺจยา น-อารมฺมเณ ตีณิ. (สํขิตฺตํ.)}
\csromanheader{2}{นเหตุปจฺจยา อารมฺมเณ เทฺว. (สํขิตฺตํ.)}
\prose{(สหชาตวาโรปิ ปจฺจยวาโรปิ นิสฺสยวาโรปิ สํสฏฺฐวาโรปิ สมฺปยุตฺตวาโรปิ ปฏิจฺจวารสทิสา วิตฺถาเรตพฺพา.)}

\csromanheader{2}{3. เนวเสกฺขนาเสกฺขปท 7. ปญฺหาวาร}
\csromanheader{2}{\textbf{ปจฺจยจตุกฺก-เหตฺวาทิ}}
\proseitem{351}{เหตุ เนวเสกฺขนาเสกฺโข ธมฺโม เหตุสฺส เนวเสกฺขนาเสกฺขสฺส ธมฺมสฺส เหตุปจฺจเยน ปจฺจโย.\csromanspacer{} ตีณิ.}

\prose{เหตุ เนวเสกฺขนาเสกฺโข ธมฺโม เหตุสฺส เนวเสกฺขนาเสกฺขสฺส ธมฺมสฺส อารมฺมณปจฺจเยน ปจฺจโย.\csromanspacer{} นว.}

\prose{เหตุ เนวเสกฺขนาเสกฺโข ธมฺโม เหตุสฺส เนวเสกฺขนาเสกฺขสฺส ธมฺมสฺส อธิปติปจฺจเยน ปจฺจโย.\csromanspacer{} อารมฺมณาธิปติ, สหชาตาธิปติ.\csromanspacer{} นว ฯเปฯ.}

\prose{นเหตุ เนวเสกฺขนาเสกฺโข ธมฺโม นเหตุสฺส เนวเสกฺขนาเสกฺขสฺส ธมฺมสฺส ปุเรชาตปจฺจเยน ปจฺจโย. (สํขิตฺตํ.)}

\proseitem{352}{เหตุยา ตีณิ, อารมฺมเณ นว, อธิปติยา นว, อนนฺตเร นว, สมนนฺตเร นว, สหชาเต นว, อญฺญมญฺเญ นว, นิสฺสเย นว,}

\vspace{\tipitakaparskip}
\csromancenter{เหตุทุก 12.\csromanspacer{} ปริตฺตตฺติก อุปนิสฺสเย นว, ปุเรชาเต ตีณิ, ปจฺฉาชาเต ตีณิ, อาเสวเน นว, กมฺเม ตีณิ, วิปาเก นว, อาหาเร ตีณิ, อินฺทฺริเย นว, ฌาเน ตีณิ, มคฺเค นว, สมฺปยุตฺเต นว, วิปฺปยุตฺเต ปญฺจ, อตฺถิยา นว, นตฺถิยา นว, วิคเต นว, อวิคเต นว. (สํขิตฺตํ.)}
\csromanheader{2}{\textbf{ปจฺจนียุทฺธาร}}
\proseitem{353}{\csromanfolio{275}เหตุ เนวเสกฺขนาเสกฺโข ธมฺโม เหตุสฺส เนวเสกฺขนาเสกฺขสฺส ธมฺมสฺส อารมฺมณปจฺจเยน ปจฺจโย.\csromanspacer{} สหชาตปจฺจเยน ปจฺจโย.\csromanspacer{} อุปนิสฺสยปจฺจเยน ปจฺจโย. (สํขิตฺตํ.)}

\proseitem{354}{นเหตุยา นว, น-อารมฺมเณ นว. (สํขิตฺตํ.)}

\csromanheader{2}{เหตุปจฺจยา น-อารมฺมเณ ตีณิ. (สํขิตฺตํ.)}
\csromanheader{2}{นเหตุปจฺจยา อารมฺมเณ นว. (สํขิตฺตํ.)}
\prose{(ยถา กุสลตฺติเก ปญฺหาวารสฺส อนุโลมมฺปิ ปจฺจนียมฺปิ อนุโลมปจฺจนียมฺปิ ปจฺจนียานุโลมมฺปิ คณิตํ, เอวํ คเณตพฺพํ.)}

\nitthitam{เหตุทุกเสกฺขตฺติกํ นิฏฺฐิตํ.}
\csromanmatikaanchor{mtk.43}
\csromantocmarkref{section}{12. ปริตฺตตฺติก}{mtk.43}
\bookmark[dest={mtk.43},level=1]{12. ปริตฺตตฺติก}
\csromanheader{1}{1. เหตุทุก 12. ปริตฺตตฺติก}
\csromanheader{2}{1. ปริตฺตปท 1. ปฏิจฺจวาร}
\csromanheader{2}{\textbf{ปจฺจยจตุกฺก-เหตุ}}
\proseitem{355}{เหตุํ ปริตฺตํ ธมฺมํ ปฏิจฺจ เหตุ ปริตฺโต ธมฺโม อุปฺปชฺชติ เหตุปจฺจยา.\csromanspacer{} ตีณิ.}

\prose{นเหตุํ ปริตฺตํ ธมฺมํ ปฏิจฺจ นเหตุ ปริตฺโต ธมฺโม อุปฺปชฺชติ เหตุปจฺจยา.\csromanspacer{} ตีณิ.}

\prose{เหตุํ ปริตฺตญฺจ นเหตุํ ปริตฺตญฺจ ธมฺมํ ปฏิจฺจ เหตุ ปริตฺโต ธมฺโม อุปฺปชฺชติ เหตุปจฺจยา.\csromanspacer{} ตีณิ. (สํขิตฺตํ.)}

\gambhira{ธมฺมานุโลม ทุกติกปฏฺฐานปาฬิ}
\proseitem{356}{\csromanfolio{276}เหตุยา นว, อารมฺมเณ นว, อธิปติยา นว ฯเปฯ อวิคเต นว. (สํขิตฺตํ.)}

\csromanheader{2}{\textbf{นเหตุ น-อารมฺมณาทิ}}
\proseitem{357}{นเหตุํ ปริตฺตํ ธมฺมํ ปฏิจฺจ นเหตุ ปริตฺโต ธมฺโม อุปฺปชฺชติ นเหตุปจฺจยา.\csromanspacer{}\csromanspacer{}นเหตุํ ปริตฺตํ ธมฺมํ ปฏิจฺจ เหตุ ปริตฺโต ธมฺโม อุปฺปชฺชติ นเหตุปจฺจยา. (2)}

\prose{เหตุํ ปริตฺตํ ธมฺมํ ปฏิจฺจ นเหตุ ปริตฺโต ธมฺโม อุปฺปชฺชติ น-อารมฺมณปจฺจยา. (1)}

\prose{นเหตุํ ปริตฺตํ ธมฺมํ ปฏิจฺจ นเหตุ ปริตฺโต ธมฺโม อุปฺปชฺชติ น-อารมฺมณปจฺจยา. (1)}

\prose{เหตุํ ปริตฺตญฺจ นเหตุํ ปริตฺตญฺจ ธมฺมํ ปฏิจฺจ นเหตุ ปริตฺโต ธมฺโม อุปฺปชฺชติ น-อารมฺมณปจฺจยา. (1)}

\prose{เหตุํ ปริตฺตํ ธมฺมํ ปฏิจฺจ เหตุ ปริตฺโต ธมฺโม อุปฺปชฺชติ น-อธิปติปจฺจยา. (สํขิตฺตํ.)}

\proseitem{358}{นเหตุยา เทฺว, น-อารมฺมเณ ตีณิ, น-อธิปติยา นว, น-อนนฺตเร ตีณิ, นสมนนฺตเร ตีณิ, น-อญฺญมญฺเญ ตีณิ, น-อุปนิสฺสเย ตีณิ, นปุเรชาเต นว, นปจฺฉาชาเต นว, น-อาเสวเน นว, นกมฺเม ตีณิ, นวิปาเก นว, น-อาหาเร เอกํ, น-อินฺทฺริเย เอกํ, นฌาเน เอกํ, นมคฺเค เอกํ, นสมฺปยุตฺเต ตีณิ, นวิปฺปยุตฺเต นว, โนนตฺถิยา ตีณิ, โนวิคเต ตีณิ. (สํขิตฺตํ.)}

\csromanheader{2}{เหตุปจฺจยา น-อารมฺมเณ ตีณิ. (สํขิตฺตํ.)}
\csromanheader{2}{นเหตุปจฺจยา อารมฺมเณ เทฺว. (สํขิตฺตํ.)}
\prose{(สหชาตวาโรปิ ปจฺจยวาโรปิ นิสฺสยวาโรปิ สํสฏฺฐวาโรปิ สมฺปยุตฺตวาโรปิ ปฏิจฺจวารสทิสา วิตฺถาเรตพฺพา.)}

\csromanheader{2}{1. เหตุทุก 12. ปริตฺตตฺติก \textbf{1. ปริตฺตปท 7. ปญฺหาวาร}}
\csromanheader{2}{\textbf{ปจฺจยจตุกฺก-เหตุ อารมฺมณ}}
\proseitem{359}{\csromanfolio{277}เหตุ ปริตฺโต ธมฺโม เหตุสฺส ปริตฺตสฺส ธมฺมสฺส เหตุปจฺจเยน ปจฺจโย.\csromanspacer{} ตีณิ.}

\prose{เหตุ ปริตฺโต ธมฺโม เหตุสฺส ปริตฺตสฺส ธมฺมสฺส อารมฺมณปจฺจเยน ปจฺจโย. (สํขิตฺตํ.)}

\proseitem{360}{เหตุยา ตีณิ, อารมฺมเณ นว, อธิปติยา นว, อนนฺตเร นว, สมนนฺตเร นว ฯเปฯ อวิคเต นว. (สํขิตฺตํ.)}

\csromanheader{2}{\textbf{ปจฺจนียุทฺธาร}}
\proseitem{361}{เหตุ ปริตฺโต ธมฺโม เหตุสฺส ปริตฺตสฺส ธมฺมสฺส อารมฺมณปจฺจเยน ปจฺจโย.\csromanspacer{} สหชาตปจฺจเยน ปจฺจโย.\csromanspacer{} อุปนิสฺสยปจฺจเยน ปจฺจโย. (สํขิตฺตํ.)}

\proseitem{362}{นเหตุยา นว, น-อารมฺมเณ นว. (สํขิตฺตํ.)}

\csromanheader{2}{เหตุปจฺจยา น-อารมฺมเณ ตีณิ. (สํขิตฺตํ.)}
\csromanheader{2}{นเหตุปจฺจยา อารมฺมเณ นว. (สํขิตฺตํ.)}
\prose{(ยถา กุสลตฺติเก ปญฺหาวารสฺส อนุโลมมฺปิ ปจฺจนียมฺปิ อนุโลมปจฺจนียมฺปิ ปจฺจนียานุโลมมฺปิ คณิตํ, เอวํ คเณตพฺพํ.)}

\csromanheader{2}{2. มหคฺคตปท 1. ปฏิจฺจวาร}
\csromanheader{2}{\textbf{ปจฺจยจตุกฺก-เหตุ}}
\proseitem{363}{เหตุํ มหคฺคตํ ธมฺมํ ปฏิจฺจ เหตุ มหคฺคโต ธมฺโม อุปฺปชฺชติ เหตุปจฺจยา.\csromanspacer{} ตีณิ.}

\gambhira{ธมฺมานุโลม ทุกติกปฏฺฐานปาฬิ}
\prose{\csromanfolio{278}นเหตุํ มหคฺคตํ ธมฺมํ ปฏิจฺจ นเหตุ มหคฺคโต ธมฺโม อุปฺปชฺชติ เหตุปจฺจยา.\csromanspacer{} ตีณิ.}

\prose{เหตุํ มหคฺคตญฺจ นเหตุํ มหคฺคตญฺจ ธมฺมํ ปฏิจฺจ เหตุ มหคฺคโต ธมฺโม อุปฺปชฺชติ เหตุปจฺจยา.\csromanspacer{} ตีณิ. (สํขิตฺตํ.)}

\proseitem{364}{เหตุยา นว, อารมฺมเณ นว, อธิปติยา นว, อนนฺตเร นว, สมนนฺตเร นว, สหชาเต นว, อญฺญมญฺเญ นว, นิสฺสเย นว, อุปนิสฺสเย นว, ปุเรชาเต นว, อาเสวเน นว, กมฺเม นว, วิปาเก นว, อาหาเร นว ฯเปฯ อวิคเต นว. (สํขิตฺตํ.)}

\csromanheader{2}{\textbf{น-อธิปติ}}
\proseitem{365}{เหตุํ มหคฺคตํ ธมฺมํ ปฏิจฺจ เหตุ มหคฺคโต ธมฺโม อุปฺปชฺชติ น-อธิปติปจฺจยา. (สํขิตฺตํ.)}

\vspace{\tipitakaparskip}
\csromancenter{\textbf{น-อธิปติ}ยา นว, นปุเรชาเต นว, นปจฺฉาชาเต นว, น-อาเสวเน นว, นกมฺเม ตีณิ, นวิปาเก นว, นวิปฺปยุตฺเต นว. (สํขิตฺตํ.)}
\csromanheader{2}{เหตุปจฺจยา น-อธิปติยา นว. (สํขิตฺตํ.)}
\csromanheader{2}{\textbf{น-อธิปติ}ปจฺจยา เหตุยา นว. (สํขิตฺตํ.)}
\prose{(สหชาตวาโรปิ ปจฺจยวาโรปิ นิสฺสยวาโรปิ สํสฏฺฐวาโรปิ สมฺปยุตฺตวาโรปิ ปฏิจฺจวารสทิสา วิตฺถาเรตพฺพา.)}

\csromanheader{2}{2. มหคฺคตปท 7. ปญฺหาวาร}
\csromanheader{2}{\textbf{ปจฺจยจตุกฺก-เหตฺวาทิ}}
\proseitem{367}{เหตุ มหคฺคโต ธมฺโม เหตุสฺส มหคฺคตสฺส ธมฺมสฺส เหตุปจฺจเยน ปจฺจโย.\csromanspacer{} ตีณิ.}

\prose{เหตุ มหคฺคโต ธมฺโม เหตุสฺส มหคฺคตสฺส ธมฺมสฺส อารมฺมณปจฺจเยน ปจฺจโย.\csromanspacer{} ตีณิ.}

\prose{นเหตุ มหคฺคโต ธมฺโม นเหตุสฺส มหคฺคตสฺส ธมฺมสฺส อารมฺมณปจฺจเยน ปจฺจโย.\csromanspacer{} ตีณิ.}

\csromanheader{2}{1. เหตุทุก 12. ปริตฺตตฺติก}
\prose{\csromanfolio{279}เหตุ มหคฺคโต จ นเหตุ มหคฺคโต จ ธมฺมา เหตุสฺส มหคฺคตสฺส ธมฺมสฺส อารมฺมณปจฺจเยน ปจฺจโย.\csromanspacer{} ตีณิ.}

\prose{เหตุ มหคฺคโต ธมฺโม เหตุสฺส มหคฺคตสฺส ธมฺมสฺส อธิปติปจฺจเยน ปจฺจโย.\csromanspacer{} สหชาตาธิปติ.\csromanspacer{} ตีณิ.}

\prose{นเหตุ มหคฺคโต ธมฺโม นเหตุสฺส มหคฺคตสฺส ธมฺมสฺส อธิปติปจฺจเยน ปจฺจโย.\csromanspacer{} สหชาตาธิปติ.\csromanspacer{} ตีณิ.}

\prose{เหตุ มหคฺคโต ธมฺโม เหตุสฺส มหคฺคตสฺส ธมฺมสฺส อนนฺตรปจฺจเยน ปจฺจโย. (สํขิตฺตํ.)}

\proseitem{368}{เหตุยา ตีณิ, อารมฺมเณ นว, อธิปติยา ฉ, อนนฺตเร นว, สมนนฺตเร นว, สหชาเต นว, อญฺญมญฺเญ นว, นิสฺสเย นว, อุปนิสฺสเย นว, อาเสวเน นว, กมฺเม ตีณิ, วิปาเก นว, อาหาเร ตีณิ, อินฺทฺริเย นว, ฌาเน ตีณิ, มคฺเค นว, สมฺปยุตฺเต นว, อตฺถิยา นว ฯเปฯ อวิคเต นว. (สํขิตฺตํ.)}

\csromanheader{2}{\textbf{ปจฺจนียุทฺธาร}}
\proseitem{369}{เหตุ มหคฺคโต ธมฺโม เหตุสฺส มหคฺคตสฺส ธมฺมสฺส อรมฺมณปจฺจเยน ปจฺจโย.\csromanspacer{} สหชาตปจฺจเยน ปจฺจโย.\csromanspacer{} อุปนิสฺสยปจฺจเยน ปจฺจโย. (สํขิตฺตํ.)}

\proseitem{370}{นเหตุยา นว, น-อารมฺมเณ นว. (สํขิตฺตํ.)}

\csromanheader{2}{เหตุปจฺจยา น-อารมฺมเณ ตีณิ. (สํขิตฺตํ.)}
\csromanheader{2}{นเหตุปจฺจยา อารมฺมเณ นว. (สํขิตฺตํ.)}
\prose{(ยถา กุสลตฺติเก ปญฺหาวารสฺส อนุโลมมฺปิ ปจฺจนียมฺปิ อนุโลมปจฺจนียมฺปิ ปจฺจนียานุโลมมฺปิ คณิตํ, เอวํ คเณตพฺพํ.)}

\csromanheader{2}{3. อปฺปมาณปท 1. ปฏิจฺจวาร}
\csromanheader{2}{\textbf{ปจฺจยจตุกฺก-เหตุ}}
\proseitem{371}{เหตุํ อปฺปมาณํ ธมฺมํ ปฏิจฺจ เหตุ อปฺปมาโณ ธมฺโม อุปฺปชฺชติ เหตุปจฺจยา.\csromanspacer{} ตีณิ.}

\gambhira{ธมฺมานุโลม ทุกติกปฏฺฐานปาฬิ}
\prose{\csromanfolio{280}นเหตุํ อปฺปมาณํ ธมฺมํ ปฏิจฺจ นเหตุ อปฺปมาโณ ธมฺโม อุปฺปชฺชติ เหตุปจฺจยา.\csromanspacer{} ตีณิ.}

\prose{เหตุํ อปฺปมาณญฺจ นเหตุํ อปฺปมาณญฺจ ธมฺมํ ปฏิจฺจ เหตุ อปฺปมาโณ ธมฺโม อุปฺปชฺชติ เหตุปจฺจยา.\csromanspacer{} ตีณิ. (สํขิตฺตํ.)}

\proseitem{372}{เหตุยา นว, อารมฺมเณ นว, อธิปติยา นว ฯเปฯ อุปนิสฺสเย นว ฯเปฯ กมฺเม นว, วิปาเก นว ฯเปฯ อวิคเต นว. (สํขิตฺตํ.)}

\csromanheader{2}{\textbf{น-อธิปติ}}
\proseitem{373}{เหตุํ อปฺปมาณํ ธมฺมํ ปฏิจฺจ เหตุ อปฺปมาโณ ธมฺโม อุปฺปชฺชติ น-อธิปติยา. (สํขิตฺตํ.)}

\vspace{\tipitakaparskip}
\csromancenter{\textbf{น-อธิปติ}ยา ฉ, นปุเรชาเต นว, นปจฺฉาชาเต นว, น-อาเสวเน นว, นกมฺเม ตีณิ, นวิปาเก นว, นวิปฺปยุตฺเต นว. (สํขิตฺตํ.)}
\csromanheader{2}{เหตุปจฺจยา น-อธิปติยา ฉ. (สํขิตฺตํ.)}
\csromanheader{2}{\textbf{น-อธิปติ}ปจฺจยา เหตุยา ฉ. (สํขิตฺตํ.)}
\prose{(สหชาตวาโรปิ ปจฺจยวาโรปิ นิสฺสยวาโรปิ สํสฏฺฐวาโรปิ สมฺปยุตฺตวาโรปิ ปฏิจฺจวารสทิสา วิตฺถาเรตพฺพา.)}

\csromanheader{2}{3. อปฺปมาณปท 7. ปญฺหาวาร}
\csromanheader{2}{\textbf{ปจฺจยจตุกฺก-เหตฺวาทิ}}
\proseitem{375}{เหตุ อปฺปมาโณ ธมฺโม เหตุสฺส อปฺปมาณสฺส ธมฺมสฺส เหตุปจฺจเยน ปจฺจโย.\csromanspacer{} ตีณิ.}

\prose{นเหตุ อปฺปมาโณ ธมฺโม นเหตุสฺส อปฺปมาณสฺส ธมฺมสฺส อารมฺมณปจฺจเยน ปจฺจโย.\csromanspacer{} ตีณิ.}

\proseitem{376}{เหตุ อปฺปมาโณ ธมฺโม เหตุสฺส อปฺปมาณสฺส ธมฺมสฺส อธิปติปจฺจเยน ปจฺจโย.\csromanspacer{} สหชาตาธิปติ.\csromanspacer{} ตีณิ.}

\csromanheader{2}{1. เหตุทุก 12. ปริตฺตตฺติก}
\prose{\csromanfolio{281}นเหตุ อปฺปมาโณ ธมฺโม นเหตุสฺส อปฺปมาณสฺส ธมฺมสฺส อธิปติปจฺจเยน ปจฺจโย.\csromanspacer{} อารมฺมณาธิปติ, สหชาตาธิปติ.\csromanspacer{} ตีณิ.}

\proseitem{377}{เหตุ อปฺปมาโณ ธมฺโม เหตุสฺส อปฺปมาณสฺส ธมฺมสฺส อนนฺตรปจฺจเยน ปจฺจโย.\csromanspacer{} นว ฯเปฯ.}

\prose{เหตุ อปฺปมาโณ ธมฺโม เหตุสฺส อปฺปมาณสฺส ธมฺมสฺส อุปนิสฺสยปจฺจเยน ปจฺจโย.\csromanspacer{} อนนฺตรูปนิสฺสโย, ปกตูปนิสฺสโย.\csromanspacer{} นว. (สํขิตฺตํ.)}

\proseitem{378}{เหตุยา ตีณิ, อารมฺมเณ ตีณิ, อธิปติยา ฉ, อนนฺตเร นว, สมนนฺตเร นว, สหชาเต นว, อญฺญมญฺเญ นว, นิสฺสเย นว, อุปนิสฺสเย นว, กมฺเม ตีณิ, วิปาเก นว, อาหาเร ตีณิ, อินฺทฺริเย นว, ฌาเน ตีณิ, มคฺเค นว, สมฺปยุตฺเต นว, อตฺถิยา นว, นตฺถิยา นว, วิคเต นว, อวิคเต นว. (สํขิตฺตํ.)}

\csromanheader{2}{\textbf{ปจฺจนียุทฺธาร}}
\proseitem{379}{เหตุ อปฺปมาโณ ธมฺโม เหตุสฺส อปฺปมาณสฺส ธมฺมสฺส สหชาตปจฺจเยน ปจฺจโย.\csromanspacer{} อุปนิสฺสยปจฺจเยน ปจฺจโย. (สํขิตฺตํ.)}

\proseitem{380}{นเหตุยา นว, น-อารมฺมเณ นว. (สํขิตฺตํ.)}

\csromanheader{2}{เหตุปจฺจยา น-อารมฺมเณ ตีณิ. (สํขิตฺตํ.)}
\csromanheader{2}{นเหตุปจฺจยา อารมฺมเณ ตีณิ. (สํขิตฺตํ.)}
\prose{(ยถา กุสลตฺติเก ปญฺหาวารสฺส อนุโลมมฺปิ ปจฺจนียมฺปิ อนุโลมปจฺจนียมฺปิ ปจฺจนียานุโลมมฺปิ คณิตํ, เอวํ คเณตพฺพํ.)}

\nitthitam{เหตุทุกปริตฺตตฺติกํ นิฏฺฐิตํ.}
\csromanmatikaanchor{mtk.44}
\csromantocmarkref{section}{13. ปริตฺตารมฺมณตฺติก}{mtk.44}
\bookmark[dest={mtk.44},level=1]{13. ปริตฺตารมฺมณตฺติก}
\csromanheader{1}{ธมฺมานุโลม ทุกติกปฏฺฐานปาฬิ \textbf{1. เหตุทุก 13. ปริตฺตารมฺมณตฺติก}}
\csromanheader{2}{1. ปริตฺตารมฺมณปท 1. ปฏิจฺจวาร}
\csromanheader{2}{\textbf{ปจฺจยจตุกฺก-เหตุ}}
\proseitem{381}{\csromanfolio{282}เหตุํ ปริตฺตารมฺมณํ ธมฺมํ ปฏิจฺจ เหตุ ปริตฺตารมฺมโณ ธมฺโม อุปฺปชฺชติ เหตุปจฺจยา.\csromanspacer{} ตีณิ.}

\prose{นเหตุํ ปริตฺตารมฺมณํ ธมฺมํ ปฏิจฺจ นเหตุ ปริตฺตารมฺมโณ ธมฺโม อุปฺปชฺชติ เหตุปจฺจยา.\csromanspacer{} ตีณิ.}

\prose{เหตุํ ปริตฺตารมฺมณญฺจ นเหตุํ ปริตฺตารมฺมณญฺจ ธมฺมํ ปฏิจฺจ เหตุ ปริตฺตารมฺมโณ ธมฺโม อุปฺปชฺชติ เหตุปจฺจยา.\csromanspacer{} ตีณิ. (สํขิตฺตํ.)}

\proseitem{382}{เหตุยา นว, อารมฺมเณ นว, อธิปติยา นว, อนนฺตเร นว, สมนนฺตเร นว, สหชาเต นว, อญฺญมญฺเญ นว, นิสฺสเย นว, อุปนิสฺสเย นว, ปุเรชาเต นว, อาเสวเน นว, กมฺเม นว, วิปาเก นว ฯเปฯ อวิคเต นว. (สํขิตฺตํ.)}

\csromanheader{2}{\textbf{นเหตุ น-อธิปติ}}
\proseitem{383}{นเหตุํ ปริตฺตารมฺมณํ ธมฺมํ ปฏิจฺจ นเหตุ ปริตฺตารมฺมโณ ธมฺโม อุปฺปชฺชติ นเหตุปจฺจยา.\csromanspacer{}\csromanspacer{}นเหตุํ ปริตฺตารมฺมณํ ธมฺมํ ปฏิจฺจ เหตุ ปริตฺตารมฺมโณ ธมฺโม อุปฺปชฺชติ นเหตุปจฺจยา. (2)}

\prose{เหตุํ ปริตฺตารมฺมณํ ธมฺมํ ปฏิจฺจ เหตุ ปริตฺตารมฺมโณ ธมฺโม อุปฺปชฺชติ น-อธิปติปจฺจยา. (สํขิตฺตํ.)}

\proseitem{384}{นเหตุยา เทฺว, น-อธิปติยา นว, นปุเรชาเต นว, นปจฺฉาชาเต นว, น-อาเสวเน นว, นกมฺเม ตีณิ, นวิปาเก นว, นฌาเน เอกํ, นมคฺเค เอกํ, นวิปฺปยุตฺเต นว. (สํขิตฺตํ.)}

\csromanheader{2}{เหตุปจฺจยา น-อธิปติยา นว. (สํขิตฺตํ.)}
\csromanheader{2}{นเหตุปจฺจยา อารมฺมเณ เทฺว. (สํขิตฺตํ.)}
\prose{(สหชาตวาโรปิ ปจฺจยวาโรปิ นิสฺสยวาโรปิ สํสฏฺฐวาโรปิ สมฺปยุตฺตวาโรปิ ปฏิจฺจวารสทิสา วิตฺถาเรตพฺพา.)}

\vspace{\tipitakaparskip}
\csromancenter{เหตุทุก 13.\csromanspacer{} ปริตฺตารมฺมณตฺติก \textbf{1.}\csromanspacer{}\textbf{ ปริตฺตารมฺมณปท 7.}\csromanspacer{}\textbf{ ปญฺหาวาร}}
\csromanheader{2}{\textbf{ปจฺจยจตุกฺก-เหตฺวาทิ}}
\proseitem{385}{\csromanfolio{283}เหตุ ปริตฺตารมฺมโณ ธมฺโม เหตุสฺส ปริตฺตารมฺมณสฺส ธมฺมสฺส เหตุปจฺจเยน ปจฺจโย.\csromanspacer{} ตีณิ.}

\prose{เหตุ ปริตฺตารมฺมโณ ธมฺโม เหตุสฺส ปริตฺตารมฺมณสฺส ธมฺมสฺส อารมฺมณปจฺจเยน ปจฺจโย.\csromanspacer{} นว.}

\prose{เหตุ ปริตฺตารมฺมโณ ธมฺโม เหตุสฺส ปริตฺตารมฺมณสฺส ธมฺมสฺส อธิปติปจฺจเยน ปจฺจโย. (สํขิตฺตํ.)}

\proseitem{368}{เหตุยา ตีณิ, อารมฺมเณ นว, อธิปติยา นว, อนนฺตเร นว, สมนนฺตเร นว, สหชาเต นว, อญฺญมญฺเญ นว, นิสฺสเย นว, อุปนิสฺสเย นว, อาเสวเน นว, กมฺเม ตีณิ, วิปาเก นว, อาหาเร ตีณิ, อินฺทฺริเย นว, ฌาเน ตีณิ, มคฺเค นว, สมฺปยุตฺเต นว, อตฺถิยา นว, นตฺถิยา นว, วิคเต นว, อวิคเต นว. (สํขิตฺตํ.)}

\csromanheader{2}{\textbf{ปจฺจนียุทฺธาร}}
\proseitem{387}{เหตุ ปริตฺตารมฺมโณ ธมฺโม เหตุสฺส ปริตฺตารมฺมณสฺส ธมฺมสฺส อารมฺมณปจฺจเยน ปจฺจโย.\csromanspacer{} สหชาตปจฺจเยน ปจฺจโย.\csromanspacer{} อุปนิสฺสยปจฺจเยน ปจฺจโย. (สํขิตฺตํ.)}

\proseitem{388}{นเหตุยา นว, น-อารมฺมเณ นว. (สํขิตฺตํ.)}

\csromanheader{2}{เหตุปจฺจยา น-อารมฺมเณ ตีณิ. (สํขิตฺตํ.)}
\csromanheader{2}{นเหตุปจฺจยา อารมฺมเณ นว. (สํขิตฺตํ.)}
\prose{(ยถา กุสลตฺติเก ปญฺหาวารสฺส อนุโลมมฺปิ ปจฺจนียมฺปิ อนุโลมปจฺจนียมฺปิ ปจฺจนียานุโลมมฺปิ คณิตํ, เอวํ คเณตพฺพํ.)}

\csromanheader{2}{ธมฺมานุโลม ทุกติกปฏฺฐานปาฬิ \textbf{2. มหคฺคตารมฺมณปท 1. ปฏิจฺจวาร}}
\csromanheader{2}{\textbf{ปจฺจยจตุกฺก-เหตุ}}
\proseitem{389}{\csromanfolio{284}เหตุํ มหคฺคตารมฺมณํ ธมฺมํ ปฏิจฺจ เหตุ มหคฺคตารมฺมโณ ธมฺโม อุปฺปชฺชติ เหตุปจฺจยา.\csromanspacer{} ตีณิ.}

\prose{นเหตุํ มหคฺคตารมฺมณํ ธมฺมํ ปฏิจฺจ นเหตุ มหคฺคตารมฺมโณ ธมฺโม อุปฺปชฺชติ เหตุปจฺจยา ตีณิ.}

\prose{เหตุํ มหคฺคตารมฺมณญฺจ นเหตุํ มหคฺคตารมฺมณญฺจ ธมฺมํ ปฏิจฺจ เหตุ มหคฺคตารมฺมโณ ธมฺโม อุปฺปชฺชติ เหตุปจฺจยา.\csromanspacer{} ตีณิ. (สํขิตฺตํ.)}

\proseitem{390}{เหตุยา นว, อารมฺมเณ นว, อธิปติยา นว ฯเปฯ กมฺเม นว, วิปาเก นว, อาหาเร นว ฯเปฯ อวิคเต นว. (สํขิตฺตํ.)}

\csromanheader{2}{\textbf{นเหตุ น-อธิปติ}}
\proseitem{391}{นเหตุํ มหคฺคตารมฺมณํ ธมฺมํ ปฏิจฺจ นเหตุ มหคฺคตารมฺมโณ ธมฺโม อุปฺปชฺชติ นเหตุปจฺจยา.\csromanspacer{}\csromanspacer{}นเหตุํ มหคฺคตารมฺมณํ ธมฺมํ ปฏิจฺจ เหตุ มหคฺคตารมฺมโณ ธมฺโม อุปฺปชฺชติ นเหตุปจฺจยา. (2)}

\prose{เหตุํ มหคฺคตารมฺมณํ ธมฺมํ ปฏิจฺจ เหตุ มหคฺคตารมฺมโณ ธมฺโม อุปฺปชฺชติ น-อธิปติปจฺจยา. (สํขิตฺตํ.)}

\proseitem{392}{นเหตุยา เทฺว, น-อธิปติยา นว, นปุเรชาเต นว, นปจฺฉาชาเต นว, น-อาเสวเน นว, นกมฺเม ตีณิ, นวิปาเก นว, นมคฺเค เอกํ, นวิปฺปยุตฺเต นว. (สํขิตฺตํ.)}

\csromanheader{2}{เหตุปจฺจยา น-อธิปติยา นว. (สํขิตฺตํ.)}
\csromanheader{2}{นเหตุปจฺจยา อารมฺมเณ เทฺว. (สํขิตฺตํ.)}
\prose{(สหชาตวาโรปิ ปจฺจยวาโรปิ นิสฺสยวาโรปิ สํสฏฺฐวาโรปิ สมฺปยุตฺตวาโรปิ ปฏิจฺจวารสทิสา วิตฺถาเรตพฺพา.)}

\vspace{\tipitakaparskip}
\csromancenter{เหตุทุก 13.\csromanspacer{} ปริตฺตารมฺมณตฺติก \textbf{2.}\csromanspacer{}\textbf{ มหคฺคตารมฺมณปท 7.}\csromanspacer{}\textbf{ ปญฺหาวาร}}
\csromanheader{2}{\textbf{ปจฺจยจตุกฺก-เหตฺวาทิ}}
\proseitem{393}{\csromanfolio{285}เหตุ มหคฺคตารมฺมโณ ธมฺโม เหตุสฺส มหคฺคตารมฺมณสฺส ธมฺมสฺส เหตุปจฺจเยน ปจฺจโย.\csromanspacer{} ตีณิ.}

\prose{เหตุ มหคฺคตารมฺมโณ ธมฺโม เหตุสฺส มหคฺคตารมฺมณสฺส ธมฺมสฺส อารมฺมณปจฺจเยน ปจฺจโย.\csromanspacer{} นว.}

\prose{เหตุ มหคฺคตารมฺมโณ ธมฺโม เหตุสฺส มหคฺคตารมฺมณสฺส ธมฺมสฺส อธิปติปจฺจเยน ปจฺจโย.\csromanspacer{} อารมฺมณาธิปติ, สหชาตาธิปติ.\csromanspacer{} ตีณิ.}

\prose{นเหตุ มหคฺคตารมฺมโณ ธมฺโม นเหตุสฺส มหคฺคตารมฺมณสฺส ธมฺมสฺส อธิปติปจฺจเยน ปจฺจโย.\csromanspacer{} อารมฺมณาธิปติ, สหชาตาธิปติ.\csromanspacer{} ตีณิ.}

\prose{เหตุ มหคฺคตารมฺมโณ จ นเหตุ มหคฺคตารมฺมโณ จ ธมฺมา เหตุสฺส มหคฺคตารมฺมณสฺส ธมฺมสฺส อธิปติปจฺจเยน ปจฺจโย.\csromanspacer{} อารมฺมณาธิปติ.\csromanspacer{} ตีณิ. (สํขิตฺตํ.)}

\proseitem{394}{เหตุยา ตีณิ, อารมฺมเณ นว, อธิปติยา นว ฯเปฯ นิสฺสเย นว, อุปนิสฺสเย นว, อาเสวเน นว, กมฺเม ตีณิ, วิปาเก นว, อาหาเร ตีณิ ฯเปฯ อวิคเต นว. (สํขิตฺตํ.)}

\csromanheader{2}{\textbf{ปจฺจนียุทฺธาร}}
\proseitem{395}{เหตุ มหคฺคตารมฺมโณ ธมฺโม เหตุสฺส มหคฺคตารมฺมณสฺส ธมฺมสฺส อารมฺมณปจฺจเยน ปจฺจโย.\csromanspacer{} สหชาตปจฺจเยน ปจฺจโย.\csromanspacer{} อุปนิสฺสยปจฺจเยน ปจฺจโย. (สํขิตฺตํ.)}

\proseitem{396}{นเหตุยา นว, น-อารมฺมเณ นว. (สํขิตฺตํ.)}

\csromanheader{2}{เหตุปจฺจยา น-อารมฺมเณ ตีณิ. (สํขิตฺตํ.)}
\csromanheader{2}{นเหตุปจฺจยา อารมฺมเณ นว. (สํขิตฺตํ.)}
\prose{(ยถา กุสลตฺติเก ปญฺหาวารสฺส อนุโลมมฺปิ ปจฺจนียมฺปิ อนุโลมปจฺจนียมฺปิ ปจฺจนียานุโลมมฺปิ คณิตํ, เอวํ คเณตพฺพํ.)}

\csromanheader{2}{ธมฺมานุโลม ทุกติกปฏฺฐานปาฬิ \textbf{3. อปฺปามาณารมฺมณปท 1. ปฏิจฺจวาร}}
\csromanheader{2}{\textbf{ปจฺจยจตุกฺก-เหตุ}}
\proseitem{397}{\csromanfolio{286}เหตุํ อปฺปมาณารมฺมณํ ธมฺมํ ปฏิจฺจ เหตุ อปฺปมาณารมฺมโณ ธมฺโม อุปฺปชฺชติ เหตุปจฺจยา.\csromanspacer{} ตีณิ.}

\prose{\textbf{นเหตุ}ํ อปฺปมาณารมฺมณํ ธมฺมํ ปฏิจฺจ นเหตุ อปฺปมาณารมฺมโณ ธมฺโม อุปฺปชฺชติ เหตุปจฺจยา.\csromanspacer{} ตีณิ.}

\prose{เหตุํ อปฺปมาณารมฺมณญฺจ นเหตุํ อปฺปมาณารมฺมณญฺจ ธมฺมํ ปฏิจฺจ เหตุ อปฺปมาณารมฺมโณ ธมฺโม อุปฺปชฺชติ เหตุปจฺจยา.\csromanspacer{} ตีณิ. (สํขิตฺตํ.)}

\proseitem{398}{เหตุยา นว, อารมฺมเณ นว, อธิปติยา นว, อนนฺตเร นว, สมนนฺตเร นว, สหชาเต นว, อญฺญมญฺเญ นว, นิสฺสเย นว, อุปนิสฺสเย นว, ปุเรชาเต นว, อาเสวเน นว, กมฺเม นว, วิปาเก นว ฯเปฯ อวิคเต นว. (สํขิตฺตํ.)}

\csromanheader{2}{\textbf{นเหตุ}}
\proseitem{399}{\textbf{นเหตุ}ํ อปฺปมาณารมฺมณํ ธมฺมํ ปฏิจฺจ นเหตุ อปฺปมาณารมฺมโณ ธมฺโม อุปฺปชฺชติ นเหตุปจฺจยา. (สํขิตฺตํ.)}

\vspace{\tipitakaparskip}
\csromancenter{\textbf{นเหตุ}ยา เอกํ, น-อธิปติยา นว, นปุเรชาเต นว, นปจฺฉาชาเต นว, น-อาเสวเน นว, นกมฺเม ตีณิ, นวิปาเก นว, นมคฺเค เอกํ, นวิปฺปยุตฺเต นว. (สํขิตฺตํ.)}
\csromanheader{2}{เหตุปจฺจยา น-อธิปติยา นว. (สํขิตฺตํ.)}
\csromanheader{2}{\textbf{นเหตุ}ปจฺจยา อารมฺมเณ เอกํ. (สํขิตฺตํ.)}
\prose{(สหชาตวาโรปิ ปจฺจยวาโรปิ นิสฺสยวาโรปิ สํสฏฺฐวาโรปิ สมฺปยุตฺตวาโรปิ ปฏิจฺจวารสทิสา วิตฺถาเรตพฺพา.)}

\vspace{\tipitakaparskip}
\csromancenter{เหตุทุก 13.\csromanspacer{} ปริตฺตารมฺมณตฺติก \textbf{3.}\csromanspacer{}\textbf{ อปฺปมาณารมฺมณปท 7.}\csromanspacer{}\textbf{ ปญฺหาวาร}}
\csromanheader{2}{\textbf{ปจฺจยจตุกฺก-เหตฺวาทิ}}
\proseitem{401}{\csromanfolio{287}เหตุ อปฺปมาณารมฺมโณ ธมฺโม เหตุสฺส อปฺปมาณารมฺมณสฺส ธมฺมสฺส เหตุปจฺจเยน ปจฺจโย.\csromanspacer{} ตีณิ.}

\prose{เหตุ อปฺปมาณารมฺมโณ ธมฺโม เหตุสฺส อปฺปมาณารมฺมณสฺส ธมฺมสฺส อารมฺมณปจฺจเยน ปจฺจโย.\csromanspacer{} นว.}

\prose{เหตุ อปฺปมาณารมฺมโณ ธมฺโม เหตุสฺส อปฺปมาณารมฺมณสฺส ธมฺมสฺส อธิปติปจฺจเยน ปจฺจโย.\csromanspacer{} อารมฺมณาธิปติ, สหชาตาธิปติ.\csromanspacer{} ตีณิ.}

\prose{นเหตุ อปฺปมาณารมฺมโณ ธมฺโม นเหตุสฺส อปฺปมาณารมฺมณสฺส ธมฺมสฺส อธิปติปจฺจเยน ปจฺจโย.\csromanspacer{} อารมฺมณาธิปติ, สหชาตาธิปติ.\csromanspacer{} ตีณิ.}

\prose{เหตุ อปฺปมาณารมฺมโณ จ นเหตุ อปฺปมาณารมฺมโณ จ ธมฺมา เหตุสฺส อปฺปมาณารมฺมณสฺส ธมฺมสฺส อธิปติปจฺจเยน ปจฺจโย.\csromanspacer{} อารมฺมณาธิปติ.\csromanspacer{} ตีณิ. (สํขิตฺตํ.)}

\proseitem{402}{เหตุยา ตีณิ, อารมฺมเณ นว, อธิปติยา นว, อนนฺตเร นว, สมนนฺตเร นว, สหชาเต นว, อญฺญมญฺเญ นว, นิสฺสเย นว, อุปนิสฺสเย นว, อาเสวเน นว, กมฺเม ตีณิ, วิปาเก นว, อาหาเร ตีณิ, อินฺทฺริเร นว, ฌาเน ตีณิ, มคฺเค นว, สมฺปยุตฺเต นว, อตฺถิยา นว ฯเปฯ อวิคเต นว. (สํขิตฺตํ.)}

\csromanheader{2}{\textbf{ปจฺจนียุทฺธาร}}
\proseitem{403}{เหตุ อปฺปมาณารมฺมโณ ธมฺโม เหตุสฺส อปฺปมาณารมฺมณสฺส ธมฺมสฺส อารมฺมณปจฺจเยน ปจฺจโย.\csromanspacer{} สหชาตปจฺจเยน ปจฺจโย.\csromanspacer{} อุปนิสฺสยปจฺจเยน ปจฺจโย. (สํขิตฺตํ.)}

\proseitem{404}{นเหตุยา นว, น-อารมฺมเณ นว. (สํขิตฺตํ.)}

\csromanheader{2}{เหตุปจฺจยา น-อารมฺมเณ ตีณิ. (สํขิตฺตํ.)}
\gambhira{ธมฺมานุโลม ทุกติกปฏฺฐานปาฬิ}
\csromanheader{2}{นเหตุปจฺจยา อารมฺมเณ นว. (สํขิตฺตํ.)}
\prose{\csromanfolio{288}(ยถา กุสลตฺติเก ปญฺหาวารสฺส อนุโลมมฺปิ ปจฺจนียมฺปิ อนุโลมปจฺจนียมฺปิ ปจฺจนียานุโลมมฺปิ คณิตํ, เอวํ คเณตพฺพํ.)}

\nitthitam{เหตุทุกปริตฺตารมฺมณตฺติกํ นิฏฺฐิตํ.}
\csromanmatikaanchor{mtk.45}
\csromanmatikaanchor{mtk.118}
\csromantocmarkref{section}{14. หินตฺติก}{mtk.45}
\bookmark[dest={mtk.45},level=1]{14. หินตฺติก}
\csromanheader{6}{1. เหตุทุก 14. หีนตฺติก}
\csromanheader{2}{1. หีนปท 1. ปฏิจฺจวาร}
\csromanheader{2}{\textbf{ปจฺจยจตุกฺก-เหตุ}}
\proseitem{405}{เหตุํ หีนํ ธมฺมํ ปฏิจฺจ เหตุ หีโน ธมฺโม อุปฺปชฺชติ เหตุปจฺจยา.\csromanspacer{} ตีณิ.}

\prose{นเหตุํ หีนํ ธมฺมํ ปฏิจฺจ นเหตุ หีโน ธมฺโม อุปฺปชฺชติ เหตุปจฺจยา.\csromanspacer{} ตีณิ.}

\prose{เหตุํ หีนญฺจ นเหตุํ หีนญฺจ ธมฺมํ ปฏิจฺจ เหตุ หีโน ธมฺโม อุปฺปชฺชติ เหตุปจฺจยา.\csromanspacer{} ตีณิ. (สํขิตฺตํ.)}

\proseitem{406}{เหตุยา นว, อารมฺมเณ นว, อธิปติยา นว, อนนฺตเร นว, สมนนฺตเร นว, สหชาเต นว, อญฺญมญฺเญ นว, นิสฺสเย นว, อุปนิสฺสเย นว, ปุเรชาเต นว, อาเสวเน นว, กมฺเม นว, อาหาเร นว ฯเปฯ อวิคเต นว. (สํขิตฺตํ.)}

\csromanheader{2}{\textbf{นเหตุ น-อธิปติ}}
\proseitem{407}{นเหตุํ หีนํ ธมฺมํ ปฏิจฺจ เหตุ หีโน ธมฺโม อุปฺปชฺชติ นเหตุปจฺจยา. (1)}

\prose{เหตุํ หีนํ ธมฺมํ ปฏิจฺจ เหตุ หีโน ธมฺโม อุปฺปชฺชติ น-อธิปติปจฺจยา.\csromanspacer{} ตีณิ.}

\prose{นเหตุํ หีนํ ธมฺมํ ปฏิจฺจ นเหตุ หีโน ธมฺโม อุปฺปชฺชติ น-อธิปติปจฺจยา.\csromanspacer{} ตีณิ.}

\csromanheader{2}{1. เหตุทุก 14. หีนตฺติก}
\prose{\csromanfolio{289}เหตุํ หีนญฺจ นเหตุํ หีนญฺจ ธมฺมํ ปฏิจฺจ เหตุ หีโน ธมฺโม อุปฺปชฺชติ น-อธิปติปจฺจยา.\csromanspacer{} ตีณิ. (สํขิตฺตํ.)}

\proseitem{408}{นเหตุยา เอกํ, น-อธิปติยา นว, นปุเรชาเต นว, นปจฺฉาชาเต นว, น-อาเสวเน นว, นกมฺเม ตีณิ, นวิปาเก นว, นวิปฺปยุตฺเต นว. (สํขิตฺตํ.)}

\csromanheader{2}{เหตุปจฺจยา น-อธิปติยา นว. (สํขิตฺตํ.)}
\csromanheader{2}{นเหตุปจฺจยา อารมฺมเณ เอกํ. (สํขิตฺตํ.)}
\prose{(สหชาตวาโรปิ ปจฺจยวาโรปิ นิสฺสยวาโรปิ สํสฏฺฐวาโรปิ สมฺปยุตฺตวาโรปิ ปฏิจฺจวารสทิสา วิตฺถาเรตพฺพา.)}

\csromanheader{2}{1. หีนปท 7. ปญฺหาวาร}
\csromanheader{2}{\textbf{ปจฺจยจตุกฺก-เหตฺวาทิ}}
\proseitem{409}{เหตุ หีโน ธมฺโม เหตุสฺส หีนสฺส ธมฺมสฺส เหตุปจฺจเยน ปจฺจโย.\csromanspacer{} ตีณิ.}

\prose{เหตุ หีโน ธมฺโม เหตุสฺส หีนสฺส ธมฺมสฺส อารมฺมณปจฺจเยน ปจฺจโย ฯเปฯ อธิปติปจฺจเยน ปจฺจโย.\csromanspacer{} อารมฺมณาธิปติ, สหชาตาธิปติ. (สํขิตฺตํ.)}

\proseitem{410}{เหตุยา ตีณิ, อารมฺมเณ นว, อธิปติยา นว, อนนฺตเร นว, สมนนฺตเร นว, สหชาเต นว, อญฺญมญฺเญ นว, นิสฺสเย นว, อุปนิสฺสเย นว, อาเสวเน นว, กมฺเม ตีณิ, อาหาเร ตีณิ, อินฺทฺริเย ฌาเน มคฺเค ตีณิ, สมฺปยุตฺเต นว ฯเปฯ อวิคเต นว. (สํขิตฺตํ.)}

\csromanheader{2}{\textbf{ปจฺจนียุทฺธาร}}
\proseitem{411}{เหตุ หีโน ธมฺโม เหตุสฺส หีนสฺส ธมฺมสฺส อารมฺมณปจฺจเยน ปจฺจโย.\csromanspacer{} สหชาตปจฺจเยน ปจฺจโย.\csromanspacer{} อุปนิสฺสยปจฺจเยน ปจฺจโย. (สํขิตฺตํ.)}

\gambhira{ธมฺมานุโลม ทุกติกปฏฺฐานปาฬิ}
\proseitem{412}{\csromanfolio{290}\textbf{นเหตุ}ยา นว, น-อารมฺมเณ นว. (สํขิตฺตํ.)}

\csromanheader{2}{เหตุปจฺจยา น-อารมฺมเณ ตีณิ. (สํขิตฺตํ.)}
\csromanheader{2}{\textbf{นเหตุ}ปจฺจยา อารมฺมเณ นว. (สํขิตฺตํ.)}
\prose{(ยถา กุสลตฺติเก ปญฺหาวารสฺส อนุโลมมฺปิ ปจฺจนียมฺปิ อนุโลมปจฺจนียมฺปิ ปจฺจนียานุโลมมฺปิ คณิตํ, เอวํ คเณตพฺพํ.)}

\csromanheader{2}{2. มชฺฌิมปท 1. ปฏิจฺจวาร}
\csromanheader{2}{\textbf{ปจฺจยจตุกฺก-เหตุ}}
\proseitem{413}{เหตุํ มชฺฌิมํ ธมฺมํ ปฏิจฺจ เหตุ มชฺฌิโม ธมฺโม อุปฺปชฺชติ เหตุปจฺจยา. (สํขิตฺตํ.)}

\proseitem{414}{เหตุยา นว, อารมฺมเณ นว ฯเปฯ อวิคเต นว. (สํขิตฺตํ.)}

\csromanheader{2}{\textbf{นเหตุ}}
\vspace{\tipitakaparskip}
\csromancenter{\textbf{นเหตุ}ํ มชฺฌิมํ ธมฺมํ ปฏิจฺจ นเหตุ มชฺฌิโม ธมฺโม อุปฺปชฺชติ นเหตุปจฺจยา. (สํขิตฺตํ.)}
\csromancenter{\textbf{นเหตุ}ยา เอกํ, น-อารมฺมเณ ตีณิ, น-อธิปติยา นว ฯเปฯ โนวิคเต ตีณิ. (สํขิตฺตํ.)}
\csromanheader{2}{เหตุปจฺจยา น-อารมฺมเณ ตีณิ. (สํขิตฺตํ.)}
\csromanheader{2}{\textbf{นเหตุ}ปจฺจยา อารมฺมเณ เอกํ. (สํขิตฺตํ.)}
\prose{(สหชาตวาโรปิ ปจฺจยวาโรปิ นิสฺสยวาโรปิ สํสฏฺฐวาโรปิ สมฺปยุตฺตวาโรปิ ปฏิจฺจวารสทิสา วิตฺถาเรตพฺพา.)}

\csromanheader{2}{1. เหตุทุก 14. หีนตฺติก \textbf{2. มชฺฌิมปท 7. ปญฺหาวาร}}
\csromanheader{2}{\textbf{ปจฺจยจตุกฺก-เหตฺวาทิ}}
\proseitem{417}{\csromanfolio{291}เหตุ มชฺฌิโม ธมฺโม เหตุสฺส มชฺฌิมสฺส ธมฺมสฺส เหตุปจฺจเยน ปจฺจโย.\csromanspacer{} ตีณิ.}

\prose{เหตุ มชฺฌิโม ธมฺโม เหตุสฺส มชฺฌิมสฺส ธมฺมสฺส อารมฺมณปจฺจเยน ปจฺจโย.\csromanspacer{} อธิปติปจฺจเยน ปจฺจโย.\csromanspacer{} อารมฺมณาธิปติ, สหชาตาธิปติ. (สํขิตฺตํ.)}

\proseitem{418}{เหตุยา ตีณิ, อารมฺมเณ นว, อธิปติยา นว, อนนฺตเร นว, สมนนฺตเร นว ฯเปฯ กมฺเม ตีณิ, วิปาเก นว, อาหาเร ตีณิ ฯเปฯ อวิคเต นว. (สํขิตฺตํ.)}

\csromanheader{2}{\textbf{ปจฺจนียุทฺธาร}}
\proseitem{419}{เหตุ มชฺฌิโม ธมฺโม เหตุสฺส มชฺฌิมสฺส ธมฺมสฺส อารมฺมณปจฺจเยน ปจฺจโย.\csromanspacer{} สหชาตปจฺจเยน ปจฺจโย.\csromanspacer{} อุปนิสฺสยปจฺจเยน ปจฺจโย. (สํขิตฺตํ.)}

\proseitem{420}{นเหตุยา นว, น-อารมฺมเณ นว. (สํขิตฺตํ.)}

\csromanheader{2}{เหตุปจฺจยา น-อารมฺมเณ ตีณิ. (สํขิตฺตํ.)}
\csromanheader{2}{นเหตุปจฺจยา อารมฺมเณ นว. (สํขิตฺตํ.)}
\prose{(ยถา กุสลตฺติเก ปญฺหาวารสฺส อนุโลมมฺปิ ปจฺจนียมฺปิ อนุโลมปจฺจนียมฺปิ ปจฺจนียานุโลมมฺปิ คณิตํ, เอวํ คเณตพฺพํ.)}

\csromanheader{2}{3. ปณีตปท 1. ปฏิจฺจวาร}
\csromanheader{2}{\textbf{ปจฺจยจตุกฺก-เหตุ}}
\proseitem{421}{เหตุํ ปณีตํ ธมฺมํ ปฏิจฺจ เหตุ ปณีโต ธมฺโม อุปฺปชฺชติ เหตุปจฺจยา.\csromanspacer{} ตีณิ.}

\gambhira{ธมฺมานุโลม ทุกติกปฏฺฐานปาฬิ}
\prose{\csromanfolio{292}นเหตุํ ปณีตํ ธมฺมํ ปฏิจฺจ นเหตุ ปณีโต ธมฺโม อุปฺปชฺชติ เหตุปจฺจยา.\csromanspacer{} ตีณิ.}

\prose{เหตุํ ปณีตญฺจ นเหตุํ ปณีตญฺจ ธมฺมํ ปฏิจฺจ เหตุ ปณีโต ธมฺโม อุปฺปชฺชติ เหตุปจฺจยา.\csromanspacer{} ตีณิ. (สํขิตฺตํ.)}

\proseitem{422}{เหตุยา นว, อารมฺมเณ นว, อธิปติยา นว ฯเปฯ กมฺเม นว, วิปาเก นว, อาหาเร นว ฯเปฯ อวิคเต นว. (สํขิตฺตํ.)}

\csromanheader{2}{\textbf{น-อธิปติ}}
\proseitem{423}{เหตุํ ปณีตํ ธมฺมํ ปฏิจฺจ เหตุ ปณีโต ธมฺโม อุปฺปชฺชติ น-อธิปติปจฺจยา. (สํขิตฺตํ.)}

\vspace{\tipitakaparskip}
\csromancenter{\textbf{น-อธิปติ}ยา ฉ, นปุเรชาเต นว, นปจฺฉาชาเต นว, น-อาเสวเน นว, นกมฺเม ตีณิ, นวิปาเก นว, นวิปฺปยุตฺเต นว. (สํขิตฺตํ.)}
\csromanheader{2}{เหตุปจฺจยา น-อธิปติยา ฉ. (สํขิตฺตํ.)}
\csromanheader{2}{\textbf{น-อธิปติ}ปจฺจยา เหตุยา ฉ. (สํขิตฺตํ.)}
\prose{(สหชาตวาโรปิ ปจฺจยวาโรปิ นิสฺสยวาโรปิ สํสฏฺฐวาโรปิ สมฺปยุตฺตวาโรปิ ปฏิจฺจวารสทิสา วิตฺถาเรตพฺพา.)}

\csromanheader{2}{3. ปณีตปท 7. ปญฺหาวาร}
\csromanheader{2}{\textbf{ปจฺจยจตุกฺก-เหตฺวาทิ}}
\proseitem{425}{เหตุ ปณีโต ธมฺโม เหตุสฺส ปณีตสฺส ธมฺมสฺส เหตุปจฺจเยน ปจฺจโย.\csromanspacer{} ตีณิ.}

\prose{นเหตุ ปณีโต ธมฺโม นเหตุสฺส ปณีตสฺส ธมฺมสฺส อารมฺมณปจฺจเยน ปจฺจโย.\csromanspacer{} ตีณิ.}

\prose{เหตุ ปณีโต ธมฺโม เหตุสฺส ปณีตสฺส ธมฺมสฺส อธิปติปจฺจเยน ปจฺจโย.\csromanspacer{} อารมฺมณาธิปติ, สหชาตาธิปติ.\csromanspacer{} ตีณิ. (สํขิตฺตํ.)}

\csromanmatikaanchor{mtk.46}
\csromantocmarkref{section}{15. มิจฺฉตฺตนิยตตฺติก}{mtk.46}
\bookmark[dest={mtk.46},level=1]{15. มิจฺฉตฺตนิยตตฺติก}
\csromanheader{1}{1. เหตุทุก 15. มิจฺฉตฺตนิยตตฺติก}
\proseitem{426}{\csromanfolio{293}เหตุยา ตีณิ, อารมฺมเณ ตีณิ, อธิปติยา ฉ, อนนฺตเร นว, สมนนฺตเร นว, สหชาเต นว, อญฺญมญฺเญ นว, นิสฺสเย นว, อุปนิสฺสเย นว, กมฺเม ตีณิ, วิปาเก นว, อาหาเร ตีณิ, อินฺทฺริเย นว, ฌาเน ตีณิ, มคฺเค นว, สมฺปยุตฺเต นว, อตฺถิยา นว ฯเปฯ อวิคเต นว. (สํขิตฺตํ.)}

\csromanheader{2}{\textbf{ปจฺจนียุทฺธาร}}
\proseitem{427}{เหตุ ปณีโต ธมฺโม เหตุสฺส ปณีตสฺส ธมฺมสฺส สหชาตปจฺจเยน ปจฺจโย.\csromanspacer{} อุปนิสฺสยปจฺจเยน ปจฺจโย. (สํขิตฺตํ.)}

\proseitem{428}{นเหตุยา นว, น-อารมฺมเณ นว. (สํขิตฺตํ.)}

\csromanheader{2}{เหตุปจฺจยา น-อารมฺมเณ ตีณิ. (สํขิตฺตํ.)}
\csromanheader{2}{นเหตุปจฺจยา อารมฺมเณ ตีณิ. (สํขิตฺตํ.)}
\prose{(ยถา กุสลตฺติเก ปญฺหาวารสฺส อนุโลมมฺปิ ปจฺจนียมฺปิ อนุโลมปจฺจนียมฺปิ ปจฺจนียานุโลมมฺปิ คณิตํ, เอวํ คเณตพฺพํ.)}

\nitthitam{เหตุทุกหีนตฺติกํ นิฏฺฐิตํ.}
\csromanheader{2}{1. เหตุทุก 15. มิจฺฉตฺตนิยตตฺติก}
\csromanheader{2}{1. มิจฺฉตฺตนิยตปท 1. ปฏิจฺจวาร}
\csromanheader{2}{\textbf{ปจฺจยจตุกฺก-เหตุ อารมฺมณ}}
\proseitem{429}{เหตุํ มิจฺฉตฺตนิยตํ ธมฺมํ ปฏิจฺจ เหตุ มิจฺฉตฺตนิยโต ธมฺโม อุปฺปชฺชติ เหตุปจฺจยา.\csromanspacer{} ตีณิ.}

\prose{นเหตุํ มิจฺฉตฺตนิยตํ ธมฺมํ ปฏิจฺจ นเหตุ มิจฺฉตฺตนิยโต ธมฺโม อุปฺปชฺชติ เหตุปจฺจยา.\csromanspacer{} ตีณิ.}

\prose{เหตุํ มิจฺฉตฺตนิยตญฺจ นเหตุํ มิจฺฉตฺตนิยตญฺจ ธมฺมํ ปฏิจฺจ เหตุ มิจฺฉตฺตนิยโต ธมฺโม อุปฺปชฺชติ เหตุปจฺจยา.\csromanspacer{} ตีณิ.}

\gambhira{ธมฺมานุโลม ทุกติกปฏฺฐานปาฬิ}
\prose{\csromanfolio{294}เหตุํ มิจฺฉตฺตนิยตํ ธมฺมํ ปฏิจฺจ เหตุ มิจฺฉตฺตนิยโต ธมฺโม อุปฺปชฺชติ อารมฺมณปจฺจยา. (สํขิตฺตํ.)}

\proseitem{430}{เหตุยา นว, อารมฺมเณ นว, อธิปติยา นว, อนนฺตเร นว, สมนนฺตเร นว, สหชาเต นว, อญฺญมญฺเญ นว, นิสฺสเย นว, อุปนิสฺสเย นว, ปุเรชาเต นว, อาเสวเน นว, กมฺเม นว, อาหาเร นว ฯเปฯ อวิคเต นว. (สํขิตฺตํ.)}

\csromanheader{2}{\textbf{น-อธิปติ}}
\proseitem{431}{เหตุํ มิจฺฉตฺตนิยตํ ธมฺมํ ปฏิจฺจ นเหตุ มิจฺฉตฺตนิยโต ธมฺโม อุปฺปชฺชติ น-อธิปติปจฺจยา. (สํขิตฺตํ.)}

\vspace{\tipitakaparskip}
\csromancenter{\textbf{น-อธิปติ}ยา ตีณิ, นปจฺฉาชาเต นว, นกมฺเม ตีณิ, นวิปาเก นว. (สํขิตฺตํ.)}
\csromanheader{2}{เหตุปจฺจยา น-อธิปติยา ตีณิ. (สํขิตฺตํ.)}
\csromanheader{2}{\textbf{น-อธิปติ}ปจฺจยา เหตุยา ตีณิ. (สํขิตฺตํ.)}
\prose{(สหชาตวาโรปิ ปจฺจยวาโรปิ นิสฺสยวาโรปิ สํสฏฺฐวาโรปิ สมฺปยุตฺตวาโรปิ ปฏิจฺจวารสทิสา วิตฺถาเรตพฺพา.)}

\csromanheader{2}{1. มิจฺฉตฺตนิยตปท 7. ปญฺหาวาร}
\csromanheader{2}{\textbf{ปจฺจยจตุกฺก-เหตุ อธิปติ}}
\proseitem{433}{เหตุ มิจฺฉตฺตนิยโต ธมฺโม เหตุสฺส มิจฺฉตฺตนิยตสฺส ธมฺมสฺส เหตุปจฺจเยน ปจฺจโย.\csromanspacer{} ตีณิ.}

\prose{นเหตุ มิจฺฉตฺตนิยโต ธมฺโม นเหตุสฺส มิจฺฉตฺตนิยตสฺส ธมฺมสฺส อธิปติปจฺจเยน ปจฺจโย.\csromanspacer{} สหชาตาธิปติ.\csromanspacer{} ตีณิ. (สํขิตฺตํ.)}

\proseitem{434}{เหตุยา ตีณิ, อธิปติยา ตีณิ, สหชาเต นว, อญฺญมญฺเญ นว, นิสฺสเย นว, อุปนิสฺสเย นว, กมฺเม ตีณิ, อาหาเร ตีณิ, อินฺทฺริเย ตีณิ, ฌาเน ตีณิ, มคฺเค ตีณิ, สมฺปยุตฺเต นว, อตฺถิยา นว, อวิคเต นว. (สํขิตฺตํ.)}

\csromanheader{2}{1. เหตุทุก 15. มิจฺฉตฺตนิยตตฺติก \textbf{ปจฺจนียุทฺธาร}}
\proseitem{435}{\csromanfolio{295}เหตุ มิจฺฉตฺตนิยโต ธมฺโม เหตุสฺส มิจฺฉตฺตนิยตสฺส ธมฺมสฺส สหชาตปจฺจเยน ปจฺจโย.\csromanspacer{} อุปนิสฺสยปจฺจเยน ปจฺจโย. (สํขิตฺตํ.)}

\proseitem{436}{นเหตุยา นว, น-อารมฺมเณ นว. (สํขิตฺตํ.)}

\csromanheader{2}{เหตุปจฺจยา น-อธิปติยา ตีณิ. (สํขิตฺตํ.)}
\csromanheader{2}{นเหตุยา อธิปติยา ตีณิ. (สํขิตฺตํ.)}
\prose{(ยถา กุสลตฺติเก ปญฺหาวารสฺส อนุโลมมฺปิ ปจฺจนียมฺปิ อนุโลมปจฺจนียมฺปิ ปจฺจนียานุโลมมฺปิ คณิตํ, เอวํ คเณตพฺพํ.)}

\csromanheader{2}{2. สมฺมตฺตนิยตปท 1. ปฏิจฺจวาร}
\csromanheader{2}{\textbf{ปจฺจยจตุกฺก-เหตุ}}
\proseitem{437}{เหตุํ สมฺมตฺตนิยตํ ธมฺมํ ปฏิจฺจ เหตุ สมฺมตฺตนิยโต ธมฺโม อุปฺปชฺชติ เหตุปจฺจยา.\csromanspacer{} ตีณิ.}

\prose{นเหตุํ สมฺมตฺตนิยตํ ธมฺมํ ปฏิจฺจ นเหตุ สมฺมตฺตนิยโต ธมฺโม อุปฺปชฺชติ เหตุปจฺจยา.\csromanspacer{} ตีณิ.}

\prose{เหตุํ สมฺมตฺตนิยตญฺจ นเหตุํ สมฺมตฺตนิยตญฺจ ธมฺมํ ปฏิจฺจ เหตุ สมฺมตฺตนิยโต ธมฺโม อุปฺปชฺชติ เหตุปจฺจยา.\csromanspacer{} ตีณิ. (สํขิตฺตํ.)}

\proseitem{438}{เหตุยา นว, อารมฺมเณ นว, อธิปติยา นว ฯเปฯ ปุเรชาเต นว, อาเสวเน นว, กมฺเม นว, อาหาเร นว ฯเปฯ อวิคเต นว. (สํขิตฺตํ.)}

\csromanheader{2}{\textbf{น-อธิปติ}}
\proseitem{439}{เหตุํ สมฺมตฺตนิยตํ ธมฺมํ ปฏิจฺจ เหตุ สมฺมตฺตนิยโต ธมฺโม อุปฺปชฺชติ น-อธิปติปจฺจยา. (สํขิตฺตํ.)}

\vspace{\tipitakaparskip}
\csromancenter{\textbf{น-อธิปติ}ยา ฉ, นปุเรชาเต นว, นปจฺฉาชาเต นว, นกมฺเม ตีณิ, นวิปาเก นว, นวิปฺปยุตฺเต นว. (สํขิตฺตํ.)}
\gambhira{ธมฺมานุโลม ทุกติกปฏฺฐานปาฬิ}
\csromanheader{2}{เหตุปจฺจยา น-อธิปติยา ฉ. (สํขิตฺตํ.)}
\csromanheader{2}{น-อธิปติปจฺจยา เหตุยา ฉ. (สํขิตฺตํ.)}
\prose{\csromanfolio{296}(สหชาตวาโรปิ ปจฺจยวาโรปิ นิสฺสยวาโรปิ สํสฏฺฐวาโรปิ สมฺปยุตฺตวาโรปิ ปฏิจฺจวารสทิสา วิตฺถาเรตพฺพา.)}

\csromanheader{2}{2. สมฺมตฺตนิยตปท 7. ปญฺหาวาร}
\csromanheader{2}{\textbf{ปจฺจยจตุกฺก-เหตุ อธิปติ}}
\proseitem{441}{เหตุ สมฺมตฺตนิยโต ธมฺโม เหตุสฺส สมฺมตฺตนิยตสฺส ธมฺมสฺส เหตุปจฺจเยน ปจฺจโย.\csromanspacer{} ตีณิ.}

\prose{เหตุ สมฺมตฺตนิยโต ธมฺโม เหตุสฺส สมฺมตฺตนิยตสฺส ธมฺมสฺส อธิปติปจฺจเยน ปจฺจโย.\csromanspacer{} สหชาตาธิปติ.\csromanspacer{} ตีณิ.}

\prose{นเหตุ สมฺมตฺตนิยโต ธมฺโม นเหตุสฺส สมฺมตฺตนิยตสฺส ธมฺมสฺส อธิปติปจฺจเยน ปจฺจโย.\csromanspacer{} สหชาตาธิปติ.\csromanspacer{} ตีณิ. (สํขิตฺตํ.)}

\proseitem{442}{เหตุยา ตีณิ, อธิปติยา ฉ, สหชาเต นว, อญฺญมญฺเญ นว, นิสฺสเย นว, อุปนิสฺสเย นว, กมฺเม ตีณิ, อาหาเร ตีณิ, อินฺทฺริเย นว, ฌาเน ตีณิ, มคฺเค นว, สมฺปยุตฺเต นว, อตฺถิยา นว, อวิคเต นว. (สํขิตฺตํ.)}

\csromanheader{2}{\textbf{ปจฺจนียุทฺธาร}}
\proseitem{443}{เหตุ สมฺมตฺตนิยโต ธมฺโม เหตุสฺส สมฺมตฺตนิยตสฺส ธมฺมสฺส สหชาตปจฺจเยน ปจฺจโย.\csromanspacer{} อุปนิสฺสยปจฺจเยน ปจฺจโย. (สํขิตฺตํ.)}

\proseitem{444}{นเหตุยา นว, น-อารมฺมเณ นว. (สํขิตฺตํ.)}

\csromanheader{2}{เหตุปจฺจยา น-อารมฺมเณ ตีณิ. (สํขิตฺตํ.)}
\csromanheader{2}{นเหตุปจฺจยา อธิปติยา ตีณิ. (สํขิตฺตํ.)}
\prose{(ยถา กุสลตฺติเก ปญฺหาวารสฺส อนุโลมมฺปิ ปจฺจนียมฺปิ อนุโลมปจฺจนียมฺปิ ปจฺจนียานุโลมมฺปิ คณิตํ, เอวํ คเณตพฺพํ.)}

\vspace{\tipitakaparskip}
\csromancenter{เหตุทุก 15.\csromanspacer{} มิจฺฉตฺตนิยตตฺติก \textbf{3.}\csromanspacer{}\textbf{ อนิยตปท 1.}\csromanspacer{}\textbf{ ปฏิจฺจวาร}}
\csromanheader{2}{\textbf{ปจฺจยจตุกฺก-เหตุ}}
\proseitem{445}{\csromanfolio{297}เหตุํ อนิยตํ ธมฺมํ ปฏิจฺจ เหตุ อนิยโต ธมฺโม อุปฺปชฺชติ เหตุปจฺจยา.\csromanspacer{} ตีณิ.}

\prose{นเหตุํ อนิยตํ ธมฺมํ ปฏิจฺจ นเหตุ อนิยโต ธมฺโม อุปฺปชฺชติ เหตุปจฺจยา.\csromanspacer{} ตีณิ.}

\prose{เหตุํ อนิยตญฺจ นเหตุํ อนิยตญฺจ ธมฺมํ ปฏิจฺจ เหตุ อนิยโต ธมฺโม อุปฺปชฺชติ เหตุปจฺจยา.\csromanspacer{} ตีณิ. (สํขิตฺตํ.)}

\proseitem{446}{เหตุยา นว, อารมฺมเณ นว, อธิปติยา นว ฯเปฯ กมฺเม นว, วิปาเก นว ฯเปฯ อวิคเต นว. (สํขิตฺตํ.)}

\csromanheader{2}{\textbf{นเหตุ น-อารมฺมณ}}
\proseitem{447}{นเหตุํ อนิยตํ ธมฺมํ ปฏิจฺจ นเหตุ อนิยโต ธมฺโม อุปฺปชฺชติ นเหตุปจฺจยา.\csromanspacer{}\csromanspacer{}นเหตุํ อนิยตํ ธมฺมํ ปฏิจฺจ เหตุ อนิยโต ธมฺโม อุปฺปชฺชติ นเหตุปจฺจยา. (2)}

\prose{เหตุํ อนิยตํ ธมฺมํ ปฏิจฺจ นเหตุ อนิยโต ธมฺโม อุปฺปชฺชติ น-อารมฺมณปจฺจยา. (1)}

\prose{นเหตุํ อนิยตํ ธมฺมํ ปฏิจฺจ นเหตุ อนิยโต ธมฺโม อุปฺปชฺชติ น-อารมฺมณปจฺจยา. (1)}

\prose{เหตุํ อนิยตญฺจ นเหตุํ อนิยตญฺจ ธมฺมํ ปฏิจฺจ นเหตุ อนิยโต ธมฺโม อุปฺปชฺชติ น-อารมฺมณปจฺจยา. (1) (สํขิตฺตํ.)}

\proseitem{448}{นเหตุยา เทฺว, น-อารมฺมเณ ตีณิ, น-อธิปติยา นว, น-อนนฺตเร ตีณิ, นสมนนฺตเร ตีณิ, น-อญฺญมญฺเญ ตีณิ, น-อุปนิสฺสเย ตีณิ, นปุเรชาเต นว, นปจฺฉาชาเต นว, น-อาเสวเน นว, นกมฺเม ตีณิ, นวิปาเก นว, น-อาหาเร เอกํ, น-อินฺทฺริเย เอกํ, นฌาเน เอกํ, นมคฺเค เอกํ, นสมฺปยุตฺเต ตีณิ, นวิปฺปยุตฺเต นว, โนนตฺถิยา ตีณิ, โนวิคเต ตีณิ. (สํขิตฺตํ.)}

\gambhira{ธมฺมานุโลม ทุกติกปฏฺฐานปาฬิ}
\csromanheader{2}{เหตุปจฺจยา น-อารมฺมเณ ตีณิ. (สํขิตฺตํ.)}
\csromanheader{2}{นเหตุปจฺจยา อารมฺมเณ เทฺว. (สํขิตฺตํ.)}
\prose{\csromanfolio{298}(สหชาตวาโรปิ ปจฺจยวาโรปิ นิสฺสยวาโรปิ สํสฏฺฐวาโรปิ สมฺปยุตฺตวาโรปิ ปฏิจฺจวารสทิสา วิตฺถาเรตพฺพา.)}

\csromanheader{2}{3. อนิยตปท 7. ปญฺหาวาร}
\csromanheader{2}{\textbf{ปจฺจยจตุกฺก-เหตฺวาทิ}}
\proseitem{449}{เหตุ อนิยโต ธมฺโม เหตุสฺส อนิยตสฺส ธมฺมสฺส เหตุปจฺจเยน ปจฺจโย.\csromanspacer{} ตีณิ.}

\prose{เหตุ อนิยโต ธมฺโม เหตุสฺส อนิยตสฺส ธมฺมสฺส อารมฺมณปจฺจเยน ปจฺจโย.\csromanspacer{} นว.}

\prose{เหตุ อนิยโต ธมฺโม เหตุสฺส อนิยตสฺส ธมฺมสฺส อธิปติปจฺจเยน ปจฺจโย.\csromanspacer{} อารมฺมณาธิปติ, สหชาตาธิปติ.\csromanspacer{} ตีณิ.}

\prose{นเหตุ อนิยโต ธมฺโม นเหตุสฺส อนิยตสฺส ธมฺมสฺส อธิปติปจฺจเยน ปจฺจโย.\csromanspacer{} อารมฺมณาธิปติ, สหชาตาธิปติ.\csromanspacer{} ตีณิ.}

\prose{เหตุ อนิยโต จ นเหตุ อนิยโต จ ธมฺมา เหตุสฺส อนิยตสฺส ธมฺมสฺส อธิปติปจฺจเยน ปจฺจโย.\csromanspacer{} อารมฺมณาธิปติ, ตีณิ. (สํขิตฺตํ.)}

\proseitem{450}{เหตุยา ตีณิ, อารมฺมเณ นว, อธิปติยา นว, อนนฺตเร นว, สมนนฺตเร นว, สหชาเต นว, อญฺญมญฺเญ นว, นิสฺสเย นว, อุปนิสฺสเย นว, ปุเรชาเต ตีณิ, ปจฺฉาชาเต ตีณิ, อาเสวเน นว, กมฺเม ตีณิ, วิปาเก นว, อาหาเร ตีณิ, อินฺทฺริเย นว, ฌาเน ตีณิ, มคฺเค นว, สมฺปยุตฺเต นว, วิปฺปยุตฺเต ปญฺจ, อตฺถิยา นว, นตฺถิยา นว, วิคเต นว, อวิคเต นว.}

\csromanheader{2}{1. เหตุทุก 16. มคฺคารมฺมณตฺติก \textbf{ปจฺจนียุทฺธาร}}
\proseitem{451}{\csromanfolio{299}เหตุ อนิยโต ธมฺโม เหตุสฺส อนิยตสฺส ธมฺมสฺส อารมฺมณปจฺจเยน ปจฺจโย.\csromanspacer{} สหชาตปจฺจเยน ปจฺจโย.\csromanspacer{} อุปนิสฺสยปจฺจเยน ปจฺจโย. (สํขิตฺตํ.)}

\proseitem{452}{นเหตุยา นว, น-อารมฺมเณ นว. (สํขิตฺตํ.)}

\csromanheader{2}{เหตุปจฺจยา น-อารมฺมเณ ตีณิ. (สํขิตฺตํ.)}
\csromanheader{2}{นเหตุปจฺจยา อารมฺมเณ นว. (สํขิตฺตํ.)}
\prose{(ยถา กุสลตฺติเก ปญฺหาวารสฺส อนุโลมมฺปิ ปจฺจนียมฺปิ อนุโลมปจฺจนียมฺปิ ปจฺจนียานุโลมมฺปิ คณิตํ, เอวํ คเณตพฺพํ.)}

\nitthitam{เหตุทุกมิจฺฉตฺตนิยตตฺติกํ นิฏฺฐิตํ.}
\csromanmatikaanchor{mtk.47}
\csromantocmarkref{section}{16. มคฺคารมฺมณตฺติก}{mtk.47}
\bookmark[dest={mtk.47},level=1]{16. มคฺคารมฺมณตฺติก}
\csromanheader{1}{1. เหตุทุก 16. มคฺคารมฺมณตฺติก}
\csromanheader{2}{1. มคฺคารมฺมณปท 1. ปฏิจฺจวาร}
\csromanheader{2}{\textbf{ปจฺจยจตุกฺก-เหตุ}}
\proseitem{453}{เหตุํ มคฺคารมฺมณํ ธมฺมํ ปฏิจฺจ เหตุ มคฺคารมฺมโณ ธมฺโม อุปฺปชฺชติ เหตุปจฺจยา.\csromanspacer{} ตีณิ.}

\prose{นเหตุํ มคฺคารมฺมณํ ธมฺมํ ปฏิจฺจ นเหตุ มคฺคารมฺมโณ ธมฺโม อุปฺปชฺชติ เหตุปจฺจยา.\csromanspacer{} ตีณิ.}

\prose{เหตุํ มคฺคารมฺมณญฺจ นเหตุํ มคฺคารมฺมณญฺจ ธมฺมํ ปฏิจฺจ เหตุ มคฺคารมฺมโณ ธมฺโม อุปฺปชฺชติ เหตุปจฺจยา.\csromanspacer{} ตีณิ. (สํขิตฺตํ.)}

\proseitem{454}{เหตุยา นว, อารมฺมเณ นว, อธิปติยา นว ฯเปฯ กมฺเม นว ฯเปฯ อวิคเต นว. (สํขิตฺตํ.)}

\gambhira{ธมฺมานุโลม ทุกติกปฏฺฐานปาฬิ \textbf{นเหตุ}}
\vspace{\tipitakaparskip}
\csromancenter{\textbf{นเหตุ}ํ มคฺคารมฺมณํ ธมฺมํ ปฏิจฺจ นเหตุ มคฺคารมฺมโณ ธมฺโม อุปฺปชฺชติ นเหตุปจฺจยา. (สํขิตฺตํ.)}
\csromancenter{\textbf{นเหตุ}ยา เอกํ, น-อธิปติยา นว, นปุเรชาเต นว, นปจฺฉาชาเต นว, น-อาเสวเน นว, นกมฺเม ตีณิ, นวิปาเก นว, นมคฺเค เอกํ, นวิปฺปยุตฺเต นว. (สํขิตฺตํ.)}
\csromanheader{2}{เหตุปจฺจยา น-อธิปติยา นว. (สํขิตฺตํ.)}
\csromanheader{2}{\textbf{นเหตุ}ปจฺจยา อารมฺมเณ เอกํ. (สํขิตฺตํ.)}
\prose{\csromanfolio{300}(สหชาตวาโรปิ ปจฺจยวาโรปิ นิสฺสยวาโรปิ สํสฏฺฐวาโรปิ สมฺปยุตฺตวาโรปิ ปฏิจฺจวารสทิสา วิตฺถาเรตพฺพา.)}

\csromanheader{2}{1. มคฺคารมฺมณปท 7. ปญฺหาวาร}
\csromanheader{2}{\textbf{ปจฺจยจตุกฺก-เหตฺวาทิ}}
\proseitem{457}{เหตุ มคฺคารมฺมโณ ธมฺโม เหตุสฺส มคฺคารมฺมณสฺส ธมฺมสฺส เหตุปจฺจเยน ปจฺจโย.\csromanspacer{} ตีณิ.}

\prose{เหตุ มคฺคารมฺมโณ ธมฺโม เหตุสฺส มคฺคารมฺมณสฺส ธมฺมสฺส อธิปติปจฺจเยน ปจฺจโย.\csromanspacer{} สหชาตาธิปติ.\csromanspacer{} ตีณิ.}

\vspace{\tipitakaparskip}
\csromancenter{\textbf{นเหตุ} มคฺคารมฺมโณ ธมฺโม นเหตุสฺส มคฺคารมฺมณสฺส ธมฺมสฺส อธิปติปจฺจเยน ปจฺจโย.\csromanspacer{} สหชาตาธิปติ.\csromanspacer{} ตีณิ.}
\prose{เหตุ มคฺคารมฺมโณ ธมฺโม เหตุสฺส มคฺคารมฺมณสฺส ธมฺมสฺส อนนฺตรปจฺจเยน ปจฺจโย. (สํขิตฺตํ.)}

\proseitem{458}{เหตุยา ตีณิ, อธิปติยา ฉ, อนนฺตเร นว, สมนนฺตเร นว, สหชาเต นว, อญฺญมญฺเญ นว, นิสฺสเย นว, อุปนิสฺสเย นว, อาเสวเน นว, กมฺเม ตีณิ, อาหาเร ตีณิ, อินฺทฺริเย นว, ฌาเน ตีณิ, มคฺเค นว, สมฺปยุตฺเต นว, อตฺถิยา นว, นตฺถิยา นว, วิคเต นว, อวิคเต นว. (สํขิตฺตํ.)}

\csromanheader{2}{1. เหตุทุก 16. มคฺคารมฺมณตฺติก \textbf{ปจฺจนียุทฺธาร}}
\proseitem{459}{\csromanfolio{301}เหตุ มคฺคารมฺมโณ ธมฺโม เหตุสฺส มคฺคารมฺมณสฺส ธมฺมสฺส สหชาตปจฺจเยน ปจฺจโย.\csromanspacer{} อุปนิสฺสยปจฺจเยน ปจฺจโย. (สํขิตฺตํ.)}

\proseitem{460}{นเยตุยา นว, น-อารมฺมเณ นว. (สํขิตฺตํ.)}

\csromanheader{2}{เหตุปจฺจยา น-อารมฺมเณ ตีณิ. (สํขิตฺตํ.)}
\csromanheader{2}{นเหตุปจฺจยา อธิปติยา ตีณิ. (สํขิตฺตํ.)}
\prose{(ยถา กุสลตฺติเก ปญฺหาวารสฺส อนุโลมมฺปิ ปจฺจนียมฺปิ อนุโลมปจฺจนียมฺปิ ปจฺจนียานุโลมมฺปิ คณิตํ, เอวํ คเณตพฺพํ.)}

\csromanheader{2}{2. มคฺคเหตุกปท 1. ปฏิจฺจวาร}
\csromanheader{2}{\textbf{ปจฺจยจตุกฺก-เหตุ}}
\proseitem{461}{เหตุํ มคฺคเหตุกํ ธมฺมํ ปฏิจฺจ เหตุ มคฺคเหตุโก ธมฺโม อุปฺปชฺชติ เหตุปจฺจยา.\csromanspacer{} ตีณิ.}

\prose{นเหตุํ มคฺคเหตุกํ ธมฺมํ ปฏิจฺจ นเหตุ มคฺคเหตุโก ธมฺโม อุปฺปชฺชติ เหตุปจฺจยา.\csromanspacer{} ตีณิ.}

\prose{เหตุํ มคฺคเหตุกญฺจ นเหตุํ มคฺคเหตุกญฺจ ธมฺมํ ปฏิจฺจ เหตุ มคฺคเหตุโก ธมฺโม อุปฺปชฺชติ เหตุปจฺจยา.\csromanspacer{} ตีณิ. (สํขิตฺตํ.)}

\proseitem{462}{เหตุยา นว, อารมฺมเณ นว, อธิปติยา นว ฯเปฯ อาเสวเน นว, กมฺเม นว, อาหาเร นว, อินฺทฺริเย นว ฯเปฯ อวิคเต นว. (สํขิตฺตํ.)}

\csromanheader{2}{\textbf{น-อธิปติ}}
\proseitem{463}{เหตุํ มคฺคเหตุกํ ธมฺมํ ปฏิจฺจ เหตุ มคฺคเหตุโก ธมฺโม อุปฺปชฺชติ น-อธิปติปจฺจยา. (สํขิตฺตํ.)}

\vspace{\tipitakaparskip}
\csromancenter{\textbf{น-อธิปติ}ยา ฉ, นปุเรชาเต นว, นปจฺฉาชาเต นว, นกมฺเม ตีณิ, นวิปาเก นว, นวิปฺปยุตฺเต นว. (สํขิตฺตํ.)}
\gambhira{ธมฺมานุโลม ทุกติกปฏฺฐานปาฬิ}
\csromanheader{2}{เหตุปจฺจยา น-อธิปติยา ฉ. (สํขิตฺตํ.)}
\csromanheader{2}{น-อธิปติปจฺจยา เหตุยา ฉ. (สํขิตฺตํ.)}
\prose{\csromanfolio{302}(สหชาตวาโรปิ ปจฺจยวาโรปิ นิสฺสยวาโรปิ สํสฏฺฐวาโรปิ สมฺปยุตฺตวาโรปิ ปฏิจฺจวารสทิสา วิตฺถาเรตพฺพา.)}

\csromanheader{2}{2. มคฺคเหตุกปท 7. ปญฺหาวาร}
\csromanheader{2}{\textbf{ปจฺจยจตุกฺก-เหตุ อธิปติ}}
\proseitem{465}{เหตุ มคฺคเหตุโก ธมฺโม เหตุสฺส มคฺคเหตุกสฺส ธมฺมสฺส เหตุปจฺจเยน ปจฺจโย.\csromanspacer{} ตีณิ.}

\prose{เหตุ มคฺคเหตุโก ธมฺโม เหตุสฺส มคฺคเหตุกสฺส ธมฺมสฺส อธิปติปจฺจเยน ปจฺจโย.\csromanspacer{} สหชาตาธิปติ.\csromanspacer{} ตีณิ.}

\prose{นเหตุ มคฺคเหตุโก ธมฺโม นเหตุสฺส มคฺคเหตุกสฺส ธมฺมสฺส อธิปติปจฺจเยน ปจฺจโย.\csromanspacer{} สหชาตาธิปติ.\csromanspacer{} ตีณิ. (สํขิตฺตํ.)}

\proseitem{466}{เหตุยา ตีณิ, อธิปติยา ฉ, สหชาเต นว, อญฺญมญฺเญ นว, นิสฺสเย นว, อุปนิสฺสเย นว, กมฺเม ตีณิ, อาหาเร ตีณิ, อินฺทฺริเย นว, ฌาเน ตีณิ, มคฺเค นว, สมฺปยุตฺเต นว, อตฺถิยา นว, อวิคเต นว. (สํขิตฺตํ.)}

\csromanheader{2}{\textbf{ปจฺจนียุทฺธาร}}
\proseitem{467}{เหตุ มคฺคเหตุโก ธมฺโม เหตุสฺส มคฺคเหตุกสฺส ธมฺมสฺส สหชาตปจฺจเยน ปจฺจโย.\csromanspacer{} อุปนิสฺสยปจฺจเยน ปจฺจโย. (สํขิตฺตํ.)}

\proseitem{468}{นเหตุยา นว, น-อารมฺมเณ นว. (สํขิตฺตํ.)}

\csromanheader{2}{เหตุปจฺจยา น-อารมฺมเณ ตีณิ. (สํขิตฺตํ.)}
\csromanheader{2}{นเหตุปจฺจยา อธิปติยา ตีณิ. (สํขิตฺตํ.)}
\prose{(ยถา กุสลตฺติเก ปญฺหาวารสฺส อนุโลมมฺปิ ปจฺจนียมฺปิ อนุโลมปจฺจนียมฺปิ ปจฺจนียานุโลมมฺปิ คณิตํ, เอวํ คเณตพฺพํ.)}

\vspace{\tipitakaparskip}
\csromancenter{เหตุทุก 16.\csromanspacer{} มคฺคารมฺมณตฺติก \textbf{3.}\csromanspacer{}\textbf{ มคฺคาธิปติปท 1.}\csromanspacer{}\textbf{ ปฏิจฺจวาร}}
\csromanheader{2}{\textbf{ปจฺจยจตุกฺก-เหตุ}}
\proseitem{469}{\csromanfolio{303}เหตุํ มคฺคาธิปติํ ธมฺมํ ปฏิจฺจ เหตุ มคฺคาธิปติ ธมฺโม อุปฺปชฺชติ เหตุปจฺจยา.\csromanspacer{} ตีณิ.}

\prose{นเหตุํ มคฺคาธิปติํ ธมฺมํ ปฏิจฺจ นเหตุ มคฺคาธิปติ ธมฺโม อุปฺปชฺชติ เหตุปจฺจยา.\csromanspacer{} ตีณิ.}

\prose{เหตุํ มคฺคาธิปติญฺจ นเหตุํ มคฺคาธิปติญฺจ ธมฺมํ ปฏิจฺจ เหตุ มคฺคาธิปติ ธมฺโม อุปฺปชฺชติ เหตุปจฺจยา.\csromanspacer{} ตีณิ. (สํขิตฺตํ.)}

\proseitem{470}{เหตุยา นว, อารมฺมเณ นว ฯเปฯ กมฺเม นว, อาหาเร นว ฯเปฯ อวิคเต นว. (สํขิตฺตํ.)}

\csromanheader{2}{\textbf{น-อธิปติ}}
\proseitem{471}{เหตุํ มคฺคาธิปติํ ธมฺมํ ปฏิจฺจ เหตุ มคฺคาธิปติ ธมฺโม อุปฺปชฺชติ น-อธิปติปจฺจยา. (สํขิตฺตํ.)}

\vspace{\tipitakaparskip}
\csromancenter{\textbf{น-อธิปติ}ยา นว, นปุเรชาเต นว, นปจฺฉาชาเต นว, น-อาเสวเน นว, นกมฺเม ตีณิ, นวิปาเก นว, นวิปฺปยุตฺเต นว. (สํขิตฺตํ.)}
\csromanheader{2}{เหตุปจฺจยา น-อธิปติยา นว. (สํขิตฺตํ.)}
\csromanheader{2}{\textbf{น-อธิปติ}ปจฺจยา เหตุยา นว. (สํขิตฺตํ.)}
\prose{(สหชาตวาโรปิ ปจฺจยวาโรปิ นิสฺสยวาโรปิ สํสฏฺฐวาโรปิ สมฺปยุตฺตวาโรปิ ปฏิจฺจวารสทิสา วิตฺถาเรตพฺพา.)}

\csromanheader{2}{3. มคฺคาธิปติปท 7. ปญฺหาวาร}
\csromanheader{2}{\textbf{ปจฺจยจตุกฺก-เหตฺวาทิ}}
\proseitem{473}{เหตุ มคฺคาธิปติ ธมฺโม เหตุสฺส มคฺคาธิปติสฺส ธมฺมสฺส เหตุปจฺจเยน ปจฺจโย.\csromanspacer{} ตีณิ.}

\gambhira{ธมฺมานุโลม ทุกติกปฏฺฐานปาฬิ}
\prose{\csromanfolio{304}เหตุ มคฺคาธิปติ ธมฺโม เหตุสฺส มคฺคาธิปติสฺส ธมฺมสฺส อารมฺมณปจฺจเยน ปจฺจโย.\csromanspacer{} นว.}

\prose{เหตุ มคฺคาธิปติ ธมฺโม เหตุสฺส มคฺคาธิปติสฺส ธมฺมสฺส อธิปติปจฺจเยน ปจฺจโย.\csromanspacer{} อารมฺมณาธิปติ, สหชาตาธิปติ.\csromanspacer{} ตีณิ.}

\prose{นเหตุ มคฺคาธิปติ ธมฺโม นเหตุสฺส มคฺคาธิปติสฺส ธมฺมสฺส อธิปติปจฺจเยน ปจฺจโย.\csromanspacer{} อารมฺมณาธิปติ, สหชาตาธิปติ.\csromanspacer{} ตีณิ.}

\prose{เหตุ มคฺคาธิปติ จ นเหตุ มคฺคาธิปติ จ ธมฺมา เหตุสฺส มคฺคาธิปติสฺส ธมฺมสฺส อธิปติปจฺจเยน ปจฺจโย.\csromanspacer{} อารมฺมณาธิปติ.\csromanspacer{} ตีณิ. (สํขิตฺตํ.)}

\proseitem{474}{เหตุยา ตีณิ, อารมฺมเณ นว, อธิปติยา นว, อนนฺตเร นว, สมนนฺตเร นว, สหชาเต นว, อญฺญมญฺเญ นว, นิสฺสเย นว, อุปนิสฺสเย นว, อาเสวเน นว, กมฺเม ตีณิ, อาหาเร ตีณิ, อินฺทฺริเย นว, ฌาเน ตีณิ, มคฺเค นว, สมฺปยุตฺเต นว, อตฺถิยา นว ฯเปฯ อวิคเต นว. (สํขิตฺตํ.)}

\csromanheader{2}{\textbf{ปจฺจนียุทฺธาร}}
\proseitem{475}{เหตุ มคฺคาธิปติ ธมฺโม เหตุสฺส มคฺคาธิปติสฺส ธมฺมสฺส อารมฺมณปจฺจเยน ปจฺจโย.\csromanspacer{} สหชาตปจฺจเยน ปจฺจโย.\csromanspacer{} อุปนิสฺสยปจฺจเยน ปจฺจโย. (สํขิตฺตํ.)}

\proseitem{476}{นเหตุยา นว, น-อารมฺมเณ นว. (สํขิตฺตํ.)}

\csromanheader{2}{เหตุปจฺจยา น-อารมฺมเณ ตีณิ. (สํขิตฺตํ.)}
\csromanheader{2}{นเหตุปจฺจยา อารมฺมเณ นว. (สํขิตฺตํ.)}
\prose{(ยถา กุสลตฺติเก ปญฺหาวารสฺส อนุโลมมฺปิ ปจฺจนียมฺปิ อนุโลมปจฺจนียมฺปิ ปจฺจนียานุโลมมฺปิ คณิตํ, เอวํ คเณตพฺพํ.)}

\nitthitam{เหตุทุกมคฺคารมฺมณตฺติกํ นิฏฺฐิตํ.}
\csromanmatikaanchor{mtk.48}
\csromantocmarkref{section}{17. อุปฺปนฺนตฺติก}{mtk.48}
\bookmark[dest={mtk.48},level=1]{17. อุปฺปนฺนตฺติก}
\csromanheader{1}{1. เหตุทุก 17. อุปฺปนฺนตฺติก \textbf{1. เหตุทุก 17. อุปฺปนฺนตฺติก}}
\csromanheader{2}{1. อุปฺปนฺนปท 7. ปญฺหาวาร}
\csromanheader{2}{\textbf{ปจฺจยจตุกฺก-เหตุ อุปนิสฺสย}}
\proseitem{477}{\csromanfolio{305}เหตุ อุปฺปนฺโน ธมฺโม เหตุสฺส อุปฺปนฺนสฺส ธมฺมสฺส เหตุปจฺจเยน ปจฺจโย.\csromanspacer{} ตีณิ.}

\prose{นเหตุ อุปฺปนฺโน ธมฺโม นเหตุสฺส อุปฺปนฺนสฺส ธมฺมสฺส อารมฺมณปจฺจเยน ปจฺจโย.\csromanspacer{} ตีณิ.}

\prose{เหตุ อุปฺปนฺโน ธมฺโม เหตุสฺส อุปฺปนฺนสฺส ธมฺมสฺส อธิปติปจฺจเยน ปจฺจโย.\csromanspacer{} สหชาตาธิปติ.\csromanspacer{} ตีณิ.}

\prose{นเหตุ อุปฺปนฺโน ธมฺโม นเหตุสฺส อุปฺปนฺนสฺส ธมฺมสฺส อธิปติปจฺจเยน ปจฺจโย.\csromanspacer{} อารมฺมณาธิปติ, สหชาตาธิปติ.\csromanspacer{} ตีณิ.}

\prose{เหตุ อุปฺปนฺโน ธมฺโม เหตุสฺส อุปฺปนฺนสฺส ธมฺมสฺส สหชาตปจฺจเยน ปจฺจโย. (สํขิตฺตํ.)}

\prose{นเหตุ อุปฺปนฺโน ธมฺโม นเหตุสฺส อุปฺปนฺนสฺส ธมฺมสฺส อุปนิสฺสยปจฺจเยน ปจฺจโย.\csromanspacer{} \textbf{อารมฺมณูปนิสฺสโย, ปกตูปนิสฺสโย ฯเปฯ.}\csromanspacer{} ปกตูปนิสฺสโย\csromandash{}อุปฺปนฺนํ อุตุํ อุปนิสฺสาย ฌานํ อุปฺปาเทติ, วิปสฺสนํ.\csromanspacer{} มคฺคํ.\csromanspacer{} อภิญฺญํ.\csromanspacer{} สมาปตฺติํ อุปฺปาเทติ. (สํขิตฺตํ.)}

\proseitem{478}{เหตุยา ตีณิ, อารมฺมเณ ตีณิ, อธิปติยา ฉ, สหชาเต นว, อญฺญมญฺเญ นว, นิสฺสเย นว, อุปนิสฺสเย ตีณิ, ปุเรชาเต ตีณิ, ปจฺฉาชาเต ตีณิ, กมฺเม ตีณิ, วิปาเก นว, อาหาเร ตีณิ, อินฺทฺริเย นว, ฌาเน ตีณิ, มคฺเค นว, สมฺปยุตฺเต นว, วิปฺปยุตฺเต ปญฺจ, อตฺถิยา นว, อวิคเต นว. (สํขิตฺตํ.)}

\csromanheader{2}{\textbf{ปจฺจนียุทฺธาร}}
\proseitem{479}{เหตุ อุปฺปนฺโน ธมฺโม เหตุสฺส อุปฺปนฺนสฺส ธมฺมสฺส สหชาตปจฺจเยน ปจฺจโย. (สํขิตฺตํ.)}

\gambhira{ธมฺมานุโลม ทุกติกปฏฺฐานปาฬิ}
\proseitem{480}{\csromanfolio{306}นเหตุยา นว, น-อารมฺมเณ นว. (สํขิตฺตํ.)}

\csromanheader{2}{เหตุปจฺจยา น-อารมฺมเณ ตีณิ. (สํขิตฺตํ.)}
\prose{นเหตุปจฺจยา อารมฺมเณ ตีณิ. (สํขิตฺตํ.) (ยถา กุสลตฺติเก ปญฺหาวารสฺส อนุโลมมฺปิ ปจฺจนียมฺปิ อนุโลมปจฺจนียมฺปิ ปจฺจนียานุโลมมฺปิ คณิตํ, เอวํ คเณตพฺพํ.\csromanspacer{} อิมมฺหิ ทุกติเก ปฏิจฺจวารมฺปิ สหชาตวารมฺปิ ปจฺจยวารมฺปิ นิสฺสยวารมฺปิ สํสฏฺฐวารมฺปิ สมฺปยุตฺตวารมฺปิ อนุปฺปนฺนมฺปิ อุปฺปาทีปิ นตฺถิ.)}

\nitthitam{เหตุทุก-อุปฺปนฺนตฺติกํ นิฏฺฐิตํ.}
\csromanmatikaanchor{mtk.49}
\csromantocmarkref{section}{18. อตีตตฺติก}{mtk.49}
\bookmark[dest={mtk.49},level=1]{18. อตีตตฺติก}
\csromanheader{1}{1. เหตุทุก 18. อตีตตฺติก}
\csromanheader{2}{3. ปจฺจุปฺปนฺนปท 7. ปญฺหาวาร}
\csromanheader{2}{\textbf{ปจฺจยจตุกฺก-เหตฺวาทิ}}
\proseitem{481}{เหตุ ปจฺจุปฺปนฺโน ธมฺโม เหตุสฺส ปจฺจุปฺปนฺนสฺส ธมฺมสฺส เหตุปจฺจเยน ปจฺจโย.\csromanspacer{} ตีณิ.}

\prose{นเหตุ ปจฺจุปฺปนฺโน ธมฺโม นเหตุสฺส ปจฺจุปฺปนฺนสฺส ธมฺมสฺส อารมฺมณปจฺจเยน ปจฺจโย.\csromanspacer{} ตีณิ.}

\prose{เหตุ ปจฺจุปฺปนฺโน ธมฺโม เหตุสฺส ปจฺจุปฺปนฺนสฺส ธมฺมสฺส อธิปติปจฺจเยน ปจฺจโย.\csromanspacer{} สหชาตาธิปติ.\csromanspacer{} ตีณิ.}

\prose{นเหตุ ปจฺจุปฺปนฺโน ธมฺโม นเหตุสฺส ปจฺจุปฺปนฺนสฺส ธมฺมสฺส อธิปติปจฺจเยน ปจฺจโย.\csromanspacer{} อารมฺมณาธิปติ, สหชาตาธิปติ.\csromanspacer{} ตีณิ.}

\prose{เหตุ ปจฺจุปฺปนฺโน ธมฺโม เหตุสฺส ปจฺจุปฺปนฺนสฺส ธมฺมสฺส สหชาตปจฺจเยน ปจฺจโย. (สํขิตฺตํ.)}

\proseitem{482}{เหตุยา ตีณิ, อารมฺมเณ ตีณิ, อธิปติยา ฉ, สหชาเต นว, อญฺญมญฺเญ นว, นิสฺสเย นว, อุปนิสฺสเย นว, ปุเรชาเต ตีณิ, ปจฺฉาชาเต ตีณิ, กมฺเม ตีณิ, วิปาเก นว, อาหาเร ตีณิ, อินฺทฺริเย}

\vspace{\tipitakaparskip}
\csromancenter{เหตุทุก 19.\csromanspacer{} อตีตารมฺมณตฺติก นว, ฌาเน ตีณิ, มคฺเค นว, สมฺปยุตฺเต นว, วิปฺปยุตฺเต ปญฺจ, อตฺถิยา นว, อวิคเต นว. (สํขิตฺตํ.)}
\csromanheader{2}{\textbf{ปจฺจนียุทฺธาร}}
\proseitem{483}{\csromanfolio{307}เหตุ ปจฺจุปฺปนฺโน ธมฺโม เหตุสฺส ปจฺจุปฺปนฺนสฺส ธมฺมสฺส สหชาตปจฺจเยน ปจฺจโย. (สํขิตฺตํ.)}

\proseitem{484}{นเหตุยา นว, น-อารมฺมเณ นว. (สํขิตฺตํ.)}

\csromanheader{2}{เหตุปจฺจยา น-อารมฺมเณ ตีณิ. (สํขิตฺตํ.)}
\prose{นเหตุปจฺจยา อารมฺมเณ ตีณิ. (สํขิตฺตํ.) (ยถา กุสลตฺติเก ปญฺหาวารสฺส อนุโลมมฺปิ ปจฺจนียมฺปิ อนุโลมปจฺจนียมฺปิ ปจฺจนียานุโลมมฺปิ คณิตํ, เอวํ คเณตพฺพํ.\csromanspacer{} อิมมฺหิ ทุกติเก ปฏิจฺจวารมฺปิ สหชาตวารมฺปิ ปจฺจยวารมฺปิ นิสฺสยวารมฺปิ สํสฏฺฐวารมฺปิ สมฺปยุตฺตวารมฺปิ อตีตมฺปิ อนาคตมฺปิ นตฺถิ.)}

\nitthitam{เหตุทุก-อตีตตฺติกํ นิฏฺฐิตํ.}
\csromanmatikaanchor{mtk.50}
\csromantocmarkref{section}{19. อตีตารมฺมณตฺติก}{mtk.50}
\bookmark[dest={mtk.50},level=1]{19. อตีตารมฺมณตฺติก}
\csromanheader{1}{1. เหตุทุก 19. อตีตารมฺมณตฺติก}
\csromanheader{2}{1. อตีตารมฺมณปท 1. ปฏิจฺจวาร}
\csromanheader{2}{\textbf{ปจฺจยจตุกฺก-เหตุ}}
\proseitem{485}{เหตุํ อตีตารมฺมณํ ธมฺมํ ปฏิจฺจ เหตุ อตีตารมฺมโณ ธมฺโม อุปฺปชฺชติ เหตุปจฺจยา.\csromanspacer{} ตีณิ.}

\prose{นเหตุํ อตีตารมฺมณํ ธมฺมํ ปฏิจฺจ นเหตุ อตีตารมฺมโณ ธมฺโม อุปฺปชฺชติ เหตุปจฺจยา.\csromanspacer{} ตีณิ.}

\prose{เหตุํ อตีตารมฺมณญฺจ นเหตุํ อตีตารมฺมณญฺจ ธมฺมํ ปฏิจฺจ เหตุ อตีตารมฺมโณ ธมฺโม อุปฺปชฺชติ เหตุปจฺจยา.\csromanspacer{} ตีณิ. (สํขิตฺตํ.)}

\gambhira{ธมฺมานุโลม ทุกติกปฏฺฐานปาฬิ}
\proseitem{486}{\csromanfolio{308}เหตุยา นว, อารมฺมเณ นว, อธิปติยา นว, อนนฺตเร นว, สมนนฺตเร นว ฯเปฯ กมฺเม นว ฯเปฯ อวิคเต นว. (สํขิตฺตํ.)}

\csromanheader{2}{\textbf{นเหตุ น-อธิปติ}}
\proseitem{487}{นเหตุํ อตีตารมฺมณํ ธมฺมํ ปฏิจฺจ นเหตุ อตีตารมฺมโณ ธมฺโม อุปฺปชฺชติ นเหตุปจฺจยา.\csromanspacer{}\csromanspacer{}นเหตุํ อตีตารมฺมณํ ธมฺมํ ปฏิจฺจ เหตุ อตีตารมฺมโณ ธมฺโม อุปฺปชฺชติ นเหตุปจฺจยา. (2)}

\prose{เหตุํ อตีตารมฺมณํ ธมฺมํ ปฏิจฺจ เหตุ อตีตารมฺมโณ ธมฺโม อุปฺปชฺชติ น-อธิปติปจฺจยา. (สํขิตฺตํ.)}

\proseitem{488}{นเหตุยา เทฺว, น-อธิปติยา นว, นปุเรชาเต นว, นปจฺฉาชาเต นว, น-อาเสวเน นว, นกมฺเม ตีณิ, นวิปาเก นว, นมคฺเค เอกํ, นวิปฺปยุตฺเต นว. (สํขิตฺตํ.)}

\csromanheader{2}{เหตุปจฺจยา น-อธิปติยา นว. (สํขิตฺตํ.)}
\csromanheader{2}{นเหตุปจฺจยา อารมฺมเณ เทฺว. (สํขิตฺตํ.)}
\prose{(สหชาตวาโรปิ ปจฺจยวาโรปิ นิสฺสยวาโรปิ สํสฏฺฐวาโรปิ สมฺปยุตฺตวาโรปิ ปฏิจฺจวารสทิสา วิตฺถาเรตพฺพา.)}

\csromanheader{2}{1. อตีตารมฺมณปท 7. ปญฺหาวาร}
\csromanheader{2}{\textbf{ปจฺจยจตุกฺก-เหตฺวาทิ}}
\proseitem{489}{เหตุ อตีตารมฺมโณ ธมฺโม เหตุสฺส อตีตารมฺมณสฺส ธมฺมสฺส เหตุปจฺจเยน ปจฺจโย.\csromanspacer{} ตีณิ.}

\prose{เหตุ อตีตารมฺมโณ ธมฺโม เหตุสฺส อตีตารมฺมณสฺส ธมฺมสฺส อารมฺมณปจฺจเยน ปจฺจโย.\csromanspacer{} นว.}

\prose{เหตุ อตีตารมฺมโณ ธมฺโม เหตุสฺส อตีตารมฺมณสฺส ธมฺมสฺส อธิปติปจฺจเยน ปจฺจโย.\csromanspacer{} อารมฺมณาธิปติ, สหชาตาธิปติ.\csromanspacer{} ตีณิ.}

\csromanheader{2}{1. เหตุทุก 19. อตีตารมฺมณตฺติก}
\prose{\csromanfolio{309}นเหตุ อตีตารมฺมโณ ธมฺโม นเหตุสฺส อตีตารมฺมณสฺส ธมฺมสฺส อธิปติปจฺจเยน ปจฺจโย.\csromanspacer{} อารมฺมณาธิปติ, สหชาตาธิปติ.\csromanspacer{} ตีณิ.}

\prose{เหตุ อตีตารมฺมโณ จ นเหตุ อตีตารมฺมโณ จ ธมฺมา เหตุสฺส อตีตารมฺมณสฺส ธมฺมสฺส อธิปติปจฺจเยน ปจฺจโย.\csromanspacer{} อารมฺมณาธิปติ.\csromanspacer{} ตีณิ. (สํขิตฺตํ.)}

\proseitem{490}{เหตุยา ตีณิ, อารมฺมเณ นว, อธิปติยา นว, อนนฺตเร นว, สมนนฺตเร นว, สหชาเต นว, อญฺญมญฺเญ นว, นิสฺสเย นว, อุปนิสฺสเย นว, อาเสวเน นว, กมฺเม ตีณิ, วิปาเก นว, อาหาเร ตีณิ, อินฺทฺริเย นว, ฌาเน ตีณิ, มคฺเค นว, สมฺปยุตฺเต นว, อตฺถิยา นว ฯเปฯ อวิคเต นว. (สํขิตฺตํ.)}

\csromanheader{2}{\textbf{ปจฺจนียุทฺธาร}}
\proseitem{491}{เหตุ อตีตารมฺมโณ ธมฺโม เหตุสฺส อตีตารมฺมณสฺส ธมฺมสฺส อารมฺมณปจฺจเยน ปจฺจโย.\csromanspacer{} สหชาตปจฺจเยน ปจฺจโย.\csromanspacer{} อุปนิสฺสยปจฺจเยน ปจฺจโย. (สํขิตฺตํ.)}

\proseitem{492}{นเหตุยา นว, น-อารมฺมเณ นว. (สํขิตฺตํ.)}

\csromanheader{2}{เหตุปจฺจยา น-อารมฺมเณ ตีณิ. (สํขิตฺตํ.)}
\csromanheader{2}{นเหตุปจฺจยา อารมฺมเณ นว. (สํขิตฺตํ.)}
\prose{(ยถา กุสลตฺติเก ปญฺหาวารสฺส อนุโลมมฺปิ ปจฺจนียมฺปิ อนุโลมปจฺจนียมฺปิ ปจฺจนียานุโลมมฺปิ คณิตํ, เอวํ คเณตพฺพํ.)}

\csromanheader{2}{2. อนาคตารมฺมณปท 1. ปฏิจฺจวาร}
\csromanheader{2}{\textbf{ปจฺจยจตุกฺก-เหตุ}}
\proseitem{493}{เหตุํ อนาคตารมฺมณํ ธมฺมํ ปฏิจฺจ เหตุ อนาคตารมฺมโณ ธมฺโม อุปฺปชฺชติ เหตุปจฺจยา.\csromanspacer{} ตีณิ.}

\gambhira{ธมฺมานุโลม ทุกติกปฏฺฐานปาฬิ}
\prose{\csromanfolio{310}นเหตุํ อนาคตารมฺมณํ ธมฺมํ ปฏิจฺจ นเหตุ อนาคตารมฺมโณ ธมฺโม อุปฺปชฺชติ เหตุปจฺจยา.\csromanspacer{} ตีณิ.}

\prose{เหตุํ อนาคตารมฺมณญฺจ นเหตุํ อนาคตารมฺมณญฺจ ธมฺมํ ปฏิจฺจ เหตุ อนาคตารมฺมโณ ธมฺโม อุปฺปชฺชติ เหตุปจฺจยา.\csromanspacer{} ตีณิ. (สํขิตฺตํ.)}

\proseitem{494}{เหตุยา นว, อารมฺมเณ นว ฯเปฯ กมฺเม นว ฯเปฯ อวิคเต นว. (สํขิตฺตํ.)}

\csromanheader{2}{\textbf{นเหตุ น-อธิปติ}}
\proseitem{495}{นเหตุํ อนาคตารมฺมณํ ธมฺมํ ปฏิจฺจ นเหตุ อนาคตารมฺมโณ ธมฺโม อุปฺปชฺชติ นเหตุปจฺจยา.\csromanspacer{}\csromanspacer{}นเหตุํ อนาคตารมฺมณํ ธมฺมํ ปฏิจฺจ เหตุ อนาคตารมฺมโณ ธมฺโม อุปฺปชฺชติ นเหตุปจฺจยา. (2)}

\prose{เหตุํ อนาคตารมฺมณํ ธมฺมํ ปฏิจฺจ เหตุ อนาคตารมฺมโณ ธมฺโม อุปฺปชฺชติ น-อธิปติปจฺจยา. (สํขิตฺตํ.)}

\proseitem{496}{นเหตุยา เทฺว, น-อธิปติยา นว, นปุเรชาเต นว, นปจฺฉาชาเต นว, น-อาเสวเน นว, นกมฺเม ตีณิ, นวิปาเก นว, นมคฺเค เอกํ, นวิปฺปยุตฺเต นว. (สํขิตฺตํ.)}

\csromanheader{2}{เหตุปจฺจยา น-อธิปติยา นว. (สํขิตฺตํ.)}
\csromanheader{2}{นเหตุปจฺจยา อารมฺมเณ เทฺว. (สํขิตฺตํ.)}
\prose{(สหชาตวาโรปิ ปจฺจยวาโรปิ นิสฺสยวาโรปิ สํสฏฺฐวาโรปิ สมฺปยุตฺตวาโรปิ ปฏิจฺจวารสทิสา วิตฺถาเรตพฺพา.)}

\csromanheader{2}{2. อนาคตารมฺมณปท 7. ปญฺหาวาร}
\csromanheader{2}{\textbf{ปจฺจยจตุกฺก-เหตฺวาทิ}}
\proseitem{497}{เหตุ อนาคตารมฺมโณ ธมฺโม เหตุสฺส อนาคตารมฺมณสฺส ธมฺมสฺส เหตุปจฺจเยน ปจฺจโย.\csromanspacer{} ตีณิ.}

\prose{เหตุ อนาคตารมฺมโณ ธมฺโม เหตุสฺส อนาคตารมฺมณสฺส ธมฺมสฺส อารมฺมณปจฺจเยน ปจฺจโย.\csromanspacer{} นว.}

\csromanheader{2}{1. เหตุทุก 19. อตีตารมฺมณตฺติก}
\prose{\csromanfolio{311}เหตุ อนาคตารมฺมโณ ธมฺโม เหตุสฺส อนาคตารมฺมณสฺส ธมฺมสฺส อธิปติปจฺจเยน ปจฺจโย.\csromanspacer{} อารมฺมณาธิปติ, สหชาตาธิปติ.\csromanspacer{} ตีณิ.}

\prose{นเหตุํ อนาคตารมฺมโณ ธมฺโม นเหตุสฺส อนาคตารมฺมณสฺส ธมฺมสฺส อธิปติปจฺจเยน ปจฺจโย.\csromanspacer{} อารมฺมณาธิปติ, สหชาตาธิปติ.\csromanspacer{} ตีณิ.}

\prose{เหตุ อนาคตารมฺมโณ จ นเหตุ อนาคตารมฺมโณ จ ธมฺมา เหตุสฺส อนาคตารมฺมณสฺส ธมฺมสฺส อธิปติปจฺจเยน ปจฺจโย.\csromanspacer{} อารมฺมณาธิปติ.\csromanspacer{} ตีณิ. (สํขิตฺตํ.)}

\proseitem{498}{เหตุยา ตีณิ, อารมฺมเณ นว, อธิปติยา นว, อนนฺตเร นว, สมนนฺตเร นว, สหชาเต นว, อญฺญมญฺเญ นว, นิสฺสเย นว, อุปนิสฺสเย นว, อาเสวเน นว, กมฺเม ตีณิ, วิปาเก นว, อาหาเร ตีณิ ฯเปฯ อวิคเต นว. (สํขิตฺตํ.)}

\csromanheader{2}{\textbf{ปจฺจนียุทฺธาร}}
\proseitem{499}{เหตุ อนาคตารมฺมโณ ธมฺโม เหตุสฺส อนาคตารมฺมณสฺส ธมฺมสฺส อารมฺมณปจฺจเยน ปจฺจโย.\csromanspacer{} สหชาตปจฺจเยน ปจฺจโย.\csromanspacer{} อุปนิสฺสยปจฺจเยน ปจฺจโย. (สํขิตฺตํ.)}

\proseitem{500}{นเหตุยา นว, น-อารมฺมเณ นว. (สํขิตฺตํ.)}

\csromanheader{2}{เหตุปจฺจยา น-อารมฺมเณ ตีณิ. (สํขิตฺตํ.)}
\csromanheader{2}{นเหตุปจฺจยา อารมฺมเณ นว. (สํขิตฺตํ.)}
\prose{(ยถา กุสลตฺติเก ปญฺหาวารสฺส อนุโลมมฺปิ ปจฺจนียมฺปิ อนุโลมปจฺจนียมฺปิ ปจฺจนียานุโลมมฺปิ คณิตํ, เอวํ คเณตพฺพํ.)}

\csromanheader{2}{ธมฺมานุโลม ทุกติกปฏฺฐานปาฬิ \textbf{3. ปจฺจุปฺปนฺนารมฺมณปท 1. ปฏิจฺจวาร}}
\csromanheader{2}{\textbf{ปจฺจยจตุกฺก-เหตุ}}
\proseitem{501}{\csromanfolio{312}เหตุํ ปจฺจุปฺปนฺนารมฺมณํ ธมฺมํ ปฏิจฺจ เหตุ ปจฺจุปฺปนฺนารมฺมโณ ธมฺโม อุปฺปชฺชติ เหตุปจฺจยา.\csromanspacer{} ตีณิ.}

\prose{นเหตุํ ปจฺจุปฺปนฺนารมฺมณํ ธมฺมํ ปฏิจฺจ นเหตุ ปจฺจุปฺปนฺนารมฺมโณ ธมฺโม อุปฺปชฺชติ เหตุปจฺจยา.\csromanspacer{} ตีณิ.}

\prose{เหตุํ ปจฺจุปฺปนฺนารมฺมณญฺจ นเหตุํ ปจฺจุปฺปนฺนารมฺมณญฺจ ธมฺมํ ปฏิจฺจ เหตุ ปจฺจุปฺปนฺนารมฺมโณ ธมฺโม อุปฺปชฺชติ เหตุปจฺจยา.\csromanspacer{} ตีณิ. (สํขิตฺตํ.)}

\proseitem{502}{เหตุยา นว, อารมฺมเณ นว, อธิปติยา นว ฯเปฯ กมฺเม นว ฯเปฯ อวิคเต นว. (สํขิตฺตํ.)}

\csromanheader{2}{\textbf{นเหตุ น-อธิปติ}}
\proseitem{503}{นเหตุํ ปจฺจุปฺปนฺนารมฺมณํ ธมฺมํ ปฏิจฺจ นเหตุ ปจฺจุปฺปนฺนารมฺมโณ ธมฺโม อุปฺปชฺชติ นเหตุปจฺจยา.\csromanspacer{}\csromanspacer{}นเหตุํ ปจฺจุปฺปนฺนารมฺมณํ ธมฺมํ ปฏิจฺจ เหตุ ปจฺจุปฺปนฺนารมฺมโณ ธมฺโม อุปฺปชฺชติ นเหตุปจฺจยา. (2)}

\prose{เหตุํ ปจฺจุปฺปนฺนารมฺมณํ ธมฺมํ ปฏิจฺจ เหตุ ปจฺจุปฺปนฺนารมฺมโณ ธมฺโม อุปฺปชฺชติ น-อธิปติปจฺจยา. (สํขิตฺตํ.)}

\proseitem{504}{นเหตุยา เทฺว, น-อธิปติยา นว, นปุเรชาเต นว, นปจฺฉาชาเต นว, น-อาเสวเน นว, นกมฺเม ตีณิ, นวิปาเก นว, นฌาเน เอกํ, นมคฺเค เอกํ, นวิปฺปยุตฺเต นว. (สํขิตฺตํ.)}

\csromanheader{2}{เหตุปจฺจยา น-อธิปติยา นว. (สํขิตฺตํ.)}
\csromanheader{2}{นเหตุปจฺจยา อารมฺมเณ เทฺว. (สํขิตฺตํ.)}
\prose{(สหชาตวาโรปิ ปจฺจยวาโรปิ นิสฺสยวาโรปิ สํสฏฺฐวาโรปิ สมฺปยุตฺตวาโรปิ ปฏิจฺจวารสทิสา วิตฺถาเรตพฺพา.)}

\vspace{\tipitakaparskip}
\csromancenter{เหตุทุก 19.\csromanspacer{} อตีตารมฺมณตฺติก \textbf{3.}\csromanspacer{}\textbf{ ปจฺจุปฺปนฺนารมฺมณปท 7.}\csromanspacer{}\textbf{ ปญฺหาวาร}}
\csromanheader{2}{\textbf{ปจฺจยจตุกฺก-เหตฺวาทิ}}
\proseitem{505}{\csromanfolio{313}เหตุ ปจฺจุปฺปนฺนารมฺมโณ ธมฺโม เหตุสฺส ปจฺจุปฺปนฺนารมฺมณสฺส ธมฺมสฺส เหตุปจฺจเยน ปจฺจโย.\csromanspacer{} ตีณิ.}

\prose{เหตุ ปจฺจุปฺปนฺนารมฺมโณ ธมฺโม เหตุสฺส ปจฺจุปฺปนฺนารมฺมณสฺส ธมฺมสฺส อารมฺมณปจฺจเยน ปจฺจโย.\csromanspacer{} นว.}

\prose{เหตุ ปจฺจุปฺปนฺนารมฺมโณ ธมฺโม เหตุสฺส ปจฺจุปฺปนฺนารมฺมณสฺส ธมฺมสฺส อธิปติปจฺจเยน ปจฺจโย.\csromanspacer{} อารมฺมณาธิปติ, สหชาตาธิปติ.\csromanspacer{} ตีณิ.}

\prose{นเหตุ ปจฺจุปฺปนฺนารมฺมโณ ธมฺโม นเหตุสฺส ปจฺจุปฺปนฺนารมฺมณสฺส ธมฺมสฺส อธิปติปจฺจเยน ปจฺจโย.\csromanspacer{} อารมฺมณาธิปติ, สหชาตาธิปติ.\csromanspacer{} ตีณิ.}

\prose{เหตุ ปจฺจุปฺปนฺนารมฺมโณ จ นเหตุ ปจฺจุปฺปนฺนารมฺมโณ จ ธมฺมา เหตุสฺส ปจฺจุปฺปนฺนารมฺมณสฺส ธมฺมสฺส อธิปติปจฺจเยน ปจฺจโย.\csromanspacer{} อารมฺมณาธิปติ.\csromanspacer{} ตีณิ. (สํขิตฺตํ.)}

\proseitem{506}{เหตุยา ตีณิ, อารมฺมเณ นว, อธิปติยา นว, อนนฺตเร นว, สมนนฺตเร นว, สหชาเต นว, อญฺญมญฺเญ นว, นิสฺสเย นว, อุปนิสฺสเย นว, อาเสวเน นว, กมฺเม ตีณิ, วิปาเก นว, อาหาเร ตีณิ, อินฺทฺริเย นว ฯเปฯ อวิคเต นว. (สํขิตฺตํ.)}

\csromanheader{2}{\textbf{ปจฺจนียุทฺธาร}}
\proseitem{507}{เหตุ ปจฺจุปฺปนฺนารมฺมโณ ธมฺโม เหตุสฺส ปจฺจุปฺปนฺนารมฺมณสฺส ธมฺมสฺส อารมฺมณปจฺจเยน ปจฺจโย.\csromanspacer{} สหชาตปจฺจเยน ปจฺจโย.\csromanspacer{} อุปนิสฺสยปจฺจเยน ปจฺจโย. (สํขิตฺตํ.)}

\proseitem{508}{นเหตุยา นว, น-อารมฺมเณ นว. (สํขิตฺตํ.)}

\csromanheader{2}{เหตุปจฺจยา น-อารมฺมเณ ตีณิ. (สํขิตฺตํ.)}
\csromanheader{2}{นเหตุปจฺจยา อารมฺมเณ นว. (สํขิตฺตํ.)}
\gambhira{ธมฺมานุโลม ทุกติกปฏฺฐานปาฬิ}
\prose{\csromanfolio{314}(ยถา กุสลตฺติเก ปญฺหาวารสฺส อนุโลมมฺปิ ปจฺจนียมฺปิ อนุโลมปจฺจนียมฺปิ ปจฺจนียานุโลมมฺปิ คณิตํ, เอวํ คเณตพฺพํ.)}

\nitthitam{เหตุทุก-อตีตารมฺมณตฺติกํ นิฏฺฐิตํ.}
\csromanmatikaanchor{mtk.51}
\csromantocmarkref{section}{20. อชฺฌตฺตตฺติก}{mtk.51}
\bookmark[dest={mtk.51},level=1]{20. อชฺฌตฺตตฺติก}
\csromanheader{1}{1. เหตุทุก 20. อชฺฌตฺตตฺติก}
\csromanheader{2}{1. อชฺฌตฺตปท 1. ปฏิจฺจวาร}
\csromanheader{2}{\textbf{ปจฺจยจตุกฺก-เหตุ}}
\proseitem{509}{เหตุํ อชฺฌตฺตํ ธมฺมํ ปฏิจฺจ เหตุ อชฺฌตฺโต ธมฺโม อุปฺปชฺชติ เหตุปจฺจยา.\csromanspacer{} ตีณิ.}

\prose{นเหตุํ อชฺฌตฺตํ ธมฺมํ ปฏิจฺจ นเหตุ อชฺฌตฺโต ธมฺโม อุปฺปชฺชติ เหตุปจฺจยา.\csromanspacer{} ตีณิ.}

\prose{เหตุํ อชฺฌตฺตญฺจ นเหตุํ อชฺฌตฺตญฺจ ธมฺมํ ปฏิจฺจ เหตุ อชฺฌตฺโต ธมฺโม อุปฺปชฺชติ เหตุปจฺจยา.\csromanspacer{} ตีณิ. (สํขิตฺตํ.)}

\proseitem{510}{เหตุยา นว, อารมฺมเณ นว, อธิปติยา นว ฯเปฯ กมฺเม นว, วิปาเก นว, อาหาเร นว ฯเปฯ อวิคเต นว. (สํขิตฺตํ.)}

\csromanheader{2}{\textbf{นเหตฺวาทิ}}
\proseitem{511}{นเหตุํ อชฺฌตฺตํ ธมฺมํ ปฏิจฺจ นเหตุ อชฺฌตฺโต ธมฺโม อุปฺปชฺชติ นเหตุปจฺจยา.\csromanspacer{}\csromanspacer{}นเหตุํ อชฺฌตฺตํ ธมฺมํ ปฏิจฺจ เหตุ อชฺฌตฺโต ธมฺโม อุปฺปชฺชติ นเหตุปจฺจยา. (2)}

\prose{เหตุํ อชฺฌตฺตํ ธมฺมํ ปฏิจฺจ นเหตุ อชฺฌตฺโต ธมฺโม อุปฺปชฺชติ น-อารมฺมณปจฺจยา. (1)}

\prose{นเหตุํ อชฺฌตฺตํ ธมฺมํ ปฏิจฺจ นเหตุ อชฺฌตฺโต ธมฺโม อุปฺปชฺชติ น-อารมฺมณปจฺจยา. (1)}

\prose{เหตุํ อชฺฌตฺตญฺจ นเหตุํ อชฺฌตฺตญฺจ ธมฺมํ ปฏิจฺจ นเหตุ อชฺฌตฺโต ธมฺโม อุปฺปชฺชติ น-อารมฺมณปจฺจยา. (1)}

\csromanheader{2}{1. เหตุทุก 20. อชฺฌตฺตตฺติก}
\prose{\csromanfolio{315}เหตุํ อชฺฌตฺตํ ธมฺมํ ปฏิจฺจ เหตุ อชฺฌตฺโต ธมฺโม อุปฺปชฺชติ น-อธิปติปจฺจยา. (สํขิตฺตํ.)}

\proseitem{512}{นเหตุยา เทฺว, น-อารมฺมเณ ตีณิ, น-อธิปติยา นว, น-อนนฺตเร ตีณิ, นสมนนฺตเร ตีณิ, น-อญฺญมญฺเญ ตีณิ, น-อุปนิสฺสเย ตีณิ, นปุเรชาเต นว, นปจฺฉาชาเต นว, น-อาเสวเน นว, นกมฺเม ตีณิ, นวิปาเก นว, น-อาหาเร เอกํ, น-อินฺทฺริเย เอกํ, นฌาเน เอกํ, นมคฺเค เอกํ, นสมฺปยุตฺเต ตีณิ, นวิปฺปยุตฺเต นว, โนนตฺถิยา ตีณิ, โนวิคเต ตีณิ. (สํขิตฺตํ.)}

\csromanheader{2}{เหตุปจฺจยา น-อารมฺมเณ ตีณิ. (สํขิตฺตํ.)}
\csromanheader{2}{นเหตุปจฺจยา อารมฺมเณ เทฺว. (สํขิตฺตํ.)}
\prose{(สหชาตวาโรปิ ปจฺจยวาโรปิ นิสฺสยวาโรปิ สํสฏฺฐวาโรปิ สมฺปยุตฺตวาโรปิ ปฏิจฺจวารสทิสา วิตฺถาเรตพฺพา.)}

\csromanheader{2}{1. อชฺฌตฺตปท 7. ปญฺหาวาร}
\csromanheader{2}{\textbf{ปจฺจยจตุกฺก-เหตฺวาทิ}}
\proseitem{513}{เหตุ อชฺฌตฺโต ธมฺโม เหตุสฺส อชฺฌตฺตสฺส ธมฺมสฺส เหตุปจฺจเยน ปจฺจโย.\csromanspacer{} ตีณิ.}

\prose{เหตุ อชฺฌตฺโต ธมฺโม เหตุสฺส อชฺฌตฺตสฺส ธมฺมสฺส อารมฺมณปจฺจเยน ปจฺจโย.\csromanspacer{} นว.}

\prose{เหตุ อชฺฌตฺโต ธมฺโม เหตุสฺส อชฺฌตฺตสฺส ธมฺมสฺส อธิปติปจฺจเยน ปจฺจโย.\csromanspacer{} อารมฺมณาธิปติ, สหชาตาธิปติ.\csromanspacer{} ตีณิ.}

\prose{นเหตุ อชฺฌตฺโต ธมฺโม นเหตุสฺส อชฺฌตฺตสฺส ธมฺมสฺส อธิปติปจฺจเยน ปจฺจโย.\csromanspacer{} อารมฺมณาธิปติ, สหชาตาธิปติ.\csromanspacer{} ตีณิ.}

\prose{เหตุ อชฺฌตฺโต จ นเหตุ อชฺฌตฺโต จ ธมฺมา เหตุสฺส อชฺฌตฺตสฺส ธมฺมสฺส อธิปติปจฺจเยน ปจฺจโย.\csromanspacer{} อารมฺมณาธิปติ.\csromanspacer{} ตีณิ. (สํขิตฺตํ.)}

\gambhira{ธมฺมานุโลม ทุกติกปฏฺฐานปาฬิ}
\proseitem{514}{\csromanfolio{316}เหตุยา ตีณิ, อารมฺมเณ นว, อธิปติยา นว, อนนฺตเร นว, สมนนฺตเร นว, สหชาเต นว, อญฺญมญฺเญ นว, นิสฺสเย นว, อุปนิสฺสเย นว, ปุเรชาเต ตีณิ, ปจฺฉาชาเต ตีณิ, อาเสวเน นว, กมฺเม ตีณิ, วิปาเก นว, อาหาเร ตีณิ, อินฺทฺริเย นว, ฌาเน ตีณิ, มคฺเค นว, สมฺปยุตฺเต นว, วิปฺปยุตฺเต ปญฺจ, อตฺถิยา นว, นตฺถิยา นว, วิคเต นว, อวิคเต นว. (สํขิตฺตํ.)}

\csromanheader{2}{\textbf{ปจฺจนียุทฺธาร}}
\proseitem{515}{เหตุ อชฺฌตฺโต ธมฺโม เหตุสฺส อชฺฌตฺตสฺส ธมฺมสฺส อารมฺมณปจฺจเยน ปจฺจโย.\csromanspacer{} สหชาตปจฺจเยน ปจฺจโย.\csromanspacer{} อุปนิสฺสยปจฺจเยน ปจฺจโย. (สํขิตฺตํ.)}

\proseitem{516}{\textbf{นเหตุ} นว, น-อารมฺมเณ นว. (สํขิตฺตํ.)}

\csromanheader{2}{เหตุปจฺจยา น-อารมฺมเณ ตีณิ. (สํขิตฺตํ.)}
\csromanheader{2}{\textbf{นเหตุ}ปจฺจยา อารมฺมเณ นว. (สํขิตฺตํ.)}
\prose{(ยถา กุสลตฺติเก ปญฺหาวารสฺส อนุโลมมฺปิ ปจฺจนียมฺปิ อนุโลมปจฺจนียมฺปิ ปจฺจนียานุโลมมฺปิ คณิตํ, เอวํ คเณตพฺพํ.)}

\csromanheader{2}{2. พหิทฺธาปท 1. ปฏิจฺจวาร}
\csromanheader{2}{\textbf{ปจฺจยจตุกฺก-เหตุ}}
\proseitem{517}{เหตุํ พหิทฺธา ธมฺมํ ปฏิจฺจ เหตุ พหิทฺธา ธมฺโม อุปฺปชฺชติ เหตุปจฺจยา. (สํขิตฺตํ.)}

\proseitem{518}{เหตุยา นว, อารมฺมเณ นว ฯเปฯ กมฺเม นว, วิปาเก นว ฯเปฯ อวิคเต นว. (สํขิตฺตํ.)}

\csromanheader{2}{\textbf{นเหตุ}}
\vspace{\tipitakaparskip}
\csromancenter{\textbf{นเหตุ}ํ พหิทฺธา ธมฺมํ ปฏิจฺจ นเหตุ พหิทฺธา ธมฺโม อุปฺปชฺชติ นเหตุปจฺจยา.\csromanspacer{}\csromanspacer{}\textbf{นเหตุ}ํ พหิทฺธา ธมฺมํ ปฏิจฺจ เหตุ พหิทฺธา ธมฺโม อุปฺปชฺชติ นเหตุปจฺจยา. (2) (สํขิตฺตํ.)}
\csromanheader{2}{1. เหตุทุก 20. อชฺฌตฺตตฺติก}
\proseitem{520}{\csromanfolio{317}นเหตุยา เทฺว, น-อารมฺมเณ ตีณิ, น-อธิปติยา นว, น-อนนฺตเร ตีณิ, นสมนนฺตเร ตีณิ, น-อญฺญมญฺเญ ตีณิ, น-อุปนิสฺสเย ตีณิ, นปุเรชาเต นว, นปจฺฉาชาเต นว, น-อาเสวเน นว, นกมฺเม ตีณิ, นวิปาเก นว, น-อาหาเร เอกํ, น-อินฺทฺริเย เอกํ, นฌาเน เอกํ, นมคฺเค เอกํ, นสมฺปยุตฺเต ตีณิ, นวิปฺปยุตฺเต นว, โนนตฺถิยา ตีณิ, โนวิคเต ตีณิ. (สํขิตฺตํ.)}

\csromanheader{2}{เหตุปจฺจยา น-อารมฺมเณ ตีณิ. (สํขิตฺตํ.)}
\csromanheader{2}{นเหตุปจฺจยา อารมฺมเณ เทฺว. (สํขิตฺตํ.)}
\prose{(สหชาตวาโรปิ ปจฺจยวาโรปิ นิสฺสยวาโรปิ สํสฏฺฐวาโรปิ สมฺปยุตฺตวาโรปิ ปฏิจฺจวารสทิสา วิตฺถาเรตพฺพา.)}

\csromanheader{2}{2. พหิทฺธาปท 7. ปญฺหาวาร}
\csromanheader{2}{\textbf{ปจฺจยจตุกฺก-เหตฺวาทิ}}
\proseitem{521}{เหตุ พหิทฺธา ธมฺโม เหตุสฺส พหิทฺธา ธมฺมสฺส เหตุปจฺจเยน ปจฺจโย.\csromanspacer{} ตีณิ.}

\prose{เหตุ พหิทฺธา ธมฺโม เหตุสฺส พหิทฺธา ธมฺมสฺส อารมฺมณปจฺจเยน ปจฺจโย.\csromanspacer{} นว.}

\prose{เหตุ พหิทฺธา ธมฺโม เหตุสฺส พหิทฺธา ธมฺมสฺส อธิปติปจฺจเยน ปจฺจโย.\csromanspacer{} อารมฺมณาธิปติ, สหชาตาธิปติ.\csromanspacer{} ตีณิ.}

\prose{นเหตุ พหิทฺธา ธมฺโม นเหตุสฺส พหิทฺธา ธมฺมสฺส อธิปติปจฺจเยน ปจฺจโย.\csromanspacer{} อารมฺมณาธิปติ, สหชาตาธิปติ.\csromanspacer{} ตีณิ.}

\prose{เหตุ พหิทฺธา จ นเหตุ พหิทฺธา จ ธมฺมา เหตุสฺส พหิทฺธา ธมฺมสฺส อธิปติปจฺจเยน ปจฺจโย.\csromanspacer{} อารมฺมณาธิปติ.\csromanspacer{} ตีณิ. (สํขิตฺตํ.)}

\proseitem{522}{เหตุยา ตีณิ, อารมฺมเณ นว, อธิปติยา นว ฯเปฯ ปุเรชาเต ตีณิ, อาเสวเน นว, กมฺเม ตีณิ, วิปาเก นว ฯเปฯ อวิคเต นว. (สํขิตฺตํ.)}

\csromanheader{2}{ธมฺมานุโลม ทุกติกปฏฺฐานปาฬิ \textbf{ปจฺจนียุทฺธาร}}
\proseitem{523}{\csromanfolio{318}เหตุ พหิทฺธา ธมฺโม เหตุสฺส พหิทฺธา ธมฺมสฺส อารมฺมณปจฺจเยน ปจฺจโย.\csromanspacer{} สหชาตปจฺจเยน ปจฺจโย.\csromanspacer{} อุปนิสฺสยปจฺจเยน ปจฺจโย. (สํขิตฺตํ.)}

\proseitem{524}{นเหตุยา นว, น-อารมฺมเณ นว. (สํขิตฺตํ.)}

\csromanheader{2}{เหตุปจฺจยา น-อารมฺมเณ ตีณิ. (สํขิตฺตํ.)}
\csromanheader{2}{นเหตุปจฺจยา อารมฺมเณ นว. (สํขิตฺตํ.)}
\prose{(ยถา กุสลตฺติเก ปญฺหาวารสฺส อนุโลมมฺปิ ปจฺจนียมฺปิ อนุโลมปจฺจนียมฺปิ ปจฺจนียานุโลมมฺปิ คณิตํ, เอวํ คเณตพฺพํ.\csromanspacer{} อชฺฌตฺตพหิทฺธา น ลพฺภนฺติ.)}

\nitthitam{เหตุทุก-อชฺฌตฺตตฺติกํ นิฏฺฐิตํ.}
\csromanmatikaanchor{mtk.52}
\csromantocmarkref{section}{21. อชฺฌตฺตารมฺมณตฺติก}{mtk.52}
\bookmark[dest={mtk.52},level=1]{21. อชฺฌตฺตารมฺมณตฺติก}
\csromanheader{1}{1. เหตุทุก 21. อชฺฌตฺตารมฺมณตฺติก}
\csromanheader{2}{1. อชฺฌตฺตารมฺมณปท 1. ปฏิจฺจวาร}
\csromanheader{2}{\textbf{ปจฺจยจตุกฺก-เหตุ}}
\proseitem{525}{เหตุํ อชฺฌตฺตารมฺมณํ ธมฺมํ ปฏิจฺจ เหตุ อชฺฌตฺตารมฺมโณ ธมฺโม อุปฺปชฺชติ เหตุปจฺจยา.\csromanspacer{} ตีณิ.}

\prose{นเหตุํ อชฺฌตฺตารมฺมณํ ธมฺมํ ปฏิจฺจ นเหตุ อชฺฌตฺตารมฺมโณ ธมฺโม อุปฺปชฺชติ เหตุปจฺจยา.\csromanspacer{} ตีณิ.}

\prose{เหตุํ อชฺฌตฺตารมฺมณญฺจ นเหตุํ อชฺฌตฺตารมฺมณญฺจ ธมฺมํ ปฏิจฺจ เหตุ อชฺฌตฺตารมฺมโณ ธมฺโม อุปฺปชฺชติ เหตุปจฺจยา.\csromanspacer{} ตีณิ. (สํขิตฺตํ.)}

\proseitem{526}{เหตุยา นว, อารมฺมเณ นว, อธิปติยา นว ฯเปฯ อวิคเต นว. (สํขิตฺตํ.)}

\csromanheader{2}{1. เหตุทุก 21. อชฺฌตฺตารมฺมณตฺติก \textbf{นเหตุ น-อธิปติ}}
\proseitem{527}{\csromanfolio{319}นเหตุํ อชฺฌตฺตารมฺมณํ ธมฺมํ ปฏิจฺจ นเหตุ อชฺฌตฺตารมฺมโณ ธมฺโม อุปฺปชฺชติ นเหตุปจฺจยา.\csromanspacer{}\csromanspacer{}นเหตุํ อชฺฌตฺตารมฺมณํ ธมฺมํ ปฏิจฺจ เหตุ อชฺฌตฺตารมฺมโณ ธมฺโม อุปฺปชฺชติ นเหตุปจฺจยา. (2)}

\prose{เหตุํ อชฺฌตฺตารมฺมณํ ธมฺมํ ปฏิจฺจ เหตุ อชฺฌตฺตารมฺมโณ ธมฺโม อุปฺปชฺชติ น-อธิปติปจฺจยา. (สํขิตฺตํ.)}

\proseitem{528}{นเหตุยา เทฺว, น-อธิปติยา นว, นปุเรชาเต นว, นปจฺฉาชาเต นว, น-อาเสวเน นว, นกมฺเม ตีณิ, นวิปาเก นว, นฌาเน เอกํ, นมคฺเค เอกํ, นวิปฺปยุตฺเต นว. (สํขิตฺตํ.)}

\csromanheader{2}{เหตุปจฺจยา น-อธิปติยา นว. (สํขิตฺตํ.)}
\csromanheader{2}{นเหตุปจฺจยา อารมฺมเณ เทฺว. (สํขิตฺตํ.)}
\prose{(สหชาตวาโรปิ ปจฺจยวาโรปิ นิสฺสยวาโรปิ สํสฏฺฐวาโรปิ สมฺปยุตฺตวาโรปิ ปฏิจฺจวารสทิสา วิตฺถาเรตพฺพา.)}

\csromanheader{2}{1. อชฺฌตฺตารมฺมณปท 7. ปญฺหาวาร}
\csromanheader{2}{\textbf{ปจฺจยจตุกฺก-เหตุ}}
\proseitem{529}{เหตุ อชฺฌตฺตารมฺมโณ ธมฺโม เหตุสฺส อชฺฌตฺตารมฺมณสฺส ธมฺมสฺส เหตุปจฺจเยน ปจฺจโย. (สํขิตฺตํ.)}

\proseitem{530}{เหตุยา ตีณิ, อารมฺมเณ นว, อธิปติยา นว ฯเปฯ อวิคเต นว. (สํขิตฺตํ.)}

\csromanheader{2}{\textbf{ปจฺจนียุทฺธาร}}
\proseitem{531}{เหตุ อชฺฌตฺตารมฺมโณ ธมฺโม เหตุสฺส อชฺฌตฺตารมฺมณสฺส ธมฺมสฺส อารมฺมณปจฺจเยน ปจฺจโย.\csromanspacer{} สหชาตปจฺจเยน ปจฺจโย.\csromanspacer{} อุปนิสฺสยปจฺจเยน ปจฺจโย. (สํขิตฺตํ.)}

\gambhira{ธมฺมานุโลม ทุกติกปฏฺฐานปาฬิ}
\proseitem{532}{\csromanfolio{320}\textbf{นเหตุ}ยา นว, น-อารมฺมเณ นว. (สํขิตฺตํ.)}

\csromanheader{2}{เหตุปจฺจยา น-อารมฺมเณ ตีณิ. (สํขิตฺตํ.)}
\csromanheader{2}{\textbf{นเหตุ}ปจฺจยา อารมฺมเณ นว. (สํขิตฺตํ.)}
\prose{(ยถา กุสลตฺติเก ปญฺหาวารสฺส อนุโลมมฺปิ ปจฺจนียมฺปิ อนุโลมปจฺจนียมฺปิ ปจฺจนียานุโลมมฺปิ คณิตํ, เอวํ คเณตพฺพํ.)}

\csromanheader{2}{2. พหิทฺธารมฺมณปท 1. ปฏิจฺจวาร}
\csromanheader{2}{\textbf{ปจฺจยจตุกฺก-เหตุ}}
\proseitem{533}{เหตุํ พหิทฺธารมฺมณํ ธมฺมํ ปฏิจฺจ เหตุ พหิทฺธารมฺมโณ ธมฺโม อุปฺปชฺชติ เหตุปจฺจยา. (สํขิตฺตํ.)}

\proseitem{534}{เหตุยา นว, อารมฺมเณ นว ฯเปฯ อวิคเต นว. (สํขิตฺตํ.)}

\csromanheader{2}{\textbf{นเหตุ}}
\vspace{\tipitakaparskip}
\csromancenter{\textbf{นเหตุ}ํ พหิทฺธารมฺมณํ ธมฺมํ ปฏิจฺจ นเหตุ พหิทฺธารมฺมโณ ธมฺโม อุปฺปชฺชติ นเหตุปจฺจยา. (สํขิตฺตํ.)}
\csromancenter{\textbf{นเหตุ}ยา เทฺว, น-อธิปติยา นว ฯเปฯ นปุเรชาเต นว, นปจฺฉาชาเต นว, น-อาเสวเน นว, นกมฺเม ตีณิ, นวิปาเก นว, นฌาเน เอกํ, นมคฺเค เอกํ, นวิปฺปยุตฺเต นว. (สํขิตฺตํ.)}
\csromanheader{2}{เหตุปจฺจยา น-อธิปติยา นว. (สํขิตฺตํ.)}
\csromanheader{2}{\textbf{นเหตุ}ปจฺจยา อารมฺมเณ เทฺว. (สํขิตฺตํ.)}
\prose{(สหชาตวาโรปิ ปจฺจยวาโรปิ นิสฺสยวาโรปิ สํสฏฺฐวาโรปิ สมฺปยุตฺตวาโรปิ ปฏิจฺจวารสทิสา วิตฺถาเรตพฺพา.)}

\csromanheader{2}{2. พหิทฺธารมฺมณปท 7. ปญฺหาวาร}
\csromanheader{2}{\textbf{ปจฺจยจตุกฺก-เหตฺวาทิ}}
\proseitem{537}{เหตุ พหิทฺธารมฺมโณ ธมฺโม เหตุสฺส พหิทฺธารมฺมณสฺส ธมฺมสฺส เหตุปจฺจเยน ปจฺจโย.\csromanspacer{} ตีณิ.}

\csromanmatikaanchor{mtk.53}
\csromantocmarkref{section}{22. สนิทสฺสนสปฺปฏิฆตฺติก}{mtk.53}
\bookmark[dest={mtk.53},level=1]{22. สนิทสฺสนสปฺปฏิฆตฺติก}
\csromanheader{1}{1. เหตุทุก 22. สนิทสฺสนสปฺปฏิฆตฺติก}
\prose{\csromanfolio{321}เหตุ พหิทฺธารมฺมโณ ธมฺโม เหตุสฺส พหิทฺธารมฺมณสฺส ธมฺมสฺส อารมฺมณปจฺจเยน ปจฺจโย.\csromanspacer{} นว.}

\prose{เหตุ พหิทฺธารมฺมโณ ธมฺโม เหตุสฺส พหิทฺธารมฺมณสฺส ธมฺมสฺส อธิปติปจฺจเยน ปจฺจโย.\csromanspacer{} สหชาตาธิปติ.\csromanspacer{} ตีณิ.}

\prose{นเหตุ พหิทฺธารมฺมโณ ธมฺโม นเหตุสฺส พหิทฺธารมฺมณสฺส ธมฺมสฺส อธิปติปจฺจเยน ปจฺจโย.\csromanspacer{} สหชาตาธิปติ.\csromanspacer{} ตีณิ. (สํขิตฺตํ.)}

\proseitem{538}{เหตุยา ตีณิ, อารมฺมเณ นว, อธิปติยา ฉ, กมฺเม ตีณิ, วิปาเก นว, อาหาเร ตีณิ ฯเปฯ อวิคเต นว. (สํขิตฺตํ.)}

\csromanheader{2}{\textbf{ปจฺจนียุทฺธาร}}
\proseitem{539}{เหตุ พหิทฺธารมฺมโณ ธมฺโม เหตุสฺส พหิทฺธารมฺมณสฺส ธมฺมสฺส อารมฺมณปจฺจเยน ปจฺจโย.\csromanspacer{} สหชาตปจฺจเยน ปจฺจโย.\csromanspacer{} อุปนิสฺสยปจฺจเยน ปจฺจโย. (สํขิตฺตํ.)}

\proseitem{540}{นเหตุยา นว, น-อารมฺมเณ นว. (สํขิตฺตํ.)}

\csromanheader{2}{เหตุปจฺจยา น-อารมฺมเณ ตีณิ. (สํขิตฺตํ.)}
\csromanheader{2}{นเหตุปจฺจยา อารมฺมเณ นว. (สํขิตฺตํ.)}
\prose{(ยถา กุสลตฺติเก ปญฺหาวารสฺส อนุโลมมฺปิ ปจฺจนียมฺปิ อนุโลมปจฺจนียมฺปิ ปจฺจนียานุโลมมฺปิ คณิตํ, เอวํ คเณตพฺพํ.)}

\nitthitam{เหตุทุก-อชฺฌตฺตารมฺมณตฺติกํ นิฏฺฐิตํ.}
\csromanheader{2}{1. เหตุทุก 22. สนิทสฺสนสปฺปฏิฆตฺติก}
\csromanheader{2}{1. อนิทสฺสนสปฺปฏิฆปท 1. ปฏิจฺจวาร}
\csromanheader{2}{\textbf{ปจฺจยจตุกฺก-เหตฺวาทิ}}
\proseitem{541}{นเหตุํ อนิทสฺสนสปฺปฏิฆํ ธมฺมํ ปฏิจฺจ นเหตุ อนิทสฺสนสปฺปฏิโฆ ธมฺโม อุปฺปชฺชติ เหตุปจฺจยา.\csromanspacer{} อธิปติปจฺจยา.\csromanspacer{} สหชาตปจฺจยา.}

\vspace{\tipitakaparskip}
\csromancenter{ธมฺมานุโลม ทุกติกปฏฺฐานปาฬิ อญฺญมญฺญปจฺจยา.\csromanspacer{} นิสฺสยปจฺจยา.\csromanspacer{} กมฺมปจฺจยา.\csromanspacer{} วิปากปจฺจยา.\csromanspacer{} อาหารปจฺจยา.\csromanspacer{} อินฺทฺริยปจฺจยา.\csromanspacer{} ฌานปจฺจยา.\csromanspacer{} มคฺคปจฺจยา.\csromanspacer{} วิปฺปยุตฺตปจฺจยา.\csromanspacer{} อตฺถิปจฺจยา.\csromanspacer{} อวิคตปจฺจยา.}
\csromanheader{2}{\textbf{สุทฺธ}}
\proseitem{542}{\csromanfolio{322}เหตุยา เอกํ, อธิปติยา เอกํ, สหชาเต เอกํ, อญฺญมญฺเญ เอกํ, นิสฺสเย เอกํ, กมฺเม เอกํ, วิปาเก เอกํ, มคฺเค เอกํ, วิปฺปยุตฺเต เอกํ, อตฺถิยา เอกํ, อวิคเต เอกํ. (สํขิตฺตํ.)}

\proseitem{543}{นเหตุยา เอกํ, น-อารมฺมเณ เอกํ, น-อธิปติยา เอกํ, (สพฺเพ ปจฺจยา กาตพฺพา.) ฯเปฯ โนวิคเต เอกํ. (สํขิตฺตํ.)}

\csromanheader{2}{1. อนิทสฺสนสปฺปฏิฆปท 7. ปญฺหาวาร}
\csromanheader{2}{\textbf{ปจฺจยจตุกฺก-สหชาตาทิ}}
\proseitem{544}{นเหตุ อนิทสฺสนสปฺปฏิโฆ ธมฺโม นเหตุสฺส อนิทสฺสนสปฺปฏิฆสฺส ธมฺมสฺส สหชาตปจฺจเยน ปจฺจโย.\csromanspacer{} อญฺญมญฺญปจฺจเยน ปจฺจโย.\csromanspacer{} นิสฺสยปจฺจเยน ปจฺจโย.\csromanspacer{} อตฺถิปจฺจเยน ปจฺจโย.\csromanspacer{} อวิคตปจฺจเยน ปจฺจโย. (สพฺพตฺถ เอกํ.)}

\csromanheader{2}{2. อนิทสฺสน-อปฺปฏิฆปท 1. ปฏิจฺจวาร}
\csromanheader{2}{\textbf{ปจฺจยจตุกฺก-เหตุ}}
\proseitem{545}{เหตุํ อนิทสฺสน-อปฺปฏิฆํ ธมฺมํ ปฏิจฺจ เหตุ อนิทสฺสน-อปฺปฏิโฆ ธมฺโม อุปฺปชฺชติ เหตุปจฺจยา.\csromanspacer{} ตีณิ.}

\prose{นเหตุํ อนิทสฺสน-อปฺปฏิฆํ ธมฺมํ ปฏิจฺจ นเหตุ อนิทสฺสน-อปฺปฏิโฆ ธมฺโม อุปฺปชฺชติ เหตุปจฺจยา.\csromanspacer{} ตีณิ.}

\prose{เหตุํ อนิทสฺสน-อปฺปฏิฆญฺจ นเหตุํ อนิทสฺสน-อปฺปฏิฆญฺจ ธมฺมํ ปฏิจฺจ เหตุ อนิทสฺสน-อปฺปฏิโฆ ธมฺโม อุปฺปชฺชติ เหตุปจฺจยา.\csromanspacer{} ตีณิ. (สํขิตฺตํ.)}

\csromanheader{2}{1. เหตุทุก 22. สนิทสฺสนสปฺปฏิฆตฺติก}
\proseitem{546}{\csromanfolio{323}เหตุยา นว, อารมฺมเณ นว ฯเปฯ กมฺเม นว, วิปาเก นว, อาหาเร นว ฯเปฯ อวิคเต นว. (สํขิตฺตํ.)}

\csromanheader{2}{\textbf{นเหตุ น-อธิปติ}}
\proseitem{547}{นเหตุํ อนิทสฺสน-อปฺปฏิฆํ ธมฺมํ ปฏิจฺจ นเหตุ อนิทสฺสน-อปฺปฏิโฆ ธมฺโม อุปฺปชฺชติ นเหตุปจฺจยา.\csromanspacer{}\csromanspacer{}นเหตุํ อนิทสฺสน-อปฺปฏิฆํ ธมฺมํ ปฏิจฺจ เหตุ อนิทสฺสน-อปฺปฏิโฆ ธมฺโม อุปฺปชฺชติ นเหตุปจฺจยา. (2)}

\prose{เหตุํ อนิทสฺสน-อปฺปฏิฆํ ธมฺมํ ปฏิจฺจ นเหตุ อนิทสฺสน-อปฺปฏิโฆ ธมฺโม อุปฺปชฺชติ น-อารมฺมณปจฺจยา.\csromanspacer{} ตีณิ.}

\prose{เหตุํ อนิทสฺสน-อปฺปฏิฆํ ธมฺมํ ปฏิจฺจ เหตุ อนิทสฺสน-อปฺปฏิโฆ ธมฺโม อุปฺปชฺชติ น-อธิปติปจฺจยา. (สํขิตฺตํ.)}

\proseitem{548}{นเหตุยา เทฺว, น-อารมฺมเณ ตีณิ, น-อธิปติยา นว ฯเปฯ นปุเรชาเต นว, นปจฺฉาชาเต นว, น-อาเสวเน นว, นกมฺเม ตีณิ, นวิปาเก นว, น-อาหาเร เอกํ, นฌาเน เอกํ, นมคฺเค เอกํ, นสมฺปยุตฺเต ตีณิ, นวิปฺปยุตฺเต นว, โนนตฺถิยา ตีณิ, โนวิคเต ตีณิ. (สํขิตฺตํ.)}

\csromanheader{2}{เหตุปจฺจยา น-อารมฺมเณ ตีณิ. (สํขิตฺตํ.)}
\csromanheader{2}{นเหตุปจฺจยา อารมฺมเณ เทฺว. (สํขิตฺตํ.)}
\prose{(สหชาตวาโรปิ ปจฺจยวาโรปิ นิสฺสยวาโรปิ สํสฏฺฐวาโรปิ สมฺปยุตฺตวาโรปิ ปฏิจฺจวารสทิสา วิตฺถาเรตพฺพา.)}

\csromanheader{2}{2. อนิทสฺสน-อปฺปฏิฆปท 7. ปญฺหาวาร}
\csromanheader{2}{\textbf{ปจฺจยจตุกฺก-เหตฺวาทิ}}
\proseitem{549}{เหตุ อนิทสฺสน-อปฺปฏิโฆ ธมฺโม เหตุสฺส อนิทสฺสน-อปฺปฏิฆสฺส ธมฺมสฺส เหตุปจฺจเยน ปจฺจโย.\csromanspacer{} ตีณิ.}

\prose{เหตุ อนิทสฺสน-อปฺปฏิโฆ ธมฺโม เหตุสฺส อนิทสฺสน-อปฺปฏิฆสฺส ธมฺมสฺส อารมฺมณปจฺจเยน ปจฺจโย.\csromanspacer{} นว.}

\gambhira{ธมฺมานุโลม ทุกติกปฏฺฐานปาฬิ}
\prose{\csromanfolio{324}เหตุ อนิทสฺสน-อปฺปฏิโฆ ธมฺโม เหตุสฺส อนิทสฺสน-อปฺปฏิฆสฺส ธมฺมสฺส อธิปติปจฺจเยน ปจฺจโย.\csromanspacer{} อารมฺมณาธิปติ, สหชาตาธิปติ.\csromanspacer{} ตีณิ.}

\prose{นเหตุ อนิทสฺสน-อปฺปฏิโฆ ธมฺโม นเหตุสฺส อนิทสฺสน-อปฺปฏิฆสฺส ธมฺมสฺส อธิปติปจฺจเยน ปจฺจโย.\csromanspacer{} อารมฺมณาธิปติ, สหชาตาธิปติ.\csromanspacer{} ตีณิ.}

\prose{เหตุ อนิทสฺสน-อปฺปฏิโฆ จ นเหตุ อนิทสฺสน-อปฺปฏิโฆ จ ธมฺมา เหตุสฺส อนิทสฺสน-อปฺปฏิฆสฺส ธมฺมสฺส อธิปติปจฺจเยน ปจฺจโย.\csromanspacer{} อารมฺมณาธิปติ.\csromanspacer{} ตีณิ. (สํขิตฺตํ.)}

\proseitem{550}{เหตุยา ตีณิ, อารมฺมเณ นว, อธิปติยา นว, อนนฺตเร นว, สมนนฺตเร นว, สหชาเต นว, อญฺญมญฺเญ นว, นิสฺสเย นว, อุปนิสฺสเย นว, ปุเรชาเต ตีณิ, ปจฺฉาชาเต ตีณิ, อาเสวเน นว, กมฺเม ตีณิ, วิปาเก นว, อาหาเร ตีณิ, อินฺทฺริเย นว, ฌาเน ตีณิ, มคฺเค นว, สมฺปยุตฺเต นว, วิปฺปยุตฺเต ปญฺจ, อตฺถิยา นว, นตฺถิยา นว, วิคเต นว, อวิคเต นว. (สํขิตฺตํ.)}

\csromanheader{2}{\textbf{ปจฺจนียุทฺธาร}}
\proseitem{551}{เหตุ อนิทสฺสน-อปฺปฏิโฆ ธมฺโม เหตุสฺส อนิทสฺสน-อปฺปฏิฆสฺส ธมฺมสฺส อารมฺมณปจฺจเยน ปจฺจโย.\csromanspacer{} สหชาตปจฺจเยนปจฺจโย.\csromanspacer{} อุปนิสฺสยปจฺจเยน ปจฺจโย. (สํขิตฺตํ.)}

\proseitem{552}{นเหตุยา นว, น-อารมฺมเณ นว. (สํขิตฺตํ.)}

\csromanheader{2}{เหตุปจฺจยา น-อารมฺมเณ ตีณิ. (สํขิตฺตํ.)}
\csromanheader{2}{นเหตุปจฺจยา อารมฺมเณ นว. (สํขิตฺตํ.)}
\prose{(ยถา กุสลตฺติเก ปญฺหาวารสฺส อนุโลมมฺปิ ปจฺจนียมฺปิ อนุโลมปจฺจนียมฺปิ ปจฺจนียานุโลมมฺปิ คณิตํ, เอวํ คเณตพฺพํ.)}

\nitthitam{เหตุทุกสนิทสฺสนสปฺปฏิฆตฺติกํ นิฏฺฐิตํ.}
\csromanmatikaanchor{mtk.54}
\csromanmatikaanchor{mtk.55}
\csromanmatikaanchor{mtk.58}
\csromantocmarknopage{chapter}{2. สเหตุกทุก}{mtk.54}
\bookmark[dest={mtk.54},level=0]{2. สเหตุกทุก}
\csromantocmarkref{section}{1. กุสลตฺติก ...}{mtk.55}
\bookmark[dest={mtk.55},level=1]{1. กุสลตฺติก ...}
\csromanheader{5}{2. สเหตุกทุก 1. กุสลตฺติก \textbf{2. สเหตุกทุก 1. กุสลตฺติก}}
\csromanheader{2}{1. กุสลปท 1. ปฏิจฺจวาร}
\csromanheader{2}{\textbf{ปจฺจยจตุกฺก-เหตุ อารมฺมณ}}
\proseitem{1}{\csromanfolio{325}สเหตุกํ กุสลํ ธมฺมํ ปฏิจฺจ สเหตุโก กุสโล ธมฺโม อุปฺปชฺชติ เหตุปจฺจยา. (1)}

\prose{สเหตุกํ กุสลํ ธมฺมํ ปฏิจฺจ สเหตุโก กุสโล ธมฺโม อุปฺปชฺชติ อารมฺมณปจฺจยา. (1) (สํขิตฺตํ.)}

\proseitem{2}{เหตุยา เอกํ, อารมฺมเณ เอกํ, อธิปติยา เอกํ, อนนฺตเร เอกํ, สมนนฺตเร เอกํ, สหชาเต เอกํ, อญฺญมญฺเญ เอกํ, นิสฺสเย เอกํ, อุปนิสฺสเย เอกํ, ปุเรชาเต เอกํ, อาเสวเน เอกํ, กมฺเม เอกํ, อาหาเร เอกํ ฯเปฯ อวิคเต เอกํ. (สํขิตฺตํ.)}

\csromanheader{2}{\textbf{ปจฺจนีย-น-อธิปติ}}
\proseitem{3}{สเหตุกํ กุสลํ ธมฺมํ ปฏิจฺจ สเหตุโก กุสโล ธมฺโม อุปฺปชฺชติ น-อธิปติปจฺจยา. (สํขิตฺตํ.)}

\proseitem{4}{น-อธิปติยา เอกํ, นปุเรชาเต เอกํ, นปจฺฉาชาเต เอกํ, น-อาเสวเน เอกํ, นกมฺเม เอกํ, นวิปาเก เอกํ, นวิปฺปยุตฺเต เอกํ. (สํขิตฺตํ.)}

\csromanheader{2}{เหตุปจฺจยา น-อธิปติยา เอกํ. (สํขิตฺตํ.)}
\csromanheader{2}{น-อธิปติปจฺจยา เหตุยา เอกํ. (สํขิตฺตํ.)}
\prose{(สหชาตวาโรปิ ปจฺจยวาโรปิ นิสฺสยวาโรปิ สํสฏฺฐวาโรปิ สมฺปยุตฺตวาโรปิ ปฏิจฺจวารสทิสา วิตฺถาเรตพฺพา.)}

\csromanheader{2}{1. กุสลปท 7. ปญฺหาวาร}
\csromanheader{2}{\textbf{ปจฺจยจตุกฺก\csromandash{}เหตุ อารมฺมณ}}
\proseitem{5}{สเหตุโก กุสโล ธมฺโม สเหตุกสฺส กุสลสฺส ธมฺมสฺส เหตุปจฺจเยน ปจฺจโย. (1)}

\gambhira{ธมฺมานุโลม ทุกติกปฏฺฐานปาฬิ}
\prose{\csromanfolio{326}สเหตุโก กุสโล ธมฺโม สเหตุกสฺส กุสลสฺส ธมฺมสฺส อารมฺมณปจฺจเยน ปจฺจโย. (1) (สํขิตฺตํ.)}

\proseitem{6}{เหตุยา เอกํ, อารมฺมเณ เอกํ, อธิปติยา เอกํ, อนนฺตเร เอกํ, สมนนฺตเร เอกํ, สหชาเต เอกํ, อญฺญมญฺเญ เอกํ, นิสฺสเย เอกํ, อุปนิสฺสเย เอกํ, อาเสวเน เอกํ, กมฺเม เอกํ, อาหาเร เอกํ, อินฺทฺริเย เอกํ, ฌาเน เอกํ, มคฺเค เอกํ, สมฺปยุตฺเต เอกํ, อตฺถิยา เอกํ, นตฺถิยา เอกํ, วิคเต เอกํ, อวิคเต เอกํ. (สํขิตฺตํ.)}

\csromanheader{2}{\textbf{ปจฺจนียุทฺธาร}}
\proseitem{7}{สเหตุโก กุสโล ธมฺโม สเหตุกสฺส กุสลสฺส ธมฺมสฺส อารมฺมณปจฺจเยน ปจฺจโย.\csromanspacer{} สหชาตปจฺจเยน ปจฺจโย.\csromanspacer{} อุปนิสฺสยปจฺจเยน ปจฺจโย.}

\proseitem{8}{นเหตุยา เอกํ, น-อารมฺมเณ เอกํ. (สํขิตฺตํ.)}

\csromanheader{2}{เหตุปจฺจยา น-อารมฺมเณ เอกํ. (สํขิตฺตํ.)}
\csromanheader{2}{นเหตุปจฺจยา อารมฺมเณ เอกํ. (สํขิตฺตํ.)}
\prose{(ยถา กุสลตฺติเก ปญฺหาวารสฺส อนุโลมมฺปิ ปจฺจนียมฺปิ อนุโลมปจฺจนียมฺปิ ปจฺจนียานุโลมมฺปิ คณิตํ, เอวํ คเณตพฺพํ.)}

\csromanheader{2}{2. อกุสลปท 1. ปฏิจฺจวาร}
\csromanheader{2}{\textbf{ปจฺจยจตุกฺก-เหตุ อารมฺมณ}}
\proseitem{9}{สเหตุกํ อกุสลํ ธมฺมํ ปฏิจฺจ สเหตุโก อกุสโล ธมฺโม อุปฺปชฺชติ เหตุปจฺจยา. (1)}

\prose{อเหตุกํ อกุสลํ ธมฺมํ ปฏิจฺจ สเหตุโก อกุสโล ธมฺโม อุปฺปชฺชติ เหตุปจฺจยา. (1)}

\prose{สเหตุกํ อกุสลญฺจ อเหตุกํ อกุสลญฺจ ธมฺมํ ปฏิจฺจ สเหตุโก อกุสโล ธมฺโม อุปฺปชฺชติ เหตุปจฺจยา. (1)}

\csromanheader{2}{2. สเหตุกทุก 1. กุสลตฺติก}
\prose{\csromanfolio{327}สเหตุกํ อกุสลํ ธมฺมํ ปฏิจฺจ สเหตุโก อกุสโล ธมฺโม อุปฺปชฺชติ อารมฺมณปจฺจยา.\csromanspacer{}\csromanspacer{}สเหตุกํ อกุสลํ ธมฺมํ ปฏิจฺจ อเหตุโก อกุสโล ธมฺโม อุปฺปชฺชติ อารมฺมณปจฺจยา.\csromanspacer{}\csromanspacer{}สเหตุกํ อกุสลํ ธมฺมํ ปฏิจฺจ สเหตุโก อกุสโล จ อเหตุโก อกุสโล จ ธมฺมา อุปฺปชฺชนฺติ อารมฺมณปจฺจยา. (3)}

\prose{อเหตุกํ อกุสลํ ธมฺมํ ปฏิจฺจ สเหตุโก อกุสโล ธมฺโม อุปฺปชฺชติ อารมฺมณปจฺจยา. (1)}

\prose{สเหตุกํ อกุสลญฺจ อเหตุกํ อกุสลญฺจ ธมฺมํ ปฏิจฺจ สเหตุโก อกุสโล ธมฺโม อุปฺปชฺชติ อารมฺมณปจฺจยา. (1)}

\prose{สเหตุกํ อกุสลํ ธมฺมํ ปฏิจฺจ สเหตุโก อกุสโล ธมฺโม อุปฺปชฺชติ อธิปติปจฺจยา. (สํขิตฺตํ.)}

\proseitem{10}{เหตุยา ตีณิ, อารมฺมเณ ปญฺจ, อธิปติยา เอกํ, อนนฺตเร ปญฺจ, สมนนฺตเร ปญฺจ, สหชาเต ปญฺจ, อญฺญมญฺเญ ปญฺจ, นิสฺสเย ปญฺจ, อุปนิสฺสเย ปญฺจ, ปุเรชาเต ปญฺจ, อาเสวเน ปญฺจ, กมฺเม ปญฺจ, อาหาเร ปญฺจ, อินฺทฺริเย ปญฺจ, ฌาเน ปญฺจ, มคฺเค ปญฺจ, สมฺปยุตฺเต ปญฺจ, วิปฺปยุตฺเต ปญฺจ, อตฺถิยา ปญฺจ, นตฺถิยา ปญฺจ, วิคเต ปญฺจ, อวิคเต ปญฺจ. (สํขิตฺตํ.)}

\csromanheader{2}{\textbf{นเหตุ น-อธิปติ}}
\proseitem{11}{สเหตุกํ อกุสลํ ธมฺมํ ปฏิจฺจ อเหตุโก อกุสโล ธมฺโม อุปฺปชฺชติ นเหตุปจฺจยา. (1)}

\prose{สเหตุกํ อกุสลํ ธมฺมํ ปฏิจฺจ สเหตุโก อกุสโล ธมฺโม อุปฺปชฺชติ น-อธิปติปจฺจยา. (สํขิตฺตํ.)}

\proseitem{12}{นเหตุยา เอกํ, น-อธิปติยา ปญฺจ, นปุเรชาเต ปญฺจ, นปจฺฉาชาเต ปญฺจ, น-อาเสวเน ปญฺจ, นกมฺเม ตีณิ, นวิปาเก ปญฺจ, นวิปฺปยุตฺเต ปญฺจ. (สํขิตฺตํ.)}

\csromanheader{2}{เหตุปจฺจยา น-อธิปติยา ตีณิ. (สํขิตฺตํ.)}
\csromanheader{2}{นเหตุปจฺจยา อารมฺมเณ เอกํ. (สํขิตฺตํ.)}
\prose{(สหชาตวาโรปิ ปจฺจยวาโรปิ นิสฺสยวาโรปิ สํสฏฺฐวาโรปิ สมฺปยุตฺตวาโรปิ ปฏิจฺจวารสทิสา วิตฺถาเรตพฺพา.)}

\csromanheader{2}{ธมฺมานุโลม ทุกติกปฏฺฐานปาฬิ \textbf{2. อกุสลปท 7. ปญฺหาวาร}}
\csromanheader{2}{\textbf{ปจฺจยจตุกฺก-เหตุ อารมฺมณาทิ}}
\proseitem{13}{\csromanfolio{328}สเหตุโก อกุสโล ธมฺโม สเหตุกสฺส อกุสลสฺส ธมฺมสฺส เหตุปจฺจเยน ปจฺจโย.\csromanspacer{} เทฺว.}

\prose{สเหตุโก อกุสโล ธมฺโม สเหตุกสฺส อกุสลสฺส ธมฺมสฺส อารมฺมณปจฺจเยน ปจฺจโย.\csromanspacer{} ตีณิ.}

\prose{อเหตุโก อกุสโล ธมฺโม อเหตุกสฺส อกุสลสฺส ธมฺมสฺส อารมฺมณปจฺจเยน ปจฺจโย.\csromanspacer{} ตีณิ.}

\prose{สเหตุโก อกุสโล จ อเหตุโก อกุสโล จ ธมฺมา สเหตุกสฺส อกุสลสฺส ธมฺมสฺส อารมฺมณปจฺจเยน ปจฺจโย.\csromanspacer{} ตีณิ.}

\prose{สเหตุโก อกุสโล ธมฺโม สเหตุกสฺส อกุสลสฺส ธมฺมสฺส อธิปติปจฺจเยน ปจฺจโย.\csromanspacer{} อารมฺมณาธิปติ, สหชาตาธิปติ. (สํขิตฺตํ.)}

\proseitem{14}{เหตุยา เทฺว, อารมฺมเณ นว, อธิปติยา เอกํ, อนนฺตเร นว, สมนนฺตเร นว, สหชาเต ปญฺจ, อญฺญมญฺเญ ปญฺจ, นิสฺสเย ปญฺจ, อุปนิสฺสเย นว, อาเสวเน นว, กมฺเม ตีณิ, อาหาเร ตีณิ, อินฺทฺริเย ตีณิ, ฌาเน ตีณิ, มคฺเค ตีณิ, สมฺปยุตฺเต ปญฺจ, อตฺถิยา ปญฺจ, นตฺถิยา นว, วิคเต นว, อวิคเต ปญฺจ. (สํขิตฺตํ.)}

\csromanheader{2}{\textbf{ปจฺจนียุทฺธาร}}
\proseitem{15}{สเหตุโก อกุสโล ธมฺโม สเหตุกสฺส อกุสลสฺส ธมฺมสฺส อารมฺมณปจฺจเยน ปจฺจโย.\csromanspacer{} สหชาตปจฺจเยน ปจฺจโย.\csromanspacer{} อุปนิสฺสยปจฺจเยน ปจฺจโย. (สํขิตฺตํ.)}

\proseitem{16}{นเหตุยา นว, น-อารมฺมเณ นว. (สํขิตฺตํ.)}

\csromanheader{2}{เหตุปจฺจยา น-อารมฺมเณ เทฺว. (สํขิตฺตํ.)}
\csromanheader{2}{นเหตุปจฺจยา อารมฺมเณ นว. (สํขิตฺตํ.)}
\prose{(ยถา กุสลตฺติเก ปญฺหาวารสฺส อนุโลมมฺปิ ปจฺจนียมฺปิ อนุโลมปจฺจนียมฺปิ ปจฺจนียานุโลมมฺปิ คณิตํ, เอวํ คเณตพฺพํ.)}

\csromanheader{2}{2. สเหตุกทุก 1. กุสลตฺติก \textbf{3. อพฺยากตปท 1. ปฏิจฺจวาร}}
\csromanheader{2}{\textbf{ปจฺจยจตุกฺก-เหตุ}}
\proseitem{17}{\csromanfolio{329}สเหตุกํ อพฺยากตํ ธมฺมํ ปฏิจฺจ สเหตุโก อพฺยากโต ธมฺโม อุปฺปชฺชติ เหตุปจฺจยา.\csromanspacer{} ตีณิ.}

\prose{อเหตุกํ อพฺยากตํ ธมฺมํ ปฏิจฺจ อเหตุโก อพฺยากโต ธมฺโม อุปฺปชฺชติ เหตุปจฺจยา.\csromanspacer{} ตีณิ.}

\prose{สเหตุกํ อพฺยากตญฺจ อเหตุกํ อพฺยากตญฺจ ธมฺมํ ปฏิจฺจ สเหตุโก อพฺยากโต ธมฺโม อุปฺปชฺชติ เหตุปจฺจยา.\csromanspacer{} ตีณิ. (สํขิตฺตํ.)}

\proseitem{18}{เหตุยา นว, อารมฺมเณ จตฺตาริ, อธิปติยา ปญฺจ, อนนฺตเร จตฺตาริ, สมนนฺตเร จตฺตาริ, สหชาเต นว, อญฺญมญฺเญ ฉ ฯเปฯ ปุเรชาเต เทฺว, อาเสวเน เทฺว, กมฺเม นว, วิปาเก นว ฯเปฯ วิปฺปยุตฺเต นว ฯเปฯ อวิคเต นว. (สํขิตฺตํ.)}

\csromanheader{2}{\textbf{นเหตุ น-อารมฺมณ}}
\proseitem{19}{อเหตุกํ อพฺยากตํ ธมฺมํ ปฏิจฺจ อเหตุโก อพฺยากโต ธมฺโม อุปฺปชฺชติ นเหตุปจฺจยา. (1)}

\prose{สเหตุกํ อพฺยากตํ ธมฺมํ ปฏิจฺจ อเหตุโก อพฺยากโต ธมฺโม อุปฺปชฺชติ น-อารมฺมณปจฺจยา. (1)}

\prose{อเหตุกํ อพฺยากตํ ธมฺมํ ปฏิจฺจ อเหตุโก อพฺยากโต ธมฺโม อุปฺปชฺชติ น-อารมฺมณปจฺจยา. (1)}

\prose{สเหตุกํ อพฺยากตญฺจ อเหตุกํ อพฺยากตญฺจ ธมฺมํ ปฏิจฺจ อเหตุโก อพฺยากโต ธมฺโม อุปฺปชฺชติ น-อารมฺมณปจฺจยา. (1)}

\prose{สเหตุกํ อพฺยากตํ ธมฺมํ ปฏิจฺจ สเหตุโก อพฺยากโต ธมฺโม อุปฺปชฺชติ น-อธิปติปจฺจยา. (สํขิตฺตํ.)}

\proseitem{20}{นเหตุยา เอกํ, น-อารมฺมเณ ตีณิ, น-อธิปติยา นว, น-อนนฺตเร ตีณิ, นสมนนฺตเร ตีณิ, น-อญฺญมญฺเญ ตีณิ, น-อุปนิสฺสเย ตีณิ, นปุเรชาเต นว, นปจฺฉาชาเต นว, น-อาเสวเน นว, นกมฺเม เทฺว,}

\vspace{\tipitakaparskip}
\csromancenter{ธมฺมานุโลม ทุกติกปฏฺฐานปาฬิ นวิปาเก ปญฺจ, น-อาหาเร เอกํ, น-อินฺทฺริเย เอกํ, นฌาเน เอกํ, นมคฺเค เอกํ, นสมฺปยุตฺเต ตีณิ, นวิปฺปยุตฺเต เทฺว, โนนตฺถิยา ตีณิ, โนวิคเต ตีณิ. (สํขิตฺตํ.)}
\csromanheader{2}{เหตุปจฺจยา น-อารมฺมเณ ตีณิ. (สํขิตฺตํ.)}
\csromanheader{2}{นเหตุปจฺจยา อารมฺมเณ เอกํ. (สํขิตฺตํ.)}
\prose{\csromanfolio{330}(สหชาตวาโรปิ ปจฺจยวาโรปิ นิสฺสยวาโรปิ สํสฏฺฐวาโรปิ สมฺปยุตฺตวาโรปิ ปฏิจฺจวารสทิสา วิตฺถาเรตพฺพา.)}

\csromanheader{2}{3. อพฺยากตปท 7. ปญฺหาวาร}
\csromanheader{2}{\textbf{ปจฺจยจตุกฺก-เหตุ อารมฺมณ}}
\proseitem{21}{สเหตุโก อพฺยากโต ธมฺโม สเหตุกสฺส อพฺยากตสฺส ธมฺมสฺส เหตุปจฺจเยน ปจฺจโย.\csromanspacer{} ตีณิ.}

\prose{สเหตุโก อพฺยากโต ธมฺโม สเหตุกสฺส อพฺยากตสฺส ธมฺมสฺส อารมฺมณปจฺจเยน ปจฺจโย. (สํขิตฺตํ.)}

\proseitem{22}{เหตุยา ตีณิ, อารมฺมเณ จตฺตาริ, อธิปติยา จตฺตาริ, อนนฺตเร จตฺตาริ, สมนนฺตเร จตฺตาริ, สหชาเต สตฺต, อญฺญมญฺเญ ฉ, นิสฺสเย สตฺต, อุปนิสฺสเย จตฺตาริ, ปุเรชาเต เทฺว, ปจฺฉาชาเต เทฺว, อาเสวเน เทฺว, กมฺเม จตฺตาริ, วิปาเก จตฺตาริ, อาหาเร จตฺตาริ อินฺทฺริเย จตฺตาริ, ฌาเน จตฺตาริ, มคฺเค ตีณิ, สมฺปยุตฺเต เทฺว, วิปฺปยุตฺเต ตีณิ, อตฺถิยา สตฺต ฯเปฯ อวิคเต สตฺต. (สํขิตฺตํ.)}

\csromanheader{2}{\textbf{ปจฺจนียุทฺธาร}}
\proseitem{23}{สเหตุโก อพฺยากโต ธมฺโม สเหตุกสฺส อพฺยากตสฺส ธมฺมสฺส อารมฺมณปจฺจเยน ปจฺจโย.\csromanspacer{} สหชาตปจฺจเยน ปจฺจโย.\csromanspacer{} อุปนิสฺสยปจฺจเยน ปจฺจโย. (สํขิตฺตํ.)}

\proseitem{24}{นเหตุยา สตฺต, น-อารมฺมเณ สตฺต. (สํขิตฺตํ.)}

\csromanheader{2}{เหตุปจฺจยา น-อารมฺมเณ ตีณิ. (สํขิตฺตํ.)}
\csromanheader{2}{นเหตุปจฺจยา อารมฺมเณ จตฺตาริ. (สํขิตฺตํ.)}
\csromanmatikaanchor{mtk.56}
\csromanmatikaanchor{mtk.57}
\csromantocmarknopage{subsection}{3. เหตุสมฺปยุตฺตทุก}{mtk.56}
\bookmark[dest={mtk.56},level=2]{3. เหตุสมฺปยุตฺตทุก}
\csromantocmarkref{paragraph}{1. กุสลตฺติก ...}{mtk.57}
\bookmark[dest={mtk.57},level=4]{1. กุสลตฺติก ...}
\csromantocmarknopage{subparagraph}{4. เหตุ สเหตุกทุก}{mtk.58}
\bookmark[dest={mtk.58},level=5]{4. เหตุ สเหตุกทุก}
\csromanheader{4}{3. เหตุสมฺปยุตฺตทุก 1. กุสลตฺติก}
\prose{\csromanfolio{331}(ยถา กุสลตฺติเก ปญฺหาวารสฺส อนุโลมมฺปิ ปจฺจนียมฺปิ อนุโลมปจฺจนียมฺปิ ปจฺจนียานุโลมมฺปิ คณิตํ, เอวํ คเณตพฺพํ.)}

\nitthitam{สเหตุกทุกกุสลตฺติกํ นิฏฺฐิตํ.}
\csromanheader{2}{3. เหตุสมฺปยุตฺตทุก 1. กุสลตฺติก}
\csromanheader{2}{1. กุสลปท 1. ปฏิจฺจวาร}
\csromanheader{2}{\textbf{ปจฺจยจตุกฺก-เหตุ}}
\proseitem{25}{เหตุสมฺปยุตฺตํ กุสลํ ธมฺมํ ปฏิจฺจ เหตุสมฺปยุตฺโต กุสโล ธมฺโม อุปฺปชฺชติ เหตุปจฺจยา. (1) (สํขิตฺตํ.)}

\proseitem{26}{เหตุยา เอกํ, อารมฺมเณ เอกํ, อธิปติยา เอกํ ฯเปฯ อวิคเต เอกํ. (สํขิตฺตํ.)}

\csromanheader{2}{\textbf{ปจฺจนีย}}
\proseitem{27}{น-อธิปติยา เอกํ, นปุเรชาเต เอกํ, นปจฺฉาชาเต เอกํ, น-อาเสวเน เอกํ, นกมฺเม เอกํ, นวิปาเก เอกํ, นวิปฺปยุตฺเต เอกํ. (สํขิตฺตํ.)}

\csromanheader{2}{เหตุปจฺจยา น-อธิปติยา เอกํ. (สํขิตฺตํ.)}
\csromanheader{2}{น-อธิปติปจฺจยา เหตุยา เอกํ. (สํขิตฺตํ.)}
\prose{(สหชาตวาโรปิ ปจฺจยวาโรปิ นิสฺสยวาโรปิ สํสฏฺฐวาโรปิ สมฺปยุตฺตวาโรปิ ปฏิจฺจวารสทิสา วิตฺถาเรตพฺพา.)}

\csromanheader{2}{1. กุสลปท 7. ปญฺหาวาร}
\csromanheader{2}{\textbf{ปจฺจยจตุกฺก-เหตุ}}
\proseitem{28}{เหตุสมฺปยุตฺโต กุสโล ธมฺโม เหตุสมฺปยุตฺตสฺส กุสลสฺส ธมฺมสฺส เหตุปจฺจเยน ปจฺจโย. (สํขิตฺตํ.)}

\prose{เหตุยา เอกํ, อารมฺมเณ เอกํ, อธิปติยา เอกํ, อนนฺตเร เอกํ, สมนนฺตเร เอกํ ฯเปฯ กมฺเม เอกํ ฯเปฯ อวิคเต เอกํ. (สํขิตฺตํ.)}

\csromanheader{2}{ธมฺมานุโลม ทุกติกปฏฺฐานปาฬิ \textbf{ปจฺจนียุทฺธาร}}
\proseitem{29}{\csromanfolio{332}เหตุสมฺปยุตฺโต กุสโล ธมฺโม เหตุสมฺปยุตฺตสฺส กุสลสฺส ธมฺมสฺส อารมฺมณปจฺจเยน ปจฺจโย.\csromanspacer{} สหชาตปจฺจเยน ปจฺจโย.\csromanspacer{} อุปนิสฺสยปจฺจเยน ปจฺจโย.}

\proseitem{30}{นเหตุยา เอกํ, น-อารมฺมเณ เอกํ. (สํขิตฺตํ.)}

\csromanheader{2}{เหตุปจฺจยา น-อารมฺมเณ เอกํ. (สํขิตฺตํ.)}
\csromanheader{2}{นเหตุปจฺจยา อารมฺมเณ เอกํ. (สํขิตฺตํ.)}
\prose{(ยถา กุสลตฺติเก ปญฺหาวารสฺส อนุโลมมฺปิ ปจฺจนียมฺปิ อนุโลมปจฺจนียมฺปิ ปจฺจนียานุโลมมฺปิ คณิตํ, เอวํ คเณตพฺพํ.)}

\csromanheader{2}{2. อกุสลปท 1. ปฏิจฺจวาร}
\csromanheader{2}{\textbf{ปจฺจยจตุกฺก-เหตุ อารมฺมณ}}
\proseitem{31}{เหตุสมฺปยุตฺตํ อกุสลํ ธมฺมํ ปฏิจฺจ เหตุสมฺปยุตฺโต อกุสโล ธมฺโม อุปฺปชฺชติ เหตุปจฺจยา. (1)}

\prose{เหตุวิปฺปยุตฺตํ อกุสลํ ธมฺมํ ปฏิจฺจ เหตุสมฺปยุตฺโต อกุสโล ธมฺโม อุปฺปชฺชติ เหตุปจฺจยา. (1)}

\prose{เหตุสมฺปยุตฺตํ อกุสลญฺจ เหตุวิปฺปยุตฺตํ อกุสลญฺจ ธมฺมํ ปฏิจฺจ เหตุสมฺปยุตฺโต อกุสโล ธมฺโม อุปฺปชฺชติ เหตุปจฺจยา. (1)}

\prose{เหตุสมฺปยุตฺตํ อกุสลํ ธมฺมํ ปฏิจฺจ เหตุสมฺปยุตฺโต อกุสโล ธมฺโม อุปฺปชฺชติ อารมฺมณปจฺจยา.\csromanspacer{} ตีณิ.}

\prose{เหตุวิปฺปยุตฺตํ อกุสลํ ธมฺมํ ปฏิจฺจ เหตุสมฺปยุตฺโต อกุสโล ธมฺโม อุปฺปชฺชติ อารมฺมณปจฺจยา. (1)}

\prose{เหตุสมฺปยุตฺตํ อกุสลญฺจ เหตุวิปฺปยุตฺตํ อกุสลญฺจ ธมฺมํ ปฏิจฺจ เหตุสมฺปยุตฺโต อกุสโล ธมฺโม อุปฺปชฺชติ อารมฺมณปจฺจยา. (1)}

\prose{เหตุสมฺปยุตฺตํ อกุสลํ ธมฺมํ ปฏิจฺจ เหตุสมฺปยุตฺโต อกุสโล ธมฺโม อุปฺปชฺชติ อธิปติปจฺจยา. (สํขิตฺตํ.)}

\csromanheader{2}{3. เหตุสมฺปยุตฺตทุก 1. กุสลตฺติก}
\proseitem{32}{\csromanfolio{333}เหตุยา ตีณิ, อารมฺมเณ ปญฺจ, อธิปติยา เอกํ, อนนฺตเร ปญฺจ, สมนนฺตเร ปญฺจ, สหชาเต ปญฺจ, อญฺญมญฺเญ ปญฺจ, นิสฺสเย ปญฺจ, อุปนิสฺสเย ปญฺจ, ปุเรชาเต ปญฺจ, อาเสวเน ปญฺจ, กมฺเม ปญฺจ, อาหาเร ปญฺจ, อินฺทฺริเย ปญฺจ, ฌาเน ปญฺจ, มคฺเค ปญฺจ, สมฺปยุตฺเต ปญฺจ, วิปฺปยุตฺเต ปญฺจ, อตฺถิยา ปญฺจ ฯเปฯ อวิคเต ปญฺจ. (สํขิตฺตํ.)}

\csromanheader{2}{\textbf{นเหตุ น-อธิปติ}}
\proseitem{33}{เหตุสมฺปยุตฺตํ อกุสลํ ธมฺมํ ปฏิจฺจ เหตุวิปฺปยุตฺโต อกุสโล ธมฺโม อุปฺปชฺชติ นเหตุปจฺจยา. (1)}

\prose{เหตุสมฺปยุตฺตํ อกุสลํ ธมฺมํ ปฏิจฺจ เหตุสมฺปยุตฺโต อกุสโล ธมฺโม อุปฺปชฺชติ น-อธิปติปจฺจยา. (สํขิตฺตํ.)}

\proseitem{34}{นเหตุยา เอกํ, น-อธิปติยา ปญฺจ, นปุเรชาเต ปญฺจ, นปจฺฉาชาเต ปญฺจ, น-อาเสวเน ปญฺจ, นกมฺเม ตีณิ, นวิปาเก ปญฺจ, นวิปฺปยุตฺเต ปญฺจ. (สํขิตฺตํ.)}

\csromanheader{2}{เหตุปจฺจยา น-อธิปติยา ตีณิ. (สํขิตฺตํ.)}
\csromanheader{2}{นเหตุปจฺจยา อารมฺมเณ เอกํ. (สํขิตฺตํ.)}
\prose{(สหชาตวาโรปิ ปจฺจยวาโรปิ นิสฺสยวาโรปิ สํสฏฺฐวาโรปิ สมฺปยุตฺตวาโรปิ ปฏิจฺจวารสทิสา วิตฺถาเรตพฺพา.)}

\csromanheader{2}{2. อกุสลปท 7. ปญฺหาวาร}
\csromanheader{2}{\textbf{ปจฺจยจตุกฺก-เหตุ อารมฺมณ}}
\proseitem{35}{เหตุสมฺปยุตฺโต อกุสโล ธมฺโม เหตุสมฺปยุตฺตสฺส อกุสลสฺส ธมฺมสฺส เหตุปจฺจเยน ปจฺจโย.\csromanspacer{} เทฺว.}

\prose{เหตุสมฺปยุตฺโต อกุสโล ธมฺโม เหตุสมฺปยุตฺตสฺส อกุสลสฺส ธมฺมสฺส อารมฺมณปจฺจเยน ปจฺจโย.\csromanspacer{} ตีณิ.}

\prose{เหตุวิปฺปยุตฺโต อกุสโล ธมฺโม เหตุวิปฺปยุตฺตสฺส อกุสลสฺส ธมฺมสฺส อารมฺมณปจฺจเยน ปจฺจโย.\csromanspacer{} ตีณิ.}

\gambhira{ธมฺมานุโลม ทุกติกปฏฺฐานปาฬิ}
\prose{\csromanfolio{334}เหตุสมฺปยุตฺโต อกุสโล จ เหตุวิปฺปยุตฺโต อกุสโล จ ธมฺมา เหตุสมฺปยุตฺตสฺส อกุสลสฺส ธมฺมสฺส อารมฺมณปจฺจเยน ปจฺจโย.\csromanspacer{} ตีณิ. (สํขิตฺตํ.)}

\proseitem{36}{เหตุยา เทฺว, อารมฺมเณ นว, อธิปติยา เอกํ, อนนฺตเร นว, สมนนฺตเร นว, สหชาเต ปญฺจ, อญฺญมญฺเญ ปญฺจ, นิสฺสเย ปญฺจ, อุปนิสฺสเย นว, อาเสวเน นว, กมฺเม ตีณิ, อาหาเร ตีณิ, อินฺทฺริเย ตีณิ, ฌาเน ตีณิ, มคฺเค ตีณิ, สมฺปยุตฺเต ปญฺจ, อตฺถิยา ปญฺจ ฯเปฯ อวิคเต ปญฺจ. (สํขิตฺตํ.)}

\csromanheader{2}{\textbf{ปจฺจนียุทฺธาร}}
\proseitem{37}{เหตุสมฺปยุตฺโต อกุสโล ธมฺโม เหตุสมฺปยุตฺตสฺส อกุสลสฺส ธมฺมสฺส อารมฺมณปจฺจเยน ปจฺจโย.\csromanspacer{} สหชาตปจฺจเยน ปจฺจโย.\csromanspacer{} อุปนิสฺสยปจฺจเยน ปจฺจโย. (สํขิตฺตํ.)}

\proseitem{38}{นเหตุยา นว, น-อารมฺมเณ นว. (สํขิตฺตํ.)}

\csromanheader{2}{เหตุปจฺจยา น-อารมฺมเณ เทฺว. (สํขิตฺตํ.)}
\csromanheader{2}{นเหตุปจฺจยา อารมฺมเณ นว. (สํขิตฺตํ.)}
\prose{(ยถา กุสลตฺติเก ปญฺหาวารสฺส อนุโลมมฺปิ ปจฺจนียมฺปิ อนุโลมปจฺจนียมฺปิ ปจฺจนียานุโลมมฺปิ คณิตํ, เอวํ คเณตพฺพํ.)}

\csromanheader{2}{3. อพฺยากตปท 1. ปฏิจฺจวาร}
\csromanheader{2}{\textbf{ปจฺจยจตุกฺก-เหตุ}}
\proseitem{39}{เหตุสมฺปยุตฺตํ อพฺยากตํ ธมฺมํ ปฏิจฺจ เหตุสมฺปยุตฺโต อพฺยากโต ธมฺโม อุปฺปชฺชติ เหตุปจฺจยา.\csromanspacer{} ตีณิ.}

\prose{เหตุวิปฺปยุตฺตํ อพฺยากตํ ธมฺมํ ปฏิจฺจ เหตุวิปฺปยุตฺโต อพฺยากโต ธมฺโม อุปฺปชฺชติ เหตุปจฺจยา.\csromanspacer{} ตีณิ.}

\prose{เหตุสมฺปยุตฺตํ อพฺยากตญฺจ เหตุวิปฺปยุตฺตํ อพฺยากตญฺจ ธมฺมํ ปฏิจฺจ เหตุสมฺปยุตฺโต อพฺยากโต ธมฺโม อุปฺปชฺชติ เหตุปจฺจยา.\csromanspacer{} ตีณิ. (สํขิตฺตํ.)}

\csromanheader{2}{3. เหตุสมฺปยุตฺตทุก 1. กุสลตฺติก}
\proseitem{40}{\csromanfolio{335}เหตุยา นว, อารมฺมเณ จตฺตาริ, อธิปติยา ปญฺจ, อนนฺตเร จตฺตาริ, สมนนฺตเร จตฺตาริ, สหชาเต นว, อญฺญมญฺเญ ฉ ฯเปฯ ปุเรชาเต เทฺว, อาเสวเน เทฺว, กมฺเม นว, วิปาเก นว ฯเปฯ อวิคเต นว. (สํขิตฺตํ.)}

\csromanheader{2}{\textbf{นเหตุ น-อารมฺมณ}}
\proseitem{41}{เหตุวิปฺปยุตฺตํ อพฺยากตํ ธมฺมํ ปฏิจฺจ เหตุวิปฺปยุตฺโต อพฺยากโต ธมฺโม อุปฺปชฺชติ นเหตุปจฺจยา. (1)}

\prose{เหตุสมฺปยุตฺตํ อพฺยากตํ ธมฺมํ ปฏิจฺจ เหตุวิปฺปยุตฺโต อพฺยากโต ธมฺโม อุปฺปชฺชติ น-อารมฺมณปจฺจยา. (สํขิตฺตํ.)}

\proseitem{42}{นเหตุยา เอกํ, น-อารมฺมเณ ตีณิ, น-อธิปติยา นว, น-อนนฺตเร ตีณิ, นสมนนฺตเร ตีณิ, น-อญฺญมญฺเญ ตีณิ, น-อุปนิสฺสเย ตีณิ, นปุเรชาเต นว, นปจฺฉาชาเต นว, น-อาเสวเน นว, นกมฺเม เทฺว, นวิปาเก ปญฺจ, น-อาหาเร เอกํ, น-อินฺทฺริเย เอกํ, นฌาเน เอกํ, นมคฺเค เอกํ, นสมฺปยุตฺเต ตีณิ, นวิปฺปยุตฺเต เทฺว, โนนตฺถิยา ตีณิ, โนวิคเต ตีณิ. (สํขิตฺตํ.)}

\csromanheader{2}{เหตุปจฺจยา น-อารมฺมเณ ตีณิ. (สํขิตฺตํ.)}
\csromanheader{2}{นเหตุปจฺจยา อารมฺมเณ เอกํ. (สํขิตฺตํ.)}
\prose{(สหชาตวาโรปิ ปจฺจยวาโรปิ นิสฺสยวาโรปิ สํสฏฺฐวาโรปิ สมฺปยุตฺตวาโรปิ ปฏิจฺจวารสทิสา วิตฺถาเรตพฺพา.)}

\csromanheader{2}{3. อพฺยากตปท 7. ปญฺหาวาร}
\csromanheader{2}{\textbf{ปจฺจยจตุกฺก-เหตุ อารมฺมณ}}
\proseitem{43}{เหตุสมฺปยุตฺโต อพฺยากโต ธมฺโม เหตุสมฺปยุตฺตสฺส อพฺยากตสฺส ธมฺมสฺส เหตุปจฺจเยน ปจฺจโย.\csromanspacer{} ตีณิ.}

\prose{เหตุสมฺปยุตฺโต อพฺยากโต ธมฺโม เหตุสมฺปยุตฺตสฺส อพฺยากตสฺส ธมฺมสฺส อารมฺมณปจฺจเยน ปจฺจโย. (สํขิตฺตํ.)}

\gambhira{ธมฺมานุโลม ทุกติกปฏฺฐานปาฬิ}
\proseitem{44}{\csromanfolio{336}เหตุยา ตีณิ, อารมฺมเณ จตฺตาริ, อธิปติยา จตฺตาริ, อนนฺตเร จตฺตาริ, กมฺเม จตฺตาริ, วิปาเก จตฺตาริ ฯเปฯ อวิคเต สตฺต. (สํขิตฺตํ.)}

\csromanheader{2}{\textbf{ปจฺจนียุทฺธาร}}
\proseitem{45}{เหตุสมฺปยุตฺโต อพฺยากโต ธมฺโม เหตุสมฺปยุตฺตสฺส อพฺยากตสฺส ธมฺมสฺส อารมฺมณปจฺจเยน ปจฺจโย.\csromanspacer{} สหชาตปจฺจเยน ปจฺจโย.\csromanspacer{} อุปนิสฺสยปจฺจเยน ปจฺจโย.}

\proseitem{46}{นเหตุยา สตฺต, น-อารมฺมเณ สตฺต. (สํขิตฺตํ.)}

\csromanheader{2}{เหตุปจฺจยา น-อารมฺมเณ ตีณิ. (สํขิตฺตํ.)}
\csromanheader{2}{นเหตุปจฺจยา อารมฺมเณ จตฺตาริ. (สํขิตฺตํ.)}
\prose{(ยถา กุสลตฺติเก ปญฺหาวารสฺส อนุโลมมฺปิ ปจฺจนียมฺปิ อนุโลมปจฺจนียมฺปิ ปจฺจนียานุโลมมฺปิ คณิตํ, เอวํ คเณตพฺพํ.)}

\nitthitam{เหตุสมฺปยุตฺตทุกกุสลตฺติกํ นิฏฺฐิตํ.}
\csromanmatikaanchor{mtk.59}
\csromantocmarkref{subparagraph}{1. กุสลตฺติก ...}{mtk.59}
\bookmark[dest={mtk.59},level=5]{1. กุสลตฺติก ...}
\csromanheader{6}{4. เหตุสเหตุกทุก 1. กุสลตฺติก}
\csromanheader{2}{1. กุสลปท 1. ปฏิจฺจวาร}
\csromanheader{2}{\textbf{ปจฺจยจตุกฺก-เหตุ}}
\proseitem{47}{เหตุญฺเจว สเหตุกญฺจ กุสลํ ธมฺมํ ปฏิจฺจ เหตุ เจว สเหตุโก จ กุสโล ธมฺโม อุปฺปชฺชติ เหตุปจฺจยา.\csromanspacer{} ตีณิ.}

\prose{สเหตุกญฺเจว น จ เหตุํ กุสลํ ธมฺมํ ปฏิจฺจ สเหตุโก เจว น จ เหตุ กุสโล ธมฺโม อุปฺปชฺชติ เหตุปจฺจยา.\csromanspacer{} ตีณิ.}

\prose{เหตุญฺเจว สเหตุกญฺจ กุสลญฺจ สเหตุกญฺเจว น จ เหตุํ กุสลญฺจ ธมฺมํ ปฏิจฺจ เหตุ เจว สเหตุโก จ กุสโล ธมฺโม อุปฺปชฺชติ เหตุปจฺจยา.\csromanspacer{} ตีณิ. (สํขิตฺตํ.)}

\csromanheader{2}{4. เหตุสเหตุกทุก 1. กุสลตฺติก}
\proseitem{48}{\csromanfolio{337}เหตุยา นว, อารมฺมเณ นว ฯเปฯ กมฺเม นว, อาหาเร นว ฯเปฯ อวิคเต นว. (สํขิตฺตํ.)}

\csromanheader{2}{\textbf{น-อธิปติ}}
\proseitem{49}{เหตุญฺเจว สเหตุกญฺจ กุสลํ ธมฺมํ ปฏิจฺจ เหตุ เจว สเหตุโก จ กุสโล ธมฺโม อุปฺปชฺชติ น-อธิปติปจฺจยา. (สํขิตฺตํ.)}

\vspace{\tipitakaparskip}
\csromancenter{\textbf{น-อธิปติ}ยา นว, นปุเรชาเต นว, นปจฺฉาชาเต นว, น-อาเสวเน นว, นกมฺเม ตีณิ, นวิปาเก นว, นวิปฺปยุตฺเต นว. (สํขิตฺตํ.)}
\csromanheader{2}{เหตุปจฺจยา น-อธิปติยา นว. (สํขิตฺตํ.)}
\csromanheader{2}{\textbf{น-อธิปติ}ปจฺจยา เหตุยา นว. (สํขิตฺตํ.)}
\prose{(สหชาตวาโรปิ ปจฺจยวาโรปิ นิสฺสยวาโรปิ สํสฏฺฐวาโรปิ สมฺปยุตฺตวาโรปิ ปฏิจฺจวารสทิสา วิตฺถาเรตพฺพา.)}

\csromanheader{2}{1. กุสลปท 7. ปญฺหาวาร}
\csromanheader{2}{\textbf{ปจฺจยจตุกฺก-เหตุ}}
\proseitem{51}{เหตุ เจว สเหตุโก จ กุสโล ธมฺโม เหตุสฺส เจว สเหตุกสฺส จ กุสลสฺส ธมฺมสฺส เหตุปจฺจเยน ปจฺจโย. (สํขิตฺตํ.)}

\proseitem{52}{เหตุยา ตีณิ, อารมฺมเณ นว, อธิปติยา นว, อนนฺตเร นว, สมนนฺตเร นว, สหชาเต นว, อญฺญมญฺเญ นว, นิสฺสเย นว, อุปนิสฺสเย นว, อาเสวเน นว, กมฺเม ตีณิ, อาหาเร ตีณิ ฯเปฯ อวิคเต นว. (สํขิตฺตํ.)}

\csromanheader{2}{\textbf{ปจฺจนียุทฺธาร}}
\proseitem{53}{เหตุ เจว สเหตุโก จ กุสโล ธมฺโม เหตุสฺส เจว สเหตุกสฺส จ กุสลสฺส ธมฺมสฺส อารมฺมณปจฺจเยน ปจฺจโย.\csromanspacer{} สหชาตปจฺจเยน ปจฺจโย.\csromanspacer{} อุปนิสฺสยปจฺจเยน ปจฺจโย. (สํขิตฺตํ.)}

\gambhira{ธมฺมานุโลม ทุกติกปฏฺฐานปาฬิ}
\proseitem{54}{\csromanfolio{338}นเหตุยา นว, น-อารมฺมเณ นว. (สํขิตฺตํ.)}

\csromanheader{2}{เหตุปจฺจยา น-อารมฺมเณ ตีณิ. (สํขิตฺตํ.)}
\csromanheader{2}{นเหตุปจฺจยา อารมฺมเณ นว. (สํขิตฺตํ.)}
\prose{(ยถา กุสลตฺติเก ปญฺหาวารสฺส อนุโลมมฺปิ ปจฺจนียมฺปิ อนุโลมปจฺจนียมฺปิ ปจฺจนียานุโลมมฺปิ คณิตํ, เอวํ คเณตพฺพํ.)}

\csromanheader{2}{2. อกุสลปท 1. ปฏิจฺจวาร}
\csromanheader{2}{\textbf{ปจฺจยจตุกฺก-เหตุ}}
\proseitem{55}{เหตุญฺเจว สเหตุกญฺจ อกุสลํ ธมฺมํ ปฏิจฺจ เหตุ เจว สเหตุโก จ อกุสโล ธมฺโม อุปฺปชฺชติ เหตุปจฺจยา. (สํขิตฺตํ.)}

\proseitem{56}{เหตุยา นว, อารมฺมเณ นว ฯเปฯ กมฺเม นว, อาหาเร นว ฯเปฯ อวิคเต นว. (สํขิตฺตํ.)}

\csromanheader{2}{\textbf{น-อธิปติ}}
\proseitem{57}{เหตุญฺเจว สเหตุกญฺจ อกุสลํ ธมฺมํ ปฏิจฺจ เหตุ เจว สเหตุโก จ อกุสโล ธมฺโม อุปฺปชฺชติ น-อธิปติปจฺจยา.\csromanspacer{} นว. (สํขิตฺตํ.)}

\vspace{\tipitakaparskip}
\csromancenter{\textbf{น-อธิปติ}ยา นว, นปุเรชาเต นว, นปจฺฉาชาเต นว, น-อาเสวเน นว, นกมฺเม ตีณิ, นวิปาเก นว, นวิปฺปยุตฺเต นว. (สํขิตฺตํ.)}
\csromanheader{2}{เหตุปจฺจยา น-อธิปติยา นว. (สํขิตฺตํ.)}
\csromanheader{2}{\textbf{น-อธิปติ}ปจฺจยา เหตุยา นว. (สํขิตฺตํ.)}
\prose{(สหชาตวาโรปิ ปจฺจยวาโรปิ นิสฺสยวาโรปิ สํสฏฺฐวาโรปิ สมฺปยุตฺตวาโรปิ ปฏิจฺจวารสทิสา วิตฺถาเรตพฺพา.)}

\csromanheader{2}{2. อกุสลปท 7. ปญฺหาวาร}
\csromanheader{2}{\textbf{ปจฺจยจตุกฺก-เหตุ}}
\proseitem{59}{เหตุ เจว สเหตุโก จ อกุสโล ธมฺโม เหตุสฺส เจว สเหตุกสฺส จ อกุสลสฺส ธมฺมสฺส เหตุปจฺจเยน ปจฺจโย. (สํขิตฺตํ.)}

\csromanheader{2}{4. เหตุสเหตุกทุก 1. กุสลตฺติก}
\proseitem{60}{\csromanfolio{339}เหตุยา ตีณิ, อารมฺมเณ นว ฯเปฯ อาเสวเน นว, กมฺเม ตีณิ, อาหาเร ตีณิ, อินฺทฺริเย ฌาเน มคฺเค ตีณิ, สมฺปยุตฺเต นว ฯเปฯ อวิคเต นว. (สํขิตฺตํ.)}

\csromanheader{2}{\textbf{ปจฺจนียุทฺธาร}}
\proseitem{61}{เหตุ เจว สเหตุโก จ อกุสโล ธมฺโม เหตุสฺส เจว สเหตุกสฺส จ อกุสลสฺส ธมฺมสฺส อารมฺมณปจฺจเยน ปจฺจโย.\csromanspacer{} สหชาตปจฺจเยน ปจฺจโย.\csromanspacer{} อุปนิสฺสยปจฺจเยน ปจฺจโย.}

\proseitem{62}{นเหตุยา นว, น-อารมฺมเณ นว ฯเปฯ โน-อวิคเต นว. (สํขิตฺตํ.)}

\csromanheader{2}{เหตุปจฺจยา น-อารมฺมเณ ตีณิ. (สํขิตฺตํ.)}
\csromanheader{2}{นเหตุปจฺจยา อารมฺมเณ นว. (สํขิตฺตํ.)}
\prose{(ยถา กุสลตฺติเก ปญฺหาวารสฺส อนุโลมมฺปิ ปจฺจนียมฺปิ อนุโลมปจฺจนียมฺปิ ปจฺจนียานุโลมมฺปิ คณิตํ, เอวํ คเณตพฺพํ.)}

\csromanheader{2}{3. อพฺยากตปท 1. ปฏิจฺจวาร}
\csromanheader{2}{\textbf{ปจฺจยจตุกฺก-เหตุ}}
\proseitem{63}{เหตุญฺเจว สเหตุกญฺจ อพฺยากตํ ธมฺมํ ปฏิจฺจ เหตุ เจว สเหตุโก จ อพฺยากโต ธมฺโม อุปฺปชฺชติ เหตุปจฺจยา. (สํขิตฺตํ.)}

\proseitem{64}{เหตุยา นว, อารมฺมเณ นว ฯเปฯ กมฺเม นว, วิปาเก นว ฯเปฯ อวิคเต นว. (สํขิตฺตํ.)}

\csromanheader{2}{\textbf{น-อธิปติ}}
\proseitem{65}{เหตุญฺเจว สเหตุกญฺจ อพฺยากตํ ธมฺมํ ปฏิจฺจ เหตุ เจว สเหตุโก จ อพฺยากโต ธมฺโม อุปฺปชฺชติ น-อธิปติปจฺจยา. (สํขิตฺตํ.)}

\vspace{\tipitakaparskip}
\csromancenter{\textbf{น-อธิปติ}ยา นว, นปุเรชาเต นว, นปจฺฉาชาเต นว, น-อาเสวเน นว, นกมฺเม ตีณิ, นวิปาเก นว, นวิปฺปยุตฺเต นว. (สํขิตฺตํ.)}
\csromanheader{2}{เหตุปจฺจยา น-อธิปติยา นว. (สํขิตฺตํ.)}
\csromanheader{2}{\textbf{น-อธิปติ}ปจฺจยา เหตุยา นว. (สํขิตฺตํ.)}
\prose{(สหชาตวาโรปิ ปจฺจยวาโรปิ นิสฺสยวาโรปิ สํสฏฺฐวาโรปิ สมฺปยุตฺตวาโรปิ ปฏิจฺจวารสทิสา วิตฺถาเรตพฺพา.)}

\csromanheader{2}{ธมฺมานุโลม ทุกติกปฏฺฐานปาฬิ \textbf{3. อพฺยากตปท 7. ปญฺหาวาร}}
\csromanheader{2}{\textbf{ปจฺจยจตุกฺก-เหตุ}}
\proseitem{67}{\csromanfolio{340}เหตุ เจว สเหตุโก จ อพฺยากโต ธมฺโม เหตุสฺส เจว สเหตุกสฺส จ อพฺยากตสฺส ธมฺมสฺส เหตุปจฺจเยน ปจฺจโย. (สํขิตฺตํ.)}

\proseitem{68}{เหตุยา ตีณิ, อารมฺมเณ นว, อธิปติยา นว, อนนฺตเร นว, สมนนฺตเร นว, สหชาเต นว, อญฺญมญฺเญ นว, นิสฺสเย นว, อุปนิสฺสเย นว, อาเสวเน นว, กมฺเม ตีณิ, วิปาเก นว, อาหาเร ตีณิ ฯเปฯ อวิคเต นว. (สํขิตฺตํ.)}

\csromanheader{2}{\textbf{ปจฺจนียุทฺธาร}}
\proseitem{69}{เหตุ เจว สเหตุโก จ อพฺยากโต ธมฺโม เหตุสฺส เจว สเหตุกสฺส จ อพฺยากตสฺส ธมฺมสฺส อารมฺมณปจฺจเยน ปจฺจโย.\csromanspacer{} สหชาตปจฺจเยน ปจฺจโย.\csromanspacer{} อุปนิสฺสยปจฺจเยน ปจฺจโย. (สํขิตฺตํ.)}

\proseitem{70}{นเหตุยา นว, น-อารมฺมเณ นว. (สํขิตฺตํ.)}

\csromanheader{2}{เหตุปจฺจยา น-อารมฺมเณ ตีณิ. (สํขิตฺตํ.)}
\csromanheader{2}{นเหตุปจฺจยา อารมฺมเณ นว. (สํขิตฺตํ.)}
\prose{(ยถา กุสลตฺติเก ปญฺหาวารสฺส อนุโลมมฺปิ ปจฺจนียมฺปิ อนุโลมปจฺจนียมฺปิ ปจฺจนียานุโลมมฺปิ คณิตํ, เอวํ คเณตพฺพํ.)}

\nitthitam{เหตุสเหตุกทุกกุสลตฺติกํ นิฏฺฐิตํ.}
\csromanmatikaanchor{mtk.60}
\csromantocmarkref{subparagraph}{5. เหตุ เหตุสมฺปยุตฺตทุก 1. กุสลตฺติก ...}{mtk.60}
\bookmark[dest={mtk.60},level=5]{5. เหตุ เหตุสมฺปยุตฺตทุก 1. กุสลตฺติก ...}
\csromanheader{6}{5. เหตุเหตุสมฺปยุตฺตทุก 1. กุสลตฺติก}
\csromanheader{2}{1. กุสลปท 1. ปฏิจฺจวาร}
\csromanheader{2}{\textbf{ปจฺจยจตุกฺก-เหตุ}}
\proseitem{71}{เหตุญฺเจว เหตุสมฺปยุตฺตญฺจ กุสลํ ธมฺมํ ปฏิจฺจ เหตุ เจว เหตุสมฺปยุตฺโต จ กุสโล ธมฺโม อุปฺปชฺชติ เหตุปจฺจยา.\csromanspacer{}\csromanspacer{}}

\vspace{\tipitakaparskip}
\csromancenter{เหตุสเหตุกสมฺปยุตฺตทุก 1.\csromanspacer{} กุสลตฺติก เหตุญฺเจว เหตุสมฺปยุตฺตญฺจ กุสลํ ธมฺมํ ปฏิจฺจ เหตุสมฺปยุตฺโต เจว น จ เหตุ กุสโล ธมฺโม อุปฺปชฺชติ เหตุปจฺจยา. (สํขิตฺตํ.)}
\proseitem{72}{\csromanfolio{341}เหตุยา นว, อารมฺมเณ นว ฯเปฯ อวิคเต นว. (สํขิตฺตํ.)}

\csromanheader{2}{\textbf{น-อธิปติ}}
\prosecont{เหตุญฺเจว เหตุสมฺปยุตฺตญฺจ กุสลํ ธมฺมํ ปฏิจฺจ เหตุ เจว เหตุสมฺปยุตฺโต จ กุสโล ธมฺโม อุปฺปชฺชติ น-อธิปติปจฺจยา. (สํขิตฺตํ.)}

\vspace{\tipitakaparskip}
\csromancenter{\textbf{น-อธิปติ}ยา นว, นปุเรชาเต นว, นปจฺฉาชาเต นว, น-อาเสวเน นว, นกมฺเม ตีณิ, นวิปาเก นว, นวิปฺปยุตฺเต นว. (สํขิตฺตํ.)}
\csromanheader{2}{เหตุปจฺจยา น-อธิปติยา นว. (สํขิตฺตํ.)}
\csromanheader{2}{อธิปติปจฺจยา เหตุยา นว. (สํขิตฺตํ.)}
\prose{(สหชาตวาโรปิ ปจฺจยวาโรปิ นิสฺสยวาโรปิ สํสฏฺฐวาโรปิ สมฺปยุตฺตวาโรปิ ปฏิจฺจวารสทิสา วิตฺถาเรตพฺพา.)}

\csromanheader{2}{1. กุสลปท 7. ปญฺหาวาร}
\csromanheader{2}{\textbf{ปจฺจยจตุกฺก-เหตุ}}
\proseitem{75}{เหตุ เจว เหตุสมฺปยุตฺโต จ กุสโล ธมฺโม เหตุสฺส เจว เหตุสมฺปยุตฺตสฺส จ กุสลสฺส ธมฺมสฺส เหตุปจฺจเยน ปจฺจโย. (สํขิตฺตํ.)}

\proseitem{76}{เหตุยา ตีณิ, อารมฺมเณ นว, อธิปติยา นว, อนนฺตเร นว ฯเปฯ อาเสวเน นว, กมฺเม ตีณิ, อาหาเร ตีณิ ฯเปฯ อวิคเต นว. (สํขิตฺตํ.)}

\csromanheader{2}{\textbf{ปจฺจนียุทฺธาร}}
\proseitem{77}{เหตุ เจว เหตุสมฺปยุตฺโต จ กุสโล ธมฺโม เหตุสฺส เจว เหตุสมฺปยุตฺตสฺส จ กุสลสฺส ธมฺมสฺส อารมฺมณปจฺจเยน ปจฺจโย.\csromanspacer{} สหชาตปจฺจเยน ปจฺจโย.\csromanspacer{} อุปนิสฺสยปจฺจเยน ปจฺจโย. (สํขิตฺตํ.)}

\gambhira{ธมฺมานุโลม ทุกติกปฏฺฐานปาฬิ}
\proseitem{78}{\csromanfolio{342}นเหตุยา นว, น-อารมฺมเณ นว. (สํขิตฺตํ.)}

\csromanheader{2}{เหตุปจฺจยา น-อารมฺมเณ ตีณิ. (สํขิตฺตํ.)}
\csromanheader{2}{นเหตุปจฺจยา อารมฺมเณ นว. (สํขิตฺตํ.)}
\prose{(ยถา กุสลตฺติเก ปญฺหาวารสฺส อนุโลมมฺปิ ปจฺจนียมฺปิ อนุโลมปจฺจนียมฺปิ ปจฺจนียานุโลมมฺปิ คณิตํ, เอวํ คเณตพฺพํ.)}

\csromanheader{2}{2. อกุสลปท 1. ปฏิจฺจวาร}
\csromanheader{2}{\textbf{ปจฺจยจตุกฺก-เหตุ}}
\proseitem{79}{เหตุญฺเจว เหตุสมฺปยุตฺตญฺจ อกุสลํ ธมฺมํ ปฏิจฺจ เหตุ เจว เหตุสมฺปยุตฺโต จ อกุสโล ธมฺโม อุปฺปชฺชติ เหตุปจฺจยา. (สํขิตฺตํ.)}

\proseitem{80}{เหตุยา นว, อารมฺมเณ นว ฯเปฯ อวิคเต นว. (สํขิตฺตํ.)}

\csromanheader{2}{\textbf{น-อธิปติ}}
\proseitem{81}{เหตุญฺเจว เหตุสมฺปยุตฺตญฺจ อกุสลํ ธมฺมํ ปฏิจฺจ เหตุ เจว เหตุสมฺปยุตฺโต จ อกุสโล ธมฺโม อุปฺปชฺชติ น-อธิปติปจฺจยา. (สํขิตฺตํ.)}

\vspace{\tipitakaparskip}
\csromancenter{\textbf{น-อธิปติ}ยา นว, นปุเรชาเต นว, นปจฺฉาชาเต นว, น-อาเสวเน นว, นกมฺเม ตีณิ, นวิปาเก นว, นวิปฺปยุตฺเต นว. (สํขิตฺตํ.)}
\csromanheader{2}{เหตุปจฺจยา น-อธิปติยา ตีณิ. (สํขิตฺตํ.)}
\csromanheader{2}{\textbf{น-อธิปติ}ปจฺจยา เหตุยา นว. (สํขิตฺตํ.)}
\prose{(สหชาตวาโรปิ ปฺจฺจยวาโรปิ นิสฺสยวาโรปิ สํสฏฺฐวาโรปิ สมฺปยุตฺตวาโรปิ ปฏิจฺจวารสทิสา วิตฺถาเรตพฺพา.)}

\csromanheader{2}{2. อกุสลปท 7. ปญฺหาวาร}
\csromanheader{2}{\textbf{ปจฺจยจตุกฺก-เหตุ}}
\proseitem{83}{เหตุ เจว เหตุสมฺปยุตฺโต จ อกุสโล ธมฺโม เหตุสฺส เจว เหตุสมฺปยุตฺตสฺส จ อกุสลสฺส ธมฺมสฺส เหตุปจฺจเยน ปจฺจโย. (สํขิตฺตํ.)}

\csromanheader{2}{5. เหตุสเหตุกสมฺปยุตฺตทุก 1. กุสลตฺติก}
\prose{\csromanfolio{343}เหตุยา ตีณิ, อารมฺมเณ นว, อธิปติยา นว ฯเปฯ กมฺเม ตีณิ, อาหาเร ตีณิ ฯเปฯ อวิคเต นว. (สํขิตฺตํ.)}

\csromanheader{2}{\textbf{ปจฺจนียุทฺธาร}}
\proseitem{84}{เหตุ เจว เหตุสมฺปยุตฺโต จ อกุสโล ธมฺโม เหตุสฺส เจว เหตุสมฺปยุตฺตสฺส จ อกุสลสฺส ธมฺมสฺส อารมฺมณปจฺจเยน ปจฺจโย.\csromanspacer{} สหชาตปจฺจเยน ปจฺจโย.\csromanspacer{} อุปนิสฺสยปจฺจเยน ปจฺจโย. (สํขิตฺตํ.)}

\proseitem{85}{นเหตุยา นว, น-อารมฺมเณ นว. (สํขิตฺตํ.)}

\csromanheader{2}{เหตุปจฺจยา น-อารมฺมเณ ตีณิ. (สํขิตฺตํ.)}
\csromanheader{2}{นเหตุปจฺจยา อารมฺมเณ นว. (สํขิตฺตํ.)}
\prose{(ยถา กุสลตฺติเก ปญฺหาวารสฺส อนุโลมมฺปิ ปจฺจนียมฺปิ อนุโลมปจฺจนียมฺปิ ปจฺจนียานุโลมมฺปิ คณิตํ, เอวํ คเณตพฺพํ.)}

\csromanheader{2}{3. อพฺยากตปท 1. ปฏิจฺจวาร}
\csromanheader{2}{\textbf{ปจฺจยจตุกฺก-เหตุ}}
\proseitem{86}{เหตุญฺเจว เหตุสมฺปยุตฺตญฺจ อพฺยากตํ ธมฺมํ ปฏิจฺจ เหตุ เจว เหตุสมฺปยุตฺโต จ อพฺยากโต ธมฺโม อุปฺปชฺชติ เหตุปจฺจยา. (สํขิตฺตํ.)}

\prose{เหตุยา นว, อารมฺมเณ นว ฯเปฯ กมฺเม นว, วิปาเก นว ฯเปฯ อวิคเต นว. (สํขิตฺตํ.)}

\csromanheader{2}{\textbf{น-อธิปติ}}
\proseitem{87}{เหตุญฺเจว เหตุสมฺปยุตฺตญฺจ อพฺยากตํ ธมฺมํ ปฏิจฺจ เหตุ เจว เหตุสมฺปยุตฺโต จ อพฺยากโต ธมฺโม อุปฺปชฺชติ น-อธิปติปจฺจยา. (สํขิตฺตํ.)}

\vspace{\tipitakaparskip}
\csromancenter{\textbf{น-อธิปติ}ยา นว, นปุเรชาเต นว, นปจฺฉาชาเต นว, น-อาเสวเน นว, นกมฺเม ตีณิ, นวิปาเก นว, นวิปฺปยุตฺเต นว. (สํขิตฺตํ.)}
\gambhira{ธมฺมานุโลม ทุกติกปฏฺฐานปาฬิ}
\csromanheader{2}{เหตุปจฺจยา น-อธิปติยา นว. (สํขิตฺตํ.)}
\csromanheader{2}{น-อธิปติปจฺจยา เหตุยา นว. (สํขิตฺตํ.)}
\prose{\csromanfolio{344}(สหชาตวาโรปิ ปจฺจยวาโรปิ นิสฺสยวาโรปิ สํสฏฺฐวาโรปิ สมฺปยุตฺตวาโรปิ ปฏิจฺจวารสทิสา วิตฺถาเรตพฺพา.)}

\csromanheader{2}{3. อพฺยากตปท 7. ปญฺหาวาร}
\csromanheader{2}{\textbf{ปจฺจยจตุกฺก-เหตุ}}
\proseitem{89}{เหตุ เจว เหตุสมฺปยุตฺโต จ อพฺยากโต ธมฺโม เหตุสฺส เจว เหตุสมฺปยุตฺตสฺส จ อพฺยากตสฺส ธมฺมสฺส เหตุปจฺจเยน ปจฺจโย. (สํขิตฺตํ.)}

\proseitem{90}{เหตุยา ตีณิ, อารมฺมเณ นว, อธิปติยา นว ฯเปฯ อาเสวเน นว, กมฺเม ตีณิ, วิปาเก นว, อาหาเร ตีณิ, อินฺทฺริเย นว, ฌาเน ตีณิ, มคฺเค สมฺปยุตฺเต นว ฯเปฯ อวิคเต นว. (สํขิตฺตํ.)}

\csromanheader{2}{\textbf{ปจฺจนียุทฺธาร}}
\proseitem{91}{เหตุ เจว เหตุสมฺปยุตฺโต จ อพฺยากโต ธมฺโม เหตุสฺส เจว เหตุสมฺปยุตฺตสฺส จ อพฺยากตสฺส ธมฺมสฺส อารมฺมณปจฺจเยน ปจฺจโย.\csromanspacer{} สหชาตปจฺจเยน ปจฺจโย.\csromanspacer{} อุปนิสฺสยปจฺจเยน ปจฺจโย. (สํขิตฺตํ.)}

\proseitem{92}{นเหตุยา นว, น-อารมฺมเณ นว. (สํขิตฺตํ.)}

\csromanheader{2}{เหตุปจฺจยา น-อารมฺมเณ ตีณิ. (สํขิตฺตํ.)}
\csromanheader{2}{นเหตุปจฺจยา อารมฺมเณ นว. (สํขิตฺตํ.)}
\prose{(ยถา กุสลตฺติเก ปญฺหาวารสฺส อนุโลมมฺปิ ปจฺจนียมฺปิ อนุโลมปจฺจนียมฺปิ ปจฺจนียานุโลมมฺปิ คณิตํ, เอวํ คเณตพฺพํ.)}

\nitthitam{เหตุเหตุสมฺปยุตฺตทุกกุสลตฺติกํ นิฏฺฐิตํ.}
\vspace{\tipitakaparskip}
\csromancenter{นเหตุสเหตุกทุก 1.\csromanspacer{} กุสลตฺติก \textbf{6.}\csromanspacer{}\textbf{ นเหตุสเหตุกทุก 1.}\csromanspacer{}\textbf{ กุสลตฺติก}}
\csromanheader{2}{1-2. กุลลากุสลปท 1-7. ปฏิจฺจาทิวาร}
\csromanheader{2}{\textbf{ปจฺจยจตุกฺก-เหตุ}}
\proseitem{93}{\csromanfolio{345}นเหตุํ สเหตุกํ กุสลํ ธมฺมํ ปฏิจฺจ นเหตุ สเหตุโก กุสโล ธมฺโม อุปฺปชฺชติ เหตุปจฺจยา. (สํขิตฺตํ.)}

\proseitem{94}{\textbf{เหตุ}ยา เอกํ, อารมฺมเณ เอกํ, (สพฺพตฺถ เอกํ.) อวิคเต เอกํ.}

\prose{น-อธิปติยา เอกํ, นปุเรชาเต เอกํ ฯเปฯ นวิปาเก เอกํ, นวิปฺปยุตฺเต เอกํ. (ปญฺหาวาเรปิ สพฺพตฺถ เอกํ.)}

\proseitem{95}{นเหตุํ สเหตุกํ อกุสลํ ธมฺมํ ปฏิจฺจ นเหตุ สเหตุโก อกุสโล ธมฺโม อุปฺปชฺชติ เหตุปจฺจยา. (สํขิตฺตํ.)}

\proseitem{96}{\textbf{เหตุ}ยา เอกํ, อารมฺมเณ เอกํ, (สพฺพตฺถ เอกํ.) อวิคเต เอกํ. (\textbf{เหตุ}ปจฺจโย นตฺถิ.\csromanspacer{} ปญฺหาวาเรปิ สพฺพตฺถ เอกํ, ปญฺหาวาเรปิ อิเมสํ ติณฺณนฺนํ เหตุปจฺจโย นตฺถิ.)}

\csromanmatikaanchor{mtk.61}
\csromanmatikaanchor{mtk.62}
\csromantocmarknopage{subparagraph}{6. นเหตุ สเหตุกทุก}{mtk.61}
\bookmark[dest={mtk.61},level=5]{6. นเหตุ สเหตุกทุก}
\csromantocmarkref{subparagraph}{1. กุสลตฺติก ...}{mtk.62}
\bookmark[dest={mtk.62},level=5]{1. กุสลตฺติก ...}
\csromanheader{6}{3. อพฺยากตปท 1. ปฏิจฺจวาร}
\csromanheader{2}{\textbf{ปจฺจยจตุกฺก-เหตุ}}
\proseitem{97}{นเหตุํ สเหตุกํ อพฺยากตํ ธมฺมํ ปฏิจฺจ นเหตุ สเหตุโก อพฺยากโต ธมฺโม อุปฺปชฺชติ เหตุปจฺจยา.\csromanspacer{}\csromanspacer{}นเหตุํ สเหตุกํ อพฺยากตํ ธมฺมํ ปฏิจฺจ นเหตุ อเหตุโก อพฺยากโต ธมฺโม อุปฺปชฺชติ เหตุปจฺจยา.\csromanspacer{}\csromanspacer{}นเหตุํ สเหตุกํ อพฺยากตํ ธมฺมํ ปฏิจฺจ นเหตุ สเหตุโก อพฺยากโต จ นเหตุ อเหตุโก อพฺยากโต จ ธมฺมา อุปฺปชฺชนฺติ เหตุปจฺจยา. (3)}

\prose{นเหตุํ อเหตุกํ อพฺยากตํ ธมฺมํ ปฏิจฺจ นเหตุ อเหตุโก อพฺยากโต ธมฺโม อุปฺปชฺชติ เหตุปจฺจยา.\csromanspacer{} ตีณิ.}

\gambhira{ธมฺมานุโลม ทุกติกปฏฺฐานปาฬิ}
\prose{\csromanfolio{346}\textbf{นเหตุ}ํ สเหตุกํ อพฺยากตญฺจ นเหตุํ อเหตุกํ อพฺยากตญฺจ ธมฺมํ ปฏิจฺจ นเหตุ สเหตุโก อพฺยากโต ธมฺโม อุปฺปชฺชติ เหตุปจฺจยา.\csromanspacer{} ตีณิ. (สํขิตฺตํ.)}

\proseitem{98}{เหตุยา นว, อารมฺมเณ จตฺตาริ, อธิปติยา ปญฺจ ฯเปฯ ปุเรชาเต อาเสวเน เทฺว, กมฺเม นว, วิปาเก นว ฯเปฯ อวิคเต นว. (สํขิตฺตํ.)}

\csromanheader{2}{\textbf{นเหตุ}}
\vspace{\tipitakaparskip}
\csromancenter{\textbf{นเหตุ}ํ อเหตุกํ อพฺยากตํ ธมฺมํ ปฏิจฺจ นเหตุ อเหตุโก อพฺยากโต ธมฺโม อุปฺปชฺชติ นเหตุปจฺจยา. (สํขิตฺตํ.)}
\csromancenter{\textbf{นเหตุ}ยา เอกํ, น-อารมฺมเณ ตีนิ, น-อธิปติยา นว, น-อนนฺตเร ตีณิ, นสมนนฺตเร ตีณิ, น-อญฺญมญฺเญ ตีณิ, น-อุปนิสฺสเย ตีณิ, นปุเรชาเต นว, นปจฺฉาชาเต นว, น-อาเสวเน นว, นกมฺเม เทฺว, นวิปาเก ปญฺจ, น-อาหาเร เอกํ, น-อินฺทฺริเย เอกํ, นฌาเน เอกํ, นมคฺเค เอกํ, นสมฺปยุตฺเต ตีณิ, นวิปฺปยุตฺเต เทฺว, โนนตฺถิยา ตีณิ, โนวิคเต ตีณิ. (สํขิตฺตํ.)}
\csromanheader{2}{เหตุปจฺจยา น-อารมฺมเณ ตีณิ. (สํขิตฺตํ.)}
\csromanheader{2}{\textbf{นเหตุ}ปจฺจยา อารมฺมเณ เอกํ. (สํขิตฺตํ.)}
\prose{(สหชาตวาโรปิ ปจฺจยวาโรปิ นิสฺสยวาโรปิ สํสฏฺฐวาโรปิ สมฺปยุตฺตวาโรปิ ปฏิจฺจวารสทิสา วิตฺถาเรตพฺพา.)}

\csromanheader{2}{3. อพฺยากตปท 7. ปญฺหาวาร}
\csromanheader{2}{\textbf{ปจฺจยจตุกฺก-อารมฺมณ}}
\vspace{\tipitakaparskip}
\csromancenter{\textbf{นเหตุ} สเหตุโก อพฺยากโต ธมฺโม นเหตุสฺส สเหตุกสฺส อพฺยากตสฺส ธมฺมสฺส อารมฺมณปจฺจเยน ปจฺจโย. (สํขิตฺตํ.)}
\proseitem{102}{อารมฺมเณ จตฺตาริ, อธิปติยา จตฺตาริ, อนนฺตเร จตฺตาริ, สมนนฺตเร จตฺตาริ, สหชาเต สตฺต, อญฺญมญฺเญ ฉ, นิสฺสเย สตฺต,}

\vspace{\tipitakaparskip}
\csromancenter{สปฺปจฺจยทุก 1.\csromanspacer{} กุสลตฺติก อุปนิสฺสเย จตฺตาริ, ปุเรชาเต เทฺว, ปจฺฉาชาเต เทฺว, อาเสวเน เทฺว, กมฺเม จตฺตาริ, วิปาเก จตฺตาริ, อาหาเร จตฺตาริ, อินฺทฺริเย จตฺตาริ, ฌาเน จตฺตาริ, มคฺเค ตีณิ, สมฺปยุตฺเต เทฺว, วิปฺปยุตฺเต ตีณิ, อตฺถิยา สตฺต, นตฺถิยา จตฺตาริ, วิคเต จตฺตาริ, อวิคเต สตฺต. (สํขิตฺตํ.)}
\csromanheader{2}{\textbf{ปจฺจนียุทฺธาร}}
\proseitem{103}{\csromanfolio{347}นเหตุ สเหตุโก อพฺยากโต ธมฺโม นเหตุสฺส สเหตุกสฺส อพฺยากตสฺส ธมฺมสฺส อารมฺมณปจฺจเยน ปจฺจโย.\csromanspacer{} สหชาตปจฺจเยน ปจฺจโย.\csromanspacer{} อุปนิสฺสยปจฺจเยน ปจฺจโย. (สํขิตฺตํ.)}

\proseitem{104}{นเหตุยา สตฺต, น-อารมฺมเณ สตฺต. (สํขิตฺตํ.)}

\csromanheader{2}{อารมฺมณปจฺจยา นเหตุยา จตฺตาริ. (สํขิตฺตํ.)}
\csromanheader{2}{นเหตุปจฺจยา อารมฺมเณ จตฺตาริ. (สํขิตฺตํ.)}
\prose{(ยถา กุสลตฺติเก ปญฺหาวารสฺส อนุโลมมฺปิ ปจฺจนียมฺปิ อนุโลมปจฺจนียมฺปิ ปจฺจนียานุโลมมฺปิ คณิตํ, เอวํ คเณตพฺพํ.)}

\nitthitam{นเหตุสเหตุกทุกกุสลตฺติกํ นิฏฺฐิตํ.}
\nitthitam{เหตุโคจฺฉกํ นิฏฺฐิตํ.}
\csromanmatikaanchor{mtk.63}
\csromanmatikaanchor{mtk.64}
\csromantocmarknopage{subparagraph}{7. สปฺปจฺจยทุก}{mtk.63}
\bookmark[dest={mtk.63},level=5]{7. สปฺปจฺจยทุก}
\csromantocmarkref{subparagraph}{1. กุสลตฺติก ...}{mtk.64}
\bookmark[dest={mtk.64},level=5]{1. กุสลตฺติก ...}
\csromanheader{6}{7. สปฺปจฺจยทุก 1. กุสลตฺติก}
\csromanheader{2}{1-2. กุสลากุสลปท 1-7. ปฏิจฺจาทิวาร}
\csromanheader{2}{\textbf{ปจฺจยจตุกฺก-เหตุ}}
\proseitem{1}{สปฺปจฺจยํ กุสลํ ธมฺมํ ปฏิจฺจ สปฺปจฺจโย กุสโล ธมฺโม อุปฺปชฺชติ เหตุปจฺจยา. (1) (สํขิตฺตํ.)}

\proseitem{2}{เหตุยา เอกํ, อารมฺมเณ เอกํ, อธิปติยา เอกํ ฯเปฯ อวิคเต เอกํ. (สํขิตฺตํ.)}

\csromanheader{2}{\textbf{น-อธิปติ}}
\proseitem{3}{สปฺปจฺจยํ กุสลํ ธมฺมํ ปฏิจฺจ สปฺปจฺจโย กุสโล ธมฺโม อุปฺปชฺชติ น-อธิปติปจฺจยา. (สํขิตฺตํ.)}

\gambhira{ธมฺมานุโลม ทุกติกปฏฺฐานปาฬิ}
\proseitem{4}{\csromanfolio{348}น-อธิปติยา เอกํ, นปุเรชาเต เอกํ, นปจฺฉาชาเต เอกํ, น-อาเสวเน เอกํ, นกมฺเม เอกํ, นวิปาเก เอกํ, นวิปฺปยุตฺเต เอกํ. (สํขิตฺตํ.)}

\prose{(สหชาตวาโรปิ ปจฺจยวาโรปิ นิสฺสยวาโรปิ สํสฏฺฐวาโรปิ สมฺปยุตฺตวาโรปิ ปฏิจฺจวารสทิสา วิตฺถาเรตพฺพา.)}

\csromanheader{2}{\textbf{ปจฺจยจตุกฺก-เหตุ}}
\proseitem{5}{สปฺปจฺจโย กุสโล ธมฺโม สปฺปจฺจยสฺส กุสลสฺส ธมฺมสฺส เหตุปจฺจเยน ปจฺจโย. (สํขิตฺตํ.)}

\proseitem{6}{เหตุยา เอกํ, อารมฺมเณ เอกํ ฯเปฯ อวิคเต เอกํ. (สํขิตฺตํ.)}

\proseitem{7}{สปฺปจฺจยํ อกุสลํ ธมฺมํ ปฏิจฺจ สปฺปจฺจโย อกุสโล ธมฺโม อุปฺปชฺชติ เหตุปจฺจยา. (1) (สํขิตฺตํ.)}

\proseitem{8}{เหตุยา เอกํ, อารมฺมเณ เอกํ ฯเปฯ อวิคเต เอกํ. (สหชาตวาเรปิ ฯเปฯ ปญฺหาวาเรปิ สพฺพตฺถ เอกํ.)}

\csromanheader{2}{3. อพฺยากตปท 1-7. ปฏิจฺจาทิวาร}
\csromanheader{2}{\textbf{ปจฺจยจตุกฺก-เหตุ}}
\proseitem{9}{สปฺปจฺจยํ อพฺยากตํ ธมฺมํ ปฏิจฺจ สปฺปจฺจโย อพฺยากโต ธมฺโม อุปฺปชฺชติ เหตุปจฺจยา. (สํขิตฺตํ.)}

\proseitem{10}{เหตุยา เอกํ, อารมฺมเณ เอกํ ฯเปฯ อวิคเต เอกํ. (สํขิตฺตํ.)}

\prose{นเหตุยา เอกํ, น-อารมฺมเณ เอกํ, น-อธิปติยา เอกํ. (สพฺพตฺถ เอกํ, สํขิตฺตํ.)}

\csromanheader{2}{\textbf{ปจฺจยจตุกฺก-เหตฺวาทิ}}
\proseitem{11}{สปฺปจฺจโย อพฺยากโต ธมฺโม สปฺปจฺจยสฺส อพฺยากตสฺส ธมฺมสฺส เหตุปจฺจเยน ปจฺจโย. (1)}

\prose{สปฺปจฺจโย อพฺยากโต ธมฺโม สปฺปจฺจยสฺส อพฺยากตสฺส ธมฺมสฺส อารมฺมณปจฺจเยน ปจฺจโย. (1)}

\csromanheader{2}{8. สงฺขตทุก 1. กุสลตฺติก}
\prose{\csromanfolio{349}อปฺปจฺจโย อพฺยากโต ธมฺโม สปฺปจฺจยสฺส อพฺยากตสฺส ธมฺมสฺส อารมฺมณปจฺจเยน ปจฺจโย. (1)}

\prose{สปฺปจฺจโย อพฺยากโต ธมฺโม สปฺปจฺจยสฺส อพฺยากตสฺส ธมฺมสฺส อุปนิสฺสยปจฺจเยน ปจฺจโย. (1)}

\prose{อปฺปจฺจโย อพฺยากโต ธมฺโม สปฺปจฺจยสฺส อพฺยากตสฺส ธมฺมสฺส อุปนิสฺสยปจฺจเยน ปจฺจโย.\csromanspacer{} อารมฺมณูปนิสฺสโย. (1) (สํขิตฺตํ.)}

\prose{เหตุยา เอกํ, อารมฺมเณ เทฺว, อธิปติยา เทฺว, อนนฺตเร เอกํ ฯเปฯ นิสฺสเย เอกํ, อุปนิสฺสเย เทฺว, ปุเรชาเต เอกํ, ปจฺฉาชาเต เอกํ, (สพฺพตฺถ เอกํ.) อวิคเต เอกํ. (สํขิตฺตํ.)}

\csromanheader{2}{\textbf{ปจฺจนียุทฺธาร}}
\proseitem{12}{สปฺปจฺจโย อพฺยากโต ธมฺโม สปฺปจฺจยสฺส อพฺยากตสฺส ธมฺมสฺส อารมฺมณปจฺจเยน ปจฺจโย.\csromanspacer{} สหชาตปจฺจเยน ปจฺจโย.\csromanspacer{} อุปนิสฺสยปจฺจเยน ปจฺจโย.}

\prose{อปฺปจฺจโย อพฺยากโต ธมฺโม สปฺปจฺจยสฺส อพฺยากตสฺส ธมฺมสฺส อารมฺมณปจฺจเยน ปจฺจโย.\csromanspacer{} อุปนิสฺสยปจฺจเยน ปจฺจโย. (สํขิตฺตํ.)}

\proseitem{13}{นเหตุยา เทฺว, น-อารมฺมเณ เอกํ. (สํขิตฺตํ.)}

\csromanheader{2}{เหตุปจฺจยา น-อารมฺมเณ เอกํ. (สํขิตฺตํ.)}
\csromanheader{2}{นเหตุปจฺจยา อารมฺมเณ เทฺว. (สํขิตฺตํ.)}
\nitthitam{สปฺปจฺจยทุกกุสลตฺติกํ นิฏฺฐิตํ.}
\csromanheader{2}{8. สงฺขตทุก 1. กุสลตฺติก}
\proseitem{14}{สงฺขตํ กุสลํ ธมฺมํ ปฏิจฺจ สงฺขโต กุสโล ธมฺโม อุปฺปชฺชติ เหตุปจฺจยา. (สํขิตฺตํ.\csromanspacer{} สปฺปจฺจยสทิสํ วิตฺถาเรตพฺพํ.)}

\nitthitam{สงฺขตทุกกุสลตฺติกํ นิฏฺฐิตํ.}
\csromanheader{2}{ธมฺมานุโลม ทุกติกปฏฺฐานปาฬิ \textbf{9. สนิทสฺสนทุก 1. กุสลตฺติก}}
\csromanheader{2}{1. กุสลปท 1. ปฏิจฺจวาร}
\csromanheader{2}{\textbf{ปจฺจยจตุกฺก-เหตุ}}
\proseitem{15}{\csromanfolio{350}อนิทสฺสนํ กุสลํ ธมฺมํ ปฏิจฺจ อนิทสฺสโน กุสโล ธมฺโม อุปฺปชฺชติ เหตุปจฺจยา. (สํขิตฺตํ.)}

\proseitem{16}{เหตุยา เอกํ, อารมฺมเณ เอกํ ฯเปฯ อวิคเต เอกํ. (สหชาตวาเรปิ ฯเปฯ ปญฺหาวาเรปิ สพฺพตฺถ เอกํ.)}

\csromanheader{2}{2. อกุสลปท-เหตุ}
\proseitem{17}{อนิทสฺสนํ อกุสลํ ธมฺมํ ปฏิจฺจ อนิทสฺสโน อกุสโล ธมฺโม อุปฺปชฺชติ เหตุปจฺจยา. (สํขิตฺตํ.)}

\proseitem{18}{เหตุยา เอกํ, อารมฺมเณ เอกํ ฯเปฯ อวิคเต เอกํ. (สหชาตวาเรปิ ฯเปฯ ปญฺหาวาเรปิ สพฺพตฺถ เอกํ.)}

\csromanheader{2}{3. อพฺยากตปท-เหตุ}
\proseitem{19}{อนิทสฺสนํ อพฺยากตํ ธมฺมํ ปฏิจฺจ อนิทสฺสโน อพฺยากโต ธมฺโม อุปฺปชฺชติ เหตุปจฺจยา. (สํขิตฺตํ.)}

\proseitem{20}{เหตุยา ตีณิ, อารมฺมเณ เอกํ, อธิปติยา ตีณิ ฯเปฯ อวิคเต ตีณิ. (สํขิตฺตํ.)}

\prose{นเหตุยา ตีณิ, (สพฺพตฺถ ตีณิ.) นกมฺเม ตีณิ, นวิปาเก ตีณิ, น-อาหาเร ตีณิ ฯเปฯ โนวิคเต ตีณิ. (สํขิตฺตํ.)}

\prose{(สหชาตวาโรปิ ปจฺจยวาโรปิ นิสฺสยวาโรปิ สํสฏฺฐวาโรปิ สมฺปยุตฺตวาโรปิ ปฏิจฺจวารสทิสา วิตฺถาเรตพฺพา.)}

\csromanheader{2}{3. อพฺยากตปท 7. ปญฺหาวาร}
\csromanheader{2}{\textbf{ปจฺจยจตุกฺก-เหตฺวาทิ}}
\proseitem{21}{อนิทสฺสโน อพฺยากโต ธมฺโม อนิทสฺสนสฺส อพฺยากตสฺส ธมฺมสฺส เหตุปจฺจเยน ปจฺจโย.\csromanspacer{} ตีณิ.}

\csromanheader{2}{10. สปฺปฏิฆทุก 1. กุสลตฺติก}
\prose{\csromanfolio{351}สนิทสฺสโน อพฺยากโต ธมฺโม อนิทสฺสนสฺส อพฺยากตสฺส ธมฺมสฺส อารมฺมณปจฺจเยน ปจฺจโย. (1)}

\prose{อนิทสฺสโน อพฺยากโต ธมฺโม อนิทสฺสนสฺส อพฺยากตสฺส ธมฺมสฺส อารมฺมณปจฺจเยน ปจฺจโย. (1)}

\prose{สนิทสฺสโน อพฺยากโต ธมฺโม อนิทสฺสนสฺส อพฺยากตสฺส ธมฺมสฺส อธิปติปจฺจเยน ปจฺจโย.\csromanspacer{} อารมฺมณาธิปติ. (1)}

\prose{อนิทสฺสโน อพฺยากโต ธมฺโม อนิทสฺสนสฺส อพฺยากตสฺส ธมฺมสฺส อธิปติปจฺจเยน ปจฺจโย.\csromanspacer{} อารมฺมณาธิปติ, สหชาตาธิปติ.\csromanspacer{}\csromanspacer{}อนิทสฺสโน อพฺยากโต ธมฺโม สนิทสฺสนสฺส อพฺยากตสฺส จ อนิทสฺสนสฺส อพฺยากตสฺส จ ธมฺมสฺส อธิปติปจฺจเยน ปจฺจโย.\csromanspacer{} สหชาตาธิปติ. (2)}

\prose{สนิทสฺสโน อพฺยากโต ธมฺโม อนิทสฺสนสฺส อพฺยากตสฺส ธมฺมสฺส อุปนิสฺสยปจฺจเยน ปจฺจโย.\csromanspacer{} ปกตูปนิสฺสโย. (1)}

\prose{อนิทสฺสโน อพฺยากโต ธมฺโม อนิทสฺสนสฺส อพฺยากตสฺส ธมฺมสฺส อุปนิสฺสยปจฺจเยน ปจฺจโย.\csromanspacer{} อารมฺมณูปนิสฺสโย, อนนฺตรูปนิสฺสโย, ปกตูปนิสฺสโย ฯเปฯ. (1) (สํขิตฺตํ.)}

\proseitem{22}{ปุเรชาเต ตีณิ, ปจฺฉาชาเต ตีณิ, อาเสวเน เอกํ, กมฺเม ตีณิ, วิปาเก ตีณิ, (สพฺพตฺถ ตีณิ.) สมฺปยุตฺเต เอกํ, วิปฺปยุตฺเต ตีณิ, อตฺถิยา ปญฺจ, นตฺถิยา เอกํ ฯเปฯ อวิคเต ปญฺจ. (สํขิตฺตํ.)}

\nitthitam{สนิทสฺสนทุกกุสลตฺติกํ นิฏฺฐิตํ.}
\vspace{\tipitakaparskip}
\csromancenter{(อปฺปจฺจยมฺปิ อสงฺขตมฺปิ สนิทสฺสนมฺปิ น ลพฺภติ.)}
\csromanheader{2}{10. สปฺปฏิฆทุก 1. กุสลตฺติก}
\csromanheader{2}{1-2. กุสลากุสลปท 1. ปฏิจฺจวาร}
\csromanheader{2}{\textbf{ปจฺจยจตุกฺก-เหตุ}}
\proseitem{23}{อปฺปฏิฆํ กุสลํ ธมฺมํ ปฏิจฺจ อปฺปฏิโฆ กุสโล ธมฺโม อุปฺปชฺชติ เหตุปจฺจยา. (สํขิตฺตํ.)}

\gambhira{ธมฺมานุโลม ทุกติกปฏฺฐานปาฬิ}
\proseitem{24}{\csromanfolio{352}เหตุยา เอกํ, อารมฺมเณ เอกํ ฯเปฯ วิปฺปยุตฺเต เอกํ ฯเปฯ อวิคเต เอกํ. (สํขิตฺตํ.) (สหชาตวาเรปิ ฯเปฯ ปญฺหาวาเรปิ สพฺพตฺถ เอกํ.)}

\proseitem{25}{อปฺปฏิฆํ อกุสลํ ธมฺมํ ปฏิจฺจ อปฺปฏิโฆ อกุสโล ธมฺโม อุปฺปชฺชติ เหตุปจฺจยา. (สํขิตฺตํ.)}

\csromanheader{2}{เหตุยา เอกํ ฯเปฯ อวิคเต เอกํ. (ปญฺหาวาเรปิ สพฺพตฺถ เอกํ.)}
\csromanheader{2}{3. อพฺยากตปท 1. ปฏิจฺจวาร}
\csromanheader{2}{\textbf{ปจฺจยจตุกฺก-เหตุ}}
\proseitem{26}{สปฺปฏิฆํ อพฺยากตํ ธมฺมํ ปฏิจฺจ สปฺปฏิโฆ อพฺยากโต ธมฺโม อุปฺปชฺชติ เหตุปจฺจยา.\csromanspacer{}\csromanspacer{}สปฺปฏิฆํ อพฺยากตํ ธมฺมํ ปฏิจฺจ อปฺปฏิโฆ อพฺยากโต ธมฺโม อุปฺปชฺชติ เหตุปจฺจยา.\csromanspacer{}\csromanspacer{}สปฺปฏิฆํ อพฺยากตํ ธมฺมํ ปฏิจฺจ สปฺปฏิโฆ อพฺยากโต จ อปฺปฏิโฆ อพฺยากโต จ ธมฺมา อุปฺปชฺชนฺติ เหตุปจฺจยา.\csromanspacer{} นว.}

\prose{อปฺปฏิฆํ อพฺยากตํ ธมฺมํ ปฏิจฺจ อปฺปฏิโฆ อพฺยากโต ธมฺโม อุปฺปชฺชติ อารมฺมณปจฺจยา. (สํขิตฺตํ.)}

\proseitem{27}{เหตุยา นว, อารมฺมเณ เอกํ, อธิปติยา นว ฯเปฯ อญฺญมญฺเญ ฉ, นิสฺสเย นว, อุปนิสฺสเย เอกํ, ปุเรชาเต อาเสวเน เอกํ, กมฺเม นว, วิปาเก นว ฯเปฯ อวิคเต นว. (สํขิตฺตํ.)}

\csromanheader{2}{\textbf{นเหตุ}}
\proseitem{28}{สปฺปฏิฆํ อพฺยากตํ ธมฺมํ ปฏิจฺจ สปฺปฏิโฆ อพฺยากโต ธมฺโม อุปฺปชฺชติ นเหตุปจฺจยา. (สํขิตฺตํ.)}

\vspace{\tipitakaparskip}
\csromancenter{\textbf{นเหตุ}ยา นว ฯเปฯ โนวิคเต นว. (สพฺพตฺถ นว.) (สํขิตฺตํ.)}
\csromanheader{2}{3. อพฺยากตปท 7. ปญฺหาวาร}
\csromanheader{2}{\textbf{ปจฺจยจตุกฺก-เหตุ}}
\proseitem{30}{อปฺปฏิโฆ อพฺยากโต ธมฺโม อปฺปฏิฆสฺส อพฺยากตสฺส ธมฺมสฺส เหตุปจฺจเยน ปจฺจโย. (3) (สํขิตฺตํ.)}

\csromanheader{2}{11. รูปีทุก 1. กุสลตฺติก}
\proseitem{31}{\csromanfolio{353}เหตุยา ตีณิ, อารมฺมเณ เทฺว, อธิปติยา ตีณิ, อนนฺตเร เอกํ, สมนนฺตเร เอกํ, สหชาเต นว, อญฺญมญฺเญ ฉ, นิสฺสเย นว, อุปนิสฺสเย เทฺว, ปุเรชาเต ตีณิ, ปจฺฉาชาเต ตีณิ, อาเสวเน เอกํ, กมฺเม ตีณิ, วิปาเก ตีณิ, อาหาเร ตีณิ, อินฺทฺริเย ปญฺจ, ฌาเน ตีณิ, มคฺเค ตีณิ, สมฺปยุตฺเต เอกํ, วิปฺปยุตฺเต จตฺตาริ, อตฺถิยา นว, นตฺถิยา เอกํ, วิคเต เอกํ, อวิคเต นว.}

\csromanheader{2}{\textbf{ปจฺจนียุทฺธาร}}
\proseitem{32}{สปฺปฏิโฆ อพฺยากโต ธมฺโม สปฺปฏิฆสฺส อพฺยากตสฺส ธมฺมสฺส สหชาตปจฺจเยน ปจฺจโย. (สํขิตฺตํ.)}

\proseitem{33}{นเหตุยา นว, น-อารมฺมเณ นว. (สํขิตฺตํ.)}

\csromanheader{2}{เหตุปจฺจยา น-อารมฺมเณ ตีณิ. (สํขิตฺตํ.)}
\csromanheader{2}{นเหตุปจฺจยา อารมฺมเณ เทฺว. (สํขิตฺตํ.)}
\prose{(ยถา กุสลตฺติเก ปญฺหาวารสฺส อนุโลมมฺปิ ปจฺจนียมฺปิ อนุโลมปจฺจนียมฺปิ ปจฺจนียานุโลมมฺปิ คณิตํ, เอวํ คเณตพฺพํ.)}

\nitthitam{สปฺปฏิฆทุกกุสลตฺติกํ นิฏฺฐิตํ.}
\csromanheader{2}{11. รูปีทุก 1. กุสลตฺติก}
\csromanheader{2}{1-2. กุสลากุสลปท 1-7. ปฏิจฺจาทิวาร}
\csromanheader{2}{\textbf{ปจฺจยจตุกฺก-เหตุ}}
\proseitem{34}{อรูปิํ กุสลํ ธมฺมํ ปฏิจฺจ อรูปี กุสโล ธมฺโม อุปฺปชฺชติ เหตุปจฺจยา. (สํขิตฺตํ.)}

\proseitem{35}{เหตุยา เอกํ, อารมฺมเณ เอกํ. (สํขิตฺตํ.\csromanspacer{} สหชาตวาเรปิ ฯเปฯ ปญฺหาวาเรปิ สพฺพตฺถ เอกํ.)}

\proseitem{36}{อรูปิํ อกุสลํ ธมฺมํ ปฏิจฺจ อรูปี อกุสโล ธมฺโม อุปฺปชฺชติ เหตุปจฺจยา. (สํขิตฺตํ.)}

\proseitem{37}{เหตุยา เอกํ, อารมฺมเณ เอกํ ฯเปฯ อวิคเต เอกํ. (สํขิตฺตํ.\csromanspacer{} สหชาตวาเรปิ ฯเปฯ ปญฺหาวาเรปิ สพฺพตฺถ เอกํ.)}

\csromanheader{2}{ธมฺมานุโลม ทุกติกปฏฺฐานปาฬิ \textbf{3. อพฺยากตปท 1-7. ปฏิจฺจาทิวาร}}
\csromanheader{2}{\textbf{ปจฺจยจตุกฺก-เหตุ อารมฺมณ}}
\proseitem{38}{\csromanfolio{354}รูปิํ อพฺยากตํ ธมฺมํ ปฏิจฺจ รูปี อพฺยากโต ธมฺโม อุปฺปชฺชติ เหตุปจฺจยา.\csromanspacer{} ตีณิ.}

\prose{อรูปิํ อพฺยากตํ ธมฺมํ ปฏิจฺจ อรูปี อพฺยากโต ธมฺโม อุปฺปชฺชติ เหตุปจฺจยา.\csromanspacer{} ตีณิ.}

\prose{รูปิํ อพฺยากตญฺจ อรูปิํ อพฺยากตญฺจ ธมฺมํ ปฏิจฺจ รูปี อพฺยากโต ธมฺโม อุปฺปชฺชติ เหตุปจฺจยา.\csromanspacer{} ตีณิ.}

\prose{รูปิํ อพฺยากตํ ธมฺมํ ปฏิจฺจ อรูปี อพฺยากโต ธมฺโม อุปฺปชฺชติ อารมฺมณปจฺจยา. (สํขิตฺตํ.)}

\proseitem{39}{เหตุยา นว, อารมฺมเณ ตีณิ, อธิปติยา ปญฺจ, อนนฺตเร ตีณิ, สมนนฺตเร ตีณิ ฯเปฯ อญฺญมญฺเญ ฉ ฯเปฯ ปุเรชาเต เอกํ, อาเสวเน เอกํ, กมฺเม นว, วิปาเก นว ฯเปฯ อวิคเต นว. (สํขิตฺตํ.)}

\proseitem{40}{นเหตุยา นว, น-อารมฺมเณ ตีณิ, น-อธิปติยา นว ฯเปฯ นกมฺเม เทฺว, นวิปาเก ปญฺจ, น-อาหาเร เอกํ, น-อินฺทฺริเย เอกํ, นฌาเน เทฺว, นมคฺเค นว, นสมฺปยุตฺเต ตีณิ, นวิปฺปยุตฺเต เทฺว, โนนตฺถิยา ตีณิ, โนวิคเต ตีณิ. (สํขิตฺตํ.\csromanspacer{} ปจฺจนียํ.)}

\csromanheader{2}{\textbf{ปจฺจยจตุกฺก-เหตุ}}
\proseitem{41}{อรูปี อพฺยากโต ธมฺโม อรูปิสฺส อพฺยากตสฺส ธมฺมสฺส เหตุปจฺจเยน ปจฺจโย. (สํขิตฺตํ.)}

\proseitem{42}{เหตุยา ตีณิ, อารมฺมเณ เทฺว, อธิปติยา ตีณิ, อนนฺตเร เอกํ, สมนนฺตเร เอกํ, สหชาเต สตฺต, อญฺญมญฺเญ ฉ, นิสฺสเย สตฺต, อุปนิสฺสเย เทฺว, ปุเรชาเต เอกํ, ปจฺฉาชาเต เอกํ, อาเสวเน เอกํ, กมฺเม ตีณิ, วิปาเก ตีณิ, อาหาเร จตฺตาริ, อินฺทฺริเย ฉ, ฌาเน ตีณิ, มคฺเค ตีณิ, สมฺปยุตฺเต เอกํ, วิปฺปยุตฺเต เทฺว ฯเปฯ อวิคเต สตฺต. (สํขิตฺตํ.)}

\prose{นเหตุยา สตฺต, น-อารมฺมเณ สตฺต. (สํขิตฺตํ.)}

\csromanheader{2}{12. โลกิยทุก 1. กุสลตฺติก}
\csromanheader{2}{เหตุปจฺจยา น-อารมฺมเณ ตีณิ. (สํขิตฺตํ.)}
\csromanheader{2}{นเหตุปจฺจยา อารมฺมเณ เทฺว. (สํขิตฺตํ.)}
\prose{\csromanfolio{355}(ยถา กุสลตฺติเก ปญฺหาวารสฺส อนุโลมมฺปิ ปจฺจนียมฺปิ อนุโลมปจฺจนียมฺปิ ปจฺจนียานุโลมมฺปิ คณิตํ, เอวํ คเณตพฺพํ.)}

\nitthitam{รูปีทุกกุสลตฺติกํ นิฏฺฐิตํ.}
\csromanheader{2}{12. โลกิยทุก 1. กุสลตฺติก}
\csromanheader{2}{1. กุสลปท 1-7. ปฏิจฺจาทิวาร}
\csromanheader{2}{\textbf{ปจฺจยจตุกฺก-เหตุ}}
\proseitem{43}{โลกิยํ กุสลํ ธมฺมํ ปฏิจฺจ โลกิโย กุสโล ธมฺโม อุปฺปชฺชติ เหตุปจฺจยา. (1)}

\prose{โลกุตฺตรํ กุสลํ ธมฺมํ ปฏิจฺจ โลกุตฺตโร กุสโล ธมฺโม อุปฺปชฺชติ เหตุปจฺจยา. (1)}

\prose{โลกิยํ กุสลํ ธมฺมํ ปฏิจฺจ โลกิโย กุสโล ธมฺโม อุปฺปชฺชติ อารมฺมณปจฺจยา. (1)}

\prose{โลกุตฺตรํ กุสลํ ธมฺมํ ปฏิจฺจ โลกุตฺตโร กุสโล ธมฺโม อุปฺปชฺชติ อารมฺมณปจฺจยา. (1) (สํขิตฺตํ.)}

\proseitem{44}{เหตุยา เทฺว, อารมฺมเณ เทฺว, อธิปติยา เทฺว ฯเปฯ กมฺเม เทฺว ฯเปฯ อวิคเต เทฺว. (สํขิตฺตํ.\csromanspacer{} อนุโลมํ.)}

\proseitem{45}{น-อธิปติยา เทฺว ฯเปฯ น-อาเสวเน เอกํ ฯเปฯ นวิปฺปยุตฺเต เทฺว. (สํขิตฺตํ.)}

\prose{(สหชาตวาโรปิ ปจฺจยวาโรปิ นิสฺสยวาโรปิ สํสฏฺฐวาโรปิ สมฺปยุตฺตวาโรปิ ปฏิจฺจวารสทิสา วิตฺถาเรตพฺพา.)}

\csromanheader{2}{\textbf{ปจฺจยจตุกฺก-เหตฺวาทิ}}
\proseitem{46}{โลกิโย กุสโล ธมฺโม โลกิยสฺส กุสลสฺส ธมฺมสฺส เหตุปจฺจเยน ปจฺจโย. (1)}

\gambhira{ธมฺมานุโลม ทุกติกปฏฺฐานปาฬิ}
\prose{\csromanfolio{356}โลกุตฺตโร กุสโล ธมฺโม โลกุตฺตรสฺส กุสลสฺส ธมฺมสฺส เหตุปจฺจเยน ปจฺจโย. (1)}

\prose{โลกิโย กุสโล ธมฺโม โลกิยสฺส กุสลสฺส ธมฺมสฺส อารมฺมณปจฺจเยน ปจฺจโย. (1)}

\prose{โลกุตฺตโร กุสโล ธมฺโม โลกิยสฺส กุสลสฺส ธมฺมสฺส อารมฺมณปจฺจเยน ปจฺจโย. (1)}

\prose{โลกิโย กุสโล ธมฺโม โลกิยสฺส กุสลสฺส ธมฺมสฺส อธิปติปจฺจเยน ปจฺจโย.\csromanspacer{} อารมฺมณาธิปติ, สหชาตาธิปติ. (1)}

\prose{โลกุตฺตโร กุสโล ธมฺโม โลกุตฺตรสฺส กุสลสฺส ธมฺมสฺส อธิปติปจฺจเยน ปจฺจโย.\csromanspacer{} สหชาตาธิปติ.\csromanspacer{}\csromanspacer{}โลกุตฺตโร กุสโล ธมฺโม โลกิยสฺส กุสลสฺส ธมฺมสฺส อธิปติปจฺจเยน ปจฺจโย.\csromanspacer{} อารมฺมณาธิปติ. (2) (สํขิตฺตํ.)}

\proseitem{47}{เหตุยา เทฺว, อารมฺมเณ เทฺว, อธิปติยา ตีณิ, อนนฺตเร เทฺว, สมนนฺตเร เทฺว, สหชาเต เทฺว ฯเปฯ อุปนิสฺสเย จตฺตาริ, อาเสวเน เทฺว, กมฺเม เทฺว, อาหาเร เทฺว ฯเปฯ อวิคเต เทฺว. (สํขิตฺตํ.)}

\csromanheader{2}{2. อกุสลปท 1-7. ปฏิจฺจาทิวาร}
\csromanheader{2}{\textbf{ปจฺจยจตุกฺก-เหตุ}}
\proseitem{48}{โลกิยํ อกุสลํ ธมฺมํ ปฏิจฺจ โลกิโย อกุสโล ธมฺโม อุปฺปชฺชติ เหตุปจฺจยา. (สํขิตฺตํ.)}

\proseitem{49}{เหตุยา เอกํ, อารมฺมเณ เอกํ ฯเปฯ อวิคเต เอกํ. (สํขิตฺตํ.)}

\prose{(สหชาตวาโรปิ ฯเปฯ ปญฺหาวาโรปิ สพฺพตฺถ สทิสา.)}

\csromanheader{2}{3. อพฺยากตปท 1-7. ปฏิจฺจาทิวาร}
\csromanheader{2}{\textbf{ปจฺจยจตุกฺก-เหตุ}}
\proseitem{50}{โลกิยํ อพฺยากตํ ธมฺมํ ปฏิจฺจ โลกิโย อพฺยากโต ธมฺโม อุปฺปชฺชติ เหตุปจฺจยา. (1)}

\csromanheader{2}{12. โลกิยทุก 1. กุสลตฺติก}
\prose{\csromanfolio{357}โลกุตฺตรํ อพฺยากตํ ธมฺมํ ปฏิจฺจ โลกุตฺตโร อพฺยากโต ธมฺโม อุปฺปชฺชติ เหตุปจฺจยา.\csromanspacer{} ตีณิ.}

\prose{โลกิยํ อพฺยากตญฺจ โลกุตฺตรํ อพฺยากตญฺจ ธมฺมํ ปฏิจฺจ โลกิโย อพฺยากโต ธมฺโม อุปฺปชฺชติ เหตุปจฺจยา. (1) (สํขิตฺตํ.)}

\proseitem{51}{เหตุยา ปญฺจ, อารมฺมเณ เทฺว ฯเปฯ อาเสวเน เอกํ, กมฺเม ปญฺจ, วิปาเก ปญฺจ ฯเปฯ อวิคเต ปญฺจ. (สํขิตฺตํ.\csromanspacer{} อนุโลมํ.)}

\proseitem{52}{นเหตุยา เอกํ, น-อารมฺมเณ ตีณิ, น-อธิปติยา เทฺว ฯเปฯ นปุเรชาเต จตฺตาริ, นปจฺฉาชาเต ปญฺจ ฯเปฯ นกมฺเม เอกํ, นวิปาเก เอกํ, น-อาหาเร เอกํ ฯเปฯ นมคฺเค เอกํ, นสมฺปยุตฺเต ตีณิ, นวิปฺปยุตฺเต เทฺว, โนนตฺถิยา ตีณิ, โนวิคเต ตีณิ. (สํขิตฺตํ.\csromanspacer{} ปจฺจนียํ.)}

\prose{(สหชาตวาโรปิ ปจฺจยวาโรปิ นิสฺสยวาโรปิ สํสฏฺฐวาโรปิ สมฺปยุตฺตวาโรปิ ปฏิจฺจวารสทิสา วิตฺถาเรตพฺพา.)}

\csromanheader{2}{\textbf{ปจฺจยจตุกฺก-เหตุ}}
\proseitem{53}{โลกิโย อพฺยากโต ธมฺโม โลกิยสฺส อพฺยากตสฺส ธมฺมสฺส เหตุปจฺจเยน ปจฺจโย. (1)}

\prose{โลกุตฺตโร อพฺยากโต ธมฺโม โลกุตฺตรสฺส อพฺยากตสฺส ธมฺมสฺส เหตุปจฺจเยน ปจฺจโย.\csromanspacer{} ตีณิ. (สํขิตฺตํ.)}

\proseitem{54}{เหตุยา จตฺตาริ, อารมฺมเณ ตีณิ, อธิปติยา จตฺตาริ, อนนฺตเร จตฺตาริ, สหชาเต ปญฺจ, อญฺญมญฺเญ เทฺว, นิสฺสเย สตฺต, อุปนิสฺสเย จตฺตาริ, ปุเรชาเต เทฺว, ปจฺฉาชาเต เทฺว, อาเสวเน เอกํ, กมฺเม จตฺตาริ, วิปาเก จตฺตาริ, อาหาเร จตฺตาริ ฯเปฯ มคฺเค จตฺตาริ, สมฺปยุตฺเต เทฺว, วิปฺปยุตฺเต ตีณิ, อตฺถิยา สตฺต ฯเปฯ อวิคเต สตฺต. (สํขิตฺตํ.)}

\prose{(ยถา กุสลตฺติเก ปญฺหาวารสฺส อนุโลมมฺปิ ปจฺจนียมฺปิ อนุโลมปจฺจนียมฺปิ ปจฺจนียานุโลมมฺปิ คณิตํ, เอวํ คเณตพฺพํ.)}

\nitthitam{โลกิยทุกกุสลตฺติกํ นิฏฺฐิตํ.}
\csromanheader{2}{ธมฺมานุโลม ทุกติกปฏฺฐานปาฬิ \textbf{13. เกนจิวิญฺเญยฺยทุก 1. กุสลตฺติก}}
\csromanheader{2}{1. กุสลปท 1-7. ปฏิจฺจาทิวาร}
\csromanheader{2}{\textbf{ปจฺจยจตุกฺก-เหตุ}}
\proseitem{55}{\csromanfolio{358}เกนจิ วิญฺเญยฺยํ กุสลํ ธมฺมํ ปฏิจฺจ เกนจิ วิญฺเญยฺโย กุสโล ธมฺโม อุปฺปชฺชติ เหตุปจฺจยา.\csromanspacer{}\csromanspacer{}เกนจิ วิญฺเญยฺยํ กุสลํ ธมฺมํ ปฏิจฺจ เกนจิ นวิญฺเญยฺโย กุสโล ธมฺโม อุปฺปชฺชติ เหตุปจฺจยา.\csromanspacer{}\csromanspacer{}เกนจิ วิญฺเญยฺยํ กุสลํ ธมฺมํ ปฏิจฺจ เกนจิ วิญฺเญยฺโย กุสโล จ เกนจิ นวิญฺเญยฺโย กุสโล จ ธมฺมา อุปฺปชฺชนฺติ เหตุปจฺจยา. (3) (สํขิตฺตํ.)}

\proseitem{56}{เหตุยา นว, อารมฺมเณ นว ฯเปฯ อวิคเต นว. (สํขิตฺตํ.)}

\csromanheader{2}{\textbf{ปจฺจนีย-น-อธิปติ}}
\proseitem{57}{เกนจิ วิญฺเญยฺยํ กุสลํ ธมฺมํ ปฏิจฺจ เกนจิ วิญฺเญยฺโย กุสโล ธมฺโม อุปฺปชฺชติ น-อธิปติปจฺจยา. (สํขิตฺตํ.)}

\proseitem{58}{น-อธิปติยา นว, นปุเรชาเต นว, นปจฺฉาชาเต นว, น-อาเสวเน นว, นกมฺเม นว, นวิปาเก นว ฯเปฯ นวิปฺปยุตฺเต นว. (สํขิตฺตํ.)}

\prose{(สหชาตวาโรปิ ปจฺจยวาโรปิ นิสฺสยวาโรปิ สํสฏฺฐวาโรปิ สมฺปยุตฺตวาโรปิ ปฏิจฺจวารสทิสา วิตฺถาเรตพฺพา.)}

\csromanheader{2}{\textbf{ปจฺจยจตุกฺก-เหตุ}}
\proseitem{59}{เกนจิ วิญฺเญยฺโย กุสโล ธมฺโม เกนจิ วิญฺเญยฺยสฺส กุสลสฺส ธมฺมสฺส เหตุปจฺจเยน ปจฺจโย. (สํขิตฺตํ.)}

\proseitem{60}{เหตุยา นว, อารมฺมเณ นว, อธิปติยา นว, อนนฺตเร นว, สมนนฺตเร นว ฯเปฯ กมฺเม นว ฯเปฯ อวิคเต นว. (สํขิตฺตํ.)}

\prose{นเหตุยา นว, น-อารมฺมเณ นว. (สํขิตฺตํ.)}

\csromanheader{2}{เหตุปจฺจยา น-อารมฺมเณ นว. (สํขิตฺตํ.)}
\csromanheader{2}{นเหตุปจฺจยา อารมฺมเณ นว. (สํขิตฺตํ.)}
\csromanmatikaanchor{mtk.65}
\csromanmatikaanchor{mtk.66}
\csromantocmarknopage{subparagraph}{14. อาสวทุก}{mtk.65}
\bookmark[dest={mtk.65},level=5]{14. อาสวทุก}
\csromantocmarkref{subparagraph}{1. กุสลตฺติก ...}{mtk.66}
\bookmark[dest={mtk.66},level=5]{1. กุสลตฺติก ...}
\csromanheader{6}{14. อาสวทุก 1. กุสลตฺติก}
\prose{\csromanfolio{359}(ยถา กุสลตฺติเก ปญฺหาวารสฺส อนุโลมมฺปิ ปจฺจนียมฺปิ อนุโลมปจฺจนียมฺปิ ปจฺจนียานุโลมมฺปิ คณิตํ, เอวํ คเณตพฺพํ.\csromanspacer{} เกนจิ วิญฺเญยฺยํ อกุสลมฺปิ เกนจิ วิญฺเญยฺยํ อพฺยากตมฺปิ เกนจิ วิญฺเญยฺยกุสลสทิสํ วิตฺถาเรตพฺพํ.)}

\nitthitam{เกนจิ วิญฺเญยฺยทุกกุสลตฺติกํ นิฏฺฐิตํ.}
\nitthitam{จูฬนฺตรทุกํ นิฏฺฐิตํ.}
\csromanheader{2}{14. อาสวทุก 1. กุสลตฺติก}
\csromanheader{2}{1. กุสลปท 1-7. ปฏิจฺจาทิวาร-เหตุ}
\proseitem{1}{โน-อสวํ กุสลํ ธมฺมํ ปฏิจฺจ โน-อสโว กุสโล ธมฺโม อุปฺปชฺชติ เหตุปจฺจยา. (สํขิตฺตํ.)}

\proseitem{2}{เหตุยา เอกํ, อารมฺมเณ เอกํ ฯเปฯ กมฺเม เอกํ ฯเปฯ อวิคเต เอกํ. (สํขิตฺตํ.)}

\prose{(สหชาตวาเรปิ ฯเปฯ ปญฺหาวาเรปิ สพฺพตฺถ เอกํ.)}

\csromanheader{2}{2. อกุสลปท-เหตุ}
\proseitem{3}{อาสวํ อกุสลํ ธมฺมํ ปฏิจฺจ อาสโว อกุสโล ธมฺโม อุปฺปชฺชติ เหตุปจฺจยา. (สํขิตฺตํ.)}

\proseitem{4}{เหตุยา นว, อารมฺมเณ นว ฯเปฯ อวิคเต นว. (สํขิตฺตํ.)}

\prose{นเหตุยา เอกํ, น-อธิปติยา นว ฯเปฯ นกมฺเม ตีณิ, นวิปาเก นว, นวิปฺปยุตฺเต นว. (สํขิตฺตํ.)}

\prose{(สหชาตวาโรปิ ฯเปฯ สมฺปยุตฺตวาโรปิ สพฺพตฺถ วิตฺถาเรตพฺโพ.)}

\proseitem{5}{อาสโว อกุสโล ธมฺโม อาสวสฺส อกุสลสฺส ธมฺมสฺส เหตุปจฺจเยน ปจฺจโย.}

\proseitem{6}{เหตุยา สตฺต, อารมฺมเณ นว, อธิปติยา นว, (มชฺเฌ ติณฺณํ สหชาตาธิปติ ลพฺภติ.) อนนฺตเร นว ฯเปฯ อุปนิสฺสเย นว,}

\vspace{\tipitakaparskip}
\csromancenter{ธมฺมานุโลม ทุกติกปฏฺฐานปาฬิ อาเสวเน นว, กมฺเม ตีณิ ฯเปฯ ฌาเน ตีณิ, มคฺเค นว ฯเปฯ อวิคเต นว. (สํขิตฺตํ.)}
\csromanheader{2}{3. อพฺยากตปท-เหตุ}
\proseitem{7}{\csromanfolio{360}โน-อสวํ อพฺยากตํ ธมฺมํ ปฏิจฺจ โน-อสโว อพฺยากโต ธมฺโม อุปฺปชฺชติ เหตุปจฺจยา. (สํขิตฺตํ.)}

\proseitem{8}{เหตุยา เอกํ, อารมฺมเณ เอกํ ฯเปฯ อวิคเต เอกํ. (สํขิตฺตํ.)}

\prose{(สหชาตวาเรปิ ปจฺจยวาเรปิ ฯเปฯ ปญฺหาวาเรปิ สพฺพตฺถ เอกํ.)}

\nitthitam{อาสวทุกกุสลตฺติกํ นิฏฺฐิตํ.}
\csromanheader{2}{15. สาสวทุก 1. กุสลตฺติก}
\csromanheader{2}{1. กุสลปท-เหตุ}
\proseitem{9}{สาสวํ กุสลํ ธมฺมํ ปฏิจฺจ สาสโว กุสโล ธมฺโม อุปฺปชฺชติ เหตุปจฺจยา. (1)}

\prose{อนาสวํ กุสลํ ธมฺมํ ปฏิจฺจ อนาสโว กุสโล ธมฺโม อุปฺปชฺชติ เหตุปจฺจยา. (1) (สํขิตฺตํ.)}

\proseitem{10}{เหตุยา เทฺว, อารมฺมเณ เทฺว, อธิปติยา เทฺว ฯเปฯ อวิคเต เทฺว. (สํขิตฺตํ.\csromanspacer{} อนุโลมํ.)}

\proseitem{11}{น-อธิปติยา เทฺว ฯเปฯ น-อาเสวเน เอกํ ฯเปฯ นวิปฺปยุตฺเต เทฺว. (สํขิตฺตํ.\csromanspacer{} ปจฺจนียํ.) (สหชาตวาเรปิ ฯเปฯ สมฺปยุตฺตวาเรปิ สพฺพตฺถ เทฺว.)}

\csromanheader{2}{\textbf{ปญฺหาวาร-เหตุ}}
\proseitem{12}{สาสโว กุสโล ธมฺโม สาสวสฺส กุสลสฺส ธมฺมสฺส เหตุปจฺจเยน ปจฺจโย. (1)}

\prose{อนาสโว กุสโล ธมฺโม อนาสวสฺส กุสลสฺส ธมฺมสฺส เหตุปจฺจเยน ปจฺจโย. (1) (สํขิตฺตํ.)}

\csromanheader{2}{15. สาสวทุก 1. กุสลตฺติก}
\proseitem{13}{\csromanfolio{361}เหตุยา เทฺว, อารมฺมเณ เทฺว, อธิปติยา ตีณิ, อนนฺตเร เทฺว ฯเปฯ นิสฺสเย เทฺว, อุปนิสฺสเย จตฺตาริ, อาเสวเน เทฺว ฯเปฯ อวิคเต เทฺว. (สํขิตฺตํ.)}

\prose{(ยถา กุสลตฺติเก ปญฺหาวารสฺส อนุโลมมฺปิ ปจฺจนียมฺปิ อนุโลมปจฺจนียมฺปิ ปจฺจนียานุโลมมฺปิ คณิตํ, เอวํ คเณตพฺพํ.)}

\csromanheader{2}{2. อกุสลปท-เหตุ}
\proseitem{14}{สาสวํ อกุสลํ ธมฺมํ ปฏิจฺจ สาสโว อกุสโล ธมฺโม อุปฺปชฺชติ เหตุปจฺจยา. (สํขิตฺตํ.)}

\proseitem{15}{เหตุยา เอกํ, อารมฺมเณ เอกํ ฯเปฯ อวิคเต เอกํ. (สํขิตฺตํ.)}

\prose{(สหชาตวาเรปิ ฯเปฯ ปญฺหาวาเรปิ สพฺพตฺถ เอกํ.)}

\csromanheader{2}{3. อพฺยากตปท-เหตุ}
\proseitem{16}{สาสวํ อพฺยากตํ ธมฺมํ ปฏิจฺจ สาสโว อพฺยากโต ธมฺโม อุปฺปชฺชติ เหตุปจฺจยา. (1)}

\prose{อนาสวํ อพฺยากตํ ธมฺมํ ปฏิจฺจ อนาสโว อพฺยากโต ธมฺโม อุปฺปชฺชติ เหตุปจฺจยา.\csromanspacer{}\csromanspacer{}อนาสวํ อพฺยากตํ ธมฺมํ ปฏิจฺจ สาสโว อพฺยากโต ธมฺโม อุปฺปชฺชติ เหตุปจฺจยา.\csromanspacer{}\csromanspacer{}อนาสวํ อพฺยากตํ ธมฺมํ ปฏิจฺจ สาสโว อพฺยากโต จ อนาสโว อพฺยากโต จ ธมฺมา อุปฺปชฺชนฺติ เหตุปจฺจยา. (3)}

\prose{สาสวํ อพฺยากตญฺจ อนาสวํ อพฺยากตญฺจ ธมฺมํ ปฏิจฺจ สาสโว อพฺยากโต ธมฺโม อุปฺปชฺชติ เหตุปจฺจยา. (1) (สํขิตฺตํ.)}

\proseitem{17}{เหตุยา ปญฺจ, อารมฺมเณ เทฺว, อธิปติยา ปญฺจ ฯเปฯ อาเสวเน เอกํ ฯเปฯ วิปาเก ปญฺจ ฯเปฯ อวิคเต ปญฺจ. (สํขิตฺตํ.)}

\prose{นเหตุยา เอกํ, น-อารมฺมเณ ตีณิ, น-อธิปติยา เทฺว, นปุเรชาเต จตฺตาริ, นปจฺฉาชาเต น-อาเสวเน ปญฺจ, นกมฺเม ฯเปฯ นมคฺเค เอกํ ฯเปฯ นวิปฺปยุตฺเต เทฺว ฯเปฯ โนวิคเต ตีณิ. (สํขิตฺตํ.)}

\prose{(สหชาตวาโรปิ ฯเปฯ สมฺปยุตฺตวาโรปิ สพฺพตฺถ วิตฺถาเรตพฺโพ.)}

\csromanheader{2}{ธมฺมานุโลม ทุกติกปฏฺฐานปาฬิ \textbf{ปญฺหาวาร-เหตุ}}
\proseitem{18}{\csromanfolio{362}สาสโว อพฺยากโต ธมฺโม สาสวสฺส อพฺยากตสฺส ธมฺมสฺส เหตุปจฺจเยน ปจฺจโย. (1)}

\prose{อนาสโว อพฺยากโต ธมฺโม อนาสวสฺส อพฺยากตสฺส ธมฺมสฺส เหตุปจฺจเยน ปจฺจโย. (1) (สํขิตฺตํ.)}

\proseitem{19}{เหตุยา จตฺตาริ, อารมฺมเณ ตีณิ, อธิปติยา อนนฺตเร สมนนฺตเร จตฺตาริ, สหชาเต ปญฺจ, อญฺญมญฺเญ เทฺว, นิสฺสเย สตฺต, อุปนิสฺสเย สตฺตาริ, ปุเรชาเต ปจฺฉาชาเต เทฺว, อาเสวเน เอกํ, กมฺเม ฯเปฯ มคฺเค จตฺตาริ, สมฺปยุตฺเต เทฺว, วิปฺปยุตฺเต ตีณิ ฯเปฯ อวิคเต สตฺต. (สํขิตฺตํ.)}

\prose{(ยถา กุสลตฺติเก ปญฺหาวารสฺส อนุโลมมฺปิ ปจฺจนียมฺปิ อนุโลมปจฺจนียมฺปิ ปจฺจนียานุโลมมฺปิ คณิตํ, เอวํ คเณตพฺพํ.)}

\nitthitam{สาสวทุกกุสลตฺติกํ นิฏฺฐิตํ.}
\csromanheader{2}{16. อาสวสมฺปยุตฺตทุก 1. กุสลตฺติก}
\csromanheader{2}{1. กุสลปท 1-7. ปฏิจฺจาทิวาร-เหตุ}
\proseitem{20}{อาสววิปฺปยุตฺตํ กุสลํ ธมฺมํ ปฏิจฺจ อาสววิปฺปยุตฺโต กุสโล ธมฺโม อุปฺปชฺชติ เหตุปจฺจยา. (สํขิตฺตํ.)}

\proseitem{21}{เหตุยา เอกํ, อารมฺมเณ เอกํ ฯเปฯ อวิคเต เอกํ. (สํขิตฺตํ.\csromanspacer{} สหชาตวาเรปิ ฯเปฯ ปญฺหาวาเรปิ สพฺพตฺถ เอกํ.)}

\csromanheader{2}{2. อกุสลปท-เหตุ}
\proseitem{22}{อาสวสมฺปยุตฺตํ อกุสลํ ธมฺมํ ปฏิจฺจ อาสวสมฺปยุตฺโต อกุสโล ธมฺโม อุปฺปชฺชติ เหตุปจฺจยา.\csromanspacer{} ตีณิ.}

\prose{อาสววิปฺปยุตฺตํ อกุสลํ ธมฺมํ ปฏิจฺจ อาสวสมฺปยุตฺโต อกุสโล ธมฺโม อุปฺปชฺชติ เหตุปจฺจยา. (1)}

\csromanheader{2}{16. อาสวสมฺปยุตฺตทุก 1. กุสลตฺติก}
\prose{\csromanfolio{363}อาสวสมฺปยุตฺตํ อกุสลญฺจ อาสววิปฺปยุตฺตํ อกุสลญฺจ ธมฺมํ ปฏิจฺจ อาสวสมฺปยุตฺโต อกุสโล ธมฺโม อุปฺปชฺชติ เหตุปจฺจยา. (1) (สํขิตฺตํ.)}

\proseitem{23}{เหตุยา ปญฺจ, อารมฺมเณ ปญฺจ ฯเปฯ อวิคเต ปญฺจ. (สํขิตฺตํ.)}

\prose{(สหชาตวาเรปิ ฯเปฯ สมฺปยุตฺตวาเรปิ สพฺพตฺถ ปญฺจ.)}

\prose{นเหตุยา เอกํ, น-อธิปติยา ปญฺจ, นปุเรชาเต ปญฺจ ฯเปฯ นกมฺเม ตีณิ, นวิปาเก ปญฺจ, นวิปฺปยุตฺเต ปญฺจ. (สํขิตฺตํ.)}

\proseitem{24}{อาสวสมฺปยุตฺโต อกุสโล ธมฺโม อาสวสมฺปยุตฺตสฺส อกุสลสฺส ธมฺมสฺส เหตุปจฺจเยน ปจฺจโย.\csromanspacer{} ตีณิ.}

\prose{อาสววิปฺปยุตฺโต อกุสโล ธมฺโม อาสวสมฺปยุตฺตสฺส อกุสลสฺส ธมฺมสฺส เหตุปจฺจเยน ปจฺจโย. (1)}

\prose{อาสวสมฺปยุตฺโต อกุสโล จ อาสววิปฺปยุตฺโต อกุสโล จ ธมฺมา อาสวสมฺปยุตฺตสฺส อกุสลสฺส ธมฺมสฺส เหตุปจฺจเยน ปจฺจโย. (1) (สํขิตฺตํ.)}

\proseitem{25}{เหตุยา ปญฺจ, อารมฺมเณ นว, อธิปติยา เอกํ, อนนฺตเร สมนนฺตเร นว, สหชาเต อญฺญมญฺเญ นิสฺสเย ปญฺจ, อุปนิสฺสเย อาเสวเน นว, กมฺเม อาหาเร อินฺทฺริเย ฌาเน มคฺเค ตีณิ, สมฺปยุตฺเต ปญฺจ ฯเปฯ อวิคเต ปญฺจ. (สํขิตฺตํ.)}

\prose{(ยถา กุสลตฺติเก ปญฺหาวารสฺส อนุโลมมฺปิ ปจฺจนียมฺปิ อนุโลมปจฺจนียมฺปิ ปจฺจนียานุโลมมฺปิ คณิตํ, เอวํ คเณตพฺพํ.)}

\csromanheader{2}{3. อพฺยากตปท-เหตุ}
\proseitem{26}{อาสววิปฺปยุตฺตํ อพฺยากตํ ธมฺมํ ปฏิจฺจ อาสววิปฺปยุตฺโต อพฺยากโต ธมฺโม อุปฺปชฺชติ เหตุปจฺจยา. (สํขิตฺตํ.)}

\proseitem{27}{เหตุยา เอกํ, อารมฺมเณ เอกํ ฯเปฯ อวิคเต เอกํ. (สํขิตฺตํ.)}

\prose{นเหตุยา เอกํ, น-อารมฺมเณ เอกํ, น-อธิปติยา เอกํ ฯเปฯ โนวิคเต เอกํ. (สํขิตฺตํ.\csromanspacer{} สหชาตวาเรปิ ฯเปฯ สมฺปยุตฺตวาเรปิ สพฺพตฺถ เอกํ.)}

\gambhira{ธมฺมานุโลม ทุกติกปฏฺฐานปาฬิ}
\proseitem{28}{\csromanfolio{364}อาสววิปฺปยุตฺโต อพฺยากโต ธมฺโม อาสววิปฺปยุตฺตสฺส อพฺยากตสฺส ธมฺมสฺส เหตุปจฺจเยน ปจฺจโย. (สํขิตฺตํ.)}

\proseitem{29}{เหตุยา เอกํ, อารมฺมเณ เอกํ ฯเปฯ อวิคเต เอกํ. (สํขิตฺตํ.)}

\prose{(ยถา กุสลตฺติเก ปญฺหาวารสฺส อนุโลมมฺปิ ปจฺจนียมฺปิ อนุโลมปจฺจนียมฺปิ ปจฺจนียานุโลมมฺปิ คณิตํ, เอวํ คเณตพฺพํ.)}

\nitthitam{อาสวสมฺปยุตฺตทุกกุสลตฺติกํ นิฏฺฐิตํ.}
\csromanheader{2}{17. อาสวสาสวทุก 1. กุสลตฺติก}
\csromanheader{2}{1. กุสลปท 1-7. ปฏิจฺจาทิวาร-เหตุ}
\proseitem{30}{สาสวญฺเจว โน จ อาสวํ กุสลํ ธมฺมํ ปฏิจฺจ สาสโว เจว โน จ อาสโว กุสโล ธมฺโม อุปฺปชฺชติ เหตุปจฺจยา. (สํขิตฺตํ.)}

\proseitem{31}{เหตุยา เอกํ, อารมฺมเณ เอกํ ฯเปฯ อวิคเต เอกํ. (สํขิตฺตํ.)}

\prose{(สหชาตวาเรปิ ฯเปฯ ปญฺหาวาเรปิ สพฺพตฺถ เอกํ.)}

\csromanheader{2}{2. อกุสลปท-เหตุ}
\proseitem{32}{อาสวญฺเจว สาสวญฺจ อกุสลํ ธมฺมํ ปฏิจฺจ อาสโว เจว สาสโว จ อกุสโล ธมฺโม อุปฺปชฺชติ เหตุปจฺจยา.\csromanspacer{} ตีณิ.}

\prose{สาสวญฺเจว โน จ อาสวํ อกุสลํ ธมฺมํ ปฏิจฺจ สาสโว เจว โน จ อาสโว อกุสโล ธมฺโม อุปฺปชฺชติ เหตุปจฺจยา.\csromanspacer{} ตีณิ.}

\prose{อาสวญฺเจว สาสวญฺจ อกุสลญฺจ สาสวญฺเจว โน จ อาสวํ อกุสลญฺจ ธมฺมํ ปฏิจฺจ อาสโว เจว สาสโว จ อกุสโล ธมฺโม อุปฺปชฺชติ เหตุปจฺจยา.\csromanspacer{} ตีณิ. (สํขิตฺตํ.)}

\proseitem{33}{เหตุยา นว, อารมฺมเณ นว ฯเปฯ อวิคเต นว. (สํขิตฺตํ.)}

\csromanheader{2}{\textbf{ปจฺจนีย-นเหตุ}}
\proseitem{34}{สาสวญฺเจว โน จ อาสวํ อกุสลํ ธมฺมํ ปฏิจฺจ อาสโว เจว สาสโว จ อกุสโล ธมฺโม อุปฺปชฺชติ นเหตุปจฺจยา. (สํขิตฺตํ.)}

\csromanheader{2}{17. อาสวสาสวทุก 1. กุสลตฺติก}
\proseitem{35}{\csromanfolio{365}นเหตุยา เอกํ, น-อธิปติยา นว, นปุเรชาเต นว, นปจฺฉาชาเต นว, น-อาเสวเน นว, นกมฺเม ตีณิ, นวิปาเก นว, นวิปฺปยุตฺเต นว. (สํขิตฺตํ.)}

\prose{(สหชาตวาโรปิ ปจฺจยวาโรปิ นิสฺสยวาโรปิ สํสฏฺฐวาโรปิ สมฺปยุตฺตวาโรปิ วิตฺถาเรตพฺโพ.)}

\csromanheader{2}{7. ปญฺหาวาร-เหตุ}
\proseitem{36}{อาสโว เจว สาสโวจ อกุสโล ธมฺโม อาสวสฺส เจว สาสวสฺส จ อกุสลสฺส ธมฺมสฺส เหตุปจฺจเยน ปจฺจโย. (สํขิตฺตํ.)}

\proseitem{37}{เหตุยา สตฺต, อารมฺมเณ นว, อธิปติยา นว ฯเปฯ อุปนิสฺสเย อาเสวเน นว, กมฺเม อาหาเร อินฺทฺริเย ฌาเน ตีณิ, มคฺเค สมฺปยุตฺเต นว ฯเปฯ อวิคเต นว. (สํขิตฺตํ.)}

\csromanheader{2}{\textbf{ปจฺจนียุทฺธาร}}
\proseitem{38}{อาสโว เจว สาสโว จ อกุสโล ธมฺโม อาสวสฺส เจว สาสวสฺส จ อกุสลสฺส ธมฺมสฺส อารมฺมณปจฺจเยน ปจฺจโย.\csromanspacer{} สหชาตปจฺจเยน ปจฺจโย.\csromanspacer{} อุปนิสฺสยปจฺจเยน ปจฺจโย. (สํขิตฺตํ.)}

\proseitem{39}{นเหตุยา นว, น-อารมฺมเณ นว. (สํขิตฺตํ.)}

\csromanheader{2}{เหตุปจฺจยา น-อารมฺมเณ สตฺต. (สํขิตฺตํ.)}
\csromanheader{2}{นเหตุปจฺจยา อารมฺมเณ นว. (สํขิตฺตํ.)}
\prose{(ยถา กุสลตฺติเก ปญฺหาวารสฺส อนุโลมมฺปิ ปจฺจนียมฺปิ อนุโลมปจฺจนียมฺปิ ปจฺจนียานุโลมมฺปิ คณิตํ, เอวํ คเณตพฺพํ.)}

\csromanheader{2}{3. อพฺยากตปท-เหตุ}
\proseitem{40}{สาสวญฺเจว โน จ อาสวํ อพฺยากตํ ธมฺมํ ปฏิจฺจ สาสโว เจว โน จ อาสโว อพฺยากโต ธมฺโม อุปฺปชฺชติ เหตุปจฺจยา. (สํขิตฺตํ.)}

\proseitem{41}{เหตุยา เอกํ, อารมฺมเณ เอกํ ฯเปฯ อวิคเต เอกํ. (สํขิตฺตํ.)}

\prose{(สหชาตวาเรปิ ฯเปฯ ปญฺหาวาเรปิ สพฺพตฺถ เอกํ.)}

\nitthitam{อาสวสาสวทุกกุสลตฺติกํ นิฏฺฐิตํ.}
\csromanheader{2}{ธมฺมานุโลม ทุกติกปฏฺฐานปาฬิ \textbf{18. อาสว-อาสวสมฺปยุตฺตทุก 1. กุสลตฺติก}}
\csromanheader{2}{1. ปฏิจฺจวาร-อกุสลปท-เหตุ}
\proseitem{42}{\csromanfolio{366}อาสวญฺเจว อาสวสมฺปยุตฺตญฺจ อกุสลํ ธมฺมํ ปฏิจฺจ อาสโว เจว อาสวสมฺปยุตฺโต จ อกุสโล ธมฺโม อุปฺปชฺชติ เหตุปจฺจยา.\csromanspacer{} ตีณิ.}

\prose{อาสวสมฺปยุตฺตญฺเจว โน จ อาสวํ อกุสลํ ธมฺมํ ปฏิจฺจ อาสวสมฺปยุตฺโต เจว โน จ อาสโว อกุสโล ธมฺโม อุปฺปชฺชติ เหตุปจฺจยา.\csromanspacer{} ตีณิ.}

\prose{อาสวญฺเจว อาสวสมฺปยุตฺตญฺจ อกุสลญฺจ อาสวสมฺปยุตฺตญฺเจว โน จ อาสวํ อกุสลญฺจ ธมฺมํ ปฏิจฺจ อาสโว เจว อาสวสมฺปยุตฺโต จ อกุสโล ธมฺโม อุปฺปชฺชติ เหตุปจฺจยา.\csromanspacer{} ตีณิ. (สํขิตฺตํ.)}

\proseitem{43}{เหตุยา นว, อารมฺมเณ นว ฯเปฯ อวิคเต นว. (สํขิตฺตํ.)}

\csromanheader{2}{\textbf{ปจฺจนีย-น-อธิปติ}}
\proseitem{44}{อาสวญฺเจว อาสวสมฺปยุตฺตญฺจ อกุสลํ ธมฺมํ ปฏิจฺจ อาสโว เจว อาสวสมฺปยุตฺโต จ อกุสโล ธมฺโม อุปฺปชฺชติ น-อธิปติปจฺจยา.}

\proseitem{45}{น-อธิปติยา นว, นปุเรชาเต นว, นปจฺฉาชาเต นว, น-อาเสวเน นว, นกมฺเม ตีณิ, นวิปาเก นว, นวิปฺปยุตฺเต นว. (สํขิตฺตํ.)}

\prose{(สหชาตวาเรปิ สมฺปยุตฺตวาเรปิ สพฺพตฺถ นว.)}

\csromanheader{2}{7. ปญฺหาวาร ปจฺจยจตุกฺก-เหตุ}
\proseitem{46}{อาสโว เจว อาสวสมฺปยุตฺโต จ อกุสโล ธมฺโม อาสวสฺส เจว อาสวสมฺปยุตฺตสฺส จ อกุสลสฺส ธมฺมสฺส เหตุปจฺจเยน ปจฺจโย. (สํขิตฺตํ.)}

\proseitem{47}{เหตุยา จตฺตาริ, อารมฺมเณ นว, อธิปติยา นว ฯเปฯ อวิคเต นว. (สํขิตฺตํ.)}

\csromanheader{2}{19. อาสววิปฺปยุตฺตสาสวทุก 1. กุสลตฺติก}
\prose{\csromanfolio{367}นเหตุยา นว, น-อารมฺมเณ นว. (สํขิตฺตํ.)}

\csromanheader{2}{เหตุปจฺจยา น-อารมฺมเณ จตฺตาริ. (สํขิตฺตํ.)}
\csromanheader{2}{นเหตุปจฺจยา อารมฺมเณ นว. (สํขิตฺตํ.)}
\prose{(ยถา กุสลตฺติเก ปญฺหาวารสฺส อนุโลมมฺปิ ปจฺจนียมฺปิ อนุโลมปจฺจนียมฺปิ ปจฺจนียานุโลมมฺปิ คณิตํ, เอวํ คเณตพฺพํ.)}

\nitthitam{อาสว-อาสวสมฺปยุตฺตทุกํ นิฏฺฐิตํ.}
\csromanheader{2}{19. อาสววิปฺปยุตฺตสาสวทุก 1. กุสลตฺติก}
\csromanheader{2}{\textbf{กุสลปท-เหตุ}}
\proseitem{48}{อาสววิปฺปยุตฺตํ สาสวํ กุสลํ ธมฺมํ ปฏิจฺจ อาสววิปฺปยุตฺโต สาสโว กุสโล ธมฺโม อุปฺปชฺชติ เหตุปจฺจยา.\csromanspacer{}\csromanspacer{}อาสววิปฺปยุตฺตํ อนาสวํ กุสลํ ธมฺมํ ปฏิจฺจ อาสววิปฺปยุตฺโต อนาสโว กุสโล ธมฺโม อุปฺปชฺชติ เหตุปจฺจยา. (2) (สํขิตฺตํ.)}

\proseitem{49}{เหตุยา เทฺว, อารมฺมเณ เทฺว ฯเปฯ อวิคเต เทฺว. (สํขิตฺตํ.)}

\prose{(สหชาตวาโรปิ ฯเปฯ ปญฺหาวาโรปิ วิตฺถาเรตพฺโพ.)}

\csromanheader{2}{\textbf{อพฺยากตปท-เหตุ}}
\proseitem{50}{อาสววิปฺปยุตฺตํ สาสวํ อพฺยากตํ ธมฺมํ ปฏิจฺจ อาสววิปฺปยุตฺโต สาสโว อพฺยากโต ธมฺโม อุปฺปชฺชติ เหตุปจฺจยา. (1)}

\prose{อาสววิปฺปยุตฺตํ อนาสวํ อพฺยากตํ ธมฺมํ ปฏิจฺจ อาสววิปฺปยุตฺโต อนาสโว อพฺยากโต ธมฺโม อุปฺปชฺชติ เหตุปจฺจยา.\csromanspacer{}\csromanspacer{}อาสววิปฺปยุตฺตํ อนาสวํ อพฺยากตํ ธมฺมํ ปฏิจฺจ อาสววิปฺปยุตฺโต สาสโว อพฺยากโต ธมฺโม อุปฺปชฺชติ เหตุปจฺจยา.\csromanspacer{}\csromanspacer{}อาสววิปฺปยุตฺตํ อนาสวํ อพฺยากตํ ธมฺมํ ปฏิจฺจ อาสววิปฺปยุตฺโต สาสโว อพฺยากโต จ อาสววิปฺปยุตฺโต อนาสโว อพฺยากโต จ ธมฺมา อุปฺปชฺชนฺติ เหตุปจฺจยา. (3)}

\prose{อาสววิปฺปยุตฺตํ สาสวํ อพฺยากตญฺจ อาสววิปฺปยุตฺตํ อนาสวํ อพฺยากตญฺจ ธมฺมํ ปฏิจฺจ อาสววิปฺปยุตฺโต สาสโว อพฺยากโต ธมฺโม อุปฺปชฺชติ เหตุปจฺจยา. (1) (สํขิตฺตํ.)}

\gambhira{ธมฺมานุโลม ทุกติกปฏฺฐานปาฬิ}
\proseitem{51}{\csromanfolio{368}เหตุยา ปญฺจ, อารมฺมเณ เทฺว, อธิปติยา ปญฺจ ฯเปฯ อาเสวเน เอกํ, กมฺเม ปญฺจ, วิปาเก ปญฺจ ฯเปฯ อวิคเต ปญฺจ. (สํขิตฺตํ.)}

\csromanheader{2}{\textbf{ปจฺจนีย-นเหตุ}}
\proseitem{52}{อาสววิปฺปยุตฺตํ สาสวํ อพฺยากตํ ธมฺมํ ปฏิจฺจ อาสววิปฺปยุตฺโต สาสโว อพฺยากโต ธมฺโม อุปฺปชฺชติ นเหตุปจฺจยา. (สํขิตฺตํ.)}

\proseitem{53}{นเหตุยา เอกํ, น-อารมฺมเณ ตีณิ, น-อธิปติยา เทฺว ฯเปฯ นปุเรชาเต จตฺตาริ, นปจฺฉาชาเต น-อาเสวเน ปญฺจ, นกมฺเม นวิปาเก ฯเปฯ นมคฺเค เอกํ ฯเปฯ นวิปฺปยุตฺเต เทฺว ฯเปฯ โนวิคเต ตีณิ. (สํขิตฺตํ.)}

\prose{(สหชาตวาโรปิ ปจฺจยวาโรปิ นิสฺสยวาโรปิ สํสฏฺฐวาโรปิ สมฺปยุตฺตวาโรปิ ปฏิจฺจวารสทิสา วิตฺถาเรตพฺพา.)}

\csromanheader{2}{7. ปญฺหาวาร ปจฺจยจตุกฺก}
\csromanheader{2}{\textbf{เหตุ อารมฺมณ}}
\proseitem{54}{อาสววิปฺปยุตฺโต สาสโว อพฺยากโต ธมฺโม อาสววิปฺปยุตฺตสฺส สาสวสฺส อพฺยากตสฺส ธมฺมสฺส เหตุปจฺจเยน ปจฺจโย. (1)}

\prose{อาสววิปฺปยุตฺโต อนาสโว อพฺยากโต ธมฺโม อาสววิปฺปยุตฺตสฺส อนาสวสฺส อพฺยากตสฺส ธมฺมสฺส เหตุปจฺจเยน ปจฺจโย.\csromanspacer{}\csromanspacer{}อาสววิปฺปยุตฺโต อนาสโว อพฺยากโต ธมฺโม อาสววิปฺปยุตฺตสฺส สาสวสฺส อพฺยากตสฺส ธมฺมสฺส เหตุปจฺจเยน ปจฺจโย.\csromanspacer{}\csromanspacer{}อาสววิปฺปยุตฺโต อนาสโว อพฺยากโต ธมฺโม อาสววิปฺปยุตฺตสฺส สาสวสฺส อพฺยากตสฺส จ อาสววิปฺปยุตฺตสฺส อนาสวสฺส อพฺยากตสฺส จ ธมฺมสฺส เหตุปจฺจเยน ปจฺจโย. (3)}

\prose{อาสววิปฺปยุตฺโต สาสโว อพฺยากโต ธมฺโม อาสววิปฺปยุตฺตสฺส สาสวสฺส อพฺยากตสฺส ธมฺมสฺส อารมฺมณปจฺจเยน ปจฺจโย. (1)}

\csromanmatikaanchor{mtk.67}
\csromanmatikaanchor{mtk.68}
\csromantocmarknopage{subparagraph}{20. สญฺโญชนทุก}{mtk.67}
\bookmark[dest={mtk.67},level=5]{20. สญฺโญชนทุก}
\csromantocmarkref{subparagraph}{1. กุสลตฺติก ...}{mtk.68}
\bookmark[dest={mtk.68},level=5]{1. กุสลตฺติก ...}
\csromanheader{6}{20. สญฺโญชนทุก 1. กุสลตฺติก}
\prose{\csromanfolio{369}อาสววิปฺปยุตฺโต อนาสโว อพฺยากโต ธมฺโม อาสววิปฺปยุตฺตสฺส อนาสวสฺส อพฺยากตสฺส ธมฺมสฺส อารมฺมณปจฺจเยน ปจฺจโย.\csromanspacer{}\csromanspacer{}อาสววิปฺปยุตฺโต อนาสโว อพฺยากโต ธมฺโม อาสววิปฺปยุตฺตสฺส สาสวสฺส อพฺยากตสฺส ธมฺมสฺส อารมฺมณปจฺจเยน ปจฺจโย. (2) (สํขิตฺตํ.)}

\proseitem{55}{เหตุยา จตฺตาริ, อารมฺมเณ ตีณิ, อธิปติยา จตฺตาริ, อนนฺตเร จตฺตาริ, สมนนฺตเร จตฺตาริ, สหชาเต ปญฺจ, อญฺญมญฺเญ เทฺว, นิสฺสเย สตฺต, อุปนิสฺสเย จตฺตาริ, ปุเรชาเต เทฺว, อาเสวเน เอกํ, กมฺเม จตฺตาริ, วิปาเก จตฺตาริ, อาหาเร จตฺตาริ ฯเปฯ อวิคเต สตฺต. (สํขิตฺตํ.)}

\csromanheader{2}{\textbf{ปจฺจนียุทฺธาร}}
\proseitem{56}{อาสววิปฺปยุตฺโต สาสโว อพฺยากโต ธมฺโม อาสววิปฺปยุตฺตสฺส สาสวสฺส อพฺยากตสฺส ธมฺมสฺส อารมฺมณปจฺจเยน ปจฺจโย.\csromanspacer{} สหชาตปจฺจเยน ปจฺจโย.\csromanspacer{} อุปนิสฺสยปจฺจเยน ปจฺจโย.\csromanspacer{} ปุเรชาตปจฺจเยน ปจฺจโย.\csromanspacer{} ปจฺฉาชาตปจฺจเยน ปจฺจโย.\csromanspacer{} อาหารปจฺจเยน ปจฺจโย.\csromanspacer{} อินฺทฺริยปจฺจเยน ปจฺจโย. (สํขิตฺตํ.)}

\proseitem{57}{นเหตุยา สตฺต, น-อารมฺมเณ สตฺต. (สํขิตฺตํ.)}

\csromanheader{2}{เหตุปจฺจยา น-อารมฺมเณ จตฺตาริ. (สํขิตฺตํ.)}
\csromanheader{2}{นเหตุปจฺจยา อารมฺมเณ ตีณิ. (สํขิตฺตํ.)}
\prose{(ยถา กุสลตฺติเก ปญฺหาวารสฺส อนุโลมมฺปิ ปจฺจนียมฺปิ อนุโลมปจฺจนียมฺปิ ปจฺจนียานุโลมมฺปิ คณิตํ, เอวํ คเณตพฺพํ.)}

\nitthitam{อาสววิปฺปยุตฺตสาสวทุกกุสลตฺติกํ นิฏฺฐิตํ.}
\nitthitam{อาสวโคจฺฉกํ นิฏฺฐิตํ.}
\csromanheader{2}{20. สญฺโญชนทุก 1. กุสลตฺติก}
\csromanheader{2}{1. ปฏิจฺจวาร ปจฺจยจตุกฺก-เหตุ}
\proseitem{1}{โนสญฺโญชนํ กุสลํ ธมฺมํ ปฏิจฺจ โนสญฺโญชโน กุสโล ธมฺโม อุปฺปชฺชติ เหตุปจฺจยา. (สํขิตฺตํ.)}

\gambhira{ธมฺมานุโลม ทุกติกปฏฺฐานปาฬิ}
\proseitem{2}{\csromanfolio{370}เหตุยา เอกํ, อารมฺมเณ เอกํ ฯเปฯ อวิคเต เอกํ. (สํขิตฺตํ.) (สหชาตวาเรปิ ฯเปฯ ปญฺหาวาเรปิ สพฺพตฺถ เอกํ.)}

\proseitem{3}{สญฺโญชนํ อกุสลํ ธมฺมํ ปฏิจฺจ สญฺโญชโน อกุสโล ธมฺโม อุปฺปชฺชติ เหตุปจฺจยา.\csromanspacer{}\csromanspacer{}สญฺโญชนํ อกุสลํ ธมฺมํ ปฏิจฺจ โนสญฺโญชโน อกุสโล ธมฺโม อุปฺปชฺชติ เหตุปจฺจยา.\csromanspacer{}\csromanspacer{}สญฺโญชนํ อกุสลํ ธมฺมํ ปฏิจฺจ สญฺโญชโน อกุสโล จ โนสญฺโญชโน อกุสโล จ ธมฺมา อุปฺปชฺชนฺติ เหตุปจฺจยา. (3) (สํขิตฺตํ.)}

\proseitem{4}{เหตุยา นว, อารมฺมเณ นว, อธิปติยา นว ฯเปฯ อวิคเต นว. (สํขิตฺตํ.)}

\prose{นเหตุยา ตีณิ, น-อธิปติยา นว ฯเปฯ นกมฺเม ตีณิ ฯเปฯ นวิปฺปยุตฺเต นว. (สํขิตฺตํ.)}

\prose{(สหชาตวาโรปิ ฯเปฯ สมฺปยุตฺตวาโรปิ ปฏิจฺจวารสทิสา วิตฺถาเรตพฺพา.)}

\csromanheader{2}{\textbf{ปญฺหาวาร ปจฺจยจตุกฺก-เหตุ}}
\proseitem{5}{สญฺโญชโน อกุสโล ธมฺโม สญฺโญชนสฺส อกุสลสฺส ธมฺมสฺส เหตุปจฺจเยน ปจฺจโย. (สํขิตฺตํ.)}

\proseitem{6}{เหตุยา ตีณิ, อารมฺมเณ นว, อธิปติยา นว ฯเปฯ อุปนิสฺสเย อาเสวเน นว, กมฺเม ตีณิ, อาหาเร อินฺทฺริเย ฌาเน ตีณิ, มคฺเค สมฺปยุตฺเต นว ฯเปฯ อวิคเต นว. (สํขิตฺตํ.)}

\proseitem{7}{นเหตุยา นว, น-อารมฺมเณ นว. (สํขิตฺตํ.)}

\csromanheader{2}{เหตุปจฺจยา น-อารมฺมเณ ตีณิ. (สํขิตฺตํ.)}
\csromanheader{2}{นเหตุปจฺจยา อารมฺมเณ นว. (สํขิตฺตํ.)}
\prose{(ยถา กุสลตฺติเก ปญฺหาวารสฺส อนุโลมมฺปิ ปจฺจนียมฺปิ อนุโลมปจฺจนียมฺปิ ปจฺจนียานุโลมมฺปิ คณิตํ, เอวํ คเณตพฺพํ.)}

\csromanheader{2}{\textbf{อพฺยากตปท-เหตุ}}
\proseitem{8}{โนสญฺโญชนํ อพฺยากตํ ธมฺมํ ปฏิจฺจ โนสญฺโญชโน อพฺยากโต ธมฺโม อุปฺปชฺชติ เหตุปจฺจยา. (สํขิตฺตํ.)}

\csromanheader{2}{21. สญฺโญชนิยทุก 1. กุสลตฺติก}
\proseitem{9}{\csromanfolio{371}เหตุยา เอกํ, อารมฺมเณ เอกํ ฯเปฯ อวิคเต เอกํ. (สํขิตฺตํ.)}

\vspace{\tipitakaparskip}
\csromancenter{(สหชาตวาเรปิ ฯเปฯ ปญฺหาวาเรปิ สพฺพตฺถ เอกํ.)}
\csromanheader{2}{21. สญฺโญชนิยทุก 1. กุสลตฺติก}
\csromanheader{2}{1. ปฏิจฺจวาร ปจฺจยจตุกฺก-เหตุ}
\proseitem{10}{สญฺโญชนิยํ กุสลํ ธมฺมํ ปฏิจฺจ สญฺโญชนิโย กุสโล ธมฺโม อุปฺปชฺชติ เหตุปจฺจยา. (1)}

\prose{อสญฺโญชนิยํ กุสลํ ธมฺมํ ปฏิจฺจ อสญฺโญชนิโย กุสโล ธมฺโม อุปฺปชฺชติ เหตุปจฺจยา. (1) (สํขิตฺตํ.)}

\proseitem{11}{เหตุยา เทฺว, อารมฺมเณ เทฺว. (สํขิตฺตํ.)}

\prose{(ยถา จูฬนฺตรทุเก โลกิยทุกคมนํ, เอวํ อิมมฺปิ ญาตพฺพํ.\csromanspacer{} สหชาตวาโรปิ ฯเปฯ ปญฺหาวาโรปิ วิตฺถาเรตพฺพา.)}

\proseitem{12}{สญฺโญชนิยํ อกุสลํ ธมฺมํ ปฏิจฺจ สญฺโญชนิโย อกุสโล ธมฺโม อุปฺปชฺชติ เหตุปจฺจยา. (สํขิตฺตํ.)}

\proseitem{13}{เหตุยา เอกํ, อารมฺมเณ เอกํ ฯเปฯ อวิคเต เอกํ. (สํขิตฺตํ.) (สหชาตวาเรปิ ฯเปฯ ปญฺหาวาเรปิ สพฺพตฺถ เอกํ.)}

\csromanheader{2}{\textbf{อพฺยากตปท-เหตุ}}
\proseitem{14}{สญฺโญชนิยํ อพฺยากตํ ธมฺมํ ปฏิจฺจ สญฺโญชนิโย อพฺยากโต ธมฺโม อุปฺปชฺชติ เหตุปจฺจยา. (1)}

\prose{อสญฺโญชนิยํ อพฺยากตํ ธมฺมํ ปฏิจฺจ อสญฺโญชนิโย อพฺยากโต ธมฺโม อุปฺปชฺชติ เหตุปจฺจยา.\csromanspacer{}\csromanspacer{}อสญฺโญชนิยํ อพฺยากตํ ธมฺมํ ปฏิจฺจ สญฺโญชนิโย อพฺยากโต ธมฺโม อุปฺปชฺชติ เหตุปจฺจยา.\csromanspacer{}\csromanspacer{}อสญฺโญชนิยํ อพฺยากตํ ธมฺมํ ปฏิจฺจ สญฺโญชนิโย อพฺยากโต จ อสญฺโญชนิโย อพฺยากโต จ ธมฺโม อุปฺปชฺชนฺติ เหตุปจฺจยา. (3)}

\prose{สญฺโญชนิยํ อพฺยากตญฺจ อสญฺโญชนิยํ อพฺยากตญฺจ ธมฺมํ ปฏิจฺจ สญฺโญชนิโย อพฺยากโต ธมฺโม อุปฺปชฺชติ เหตุปจฺจยา. (1)}

\gambhira{ธมฺมานุโลม ทุกติกปฏฺฐานปาฬิ}
\proseitem{15}{\csromanfolio{372}เหตุยา ปญฺจ, อารมฺมเณ เทฺว, อธิปติยา ปญฺจ ฯเปฯ อวิคเต ปญฺจ. (สํขิตฺตํ.)}

\prose{จูฬนฺตรทุเก โลกิยทุกสทิสํ. (สหชาตวาโรปิ ฯเปฯ สมฺปยุตฺตวาโรปิ วิตฺถาเรตพฺพา.)}

\csromanheader{2}{\textbf{ปญฺหาวาร-เหตุ}}
\proseitem{16}{สญฺโญชนิโย อพฺยากโต ธมฺโม สญฺโญชนิยสฺส อพฺยากตสฺส ธมฺมสฺส เหตุปจฺจเยน ปจฺจโย. (1)}

\prose{อสญฺโญชนิโย อพฺยากโต ธมฺโม อสญฺโญชนิยสฺส อพฺยากตสฺส ธมฺมสฺส เหตุปจฺจเยน ปจฺจโย.\csromanspacer{} ตีณิ. (สํขิตฺตํ.)}

\proseitem{17}{เหตุยา จตฺตาริ, อารมฺมเณ ตีณิ, อธิปติยา จตฺตาริ ฯเปฯ อวิคเต สตฺต. (สํขิตฺตํ.)}

\csromanheader{2}{\textbf{ปจฺจนียุทฺธาร}}
\proseitem{18}{สญฺโญชนิโย อพฺยากโต ธมฺโม สญฺโญชนิยสฺส อพฺยากตสฺส ธมฺมสฺส อารมฺมณปจฺจเยน ปจฺจโย.\csromanspacer{} สหชาตปจฺจเยน ปจฺจโย.\csromanspacer{} อุปนิสฺสยปจฺจเยน ปจฺจโย.\csromanspacer{} ปุเรชาตปจฺจเยน ปจฺจโย.\csromanspacer{} ปจฺฉาชาตปจฺจเยน ปจฺจโย.\csromanspacer{} อาหารปจฺจเยน ปจฺจโย.\csromanspacer{} อินฺทฺริยปจฺจเยน ปจฺจโย. (สํขิตฺตํ.)}

\proseitem{19}{นเหตุยา สตฺต, น-อารมฺมเณ สตฺต. (สํขิตฺตํ.)}

\csromanheader{2}{เหตุปจฺจยา น-อารมฺมเณ จตฺตาริ. (สํขิตฺตํ.)}
\csromanheader{2}{นเหตุปจฺจยา อารมฺมเณ ตีณิ. (สํขิตฺตํ.)}
\prose{(ยถา กุสลตฺติเก ปญฺหาวารสฺส อนุโลมมฺปิ ปจฺจนียมฺปิ อนุโลมปจฺจนียมฺปิ ปจฺจนียานุโลมมฺปิ คณิตํ, เอวํ คเณตพฺพํ.)}

\csromanheader{2}{22. สญฺโญชนสมฺปยุตฺตทุก 1. กุสลตฺติก}
\csromanheader{2}{1. ปฏิจฺจวาร ปจฺจยจตุกฺก-เหตุ}
\proseitem{20}{สญฺโญชนวิปฺปยุตฺตํ กุสลํ ธมฺมํ ปฏิจฺจ สญฺโญชนวิปฺปยุตฺโต กุสโล ธมฺโม อุปฺปชฺชติ เหตุปจฺจยา. (สํขิตฺตํ.)}

\csromanheader{2}{22. สญฺโญชนสมฺปยุตฺตทุก 1. กุสลตฺติก}
\prose{\csromanfolio{373}เหตุยา เอกํ, อารมฺมเณ เอกํ ฯเปฯ อวิคเต เอกํ. (สํขิตฺตํ.) (สหชาตวาเรปิ ฯเปฯ ปญฺหาวาเรปิ สพฺพตฺถ เอกํ.)}

\proseitem{21}{สญฺโญชนสมฺปยุตฺตํ อกุสลํ ธมฺมํ ปฏิจฺจ สญฺโญชนสมฺปยุตฺโต อกุสโล ธมฺโม อุปฺปชฺชติ เหตุปจฺจยา.\csromanspacer{} ตีณิ.}

\prose{สญฺโญชนสมฺปยุตฺตํ อกุสลํ ธมฺมํ ปฏิจฺจ สญฺโญชนสมฺปยุตฺโต อกุสโล ธมฺโม อุปฺปชฺชติ อารมฺมณปจฺจยา. (สํขิตฺตํ.)}

\proseitem{22}{เหตุยา ตีณิ, อารมฺมเณ ปญฺจ, อธิปติยา เอกํ ฯเปฯ อวิคเต ปญฺจ. (สํขิตฺตํ.\csromanspacer{} อนุโลมํ.)}

\csromanheader{2}{\textbf{ปจฺจนีย-นเหตุ}}
\proseitem{23}{สญฺโญชนสมฺปยุตฺตํ อกุสลํ ธมฺมํ ปฏิจฺจ สญฺโญชนสมฺปยุตฺโต อกุสโล ธมฺโม อุปฺปชฺชติ นเหตุปจฺจยา.\csromanspacer{} วิจิกิจฺฉาสหคเต ขนฺเธ ปฏิจฺจ วิจิกิจฺฉาสหคโต โมโห.\csromanspacer{}\csromanspacer{}สญฺโญชนสมฺปยุตฺตํ อกุสลํ ธมฺมํ ปฏิจฺจ สญฺโญชนวิปฺปยุตฺโต อกุสโล ธมฺโม อุปฺปชฺชติ นเหตุปจฺจยา.\csromanspacer{} อุทฺธจฺจสหคเต ขนฺเธ ปฏิจฺจ อุทฺธจฺจสหคโต โมโห. (2) (สํขิตฺตํ.)}

\proseitem{24}{นเหตุยา เทฺว, น-อธิปติยา ปญฺจ, นปุเรชาเต ปญฺจ ฯเปฯ นกมฺเม ตีณิ ฯเปฯ นวิปฺปยุตฺเต ปญฺจ. (สํขิตฺตํ.\csromanspacer{} ปจฺจนียํ.)}

\csromanheader{2}{\textbf{ปญฺหาวาร-เหตุ}}
\proseitem{25}{สญฺโญชนสมฺปยุตฺโต อกุสโล ธมฺโม สญฺโญชนสมฺปยุตฺตสฺส อกุสลสฺส ธมฺมสฺส เหตุปจฺจเยน ปจฺจโย. (1)}

\prose{สญฺโญชนวิปฺปยุตฺโต อกุสโล ธมฺโม สญฺโญชนสมฺปยุตฺตสฺส อกุสลสฺส ธมฺมสฺส อารมฺมณปจฺจเยน ปจฺจโย. (1)}

\prose{สญฺโญชนสมฺปยุตฺโต อกุสโล ธมฺโม สญฺโญชนสมฺปยุตฺตสฺส อกุสลสฺส ธมฺมสฺส อารมฺมณปจฺจเยน ปจฺจโย.\csromanspacer{} นว.}

\prose{สญฺโญชนสมฺปยุตฺโต อกุสโล ธมฺโม สญฺโญชนสมฺปยุตฺตสฺส อกุสลสฺส ธมฺมสฺส อธิปติปจฺจเยน ปจฺจโย. (1) (สํขิตฺตํ.)}

\gambhira{ธมฺมานุโลม ทุกติกปฏฺฐานปาฬิ}
\proseitem{26}{\csromanfolio{374}อธิปติยา เอกํ, อนนฺตเร สมนนฺตเร นว, สหชาเต ปญฺจ, อุปนิสฺสเย อาเสวเน นว, กมฺเม ตีณิ, มคฺเค ตีณิ, สมฺปยุตฺเต ปญฺจ, อตฺถิยา ปญฺจ. (สํขิตฺตํ.)}

\csromanheader{2}{\textbf{อพฺยากตปท-เหตุ}}
\proseitem{27}{สญฺโญชนวิปฺปยุตฺตํ อพฺยากตํ ธมฺมํ ปฏิจฺจ สญฺโญชนวิปฺปยุตฺโต อพฺยากโต ธมฺโม อุปฺปชฺชติ เหตุปจฺจยา. (สํขิตฺตํ.)}

\prose{เหตุยา เอกํ, อารมฺมเณ เอกํ ฯเปฯ อวิคเต เอกํ. (สํขิตฺตํ.)}

\vspace{\tipitakaparskip}
\csromancenter{(สหชาตวาเรปิ ฯเปฯ ปญฺหาวาเรปิ สพฺพตฺถ เอกํ.)}
\csromanheader{2}{23. สญฺโญชนสญฺโญชนิยทุก 1. กุสลตฺติก}
\csromanheader{2}{1-7. ปฏิจฺจาทิวาร ปจฺจยจตุกฺก-เหตุ}
\proseitem{28}{สญฺโญชนิยญฺเจว โน จ สญฺโญชนํ กุสลํ ธมฺมํ ปฏิจฺจ สญฺโญชนิโย เจว โน จ สญฺโญชโน กุสโล ธมฺโม อุปฺปชฺชติ เหตุปจฺจยา. (สํขิตฺตํ.)}

\proseitem{29}{เหตุยา เอกํ, อารมฺมเณ เอกํ ฯเปฯ อวิคเต เอกํ. (สํขิตฺตํ.) (สหชาตวาเรปิ ฯเปฯ ปญฺหาวาเรปิ สพฺพตฺถ เอกํ.)}

\proseitem{30}{สญฺโญชนญฺเจว สญฺโญชนิยญฺจ อกุสลํ ธมฺมํ ปฏิจฺจ สญฺโญชโน เจว สญฺโญชนิโย จ อกุสโล ธมฺโม อุปฺปชฺชติ เหตุปจฺจยา.\csromanspacer{} ตีณิ. (สํขิตฺตํ.)}

\proseitem{31}{เหตุยา นว, อารมฺมเณ นว, อธิปติยา นว ฯเปฯ อวิคเต นว. (สํขิตฺตํ.)}

\prose{นเหตุยา ตีณิ, น-อธิปติยา นว, นปุเรชาเต นว, นปจฺฉาชาเต นว, น-อาเสวเน นว, นกมฺเม ตีณิ, นวิปาเก นว, นวิปฺปยุตฺเต นว. (สํขิตฺตํ.)}

\prose{(สหชาตวาโรปิ ฯเปฯ สมฺปยุตฺตวาโรปิ ปฏิจฺจวารสทิสา.)}

\csromanheader{2}{24. สญฺโญชนสญฺโญชนสมฺปยุตฺตทุก 1. กุสลตฺติก}
\proseitem{32}{\csromanfolio{375}สญฺโญชโน เจว สญฺโญชนิโย จ อกุสโล ธมฺโม สญฺโญชนสฺส เจว สญฺโญชนิยสฺส จ อกุสลสฺส ธมฺมสฺส เหตุปจฺจเยน ปจฺจโย. (สํขิตฺตํ.)}

\proseitem{33}{เหตุยา ตีณิ, อารมฺมเณ นว, อธิปติยา นว ฯเปฯ อุปนิสฺสเย อาเสวเน นว, กมฺเม อาหาเร อินฺทฺริเย ฌาเน ตีณิ, มคฺเค สมฺปยุตฺเต นว ฯเปฯ อวิคเต นว. (สํขิตฺตํ.)}

\proseitem{34}{นเหตุยา นว, น-อารมฺมเณ นว. (สํขิตฺตํ.)}

\csromanheader{2}{เหตุปจฺจยา น-อารมฺมเณ ตีณิ. (สํขิตฺตํ.)}
\csromanheader{2}{นเหตุปจฺจยา อารมฺมเณ นว. (สํขิตฺตํ.)}
\prose{(ยถา กุสลตฺติเก ปญฺหาวารสฺส อนุโลมมฺปิ ปจฺจนียมฺปิ อนุโลมปจฺจนียมฺปิ ปจฺจนียานุโลมมฺปิ คณิตํ, เอวํ คเณตพฺพํ.)}

\proseitem{35}{สญฺโญชนิยญฺเจว โน จ สญฺโญชนํ อพฺยากตํ ธมฺมํ ปฏิจฺจ สญฺโญชนิโย เจว โน จ สญฺโญชโน อพฺยากโต ธมฺโม อุปฺปชฺชติ เหตุปจฺจยา. (สํขิตฺตํ.)}

\csromanheader{2}{36. เหตุยา เอกํ อารมฺมเณ เอกํ ฯเปฯ อวิคเต เอกํ. (สํขิตฺตํ.)}
\vspace{\tipitakaparskip}
\csromancenter{(สหชาตวาเรปิ ฯเปฯ ปญฺหาวาเรปิ สพฺพตฺถ เอกํ.)}
\csromanheader{2}{24. สญฺโญชนสญฺโญชนสมฺปยุตฺตทุก 1. กุสลตฺติก}
\csromanheader{2}{1-7. ปฏิจฺจาทิวาร ปจฺจยจตุกฺก-เหตุ}
\proseitem{37}{สญฺโญชนญฺเจว สญฺโญชนสมฺปยุตฺตญฺจ อกุสลํ ธมฺมํ ปฏิจฺจ สญฺโญชโน เจว สญฺโญชนสมฺปยุตฺโต จ อกุสโล ธมฺโม อุปฺปชฺชติ เหตุปจฺจยา.\csromanspacer{} ตีณิ.}

\prose{สญฺโญชนสมฺปยุตฺตญฺเจว โน จ สญฺโญชนํ อกุสลํ ธมฺมํ ปฏิจฺจ สญฺโญชนสมฺปยุตฺโต เจว โน จ สญฺโญชโน อกุสโล ธมฺโม อุปฺปชฺชติ เหตุปจฺจยา.\csromanspacer{} ตีณิ.}

\gambhira{ธมฺมานุโลม ทุกติกปฏฺฐานปาฬิ}
\prose{\csromanfolio{376}สญฺโญชนญฺเจว สญฺโญชนสมฺปยุตฺตํ อกุสลญฺจ สญฺโญชนสมฺปยุตฺตญฺเจว โน จ สญฺโญชนํ อกุสลญฺจ ธมฺมํ ปฏิจฺจ สญฺโญชโน เจว สญฺโญชนสมฺปยุตฺโต จ อกุสโล ธมฺโม อุปฺปชฺชติ เหตุปจฺจยา.\csromanspacer{} ตีณิ. (สํขิตฺตํ.)}

\proseitem{38}{เหตุยา นว, อารมฺมเณ นว, อธิปติยา นว ฯเปฯ อวิคเต นว. (สํขิตฺตํ.)}

\prose{นเหตุยา ตีณิ, น-อธิปติยา นว, นปุเรชาเต นว, นปจฺฉาชาเต นว, น-อาเสวเน นว, นกมฺเม ตีณิ, นวิปาเก นว, นวิปฺปยุตฺเต นว. (สํขิตฺตํ.)}

\prose{(สหชาตวาโรปิ สมฺปยุตฺตวาโรปิ ปฏิจฺจวารสทิสา วิตฺถาเรตพฺพา.)}

\proseitem{39}{สญฺโญชโน เจว สญฺโญชนสมฺปยุตฺโต จ อกุสโล ธมฺโม สญฺโญชนสฺส เจว สญฺโญชนสมฺปยุตฺตสฺส จ อกุสลสฺส ธมฺมสฺส เหตุปจฺจเยน ปจฺจโย. (สํขิตฺตํ.)}

\proseitem{40}{เหตุยา ตีณิ, อารมฺมเณ นว, อธิปติยา นว ฯเปฯ อุปนิสฺสเย อาเสวเน นว, กมฺเม อาหาเร อินฺทฺริเย ฌาเน ตีณิ, มคฺเค สมฺปยุตฺเต นว ฯเปฯ อวิคเต นว. (สํขิตฺตํ.)}

\proseitem{41}{นเหตุยา นว, น-อารมฺมเณ นว, น-อธิปติยา นว. (สํขิตฺตํ.)}

\csromanheader{2}{เหตุปจฺจยา น-อารมฺมเณ ตีณิ. (สํขิตฺตํ.)}
\csromanheader{2}{นเหตุปจฺจยา อารมฺมเณ นว. (สํขิตฺตํ.)}
\prose{(ยถา กุสลตฺติเก ปญฺหาวารสฺส อนุโลมมฺปิ ปจฺจนียมฺปิ อนุโลมปจฺจนียมฺปิ ปจฺจนียานุโลมมฺปิ คณิตํ, เอวํ คเณตพฺพํ.)}

\csromanheader{2}{25. สญฺโญชนวิปฺปยุตฺตสญฺโญชนิยทุก 1. กุสลตฺติก}
\csromanheader{2}{1-7. ปฏิจฺจาทิวาร ปจฺจยจตุกฺก-เหตุ}
\proseitem{42}{สญฺโญชนวิปฺปยุตฺตํ สญฺโญชนิยํ กุสลํ ธมฺมํ ปฏิจฺจ สญฺโญชนวิปฺปยุตฺโต สญฺโญชนิโย กุสโล ธมฺโม อุปฺปชฺชติ เหตุปจฺจยา. (1)}

\csromanheader{2}{25. สญฺโญชนวิปฺปยุตฺตสญฺโญชนิยทุก 1. กุสลตฺติก}
\prose{\csromanfolio{377}สญฺโญชนวิปฺปยุตฺตํ อสญฺโญชนิยํ กุสลํ ธมฺมํ ปฏิจฺจ สญฺโญชนวิปฺปยุตฺโต อสญฺโญชนิโย กุสโล ธมฺโม อุปฺปชฺชติ เหตุปจฺจยา. (1)}

\proseitem{43}{เหตุยา เทฺว, อารมฺมเณ เทฺว ฯเปฯ อวิคเต เทฺว. (สํขิตฺตํ.)}

\prose{(ยถา จูฬนฺตรทุเก โลกิยทุกสทิสํ.\csromanspacer{} สหชาตวาโรปิ ฯเปฯ ปญฺหาวาโรปิ สพฺพตฺถ วิตฺถาเรตพฺพา.)}

\proseitem{44}{สญฺโญชนวิปฺปยุตฺโต สญฺโญชนิโย อกุสโล ธมฺโม สญฺโญชนวิปฺปยุตฺตสฺส สญฺโญชนิยสฺส อกุสลสฺส ธมฺมสฺส อารมฺมณปจฺจเยน ปจฺจโย. (สํขิตฺตํ.)}

\proseitem{45}{อารมฺมเณ เอกํ. (สพฺพตฺถ เอกํ, สํขิตฺตํ.)}

\proseitem{46}{สญฺโญชนวิปฺปยุตฺตํ สญฺโญชนิยํ อพฺยากตํ ธมฺมํ ปฏิจฺจ สญฺโญชนวิปฺปยุตฺโต สญฺโญชนิโย อพฺยากโต ธมฺโม อุปฺปชฺชติ เหตุปจฺจยา. (1)}

\prose{สญฺโญชนวิปฺปยุตฺตํ อสญฺโญชนิยํ อพฺยากตํ ธมฺมํ ปฏิจฺจ สญฺโญชนวิปฺปยุตฺโต อสญฺโญชนิโย อพฺยากโต ธมฺโม อุปฺปชฺชติ เหตุปจฺจยา.\csromanspacer{}\csromanspacer{}สญฺโญชนวิปฺปยุตฺตํ อสญฺโญชนิยํ อพฺยากตํ ธมฺมํ ปฏิจฺจ สญฺโญชนวิปฺปยุตฺโต สญฺโญชนิโย อพฺยากโต ธมฺโม อุปฺปชฺชติ เหตุปจฺจยา.\csromanspacer{}\csromanspacer{}สญฺโญชนวิปฺปยุตฺตํ อสญฺโญชนิยํ อพฺยากตํ ธมฺมํ ปฏิจฺจ สญฺโญชนวิปฺปยุตฺโต สญฺโญชนิโย อพฺยากโต จ สญฺโญชนวิปฺปยุตฺโต อสญฺโญชนิโย อพฺยากโต จ ธมฺมา อุปฺปชฺชนฺติ เหตุปจฺจยา. (3)}

\prose{สญฺโญชนวิปฺปยุตฺตํ สญฺโญชนิยํ อพฺยากตญฺจ สญฺโญชนวิปฺปยุตฺตํ อสญฺโญชนิยํ อพฺยากตญฺจ ธมฺมํ ปฏิจฺจ สญฺโญชนวิปฺปยุตฺโต สญฺโญชนิโย อพฺยากโต ธมฺโม อุปฺปชฺชติ เหตุปจฺจยา. (1) (สํขิตฺตํ.)}

\proseitem{47}{เหตุยา ปญฺจ, อารมฺมเณ เทฺว, อธิปติยา ปญฺจ ฯเปฯ อาเสวเน เอกํ, กมฺเม ปญฺจ, วิปาเก ปญฺจ ฯเปฯ อวิคเต ปญฺจ. (สํขิตฺตํ.)}

\gambhira{ธมฺมานุโลม ทุกติกปฏฺฐานปาฬิ}
\prose{\csromanfolio{378}นเหตุยา เอกํ, น-อารมฺมเณ ตีณิ, น-อธิปติยา เทฺว, นปุเรชาเต จตฺตาริ, นปจฺฉาชาเต ปญฺจ, น-อาเสวเน ปญฺจ, นกมฺเม เอกํ, นวิปาเก เอกํ, นวิปฺปยุตฺเต เทฺว ฯเปฯ โนวิคเต ตีณิ. (สํขิตฺตํ.) (สหชาตวาโรปิ ฯเปฯ สมฺปยุตฺตวาโรปิ ปฏิจฺจวารสทิสา.)}

\proseitem{48}{สญฺโญชนวิปฺปยุตฺโต สญฺโญชนิโย อพฺยากโต ธมฺโม สญฺโญชนวิปฺปยุตฺตสฺส สญฺโญชนิยสฺส อพฺยากตสฺส ธมฺมสฺส เหตุปจฺจเยน ปจฺจโย. (สํขิตฺตํ.)}

\proseitem{49}{เหตุยา จตฺตาริ, อารมฺมเณ ตีณิ, อธิปติยา จตฺตาริ, อนนฺตเร จตฺตาริ ฯเปฯ อวิคเต สตฺต. (สํขิตฺตํ.)}

\proseitem{50}{นเหตุยา สตฺต, น-อารมฺมเณ สตฺต. (สํขิตฺตํ.)}

\csromanheader{2}{เหตุปจฺจยา น-อารมฺมเณ จตฺตาริ. (สํขิตฺตํ.)}
\csromanheader{2}{นเหตุปจฺจยา อารมฺมเณ ตีณิ. (สํขิตฺตํ.)}
\prose{(ยถา กุสลตฺติเก ปญฺหาวารสฺส อนุโลมมฺปิ ปจฺจนียมฺปิ อนุโลมปจฺจนียมฺปิ ปจฺจนียานุโลมมฺปิ คณิตํ, เอวํ คเณตพฺพํ.)}

\nitthitam{สญฺโญชนโคจฺฉกกุสลตฺติกํ นิฏฺฐิตํ.}
\csromanmatikaanchor{mtk.69}
\csromanmatikaanchor{mtk.70}
\csromantocmarknopage{subparagraph}{26. คนฺถทุก}{mtk.69}
\bookmark[dest={mtk.69},level=5]{26. คนฺถทุก}
\csromantocmarkref{subparagraph}{1. กุสลตฺติก ...}{mtk.70}
\bookmark[dest={mtk.70},level=5]{1. กุสลตฺติก ...}
\csromanheader{6}{26. คนฺถทุก 1. กุสลตฺติก}
\csromanheader{2}{1-7. ปฏิจฺจาทิวาร ปจฺจยจตุกฺก}
\proseitem{1}{โนคนฺถํ กุสลํ ธมฺมํ ปฏิจฺจ โนคนฺโถ กุสโล ธมฺโม อุปฺปชฺชติ เหตุปจฺจยา. (สํขิตฺตํ.)}

\proseitem{2}{เหตุยา เอกํ, อารมฺมเณ เอกํ ฯเปฯ อวิคเต เอกํ. (สํขิตฺตํ.)}

\vspace{\tipitakaparskip}
\csromancenter{(สหชาตวาเรปิ ฯเปฯ ปญฺหาวาเรปิ สพฺพตฺถ เอกํ.)}
\csromanheader{2}{26. คนฺถทุก 1. กุสลตฺติก}
\proseitem{3}{\csromanfolio{379}คนฺถํ อกุสลํ ธมฺมํ ปฏิจฺจ คนฺโถ อกุสโล ธมฺโม อุปฺปชฺชติ เหตุปจฺจยา.\csromanspacer{}\csromanspacer{}คนฺถํ อกุสลํ ธมฺมํ ปฏิจฺจ โนคนฺโถ อกุสโล ธมฺโม อุปฺปชฺชติ เหตุปจฺจยา.\csromanspacer{}\csromanspacer{}คนฺถํ อกุสลํ ธมฺมํ ปฏิจฺจ คนฺโถ อกุสโล จ โนคนฺโถ อกุสโล จ ธมฺมา อุปฺปชฺชนฺติ เหตุปจฺจยา. (3)}

\prose{โนคนฺถํ อกุสลํ ธมฺมํ ปฏิจฺจ โนคนฺโถ อกุสโล ธมฺโม อุปฺปชฺชติ เหตุปจฺจยา.\csromanspacer{} ตีณิ.}

\prose{คนฺถํ อกุสลญฺจ โนคนฺถํ อกุสลญฺจ ธมฺมํ ปฏิจฺจ คนฺโถ อกุสโล ธมฺโม อุปฺปชฺชติ เหตุปจฺจยา.\csromanspacer{} ตีณิ. (สํขิตฺตํ.)}

\proseitem{4}{เหตุยา นว, อารมฺมเณ นว ฯเปฯ อวิคเต นว. (สํขิตฺตํ.)}

\prose{นเหตุยา เอกํ, น-อธิปติยา นว, นปุเรชาเต นว, นปจฺฉาชาเต นว, น-อาเสวเน นว, นกมฺเม ตีณิ, นวิปาเก นว, นวิปฺปยุตฺเต นว. (สํขิตฺตํ.)}

\prose{(สหชาตวาโรปิ ฯเปฯ สมฺปยุตฺตวาโรปิ ปฏิจฺจวารสทิสา วิตฺถาเรตพฺพา.)}

\proseitem{5}{คนฺโถ อกุสโล ธมฺโม คนฺถสฺส อกุสลสฺส ธมฺมสฺส เหตุปจฺจเยน ปจฺจโย. (สํขิตฺตํ.)}

\proseitem{6}{เหตุยา นว, อารมฺมเณ นว, อธิปติยา นว ฯเปฯ กมฺเม อาหาเร อินฺทฺริเย ฌาเน ตีณิ ฯเปฯ อวิคเต นว. (สํขิตฺตํ.)}

\proseitem{7}{นเหตุยา นว, น-อารมฺมเณ นว. (สํขิตฺตํ.)}

\csromanheader{2}{เหตุปจฺจยา น-อารมฺมเณ นว. (สํขิตฺตํ.)}
\csromanheader{2}{นเหตุปจฺจยา อารมฺมเณ นว. (สํขิตฺตํ.)}
\prose{(ยถา กุสลตฺติเก ปญฺหาวารสฺส อนุโลมมฺปิ ปจฺจนียมฺปิ อนุโลมปจฺจนียมฺปิ ปจฺจนียานุโลมมฺปิ คณิตํ, เอวํ คเณตพฺพํ.)}

\proseitem{8}{โนคนฺถํ อพฺยากตํ ธมฺมํ ปฏิจฺจ โนคนฺโถ อพฺยากโต ธมฺโม อุปฺปชฺชติ เหตุปจฺจยา. (สํขิตฺตํ.)}

\csromanheader{2}{9. เหตุยา เอกํ ฯเปฯ อวิคเต เอกํ. (สํขิตฺตํ.)}
\vspace{\tipitakaparskip}
\csromancenter{(สหชาตวาเรปิ ฯเปฯ ปญฺหาวาเรปิ สพฺพตฺถ เอกํ.)}
\csromanheader{2}{ธมฺมานุโลม ทุกติกปฏฺฐานปาฬิ \textbf{27. คนฺถนิยทุก 1. กุสลตฺติก}}
\csromanheader{2}{1-7. ปฏิจฺจาทิวาร ปจฺจยจตุกฺก-เหตุ}
\proseitem{10}{\csromanfolio{380}คนฺถนิยํ กุสลํ ธมฺมํ ปฏิจฺจ คนฺถนิโย กุสโล ธมฺโม อุปฺปชฺชติ เหตุปจฺจยา. (1)}

\prose{อคนฺถนิยํ กุสลํ ธมฺมํ ปฏิจฺจ อคนฺถนิโย กุสโล ธมฺโม อุปฺปชฺชติ เหตุปจฺจยา. (1) (สํขิตฺตํ.)}

\proseitem{11}{เหตุยา เทฺว, อารมฺมเณ เทฺว ฯเปฯ อวิคเต เทฺว. (สํขิตฺตํ.\csromanspacer{} โลกิยโลกุตฺตรคมนสทิสํ.\csromanspacer{} สหชาตวาโรปิ ฯเปฯ ปญฺหาวาโรปิ วิตฺถาเรตพฺพา.)}

\proseitem{12}{คนฺถนิยํ อกุสลํ ธมฺมํ ปฏิจฺจ คนฺถนิโย อกุสโล ธมฺโม อุปฺปชฺชติ เหตุปจฺจยา. (สํขิตฺตํ.)}

\proseitem{13}{เหตุยา เอกํ, อารมฺมเณ เอกํ ฯเปฯ อวิคเต เอกํ. (สํขิตฺตํ.) (สหชาตวาเรปิ ฯเปฯ ปญฺหาวาเรปิ สพฺพตฺถ เอกํ.)}

\proseitem{14}{คนฺถนิยํ อพฺยากตํ ธมฺมํ ปฏิจฺจ คนฺถนิโย อพฺยากโต ธมฺโม อุปฺปชฺชติ เหตุปจฺจยา. (1)}

\prose{อคนฺถนิยํ อพฺยากตํ ธมฺมํ ปฏิจฺจ อคนฺถนิโย อพฺยากโต ธมฺโม อุปฺปชฺชติ เหตุปจฺจยา.\csromanspacer{} ตีณิ. (3)}

\prose{คนฺถนิยํ อพฺยากตญฺจ อคนฺถนิยํ อพฺยากตญฺจ ธมฺมํ ปฏิจฺจ คนฺถนิโย อพฺยากโต ธมฺโม อุปฺปชฺชติ เหตุปจฺจยา. (1) (สํขิตฺตํ.)}

\proseitem{15}{เหตุยา ปญฺจ, อารมฺมเณ เทฺว ฯเปฯ อวิคเต ปญฺจ. (สํขิตฺตํ.\csromanspacer{} โลกิยโลกุตฺตรคมนสทิสํ.\csromanspacer{} สหชาตวาโรปิ ฯเปฯ ปญฺหาวาโรปิ วิตฺถาเรตพฺพา.)}

\csromanheader{2}{28. คนฺถสมฺปยุตฺตทุก 1. กุสลตฺติก}
\csromanheader{2}{1-7. ปฏิจฺจาทิวาร ปจฺจยจตุกฺก-เหตุ}
\proseitem{16}{คนฺถวิปฺปยุตฺตํ กุสลํ ธมฺมํ ปฏิจฺจ คนฺถวิปฺปยุตฺโต กุสโล ธมฺโม อุปฺปชฺชติ เหตุปจฺจยา. (สํขิตฺตํ.)}

\csromanheader{2}{29. คนฺถคนฺถนิยทุก 1. กุสลตฺติก}
\proseitem{17}{\csromanfolio{381}เหตุยา เอกํ, อารมฺมเณ เอกํ ฯเปฯ อวิคเต เอกํ. (สํขิตฺตํ.) (สหชาตวาเรปิ ฯเปฯ ปญฺหาวาเรปิ สพฺพตฺถ เอกํ.)}

\proseitem{18}{คนฺถสมฺปยุตฺตํ อกุสลํ ธมฺมํ ปฏิจฺจ คนฺถสมฺปยุตฺโต อกุสโล ธมฺโม อุปฺปชฺชติ เหตุปจฺจยา.\csromanspacer{} ตีณิ.}

\prose{คนฺถวิปฺปยุตฺตํ อกุสลํ ธมฺมํ ปฏิจฺจ คนฺถวิปฺปยุตฺโต อกุสโล ธมฺโม อุปฺปชฺชติ เหตุปจฺจยา.\csromanspacer{} เทฺว.}

\prose{คนฺถสมฺปยุตฺตํ อกุสลญฺจ คนฺถวิปฺปยุตฺตํ อกุสลญฺจ ธมฺมํ ปฏิจฺจ คนฺถสมฺปยุตฺโต อกุสโล ธมฺโม อุปฺปชฺชติ เหตุปจฺจยา. (1) (สํขิตฺตํ.)}

\proseitem{19}{เหตุยา ฉ, อารมฺมเณ ฉ, อธิปติยา ปญฺจ ฯเปฯ อวิคเต ฉ. (สํขิตฺตํ.)}

\prose{นเหตุยา เอกํ, น-อธิปติยา ฉ, นปุเรชาเต ฉ ฯเปฯ นกมฺเม จตฺตาริ, นวิปฺปยุตฺเต ฉ. (สํขิตฺตํ.) (สหชาตวาโรปิ ฯเปฯ สมฺปยุตฺตวาโรปิ วิตฺถาเรตพฺพา.)}

\proseitem{20}{คนฺถสมฺปยุตฺโต อกุสโล ธมฺโม คนฺถสมฺปยุตฺตสฺส อกุสลสฺส ธมฺมสฺส เหตุปจฺจเยน ปจฺจโย. (สํขิตฺตํ.)}

\proseitem{21}{เหตุยา ฉ, อารมฺมเณ นว, อธิปติยา นว, อนนฺตเร นว, สหชาเต ฉ ฯเปฯ นิสฺสเย ฉ, อุปนิสฺสเย นว, อาเสวเน นว, กมฺเม จตฺตาริ ฯเปฯ มคฺเค จตฺตาริ, สมฺปยุตฺเต ฉ ฯเปฯ อวิคเต ฉ. (สํขิตฺตํ.)}

\proseitem{22}{คนฺถวิปฺปยุตฺตํ อพฺยากตํ ธมฺมํ ปฏิจฺจ คนฺถวิปฺปยุตฺโต อพฺยากโต ธมฺโม อุปฺปชฺชติ เหตุปจฺจยา. (สํขิตฺตํ.)}

\proseitem{23}{เหตุยา เอกํ, อารมฺมเณ เอกํ ฯเปฯ อวิคเต เอกํ. (สํขิตฺตํ.)}

\vspace{\tipitakaparskip}
\csromancenter{(สหชาตวาโรปิ ฯเปฯ ปญฺหาวาเรปิ สพฺพตฺถ เอกํ.)}
\csromanheader{2}{29. คนฺถคนฺถนิยทุก 1. กุสลตฺติก}
\csromanheader{2}{1-7. ปฏิจฺจาทิวาร ปจฺจยจตุกฺก-เหตุ}
\proseitem{24}{คนฺถนิยญฺเจว โน จ คนฺถํ กุสลํ ธมฺมํ ปฏิจฺจ คนฺถนิโย เจว โน จ คนฺโถ กุสโล ธมฺโม อุปฺปชฺชติ เหตุปจฺจยา. (1) (สํขิตฺตํ.)}

\gambhira{ธมฺมานุโลม ทุกติกปฏฺฐานปาฬิ}
\proseitem{25}{\csromanfolio{382}เหตุยา เอกํ, อารมฺมเณ เอกํ, อวิคเต เอกํ. (สหชาตวาเรปิ ฯเปฯ ปญฺหาวาเรปิ สพฺพตฺถ เอกํ.)}

\proseitem{26}{คนฺถญฺเจว คนฺถนิยญฺจ อกุสลํ ธมฺมํ ปฏิจฺจ คนฺโถ เจว คนฺถนิโย จ อกุสโล ธมฺโม อุปฺปชฺชติ เหตุปจฺจยา.\csromanspacer{} ตีณิ.}

\prose{คนฺถนิยญฺเจว โน จ คนฺถํ อกุสลํ ธมฺมํ ปฏิจฺจ คนฺถนิโย เจว โน จ คนฺโถ อกุสโล ธมฺโม อุปฺปชฺชติ เหตุปจฺจยา.\csromanspacer{} ตีณิ.}

\prose{คนฺถญฺเจว คนฺถนิยํ อกุสลญฺจ คนฺถนิยญฺเจว โน จ คนฺถํ อกุสลญฺจ ธมฺมํ ปฏิจฺจ คนฺโถ เจว คนฺถนิโย จ อกุสโล ธมฺโม อุปฺปชฺชติ เหตุปจฺจยา.\csromanspacer{} ตีณิ. (สํขิตฺตํ.)}

\proseitem{27}{เหตุยา นว, อารมฺมเณ นว, อธิปติยา นว ฯเปฯ อวิคเต นว. (สํขิตฺตํ.)}

\prose{นเหตุยา เอกํ, น-อธิปติยา นว ฯเปฯ นกมฺเม ตีณิ ฯเปฯ นวิปฺปยุตฺเต นว. (สํขิตฺตํ.)}

\prose{(สหชาตวารมฺปิ ฯเปฯ สมฺปยุตฺตวารมฺปิ ปฏิจฺจวารสทิสํ วิตฺถาเรตพฺพํ.)}

\proseitem{28}{คนฺโถ เจว คนฺถนิโย จ อกุสโล ธมฺโม คนฺถสฺส เจว คนฺถนิยสฺส จ อกุสลสฺส ธมฺมสฺส เหตุปจฺจเยน ปจฺจโย. (สํขิตฺตํ.)}

\proseitem{29}{เหตุยา นว, อารมฺมเณ นว, อธิปติยา นว ฯเปฯ กมฺเม อาหาเร อินฺทฺริเย ฌาเน ตีณิ ฯเปฯ อวิคเต นว. (สํขิตฺตํ.)}

\proseitem{30}{นเหตุยา นว, น-อารมฺมเณ นว. (สํขิตฺตํ.)}

\csromanheader{2}{เหตุปจฺจยา น-อารมฺมเณ นว. (สํขิตฺตํ.)}
\csromanheader{2}{นเหตุปจฺจยา อารมฺมเณ นว. (สํขิตฺตํ.)}
\prose{(ยถา กุสลตฺติเก ปญฺหาวารสฺส อนุโลมมฺปิ ปจฺจนียมฺปิ อนุโลมปจฺจนียมฺปิ ปจฺจนียานุโลมมฺปิ คณิตํ, เอวํ คเณตพฺพํ.)}

\proseitem{31}{คนฺถนิยญฺเจว โน จ คนฺถํ อพฺยากตํ ธมฺมํ ปฏิจฺจ คนฺถนิโย เจว โน จ คนฺโถ อพฺยากโต ธมฺโม อุปฺปชฺชติ เหตุปจฺจยา. (สํขิตฺตํ.)}

\proseitem{32}{เหตุยา เอกํ, อารมฺมเณ เอกํ ฯเปฯ อวิคเต เอกํ. (สํขิตฺตํ.) (สหชาตวาเรปิ ฯเปฯ ปญฺหาวาเรปิ สพฺพตฺถ เอกํ.)}

\vspace{\tipitakaparskip}
\csromancenter{คนฺถวิปฺปยุตฺตคนฺถนิยทุก 1.\csromanspacer{} กุสลตฺติก \textbf{30.}\csromanspacer{}\textbf{ คนฺถคนฺถสมฺปยุตฺตทุก 1.}\csromanspacer{}\textbf{ กุสลตฺติก}}
\csromanheader{2}{1-7. ปฏิจฺจาทิวาร ปจฺจยจตุกฺก-เหตุ}
\proseitem{33}{\csromanfolio{383}คนฺถญฺเจว คนฺถสมฺปยุตฺตญฺจ อกุสลํ ธมฺมํ ปฏิจฺจ คนฺโถ เจว คนฺถสมฺปยุตฺโต จ อกุสโล ธมฺโม อุปฺปชฺชติ เหตุปจฺจยา. (สํขิตฺตํ.)}

\proseitem{34}{เหตุยา นว, อารมฺมเณ นว ฯเปฯ อวิคเต นว. (สํขิตฺตํ.)}

\prose{น-อธิปติยา นว, นปุเรชาเต นว, นกมฺเม ตีณิ ฯเปฯ นวิปฺปยุตฺเต นว. (สํขิตฺตํ.) (สหชาตวาโรปิ ฯเปฯ สมฺปยุตฺตวาโรปิ ปฏิจฺจวารสทิสา.)}

\proseitem{35}{คนฺโถ เจว คนฺถสมฺปยุตฺโต จ อกุสโล ธมฺโม คนฺถสฺส เจว คนฺถสมฺปยุตฺตสฺส จ อกุสลสฺส ธมฺมสฺส เหตุปจฺจเยน ปจฺจโย. (สํขิตฺตํ.)}

\proseitem{36}{เหตุยา นว, อารมฺมเณ นว ฯเปฯ กมฺเม อาหาเร อินฺทฺริเย ฌาเน ตีณิ ฯเปฯ อวิคเต นว. (สํขิตฺตํ.)}

\proseitem{37}{นเหตุยา นว, น-อารมฺมเณ นว. (สํขิตฺตํ.)}

\csromanheader{2}{เหตุปจฺจยา น-อารมฺมเณ นว. (สํขิตฺตํ.)}
\csromanheader{2}{นเหตุปจฺจยา อารมฺมเณ นว. (สํขิตฺตํ.)}
\prose{(ยถา กุสลตฺติเก ปญฺหาวารสฺส อนุโลมมฺปิ ปจฺจนียมฺปิ อนุโลมปจฺจนียมฺปิ ปจฺจนียานุโลมมฺปิ คณิตํ, เอวํ คเณตพฺพํ.)}

\csromanheader{2}{31. คนฺถวิปฺปยุตฺตคนฺถนิยทุก 1. กุสลตฺติก}
\csromanheader{2}{1-7. ปฏิจฺจาทิวาร ปจฺจยจตุกฺก-เหตุ}
\proseitem{38}{คนฺถวิปฺปยุตฺตํ คนฺถนิยํ กุสลํ ธมฺมํ ปฏิจฺจ อนฺถวิปฺปยุตฺโต คนฺถนิโย กุสโล ธมฺโม อุปฺปชฺชติ เหตุปจฺจยา. (1)}

\prose{คนฺถวิปฺปยุตฺตํ อคนฺถนิยํ กุสลํ ธมฺมํ ปฏิจฺจ คนฺถวิปฺปยุตฺโต อคนฺถนิโย กุสโล ธมฺโม อุปฺปชฺชติ เหตุปจฺจยา. (1) (สํขิตฺตํ.)}

\gambhira{ธมฺมานุโลม ทุกติกปฏฺฐานปาฬิ}
\proseitem{39}{\csromanfolio{384}เหตุยา เทฺว, อารมฺมเณ เทฺว ฯเปฯ อวิคเต เทฺว. (สํขิตฺตํ.) (สหชาตวาโรปิ ฯเปฯ ปญฺหาวาโรปิ สพฺพตฺถ วิตฺถาเรตพฺพา.)}

\proseitem{40}{คนฺถวิปฺปยุตฺตํ คนฺถนิยํ อกุสลํ ธมฺมํ ปฏิจฺจ คนฺถวิปฺปยุตฺโต คนฺถนิโย อกุสโล ธมฺโม อุปฺปชฺชติ เหตุปจฺจยา. (สํขิตฺตํ.)}

\proseitem{41}{เหตุยา เอกํ ฯเปฯ อวิคเต เอกํ. (สํขิตฺตํ.) (สหชาตวาเรปิ ฯเปฯ ปญฺหาวาเรปิ สพฺพตฺถ เอกํ.)}

\proseitem{42}{คนฺถวิปฺปยุตฺตํ คนฺถนิยํ อพฺยากตํ ธมฺมํ ปฏิจฺจ คนฺถวิปฺปยุตฺโต อคนฺถนิโย อพฺยากโต ธมฺโม อุปฺปชฺชติ เหตุปจฺจยา. (1)}

\prose{คนฺถวิปฺปยุตฺตํ อคนฺถนิยํ อพฺยากตํ ธมฺมํ ปฏิจฺจ คนฺถวิปฺปยุตฺโต อคนฺถนิโย อพฺยากโต ธมฺโม อุปฺปชฺชติ เหตุปจฺจยา.\csromanspacer{} ตีณิ.}

\prose{คนฺถวิปฺปยุตฺตํ คนฺถนิยํ อพฺยากตญฺจ คนฺถวิปฺปยุตฺตํ อคนฺถนิยํ อพฺยากตญฺจ ธมฺมํ ปฏิจฺจ คนฺถวิปฺปยุตฺโต คนฺถนิโย อพฺยากโต ธมฺโม อุปฺปชฺชติ เหตุปจฺจยา. (1) (สํขิตฺตํ.)}

\proseitem{43}{เหตุยา ปญฺจ, อารมฺมเณ เทฺว ฯเปฯ วิปาเก ปญฺจ ฯเปฯ อวิคเต ปญฺจ. (สํขิตฺตํ.\csromanspacer{} สหชาตวารมฺปิ ฯเปฯ สมฺปยุตฺตวารมฺปิ ปฏิจฺจวารสทิสํ.)}

\proseitem{44}{คนฺถวิปฺปยุตฺโต คนฺถนิโย อพฺยากโต ธมฺโม คนฺถวิปฺปยุตฺตสฺส คนฺถนิยสฺส อพฺยากตสฺส ธมฺมสฺส เหตุปจฺจเยน ปจฺจโย. (1)}

\prose{คนฺถวิปฺปยุตฺโต อคนฺถนิโย อพฺยากโต ธมฺโม คนฺถวิปฺปยุตฺตสฺส อคนฺถนิยสฺส อพฺยากตสฺส ธมฺมสฺส เหตุปจฺจเยน ปจฺจโย.\csromanspacer{} ตีณิ. (สํขิตฺตํ.)}

\proseitem{45}{เหตุยา จตฺตาริ, อารมฺมเณ ตีณิ, อธิปติยา จตฺตาริ ฯเปฯ อวิคเต สตฺต. (สํขิตฺตํ.)}

\proseitem{46}{นเหตุยา สตฺต, น-อารมฺมเณ สตฺต. (สํขิตฺตํ.)}

\csromanmatikaanchor{mtk.71}
\csromanmatikaanchor{mtk.72}
\csromantocmarknopage{subparagraph}{44. นีวรณทุก}{mtk.71}
\bookmark[dest={mtk.71},level=5]{44. นีวรณทุก}
\csromantocmarkref{subparagraph}{1. กุสลตฺติก ...}{mtk.72}
\bookmark[dest={mtk.72},level=5]{1. กุสลตฺติก ...}
\csromanheader{6}{44. นีวรณทุก 1. กุสลตฺติก}
\csromanheader{2}{เหตุปจฺจยา น-อารมฺมเณ จตฺตาริ. (สํขิตฺตํ.)}
\csromanheader{2}{นเหตุปจฺจยา อารมฺมเณ ตีณิ. (สํขิตฺตํ.)}
\prose{\csromanfolio{385}(ยถา กุสลตฺติเก ปญฺหาวารสฺส อนุโลมมฺปิ ปจฺจนียมฺปิ อนุโลมปจฺจนียมฺปิ ปจฺจนียานุโลมมฺปิ คณิตํ, เอวํ คเณตพฺพํ.)}

\nitthitam{คนฺถโคจฺฉกกุสลตฺติกํ นิฏฺฐิตํ.}
\vspace{\tipitakaparskip}
\csromancenter{(โอฆโคจฺฉกมฺปิ โยคโคจฺฉกมฺปิ อาสวโคจฺฉก- กุสลตฺติกสทิสํ.)}
\csromanheader{2}{44. นีวรณทุก 1. กุสลตฺติก}
\csromanheader{2}{1-7. ปฏิจฺจาทิวาร ปจฺจยจตุกฺก-เหตุ}
\proseitem{1}{โนนีวรณํ กุสลํ ธมฺมํ ปฏิจฺจ โนนีวรโณ กุสโล ธมฺโม อุปฺปชฺชติ เหตุปจฺจยา. (สํขิตฺตํ.)}

\proseitem{2}{เหตุยา เอกํ, อารมฺมเณ เอกํ ฯเปฯ อวิคเต เอกํ. (สํขิตฺตํ.)}

\prose{(สหชาตวาเรปิ ฯเปฯ ปญฺหาวาเรปิ สพฺพตฺถ เอกํ.)}

\proseitem{3}{นีวรณํ อกุสลํ ธมฺมํ ปฏิจฺจ นีวรโณ อกุสโล ธมฺโม อุปฺปชฺชติ เหตุปจฺจยา.\csromanspacer{}\csromanspacer{}นีวรณํ อกุสลํ ธมฺมํ ปฏิจฺจ โนนีวรโณ อกุสโล ธมฺโม อุปฺปชฺชติ เหตุปจฺจยา.\csromanspacer{}\csromanspacer{}นีวรณํ อกุสลํ ธมฺมํ ปฏิจฺจ นีวรโณ อกุสโล จ โนนีวรโณ อกุสโล จ ธมฺมา อุปฺปชฺชนฺติ เหตุปจฺจยา. (3) (สํขิตฺตํ.)}

\proseitem{4}{เหตุยา นว, อารมฺมเณ นว ฯเปฯ อวิคเต นว. (สํขิตฺตํ.)}

\csromanheader{2}{\textbf{ปจฺจนีย-นเหตุ}}
\proseitem{5}{นีวรณํ อกุสลํ ธมฺมํ ปฏิจฺจ นีวรโณ อกุสโล ธมฺโม อุปฺปชฺชติ นเหตุปจฺจยา.\csromanspacer{} วิจิกิจฺฉานีวรณํ อุทฺธจฺจนีวรณํ ปฏิจฺจ อวิชฺชานีวรณํ. (1)}

\prose{โนนีวรณํ อกุสลํ ธมฺมํ ปฏิจฺจ นีวรโณ อกุสโล ธมฺโม อุปฺปชฺชติ นเหตุปจฺจยา.\csromanspacer{} วิจิกิจฺฉาสหคเต อุทฺธจฺจสหคเต ขนฺเธ ปฏิจฺจ อวิชฺชานีวรณํ. (1)}

\gambhira{ธมฺมานุโลม ทุกติกปฏฺฐานปาฬิ}
\prose{\csromanfolio{386}นีวรณญฺจ โนนีวรณญฺจ อกุสลํ ธมฺมํ ปฏิจฺจ นีวรโณ อกุสโล ธมฺโม อุปฺปชฺชติ นเหตุปจฺจยา.\csromanspacer{} วิจิกิจฺฉานีวรณญฺจ อุทฺธจฺจนีวรณญฺจ สมฺปยุตฺตเก จ ขนฺเธ ปฏิจฺจ อวิชฺชานีวรณํ. (1) (สํขิตฺตํ.)}

\proseitem{6}{นเหตุยา ตีณิ, น-อธิปติยา นว, ปุเรชาเต นว ฯเปฯ นกมฺเม ตีณิ ฯเปฯ นวิปฺปยุตฺเต นว. (สํขิตฺตํ.\csromanspacer{} สหชาตวาเรปิ สมฺปยุตฺตวาเรปิ สพฺพตฺถ นว.)}

\proseitem{7}{นีวรโณ อกุสโล ธมฺโม นีวรณสฺส อกุสลสฺส ธมฺมสฺส เหตุปจฺจเยน ปจฺจโย. (สํขิตฺตํ.)}

\proseitem{8}{เหตุยา ตีณิ, อารมฺมเณ นว, อธิปติยา นว ฯเปฯ อาเสวเน นว, กมฺเม อาหาเร อินฺทฺริเย ฌาเน มคฺเค ตีณิ, สมฺปยุตฺเต นว ฯเปฯ อวิคเต นว. (สํขิตฺตํ.)}

\prose{นเหตุยา นว, น-อารมฺมเณ นว. (สํขิตฺตํ.)}

\csromanheader{2}{เหตุปจฺจยา น-อารมฺมเณ ตีณิ. (สํขิตฺตํ.)}
\csromanheader{2}{นเหตุปจฺจยา อารมฺมเณ นว. (สํขิตฺตํ.)}
\prose{(ยถา กุสลตฺติเก ปญฺหาวารสฺส อนุโลมมฺปิ ปจฺจนียมฺปิ อนุโลมปจฺจนียมฺปิ ปจฺจนียานุโลมมฺปิ คณิตํ, เอวํ คเณตพฺพํ.)}

\proseitem{9}{โนนีวรณํ อพฺยากตํ ธมฺมํ ปฏิจฺจ โนนีวรโณ อพฺยากโต ธมฺโม อุปฺปชฺชติ เหตุปจฺจยา. (สํขิตฺตํ.)}

\proseitem{10}{เหตุยา เอกํ, อารมฺมเณ เอกํ ฯเปฯ อวิคเต เอกํ. (สํขิตฺตํ.)}

\vspace{\tipitakaparskip}
\csromancenter{(สหชาตวาเรปิ ฯเปฯ ปญฺหาวาเรปิ สพฺพตฺถ เอกํ.)}
\csromanheader{2}{45. นีวรณิยทุก 1. กุสลตฺติก}
\csromanheader{2}{1. ปฏิจฺจวาร ปจฺจยจตุกฺก-เหตุ}
\proseitem{11}{นีวรณิยํ กุสลํ ธมฺมํ ปฏิจฺจ นีวรณิโย กุสโล ธมฺโม อุปฺปชฺชติ เหตุปจฺจยา. (1)}

\csromanheader{2}{46. นีวรณสมฺปยุตฺตทุก 1. กุสลตฺติก}
\prose{\csromanfolio{387}อนีวรณิยํ กุสลํ ธมฺมํ ปฏิจฺจ อนีวรณิโย กุสโล ธมฺโม อุปฺปชฺชติ เหตุปจฺจยา. (1) (สํขิตฺตํ.)}

\proseitem{12}{เหตุยา เทฺว, อารมฺมเณ เทฺว ฯเปฯ อวิคเต เทฺว. (สํขิตฺตํ.)}

\prose{(สหชาตวาโรปิ ฯเปฯ ปญฺหาวาโรปิ วิตฺถาเรตพฺพา.)}

\proseitem{13}{นีวรณิยํ อกุสลํ ธมฺมํ ปฏิจฺจ นีวรณิโย อกุสโล ธมฺโม อุปฺปชฺชติ เหตุปจฺจยา. (สํขิตฺตํ.)}

\proseitem{14}{เหตุยา เอกํ, อารมฺมเณ เอกํ ฯเปฯ อวิคเต เอกํ. (สํขิตฺตํ.) (สหชาตวาเรปิ ฯเปฯ ปญฺหาวาเรปิ สพฺพตฺถ เอกํ.)}

\proseitem{15}{นีวรณิยํ อพฺยากตํ ธมฺมํ ปฏิจฺจ นีวรณิโย อพฺยากโต ธมฺโม อุปฺปชฺชติ เหตุปจฺจยา. (1)}

\prose{อนีวรณิยํ อพฺยากตํ ธมฺมํ ปฏิจฺจ อนีวรณิโย อพฺยากโต ธมฺโม อุปฺปชฺชติ เหตุปจฺจยา.\csromanspacer{} ตีณิ.}

\prose{นีวรณิยํ อพฺยากตญฺจ อนีวรณิยํ อพฺยากตญฺจ ธมฺมํ ปฏิจฺจ นีวรณิโย อพฺยากโต ธมฺโม อุปฺปชฺชติ เหตุปจฺจยา. (1) (สํขิตฺตํ.)}

\proseitem{16}{เหตุยา ปญฺจ, อารมฺมเณ เทฺว ฯเปฯ อาเสวเน เอกํ ฯเปฯ วิปาเก ปญฺจ ฯเปฯ อวิคเต ปญฺจ. (สํขิตฺตํ.\csromanspacer{} สหชาตวาโรปิ ฯเปฯ ปญฺหาวาโรปิ วิตฺถาเรตพฺพา.)}

\csromanheader{2}{46. นีวรณสมฺปยุตฺตทุก 1. กุสลตฺติก}
\csromanheader{2}{1. ปฏิจฺจวาร ปจฺจยจตุกฺก-เหตุ}
\proseitem{17}{นีวรณวิปฺปยุตฺตํ กุสลํ ธมฺมํ ปฏิจฺจ นีวรณวิปฺปยุตฺโต กุสโล ธมฺโม อุปฺปชฺชติ เหตุปจฺจยา. (สํขิตฺตํ.)}

\proseitem{18}{เหตุยา เอกํ, อารมฺมเณ เอกํ ฯเปฯ อวิคเต เอกํ. (สํขิตฺตํ.)}

\vspace{\tipitakaparskip}
\csromancenter{(สหชาตวาเรปิ ฯเปฯ ปญฺหาวาเรปิ สพฺพตฺถ เอกํ.)}
\gambhira{ธมฺมานุโลม ทุกติกปฏฺฐานปาฬิ}
\proseitem{19}{\csromanfolio{388}นีวรณสมฺปยุตฺตํ อกุสลํ ธมฺมํ ปฏิจฺจ นีวรณสมฺปยุตฺโต อกุสโล ธมฺโม อุปฺปชฺชติ เหตุปจฺจยา. (สํขิตฺตํ.)}

\proseitem{20}{เหตุยา เอกํ, อารมฺมเณ เอกํ ฯเปฯ อวิคเต เอกํ. (สํขิตฺตํ.) (สหชาตวาเรปิ ฯเปฯ ปญฺหาวาเรปิ สพฺพตฺถ เอกํ.)}

\proseitem{21}{นีวรณวิปฺปยุตฺตํ อพฺยากตํ ธมฺมํ ปฏิจฺจ นีวรณวิปฺปยุตฺโต อพฺยากโต ธมฺโม อุปฺปชฺชติ เหตุปจฺจยา. (สํขิตฺตํ.)}

\proseitem{22}{เหตุยา เอกํ, อารมฺมเณ เอกํ ฯเปฯ อวิคเต เอกํ. (สํขิตฺตํ.)}

\vspace{\tipitakaparskip}
\csromancenter{(สหชาตวาเรปิ ฯเปฯ ปญฺหาวาเรปิ สพฺพตฺถ เอกํ.)}
\csromanheader{2}{47. นีวรณนีวรณิยทุก 1. กุสลตฺติก}
\csromanheader{2}{1. ปฏิจฺจวาร ปจฺจยจตุกฺก-เหตุ}
\proseitem{23}{นีวรณิยญฺเจว โน จ นีวรณํ กุสลํ ธมฺมํ ปฏิจฺจ นีวรณิโย เจว โน จ นีวรโณ กุสโล ธมฺโม อุปฺปชฺชติ เหตุปจฺจยา. (สํขิตฺตํ.)}

\proseitem{24}{เหตุยา เอกํ, อารมฺมเณ เอกํ ฯเปฯ อวิคเต เอกํ. (สํขิตฺตํ.) (สหชาตวาเรปิ ฯเปฯ ปญฺหาวาเรปิ สพฺพตฺถ เอกํ.)}

\proseitem{25}{นีวรณญฺเจว นีวรณิยญฺจ อกุสลํ ธมฺมํ ปฏิจฺจ นีวรโณ เจว นีวรณิโย จ อกุสโล ธมฺโม อุปฺปชฺชติ เหตุปจฺจยา.\csromanspacer{} ตีณิ.}

\prose{นีวรณิยญฺเจว โน จ นีวรณํ อกุสลํ ธมฺมํ ปฏิจฺจ นีวรณิโย เจว โน จ นีวรโณ อกุสโล ธมฺโม อุปฺปชฺชติ เหตุปจฺจยา.\csromanspacer{} ตีณิ.}

\prose{นีวรณญฺเจว นีวรณิยํ อกุสลญฺจ นีวรณิยญฺเจว โน จ นีวรณํ อกุสลญฺจ ธมฺมํ ปฏิจฺจ นีวรโณ จ นีวรณิโย จ อกุสโล ธมฺโม อุปฺปชฺชติ เหตุปจฺจยา.\csromanspacer{} ตีณิ. (สํขิตฺตํ.)}

\proseitem{26}{เหตุยา นว, อารมฺมเณ นว ฯเปฯ อวิคเต นว. (สํขิตฺตํ.)}

\vspace{\tipitakaparskip}
\csromancenter{(สหชาตวาเรปิ ฯเปฯ ปญฺหาวาเรปิ สพฺพตฺถ นว.)}
\csromanheader{2}{49. นีวรณวิปฺปยุตฺตนีวรณิยทุก 1. กุสลตฺติก}
\proseitem{27}{\csromanfolio{389}นีวรณิยญฺเจว โน จ นีวรณํ อพฺยากตํ ธมฺมํ ปฏิจฺจ นีวรณิโย เจว โน จ นีวรโณ อพฺยากโต ธมฺโม อุปฺปชฺชติ เหตุปจฺจยา. (สํขิตฺตํ.)}

\proseitem{28}{เหตุยา เอกํ, อารมฺมเณ เอกํ ฯเปฯ อวิคเต เอกํ. (สํขิตฺตํ.)}

\vspace{\tipitakaparskip}
\csromancenter{(สหชาตวาเรปิ ฯเปฯ ปญฺหาวาเรปิ สพฺพตฺถ เอกํ.)}
\csromanheader{2}{48. นีวรณนีวรณสมฺปยุตฺตทุก 1. กุสลตฺติก}
\csromanheader{2}{1. ปฏิจฺจวาร ปจฺจยจตุกฺก-เหตุ}
\proseitem{29}{นีวรณญฺเจว นีวรณสมฺปยุตฺตญฺจ อกุสลํ ธมฺมํ ปฏิจฺจ นีวรโณ เจว นีวรณสมฺปยุตฺโต จ อกุสโล ธมฺโม อุปฺปชฺชติ เหตุปจฺจยา. (สํขิตฺตํ.)}

\proseitem{30}{เหตุยา นว, อารมฺมเณ นว ฯเปฯ อวิคเต นว. (สํขิตฺตํ.)}

\vspace{\tipitakaparskip}
\csromancenter{(สหชาตวาเรปิ ฯเปฯ ปญฺหาวาเรปิ สพฺพตฺถ นว.)}
\csromanheader{2}{49. นีวรณวิปฺปยุตฺตนีวรณิยทุก 1. กุสลตฺติก}
\csromanheader{2}{1. ปฏิจฺจวาร ปจฺจยจตุกฺก-เหตุ}
\proseitem{31}{นีวรณวิปฺปยุตฺตํ นีวรณิยํ กุสลํ ธมฺมํ ปฏิจฺจ นีวรณวิปฺปยุตฺโต นีวรณิโย กุสโล ธมฺโม อุปฺปชฺชติ เหตุปจฺจยา. (1)}

\prose{นีวรณวิปฺปยุตฺตํ อนีวรณิยํ กุสลํ ธมฺมํ ปฏิจฺจ นีวรณวิปฺปยุตฺโต อนีวรณิโย กุสโล ธมฺโม อุปฺปชฺชติ เหตุปจฺจยา. (1) (สํขิตฺตํ.)}

\proseitem{32}{เหตุยา เทฺว, อารมฺมเณ เทฺว ฯเปฯ อวิคเต เทฺว. (สํขิตฺตํ.) (สหชาตวาเรปิ ฯเปฯ ปญฺหาวาเรปิ วิตฺถาเรตพฺพา.)}

\proseitem{33}{นีวรณวิปฺปยุตฺตํ นีวรณิยํ อพฺยากตํ ธมฺมํ ปฏิจฺจ นีวรณวิปฺปยุตฺโต นีวรณิโย อพฺยากโต ธมฺโม อุปฺปชฺชติ เหตุปจฺจยา. (1)}

\gambhira{ธมฺมานุโลม ทุกติกปฏฺฐานปาฬิ}
\prose{\csromanfolio{390}นีวรณวิปฺปยุตฺตํ อนีวรณิยํ อพฺยากตํ ธมฺมํ ปฏิจฺจ นีวรณวิปฺปยุตฺโต อนีวรณิโย อพฺยากโต ธมฺโม อุปฺปชฺชติ เหตุปจฺจยา.\csromanspacer{} ตีณิ.}

\prose{นีวรณวิปฺปยุตฺตํ นีวรณิยํ อพฺยากตญฺจ นีวรณวิปฺปยุตฺตํ อนีวรณิยํ อพฺยากตญฺจ ธมฺมํ ปฏิจฺจ นีวรณวิปฺปยุตฺโต นีวรณิโย อพฺยากโต ธมฺโม อุปฺปชฺชติ เหตุปจฺจยา. (1) (สํขิตฺตํ.)}

\proseitem{34}{เหตุยา ปญฺจ, อารมฺมเณ เทฺว ฯเปฯ วิปาเก ปญฺจ. (สํขิตฺตํ.\csromanspacer{} สหชาตวาเรปิ ฯเปฯ ปญฺหาวาเรปิ สพฺพตฺถ วิตฺถาเรตพฺพํ.)}

\nitthitam{นีวรณโคจฺฉกกุสลตฺติกํ นิฏฺฐิตํ.}
\csromanmatikaanchor{mtk.73}
\csromanmatikaanchor{mtk.74}
\csromantocmarknopage{subparagraph}{50. ปรามาสทุก}{mtk.73}
\bookmark[dest={mtk.73},level=5]{50. ปรามาสทุก}
\csromantocmarkref{subparagraph}{1. กุสลตฺติก ...}{mtk.74}
\bookmark[dest={mtk.74},level=5]{1. กุสลตฺติก ...}
\csromanheader{6}{50. ปรามาสทุก 1. กุสลตฺติก}
\csromanheader{2}{1-7. ปฏิจฺจาทิวาร ปจฺจยจตุกฺก-เหตุ}
\proseitem{1}{โนปรามาสํ กุสลํ ธมฺมํ ปฏิจฺจ โนปรามาโส กุสโล ธมฺโม อุปฺปชฺชติ เหตุปจฺจยา. (สํขิตฺตํ.)}

\proseitem{2}{เหตุยา เอกํ, อารมฺมเณ เอกํ ฯเปฯ อวิคเต เอกํ. (สํขิตฺตํ.) (สหชาตวาเรปิ ฯเปฯ ปญฺหาวาเรปิ สพฺพตฺถ เอกํ.)}

\proseitem{3}{ปรามาสํ อกุสลํ ธมฺมํ ปฏิจฺจ โนปรามาโส อกุสโล ธมฺโม อุปฺปชฺชติ เหตุปจฺจยา. (1)}

\prose{โนปรามาสํ อกุสลํ ธมฺมํ ปฏิจฺจ โนปรามาโส อกุสโล ธมฺโม อุปฺปชฺชติ เหตุปจฺจยา.\csromanspacer{} ตีณิ.}

\prose{ปรามาสํ อกุสลญฺจ โนปรามาสํ อกุสลญฺจ ธมฺมํ ปฏิจฺจ โนปรามาโส อกุสโล ธมฺโม อุปฺปชฺชติ เหตุปจฺจยา. (1) (สํขิตฺตํ.)}

\proseitem{4}{เหตุยา ปญฺจ, อารมฺมเณ ปญฺจ ฯเปฯ อวิคเต ปญฺจ. (สํขิตฺตํ.)}

\prose{นเหตุยา เอกํ, น-อธิปติยา ปญฺจ, นปุเรชาเต ปญฺจ, นปจฺฉาชาเต ปญฺจ ฯเปฯ. (สํขิตฺตํ.\csromanspacer{} สหชาตวารมฺปิ ฯเปฯ สมฺปยุตฺตวารมฺปิ วิตฺถาเรตพฺพํ.)}

\csromanheader{2}{51. ปรามฏฺฐทุก 1. กุสลตฺติก}
\proseitem{5}{\csromanfolio{391}โนปรามาโส อกุสโล ธมฺโม โนปรามาสสฺส อกุสลสฺส ธมฺมสฺส เหตุปจฺจเยน ปจฺจโย. (สํขิตฺตํ.)}

\proseitem{6}{เหตุยา ตีณิ, อารมฺมเณ นว, อธิปติยา นว ฯเปฯ อวิคเต ปญฺจ. (สํขิตฺตํ.)}

\prose{(ยถา กุสลตฺติเก ปญฺหาวารสฺส อนุโลมมฺปิ ปจฺจนียมฺปิ อนุโลมปจฺจนียมฺปิ ปจฺจนียานุโลมมฺปิ คณิตํ, เอวํ คเณตพฺพํ.)}

\proseitem{7}{โนปรามาสํ อพฺยากตํ ธมฺมํ ปฏิจฺจ โนปรามาโส อพฺยากโต ธมฺโม อุปฺปชฺชติ เหตุปจฺจยา. (สํขิตฺตํ.)}

\proseitem{8}{เหตุยา เอกํ, อารมฺมเณ เอกํ ฯเปฯ อวิคเต เอกํ. (สํขิตฺตํ.)}

\vspace{\tipitakaparskip}
\csromancenter{(สหชาตวาเรปิ ฯเปฯ ปญฺหาวาเรปิ สพฺพตฺถ เอกํ.)}
\csromanheader{2}{51. ปรามฏฺฐทุก 1. กุสลตฺติก}
\csromanheader{2}{1. ปฏิจฺจวาร ปจฺจยจตุกฺก-เหตุ}
\proseitem{9}{ปรามฏฺฐํ กุสลํ ธมฺมํ ปฏิจฺจ ปรามฏฺโฐ กุสโล ธมฺโม อุปฺปชฺชติ เหตุปจฺจยา. (1)}

\prose{อปรามฏฺฐํ กุสลํ ธมฺมํ ปฏิจฺจ อปรามฏฺโฐ กุสโล ธมฺโม อุปฺปชฺชติ เหตุปจฺจยา. (1) (สํขิตฺตํ.)}

\proseitem{10}{เหตุยา เทฺว, อารมฺมเณ เทฺว ฯเปฯ อวิคเต เทฺว. (สํขิตฺตํ.)}

\prose{(สหชาตวาโรปิ ฯเปฯ ปญฺหาวาโรปิ วิตฺถาเรตพฺพา.)}

\proseitem{11}{ปรามฏฺฐํ อกุสลํ ธมฺมํ ปฏิจฺจ ปรามฏฺโฐ อกุสโล ธมฺโม อุปฺปชฺชติ เหตุปจฺจยา. (สํขิตฺตํ.)}

\proseitem{12}{เหตุยา เอกํ, อารมฺมเณ เอกํ ฯเปฯ อวิคเต เอกํ. (สํขิตฺตํ.)}

\vspace{\tipitakaparskip}
\csromancenter{(สหชาตวาเรปิ ฯเปฯ ปญฺหาวาเรปิ สพฺพตฺถ เอกํ.)}
\gambhira{ธมฺมานุโลม ทุกติกปฏฺฐานปาฬิ}
\proseitem{13}{\csromanfolio{392}ปรามฏฺฐํ อพฺยากตํ ธมฺมํ ปฏิจฺจ ปรามฏฺโฐ อพฺยากโต ธมฺโม อุปฺปชฺชติ เหตุปจฺจยา. (1)}

\prose{อปรามฏฺฐํ อพฺยากตํ ธมฺมํ ปฏิจฺจ อปรามฏฺโฐ อพฺยากโต ธมฺโม อุปฺปชฺชติ เหตุปจฺจยา.\csromanspacer{} ตีณิ.}

\prose{ปรามฏฺฐํ อพฺยากตญฺจ อปรามฏฺฐํ อพฺยากตญฺจ ธมฺมํ ปฏิจฺจ ปรามฏฺโฐ อพฺยากโต ธมฺโม อุปฺปชฺชติ เหตุปจฺจยา. (1) (สํขิตฺตํ.)}

\proseitem{14}{เหตุยา ปญฺจ, อารมฺมเณ เทฺว ฯเปฯ อาเสวเน เอกํ ฯเปฯ อวิคเต ปญฺจ. (สํขิตฺตํ.\csromanspacer{} สหชาตวารมฺปิ ฯเปฯ ปญฺหาวารมฺปิ วิตฺถาเรตพฺพํ.)}

\csromanheader{2}{52. ปรามาสสมฺปยุตฺตทุก 1. กุสลตฺติก}
\csromanheader{2}{1-7. ปฏิจฺจาทิวาร ปจฺจยจตุกฺก-เหตุ}
\proseitem{15}{ปรามาสวิปฺปยุตฺตํ กุสลํ ธมฺมํ ปฏิจฺจ ปรามาสวิปฺปยุตฺโต กุสโล ธมฺโม อุปฺปชฺชติ เหตุปจฺจยา. (สํขิตฺตํ.)}

\proseitem{16}{เหตุยา เอกํ, อารมฺมเณ เอกํ ฯเปฯ อวิคเต เอกํ. (สํขิตฺตํ.)}

\prose{(สหชาตวาเรปิ ฯเปฯ ปญฺหาวาเรปิ สพฺพตฺถ เอกํ.)}

\proseitem{17}{ปรามาสสมฺปยุตฺตํ อกุสลํ ธมฺมํ ปฏิจฺจ ปรามาสสมฺปยุตฺโต อกุสโล ธมฺโม อุปฺปชฺชติ เหตุปจฺจยา. (1)}

\prose{ปรามาสวิปฺปยุตฺตํ อกุสลํ ธมฺมํ ปฏิจฺจ ปรามาสวิปฺปยุตฺโต อกุสโล ธมฺโม อุปฺปชฺชติ เหตุปจฺจยา. (1) (สํขิตฺตํ.)}

\proseitem{18}{เหตุยา เทฺว, อารมฺมเณ เทฺว, (สพฺพตฺถ เทฺว.) ฯเปฯ อวิคเต เทฺว. (สํขิตฺตํ.)}

\prose{นเหตุยา เอกํ, น-อธิปติยา เทฺว, (สพฺพตฺถ เทฺว.) ฯเปฯ นวิปฺปยุตฺเต เทฺว. (สํขิตฺตํ.\csromanspacer{} สหชาตวาโรปิ ฯเปฯ สมฺปยุตฺตวาโรปิ วิตฺถาเรตพฺพา.)}

\proseitem{19}{ปรามาสสมฺปยุตฺโต อกุสโล ธมฺโม ปรามาสสมฺปยุตฺตสฺส อกุสลสฺส ธมฺมสฺส เหตุปจฺจเยน ปจฺจโย. (1)}

\csromanheader{2}{53. ปรามาสปรามฏฺฐทุก 1. กุสลตฺติก}
\prose{\csromanfolio{393}ปรามาสวิปฺปยุตฺโต อกุสโล ธมฺโม ปรามาสวิปฺปยุตฺตสฺส อกุสลสฺส ธมฺมสฺส เหตุปจฺจเยน ปจฺจโย. (1) (สํขิตฺตํ.)}

\proseitem{20}{เหตุยา เทฺว, อารมฺมเณ จตฺตาริ, อธิปติยา จตฺตาริ, (มชฺฌิเม จ อนฺเต จ อารมฺมณาธิปติ.) อนนฺตเร เทฺว ฯเปฯ สหชาเต เทฺว ฯเปฯ อุปนิสฺสเย จตฺตาริ, อาเสวเน เทฺว, กมฺเม เทฺว, อาหาเร เทฺว ฯเปฯ มคฺเค เทฺว ฯเปฯ อวิคเต เทฺว. (สํขิตฺตํ.)}

\proseitem{21}{ปรามาสวิปฺปยุตฺตํ อพฺยากตํ ธมฺมํ ปฏิจฺจ ปรามาสวิปฺปยุตฺโต อพฺยากโต ธมฺโม อุปฺปชฺชติ เหตุปจฺจยา. (สํขิตฺตํ.)}

\proseitem{22}{เหตุยา เอกํ, อารมฺมเณ เอกํ ฯเปฯ อวิคเต เอกํ. (สํขิตฺตํ.)}

\vspace{\tipitakaparskip}
\csromancenter{(สหชาตวาเรปิ ฯเปฯ ปญฺหาวาเรปิ สพฺพตฺถ เอกํ.)}
\csromanheader{2}{53. ปรามาสปรามฏฺฐทุก 1. กุสลตฺติก}
\csromanheader{2}{1-7. ปฏิจฺจาทิวาร ปจฺจยจตุกฺก}
\proseitem{23}{ปรามฏฺฐญฺเจว โน จ ปรามาสํ กุสลํ ธมฺมํ ปฏิจฺจ ปรามฏฺโฐ เจว โน จ ปรามาโส กุสโล ธมฺโม อุปฺปชฺชติ เหตุปจฺจยา. (สํขิตฺตํ.)}

\proseitem{24}{เหตุยา เอกํ, อารมฺมเณ เอกํ ฯเปฯ อวิคเต เอกํ. (สํขิตฺตํ.) (สหชาตวาเรปิ ฯเปฯ ปญฺหาวาเรปิ สพฺพตฺถ เอกํ.)}

\proseitem{25}{ปรามาสญฺเจว ปรามฏฺฐญฺจ อกุสลํ ธมฺมํ ปฏิจฺจ ปรามฏฺโฐ เจว โน จ ปรามาโส อกุสโล ธมฺโม อุปฺปชฺชติ เหตุปจฺจยา.\csromanspacer{} เอกํ.}

\prose{ปรามฏฺฐญฺเจว โน จ ปรามาสํ อกุสลํ ธมฺมํ ปฏิจฺจ ปรามฏฺโฐ เจว โน จ ปรามาโส อกุสโล ธมฺโม อุปฺปชฺชติ เหตุปจฺจยา.\csromanspacer{} ตีณิ.}

\prose{ปรามาสญฺเจว ปรามฏฺฐํ อกุสลญฺจ ปรามฏฺฐญฺเจว โน จ ปรามาสํ อกุสลญฺจ ธมฺมํ ปฏิจฺจ ปรามฏฺโฐ เจว โน จ ปรามาโส อกุสโล ธมฺโม อุปฺปชฺชติ เหตุปจฺจยา.\csromanspacer{} เอกํ. (สํขิตฺตํ.)}

\csromanhangingbegin
\hangingheaditem{26}{เหตุยา ปญฺจ, อารมฺมเณ ปญฺจ ฯเปฯ อวิคเต ปญฺจ. (สํขิตฺตํ.)}
\hangingbody{(สหชาตวารมฺปิ ฯเปฯ สมฺปยุตฺตวารมฺปิ วิตฺถาเรตพฺพํ.)}
\csromanhangingend

\gambhira{ธมฺมานุโลม ทุกติกปฏฺฐานปาฬิ}
\proseitem{27}{\csromanfolio{394}ปรามฏฺโฐ เจว โน จ ปรามาโส อกุสโล ธมฺโม ปรามฏฺฐสฺส เจว โน จ ปรามาสสฺส อกุสลสฺส ธมฺมสฺส เหตุปจฺจเยน ปจฺจโย. (สํขิตฺตํ.)}

\proseitem{28}{เหตุยา ตีณิ, อารมฺมเณ นว, อธิปติยา นว ฯเปฯ อวิคเต ปญฺจ. (สํขิตฺตํ.)}

\prose{(ยถา กุสลตฺติเก ปญฺหาวารสฺส อนุโลมมฺปิ ปจฺจนียมฺปิ อนุโลมปจฺจนียมฺปิ ปจฺจนียานุโลมมฺปิ คณิตํ, เอวํ คเณตพฺพํ.)}

\proseitem{29}{ปรามฏฺฐญฺเจว โน จ ปรามาสํ อพฺยากตํ ธมฺมํ ปฏิจฺจ ปรามฏฺโฐ เจว โน จ ปรามาโส อพฺยากโต ธมฺโม อุปฺปชฺชติ เหตุปจฺจยา. (สํขิตฺตํ.)}

\proseitem{30}{เหตุยา เอกํ, อารมฺมเณ เอกํ ฯเปฯ อวิคเต เอกํ. (สํขิตฺตํ.)}

\vspace{\tipitakaparskip}
\csromancenter{(สหชาตวาเรปิ ฯเปฯ ปญฺหาวาเรปิ สพฺพตฺถ เอกํ.)}
\csromanheader{2}{54. ปรามาสวิปฺปยุตฺตปรามฏฺฐทุก 1. กุสลตฺติก}
\csromanheader{2}{1. ปฏิจฺจวาร ปจฺจยจตุกฺก-เหตุ}
\proseitem{31}{ปรามาสวิปฺปยุตฺตํ ปรามฏฺฐํ กุสลํ ธมฺมํ ปฏิจฺจ ปรามาสวิปฺปยุตฺโต ปรามฏฺโฐ กุสโล ธมฺโม อุปฺปชฺชติ เหตุปจฺจยา. (1)}

\prose{ปรามาสวิปฺปยุตฺตํ อปรามฏฺฐํ กุสลํ ธมฺมํ ปฏิจฺจ ปรามาสวิปฺปยุตฺโต อปรามฏฺโฐ กุสโล ธมฺโม อุปฺปชฺชติ เหตุปจฺจยา. (1) (สํขิตฺตํ.)}

\proseitem{32}{เหตุยา เทฺว, อารมฺมเณ เทฺว ฯเปฯ อวิคเต เทฺว. (สํขิตฺตํ.\csromanspacer{} โลกิยโลกุตฺตรทุกสทิสํ.\csromanspacer{} สหชาตวาโรปิ ฯเปฯ ปญฺหาวาโรปิ วิตฺถาเรตพฺพา.)}

\proseitem{33}{ปรามาสวิปฺปยุตฺตํ ปรามฏฺฐํ อกุสลํ ธมฺมํ ปฏิจฺจ ปรามาสวิปฺปยุตฺโต ปรามฏฺโฐ อกุสโล ธมฺโม อุปฺปชฺชติ เหตุปจฺจยา. (สพฺพตฺถ เอกํ.)}

\proseitem{34}{ปรามาสวิปฺปยุตฺตํ ปรามฏฺฐํ อพฺยากตํ ธมฺมํ ปฏิจฺจ ปรามาสวิปฺปยุตฺโต ปรามฏฺโฐ อพฺยากโต ธมฺโม อุปฺปชฺชติ เหตุปจฺจยา. (1)}

\csromanmatikaanchor{mtk.75}
\csromanmatikaanchor{mtk.76}
\csromantocmarknopage{subparagraph}{55. สารมฺมณทุก}{mtk.75}
\bookmark[dest={mtk.75},level=5]{55. สารมฺมณทุก}
\csromantocmarkref{subparagraph}{1. กุสลตฺติก ...}{mtk.76}
\bookmark[dest={mtk.76},level=5]{1. กุสลตฺติก ...}
\csromanheader{6}{55. สารมฺมณทุก 1. กุสลตฺติก}
\prose{\csromanfolio{395}ปรามาสวิปฺปยุตฺตํ อปรามฏฺฐํ อพฺยากตํ ธมฺมํ ปฏิจฺจ ปรามาสวิปฺปยุตฺโต อปรามฏฺโฐ อพฺยากโต ธมฺโม อุปฺปชฺชติ เหตุปจฺจยา.\csromanspacer{} ตีณิ.}

\prose{ปรามาสวิปฺปยุตฺตํ ปรามฏฺฐํ อพฺยากตญฺจ ปรามาสวิปฺปยุตฺตํ อปรามฏฺฐํ อพฺยากตญฺจ ธมฺมํ ปฏิจฺจ ปรามาสวิปฺปยุตฺโต ปรามฏฺโฐ อพฺยากโต ธมฺโม อุปฺปชฺชติ เหตุปจฺจยา. (1) (สํขิตฺตํ.)}

\proseitem{35}{เหตุยา ปญฺจ, อารมฺมเณ เทฺว ฯเปฯ อวิคเต ปญฺจ. (สํขิตฺตํ.\csromanspacer{} โลกิยทุกสทิสํ.\csromanspacer{} สหชาตวารมฺปิ ฯเปฯ ปญฺหาวารมฺปิ วิตฺถาเรตพฺพํ.)}

\nitthitam{ปรามาสโคจฺฉกกุสลตฺติกํ นิฏฺฐิตํ.}
\csromanheader{2}{55. สารมฺมณทุก 1. กุสลตฺติก}
\csromanheader{2}{1-7. ปฏิจฺจาทิวาร ปจฺจยจตุกฺก-เหตุ}
\proseitem{1}{สารมฺมณํ กุสลํ ธมฺมํ ปฏิจฺจ สารมฺมโณ กุสโล ธมฺโม อุปฺปชฺชติ เหตุปจฺจยา. (ยถา สปฺปจฺจยํ กุสลํ, เอวํ วิตฺถาเรตพฺพํ.)}

\proseitem{2}{สารมฺมณํ อกุสลํ ธมฺมํ ปฏิจฺจ สารมฺมโณ อกุสโล ธมฺโม อุปฺปชฺชติ เหตุปจฺจยา. (ยถา สปฺปจฺจยํ อกุสลํ, เอวํ กาตพฺพํ.)}

\proseitem{3}{สารมฺมณํ อพฺยากตํ ธมฺมํ ปฏิจฺจ สารมฺมโณ อพฺยากโต ธมฺโม อุปฺปชฺชติ เหตุปจฺจยา.\csromanspacer{} ตีณิ.}

\prose{อนารมฺมณํ อพฺยากตํ ธมฺมํ ปฏิจฺจ อนารมฺมโณ อพฺยากโต ธมฺโม อุปฺปชฺชติ เหตุปจฺจยา.\csromanspacer{} ตีณิ.}

\prose{สารมฺมณํ อพฺยากตญฺจ อนารมฺมณํ อพฺยากตญฺจ ธมฺมํ ปฏิจฺจ สารมฺมโณ อพฺยากโต ธมฺโม อุปฺปชฺชติ เหตุปจฺจยา.\csromanspacer{} ตีณิ. (สํขิตฺตํ.)}

\proseitem{4}{เหตุยา นว, อารมฺมเณ ตีณิ, อธิปติยา ปญฺจ ฯเปฯ อญฺญมญฺเญ ฉ ฯเปฯ ปุเรชาเต เอกํ, อาเสวเน เอกํ. (สํขิตฺตํ.)}

\prose{นเหตุยา นว, น-อารมฺมเณ ตีณิ, น-อธิปติยา นว ฯเปฯ นกมฺเม เทฺว, นวิปาเก ปญฺจ, น-อาหาเร เอกํ, น-อินฺทฺริเย เอกํ, นฌาเน เทฺว,}

\vspace{\tipitakaparskip}
\csromancenter{ธมฺมานุโลม ทุกติกปฏฺฐานปาฬิ นมคฺเค นว, นสมฺปยุตฺเต ตีณิ, นวิปฺปยุตฺเต เทฺว, โนนตฺถิยา ตีณิ, โนวิคเต ตีณิ. (สํขิตฺตํ.)}
\csromanheader{2}{\textbf{เหตุ อารมฺมณ อธิปติ}}
\proseitem{5}{\csromanfolio{396}สารมฺมโณ อพฺยากโต ธมฺโม สารมฺมณสฺส อพฺยากตสฺส ธมฺมสฺส เหตุปจฺจเยน ปจฺจโย.\csromanspacer{} ตีณิ.}

\proseitem{6}{สารมฺมโณ อพฺยากโต ธมฺโม สารมฺมณสฺส อพฺยากตสฺส ธมฺมสฺส อารมฺมณปจฺจเยน ปจฺจโย. (1)}

\prose{อนารมฺมโณ อพฺยากโต ธมฺโม สารมฺมณสฺส อพฺยากตสฺส ธมฺมสฺส อารมฺมณปจฺจเยน ปจฺจโย. (1)}

\proseitem{7}{สารมฺมโณ อพฺยากโต ธมฺโม สารมฺมณสฺส อพฺยากตสฺส ธมฺมสฺส อธิปติปจฺจเยน ปจฺจโย.\csromanspacer{} อารมฺมณาธิปติ, สหชาตาธิปติ.\csromanspacer{} ตีณิ.\csromanspacer{} อนารมฺมโณ อพฺยากโต ธมฺโม สารมฺมณสฺส อพฺยากตสฺส ธมฺมสฺส อธิปติปจฺจเยน ปจฺจโย.\csromanspacer{} อารมฺมณาธิปติ. (สํขิตฺตํ.)}

\proseitem{8}{เหตุยา ตีณิ, อารมฺมเณ เทฺว, อธิปติยา จตฺตาริ, อนนฺตเร เอกํ ฯเปฯ สหชาเต สตฺต, อญฺญมญฺเญ ฉ, นิสฺสเย สตฺต, อุปนิสฺสเย เทฺว, ปุเรชาเต เอกํ, ปจฺฉาชาเต เอกํ, อาเสวเน เอกํ, กมฺเม ตีณิ, วิปาเก ตีณิ, สมฺปยุตฺเต เอกํ, วิปฺปยุตฺเต เทฺว. (สํขิตฺตํ.)}

\csromanheader{2}{56. จิตฺตทุก 1. กุสลตฺติก}
\csromanheader{2}{1-7. ปฏิจฺจาทิวาร ปจฺจยจตุกฺก-เหตุ}
\proseitem{9}{จิตฺตํ กุสลํ ธมฺมํ ปฏิจฺจ โนจิตฺโต กุสโล ธมฺโม อุปฺปชฺชติ เหตุปจฺจยา. (1)}

\prose{โนจิตฺตํ กุสลํ ธมฺมํ ปฏิจฺจ โนจิตฺโต กุสโล ธมฺโม อุปฺปชฺชติ เหตุปจฺจยา.\csromanspacer{}\csromanspacer{}โนจิตฺตํ กุสลํ ธมฺมํ ปฏิจฺจ จิตฺโต กุสโล ธมฺโม อุปฺปชฺชติ เหตุปจฺจยา.\csromanspacer{}\csromanspacer{}โนจิตฺตํ กุสลํ ธมฺมํ ปฏิจฺจ จิตฺโต กุสโล จ โนจิตฺโต กุสโล จ ธมฺมา อุปฺปชฺชนฺติ เหตุปจฺจยา. (3)}

\csromanheader{2}{56. จิตฺตทุก 1. กุสลตฺติก}
\prose{\csromanfolio{397}จิตฺตํ กุสลญฺจ โนจิตฺตํ กุสลญฺจ ธมฺมํ ปฏิจฺจ โนจิตฺโต กุสโล ธมฺโม อุปฺปชฺชติ เหตุปจฺจยา. (1) (สํขิตฺตํ.)}

\proseitem{10}{เหตุยา ปญฺจ, อารมฺมเณ ปญฺจ. (สพฺพตฺถ ปญฺจ.) อวิคเต ปญฺจ. (สํขิตฺตํ.)}

\prose{น-อธิปติยา ปญฺจ, นปุเรชาเต ปญฺจ, นปจฺฉาชาเต ปญฺจ, น-อาเสวเน ปญฺจ, นกมฺเม ตีณิ, นวิปาเก ปญฺจ, นวิปฺปยุตฺเต ปญฺจ. (สํขิตฺตํ.) (สหชาตวารมฺปิ ฯเปฯ สมฺปยุตฺตวารมฺปิ วิตฺถาเรตพฺพํ.)}

\proseitem{11}{โนจิตฺโต กุสโล ธมฺโม โนจิตฺตสฺส กุสลสฺส ธมฺมสฺส เหตุปจฺจเยน ปจฺจโย.\csromanspacer{}\csromanspacer{}โนจิตฺโต กุสโล ธมฺโม จิตฺตสฺส กุสลสฺส ธมฺมสฺส เหตุปจฺจเยน ปจฺจโย.\csromanspacer{} โนจิตฺโต กุสโล ธมฺโม จิตฺตสฺส กุสลสฺส จ โนจิตฺตสฺส กุสลสฺส จ ธมฺมสฺส เหตุปจฺจเยน ปจฺจโย. (3) (สํขิตฺตํ.)}

\proseitem{12}{เหตุยา ตีณิ, อารมฺมเณ นว, อธิปติยา นว, อนนฺตเร นว ฯเปฯ สหชาเต อญฺญมญฺเญ นิสฺสเย ปญฺจ, อุปนิสฺสเย อาเสวเน นว, กมฺเม ตีณิ, อาหาเร อินฺทฺริเย ปญฺจ, ฌาเน มคฺเค ตีณิ, สมฺปยุตฺเต ปญฺจ ฯเปฯ อวิคเต ปญฺจ. (สํขิตฺตํ.)}

\prose{นเหตุยา นว, น-อารมฺมเณ นว. (สํขิตฺตํ.)}

\csromanheader{2}{เหตุปจฺจยา น-อารมฺมเณ ตีนิ. (สํขิตฺตํ.)}
\csromanheader{2}{นเหตุปจฺจยา อารมฺมเณ นว. (สํขิตฺตํ.)}
\prose{(ยถา กุสลตฺติเก ปญฺหาวารสฺส อนุโลมมฺปิ ปจฺจนียมฺปิ อนุโลมปจฺจนียมฺปิ ปจฺจนียานุโลมมฺปิ คณิตํ, เอวํ คเณตพฺพํ.)}

\proseitem{13}{จิตฺตํ อกุสลํ ธมฺมํ ปฏิจฺจ โนจิตฺโต อกุสโล ธมฺโม อุปฺปชฺชติ เหตุปจฺจยา. (1)}

\prose{โนจิตฺตํ อกุสลํ ธมฺมํ ปฏิจฺจ โนจิตฺโต อกุสโล ธมฺโม อุปฺปชฺชติ เหตุปจฺจยา.\csromanspacer{}\csromanspacer{}โนจิตฺตํ อกุสลํ ธมฺมํ ปฏิจฺจ จิตฺโต อกุสโล ธมฺโม อุปฺปชฺชติ เหตุปจฺจยา.\csromanspacer{}\csromanspacer{}โนจิตฺตํ อกุสลํ ธมฺมํ ปฏิจฺจ จิตฺโต อกุสโล จ โนจิตฺโต อกุสโล จ ธมฺมา อุปฺปชฺชนฺติ เหตุปจฺจโย. (3)}

\gambhira{ธมฺมานุโลม ทุกติกปฏฺฐานปาฬิ}
\prose{\csromanfolio{398}จิตฺตํ อกุสลญฺจ โนจิตฺตํ อกุสลญฺจ ธมฺมํ ปฏิจฺจ โนจิตฺโต อกุสโล ธมฺโม อุปฺปชฺชติ เหตุปจฺจยา. (1) (สํขิตฺตํ.)}

\proseitem{14}{เหตุยา ปญฺจ, อารมฺมเณ ปญฺจ ฯเปฯ อวิคเต ปญฺจ. (สํขิตฺตํ.\csromanspacer{} จิตฺตทุกกุสลสทิสํ, นเหตุยาปิ กตฺตพฺพํ.\csromanspacer{} สหชาตวารมฺปิ ฯเปฯ ปญฺหาวารมฺปิ กาตพฺพํ.)}

\proseitem{15}{จิตฺตํ อพฺยากตํ ธมฺมํ ปฏิจฺจ โนจิตฺโต อพฺยากโต ธมฺโม อุปฺปชฺชติ เหตุปจฺจยา. (1)}

\prose{โนจิตฺตํ อพฺยากตํ ธมฺมํ ปฏิจฺจ โนจิตฺโต อพฺยากโต ธมฺโม อุปฺปชฺชติ เหตุปจฺจยา.\csromanspacer{}\csromanspacer{}โนจิตฺตํ อพฺยากตํ ธมฺมํ ปฏิจฺจ จิตฺโต อพฺยากโต ธมฺโม อุปฺปชฺชติ เหตุปจฺจยา.\csromanspacer{}\csromanspacer{}โนจิตฺตํ อพฺยากตํ ธมฺมํ ปฏิจฺจ จิตฺโต อพฺยากโต จ โนจิตฺโต อพฺยากโต จ ธมฺมา อุปฺปชฺชนฺติ เหตุปจฺจยา. (3)}

\prose{จิตฺตํ อพฺยากตญฺจ โนจิตฺตํ อพฺยากตญฺจ ธมฺมํ ปฏิจฺจ โนจิตฺโต อพฺยากโต ธมฺโม อุปฺปชฺชติ เหตุปจฺจยา. (1) (สํขิตฺตํ.)}

\proseitem{16}{เหตุยา ปญฺจ, อารมฺมเณ ปญฺจ ฯเปฯ วิปาเก ปญฺจ ฯเปฯ อวิคเต ปญฺจ. (สํขิตฺตํ.)}

\prose{นเหตุยา ปญฺจ, น-อารมฺมเณ ตีณิ, น-อธิปติยา ปญฺจ ฯเปฯ นปุเรชาเต นปจฺฉาชาเต น-อาเสวเน ปญฺจ, นกมฺเม ตีณิ, นวิปาเก ปญฺจ, น-อาหาเร เอกํ, น-อินฺทฺริเย เอกํ, นฌาเน ปญฺจ, นมคฺเค ปญฺจ, นสมฺปยุตฺเต ตีณิ, นวิปฺปยุตฺเต ปญฺจ ฯเปฯ โนวิคเต ตีณิ. (สํขิตฺตํ.)}

\prose{(สหชาตวารมฺปิ ฯเปฯ สมฺปยุตฺตวารมฺปิ วิตฺถาเรตพฺพํ.)}

\proseitem{17}{โนจิตฺโต อพฺยากโต ธมฺโม โนจิตฺตสฺส อพฺยากตสฺส ธมฺมสฺส เหตุปจฺจเยน ปจฺจโย.\csromanspacer{} ตีณิ. (จิตฺตทุกกุสลสทิสํ วิตฺถาเรตพฺพํ.)}

\proseitem{18}{เหตุยา ตีณิ, อารมฺมเณ นว, อธิปติยา นว, อนนฺตเร นว ฯเปฯ สหชาเต ปญฺจ, ปุเรชาเต ตีณิ, อาเสวเน นว, กมฺเม}

\vspace{\tipitakaparskip}
\csromancenter{เจตสิกทุก 1.\csromanspacer{} กุสลตฺติก ตีณิ, วิปาเก ปญฺจ, วิปฺปยุตฺเต ปญฺจ, อตฺถิยา ปญฺจ, นตฺถิยา นว ฯเปฯ อวิคเต ปญฺจ. (สํขิตฺตํ.)}
\csromanheader{2}{57. เจตสิกทุก 1. กุสลตฺติก}
\csromanheader{2}{1-7. ปฏิจฺจาทิวาร ปจฺจยจตุกฺก-เหตุ}
\proseitem{19}{\csromanfolio{399}เจตสิกํ กุสลํ ธมฺมํ ปฏิจฺจ เจตสิโก กุสโล ธมฺโม อุปฺปชฺชติ เหตุปจฺจยา.\csromanspacer{}\csromanspacer{}เจตสิกํ กุสลํ ธมฺมํ ปฏิจฺจ อเจตสิโก กุสโล ธมฺโม อุปฺปชฺชติ เหตุปจฺจยา.\csromanspacer{}\csromanspacer{}เจตสิกํ กุสลํ ธมฺมํ ปฏิจฺจ เจตสิโก กุสโล จ อเจตสิโก กุสโล จ ธมฺมา อุปฺปชฺชนฺติ เหตุปจฺจยา. (3)}

\prose{อเจตสิกํ กุสลํ ธมฺมํ ปฏิจฺจ เจตสิโก กุสโล ธมฺโม อุปฺปชฺชติ เหตุปจฺจยา. (1)}

\prose{เจตสิกํ กุสลญฺจ อเจตสิกํ กุสลญฺจ ธมฺมํ ปฏิจฺจ เจตสิโก กุสโล ธมฺโม อุปฺปชฺชติ เหตุปจฺจยา. (1) (สํขิตฺตํ.)}

\proseitem{20}{เหตุยา ปญฺจ, อารมฺมเณ ปญฺจ, อธิปติยา ปญฺจ ฯเปฯ อวิคเต ปญฺจ.}

\prose{น-อธิปติยา ปญฺจ, นปุเรชาเต ปญฺจ ฯเปฯ นกมฺเม ตีณิ, นวิปาเก ปญฺจ ฯเปฯ นวิปฺปยุตฺเต ปญฺจ. (สํขิตฺตํ.\csromanspacer{} สหชาตวาราทิ วิตฺถาเรตพฺโพ.)}

\proseitem{21}{เจตสิโก กุสโล ธมฺโม เจตสิกสฺส กุสลสฺส ธมฺมสฺส เหตุปจฺจเยน ปจฺจโย. (สํขิตฺตํ.)}

\proseitem{22}{เหตุยา ตีณิ, อารมฺมเณ นว, อธิปติยา นว, (อาทิมฺหิ ตีณิ, สหชาตาธิปติ, มชฺฌิเมสุ ตีสุ มชฺฌิมานุโลมิกาเยว ปญฺหา, สหชาตาธิปติ,) อนนฺตเร นว ฯเปฯ สหชาเต ปญฺจ ฯเปฯ อุปนิสฺสเย นว, อาเสวเน นว, กมฺเม ตีณิ, อาหาเร ปญฺจ, อินฺทฺริเย ปญฺจ, ฌาเน ตีณิ, มคฺเค ตีณิ, สมฺปยุตฺเต ปญฺจ, อตฺถิยา ปญฺจ, นตฺถิยา นว ฯเปฯ อวิคเต ปญฺจ. (สํขิตฺตํ.)}

\gambhira{ธมฺมานุโลม ทุกติกปฏฺฐานปาฬิ}
\proseitem{23}{\csromanfolio{400}เจตสิกํ อกุสลํ ธมฺมํ ปฏิจฺจ เจตสิโก อกุสโล ธมฺโม อุปฺปชฺชติ เหตุปจฺจยา.\csromanspacer{}\csromanspacer{}เจตสิกํ อกุสลํ ธมฺมํ ปฏิจฺจ อเจตสิโก อกุสโล ธมฺโม อุปฺปชฺชติ เหตุปจฺจยา.\csromanspacer{}\csromanspacer{}เจตสิกํ อกุสลํ ธมฺมํ ปฏิจฺจ เจตสิโก อกุสโล จ อเจตสิโก อกุสโล จ ธมฺมา อุปฺปชฺชนฺติ เหตุปจฺจยา. (3)}

\prose{อเจตสิกํ อกุสลํ ธมฺมํ ปฏิจฺจ เจตสิโก อกุสโล ธมฺโม อุปฺปชฺชติ เหตุปจฺจยา. (1)}

\prose{เจตสิกํ อกุสลญฺจ อเจตสิกํ อกุสลญฺจ ธมฺมํ ปฏิจฺจ เจตสิโก อกุสโล ธมฺโม อุปฺปชฺชติ เหตุปจฺจยา. (1) (สํขิตฺตํ.)}

\proseitem{24}{เหตุยา ปญฺจ, อารมฺมเณ ปญฺจ ฯเปฯ อวิคเต ปญฺจ. (สํขิตฺตํ.) (ยถา กุสลนยํ เอวํ นเหตุปจฺจยมฺปิ กาตพฺพํ, สหชาตวารมฺปิ ฯเปฯ ปญฺหาวารมฺปิ วิตฺถาเรตพฺพํ.)}

\proseitem{25}{เจตสิกํ อพฺยากตํ ธมฺมํ ปฏิจฺจ เจตสิโก อพฺยากโต ธมฺโม อุปฺปชฺชติ เหตุปจฺจยา.\csromanspacer{}\csromanspacer{}เจตสิกํ อพฺยากตํ ธมฺมํ ปฏิจฺจ อเจตสิโก อพฺยากโต ธมฺโม อุปฺปชฺชติ เหตุปจฺจยา.\csromanspacer{}\csromanspacer{}เจตสิกํ อพฺยากตํ ธมฺมํ ปฏิจฺจ เจตสิโก อพฺยากโต จ อเจตสิโก อพฺยากโต จ ธมฺมา อุปฺปชฺชนฺติ เหตุปจฺจยา. (3)}

\prose{อเจตสิกํ อพฺยากตํ ธมฺมํ ปฏิจฺจ อเจตสิโก อพฺยากโต ธมฺโม อุปฺปชฺชติ เหตุปจฺจยา.\csromanspacer{} ตีณิ.}

\prose{เจตสิกํ อพฺยากตญฺจ อเจตสิกํ อพฺยากตญฺจ ธมฺมํ ปฏิจฺจ เจตสิโก อพฺยากโต ธมฺโม อุปฺปชฺชติ เหตุปจฺจยา.\csromanspacer{} ตีณิ. (สํขิตฺตํ.)}

\proseitem{26}{เหตุยา นว, อารมฺมเณ นว, อธิปติยา นว ฯเปฯ ปุเรชาเต อาเสวเน ปญฺจ, กมฺเม นว, วิปาเก นว ฯเปฯ อวิคเต นว. (สํขิตฺตํ.)}

\prose{นเหตุยา นว, น-อารมฺมเณ ตีณิ, น-อธิปติยา นว, (สพฺเพ กตพฺพา) ฯเปฯ นกมฺเม จตฺตาริ, นวิปาเก นว, น-อาหาเร เอกํ, น-อินฺทฺริเย เอกํ, นฌาเน ฉ, นมคฺเค นว, นสมฺปยุตฺเต ตีณิ, นวิปฺปยุตฺเต ฉ ฯเปฯ โนวิคเต ตีณิ. (สํขิตฺตํ.\csromanspacer{} สหชาตวาราทิ วิตฺถาเรตพฺโพ.)}

\csromanheader{2}{58. จิตฺตสมฺปยุตฺตทุก 1. กุสลตฺติก}
\proseitem{27}{\csromanfolio{401}เจตสิโก อพฺยากโต ธมฺโม เจตสิกสฺส อพฺยากตสฺส ธมฺมสฺส เหตุปจฺจเยน ปจฺจโย. (สํขิตฺตํ.)}

\proseitem{28}{เหตุยา ตีณิ, อารมฺมเณ นว, อธิปติยา นว, (เอตฺถ ฉสุ สหชาตาธิปติ.) อนนฺตเร นว ฯเปฯ สหชาเต นว, อญฺญมญฺเญ นว, นิสฺสเย นว, อุปนิสฺสเย นว, ปุเรชาเต ตีณิ, อาเสวเน นว, กมฺเม ตีณิ, วิปาเก นว, อาหาเร นว, อินฺทฺริเย นว, ฌาเน ตีณิ, มคฺเค ตีณิ, สมฺปยุตฺเต ปญฺจ, วิปฺปยุตฺเต ปญฺจ, อตฺถิยา นว ฯเปฯ อวิคเต นว. (สํขิตฺตํ.)}

\csromanheader{2}{58. จิตฺตสมฺปยุตฺตทุก 1. กุสลตฺติก}
\csromanheader{2}{1-7. ปฏิจฺจาทิวาร ปจฺจยจตุกฺก-เหตุ}
\proseitem{29}{จิตฺตสมฺปยุตฺตํ กุสลํ ธมฺมํ ปฏิจฺจ จิตฺตสมฺปยุตฺโต กุสโล ธมฺโม อุปฺปชฺชติ เหตุปจฺจยา. (สํขิตฺตํ.)}

\proseitem{30}{เหตุยา เอกํ, อารมฺมเณ เอกํ ฯเปฯ อวิคเต เอกํ. (สํขิตฺตํ.) (สหชาตวาเรปิ ฯเปฯ ปญฺหาวาเรปิ สพฺพตฺถ เอกํ.)}

\proseitem{31}{จิตฺตสมฺปยุตฺตํ อกุสลํ ธมฺมํ ปฏิจฺจ จิตฺตสมฺปยุตฺโต อกุสโล ธมฺโม อุปฺปชฺชติ เหตุปจฺจยา. (สํขิตฺตํ.)}

\proseitem{32}{เหตุยา เอกํ, อารมฺมเณ เอกํ ฯเปฯ อวิคเต เอกํ. (สํขิตฺตํ.) (สหชาตวาเรปิ ฯเปฯ ปญฺหาวาเรปิ สพฺพตฺถ เอกํ.)}

\csromanheader{2}{\textbf{อพฺยากตปท-เหตุ}}
\proseitem{33}{จิตฺตสมฺปยุตฺตํ อพฺยากตํ ธมฺมํ ปฏิจฺจ จิตฺตสมฺปยุตฺโต อพฺยากโต ธมฺโม อุปฺปชฺชติ เหตุปจฺจยา.\csromanspacer{}\csromanspacer{}จิตฺตสมฺปยุตฺตํ อพฺยากตํ ธมฺมํ ปฏิจฺจ จิตฺตวิปฺปยุตฺโต อพฺยากโต ธมฺโม อุปฺปชฺชติ เหตุปจฺจยา.\csromanspacer{}\csromanspacer{}จิตฺตสมฺปยุตฺตํ อพฺยากตํ ธมฺมํ ปฏิจฺจ จิตฺตสมฺปยุตฺโต อพฺยากโต จ จิตฺตวิปฺปยุตฺโต อพฺยากโต จ ธมฺมา อุปฺปชฺชนฺติ เหตุปจฺจยา. (3)}

\prose{จิตฺตวิปฺปยุตฺตํ อพฺยากตํ ธมฺมํ ปฏิจฺจ จิตฺตวิปฺปยุตฺโต อพฺยากโต ธมฺโม อุปฺปชฺชติ เหตุปจฺจยา.\csromanspacer{}\csromanspacer{}จิตฺตวิปฺปยุตฺตํ อพฺยากตํ ธมฺมํ ปฏิจฺจ จิตฺตธมฺมานุโลม \csromanfolio{402}ทุกติกปฏฺฐานปาฬิ สมฺปยุตฺโต อพฺยากโต ธมฺโม อุปฺปชฺชติ เหตุปจฺจยา.\csromanspacer{}\csromanspacer{}จิตฺตวิปฺปยุตฺตํ อพฺยากตํ ธมฺมํ ปฏิจฺจ จิตฺตสมฺปยุตฺโต อพฺยากโต จ จิตฺตวิปฺปยุตฺโต อพฺยากโต จ ธมฺมา อุปฺปชฺชนฺติ เหตุปจฺจยา. (3)}

\prose{จิตฺตสมฺปยุตฺตํ อพฺยากตญฺจ จิตฺตวิปฺปยุตฺตํ อพฺยากตญฺจ ธมฺมํ ปฏิจฺจ จิตฺตสมฺปยุตฺโต อพฺยากโต ธมฺโม อุปฺปชฺชติ เหตุปจฺจยา.\csromanspacer{}\csromanspacer{}จิตฺตสมฺปยุตฺตํ อพฺยากตญฺจ จิตฺตวิปฺปยุตฺตํ อพฺยากตญฺจ ธมฺมํ ปฏิจฺจ จิตฺตวิปฺปยุตฺโต อพฺยากโต ธมฺโม อุปฺปชฺชติ เหตุปจฺจยา.\csromanspacer{}\csromanspacer{}จิตฺตสมฺปยุตฺตํ อพฺยากตญฺจ จิตฺตวิปฺปยุตฺตํ อพฺยากตญฺจ ธมฺมํ ปฏิจฺจ จิตฺตสมฺปยุตฺโต อพฺยากโต จ จิตฺตวิปฺปยุตฺโต อพฺยากโต จ ธมฺมา อุปฺปชฺชนฺติ เหตุปจฺจยา. (3) (สํขิตฺตํ.)}

\proseitem{34}{เหตุยา นว, อารมฺมเณ ตีณิ, อธิปติยา ปญฺจ ฯเปฯ อญฺญมญฺเญ ฉ ฯเปฯ ปุเรชาเต เอกํ, อาเสวเน เอกํ, กมฺเม นว, วิปาเก นว ฯเปฯ อวิคเต นว. (สํขิตฺตํ.)}

\prose{นเหตุยา นว, น-อารมฺมเณ ตีณิ, น-อธิปติยา นว ฯเปฯ นปุเรชาเต นว ฯเปฯ นกมฺเม เทฺว, นวิปาเก ปญฺจ, น-อาหาเร เอกํ, น-อินฺทฺริเย เอกํ, นฌาเน เทฺว, นมคฺเค นว, นสมฺปยุตฺเต ตีณิ, นวิปฺปยุตฺเต เทฺว. (สํขิตฺตํ.\csromanspacer{} สหชาตวาราทิ วิตฺถาเรตพฺโพ.)}

\csromanheader{2}{\textbf{เหตุ-อารมฺมณ}}
\proseitem{35}{จิตฺตสมฺปยุตฺโต อพฺยากโต ธมฺโม จิตฺตสมฺปยุตฺตสฺส อพฺยากตสฺส ธมฺมสฺส เหตุปจฺจเยน ปจฺจโย.\csromanspacer{} ตีณิ.}

\prose{จิตฺตสมฺปยุตฺโต อพฺยากโต ธมฺโม จิตฺตสมฺปยุตฺตสฺส อพฺยากตสฺส ธมฺมสฺส อารมฺมณปจฺจเยน ปจฺจโย. (1)}

\prose{จิตฺตวิปฺปยุตฺโต อพฺยากโต ธมฺโม จิตฺตสมฺปยุตฺตสฺส อพฺยากตสฺส ธมฺมสฺส อารมฺมณปจฺจเยน ปจฺจโย. (1) (สํขิตฺตํ.)}

\proseitem{36}{เหตุยา ตีณิ, อารมฺมเณ เทฺว, อธิปติยา จตฺตาริ, อนนฺตเร เอกํ ฯเปฯ สหชาเต สตฺต, อญฺญมญฺเญ ฉ, นิสฺสเย สตฺต, อุปนิสฺสเย เทฺว, ปุเรชาเต เอกํ, ปจฺฉาชาเต เอกํ, อาเสวเน เอกํ, กมฺเม ตีณิ, วิปาเก ตีณิ, อาหาเร จตฺตาริ, อินฺทฺริเย ฉ, ฌาเน ตีณิ, มคฺเค ตีณิ, สมฺปยุตฺเต เอกํ, วิปฺปยุตฺเต เทฺว. (สํขิตฺตํ.)}

\vspace{\tipitakaparskip}
\csromancenter{จิตฺตสมฺปยุตฺตทุก 1.\csromanspacer{} กุสลตฺติก \textbf{59.}\csromanspacer{}\textbf{ จิตฺตสํสฏฺฐทุก 1.}\csromanspacer{}\textbf{ กุสลตฺติก}}
\csromanheader{2}{1. ปฏิจฺจวาร ปจฺจยจตุกฺก-เหตุ}
\proseitem{37}{\csromanfolio{403}จิตฺตสํสฏฺฐํ กุสลํ ธมฺมํ ปฏิจฺจ จิตฺตสํสฏฺโฐ กุสโล ธมฺโม อุปฺปชฺชติ เหตุปจฺจยา. (สํขิตฺตํ.)}

\proseitem{38}{เหตุยา เอกํ, อารมฺมเณ เอกํ ฯเปฯ อวิคเต เอกํ. (สํขิตฺตํ.) (สหชาตวาเรปิ ฯเปฯ ปญฺหาวาเรปิ สพฺพตฺถ เอกํ.)}

\proseitem{39}{จิตฺตสํสฏฺฐํ อกุสลํ ธมฺมํ ปฏิจฺจ จิตฺตสํสฏฺโฐ อกุสโล ธมฺโม อุปฺปชฺชติ เหตุปจฺจยา. (สํขิตฺตํ.)}

\proseitem{40}{เหตุยา เอกํ, อารมฺมเณ เอกํ ฯเปฯ อวิคเต เอกํ. (สํขิตฺตํ.) (สหชาตวาเรปิ ฯเปฯ ปญฺหาวาเรปิ สพฺพตฺถ เอกํ.)}

\proseitem{41}{จิตฺตสํสฏฺฐํ อพฺยากตํ ธมฺมํ ปฏิจฺจ จิตฺตสํสฏฺโฐ อพฺยากโต ธมฺโม อุปฺปชฺชติ เหตุปจฺจยา.\csromanspacer{} ตีณิ.}

\prose{โนจิตฺตสํสฏฺฐํ อพฺยากตํ ธมฺมํ ปฏิจฺจ โนจิตฺตสํสฏฺโฐ อพฺยากโต ธมฺโม อุปฺปชฺชติ เหตุปจฺจยา.\csromanspacer{} ตีณิ.}

\prose{จิตฺตสํสฏฺฐํ อพฺยากตญฺจ โนจิตฺตสํสฏฺฐํ อพฺยากตญฺจ ธมฺมํ ปฏิจฺจ จิตฺตสํสฏฺโฐ อพฺยากโต ธมฺโม อุปฺปชฺชติ เหตุปจฺจยา.\csromanspacer{} ตีณิ. (สํขิตฺตํ.)}

\proseitem{42}{เหตุยา นว, อารมฺมเณ ตีณิ ฯเปฯ อวิคเต นว. (สํขิตฺตํ.)}

\prose{(ยถา จิตฺตสมฺปยุตฺตทุกํ อพฺยากตสทิสํ.\csromanspacer{} สหชาตวารมฺปิ ฯเปฯ ปญฺหาวารมฺปิ วิตฺถาเรตพฺพํ.)}

\csromanheader{2}{60. จิตฺตสมุฏฺฐานทุก 1. กุสลตฺติก}
\csromanheader{2}{1-7. ปฏิจฺจาทิวาร ปจฺจยจตุกฺก-เหตุ}
\proseitem{43}{จิตฺตสมุฏฺฐานํ กุสลํ ธมฺมํ ปฏิจฺจ จิตฺตสมุฏฺฐาโน กุสโล ธมฺโม อุปฺปชฺชติ เหตุปจฺจยา.\csromanspacer{}\csromanspacer{}จิตฺตสมุฏฺฐานํ กุสลํ ธมฺมํ ปฏิจฺจ โนจิตฺตสมุฏฺฐาโน กุสโล ธมฺโม อุปฺปชฺชติ เหตุปจฺจยา.\csromanspacer{}\csromanspacer{}จิตฺตสมุฏฺฐานํ กุสลํ ธมฺมํ ปฏิจฺจ}

\vspace{\tipitakaparskip}
\csromancenter{ธมฺมานุโลม ทุกติกปฏฺฐานปาฬิ จิตฺตสมุฏฺฐาโน กุสโล จ โนจิตฺตสมุฏฺฐาโน กุสโล จ ธมฺมา อุปฺปชฺชนฺติ เหตุปจฺจยา. (3)}
\prose{\csromanfolio{404}โนจิตฺตสมุฏฺฐานํ กุสลํ ธมฺมํ ปฏิจฺจ จิตฺตสมุฏฺฐาโน กุสโล ธมฺโม อุปฺปชฺชติ เหตุปจฺจยา. (1)}

\prose{จิตฺตสมุฏฺฐานํ กุสลญฺจ โนจิตฺตสมุฏฺฐานํ กุสลญฺจ ธมฺมํ ปฏิจฺจ จิตฺตสมุฏฺฐาโน กุสโล ธมฺโม อุปฺปชฺชติ เหตุปจฺจยา. (1) (สํขิตฺตํ.)}

\proseitem{44}{เหตุยา ปญฺจ, อารมฺมเณ ปญฺจ, (สพฺพตฺถ ปญฺจ.) อวิคเต ปญฺจ. (สํขิตฺตํ.)}

\prose{น-อธิปติยา ปญฺจ, นปุเรชาเต ปญฺจ ฯเปฯ นกมฺเม ตีณิ, นวิปาเก ปญฺจ ฯเปฯ นวิปฺปยุตฺเต ปญฺจ. (สํขิตฺตํ.\csromanspacer{} สหชาตวาราทิ วิตฺถาเรตพฺโพ.)}

\proseitem{45}{จิตฺตสมุฏฺฐาโน กุสโล ธมฺโม จิตฺตสมุฏฺฐานสฺส กุสลสฺส ธมฺมสฺส เหตุปจฺจเยน ปจฺจโย. (สํขิตฺตํ.)}

\proseitem{46}{เหตุยา ตีณิ, อารมฺมเณ นว, อธิปติยา นว, อนนฺตเร นว, สมนนฺตเร นว, สหชาเต ปญฺจ ฯเปฯ อุปนิสฺสเย นว, อาเสวเน นว, กมฺเม ตีณิ, อาหาเร ปญฺจ, อินฺทฺริเย ปญฺจ, ฌาเน ตีณิ, มคฺเค ตีณิ, สมฺปยุตฺเต ปญฺจ, อตฺถิยา ปญฺจ ฯเปฯ อวิคเต ปญฺจ. (สํขิตฺตํ.)}

\proseitem{47}{จิตฺตสมุฏฺฐานํ อกุสลํ ธมฺมํ ปฏิจฺจ จิตฺตสมุฏฺฐาโน อกุสโล ธมฺโม อุปฺปชฺชติ เหตุปจฺจยา.\csromanspacer{} ตีณิ.}

\prose{โนจิตฺตสมุฏฺฐานํ อกุสลํ ธมฺมํ ปฏิจฺจ จิตฺตสมุฏฺฐาโน อกุสโล ธมฺโม อุปฺปชฺชติ เหตุปจฺจยา. (1)}

\prose{จิตฺตสมุฏฺฐานํ อกุสลญฺจ โนจิตฺตสมุฏฺฐานํ อกุสลญฺจ ธมฺมํ ปฏิจฺจ จิตฺตสมุฏฺฐาโน อกุสโล ธมฺโม อุปฺปชฺชติ เหตุปจฺจยา. (1) (สํขิตฺตํ.)}

\proseitem{48}{เหตุยา ปญฺจ, อารมฺมเณ ปญฺจ ฯเปฯ อวิคเต ปญฺจ. (สํขิตฺตํ.)}

\prose{(กุสลตฺติกสทิสํ สหชาตวารมฺปิ ฯเปฯ ปญฺหาวารมฺปิ วิตฺถาเรตพฺพํ.)}

\proseitem{49}{จิตฺตสมุฏฺฐานํ อพฺยากตํ ธมฺมํ ปฏิจฺจ จิตฺตสมุฏฺฐาโน อพฺยากโต ธมฺโม อุปฺปชฺชติ เหตุปจฺจยา.\csromanspacer{} ตีณิ.}

\prose{โนจิตฺตสมุฏฺฐานํ อพฺยากตํ ธมฺมํ ปฏิจฺจ โนจิตฺตสมุฏฺฐาโน อพฺยากโต ธมฺโม อุปฺปชฺชติ เหตุปจฺจยา.\csromanspacer{} ตีณิ.}

\csromanheader{2}{61. จิตฺตสหภูทุก 1. กุสลตฺติก}
\prose{\csromanfolio{405}จิตฺตสมุฏฺฐานํ อพฺยากตญฺจ โนจิตฺตสมุฏฺฐานํ อพฺยากตญฺจ ธมฺมํ ปฏิจฺจ จิตฺตสมุฏฺฐาโน อพฺยากโต ธมฺโม อุปฺปชฺชติ เหตุปจฺจยา.\csromanspacer{} ตีณิ. (สํขิตฺตํ.)}

\proseitem{50}{เหตุยา นว, อารมฺมเณ นว, อธิปติยา ปญฺจ, (สพฺพตฺถ นว.) ปุเรชาเต ปญฺจ, อาเสวเน ปญฺจ ฯเปฯ อวิคเต นว. (สํขิตฺตํ.)}

\prose{นเหตุยา นว, น-อารมฺมเณ ฉ, น-อธิปติยา นว ฯเปฯ นกมฺเม ตีณิ, นวิปาเก ฉ, น-อาหาเร เอกํ, น-อินฺทฺริเย เอกํ, นฌาเน ฉ, นมคฺเค นว, นสมฺปยุตฺเต ฉ, นวิปฺปยุตฺเต ฉ, โนนตฺถิยา ฉ, โนวิคเต ฉ. (สํขิตฺตํ.\csromanspacer{} สหชาตวารมฺปิ ฯเปฯ สมฺปยุตฺตวารมฺปิ วิตฺถาเรตพฺพํ.)}

\proseitem{51}{จิตฺตสมุฏฺฐาโน อพฺยากโต ธมฺโม จิตฺตสมุฏฺฐานสฺส อพฺยากตสฺส ธมฺมสฺส เหตุปจฺจเยน ปจฺจโย.\csromanspacer{} ตีณิ.}

\prose{จิตฺตสมุฏฺฐาโน อพฺยากโต ธมฺโม จิตฺตสมุฏฺฐานสฺส อพฺยากตสฺส ธมฺมสฺส อารมฺมณปจฺจเยน ปจฺจโย. (สํขิตฺตํ.)}

\proseitem{52}{เหตุยา ตีณิ, อารมฺมเณ นว, อธิปติยา นว, อนนฺตเร นว, (สพฺพตฺถ นว.) ปุเรชาเต นว, ปจฺฉาชาเต นว, อาเสวเน นว, กมฺเม ตีณิ, วิปาเก นว, อาหาเร นว, (จิตฺตสมุฏฺฐานมูลํ โนจิตฺตสมุฏฺฐานสฺส กพฬีกาโร อาหาโร กาตพฺโพ.\csromanspacer{} โนจิตฺตสมุฏฺฐาโน จิตฺตสมุฏฺฐานสฺส กพฬีกาโร อาหาโร ฆฏเน มชฺเฌ กพฬีกาโร อาหาโร.) อินฺทฺริเย นว, (รูปชีวตินฺทฺริยํ เอกํ) ฌาเน ตีณิ, มคฺเค ตีณิ, สมฺปยุตฺเต ปญฺจ, วิปฺปยุตฺเต นว, อตฺถิยา นว, นตฺถิยา นว, วิคเต นว, อวิคเต นว. (สํขิตฺตํ.)}

\prose{(ยถา กุสลตฺติเก ปญฺหาวารสฺส อนุโลมมฺปิ ปจฺจนียมฺปิ อนุโลมปจฺจนียมฺปิ ปจฺจนียานุโลมมฺปิ คณิตํ, เอวํ คเณตพฺพํ.)}

\csromanheader{2}{61. จิตฺตสหภูทุก 1. กุสลตฺติก}
\csromanheader{2}{1. ปฏิจฺจวาร ปจฺจยจตุกฺก-เหตุ}
\proseitem{53}{จิตฺตสหภุํ กุสลํ ธมฺมํ ปฏิจฺจ จิตฺตสหภู กุสโล ธมฺโม อุปฺปชฺชติ เหตุปจฺจยา.\csromanspacer{}\csromanspacer{}จิตฺตสหภุํ กุสลํ ธมฺมํ ปฏิจฺจ โนจิตฺตสหภู กุสโล}

\vspace{\tipitakaparskip}
\csromancenter{ธมฺมานุโลม ทุกติกปฏฺฐานปาฬิ ธมฺโม อุปฺปชฺชติ เหตุปจฺจยา.\csromanspacer{}\csromanspacer{}จิตฺตสหภุํ กุสลํ ธมฺมํ ปฏิจฺจ จิตฺตสหภู กุสโล จ โนจิตฺตสหภู กุสโล จ ธมฺมา อุปฺปชฺชนฺติ เหตุปจฺจยา. (3)}
\prose{\csromanfolio{406}โนจิตฺตสหภุํ กุสลํ ธมฺมํ ปฏิจฺจ จิตฺตสหภู กุสโล ธมฺโม อุปฺปชฺชติ เหตุปจฺจยา. (1)}

\prose{จิตฺตสหภุํ กุสลญฺจ โนจิตฺตสหภุํ กุสลญฺจ ธมฺมํ ปฏิจฺจ จิตฺตสหภู กุสโล ธมฺโม อุปฺปชฺชติ เหตุปจฺจยา. (1) (สํขิตฺตํ.)}

\proseitem{54}{เหตุยา ปญฺจ, อารมฺมเณ ปญฺจ ฯเปฯ อวิคเต ปญฺจ. (สํขิตฺตํ.) (สหชาตวารมฺปิ ฯเปฯ ปญฺหาวารมฺปิ วิตฺถาเรตพฺพํ.)}

\proseitem{55}{จิตฺตสหภุํ อกุสลํ ธมฺมํ ปฏิจฺจ จิตฺตสหภู อกุสโล ธมฺโม อุปฺปชฺชติ เหตุปจฺจยา.\csromanspacer{} ตีณิ.}

\prose{โนจิตฺตสหภุํ อกุสลํ ธมฺมํ ปฏิจฺจ จิตฺตสหภู อกุสโล ธมฺโม อุปฺปชฺชติ เหตุปจฺจยา. (1)}

\prose{จิตฺตสหภุํ อกุสลญฺจ โนจิตฺตสหภุํ อกุสลญฺจ ธมฺมํ ปฏิจฺจ จิตฺตสหภู อกุสโล ธมฺโม อุปฺปชฺชติ เหตุปจฺจยา. (1) (สํขิตฺตํ.)}

\proseitem{56}{เหตุยา ปญฺจ, อารมฺมเณ ปญฺจ ฯเปฯ อวิคเต ปญฺจ. (สํขิตฺตํ.) (สหชาตวารมฺปิ ฯเปฯ ปญฺหาวารมฺปิ วิตฺถาเรตพฺพํ.)}

\proseitem{57}{จิตฺตสหภุํ อพฺยากตํ ธมฺมํ ปฏิจฺจ จิตฺตสหภู อพฺยากโต ธมฺโม อุปฺปชฺชติ เหตุปจฺจยา.\csromanspacer{} ตีณิ.}

\prose{โนจิตฺตสหภุํ อพฺยากตํ ธมฺมํ ปฏิจฺจ โนจิตฺตสหภู อพฺยากโต ธมฺโม อุปฺปชฺชติ เหตุปจฺจยา.\csromanspacer{} ตีณิ.}

\prose{จิตฺตสหภุํ อพฺยากตญฺจ โนจิตฺตสหภุํ อพฺยากตญฺจ ธมฺมํ ปฏิจฺจ จิตฺตสหภู อพฺยากโต ธมฺโม อุปฺปชฺชติ เหตุปจฺจยา.\csromanspacer{} ตีณิ. (สํขิตฺตํ.)}

\proseitem{58}{เหตุยา นว, อารมฺมเณ นว ฯเปฯ อวิคเต นว. (สํขิตฺตํ.\csromanspacer{} ยถา เจตสิกทุกมูลา ตีณิ คมนา, เอวํ อิเมปิ ตีณิ คมนา กาตพฺพา.\csromanspacer{} สหชาตวารมฺปิ ฯเปฯ ปญฺหาวารมฺปิ วิตฺถาเรตพฺพํ.)}

\vspace{\tipitakaparskip}
\csromancenter{จิตฺตานุปริวตฺติทุก 1.\csromanspacer{} กุสลตฺติก \textbf{62.}\csromanspacer{}\textbf{ จิตฺตานุปริวตฺติทุก 1.}\csromanspacer{}\textbf{ กุสลตฺติก}}
\csromanheader{2}{1. ปฏิจฺจวาร ปจฺจยจตุกฺก-เหตุ}
\proseitem{59}{\csromanfolio{407}จิตฺตานุปริวตฺติํ กุสลํ ธมฺมํ ปฏิจฺจ จิตฺตานุปริวตฺตี กุสโล ธมฺโม อุปฺปชฺชติ เหตุปจฺจยา.\csromanspacer{} ตีณิ.}

\prose{โนจิตฺตานุปริวตฺติํ กุสลํ ธมฺมํ ปฏิจฺจ จิตฺตานุปริวตฺตี กุสโล ธมฺโม อุปฺปชฺชติ เหตุปจฺจยา. (1)}

\prose{จิตฺตานุปริวตฺติํ กุสลญฺจ โนจิตฺตานุปริวตฺติํ กุสลญฺจ ธมฺมํ ปฏิจฺจ จิตฺตานุปริวตฺตี กุสโล ธมฺโม อุปฺปชฺชติ เหตุปจฺจยา. (1) (สํขิตฺตํ.)}

\proseitem{60}{เหตุยา ปญฺจ, อารมฺมเณ ปญฺจ ฯเปฯ อวิคเต ปญฺจ. (สํขิตฺตํ.) (สหชาตวารมฺปิ ฯเปฯ ปญฺหาวารมฺปิ วิตฺถาเรตพฺพํ.)}

\proseitem{61}{จิตฺตานุปริวตฺติํ อกุสลํ ธมฺมํ ปฏิจฺจ จิตฺตานุปริวตฺตี อกุสโล ธมฺโม อุปฺปชฺชติ เหตุปจฺจยา.\csromanspacer{} ตีณิ.}

\prose{โนจิตฺตานุปริวตฺติํ อกุสลํ ธมฺมํ ปฏิจฺจ จิตฺตานุปริวตฺตี อกุสโล ธมฺโม อุปฺปชฺชติ เหตุปจฺจยา. (1)}

\prose{จิตฺตานุปริวตฺติํ อกุสลญฺจ โนจิตฺตานุปริวตฺติํ อกุสลญฺจ ธมฺมํ ปฏิจฺจ จิตฺตานุปริวตฺตี อกุสโล ธมฺโม อุปฺปชฺชติ เหตุปจฺจยา. (1) (สํขิตฺตํ.)}

\proseitem{62}{เหตุยา ปญฺจ, อารมฺมเณ ปญฺจ ฯเปฯ อวิคเต ปญฺจ. (สํขิตฺตํ.) (สหชาตวารมฺปิ ฯเปฯ ปญฺหาวารมฺปิ วิตฺถาเรตพฺพํ.)}

\proseitem{63}{จิตฺตานุปริวตฺติํ อพฺยากตํ ธมฺมํ ปฏิจฺจ จิตฺตานุปริวตฺตี อพฺยากโต ธมฺโม อุปฺปชฺชติ เหตุปจฺจยา. (สํขิตฺตํ.)}

\proseitem{64}{เหตุยา นว, อารมฺมเณ นว ฯเปฯ อวิคเต นว. (สํขิตฺตํ.\csromanspacer{} เจตสิกทุกสทิสํ.\csromanspacer{} สหชาตวารมฺปิ ฯเปฯ ปญฺหาวารมฺปิ วิตฺถาเรตพฺพํ.)}

\vspace{\tipitakaparskip}
\csromancenter{ธมฺมานุโลม ทุกติกปฏฺฐานปาฬิ \textbf{63.}\csromanspacer{}\textbf{ จิตฺตสํสฏฺฐสมุฏฺฐานทุก 1.}\csromanspacer{}\textbf{ กุสลตฺติก}}
\csromanheader{2}{1-7. ปฏิจฺจาทิวาร ปจฺจยจตุกฺก-เหตุ}
\proseitem{65}{\csromanfolio{408}จิตฺตสํสฏฺฐสมุฏฺฐานํ กุสลํ ธมฺมํ ปฏิจฺจ จิตฺตสํสฏฺฐสมุฏฺฐาโน กุสโล ธมฺโม อุปฺปชฺชติ เหตุปจฺจยา.\csromanspacer{}\csromanspacer{}จิตฺตสํสฏฺฐสมุฏฺฐานํ กุสลํ ธมฺมํ ปฏิจฺจ โนจิตฺตสํสฏฺฐสมุฏฺฐาโน กุสโล ธมฺโม อุปฺปชฺชติ เหตุปจฺจยา.\csromanspacer{}\csromanspacer{}จิตฺตสํสฏฺฐสมุฏฺฐานํ กุสลํ ธมฺมํ ปฏิจฺจ จิตฺตสํสฏฺฐสมุฏฺฐาโน กุสโล จ โนจิตฺตสํสฏฺฐสมุฏฺฐาโน กุสโล จ ธมฺมา อุปฺปชฺชนฺติ เหตุปจฺจยา. (3)}

\prose{โนจิตฺตสํสฏฺฐสมุฏฺฐานํ กุสลํ ธมฺมํ ปฏิจฺจ จิตฺตสํสฏฺฐสมุฏฺฐาโน กุสโล ธมฺโม อุปฺปชฺชติ เหตุปจฺจยา. (1)}

\prose{จิตฺตสํสฏฺฐสมุฏฺฐานํ กุสลญฺจ โนจิตฺตสํสฏฺฐสมุฏฺฐานํ กุสลญฺจ ธมฺมํ ปฏิจฺจ จิตฺตสํสฏฺฐสมุฏฺฐาโน กุสโล ธมฺโม อุปฺปชฺชติ เหตุปจฺจยา. (1) (สํขิตฺตํ.)}

\proseitem{66}{เหตุยา ปญฺจ, อารมฺมเณ ปญฺจ ฯเปฯ อวิคเต ปญฺจ. (สํขิตฺตํ.\csromanspacer{} ยถา มหนฺตรทุเก เจตสิกทุกกุสลสทิสํ, ตตฺตกา เอว ปญฺหา.\csromanspacer{} สหชาตวารมฺปิ ฯเปฯ ปญฺหาวารมฺปิ วิตฺถาเรตพฺพํ.)}

\proseitem{67}{จิตฺตสํสฏฺฐสมุฏฺฐานํ อกุสลํ ธมฺมํ ปฏิจฺจ จิตฺตสํสฏฺฐสมุฏฺฐาโน อกุสโล ธมฺโม อุปฺปชฺชติ เหตุปจฺจยา.\csromanspacer{} ตีณิ.}

\prose{โนจิตฺตสํสฏฺฐสมุฏฺฐานํ อกุสลํ ธมฺมํ ปฏิจฺจ จิตฺตสํสฏฺฐสมุฏฺฐาโน อกุสโล ธมฺโม อุปฺปชฺชติ เหตุปจฺจยา. (1)}

\prose{จิตฺตสํสฏฺฐสมุฏฺฐานํ อกุสลญฺจ โนจิตฺตสํสฏฺฐสมุฏฺฐานํ อกุสลญฺจ ธมฺมํ ปฏิจฺจ จิตฺตสํสฏฺฐสมุฏฺฐาโน อกุสโล ธมฺโม อุปฺปชฺชติ เหตุปจฺจยา. (1) (สํขิตฺตํ.)}

\proseitem{68}{เหตุยา ปญฺจ, อารมฺมเณ ปญฺจ ฯเปฯ อวิคเต ปญฺจ. (สํขิตฺตํ.\csromanspacer{} เจตสิกทุก-อกุสลสทิสํ.\csromanspacer{} สหชาตวารมฺปิ ฯเปฯ ปญฺหาวารมฺปิ วิตฺถาเรตพฺพํ.)}

\proseitem{69}{จิตฺตสํสฏฺฐสมุฏฺฐานํ อพฺยากตํ ธมฺมํ ปฏิจฺจ จิตฺตสํสฏฺฐสมุฏฺฐาโน อพฺยากโต ธมฺโม อุปฺปชฺชติ เหตุปจฺจยา.\csromanspacer{} ตีณิ.}

\csromanheader{2}{64. จิตฺตสํสฏฺฐสมุฏฺฐานสหภูทุก 1. กุสลตฺติก}
\prose{\csromanfolio{409}โนจิตฺตสํสฏฺฐสมุฏฺฐานํ อพฺยากตํ ธมฺมํ ปฏิจฺจ โนจิตฺตสํสฏฺฐสมุฏฺฐาโน อพฺยากโต ธมฺโม อุปฺปชฺชติ เหตุปจฺจยา.\csromanspacer{} ตีณิ.}

\prose{จิตฺตสํสฏฺฐสมุฏฺฐานํ อพฺยากตญฺจ โนจิตฺตสํสฏฺฐสมุฏฺฐานํ อพฺยากตญฺจ ธมฺมํ ปฏิจฺจ จิตฺตสํสฏฺฐสมุฏฺฐาโน อพฺยากโต ธมฺโม อุปฺปชฺชติ เหตุปจฺจยา.\csromanspacer{} ตีณิ. (สํขิตฺตํ.)}

\proseitem{70}{เหตุยา นว, อารมฺมเณ นว, อธิปติยา นว ฯเปฯ ปุเรชาเต ปญฺจ, อาเสวเน ปญฺจ, กมฺเม นว ฯเปฯ อวิคเต นว. (สํขิตฺตํ.)}

\prose{นเหตุยา นว, น-อารมฺมเณ ตีณิ ฯเปฯ นกมฺเม จตฺตาริ, นวิปาเก นว, น-อาหาเร เอกํ, น-อินฺทฺริเย เอกํ, นฌาเน ฉ, นมคฺเค นว, นสมฺปยุตฺเต ตีณิ, นวิปฺปยุตฺเต ฉ, โนนตฺถิยา ตีณิ, โนวิคเต ตีณิ. (สํขิตฺตํ.\csromanspacer{} สหชาตวารมฺปิ ฯเปฯ สมฺปยุตฺตวารมฺปิ วิตฺถาเรตพฺพํ.)}

\proseitem{71}{จิตฺตสํสฏฺฐสมุฏฺฐาโน อพฺยากโต ธมฺโม จิตฺตสํสฏฺฐสมุฏฺฐานสฺส อพฺยากตสฺส ธมฺมสฺส เหตุปจฺจเยน ปจฺจโย. (สํขิตฺตํ.)}

\proseitem{72}{เหตุยา ตีณิ, อารมฺมเณ นว, อธิปติยา นว, (เอตฺถ ฉสุ สหชาตาธิปติ.) อนนฺตเร นว, (สพฺพตฺถ นว.) ปุเรชาเต ตีณิ, ปจฺฉาชาเต ตีณิ, อาเสวเน นว, กมฺเม ตีณิ, วิปาเก นว, อาหาเร นว, อินฺทฺริเย นว, ฌาเน ตีณิ, มคฺเค ตีณิ, สมฺปยุตฺเต ปญฺจ, วิปฺปยุตฺเต ปญฺจ, อตฺถิยา นว ฯเปฯ อวิคเต นว. (สํขิตฺตํ.)}

\prose{นเหตุยา นว, น-อารมฺมเณ นว. (สํขิตฺตํ.)}

\csromanheader{2}{เหตุปจฺจยา น-อารมฺมเณ ตีณิ. (สํขิตฺตํ.)}
\csromanheader{2}{นเหตุปจฺจยา อารมฺมเณ นว. (สํขิตฺตํ.)}
\prose{(ยถา กุสลตฺติเก ปญฺหาวารสฺส อนุโลมมฺปิ ปจฺจนียมฺปิ อนุโลมปจฺจนียมฺปิ ปจฺจนียานุโลมมฺปิ คณิตํ, เอวํ คเณตพฺพํ.)}

\csromanheader{2}{64. จิตฺตสํสฏฺฐสมุฏฺฐานสหภูทุก 1. กุสลตฺติก}
\csromanheader{2}{1. ปฏิจฺจวาร ปจฺจยจตุกฺก-เหตุ}
\proseitem{73}{จิตฺตสํสฏฺฐสมุฏฺฐานสหภุํ กุสลํ ธมฺมํ ปฏิจฺจ จิตฺตสํสฏฺฐสมุฏฺฐานสหภู กุสโล ธมฺโม อุปฺปชฺชติ เหตุปจฺจยา.\csromanspacer{} ตีณิ.}

\gambhira{ธมฺมานุโลม ทุกติกปฏฺฐานปาฬิ}
\prose{\csromanfolio{410}โนจิตฺตสํสฏฺฐสมุฏฺฐานสหภุํ กุสลํ ธมฺมํ ปฏิจฺจ จิตฺตสํสฏฺฐสมุฏฺฐานสหภู กุสโล ธมฺโม อุปฺปชฺชติ เหตุปจฺจยา. (1)}

\prose{จิตฺตสํสฏฺฐสมุฏฺฐานสหภุํ กุสลญฺจ โนจิตฺตสํสฏฺฐสมุฏฺฐานสหภุํ กิสลญฺจ ธมฺมํ ปฏิจฺจ จิตฺตสํสฏฺฐสมุฏฺฐานสหภู กุสโล ธมฺโม อุปฺปชฺชติ เหตุปจฺจยา. (1) (สํขิตฺตํ.)}

\proseitem{74}{เหตุยา ปญฺจ, อารมฺมเณ ปญฺจ ฯเปฯ อวิคเต ปญฺจ. (สํขิตฺตํ.) (สหชาตวารมฺปิ ฯเปฯ ปญฺหาวารมฺปิ วิตฺถาเรตพฺพํ.)}

\proseitem{75}{จิตฺตสํสฏฺฐสมุฏฺฐานสหภุํ อกุสลํ ธมฺมํ ปฏิจฺจ จิตฺตสํสฏฺฐสมุฏฺฐานสหภู อกุสโล ธมฺโม อุปฺปชฺชติ เหตุปจฺจยา.\csromanspacer{} ตีณิ.}

\prose{โนจิตฺตสํสฏฺฐสมุฏฺฐานสหภุํ อกุสลํ ธมฺมํ ปฏิจฺจ จิตฺตสํสฏฺฐสมุฏฺฐานสหภู อกุสโล ธมฺโม อุปฺปชฺชติ เหตุปจฺจยา. (1)}

\prose{จิตฺตสํสฏฺฐสมุฏฺฐานสหภุํ อกุสลญฺจ โนจิตฺตสํสฏฺฐสมุฏฺฐานสหภุํ อกุสลญฺจ ธมฺมํ ปฏิจฺจ จิตฺตสํสฏฺฐสมุฏฺฐานสหภู อกุสโล ธมฺโม อุปฺปชฺชติ เหตุปจฺจยา. (1) (สํขิตฺตํ.)}

\proseitem{76}{เหตุยา ปญฺจ, อารมฺมเณ ปญฺจ ฯเปฯ อวิคเต ปญฺจ. (สํขิตฺตํ.) (สหชาตวารมฺปิ ฯเปฯ ปญฺหาวารมฺปิ วิตฺถาเรตพฺพํ.)}

\proseitem{77}{จิตฺตสํสฏฺฐสมุฏฺฐานสหภุํ อพฺยากตํ ธมฺมํ ปฏิจฺจ จิตฺตสํสฏฺฐสมุฏฺฐานสหภู อพฺยากโต ธมฺโม อุปฺปชฺชติ เหตุปจฺจยา. (สํขิตฺตํ.)}

\proseitem{78}{เหตุยา นว, อารมฺมเณ นว ฯเปฯ อวิคเต นว. (สํขิตฺตํ.\csromanspacer{} จิตฺตสํสฏฺฐสมุฏฺฐานทุก-อพฺยากตสทิสํ.\csromanspacer{} สหชาตวารมฺปิ ฯเปฯ ปญฺหาวารมฺปิ วิตฺถาเรตพฺพํ.)}

\csromanheader{2}{65. จิตฺตสํสฏฺฐสมุฏฺฐานานุปริวตฺติทุก 1. กุสลตฺติก}
\csromanheader{2}{1. ปฏิจฺจวาร ปจฺจยจตุกฺก-เหตุ}
\proseitem{79}{จิตฺตสํสฏฺฐสมุฏฺฐานานุปริวตฺติํ กุสลํ ธมฺมํ ปฏิจฺจ จิตฺตสํสฏฺฐสมุฏฺฐานานุปริวตฺตี กุสโล ธมฺโม อุปฺปชฺชติ เหตุปจฺจยา.\csromanspacer{} ตีณิ.}

\csromanheader{2}{66. อชฺฌตฺติกทุก 1. กุสลตฺติก}
\prose{\csromanfolio{411}โนจิตฺตสํสฏฺฐสมุฏฺฐานานุปริวตฺติํ กุสลํ ธมฺมํ ปฏิจฺจ จิตฺตสํสฏฺฐสมุฏฺฐานานุปริวตฺตี กุสโล ธมฺโม อุปฺปชฺชติ เหตุปจฺจยา. (1)}

\prose{จิตฺตสํสฏฺฐสมุฏฺฐานานุปริวตฺติํ กุสลญฺจ โนจิตฺตสํสฏฺฐสมุฏฺฐานานุปริวตฺติํ กุสลญฺจ ธมฺมํ ปฏิจฺจ จิตฺตสํสฏฺฐสมุฏฺฐานานุปริวตฺตี กุสโล ธมฺโม อุปฺปชฺชติ เหตุปจฺจยา. (1) (สํขิตฺตํ.)}

\proseitem{80}{เหตุยา ปญฺจ, อารมฺมเณ ปญฺจ ฯเปฯ อวิคเต ปญฺจ. (สํขิตฺตํ.) (สหชาตวารมฺปิ ฯเปฯ ปญฺหาวารมฺปิ วิตฺถาเรตพฺพํ.)}

\proseitem{81}{จิตฺตสํสฏฺฐสมุฏฺฐานานุปริวตฺติํ อกุสลํ ธมฺมํ ปฏิจฺจ จิตฺตสํสฏฺฐสมุฏฺฐานานุปริวตฺตี อกุสโล ธมฺโม อุปฺปชฺชติ เหตุปจฺจยา. (สํขิตฺตํ.)}

\proseitem{82}{เหตุยา ปญฺจ, อรมฺมเณ ปญฺจ ฯเปฯ อวิคเต ปญฺจ. (สํขิตฺตํ.) (สหชาตวารมฺปิ ฯเปฯ ปญฺหาวารมฺปิ วิตฺถาเรตพฺพํ.)}

\proseitem{83}{จิตฺตสํสฏฺฐสมุฏฺฐานานุปริวตฺติํ อพฺยากตํ ธมฺมํ ปฏิจฺจ จิตฺตสํสฏฺฐสมุฏฺฐานานุปริวตฺตี อพฺยากโต ธมฺโม อุปฺปชฺชติ เหตุปจฺจยา. (สํขิตฺตํ.)}

\prose{เหตุยา นว, อารมฺมเณ นว ฯเปฯ อวิคเต นว. (สํขิตฺตํ.\csromanspacer{} จิตฺตสํสฏฺฐสมุฏฺฐานทุก-อพฺยากตสทิสํ.\csromanspacer{} สหชาตวารมฺปิ ฯเปฯ ปญฺหาวารมฺปิ วิตฺถาเรตพฺพํ.)}

\csromanheader{2}{66. อชฺฌตฺติกทุก 1. กุสลตฺติก}
\csromanheader{2}{1-7. ปฏิจฺจาทิวาร ปจฺจยจตุกฺก-เหตุ}
\proseitem{84}{อชฺฌตฺติกํ กุสลํ ธมฺมํ ปฏิจฺจ พาหิโร กุสโล ธมฺโม อุปฺปชฺชติ เหตุปจฺจยา. (1)}

\prose{พาหิรํ กุสลํ ธมฺมํ ปฏิจฺจ พาหิโร กุสโล ธมฺโม อุปฺปชฺชติ เหตุปจฺจยา.\csromanspacer{}\csromanspacer{}พาหิรํ กุสลํ ธมฺมํ ปฏิจฺจ อชฺฌตฺติโก กุสโล ธมฺโม อุปฺปชฺชติ เหตุปจฺจยา.\csromanspacer{}\csromanspacer{}พาหิรํ กุสลํ ธมฺมํ ปฏิจฺจ อชฺฌตฺติโก กุสโล จ พาหิโร กุสโล จ ธมฺมา อุปฺปชฺชนฺติ เหตุปจฺจยา. (3)}

\gambhira{ธมฺมานุโลม ทุกติกปฏฺฐานปาฬิ}
\prose{\csromanfolio{412}อชฺฌตฺติกํ กุสลญฺจ พาหิรํ กุสลญฺจ ธมฺมํ ปฏิจฺจ พาหิโร กุสโล ธมฺโม อุปฺปชฺชติ เหตุปจฺจยา. (1) (สํขิตฺตํ.)}

\proseitem{85}{เหตุยา ปญฺจ, อารมฺมเณ ปญฺจ ฯเปฯ อวิคเต ปญฺจ. (สํขิตฺตํ.\csromanspacer{} จิตฺตทุกสทิสํ.\csromanspacer{} สหชาตวารมฺปิ ฯเปฯ ปญฺหาวารมฺปิ วิตฺถาเรตพฺพํ.)}

\proseitem{86}{อชฺฌตฺติกํ อกุสลํ ธมฺมํ ปฏิจฺจ พาหิโร อกุสโล ธมฺโม อุปฺปชฺชติ เหตุปจฺจยา. (1)}

\prose{พาหิรํ อกุสลํ ธมฺมํ ปฏิจฺจ พาหิโร อกุสโล ธมฺโม อุปฺปชฺชติ เหตุปจฺจยา.\csromanspacer{}\csromanspacer{}พาหิรํ อกุสลํ ธมฺมํ ปฏิจฺจ อชฺฌตฺติโก อกุสโล ธมฺโม อุปฺปชฺชติ เหตุปจฺจยา.\csromanspacer{}\csromanspacer{}พาหิรํ อกุสลํ ธมฺมํ ปฏิจฺจ อชฺฌตฺติโก อกุสโล จ พาหิโร อกุสโล จ ธมฺมา อุปฺปชฺชนฺติ เหตุปจฺจยา. (3)}

\prose{อชฺฌตฺติกํ อกุสลญฺจ พาหิรํ อกุสลญฺจ ธมฺมํ ปฏิจฺจ พาหิโร อกุสโล ธมฺโม อุปฺปชฺชติ เหตุปจฺจยา. (1) (สํขิตฺตํ.)}

\proseitem{87}{เหตุยา ปญฺจ, อารมฺมเณ ปญฺจ ฯเปฯ อวิคเต ปญฺจ. (สํขิตฺตํ.\csromanspacer{} จิตฺตทุกสทิสํ.\csromanspacer{} สหชาตวารมฺปิ ฯเปฯ ปญฺหาวารมฺปิ วิตฺถาเรตพฺพํ.)}

\proseitem{88}{อชฺฌตฺติกํ อพฺยากตํ ธมฺมํ ปฏิจฺจ อชฺฌตฺติโก อพฺยากโต ธมฺโม อุปฺปชฺชติ เหตุปจฺจยา.\csromanspacer{} ตีณิ.}

\prose{พาหิรํ อพฺยากตํ ธมฺมํ ปฏิจฺจ พาหิโร อพฺยากโต ธมฺโม อุปฺปชฺชติ เหตุปจฺจยา.\csromanspacer{} ตีณิ.}

\prose{อชฺฌตฺติกํ อพฺยากตญฺจ พาหิรํ อพฺยากตญฺจ ธมฺมํ ปฏิจฺจ อชฺฌตฺติโก อพฺยากโต ธมฺโม อุปฺปชฺชติ เหตุปจฺจยา.\csromanspacer{} ตีณิ. (สํขิตฺตํ.)}

\proseitem{89}{เหตุยา นว, อารมฺมเณ ปญฺจ, อธิปติยา ปญฺจ ฯเปฯ อญฺญมญฺเญ ปญฺจ, นิสฺสเย นว, อุปนิสฺสเย ปญฺจ, ปุเรชาเต ปญฺจ, อาเสวเน ปญฺจ, กมฺเม ฯเปฯ มคฺเค นว, สมฺปยุตฺเต ปญฺจ, วิปฺปยุตฺเต นว ฯเปฯ อวิคเต นว. (สํขิตฺตํ.)}

\prose{นเหตุยา นว, น-อารมฺมเณ นว, น-อธิปติยา นว ฯเปฯ นกมฺเม ตีณิ, นวิปาเก ปญฺจ, น-อาหาเร เอกํ, น-อินฺทฺริเย เอกํ, นฌาเน ปญฺจ, นมคฺเค นว, นสมฺปยุตฺเต นว, นวิปฺปยุตฺเต ปญฺจ, โนนตฺถิยา นว, โนวิคเต นว. (สํขิตฺตํ.\csromanspacer{} สหชาตวารมฺปิ ฯเปฯ สมฺปยุตฺตวารมฺปิ วิตฺถาเรตพฺพํ.)}

\csromanheader{2}{67. อุปาทาทุก 1. กุสลตฺติก}
\proseitem{90}{\csromanfolio{413}พาหิโร อพฺยากโต ธมฺโม พาหิรสฺส อพฺยากตสฺส ธมฺมสฺส เหตุปจฺจเยน ปจฺจโย. (สํขิตฺตํ.)}

\prose{พาหิโร อพฺยากโต ธมฺโม พาหิรสฺส อพฺยากตสฺส ธมฺมสฺส กมฺมปจฺจเยน ปจฺจโย.\csromanspacer{} ตีณิ.}

\proseitem{91}{เหตุยา ตีณิ, อารมฺมเณ นว, อธิปติยา นว, (อชฺฌตฺติโก พาหิรสฺส คมนกาเล สหชาตาธิปติ, มชฺเฌ ตีสุปิ สหชาตาธิปติ.) อนนฺตเร นว, สหชาเต นว, อญฺญมญฺเญ ปญฺจ, นิสฺสเย นว, อุปนิสฺสเย นว, ปุเรชาเต นว, (อารมฺมณปุเรชาตมฺปิ วตฺถุปุเรชาตมฺปิ.) ปจฺฉาชาเต นว, อาเสวเน นว, กมฺเม ตีณิ, วิปาเก นว, อาหาเร นว, (ติณฺณนฺนํ กพฬีกาโร อาหาโร.) อินฺทฺริเย นว, (ติณฺณนฺนํ รูปชีวิตินฺทฺริยํ.) ฌาเน มคฺเค ตีณิ, สมฺปยุตฺเต ปญฺจ, วิปฺปยุตฺเต นว, อตฺถิยา นว ฯเปฯ อวิคเต นว. (สํขิตฺตํ.)}

\csromanheader{2}{67. อุปาทาทุก 1. กุสลตฺติก}
\csromanheader{2}{1-7. ปฏิจฺจาทิวาร ปจฺจยจตุกฺก-เหตุ}
\proseitem{92}{โน-อุปาทา กุสลํ ธมฺมํ ปฏิจฺจ โน-อุปาทา กุสโล ธมฺโม อุปฺปชฺชติ เหตุปจฺจยา. (สํขิตฺตํ.)}

\prose{เหตุยา เอกํ, อารมฺมเณ เอกํ ฯเปฯ อวิคเต เอกํ. (สํขิตฺตํ.) (สหชาตวาเรปิ ฯเปฯ ปญฺหาวาเรปิ สพฺพตฺถ เอกํ.)}

\proseitem{93}{โน-อุปาทา อกุสลํ ธมฺมํ ปฏิจฺจ โน-อุปาทา อกุสโล ธมฺโม อุปฺปชฺชติ เหตุปจฺจยา. (สํขิตฺตํ.)}

\prose{เหตุยา เอกํ, อารมฺมเณ เอกํ ฯเปฯ อวิคเต เอกํ. (สํขิตฺตํ.) (สหชาตวาเรปิ ฯเปฯ ปญฺหาวาเรปิ สพฺพตฺถ เอกํ.)}

\proseitem{94}{อุปาทา อพฺยากตํ ธมฺมํ ปฏิจฺจ โน-อุปาทา อพฺยากโต ธมฺโม อุปฺปชฺชติ เหตุปจฺจยา. (1)}

\prose{โน-อุปาทา อพฺยากตํ ธมฺมํ ปฏิจฺจ โน-อุปาทา อพฺยากโต ธมฺโม อุปฺปชฺชติ เหตุปจฺจยา.\csromanspacer{}\csromanspacer{}โน-อุปาทา อพฺยากตํ ธมฺมํ ปฏิจฺจ อุปาทา}

\vspace{\tipitakaparskip}
\csromancenter{ธมฺมานุโลม ทุกติกปฏฺฐานปาฬิ อพฺยากโต ธมฺโม อุปฺปชฺชติ เหตุปจฺจยา.\csromanspacer{}\csromanspacer{}โน-อุปาทา อพฺยากตํ ธมฺมํ ปฏิจฺจ อุปาทา อพฺยากโต จ โน-อุปาทา อพฺยากโต จ ธมฺมา อุปฺปชฺชนฺติ เหตุปจฺจยา. (3)}
\prose{\csromanfolio{414}อุปาทา อพฺยากตญฺจ โน-อุปาทา อพฺยากตญฺจ ธมฺมํ ปฏิจฺจ โน-อุปาทา อพฺยากโต ธมฺโม อุปฺปชฺชติ เหตุปจฺจยา. (1) (สํขิตฺตํ.)}

\proseitem{95}{เหตุยา ปญฺจ, อารมฺมเณ ตีณิ ฯเปฯ อญฺญมญฺเญ ปญฺจ ฯเปฯ ปุเรชาเต เอกํ, อาเสวเน เอกํ, กมฺเม ปญฺจ, วิปาเก ปญฺจ ฯเปฯ อวิคเต ปญฺจ. (สํขิตฺตํ.)}

\prose{นเหตุยา ปญฺจ, น-อารมฺมเณ ตีณิ, น-อธิปติยา ปญฺจ ฯเปฯ นปุเรชาเต ปญฺจ ฯเปฯ นกมฺเม ตีณิ, นวิปาเก ตีณิ, น-อาหาเร ตีณิ, น-อินฺทฺริเย ตีณิ, นฌาเน ตีณิ, นมคฺเค ปญฺจ, นสมฺปยุตฺเต ตีณิ, นวิปฺปยุตฺเต ตีณิ, โนนตฺถิยา ตีณิ, โนวิคเต ตีณิ. (สํขิตฺตํ.\csromanspacer{} สหชาตวารมฺปิ ฯเปฯ สมฺปยุตฺตวารมฺปิ วิตฺถาเรตพฺพํ.)}

\csromanheader{2}{\textbf{เหตุ อารมฺมณาทิ}}
\proseitem{96}{โน-อุปาทา อพฺยากโต ธมฺโม โน-อุปาทา อพฺยากตสฺส ธมฺมสฺส เหตุปจฺจเยน ปจฺจโย.\csromanspacer{} ตีณิ.}

\prose{อุปาทา อพฺยากโต ธมฺโม โน-อุปาทา อพฺยากตสฺส ธมฺมสฺส อารมฺมณปจฺจเยน ปจฺจโย. (1)}

\prose{โน-อุปาทา อพฺยากโต ธมฺโม โน-อุปาทา อพฺยากตสฺส ธมฺมสฺส อารมฺมณปจฺจเยน ปจฺจโย. (1)}

\prose{โน-อุปาทา อพฺยากโต ธมฺโม โน-อุปาทา อพฺยากตสฺส ธมฺมสฺส อธิปติปจฺจเยน ปจฺจโย.\csromanspacer{} ตีณิ. (ปฐเม เทฺวปิ อธิปตี.\csromanspacer{} ทฺวีสุ สหชาตาธิปติ.)}

\prose{อุปาทา อพฺยากโต ธมฺโม โน-อุปาทา อพฺยากตสฺส ธมฺมสฺส อุปนิสฺสยปจฺจเยน ปจฺจโย. (1)}

\prose{โน-อุปาทา อพฺยากโต ธมฺโม โน-อุปาทา อพฺยากตสฺส ธมฺมสฺส อุปนิสฺสยปจฺจเยน ปจฺจโย. (1) (สํขิตฺตํ.)}

\csromanmatikaanchor{mtk.77}
\csromanmatikaanchor{mtk.78}
\csromantocmarknopage{subparagraph}{68. อุปาทินฺนทุก}{mtk.77}
\bookmark[dest={mtk.77},level=5]{68. อุปาทินฺนทุก}
\csromantocmarkref{subparagraph}{1. กุสลตฺติก ...}{mtk.78}
\bookmark[dest={mtk.78},level=5]{1. กุสลตฺติก ...}
\csromanheader{6}{68. อุปาทินฺนทุก 1. กุสลตฺติก}
\proseitem{97}{\csromanfolio{415}อุปนิสฺสเย เทฺว, ปุเรชาเต ตีณิ, ปจฺฉาชาเต ตีณิ, อาเสวเน เอกํ, กมฺเม ตีณิ, วิปาเก ตีณิ, อาหาเร ฉ, อินฺทฺริเย สตฺต, ฌาเน ตีณิ, มคฺเค ตีณิ, สมฺปยุตฺเต เอกํ, วิปฺปยุตฺเต จตฺตาริ, อตฺถิยา นว ฯเปฯ อวิคเต นว. (สํขิตฺตํ.)}

\csromanheader{2}{68. อุปาทินฺนทุก 1. กุสลตฺติก}
\csromanheader{2}{1-7. ปฏิจฺจาทิวาร ปจฺจยจตุกฺก-เหตุ}
\proseitem{98}{อนุปาทินฺนํ กุสลํ ธมฺมํ ปฏิจฺจ อนุปาทินฺโน กุสโล ธมฺโม อุปฺปชฺชติ เหตุปจฺจยา. (สํขิตฺตํ.)}

\prose{เหตุยา เอกํ, อารมฺมเณ เอกํ ฯเปฯ อวิคเต เอกํ. (สพฺพตฺถ เอกํ.\csromanspacer{} สํขิตฺตํ.\csromanspacer{} สหชาตวาเรปิ ฯเปฯ ปญฺหาวาเรปิ เอกํ, สปฺปจฺจยกุสลสทิสํ.)}

\proseitem{99}{อนุปาทินฺนํ อกุสลํ ธมฺมํ ปฏิจฺจ อนุปาทินฺโน อกุสโล ธมฺโม อุปฺปชฺชติ เหตุปจฺจยา. (สํขิตฺตํ.)}

\prose{เหตุยา เอกํ ฯเปฯ อวิคเต เอกํ. (สํขิตฺตํ.\csromanspacer{} สหชาตวาเรปิ ฯเปฯ ปญฺหาวาเรปิ สพฺพตฺถ เอกํ.}

\proseitem{100}{อุปาทินฺนํ อพฺยากตํ ธมฺมํ ปฏิจฺจ อุปาทินฺโน อพฺยากโต ธมฺโม อุปฺปชฺชติ เหตุปจฺจยา.\csromanspacer{} ตีณิ.}

\prose{อนุปาทินฺนํ อพฺยากตํ ธมฺมํ ปฏิจฺจ อนุปาทินฺโน อพฺยากโต ธมฺโม อุปฺปชฺชติ เหตุปจฺจยา. (1)}

\prose{อุปาทินฺนํ อพฺยากตญฺจ อนุปาทินฺนํ อพฺยากตญฺจ ธมฺมํ ปฏิจฺจ อนุปาทินฺโน อพฺยากโต ธมฺโม อุปฺปชฺชติ เหตุปจฺจยา.\csromanspacer{} อุปาทินฺเน ขนฺเธ จ มหาภูเต จ ปฏิจฺจ จิตฺตสมุฏฺฐานํ รูปํ. (1) (สํขิตฺตํ.)}

\proseitem{101}{เหตุยา ปญฺจ, อารมฺมเณ เทฺว, อธิปติยา เอกํ ฯเปฯ ปุเรชาเต เทฺว, อาเสวเน เอกํ, กมฺเม ปญฺจ, วิปาเก ปญฺจ ฯเปฯ อวิคเต ปญฺจ. (สํขิตฺตํ.)}

\gambhira{ธมฺมานุโลม ทุกติกปฏฺฐานปาฬิ}
\prose{\csromanfolio{416}นเหตุยา ปญฺจ, น-อารมฺมเณ จตฺตาริ, น-อธิปติยา ปญฺจ ฯเปฯ นปุเรชาเต จตฺตาริ, นปจฺฉาชาเต น-อาเสวเน ปญฺจ, นกมฺเม เอกํ, นวิปาเก เทฺว, น-อาหาเร เทฺว, น-อินฺทฺริเย เทฺว, นฌาเน เทฺว, นมคฺเค ปญฺจ, นสมฺปยุตฺเต จตฺตาริ, นวิปฺปยุตฺเต เทฺว, โนนตฺถิยา จตฺตาริ, โนวิคเต จตฺตาริ. (สํขิตฺตํ.\csromanspacer{} สหชาตวาราทิ วิตฺถาเรตพฺโพ.)}

\csromanheader{2}{\textbf{เหตุ ปุเรชาต}}
\proseitem{102}{อุปาทินฺโน อพฺยากโต ธมฺโม อุปาทินฺนสฺส อพฺยากตสฺส ธมฺมสฺส เหตุปจฺจเยน ปจฺจโย.\csromanspacer{} ตีณิ.}

\prose{อนุปาทินฺโน อพฺยากโต ธมฺโม อนุปาทินฺนสฺส อพฺยากตสฺส ธมฺมสฺส เหตุปจฺจเยน ปจฺจโย. (1)}

\prose{อนุปาทินฺโน อพฺยากโต ธมฺโม อนุปาทินฺนสฺส อพฺยากตสฺส ธมฺมสฺส อธิปติปจฺจเยน ปจฺจโย.\csromanspacer{} อารมฺมณาธิปติ, สหชาตาธิปติ.\csromanspacer{} เอกํ.}

\prose{อุปาทินฺโน อพฺยากโต ธมฺโม อุปาทินฺนสฺส อพฺยากตสฺส ธมฺมสฺส ปุเรชาตปจฺจเยน ปจฺจโย.\csromanspacer{} อารมฺมณปุเรชาตํ, วตฺถุปุเรชาตํ. (เทฺว ปญฺหา.)}

\prose{อนุปาทินฺโน อพฺยากโต ธมฺโม อนุปาทินฺนสฺส อพฺยากตสฺส ธมฺมสฺส ปุเรชาตปจฺจเยน ปจฺจโย.\csromanspacer{} เทฺว. (อารมฺมณปุเรชาตํเยว.\csromanspacer{} ฆฏนา เทฺว, อารมฺมณปุเรชาตมฺปิ วตฺถุปุเรชาตมฺปิ.)}

\proseitem{103}{เหตุยา จตฺตาริ, อารมฺมเณ จตฺตาริ, (อุปาทินฺนมูลเก เทฺว.\csromanspacer{} อนุปาทินฺนมูลเก เทฺว.) อธิปติยา เอกํ, อนนฺตเร สมนนฺตเร จตฺตาริ, สหชาเต ปญฺจ, อญฺญมญฺเญ เทฺว, นิสฺสเย ปญฺจ, อุปนิสฺสเย จตฺตาริ, ปุเรชาเต ฉ, ปจฺฉาชาเต ฉ, อาเสวเน เอกํ, กมฺเม จตฺตาริ, วิปาเก จตฺตาริ, อาหาเร นว, อินฺทฺริเย จตฺตาริ, ฌาเน จตฺตาริ, มคฺเค จตฺตาริ, สมฺปยุตฺเต เทฺว, วิปฺปยุตฺเต ฉ, อตฺถิยา นว, นตฺถิยา จตฺตาริ, วิคเต จตฺตาริ, อวิคเต นว. (สํขิตฺตํ.)}

\prose{นเหตุยา นว, น-อารมฺมเณ นว. (สํขิตฺตํ.)}

\csromanheader{2}{เหตุปจฺจยา น-อารมฺมเณ จตฺตาริ. (สํขิตฺตํ.)}
\csromanheader{2}{นเหตุปจฺจยา อารมฺมเณ จตฺตาริ. (สํขิตฺตํ.)}
\nitthitam{มหนฺตรทุกกุสลตฺติกํ นิฏฺฐิตํ.}
\csromanmatikaanchor{mtk.79}
\csromanmatikaanchor{mtk.80}
\csromantocmarknopage{subparagraph}{69. อุปาทานทุก}{mtk.79}
\bookmark[dest={mtk.79},level=5]{69. อุปาทานทุก}
\csromantocmarkref{subparagraph}{1. กุสลตฺติก ...}{mtk.80}
\bookmark[dest={mtk.80},level=5]{1. กุสลตฺติก ...}
\csromanheader{6}{69. อุปาทานทุก 1. กุสลตฺติก \textbf{69. อุปาทานทุก 1. กุสลตฺติก}}
\csromanheader{2}{1-7. ปฏิจฺจาทิวาร ปจฺจยจตุกฺก}
\proseitem{1}{\csromanfolio{417}โน-อุปาทานํ กุสลํ ธมฺมํ ปฏิจฺจ โน-อุปาทาโน กุสโล ธมฺโม อุปฺปชฺชติ เหตุปจฺจยา. (สํขิตฺตํ.)}

\prose{เหตุยา เอกํ, อารมฺมเณ เอกํ ฯเปฯ อวิคเต เอกํ. (สํขิตฺตํ.) (สหชาตวาเรปิ ฯเปฯ ปญฺหาวาเรปิ สพฺพตฺถ เอกํ.)}

\proseitem{2}{อุปาทานํ อกุสลํ ธมฺมํ ปฏิจฺจ อุปาทาโน อกุสโล ธมฺโม อุปฺปชฺชติ เหตุปจฺจยา.\csromanspacer{} อุปาทานํ อกุสลํ ธมฺมํ ปฏิจฺจ โน-อุปาทาโน อกุสโล ธมฺโม อุปฺปชฺชติ เหตุปจฺจยา.\csromanspacer{}\csromanspacer{}อุปาทานํ อกุสลํ ธมฺมํ ปฏิจฺจ อุปาทาโน อกุสโล จ โน-อุปาทาโน อกุสโล จ ธมฺมา อุปฺปชฺชนฺติ เหตุปจฺจยา. (3)}

\prose{โน-อุปาทานํ อกุสลํ ธมฺมํ ปฏิจฺจ โน-อุปาทาโน อกุสโล ธมฺโม อุปฺปชฺชติ เหตุปจฺจยา.\csromanspacer{} ตีณิ.}

\prose{อุปาทานํ อกุสลญฺจ โน-อุปาทานํ อกุสลญฺจ ธมฺมํ ปฏิจฺจ อุปาทาโน อกุสโล ธมฺโม อุปฺปชฺชติ เหตุปจฺจยา.\csromanspacer{} ตีณิ. (สํขิตฺตํ.)}

\proseitem{3}{เหตุยา นว, อารมฺมเณ นว, (สพฺพตฺถ นว,) อวิคเต นว. (สํขิตฺตํ.)}

\proseitem{4}{โน-อุปาทานํ อกุสลํ ธมฺมํ ปฏิจฺจ โน-อุปาทาโน อกุสโล ธมฺโม อุปฺปชฺชติ นเหตุปจฺจยา. (สํขิตฺตํ.)}

\proseitem{5}{นเหตุยา เอกํ, น-อธิปติยา นว, นปุเรชาเต นว ฯเปฯ นกมฺเม ตีณิ, นวิปาเก นว, นวิปฺปยุตฺเต นว. (สํขิตฺตํ.\csromanspacer{} สหชาตวาราทิ วิตฺถาเตตพฺโพ.)}

\proseitem{6}{อุปาทาโน อกุสโล ธมฺโม อุปาทานสฺส อกุสลสฺส ธมฺมสฺส เหตุปจฺจเยน ปจฺจโย. (สํขิตฺตํ.)}

\proseitem{7}{เหตุยา นว, อารมฺมเณ นว, อธิปติยา นว, (โน-อุปาทานมูลเก ตีณิ.\csromanspacer{} สหชาตาธิปติ.) อนนฺตเร นว ฯเปฯ นิสฺสเย นว,}

\vspace{\tipitakaparskip}
\csromancenter{ธมฺมานุโลม ทุกติกปฏฺฐานปาฬิ อุปนิสฺสเย นว, อาเสวเน นว, กมฺเม ตีณิ, (โน-อุปาทานมูลเก.) อาหาเร ตีณิ, อินฺทฺริเย ตีณิ, ฌาเน ตีณิ, มคฺเค นว, สมฺปยุตฺเต นว, อตฺถิยา นว ฯเปฯ อวิคเต นว. (สํขิตฺตํ.)}
\proseitem{8}{\csromanfolio{418}โน-อุปาทานํ อพฺยากตํ ธมฺมํ ปฏิจฺจ โน-อุปาทาโน อพฺยากโต ธมฺโม อุปฺปชฺชติ เหตุปจฺจยา. (สํขิตฺตํ.)}

\proseitem{9}{เหตุยา เอกํ, อารมฺมเณ เอกํ ฯเปฯ อวิคเต เอกํ. (สํขิตฺตํ.) (สหชาตวาเรปิ ฯเปฯ ปญฺหาวาเรปิ สพฺพตฺถ เอกํ.)}

\csromanheader{2}{70. อุปาทานิยทุก 1. กุสลตฺติก}
\csromanheader{2}{1. ปฏิจฺจวาร ปจฺจยจตุกฺก-เหตุ}
\proseitem{10}{อุปาทานิยํ กุสลํ ธมฺมํ ปฏิจฺจ อุปาทานิโย กุสโล ธมฺโม อุปฺปชฺชติ เหตุปจฺจยา. (1)}

\prose{อนุปาทานิยํ กุสลํ ธมฺมํ ปฏิจฺจ อนุปาทานิโย กุสโล ธมฺโม อุปฺปชฺชติ เหตุปจฺจยา. (1) (สํขิตฺตํ.)}

\proseitem{11}{เหตุยา เทฺว, อารมฺมเณ เทฺว ฯเปฯ อวิคเต เทฺว. (สํขิตฺตํ.\csromanspacer{} โลกิยโลกุตฺตรทุกกุสลสทิสํ.\csromanspacer{} สหชาตวารมฺปิ ฯเปฯ ปญฺหาวารมฺปิ วิตฺถาเรตพฺพํ.)}

\proseitem{12}{อุปาทานิยํ อกุสลํ ธมฺมํ ปฏิจฺจ อุปาทานิโย อกุสโล ธมฺโม อุปฺปชฺชติ เหตุปจฺจยา. (สํขิตฺตํ.)}

\prose{เหตุยา เอกํ, อารมฺมเณ เอกํ ฯเปฯ อวิคเต เอกํ. (สํขิตฺตํ.) (สหชาตวาเรปิ ฯเปฯ ปญฺหาวาเรปิ สพฺพตฺถ เอกํ.)}

\proseitem{13}{อุปาทานิยํ อพฺยากตํ ธมฺมํ ปฏิจฺจ อุปาทานิโย อพฺยากโต ธมฺโม อุปฺปชฺชติ เหตุปจฺจยา. (1)}

\prose{อนุปาทานิยํ อพฺยากตํ ธมฺมํ ปฏิจฺจ อนุปาทานิโย อพฺยากโต ธมฺโม อุปฺปชฺชติ เหตุปจฺจยา.\csromanspacer{}\csromanspacer{}อนุปาทานิยํ อพฺยากตํ ธมฺมํ ปฏิจฺจ อุปาทานิโย}

\vspace{\tipitakaparskip}
\csromancenter{อุปาทานสมฺปยุตฺตทุก 1.\csromanspacer{} กุสลตฺติก อพฺยากโต ธมฺโม อุปฺปชฺชติ เหตุปจฺจยา.\csromanspacer{}\csromanspacer{}อนุปาทานิยํ อพฺยากตํ ธมฺมํ ปฏิจฺจ อุปาทานิโย อพฺยากโต จ อนุปาทานิโย อพฺยากโต จ ธมฺมา อุปฺปชฺชนฺติ เหตุปจฺจยา. (3)}
\prose{\csromanfolio{419}อุปาทานิยํ อพฺยากตญฺจ อนุปาทานิยํ อพฺยากตญฺจ ธมฺมํ ปฏิจฺจ อุปาทานิโย อพฺยากโต ธมฺโม อุปฺปชฺชติ เหตุปจฺจยา. (1) (สํขิตฺตํ.)}

\proseitem{14}{เหตุยา ปญฺจ, อารมฺมเณ เทฺว ฯเปฯ อาเสวเน เอกํ ฯเปฯ อวิคเต ปญฺจ. (สํขิตฺตํ.\csromanspacer{} โลกิยทุก-อพฺยากตสทิสํ.\csromanspacer{} สหชาตวารมฺปิ ฯเปฯ ปญฺหาวารมฺปิ วิตฺถาเรตพฺพํ.)}

\csromanheader{2}{71. อุปาทานสมฺปยุตฺตทุก 1. กุสลตฺติก}
\csromanheader{2}{1-7. ปฏิจฺจาทิวาร ปจฺจยจตุกฺก-เหตุ}
\proseitem{15}{อุปาทานวิปฺปยุตฺตํ กุสลํ ธมฺมํ ปฏิจฺจ อุปาทานวิปฺปยุตฺโต กุสโล ธมฺโม อุปฺปชฺชติ เหตุปจฺจยา. (สํขิตฺตํ.)}

\prose{เหตุยา เอกํ, อารมฺมเณ เอกํ ฯเปฯ อวิคเต เอกํ. (สํขิตฺตํ.) (สหชาตวาเรปิ ฯเปฯ ปญฺหาวาเรปิ สพฺพตฺถ เอกํ.)}

\proseitem{16}{อุปาทานสมฺปยุตฺตํ อกุสลํ ธมฺมํ ปฏิจฺจ อุปาทานสมฺปยุตฺโต อกุสโล ธมฺโม อุปฺปชฺชติ เหตุปจฺจยา.\csromanspacer{} ตีณิ.}

\prose{อุปาทานวิปฺปยุตฺตํ อกุสลํ ธมฺมํ ปฏิจฺจ อุปาทานวิปฺปยุตฺโต อกุสโล ธมฺโม อุปฺปชฺชติ เหตุปจฺจยา.\csromanspacer{} เทฺว.}

\prose{อุปาทานสมฺปยุตฺตํ อกุสลญฺจ อุปาทานวิปฺปยุตฺตํ อกุสลญฺจ ธมฺมํ ปฏิจฺจ อุปาทานสมฺปยุตฺโต อกุสโล ธมฺโม อุปฺปชฺชติ เหตุปจฺจยา. (สํขิตฺตํ.)}

\proseitem{17}{เหตุยา ฉ, อารมฺมเณ ฉ, (สพฺพตฺถ ฉ,) อวิคเต ฉ. (สํขิตฺตํ.)}

\prose{นเหตุยา เอกํ, น-อธิปติยา ฉ ฯเปฯ นกมฺเม จตฺตาริ, นวิปาเก ฉ.\csromanspacer{} นวิปฺปยุตฺเต ฉ. (สํขิตฺตํ.\csromanspacer{} สหชาตวาราทิ วิตฺถาเรตพฺโพ.)}

\gambhira{ธมฺมานุโลม ทุกติกปฏฺฐานปาฬิ}
\proseitem{18}{\csromanfolio{420}อุปาทานสมฺปยุตฺโต อกุสโล ธมฺโม อุปาทานสมฺปยุตฺตสฺส อกุสลสฺส ธมฺมสฺส เหตุปจฺจเยน ปจฺจโย.\csromanspacer{} ตีณิ.}

\prose{อุปาทานวิปฺปยุตฺโต อกุสโล ธมฺโม อุปาทานวิปฺปยุตฺตสฺส อกุสลสฺส ธมฺมสฺส เหตุปจฺจเยน ปจฺจโย.\csromanspacer{} เทฺว.}

\prose{อุปาทานสมฺปยุตฺโต อกุสโล จ อุปาทานวิปฺปยุตฺโต อกุสโล จ ธมฺมา อุปาทานสมฺปยุตฺตสฺส อกุสลสฺส ธมฺมสฺส เหตุปจฺจเยน ปจฺจโย. (สํขิตฺตํ.)}

\proseitem{19}{เหตุยา ฉ, อารมฺมเณ นว, อธิปติยา นว. (อุปาทานสมฺปยุตฺตมูลเก ตีณิ, สหชาตาธิปติ, อุปาทานวิปฺปยุตฺเต เอกํ.) อนนฺตเร นว, สมนนฺตเร นว, สหชาเต ฉ, อญฺญมญฺเญ ฉ, นิสฺสเย ฉ, อุปนิสฺสเย นว, อาเสวเน นว, กมฺเม จตฺตาริ, อาหาเร จตฺตาริ, อินฺทฺริเย จตฺตาริ, ฌาเน จตฺตาริ, มคฺเค จตฺตาริ, สมฺปยุตฺเต ฉ, อตฺถิยา ฉ, นตฺถิยา นว, วิคเต นว, อวิคเต ฉ. (สํขิตฺตํ.)}

\proseitem{20}{อุปาทานวิปฺปยุตฺตํ อพฺยากตํ ธมฺมํ ปฏิจฺจ อุปาทานวิปฺปยุตฺโต อพฺยากโต ธมฺโม อุปฺปชฺชติ เหตุปจฺจยา. (สํขิตฺตํ.)}

\prose{เหตุยา เอกํ, อารมฺมเณ เอกํ ฯเปฯ อวิคเต เอกํ. (สํขิตฺตํ.)}

\vspace{\tipitakaparskip}
\csromancenter{(สหชาตวาเรปิ ฯเปฯ ปญฺหาวาเรปิ สพฺพตฺถ เอกํ.)}
\csromanheader{2}{72. อุปาทาน-อุปาทานิยทุก 1. กุสลตฺติก}
\csromanheader{2}{1. ปฏิจฺจวาร ปจฺจยจตุกฺก-เหตุ}
\proseitem{21}{อุปาทานิยํ เจว โน จ อุปาทานํ กุสลํ ธมฺมํ ปฏิจฺจ อุปาทานิโย เจว โน จ อุปาทาโน กุสโล ธมฺโม อุปฺปชฺชติ เหตุปจฺจยา. (สํขิตฺตํ.)}

\prose{เหตุยา เอกํ, อารมฺมเณ เอกํ ฯเปฯ อวิคเต เอกํ. (สํขิตฺตํ.) (สหชาตวาเรปิ ฯเปฯ ปญฺหาวาเรปิ สพฺพตฺถ เอกํ.)}

\proseitem{22}{อุปาทานญฺเจว อุปาทานิยญฺจ อกุสลํ ธมฺมํ ปฏิจฺจ อุปาทาโน เจว อุปาทานิโย จ อกุสโล ธมฺโม อุปฺปชฺชติ เหตุปจฺจยา.\csromanspacer{} ตีณิ.}

\csromanheader{2}{73. อุปาทาน-อุปาทานสมฺปยุตฺตทุก 1. กุสลตฺติก}
\prose{\csromanfolio{421}อุปาทานิยญฺเจว โน จ อุปาทานํ อกุสลํ ธมฺมํ ปฏิจฺจ อุปาทานิโย เจว โน จ อุปาทาโน อกุสโล ธมฺโม อุปฺปชฺชติ เหตุปจฺจยา.\csromanspacer{} ตีณิ.}

\prose{อุปาทานญฺเจว อุปาทานิยํ อกุสลญฺจ อุปาทานิยญฺเจว โน จ อุปาทานํ อกุสลญฺจ ธมฺมํ ปฏิจฺจ อุปาทาโน เจว อุปาทานิโย จ อกุสโล ธมฺโม อุปฺปชฺชติ เหตุปจฺจยา.\csromanspacer{} ตีณิ. (สํขิตฺตํ.)}

\proseitem{23}{เหตุยา นว, อารมฺมเณ นว ฯเปฯ อวิคเต นว. (สํขิตฺตํ.\csromanspacer{} อิมํ คมนํ อุปาทานทุก-อกุสลสทิสํ.\csromanspacer{} สหชาตวารมฺปิ ฯเปฯ ปญฺหาวารมฺปิ วิตฺถาเรตพฺพํ.)}

\proseitem{24}{อุปาทานิยญฺเจว โน จ อุปาทานํ อพฺยากตํ ธมฺมํ ปฏิจฺจ อุปาทานิโย เจว โน จ อุปาทาโน อพฺยากโต ธมฺโม อุปฺปชฺชติ เหตุปจฺจยา. (สํขิตฺตํ.)}

\proseitem{25}{เหตุยา เอกํ, อารมฺมเณ เอกํ ฯเปฯ อวิคเต เอกํ. (สํขิตฺตํ.)}

\vspace{\tipitakaparskip}
\csromancenter{(สหชาตวาเรปิ ฯเปฯ ปญฺหาวาเรปิ สพฺพตฺถ เอกํ.)}
\csromanheader{2}{73. อุปาทาน-อุปาทานสมฺปยุตฺตทุก 1. กุสลตฺติก}
\csromanheader{2}{1. ปฏิจฺจวาร ปจฺจยจตุกฺก-เหตุ}
\proseitem{26}{อุปาทานญฺเจว อุปาทานสมฺปยุตฺตญฺจ อกุสลํ ธมฺมํ ปฏิจฺจ อุปาทาโน เจว อุปาทานสมฺปยุตฺโต จ อกุสโล ธมฺโม อุปฺปชฺชติ เหตุปจฺจยา.\csromanspacer{} ตีณิ.}

\prose{อุปาทานสมฺปยุตฺตญฺเจว โน จ อุปาทานํ อกุสลํ ธมฺมํ ปฏิจฺจ อุปาทานสมฺปยุตฺโต เจว โน จ อุปาทาโน อกุสโล ธมฺโม อุปฺปชฺชติ เหตุปจฺจยา.\csromanspacer{} ตีณิ.}

\prose{อุปาทานญฺเจว อุปาทานสมฺปยุตฺตํ อกุสลญฺจ อุปาทานสมฺปยุตฺตญฺเจว โน จ อุปาทานํ อกุสลญฺจ ธมฺมํ ปฏิจฺจ อุปาทาโน เจว อุปาทานสมฺปยุตฺโต จ อกุสโล ธมฺโม อุปฺปชฺชติ เหตุปจฺจยา.\csromanspacer{} ตีณิ. (สํขิตฺตํ.)}

\gambhira{ธมฺมานุโลม ทุกติกปฏฺฐานปาฬิ}
\proseitem{27}{\csromanfolio{422}เหตุยา นว, อารมฺมเณ นว, อธิปติยา นว ฯเปฯ อวิคเต นว. (สํขิตฺตํ.\csromanspacer{} อิมํ คมนํ อุปาทานทุก-อกุสลสทิสํ.\csromanspacer{} อิธ ปน เหตุปจฺจยา อิทํ นานํ.\csromanspacer{} สหชาตวาเรปิ ฯเปฯ ปญฺหาวาเรปิ สพฺพตฺถ นว.)}

\csromanheader{2}{74. อุปาทานวิปฺปยุตฺต-อุปาทานิยทุก 1. กุสลตฺติก}
\csromanheader{2}{1. ปฏิจฺจวาร ปจฺจยจตุกฺก-เหตุ}
\proseitem{28}{อุปาทานวิปฺปยุตฺตํ อุปาทานิยํ กุสลํ ธมฺมํ ปฏิจฺจ อุปาทานวิปฺปยุตฺโต อุปาทานิโย กุสโล ธมฺโม อุปฺปชฺชติ เหตุปจฺจยา. (1)}

\prose{อุปาทานวิปฺปยุตฺตํ อนุปาทานิยํ กุสลํ ธมฺมํ ปฏิจฺจ อุปาทานวิปฺปยุตฺโต อนุปาทานิโย กุสโล ธมฺโม อุปฺปชฺชติ เหตุปจฺจยา. (1) (สํขิตฺตํ.)}

\proseitem{29}{เหตุยา เทฺว, อารมฺมเณ เทฺว ฯเปฯ อวิคเต เทฺว. (สํขิตฺตํ.\csromanspacer{} โลกิยทุกกุสลสทิสํ.\csromanspacer{} สหชาตวารมฺปิ ฯเปฯ ปญฺหาวารมฺปิ วิตฺถาเรตพฺพํ.)}

\proseitem{30}{อุปาทานวิปฺปยุตฺตํ อุปาทานิยํ อกุสลํ ธมฺมํ ปฏิจฺจ อุปาทานวิปฺปยุตฺโต อุปาทานิโย อกุสโล ธมฺโม อุปฺปชฺชติ เหตุปจฺจยา. (สพฺพตฺถ เอกํ.)}

\prose{อุปาทานวิปฺปยุตฺตํ อุปาทานิยํ อพฺยากตํ ธมฺมํ ปฏิจฺจ อุปาทานวิปฺปยุตฺโต อุปาทานิโย อพฺยากโต ธมฺโม อุปฺปชฺชติ เหตุปจฺจยา. (1)}

\prose{อุปาทานวิปฺปยุตฺตํ อนุปาทานิยํ อพฺยากตํ ธมฺมํ ปฏิจฺจ อุปาทานวิปฺปยุตฺโต อนุปาทานิโย อพฺยากโต ธมฺโม อุปฺปชฺชติ เหตุปจฺจยา.\csromanspacer{} ตีณิ.}

\prose{อุปาทานวิปฺปยุตฺตํ อุปาทานิยํ อพฺยากตญฺจ อุปาทานวิปฺปยุตฺตํ อนุปาทานิยํ อพฺยากตญฺจ ธมฺมํ ปฏิจฺจ อุปาทานวิปฺปยุตฺโต อุปาทานิโย อพฺยากโต ธมฺโม อุปฺปชฺชติ เหตุปจฺจยา. (1) (สํขิตฺตํ.)}

\proseitem{31}{เหตุยา ปญฺจ, อารมฺมเณ เทฺว ฯเปฯ อวิคเต ปญฺจ. (สํขิตฺตํ.\csromanspacer{} โลกิยทุก-อพฺยากตสทิสํ.\csromanspacer{} สหชาตวารมฺปิ ฯเปฯ ปญฺหาวารมฺปิ สพฺพตฺถ วิตฺถาเรตพฺพํ.)}

\nitthitam{อุปาทานโคจฺฉกกุสลตฺติกํ นิฏฺฐิตํ.}
\csromanmatikaanchor{mtk.81}
\csromanmatikaanchor{mtk.82}
\csromantocmarknopage{subparagraph}{75. กิเลสทุก}{mtk.81}
\bookmark[dest={mtk.81},level=5]{75. กิเลสทุก}
\csromantocmarkref{subparagraph}{1. กุสลตฺติก ...}{mtk.82}
\bookmark[dest={mtk.82},level=5]{1. กุสลตฺติก ...}
\csromanheader{6}{75. กิเลสทุก 1. กุสลตฺติก \textbf{75. กิเลสทุก 1. กุสลตฺติก}}
\csromanheader{2}{1-7. ปฏิจฺจาทิวาร ปจฺจยจตุกฺก-เหตุ}
\proseitem{1}{\csromanfolio{423}โนกิเลสํ กุสลํ ธมฺมํ ปฏิจฺจ โนกิเลโส กุสโล ธมฺโม อุปฺปชฺชติ เหตุปจฺจยา. (สํขิตฺตํ.)}

\prose{เหตุยา เอกํ, อารมฺมเณ เอกํ ฯเปฯ อวิคเต เอกํ. (สํขิตฺตํ.)}

\prose{(สหชาตวาเรปิ ฯเปฯ ปญฺหาวาเรปิ สพฺพตฺถ เอกํ.)}

\proseitem{2}{กิเลสํ อกุสลํ ธมฺมํ ปฏิจฺจ กิเลโส อกุสโล ธมฺโม อุปฺปชฺชติ เหตุปจฺจยา.\csromanspacer{} ตีณิ.}

\prose{โนกิเลสํ อกุสลํ ธมฺมํ ปฏิจฺจ โนกิเลโส อกุสโล ธมฺโม อุปฺปชฺชติ เหตุปจฺจยา.\csromanspacer{} ตีณิ.}

\prose{กิเลสํ อกุสลญฺจ โนกิเลสํ อกุสลญฺจ ธมฺมํ ปฏิจฺจ กิเลโส อกุสโล ธมฺโม อุปฺปชฺชติ เหตุปจฺจยา.\csromanspacer{} ตีณิ. (สํขิตฺตํ.)}

\proseitem{3}{เหตุยา นว, อารมฺมเณ นว, (สพฺพตฺถ นว,) อวิคเต นว. (สํขิตฺตํ.)}

\proseitem{4}{กิเลสํ อกุสลํ ธมฺมํ ปฏิจฺจ กิเลโส อกุสโล ธมฺโม อุปฺปชฺชติ นเหตุปจฺจยา.\csromanspacer{} วิจิกิจฺฉํ ปฏิจฺจ วิจิกิจฺฉาสหคโต โมโห.\csromanspacer{} อุทฺธจฺจํ ปฏิจฺจ อุทฺธจฺจสหคโต โมโห. (1)}

\prose{โนกิเลสํ อกุสลํ ธมฺมํ ปฏิจฺจ กิเลโส อกุสโล ธมฺโม อุปฺปชฺชติ นเหตุปจฺจยา.\csromanspacer{} วิจิกิจฺฉาสหคเต อุทฺธจฺจสหคเต ขนฺเธ ปฏิจฺจ วิจิกิจฺฉาสหคโต อุทฺธจฺจสหคโต โมโห. (1)}

\prose{กิเลสํ อกุสลญฺจ โนกิเลสํ อกุสลญฺจ ธมฺมํ ปฏิจฺจ กิเลโส อกุสโล ธมฺโม อุปฺปชฺชติ นเหตุปจฺจยา.\csromanspacer{} วิจิกิจฺฉาสหคเต ขนฺเธ จ วิจิกิจฺฉญฺจ ปฏิจฺจ วิจิกิจฺฉาสหคโต โมโห.\csromanspacer{} อุทฺธจฺจสหคเต ขนฺเธ จ อุทฺธจฺจญฺจ ปฏิจฺจ อุทฺธจฺจสหคโต โมโห. (1) (สํขิตฺตํ.)}

\proseitem{5}{นเหตุยา ตีณิ, น-อธิปติยา นว, นปุเรชาเต นว, นปจฺฉาชาเต นว, น-อาเสวเน นว, นกมฺเม ตีณิ, นวิปาเก นว, นวิปฺปยุตฺเต นว. (สํขิตฺตํ.\csromanspacer{} สหชาตวาราทิ วิตฺถาเรตพฺโพ.)}

\gambhira{ธมฺมานุโลม ทุกติกปฏฺฐานปาฬิ}
\proseitem{6}{\csromanfolio{424}กิเลโส อกุสโล ธมฺโม กิเลสสฺส อกุสลสฺส ธมฺมสฺส เหตุปจฺจเยน ปจฺจโย. (สํขิตฺตํ.)}

\proseitem{7}{เหตุยา ตีณิ, อารมฺมเณ นว, อธิปติยา นว, (มชฺเฌ ตีณิ, สหชาตาธิปติ.) อนนฺตเร นว, สมนนฺตเร นว ฯเปฯ อุปนิสฺสเย นว, อาเสวเน นว, กมฺเม ตีณิ, อาหาเร ตีณิ, อินฺทฺริเย ตีณิ, ฌาเน ตีณิ, มคฺเค นว, สมฺปยุตฺเต นว ฯเปฯ อวิคเต นว. (สํขิตฺตํ.)}

\prose{โนกิเลสํ อพฺยากตํ ธมฺมํ ปฏิจฺจ โนกิเลโส อพฺยากโต ธมฺโม อุปฺปชฺชติ เหตุปจฺจยา. (สํขิตฺตํ.)}

\proseitem{8}{เหตุยา เอกํ, อารมฺมเณ เอกํ ฯเปฯ อวิคเต เอกํ. (สํขิตฺตํ.)}

\vspace{\tipitakaparskip}
\csromancenter{(สหชาตวาเรปิ ฯเปฯ ปญฺหาวาเรปิ สพฺพตฺถ เอกํ.)}
\csromanheader{2}{76. สํกิเลสิกทุก 1. กุสลตฺติก}
\csromanheader{2}{1. ปฏิจฺจวาร ปจฺจยจตุกฺก-เหตุ}
\proseitem{9}{สํกิเลสิกํ กุสลํ ธมฺมํ ปฏิจฺจ สํกิเลสิโก กุสโล ธมฺโม อุปฺปชฺชติ เหตุปจฺจยา. (1)}

\prose{อสํกิเลสิกํ กุสลํ ธมฺมํ ปฏิจฺจ อสํกิเลสิโก กุสโล ธมฺโม อุปฺปชฺชติ เหตุปจฺจยา. (1) (สํขิตฺตํ.)}

\proseitem{10}{เหตุยา เทฺว, อารมฺมเณ เทฺว ฯเปฯ อวิคเต เทฺว. (สํขิตฺตํ.\csromanspacer{} โลกิยทุกกุสลสทิสํ.\csromanspacer{} สหชาตวารมฺปิ ฯเปฯ ปญฺหาวารมฺปิ วิตฺถาเรตพฺพํ.)}

\proseitem{11}{สํกิเลสิกํ อกุสลํ ธมฺมํ ปฏิจฺจ สํกิเลสิโก อกุสโล ธมฺโม อุปฺปชฺชติ เหตุปจฺจยา. (สํขิตฺตํ.)}

\prose{เหตุยา เอกํ, อารมฺมเณ เอกํ ฯเปฯ อวิคเต เอกํ. (สํขิตฺตํ.)}

\vspace{\tipitakaparskip}
\csromancenter{(สหชาตวาเรปิ ฯเปฯ ปญฺหาวาเรปิ สพฺพตฺถ เอกํ.)}
\proseitem{12}{สํกิเลสิกํ อพฺยากตํ ธมฺมํ ปฏิจฺจ สํกิเลสิโก อพฺยากโต ธมฺโม อุปฺปชฺชติ เหตุปจฺจยา. (1)}

\csromanheader{2}{78. กิเลสสมฺปยุตฺตทุก 1. กุสลตฺติก}
\prose{\csromanfolio{425}อสํกิเลสิกํ อพฺยากตํ ธมฺมํ ปฏิจฺจ อสํกิเลสิโก อพฺยากโต ธมฺโม อุปฺปชฺชติ เหตุปจฺจยา.\csromanspacer{} ตีณิ.}

\prose{สํกิเลสิกํ อพฺยากตญฺจ อสํกิเลสิกํ อพฺยากตญฺจ ธมฺมํ ปฏิจฺจ สํกิเลสิโก อพฺยากโต ธมฺโม อุปฺปชฺชติ เหตุปจฺจยา. (สํขิตฺตํ.)}

\proseitem{13}{เหตุยา ปญฺจ, อารมฺมเณ เทฺว ฯเปฯ อวิคเต ปญฺจ. (สํขิตฺตํ.\csromanspacer{} โลกิยทุก-อพฺยากตสทิสํ.\csromanspacer{} สหชาตวาเรปิ ฯเปฯ ปญฺหาวาเรปิ สพฺพตฺถ วิตฺถาเรตพฺพํ.)}

\csromanheader{2}{77. สํกิลิฏฺฐทุก 1. กุสลตฺติก}
\csromanheader{2}{1. ปฏิจฺจวาร ปจฺจยจตุกฺก}
\proseitem{14}{อสํกิลิฏฺฐํ กุสลํ ธมฺมํ ปฏิจฺจ อสํกิลิฏฺโฐ กุสโล ธมฺโม อุปฺปชฺชติ เหตุปจฺจยา. (สํขิตฺตํ.)}

\prose{เหตุยา เอกํ, อารมฺมเณ เอกํ ฯเปฯ อวิคเต เอกํ. (สํขิตฺตํ.)}

\prose{(สหชาตวาเรปิ ฯเปฯ ปญฺหาวาเรปิ สพฺพตฺถ เอกํ.)}

\proseitem{15}{สํกิลิฏฺฐํ อกุสลํ ธมฺมํ ปฏิจฺจ สํกิลิฏฺโฐ อกุสโล ธมฺโม อุปฺปชฺชติ เหตุปจฺจยา. (สํขิตฺตํ.)}

\prose{เหตุยา เอกํ, อารมฺมเณ เอกํ ฯเปฯ อวิคเต เอกํ. (สํขิตฺตํ.)}

\prose{(สหชาตวาเรปิ ฯเปฯ ปญฺหาวาเรปิ สพฺพตฺถ เอกํ.)}

\proseitem{16}{อสํกิลิฏฺฐํ อพฺยากตํ ธมฺมํ ปฏิจฺจ อสํกิลิฏฺโฐ อพฺยากโต ธมฺโม อุปฺปชฺชติ เหตุปจฺจยา. (สํขิตฺตํ.)}

\prose{เหตุยา เอกํ, อารมฺมเณ เอกํ ฯเปฯ อวิคเต เอกํ. (สํขิตฺตํ.)}

\prose{(สหชาตวาเรปิ ฯเปฯ ปญฺหาวาเรปิ สพฺพตฺถ เอกํ.)}

\csromanheader{2}{78. กิเลสสมฺปยุตฺตทุก 1. กุสลตฺติก}
\csromanheader{2}{1. ปฏิจฺจวาร ปจฺจยจตุกฺก}
\proseitem{17}{กิเลสวิปฺปยุตฺตํ กุสลํ ธมฺมํ ปฏิจฺจ กิเลสวิปฺปยุตฺโต กุสโล ธมฺโม อุปฺปชฺชติ เหตุปจฺจยา. (สํขิตฺตํ.)}

\gambhira{ธมฺมานุโลม ทุกติกปฏฺฐานปาฬิ}
\prose{\csromanfolio{426}เหตุยา เอกํ, อารมฺมเณ เอกํ ฯเปฯ อวิคเต เอกํ. (สํขิตฺตํ.)}

\prose{(สหชาตวาเรปิ ฯเปฯ ปญฺหาวาเรปิ สพฺพตฺถ เอกํ.)}

\proseitem{18}{กิเลสสมฺปยุตฺตํ อกุสลํ ธมฺมํ ปฏิจฺจ กิเลสสมฺปยุตฺโต อกุสโล ธมฺโม อุปฺปชฺชติ เหตุปจฺจยา. (สํขิตฺตํ.)}

\prose{เหตุยา เอกํ, อารมฺมเณ เอกํ ฯเปฯ อวิคเต เอกํ. (สํขิตฺตํ.)}

\prose{(สหชาตวาเรปิ ฯเปฯ ปญฺหาวาเรปิ สพฺพตฺถ เอกํ.)}

\proseitem{19}{กิเลสวิปฺปยุตฺตํ อพฺยากตํ ธมฺมํ ปฏิจฺจ กิเลสวิปฺปยุตฺโต อพฺยากโต ธมฺโม อุปฺปชฺชติ เหตุปจฺจยา. (สํขิตฺตํ.)}

\prose{เหตุยา เอกํ, อารมฺมเณ เอกํ ฯเปฯ อวิคเต เอกํ. (สํขิตฺตํ.)}

\prose{(สหชาตวาเรปิ ฯเปฯ ปญฺหาวาเรปิ สพฺพตฺถ เอกํ.)}

\csromanheader{2}{79. สํกิเลสสํกิเลสิกทุก 1. กุสลตฺติก}
\csromanheader{2}{1. ปฏิจฺจวาร ปจฺจยจตุกฺก}
\proseitem{20}{สํกิเลสิกญฺเจว โน จ กิเลสํ กุสลํ ธมฺมํ ปฏิจฺจ สํกิเลสิโก เจว โน จ กิเลโส กุสโล ธมฺโม อุปฺปชฺชติ เหตุปจฺจยา. (สํขิตฺตํ.)}

\prose{เหตุยา เอกํ, อารมฺมเณ เอกํ ฯเปฯ อวิคเต เอกํ. (สํขิตฺตํ.)}

\prose{(สหชาตวาเรปิ ฯเปฯ ปญฺหาวาเรปิ สพฺพตฺถ เอกํ.)}

\proseitem{21}{กิเลสญฺเจว สํกิเลสิกญฺจ อกุสลํ ธมฺมํ ปฏิจฺจ กิเลโส เจว สํกิเลสิโก จ อกุสโล ธมฺโม อุปฺปชฺชติ เหตุปจฺจยา.\csromanspacer{} ตีณิ.}

\prose{สํกิเลสิกญฺเจว โน จ กิเลสํ อกุสลํ ธมฺมํ ปฏิจฺจ สํกิเลสิโก เจว โน จ กิเลโส อกุสโล ธมฺโม อุปฺปชฺชติ เหตุปจฺจยา.\csromanspacer{} ตีณิ.}

\prose{กิเลสญฺเจว สํกิเลสิกํ อกุสลญฺจ สํกิเลสิกญฺเจว โน จ กิเลสํ อกุสลญฺจ ธมฺมํ ปฏิจฺจ กิเลโส เจว สํกิเลสิโก จ อกุสโล ธมฺโม อุปฺปชฺชติ เหตุปจฺจยา.\csromanspacer{} ตีณิ. (สํขิตฺตํ.)}

\csromanheader{2}{81. กิเลสกิเลสสมฺปยุตฺตทุก 1. กุสลตฺติก}
\proseitem{22}{\csromanfolio{427}เหตุยา นว, อารมฺมเณ นว ฯเปฯ อวิคเต นว. (สํขิตฺตํ.\csromanspacer{} กิเลสทุก-อกุสลสทิสํ.\csromanspacer{} สหชาตวาเรปิ ฯเปฯ ปญฺหาวาเรปิ สพฺพตฺถ วิตฺถาเรตพฺพํ.)}

\proseitem{23}{สํกิเลสิกญฺเจว โน จ กิเลสํ อพฺยากตํ ธมฺมํ ปฏิจฺจ สํกิเลสิโก เจว โน จ กิเลโส อพฺยากโต ธมฺโม อุปฺปชฺชติ เหตุปจฺจยา. (สํขิตฺตํ.)}

\prose{เหตุยา เอกํ, อารมฺมเณ เอกํ ฯเปฯ อวิคเต เอกํ. (สํขิตฺตํ.)}

\vspace{\tipitakaparskip}
\csromancenter{(สหชาตวาเรปิ ฯเปฯ ปญฺหาวาเรปิ สพฺพตฺถ เอกํ.)}
\csromanheader{2}{80. กิเลสสํกิลิฏฺฐทุก 1. กุสลตฺติก}
\csromanheader{2}{1. ปฏิจฺจวาร ปจฺจยจตุกฺก}
\proseitem{24}{กิเลสญฺเจว สํกิลิฏฺฐญฺจ อกุสลํ ธมฺมํ ปฏิจฺจ กิเลโส เจว สํกิลิฏฺโฐ จ อกุสโล ธมฺโม อุปฺปชฺชติ เหตุปจฺจยา.\csromanspacer{} ตีณิ.}

\prose{สํกิลิฏฺฐญฺเจว โน จ กิเลสํ อกุสลํ ธมฺมํ ปฏิจฺจ สํกิลิฏฺโฐ เจว โน จ กิเลโส อกุสโล ธมฺโม อุปฺปชฺชติ เหตุปจฺจยา.\csromanspacer{} ตีณิ.}

\prose{กิเลสญฺเจว สํกิลิฏฺฐํ อกุสลญฺจ สํกิลิฏฺฐญฺเจว โน จ กิเลสํ อกุสลญฺจ ธมฺมํ ปฏิจฺจ กิเลโส เจว สํกิลิฏฺโฐ จ อกุสโล ธมฺโม อุปฺปชฺชติ เหตุปจฺจยา.\csromanspacer{} ตีณิ. (สํขิตฺตํ.)}

\proseitem{25}{เหตุยา นว, อารมฺมเณ นว ฯเปฯ อวิคเต นว. (สํขิตฺตํ.\csromanspacer{} กิเลสทุก-อกุสลสทิสํ.\csromanspacer{} สหชาตวารมฺปิ ฯเปฯ ปญฺหาวารมฺปิ สพฺพตฺถ วิตฺถาเรตพฺพํ.)}

\csromanheader{2}{81. กิเลสกิเลสสมฺปยุตฺตทุก 1. กุสลตฺติก}
\csromanheader{2}{1. ปฏิจฺจวาร ปจฺจยจตุกฺก}
\proseitem{26}{กิเลสญฺเจว กิเลสสมฺปยุตฺตญฺจ อกุสลํ ธมฺมํ ปฏิจฺจ กิเลโส เจว กิเลสสมฺปยุตฺโต จ อกุสโล ธมฺโม อุปฺปชฺชติ เหตุปจฺจยา.\csromanspacer{} ตีณิ.}

\gambhira{ธมฺมานุโลม ทุกติกปฏฺฐานปาฬิ}
\prose{\csromanfolio{428}กิเลสสมฺปยุตฺตญฺเจว โน จ กิเลสํ อกุสลํ ธมฺมํ ปฏิจฺจ กิเลสสมฺปยุตฺโต เจว โน จ กิเลโส อกุสโล ธมฺโม อุปฺปชฺชติ เหตุปจฺจยา.\csromanspacer{} ตีณิ.}

\prose{กิเลสญฺเจว กิเลสสมฺปยุตฺตํ อกุสลญฺจ กิเลสสมฺปยุตฺตญฺเจว โน จ กิเลสํ อกุสลญฺจ ธมฺมํ ปฏิจฺจ กิเลโส เจว กิเลสสมฺปยุตฺโต จ อกุสโล ธมฺโม อุปฺปชฺชติ เหตุปจฺจยา.\csromanspacer{} ตีณิ. (สํขิตฺตํ.)}

\proseitem{27}{เหตุยา นว, อารมฺมเณ นว ฯเปฯ อวิคเต นว. (สํขิตฺตํ.\csromanspacer{} กิเลสทุก-อกุสลสทิสํ.\csromanspacer{} สหชาตวาเรปิ ฯเปฯ ปญฺหาวาเรปิ สพฺพตฺถ วิตฺถาเรตพฺพํ.)}

\csromanheader{2}{82. กิเลสวิปฺปยุตฺตสํกิเลสิกทุก 1. กุสลตฺติก}
\csromanheader{2}{1. ปฏิจฺจวาร ปจฺจยจตุกฺก}
\proseitem{28}{กิเลสวิปฺปยุตฺตํ สํกิเลสิกํ กุสลํ ธมฺมํ ปฏิจฺจ กิเลสวิปฺปยุตฺโต สํกิเลสิโก กุสโล ธมฺโม อุปฺปชฺชติ เหตุปจฺจยา. (1)}

\prose{กิเลสวิปฺปยุตฺตํ อสํกิเลสิกํ กุสลํ ธมฺมํ ปฏิจฺจ กิเลสวิปฺปยุตฺโต อสํกิเลสิโก กุสโล ธมฺโม อุปฺปชฺชติ เหตุปจฺจยา. (สํขิตฺตํ.)}

\proseitem{29}{เหตุยา เทฺว, อารมฺมเณ เทฺว ฯเปฯ อวิคเต เทฺว. (สํขิตฺตํ.\csromanspacer{} โลกิยทุกกุสลสทิสํ.\csromanspacer{} สหชาตวาเรปิ ฯเปฯ ปญฺหาวาเรปิ วิตฺถาเรตพฺพํ.)}

\proseitem{30}{กิเลสวิปฺปยุตฺตํ สํกิเลสิกํ อพฺยากตํ ธมฺมํ ปฏิจฺจ กิเลสวิปยุตฺโต สํกิเลสิโก อพฺยากโต ธมฺโม อุปฺปชฺชติ เหตุปจฺจยา. (1)}

\prose{กิเลสวิปฺปยุตฺตํ อสํกิเลสิกํ อพฺยากตํ ธมฺมํ ปฏิจฺจ กิเลสวิปฺปยุตฺโต อสํกิเลสิโก อพฺยากโต ธมฺโม อุปฺปชฺชติ เหตุปจฺจยา.\csromanspacer{} ตีณิ.}

\csromanmatikaanchor{mtk.83}
\csromantocmarkref{subparagraph}{83. ทสฺสเนนปหาตพฺพทุก 1. กุสลตฺติก ...}{mtk.83}
\bookmark[dest={mtk.83},level=5]{83. ทสฺสเนนปหาตพฺพทุก 1. กุสลตฺติก ...}
\csromanheader{6}{83. ทสฺสเนนปหาตพฺพทุก 1. กุสลตฺติก}
\prose{\csromanfolio{429}กิเลสวิปฺปยุตฺตํ สํกิเลสิกํ อพฺยากตญฺจ กิเลสวิปฺปยุตฺตํ อสํกิเลสิกํ อพฺยากตญฺจ ธมฺมํ ปฏิจฺจ กิเลสวิปฺปยุตฺโต สํกิเลสิโก อพฺยากโต ธมฺโม อุปฺปชฺชติ เหตุปจฺจยา. (1) (สํขิตฺตํ.)}

\proseitem{31}{เหตุยา ปญฺจ, อารมฺมเณ เทฺว ฯเปฯ อาเสวเน เอกํ ฯเปฯ อวิคเต ปญฺจ. (สํขิตฺตํ.\csromanspacer{} โลกิยทุก-อพฺยากตสทิสํ.\csromanspacer{} สหชาตวาเรปิ ฯเปฯ ปญฺหาวาเรปิ สพฺพตฺถ วิตฺถาเรตพฺพํ.)}

\nitthitam{กิเลสโคจฺฉกกุสลตฺติกํ นิฏฺฐิตํ.}
\csromanheader{2}{83. ทสฺสเนนปหาตพฺพทุก 1. กุสลตฺติก}
\csromanheader{2}{1-7. ปฏิจฺจาทิวาร ปจฺจยจตุกฺก}
\proseitem{1}{นทสฺสเนน ปหาตพฺพํ กุสลํ ธมฺมํ ปฏิจฺจ นทสฺสเนน ปหาตพฺโพ กุสโล ธมฺโม อุปฺปชฺชติ เหตุปจฺจยา. (สํขิตฺตํ.)}

\prose{เหตุยา เอกํ, อารมฺมเณ เอกํ, อธิปติยา เอกํ ฯเปฯ อวิคเต เอกํ. (สํขิตฺตํ.\csromanspacer{} สหชาตวาเรปิ ฯเปฯ ปญฺหาวาเรปิ สพฺพตฺถ เอกํ.)}

\proseitem{2}{ทสฺสเนน ปหาตพฺพํ อกุสลํ ธมฺมํ ปฏิจฺจ ทสฺสเนน ปหาตพฺโพ อกุสโล ธมฺโม อุปฺปชฺชติ เหตุปจฺจยา. (สํขิตฺตํ.)}

\prose{เหตุยา เทฺว, อารมฺมเณ เทฺว, อธิปติยา เทฺว ฯเปฯ อวิคเต เทฺว. (สํขิตฺตํ.)}

\proseitem{3}{ทสฺสเนน ปหาตพฺพํ อกุสลํ ธมฺมํ ปฏิจฺจ ทสฺสเนน ปหาตพฺโพ อกุสโล ธมฺโม อุปฺปชฺชติ นเหตุปจฺจยา.\csromanspacer{} วิจิกิจฺฉาสหคเต ขนฺเธ ปฏิจฺจ วิจิกิจฺฉาสหคโต โมโห. (1)}

\prose{นทสฺสเนน ปหาตพฺพํ อกุสลํ ธมฺมํ ปฏิจฺจ นทสฺสเนน ปหาตพฺโพ อกุสโล ธมฺโม อุปฺปชฺชติ นเหตุปจฺจยา. (1) (สํขิตฺตํ.)}

\proseitem{4}{นเหตุยา เทฺว, น-อธิปติยา เทฺว, นปุเรชาเต เทฺว ฯเปฯ นกมฺเม เทฺว, นวิปาเก เทฺว, นวิปฺปยุตฺเต เทฺว. (สํขิตฺตํ.\csromanspacer{} สหชาตวาเรปิ ฯเปฯ สมฺปยุตฺตวาเรปิ สพฺพตฺถ วิตฺถาเรตพฺพํ.)}

\gambhira{ธมฺมานุโลม ทุกติกปฏฺฐานปาฬิ}
\proseitem{5}{\csromanfolio{430}ทสฺสเนน ปหาตพฺโพ อกุสโล ธมฺโม ทสฺสเนน ปหาตพฺพสฺส อกุสลสฺส ธมฺมสฺส เหตุปจฺจเยน ปจฺจโย. (1)}

\prose{นทสฺสเนน ปหาตพฺโพ อกุสโล ธมฺโม นทสฺสเนน ปหาตพฺพสฺส อกุสลสฺส ธมฺมสฺส เหตุปจฺจเยน ปจฺจโย. (1) (สํขิตฺตํ.)}

\proseitem{6}{เหตุยา เทฺว, อารมฺมเณ ตีณิ. (ทสฺสเน เอกํ, นทสฺสเน เทฺว,) อธิปติยา ตีณิ. (ทสฺสเนน ปหาตพฺพมูลกํ เอกํ, นทสฺสเน เทฺว, อารมฺมณาธิปติ, สหชาตาธิปติ, เอการมฺมณาธิปติ.) อนนฺตเร เทฺว, (ทสฺสนมูลกํ เอกํ, นทสฺสเน เอกํ.) สมนนฺตเร เทฺว, สหชาเต เทฺว ฯเปฯ อุปนิสฺสเย ตีณิ, อาเสวเน เทฺว, กมฺเม เทฺว, อาหาเร เทฺว ฯเปฯ สมฺปยุตฺเต เทฺว, อตฺถิยา เทฺว, นตฺถิยา เทฺว ฯเปฯ อวิคเต เทฺว. (สํขิตฺตํ.)}

\proseitem{7}{นทสฺสเนน ปหาตพฺพํ อพฺยากตํ ธมฺมํ ปฏิจฺจ นทสฺสเนน ปหาตพฺโพ อพฺยากโต ธมฺโม อุปฺปชฺชติ เหตุปจฺจยา. (สํขิตฺตํ.)}

\prose{เหตุยา เอกํ, อารมฺมเณ เอกํ ฯเปฯ อวิคเต เอกํ. (สํขิตฺตํ.)}

\vspace{\tipitakaparskip}
\csromancenter{(สหชาตวาเรปิ ฯเปฯ ปญฺหาวาเรปิ สพฺพตฺถ เอกํ.)}
\csromanheader{2}{84. ภาวนายปหาตพฺพทุก 1. กุสลตฺติก}
\csromanheader{2}{1-7. ปฏิจฺจาทิวาร ปจฺจยจตุกฺก}
\proseitem{8}{นภาวนาย ปหาตพฺพํ กุสลํ ธมฺมํ ปฏิจฺจ นภาวนาย ปหาตพฺโพ กุสโล ธมฺโม อุปฺปชฺชติ เหตุปจฺจยา. (สํขิตฺตํ.)}

\prose{เหตุยา เอกํ, อารมฺมเณ เอกํ ฯเปฯ อวิคเต เอกํ. (สํขิตฺตํ.)}

\vspace{\tipitakaparskip}
\csromancenter{(สหชาตวาเรปิ ฯเปฯ ปญฺหาวาเรปิ สพฺพตฺถ เอกํ.)}
\proseitem{9}{ภาวนาย ปหาตพฺพํ อกุสลํ ธมฺมํ ปฏิจฺจ ภาวนาย ปหาตพฺโพ อกุสโล ธมฺโม อุปฺปชฺชติ เหตุปจฺจยา. (สํขิตฺตํ.)}

\prose{เหตุยา เทฺว, อารมฺมเณ เทฺว ฯเปฯ อวิคเต เทฺว. (สํขิตฺตํ.\csromanspacer{} ทสฺสเนน ปหาตพฺพทุก-อกุสลสทิสํ.\csromanspacer{} สหชาตวาเรปิ ฯเปฯ ปญฺหาวาเรปิ สพฺพตฺถ วิตฺถาเรตพฺพํ.)}

\csromanheader{2}{85. ทสฺสเนนปหาตพฺพเหตุกทุก 1. กุสลตฺติก}
\proseitem{10}{\csromanfolio{431}นภาวนาย ปหาตพฺพํ อพฺยากตํ ธมฺมํ ปฏิจฺจ นภาวนาย ปหาตพฺโพ อพฺยากโต ธมฺโม อุปฺปชฺชติ เหตุปจฺจยา. (สํขิตฺตํ.)}

\prose{เหตุยา เอกํ, อารมฺมเณ เอกํ ฯเปฯ อวิคเต เอกํ. (สํขิตฺตํ.)}

\vspace{\tipitakaparskip}
\csromancenter{(สหชาตวาเรปิ ฯเปฯ ปญฺหาวาเรปิ สพฺพตฺถ เอกํ.)}
\csromanheader{2}{85. ทสฺสเนนปหาตพฺพเหตุกทุก 1. กุสลตฺติก}
\csromanheader{2}{1-7. ปฏิจฺจาทิวาร ปจฺจยจตุกฺก}
\proseitem{11}{นทสฺสเนน ปหาตพฺพเหตุกํ กุสลํ ธมฺมํ ปฏิจฺจ นทสฺสเนน ปหาตพฺพเหตุโก กุสโล ธมฺโม อุปฺปชฺชติ เหตุปจฺจยา. (สํขิตฺตํ.)}

\prose{เหตุยา เอกํ, อารมฺมเณ เอกํ ฯเปฯ อวิคเต เอกํ. (สํขิตฺตํ.)}

\vspace{\tipitakaparskip}
\csromancenter{(สหชาตวาเรปิ ฯเปฯ ปญฺหาวาเรปิ สพฺพตฺถ เอกํ.)}
\proseitem{12}{ทสฺสเนน ปหาตพฺพเหตุกํ อกุสลํ ธมฺมํ ปฏิจฺจ ทสฺสเนน ปหาตพฺพเหตุโก อกุสโล ธมฺโม อุปฺปชฺชติ เหตุปจฺจยา. (1)}

\prose{นทสฺสเนน ปหาตพฺพเหตุกํ อกุสลํ ธมฺมํ ปฏิจฺจ นทสฺสเนน ปหาตพฺพเหตุโก อกุสโล ธมฺโม อุปฺปชฺชติ เหตุปจฺจยา.\csromanspacer{}\csromanspacer{}นทสฺสเนน ปหาตพฺพเหตุกํ อกุสลํ ธมฺมํ ปฏิจฺจ ทสฺสเนน ปหาตพฺพเหตุโก อกุสโล ธมฺโม อุปฺปชฺชติ เหตุปจฺจยา. (2)}

\prose{ทสฺสเนน ปหาตพฺพเหตุกํ อกุสลญฺจ นทสฺสเนน ปหาตพฺพเหตุกํ อกุสลญฺจ ธมฺมํ ปฏิจฺจ ทสฺสเนน ปหาตพฺพเหตุโก อกุสโล ธมฺโม อุปฺปชฺชติ เหตุปจฺจยา. (1)}

\prose{ทสฺสเนน ปหาตพฺพเหตุกํ อกุสลํ ธมฺมํ ปฏิจฺจ ทสฺสเนน ปหาตพฺพเหตุโก อกุสโล ธมฺโม อุปฺปชฺชติ อารมฺมณปจฺจยา.\csromanspacer{} ตีณิ. (นทสฺสเน เทฺว, ฆฏเน เอกํ.) (สํขิตฺตํ.)}

\proseitem{13}{เหตุยา จตฺตาริ, อารมฺมเณ ฉ, อธิปติยา เทฺว, (สพฺพตฺถ ฉ,) อวิคเต ฉ. (สํขิตฺตํ.)}

\proseitem{14}{ทสฺสเนน ปหาตพฺพเหตุกํ อกุสลํ ธมฺมํ ปฏิจฺจ นทสฺสเนน ปหาตพฺพเหตุโก อกุสโล ธมฺโม อุปฺปชฺชติ นเหตุปจฺจยา.}

\gambhira{ธมฺมานุโลม ทุกติกปฏฺฐานปาฬิ}
\prose{\csromanfolio{432}นทสฺสเนน ปหาตพฺพเหตุกํ อกุสลํ ธมฺมํ ปฏิจฺจ นทสฺสเนน ปหาตพฺพเหตุโก อกุสโล ธมฺโม อุปฺปชฺชติ นเหตุปจฺจยา. (สํขิตฺตํ.)}

\proseitem{15}{นเหตุยา เทฺว, น-อธิปติยา ฉ ฯเปฯ นกมฺเม จตฺตาริ, นวิปาเก ฉ, นวิปฺปยุตฺเต ฉ. (สํขิตฺตํ.\csromanspacer{} สหชาตวาราทิ วิตฺถาเรตพฺโพ.)}

\proseitem{16}{ทสฺสเนน ปหาตพฺพเหตุโก อกุสโล ธมฺโม ทสฺสเนน ปหาตพฺพเหตุกสฺส อกุสลสฺส ธมฺมสฺส เหตุปจฺจเยน ปจฺจโย. (สํขิตฺตํ.)}

\proseitem{17}{เหตุยา ตีณิ, อารมฺมเณ นว, อธิปติยา ตีณิ, (ทสฺสเน เอกํ, นทสฺสเน เทฺว.) อนนฺตเร นว, สมนนฺตเร นว, สหชาเต ฉ ฯเปฯ อุปนิสฺสเย นว, อาเสวเน นว, กมฺเม จตฺตาริ, อาหาเร จตฺตาริ, อินฺทฺริเย จตฺตาริ ฯเปฯ สมฺปยุตฺเต ฉ, อตฺถิยา ฉ, นตฺถิยา นว ฯเปฯ อวิคเต ฉ. (สํขิตฺตํ.)}

\proseitem{18}{นทสฺสเนน ปหาตพฺพเหตุกํ อพฺยากตํ ธมฺมํ ปฏิจฺจ นทสฺสเนน ปหาตพฺพเหตุโก อพฺยากโต ธมฺโม อุปฺปชฺชติ เหตุปจฺจยา. (สํขิตฺตํ.)}

\prose{เหตุยา เอกํ, อารมฺมเณ เอกํ ฯเปฯ อวิคเต เอกํ. (สํขิตฺตํ.)}

\vspace{\tipitakaparskip}
\csromancenter{(สหชาตวาเรปิ ฯเปฯ ปญฺหาวาเรปิ สพฺพตฺถ เอกํ.)}
\csromanheader{2}{86. ภาวนายปหาตพฺพเหตุกทุก 1. กุสลตฺติก}
\csromanheader{2}{1. ปฏิจฺจวาร ปจฺจยจตุกฺก}
\proseitem{19}{นภาวนาย ปหาตพฺพเหตุกํ กุสลํ ธมฺมํ ปฏิจฺจ นภาวนาย ปหาตพฺพเหตุโก กุสโล ธมฺโม อุปฺปชฺชติ เหตุปจฺจยา. (สํขิตฺตํ.)}

\prose{เหตุยา เอกํ, อารมฺมเณ เอกํ ฯเปฯ อวิคเต เอกํ. (สํขิตฺตํ.)}

\vspace{\tipitakaparskip}
\csromancenter{(สหชาตวาเรปิ ฯเปฯ ปญฺหาวาเรปิ สพฺพตฺถ เอกํ.)}
\proseitem{20}{ภาวนาย ปหาตพฺพเหตุกํ อกุสลํ ธมฺมํ ปฏิจฺจ ภาวนาย ปหาตพฺพเหตุโก อกุสโล ธมฺโม อุปฺปชฺชติ เหตุปจฺจยา. (1)}

\csromanheader{2}{87. สวิตกฺกทุก 1. กุสลตฺติก}
\prose{\csromanfolio{433}นภาวนาย ปหาตพฺพเหตุกํ อกุสลํ ธมฺมํ ปฏิจฺจ นภาวนาย ปหาตพฺพเหตุโก อกุสโล ธมฺโม อุปฺปชฺชติ เหตุปจฺจยา.\csromanspacer{} ตีณิ. (สํขิตฺตํ.)}

\proseitem{21}{เหตุยา จตฺตาริ, อารมฺมเณ ฉ, อธิปติยา เทฺว ฯเปฯ อวิคเต ฉ. (สํขิตฺตํ.\csromanspacer{} ทสฺสเนน ปหาตพฺพเหตุกทุก-อกุสลสทิสํ.\csromanspacer{} สหชาตวาเรปิ ฯเปฯ ปญฺหาวาเรปิ สพฺพตฺถ วิตฺถาเรตพฺพํ.)}

\proseitem{22}{นภาวนาย ปหาตพฺพเหตุกํ อพฺยากตํ ธมฺมํ ปฏิจฺจ นภาวนาย ปหาตพฺพเหตุโก อพฺยากโต ธมฺโม อุปฺปชฺชติ เหตุปจฺจยา. (สํขิตฺตํ.)}

\prose{เหตุยา เอกํ, อารมฺมเณ เอกํ ฯเปฯ อวิคเต เอกํ. (สํขิตฺตํ.)}

\vspace{\tipitakaparskip}
\csromancenter{(สหชาตวาเรปิ ฯเปฯ ปญฺหาวาเรปิ สพฺพตฺถ เอกํ.)}
\csromanheader{2}{87. สวิตกฺกทุก 1. กุสลตฺติก}
\csromanheader{2}{1. ปฏิจฺจวาร ปจฺจยจตุกฺก}
\proseitem{23}{สวิตกฺกํ กุสลํ ธมฺมํ ปฏิจฺจ สวิตกฺโก กุสโล ธมฺโม อุปฺปชฺชติ เหตุปจฺจยา.\csromanspacer{}\csromanspacer{}สวิตกฺกํ กุสลํ ธมฺมํ ปฏิจฺจ อวิตกฺโก กุสโล ธมฺโม อุปฺปชฺชติ เหตุปจฺจยา.\csromanspacer{}\csromanspacer{}สวิตกฺกํ กุสลํ ธมฺมํ ปฏิจฺจ สวิตกฺโก กุสโล จ อวิตกฺโก กุสโล จ ธมฺมา อุปฺปชฺชนฺติ เหตุปจฺจยา. (3)}

\prose{อวิตกฺกํ กุสลํ ธมฺมํ ปฏิจฺจ อวิตกฺโก กุสโล ธมฺโม อุปฺปชฺชติ เหตุปจฺจยา.\csromanspacer{}\csromanspacer{}อวิตกฺกํ กุสลํ ธมฺมํ ปฏิจฺจ สวิตกฺโก กุสโล ธมฺโม อุปฺปชฺชติ เหตุปจฺจยา. (2)}

\prose{สวิตกฺกํ กุสลญฺจ อวิตกฺกํ กุสลญฺจ ธมฺมํ ปฏิจฺจ สวิตกฺโก กุสโล ธมฺโม อุปฺปชฺชติ เหตุปจฺจยา. (1) (สํขิตฺตํ.)}

\proseitem{24}{เหตุยา ฉ, อารมฺมเณ ฉ, (สพฺพตฺถ ฉ.) อวิคเต ฉ. (สํขิตฺตํ.)}

\prose{น-อธิปติยา ฉ, นปุเรชาเต ฉ ฯเปฯ นกมฺเม จตฺตาริ ฯเปฯ นวิปฺปยุตฺเต ฉ. (สํขิตฺตํ.\csromanspacer{} สหชาตวารมฺปิ ฯเปฯ สมฺปยุตฺตวารมฺปิ วิตฺถาเรตพฺพํ.)}

\csromanheader{2}{ธมฺมานุโลม ทุกติกปฏฺฐานปาฬิ \textbf{ปญฺหาวาร-เหตุ อารมฺมณ}}
\proseitem{25}{\csromanfolio{434}สวิตกฺโก กุสโล ธมฺโม สวิตกฺกสฺส กุสลสฺส ธมฺมสฺส เหตุปจฺจเยน ปจฺจโย.\csromanspacer{}\csromanspacer{}สวิตกฺโก กุสโล ธมฺโม อวิตกฺกสฺส กุสลสฺส ธมฺมสฺส เหตุปจฺจเยน ปจฺจโย.\csromanspacer{}\csromanspacer{}สวิตกฺโก กุสโล ธมฺโม สวิตกฺกสฺส กุสลสฺส จ อวิตกฺกสฺส กุสลสฺส จ ธมฺมสฺส เหตุปจฺจเยน ปจฺจโย. (3)}

\prose{อวิตกฺโก กุสโล ธมฺโม อวิตกฺกสฺส กุสลสฺส ธมฺมสฺส เหตุปจฺจเยน ปจฺจโย. (1)}

\prose{สวิตกฺโก กุสโล ธมฺโม สวิตกฺกสฺส กุสลสฺส ธมฺมสฺส อารมฺมณปจฺจเยน ปจฺจโย.\csromanspacer{} ตีณิ.}

\prose{อวิตกฺโก กุสโล ธมฺโม อวิตกฺกสฺส กุสลสฺส ธมฺมสฺส อารมฺมณปจฺจเยน ปจฺจโย.\csromanspacer{} ตีณิ.}

\prose{สวิตกฺโก กุสโล จ อวิตกฺโก กุสโล จ ธมฺมา สวิตกฺกสฺส กุสลสฺส ธมฺมสฺส อารมฺมณปจฺจเยน ปจฺจโย.\csromanspacer{} ตีณิ. (สํขิตฺตํ.)}

\proseitem{26}{เหตุยา จตฺตาริ, อารมฺมเณ นว, อธิปติยา นว. (เหฏฺฐา ตีสุ สหชาตาธิปติ, อวิตกฺเก เอกํ สหชาตาธิปติ.) อนนฺตเร นว, สมนนฺตเร นว, สหชาเต ฉ ฯเปฯ อุปนิสฺสเย นว, อาเสวเน นว, กมฺเม จตฺตาริ, อาหาเร จตฺตาริ, อินฺทฺริเย จตฺตาริ, ฌาเน ฉ, มคฺเค ฉ, สมฺปยุตฺเต ฉ, อตฺถิยา ฉ, นตฺถิยา นว, วิคเต นว, อวิคเต ฉ. (สํขิตฺตํ.)}

\prose{นเหตุยา นว, น-อารมฺมเณ นว. (สํขิตฺตํ.)}

\csromanheader{2}{เหตุปจฺจยา น-อารมฺมเณ จตฺตาริ. (สํขิตฺตํ.)}
\prose{นเหตุปจฺจยา อารมฺมเณ นว. (สํขิตฺตํ.) (ยถา กุสลตฺติเก ปญฺหาวารํ, เอวํ วิตฺถาเรตพฺพํ.)}

\csromanheader{2}{\textbf{อกุสลปท}}
\proseitem{27}{สวิตกฺกํ อกุสลํ ธมฺมํ ปฏิจฺจ สวิตกฺโก อกุสโล ธมฺโม อุปฺปชฺชติ เหตุปจฺจยา.\csromanspacer{}\csromanspacer{}สวิตกฺกํ อกุสลํ ธมฺมํ ปฏิจฺจ อวิตกฺโก}

\vspace{\tipitakaparskip}
\csromancenter{สวิตกฺกทุก 1.\csromanspacer{} กุสลตฺติก อกุสโล ธมฺโม อุปฺปชฺชติ เหตุปจฺจยา.\csromanspacer{}\csromanspacer{}สวิตกฺกํ อกุสลํ ธมฺมํ ปฏิจฺจ สวิตกฺโก อกุสโล จ อวิตกฺโก อกุสโล จ ธมฺมา อุปฺปชฺชนฺติ เหตุปจฺจยา. (3)}
\prose{\csromanfolio{435}อวิตกฺกํ อกุสลํ ธมฺมํ ปฏิจฺจ สวิตกฺโก อกุสโล ธมฺโม อุปฺปชฺชติ เหตุปจฺจยา. (1)}

\prose{สวิตกฺกํ อกุสลญฺจ อวิตกฺกํ อกุสลญฺจ ธมฺมํ ปฏิจฺจ สวิตกฺโก อกุสโล ธมฺโม อุปฺปชฺชติ เหตุปจฺจยา. (1) (สํขิตฺตํ.)}

\proseitem{28}{เหตุยา ปญฺจ, อารมฺมเณ ปญฺจ, (สพฺพตฺถ ปญฺจ.) ฯเปฯ อวิคเต ปญฺจ. (สํขิตฺตํ.)}

\proseitem{29}{สวิตกฺกํ อกุสลํ ธมฺมํ ปฏิจฺจ สวิตกฺโก อกุสโล ธมฺโม อุปฺปชฺชติ นเหตุปจฺจยา. (1)}

\prose{อวิตกฺกํ อกุสลํ ธมฺมํ ปฏิจฺจ สวิตกฺโก อกุสโล ธมฺโม อุปฺปชฺชติ นเหตุปจฺจยา. (1)}

\prose{สวิตกฺกํ อกุสลญฺจ อวิตกฺกํ อกุสลญฺจ ธมฺมํ ปฏิจฺจ สวิตกฺโก อกุสโล ธมฺโม อุปฺปชฺชติ นเหตุปจฺจยา. (1) (สํขิตฺตํ.)}

\proseitem{30}{นเหตุยา ตีณิ, น-อธิปติยา ปญฺจ, นกมฺเม ตีณิ ฯเปฯ นวิปฺปยุตฺเต ปญฺจ. (สํขิตฺตํ.\csromanspacer{} สหชาตวาราทิ วิตฺถาเรตพฺโพ.)}

\proseitem{31}{สวิตกฺโก อกุสโล ธมฺโม สวิตกฺกสฺส อกุสลสฺส ธมฺมสฺส เหตุปจฺจเยน ปจฺจโย. (สํขิตฺตํ.)}

\proseitem{32}{เหตุยา ตีณิ, อารมฺมเณ นว, อธิปติยา นว, อนนฺตเร นว ฯเปฯ สหชาเต ปญฺจ ฯเปฯ อุปนิสฺสเย นว, อาเสวเน นว, กมฺเม ตีณิ, อาหาเร ตีณิ ฯเปฯ ฌาเน ปญฺจ, มคฺเค ปญฺจ, สมฺปยุตฺเต ปญฺจ, อตฺถิยา ปญฺจ, นตฺถิยา นว ฯเปฯ อวิคเต ปญฺจ. (สํขิตฺตํ.)}

\csromanheader{2}{\textbf{อพฺยากตปท}}
\proseitem{33}{สวิตกฺกํ อพฺยากตํ ธมฺมํ ปฏิจฺจ สวิตกฺโก อพฺยากโต ธมฺโม อุปฺปชฺชติ เหตุปจฺจยา.\csromanspacer{} ตีณิ.}

\gambhira{ธมฺมานุโลม ทุกติกปฏฺฐานปาฬิ}
\prose{\csromanfolio{436}อวิตกฺกํ อพฺยากตํ ธมฺมํ ปฏิจฺจ อวิตกฺโก อพฺยากโต ธมฺโม อุปฺปชฺชติ เหตุปจฺจยา.\csromanspacer{} ตีณิ.}

\prose{สวิตกฺกํ อพฺยากตญฺจ อวิตกฺกํ อพฺยากตญฺจ ธมฺมํ ปฏิจฺจ สวิตกฺโก อพฺยากโต ธมฺโม อุปฺปชฺชติ เหตุปจฺจยา.\csromanspacer{} ตีณิ. (สํขิตฺตํ.)}

\proseitem{34}{เหตุยา นว, อารมฺมเณ นว ฯเปฯ ปุเรชาเต อาเสวเน ฉ ฯเปฯ อวิคเต นว. (สํขิตฺตํ.)}

\prose{นเหตุยา นว, น-อารมฺมเณ ตีณิ, น-อธิปติยา นว ฯเปฯ นปุเรชาเต นว ฯเปฯ นกมฺเม จตฺตาริ, นวิปาเก นว, น-อาหาเร เอกํ, น-อินฺทฺริเย เอกํ, นฌาเน เอกํ, นมคฺเค นว, นสมฺปยุตฺเต ตีณิ, นวิปฺปยุตฺเต ฉ ฯเปฯ โนวิคเต ตีณิ. (สํขิตฺตํ.\csromanspacer{} สหชาตวารมฺปิ ฯเปฯ สมฺปยุตฺตวารมฺปิ วิตฺถาเรตพฺพํ.)}

\proseitem{35}{สวิตกฺโก อพฺยากโต ธมฺโม สวิตกฺกสฺส อพฺยากตสฺส ธมฺมสฺส เหตุปจฺจเยน ปจฺจโย.\csromanspacer{} ตีณิ.}

\prose{อวิตกฺโก อพฺยากโต ธมฺโม อวิตกฺกสฺส อพฺยากตสฺส ธมฺมสฺส เหตุปจฺจเยน ปจฺจโย. (1)}

\prose{สวิตกฺโก อพฺยากโต ธมฺโม สวิตกฺกสฺส อพฺยากตสฺส ธมฺมสฺส อารมฺมณปจฺจเยน ปจฺจโย. (สํขิตฺตํ.)}

\proseitem{36}{เหตุยา จตฺตาริ, อารมฺมเณ นว, อธิปติยา นว, (สพฺพตฺถ นว.) อุปนิสฺสเย นว, ปุเรชาเต ตีณิ, ปจฺฉาชาเต ตีณิ, อาเสวเน นว, กมฺเม จตฺตาริ, วิปาเก นว, อาหาเร จตฺตาริ ฯเปฯ ฌาเน นว, มคฺเค นว, สมฺปยุตฺเต ฉ, วิปฺปยุตฺเต ปญฺจ, อตฺถิยา นว, นตฺถิยา นว, วิคเต นว, อวิคเต นว. (สํขิตฺตํ.)}

\prose{นเหตุยา นว, น-อารมฺมเณ นว. (สํขิตฺตํ.)}

\csromanheader{2}{เหตุปจฺจยา น-อารมฺมเณ จตฺตาริ. (สํขิตฺตํ.)}
\csromanheader{2}{นเหตุปจฺจยา อารมฺมเณ นว. (สํขิตฺตํ.)}
\prose{(ยถา กุสลตฺติเก ปญฺหาวารํ, เอวํ วิตฺถาเรตพฺพํ.)}

\csromanheader{2}{88. สวิจารทุก 1. กุสลตฺติก \textbf{88. สวิจารทุก 1. กุสลตฺติก}}
\csromanheader{2}{1-7. ปฏิจฺจาทิวาร ปจฺจยจตุกฺก}
\proseitem{37}{\csromanfolio{437}สวิจารํ กุสลํ ธมฺมํ ปฏิจฺจ สวิจาโร กุสโล ธมฺโม อุปฺปชฺชติ เหตุปจฺจยา.\csromanspacer{} ตีณิ.}

\prose{อวิจารํ กุสลํ ธมฺมํ ปฏิจฺจ อวิจาโร กุสโล ธมฺโม อุปฺปชฺชติ เหตุปจฺจยา.\csromanspacer{} เทฺว.}

\prose{สวิจารํ กุสลญฺจ อวิจารํ กุสลญฺจ ธมฺมํ ปฏิจฺจ สวิจาโร กุสโล ธมฺโม อุปฺปชฺชติ เหตุปจฺจยา. (1) (สํขิตฺตํ.)}

\proseitem{38}{เหตุยา ฉ, อารมฺมเณ ฉ ฯเปฯ อวิคเต ฉ. (สํขิตฺตํ.\csromanspacer{} สวิตกฺกทุกกุสลสทิสํ.\csromanspacer{} สหชาตวารมฺปิ ฯเปฯ ปญฺหาวารมฺปิ วิตฺถาเรตพฺพํ.)}

\proseitem{39}{สวิจารํ อกุสลํ ธมฺมํ ปฏิจฺจ สวิจาโร อกุสโล ธมฺโม อุปฺปชฺชติ เหตุปจฺจยา.\csromanspacer{} ตีณิ.}

\prose{อวิจารํ อกุสลํ ธมฺมํ ปฏิจฺจ สวิจาโร อกุสโล ธมฺโม อุปฺปชฺชติ เหตุปจฺจยา. (1)}

\prose{สวิจารํ อกุสลญฺจ อวิจารํ อกุสลญฺจ ธมฺมํ ปฏิจฺจ สวิจาโร อกุสโล ธมฺโม อุปฺปชฺชติ เหตุปจฺจยา. (1) (สํขิตฺตํ.)}

\proseitem{40}{เหตุยา ปญฺจ, อารมฺมเณ ปญฺจ ฯเปฯ อวิคเต ปญฺจ. (สํขิตฺตํ.\csromanspacer{} สวิตกฺกทุก-อกุสลสทิสํ.\csromanspacer{} สหชาตวารมฺปิ ฯเปฯ ปญฺหาวารมฺปิ วิตฺถาเรตพฺพํ.)}

\proseitem{41}{สวิจารํ อพฺยากตํ ธมฺมํ ปฏิจฺจ สวิจาโร อพฺยากโต ธมฺโม อุปฺปชฺชติ เหตุปจฺจยา.\csromanspacer{} ตีณิ.}

\prose{อวิจารํ อพฺยากตํ ธมฺมํ ปฏิจฺจ อวิจาโร อพฺยากโต ธมฺโม อุปฺปชฺชติ เหตุปจฺจยา.\csromanspacer{} ตีณิ.}

\prose{สวิจารํ อพฺยากตญฺจ อวิจารํ อพฺยากตญฺจ ธมฺมํ ปฏิจฺจ สวิจาโร อพฺยากโต ธมฺโม อุปฺปชฺชติ เหตุปจฺจยา.\csromanspacer{} ตีณิ. (สํขิตฺตํ.)}

\proseitem{42}{เหตุยา นว, อารมฺมเณ นว ฯเปฯ วิปาเก นว ฯเปฯ อวิคเต นว. (สํขิตฺตํ.\csromanspacer{} สวิตกฺกทุก-อพฺยากตสทิสํ.\csromanspacer{} สหชาตวาโรปิ ฯเปฯ สมฺปยุตฺตวาโรปิ วิตฺถาเรตพฺพา.)}

\gambhira{ธมฺมานุโลม ทุกติกปฏฺฐานปาฬิ}
\proseitem{43}{\csromanfolio{438}สวิจาโร อพฺยากโต ธมฺโม สวิจารสฺส อพฺยากตสฺส ธมฺมสฺส เหตุปจฺจเยน ปจฺจโย. (สํขิตฺตํ.)}

\proseitem{44}{เหตุยา จตฺตาริ, อารมฺมเณ นว ฯเปฯ มคฺเค จตฺตาริ ฯเปฯ อวิคเต นว. (สํขิตฺตํ.)}

\prose{นเหตุยา นว, น-อารมฺมเณ นว. (สํขิตฺตํ.)}

\csromanheader{2}{เหตุปจฺจยา น-อารมฺมเณ จตฺตาริ. (สํขิตฺตํ.)}
\csromanheader{2}{นเหตุปจฺจยา อารมฺมเณ นว. (สํขิตฺตํ.)}
\prose{(ยถา กุสลตฺติเก ปญฺหาวารํ, เอวํ วิตฺถาเรตพฺพํ.)}

\csromanheader{2}{89. สปฺปีติกทุก 1. กุสลตฺติก}
\csromanheader{2}{1-7. ปฏิจฺจาทิวาร ปจฺจยจตุกฺก}
\proseitem{45}{สปฺปีติกํ กุสลํ ธมฺมํ ปฏิจฺจ สปฺปีติโก กุสโล ธมฺโม อุปฺปชฺชติ เหตุปจฺจยา.\csromanspacer{}\csromanspacer{}สปฺปีติกํ กุสลํ ธมฺมํ ปฏิจฺจ อปฺปีติโก กุสโล ธมฺโม อุปฺปชฺชติ เหตุปจฺจยา.\csromanspacer{}\csromanspacer{}สปฺปีติกํ กุสลํ ธมฺมํ ปฏิจฺจ สปฺปีติโก กุสโล จ อปฺปีติโก กุสโล จ ธมฺมา อุปฺปชฺชนฺติ เหตุปจฺจยา. (3)}

\prose{อปฺปีติกํ กุสลํ ธมฺมํ ปฏิจฺจ อปฺปีติโก กุสโล ธมฺโม อุปฺปชฺชติ เหตุปจฺจยา.\csromanspacer{}\csromanspacer{}อปฺปีติกํ กุสลํ ธมฺมํ ปฏิจฺจ สปฺปีติโก กุสโล ธมฺโม อุปฺปชฺชติ เหตุปจฺจยา. (2)}

\prose{สปฺปีติกํ กุสลญฺจ อปฺปีติกํ กุสลญฺจ ธมฺมํ ปฏิจฺจ สปฺปีติโก กุสโล ธมฺโม อุปฺปชฺชติ เหตุปจฺจยา. (1) (สํขิตฺตํ.)}

\proseitem{46}{เหตุยา ฉ, (สพฺพตฺถ ฉ.) ฯเปฯ อวิคเต ฉ. (สํขิตฺตํ.)}

\prose{น-อธิปติยา ฉ, นปุเรชาเต ฉ ฯเปฯ นกมฺเม จตฺตาริ, นวิปฺปยุตฺเต ฉ. (สํขิตฺตํ.\csromanspacer{} สหชาตวารมฺปิ ฯเปฯ สมฺปยุตฺตวารมฺปิ วิตฺถาเรตพฺพํ.)}

\proseitem{47}{สปฺปีติโก กุสโล ธมฺโม สปฺปีติกสฺส กุสลสฺส ธมฺมสฺส เหตุปจฺจเยน ปจฺจโย.\csromanspacer{} ตีณิ.}

\csromanheader{2}{89. สปฺปีติกทุก 1. กุสลตฺติก}
\prose{\csromanfolio{439}อปฺปีติโก กุสโล ธมฺโม อปฺปีติกสฺส กุสลสฺส ธมฺมสฺส เหตุปจฺจเยน ปจฺจโย. (1)}

\prose{สปฺปีติโก กุสโล ธมฺโม สปฺปีติกสฺส กุสลสฺส ธมฺมสฺส อารมฺมณปจฺจเยน ปจฺจโย. (สํขิตฺตํ.)}

\proseitem{48}{เหตุยา จตฺตาริ, อารมฺมเณ นว, อธิปติยา นว, (จตฺตาริ สหชาตาธิปติ.) อนนฺตเร นว, สมนนฺตเร นว, สหชาเต ฉ ฯเปฯ อุปนิสฺสเย นว, อาเสวเน นว, กมฺเม จตฺตาริ, อาหาเร จตฺตาริ, อินฺทฺริเย จตฺตาริ, ฌาเน ฉ, มคฺเค จตฺตาริ, สมฺปยุตฺเต ฉ, อตฺถิยา ฉ, นตฺถิยา นว, วิคเต นว, อวิคเต ฉ. (สํขิตฺตํ.)}

\prose{นเหตุยา นว, น-อารมฺมเณ นว. (สํขิตฺตํ.)}

\csromanheader{2}{เหตุปจฺจยา น-อารมฺมเณ จตฺตาริ. (สํขิตฺตํ.)}
\prose{นเหตุปจฺจยา อารมฺมเณ นว. (สํขิตฺตํ.) (ยถา กุสลตฺติเก ปญฺหาวารํ, เอวํ วิตฺถาเรตพฺพํ.)}

\csromanheader{2}{\textbf{อกุสลปท-เหตุ}}
\proseitem{49}{สปฺปีติกํ อกุสลํ ธมฺมํ ปฏิจฺจ สปฺปีติโก อกุสโล ธมฺโม อุปฺปชฺชติ เหตุปจฺจยา.\csromanspacer{} ตีณิ.}

\prose{อปฺปีติกํ อกุสลํ ธมฺมํ ปฏิจฺจ อปฺปีติโก อกุสโล ธมฺโม อุปฺปชฺชติ เหตุปจฺจยา.\csromanspacer{} เทฺว.}

\prose{สปฺปีติกํ อกุสลญฺจ อปฺปีติกํ อกุสลญฺจ ธมฺมํ ปฏิจฺจ สปฺปีติโก อกุสโล ธมฺโม อุปฺปชฺชติ เหตุปจฺจยา. (1) (สํขิตฺตํ.)}

\proseitem{50}{เหตุยา ฉ, อารมฺมเณ ฉ ฯเปฯ อวิคเต ฉ. (สํขิตฺตํ.\csromanspacer{} กุสลสทิสํ.\csromanspacer{} สหชาตวาโรปิ ฯเปฯ ปญฺหาวาโรปิ วิตฺถาเรตพฺพา.)}

\csromanheader{2}{\textbf{อพฺยากตปท-เหตุ}}
\proseitem{51}{สปฺปีติกํ อพฺยากตํ ธมฺมํ ปฏิจฺจ สปฺปีติโก อพฺยากโต ธมฺโม อุปฺปชฺชติ เหตุปจฺจยา.\csromanspacer{} ตีณิ.}

\prose{อปฺปีติกํ อพฺยากตํ ธมฺมํ ปฏิจฺจ อปฺปีติโก อพฺยากโต ธมฺโม อุปฺปชฺชติ เหตุปจฺจยา.\csromanspacer{} ตีณิ.}

\gambhira{ธมฺมานุโลม ทุกติกปฏฺฐานปาฬิ}
\prose{\csromanfolio{440}สปฺปีติกํ อพฺยากตญฺจ อปฺปีติกํ อพฺยากตญฺจ ธมฺมํ ปฏิจฺจ สปฺปีติโก อพฺยากโต ธมฺโม อุปฺปชฺชติ เหตุปจฺจยา.\csromanspacer{} ตีณิ. (สํขิตฺตํ.)}

\proseitem{52}{เหตุยา นว, อารมฺมเณ นว ฯเปฯ ปุเรชาเต ฉ, อาเสวเน ฉ ฯเปฯ อวิคเต นว. (สํขิตฺตํ.)}

\prose{นเหตุยา นว, น-อารมฺมเณ ตีณิ, น-อธิปติยา นว ฯเปฯ นกมฺเม จตฺตาริ, นวิปาเก นว, น-อาหาเร เอกํ, น-อินฺทฺริเย เอกํ, นฌาเน เอกํ, นมคฺเค นว, นสมฺปยุตฺเต ตีณิ, นวิปฺปยุตฺเต ฉ. (สํขิตฺตํ.\csromanspacer{} สหชาตวาราทิ วิตฺถาเรตพฺโพ.)}

\proseitem{53}{สปฺปีติโก อพฺยากโต ธมฺโม สปฺปีติกสฺส อพฺยากตสฺส ธมฺมสฺส เหตุปจฺจเยน ปจฺจโย. (สํขิตฺตํ.)}

\proseitem{54}{เหตุยา จตฺตาริ, อารมฺมเณ นว, อธิปติยา นว, อนนฺตเร นว ฯเปฯ (สพฺพตฺถ นว.) อุปนิสฺสเย นว, ปุเรชาเต ตีณิ, ปจฺฉาชาเต ตีณิ, อาเสวเน นว, กมฺเม จตฺตาริ, วิปาเก นว, อาหาเร จตฺตาริ, อินฺทฺริเย จตฺตาริ, ฌาเน นว, มคฺเค จตฺตาริ, สมฺปยุตฺเต ฉ, วิปฺปยุตฺเต ปญฺจ, อตฺถิยา นว. (สํขิตฺตํ.)}

\csromanheader{2}{90-91. ปีติสหคตาทิทุก 1. กุสลตฺติก}
\csromanheader{2}{1. ปฏิจฺจวาร ปจฺจยจตุกฺก}
\proseitem{55}{ปีติสหคตํ กุสลํ ธมฺมํ ปฏิจฺจ ฯเปฯ. (กุสลมฺปิ อกุสลมฺปิ อพฺยากตมฺปิ สปฺปีติกทุกสทิสํ.)}

\proseitem{56}{สุขสหคตํ กุสลํ ธมฺมํ ปฏิจฺจ ฯเปฯ. (กุสลมฺปิ อกุสลมฺปิ อพฺยากตมฺปิ สปฺปีติกทุกสทิสํ.\csromanspacer{} อกุสลํ ธมฺมํ ปฏิจฺจ ฯเปฯ.\csromanspacer{} ปจฺจนีเย, นเหตุยา เอกํ.\csromanspacer{}\csromanspacer{}อพฺยากตํ ธมฺมํ ปฏิจฺจ ฯเปฯ.\csromanspacer{} ปจฺจนีเย, นเหตุยา นว ฯเปฯ นฌาเน ฉ, กาตพฺพา.\csromanspacer{}\csromanspacer{}ปจฺจนีเย ปญฺหาวาเร กุสลากุสเล อินฺทฺริเย ฌาเน ฉ, ปญฺหาวาเร อพฺยากเต นว.)}

\vspace{\tipitakaparskip}
\csromancenter{กามาวจรทุก 1.\csromanspacer{} กุสลตฺติก \textbf{92.}\csromanspacer{}\textbf{ อุเปกฺขาสหคตทุก 1.}\csromanspacer{}\textbf{ กุสลตฺติก}}
\csromanheader{2}{1. ปฏิจฺจวาร ปจฺจยจตุกฺก}
\proseitem{57}{\csromanfolio{441}อุเปกฺขาสหคตํ กุสลํ ธมฺมํ ปฏิจฺจ ฯเปฯ.\csromanspacer{} ฉ. (สปฺปีติกทุกสทิสํ. “อุเปกฺขา”ติ นานา-อุเปกฺขา อพฺยากเต.\csromanspacer{} ปจฺจนีเย นเหตุยา นว, นฌาเน ฉ.)}

\proseitem{58}{อุเปกฺขาสหคตํ อกุสลํ ธมฺมํ ปฏิจฺจ อุเปกฺขาสหคโต อกุสโล ธมฺโม อุปฺปชฺชติ เหตุปจฺจยา. (สํขิตฺตํ.)}

\csromanheader{2}{เหตุยา ฉ. (สพฺพตฺถ ฉ. สํขิตฺตํ.)}
\proseitem{59}{อุเปกฺขาสหคตํ อกุสลํ ธมฺมํ ปฏิจฺจ อุเปกฺขาสหคโต อกุสโล ธมฺโม อุปชฺชติ นเหตุปจฺจยา. (1)}

\prose{น-อุเปกฺขาสหคตํ อกุสลํ ธมฺมํ ปฏิจฺจ อุเปกฺขาสหคโต อกุสโล ธมฺโม อุปฺปชฺชติ นเหตุปจฺจยา. (1)}

\prose{อุเปกฺขาสหคตญฺจ น-อุเปกฺขาสหคตญฺจ ธมฺมํ ปฏิจฺจ อุเปกฺขาสหคโต อกุสโล ธมฺโม อุปฺปชฺชติ นเหตุปจฺจยา. (1) (สํขิตฺตํ.)}

\proseitem{60}{นเหตุยา ตีณิ, น-อธิปติยา ฉ ฯเปฯ นวิปฺปยุตฺเต ฉ. (สํขิตฺตํ.\csromanspacer{} เอวํ สพฺพตฺถ อกุสลํ วิตฺถาเรตพฺพํ.\csromanspacer{} สปฺปีติกทุกสทิสํ.)}

\proseitem{61}{อุเปกฺขาสหคตํ อพฺยากตํ ธมฺมํ ปฏิจฺจ อุเปกฺขาสหคโต อพฺยากโต ธมฺโม อุปฺปชฺชติ เหตุปจฺจยา นว ปญฺหา. (สปฺปีติกทุก-อพฺยากตสทิสํ.\csromanspacer{} ปญฺหาวาเร กุสลากุสเล อินฺทฺริเย ฌาเน ฉ, อพฺยากเต นว.)}

\csromanheader{2}{93. กามาวจรทุก 1. กุสลตฺติก}
\csromanheader{2}{1-7. ปฏิจฺจาทิวาร ปจฺจยจตุกฺก}
\proseitem{62}{กามาวจรํ กุสลํ ธมฺมํ ปฏิจฺจ กามาวจโร กุสโล ธมฺโม อุปฺปชฺชติ เหตุปจฺจยา. (1)}

\prose{นกามาวจรํ กุสลํ ธมฺมํ ปฏิจฺจ นกามาวจโร กุสโล ธมฺโม อุปฺปชฺชติ เหตุปจฺจยา. (1) (สํขิตฺตํ.)}

\gambhira{ธมฺมานุโลม ทุกติกปฏฺฐานปาฬิ}
\proseitem{63}{\csromanfolio{442}เหตุยา เทฺว, อารมฺมเณ เทฺว, (สพฺพตฺถ เทฺว.) อวิคเต เทฺว. (สํขิตฺตํ.)}

\prose{น-อธิปติยา เทฺว ฯเปฯ นวิปฺปยุตฺเต เทฺว. (สํขิตฺตํ.) (สหชาตวาราทิ วิตฺถาเรตพฺโพ.)}

\proseitem{64}{กามาวจโร กุสโล ธมฺโม กามาวจรสฺส กุสลสฺส ธมฺมสฺส เหตุปจฺจเยน ปจฺจโย. (1)}

\prose{นกามาวจโร กุสโล ธมฺโม นกามาวจรสฺส กุสลสฺส ธมฺมสฺส เหตุปจฺจเยน ปจฺจโย. (1) (สํขิตฺตํ.)}

\proseitem{65}{เหตุยา เทฺว, อารมฺมเณ จตฺตาริ, อธิปติยา ตีณิ, (กามาวจเร เอกํ, นกามาวจเร เทฺว.) อนนฺตเร ตีณิ, (กามาวจเร เทฺว, นกามาวจเร เอกํ.) ฯเปฯ สหชาเต เทฺว ฯเปฯ อุปนิสฺสเย จตฺตาริ, อาเสวเน ตีณิ, กมฺเม เทฺว, อาหาเร เทฺว ฯเปฯ นตฺถิยา ตีณิ ฯเปฯ อวิคเต เทฺว. (สํขิตฺตํ.)}

\proseitem{66}{กามาวจรํ อกุสลํ ธมฺมํ ปฏิจฺจ กามาวจโร อกุสโล ธมฺโม อุปฺปชฺชติ เหตุปจฺจยา. (สํขิตฺตํ.)}

\prose{เหตุยา เอกํ, อารมฺมเณ เอกํ ฯเปฯ อวิคเต เอกํ. (สํขิตฺตํ.) (สหชาตวาเรปิ ฯเปฯ ปญฺหาวาเรปิ สพฺพตฺถ เอกํ.)}

\proseitem{67}{กามาวจรํ อพฺยากตํ ธมฺมํ ปฏิจฺจ กามาวจโร อพฺยากโต ธมฺโม อุปฺปชฺชติ เหตุปจฺจยา.\csromanspacer{} ตีณิ.}

\prose{นกามาวจรํ อพฺยากตํ ธมฺมํ ปฏิจฺจ นกามาวจโร อพฺยากโต ธมฺโม อุปฺปชฺชติ เหตุปจฺจยา.\csromanspacer{} ตีณิ.}

\prose{กามาวจรํ อพฺยากตญฺจ นกามาวจรํ อพฺยากตญฺจ ธมฺมํ ปฏิจฺจ กามาวจโร อพฺยากโต ธมฺโม อุปฺปชฺชติ เหตุปจฺจยา.\csromanspacer{} ตีณิ. (สํขิตฺตํ.)}

\proseitem{68}{เหตุยา นว, อารมฺมเณ จตฺตาริ, อธิปติยา ปญฺจ ฯเปฯ อญฺญมญฺเญ ฉ ฯเปฯ ปุเรชาเต เทฺว, อาเสวเน เทฺว ฯเปฯ อวิคเต นว. (สํขิตฺตํ.)}

\csromanheader{2}{94. รูปาวจรทุก 1. กุสลตฺติก}
\prose{\csromanfolio{443}นเหตุยา เอกํ, น-อารมฺมเณ ตีณิ, น-อธิปติยา นว ฯเปฯ นกมฺเม เทฺว, นวิปาเก ปญฺจ, น-อาหาเร เอกํ, น-อินฺทฺริเย เอกํ, นฌาเน เอกํ, นมคฺเค เอกํ, นสมฺปยุตฺเต ตีณิ, นวิปฺปยุตฺเต เทฺว, โนนตฺถิยา ตีณิ, โนวิคเต ตีณิ. (สํขิตฺตํ.\csromanspacer{} สหชาตวาราทิ วิตฺถาเรตพฺโพ.)}

\proseitem{69}{กามาวจโร อพฺยากโต ธมฺโม กามาวจรสฺส อพฺยากตสฺส ธมฺมสฺส เหตุปจฺจเยน ปจฺจโย. (1)}

\prose{นกามาวจโร อพฺยากโต ธมฺโม นกามาวจรสฺส อพฺยากตสฺส ธมฺมสฺส เหตุปจฺจเยน ปจฺจโย.\csromanspacer{} ตีณิ. (สํขิตฺตํ.)}

\proseitem{70}{เหตุยา จตฺตาริ, อารมฺมเณ จตฺตาริ, อธิปติยา จตฺตาริ, (กามาวจเร เอกํ, นกามาวจเร ตีณิ, กามาวจเร สหชาตาธิปติเยว.) อนนฺตเร จตฺตาริ, สมนนฺตเร จตฺตาริ, สหชาเต สตฺต, อญฺญมญฺเญ ฉ, นิสฺสเย สตฺต, อุปนิสฺสเย จตฺตาริ, ปุเรชาเต เทฺว, ปจฺฉาชาเต เทฺว, อาเสวเน ตีณิ, กมฺเม จตฺตาริ, วิปาเก จตฺตาริ, อาหาเร จตฺตาริ ฯเปฯ สมฺปยุตฺเต เทฺว, วิปฺปยุตฺเต ตีณิ, อตฺถิยา สตฺต, นตฺถิยา จตฺตาริ ฯเปฯ อวิคเต สตฺต. (สํขิตฺตํ.)}

\csromanheader{2}{94. รูปาวจรทุก 1. กุสลตฺติก}
\csromanheader{2}{1-7. ปฏิจฺจาทิวาร ปจฺจยจตุกฺก}
\proseitem{71}{รูปาวจรํ กุสลํ ธมฺมํ ปฏิจฺจ รูปาวจโร กุสโล ธมฺโม อุปฺปชฺชติ เหตุปจฺจยา. (1)}

\prose{นรูปาวจรํ กุสลํ ธมฺมํ ปฏิจฺจ นรูปาวจโร กุสโล ธมฺโม อุปฺปชฺชติ เหตุปจฺจยา. (1) (สํขิตฺตํ.)}

\proseitem{72}{เหตุยา เทฺว, อารมฺมเณ เทฺว. (สพฺพตฺถ เทฺว, สํขิตฺตํ.)}

\prose{น-อธิปติยา เทฺว ฯเปฯ นปุเรชาเต เอกํ, น-อาเสวเน เอกํ ฯเปฯ นวิปฺปยุตฺเต เอกํ. (สํขิตฺตํ.\csromanspacer{} สหชาตวาราทิ วิตฺถาเรตพฺโพ.)}

\proseitem{73}{รูปาวจโร กุสโล ธมฺโม รูปาวจรสฺส กุสลสฺส ธมฺมสฺส เหตุปจฺจเยน ปจฺจโย. (1)}

\gambhira{ธมฺมานุโลม ทุกติกปฏฺฐานปาฬิ}
\prose{\csromanfolio{444}นรูปาวจโร กุสโล ธมฺโม นรูปาวจรสฺส กุสลสฺส ธมฺมสฺส เหตุปจฺจเยน ปจฺจโย. (1) (สํขิตฺตํ.)}

\proseitem{74}{เหตุยา เทฺว, อารมฺมเณ จตฺตาริ, อธิปติยา ตีณิ, (รูปาวจเร เอกํ, สหชาตาธิปติเยว, นรูปาวจเร เทฺว.) อนนฺตเร ตีณิ, (รูปาวจเร เอกํ, นรูปาวจเร เทฺว.) สมนนฺตเร ตีณิ, สหชาเต เทฺว ฯเปฯ อุปนิสฺสเย จตฺตาริ, อาเสวเน ตีณิ, กมฺเม เทฺว ฯเปฯ อตฺถิยา เทฺว, นตฺถิยา ตีณิ ฯเปฯ อวิคเต เทฺว. (สํขิตฺตํ.)}

\proseitem{75}{นรูปาวจรํ อกุสลํ ธมฺมํ ปฏิจฺจ นรูปาวจโร อกุสโล ธมฺโม อุปฺปชฺชติ เหตุปจฺจยา. (สพฺพตฺถ เอกํ.)}

\prose{รูปาวจรํ อพฺยากตํ ธมฺมํ ปฏิจฺจ รูปาวจโร อพฺยากโต ธมฺโม อุปฺปชฺชติ เหตุปจฺจยา.\csromanspacer{} ตีณิ.}

\prose{นรูปาวจรํ อพฺยากตํ ธมฺมํ ปฏิจฺจ นรูปาวจโร อพฺยากโต ธมฺโม อุปฺปชฺชติ เหตุปจฺจยา.\csromanspacer{} ตีณิ.}

\prose{รูปาวจรํ อพฺยากตญฺจ นรูปวจรํ อพฺยากตญฺจ ธมฺมํ ปฏิจฺจ รูปาวจโร อพฺยากโต ธมฺโม อุปฺปชฺชติ เหตุปจฺจยา.\csromanspacer{} ตีณิ. (สํขิตฺตํ.)}

\proseitem{76}{เหตุยา นว, อารมฺมเณ จตฺตาริ, อธิปติยา ปญฺจ ฯเปฯ ปุเรชาเต อาเสวเน เทฺว ฯเปฯ อวิคเต นว. (สํขิตฺตํ.\csromanspacer{} ยถา กามาวจรทุก-อพฺยากตสทิสํ, เอตฺตกาเยว ปญฺหา เหฏฺฐุปริกํ ปริวตฺตนฺติ.\csromanspacer{} สหชาตวารมฺปิ ฯเปฯ ปญฺหาวารมฺปิ สพฺพตฺถ วิตฺถาเรตพฺพํ.)}

\csromanheader{2}{95. อรูปาวจรทุก 1. กุสลตฺติก}
\csromanheader{2}{1-7. ปฏิจฺจาทิวาร ปจฺจยจตุกฺก}
\proseitem{77}{อรูปาวจรํ กุสลํ ธมฺมํ ปฏิจฺจ อรูปาวจโร กุสโล ธมฺโม อุปฺปชฺชติ เหตุปจฺจยา. (1)}

\prose{น-อรูปาวจรํ กุสลํ ธมฺมํ ปฏิจฺจ น-อรูปาวจโร กุสโล ธมฺโม อุปฺปชฺชติ เหตุปจฺจยา. (1) (สํขิตฺตํ.)}

\csromanheader{2}{95. อรูปาวจรทุก 1. กุสลตฺติก}
\proseitem{78}{\csromanfolio{445}เหตุยา เทฺว, อารมฺมเณ เทฺว ฯเปฯ อวิคเต เทฺว. (สํขิตฺตํ.)}

\prose{น-อธิปติยา เทฺว ฯเปฯ น-อาเสวเน เอกํ ฯเปฯ นวิปฺปยุตฺเต เทฺว. (สํขิตฺตํ.\csromanspacer{} สหชาตวาราทิ วิตฺถาเรตพฺโพ.)}

\proseitem{79}{อรูปาวจโร กุสโล ธมฺโม อรูปาวจรสฺส กุสลสฺส ธมฺมสฺส เหตุปจฺจเยน ปจฺจโย. (1)}

\prose{น-อรูปาวจโร กุสโล ธมฺโม น-อรูปาวจรสฺส กุสลสฺส ธมฺมสฺส เหตุปจฺจเยน ปจฺจโย. (1) (สํขิตฺตํ.)}

\proseitem{80}{เหตุยา เทฺว, อารมฺมเณ ตีณิ, อธิปติยา ตีณิ, อนนฺตเร ตีณิ ฯเปฯ สหชาเต เทฺว ฯเปฯ อุปนิสฺสเย จตฺตาริ, อาเสวเน ตีณิ, กมฺเม เทฺว ฯเปฯ อตฺถิยา เทฺว, นตฺถิยา ตีณิ ฯเปฯ อวิคเต เทฺว. (สํขิตฺตํ.)}

\proseitem{81}{น-อรูปาวจรํ อกุสลํ ธมฺมํ ปฏิจฺจ น-อรูปาวจโร อกุสโล ธมฺโม อุปฺปชฺชติ เหตุปจฺจยา.\csromanspacer{} เอกํ. (สพฺพตฺถ เอกํ.\csromanspacer{} สํขิตฺตํ.)}

\prose{อรูปาวจรํ อพฺยากตํ ธมฺมํ ปฏิจฺจ อรูปาวจโร อพฺยากโต ธมฺโม อุปฺปชฺชติ เหตุปจฺจยา.\csromanspacer{} ตีณิ.}

\prose{น-อรูปาวจรํ อพฺยากตํ ธมฺมํ ปฏิจฺจ น-อรูปาวจโร อพฺยากโต ธมฺโม อุปฺปชฺชติ เหตุปจฺจยา. (1)}

\prose{อรูปาวจรํ อพฺยากตญฺจ น-อรูปาวจรํ อพฺยากตญฺจ ธมฺมํ ปฏิจฺจ น-อรูปาวจโร อพฺยากโต ธมฺโม อุปฺปชฺชติ เหตุปจฺจยา. (1) (สํขิตฺตํ.)}

\proseitem{82}{เหตุยา ปญฺจ, อารมฺมเณ เทฺว, อธิปติยา ปญฺจ ฯเปฯ อวิคเต ปญฺจ. (สํขิตฺตํ.)}

\prose{นเหตุยา เอกํ, น-อารมฺมเณ ตีณิ, น-อธิปติยา เทฺว ฯเปฯ นปุเรชาเต จตฺตาริ, นปจฺฉาชาเต น-อาเสวเน ปญฺจ, นกมฺเม เทฺว, นวิปาเก ปญฺจ, น-อาหาเร เอกํ ฯเปฯ นวิปฺปยุตฺเต เทฺว ฯเปฯ โนวิคเต ตีณิ. (สํขิตฺตํ.\csromanspacer{} สหชาตวาราทิ วิตฺถาเรตพฺโพ.)}

\proseitem{83}{อรูปาวจโร อพฺยากโต ธมฺโม อรูปาวจรสฺส อพฺยากตสฺส ธมฺมสฺส เหตุปจฺจเยน ปจฺจโย.\csromanspacer{} ตีณิ.}

\prose{น-อรูปาวจโร อพฺยากโต ธมฺโม น-อรูปาวจรสฺส อพฺยากตสฺส ธมฺมสฺส เหตุปจฺจเยน ปจฺจโย. (1) (สํขิตฺตํ.)}

\gambhira{ธมฺมานุโลม ทุกติกปฏฺฐานปาฬิ}
\proseitem{84}{\csromanfolio{446}เหตุยา จตฺตาริ, อารมฺมเณ ตีณิ, (อรูปาวจรมูเล เทฺว, น-อรูปาวจเร เอกํ.) อธิปติยา จตฺตาริ, (อรูปาวจรมูเล ตีณิ, น-อรูเป เอกํ.) อนนฺตเร จตฺตาริ ฯเปฯ สหชาเต ปญฺจ, อญฺญมญฺเญ เทฺว, นิสฺสเย สตฺต, อุปนิสฺสเย จตฺตาริ, ปุเรชาเต เทฺว, ปจฺฉาชาเต เทฺว, อาเสวเน ตีณิ, กมฺเม จตฺตาริ, วิปาเก เทฺว ฯเปฯ สมฺปยุตฺเต เทฺว, วิปฺปยุตฺเต ตีณิ, อตฺถิยา สตฺต, นตฺถิยา จตฺตาริ ฯเปฯ อวิคเต สตฺต. (สํขิตฺตํ.)}

\csromanheader{2}{96. ปริยาปนฺนทุก 1. กุสลตฺติก}
\csromanheader{2}{1. ปฏิจฺจวาร ปจฺจยจตุกฺก}
\proseitem{85}{ปริยาปนฺนํ กุสลํ ธมฺมํ ปฏิจฺจ ปริยาปนฺโน กุสโล ธมฺโม อุปฺปชฺชติ เหตุปจฺจยา. (1)}

\prose{อปริยาปนฺนํ กุสลํ ธมฺมํ ปฏิจฺจ อปริยาปนฺโน กุสโล ธมฺโม อุปฺปชฺชติ เหตุปจฺจยา. (1) (สํขิตฺตํ.)}

\proseitem{86}{เหตุยา เทฺว, อารมฺมเณ เทฺว ฯเปฯ อวิคเต เทฺว. (สํขิตฺตํ.\csromanspacer{} โลกิยทุกกุสลสทิสํ.\csromanspacer{} สหชาตวาโรปิ ฯเปฯ ปญฺหาวาโรปิ วิตฺถาเรตพฺพา.)}

\proseitem{87}{ปริยาปนฺนํ อกุสลํ ธมฺมํ ปฏิจฺจ ปริยาปนฺโน อกุสโล ธมฺโม อุปฺปชฺชติ เหตุปจฺจยา. (สํขิตฺตํ.)}

\prose{เหตุยา เอกํ, อารมฺมเณ เอกํ ฯเปฯ อวิคเต เอกํ. (สํขิตฺตํ.) (สหชาตวาเรปิ ฯเปฯ ปญฺหาวาเรปิ สพฺพตฺถ เอกํ.)}

\proseitem{88}{ปริยาปนฺนํ อพฺยากตํ ธมฺมํ ปฏิจฺจ ปริยาปนฺโน อพฺยากโต ธมฺโม อุปฺปชฺชติ เหตุปจฺจยา. (1)}

\prose{อปริยาปนฺนํ อพฺยากตํ ธมฺมํ ปฏิจฺจ อปริยาปนฺโน อพฺยากโต ธมฺโม อุปฺปชฺชติ เหตุปจฺจยา.\csromanspacer{} ตีณิ.}

\prose{ปริยาปนฺนํ อพฺยากตญฺจ อปริยาปนฺนํ อพฺยากตญฺจ ธมฺมํ ปฏิจฺจ ปริยาปนฺโน อพฺยากโต ธมฺโม อุปฺปชฺชติ เหตุปจฺจยา. (1) (สํขิตฺตํ.)}

\csromanheader{2}{98. นิยตทุก 1. กุสลตฺติก}
\proseitem{89}{\csromanfolio{447}เหตุยา ปญฺจ, อารมฺมเณ เทฺว ฯเปฯ อวิคเต ปญฺจ. (สํขิตฺตํ.\csromanspacer{} โลกิยทุก-อพฺยากตสทิสํ.\csromanspacer{} สหชาตวารมฺปิ ฯเปฯ ปญฺหาวารมฺปิ สพฺพตฺถ วิตฺถาเรตพฺพํ.)}

\csromanheader{2}{97. นิยฺยานิกทุก 1. กุสลตฺติก}
\csromanheader{2}{1. ปฏิจฺจวาร ปจฺจยจตุกฺก}
\proseitem{90}{นิยฺยานิกํ กุสลํ ธมฺมํ ปฏิจฺจ นิยฺยานิโก กุสโล ธมฺโม อุปฺปชฺชติ เหตุปจฺจยา. (1)}

\prose{อนิยฺยานิกํ กุสลํ ธมฺมํ ปฏิจฺจ อนิยฺยานิโก กุสโล ธมฺโม อุปฺปชฺชติ เหตุปจฺจยา. (1) (สํขิตฺตํ.)}

\proseitem{91}{เหตุยา เทฺว, อารมฺมเณ เทฺว ฯเปฯ อวิคเต เทฺว. (สํขิตฺตํ.\csromanspacer{} โลกิยทุกสทิสํ.\csromanspacer{} สหชาตวารมฺปิ ฯเปฯ ปญฺหาวารมฺปิ วิตฺถาเรตพฺพํ.)}

\proseitem{92}{อนิยฺยานิกํ อกุสลํ ธมฺมํ ปฏิจฺจ อนิยฺยานิโก อกุสโล ธมฺโม อุปฺปชฺชติ เหตุปจฺจยา. (1) (สํขิตฺตํ.)}

\prose{เหตุยา เอกํ, อารมฺมเณ เอกํ ฯเปฯ อวิคเต เอกํ. (สํขิตฺตํ.) (สหชาตวาเรปิ ฯเปฯ ปญฺหาวาเรปิ สพฺพตฺถ เอกํ.)}

\proseitem{93}{อนิยฺยานิกํ อพฺยากตํ ธมฺมํ ปฏิจฺจ อนิยฺยานิโก อพฺยากโต ธมฺโม อุปฺปชฺชติ เหตุปจฺจยา. (1) (สํขิตฺตํ.)}

\prose{เหตุยา เอกํ, อารมฺมเณ เอกํ ฯเปฯ อวิคเต เอกํ. (สํขิตฺตํ.)}

\vspace{\tipitakaparskip}
\csromancenter{(สหชาตวาเรปิ ฯเปฯ ปญฺหาวาเรปิ สพฺพตฺถ เอกํ.)}
\csromanheader{2}{98. นิยตทุก 1. กุสลตฺติก}
\csromanheader{2}{1-7. ปฏิจฺจาทิวาร ปจฺจยจตุกฺก}
\proseitem{94}{นิยตํ กุสลํ ธมฺมํ ปฏิจฺจ นิยโต กุสโล ธมฺโม อุปฺปชฺชติ เหตุปจฺจยา. (1)}

\gambhira{ธมฺมานุโลม ทุกติกปฏฺฐานปาฬิ}
\prose{\csromanfolio{448}อนิยตํ กุสลํ ธมฺมํ ปฏิจฺจ อนิยโต กุสโล ธมฺโม อุปฺปชฺชติ เหตุปจฺจยา. (1) (สํขิตฺตํ.)}

\proseitem{95}{เหตุยา เทฺว, อารมฺมเณ เทฺว ฯเปฯ อวิคเต เทฺว. (สํขิตฺตํ.\csromanspacer{} โลกิยทุกกุสลสทิสํ.\csromanspacer{} สหชาตวารมฺปิ ฯเปฯ ปญฺหาวารมฺปิ วิตฺถาเรตพฺพํ.)}

\proseitem{96}{นิยตํ อกุสลํ ธมฺมํ ปฏิจฺจ นิยโต อกุสโล ธมฺโม อุปฺปชฺชติ เหตุปจฺจยา. (1)}

\prose{อนิยตํ อกุสลํ ธมฺมํ ปฏิจฺจ อนิยโต อกุสโล ธมฺโม อุปฺปชฺชติ เหตุปจฺจยา. (1) (สํขิตฺตํ.)}

\proseitem{97}{เหตุยา เทฺว, อารมฺมเณ เทฺว, (สพฺพตฺถ เทฺว.) อวิคเต เทฺว. (สํขิตฺตํ.)}

\prose{นเหตุยา เอกํ, น-อาธิปติยา เทฺว, นปุเรชาเต เอกํ ฯเปฯ นกมฺเม เทฺว ฯเปฯ นวิปฺปยุตฺเต เอกํ. (สํขิตฺตํ.\csromanspacer{} สหชาตวาราทิ วิตฺถาเรตพฺโพ.)}

\proseitem{98}{นิยโต อกุสโล ธมฺโม นิยตสฺส อกุสลสฺส ธมฺมสฺส เหตุปจฺจเยน ปจฺจโย. (1)}

\prose{อนิยโต อกุสโล ธมฺโม อนิยตสฺส อกุสลสฺส ธมฺมสฺส เหตุปจฺจเยน ปจฺจโย. (1) (สํขิตฺตํ.)}

\proseitem{99}{เหตุยา เทฺว, อารมฺมเณ ตีณิ, อธิปติยา เทฺว, (นิยเต สหชาตาธิปติ.\csromanspacer{} ทุติเย อารมฺมณาธิปติ, สหชาตาธิปติ.) อนนฺตเร เทฺว ฯเปฯ นิสฺสเย เทฺว, อุปนิสฺสเย ตีณิ, อาเสวเน เทฺว, กมฺเม เทฺว ฯเปฯ อวิคเต เทฺว. (สํขิตฺตํ.)}

\proseitem{100}{อนิยตํ อพฺยากตํ ธมฺมํ ปฏิจฺจ อนิยโต อพฺยากโต ธมฺโม อุปฺปชฺชติ เหตุปจฺจยา. (สํขิตฺตํ.)}

\prose{เหตุยา เอกํ, อารมฺมเณ เอกํ ฯเปฯ อวิคเต เอกํ. (สํขิตฺตํ.\csromanspacer{} สหชาตวาเรปิ ฯเปฯ ปญฺหาวาเรปิ สพฺพตฺถ เอกํ.)}

\csromanheader{2}{100. สรณทุก 1. กุสลตฺติก \textbf{99. ส-อุตฺตรทุก 1. กุสลตฺติก}}
\csromanheader{2}{1. ปฏิจฺจวาร ปจฺจยจตุกฺก}
\proseitem{101}{\csromanfolio{449}ส-อุตฺตรํ กุสลํ ธมฺมํ ปฏิจฺจ ส-อุตฺตโร กุสโล ธมฺโม อุปฺปชฺชติ เหตุปจฺจยา. (1)}

\prose{อนุตฺตรํ กุสลํ ธมฺมํ ปฏิจฺจ อนุตฺตโร กุสโล ธมฺโม อุปฺปชฺชติ เหตุปจฺจยา. (1) (สํขิตฺตํ.)}

\proseitem{102}{เหตุยา เทฺว, อารมฺมเณ เทฺว ฯเปฯ อวิคเต เทฺว. (สํขิตฺตํ.\csromanspacer{} สหชาตวาโรปิ ฯเปฯ ปญฺหาวาโรปิ สพฺพตฺถ วิตฺถาเรตพฺพา.)}

\proseitem{103}{ส-อุตฺตรํ อกุสลํ ธมฺมํ ปฏิจฺจ ส-อุตฺตโร อกุสโล ธมฺโม อุปฺปชฺชติ เหตุปจฺจยา. (สํขิตฺตํ.)}

\prose{เหตุยา เอกํ, อารมฺมเณ เอกํ ฯเปฯ อวิคเต เอกํ. (สํขิตฺตํ.\csromanspacer{} สหชาตวาเรปิ ฯเปฯ ปญฺหาวาเรปิ สพฺพตฺถ เอกํ.)}

\prose{ส-อุตฺตรํ อพฺยากตํ ธมฺมํ ปฏิจฺจ ส-อุตฺตโร อพฺยากโต ธมฺโม อุปฺปชฺชติ เหตุปจฺจยา. (1)}

\prose{อนุตฺตรํ อพฺยากตํ ธมฺมํ ปฏิจฺจ อนุตฺตโร อพฺยากโต ธมฺโม อุปฺปชฺชติ เหตุปจฺจยา.\csromanspacer{} ตีณิ.}

\prose{ส-อุตฺตรํ อพฺยากตญฺจ อนุตฺตรํ อพฺยากตญฺจ ธมฺมํ ปฏิจฺจ ส-อุตฺตโร อพฺยากโต ธมฺโม อุปฺปชฺชติ เหตุปจฺจยา. (1) (สํขิตฺตํ.)}

\proseitem{104}{เหตุยา ปญฺจ, อารมฺมเณ เทฺว ฯเปฯ อวิคเต ปญฺจ. (สํขิตฺตํ.\csromanspacer{} สหชาตวารมฺปิ ฯเปฯ ปญฺหาวารมฺปิ สพฺพตฺถ วิตฺถาเรตพฺพํ.\csromanspacer{} โลกิยทุก-อพฺยากตสทิสํ.)}

\csromanmatikaanchor{mtk.84}
\csromanmatikaanchor{mtk.85}
\csromantocmarknopage{subparagraph}{100. สรณทุก}{mtk.84}
\bookmark[dest={mtk.84},level=5]{100. สรณทุก}
\csromantocmarkref{subparagraph}{1. กุสลตฺติก}{mtk.85}
\bookmark[dest={mtk.85},level=5]{1. กุสลตฺติก}
\csromanheader{6}{100. สรณทุก 1. กุสลตฺติก}
\csromanheader{2}{1-7. ปฏิจฺจาทิวาร ปจฺจยจตุกฺก}
\proseitem{105}{อรณํ กุสลํ ธมฺมํ ปฏิจฺจ อรโณ กุสโล ธมฺโม อุปฺปชฺชติ เหตุปจฺจยา. (สํขิตฺตํ.)}

\gambhira{ธมฺมานุโลม ทุกติกปฏฺฐานปาฬิ}
\prose{\csromanfolio{450}เหตุยา เอกํ, อารมฺมเณ เอกํ ฯเปฯ อวิคเต เอกํ. (สํขิตฺตํ.\csromanspacer{} สหชาตวาเรปิ ฯเปฯ ปญฺหาวาเรปิ สพฺพตฺถ เอกํ.)}

\proseitem{106}{สรณํ อกุสลํ ธมฺมํ ปฏิจฺจ สรโณ อกุสโล ธมฺโม อุปฺปชฺชติ เหตุปจฺจยา. (สํขิตฺตํ.)}

\prose{เหตุยา เอกํ, อารมฺมเณ เอกํ ฯเปฯ อวิคเต เอกํ. (สํขิตฺตํ.\csromanspacer{} สหชาตวาเรปิ ฯเปฯ ปญฺหาวาเรปิ สพฺพตฺถ เอกํ.)}

\proseitem{107}{อรณํ อพฺยากตํ ธมฺมํ ปฏิจฺจ อรโณ อพฺยากโต ธมฺโม อุปฺปชฺชติ เหตุปจฺจยา. (สํขิตฺตํ.)}

\prose{เหตุยา เอกํ, อารมฺมเณ เอกํ ฯเปฯ อวิคเต เอกํ. (สํขิตฺตํ.\csromanspacer{} สหชาตวาเรปิ ฯเปฯ สมฺปยุตฺตวาเรปิ สพฺพตฺถ เอกํ.)}

\proseitem{108}{อรโณ อพฺยากโต ธมฺโม อรณสฺส อพฺยากตสฺส ธมฺมสฺส เหตุปจฺจเยน ปจฺจโย. (สํขิตฺตํ.)}

\prose{เหตุยา เอกํ, อารมฺมเณ เอกํ ฯเปฯ อวิคเต เอกํ. (สํขิตฺตํ.) (ยถา กุสลตฺติเก ปญฺหาวารํ, เอวํ วิตฺถาเรตพฺพํ.)}

\nitthitam{ปิฏฺฐิทุกกุสลตฺติกํ นิฏฺฐิตํ.}
\csromanmatikaanchor{mtk.86}
\csromantocmarkref{subparagraph}{2. เวทนาตฺติก}{mtk.86}
\bookmark[dest={mtk.86},level=5]{2. เวทนาตฺติก}
\csromanheader{6}{100. สรณทุก 2. เวทนาตฺติก}
\csromanheader{2}{1. ปฏิจฺจวาร ปจฺจยจตุกฺก}
\proseitem{1}{สรณํ สุขาย เวทนาย สมฺปยุตฺตํ ธมฺมํ ปฏิจฺจ สรโณ สุขาย เวทนาย สมฺปยุตฺโต ธมฺโม อุปฺปชฺชติ เหตุปจฺจยา. (1)}

\prose{อรณํ สุขาย เวทนาย สมฺปยุตฺตํ ธมฺมํ ปฏิจฺจ อรโณ สุขาย เวทนาย สมฺปยุตฺโต ธมฺโม อุปฺปชฺชติ เหตุปจฺจยา. (1) (สํขิตฺตํ.)}

\proseitem{2}{เหตุยา เทฺว, อารมฺมเณ เทฺว ฯเปฯ กมฺเม เทฺว, วิปาเก เอกํ ฯเปฯ อวิคเต เทฺว. (สํขิตฺตํ.\csromanspacer{} สหชาตวาโรปิ ฯเปฯ ปญฺหาวาโรปิ สพฺพตฺถ วิตฺถาเรตพฺพา.)}

\csromanmatikaanchor{mtk.87}
\csromantocmarkref{subparagraph}{3. วิปากตฺติก}{mtk.87}
\bookmark[dest={mtk.87},level=5]{3. วิปากตฺติก}
\csromanheader{6}{100. สรณทุก 3. วิปากตฺติก}
\proseitem{3}{\csromanfolio{451}สรณํ ทุกฺขาย เวทนาย สมฺปยุตฺตํ ธมฺมํ ปฏิจฺจ สรโณ ทุกฺขาย เวทนาย สมฺปยุตฺโต ธมฺโม อุปฺปชฺชติ เหตุปจฺจยา. (1)}

\prose{สรณํ ทุกฺขาย เวทนาย สมฺปยุตฺตํ ธมฺมํ ปฏิจฺจ สรโณ ทุกฺขาย เวทนาย สมฺปยุตฺโต ธมฺโม อุปฺปชฺชติ อารมฺมณปจฺจยา. (1) (สํขิตฺตํ.)}

\proseitem{4}{เหตุยา เอกํ, อารมฺมเณ เทฺว ฯเปฯ อวิคเต เทฺว. (สํขิตฺตํ.\csromanspacer{} สหชาตวารมฺปิ ฯเปฯ ปญฺหาวารมฺปิ วิตฺถาเรตพฺพํ.)}

\proseitem{5}{สรณํ อทุกฺขมสุขาย เวทนาย สมฺปยุตฺตํ ธมฺมํ ปฏิจฺจ สรโณ อทุกฺขมสุขาย เวทนาย สมฺปยุตฺโต ธมฺโม อุปฺปชฺชติ เหตุปจฺจยา. (1)}

\prose{อรณํ อทุกฺขมสุขาย เวทนาย สมฺปยุตฺตํ ธมฺมํ ปฏิจฺจ อรโณ อทุกฺขมสุขาย เวทนาย สมฺปยุตฺโต ธมฺโม อุปฺปชฺชติ เหตุปจฺจยา. (สํขิตฺตํ.)}

\proseitem{6}{เหตุยา เทฺว, อารมฺมเณ เทฺว ฯเปฯ อวิคเต เทฺว. (สํขิตฺตํ.\csromanspacer{} สหชาตวารมฺปิ ฯเปฯ ปญฺหาวารมฺปิ สพฺพตฺถ วิตฺถาเรตพฺพํ.)}

\nitthitam{สรณทุกเวทนาตฺติกํ นิฏฺฐิตํ.}
\csromanheader{2}{100. สรณทุก 3. วิปากตฺติก}
\csromanheader{2}{1. ปฏิจฺจวาร ปจฺจยจตุกฺก}
\proseitem{7}{อรณํ วิปากํ ธมฺมํ ปฏิจฺจ อรโณ วิปาโก ธมฺโม อุปฺปชฺชติ เหตุปจฺจยา. (สํขิตฺตํ.)}

\prose{เหตุยา เอกํ, อารมฺมเณ เอกํ ฯเปฯ อวิคเต เอกํ. (สํขิตฺตํ.\csromanspacer{} สหชาตวาเรปิ ฯเปฯ ปญฺหาวาเรปิ สพฺพตฺถ เอกํ.)}

\proseitem{8}{สรณํ วิปากธมฺมธมฺมํ ปฏิจฺจ สรโณ วิปากธมฺมธมฺโม อุปฺปชฺชติ เหตุปจฺจยา. (1)}

\prose{อรณํ วิปากธมฺมธมฺมํ ปฏิจฺจ อรโณ วิปากธมฺมธมฺโม อุปฺปชฺชติ เหตุปจฺจยา. (1) (สํขิตฺตํ.)}

\gambhira{ธมฺมานุโลม ทุกติกปฏฺฐานปาฬิ}
\proseitem{9}{\csromanfolio{452}เหตุยา เทฺว, อารมฺมเณ เทฺว ฯเปฯ อวิคเต เทฺว. (สํขิตฺตํ.\csromanspacer{} สหชาตวารมฺปิ ฯเปฯ ปญฺหาวารมฺปิ วิตฺถาเรตพฺพํ.)}

\prose{อรณํ เนววิปากนวิปากธมฺมธมฺมํ ปฏิจฺจ อรโณ เนววิปากนวิปากธมฺมธมฺโม อุปฺปชฺชติ เหตุปจฺจยา. (สํขิตฺตํ.)}

\proseitem{10}{เหตุยา เอกํ, อารมฺมเณ เอกํ ฯเปฯ อวิคเต เอกํ. (สํขิตฺตํ.\csromanspacer{} สหชาตวาเรปิ ฯเปฯ ปญฺหาวาเรปิ สพฺพตฺถ เอกํ.)}

\csromanmatikaanchor{mtk.88}
\csromantocmarkref{subparagraph}{4. อุปทินฺนตฺติก}{mtk.88}
\bookmark[dest={mtk.88},level=5]{4. อุปทินฺนตฺติก}
\csromanheader{6}{100. สรณทุก 4. อุปาทินฺนตฺติก}
\csromanheader{2}{1. ปฏิจฺจวาร ปจฺจยจตุกฺก}
\proseitem{11}{อรณํ อุปาทินฺนุปาทานิยํ ธมฺมํ ปฏิจฺจ อรโณ อุปาทินฺนุปาทานิโย ธมฺโม อุปฺปชฺชติ เหตุปจฺจยา. (สํขิตฺตํ.)}

\prose{เหตุยา เอกํ, อารมฺมเณ เอกํ ฯเปฯ อวิคเต เอกํ. (สํขิตฺตํ.\csromanspacer{} สหชาตวาเรปิ ฯเปฯ ปญฺหาวาเรปิ สพฺพตฺถ เอกํ.)}

\proseitem{12}{สรณํ อนุปาทินฺนุปาทานิยํ ธมฺมํ ปฏิจฺจ สรโณ อนุปาทินฺนุปาทานิโย ธมฺโม อุปฺปชฺชติ เหตุปจฺจยา.\csromanspacer{} ตีณิ.}

\prose{อรณํ อนุปาทินฺนุปาทานิยํ ธมฺมํ ปฏิจฺจ อรโณ อนุปาทินฺนุปาทานิโย ธมฺโม อุปฺปชฺชติ เหตุปจฺจยา. (1)}

\prose{สรณํ อนุปาทินฺนุปาทานิยญฺจ อรณํ อนุปาทินฺนุปาทานิยญฺจ ธมฺมํ ปฏิจฺจ อรโณ อนุปาทินฺนุปาทานิโย ธมฺโม อุปฺปชฺชติ เหตุปจฺจยา. (1) (สํขิตฺตํ.)}

\proseitem{13}{เหตุยา ปญฺจ, อารมฺมเณ เทฺว ฯเปฯ อวิคเต ปญฺจ. (สํขิตฺตํ.\csromanspacer{} สหชาตวารมฺปิ ฯเปฯ ปญฺหาวารมฺปิ สพฺพตฺถ วิตฺถาเรตพฺพํ.)}

\proseitem{14}{อรณํ อนุปาทินฺน-อนุปาทานิยํ ธมฺมํ ปฏิจฺจ อรโณ อนุปาทินฺน-อนุปาทานิโย ธมฺโม อุปฺปชฺชติ เหตุปจฺจยา. (สํขิตฺตํ.)}

\prose{เหตุยา เอกํ, อารมฺมเณ เอกํ ฯเปฯ อวิคเต เอกํ. (สํขิตฺตํ.\csromanspacer{} สหชาตวาเรปิ ฯเปฯ ปญฺหาวาเรปิ สพฺพตฺถ เอกํ.)}

\vspace{\tipitakaparskip}
\csromancenter{สรณทุก 6.\csromanspacer{} วิตกฺกตฺติก \textbf{100.}\csromanspacer{}\textbf{ สรณทุก 5.}\csromanspacer{}\textbf{ สํกิลิฏฺฐตฺติก}}
\csromanheader{2}{1. ปฏิจฺจวาร ปจฺจยจตุกฺก}
\csromanmatikaanchor{mtk.89}
\csromantocmarkref{subparagraph}{5-21. สํกิลิฏฺฐตฺติกาทิ}{mtk.89}
\bookmark[dest={mtk.89},level=5]{5-21. สํกิลิฏฺฐตฺติกาทิ}
\proseitem{15}{\csromanfolio{453}สรณํ สํกิลิฏฺฐสํกิเลสิกํ ธมฺมํ ปฏิจฺจ สรโณ สํกิลิฏฺฐสํกิเลสิโก ธมฺโม อุปฺปชฺชติ เหตุปจฺจยา. (สํขิตฺตํ.)}

\prose{เหตุยา เอกํ, อารมฺมเณ เอกํ ฯเปฯ อวิคเต เอกํ. (สํขิตฺตํ.\csromanspacer{} สหชาตวาเรปิ ฯเปฯ ปญฺหาวาเรปิ สพฺพตฺถ เอกํ.)}

\proseitem{16}{อรณํ อสํกิลิฏฺฐสํกิเลสิกํ ธมฺมํ ปฏิจฺจ อรโณ อสํกิลิฏฺฐสํกิเลสิโก ธมฺโม อุปฺปชฺชติ เหตุปจฺจยา. (สํขิตฺตํ.)}

\prose{เหตุยา เอกํ, อารมฺมเณ เอกํ ฯเปฯ อวิคเต เอกํ. (สํขิตฺตํ.\csromanspacer{} สหชาตวาเรปิ ฯเปฯ ปญฺหาวาเรปิ สพฺพตฺถ เอกํ.)}

\proseitem{17}{อรณํ อสํกิลิฏฺฐ-อสํกิเลสิกํ ธมฺมํ ปฏิจฺจ อรโณ อสํกิลิฏฺฐ-อสํกิเลสิโก ธมฺโม อุปฺปชฺชติ เหตุปจฺจยา. (สํขิตฺตํ.)}

\prose{เหตุยา เอกํ, อารมฺมเณ เอกํ ฯเปฯ อวิคเต เอกํ. (สํขิตฺตํ.\csromanspacer{} สหชาตวาเรปิ ฯเปฯ ปญฺหาวาเรปิ สพฺพตฺถ เอกํ.)}

\csromanheader{2}{100. สรณทุก 6. วิตกฺกตฺติก}
\csromanheader{2}{1. ปฏิจฺจวาร ปจฺจยจตุกฺก}
\proseitem{18}{สรณํ สวิตกฺกสวิจารํ ธมฺมํ ปฏิจฺจ สรโณ สวิตกฺกสวิจาโร ธมฺโม อุปฺปชฺชติ เหตุปจฺจยา. (1)}

\prose{อรณํ สวิตกฺกสวิจารํ ธมฺมํ ปฏิจฺจ อรโณ สวิตกฺกสวิจาโร ธมฺโม อุปฺปชฺชติ เหตุปจฺจยา. (1) (สํขิตฺตํ.)}

\proseitem{19}{เหตุยา เทฺว, อารมฺมเณ เทฺว ฯเปฯ อวิคเต เทฺว. (สํขิตฺตํ.\csromanspacer{} สหชาตวารมฺปิ ฯเปฯ ปญฺหาวารมฺปิ วิตฺถาเรตพฺพํ.)}

\proseitem{20}{อรณํ อวิตกฺกวิจารมตฺตํ ธมฺมํ ปฏิจฺจ อรโณ อวิตกฺกวิจารมตฺโต ธมฺโม อุปฺปชฺชติ เหตุปจฺจยา. (สํขิตฺตํ.)}

\gambhira{ธมฺมานุโลม ทุกติกปฏฺฐานปาฬิ}
\prose{\csromanfolio{454}เหตุยา เอกํ ฯเปฯ อวิคเต เอกํ. (สํขิตฺตํ.\csromanspacer{} สหชาตวาเรปิ ฯเปฯ ปญฺหาวาเรปิ สพฺพตฺถ เอกํ.)}

\proseitem{21}{อรณํ อวิตกฺก-อวิจารํ ธมฺมํ ปฏิจฺจ อรโณ อวิตกฺก-อวิจาโร ธมฺโม อุปฺปชฺชติ เหตุปจฺจยา. (สํขิตฺตํ.)}

\prose{เหตุยา เอกํ, (สพฺพตฺถ เอกํ) ฯเปฯ อวิคเต เอกํ. (สํขิตฺตํ.)}

\csromanheader{2}{100. สรณทุก 7. ปีติตฺติก}
\csromanheader{2}{1. ปฏิจฺจวาร ปจฺจยจตุกฺก}
\proseitem{22}{สรณํ ปีติสหคตํ ธมฺมํ ปฏิจฺจ สรโณ ปีติสหคโต ธมฺโม อุปฺปชฺชติ เหตุปจฺจยา. (1)}

\prose{อรณํ ปีติสหคตํ ธมฺมํ ปฏิจฺจ อรโณ ปีติสหคโต ธมฺโม อุปฺปชฺชติ เหตุปจฺจยา. (1) (สํขิตฺตํ.)}

\proseitem{23}{เหตุยา เทฺว, อารมฺมเณ เทฺว ฯเปฯ อวิคเต เทฺว. (สํขิตฺตํ.\csromanspacer{} สหชาตวาโรปิ ฯเปฯ ปญฺหาวาโรปิ สพฺพตฺถ วิตฺถาเรตพฺพา.)}

\proseitem{24}{สรณํ สุขสหคตํ ธมฺมํ ปฏิจฺจ สรโณ สุขสหคโต ธมฺโม อุปฺปชฺชติ เหตุปจฺจยา. (1)}

\prose{อรณํ สุขสหคตํ ธมฺมํ ปฏิจฺจ อรโณ สุขสหคโต ธมฺโม อุปฺปชฺชติ เหตุปจฺจยา. (1) (สํขิตฺตํ.)}

\proseitem{25}{เหตุยา เทฺว, อารมฺมเณ เทฺว ฯเปฯ อวิคเต เทฺว. (สํขิตฺตํ.\csromanspacer{} สหชาตวาโรปิ ฯเปฯ ปญฺหาวาโรปิ วิตฺถาเรตพฺพา.)}

\proseitem{26}{สรณํ อุเปกฺขาสหคตํ ธมฺมํ ปฏิจฺจ สรโณ อุเปกฺขาสหคโต ธมฺโม อุปฺปชฺชติ เหตุปจฺจยา. (1)}

\prose{อรณํ อุเปกฺขาสหคตํ ธมฺมํ ปฏิจฺจ อรโณ อุเปกฺขาสหคโต ธมฺโม อุปฺปชฺชติ เหตุปจฺจยา. (1) (สํขิตฺตํ.)}

\proseitem{27}{เหตุยา เทฺว, อารมฺมเณ เทฺว ฯเปฯ อวิคเต เทฺว. (สํขิตฺตํ.\csromanspacer{} สหชาตวาเรปิ ฯเปฯ ปญฺหาวาเรปิ วิตฺถาเรตพฺโพ.)}

\vspace{\tipitakaparskip}
\csromancenter{สรณทุก 9.\csromanspacer{} ทสฺสเนนปหาตพฺพเหตุกตฺติก \textbf{100.}\csromanspacer{}\textbf{ สรณทุก 8.}\csromanspacer{}\textbf{ ทสฺสเนนปหาตพฺพตฺติก}}
\csromanheader{2}{1. ปฏิจฺจวาร ปจฺจยจตุกฺก}
\proseitem{28}{\csromanfolio{455}สรณํ ทสฺสเนน ปหาตพฺพํ ธมฺมํ ปฏิจฺจ สรโณ ทสฺสเนน ปหาตพฺโพ ธมฺโม อุปฺปชฺชติ เหตุปจฺจยา. (สํขิตฺตํ.)}

\prose{เหตุยา เอกํ, อารมฺมเณ เอกํ ฯเปฯ อวิคเต เอกํ. (สํขิตฺตํ.\csromanspacer{} สหชาตวาเรปิ ฯเปฯ ปญฺหาวาเรปิ สพฺพตฺถ เอกํ.)}

\proseitem{29}{สรณํ ภาวนาย ปหาตพฺพํ ธมฺมํ ปฏิจฺจ สรโณ ภาวนาย ปหาตพฺโพ ธมฺโม อุปฺปชฺชติ เหตุปจฺจยา. (สํขิตฺตํ.)}

\prose{เหตุยา เอกํ, อารมฺมเณ เอกํ ฯเปฯ อวิคเต เอกํ. (สํขิตฺตํ.\csromanspacer{} สหชาตวาเรปิ ฯเปฯ ปญฺหาวาเรปิ สพฺพตฺถ เอกํ.)}

\proseitem{30}{อรณํ เนวทสฺสเนน นภาวนาย ปหาตพฺพํ ธมฺมํ ปฏิจฺจ อรโณ เนวทสฺสเนน นภาวนาย ปหาตพฺโพ ธมฺโม อุปฺปชฺชติ เหตุปจฺจยา.}

\prose{เหตุยา เอกํ, (สพฺพตฺถ เอกํ) ฯเปฯ อวิคเต เอกํ. (สํขิตฺตํ.)}

\csromanheader{2}{100. สรณทุก 9. ทสฺสเนนปหาตพฺพเหตุกตฺติก}
\csromanheader{2}{1. ปฏิจฺจวาร ปจฺจยจตุกฺก}
\proseitem{31}{สรณํ ทสฺสเนน ปหาตพฺพเหตุกํ ธมฺมํ ปฏิจฺจ สรโณ ทสฺสเนน ปหาตพฺพเหตุโก ธมฺโม อุปฺปชฺชติ เหตุปจฺจยา. (สํขิตฺตํ.)}

\prose{เหตุยา เอกํ, อารมฺมเณ เอกํ ฯเปฯ อวิคเต เอกํ. (สํขิตฺตํ.\csromanspacer{} สหชาตวาเรปิ ฯเปฯ ปญฺหาวาเรปิ สพฺพตฺถ เอกํ.)}

\proseitem{32}{สรณํ ภาวนาย ปหาตพฺพเหตุกํ ธมฺมํ ปฏิจฺจ สรโณ ภาวนาย ปหาตพฺพเหตุโก ธมฺโม อุปฺปชฺชติ เหตุปจฺจยา. (สํขิตฺตํ.)}

\prose{เหตุยา เอกํ, อารมฺมเณ เอกํ ฯเปฯ อวิคเต เอกํ. (สํขิตฺตํ.\csromanspacer{} สหชาตวาเรปิ ฯเปฯ ปญฺหาวาเรปิ สพฺพตฺถ เอกํ.)}

\gambhira{ธมฺมานุโลม ทุกติกปฏฺฐานปาฬิ}
\proseitem{33}{\csromanfolio{456}สรณํ เนวทสฺสเนน นภาวนาย ปหาตพฺพเหตุกํ ธมฺมํ ปฏิจฺจ อรโณ เนวทสฺสเนน นภาวนาย ปหาตพฺพเหตุโก ธมฺโม อุปฺปชฺชติ เหตุปจฺจยา. (สํขิตฺตํ.)}

\prose{เหตุยา ตีณิ, อารมฺมเณ เอกํ. (สพฺพตฺถ วิตฺถาเรตพฺพํ.)}

\csromanheader{2}{100. สรณทุก 10. อาจยคามิตฺติก}
\csromanheader{2}{1. ปฏิจฺจวาร ปจฺจยจตุกฺก}
\proseitem{34}{สรณํ อาจยคามิํ ธมฺมํ ปฏิจฺจ สรโณ อาจยคามี ธมฺโม อุปฺปชฺชติ เหตุปจฺจยา. (1)}

\prose{อรณํ อาจยคามิํ ธมฺมํ ปฏิจฺจ อรโณ อาจยคามี ธมฺโม อุปฺปชฺชติ เหตุปจฺจยา. (1) (สํขิตฺตํ.)}

\proseitem{35}{เหตุยา เทฺว, อารมฺมเณ เทฺว ฯเปฯ อวิคเต เทฺว. (สํขิตฺตํ.)}

\prose{(สหชาตวาเรปิ ฯเปฯ ปญฺหาวาเรปิ สพฺพตฺถ วิตฺถาเรตพฺโพ.)}

\proseitem{36}{อรณํ อปจยคามิํ ธมฺมํ ปฏิจฺจ อรโณ อปจยคามี ธมฺโม อุปฺปชฺชติ เหตุปจฺจยา. (สํขิตฺตํ.)}

\prose{เหตุยา เอกํ, อารมฺมเณ เอกํ ฯเปฯ อวิคเต เอกํ. (สํขิตฺตํ.)}

\prose{(สหชาตวาเรปิ ฯเปฯ ปญฺหาวาเรปิ สพฺพตฺถ เอกํ.)}

\proseitem{37}{อรณํ เนวาจยคามินาปจยคามิํ ธมฺมํ ปฏิจฺจ อรโณ เนวาจยคามินาปจยคามี ธมฺโม อุปฺปชฺชติ เหตุปจฺจยา. (สํขิตฺตํ.)}

\prose{เหตุยา เอกํ, อารมฺมเณ เอกํ ฯเปฯ อวิคเต เอกํ. (สํขิตฺตํ.)}

\prose{(สหชาตวาเรปิ ฯเปฯ ปญฺหาวาเรปิ สพฺพตฺถ เอกํ.)}

\csromanheader{2}{100. สรณทุก 12. ปริตฺตตฺติก \textbf{100. สรณทุก 11. เสกฺขตฺติก}}
\csromanheader{2}{1. ปฏิจฺจวาร ปจฺจยจตุกฺก}
\proseitem{38}{\csromanfolio{457}อรณํ เสกฺขํ ธมฺมํ ปฏิจฺจ อรโณ เสกฺโข ธมฺโม อุปฺปชฺชติ เหตุปจฺจยา. (สํขิตฺตํ.)}

\prose{เหตุยา เอกํ, อารมฺมเณ เอกํ ฯเปฯ อวิคเต เอกํ. (สํขิตฺตํ.)}

\prose{(สหชาตวาเรปิ ฯเปฯ ปญฺหาวาเรปิ สพฺพตฺถ เอกํ.)}

\proseitem{39}{อรณํ อเสกฺขํ ธมฺมํ ปฏิจฺจ อรโณ อเสกฺโข ธมฺโม อุปฺปชฺชติ เหตุปจฺจยา. (สํขิตฺตํ.)}

\prose{เหตุยา เอกํ, อารมฺมเณ เอกํ ฯเปฯ อวิคเต เอกํ. (สํขิตฺตํ.)}

\prose{(สหชาตวาเรปิ ฯเปฯ ปญฺหาวาเรปิ สพฺพตฺถ เอกํ.)}

\proseitem{40}{สรณํ เนวเสกฺขนาเสกฺขํ ธมฺมํ ปฏิจฺจ สรโณ เนวเสกฺขนาเสกฺโข ธมฺโม อุปฺปชฺชติ เหตุปจฺจยา.\csromanspacer{} ตีณิ.}

\prose{อรณํ เนวเสกฺขนาเสกฺขํ ธมฺมํ ปฏิจฺจ อรโณ เนวเสกฺขนาเสกฺโข ธมฺโม อุปฺปชฺชติ เหตุปจฺจยา. (1)}

\prose{สรณํ เนวเสกฺขนาเสกฺขญฺจ อรณํ เนวเสกฺขนาเสกฺขญฺจ ธมฺมํ ปฏิจฺจ อรโณ เนวเสกฺขนาเสกฺโข ธมฺโม อุปฺปชฺชติ เหตุปจฺจยา. (1) (สํขิตฺตํ.)}

\proseitem{41}{เหตุยา ปญฺจ, อารมฺมเณ เทฺว, อธิปติยา ปญฺจ ฯเปฯ อวิคเต ปญฺจ. (สํขิตฺตํ.\csromanspacer{} สหชาตวาเรปิ ฯเปฯ ปญฺหาวาเรปิ สพฺพตฺถ วิตฺถาเรตพฺโพ.)}

\csromanheader{2}{100. สรณทุก 12. ปริตฺตตฺติก}
\csromanheader{2}{1. ปฏิจฺจวาร ปจฺจยจตุกฺก}
\proseitem{42}{สรณํ ปริตฺตํ ธมฺมํ ปฏิจฺจ สรโณ ปริตฺโต ธมฺโม อุปฺปชฺชติ เหตุปจฺจยา.\csromanspacer{} ตีณิ.}

\prose{อรณํ ปริตฺตํ ธมฺมํ ปฏิจฺจ อรโณ ปริตฺโต ธมฺโม อุปฺปชฺชติ เหตุปจฺจยา. (1)}

\gambhira{ธมฺมานุโลม ทุกติกปฏฺฐานปาฬิ}
\prose{\csromanfolio{458}สรณํ ปริตฺตญฺจ อรณํ ปริตฺตญฺจ ธมฺมํ ปฏิจฺจ อรโณ ปริตฺโต ธมฺโม อุปฺปชฺชติ เหตุปจฺจยา. (1) (สํขิตฺตํ.)}

\prose{เหตุยา ปญฺจ, อารมฺมเณ เทฺว ฯเปฯ อวิคเต ปญฺจ. (สํขิตฺตํ.)}

\prose{(สหชาตวาเรปิ ฯเปฯ ปญฺหาวาเรปิ สพฺพตฺถ วิตฺถาเรตพฺโพ.)}

\proseitem{43}{อรณํ มหคฺคตํ ธมฺมํ ปฏิจฺจ อรโณ มหคฺคโต ธมฺโม อุปฺปชฺชติ เหตุปจฺจยา. (สํขิตฺตํ.)}

\prose{เหตุยา เอกํ, อารมฺมเณ เอกํ ฯเปฯ อวิคเต เอกํ. (สํขิตฺตํ.)}

\prose{(สหชาตวาเรปิ ฯเปฯ ปญฺหาวาเรปิ สพฺพตฺถ เอกํ.)}

\proseitem{44}{อรณํ อปฺปมาณํ ธมฺมํ ปฏิจฺจ อรโณ อปฺปมาโณ ธมฺโม อุปฺปชฺชติ เหตุปจฺจยา. (สํขิตฺตํ.)}

\prose{เหตุยา เอกํ, อารมฺมเณ เอกํ ฯเปฯ อวิคเต เอกํ. (สํขิตฺตํ.)}

\prose{(สหชาตวาเรปิ ฯเปฯ ปญฺหาวาเรปิ สพฺพตฺถ เอกํ.)}

\csromanheader{2}{100. สรณทุก 13. ปริตฺตารมฺมณตฺติก}
\csromanheader{2}{1. ปฏิจฺจวาร ปจฺจยจตุกฺก}
\proseitem{45}{สรณํ ปริตฺตารมฺมณํ ธมฺมํ ปฏิจฺจ สรโณ ปริตฺตารมฺมโณ ธมฺโม อุปฺปชฺชติ เหตุปจฺจยา. (1)}

\prose{อรณํ ปริตฺตารมฺมณํ ธมฺมํ ปฏิจฺจ อรโณ ปริตฺตารมฺมโณ ธมฺโม อุปฺปชฺชติ เหตุปจฺจยา. (1) (สํขิตฺตํ.)}

\prose{เหตุยา เทฺว, อารมฺมเณ เทฺว ฯเปฯ อวิคเต เทฺว. (สํขิตฺตํ.)}

\prose{(สหชาตวาโรปิ ฯเปฯ ปญฺหาวาโรปิ สพฺพตฺถ วิตฺถาเรตพฺพา.)}

\proseitem{46}{สรณํ มหคฺคตารมฺมณํ ธมฺมํ ปฏิจฺจ สรโณ มหคฺคตารมฺมโณ ธมฺโม อุปฺปชฺชติ เหตุปจฺจยา. (1)}

\prose{อรณํ มหคฺคตารมฺมณํ ธมฺมํ ปฏิจฺจ อรโณ มหคฺคตารมฺมโณ ธมฺโม อุปฺปชฺชติ เหตุปจฺจยา. (1) (สํขิตฺตํ.)}

\csromanheader{2}{100. สรณทุก 15. มิจฺฉตฺตตฺติก}
\prose{\csromanfolio{459}เหตุยา เทฺว, อารมฺมเณ เทฺว ฯเปฯ อวิคเต เทฺว. (สํขิตฺตํ.)}

\prose{(สหชาตวาโรปิ ฯเปฯ ปญฺหาวาโรปิ วิตฺถาเรตพฺพา.)}

\proseitem{47}{อรณํ อปฺปมาณารมฺมณํ ธมฺมํ ปฏิจฺจ อรโณ อปฺปมาณารมฺมโณ ธมฺโม อุปฺปชฺชติ เหตุปจฺจยา. (สํขิตฺตํ.)}

\prose{เหตุยา เอกํ, อารมฺมเณ เอกํ ฯเปฯ อวิคเต เอกํ. (สํขิตฺตํ.)}

\prose{(สหชาตวาเรปิ ฯเปฯ ปญฺหาวาเรปิ สพฺพตฺถ เอกํ.)}

\csromanheader{2}{100. สรณทุก 14. หีนตฺติก}
\csromanheader{2}{1. ปฏิจฺจวาร ปจฺจยจตุกฺก}
\proseitem{48}{สรณํ หีนํ ธมฺมํ ปฏิจฺจ สรโณ หีโน ธมฺโม อุปฺปชฺชติ เหตุปจฺจยา.}

\prose{เหตุยา เอกํ, อารมฺมเณ เอกํ ฯเปฯ อวิคเต เอกํ. (สํขิตฺตํ.)}

\prose{(สหชาตวาเรปิ ฯเปฯ ปญฺหาวาเรปิ สพฺพตฺถ เอกํ.)}

\proseitem{49}{อรณํ มชฺฌิมํ ธมฺมํ ปฏิจฺจ อรโณ มชฺฌิโม ธมฺโม อุปฺปชฺชติ เหตุปจฺจยา.}

\prose{เหตุยา เอกํ, อารมฺมเณ เอกํ ฯเปฯ อวิคเต เอกํ. (สํขิตฺตํ.)}

\prose{(สหชาตวาเรปิ ฯเปฯ ปญฺหาวาเรปิ สพฺพตฺถ เอกํ.)}

\proseitem{50}{อรณํ ปณีตํ ธมฺมํ ปฏิจฺจ อรโณ ปณีโต ธมฺโม อุปฺปชฺชติ เหตุปจฺจยา. (สํขิตฺตํ.)}

\prose{เหตุยา เอกํ, อารมฺมเณ เอกํ ฯเปฯ อวิคเต เอกํ. (สํขิตฺตํ.)}

\prose{(สหชาตวาเรปิ ฯเปฯ ปญฺหาวาเรปิ สพฺพตฺถ เอกํ.)}

\csromanheader{2}{100. สรณทุก 15. มิจฺฉตฺตตฺติก}
\csromanheader{2}{1. ปฏิจฺจวาร ปจฺจยจตุกฺก}
\proseitem{51}{สรณํ มิจฺฉตฺตนิยตํ ธมฺมํ ปฏิจฺจ สรโณ มิจฺฉตฺตนิยโต ธมฺโม อุปฺปชฺชติ เหตุปจฺจยา. (สํขิตฺตํ.)}

\gambhira{ธมฺมานุโลม ทุกติกปฏฺฐานปาฬิ}
\prose{\csromanfolio{460}เหตุยา เอกํ, อารมฺมเณ เอกํ ฯเปฯ อวิคเต เอกํ. (สํขิตฺตํ.)}

\prose{(สหชาตวาเรปิ ฯเปฯ ปญฺหาวาเรปิ สพฺพตฺถ เอกํ.)}

\proseitem{52}{อรณํ สมฺมตฺตนิยตํ ธมฺมํ ปฏิจฺจ อรโณ สมฺมตฺตนิยโต ธมฺโม อุปฺปชฺชติ เหตุปจฺจยา. (สํขิตฺตํ.)}

\prose{เหตุยา เอกํ, อารมฺมเณ เอกํ ฯเปฯ อวิคเต เอกํ. (สํขิตฺตํ.)}

\prose{(สหชาตวาเรปิ ฯเปฯ ปญฺหาวาเรปิ สพฺพตฺถ เอกํ.)}

\proseitem{53}{สรณํ อนิยตํ ธมฺมํ ปฏิจฺจ สรโณ อนิยโต ธมฺโม อุปฺปชฺชติ เหตุปจฺจยา.\csromanspacer{} ตีณิ.}

\prose{อรณํ อนิยตํ ธมฺมํ ปฏิจฺจ อรโณ อนิยโต ธมฺโม อุปฺปชฺชติ เหตุปจฺจยา. (1)}

\prose{สรณํ อนิยตญฺจ อรณํ อนิยตญฺจ ธมฺมํ ปฏิจฺจ อรโณ อนิยโต ธมฺโม อุปฺปชฺชติ เหตุปจฺจยา. (1) (สํขิตฺตํ.)}

\proseitem{54}{เหตุยา ปญฺจ, อารมฺมเณ เทฺว ฯเปฯ อวิคเต ปญฺจ. (สํขิตฺตํ.)}

\prose{(สหชาตวาโรปิ ฯเปฯ ปญฺหาวาโรปิ สพฺพตฺถ วิตฺถาเรตพฺพา.)}

\csromanheader{2}{100. สรณทุก 16. มคฺคารมฺมณตฺติก}
\csromanheader{2}{1. ปฏิจฺจวาร ปจฺจยจตุกฺก}
\proseitem{55}{อรณํ มคฺคารมฺมณํ ธมฺมํ ปฏิจฺจ อรโณ มคฺคารมฺมโณ ธมฺโม อุปฺปชฺชติ เหตุปจฺจยา. (สํขิตฺตํ.)}

\prose{เหตุยา เอกํ, อารมฺมเณ เอกํ ฯเปฯ อวิคเต เอกํ. (สํขิตฺตํ.)}

\prose{(สหชาตวาเรปิ ฯเปฯ ปญฺหาวาเรปิ สพฺพตฺถ เอกํ.)}

\proseitem{56}{อรณํ มคฺคเหตุกํ ธมฺมํ ปฏิจฺจ อรโณ มคฺคเหตุโก ธมฺโม อุปฺปชฺชติ เหตุปจฺจยา. (สํขิตฺตํ.)}

\prose{เหตุยา เอกํ, อารมฺมเณ เอกํ ฯเปฯ อวิคเต เอกํ. (สํขิตฺตํ.)}

\prose{(สหชาตวาเรปิ ฯเปฯ ปญฺหาวาเรปิ สพฺพตฺถ เอกํ.)}

\csromanheader{2}{100. สรณทุก 18. อตีตตฺติก}
\proseitem{57}{\csromanfolio{461}อรณํ มคฺคาธิปติํ ธมฺมํ ปฏิจฺจ อรโณ มคฺคาธิปติ ธมฺโม อุปฺปชฺชติ เหตุปจฺจยา. (สํขิตฺตํ.)}

\prose{เหตุยา เอกํ, อารมฺมเณ เอกํ ฯเปฯ อวิคเต เอกํ. (สํขิตฺตํ.)}

\prose{(สหชาตวาเรปิ ฯเปฯ ปญฺหาวาเรปิ สพฺพตฺถ เอกํ.)}

\csromanheader{2}{100. สรณทุก 17. อุปฺปนฺนตฺติก}
\csromanheader{2}{7. ปญฺหาวาร ปจฺจยจตุกฺก}
\proseitem{58}{สรโณ อุปฺปนฺโน ธมฺโม สรณสฺส อุปฺปนฺนสฺส ธมฺมสฺส เหตุปจฺจเยน ปจฺจโย.\csromanspacer{}\csromanspacer{}สรโณ อุปฺปนฺโน ธมฺโม อรณสฺส อุปฺปนฺนสฺส ธมฺมสฺส เหตุปจฺจเยน ปจฺจโย.\csromanspacer{}\csromanspacer{}สรโณ อุปฺปนฺโน ธมฺโม สรณสฺส อุปฺปนฺนสฺส จ อรณสฺส อุปฺปนฺนสฺส จ ธมฺมสฺส เหตุปจฺจเยน ปจฺจโย. (3)}

\prose{อรโณ อุปฺปนฺโน ธมฺโม อรณสฺส อุปฺปนฺนสฺส ธมฺมสฺส เหตุปจฺจเยน ปจฺจโย. (1)}

\prose{อรโณ อุปฺปนฺโน ธมฺโม อรณสฺส อุปฺปนฺนสฺส ธมฺมสฺส อารมฺมณปจฺจเยน ปจฺจโย. (สํขิตฺตํ.)}

\csromanhangingbegin
\hangingheaditem{59}{เหตุยา จตฺตาริ, อารมฺมเณ เทฺว ฯเปฯ อวิคเต สตฺต. (สํขิตฺตํ.)}
\hangingbody{(ยถา กุสลตฺติเก ปญฺหาวารํ, เอวํ วิตฺถาเรตพฺพํ.)}
\csromanhangingend

\csromanheader{2}{100. สรณทุก 18. อตีตตฺติก}
\csromanheader{2}{7. ปญฺหาวาร ปจฺจยจตุกฺก}
\proseitem{60}{สรโณ ปจฺจุปฺปนฺโน ธมฺโม สรณสฺส ปจฺจุปฺปนฺนสฺส ธมฺมสฺส เหตุปจฺจเยน ปจฺจโย.\csromanspacer{} ตีณิ.}

\gambhira{ธมฺมานุโลม ทุกติกปฏฺฐานปาฬิ}
\prose{\csromanfolio{462}อรโณ ปจฺจุปฺปนฺโน ธมฺโม อรณสฺส ปจฺจุปฺปนฺนสฺส ธมฺมสฺส เหตุปจฺจเยน ปจฺจโย. (1) (สํขิตฺตํ.)}

\prose{เหตุยา จตฺตาริ, อารมฺมเณ เทฺว ฯเปฯ อวิคเต สตฺต. (สํขิตฺตํ.)}

\prose{(ยถา กุสลตฺติเก ปญฺหาวารํ, เอวํ วิตฺถาเรตพฺพํ.)}

\csromanheader{2}{100. สรณทุก 19. อตีตารมฺมณตฺติก}
\csromanheader{2}{1. ปฏิจฺจวาร ปจฺจยจตุกฺก}
\proseitem{61}{สรณํ อตีตารมฺมณํ ธมฺมํ ปฏิจฺจ สรโณ อตีตารมฺมโณ ธมฺโม อุปฺปชฺชติ เหตุปจฺจยา. (1)}

\prose{อรณํ อตีตารมฺมณํ ธมฺมํ ปฏิจฺจ อรโณ อตีตารมฺมโณ ธมฺโม อุปฺปชฺชติ เหตุปจฺจยา. (1) (สํขิตฺตํ.)}

\prose{เหตุยา เทฺว, อารมฺมเณ เทฺว ฯเปฯ วิปาเก เอกํ ฯเปฯ อวิคเต เทฺว. (สํขิตฺตํ.\csromanspacer{} สหชาตวาเรปิ ฯเปฯ ปญฺหาวาเรปิ สพฺพตฺถ วิตฺถาเรตพฺโพ.)}

\proseitem{62}{สรณํ อนาคตารมฺมณํ ธมฺมํ ปฏิจฺจ สรโณ อนาคตารมฺมโณ ธมฺโม อุปฺปชฺชติ เหตุปจฺจยา. (1)}

\prose{อรณํ อนาคตารมฺมณํ ธมฺมํ ปฏิจฺจ อรโณ อนาคตารมฺมโณ ธมฺโม อุปฺปชฺชติ เหตุปจฺจยา. (1) (สํขิตฺตํ.)}

\prose{เหตุยา เทฺว, อารมฺมเณ เทฺว ฯเปฯ วิปาเก เอกํ ฯเปฯ อวิคเต เทฺว. (สํขิตฺตํ.\csromanspacer{} สหชาตวาเรปิ ฯเปฯ ปญฺหาวาเรปิ สพฺพตฺถ วิตฺถาเรตพฺโพ.)}

\proseitem{63}{สรณํ ปจฺจุปฺปนฺนารมฺมณํ ธมฺมํ ปฏิจฺจ สรโณ ปจฺจุปฺปนฺนารมฺมโณ ธมฺโม อุปฺปชฺชติ เหตุปจฺจยา. (1)}

\prose{อรณํ ปจฺจุปฺปนฺนารมฺมณํ ธมฺมํ ปฏิจฺจ อรโณ ปจฺจุปฺปนฺนารมฺมโณ ธมฺโม อุปฺปชฺชติ เหตุปจฺจยา. (1) (สํขิตฺตํ.)}

\prose{เหตุยา เทฺว, อารมฺมเณ เทฺว ฯเปฯ วิปาเก เอกํ ฯเปฯ อวิคเต เทฺว. (สํขิตฺตํ.\csromanspacer{} สหชาตวาเรปิ ฯเปฯ ปญฺหาวาเรปิ สพฺพตฺถ วิตฺถาเรตพฺโพ.)}

\vspace{\tipitakaparskip}
\csromancenter{สรณทุก 21.\csromanspacer{} อชฺฌตฺตารมฺมณตฺติก \textbf{100.}\csromanspacer{}\textbf{ สรณทุก 20.}\csromanspacer{}\textbf{ อชฺฌตฺตตฺติก}}
\csromanheader{2}{1. ปฏิจฺจวาร ปจฺจยจตุกฺก}
\proseitem{64}{\csromanfolio{463}สรณํ อชฺฌตฺตํ ธมฺมํ ปฏิจฺจ สรโณ อชฺฌตฺโต ธมฺโม อุปฺปชฺชติ เหตุปจฺจยา.\csromanspacer{} ตีณิ.}

\prose{อรณํ อชฺฌตฺตํ ธมฺมํ ปฏิจฺจ อรโณ อชฺฌตฺโต ธมฺโม อุปฺปชฺชติ เหตุปจฺจยา. (1)}

\prose{สรณํ อชฺฌตฺตญฺจ อรณํ อชฺฌตฺตญฺจ ธมฺมํ ปฏิจฺจ อรโณ อชฺฌตฺโต ธมฺโม อุปฺปชฺชติ เหตุปจฺจยา. (1) (สํขิตฺตํ.)}

\proseitem{65}{เหตุยา ปญฺจ, อรมฺมเณ เทฺว ฯเปฯ อวิคเต ปญฺจ. (สํขิตฺตํ.)}

\prose{(สหชาตวาเรปิ ฯเปฯ ปญฺหาวาเรปิ สพฺพตฺถ วิตฺถาเรตพฺโพ.)}

\proseitem{66}{สรณํ พหิทฺธา ธมฺมํ ปฏิจฺจ สรโณ พหิทฺธา ธมฺโม อุปฺปชฺชติ เหตุปจฺจยา.\csromanspacer{} ตีณิ.}

\prose{อรณํ พหิทฺธา ธมฺมํ ปฏิจฺจ อรโณ พหิทฺธา ธมฺโม อุปฺปชฺชติ เหตุปจฺจยา.}

\prose{สรณํ พหิทฺธา จ อรณํ หิทฺธา จ ธมฺมํ ปฏิจฺจ อรโณ พหิทฺธา ธมฺโม อุปฺปชฺชติ เหตุปจฺจยา. (1) (สํขิตฺตํ.)}

\prose{เหตุยา ปญฺจ, อารมฺมเณ เทฺว ฯเปฯ อวิคเต ปญฺจ. (สํขิตฺตํ.)}

\prose{(สหชาตวาเรปิ ฯเปฯ ปญฺหาวาเรปิ สพฺพตฺถ วิตฺถาเรตพฺโพ.)}

\csromanheader{2}{100. สรณทุก 21. อชฺฌตฺตารมฺมณตฺติก}
\csromanheader{2}{1. ปฏิจฺจวาร ปจฺจยจตุกฺก}
\proseitem{67}{สรณํ อชฺฌตฺตารมฺมณํ ธมฺมํ ปฏิจฺจ สรโณ อชฺฌตฺตารมฺมโณ ธมฺโม อุปฺปชฺชติ เหตุปจฺจยา. (1)}

\prose{อรณํ อชฺฌตฺตารมฺมณํ ธมฺมํ ปฏิจฺจ อรโณ อชฺฌตฺตารมฺมโณ ธมฺโม อุปฺปชฺชติ เหตุปจฺจยา. (1) (สํขิตฺตํ.)}

\gambhira{ธมฺมานุโลม ทุกติกปฏฺฐานปาฬิ}
\prose{\csromanfolio{464}เหตุยา เทฺว, อารมฺมเณ เทฺว ฯเปฯ อวิคเต เทฺว. (สํขิตฺตํ.)}

\prose{(สหชาตวาเรปิ ฯเปฯ ปญฺหาวาเรปิ สพฺพตฺถ วิตฺถาเรตพฺโพ.)}

\proseitem{68}{สรณํ พหิทฺธารมฺมณํ ธมฺมํ ปฏิจฺจ สรโณ พหิทฺธารมฺมโณ ธมฺโม อุปฺปชฺชติ เหตุปจฺจยา. (1)}

\prose{อรณํ พหิทฺธารมฺมณํ ธมฺมํ ปฏิจฺจ อรโณ พหิทฺธารมฺมโณ ธมฺโม อุปฺปชฺชติ เหตุปจฺจยา. (1) (สํขิตฺตํ.)}

\prose{เหตุยา เทฺว, อารมฺมเณ เทฺว ฯเปฯ อวิคเต เทฺว. (สํขิตฺตํ.)}

\prose{(สหชาตวาเรปิ ฯเปฯ ปญฺหาวาเรปิ สพฺพตฺถ วิตฺถาเรตพฺโพ.)}

\csromanmatikaanchor{mtk.90}
\csromantocmarkref{subparagraph}{22. สนิทสฺสนตฺติก}{mtk.90}
\bookmark[dest={mtk.90},level=5]{22. สนิทสฺสนตฺติก}
\csromanheader{6}{100. สรณทุก 22. สนิทสฺสนตฺติก}
\csromanheader{2}{1-7. ปฏิจฺจาทิวาร ปจฺจยจตุกฺก}
\proseitem{69}{อรณํ อนิทสฺสนสปฺปฏิฆํ ธมฺมํ ปฏิจฺจ อรโณ อนิทสฺสนสปฺปฏิโฆ ธมฺโม อุปฺปชฺชติ เหตุปจฺจยา. (สํขิตฺตํ.)}

\prose{เหตุยา เอกํ, อธิปติยา เอกํ ฯเปฯ อวิคเต เอกํ. (สํขิตฺตํ.)}

\prose{(สหชาตวาเรปิ ฯเปฯ ปญฺหาวาเรปิ สพฺพตฺถ เอกํ.)}

\proseitem{70}{สรณํ อนิทสฺสน-อปฺปฏิฆํ ธมฺมํ ปฏิจฺจ สรโณ อนิทสฺสน-อปฺปฏิโฆ ธมฺโม อุปฺปชฺชติ เหตุปจฺจยา.\csromanspacer{} ตีณิ.}

\prose{อรณํ อนิทสฺสน-อปฺปฏิฆํ ธมฺมํ ปฏิจฺจ อรโณ อนิทสฺสน-อปฺปฏิโฆ ธมฺโม อุปฺปชฺชติ เหตุปจฺจยา. (1)}

\prose{สรณํ อนิทสฺสน-อปฺปฏิฆญฺจ อรณํ อนิทสฺสน-อปฺปฏิฆญฺจ ธมฺมํ ปฏิจฺจ อรโณ อนิทสฺสน-อปฺปฏิโฆ ธมฺโม อุปฺปชฺชติ เหตุปจฺจยา. (1) (สํขิตฺตํ.)}

\prose{เหตุยา ปญฺจ, อารมฺมเณ เทฺว ฯเปฯ อวิคเต ปญฺจ. (สํขิตฺตํ.)}

\prose{(สหชาตวาโรปิ ฯเปฯ สมฺปยุตฺตวาโรปิ วิตฺถาเรตพฺพา.)}

\csromanheader{2}{100. สรณทุก 22. สนิทสฺสนตฺติก}
\proseitem{71}{\csromanfolio{465}อรโณ อนิทสฺสนสปฺปฏิโฆ ธมฺโม อรณสฺส อนิทสฺสนสปฺปฏิฆสฺส ธมฺมสฺส สหชาตปจฺจเยน ปจฺจโย.\csromanspacer{} เอกํ. (อญฺญมญฺเญ เอกํ, นิสฺสเย เอกํ, อตฺถิยา เอกํ, อวิคเต เอกํ.\csromanspacer{} สํขิตฺตํ.)}

\proseitem{72}{สรโณ อนิทสฺสน-อปฺปฏิโฆ ธมฺโม สรณสฺส อนิทสฺสน-อปฺปฏิฆสฺส ธมฺมสฺส เหตุปจฺจเยน ปจฺจโย.\csromanspacer{} ตีณิ.}

\prose{อรโณ อนิทสฺสน-อปฺปฏิโฆ ธมฺโม อรณสฺส อนิทสฺสน-อปฺปฏิฆสฺส ธมฺมสฺส เหตุปจฺจเยน ปจฺจโย. (1) (สํขิตฺตํ.)}

\proseitem{73}{เหตุยา จตฺตาริ, อารมฺมเณ จตฺตาริ ฯเปฯ อวิคเต สตฺต. (สํขิตฺตํ.)}

\prose{นเหตุยา สตฺต, น-อารมฺมเณ สตฺต. (สํขิตฺตํ.)}

\csromanheader{2}{เหตุปจฺจยา น-อารมฺมเณ จตฺตาริ. (สํขิตฺตํ.)}
\csromanheader{2}{นเหตุยา อารมฺมเณ จตฺตาริ. (สํขิตฺตํ.)}
\prose{(ยถา กุสลตฺติเก ปญฺหาวารสฺส อนุโลมมฺปิ ปจฺจนียมฺปิ อนุโลมปจฺจนียมฺปิ ปจฺจนียานุโลมมฺปิ คณิตํ, เอวํ คเณตพฺพํ.)}

\nitthitam{\textbf{ธมฺมานุโลเม ทุกติกปฏฺฐานํ นิฏฺฐิตํ.}}
\pitaka{\textbf{อภิธมฺมปิฏก}}
\csromanheader{2}{\textbf{ธมฺมานุโลม}}
\csromanmatikaanchor{mtk.91}
\csromantocmarknopage{chapter}{ติกทุกปฏฺฐานปาฬิ}{mtk.91}
\bookmark[dest={mtk.91},level=0]{ติกทุกปฏฺฐานปาฬิ}
\csromantocmarknopage{chapter}{1. กุสลตฺติก}{mtk.92}
\bookmark[dest={mtk.92},level=0]{1. กุสลตฺติก}
\gambhira{\textbf{ติกทุกปฏฺฐานปาฬิ}}
\namakkaram{นโม ตสฺส ภควโต อรหโต สมฺมาสมฺพุทฺธสฺส.}
\csromanmatikaanchor{mtk.93}
\csromantocmarkref{section}{1. เหตุทุก}{mtk.93}
\bookmark[dest={mtk.93},level=1]{1. เหตุทุก}
\csromanheader{1}{1. กุสลตฺติก 1. เหตุทุก}
\csromanheader{2}{1. เหตุปท 1. ปฏิจฺจวาร}
\csromanheader{2}{\textbf{ปจฺจยจตุกฺก-เหตุ อารมฺมณ}}
\proseitem{1}{\csromanfolio{467}กุสลํ เหตุํ ธมฺมํ ปฏิจฺจ กุสโล เหตุ ธมฺโม อุปฺปชฺชติ เหตุปจฺจยา. (1)}

\prose{อกุสลํ เหตุํ ธมฺมํ ปฏิจฺจ อกุสโล เหตุ ธมฺโม อุปฺปชฺชติ เหตุปจฺจยา. (1)}

\prose{อพฺยากตํ เหตุํ ธมฺมํ ปฏิจฺจ อพฺยากโต เหตุ ธมฺโม อุปฺปชฺชติ เหตุปจฺจยา. (1)}

\proseitem{2}{กุสลํ เหตุํ ธมฺมํ ปฏิจฺจ กุสโล เหตุ ธมฺโม อุปฺปชฺชติ อารมฺมณปจฺจยา. (1)}

\prose{อกุสลํ เหตุํ ธมฺมํ ปฏิจฺจ อกุสโล เหตุ ธมฺโม อุปฺปชฺชติ อารมฺมณปจฺจยา. (1)}

\prose{อพฺยากตํ เหตุํ ธมฺมํ ปฏิจฺจ อพฺยากโต เหตุ ธมฺโม อุปฺปชฺชติ อารมฺมณปจฺจยา. (1) (สํขิตฺตํ.)}

\proseitem{3}{เหตุยา ตีณิ, อารมฺมเณ ตีณิ, อธิปติยา ตีณิ, อนนฺตเร ตีณิ, สมนนฺตเร ตีณิ, สหชาเต ตีณิ, อญฺญมญฺเญ ตีณิ, นิสฺสเย}

\vspace{\tipitakaparskip}
\csromancenter{ธมฺมานุโลม ติกทุกปฏฺฐานปาฬิ ตีณิ, อุปนิสฺสเย ตีณิ, ปุเรชาเต ตีณิ, อาเสวเน ตีณิ, กมฺเม ตีณิ, วิปาเก เอกํ, อาหาเร ตีณิ ฯเปฯ อวิคเต ตีณิ. (สํขิตฺตํ.)}
\csromanheader{2}{\textbf{ปจฺจนีย-น-อธิปติ}}
\proseitem{4}{\csromanfolio{468}กุสลํ เหตุํ ธมฺมํ ปฏิจฺจ กุสโล เหตุ ธมฺโม อุปฺปชฺชติ น-อธิปติปจฺจยา. (1)}

\prose{อกุสลํ เหตุํ ธมฺมํ ปฏิจฺจ อกุสโล เหตุ ธมฺโม อุปฺปชฺชติ น-อธิปติปจฺจยา. (1)}

\prose{อพฺยากตํ เหตุํ ธมฺมํ ปฏิจฺจ อพฺยากโต เหตุ ธมฺโม อุปฺปชฺชติ น-อธิปติปจฺจยา. (1) (สํขิตฺตํ.)}

\proseitem{5}{น-อธิปติยา ตีณิ, นปุเรชาเต ตีณิ, นปจฺฉาชาเต ตีณิ, น-อาเสวเน ตีณิ, นวิปาเก ตีณิ, นวิปฺปยุตฺเต ตีณิ. (สํขิตฺตํ.)}

\csromanheader{2}{เหตุปจฺจยา น-อธิปติยา ตีณิ. (สํขิตฺตํ.)}
\csromanheader{2}{น-อธิปติปจฺจยา เหตุยา ตีณิ. (สํขิตฺตํ.)}
\prose{(สหชาตวาเรปิ ปจฺจยวาเรปิ นิสฺสยวาเรปิ สํสฏฺฐวาเรปิ สมฺปยุตฺตวาเรปิ สพฺพตฺถ ตีณิ.)}

\csromanheader{2}{7. ปญฺหาวาร ปจฺจยจตุกฺก}
\csromanheader{2}{\textbf{เหตุ-อารมฺมณาทิ}}
\proseitem{6}{กุสโล เหตุ ธมฺโม กุสลสฺส เหตุสฺส ธมฺมสฺส เหตุปจฺจเยน ปจฺจโย. (1)}

\prose{อกุสโล เหตุ ธมฺโม อกุสลสฺส เหตุสฺส ธมฺมสฺส เหตุปจฺจเยน ปจฺจโย. (1)}

\prose{อพฺยากโต เหตุ ธมฺโม อพฺยากตสฺส เหตุสฺส ธมฺมสฺส เหตุปจฺจเยน ปจฺจโย. (1)}

\proseitem{7}{กุสโล เหตุ ธมฺโม กุสลสฺส เหตุสฺส ธมฺมสฺส อารมฺมณปจฺจเยน ปจฺจโย.\csromanspacer{}\csromanspacer{}กุสโล เหตุ ธมฺโม อกุสลสฺส เหตุสฺส}

\vspace{\tipitakaparskip}
\csromancenter{กุสลตฺติก 1.\csromanspacer{} เหตุทุก ธมฺมสฺส อารมฺมณปจฺจเยน ปจฺจโย.\csromanspacer{}\csromanspacer{}กุสโล เหตุ ธมฺโม อพฺยากตสฺส เหตุสฺส ธมฺมสฺส อารมฺมณปจฺจเยน ปจฺจโย. (3)}
\prose{\csromanfolio{469}อกุสโล เหตุ ธมฺโม อกุสลสฺส เหตุสฺส ธมฺมสฺส อารมฺมณปจฺจเยน ปจฺจโย.\csromanspacer{} ตีณิ.}

\prose{อพฺยากโต เหตุ ธมฺโม อพฺยากตสฺส เหตุสฺส ธมฺมสฺส อารมฺมณปจฺจเยน ปจฺจโย.\csromanspacer{} ตีณิ.}

\proseitem{8}{กุสโล เหตุ ธมฺโม กุสลสฺส เหตุสฺส ธมฺมสฺส อธิปติปจฺจเยน ปจฺจโย.\csromanspacer{}\csromanspacer{}กุสโล เหตุ ธมฺโม อกุสลสฺส เหตุสฺส ธมฺมสฺส อธิปติปจฺจเยน ปจฺจโย.\csromanspacer{}\csromanspacer{}กุสโล เหตุ ธมฺโม อพฺยากตสฺส เหตุสฺส ธมฺมสฺส อธิปติปจฺจเยน ปจฺจโย.\csromanspacer{} ตีณิ.}

\prose{อกุสโล เหตุ ธมฺโม อกุสลสฺส เหตุสฺส ธมฺมสฺส อธิปติปจฺจเยน ปจฺจโย. (1)}

\prose{อพฺยากโต เหตุ ธมฺโม อพฺยากตสฺส เหตุสฺส ธมฺมสฺส อธิปติปจฺจเยน ปจฺจโย.\csromanspacer{}\csromanspacer{}อพฺยากโต เหตุ ธมฺโม กุสลสฺส เหตุสฺส ธมฺมสฺส อธิปติปจฺจเยน ปจฺจโย.\csromanspacer{}\csromanspacer{}อพฺยากโต เหตุ ธมฺโม อกุสลสฺส เหตุสฺส ธมฺมสฺส อธิปติปจฺจเยน ปจฺจโย.\csromanspacer{} ตีณิ.}

\csromanheader{2}{\textbf{อนนฺตราทิ}}
\proseitem{9}{กุสโล เหตุ ธมฺโม กุสลสฺส เหตุสฺส ธมฺมสฺส อนนฺตรปจฺจเยน ปจฺจโย.\csromanspacer{} เทฺว.}

\prose{อกุสโล เหตุ ธมฺโม อกุสลสฺส เหตุสฺส ธมฺมสฺส อนนฺตรปจฺจเยน ปจฺจโย.\csromanspacer{} เทฺว.}

\prose{อพฺยากโต เหตุ ธมฺโม อพฺยากตสฺส เหตุสฺส ธมฺมสฺส อนนฺตรปจฺจเยน ปจฺจโย. (1)}

\proseitem{10}{กุสโล เหตุ ธมฺโม กุสลสฺส เหตุสฺส ธมฺมสฺส สหชาตปจฺจเยน ปจฺจโย. (1)}

\prose{อกุสโล เหตุ ธมฺโม อกุสลสฺส เหตุสฺส ธมฺมสฺส สหชาตปจฺจเยน ปจฺจโย. (1)}

\gambhira{ธมฺมานุโลม ติกทุกปฏฺฐานปาฬิ}
\prose{\csromanfolio{470}อพฺยากโต เหตุ ธมฺโม อพฺยากตสฺส เหตุสฺส ธมฺมสฺส สหชาตปจฺจเยน ปจฺจโย. (1) (อญฺญมญฺญนิสฺสยปจฺจยา สหชาตปจฺจยสทิสา.)}

\csromanheader{2}{\textbf{อุปนิสฺสยาทิ}}
\proseitem{11}{กุสโล เหตุ ธมฺโม กุสลสฺส เหตุสฺส ธมฺมสฺส อุปนิสฺสยปจฺจเยน ปจฺจโย.\csromanspacer{}\csromanspacer{}กุสโล เหตุ ธมฺโม อกุสลสฺส เหตุสฺส ธมฺมสฺส อุปนิสฺสยปจฺจเยน ปจฺจโย.\csromanspacer{}\csromanspacer{}กุสโล เหตุ ธมฺโม อพฺยากตสฺส เหตุสฺส ธมฺมสฺส อุปนิสฺสยปจฺจเยน ปจฺจโย. (3)}

\prose{อกุสโล เหตุ ธมฺโม อกุสลสฺส เหตุสฺส ธมฺมสฺส อุปนิสฺสยปจฺจเยน ปจฺจโย.\csromanspacer{}\csromanspacer{}อกุสโล เหตุ ธมฺโม กุสลสฺส เหตุสฺส ธมฺมสฺส อุปนิสฺสยปจฺจเยน ปจฺจโย.\csromanspacer{}\csromanspacer{}อกุสโล เหตุ ธมฺโม อพฺยากตสฺส เหตุสฺส ธมฺมสฺส อุปนิสฺสยปจฺจเยน ปจฺจโย.\csromanspacer{} ตีณิ.}

\prose{อพฺยากโต เหตุ ธมฺโม อพฺยากตสฺส เหตุสฺส ธมฺมสฺส อุปนิสฺสยปจฺจเยน ปจฺจโย.\csromanspacer{}\csromanspacer{}อพฺยากโต เหตุ ธมฺโม กุสลสฺส เหตุสฺส ธมฺมสฺส อุปนิสฺสยปจฺจเยน ปจฺจโย.\csromanspacer{}\csromanspacer{}อพฺยากโต เหตุ ธมฺโม อกุสลสฺส เหตุสฺส ธมฺมสฺส อุปนิสฺสยปจฺจเยน ปจฺจโย. (3)}

\prose{อพฺยากโต เหตุ ธมฺโม อพฺยากตสฺส เหตุสฺส ธมฺมสฺส วิปากปจฺจเยน ปจฺจโย. (1) (สํขิตฺตํ.)}

\proseitem{12}{เหตุยา ตีณิ, อารมฺมเณ นว, อธิปติยา สตฺต, อนนฺตเร ปญฺจ, สมนนฺตเร ปญฺจ, สหชาเต ตีณิ, อญฺญมญฺเญ ตีณิ, นิสฺสเย ตีณิ, อุปนิสฺสเย นว, อาเสวเน ตีณิ, วิปาเก เอกํ, อินฺทฺริเย เทฺว, มคฺเค เทฺว, สมฺปยุตฺเต ตีณิ, อตฺถิยา ตีณิ, นตฺถิยา ปญฺจ, วิคเต ปญฺจ, อวิคเต ตีณิ. (สํขิตฺตํ.)}

\csromanheader{2}{\textbf{ปจฺจนียุทฺธาร}}
\proseitem{13}{กุสโล เหตุ ธมฺโม กุสลสฺส เหตุสฺส ธมฺมสฺส อารมฺมณปจฺจเยน ปจฺจโย.\csromanspacer{} สหชาตปจฺจเยน ปจฺจโย.\csromanspacer{} อุปนิสฺสยปจฺจเยน ปจฺจโย.\csromanspacer{}\csromanspacer{}กุสโล เหตุ ธมฺโม อกุสลสฺส เหตุสฺส ธมฺมสฺส อารมฺมณปจฺจเยน ปจฺจโย.\csromanspacer{} อุปนิสฺสยปจฺจเยน ปจฺจโย.\csromanspacer{}\csromanspacer{}}

\vspace{\tipitakaparskip}
\csromancenter{กุสลตฺติก 1.\csromanspacer{} เหตุทุก กุสโล เหตุ ธมฺโม อพฺยากตสฺส เหตุสฺส ธมฺมสฺส อารมฺมณปจฺจเยน ปจฺจโย.\csromanspacer{} อุปนิสฺสยปจฺจเยน ปจฺจโย. (3)}
\prose{\csromanfolio{471}อกุสโล เหตุ ธมฺโม อกุสลสฺส เหตุสฺส ธมฺมสฺส อารมฺมณปจฺจเยน ปจฺจโย.\csromanspacer{} สหชาตปจฺจเยน ปจฺจโย.\csromanspacer{} อุปนิสฺสยปจฺจเยน ปจฺจโย.\csromanspacer{}\csromanspacer{}อกุสโล เหตุ ธมฺโม กุสลสฺส เหตุสฺส ธมฺมสฺส อารมฺมณปจฺจเยน ปจฺจโย.\csromanspacer{} อุปนิสฺสยปจฺจเยน ปจฺจโย.\csromanspacer{}\csromanspacer{}อกุสโล เหตุ ธมฺโม อพฺยากตสฺส เหตุสฺส ธมฺมสฺส อารมฺมณปจฺจเยน ปจฺจโย.\csromanspacer{} อุปนิสฺสยปจฺจเยน ปจฺจโย. (3)}

\prose{อพฺยากโต เหตุ ธมฺโม อพฺยากตสฺส เหตุสฺส ธมฺมสฺส อารมฺมณปจฺจเยน ปจฺจโย.\csromanspacer{} สหชาตปจฺจเยน ปจฺจโย.\csromanspacer{} อุปนิสฺสยปจฺจเยน ปจฺจโย. (สํขิตฺตํ.)}

\proseitem{14}{นเหตุยา นว, น-อารมฺมเณ นว. (สํขิตฺตํ.)}

\csromanheader{2}{เหตุปจฺจยา น-อารมฺมเณ ตีณิ. (สํขิตฺตํ.)}
\csromanheader{2}{นเหตุปจฺจยา อารมฺมเณ นว. (สํขิตฺตํ.)}
\prose{(ยถา กุสลตฺติเก ปญฺหาวารสฺส อนุโลมมฺปิ ปจฺจนียมฺปิ อนุโลมปจฺจนียมฺปิ ปจฺจนียานุโลมมฺปิ คณิตํ, เอวํ คเณตพฺพํ.)}

\csromanheader{2}{2. นเหตุปท 1. ปฏิจฺจวาร}
\csromanheader{2}{\textbf{ปจฺจยจตุกฺก-เหตุ อารมฺมณ}}
\proseitem{15}{กุสลํ นเหตุํ ธมฺมํ ปฏิจฺจ กุสโล นเหตุ ธมฺโม อุปฺปชฺชติ เหตุปจฺจยา.\csromanspacer{}\csromanspacer{}กุสลํ นเหตุํ ธมฺมํ ปฏิจฺจ อพฺยากโต นเหตุ ธมฺโม อุปฺปชฺชติ เหตุปจฺจยา.\csromanspacer{}\csromanspacer{}กุสลํ นเหตุํ ธมฺมํ ปฏิจฺจ กุสโล นเหตุ จ อพฺยากโต นเหตุ จ ธมฺมา อุปฺปชฺชนฺติ เหตุปจฺจยา. (3)}

\prose{อกุสลํ นเหตุํ ธมฺมํ ปฏิจฺจ อกุสโล นเหตุ ธมฺโม อุปฺปชฺชติ เหตุปจฺจยา.\csromanspacer{}\csromanspacer{}อกุสลํ นเหตุํ ธมฺมํ ปฏิจฺจ อพฺยากโต นเหตุ ธมฺโม อุปฺปชฺชติ เหตุปจฺจยา.\csromanspacer{}\csromanspacer{}อกุสลํ นเหตุํ ธมฺมํ ปฏิจฺจ อกุสโล นเหตุ จ อพฺยากโต นเหตุ จ ธมฺมา อุปฺปชฺชนฺติ เหตุปจฺจยา. (3)}

\gambhira{ธมฺมานุโลม ติกทุกปฏฺฐานปาฬิ}
\prose{\csromanfolio{472}อพฺยากตํ นเหตุํ ธมฺมํ ปฏิจฺจ อพฺยากโต นเหตุ ธมฺโม อุปฺปชฺชติ เหตุปจฺจยา. (1)}

\prose{กุสลํ นเหตุญฺจ อพฺยากตํ นเหตุญฺจ ธมฺมํ ปฏิจฺจ อพฺยากโต นเหตุ ธมฺโม อุปฺปชฺชติ เหตุปจฺจยา. (1)}

\prose{อกุสลํ นเหตุญฺจ อพฺยากตํ นเหตุญฺจ ธมฺมํ ปฏิจฺจ อพฺยากโต นเหตุ ธมฺโม อุปฺปชฺชติ เหตุปจฺจยา. (1)}

\proseitem{16}{กุสลํ นเหตุํ ธมฺมํ ปฏิจฺจ กุสโล นเหตุ ธมฺโม อุปฺปชฺชติ อารมฺมณปจฺจยา. (1)}

\prose{อกุสลํ นเหตุํ ธมฺมํ ปฏิจฺจ อกุสโล นเหตุ ธมฺโม อุปฺปชฺชติ อารมฺมณปจฺจยา. (1)}

\prose{อพฺยากตํ นเหตุํ ธมฺมํ ปฏิจฺจ อพฺยากโต นเหตุ ธมฺโม อุปฺปชฺชติ อารมฺมณปจฺจยา. (1) (สํขิตฺตํ.)}

\proseitem{17}{เหตุยา นว, อารมฺมเณ ตีณิ, อธิปติยา นว, อนนฺตเร ตีณิ, สมนนฺตเร ตีณิ, สหชาเต นว, อญฺญมญฺเญ ตีณิ, นิสฺสเย นว, อุปนิสฺสเย ตีณิ, ปุเรชาเต ตีณิ, อาเสวเน ตีณิ, กมฺเม นว, วิปาเก เอกํ, อาหาเร นว, อินฺทฺริเย นว, ฌาเน นว, มคฺเค นว, สมฺปยุตฺเต ตีณิ, วิปฺปยุตฺเต นว, อตฺถิยา นว, นตฺถิยา ตีณิ, วิคเต ตีณิ, อวิคเต นว. (สํขิตฺตํ.)}

\csromanheader{2}{\textbf{นเหตุ น-อารมฺมณ}}
\proseitem{18}{อพฺยากตํ นเหตุํ ธมฺมํ ปฏิจฺจ อพฺยากโต นเหตุ ธมฺโม อุปฺปชฺชติ นเหตุปจฺจยา. (1)}

\prose{กุสลํ นเหตุํ ธมฺมํ ปฏิจฺจ อพฺยากโต นเหตุ ธมฺโม อุปฺปชฺชติ น-อารมฺมณปจฺจยา. (1)}

\prose{อกุสลํ นเหตุํ ธมฺมํ ปฏิจฺจ อพฺยากโต นเหตุ ธมฺโม อุปฺปชฺชติ น-อารมฺมณปจฺจยา. (1)}

\prose{อพฺยากตํ นเหตุํ ธมฺมํ ปฏิจฺจ อพฺยากโต นเหตุ ธมฺโม อุปฺปชฺชติ น-อารมฺมณปจฺจยา. (1)}

\csromanheader{2}{1. กุสลตฺติก 1. เหตุทุก}
\prose{\csromanfolio{473}กุสลํ นเหตุญฺจ อพฺยากตํ นเหตุญฺจ ธมฺมํ ปฏิจฺจ อพฺยากโต นเหตุ ธมฺโม อุปฺปชฺชติ น-อารมฺมณปจฺจยา. (1)}

\prose{อกุสลํ นเหตุญฺจ อพฺยากตํ นเหตุญฺจ ธมฺมํ ปฏิจฺจ อพฺยากโต นเหตุ ธมฺโม อุปฺปชฺชติ น-อารมฺมณปจฺจยา. (1) (สํขิตฺตํ.)}

\proseitem{19}{นเหตุยา เอกํ, น-อารมฺมเณ ปญฺจ, น-อธิปติยา นว, น-อนนฺตเร ปญฺจ, นสมนนฺตเร ปญฺจ, น-อญฺญมญฺเญ ปญฺจ, น-อุปนิสฺสเย ปญฺจ, นปุเรชาเต สตฺต, นปจฺฉาชาเต นว, น-อาเสวเน นว, นกมฺเม ตีณิ, นวิปาเก นว, น-อาหาเร เอกํ, น-อินฺทฺริเย เอกํ, นฌาเน เอกํ, นมคฺเค เอกํ, นสมฺปยุตฺเต ปญฺจ, นวิปฺปยุตฺเต ตีณิ, โนนตฺถิยา ปญฺจ, โนวิคเต ปญฺจ. (สํขิตฺตํ.)}

\csromanheader{2}{เหตุปจฺจยา น-อารมฺมเณ ปญฺจ. (สํขิตฺตํ.)}
\csromanheader{2}{นเหตุปจฺจยา อารมฺมเณ เอกํ. (สํขิตฺตํ.)}
\prose{(สหชาตวารมฺปิ ปจฺจยวารมฺปิ นิสฺสยวารมฺปิ สํสฏฺฐวารมฺปิ สมฺปยุตฺตวารมฺปิ ปฏิจฺจวารสทิสา วิตฺถาเรตพฺพา.)}

\csromanheader{2}{7. ปญฺหาวาร ปจฺจยจตุกฺก-อารมฺมณาทิ}
\proseitem{20}{กุสโล นเหตุ ธมฺโม กุสลสฺส นเหตุสฺส ธมฺมสฺส อารมฺมณปจฺจเยน ปจฺจโย.\csromanspacer{}\csromanspacer{}กุสโล นเหตุ ธมฺโม อกุสลสฺส นเหตุสฺส ธมฺมสฺส อารมฺมณปจฺจเยน ปจฺจโย.\csromanspacer{}\csromanspacer{}กุสโล นเหตุ ธมฺโม อพฺยากตสฺส นเหตุสฺส ธมฺมสฺส อารมฺมณปจฺจเยน ปจฺจโย. (3)}

\prose{อกุสโล นเหตุ ธมฺโม อกุสลสฺส นเหตุสฺส ธมฺมสฺส อารมฺมณปจฺจเยน ปจฺจโย.\csromanspacer{}\csromanspacer{}อกุสโล นเหตุ ธมฺโม กุสลสฺส นเหตุสฺส ธมฺมสฺส อารมฺมณปจฺจเยน ปจฺจโย.\csromanspacer{}\csromanspacer{}อกุสโล นเหตุ ธมฺโม อพฺยากตสฺส นเหตุสฺส ธมฺมสฺส อารมฺมณปจฺจเยน ปจฺจโย. (3)}

\prose{อพฺยากโต นเหตุ ธมฺโม อพฺยากตสฺส นเหตุสฺส ธมฺมสฺส อารมฺมณปจฺจเยน ปจฺจโย.\csromanspacer{}\csromanspacer{}อพฺยากโต นเหตุ ธมฺโม กุสลสฺส นเหตุสฺส ธมฺมสฺส อารมฺมณปจฺจเยน ปจฺจโย.\csromanspacer{}\csromanspacer{}อพฺยากโต นเหตุ ธมฺโม อกุสลสฺส นเหตุสฺส ธมฺมสฺส อารมฺมณปจฺจเยน ปจฺจโย. (3)}

\gambhira{ธมฺมานุโลม ติกทุกปฏฺฐานปาฬิ}
\proseitem{21}{\csromanfolio{474}กุสโล นเหตุ ธมฺโม กุสลสฺส นเหตุสฺส ธมฺมสฺส อธิปติปจฺจเยน ปจฺจโย.\csromanspacer{}\csromanspacer{}กุสโล นเหตุ ธมฺโม อกุสลสฺส นเหตุสฺส ธมฺมสฺส อธิปติปจฺจเยน ปจฺจโย.\csromanspacer{}\csromanspacer{}กุสโล นเหตุ ธมฺโม อพฺยากตสฺส นเหตุสฺส ธมฺมสฺส อธิปติปจฺจเยน ปจฺจโย.\csromanspacer{}\csromanspacer{}กุสโล นเหตุ ธมฺโม กุสลสฺส นเหตุสฺส จ อพฺยากตสฺส นเหตุสฺส จ ธมฺมสฺส อธิปติปจฺจเยน ปจฺจโย.\csromanspacer{} จตฺตาริ.}

\prose{อกุสโล นเหตุ ธมฺโม อกุสลสฺส นเหตุสฺส ธมฺมสฺส อธิปติปจฺจเยน ปจฺจโย.\csromanspacer{} ตีณิ.}

\prose{อพฺยากโต นเหตุ ธมฺโม อพฺยากตสฺส นเหตุสฺส ธมฺมสฺส อธิปติปจฺจเยน ปจฺจโย.\csromanspacer{}\csromanspacer{}อพฺยากโต นเหตุ ธมฺโม กุสลสฺส นเหตุสฺส ธมฺมสฺส อธิปติปจฺจเยน ปจฺจโย.\csromanspacer{} ตีณิ.}

\proseitem{22}{กุสโล นเหตุ ธมฺโม กุสลสฺส นเหตุสฺส ธมฺมสฺส อนนฺตรปจฺจเยน ปจฺจโย.\csromanspacer{}\csromanspacer{}กุสโล นเหตุ ธมฺโม อพฺยากตสฺส นเหตุสฺส ธมฺมสฺส อนนฺตรปจฺจเยน ปจฺจโย. (2)}

\prose{อกุสโล นเหตุ ธมฺโม อกุสลสฺส นเหตุสฺส ธมฺมสฺส อนนฺตรปจฺจเยน ปจฺจโย.\csromanspacer{}\csromanspacer{}อกุสโล นเหตุ ธมฺโม อพฺยากตสฺส นเหตุสฺส ธมฺมสฺส อนนฺตรปจฺจเยน ปจฺจโย. (2)}

\prose{อพฺยากโต นเหตุ ธมฺโม อพฺยากตสฺส นเหตุสฺส ธมฺมสฺส อนนฺตรปจฺจเยน ปจฺจโย.\csromanspacer{} ตีณิ ฯเปฯ.}

\csromanheader{2}{\textbf{อุปนิสฺสย}}
\proseitem{23}{กุสโล นเหตุ ธมฺโม กุสลสฺส นเหตุสฺส ธมฺมสฺส อุปนิสฺสยปจฺจเยน ปจฺจโย.\csromanspacer{}\csromanspacer{}กุสโล นเหตุ ธมฺโม อกุสลสฺส นเหตุสฺส ธมฺมสฺส อุปนิสฺสยปจฺจเยน ปจฺจโย.\csromanspacer{}\csromanspacer{}กุสโล นเหตุ ธมฺโม อพฺยากตสฺส นเหตุสฺส ธมฺมสฺส อุปนิสฺสยปจฺจเยน ปจฺจโย. (3)}

\prose{อกุสโล นเหตุ ธมฺโม อกุสลสฺส นเหตุสฺส ธมฺมสฺส อุปนิสฺสยปจฺจเยน ปจฺจโย.\csromanspacer{}\csromanspacer{}อกุสโล นเหตุ ธมฺโม กุสลสฺส นเหตุสฺส ธมฺมสฺส อุปนิสฺสยปจฺจเยน ปจฺจโย.\csromanspacer{}\csromanspacer{}อกุสโล นเหตุ ธมฺโม อพฺยากตสฺส นเหตุสฺส ธมฺมสฺส อุปนิสฺสยปจฺจเยน ปจฺจโย. (3)}

\csromanmatikaanchor{mtk.94}
\csromanmatikaanchor{mtk.95}
\csromantocmarknopage{section}{2. เวทนาตฺติก}{mtk.94}
\bookmark[dest={mtk.94},level=1]{2. เวทนาตฺติก}
\csromantocmarkref{subsection}{1. เหตุทุก}{mtk.95}
\bookmark[dest={mtk.95},level=2]{1. เหตุทุก}
\csromanheader{2}{2. เวทนาตฺติก 1. เหตุทุก}
\prose{\csromanfolio{475}อพฺยากโต นเหตุ ธมฺโม อพฺยากตสฺส นเหตุสฺส ธมฺมสฺส อุปนิสฺสยปจฺจเยน ปจฺจโย.\csromanspacer{}\csromanspacer{}อพฺยากโต นเหตุ ธมฺโม กุสลสฺส นเหตุสฺส ธมฺมสฺส อุปนิสฺสยปจฺจเยน ปจฺจโย.\csromanspacer{}\csromanspacer{}อพฺยากโต นเหตุ ธมฺโม อกุสลสฺส นเหตุสฺส ธมฺมสฺส อุปนิสฺสยปจฺจเยน ปจฺจโย. (3) (สํขิตฺตํ.)}

\proseitem{24}{อารมฺมเณ นว, อธิปติยา ทส, อนนฺตเร สตฺต, สมนนฺตเร สตฺต, สหชาเต นว, อญฺญมญฺเญ ตีณิ, นิสฺสเย เตรส, อุปนิสฺสเย นว, ปุเรชาเต ตีณิ, ปจฺฉาชาเต ตีณิ, อาเสวเน ตีณิ, กมฺเม สตฺต, วิปาเก เอกํ, อาหาเร สตฺต, อินฺทฺริเย สตฺต, ฌาเน สตฺต, มคฺเค สตฺต, สมฺปยุตฺเต ตีณิ, วิปฺปยุตฺเต ปญฺจ, อตฺถิยา เตรส, นตฺถิยา สตฺต, วิคเต สตฺต, อวิคเต เตรส. (สํขิตฺตํ.)}

\prose{นเหตุยา ปนฺนรส, น-อารมฺมเณ ปนฺนรส. (สํขิตฺตํ.)}

\csromanheader{2}{อารมฺมณปจฺจยา นเหตุยา นว. (สํขิตฺตํ.)}
\csromanheader{2}{นเหตุปจฺจยา อารมฺมเณ นว. (สํขิตฺตํ.)}
\prose{(ยถา กุสลตฺติเก ปญฺหาวารํ, เอวํ วิตฺถาเรตพฺพํ.)}

\csromanheader{2}{2. เวทนาตฺติก 1. เหตุทุก}
\csromanheader{2}{1. ปฏิจฺจาทิวาร ปจฺจยจตุกฺก}
\proseitem{25}{สุขาย เวทนาย สมฺปยุตฺตํ เหตุํ ธมฺมํ ปฏิจฺจ สุขาย เวทนาย สมฺปยุตฺโต เหตุ ธมฺโม อุปฺปชฺชติ เหตุปจฺจยา. (1)}

\prose{ทุกฺขาย เวทนาย สมฺปยุตฺตํ เหตุํ ธมฺมํ ปฏิจฺจ ทุกฺขาย เวทนาย สมฺปยุตฺโต เหตุ ธมฺโม อุปฺปชฺชติ เหตุปจฺจยา. (1)}

\prose{อทุกฺขมสุขาย เวทนาย สมฺปยุตฺตํ เหตุํ ธมฺมํ ปฏิจฺจ อทุกฺขมสุขาย เวทนาย สมฺปยุตฺโต เหตุ ธมฺโม อุปฺปชฺชติ เหตุปจฺจยา. (1) (สํขิตฺตํ.)}

\proseitem{26}{เหตุยา ตีณิ, อารมฺมเณ ตีณิ, อธิปติยา ตีณิ ฯเปฯ วิปาเก เทฺว ฯเปฯ อวิคเต ตีณิ. (สํขิตฺตํ.)}

\gambhira{ธมฺมานุโลม ติกทุกปฏฺฐานปาฬิ}
\prose{\csromanfolio{476}น-อธิปติยา ตีณิ, นปุเรชาเต เทฺว, นปจฺฉาชาเต ตีณิ, น-อาเสวเน ตีณิ, นวิปาเก ตีณิ, นวิปฺปยุตฺเต เทฺว. (สํขิตฺตํ.)}

\csromanheader{2}{เหตุปจฺจยา น-อธิปติยา ตีณิ. (สํขิตฺตํ.)}
\prose{น-อธิปติปจฺจยา เหตุยา ตีณิ. (สํขิตฺตํ.) (สหชาตวารมฺปิ ฯเปฯ สมฺปยุตฺตวารมฺปิ วิตฺถาเรตพฺพํ.)}

\csromanheader{2}{\textbf{เหตุ อารมฺมณ}}
\proseitem{27}{สุขาย เวทนาย สมฺปยุตฺโต เหตุ ธมฺโม สุขาย เวทนาย สมฺปยุตฺตสฺส เหตุสฺส ธมฺมสฺส เหตุปจฺจเยน ปจฺจโย. (1)}

\prose{ทุกฺขาย เวทนาย สมฺปยุตฺโต เหตุ ธมฺโม ทุกฺขาย เวทนาย สมฺปยุตฺตสฺส เหตุสฺส ธมฺมสฺส เหตุปจฺจเยน ปจฺจโย. (1)}

\prose{อทุกฺขมสุขาย เวทนาย สมฺปยุตฺโต เหตุ ธมฺโม อทุกฺขมสุขาย เวทนาย สมฺปยุตฺตสฺส เหตุสฺส ธมฺมสฺส เหตุปจฺจเยน ปจฺจโย. (1)}

\proseitem{28}{สุขาย เวทนาย สมฺปยุตฺโต เหตุ ธมฺโม สุขาย เวทนาย สมฺปยุตฺตสฺส เหตุสฺส ธมฺมสฺส อารมฺมณปจฺจเยน ปจฺจโย.\csromanspacer{} ตีณิ.}

\prose{ทุกฺขาย เวทนาย สมฺปยุตฺโต เหตุ ธมฺโม ทุกฺขาย เวทนาย สมฺปยุตฺตสฺส เหตุสฺส ธมฺมสฺส อารมฺมณปจฺจเยน ปจฺจโย.\csromanspacer{} ตีณิ.}

\prose{อทุกฺขมสุขาย เวทนาย สมปยุตฺโต เหตุ ธมฺโม อทุกฺขมสุขาย เวทนาย สมฺปยุตฺตสฺส เหตุสฺส ธมฺมสฺส อารมฺมณปจฺจเยน ปจฺจโย.\csromanspacer{} ตีณิ. (สํขิตฺตํ.)}

\proseitem{29}{เหตุยา ตีณิ, อารมฺมเณ นว, อธิปติยา จตฺตาริ, อนนฺตเร ฉ, สมนนฺตเร ฉ, สหชาเต ตีณิ, อญฺญมญฺเญ ตีณิ, นิสฺสเย ตีณิ, อุปนิสฺสเย นว, อาเสวเน ตีณิ, วิปาเก เทฺว, อินฺทฺริเย เทฺว, มคฺเค เทฺว, สมฺปยุตฺเต ตีณิ, อตฺถิยา ตีณิ, นตฺถิยา ฉ, วิคเต ฉ, อวิคเต ตีณิ. (สํขิตฺตํ.)}

\prose{นเหตุยา นว, น-อารมฺมเณ นว. (สํขิตฺตํ.)}

\csromanheader{2}{เหตุปจฺจยา น-อารมฺมเณ ตีณิ. (สํขิตฺตํ.)}
\csromanheader{2}{นเหตุปจฺจยา อารมฺมเณ นว. (สํขิตฺตํ.)}
\prose{(ยถา กุสลตฺติเก ปญฺหาวารํ, เอวํ วิตฺถาเรตพฺพํ.)}

\csromanheader{2}{2. เวทนาตฺติก 1. เหตุทุก \textbf{1. ปฏิจฺจวาร ปจฺจยจตุกฺก}}
\proseitem{30}{\csromanfolio{477}สุขาย เวทนาย สมฺปยุตฺตํ นเหตุํ ธมฺมํ ปฏิจฺจ สุขาย เวทนาย สมฺปยุตฺโต นเหตุ ธมฺโม อุปฺปชฺชติ เหตุปจฺจยา. (1)}

\prose{ทุกฺขาย เวทนาย สมฺปยุตฺตํ นเหตุํ ธมฺมํ ปฏิจฺจ ทุกฺขาย เวทนาย สมฺปยุตฺโต นเหตุ ธมฺโม อุปฺปชฺชติ เหตุปจฺจยา. (1)}

\prose{อทุกฺขมสุขาย เวทนาย สมฺปยุตฺตํ นเหตุํ ธมฺมํ ปฏิจฺจ อทุกฺขมสุขาย เวทนาย สมฺปยุตฺโต นเหตุ ธมฺโม อุปฺปชฺชติ เหตุปจฺจยา. (1) (สํขิตฺตํ.)}

\proseitem{31}{เหตุยา ตีณิ, อารมฺมเณ ตีณิ ฯเปฯ อวิคเต ตีณิ. (สํขิตฺตํ.)}

\prose{นเหตุยา ตีณิ, น-อธิปติยา ตีณิ, นปุเรชาเต เทฺว, นปจฺฉาชาเต ตีณิ, น-อาเสวเน ตีณิ, นกมฺเม ตีณิ, นวิปาเก ตีณิ ฯเปฯ นวิปฺปยุตฺเต เทฺว. (สํขิตฺตํ.)}

\prose{(สหชาตวารมฺปิ ฯเปฯ สมฺปยุตฺตวารมฺปิ วิตฺถาเรตพฺพํ.)}

\csromanheader{2}{7. ปญฺหาวาร ปจฺจยจตุกฺก-อารมฺมณาทิ}
\proseitem{32}{สุขาย เวทนาย สมฺปยุตฺโต นเหตุ ธมฺโม สุขาย เวทนาย สมฺปยุตฺตสฺส นเหตุสฺส ธมฺมสฺส อารมฺมณปจฺจเยน ปจฺจโย.\csromanspacer{}\csromanspacer{}สุขาย เวทนาย สมฺปยุตฺโต นเหตุ ธมฺโม ทุกฺขาย เวทนาย สมฺปยุตฺตสฺส นเหตุสฺส ธมฺมสฺส อารมฺมณปจฺจเยน ปจฺจโย.\csromanspacer{}\csromanspacer{}สุขาย เวทนาย สมฺปยุตฺโต นเหตุ ธมฺโม อทุกฺขมสุขาย เวทนาย สมฺปยุตฺตสฺส นเหตุสฺส ธมฺมสฺส อารมฺมณปจฺจเยน ปจฺจโย. (3)}

\prose{ทุกฺขาย เวทนาย สมฺปยุตฺโต นเหตุ ธมฺโม ทุกฺขาย เวทนาย สมฺปยุตฺตสฺส นเหตุสฺส ธมฺมสฺส อารมฺมณปจฺจเยน ปจฺจโย.\csromanspacer{}\csromanspacer{}ทุกฺขาย เวทนาย สมฺปยุตฺโต เนเหตุ ธมฺโม สุขาย เวทนาย สมฺปยุตฺตสฺส นเหตุสฺส ธมฺมสฺส อารมฺมณปจฺจเยน ปจฺจโย.\csromanspacer{}\csromanspacer{}ทุกฺขาย เวทนาย สมฺปยุตฺโต นเหตุ ธมฺโม อทุกฺขมสุขาย เวทนาย สมฺปยุตฺตสฺส นเหตุสฺส ธมฺมสฺส อารมฺมณปจฺจเยน ปจฺจโย. (3)}

\gambhira{ธมฺมานุโลม ติกทุกปฏฺฐานปาฬิ}
\prose{\csromanfolio{478}อทุกฺขมสุขาย เวทนาย สมฺปยุตฺโต นเหตุ ธมฺโม อทุกฺขมสุขาย เวทนาย สมฺปยุตฺตสฺส นเหตุสฺส ธมฺมสฺส อารมฺมณปจฺจเยน ปจฺจโย.\csromanspacer{}\csromanspacer{}อทุกฺขมสุขาย เวทนาย สมฺปยุตฺโต นเหตุ ธมฺโม สุขาย เวทนาย สมฺปยุตฺตสฺส นเหตุสฺส ธมฺมสฺส อารมฺมณปจฺจเยน ปจฺจโย.\csromanspacer{}\csromanspacer{}อทุกฺขมสุขาย เวทนาย สมฺปยุตฺโต นเหตุ ธมฺโม ทุกฺขาย เวทนาย สมฺปยุตฺตสฺส นเหตุสฺส ธมฺมสฺส อารมฺมณปจฺจเยน ปจฺจโย. (3)}

\proseitem{33}{สุขาย เวทนาย สมฺปยุตฺโต นเหตุ ธมฺโม สุขาย เวทนาย สมฺปยุตฺตสฺส นเหตุสฺส ธมฺมสฺส อธิปติปจฺจเยน ปจฺจโย.\csromanspacer{}\csromanspacer{}สุขาย เวทนาย สมฺปยุตฺโต นเหตุ ธมฺโม อทุกฺขมสุขาย เวทนาย สมฺปยุตฺตสฺส นเหตุสฺส ธมฺมสฺส อธิปติปจฺจเยน ปจฺจโย. (2)}

\prose{ทุกฺขาย เวทนาย สมฺปยุตฺโต นเหตุ ธมฺโม ทุกฺขาย เวทนาย สมฺปยุตฺตสฺส นเหตุสฺส ธมฺมสฺส อธิปติปจฺจเยน ปจฺจโย. (1)}

\prose{อทุกฺขมสุขาย เวทนาย สมฺปยุตฺโต นเหตุ ธมฺโม อทุกฺขมสุขาย เวทนาย สมฺปยุตฺตสฺส นเหตุสฺส ธมฺมสฺส อธิปติปจฺจเยน ปจฺจโย.\csromanspacer{}\csromanspacer{}อทุกฺขมสุขาย เวทนาย สมฺปยุตฺโต นเหตุ ธมฺโม สุขาย เวทนาย สมฺปยุตฺตสฺส นเหตุสฺส ธมฺมสฺส อธิปติปจฺจเยน ปจฺจโย. (2)}

\csromanheader{2}{\textbf{อนนฺตร อุปนิสฺสย}}
\proseitem{34}{สุขาย เวทนาย สมฺปยุตฺโต นเหตุ ธมฺโม สุขาย เวทนาย สมฺปยุตฺตสฺส นเหตุสฺส ธมฺมสฺส อนนฺตรปจฺจเยน ปจฺจโย.\csromanspacer{} เทฺว.}

\prose{ทุกฺขาย เวทนาย สมฺปยุตฺโต นเหตุ ธมฺโม ทุกฺขาย เวทนาย สมฺปยุตฺตสฺส นเหตุสฺส ธมฺมสฺส อนนฺตรปจฺจเยน ปจฺจโย.\csromanspacer{} เทฺว.}

\prose{อทุกฺขมสุขาย เวทนาย สมฺปยุตฺโต นเหตุ ธมฺโม อทุกฺขมสุขาย เวทนาย สมฺปยุตฺตสฺส นเหตุสฺส ธมฺมสฺส อนนฺตรปจฺจเยน ปจฺจโย.\csromanspacer{} ตีณิ ฯเปฯ.}

\proseitem{35}{สุขาย เวทนาย สมฺปยุตฺโต นเหตุ ธมฺโม สุขาย เวทนาย สมฺปยุตฺตสฺส นเหตุสฺส ธมฺมสฺส อุปนิสฺสยปจฺจเยน ปจฺจโย.\csromanspacer{} ตีณิ.}

\prose{ทุกฺขาย เวทนาย สมฺปยุตฺโต นเหตุ ธมฺโม ทุกฺขาย เวทนาย สมฺปยุตฺตสฺส นเหตุ ธมฺมสฺส อุปนิสฺสยปจฺจเยน ปจฺจโย.\csromanspacer{} ตีณิ.}

\csromanmatikaanchor{mtk.96}
\csromanmatikaanchor{mtk.97}
\csromantocmarknopage{paragraph}{3. วิปากตฺติก}{mtk.96}
\bookmark[dest={mtk.96},level=4]{3. วิปากตฺติก}
\csromantocmarkref{subparagraph}{1. เหตุทุก}{mtk.97}
\bookmark[dest={mtk.97},level=5]{1. เหตุทุก}
\csromanheader{5}{3. วิปากตฺติก 1. เหตุทุก}
\prose{\csromanfolio{479}อทุกฺขมสุขาย เวทนาย สมฺปยุตฺโต นเหตุ ธมฺโม อทุกฺขมสุขาย เวทนาย สมฺปยุตฺตสฺส นเหตุสฺส ธมฺมสฺส อุปนิสฺสยปจฺจเยน ปจฺจโย.\csromanspacer{} ตีณิ. (สํขิตฺตํ.)}

\proseitem{36}{อารมฺมเณ นว, อธิปติยา ปญฺจ, อนนฺตเร สตฺต, สมนนฺตเร สตฺต, สหชาเต ตีณิ, อญฺญมญฺเญ ตีณิ, นิสฺสเย ตีณิ, อุปนิสฺสเย นว, อาเสวเน ตีณิ, กมฺเม อฏฺฐ, วิปาเก ตีณิ, อาหาเร ตีณิ, อินฺทฺริเย ตีณิ, ฌาเน ตีณิ, มคฺเค ตีณิ, สมฺปยุตฺเต ตีณิ, อตฺถิยา ตีณิ, นตฺถิยา สตฺต, วิคเต สตฺต, อวิคเต ตีณิ. (สํขิตฺตํ.)}

\prose{นเหตุยา นว, น-อารมฺมเณ นว. (สํขิตฺตํ.)}

\csromanheader{2}{อารมฺมณปจฺจยา นเหตุยา นว. (สํขิตฺตํ.)}
\csromanheader{2}{นเหตุปจฺจยา อารมฺมเณ นว. (สํขิตฺตํ.)}
\prose{(ยถา กุสลตฺติเก ปญฺหาวารํ, เอวํ วิตฺถาเรตพฺพํ.)}

\csromanheader{2}{3. วิปากตฺติก 1. เหตุทุก}
\csromanheader{2}{1. ปฏิจฺจวาร ปจฺจยจตุกฺก}
\proseitem{37}{วิปากํ เหตุํ ธมฺมํ ปฏิจฺจ วิปาโก เหตุ ธมฺโม อุปฺปชฺชติ เหตุปจฺจยา. (1)}

\prose{วิปากธมฺมธมฺมํ เหตุํ ธมฺมํ ปฏิจฺจ วิปากธมฺมธมฺโม เหตุ ธมฺโม อุปฺปชฺชติ เหตุปจฺจยา. (1)}

\prose{เนววิปากนวิปากธมฺมธมฺมํ เหตุํ ธมฺมํ ปฏิจฺจ เนววิปากนวิปากธมฺมธมฺโม เหตุ ธมฺโม อุปฺปชฺชติ เหตุปจฺจยา. (1) (สํขิตฺตํ.)}

\proseitem{38}{เหตุยา ตีณิ, อารมฺมเณ ตีณิ ฯเปฯ อาเสวเน เทฺว, กมฺเม ตีณิ, วิปาเก เอกํ ฯเปฯ อวิคเต ตีณิ. (สํขิตฺตํ.)}

\prose{น-อธิปติยา ตีณิ, นปุเรชาเต ตีณิ, นปจฺฉาชาเต ตีณิ, น-อาเสวเน ตีณิ, นวิปาเก เทฺว, นวิปฺปยุตฺเต ตีณิ. (สํขิตฺตํ.)}

\prose{(สหชาตวารมฺปิ ฯเปฯ สมฺปยุตฺตวารมฺปิ ปฏิจฺจวารสทิสํ วิตฺถาเรตพฺพํ.)}

\csromanheader{2}{ธมฺมานุโลม ติกทุกปฏฺฐานปาฬิ \textbf{7. ปญฺหาวาร ปจฺจยจตุกฺก}}
\proseitem{39}{\csromanfolio{480}วิปาโก เหตุ ธมฺโม วิปากสฺส เหตุสฺส ธมฺมสฺส เหตุปจฺจเยน ปจฺจโย. (1)}

\prose{วิปากธมฺมธมฺโม เหตุ ธมฺโม วิปากธมฺมธมฺมสฺส เหตุสฺส ธมฺมสฺส เหตุปจฺจเยน ปจฺจโย. (1)}

\prose{เนววิปากนวิปากธมฺมธมฺโม เหตุ ธมฺโม เนววิปากนวิปากธมฺมธมฺมสฺส เหตุสฺส ธมฺมสฺส เหตุปจฺจเยน ปจฺจโย. (1)}

\proseitem{40}{วิปาโก เหตุ ธมฺโม วิปากสฺส เหตุสฺส ธมฺมสฺส อารมฺมณปจฺจเยน ปจฺจโย.\csromanspacer{}\csromanspacer{}วิปาโก เหตุ ธมฺโม วิปากธมฺมธมฺมสฺส เหตุสฺส ธมฺมสฺส อารมฺมณปจฺจเยน ปจฺจโย.\csromanspacer{}\csromanspacer{}วิปาโก เหตุ ธมฺโม เนววิปากนวิปากธมฺมธมฺมสฺส เหตุสฺส ธมฺมสฺส อารมฺมณปจฺจเยน ปจฺจโย. (3)}

\prose{วิปากธมฺมธมฺโม เหตุ ธมฺโม วิปากธมฺมธมฺมสฺส เหตุสฺส ธมฺมสฺส อารมฺมณปจฺจเยน ปจฺจโย.\csromanspacer{}\csromanspacer{}วิปากธมฺมธมฺโม เหตุ ธมฺโม วิปากสฺส เหตุสฺส ธมฺมสฺส อารมฺมณปจฺจเยน ปจฺจโย.\csromanspacer{}\csromanspacer{}วิปากธมฺมธมฺโม เหตุ ธมฺโม เนววิปากนวิปากธมฺมธมฺมสฺส เหตุสฺส ธมฺมสฺส อารมฺมณปจฺจเยน ปจฺจโย. (3)}

\prose{เนววิปากนวิปากธมฺมธมฺโม เหตุ ธมฺโม เนววิปากนวิปากธมฺมธมฺมสฺส เหตุสฺส ธมฺมสฺส อารมฺมณปจฺจเยน ปจฺจโย.\csromanspacer{} ตีณิ.}

\proseitem{41}{วิปาโก เหตุ ธมฺโม วิปากสฺส เหตุสฺส ธมฺมสฺส อธิปติปจฺจเยน ปจฺจโย.\csromanspacer{}\csromanspacer{}วิปาโก เหตุ ธมฺโม วิปากธมฺมธมฺมสฺส เหตุสฺส ธมฺมสฺส อธิปติปจฺจเยน ปจฺจโย.\csromanspacer{}\csromanspacer{}วิปาโก เหตุ ธมฺโม เนววิปากนวิปากธมฺมธมฺมสฺส เหตุสฺส ธมฺมสฺส อธิปติปจฺจเยน ปจฺจโย.\csromanspacer{} ตีณิ.}

\prose{วิปากธมฺมธมฺโม เหตุ ธมฺโม วิปากธมฺมธมฺมสฺส เหตุสฺส ธมฺมสฺส อธิปติปจฺจเยน ปจฺจโย.\csromanspacer{} เทฺว.}

\prose{เนววิปากนวิปากธมฺมธมฺโม เหตุ ธมฺโม เนววิปากนวิปากธมฺมธมฺมสฺส เหตุสฺส ธมฺมสฺส อธิปติปจฺจเยน ปจฺจโย. (1) (สํขิตฺตํ.)}

\csromanheader{2}{3. วิปากตฺติก 1. เหตุทุก}
\proseitem{42}{\csromanfolio{481}เหตุยา ตีณิ, อารมฺมเณ นว, อธิปติยา ฉ, อนนฺตเร ปญฺจ, สมนนฺตเร ปญฺจ, สหชาเต ตีณิ, อญฺญมญฺเญ ตีณิ, นิสฺสเย ตีณิ, อุปนิสฺสเย นว, อาเสวเน เทฺว, วิปาเก เอกํ ฯเปฯ อวิคเต ตีณิ. (สํขิตฺตํ.)}

\prose{นเหตุยา นว, น-อารมฺมเณ นว. (สํขิตฺตํ.)}

\csromanheader{2}{เหตุปจฺจยา น-อารมฺมเณ ตีณิ. (สํขิตฺตํ.)}
\prose{นเหตุปจฺจยา อารมฺมเณ นว. (สํขิตฺตํ.) (ยถา กุสลตฺติเก ปญฺหาวารํ, เอวํ วิตฺถาเรตพฺพํ.)}

\csromanheader{2}{\textbf{นเหตุปทปจฺจยจตุกฺก}}
\proseitem{43}{วิปากํ นเหตุํ ธมฺมํ ปฏิจฺจ วิปาโก นเหตุ ธมฺโม อุปฺปชฺชติ เหตุปจฺจยา.\csromanspacer{}\csromanspacer{}วิปากํ นเหตุํ ธมฺมํ ปฏิจฺจ เนววิปากนวิปากธมฺมธมฺโม นเหตุ ธมฺโม อุปฺปชฺชติ เหตุปจฺจยา.\csromanspacer{}\csromanspacer{}วิปากํ นเหตุํ ธมฺมํ ปฏิจฺจ วิปาโก นเหตุ จ เนววิปากนวิปากธมฺมธมฺโม นเหตุ จ ธมฺมา อุปฺปชฺชนฺติ เหตุปจฺจยา. (3)}

\prose{วิปากธมฺมธมฺมํ นเหตุํ ธมฺมํ ปฏิจฺจ วิปากธมฺมธมฺโม นเหตุ ธมฺโม อุปฺปชฺชติ เหตุปจฺจยา.\csromanspacer{}\csromanspacer{}วิปากธมฺมธมฺมํ นเหตุํ ธมฺมํ ปฏิจฺจ เนววิปากนวิปากธมฺมธมฺโม นเหตุ ธมฺโม อุปฺปชฺชติ เหตุปจฺจยา.\csromanspacer{}\csromanspacer{}วิปากธมฺมธมฺมํ นเหตุํ ธมฺมํ ปฏิจฺจ วิปากธมฺมธมฺโม นเหตุ จ เนววิปากนวิปากธมฺมธมฺโม นเหตุ จ ธมฺมา อุปฺปชฺชนฺติ เหตุปจฺจยา. (3)}

\prose{นววิปากนวิปากธมฺมธมฺมํ นเหตุํ ธมฺมํ ปฏิจฺจ เนววิปากนวิปากธมฺมธมฺโม นเหตุ ธมฺโม อุปฺปชฺชติ เหตุปจฺจยา.\csromanspacer{} ตีณิ.}

\prose{วิปากํ นเหตุญฺจ เนววิปากนวิปากธมฺมธมฺมํ นเหตุญฺจ ธมฺมํ ปฏิจฺจ เนววิปากนวิปากธมฺมธมฺโม นเหตุ ธมฺโม อุปฺปชฺชติ เหตุปจฺจยา.\csromanspacer{} ตีณิ.}

\prose{วิปากธมฺมธมฺมํ นเหตุญฺจ เนววิปากนวิปากธมฺมธมฺมํ นเหตุญฺจ ธมฺมํ ปฏิจฺจ เนววิปากนวิปากธมฺมธมฺโม นเหตุ ธมฺโม อุปฺปชฺชติ เหตุปจฺจยา. (สํขิตฺตํ.)}

\proseitem{44}{เหตุยา เตรส, อารมฺมเณ ปญฺจ, อธิปติยา นว, อนนฺตเร ปญฺจ, สมนนฺตเร ปญฺจ, สหชาเต เตรส, อญฺญมญฺเญ สตฺต, นิสฺสเย}

\vspace{\tipitakaparskip}
\csromancenter{ธมฺมานุโลม ติกทุกปฏฺฐานปาฬิ เตรส, อุปนิสฺสเย ปญฺจ, ปุเรชาเต ตีณิ, อาเสวเน เทฺว, กมฺเม เตรส, วิปาเก นว, อาหาเร เตรส ฯเปฯ สมฺปยุตฺเต ปญฺจ ฯเปฯ อวิคเต เตรส. (สํขิตฺตํ.)}
\csromanheader{2}{\textbf{นเหตุ น-อารมฺมณ}}
\proseitem{45}{\csromanfolio{482}วิปากํ นเหตุํ ธมฺมํ ปฏิจฺจ วิปาโก นเหตุ ธมฺโม อุปฺปชฺชติ นเหตุปจฺจยา.\csromanspacer{} นว.}

\prose{วิปากํ นเหตุํ ธมฺมํ ปฏิจฺจ เนววิปากนวิปากธมฺมธมฺโม นเหตุ ธมฺโม อุปฺปชฺชติ น-อารมฺมณปจฺจยา. (สํขิตฺตํ.)}

\proseitem{46}{นเหตุยา นว, น-อารมฺมเณ ปญฺจ, น-อธิปติยา เตรส, น-อนนฺตเร ปญฺจ, นสมนนฺตเร ปญฺจ, น-อญฺญมญฺเญ ปญฺจ, น-อุปนิสฺสเย ปญฺจ, นปุเรชาเต ทฺวาทส, นปจฺฉาชาเต เตรส, น-อาเสวเน เตรส, นกมฺเม เทฺว, นวิปาเก ปญฺจ ฯเปฯ นวิปฺปยุตฺเต ตีณิ. (สํขิตฺตํ.) (สหชาตวารมฺปิ ฯเปฯ สมฺปยุตฺตวารมฺปิ ปฏิจฺจวารสทิสํ วิตฺถาเรตพฺพํ.)}

\csromanheader{2}{7. ปญฺหาวาร ปจฺจยจตุกฺก}
\proseitem{47}{วิปาโก นเหตุ ธมฺโม วิปากสฺส นเหตุสฺส ธมฺมสฺส อารมฺมณปจฺจเยน ปจฺจโย.\csromanspacer{}\csromanspacer{}วิปาโก นเหตุ ธมฺโม วิปากธมฺมธมฺมสฺส นเหตุสฺส ธมฺมสฺส อารมฺมณปจฺจเยน ปจฺจโย.\csromanspacer{}\csromanspacer{}วิปาโก นเหตุ ธมฺโม เนววิปากนวิปากธมฺมธมฺมสฺส นเหตุสฺส ธมฺมสฺส อารมฺมณปจฺจเยน ปจฺจโย. (3)}

\prose{วิปากธมฺมธมฺโม นเหตุ ธมฺโม วิปากธมฺมธมฺมสฺส นเหตุสฺส ธมฺมสฺส อารมฺมณปจฺจเยน ปจฺจโย.\csromanspacer{} ตีณิ.}

\prose{เนววิปากนวิปากธมฺมธมฺโม นเหตุ ธมฺโม เนววิปากนวิปากธมฺมธมฺมสฺส นเหตุสฺส ธมฺมสฺส อารมฺมณปจฺจเยน ปจฺจโย.\csromanspacer{} ตีณิ.}

\proseitem{48}{วิปาโก นเหตุ ธมฺโม วิปากสฺส นเหตุสฺส ธมฺมสฺส อธิปติปจฺจเยน ปจฺจโย.\csromanspacer{} จตฺตาริ.}

\prose{วิปากธมฺมธมฺโม นเหตุ ธมฺโม วิปากธมฺมธมฺมสฺส นเหตุสฺส ธมฺมสฺส อธิปติปจฺจเยน ปจฺจโย.\csromanspacer{} ตีณิ.}

\csromanmatikaanchor{mtk.98}
\csromanmatikaanchor{mtk.99}
\csromantocmarknopage{subparagraph}{4. อุปาทินฺนตฺติก}{mtk.98}
\bookmark[dest={mtk.98},level=5]{4. อุปาทินฺนตฺติก}
\csromantocmarkref{subparagraph}{1. เหตุทุก}{mtk.99}
\bookmark[dest={mtk.99},level=5]{1. เหตุทุก}
\csromanheader{6}{4. อุปาทินฺนตฺติก 1. เหตุทุก}
\prose{\csromanfolio{483}เนววิปากนวิปากธมฺมธมฺโม นเหตุ ธมฺโม เนววิปากนวิปากธมฺมธมฺมสฺส นเหตุสฺส ธมฺมสฺส อธิปติปจฺจเยน ปจฺจโย.\csromanspacer{}\csromanspacer{}เนววิปากนวิปากธมฺมธมฺโม นเหตุ ธมฺโม วิปากสฺส นเหตุสฺส ธมฺมสฺส อธิปติปจฺจเยน ปจฺจโย.\csromanspacer{}\csromanspacer{}เนววิปากนวิปากธมฺมธมฺโม นเหตุ ธมฺโม วิปากธมฺมธมฺมสฺส นเหตุสฺส ธมฺมสฺส อธิปติปจฺจเยน ปจฺจโย.\csromanspacer{} ตีณิ.}

\proseitem{49}{วิปาโก นเหตุ ธมฺโม วิปากสฺส นเหตุสฺส ธมฺมสฺส อนนฺตรปจฺจเยน ปจฺจโย.\csromanspacer{} เทฺว.}

\prose{วิปากธมฺมธมฺโม นเหตุ ธมฺโม วิปากธมฺมธมฺมสฺส นเหตุสฺส ธมฺมสฺส อนนฺตรปจฺจเยน ปจฺจโย.\csromanspacer{} เทฺว.}

\prose{เนววิปากนวิปากธมฺมธมฺโม นเหตุ ธมฺโม เนววิปากนวิปากธมฺมธมฺมสฺส นเหตุสฺส ธมฺมสฺส อนนฺตรปจฺจเยน ปจฺจโย.\csromanspacer{} ตีณิ. (สํขิตฺตํ.)}

\proseitem{50}{อารมฺมเณ นว, อธิปติยา ทส, อนนฺตเร สตฺต, สมนนฺตเร สตฺต, สหชาเต เอกาทส, อญฺญมญฺเญ สตฺต, นิสฺสเย เตรส, อุปนิสฺสเย นว, อาเสวเน เทฺว, กมฺเม นว, วิปาเก ตีณิ, อาหาเร สตฺต, อินฺทฺริเย นว, ฌาเน สตฺต, มคฺเค สตฺต, สมฺปยุตฺเต ตีณิ, วิปฺปยุตฺเต ปญฺจ, อตฺถิยา เตรส ฯเปฯ อวิคเต เตรส. (สํขิตฺตํ.)}

\prose{นเหตุยา โสฬส, น-อารมฺมเณ โสฬส. (สํขิตฺตํ.)}

\csromanheader{2}{เหตุปจฺจยา น-อารมฺมเณ นว. (สํขิตฺตํ.)}
\csromanheader{2}{นเหตุปจฺจยา อารมฺมเณ นว. (สํขิตฺตํ.)}
\prose{(ยถา กุสลตฺติเก ปญฺหาวารํ, เอวํ วิตฺถาเรตพฺพํ.)}

\csromanheader{2}{4. อุปาทินฺนตฺติก 1. เหตุทุก}
\csromanheader{2}{1. ปฏิจฺจวาร ปจฺจยจตุกฺก}
\proseitem{51}{อุปาทินฺนุปาทานิยํ เหตุํ ธมฺมํ ปฏิจฺจ อุปาทินฺนุปาทานิโย เหตุ ธมฺโม อุปฺปชฺชติ เหตุปจฺจยา. (1)}

\prose{อนุปาทินฺนุปาทานิยํ เหตุํ ธมฺมํ ปฏิจฺจ อนุปาทินฺนุปาทานิโย เหตุ ธมฺโม อุปฺปชฺชติ เหตุปจฺจยา. (1)}

\gambhira{ธมฺมานุโลม ติกทุกปฏฺฐานปาฬิ}
\prose{\csromanfolio{484}อนุปาทินฺน-อนุปาทานิยํ เหตุํ ธมฺมํ ปฏิจฺจ อนุปาทินฺน-อนุปาทานิโย เหตุ ธมฺโม อุปฺปชฺชติ เหตุปจฺจยา. (1) (สํขิตฺตํ.)}

\proseitem{52}{เหตุยา ตีณิ, อารมฺมเณ ตีณิ, อธิปติยา เทฺว ฯเปฯ วิปาเก เทฺว. (สํขิตฺตํ.)}

\prose{น-อธิปติยา ตีณิ, นปุเรชาเต ตีณิ ฯเปฯ น-อาเสวเน ตีณิ, นวิปาเก เทฺว, นวิปฺปยุตฺเต ตีณิ. (สํขิตฺตํ.\csromanspacer{} สหชาตวารมฺปิ ฯเปฯ สมฺปยุตฺตวารมฺปิ วิตฺถาเรตพฺพํ.)}

\csromanheader{2}{7. ปญฺหาวาร ปจฺจยจตุกฺก}
\proseitem{53}{อุปาทินฺนุปาทานิโย เหตุ ธมฺโม อุปาทินฺนุปาทานิยสฺส เหตุสฺส ธมฺมสฺส เหตุปจฺจเยน ปจฺจโย. (1)}

\prose{อนุปาทินฺนุปาทานิโย เหตุ ธมฺโม อนุปาทินฺนุปาทานิยสฺส เหตุสฺส ธมฺมสฺส เหตุปจฺจเยน ปจฺจโย. (1)}

\prose{อนุปาทินฺน-อนุปาทานิโย เหตุ ธมฺโม อนุปาทินฺน-อนุปาทานิยสฺส เหตุสฺส ธมฺมสฺส เหตุปจฺจเยน ปจฺจโย. (1)}

\proseitem{54}{อุปาทินฺนุปาทานิโย เหตุ ธมฺโม อุปาทินฺนุปาทานิยสฺส เหตุสฺส ธมฺมสฺส อารมฺมณปจฺจเยน ปจฺจโย.\csromanspacer{}\csromanspacer{}อุปาทินฺนุปาทานิโย เหตุ ธมฺโม อนุปาทินฺนุปาทานิยสฺส เหตุสฺส ธมฺมสฺส อารมฺมณปจฺจเยน ปจฺจโย. (2)}

\prose{อนุปาทินฺนุปาทานิโย เหตุ ธมฺโม อนุปาทินฺนุปาทานิยสฺส เหตุสฺส ธมฺมสฺส อารมฺมณปจฺจเยน ปจฺจโย.\csromanspacer{}\csromanspacer{}อนุปาทินฺนุปาทานิโย เหตุ ธมฺโม อุปาทินฺนุปาทานิยสฺส เหตุสฺส ธมฺมสฺส อารมฺมณปจฺจเยน ปจฺจโย. (2)}

\prose{อนุปาทินฺน-อนุปาทานิโย เหตุ ธมฺโม อนุปาทินฺนุปาทานิยสฺส เหตุสฺส ธมฺมสฺส อารมฺมณปจฺจเยน ปจฺจโย. (1)}

\proseitem{55}{อุปาทินฺนุปาทานิโย เหตุ ธมฺโม อนุปาทินฺนุปาทานิยสฺส เหตุสฺส ธมฺมสฺส อธิปติปจฺจเยน ปจฺจโย. (1)}

\prose{อนุปาทินฺนุปาทานิโย เหตุ ธมฺโม อนุปาทินฺนุปาทานิยสฺส เหตุสฺส ธมฺมสฺส อธิปติปจฺจเยน ปจฺจโย. (1)}

\prose{อนุปาทินฺน-อนุปาทานิโย เหตุ ธมฺโม อนุปาทินฺน-อนุปาทานิยสฺส เหตุสฺส ธมฺมสฺส อธิปติปจฺจเยน ปจฺจโย.\csromanspacer{}\csromanspacer{}อนุปาทินฺน-อนุปาทานิโย}

\vspace{\tipitakaparskip}
\csromancenter{อุปาทินฺนตฺติก 1.\csromanspacer{} เหตุทุก เหตุ ธมฺโม อนุปาทินฺนุปาทานิยสฺส เหตุสฺส ธมฺมสฺส อธิปติปจฺจเยน ปจฺจโย. (2)}
\csromanheader{2}{\textbf{อนนฺตร อุปนิสฺสย}}
\proseitem{56}{\csromanfolio{485}อุปาทินฺนุปาทานิโย เหตุ ธมฺโม อุปาทินฺนุปาทานิยสฺส เหตุสฺส ธมฺมสฺส อนนฺตรปจฺจเยน ปจฺจโย. (1)}

\prose{อนุปาทินฺนุปาทานิโย เหตุ ธมฺโม อนุปาทินฺนุปาทานิยสฺส เหตุสฺส ธมฺมสฺส อนนฺตรปจฺจเยน ปจฺจโย.\csromanspacer{} ตีณิ.}

\prose{อนุปาทินฺน-อนุปาทานิโย เหตุ ธมฺโม อนุปาทินฺน-อนุปาทานิยสฺส เหตุสฺส ธมฺมสฺส อนนฺตรปจฺจเยน ปจฺจโย.\csromanspacer{}\csromanspacer{}อนุปาทินฺน-อนุปาทานิโย เหตุ ธมฺโม อุปาทินฺนุปาทานิยสฺส เหตุสฺส ธมฺมสฺส อนนฺตรปจฺจเยน ปจฺจโย ฯเปฯ. (2)}

\proseitem{57}{อุปาทินฺนุปาทานิโย เหตุ ธมฺโม อุปาทินฺนุปาทานิยสฺส เหตุสฺส ธมฺมสฺส อุปนิสฺสยปจฺจเยน ปจฺจโย.\csromanspacer{}\csromanspacer{}อุปาทินฺนุปาทานิโย เหตุ ธมฺโม อนุปาทินฺนุปาทานิยสฺส เหตุสฺส ธมฺมสฺส อุปนิสฺสยปจฺจเยน ปจฺจโย.\csromanspacer{}\csromanspacer{}อุปาทินฺนุปาทานิโย เหตุ ธมฺโม อนุปาทินฺน-อนุปาทานิยสฺส เหตุสฺส ธมฺมสฺส อุปนิสฺสยปจฺจเยน ปจฺจโย. (3)}

\prose{อนุปาทินฺนุปาทานิโย เหตุ ธมฺโม อนุปาทินฺนุปาทานิยสฺส เหตุสฺส ธมฺมสฺส อุปนิสฺสยปจฺจเยน ปจฺจโย.\csromanspacer{} ตีณิ.}

\prose{อนุปาทินฺน-อนุปาทานิโย เหตุ ธมฺโม อนุปาทินฺน-อนุปาทานิยสฺส เหตุสฺส ธมฺมสฺส อุปนิสฺสยปจฺจเยน ปจฺจโย.\csromanspacer{} ตีณิ. (สํขิตฺตํ.)}

\proseitem{58}{เหตุยา ตีณิ, อารมฺมเณ ปญฺจ, อธิปติยา จตฺตาริ, อนนฺตเร ฉ, สมนนฺตเร ฉ, สหชาเต ตีณิ, อญฺญมญฺเญ ตีณิ, นิสฺสเย ตีณิ, อุปนิสฺสเย นว, อาเสวเน เทฺว, วิปาเก เทฺว ฯเปฯ อวิคเต ตีณิ. (สํขิตฺตํ.)}

\prose{นเหตุยา นว, น-อารมฺมเณ นว. (สํขิตฺตํ.)}

\csromanheader{2}{เหตุปจฺจยา น-อารมฺมเณ ตีณิ. (สํขิตฺตํ.)}
\csromanheader{2}{นเหตุปจฺจยา อารมฺมเณ ปญฺจ. (สํขิตฺตํ.)}
\prose{(ยถา กุสลตฺติเก ปญฺหาวารํ, เอวํ วิตฺถาเรตพฺพํ.)}

\csromanheader{2}{ธมฺมานุโลม ติกทุกปฏฺฐานปาฬิ \textbf{นเหตุปทปจฺจยจตุกฺก}}
\proseitem{59}{\csromanfolio{486}อุปาทินฺนุปาทานิยํ นเหตุํ ธมฺมํ ปฏิจฺจ อุปาทินฺนุปาทานิโย นเหตุ ธมฺโม อุปฺปชฺชติ เหตุปจฺจยา.\csromanspacer{}\csromanspacer{}อุปาทินฺนุปาทานิยํ นเหตุํ ธมฺมํ ปฏิจฺจ อนุปาทินฺนุปาทานิโย นเหตุ ธมฺโม อุปฺปชฺชติ เหตุปจฺจยา.\csromanspacer{}\csromanspacer{}อุปาทินฺนุปาทานิยํ นเหตุํ ธมฺมํ ปฏิจฺจ อุปาทินฺนุปาทานิโย นเหตุ จ อนุปาทินฺนุปาทานิโย นเหตุ จ ธมฺมา อุปฺปชฺชนฺติ เหตุปจฺจยา. (3)}

\prose{อนุปาทินฺนุปาทานิยํ นเหตุํ ธมฺมํ ปฏิจฺจ อนุปาทินฺนุปาทานิโย นเหตุ ธมฺโม อุปฺปชฺชติ เหตุปจฺจยา. (1)}

\prose{อนุปาทินฺน-อนุปาทานิยํ นเหตุํ ธมฺมํ ปฏิจฺจ อนุปาทินฺน-อนุปาทานิโย นเหตุ ธมฺโม อุปฺปชฺชติ เหตุปจฺจยา.\csromanspacer{}\csromanspacer{}อนุปาทินฺน-อนุปาทานิยํ นเหตุํ ธมฺมํ ปฏิจฺจ อนุปาทินฺนุปาทานิโย นเหตุ ธมฺโม อุปฺปชฺชติ เหตุปจฺจยา.\csromanspacer{}\csromanspacer{}อนุปาทินฺน-อนุปาทานิยํ นเหตุํ ธมฺมํ ปฏิจฺจ อนุปาทินฺนุปาทานิโย นเหตุ จ อนุปาทินฺน-อนุปาทานิโย นเหตุ จ ธมฺมา อุปฺปชฺชนฺติ เหตุปจฺจยา. (3)}

\prose{อุปาทินฺนุปาทานิยํ นเหตุญฺจ อนุปาทินฺนุปาทานิยํ นเหตุญฺจ ธมฺมํ ปฏิจฺจ อนุปาทินฺนุปาทานิโย นเหตุ ธมฺโม อุปฺปชฺชติ เหตุปจฺจยา. (1)}

\prose{อนุปาทินฺนุปาทานิยํ นเหตุญฺจ อนุปาทินฺน-อนุปาทานิยํ นเหตุญฺจ ธมฺมํ ปฏิจฺจ อนุปาทินฺนุปาทานิโย นเหตุ ธมฺโม อุปฺปชฺชติ เหตุปจฺจยา. (1)}

\proseitem{60}{อุปาทินฺนุปาทานิยํ นเหตุํ ธมฺมํ ปฏิจฺจ อุปาทินฺนุปาทานิโย นเหตุ ธมฺโม อุปฺปชฺชติ อารมฺมณปจฺจยา. (1)}

\prose{อนุปาทินฺนุปาทานิยํ นเหตุํ ธมฺมํ ปฏิจฺจ อนุปาทินฺนุปาทานิโย นเหตุ ธมฺโม อุปฺปชฺชติ อารมฺมณปจฺจยา. (1)}

\prose{อนุปาทินฺน-อนุปาทานิยํ นเหตุํ ธมฺมํ ปฏิจฺจ อนุปาทินฺน-อนุปาทานิโย นเหตุ ธมฺโม อุปฺปชฺชติ อารมฺมณปจฺจยา. (1) (สํขิตฺตํ.)}

\proseitem{61}{เหตุยา นว, อารมฺมเณ ตีณิ, อธิปติยา ปญฺจ, อนนฺตเร ตีณิ, สมนนฺตเร ตีณิ, สหชาเต นว, อญฺญมญฺเญ ตีณิ, นิสฺสเย นว, อุปนิสฺสเย ตีณิ, ปุเรชาเต ตีณิ, อาเสวเน เทฺว, กมฺเม นว, วิปาเก นว ฯเปฯ อวิคเต นว. (สํขิตฺตํ.)}

\csromanheader{2}{5. สํกิลิฏฺฐตฺติก 1. เหตุทุก \textbf{ปจฺจนีย-นเหตุ}}
\proseitem{62}{\csromanfolio{487}อุปาทินฺนุปาทานิยํ นเหตุํ ธมฺมํ ปฏิจฺจ อุปาทินฺนุปาทานิโย นเหตุ ธมฺโม อุปฺปชฺชติ นเหตุปจฺจยา.\csromanspacer{} ตีณิ.}

\prose{อนุปาทินฺนุปาทานิยํ นเหตุํ ธมฺมํ ปฏิจฺจ อนุปาทินฺนุปาทานิโย นเหตุ ธมฺโม อุปฺปชฺชติ นเหตุปจฺจยา. (1)}

\prose{อุปาทินฺนุปาทานิยํ นเหตุญฺจ อนุปาทินฺนุปาทานิยํ นเหตุญฺจ ธมฺมํ ปฏิจฺจ อนุปาทินฺนุปาทานิโย นเหตุ ธมฺโม อุปฺปชฺชติ นเหตุปจฺจยา. (1) (สํขิตฺตํ.)}

\proseitem{63}{นเหตุยา ปญฺจ, น-อารมฺมเณ ฉ. (สํขิตฺตํ.)}

\csromanheader{2}{เหตุปจฺจยา น-อารมฺมเณ ฉ. (สํขิตฺตํ.)}
\csromanheader{2}{นเหตุปจฺจยา อารมฺมเณ เทฺว. (สํขิตฺตํ.)}
\prose{(ยถา กุสลตฺติเก ปญฺหาวารํ, เอวํ วิตฺถาเรตพฺพํ.)}

\csromanmatikaanchor{mtk.100}
\csromanmatikaanchor{mtk.101}
\csromantocmarknopage{subparagraph}{5. สํกิลิฏฺฐตฺติก}{mtk.100}
\bookmark[dest={mtk.100},level=5]{5. สํกิลิฏฺฐตฺติก}
\csromantocmarkref{subparagraph}{1. เหตุทุก}{mtk.101}
\bookmark[dest={mtk.101},level=5]{1. เหตุทุก}
\csromanheader{6}{5. สํกิลิฏฺฐตฺติก 1. เหตุทุก}
\csromanheader{2}{1. ปฏิจฺจวาร ปจฺจยจตุกฺก}
\proseitem{64}{สํกิลิฏฺฐสํกิเลสิกํ เหตุํ ธมฺมํ ปฏิจฺจ สํกิลิฏฺฐสํกิเลสิโก เหตุ ธมฺโม อุปฺปชฺชติ เหตุปจฺจยา. (1)}

\prose{อสํกิลิฏฺฐสํกิเลสิกํ เหตุํ ธมฺมํ ปฏิจฺจ อสํกิลิฏฺฐสํกิเลสิโก เหตุ ธมฺโม อุปฺปชฺชติ เหตุปจฺจยา. (1)}

\prose{อสํกิลิฏฺฐ-อสํกิเลสิกํ เหตุํ ธมฺมํ ปฏิจฺจ อสํกิลิฏฺฐ-อสํกิเลสิโก เหตุ ธมฺโม อุปฺปชฺชติ เหตุปจฺจยา. (1) (สํขิตฺตํ.)}

\proseitem{65}{เหตุยา ตีณิ, อารมฺมเณ ตีณิ ฯเปฯ อวิคเต ตีณิ. (สํขิตฺตํ.)}

\gambhira{ธมฺมานุโลม ติกทุกปฏฺฐานปาฬิ}
\prose{\csromanfolio{488}น-อธิปติยา ตีณิ, นปุเรชาเต ตีณิ, นปจฺฉาชาเต ตีณิ, น-อาเสวเน ตีณิ, นวิปาเก ตีณิ, นวิปฺปยุตฺเต ตีณิ. (สํขิตฺตํ.)}

\prose{(สหชาตวารมฺปิ ฯเปฯ สมฺปยุตฺตวารมฺปิ วิตฺถาเรตพฺพํ.)}

\csromanheader{2}{7. ปญฺหาวาร ปจฺจยจตุกฺก}
\csromanheader{2}{\textbf{เหตุ-อารมฺมณาทิ}}
\proseitem{66}{สํกิลิฏฺฐสํกิเลสิโก เหตุ ธมฺโม สํกิลิฏฺฐสํกิเลสิกสฺส เหตุสฺส ธมฺมสฺส เหตุปจฺจเยน ปจฺจโย. (1)}

\prose{อสํกิลิฏฺฐสํกิเลสิโก เหตุ ธมฺโม อสํกิลิฏฺฐสํกิเลสิกสฺส เหตุสฺส ธมฺมสฺส เหตุปจฺจเยน ปจฺจโย. (1)}

\prose{อสํกิลิฏฺฐ-อสํกิเลสิโก เหตุ ธมฺโม อสํกิลิฏฺฐ-อสํกิเลสิกสฺส เหตุสฺส ธมฺมสฺส เหตุปจฺจเยน ปจฺจโย. (1)}

\proseitem{67}{สํกิลิฏฺฐสํกิเลสิโก เหตุ ธมฺโม สํกิลิฏฺฐสํกิเลสิกสฺส เหตุสฺส ธมฺมสฺส อารมฺมณปจฺจเยน ปจฺจโย.\csromanspacer{} เทฺว.}

\prose{อสํกิลิฏฺฐสํกิเลสิโก เหตุ ธมฺโม อสํกิลิฏฺฐสํกิเลสิกสฺส เหตุสฺส ธมฺมสฺส อารมฺมณปจฺจเยน ปจฺจโย.\csromanspacer{} เทฺว.}

\prose{อสํกิลิฏฺฐ-อสํกิเลสิโก เหตุ ธมฺโม อสํกิลิฏฺฐสํกิเลสิกสฺส เหตุสฺส ธมฺมสฺส อารมฺมณปจฺจเยน ปจฺจโย. (1)}

\proseitem{68}{สํกิลิฏฺฐสํกิเลสิโก เหตุ ธมฺโม สํกิลิฏฺฐสํกิเลสิกสฺส เหตุสฺส ธมฺมสฺส อธิปติปจฺจเยน ปจฺจโย. (1)}

\prose{อสํกิลิฏฺฐสํกิเลสิโก เหตุ ธมฺโม อสํกิลิฏฺฐสํกิเลสิกสฺส เหตุสฺส ธมฺมสฺส อธิปติปจฺจเยน ปจฺจโย.\csromanspacer{}\csromanspacer{}อสํกิลิฏฺฐสํกิเลสิโก เหตุ ธมฺโม สํกิลิฏฺฐสํกิเลสิกสฺส เหตุสฺส ธมฺมสฺส อธิปติปจฺจเยน ปจฺจโย. (2)}

\prose{อสํกิลิฏฺฐ-อสํกิเลสิโก เหตุ ธมฺโม อสํกิลิฏฺฐ-อสํกิเลสิกสฺส เหตุสฺส ธมฺมสฺส อธิปติปจฺจเยน ปจฺจโย.\csromanspacer{} เทฺว.}

\csromanheader{2}{5. สํกิลิฏฺฐตฺติก 1. เหตุทุก}
\proseitem{69}{\csromanfolio{489}สํกิลิฏฺฐสํกิเลสิโก เหตุ ธมฺโม สํกิลิฏฺฐสํกิเลสิกสฺส เหตุสฺส ธมฺมสฺส อนนฺตรปจฺจเยน ปจฺจโย.\csromanspacer{} ฉ ฯเปฯ.}

\prose{สํกิลิฏฺฐสํกิเลสิโก เหตุ ธมฺโม สํกิลิฏฺฐสํกิเลสิกสฺส เหตุสฺส ธมฺมสฺส อุปนิสฺสยปจฺจเยน ปจฺจโย.\csromanspacer{}\csromanspacer{}สํกิลิฏฺฐสํกิเลสิโก เหตุ ธมฺโม อสํกิลิฏฺฐสํกิเลสิกสฺส เหตุสฺส ธมฺมสฺส อุปนิสฺสยปจฺจเยน ปจฺจโย.\csromanspacer{}\csromanspacer{}สํกิลิฏฺฐสํกิเลสิโก เหตุ ธมฺโม อสํกิลิฏฺฐ-อสํกิเลสิกสฺส เหตุสฺส ธมฺมสฺส อุปนิสฺสยปจฺจเยน ปจฺจโย. (3)}

\prose{อสํกิลิฏฺฐสํกิเลสิโก เหตุ ธมฺโม อสํกิลิฏฺฐสํกิเลสิกสฺส เหตุสฺส ธมฺมสฺส อุปนิสฺสยปจฺจเยน ปจฺจโย.\csromanspacer{} ตีณิ.}

\prose{อสํกิลิฏฺฐ-อสํกิเลสิโก เหตุ ธมฺโม อสํกิลิฏฺฐ-อสํกิเลสิกสฺส เหตุสฺส ธมฺมสฺส อุปนิสฺสยปจฺจเยน ปจฺจโย.\csromanspacer{} เทฺว. (สํขิตฺตํ.)}

\proseitem{70}{เหตุยา ตีณิ, อารมฺมเณ ปญฺจ, อธิปติยา ปญฺจ, อนนฺตเร ฉ, สมนนฺตเร ฉ, สหชาเต ตีณิ, อญฺญมญฺเญ ตีณิ, นิสฺสเย ตีณิ, อุปนิสฺสเย อฏฺฐ, อาเสวเน ตีณิ, วิปาเก เทฺว, อินฺทฺริเย เทฺว ฯเปฯ อวิคเต ตีณิ. (สํขิตฺตํ.)}

\prose{นเหตุยา อฏฺฐ, น-อารมฺมเณ อฏฺฐ. (สํขิตฺตํ.)}

\csromanheader{2}{เหตุปจฺจยา น-อารมฺมเณ ตีณิ. (สํขิตฺตํ.)}
\csromanheader{2}{นเหตุปจฺจยา อารมฺมเณ ปญฺจ. (สํขิตฺตํ.)}
\prose{(ยถา กุสลตฺติเก ปญฺหาวารํ, เอวํ วิตฺถาเรตพฺพํ.)}

\csromanheader{2}{\textbf{นเหตุปทปจฺจยจตุกฺก}}
\proseitem{71}{สํกิลิฏฺฐสํกิเลสิกํ นเหตุํ ธมฺมํ ปฏิจฺจ สํกิลิฏฺฐสํกิเลสิโก นเหตุ ธมฺโม อุปฺปชฺชติ เหตุปจฺจยา.\csromanspacer{}\csromanspacer{}สํกิลิฏฺฐสํกิเลสิกํ นเหตุํ ธมฺมํ ปฏิจฺจ อสํกิลิฏฺฐสํกิเลสิโก นเหตุ ธมฺโม อุปฺปชฺชติ เหตุปจฺจยา.\csromanspacer{}\csromanspacer{}สํกิลิฏฺฐสํกิเลสิกํ นเหตุํ ธมฺมํ ปฏิจฺจ สํกิลิฏฺฐสํกิเลสิโก นเหตุ จ อสํกิลิฏฺฐสํกิเลสิโก นเหตุ จ ธมฺมา อุปฺปชฺชนฺติ เหตุปจฺจยา. (3)}

\prose{อสํกิลิฏฺฐสํกิเลสิกํ นเหตุํ ธมฺมํ ปฏิจฺจ อสํกิลิฏฺฐสํกิเลสิโก นเหตุ ธมฺโม อุปฺปชฺชติ เหตุปจฺจยา. (1)}

\gambhira{ธมฺมานุโลม ติกทุกปฏฺฐานปาฬิ}
\prose{\csromanfolio{490}อสํกิลิฏฺฐ-อสํกิเลสิกํ นเหตุํ ธมฺมํ ปฏิจฺจ อสํกิลิฏฺฐ-อสํกิเลสิโก นเหตุ ธมฺโม อุปฺปชฺชติ เหตุปจฺจยา.\csromanspacer{} ตีณิ.}

\prose{สํกิลิฏฺฐสํกิเลสิกํ นเหตุญฺจ อสํกิลิฏฺฐสํกิเลสิกํ นเหตุญฺจ ธมฺมํ ปฏิจฺจ อสํกิลิฏฺฐสํกิเลสิโก นเหตุ ธมฺโม อุปฺปชฺชติ เหตุปจฺจยา. (1)}

\prose{อสํกิลิฏฺฐสํกิเลสิกํ นเหตุญฺจ อสํกิลิฏฺฐ-อสํกิเลสิกํ นเหตุญฺจ ธมฺมํ ปฏิจฺจ อสํกิลิฏฺฐสํกิเลสิโก นเหตุ ธมฺโม อุปฺปชฺชติ เหตุปจฺจยา. (1)}

\proseitem{72}{สํกิลิฏฺฐสํกิเลสิกํ นเหตุํ ธมฺมํ ปฏิจฺจ สํกิลิฏฺฐสํกิเลสิโก นเหตุ ธมฺโม อุปฺปชฺชติ อารมฺมณปจฺจยา. (1)}

\prose{อสํกิลิฏฺฐสํกิเลสิกํ นเหตุํ ธมฺมํ ปฏิจฺจ อสํกิลิฏฺฐสํกิเลสิโก นเหตุ ธมฺโม อุปฺปชฺชติ อารมฺมณปจฺจยา. (1)}

\prose{อสํกิลิฏฺฐ-อสํกิเลสิกํ นเหตุํ ธมฺมํ ปฏิจฺจ อสํกิลิฏฺฐ-อสํกิเลสิโก นเหตุ ธมฺโม อุปฺปชฺชติ อารมฺมณปจฺจยา. (1) (สํขิตฺตํ.)}

\proseitem{73}{เหตุยา นว, อารมฺมเณ ตีณิ, อธิปติยา นว, อนนฺตเร ตีณิ, สมนนฺตเร ตีณิ, สหชาเต นว, อญฺญมญฺเญ ตีณิ, นิสฺสเย นว, อุปนิสฺสเย ตีณิ, ปุเรชาเต ตีณิ, อาเสวเน ตีณิ, กมฺเม นว, วิปาเก ปญฺจ, อาหาเร นว ฯเปฯ อวิคเต นว. (สํขิตฺตํ.)}

\csromanheader{2}{\textbf{ปจฺจนีย-นเหตุ}}
\proseitem{74}{อสํกิลิฏฺฐสํกิเลสิกํ นเหตุํ ธมฺมํ ปฏิจฺจ อสํกิลิฏฺฐสํกิเลสิโก นเหตุ ธมฺโม อุปฺปชฺชติ นเหตุปจฺจยา. (สํขิตฺตํ.)}

\prose{นเหตุยา เอกํ, น-อารมฺมเณ ปญฺจ, น-อธิปติยา ฉ, นปุเรชาเต สตฺต ฯเปฯ นกมฺเม ตีณิ ฯเปฯ นวิปฺปยุตฺเต ตีณิ ฯเปฯ โนวิคเต ปญฺจ. (สํขิตฺตํ.)}

\prose{(สหชาตวารมฺปิ ฯเปฯ สมฺปยุตฺตวารมฺปิ วิตฺถาเรตพฺพํ.)}

\csromanheader{2}{6. วิตกฺกตฺติก 1. เหตุทุก \textbf{7. ปญฺหาวาร ปจฺจยจตุกฺก}}
\csromanheader{2}{\textbf{อารมฺมณ-อธิปติ}}
\proseitem{75}{\csromanfolio{491}สํกิลิฏฺฐสํกิเลสิโก นเหตุ ธมฺโม สํกิลิฏฺฐสํกิเลสิกสฺส นเหตุสฺส ธมฺมสฺส อารมฺมณปจฺจเยน ปจฺจโย.\csromanspacer{} เทฺว.}

\prose{อสํกิลิฏฺฐสํกิเลสิโก นเหตุ ธมฺโม อสํกิลิฏฺฐสํกิเลสิกสฺส นเหตุสฺส ธมฺมสฺส อารมฺมณปจฺจเยน ปจฺจโย.\csromanspacer{} เทฺว.}

\prose{อสํกิลิฏฺฐ-อสํกิเลสิโก นเหตุ ธมฺโม อสํกิลิฏฺฐ-อสํกิเลสิกสฺส นเหตุสฺส ธมฺมสฺส อารมฺมณปจฺจเยน ปจฺจโย.\csromanspacer{} เทฺว.}

\proseitem{76}{สํกิลิฏฺฐสํกิเลสิโก นเหตุ ธมฺโม สํกิลิฏฺฐสํกิเลสิกสฺส นเหตุสฺส ธมฺมสฺส อธิปติปจฺจเยน ปจฺจโย.\csromanspacer{} ตีณิ.}

\prose{อสํกิลิฏฺฐสํกิเลสิโก นเหตุ ธมฺโม อสํกิลิฏฺฐสํกิเลสิกสฺส นเหตุสฺส ธมฺมสฺส อธิปติเยน ปจฺจโย.\csromanspacer{} เทฺว.}

\prose{อสํกิลิฏฺฐ-อสํกิเลสิโก นเหตุ ธมฺโม อสํกิลิฏฺฐ-อสํกิเลสิกสฺส นเหตุสฺส ธมฺมสฺส อธิปติปจฺจเยน ปจฺจโย.\csromanspacer{} ตีณิ. (สํขิตฺตํ.)}

\proseitem{77}{อารมฺมเณ ฉ, อธิปติยา อฏฺฐ, อนนฺตเร สตฺต, สมนนฺตเร สตฺต, สหชาเต นว, อญฺญมญฺเญ ตีณิ, นิสฺสเย เตรส, อุปนิสฺสเย อฏฺฐ, อาเสวเน ตีณิ, กมฺเม สตฺต, วิปาเก จตฺตาริ, อาหาเร สตฺต ฯเปฯ อวิคเต เตรส. (สํขิตฺตํ.)}

\prose{นเหตุยา จุทฺทส, น-อารมฺมเณ จุทฺทส. (สํขิตฺตํ.)}

\csromanheader{2}{อารมฺมณปจฺจยา นเหตุยา ฉ. (สํขิตฺตํ.)}
\csromanheader{2}{นเหตุปจฺจยา อารมฺมเณ ฉ. (สํขิตฺตํ.)}
\prose{(ยถา กุสลตฺติเก ปญฺหาวารํ, เอวํ วิตฺถาเรตพฺพํ.)}

\csromanmatikaanchor{mtk.102}
\csromanmatikaanchor{mtk.103}
\csromantocmarknopage{subparagraph}{6. วิตกฺกตฺติก}{mtk.102}
\bookmark[dest={mtk.102},level=5]{6. วิตกฺกตฺติก}
\csromantocmarkref{subparagraph}{1. เหตุทุก}{mtk.103}
\bookmark[dest={mtk.103},level=5]{1. เหตุทุก}
\csromanheader{6}{6. วิตกฺกตฺติก 1. เหตุทุก}
\csromanheader{2}{1. ปฏิจฺจวาร ปจฺจยจตุกฺก}
\proseitem{78}{สวิตกฺกสวิจารํ เหตุํ ธมฺมํ ปฏิจฺจ สวิตกฺกสวิจาโร เหตุ ธมฺโม อุปฺปชฺชติ เหตุปจฺจยา. (1)}

\gambhira{ธมฺมานุโลม ติกทุกปฏฺฐานปาฬิ}
\prose{\csromanfolio{492}อวิตกฺกวิจารมตฺตํ เหตุํ ธมฺมํ ปฏิจฺจ อวิตกฺกวิจารมตฺโต เหตุ ธมฺโม อุปฺปชฺชติ เหตุปจฺจยา. (1)}

\prose{อวิตกฺก-อวิจารํ เหตุํ ธมฺมํ ปฏิจฺจ อวิตกฺก-อวิจาโร เหตุ ธมฺโม อุปฺปชฺชติ เหตุปจฺจยา. (1) (สํขิตฺตํ.)}

\proseitem{79}{เหตุยา ตีณิ, อารมฺมเณ ตีณิ ฯเปฯ อวิคเต ตีณิ. (สํขิตฺตํ.)}

\prose{น-อธิปติยา ตีณิ, นปุเรชาเต ตีณิ, นปจฺฉาชาเต ตีณิ, น-อาเสวเน ตีณิ, นวิปาเก ตีณิ, นวิปฺปยุตฺเต ตีณิ. (สํขิตฺตํ.)}

\prose{(ยถา สหชาตวารมฺปิ สมฺปยุตฺตวารมฺปิ, เอวํ วิตฺถาเรตพฺพํ.)}

\csromanheader{2}{7. ปญฺหาวาร ปจฺจยจตุกฺก}
\proseitem{80}{สวิตกฺกสวิจาโร เหตุ ธมฺโม สวิตกฺกสวิจารสฺส เหตุสฺส ธมฺมสฺส เหตุปจฺจเยน ปจฺจโย. (1)}

\prose{อวิตกฺกวิจารมตฺโต เหตุ ธมฺโม อวิตกฺกวิจารมตฺตสฺส เหตุสฺส ธมฺมสฺส เหตุปจฺจเยน ปจฺจโย. (1)}

\prose{อวิตกฺก-อวิจาโร เหตุ ธมฺโม อวิตกฺก-อวิจารสฺส เหตุสฺส ธมฺมสฺส เหตุปจฺจเยน ปจฺจโย. (1)}

\proseitem{81}{สวิตกฺกสวิจาโร เหตุ ธมฺโม สวิตกฺกสวิจารสฺส เหตุสฺส ธมฺมสฺส อารมฺมณปจฺจเยน ปจฺจโย.\csromanspacer{}\csromanspacer{}สวิตกฺกสวิจาโร เหตุ ธมฺโม อวิตกฺก-อวิจารสฺส เหตุสฺส ธมฺมสฺส อารมฺมณปจฺจเยน ปจฺจโย. (2)}

\prose{อวิตกฺกวิจารมตฺโต เหตุ ธมฺโม สวิตกฺกสวิจารสฺส เหตุสฺส ธมฺมสฺส อารมฺมณปจฺจเยน ปจฺจโย.\csromanspacer{}\csromanspacer{}อวิตกฺกวิจารมตฺโต เหตุ ธมฺโม อวิตกฺก-อวิจารสฺส เหตุสฺส ธมฺมสฺส อารมฺมณปจฺจเยน ปจฺจโย. (2)}

\prose{อวิตกฺก-อวิจาโร เหตุ ธมฺโม อวิตกฺก-อวิจารสฺส เหตุสฺส ธมฺมสฺส อารมฺมณปจฺจเยน ปจฺจโย.\csromanspacer{}\csromanspacer{}อวิตกฺก-อวิจาโร เหตุ ธมฺโม สวิตกฺกสวิจารสฺส เหตุสฺส ธมฺมสฺส อารมฺมณปจฺจเยน ปจฺจโย. (2)}

\proseitem{82}{สวิตกฺกสวิจาโร เหตุ ธมฺโม สวิตกฺกสวิจารสฺส เหตุสฺส ธมฺมสฺส อธิปติปจฺจเยน ปจฺจโย. (1)}

\csromanheader{2}{6. วิตกฺกตฺติก 1. เหตุทุก}
\prose{\csromanfolio{493}อวิตกฺกวิจารมตฺโต เหตุ ธมฺโม อวิตกฺกวิจารมตฺตสฺส เหตุสฺส ธมฺมสฺส อธิปติปจฺจเยน ปจฺจโย.\csromanspacer{}\csromanspacer{}อวิตกฺกวิจารมตฺโต เหตุ ธมฺโม สวิตกฺกสวิจารสฺส เหตุสฺส ธมฺมสฺส อธิปติปจฺจเยน ปจฺจโย. (2)}

\prose{อวิตกฺก-อวิจาโร เหตุ ธมฺโม อวิตกฺก-อวิจารสฺส เหตุสฺส ธมฺมสฺส อธิปติปจฺจเยน ปจฺจโย.\csromanspacer{}\csromanspacer{}อวิตกฺก-อวิจาโร เหตุ ธมฺโม สวิตกฺกสวิจารสฺส เหตุสฺส ธมฺมสฺส อธิปติปจฺจเยน ปจฺจโย. (2)}

\proseitem{83}{สวิตกฺกสวิจาโร เหตุ ธมฺโม สวิตกฺกสวิจารสฺส เหตุสฺส ธมฺมสฺส อนนฺตรปจฺจเยน ปจฺจโย.\csromanspacer{} ตีณิ.}

\prose{อวิตกฺกวิจารมตฺโต เหตุ ธมฺโม อวิตกฺกวิจารมตฺตสฺส เหตุสฺส ธมฺมสฺส อนนฺตรปจฺจเยน ปจฺจโย.\csromanspacer{} ตีณิ.}

\prose{อวิตกฺก-อวิจาโร เหตุ ธมฺโม อวิตกฺก-อวิจารสฺส เหตุสฺส ธมฺมสฺส อนนฺตรปจฺจเยน ปจฺจโย.\csromanspacer{} ตีณิ ฯเปฯ.}

\prose{สวิตกฺกสวิจาโร เหตุ ธมฺโม สวิตกฺกสวิจารสฺส เหตุสฺส ธมฺมสฺส อุปนิสฺสยปจฺจเยน ปจฺจโย.\csromanspacer{} นว. (สํขิตฺตํ.)}

\proseitem{84}{เหตุยา ตีณิ, อารมฺมเณ ฉ, อธิปติยา ปญฺจ, อนนฺตเร นว, สมนนฺตเร นว, สหชาเต ตีณิ, อญฺญมญฺเญ ตีณิ, นิสฺสเย ตีณิ, อุปนิสฺสเย นว, อาเสวเน ปญฺจ, วิปาเก ตีณิ ฯเปฯ อวิคเต ตีณิ.}

\prose{นเหตุยา นว, น-อารมฺมเณ นว. (สํขิตฺตํ.)}

\csromanheader{2}{เหตุปจฺจยา น-อารมฺมเณ ตีณิ. (สํขิตฺตํ.)}
\prose{นเหตุปจฺจยา อารมฺมเณ ฉ. (สํขิตฺตํ.) (ยถา กุสลตฺติเก ปญฺหาวารํ, เอวํ วิตฺถาเรตพฺพํ.)}

\csromanheader{2}{\textbf{นเหตุปท 1. ปฏิจฺจวาร}}
\csromanheader{2}{\textbf{ปจฺจยจตุกฺก-เหตุ}}
\proseitem{85}{สวิตกฺกสวิจารํ นเหตุํ ธมฺมํ ปฏิจฺจ สวิตกฺกสวิจาโร นเหตุ ธมฺโม อุปฺปชฺชติ เหตุปจฺจยา. (เอกํ.).\csromanspacer{} สวิตกฺกสวิจารํ นเหตุํ ธมฺมํ ปฏิจฺจ อวิตกฺกวิจารมตฺโต นเหตุ ธมฺโม อุปฺปชฺชติ เหตุปจฺจยา.\csromanspacer{} เทฺว.}

\vspace{\tipitakaparskip}
\csromancenter{ธมฺมานุโลม ติกทุกปฏฺฐานปาฬิ สวิตกฺกสวิจารํ นเหตุํ ธมฺมํ ปฏิจฺจ อวิตกฺก-อวิจาโร นเหตุ ธมฺโม อุปฺปชฺชติ เหตุปจฺจยา. (ตีณิ.).\csromanspacer{} สวิตกฺกสวิจารํ นเหตุํ ธมฺมํ ปฏิจฺจ สวิตกฺกสวิจาโร นเหตุ จ อวิตกฺก-อวิจาโร นเหตุ จ ธมฺมา อุปฺปชฺชนฺติ เหตุปจฺจยา. (จตฺตาริ.).\csromanspacer{} สวิตกฺกสวิจารํ นเหตุํ ธมฺมํ ปฏิจฺจ อวิตกฺกวิจารมตฺโต นเหตุ จ อวิตกฺก-อวิจาโร นเหตุ จ ธมฺมา อุปฺปชฺชนฺติ เหตุปจฺจยา. (ปญฺจ.).\csromanspacer{} สวิตกฺกสวิจารํ นเหตุํ ธมฺมํ ปฏิจฺจ สวิตกฺกสวิจาโร นเหตุ จ อวิตกฺกวิจารมตฺโต นเหตุ จ ธมฺมา อุปฺปชฺชนฺติ เหตุปจฺจยา. (ฉ.).\csromanspacer{} สวิตกฺกสวิจารํ นเหตุํ ธมฺมํ ปฏิจฺจ สวิตกฺกสวิจาโร นเหตุ จ อวิตกฺกวิจารมตฺโต นเหตุ จ อวิตกฺก-อวิจาโร นเหตุ จ ธมฺมา อุปฺปชฺชนฺติ เหตุปจฺจยา. (สตฺต.)}
\proseitem{86}{\csromanfolio{494}อวิตกฺกวิจารมตฺตํ นเหตุํ ธมฺมํ ปฏิจฺจ อวิตกฺกวิจารมตฺโต นเหตุ ธมฺโม อุปฺปชฺชติ เหตุปจฺจยา. (เอกํ.).\csromanspacer{} อวิตกฺกวิจารมตฺตํ นเหตุํ ธมฺมํ ปฏิจฺจ สวิตกฺกสวิจาโร นเหตุ ธมฺโม อุปฺปชฺชติ เหตุปจฺจยา. (เทฺว.).\csromanspacer{} อวิตกฺกวิจารมตฺตํ นเหตุํ ธมฺมํ ปฏิจฺจ อวิตกฺก-อวิจาโร นเหตุ ธมฺโม อุปฺปชฺชติ เหตุปจฺจยา. (ตีณิ.).\csromanspacer{} อวิตกฺกวิจารมตฺตํ นเหตุํ ธมฺมํ ปฏิจฺจ สวิตกฺกสวิจาโร นเหตุ จ อวิตกฺก-อวิจาโร นเหตุ จ ธมฺมา อุปฺปชฺชนฺติ เหตุปจฺจยา. (จตฺตาริ.).\csromanspacer{} อวิตกฺกวิจารมตฺตํ นเหตุํ ธมฺมํ ปฏิจฺจ อวิตกฺกวิจารมตฺโต นเหตุ จ อวิตกฺก-อวิจาโร นเหตุ จ ธมฺมา อุปฺปชฺชนฺติ เหตุปจฺจยา. (ปญฺจ.)}

\proseitem{87}{อวิตกฺก-อวิจารํ นเหตุํ ธมฺมํ ปฏิจฺจ อวิตกฺก-อวิจาโร นเหตุ ธมฺโม อุปฺปชฺชติ เหตุปจฺจยา. (เอกํ.).\csromanspacer{} อวิตกฺก-อวิจารํ นเหตุํ ธมฺมํ ปฏิจฺจ สวิตกฺกสวิจาโร นเหตุ ธมฺโม อุปฺปชฺชติ เหตุปจฺจยา. (เทฺว.).\csromanspacer{} อวิตกฺก-อวิจารํ นเหตุํ ธมฺมํ ปฏิจฺจ อวิตกฺกวิจารมตฺโต นเหตุ ธมฺโม อุปฺปชฺชติ เหตุปจฺจยา. (ตีณิ.).\csromanspacer{} อวิตกฺก-อวิจารํ นเหตุํ ธมฺมํ ปฏิจฺจ สวิตกฺกสวิจาโร นเหตุ จ อวิตกฺก-อวิจาโร นเหตุ จ ธมฺมา อุปฺปชฺชนฺติ เหตุปจฺจยา. (จตฺตาริ.).\csromanspacer{} อวิตกฺก-อวิจารํ นเหตุํ ธมฺมํ ปฏิจฺจ อวิตกฺกวิจารมตฺโต นเหตุ จ อวิตกฺก-อวิจาโร นเหตุ จ ธมฺมา อุปฺปชฺชนฺติ เหตุปจฺจยา. (ปญฺจ.).\csromanspacer{} อวิตกฺก-อวิจารํ นเหตุํ ธมฺมํ ปฏิจฺจ สวิตกฺกสวิจาโร นเหตุ จ อวิตกฺกวิจารมตฺโต นเหตุ จ ธมฺมา อุปฺปชฺชนฺติ เหตุปจฺจยา. (ฉ.).\csromanspacer{} อวิตกฺก-อวิจารํ นเหตุํ ธมฺมํ ปฏิจฺจ สวิตกฺกสวิจาโร นเหตุ จ}

\vspace{\tipitakaparskip}
\csromancenter{วิตกฺกตฺติก 1.\csromanspacer{} เหตุทุก อวิตกฺกวิจารมตฺโต นเหตุ จ อวิตกฺก-อวิจาโร นเหตุ จ ธมฺมา อุปฺปชฺชนฺติ เหตุปจฺจยา. (สตฺต.)}
\proseitem{88}{\csromanfolio{495}สวิตกฺกสวิจารํ นเหตุญฺจ อวิตกฺก-อวิจารํ นเหตุญฺจ ธมฺมํ ปฏิจฺจ สวิตกฺกสวิจาโร นเหตุ ธมฺโม อุปฺปชฺชติ เหตุปจฺจยา. (เอกํ.) สวิตกฺกสวิจารํ นเหตุญฺจ อวิตกฺก-อวิจารํ นเหตุญฺจ ธมฺมํ ปฏิจฺจ อวิตกฺกวิจารมตฺโต นเหตุ ธมฺโม อุปฺปชฺชติ เหตุปจฺจยา. (เทฺว.) สวิตกฺกสวิจารํ นเหตุญฺจ อวิตกฺก-อวิจารํ นเหตุญฺจ ธมฺมํ ปฏิจฺจ อวิตกฺก-อวิจาโร นเหตุ ธมฺโม อุปฺปชฺชติ เหตุปจฺจยา. (ตีณิ.) สวิตกฺกสวิจารํ นเหตุญฺจ อวิตกฺก-อวิจารํ นเหตุญฺจ ธมฺมํ ปฏิจฺจ สวิตกฺกสวิจาโร นเหตุ จ อวิตกฺก-อวิจาโร นเหตุ จ ธมฺมา อุปฺปชฺชนฺติ เหตุปจฺจยา. (จตฺตาริ.).\csromanspacer{} สวิตกฺกสวิจารํ นเหตุญฺจ อวิตกฺก-อวิจารํ นเหตุญฺจ ธมฺมํ ปฏิจฺจ อวิตกฺกวิจารมตฺโต นเหตุ จ อวิตกฺก-อวิจาโร นเหตุ จ ธมฺมา อุปฺปชฺชนฺติ เหตุปจฺจยา. (ปญฺจ.).\csromanspacer{} สวิตกฺกสวิจารํ นเหตุญฺจ อวิตกฺก-อวิจารํ นเหตุญฺจ ธมฺมํ ปฏิจฺจ สวิตกฺกสวิจาโร นเหตุ จ อวิตกฺกวิจารมตฺโต นเหตุ จ ธมฺมา อุปฺปชฺชนฺติ เหตุปจฺจยา. (ฉ.).\csromanspacer{} สวิตกฺกสวิจารํ นเหตุญฺจ อวิตกฺก-อวิจารํ นเหตุญฺจ ธมฺมํ ปฏิจฺจ สวิตกฺกสวิจาโร นเหตุ จ อวิตกฺกวิจารมตฺโต นเหตุ จ อวิตกฺก-อวิจาโร นเหตุ จ ธมฺมา อุปฺปชฺชนฺติ เหตุปจฺจยา. (สตฺต.)}

\proseitem{89}{อวิตกฺกวิจารมตฺตํ นเหตุญฺจ อวิตกฺก-อวิจารํ นเหตุญฺจ ธมฺมํ ปฏิจฺจ สวิตกฺกสวิจาโร นเหตุ ธมฺโม อุปฺปชฺชติ เหตุปจฺจยา. (เอกํ.).\csromanspacer{} อวิตกฺกวิจารมตฺตํ นเหตุญฺจ อวิตกฺก-อวิจารํ นเหตุญฺจ ธมฺมํ ปฏิจฺจ อวิตกฺกวิจารมตฺโต นเหตุ ธมฺโม อุปฺปชฺชติ เหตุปจฺจยา. (เทฺว.).\csromanspacer{} อวิตกฺกวิจารมตฺตํ นเหตุญฺจ อวิตกฺก-อวิจารํ นเหตุญฺจ ธมฺมํ ปฏิจฺจ อวิตกฺก-อวิจาโร นเหตุ ธมฺโม อุปฺปชฺชติ เหตุปจฺจยา. (ตีณิ.).\csromanspacer{} อวิตกฺกวิจารมตฺตํ นเหตุญฺจ อวิตกฺก-อวิจารํ นเหตุญฺจ ธมฺมํ ปฏิจฺจ สวิตกฺกสวิจาโร นเหตุ จ อวิตกฺก-อวิจาโร นเหตุ จ ธมฺมา อุปฺปชฺชนฺติ เหตุปจฺจยา. (จตฺตาริ.).\csromanspacer{} อวิตกฺกวิจารมตฺตํ นเหตุญฺจ อวิตกฺก-อวิจารํ นเหตุญฺจ ธมฺมํ ปฏิจฺจ อวิตกฺกวิจารมตฺโต นเหตุ จ อวิตกฺก-อวิจาโร นเหตุ จ ธมฺมา อุปฺปชฺชนฺติ เหตุปจฺจยา.\csromanspacer{} ปญฺจ.}

\gambhira{ธมฺมานุโลม ติกทุกปฏฺฐานปาฬิ}
\proseitem{90}{\csromanfolio{496}สวิตกฺกสวิจารํ นเหตุญฺจ อวิตกฺกวิจารมตฺตํ นเหตุญฺจ ธมฺมํ ปฏิจฺจ สวิตกฺกสวิจาโร นเหตุ ธมฺโม อุปฺปชฺชติ เหตุปจฺจยา. (เอกํ.).\csromanspacer{} สวิตกฺกสวิจารํ นเหตุญฺจ อวิตกฺกวิจารมตฺตํ นเหตุญฺจ ธมฺมํ ปฏิจฺจ อวิตกฺก-อวิจาโร นเหตุ ธมฺโม อุปฺปชฺชติ เหตุปจฺจยา. (เทฺว.).\csromanspacer{} สวิตกฺกสวิจารํ นเหตุญฺจ อวิตกฺกวิจารมตฺตํ นเหตุญฺจ ธมฺมํ ปฏิจฺจ สวิตกฺกสวิจาโร นเหตุ จ อวิตกฺก-อวิจาโร นเหตุ จ ธมฺมา อุปฺปชฺชนฺติ เหตุปจฺจยา. (ตีณิ.)}

\proseitem{91}{สวิตกฺกสวิจารํ นเหตุญฺจ อวิตกฺกวิจารมตฺตํ นเหตุญฺจ อวิตกฺก-อวิจารํ นเหตุญฺจ ธมฺมํ ปฏิจฺจ สวิตกฺกสวิจาโร นเหตุ ธมฺโม อุปฺปชฺชติ เหตุปจฺจยา. (เอกํ.).\csromanspacer{} สวิตกฺกสวิจารํ นเหตุญฺจ อวิตกฺกวิจารมตฺตํ นเหตุญฺจ อวิตกฺก-อวิจารํ นเหตุญฺจ ธมฺมํ ปฏิจฺจ อวิตกฺก-อวิจาโร นเหตุ ธมฺโม อุปฺปชฺชติ เหตุปจฺจยา. (เทฺว.).\csromanspacer{} สวิตกฺกสวิจารํ นเหตุญฺจ อวิตกฺกวิจารมตฺตํ นเหตุญฺจ อวิตกฺก-อวิจารํ นเหตุญฺจ ธมฺมํ ปฏิจฺจ สวิตกฺกสวิจาโร นเหตุ จ อวิตกฺก-อวิจาโร นเหตุ จ ธมฺมา อุปฺปชฺชนฺติ เหตุปจฺจยา. (ตีณิ.)}

\csromanheader{2}{\textbf{อารมฺมณ}}
\proseitem{92}{สวิตกฺกสวิจารํ นเหตุํ ธมฺมํ ปฏิจฺจ สวิตกฺกสวิจาโร นเหตุ ธมฺโม อุปฺปชฺชติ อารมฺมณปจฺจยา.\csromanspacer{}\csromanspacer{}สวิตกฺกสวิจารํ นเหตุํ ธมฺมํ ปฏิจฺจ อวิตกฺกวิจารมตฺโต นเหตุ ธมฺโม อุปฺปชฺชติ อารมฺมณปจฺจยา.\csromanspacer{}\csromanspacer{}สวิตกฺกสวิจารํ นเหตุํ ธมฺมํ ปฏิจฺจ สวิตกฺกสวิจาโร นเหตุ จ อวิตกฺกวิจารมตฺโต นเหตุ จ ธมฺมา อุปฺปชฺชนฺติ อารมฺมณปจฺจยา.\csromanspacer{} ตีณิ.}

\prose{อวิตกฺกวิจารมตฺตํ นเหตุํ ธมฺมํ ปฏิจฺจ อวิตกฺกวิจารมตฺโต นเหตุ ธมฺโม อุปฺปชฺชติ อารมฺมณปจฺจยา.\csromanspacer{} จตฺตาริ.}

\prose{อวิตกฺก-อวิจารํ นเหตุํ ธมฺมํ ปฏิจฺจ อวิตกฺก-อวิจาโร นเหตุ ธมฺโม อุปฺปชฺชติ อารมฺมณปจฺจยา.\csromanspacer{} ปญฺจ.}

\prose{สวิตกฺกสวิจารํ นเหตุญฺจ อวิตกฺก-อวิจารํ นเหตุญฺจ ธมฺมํ ปฏิจฺจ อวิตกฺก-อวิจาโร นเหตุ ธมฺโม อุปฺปชฺชติ อารมฺมณปจฺจยา.\csromanspacer{} ตีณิ.}

\prose{อวิตกฺกสวิจาโร นเหตุญฺจ อวิตกฺก-อวิจารํ นเหตุญฺจ ธมฺมํ ปฏิจฺจ สวิตกฺกสวิจาโร นเหตุ ธมฺโม อุปฺปชฺชติ อารมฺมณปจฺจยา.\csromanspacer{} จตฺตาริ.}

\csromanheader{2}{6. วิตกฺกตฺติก 1. เหตุทุก}
\prose{\csromanfolio{497}สวิตกฺกสวิจารํ นเหตุญฺจ อวิตกฺกวิจารมตฺตํ นเหตุญฺจ ธมฺมํ ปฏิจฺจ สวิตกฺกสวิจาโร นเหตุ ธมฺโม อุปฺปชฺชติ อารมฺมณปจฺจยา.\csromanspacer{} เอกํ.}

\prose{สวิตกฺกสวิจารํ นเหตุญฺจ อวิตกฺกวิจารมตฺตํ นเหตุญฺจ อวิตกฺก-อวิจารํ นเหตุญฺจ ธมฺมํ ปฏิจฺจ สวิตกฺกสวิจาโร นเหตุ ธมฺโม อุปฺปชฺชติ อารมฺมณปจฺจยา.\csromanspacer{} เอกํ. (สํขิตฺตํ.)}

\proseitem{93}{เหตุยา สตฺตติํส, อารมฺมเณ เอกวีส, อธิปติยา เตวีส, อนนฺตเร เอกวีส, สหชาเต สตฺตติํส, อญฺญมญฺเญ อฏฺฐวีส, นิสฺสเย สตฺตติํส, อุปนิสฺสเย เอกวีส, ปุเรชาเต เอกาทส, อาเสวเน เอกาทส, กมฺเม สตฺตติํส, วิปาเก อาหาเร อินฺทฺริเย ฌาเน มคฺเค สตฺตติํส, สมฺปยุตฺเต เอกวีส, วิปฺปยุตฺเต สตฺตติํส ฯเปฯ อวิคเต สตฺตติํส. (สํขิตฺตํ.)}

\prose{นเหตุยา เตตฺติํส, น-อารมฺมเณ สตฺต. (สํขิตฺตํ.)}

\csromanheader{2}{เหตุปจฺจยา น-อารมฺมเณ สตฺต. (สํขิตฺตํ.)}
\csromanheader{2}{นเหตุปจฺจยา อารมฺมเณ จุทฺทส. (สํขิตฺตํ.)}
\prose{(สหชาตวารมฺปิ ฯเปฯ สมฺปยุตฺตวารมฺปิ วิตฺถาเรตพฺพํ.)}

\csromanheader{2}{\textbf{นเหตุปท 7. ปญฺหาวาร}}
\csromanheader{2}{\textbf{ปจฺจยจตุกฺก-อารมฺมณ}}
\proseitem{94}{สวิตกฺกสวิจาโร นเหตุ ธมฺโม สวิตกฺกสวิจารสฺส นเหตุสฺส ธมฺมสฺส อารมฺมณปจฺจเยน ปจฺจโย.\csromanspacer{}\csromanspacer{}สวิตกฺกสวิจาโร นเหตุ ธมฺโม อวิตกฺกวิจารมตฺตสฺส นเหตุสฺส ธมฺมสฺส อารมฺมณปจฺจเยน ปจฺจโย.\csromanspacer{}\csromanspacer{}สวิตกฺกสวิจาโร นเหตุ ธมฺโม อวิตกฺก-อวิจารสฺส นเหตุสฺส ธมฺมสฺส อารมฺมณปจฺจเยน ปจฺจโย.\csromanspacer{}\csromanspacer{}สวิตกฺกสวิจาโร นเหตุ ธมฺโม สวิตกฺกสวิจารสฺส นเหตุสฺส จ อวิตกฺกวิจารมตฺตสฺส นเหตุสฺส จ ธมฺมสฺส อารมฺมณปจฺจเยน ปจฺจโย.\csromanspacer{} จตฺตาริ.}

\gambhira{ธมฺมานุโลม ติกทุกปฏฺฐานปาฬิ}
\proseitem{95}{\csromanfolio{498}อวิตกฺกวิจารมตฺโต นเหตุ ธมฺโม อวิตกฺกวิจารมตฺตสฺส นเหตุสฺส ธมฺมสฺส อารมฺมณปจฺจเยน ปจฺจโย.\csromanspacer{}\csromanspacer{}อวิตกฺกวิจารมตฺโต เนเหตุ ธมฺโม สวิตกฺกสวิจารสฺส นเหตุสฺส ธมฺมสฺส อารมฺมณปจฺจเยน ปจฺจโย.\csromanspacer{}\csromanspacer{}อวิตกฺกวิจารมตฺโต นเหตุ ธมฺโม อวิตกฺก-อวิจารสฺส นเหตุสฺส ธมฺมสฺส อารมฺมณปจฺจเยน ปจฺจโย.\csromanspacer{}\csromanspacer{}อวิตกฺกวิจารมตฺโต นเหตุ ธมฺโม สวิตกฺกสวิจารสฺส นเหตุสฺส จ อวิตกฺกวิจารมตฺตสฺส นเหตุสฺส จ ธมฺมสฺส อารมฺมณปจฺจเยน ปจฺจโย.\csromanspacer{} จตฺตาริ.}

\proseitem{96}{อวิตกฺก-อวิจาโร นเหตุ ธมฺโม อวิตกฺก-อวิจารสฺส นเหตุสฺส ธมฺมสฺส อารมฺมณปจฺจเยน ปจฺจโย.\csromanspacer{}\csromanspacer{}อวิตกฺก-อวิจาโร นเหตุ ธมฺโม สวิตกฺกสวิจารสฺส นเหตุสฺส ธมฺมสฺส อารมฺมณปจฺจเยน ปจฺจโย.\csromanspacer{}\csromanspacer{}อวิตกฺก-อวิจาโร นเหตุ ธมฺโม อวิตกฺกวิจารมตฺตสฺส นเหตุสฺส ธมฺมสฺส อารมฺมณปจฺจเยน ปจฺจโย.\csromanspacer{}\csromanspacer{}อวิตกฺก-อวิจาโร นเหตุ ธมฺโม อวิตกฺกวิจารมตฺตสฺส นเหตุสฺส จ อวิตกฺก-อวิจารสฺส นเหตุสฺส จ ธมฺมสฺส อารมฺมณปจฺจเยน ปจฺจโย.\csromanspacer{}\csromanspacer{}อวิตกฺก-อวิจาโร นเหตุ ธมฺโม สวิตกฺกสวิจารสฺส นเหตุสฺส จ อวิตกฺกวิจารมตฺตสฺส นเหตุสฺส จ ธมฺมสฺส อารมฺมณปจฺจเยน ปจฺจโย.\csromanspacer{} ปญฺจ.}

\proseitem{97}{อวิตกฺกวิจารมตฺโต นเหตุ จ อวิตกฺก-อวิจาโร นเหตุ จ ธมฺมา สวิตกฺกสวิจารสฺส นเหตุสฺส ธมฺมสฺส อารมฺมณปจฺจเยน ปจฺจโย.\csromanspacer{}\csromanspacer{}อวิตกฺกวิจารมตฺโต นเหตุ จ อวิตกฺก-อวิจาโร นเหตุ จ ธมฺมา อวิตกฺกวิจารมตฺตสฺส นเหตุสฺส ธมฺมสฺส อารมฺมณปจฺจเยน ปจฺจโย.\csromanspacer{}\csromanspacer{}อวิตกฺกวิจารมตฺโต นเหตุ จ อวิตกฺก-อวิจาโร นเหตุ จ ธมฺมา อวิตกฺก-อวิจารสฺส นเหตุสฺส ธมฺมสฺส อารมฺมณปจฺจเยน ปจฺจโย.\csromanspacer{}\csromanspacer{}อวิตกฺกวิจารมตฺโต นเหตุ จ อวิตกฺก-อวิจาโร นเหตุ จ ธมฺมา สวิตกฺกสวิจารสฺส นเหตุสฺส จ อวิตกฺกวิจารมตฺตสฺส นเหตุสฺส จ ธมฺมสฺส อารมฺมณปจฺจเยน ปจฺจโย.\csromanspacer{} จตฺตาริ.}

\proseitem{98}{สวิตกฺกสวิจาโร นเหตุ จ อวิตกฺกวิจารมตฺโต นเหตุ จ ธมฺมา สวิตกฺกสวิจารสฺส นเหตุสฺส ธมฺมสฺส อารมฺมณปจฺจเยน ปจฺจโย.\csromanspacer{}\csromanspacer{}สวิตกฺกสวิจาโร นเหตุ จ อวิตกฺกวิจารมตฺโต นเหตุ จ ธมฺมา อวิตกฺกวิจารมตฺตสฺส นเหตุสฺส ธมฺมสฺส อารมฺมณปจฺจเยน ปจฺจโย.\csromanspacer{}\csromanspacer{}}

\vspace{\tipitakaparskip}
\csromancenter{วิตกฺกตฺติก 1.\csromanspacer{} เหตุทุก สวิตกฺกสวิจาโร นเหตุ จ อวิตกฺกวิจารมตฺโต นเหตุ จ ธมฺมา อวิตกฺก-อวิจารสฺส นเหตุสฺส ธมฺมสฺส อารมฺมณปจฺจเยน ปจฺจโย.\csromanspacer{}\csromanspacer{}สวิตกฺกสวิจาโร นเหตุ จ อวิตกฺกวิจารมตฺโต นเหตุ จ ธมฺมา สวิตกฺกสวิจารสฺส นเหตุสฺส จ อวิตกฺกวิจารมตฺตสฺส นเหตุสฺส จ ธมฺมสฺส อารมฺมณปจฺจเยน ปจฺจโย.\csromanspacer{} จตฺตาริ.}
\csromanheader{2}{\textbf{อธิปติ}}
\proseitem{99}{\csromanfolio{499}สวิตกฺกสวิจาโร นเหตุ ธมฺโม สวิตกฺกสวิจารสฺส นเหตุสฺส ธมฺมสฺส อธิปติปจฺจเยน ปจฺจโย.\csromanspacer{} สตฺต.}

\prose{อวิตกฺกวิจารมตฺโต นเหตุ ธมฺโม อวิตกฺกวิจารมตฺตสฺส นเหตุสฺส ธมฺมสฺส อธิปติปจฺจเยน ปจฺจโย.\csromanspacer{} ปญฺจ.}

\prose{อวิตกฺก-อวิจาโร นเหตุ ธมฺโม อวิตกฺก-อวิจารสฺส นเหตุสฺส ธมฺมสฺส อธิปติปจฺจเยน ปจฺจโย.\csromanspacer{} ปญฺจ.}

\prose{อวิตกฺกวิจารมตฺโต นเหตุ จ อวิตกฺก-อวิจาโร นเหตุ จ ธมฺมา สวิตกฺกสวิจารสฺส นเหตุสฺส ธมฺมสฺส อธิปติปจฺจเยน ปจฺจโย.\csromanspacer{} ตีณิ. สวิตกฺกสวิจาโร นเหตุ จ อวิตกฺกวิจารมตฺโต นเหตุ จ ธมฺมา สวิตกฺกสวิจารสฺส นเหตุสฺส ธมฺมสฺส อธิปติปจฺจเยน ปจฺจโย.\csromanspacer{} ตีณิ.}

\csromanheader{2}{\textbf{อนนฺตราทิ}}
\proseitem{100}{สวิตกฺกสวิจาโร นเหตุ ธมฺโม สวิตกฺกสวิจารสฺส นเหตุสฺส ธมฺมสฺส อนนฺตรปจฺจเยน ปจฺจโย ฯเปฯ.\csromanspacer{} สมนนฺตรปจฺจเยน ปจฺจโย.}

\prose{สวิตกฺกสวิจาโร นเหตุ ธมฺโม สวิตกฺกสวิจารสฺส นเหตุสฺส ธมฺมสฺส อุปนิสฺสยปจฺจเยน ปจฺจโย. (สํขิตฺตํ.)}

\proseitem{101}{อารมฺมเณ เอกวีส, อธิปติยา เตวีส, อนนฺตเร ปญฺจวีส, สมนนฺตเร ปญฺจวีส, สหชาเต ติํส, อญฺญมญฺเญ อฏฺฐวีส, นิสฺสเย ติํส, อุปนิสฺสเย ปญฺจวีส, ปุเรชาเต ปญฺจ, ปจฺฉาชาเต ปญฺจ, อาเสวเน เอกวีส, กมฺเม เอกาทส, วิปาเก เอกวีส, อาหาเร เอกาทส ฯเปฯ อวิคเต ติํส. (สํขิตฺตํ.)}

\gambhira{ธมฺมานุโลม ติกทุกปฏฺฐานปาฬิ}
\prose{\csromanfolio{500}นเหตุยา ปญฺจติํส, น-อารมฺมเณ ปญฺจติํส, น-อธิปติยา ปญฺจติํส. (สํขิตฺตํ.)}

\csromanheader{2}{อารมฺมณปจฺจยา นเหตุยา เอกวีส. (สํขิตฺตํ.)}
\csromanheader{2}{นเหตุปจฺจยา อารมฺมเณ เอกวีส. (สํขิตฺตํ.)}
\prose{(ยถา กุสลตฺติเก ปญฺหาวารํ, เอวํ วิตฺถาเรตพฺพํ.)}

\csromanmatikaanchor{mtk.104}
\csromanmatikaanchor{mtk.105}
\csromantocmarknopage{subparagraph}{7. ปีติตฺติก}{mtk.104}
\bookmark[dest={mtk.104},level=5]{7. ปีติตฺติก}
\csromantocmarkref{subparagraph}{1. เหตุทุก}{mtk.105}
\bookmark[dest={mtk.105},level=5]{1. เหตุทุก}
\csromanheader{6}{7. ปีติตฺติก 1. เหตุทุก}
\csromanheader{2}{1. ปฏิจฺจวาร ปจฺจยจตุกฺก}
\proseitem{102}{ปีติสหคตํ เหตุํ ธมฺมํ ปฏิจฺจ ปีติสหคโต เหตุ ธมฺโม อุปฺปชฺชติ เหตุปจฺจยา.\csromanspacer{} ตีณิ.}

\prose{สุขสหคตํ เหตุํ ธมฺมํ ปฏิจฺจ สุขสหคโต เหตุ ธมฺโม อุปฺปชฺชติ เหตุปจฺจยา.\csromanspacer{} ตีณิ.}

\prose{อุเปกฺขาสหคตํ เหตุํ ธมฺมํ ปฏิจฺจ อุเปกฺขาสหคโต เหตุ ธมฺโม อุปฺปชฺชติ เหตุปจฺจยา.\csromanspacer{} เอกํ.}

\prose{ปีติสหคตํ เหตุญฺจ สุขสหคตํ เหตุญฺจ ธมฺมํ ปฏิจฺจ ปีติสหคโต เหตุ ธมฺโม อุปฺปชฺชติ เหตุปจฺจยา.\csromanspacer{} ตีณิ. (สํขิตฺตํ.)}

\proseitem{103}{เหตุยา ทส, อารมฺมเณ ทส ฯเปฯ อวิคเต ทส. (สํขิตฺตํ.)}

\prose{น-อธิปติยา ทส, นปุเรชาเต ทส, นปจฺฉาชาเต ทส, น-อาเสวเน ทส, นวิปาเก ทส, นวิปฺปยุตฺเต ทส. (สํขิตฺตํ.)}

\prose{(สหชาตวารมฺปิ ฯเปฯ สมฺปยุตฺตวารมฺปิ วิตฺถาเรตพฺพํ.)}

\csromanheader{2}{7. ปญฺหาวาร ปจฺจยจตุกฺก}
\csromanheader{2}{\textbf{เหตุ อารมฺมณาทิ}}
\proseitem{104}{ปีติสหคโต เหตุ ธมฺโม ปีติสหคตสฺส เหตุสฺส ธมฺมสฺส เหตุปจฺจเยน ปจฺจโย.\csromanspacer{} ตีณิ.}

\csromanheader{2}{7. ปีติตฺติก 1. เหตุทุก}
\prose{\csromanfolio{501}สุขสหคโต เหตุ ธมฺโม สุขสหคตสฺส เหตุสฺส ธมฺมสฺส เหตุปจฺจเยน ปจฺจโย.\csromanspacer{} ตีณิ.}

\prose{อุเปกฺขาสหคโต เหตุ ธมฺโม อุเปกฺขาสหคตสฺส เหตุสฺส ธมฺมสฺส เหตุปจฺจเยน ปจฺจโย.\csromanspacer{} เอกํ.}

\prose{ปีติสหคโต เหตุ จ สุขสหคโต เหตุ จ ธมฺมา ปีติสหคตสฺส เหตุสฺส ธมฺมสฺส เหตุปจฺจเยน ปจฺจโย.\csromanspacer{} ตีณิ.}

\proseitem{105}{ปีติสหคโต เหตุ ธมฺโม ปีติสหคตสฺส เหตุสฺส ธมฺมสฺส อารมฺมณปจฺจเยน ปจฺจโย.\csromanspacer{}\csromanspacer{}ปีติสหคโต เหตุ ธมฺโม สุขสหคตสฺส เหตุสฺส ธมฺมสฺส อารมฺมณปจฺจเยน ปจฺจโย.\csromanspacer{}\csromanspacer{}ปีติสหคโต เหตุ ธมฺโม อุเปกฺขาสหคตสฺส เหตุสฺส ธมฺมสฺส อารมฺมณปจฺจเยน ปจฺจโย.\csromanspacer{}\csromanspacer{}ปีติสหคโต เหตุ ธมฺโม ปีติสหคตสฺส เหตุสฺส จ สุขสหคตสฺส เหตุสฺส จ ธมฺมสฺส อารมฺมณปจฺจเยน ปจฺจโย.\csromanspacer{} จตฺตาริ.}

\prose{สุขสหคโต เหตุ ธมฺโม สุขสหคตสฺส เหตุสฺส ธมฺมสฺส อารมฺมณปจฺจเยน ปจฺจโย.\csromanspacer{}\csromanspacer{}สุขสหคโต เหตุ ธมฺโม ปีติสหคตสฺส เหตุสฺส ธมฺมสฺส อารมฺมณปจฺจเยน ปจฺจโย.\csromanspacer{}\csromanspacer{}สุขสหคโต เหตุ ธมฺโม อุเปกฺขสหคตสฺส เหตุสฺส ธมฺมสฺส อารมฺมณปจฺจเยน ปจฺจโย.\csromanspacer{}\csromanspacer{}สุขสหคโต เหตุ ธมฺโม ปีติสหคตสฺส เหตุสฺส จ สุขสหคตสฺส เหตุสฺส จ ธมฺมสฺส อารมฺมณปจฺจเยน ปจฺจโย.\csromanspacer{} จตฺตาริ.}

\prose{อุเปกฺขาสหคโต เหตุ ธมฺโม อุเปกฺขาสหคตสฺส เหตุสฺส ธมฺมสฺส อารมฺมณปจฺจเยน ปจฺจโย.\csromanspacer{}\csromanspacer{}อุเปกฺขาสหคโต เหตุ ธมฺโม ปีติสหคตสฺส เหตุสฺส ธมฺมสฺส อารมฺมณปจฺจเยน ปจฺจโย.\csromanspacer{}\csromanspacer{}อุเปกฺขาสหคโต เหตุ ธมฺโม สุขสหคตสฺส เหตุสฺส ธมฺมสฺส อารมฺมณปจฺจเยน ปจฺจโย.\csromanspacer{}\csromanspacer{}อุเปกฺขาสหคโต เหตุ ธมฺโม ปีติสหคตสฺส เหตุสฺส จ สุขสหคตสฺส เหตุสฺส จ ธมฺมสฺส อารมฺมณปจฺจเยน ปจฺจโย.\csromanspacer{} จตฺตาริ.}

\prose{ปีติสหคโต เหตุ จ สุขสหคโต เหตุ จ ธมฺมา ปีติสหคตสฺส เหตุสฺส ธมฺมสฺส อารมฺมณปจฺจเยน ปจฺจโย.\csromanspacer{} จตฺตาริ.}

\gambhira{ธมฺมานุโลม ติกทุกปฏฺฐานปาฬิ}
\prose{\csromanfolio{502}ปีติสหคโต เหตุ ธมฺโม ปีติสหคตสฺส เหตุสฺส ธมฺมสฺส อธิปติปจฺจเยน ปจฺจโย.}

\proseitem{106}{เหตุยา ทส, อารมฺมเณ โสฬส, อธิปติยา โสฬส, อนนฺตเร โสฬส, สมนนฺตเร โสฬส, สหชาเต ทส, อญฺญมญฺเญ ทส, นิสฺสเย ทส, อุปนิสฺสเย โสฬส ฯเปฯ วิปาเก ทส ฯเปฯ อวิคเต ทส. (สํขิตฺตํ.)}

\prose{นเหตุยา โสฬส, น-อารมฺมเณ โสฬส. (สํขิตฺตํ.)}

\csromanheader{2}{เหตุปจฺจยา น-อารมฺมเณ ทส. (สํขิตฺตํ.)}
\csromanheader{2}{นเหตุปจฺจยา อารมฺมเณ โสฬส. (สํขิตฺตํ.)}
\prose{(ยถา กุสลตฺติเก ปญฺหาวารํ, เอวํ วิตฺถาเรตพฺพํ.)}

\csromanheader{2}{\textbf{นเหตุปท 1. ปฏิจฺจวาร}}
\csromanheader{2}{\textbf{ปจฺจยจตุกฺก-เหตุ อารมฺมณ}}
\proseitem{107}{ปีติสหคตํ นเหตุํ ธมฺมํ ปฏิจฺจ ปีติสหคโต นเหตุ ธมฺโม อุปฺปชฺชติ เหตุปจฺจยา.\csromanspacer{}\csromanspacer{}ปีติสหคตํ นเหตุํ ธมฺมํ ปฏิจฺจ สุขสหคโต นเหตุ ธมฺโม อุปฺปชฺชติ เหตุปจฺจยา.\csromanspacer{}\csromanspacer{}ปีติสหคตํ นเหตุํ ธมฺมํ ปฏิจฺจ ปีติสหคโต นเหตุ จ สุขสหคโต นเหตุ จ ธมฺมา อุปฺปชฺชนฺติ เหตุปจฺจยา. (3)}

\prose{สุขสหคตํ นเหตุํ ธมฺมํ ปฏิจฺจ สุขสหคโต นเหตุ ธมฺโม อุปฺปชฺชติ เหตุปจฺจยา.\csromanspacer{}\csromanspacer{}สุขสหคตํ นเหตุํ ธมฺมํ ปฏิจฺจ ปีติสหคโต นเหตุ ธมฺโม อุปฺปชฺชติ เหตุปจฺจยา.\csromanspacer{}\csromanspacer{}สุขสหคตํ นเหตุํ ธมฺมํ ปฏิจฺจ ปีติสหคโต นเหตุ จ สุขสหคโต นเหตุ จ ธมฺมา อุปฺปชฺชนฺติ เหตุปจฺจยา. (3)}

\prose{อุเปกฺขาสหคตํ นเหตุํ ธมฺมํ ปฏิจฺจ อุเปกฺขาสหคโต นเหตุ ธมฺโม อุปฺปชฺชติ เหตุปจฺจยา. (1)}

\prose{ปีติสหคตํ นเหตุญฺจ สุขสหคตํ นเหตุญฺจ ธมฺมํ ปฏิจฺจ ปีติสหคโต นเหตุ ธมฺโม อุปฺปชฺชติ เหตุปจฺจยา.\csromanspacer{}\csromanspacer{}ปีติสหคตํ นเหตุญฺจ สุขสหคตํ นเหตุญฺจ ธมฺมํ ปฏิจฺจ สุขสหคโต นเหตุ}

\vspace{\tipitakaparskip}
\csromancenter{ปีติตฺติก 1.\csromanspacer{} เหตุทุก ธมฺโม อุปฺปชฺชติ เหตุปจฺจยา.\csromanspacer{}\csromanspacer{}ปีติสหคตํ นเหตุญฺจ สุขสหคตํ นเหตุญฺจ ธมฺมํ ปฏิจฺจ ปีติสหคโต นเหตุ จ สุขสหคโต นเหตุ จ ธมฺมา อุปฺปชฺชนฺติ เหตุปจฺจยา. (3)}
\proseitem{108}{\csromanfolio{503}ปีติสหคตํ นเหตุํ ธมฺมํ ปฏิจฺจ ปีติสหคโต นเหตุ ธมฺโม อุปฺปชฺชติ อารมฺมณปจฺจยา.\csromanspacer{} ตีณิ.}

\prose{สุขสหคตํ นเหตุํ ธมฺมํ ปฏิจฺจ สุขสหคโต นเหตุ ธมฺโม อุปฺปชฺชติ อารมฺมณปจฺจยา.\csromanspacer{} ตีณิ.}

\prose{อุเปกฺขาสหคตํ นเหตุํ ธมฺมํ ปฏิจฺจ อุเปกฺขาสหคโต นเหตุ ธมฺโม อุปฺปชฺชติ อารมฺมณปจฺจยา.\csromanspacer{} เอกํ.}

\prose{ปีติสหคตํ นเหตุญฺจ สุขสหคตํ นเหตุญฺจ ธมฺมํ ปฏิจฺจ ปีติสหคโต นเหตุ ธมฺโม อุปฺปชฺชติ อารมฺมณปจฺจยา.\csromanspacer{} ตีณิ. (สํขิตฺตํ.)}

\proseitem{109}{เหตุยา ทส, อารมฺมเณ ทส, อธิปติยา ทส, อนนฺตเร ทส ฯเปฯ อวิคเต ทส. (สํขิตฺตํ.)}

\prose{นเหตุยา ทส, น-อธิปติยา ทส. (สํขิตฺตํ.\csromanspacer{} สหชาตวารมฺปิ ฯเปฯ สมฺปยุตฺตวารมฺปิ วิตฺถาเรตพฺพํ.)}

\csromanheader{2}{\textbf{นเหตุปท 7. ปญฺหาวาร}}
\csromanheader{2}{\textbf{ปจฺจยจตุกฺก-อารมฺมณาทิ}}
\proseitem{110}{ปีติสหคโต นเหตุ ธมฺโม ปีติสหคตสฺส นเหตุสฺส ธมฺมสฺส อารมฺมณปจฺจเยน ปจฺจโย.\csromanspacer{}\csromanspacer{}ปีติสหคโต นเหตุ ธมฺโม สุขสหคตสฺส นเหตุสฺส ธมฺมสฺส อารมฺมณปจฺจเยน ปจฺจโย.\csromanspacer{}\csromanspacer{}ปีติสหคโต นเหตุ ธมฺโม อุเปกฺขาสหคตสฺส นเหตุสฺส ธมฺมสฺส อารมฺมณปจฺจเยน ปจฺจโย.\csromanspacer{}\csromanspacer{}ปีติสหคโต นเหตุ ธมฺโม ปีติสหคตสฺส นเหตุสฺส จ สุขสหคตสฺส นเหตุสฺส จ ธมฺมสฺส อารมฺมณปจฺจเยน ปจฺจโย. (4)}

\prose{สุขสหคโต นเหตุ ธมฺโม สุขสหคตสฺส นเหตุสฺส ธมฺมสฺส อารมฺมณปจฺจเยน ปจฺจโย.\csromanspacer{}\csromanspacer{}สุขสหคโต นเหตุ ธมฺโม ปีติสหคตสฺส นเหตุสฺส ธมฺมสฺส อารมฺมณปจฺจเยน ปจฺจโย.\csromanspacer{}\csromanspacer{}}

\vspace{\tipitakaparskip}
\csromancenter{ธมฺมานุโลม ติกทุกปฏฺฐานปาฬิ สุขสหคโต นเหตุ ธมฺโม อุเปกฺขาสหคตสฺส นเหตุสฺส ธมฺมสฺส อารมฺมณปจฺจเยน ปจฺจโย.\csromanspacer{}\csromanspacer{}สุขสหคโต นเหตุ ธมฺโม ปีติสหคตสฺส นเหตุสฺส จ สุขสหคตสฺส นเหตุสฺส จ ธมฺมสฺส อารมฺมณปจฺจเยน ปจฺจโย. (4)}
\prose{\csromanfolio{504}อุเปกฺขาสหคโต นเหตุ ธมฺโม อุเปกฺขาสหคตสฺส นเหตุสฺส ธมฺมสฺส อารมฺมณปจฺจเยน ปจฺจโย.\csromanspacer{}\csromanspacer{}อุเปกฺขาสหคโต นเหตุ ธมฺโม ปีติสหคตสฺส นเหตุสฺส ธมฺมสฺส อารมฺมณปจฺจเยน ปจฺจโย.\csromanspacer{}\csromanspacer{}อุเปกฺขาสหคโต นเหตุ ธมฺโม สุขสหคตสฺส นเหตุสฺส ธมฺมสฺส อารมฺมณปจฺจเยน ปจฺจโย.\csromanspacer{}\csromanspacer{}อุเปกฺขาสหคโต นเหตุ ธมฺโม ปีติสหคตสฺส นเหตุสฺส จ สุขสหคตสฺส นเหตุสฺส จ ธมฺมสฺส อารมฺมณปจฺจเยน ปจฺจโย. (4)}

\prose{ปีติสหคโต นเหตุ จ สุขสหคโต นเหตุ จ ธมฺมา ปีติสหคตสฺส นเหตุสฺส ธมฺมสฺส อารมฺมณปจฺจเยน ปจฺจโย.\csromanspacer{} จตฺตาริ.}

\proseitem{111}{ปีติสหคโต นเหตุ ธมฺโม ปีติสหคตสฺส นเหตุสฺส ธมฺมสฺส อธิปติปจฺจเยน ปจฺจโย.\csromanspacer{} จตฺตาริ. สุขสหคโต นเหตุ ธมฺโม สุขสหคตสฺส นเหตุสฺส ธมฺมสฺส อธิปติปจฺจเยน ปจฺจโย.\csromanspacer{} จตฺตาริ.}

\prose{อุเปกฺขาสหคโต นเหตุ ธมฺโม อุเปกฺขาสหคตสฺส นเหตุสฺส ธมฺมสฺส อธิปติปจฺจเยน ปจฺจโย.\csromanspacer{} จตฺตาริ.}

\prose{ปีติสหคโต นเหตุ จ สุขสหคโต นเหตุ จ ธมฺมา ปีติสหคตสฺส นเหตุสฺส ธมฺมสฺส อธิปติปจฺจเยน ปจฺจโย.\csromanspacer{} จตฺตาริ.}

\proseitem{112}{ปีติสหคโต นเหตุ ธมฺโม ปีติสหคตสฺส นเหตุสฺส ธมฺมสฺส อนนฺตรปจฺจเยน ปจฺจโย ฯเปฯ อุปนิสฺสยปจฺจเยน ปจฺจโย.\csromanspacer{} จตฺตาริ. สุขสหคโต นเหตุ ธมฺโม สุขสหคตสฺส นเหตุสฺส ธมฺมสฺส อุปนิสฺสยปจฺจเยน ปจฺจโย.\csromanspacer{} จตฺตาริ.}

\prose{อุเปกฺขาสหคโต นเหตุ ธมฺโม อุเปกฺขาสหคตสฺส นเหตุสฺส ธมฺมสฺส อุปนิสฺสยปจฺจเยน ปจฺจโย.\csromanspacer{} จตฺตาริ.}

\csromanmatikaanchor{mtk.106}
\csromanmatikaanchor{mtk.107}
\csromantocmarknopage{subparagraph}{8. ทสฺสเนนปหาตพฺพตฺติก}{mtk.106}
\bookmark[dest={mtk.106},level=5]{8. ทสฺสเนนปหาตพฺพตฺติก}
\csromantocmarkref{subparagraph}{1. เหตุทุก}{mtk.107}
\bookmark[dest={mtk.107},level=5]{1. เหตุทุก}
\csromanheader{6}{8. ทสฺสเนนปหาตพฺพตฺติก 1. เหตุทุก}
\prose{\csromanfolio{505}ปีติสหคโต นเหตุ จ สุขสหคโต นเหตุ จ ธมฺมา ปีติสหคตสฺส นเหตุสฺส ธมฺมสฺส อุปนิสฺสยปจฺจเยน ปจฺจโย.\csromanspacer{} จตฺตาริ. (สํขิตฺตํ.)}

\proseitem{113}{อารมฺมเณ โสฬส, อธิปติยา โสฬส, สหชาเต ทส, อญฺญมญฺเญ ทส, นิสฺสเย ทส, อุปนิสฺสเย โสฬส, อาหาเร ทส ฯเปฯ อวิคเต ทส. (สํขิตฺตํ.)}

\prose{นเหตุยา โสฬส, น-อารมฺมเณ โสฬส, น-อธิปติยา โสฬส. (สํขิตฺตํ.)}

\csromanheader{2}{อารมฺมณปจฺจยา นเหตุยา โสฬส. (สํขิตฺตํ.)}
\csromanheader{2}{นเหตุปจฺจยา อารมฺมเณ โสฬส. (สํขิตฺตํ.)}
\prose{(ยถา กุสลตฺติเก ปญฺหาวารํ, เอวํ วิตฺถาเรตพฺพํ.)}

\csromanheader{2}{8. ทสฺสเนนปหาตพฺพตฺติก 1. เหตุทุก}
\csromanheader{2}{1-7. ปฏิจฺจาทิวาร ปจฺจยจตุกฺก}
\proseitem{114}{ทสฺสเนน ปหาตพฺพํ เหตุํ ธมฺมํ ปฏิจฺจ ทสฺสเนน ปหาตพฺโพ เหตุ ธมฺโม อุปฺปชฺชติ เหตุปจฺจยา. (1)}

\prose{ภาวนาย ปหาตพฺพํ เหตุํ ธมฺมํ ปฏิจฺจ ภาวนาย ปหาตพฺโพ เหตุ ธมฺโม อุปฺปชฺชติ เหตุปจฺจยา. (1)}

\prose{เนวทสฺสเนน นภาวนาย ปหาตพฺพํ เหตุํ ธมฺมํ ปฏิจฺจ เนวทสฺสเนน นภาวนาย ปหาตพฺโพ เหตุ ธมฺโม อุปฺปชฺชติ เหตุปจฺจยา. (สํขิตฺตํ.)}

\proseitem{115}{เหตุยา ตีณิ, (สพฺพตฺถ ตีณิ.) วิปาเก เอกํ ฯเปฯ อวิคเต ตีณิ.}

\prose{น-อธิปติยา ตีณิ, นปุเรชาเต ตีณิ ฯเปฯ นวิปฺปยุตฺเต ตีณิ. (สํขิตฺตํ.\csromanspacer{} สหชาตวารมฺปิ ฯเปฯ สมฺปยุตฺตวารมฺปิ วิตฺถาเรตพฺพํ.)}

\proseitem{116}{ทสฺสเนน ปหาตพฺโพ เหตุ ธมฺโม ทสฺสเนน ปหาตพฺพสฺส เหตุสฺส ธมฺมสฺส เหตุปจฺจเยน ปจฺจโย.\csromanspacer{} ตีณิ.}

\gambhira{ธมฺมานุโลม ติกทุกปฏฺฐานปาฬิ}
\prose{\csromanfolio{506}ทสฺสเนน ปหาตพฺโพ เหตุ ธมฺโม ทสฺสเนน ปหาตพฺพสฺส เหตุสฺส ธมฺมสฺส อารมฺมณปจฺจเยน ปจฺจโย. (สํขิตฺตํ.)}

\proseitem{117}{เหตุยา ตีณิ, อารมฺมเณ อฏฺฐ, อธิปติยา ฉ, อนนฺตเร ปญฺจ, สมนนฺตเร ปญฺจ, สหชาเต ตีณิ, อญฺญมญฺเญ ตีณิ, นิสฺสเย ตีณิ, อุปนิสฺสเย อฏฺฐ, อาเสวเน ตีณิ, วิปาเก เอกํ, อินฺทฺริเย เอกํ, มคฺเค เอกํ ฯเปฯ อวิคเต ตีณิ. (สํขิตฺตํ.)}

\prose{นเหตุยา อฏฺฐ, น-อารมฺมเณ อฏฺฐ. (สํขิตฺตํ.)}

\csromanheader{2}{เหตุปจฺจยา น-อารมฺมเณ ตีณิ. (สํขิตฺตํ.)}
\csromanheader{2}{นเหตุปจฺจยา อารมฺมเณ อฏฺฐ. (สํขิตฺตํ.)}
\prose{(ยถา กุสลตฺติเก ปญฺหาวารํ, เอวํ วิตฺถาเรตพฺพํ.)}

\csromanheader{2}{\textbf{นเหตุปท 1. ปฏิจฺจวาร}}
\csromanheader{2}{\textbf{ปจฺจยจตุกฺก-เหตุ อารมฺมณ}}
\proseitem{118}{ทสฺสเนน ปหาตพฺพํ นเหตุํ ธมฺมํ ปฏิจฺจ ทสฺสเนน ปหาตพฺโพ นเหตุ ธมฺโม อุปฺปชฺชติ เหตุปจฺจยา.\csromanspacer{}\csromanspacer{}ทสฺสเนน ปหาตพฺพํ นเหตุํ ธมฺมํ ปฏิจฺจ เนวทสฺสเนน นภาวนาย ปหาตพฺโพ นเหตุ ธมฺโม อุปฺปชฺชติ เหตุปจฺจยา.\csromanspacer{}\csromanspacer{}ทสฺสเนน ปหาตพฺพํ นเหตุํ ธมฺมํ ปฏิจฺจ ทสฺสเนน ปหาตพฺโพ นเหตุ จ เนวทสฺสเนน นภาวนาย ปหาตพฺโพ นเหตุ จ ธมฺมา อุปฺปชฺชนฺติ เหตุปจฺจยา.\csromanspacer{} ตีณิ.}

\prose{ภาวนาย ปหาตพฺพํ นเหตุํ ธมฺมํ ปฏิจฺจ ภาวนาย ปหาตพฺโพ นเหตุ ธมฺโม อุปฺปชฺชติ เหตุปจฺจยา.\csromanspacer{}\csromanspacer{}ภาวนาย ปหาตพฺพํ นเหตุํ ธมฺมํ ปฏิจฺจ เนวทสฺสเนน นภาวนาย ปหาตพฺโพ นเหตุ ธมฺโม อุปฺปชฺชติ เหตุปจฺจยา.\csromanspacer{}\csromanspacer{}ภาวนาย ปหาตพฺพํ นเหตุํ ธมฺมํ ปฏิจฺจ ภาวนาย ปหาตพฺโพ นเหตุ จ เนวทสฺสเนน นภาวนาย ปหาตพฺโพ นเหตุ จ ธมฺมา อุปฺปชฺชนฺติ เหตุปจฺจยา.\csromanspacer{} ตีณิ.}

\prose{เนวทสฺสเนน นภาวนาย ปหาตพฺพํ นเหตุํ ธมฺมํ ปฏิจฺจ เนวทสฺสเนน นภาวนาย ปหาตพฺโพ นเหตุ ธมฺโม อุปฺปชฺชติ เหตุปจฺจยา. (1)}

\csromanheader{2}{8. ทสฺสเนนปหาตพฺพตฺติก 1. เหตุทุก}
\prose{\csromanfolio{507}ทสฺสเนน ปหาตพฺพํ นเหตุญฺจ เนวทสฺสเนน นภาวนาย ปหาตพฺพํ นเหตุญฺจ ธมฺมํ ปฏิจฺจ เนวทสฺสเนน นภาวนาย ปหาตพฺโพ นเหตุ ธมฺโม อุปฺปชฺชติ เหตุปจฺจยา. (1)}

\prose{ภาวนาย ปหาตพฺพํ นเหตุญฺจ เนวทสฺสเนน นภาวนาย ปหาตพฺพํ นเหตุญฺจ ธมฺมํ ปฏิจฺจ เนวทสฺสเนน นภาวนาย ปหาตพฺโพ นเหตุ ธมฺโม อุปฺปชฺชติ เหตุปจฺจยา. (1)}

\prose{ทสฺสเนน ปหาตพฺพํ นเหตุํ ธมฺมํ ปฏิจฺจ ทสฺสเนน ปหาตพฺโพ นเหตุ ธมฺโม อุปฺปชฺชติ อารมฺมณปจฺจยา. (1)}

\prose{ภาวนาย ปหาตพฺพํ นเหตุํ ธมฺมํ ปฏิจฺจ ภาวนาย ปหาตพฺโพ นเหตุ ธมฺโม อุปฺปชฺชติ อารมฺมณปจฺจยา. (1)}

\prose{เนวทสฺสเนน นภาวนาย ปหาตพฺพํ นเหตุํ ธมฺมํ ปฏิจฺจ เนวทสฺสเนน นภาวนาย ปหาตพฺโพ นเหตุ ธมฺโม อุปฺปชฺชติ อารมฺมณปจฺจยา. (สํขิตฺตํ.)}

\proseitem{119}{เหตุยา นว, อารมฺมเณ ตีณิ, อธิปติยา นว ฯเปฯ อวิคเต นว. (สํขิตฺตํ.)}

\csromanheader{2}{\textbf{นเหตุ น-อารมฺมณ}}
\proseitem{120}{เนวทสฺสเนน นภาวนาย ปหาตพฺพํ นเหตุํ ธมฺมํ ปฏิจฺจ เนวทสฺสเนน นภาวนาย ปหาตพฺโพ นเหตุ ธมฺโม อุปฺปชฺชติ นเหตุปจฺจยา.}

\prose{ทสฺสเนน ปหาตพฺพํ นเหตุํ ธมฺมํ ปฏิจฺจ เนวทสฺสเนน นภาวนาย ปหาตพฺโพ นเหตุ ธมฺโม อุปฺปชฺชติ น-อารมฺมณปจฺจยา. (1)}

\prose{ภาวนาย ปหาตพฺพํ นเหตุํ ธมฺมํ ปฏิจฺจ เนวทสฺสเนน นภาวนาย ปหาตพฺโพ นเหตุ ธมฺโม อุปฺปชฺชติ น-อารมฺมณปจฺจยา. (1)}

\prose{เนวทสฺสเนน นภาวนาย ปหาตพฺพํ นเหตุํ ธมฺมํ ปฏิจฺจ เนวทสฺสเนน นภาวนาย ปหาตพฺโพ นเหตุ ธมฺโม อุปฺปชฺชติ น-อารมฺมณปจฺจยา. (1)}

\prose{ทสฺสเนน ปหาตพฺพํ นเหตุญฺจ เนวทสฺสเนน นภาวนาย ปหาตพฺพํ นเหตุญฺจ ธมฺมํ ปฏิจฺจ เนวทสฺสเนน นภาวนาย ปหาตพฺโพ นเหตุ ธมฺโม อุปฺปชฺชติ น-อารมฺมณปจฺจยา. (1)}

\gambhira{ธมฺมานุโลม ติกทุกปฏฺฐานปาฬิ}
\prose{\csromanfolio{508}ภาวนาย ปหาตพฺพํ นเหตุญฺจ เนวทสฺสเนน นภาวนาย ปหาตพฺพํ นเหตุญฺจ ธมฺมํ ปฏิจฺจ เนวทสฺสเนน นภาวนาย ปหาตพฺโพ นเหตุ ธมฺโม อุปฺปชฺชติ น-อารมฺมณปจฺจยา. (1) (สํขิตฺตํ.)}

\proseitem{121}{นเหตุยา เอกํ, น-อารมฺมเณ ปญฺจ, น-อธิปติยา นว ฯเปฯ โนวิคเต ปญฺจ. (สํขิตฺตํ.)}

\csromanheader{2}{เหตุปจฺจยา น-อารมฺมเณ ปญฺจ. (สํขิตฺตํ.)}
\csromanheader{2}{นเหตุปจฺจยา อารมฺมเณ เอกํ. (สํขิตฺตํ.)}
\prose{(สหชาตวารมฺปิ ปจฺจยวารมฺปิ นิสฺสยวารมฺปิ สํสฏฺฐวารมฺปิ สมฺปยุตฺตวารมฺปิ ปฏิจฺจวารสทิสํ วิตฺถาเรตพฺพํ.)}

\csromanheader{2}{7. ปญฺหาวาร ปจฺจยจตุกฺก}
\proseitem{122}{ทสฺสเนน ปหาตพฺโพ นเหตุ ธมฺโม ทสฺสเนน ปหาตพฺพสฺส นเหตุสฺส ธมฺมสฺส อารมฺมณปจฺจเยน ปจฺจโย.\csromanspacer{}\csromanspacer{}ทสฺสเนน ปหาตพฺโพ เหตุ ธมฺโม เนวทสฺสเนน นภาวนาย ปหาตพฺพสฺส นเหตุสฺส ธมฺมสฺส อารมฺมณปจฺจเยน ปจฺจโย.\csromanspacer{} เทฺว.}

\prose{ภาวนาย ปหาตพฺโพ นเหตุ ธมฺโม ภาวนาย ปหาตพฺพสฺส นเหตุสฺส ธมฺมสฺส อารมฺมณปจฺจเยน ปจฺจโย.\csromanspacer{} ตีณิ.}

\prose{เนวทสฺสเนน นภาวนาย ปหาตพฺโพ นเหตุ ธมฺโม เนวทสฺสเนน นภาวนาย ปหาตพฺพสฺส นเหตุสฺส ธมฺมสฺส อารมฺมณปจฺจเยน ปจฺจโย.\csromanspacer{} ตีณิ. (สํขิตฺตํ.)}

\proseitem{123}{อารมฺมเณ อฏฺฐ, อธิปติยา ทส ฯเปฯ อวิคเต เตรส. (สํขิตฺตํ.)}

\csromanmatikaanchor{mtk.108}
\csromanmatikaanchor{mtk.109}
\csromantocmarknopage{subparagraph}{9. ทสฺสเนนปหาตพฺพเหตุกตฺติก}{mtk.108}
\bookmark[dest={mtk.108},level=5]{9. ทสฺสเนนปหาตพฺพเหตุกตฺติก}
\csromantocmarkref{subparagraph}{1. เหตุทุก}{mtk.109}
\bookmark[dest={mtk.109},level=5]{1. เหตุทุก}
\csromanheader{6}{9. ทสฺสเนนปหาตพฺพเหตุกตฺติก 1. เหตุทุก}
\prose{\csromanfolio{509}นเหตุยา จุทฺทส, น-อารมฺมเณ จุทฺทส. (สํขิตฺตํ.)}

\csromanheader{2}{อารมฺมณปจฺจยา นเหตุยา อฏฺฐ. (สํขิตฺตํ.)}
\csromanheader{2}{นเหตุปจฺจยา อารมฺมเณ อฏฺฐ. (สํขิตฺตํ.)}
\prose{(ยถา กุสลตฺติเก ปญฺหาวารํ, เอวํ วิตฺถาเรตพฺพํ.)}

\csromanheader{2}{9. ทสฺสเนนปหาตพฺพเหตุกตฺติก 1. เหตุทุก}
\csromanheader{2}{1-7. ปฏิจฺจาทิวาร ปจฺจยจตุกฺก-เหตุ}
\proseitem{124}{ทสฺสเนน ปหาตพฺพเหตุกํ เหตุํ ธมฺมํ ปฏิจฺจ ทสฺสเนน ปหาตพฺพเหตุโก เหตุ ธมฺโม อุปฺปชฺชติ เหตุปจฺจยา. (1)}

\prose{ภาวนาย ปหาตพฺพเหตุกํ เหตุํ ธมฺมํ ปฏิจฺจ ภาวนาย ปหาตพฺพเหตุโก เหตุ ธมฺโม อุปฺปชฺชติ เหตุปจฺจยา. (1)}

\prose{เนวทสฺสเนน นภาวนาย ปหาตพฺพเหตุกํ เหตุํ ธมฺมํ ปฏิจฺจ เนวทสฺสเนน นภาวนาย ปหาตพฺพเหตุโก เหตุ ธมฺโม อุปฺปชฺชติ เหตุปจฺจยา. (1) (สํขิตฺตํ.)}

\proseitem{125}{เหตุยา ตีณิ, อารมฺมเณ ตีณิ ฯเปฯ อวิคเต ตีณิ. (สํขิตฺตํ.)}

\prose{น-อธิปติยา ตีณิ, นปุเรชาเต ตีณิ ฯเปฯ นวิปฺปยุตฺเต ตีณิ. (สํขิตฺตํ.\csromanspacer{} สหชาตวารมฺปิ ฯเปฯ สมฺปยุตฺตวารมฺปิ ปฏิจฺจวารสทิสํ วิตฺถาเรตพฺพํ.)}

\proseitem{126}{ทสฺสเนน ปหาตพฺพเหตุโก เหตุ ธมฺโม ทสฺสเนน ปหาตพฺพเหตุกสฺส เหตุสฺส ธมฺมสฺส เหตุปจฺจเยน ปจฺจโย. (สํขิตฺตํ.)}

\prose{เหตุยา ตีณิ, อารมฺมเณ อฏฺฐ, อธิปติยา ฉ ฯเปฯ อวิคเต ตีณิ. (สํขิตฺตํ.)}

\gambhira{ธมฺมานุโลม ติกทุกปฏฺฐานปาฬิ}
\prose{\csromanfolio{510}นเหตุยา อฏฺฐ, น-อารมฺมเณ อฏฺฐ. (สํขิตฺตํ.)}

\csromanheader{2}{เหตุปจฺจยา น-อารมฺมเณ ตีณิ. (สํขิตฺตํ.)}
\csromanheader{2}{นเหตุปจฺจยา อารมฺมเณ อฏฺฐ. (สํขิตฺตํ.)}
\prose{(ยถา กุสลตฺติเก ปญฺหาวารํ, เอวํ วิตฺถาเรตพฺพํ.)}

\csromanheader{2}{\textbf{นเหตุปท 1-7. ปฏิจฺจาทิวาร}}
\csromanheader{2}{\textbf{ปจฺจยจตุกฺก-เหตุ}}
\proseitem{127}{ทสฺสเนน ปหาตพฺพเหตุกํ นเหตุํ ธมฺมํ ปฏิจฺจ ทสฺสเนน ปหาตพฺพเหตุโก นเหตุ ธมฺโม อุปฺปชฺชติ เหตุปจฺจยา.\csromanspacer{} ตีณิ.}

\prose{ภาวนาย ปหาตพฺพเหตุกํ นเหตุํ ธมฺมํ ปฏิจฺจ ภาวนาย ปหาตพฺพเหตุโก นเหตุ ธมฺโม อุปฺปชฺชติ เหตุปจฺจยา.\csromanspacer{} ตีณิ.}

\prose{เนวทสฺสเนน นภาวนาย ปหาตพฺพเหตุกํ นเหตุํ ธมฺมํ ปฏิจฺจ เนวทสฺสเนน นภาวนาย ปหาตพฺพเหตุโก นเหตุ ธมฺโม อุปฺปชฺชติ เหตุปจฺจยา.\csromanspacer{} ตีณิ. (สํขิตฺตํ.)}

\prose{เหตุยา นว, อารมฺมเณ ตีณิ ฯเปฯ อวิคเต นว. (สํขิตฺตํ.)}

\csromanheader{2}{\textbf{ปจฺจนีย-นเหตุ}}
\proseitem{128}{เนวทสฺสเนน นภาวนาย ปหาตพฺพเหตุกํ นเหตุํ ธมฺมํ ปฏิจฺจ เนวทสฺสเนน นภาวนาย ปหาตพฺพเหตุโก นเหตุ ธมฺโม อุปฺปชฺชติ นเหตุปจฺจยา. (1) (สํขิตฺตํ.)}

\proseitem{129}{นเหตุยา เอกํ, น-อารมฺมเณ ปญฺจ ฯเปฯ โนวิคเต ปญฺจ. (สํขิตฺตํ.)}

\csromanheader{2}{เหตุปจฺจยา น-อารมฺมเณ ปญฺจ. (สํขิตฺตํ.)}
\csromanheader{2}{นเหตุปจฺจยา อารมฺมเณ เอกํ. (สํขิตฺตํ.)}
\prose{(สหชาตวารมฺปิ ฯเปฯ สมฺปยุตฺตวารมฺปิ ปฏิจฺจวารสทิสํ วิตฺถาเรตพฺพํ.)}

\csromanheader{2}{10. อาจยคามิตฺติก 1. เหตุทุก \textbf{อารมฺมณ}}
\proseitem{130}{\csromanfolio{511}ทสฺสเนน ปหาตพฺพเหตุโก นเหตุ ธมฺโม ทสฺสเนน ปหาตพฺพเหตุกสฺส นเหตุสฺส ธมฺมสฺส อารมฺมณปจฺจเยน ปจฺจโย.\csromanspacer{} เทฺว.}

\prose{ภาวนาย ปหาตพฺพเหตุโก นเหตุ ธมฺโม ภาวนาย ปหาตพฺพเหตุกสฺส นเหตุสฺส ธมฺมสฺส อารมฺมณปจฺจเยน ปจฺจโย.\csromanspacer{} ตีณิ.}

\prose{เนวทสฺสเนน นภาวนาย ปหาตพฺพเหตุโก นเหตุ ธมฺโม เนวทสฺสเนน นภาวนาย ปหาตพฺพเหตุกสฺส นเหตุสฺส ธมฺมสฺส อารมฺมณปจฺจเยน ปจฺจโย.\csromanspacer{} ตีณิ. (สํขิตฺตํ.)}

\proseitem{131}{อารมฺมเณ อฏฺฐ, อธิปติยา ทส ฯเปฯ อวิคเต เตรส. (สํขิตฺตํ.)}

\csromanheader{2}{นเหตุยา จุทฺทส น-อารมฺมเณ จุทฺทส. (สํขิตฺตํ.)}
\csromanheader{2}{\textbf{อารมฺมณ}ปจฺจยา นเหตุยา อฏฺฐ. (สํขิตฺตํ.)}
\csromanheader{2}{นเหตุปจฺจยา อารมฺมเณ อฏฺฐ. (สํขิตฺตํ.)}
\prose{(ยถา กุสลตฺติเก ปญฺหาวารํ, เอวํ วิตฺถาเรตพฺพํ.)}

\csromanmatikaanchor{mtk.110}
\csromanmatikaanchor{mtk.111}
\csromantocmarknopage{subparagraph}{10. อาจยคามิตฺติก}{mtk.110}
\bookmark[dest={mtk.110},level=5]{10. อาจยคามิตฺติก}
\csromantocmarkref{subparagraph}{1. เหตุทุก}{mtk.111}
\bookmark[dest={mtk.111},level=5]{1. เหตุทุก}
\csromanheader{6}{10. อาจยคามิตฺติก 1. เหตุทุก}
\csromanheader{2}{1-7. ปฏิจฺจาทิวาร ปจฺจยจตุกฺก-เหตุ}
\proseitem{132}{อาจยคามิํ เหตุํ ธมฺมํ ปฏิจฺจ อาจยคามี เหตุ ธมฺโม อุปฺปชฺชติ เหตุปจฺจยา. (1)}

\prose{อปจยคามิํ เหตุํ ธมฺมํ ปฏิจฺจ อปจยคามี เหตุ ธมฺโม อุปฺปชฺชติ เหตุปจฺจยา. (1)}

\prose{เนวาจยคามินาปจยคามิํ เหตุํ ธมฺมํ ปฏิจฺจ เนวาจยคามินาปจยคามี เหตุ ธมฺโม อุปฺปชฺชติ เหตุปจฺจยา. (1) (สํขิตฺตํ.)}

\proseitem{133}{เหตุยา ตีณิ, อารมฺมเณ ตีณิ, อธิปติยา ตีณิ, (สพฺพตฺถ ตีณิ.)}

\gambhira{ธมฺมานุโลม ติกทุกปฏฺฐานปาฬิ}
\prose{\csromanfolio{512}น-อธิปติยา ตีณิ, น-อาเสวเน เทฺว, นวิปาเก ตีณิ, นวิปฺปยุตฺเต ตีณิ. (สํขิตฺตํ.)}

\csromanheader{2}{เหตุปจฺจยา น-อธิปติยา ตีณิ. (สํขิตฺตํ.)}
\csromanheader{2}{น-อธิปติปจฺจยา เหตุยา ตีณิ. (สํขิตฺตํ.)}
\prose{(สหชาตวารมฺปิ ฯเปฯ สมฺปยุตฺตวารมฺปิ ปฏิจฺจวารสทิสํ วิตฺถาเรตพฺพํ.)}

\csromanheader{2}{\textbf{เหตุ อารมฺมณาทิ}}
\proseitem{134}{อาจยคามี เหตุ ธมฺโม อาจยคามิสฺส เหตุสฺส ธมฺมสฺส เหตุปจฺจเยน ปจฺจโย. (1)}

\prose{อปจยคามี เหตุ ธมฺโม อปจยคามิสฺส เหตุสฺส ธมฺมสฺส เหตุปจฺจเยน ปจฺจโย. (1)}

\prose{เนวาจยคามินาปจยคามี เหตุ ธมฺโม เนวาจยคามินาปจยคามิสฺส เหตุสฺส ธมฺมสฺส เหตุปจฺจเยน ปจฺจโย. (1)}

\prose{อาจยคามี เหตุ ธมฺโม อาจยคามิสฺส เหตุสฺส ธมฺมสฺส อารมฺมณปจฺจเยน ปจฺจโย.\csromanspacer{}\csromanspacer{}อาจยคามี เหตุ ธมฺโม เนวาจยคามินาปจยคามิสฺส เหตุสฺส ธมฺมสฺส อารมฺมณปจฺจเยน ปจฺจโย. (2)}

\prose{อปจยคามี เหตุ ธมฺโม อาจยคามิสฺส เหตุสฺส ธมฺมสฺส อารมฺมณปจฺจเยน ปจฺจโย.\csromanspacer{}\csromanspacer{}อปจยคามี เหตุ ธมฺโม เนวาจยคามินาปจยคามิสฺส เหตุสฺส ธมฺมสฺส อารมฺมณปจฺจเยน ปจฺจโย. (2)}

\prose{เนวาจยคามินาปจยคามี เหตุ ธมฺโม เนวาจยคามินาปจยคามิสฺส เหตุสฺส ธมฺมสฺส อารมฺมณปจฺจเยน ปจฺจโย.\csromanspacer{}\csromanspacer{}เนวาจยคามินาปจยคามี เหตุ ธมฺโม อาจยคามิสฺส เหตุสฺส ธมฺมสฺส อารมฺมณปจฺจเยน ปจฺจโย. (2)}

\proseitem{135}{อาจยคามี เหตุ ธมฺโม อาจยคามิสฺส เหตุสฺส ธมฺมสฺส อธิปติปจฺจเยน ปจฺจโย.\csromanspacer{} อารมฺมณาธิปติ, สหชาตาธิปติ. (1)}

\prose{อปจยคามี เหตุ ธมฺโม อปจยคามิสฺส เหตุสฺส ธมฺมสฺส อธิปติปจฺจเยน ปจฺจโย.\csromanspacer{} สหชาตาธิปติ.\csromanspacer{}\csromanspacer{}อปจยคามี เหตุ ธมฺโม อาจยคามิสฺส เหตุสฺส ธมฺมสฺส อธิปติปจฺจเยน ปจฺจโย.\csromanspacer{}\csromanspacer{}อปจยคามี เหตุ ธมฺโม เนวาจยคามินาปจยคามิสฺส เหตุสฺส ธมฺมสฺส อธิปติปจฺจเยน ปจฺจโย.\csromanspacer{} อารมฺมณาธิปติ. (3)}

\csromanheader{2}{10. อาจยคามิตฺติก 1. เหตุทุก}
\prose{\csromanfolio{513}เนวาจยคามินาปจยคามี เหตุ ธมฺโม เนวาจยคามินาปจยคามิสฺส เหตุสฺส ธมฺมสฺส อธิปติปจฺจเยน ปจฺจโย.\csromanspacer{}\csromanspacer{}เนวาจยคามินาปจยคามี เหตุ ธมฺโม อาจยคามิสฺส เหตุสฺส ธมฺมสฺส อธิปติปจฺจเยน ปจฺจโย.}

\proseitem{136}{อาจยคามี เหตุ ธมฺโม อาจยคามิสฺส เหตุสฺส ธมฺมสฺส อนนฺตรปจฺจเยน ปจฺจโย.\csromanspacer{}\csromanspacer{}อาจยคามี เหตุ ธมฺโม อปจยคามิสฺส เหตุสฺส ธมฺมสฺส อนนฺตรปจฺจเยน ปจฺจโย.\csromanspacer{}\csromanspacer{}อาจยคามี เหตุ ธมฺโม เนวาจยคามินาปจยคามิสฺส เหตุสฺส ธมฺมสฺส อนนฺตรปจฺจเยน ปจฺจโย. (3)}

\prose{อปจยคามี เหตุ ธมฺโม เนวาจยคามินาปจยคามิสฺส เหตุสฺส ธมฺมสฺส อนนฺตรปจฺจเยน ปจฺจโย. (1)}

\prose{เนวาจยคามินาปจยคามี เหตุ ธมฺโม เนวาจยคามินาปจยคามิสฺส เหตุสฺส ธมฺมสฺส อนนฺตรปจฺจเยน ปจฺจโย. (1) (สํขิตฺตํ.)}

\proseitem{137}{เหตุยา ตีณิ, อารมฺมเณ ฉ, อธิปติยา ฉ, อนนฺตเร ปญฺจ, สมนนฺตเร ปญฺจ, สหชาเต ตีณิ, อญฺญมญฺเญ ตีณิ, นิสฺสเย ตีณิ, อุปนิสฺสเย อฏฺฐ, อาเสวเน ตีณิ, วิปาเก เอกํ ฯเปฯ อวิคเต ตีณิ. (สํขิตฺตํ.)}

\prose{นเหตุยา อฏฺฐ, น-อารมฺมเณ อฏฺฐ. (สํขิตฺตํ.)}

\csromanheader{2}{เหตุปจฺจยา น-อารมฺมเณ ตีณิ. (สํขิตฺตํ.)}
\csromanheader{2}{นเหตุปจฺจยา อารมฺมเณ ฉ. (สํขิตฺตํ.)}
\prose{(ยถา กุสลตฺติเก ปญฺหาวารํ, เอวํ วิตฺถาเรตพฺพํ.)}

\csromanheader{2}{\textbf{นเหตุปท 1. ปฏิจฺจวาร}}
\csromanheader{2}{\textbf{ปจฺจยจตุกฺก-เหตุ อารมฺมณ}}
\proseitem{138}{อาจยคามิํ นเหตุํ ธมฺมํ ปฏิจฺจ อาจยคามี นเหตุ ธมฺโม อุปฺปชฺชติ เหตุปจฺจยา.\csromanspacer{}\csromanspacer{}อาจยคามิํ นเหตุํ ธมฺมํ ปฏิจฺจ เนวาจยคามินาปจยคามี นเหตุ ธมฺโม อุปฺปชฺชติ เหตุปจฺจยา.\csromanspacer{}\csromanspacer{}อาจยคามิํ}

\vspace{\tipitakaparskip}
\csromancenter{ธมฺมานุโลม ติกทุกปฏฺฐานปาฬิ นเหตุํ ธมฺมํ ปฏิจฺจ อาจยคามี นเหตุ จ เนวาจยคามินาปจยคามี นเหตุ จ ธมฺมา อุปฺปชฺชนฺติ เหตุปจฺจยา. (3)}
\prose{\csromanfolio{514}อปจยคามิํ นเหตุํ ธมฺมํ ปฏิจฺจ อปจยคามี นเหตุ ธมฺโม อุปฺปชฺชติ เหตุปจฺจยา.\csromanspacer{}\csromanspacer{}อปจยคามิํ นเหตุํ ธมฺมํ ปฏิจฺจ เนวาจยคามินาปจยคามี นเหตุ ธมฺโม อุปฺปชฺชติ เหตุปจฺจยา.\csromanspacer{}\csromanspacer{}อปจยคามิํ นเหตุํ ธมฺมํ ปฏิจฺจ อปจยคามี นเหตุ จ เนวาจยคามินาปจยคามี นเหตุ จ ธมฺมา อุปฺปชฺชนฺติ เหตุปจฺจยา. (3)}

\prose{เนวาจยคามินาปจยคามิํ นเหตุํ ธมฺมํ ปฏิจฺจ เนวาจยคามินาปจยคามี นเหตุ ธมฺโม อุปฺปชฺชติ เหตุปจฺจยา. (1)}

\prose{อาจยคามิํ นเหตุญฺจ เนวาจยคามินาปจยคามิํ นเหตุญฺจ ธมฺมํ ปฏิจฺจ เนวาจยคามินาปจยคามี นเหตุ ธมฺโม อุปฺปชฺชติ เหตุปจฺจยา. (1)}

\prose{อปจยคามิํ นเหตุญฺจ เนวาจยคามินาปจยคามิํ นเหตุญฺจ ธมฺมํ ปฏิจฺจ เนวาจยคามินาปจยคามี นเหตุ ธมฺโม อุปฺปชฺชติ เหตุปจฺจยา. (1)}

\proseitem{139}{อาจยคามิํ นเหตุํ ธมฺมํ ปฏิจฺจ อาจยคามี นเหตุ ธมฺโม อุปฺปชฺชติ อารมฺมณปจฺจยา. (1)}

\prose{อปจยคามิํ นเหตุํ ธมฺมํ ปฏิจฺจ อปจยคามี นเหตุ ธมฺโม อุปฺปชฺชติ อารมฺมณปจฺจยา. (1)}

\prose{เนวาจยคามินาปจยคามิํ นเหตุํ ธมฺมํ ปฏิจฺจ เนวาจยคามินาปจยคามี นเหตุ ธมฺโม อุปฺปชฺชติ อารมฺมณปจฺจยา. (1) (สํขิตฺตํ.)}

\prose{เหตุยา นว, อารมฺมเณ ตีณิ, อธิปติยา นว ฯเปฯ อวิคเต นว. (สํขิตฺตํ.)}

\csromanheader{2}{\textbf{นเหตุ น-อารมฺมณ}}
\proseitem{140}{เนวาจยคามินาปจยคามิํ นเหตุํ ธมฺมํ ปฏิจฺจ เนวาจยคามินาปจยคามี นเหตุ ธมฺโม อุปฺปชฺชติ นเหตุปจฺจยา. (1)}

\prose{อาจยคามิํ นเหตุํ ธมฺมํ ปฏิจฺจ เนวาจยคามินาปจยคามี นเหตุ ธมฺโม อุปฺปชฺชติ น-อารมฺมณปจฺจยา. (สํขิตฺตํ.)}

\proseitem{141}{นเหตุยา เอกํ, น-อารมฺมเณ ปญฺจ, น-อธิปติยา ฉ ฯเปฯ โนวิคเต ปญฺจ. (สํขิตฺตํ.)}

\csromanheader{2}{10. อาจยคามิตฺติก 1. เหตุทุก}
\csromanheader{2}{เหตุปจฺจยา น-อารมฺมเณ ปญฺจ. (สํขิตฺตํ.)}
\csromanheader{2}{นเหตุปจฺจยา อารมฺมเณ เอกํ. (สํขิตฺตํ.)}
\prose{\csromanfolio{515}(สหชาตวารมฺปิ ฯเปฯ สมฺปยุตฺตวารมฺปิ ปฏิจฺจวารสทิสํ วิตฺถาเรตพฺพํ.)}

\csromanheader{2}{7. ปญฺหาวาร ปจฺจยจตุกฺก}
\proseitem{142}{อาจยคามี นเหตุ ธมฺโม อาจยคามิสฺส นเหตุสฺส ธมฺมสฺส อารมฺมณปจฺจเยน ปจฺจโย.\csromanspacer{}\csromanspacer{}อาจยคามี นเหตุ ธมฺโม เนวาจยคามินาปจยคามิสฺส นเหตุสฺส ธมฺมสฺส อารมฺมณปจฺจเยน ปจฺจโย. (2)}

\prose{อปจยคามี นเหตุ ธมฺโม อาจยคามิสฺส นเหตุสฺส ธมฺมสฺส อารมฺมณปจฺจเยน ปจฺจโย.\csromanspacer{}\csromanspacer{}อปจยคามี นเหตุ ธมฺโม เนวาจยคามินาปจยคามิสฺส นเหตุสฺส ธมฺมสฺส อารมฺมณปจฺจเยน ปจฺจโย. (2)}

\prose{เนวาจยคามินาปจยคามี นเหตุ ธมฺโม เนวาจยคามินาปจยคามิสฺส นเหตุสฺส ธมฺมสฺส อารมฺมณปจฺจเยน ปจฺจโย.\csromanspacer{}\csromanspacer{}เนวาจยคามินาปจยคามี นเหตุ ธมฺโม อาจยคามิสฺส นเหตุสฺส ธมฺมสฺส อารมฺมณปจฺจเยน ปจฺจโย.\csromanspacer{}\csromanspacer{}เนวาจยคามินาปจยคามี นเหตุ ธมฺโม อปจยคามิสฺส นเหตุสฺส ธมฺมสฺส อารมฺมณปจฺจเยน ปจฺจโย. (3)}

\proseitem{143}{อาจยคามี นเหตุ ธมฺโม อาจยคามิสฺส นเหตุสฺส ธมฺมสฺส อธิปติปจฺจเยน ปจฺจโย.\csromanspacer{} ตีณิ.}

\prose{อปจยคามี นเหตุ ธมฺโม อปจยคามิสฺส นเหตุสฺส ธมฺมสฺส อธิปติปจฺจเยน ปจฺจโย.\csromanspacer{}\csromanspacer{}อปจยคามี นเหตุ ธมฺโม อาจยคามิสฺส นเหตุสฺส ธมฺมสฺส อธิปติปจฺจเยน ปจฺจโย.\csromanspacer{}\csromanspacer{}อปจยคามี นเหตุ ธมฺโม เนวาจยคามินาปจยคามิสฺส นเหตุสฺส ธมฺมสฺส อธิปติปจฺจเยน ปจฺจโย.\csromanspacer{}\csromanspacer{}อปจยคามี นเหตุ ธมฺโม อปจยคามิสฺส นเหตุสฺส จ เนวาจยคามินาปจยคามิสฺส นเหตุสฺส จ ธมฺมสฺส อธิปติปจฺจเยน ปจฺจโย. (4)}

\proseitem{144}{เนวาจยคามินาปจยคามี นเหตุ ธมฺโม เนวาจยคามินาปจยคามิสฺส นเหตุสฺส ธมฺมสฺส อธิปติปจฺจเยน ปจฺจโย.\csromanspacer{} ตีณิ.}

\gambhira{ธมฺมานุโลม ติกทุกปฏฺฐานปาฬิ}
\prose{\csromanfolio{516}อาจยคามี นเหตุ ธมฺโม อาจยคามิสฺส นเหตุสฺส ธมฺมสฺส อนนฺตรปจฺจเยน ปจฺจโย.\csromanspacer{}\csromanspacer{}อาจยคามี นเหตุ ธมฺโม อปจยคามิสฺส นเหตุสฺส ธมฺมสฺส อนนฺตรปจฺจเยน ปจฺจโย.\csromanspacer{}\csromanspacer{}อาจยคามี นเหตุ ธมฺโม เนวาจยคามินาปจยคามิสฺส นเหตุสฺส ธมฺมสฺส อนนฺตรปจฺจเยน ปจฺจโย. (3)}

\prose{อปจยคามี นเหตุ ธมฺโม เนวาจยคามินาปจยคามิสฺส นเหตุสฺส ธมฺมสฺส อนนฺตรปจฺจเยน ปจฺจโย. (1)}

\prose{เนวาจยคามินาปจยคามี นเหตุ ธมฺโม เนวาจยคามินาปจยคามิสฺส นเหตุสฺส ธมฺมสฺส อนนฺตรปจฺจเยน ปจฺจโย.\csromanspacer{}\csromanspacer{}เนวาจยคามินาปจยคามี นเหตุ ธมฺโม อาจยคามิสฺส นเหตุสฺส ธมฺมสฺส อนนฺตรปจฺจเยน ปจฺจโย. (2) ฯเปฯ.}

\prose{อาจยคามี นเหตุ ธมฺโม อาจยคามิสฺส นเหตุสฺส ธมฺมสฺส อุปนิสฺสยปจฺจเยน ปจฺจโย. (สํขิตฺตํ.)}

\proseitem{145}{อารมฺมเณ สตฺต, อธิปติยา ทส, อนนฺตเร ฉ, สมนนฺตเร ฉ, สหชาเต นว, อญฺญมญฺเญ ตีณิ, นิสฺสเย เตรส, อุปนิสฺสเย นว ฯเปฯ อวิคเต เตรส. (สํขิตฺตํ.)}

\prose{นเหตุยา ปนฺนรส, น-อารมฺมเณ ปนฺนรส. (สํขิตฺตํ.)}

\csromanheader{2}{อารมฺมณปจฺจยา นเหตุยา สตฺต. (สํขิตฺตํ.)}
\csromanheader{2}{นเหตุปจฺจยา อารมฺมเณ สตฺต. (สํขิตฺตํ.)}
\prose{(ยถา กุสลตฺติเก ปญฺหาวารํ, เอวํ วิตฺถาเรตพฺพํ.)}

\csromanmatikaanchor{mtk.112}
\csromanmatikaanchor{mtk.113}
\csromantocmarknopage{subparagraph}{11. เสกฺขตฺติก}{mtk.112}
\bookmark[dest={mtk.112},level=5]{11. เสกฺขตฺติก}
\csromantocmarkref{subparagraph}{1. เหตุทุก}{mtk.113}
\bookmark[dest={mtk.113},level=5]{1. เหตุทุก}
\csromanheader{6}{11. เสกฺขตฺติก 1. เหตุทุก}
\csromanheader{2}{1-7. ปฏิจฺจาทิวาร ปจฺจยจตุกฺก-เหตุ}
\vspace{\tipitakaparskip}
\csromancenter{เสกฺขํ เหตุํ ธมฺมํ ปฏิจฺจ เสกฺโข เหตุ ธมฺโม อุปฺปชฺชติ เหตุปจฺจยา. (1)}
\csromanheader{2}{11. เสกฺขตฺติก 1. เหตุทุก}
\prose{\csromanfolio{517}อเสกฺขํ เหตุํ ธมฺมํ ปฏิจฺจ อเสกฺโข เหตุ ธมฺโม อุปฺปชฺชติ เหตุปจฺจยา. (1)}

\prose{เนวเสกฺขนาเสกฺขํ เหตุํ ธมฺมํ ปฏิจฺจ เนวเสกฺขนาเสกฺโข เหตุ ธมฺโม อุปฺปชฺชติ เหตุปจฺจยา. (1) (สํขิตฺตํ.)}

\proseitem{147}{เหตุยา ตีณิ, อารมฺมเณ ตีณิ, อธิปติยา ตีณิ ฯเปฯ อวิคเต ตีณิ. (สํขิตฺตํ.)}

\prose{น-อธิปติยา ตีณิ, นปุเรชาเต ตีณิ, นปจฺฉาชาเต ตีณิ, น-อาเสวเน ตีณิ, นวิปาเก เทฺว, นวิปฺปยุตฺเต ตีณิ. (สํขิตฺตํ.)}

\csromanheader{2}{เหตุปจฺจยา น-อธิปติยา ตีณิ. (สํขิตฺตํ.)}
\prose{น-อธิปติปจฺจยา เหตุยา ตีณิ. (สํขิตฺตํ.) (สหชาตวารมฺปิ ฯเปฯ สมฺปยุตฺตวารมฺปิ วิตฺถาเรตพฺพํ.)}

\csromanheader{2}{\textbf{เหตุ-อารมฺมณาทิ}}
\proseitem{148}{เสกฺโข เหตุ ธมฺโม เสกฺขสฺส เหตุสฺส ธมฺมสฺส เหตุปจฺจเยน ปจฺจโย. (1)}

\prose{อเสกฺโข เหตุ ธมฺโม อเสกฺขสฺส เหตุสฺส ธมฺมสฺส เหตุปจฺจเยน ปจฺจโย. (1)}

\prose{เนวเสกฺขนาเสกฺโข เหตุ ธมฺโม เนวเสกฺขนาเสกฺขสฺส เหตุสฺส ธมฺมสฺส เหตุปจฺจเยน ปจฺจโย. (1)}

\proseitem{149}{เสกฺโข เหตุ ธมฺโม เนวเสกฺขนาเสกฺขสฺส เหตุสฺส ธมฺมสฺส อารมฺมณปจฺจเยน ปจฺจโย. (1)}

\prose{อเสกฺโข เหตุ ธมฺโม เนวเสกฺขนาเสกฺขสฺส เหตุสฺส ธมฺมสฺส อารมฺมณปจฺจเยน ปจฺจโย. (1)}

\prose{เนวเสกฺขนาเสกฺโข เหตุ ธมฺโม เนวเสกฺขนาเสกฺขสฺส เหตุสฺส ธมฺมสฺส อารมฺมณปจฺจเยน ปจฺจโย. (1)}

\gambhira{ธมฺมานุโลม ติกทุกปฏฺฐานปาฬิ}
\proseitem{150}{\csromanfolio{518}เสกฺโข เหตุ ธมฺโม เสกฺขสฺส เหตุสฺส ธมฺมสฺส อธิปติปจฺจเยน ปจฺจโย.\csromanspacer{}\csromanspacer{}เสกฺโข เหตุ ธมฺโม เนวเสกฺขนาเสกฺขสฺส เหตุสฺส ธมฺมสฺส อธิปติปจฺจเยน ปจฺจโย. (2)}

\prose{อเสกฺโข เหตุ ธมฺโม อเสกฺขสฺส เหตุสฺส ธมฺมสฺส อธิปติปจฺจเยน ปจฺจโย.\csromanspacer{}\csromanspacer{}อเสกฺโข เหตุ ธมฺโม เนวเสกฺขนาเสกฺขสฺส เหตุสฺส ธมฺมสฺส อธิปติปจฺจเยน ปจฺจโย. (2)}

\prose{เนวเสกฺขนาเสกฺโข เหตุ ธมฺโม เนวเสกฺขนาเสกฺขสฺส เหตุสฺส ธมฺมสฺส อธิปติปจฺจเยน ปจฺจโย. (1)}

\proseitem{151}{เสกฺโข เหตุ ธมฺโม เสกฺขสฺส เหตุสฺส ธมฺมสฺส อนนฺตรปจฺจเยน ปจฺจโย.\csromanspacer{}\csromanspacer{}เสกฺโข เหตุ ธมฺโม อเสกฺขสฺส เหตุสฺส ธมฺมสฺส อนนฺตรปจฺจเยน ปจฺจโย.\csromanspacer{}\csromanspacer{}เสกฺโข เหตุ ธมฺโม เนวเสกฺขนาเสกฺขสฺส เหตุสฺส ธมฺมสฺส อนนฺตรปจฺจเยน ปจฺจโย. (3)}

\prose{อเสกฺโข เหตุ ธมฺโม อเสกฺขสฺส เหตุสฺส ธมฺมสฺส อนนฺตรปจฺจเยน ปจฺจโย.\csromanspacer{}\csromanspacer{}อเสกฺโข เหตุ ธมฺโม เนวเสกฺขนาเสกฺขสฺส เหตุสฺส ธมฺมสฺส อนนฺตรปจฺจเยน ปจฺจโย. (2)}

\prose{เนวเสกฺขนาเสกฺโข เหตุ ธมฺโม เนวเสกฺขนาเสกฺขสฺส เหตุสฺส ธมฺมสฺส อนนฺตรปจฺจเยน ปจฺจโย.\csromanspacer{} เนวเสกฺขนาเสกฺโข เหตุ ธมฺโม เสกฺขสฺส เหตุสฺส ธมฺมสฺส อนนฺตรปจฺจเยน ปจฺจโย.\csromanspacer{}\csromanspacer{}เนวเสกฺขนาเสกฺโข เหตุ ธมฺโม อเสกฺขสฺส เหตุสฺส ธมฺมสฺส อนนฺตรปจฺจเยน ปจฺจโย. (3) (สํขิตฺตํ.)}

\proseitem{152}{เหตุยา ตีณิ, อารมฺมเณ ตีณิ, อธิปติยา ปญฺจ, อนนฺตเร อฏฺฐ, สมนนฺตเร อฏฺฐ, สหชาเต ตีณิ, อญฺญมญฺเญ ตีณิ, นิสฺสเย ตีณิ, อุปนิสฺสเย อฏฺฐ ฯเปฯ อวิคเต ตีณิ. (สํขิตฺตํ.)}

\prose{นเหตุยา อฏฺฐ, น-อารมฺมเณ อฏฺฐ. (สํขิตฺตํ.)}

\csromanheader{2}{เหตุปจฺจยา น-อารมฺมเณ ตีณิ. (สํขิตฺตํ.)}
\csromanheader{2}{นเหตุปจฺจยา อารมฺมเณ ตีณิ. (สํขิตฺตํ.)}
\prose{(ยถา กุสลตฺติเก ปญฺหาวารํ, เอวํ วิตฺถาเรตพฺพํ.)}

\csromanheader{2}{11. เสกฺขตฺติก 1. เหตุทุก \textbf{นเหตุปท 1. ปฏิจฺจวาร}}
\csromanheader{2}{\textbf{ปจฺจยจตุกฺก เหตุ-อารมฺมณ}}
\proseitem{153}{\csromanfolio{519}เสกฺขํ นเหตุํ ธมฺมํ ปฏิจฺจ เสกฺโข นเหตุ ธมฺโม อุปฺปชฺชติ เหตุปจฺจยา.\csromanspacer{}\csromanspacer{}เสกฺขํ นเหตุํ ธมฺมํ ปฏิจฺจ เนวเสกฺขนาเสกฺโข นเหตุ ธมฺโม อุปฺปชฺชติ เหตุปจฺจยา.\csromanspacer{}\csromanspacer{}เสกฺขํ นเหตุํ ธมฺมํ ปฏิจฺจ เสกฺโข นเหตุ จ เนวเสกฺขนาเสกฺโข นเหตุ จ ธมฺมา อุปฺปชฺชนฺติ เหตุปจฺจยา. (3)}

\prose{อเสกฺขํ นเหตุํ ธมฺมํ ปฏิจฺจ อเสกฺโข นเหตุ ธมฺโม อุปฺปชฺชติ เหตุปจฺจยา.\csromanspacer{} ตีณิ.}

\prose{เนวเสกฺขนาเสกฺขํ นเหตุํ ธมฺมํ ปฏิจฺจ เนวเสกฺขนาเสกฺโข นเหตุ ธมฺโม อุปฺปชฺชติ เหตุปจฺจยา. (1)}

\prose{เสกฺขํ นเหตุญฺจ เนวเสกฺขนาเสกฺขํ นเหตุญฺจ ธมฺมํ ปฏิจฺจ เนวเสกฺขนาเสกฺโข นเหตุ ธมฺโม อุปฺปชฺชติ เหตุปจฺจยา. (1)}

\prose{อเสกฺขํ นเหตุญฺจ เนวเสกฺขนาเสกฺขํ นเหตุญฺจ ธมฺมํ ปฏิจฺจ เนวเสกฺขนาเสกฺโข นเหตุ ธมฺโม อุปฺปชฺชติ เหตุปจฺจยา. (1)}

\proseitem{154}{เสกฺขํ นเหตุํ ธมฺมํ ปฏิจฺจ เสกฺโข นเหตุ ธมฺโม อุปฺปชฺชติ อารมฺมณปจฺจยา. (1)}

\prose{อเสกฺขํ นเหตุํ ธมฺมํ ปฏิจฺจ อเสกฺโข นเหตุ ธมฺโม อุปฺปชฺชติ อารมฺมณปจฺจยา. (1)}

\prose{เนวเสกฺขนาเสกฺขํ นเหตุํ ธมฺมํ ปฏิจฺจ เนวเสกฺขนาเสกฺโข นเหตุ ธมฺโม อุปฺปชฺชติ อารมฺมณปจฺจยา. (1) (สํขิตฺตํ.)}

\prose{เหตุยา นว, อารมฺมเณ ตีณิ, อธิปติยา นว ฯเปฯ อาเสวเน เทฺว ฯเปฯ อวิคเต นว. (สํขิตฺตํ.)}

\csromanheader{2}{\textbf{ปจฺจนีย-นเหตุ}}
\proseitem{155}{เนวเสกฺขนาเสกฺขํ นเหตุํ ธมฺมํ ปฏิจฺจ เนวเสกฺขนาเสกฺโข นเหตุ ธมฺโม อุปฺปชฺชติ เหตุปจฺจยา. (สํขิตฺตํ.)}

\gambhira{ธมฺมานุโลม ติกทุกปฏฺฐานปาฬิ}
\prose{\csromanfolio{520}นเหตุยา เอกํ, น-อารมฺมเณ ปญฺจ, น-อธิปติยา ตีณิ ฯเปฯ โนวิคเต ปญฺจ. (สํขิตฺตํ.)}

\csromanheader{2}{เหตุปจฺจยา น-อารมฺมเณ ปญฺจ. (สํขิตฺตํ.)}
\csromanheader{2}{นเหตุปจฺจยา อารมฺมเณ เอกํ. (สํขิตฺตํ.)}
\prose{(สหชาตวารมฺปิ ฯเปฯ สมฺปยุตฺตวารมฺปิ ปฏิจฺจวารสทิสํ วิตฺถาเรตพฺพํ.)}

\csromanheader{2}{7. ปญฺหาวาร ปจฺจยจตุกฺก}
\proseitem{156}{เสกฺโข นเหตุ ธมฺโม เนวเสกฺขนาเสกฺขสฺส นเหตุสฺส ธมฺมสฺส อารมฺมณปจฺจเยน ปจฺจโย. (1)}

\prose{อเสกฺโข นเหตุ ธมฺโม เนวเสกฺขนาเสกฺขสฺส นเหตุสฺส ธมฺมสฺส อารมฺมณปจฺจเยน ปจฺจโย. (1)}

\prose{เนวเสกฺขนาเสกฺโข นเหตุ ธมฺโม เนวเสกฺขนาเสกฺขสฺส นเหตุสฺส ธมฺมสฺส อารมฺมณปจฺจเยน ปจฺจโย.\csromanspacer{}\csromanspacer{}เนวเสกฺขนาเสกฺโข นเหตุ ธมฺโม เสกฺขสฺส นเหตุสฺส ธมฺมสฺส อารมฺมณปจฺจเยน ปจฺจโย.\csromanspacer{}\csromanspacer{}เนวเสกฺขนาเสกฺโข นเหตุ ธมฺโม อเสกฺขสฺส นเหตุสฺส ธมฺมสฺส อารมฺมณปจฺจเยน ปจฺจโย. (3)}

\prose{เสกฺโข นเหตุ ธมฺโม เสกฺขสฺส นเหตุสฺส ธมฺมสฺส อธิปติปจฺจเยน ปจฺจโย.\csromanspacer{} ตีณิ.}

\prose{อเสกฺโข นเหตุ ธมฺโม อเสกฺขสฺส นเหตุสฺส ธมฺมสฺส อธิปติปจฺจเยน ปจฺจโย.\csromanspacer{} ตีณิ.}

\prose{เนวเสกฺขนาเสกฺโข นเหตุ ธมฺโม เนวเสกฺขนาเสกฺขสฺส นเหตุสฺส ธมฺมสฺส อธิปติปจฺจเยน ปจฺจโย.\csromanspacer{} ตีณิ. (สํขิตฺตํ.)}

\proseitem{157}{อารมฺมเณ ปญฺจ, อธิปติยา นว, อนนฺตเร อฏฺฐ, สมนนฺตเร อฏฺฐ.\csromanspacer{} สหชาเต นว, อญฺญมญฺเญ ตีณิ, นิสฺสเย เตรส, อุปนิสฺสเย อฏฺฐ ฯเปฯ อวิคเต เตรส. (สํขิตฺตํ.)}

\csromanmatikaanchor{mtk.114}
\csromanmatikaanchor{mtk.115}
\csromantocmarknopage{subparagraph}{12. ปริตฺตตฺติก}{mtk.114}
\bookmark[dest={mtk.114},level=5]{12. ปริตฺตตฺติก}
\csromantocmarkref{subparagraph}{1. เหตุทุก}{mtk.115}
\bookmark[dest={mtk.115},level=5]{1. เหตุทุก}
\csromanheader{6}{12. ปริตฺตตฺติก 1. เหตุทุก}
\prose{\csromanfolio{521}นเหตุยา จุทฺทส, น-อารมฺมเณ จุทฺทส. (สํขิตฺตํ.)}

\csromanheader{2}{อารมฺมณปจฺจยา นเหตุยา ปญฺจ. (สํขิตฺตํ.)}
\csromanheader{2}{นเหตุปจฺจยา อารมฺมเณ ปญฺจ. (สํขิตฺตํ.)}
\prose{(ยถา กุสลตฺติเก ปญฺหาวารํ, เอวํ วิตฺถาเรตพฺพํ.)}

\csromanheader{2}{12. ปริตฺตตฺติก 1. เหตุทุก}
\csromanheader{2}{1-7. ปฏิจฺจาทิวาร ปจฺจยจตุกฺก}
\proseitem{158}{ปริตฺตํ เหตุํ ธมฺมํ ปฏิจฺจ ปริตฺโต เหตุ ธมฺโม อุปฺปชฺชติ เหตุปจฺจยา. (1)}

\prose{มหคฺคตํ เหตุํ ธมฺมํ ปฏิจฺจ มหคฺคโต เหตุ ธมฺโม อุปฺปชฺชติ เหตุปจฺจยา. (1)}

\prose{อปฺปมาณํ เหตุํ ธมฺมํ ปฏิจฺจ อปฺปมาโณ เหตุ ธมฺโม อุปฺปชฺชติ เหตุปจฺจยา. (1) (สํขิตฺตํ.)}

\prose{เหตุยา ตีณิ, อารมฺมเณ ตีณิ ฯเปฯ อวิคเต ตีณิ. (สํขิตฺตํ.)}

\csromanheader{2}{น-อธิปติยา ตีณิ ฯเปฯ นวิปฺปยุตฺเต ตีณิ. (สํขิตฺตํ.)}
\prose{(สหชาตวารมฺปิ ฯเปฯ สมฺปยุตฺตวารมฺปิ ปฏิจฺจวารสทิสํ วิตฺถาเรตพฺพํ.)}

\csromanheader{2}{\textbf{เหตุ-อารมฺมณาทิ}}
\proseitem{159}{ปริตฺโต เหตุ ธมฺโม ปริตฺตสฺส เหตุสฺส ธมฺมสฺส เหตุปจฺจเยน ปจฺจโย. (1)}

\prose{มหคฺคโต เหตุ ธมฺโม มหคฺคตสฺส เหตุสฺส ธมฺมสฺส เหตุปจฺจเยน ปจฺจโย. (1)}

\prose{อปฺปมาโณ เหตุ ธมฺโม อปฺปมาณสฺส เหตุสฺส ธมฺมสฺส เหตุปจฺจเยน ปจฺจโย. (1)}

\prose{ปริตฺโต เหตุ ธมฺโม ปริตฺตสฺส เหตุสฺส ธมฺมสฺส อารมฺมณปจฺจเยน ปจฺจโย.\csromanspacer{}\csromanspacer{}ปริตฺโต เหตุ ธมฺโม มหคฺคตสฺส เหตุสฺส ธมฺมสฺส อารมฺมณปจฺจเยน ปจฺจโย. (2)}

\gambhira{ธมฺมานุโลม ติกทุกปฏฺฐานปาฬิ}
\prose{\csromanfolio{522}มหคฺคโต เหตุ ธมฺโม มหคฺคตสฺส เหตุสฺส ธมฺมสฺส อารมฺมณปจฺจเยน ปจฺจโย.\csromanspacer{}\csromanspacer{}มหคฺคโต เหตุ ธมฺโม ปริตฺตสฺส เหตุสฺส ธมฺมสฺส อารมฺมณปจฺจเยน ปจฺจโย. (2)}

\prose{อปฺปมาโณ เหตุ ธมฺโม ปริตฺตสฺส เหตุสฺส ธมฺมสฺส อารมฺมณปจฺจเยน ปจฺจโย.\csromanspacer{}\csromanspacer{}อปฺปมาโณ เหตุ ธมฺโม มหคฺคตสฺส เหตุสฺส ธมฺมสฺส อารมฺมณปจฺจเยน ปจฺจโย. (2)}

\proseitem{160}{ปริตฺโต เหตุ ธมฺโม ปริตฺตสฺส เหตุสฺส ธมฺมสฺส อธิปติปจฺจเยน ปจฺจโย. (1)}

\prose{มหคฺคโต เหตุ ธมฺโม มหคฺคตสฺส เหตุสฺส ธมฺมสฺส อธิปติปจฺจเยน ปจฺจโย.\csromanspacer{}\csromanspacer{}มหคฺคโต เหตุ ธมฺโม ปริตฺตสฺส เหตุสฺส ธมฺมสฺส อธิปติปจฺจเยน ปจฺจโย. (2)}

\prose{อปฺปมาโณ เหตุ ธมฺโม อปฺปมาณสฺส เหตุสฺส ธมฺมสฺส อธิปติปจฺจเยน ปจฺจโย.\csromanspacer{} เทฺว.}

\prose{ปริตฺโต เหตุ ธมฺโม ปริตฺตสฺส เหตุสฺส ธมฺมสฺส อนนฺตรปจฺจเยน ปจฺจโย.\csromanspacer{}\csromanspacer{}ปริตฺโต เหตุ ธมฺโม มหคฺคตสฺส เหตุสฺส ธมฺมสฺส อนนฺตรปจฺจเยน ปจฺจโย.\csromanspacer{}\csromanspacer{}ปริตฺโต เหตุ ธมฺโม อปฺปมาณสฺส เหตุสฺส ธมฺมสฺส อนนฺตรปจฺจเยน ปจฺจโย. (3)}

\prose{มหคฺคโต เหตุ ธมฺโม มหคฺคตสฺส เหตุสฺส ธมฺมสฺส อนนฺตรปจฺจเยน ปจฺจโย.\csromanspacer{}\csromanspacer{}มหคฺคโต เหตุ ธมฺโม ปริตฺตสฺส เหตุสฺส ธมฺมสฺส อนนฺตรปจฺจเยน ปจฺจโย.\csromanspacer{}\csromanspacer{}มหคฺคโต เหตุ ธมฺโม อปฺปมาณสฺส เหตุสฺส ธมฺมสฺส อนนฺตรปจฺจเยน ปจฺจโย. (3)}

\prose{อปฺปมาโณ เหตุ ธมฺโม อปฺปมาณสฺส เหตุสฺส ธมฺมสฺส อนนฺตรปจฺจเยน ปจฺจโย.\csromanspacer{}\csromanspacer{}อปฺปมาโณ เหตุ ธมฺโม ปริตฺตสฺส เหตุสฺส ธมฺมสฺส อนนฺตรปจฺจเยน ปจฺจโย.\csromanspacer{}\csromanspacer{}อปฺปมาโณ เหตุ ธมฺโม มหคฺคตสฺส เหตุสฺส ธมฺมสฺส อนนฺตรปจฺจเยน ปจฺจโย. (3) ฯเปฯ.}

\prose{ปริตฺโต เหตุ ธมฺโม ปริตฺตสฺส เหตุสฺส ธมฺมสฺส อุปนิสฺสยปจฺจเยน ปจฺจโย. (สํขิตฺตํ.)}

\csromanheader{2}{12. ปริตฺตตฺติก 1. เหตุทุก}
\proseitem{161}{\csromanfolio{523}เหตุยา ตีณิ, อารมฺมเณ ฉ, อธิปติยา ปญฺจ, อนนฺตเร นว, สมนนฺตเร นว, สหชาเต ตีณิ, อญฺญมญฺเญ ตีณิ, นิสฺสเย ตีณิ, อุปนิสฺสเย นว ฯเปฯ อวิคเต ตีณิ. (สํขิตฺตํ.)}

\prose{นเหตุยา นว, น-อารมฺมเณ นว. (สํขิตฺตํ.)}

\csromanheader{2}{เหตุปจฺจยา น-อารมฺมเณ ตีณิ. (สํขิตฺตํ.)}
\csromanheader{2}{นเหตุปจฺจยา อารมฺมเณ ฉ. (สํขิตฺตํ.)}
\prose{(ยถา กุสลตฺติเก ปญฺหาวารํ, เอวํ วิตฺถาเรตพฺพํ.)}

\csromanheader{2}{\textbf{นเหตุปท 1. ปฏิจฺจวาร}}
\csromanheader{2}{\textbf{ปจฺจยจตุกฺก-เหตุ}}
\proseitem{162}{ปริตฺตํ นเหตุํ ธมฺมํ ปฏิจฺจ ปริตฺโต นเหตุ ธมฺโม อุปฺปชฺชติ เหตุปจฺจยา.\csromanspacer{} ตีณิ.}

\prose{มหคฺคตํ นเหตุํ ธมฺมํ ปฏิจฺจ มหคฺคโต นเหตุ ธมฺโม อุปฺปชฺชติ เหตุปจฺจยา.\csromanspacer{}\csromanspacer{}มหคฺคตํ นเหตุํ ธมฺมํ ปฏิจฺจ ปริตฺโต นเหตุ ธมฺโม อุปฺปชฺชติ เหตุปจฺจยา.\csromanspacer{}\csromanspacer{}มหคฺคตํ นเหตุํ ธมฺมํ ปฏิจฺจ ปริตฺโต นเหตุ จ มหคฺคโต นเหตุ จ ธมฺมา อุปฺปชฺชนฺติ เหตุปจฺจยา. (3)}

\prose{อปฺปมาณํ นเหตุํ ธมฺมํ ปฏิจฺจ อปฺปมาโณ นเหตุ ธมฺโม อุปฺปชฺชติ เหตุปจฺจยา.\csromanspacer{}\csromanspacer{}อปฺปมาณํ นเหตุํ ธมฺมํ ปฏิจฺจ ปริตฺโต นเหตุ ธมฺโม อุปฺปชฺชติ เหตุปจฺจยา.\csromanspacer{}\csromanspacer{}อปฺปมาณํ นเหตุํ ธมฺมํ ปฏิจฺจ ปริตฺโต นเหตุ จ อปฺปมาโณ นเหตุ จ ธมฺมา อุปฺปชฺชนฺติ เหตุปจฺจยา. (3)}

\prose{ปริตฺตํ นเหตุญฺจ มหคฺคตํ นเหตุญฺจ ธมฺมํ ปฏิจฺจ ปริตฺโต นเหตุ ธมฺโม อุปฺปชฺชติ เหตุปจฺจยา.\csromanspacer{} ตีณิ.}

\prose{ปริตฺตํ นเหตุญฺจ อปฺปมาณํ นเหตุญฺจ ธมฺมํ ปฏิจฺจ ปริตฺโต นเหตุ ธมฺโม อุปฺปชฺชติ เหตุปจฺจยา. (1) (สํขิตฺตํ.)}

\proseitem{163}{เหตุยา เตรส, อารมฺมเณ ปญฺจ, อธิปติยา นว ฯเปฯ วิคเต ปญฺจ, อวิคเต เตรส. (สํขิตฺตํ.)}

\csromanheader{2}{ธมฺมานุโลม ติกทุกปฏฺฐานปาฬิ \textbf{ปจฺจนีย-นเหตุ}}
\proseitem{164}{\csromanfolio{524}ปริตฺตํ นเหตุํ ธมฺมํ ปฏิจฺจ ปริตฺโต นเหตุ ธมฺโม อุปฺปชฺชติ นเหตุปจฺจยา. (1) (สํขิตฺตํ.)}

\prose{นเหตุยา เอกํ, น-อารมฺมเณ ปญฺจ, น-อธิปติยา ทส ฯเปฯ โนวิคเต ปญฺจ. (สํขิตฺตํ.)}

\csromanheader{2}{เหตุปจฺจยา น-อารมฺมเณ ปญฺจ. (สํขิตฺตํ.)}
\csromanheader{2}{นเหตุปจฺจยา อารมฺมเณ เอกํ. (สํขิตฺตํ.)}
\prose{(สหชาตวารมฺปิ ฯเปฯ สมฺปยุตฺตวารมฺปิ ปฏิจฺจวารสทิสํ วิตฺถาเรตพฺพํ.)}

\csromanheader{2}{7. ปญฺหาวาร ปจฺจยจตุกฺก}
\proseitem{165}{ปริตฺโต นเหตุ ธมฺโม ปริตฺตสฺส นเหตุสฺส ธมฺมสฺส อารมฺมณปจฺจเยน ปจฺจโย.\csromanspacer{}\csromanspacer{}ปริตฺโต นเหตุ ธมฺโม มหคฺคตสฺส นเหตุสฺส ธมฺมสฺส อารมฺมณปจฺจเยน ปจฺจโย. (2)}

\prose{มหคฺคโต นเหตุ ธมฺโม มหคฺคตสฺส นเหตุสฺส ธมฺมสฺส อารมฺมณปจฺจเยน ปจฺจโย.\csromanspacer{}\csromanspacer{}มหคฺคโต นเหตุ ธมฺโม ปริตฺตสฺส นเหตุสฺส ธมฺมสฺส อารมฺมณปจฺจเยน ปจฺจโย. (2)}

\prose{อปฺปมาโณ นเหตุ ธมฺโม อปฺปมาณสฺส นเหตุสฺส ธมฺมสฺส อารมฺมณปจฺจเยน ปจฺจโย.\csromanspacer{}\csromanspacer{}อปฺปมาโณ นเหตุ ธมฺโม ปริตฺตสฺส นเหตุสฺส ธมฺมสฺส อารมฺมณปจฺจเยน ปจฺจโย.\csromanspacer{}\csromanspacer{}อปฺปมาโณ นเหตุ ธมฺโม มหคฺคตสฺส นเหตุสฺส ธมฺมสฺส อารมฺมณปจฺจเยน ปจฺจโย. (3)}

\prose{ปริตฺโต นเหตุ ธมฺโม ปริตฺตสฺส นเหตุสฺส ธมฺมสฺส อธิปติปจฺจเยน ปจฺจโย. (สํขิตฺตํ.)}

\proseitem{166}{อารมฺมเณ สตฺต, อธิปติยา สตฺต, อนนฺตเร นว, สมนนฺตเร นว, สหชาเต เอกาทส, อญฺญมญฺเญ สตฺต, นิสฺสเย เตรส, อุปนิสฺสเย นว ฯเปฯ อวิคเต เตรส. (สํขิตฺตํ.)}

\csromanmatikaanchor{mtk.116}
\csromanmatikaanchor{mtk.117}
\csromantocmarknopage{subparagraph}{13. ปริตฺตารมฺมณตฺติก}{mtk.116}
\bookmark[dest={mtk.116},level=5]{13. ปริตฺตารมฺมณตฺติก}
\csromantocmarkref{subparagraph}{1. เหตุทุก}{mtk.117}
\bookmark[dest={mtk.117},level=5]{1. เหตุทุก}
\csromantocmarknopage{subparagraph}{14. หีนตฺติก}{mtk.118}
\bookmark[dest={mtk.118},level=5]{14. หีนตฺติก}
\csromanheader{6}{13. ปริตฺตารมฺมณตฺติก 1. เหตุทุก}
\prose{\csromanfolio{525}นเหตุยา ปนฺนรส, น-อารมฺมเณ ปนฺนรส. (สํขิตฺตํ.)}

\csromanheader{2}{อารมฺมณปจฺจยา นเหตุยา สตฺต. (สํขิตฺตํ.)}
\csromanheader{2}{นเหตุปจฺจยา อารมฺมเณ สตฺต. (สํขิตฺตํ.)}
\prose{(ยถา กุสลตฺติเก ปญฺหาวารํ, เอวํ วิตฺถาเรตพฺพํ.)}

\csromanheader{2}{13. ปริตฺตารมฺมณตฺติก 1. เหตุทุก}
\csromanheader{2}{1-7. ปฏิจฺจาทิวาร ปจฺจยจตุกฺก-เหตุ}
\proseitem{167}{ปริตฺตารมฺมณํ เหตุํ ธมฺมํ ปฏิจฺจ ปริตฺตารมฺมโณ เหตุ ธมฺโม อุปฺปชฺชติ เหตุปจฺจยา. (1)}

\prose{มหคฺคตารมฺมณํ เหตุํ ธมฺมํ ปฏิจฺจ มหคฺคตารมฺมโณ เหตุ ธมฺโม อุปฺปชฺชติ เหตุปจฺจยา. (1)}

\prose{อปฺปมาณารมฺมณํ เหตุํ ธมฺมํ ปฏิจฺจ อปฺปมาณารมฺมโณ เหตุ ธมฺโม อุปฺปชฺชติ เหตุปจฺจยา. (1) (สํขิตฺตํ.)}

\prose{เหตุยา ตีณิ, อารมฺมเณ ตีณิ ฯเปฯ วิคเต ตีณิ, อวิคเต ตีณิ. (สํขิตฺตํ.)}

\prose{น-อธิปติยา ตีณิ, นปุเรชาเต ตีณิ, นปจฺฉาชาเต ตีณิ, น-อาเสวเน ตีณิ, นวิปาเก ตีณิ, นวิปฺปยุตฺเต ตีณิ. (สํขิตฺตํ.)}

\prose{(สหชาตวารมฺปิ ฯเปฯ สมฺปยุตฺตวารมฺปิ ปฏิจฺจวารสทิสํ วิตฺถาเรตพฺพํ.)}

\csromanheader{2}{\textbf{เหตุ-อารมฺมณ}}
\proseitem{168}{ปริตฺตารมฺมโณ เหตุ ธมฺโม ปริตฺตารมฺมณสฺส เหตุสฺส ธมฺมสฺส เหตุปจฺจเยน ปจฺจโย. (1)}

\prose{มหคฺคตารมฺมโณ เหตุ ธมฺโม มหคฺคตารมฺมณสฺส เหตุสฺส ธมฺมสฺส เหตุปจฺจเยน ปจฺจโย. (1)}

\prose{อปฺปมาณารมฺมโณ เหตุ ธมฺโม อปฺปมาณารมฺมณสฺส เหตุสฺส ธมฺมสฺส เหตุปจฺจเยน ปจฺจโย. (1)}

\gambhira{ธมฺมานุโลม ติกทุกปฏฺฐานปาฬิ}
\prose{\csromanfolio{526}ปริตฺตารมฺมโณ เหตุ ธมฺโม ปริตฺตารมฺมณสฺส เหตุสฺส ธมฺมสฺส อารมฺมณปจฺจเยน ปจฺจโย.\csromanspacer{}\csromanspacer{}ปริตฺตารมฺมโณ เหตุ ธมฺโม มหคฺคตารมฺมณสฺส เหตุสฺส ธมฺมสฺส อารมฺมณปจฺจเยน ปจฺจโย. (2)}

\prose{มหคฺคตารมฺมโณ เหตุ ธมฺโม มหคฺคตารมฺมณสฺส เหตุสฺส ธมฺมสฺส อารมฺมณปจฺจเยน ปจฺจโย.\csromanspacer{}\csromanspacer{}มหคฺคตารมฺมโณ เหตุ ธมฺโม ปริตฺตารมฺมณสฺส เหตุสฺส ธมฺมสฺส อารมฺมณปจฺจเยน ปจฺจโย. (2)}

\prose{อปฺปมาณารมฺมโณ เหตุ ธมฺโม อปฺปมาณารมฺมณสฺส เหตุสฺส ธมฺมสฺส อารมฺมณปจฺจเยน ปจฺจโย.\csromanspacer{}\csromanspacer{}อปฺปมาณารมฺมโณ เหตุ ธมฺโม ปริตฺตารมฺมณสฺส เหตุสฺส ธมฺมสฺส อารมฺมณปจฺจเยน ปจฺจโย.\csromanspacer{}\csromanspacer{}อปฺปมาณารมฺมโณ เหตุ ธมฺโม มหคฺคตารมฺมณสฺส เหตุสฺส ธมฺมสฺส อารมฺมณปจฺจเยน ปจฺจโย. (3) (สํขิตฺตํ.)}

\proseitem{169}{เหตุยา ตีณิ, อารมฺมเณ สตฺต, อธิปติยา สตฺต, อนนฺตเร นว, สมนนฺตเร นว, อุปนิสฺสเย นว ฯเปฯ อวิคเต ตีณิ. (สํขิตฺตํ.)}

\prose{นเหตุยา นว, น-อารมฺมเณ นว. (สํขิตฺตํ.)}

\csromanheader{2}{เหตุปจฺจยา น-อารมฺมเณ ตีณิ. (สํขิตฺตํ.)}
\csromanheader{2}{นเหตุปจฺจยา อารมฺมเณ สตฺต. (สํขิตฺตํ.)}
\prose{(ยถา กุสลตฺติเก ปญฺหาวารํ, เอวํ วิตฺถาเรตพฺพํ.)}

\csromanheader{2}{\textbf{นเหตุปท 1-7. ปฏิจฺจาทิวาร}}
\csromanheader{2}{\textbf{ปจฺจยจตุกฺก-เหตุ}}
\proseitem{170}{ปริตฺตารมฺมณํ นเหตุํ ธมฺมํ ปฏิจฺจ ปริตฺตารมฺมโณ นเหตุ ธมฺโม อุปฺปชฺชติ เหตุปจฺจยา. (สํขิตฺตํ.)}

\prose{เหตุยา ตีณิ, อารมฺมเณ ตีณิ, อธิปติยา ตีณิ ฯเปฯ อวิคเต ตีณิ. (สํขิตฺตํ.)}

\csromanheader{2}{13. ปริตฺตารมฺมณตฺติก 1. เหตุทุก}
\proseitem{171}{\csromanfolio{527}ปริตฺตารมฺมณํ นเหตุํ ธมฺมํ ปฏิจฺจ ปริตฺตารมฺมโณ นเหตุ ธมฺโม อุปฺปชฺชติ นเหตุปจฺจยา. (สํขิตฺตํ.)}

\prose{นเหตุยา ตีณิ, น-อธิปติยา ตีณิ ฯเปฯ นวิปฺปยุตฺเต ตีณิ. (สํขิตฺตํ.)}

\csromanheader{2}{เหตุปจฺจยา น-อธิปติยา ตีณิ. (สํขิตฺตํ.)}
\csromanheader{2}{นเหตุปจฺจยา อารมฺมเณ ตีณิ. (สํขิตฺตํ.)}
\prose{(สหชาตวารมฺปิ ฯเปฯ สมฺปยุตฺตวารมฺปิ ปฏิจฺจวารสทิสํ วิตฺถาเรตพฺพํ.)}

\csromanheader{2}{\textbf{อารมฺมณ-อธิปติ}}
\proseitem{172}{ปริตฺตารมฺมโณ นเหตุ ธมฺโม ปริตฺตารมฺมณสฺส นเหตุสฺส ธมฺมสฺส อารมฺมณปจฺจเยน ปจฺจโย.\csromanspacer{}\csromanspacer{}ปริตฺตารมฺมโณ นเหตุ ธมฺโม มหคฺคตารมฺมณสฺส นเหตุสฺส ธมฺมสฺส อารมฺมณปจฺจเยน ปจฺจโย. (2)}

\prose{มหคฺคตารมฺมโณ นเหตุ ธมฺโม มหคฺคตารมฺมณสฺส นเหตุสฺส ธมฺมสฺส อารมฺมณปจฺจเยน ปจฺจโย.\csromanspacer{}\csromanspacer{}มหคฺคตารมฺมโณ นเหตุ ธมฺโม ปริตฺตารมฺมณสฺส นเหตุสฺส ธมฺมสฺส อารมฺมณปจฺจเยน ปจฺจโย. (2)}

\prose{อปฺปมาณารมฺมโณ นเหตุ ธมฺโม อปฺปมาณารมฺมณสฺส นเหตุสฺส ธมฺมสฺส อารมฺมณปจฺจเยน ปจฺจโย.\csromanspacer{}\csromanspacer{}อปฺปมาณารมฺมโณ นเหตุ ธมฺโม ปริตฺตารมฺมณสฺส นเหตุสฺส ธมฺมสฺส อารมฺมณปจฺจเยน ปจฺจโย.\csromanspacer{}\csromanspacer{}อปฺปมาณารมฺมโณ นเหตุ ธมฺโม มหคฺคตารมฺมณสฺส นเหตุสฺส ธมฺมสฺส อารมฺมณปจฺจเยน ปจฺจโย. (3)}

\prose{ปริตฺตารมฺมโณ นเหตุ ธมฺโม ปริตฺตารมฺมณสฺส นเหตุสฺส ธมฺมสฺส อธิปติปจฺจเยน ปจฺจโย. (สํขิตฺตํ.)}

\proseitem{173}{อารมฺมเณ สตฺต, อธิปติยา สตฺต, อนนฺตเร นว, สมนนฺตเร นว, อุปนิสฺสเย นว ฯเปฯ อวิคเต ตีณิ. (สํขิตฺตํ.)}

\prose{นเหตุยา นว, น-อารมฺมเณ นว. (สํขิตฺตํ.)}

\csromanheader{2}{อารมฺมณปจฺจยา นเหตุยา สตฺต. (สํขิตฺตํ.)}
\csromanheader{2}{นเหตุปจฺจยา อารมฺมเณ สตฺต. (สํขิตฺตํ.)}
\prose{(ยถา กุสลตฺติเก ปญฺหาวารํ, เอวํ วิตฺถาเรตพฺพํ.)}

\csromanmatikaanchor{mtk.119}
\csromantocmarkref{subparagraph}{1. เหตุทุก}{mtk.119}
\bookmark[dest={mtk.119},level=5]{1. เหตุทุก}
\csromanheader{6}{ธมฺมานุโลม ติกทุกปฏฺฐานปาฬิ \textbf{14. หีนตฺติก 1. เหตุทุก}}
\csromanheader{2}{1-7. ปฏิจฺจาทิวาร ปจฺจยจตุกฺก-เหตุ}
\proseitem{174}{\csromanfolio{528}หีนํ เหตุํ ธมฺมํ ปฏิจฺจ หีโน เหตุ ธมฺโม อุปฺปชฺชติ เหตุปจฺจยา.}

\prose{มชฺฌิมํ เหตุํ ธมฺมํ ปฏิจฺจ มชฺฌิโม เหตุ ธมฺโม อุปฺปชฺชติ เหตุปจฺจยา.}

\prose{ปณีตํ เหตุํ ธมฺมํ ปฏิจฺจ ปณีโต เหตุ ธมฺโม อุปฺปชฺชติ เหตุปจฺจยา. (3) (สํขิตฺตํ.)}

\proseitem{175}{เหตุยา ตีณิ, อารมฺมเณ ตีณิ ฯเปฯ อวิคเต ตีณิ. (สํขิตฺตํ.)}

\prose{น-อธิปติยา ตีณิ, นปุเรชาเต ตีณิ, นปจฺฉาชาเต ตีณิ, น-อาเสวเน ตีณิ, นวิปาเก ตีณิ, นวิปฺปยุตฺเต ตีณิ. (สํขิตฺตํ.)}

\prose{(สหชาตวารมฺปิ ฯเปฯ สมฺปยุตฺตวารมฺปิ ปฏิจฺจวารสทิสํ วิตฺถาเรตพฺพํ.)}

\csromanheader{2}{\textbf{เหตุ-อารมฺมณาทิ}}
\proseitem{176}{หีโน เหตุ ธมฺโม หีนสฺส เหตุสฺส ธมฺมสฺส เหตุปจฺจเยน ปจฺจโย. (1)}

\prose{มชฺฌิโม เหตุ ธมฺโม มชฺฌิมสฺส เหตุสฺส ธมฺมสฺส เหตุปจฺจเยน ปจฺจโย. (1)}

\prose{ปณีโต เหตุ ธมฺโม ปณีตสฺส เหตุสฺส ธมฺมสฺส เหตุปจฺจเยน ปจฺจโย. (1)}

\prose{หีโน เหตุ ธมฺโม หีนสฺส เหตุสฺส ธมฺมสฺส อารมฺมณปจฺจเยน ปจฺจโย.\csromanspacer{}\csromanspacer{}หีโน เหตุ ธมฺโม มชฺฌิมสฺส เหตุสฺส ธมฺมสฺส อารมฺมณปจฺจเยน ปจฺจโย. (2)}

\prose{มชฺฌิโม เหตุ ธมฺโม มชฺฌิมสฺส เหตุสฺส ธมฺมสฺส อารมฺมณปจฺจเยน ปจฺจโย.\csromanspacer{}\csromanspacer{}มชฺฌิโม เหตุ ธมฺโม หีนสฺส เหตุสฺส ธมฺมสฺส อารมฺมณปจฺจเยน ปจฺจโย. (2)}

\prose{ปณีโต เหตุ ธมฺโม มชฺฌิมสฺส เหตุสฺส ธมฺมสฺส อารมฺมณปจฺจเยน ปจฺจโย. (1)}

\csromanheader{2}{14. หีนตฺติก 1. เหตุทุก}
\proseitem{177}{\csromanfolio{529}หีโน เหตุ ธมฺโม หีนสฺส เหตุสฺส ธมฺมสฺส อธิปติปจฺจเยน ปจฺจโย. (1)}

\prose{มชฺฌิโม เหตุ ธมฺโม มชฺฌิมสฺส เหตุสฺส ธมฺมสฺส อธิปติปจฺจเยน ปจฺจโย.\csromanspacer{}\csromanspacer{}มชฺฌิโม เหตุ ธมฺโม หีนสฺส เหตุสฺส ธมฺมสฺส อธิปติปจฺจเยน ปจฺจโย. (2)}

\prose{ปณีโต เหตุ ธมฺโม ปณีตสฺส เหตุสฺส ธมฺมสฺส อธิปติปจฺจเยน ปจฺจโย.\csromanspacer{}\csromanspacer{}ปณีโต เหตุ ธมฺโม มชฺฌิมสฺส เหตุสฺส ธมฺมสฺส อธิปติปจฺจเยน ปจฺจโย. (2)}

\prose{หีโน เหตุ ธมฺโม หีนสฺส เหตุสฺส ธมฺมสฺส อนนฺตรปจฺจเยน ปจฺจโย.\csromanspacer{} เทฺว.}

\prose{มชฺฌิโม เหตุ ธมฺโม มชฺฌิมสฺส เหตุสฺส ธมฺมสฺส อนนฺตรปจฺจเยน ปจฺจโย.\csromanspacer{} เทฺว.}

\prose{ปณีโต เหตุ ธมฺโม ปณีตสฺส เหตุสฺส ธมฺมสฺส อนนฺตรปจฺจเยน ปจฺจโย.\csromanspacer{} เทฺว. (สํขิตฺตํ.)}

\proseitem{178}{เหตุยา ตีณิ, อารมฺมเณ ปญฺจ, อธิปติยา ปญฺจ ฯเปฯ อวิคเต ตีณิ. (สํขิตฺตํ.)}

\prose{นเหตุยา อฏฺฐ, น-อารมฺมเณ อฏฺฐ. (สํขิตฺตํ.)}

\csromanheader{2}{เหตุปจฺจยา น-อารมฺมเณ ตีณิ. (สํขิตฺตํ.)}
\csromanheader{2}{นเหตุปจฺจยา อารมฺมเณ ปญฺจ. (สํขิตฺตํ.)}
\prose{(ยถา กุสลตฺติเก ปญฺหาวารํ, เอวํ วิตฺถาเรตพฺพํ.)}

\csromanheader{2}{\textbf{นเหตุปท-เหตุ}}
\proseitem{179}{หีนํ นเหตุํ ธมฺมํ ปฏิจฺจ หีโน นเหตุ ธมฺโม อุปฺปชฺชติ เหตุปจฺจยา.}

\prose{(ยถา สํกิลิฏฺฐตฺติกเหตุทุกํ, เอวํ วิตฺถาเรตพฺพํ.)}

\csromanmatikaanchor{mtk.120}
\csromanmatikaanchor{mtk.121}
\csromantocmarknopage{subparagraph}{15. มิจฺฉตฺตนิยตตฺติก}{mtk.120}
\bookmark[dest={mtk.120},level=5]{15. มิจฺฉตฺตนิยตตฺติก}
\csromantocmarkref{subparagraph}{1. เหตุทุก}{mtk.121}
\bookmark[dest={mtk.121},level=5]{1. เหตุทุก}
\csromanheader{6}{ธมฺมานุโลม ติกทุกปฏฺฐานปาฬิ \textbf{15. มิจฺฉตฺตนิยตตฺติก 1. เหตุทุก}}
\csromanheader{2}{1-7. ปฏิจฺจาทิวาร ปจฺจยจตุกฺก-เหตุ}
\proseitem{180}{\csromanfolio{530}มิจฺฉตฺตนิยตํ เหตุํ ธมฺมํ ปฏิจฺจ มิจฺฉตฺตนิยโต เหตุ ธมฺโม อุปฺปชฺชติ เหตุปจฺจยา. (1)}

\prose{สมฺมตฺตนิยตํ เหตุํ ธมฺมํ ปฏิจฺจ สมฺมตฺตนิยโต เหตุ ธมฺโม อุปฺปชฺชติ เหตุปจฺจยา. (1)}

\prose{อนิยตํ เหตุํ ธมฺมํ ปฏิจฺจ อนิยโต เหตุ ธมฺโม อุปฺปชฺชติ เหตุปจฺจยา. (1) (สํขิตฺตํ.)}

\proseitem{181}{เหตุยา ตีณิ, อารมฺมเณ ตีณิ, อธิปติยา ตีณิ ฯเปฯ วิปาเก เอกํ ฯเปฯ อวิคเต ตีณิ. (สํขิตฺตํ.)}

\prose{น-อธิปติยา เทฺว, นปุเรชาเต เทฺว, นปจฺฉาชาเต ตีณิ, น-อาเสวเน เอกํ, นวิปาเก ตีณิ, นวิปฺปยุตฺเต เทฺว. (สํขิตฺตํ.)}

\prose{(สหชาตวารมฺปิ ฯเปฯ สมฺปยุตฺตวารมฺปิ ปฏิจฺจวารสทิสํ วิตฺถาเรตพฺพํ.)}

\csromanheader{2}{\textbf{เหตุ อารมฺมณาทิ}}
\proseitem{182}{มิจฺฉตฺตนิยโต เหตุ ธมฺโม มิจฺฉตฺตนิยตสฺส เหตุสฺส ธมฺมสฺส เหตุปจฺจเยน ปจฺจโย. (1)}

\prose{สมฺมตฺตนิยโต เหตุ ธมฺโม สมฺมตฺตนิยตสฺส เหตุสฺส ธมฺมสฺส เหตุปจฺจเยน ปจฺจโย. (1)}

\prose{อนิยโต เหตุ ธมฺโม อนิยตสฺส เหตุสฺส ธมฺมสฺส เหตุปจฺจเยน ปจฺจโย. (1)}

\prose{มิจฺฉตฺตนิยโต เหตุ ธมฺโม อนิยตสฺส เหตุสฺส ธมฺมสฺส อารมฺมณปจฺจเยน ปจฺจโย. (1)}

\prose{สมฺมตฺตนิยโต เหตุ ธมฺโม อนิยตสฺส เหตุสฺส ธมฺมสฺส อารมฺมณปจฺจเยน ปจฺจโย. (1)}

\prose{อนิยโต เหตุ ธมฺโม อนิยตสฺส เหตุสฺส ธมฺมสฺส อารมฺมณปจฺจเยน ปจฺจโย. (1)}

\csromanheader{2}{15. มิจฺฉตฺตนิยตตฺติก 1. เหตุทุก}
\prose{\csromanfolio{531}อนิยโต เหตุ ธมฺโม มิจฺฉตฺตนิยตสฺส เหตุสฺส ธมฺมสฺส อารมฺมณปจฺจเยน ปจฺจโย. (1)}

\prose{สมฺมตฺตนิยโต เหตุ ธมฺโม สมฺมตฺตนิยตสฺส เหตุสฺส ธมฺมสฺส อธิปติปจฺจเยน ปจฺจโย.\csromanspacer{}\csromanspacer{}สมฺมตฺตนิยโต เหตุ ธมฺโม อนิยตสฺส เหตุสฺส ธมฺมสฺส อธิปติปจฺจเยน ปจฺจโย. (2)}

\prose{อนิยโต เหตุ ธมฺโม อนิยตสฺส เหตุสฺส ธมฺมสฺส อธิปติปจฺจเยน ปจฺจโย. (1)}

\csromanheader{2}{\textbf{อนนฺตราทิ}}
\proseitem{183}{มิจฺฉตฺตนิยโต เหตุ ธมฺโม อนิยตสฺส เหตุสฺส ธมฺมสฺส อนนฺตรปจฺจเยน ปจฺจโย. (1)}

\prose{สมฺมตฺตนิยโต เหตุ ธมฺโม อนิยตสฺส เหตุสฺส ธมฺมสฺส อนนฺตรปจฺจเยน ปจฺจโย. (1)}

\prose{อนิยโต เหตุ ธมฺโม อนิยตสฺส เหตุสฺส ธมฺมสฺส อนนฺตรปจฺจเยน ปจฺจโย.\csromanspacer{}\csromanspacer{}อนิยโต เหตุ ธมฺโม มิจฺฉตฺตนิยตสฺส เหตุสฺส ธมฺมสฺส อนนฺตรปจฺจเยน ปจฺจโย.\csromanspacer{}\csromanspacer{}อนิยโต เหตุ ธมฺโม สมฺมตฺตนิยตสฺส เหตุสฺส ธมฺมสฺส อนนฺตรปจฺจเยน ปจฺจโย. (3) ฯเปฯ.}

\proseitem{184}{มิจฺฉตฺตนิยโต เหตุ ธมฺโม มิจฺฉตฺตนิยตสฺส เหตุสฺส ธมฺมสฺส อุปนิสฺสยปจฺจเยน ปจฺจโย.\csromanspacer{}\csromanspacer{}มิจฺฉตฺตนิยโต เหตุ ธมฺโม อนิยตสฺส เหตุสฺส ธมฺมสฺส อุปนิสฺสยปจฺจเยน ปจฺจโย. (2)}

\prose{สมฺมตฺตนิยโต เหตุ ธมฺโม สมฺมตฺตนิยตสฺส เหตุสฺส ธมฺมสฺส อุปนิสฺสยปจฺจเยน ปจฺจโย.\csromanspacer{}\csromanspacer{}สมฺมตฺตนิยโต เหตุ ธมฺโม อนิยตสฺส เหตุสฺส ธมฺมสฺส อุปนิสฺสยปจฺจเยน ปจฺจโย. (2)}

\prose{อนิยโต เหตุ ธมฺโม อนิยตสฺส เหตุสฺส ธมฺมสฺส อุปนิสฺสยปจฺจเยน ปจฺจโย.\csromanspacer{}\csromanspacer{}อนิยโต เหตุ ธมฺโม มิจฺฉตฺตนิยตสฺส เหตุสฺส ธมฺมสฺส อุปนิสฺสยปจฺจเยน ปจฺจโย.\csromanspacer{}\csromanspacer{}อนิยโต เหตุ ธมฺโม สมฺมตฺตนิยตสฺส เหตุสฺส ธมฺมสฺส อุปนิสฺสยปจฺจเยน ปจฺจโย. (3) (สํขิตฺตํ.)}

\gambhira{ธมฺมานุโลม ติกทุกปฏฺฐานปาฬิ}
\proseitem{185}{\csromanfolio{532}เหตุยา ตีณิ, อารมฺมเณ จตฺตาริ, อธิปติยา ตีณิ, อนนฺตเร ปญฺจ, สมนนฺตเร ปญฺจ, สหชาเต ตีณิ ฯเปฯ อุปนิสฺสเย สตฺต ฯเปฯ อวิคเต ตีณิ. (สํขิตฺตํ.)}

\prose{นเหตุยา สตฺต, น-อารมฺมเณ สตฺต. (สํขิตฺตํ.)}

\csromanheader{2}{เหตุปจฺจยา น-อารมฺมเณ ตีณิ. (สํขิตฺตํ.)}
\csromanheader{2}{นเหตุปจฺจยา อารมฺมเณ จตฺตาริ. (สํขิตฺตํ.)}
\prose{(ยถา กุสลตฺติเก ปญฺหาวารํ, เอวํ วิตฺถาเรตพฺพํ.)}

\csromanheader{2}{\textbf{นเหตุปท 1-7. ปฏิจฺจาทิวาร}}
\csromanheader{2}{\textbf{ปจฺจยจตุกฺก-เหตุ}}
\proseitem{186}{มิจฺฉตฺตนิยตํ นเหตุํ ธมฺมํ ปฏิจฺจ มิจฺฉตฺตนิยโต นเหตุ ธมฺโม อุปฺปชฺชติ เหตุปจฺจยา.\csromanspacer{} ตีณิ.}

\prose{สมฺมตฺตนิยตํ นเหตุํ ธมฺมํ ปฏิจฺจ สมฺมตฺตนิยโต นเหตุ ธมฺโม อุปฺปชฺชติ เหตุปจฺจยา.\csromanspacer{}\csromanspacer{}สมฺมตฺตนิยตํ นเหตุํ ธมฺมํ ปฏิจฺจ อนิยโต นเหตุ ธมฺโม อุปฺปชฺชติ เหตุปจฺจยา.\csromanspacer{}\csromanspacer{}สมฺมตฺตนิยตํ นเหตุํ ธมฺมํ ปฏิจฺจ สมฺมตฺตนิยโต นเหตุ จ อนิยโต นเหตุ จ ธมฺมา อุปฺปชฺชนฺติ เหตุปจฺจยา. (3)}

\prose{อนิยตํ นเหตุํ ธมฺมํ ปฏิจฺจ อนิยโต นเหตุ ธมฺโม อุปฺปชฺชติ เหตุปจฺจยา. (1)}

\prose{มิจฺฉตฺตนิยตํ นเหตุญฺจ อนิยตํ นเหตุญฺจ ธมฺมํ ปฏิจฺจ อนิยโต นเหตุ ธมฺโม อุปฺปชฺชติ เหตุปจฺจยา. (1)}

\prose{สมฺมตฺตนิยตํ นเหตุญฺจ อนิยตํ นเหตุญฺจ ธมฺมํ ปฏิจฺจ อนิยโต นเหตุ ธมฺโม อุปฺปชฺชติ เหตุปจฺจยา. (1) (สํขิตฺตํ.) เหตุยา นว, อารมฺมเณ ตีณิ, อธิปติยา นว ฯเปฯ วิปาเก เอกํ ฯเปฯ อวิคเต นว. (สํขิตฺตํ.)}

\csromanheader{2}{15. มิจฺฉตฺตนิยตตฺติก 1. เหตุทุก}
\csromanheader{2}{\textbf{ปจฺจนีย-นเหตุ}}
\proseitem{188}{\csromanfolio{533}อนิยตํ นเหตุํ ธมฺมํ ปฏิจฺจ อนิยโต นเหตุ ธมฺโม อุปฺปชฺชติ นเหตุปจฺจยา. (สํขิตฺตํ.)}

\prose{นเหตุยา เอกํ, น-อารมฺมเณ ปญฺจ, น-อธิปติยา ตีณิ ฯเปฯ โนวิคเต ปญฺจ. (สํขิตฺตํ.)}

\csromanheader{2}{เหตุปจฺจยา น-อารมฺมเณ ปญฺจ. (สํขิตฺตํ.)}
\csromanheader{2}{นเหตุปจฺจยา อารมฺมเณ เอกํ. (สํขิตฺตํ.)}
\prose{(สหชาตวารมฺปิ ฯเปฯ สมฺปยุตฺตวารมฺปิ ปฏิจฺจวารสทิสํ วิตฺถาเรตพฺพํ.)}

\csromanheader{2}{\textbf{อารมฺมณ อธิปติ อนนฺตร}}
\proseitem{189}{มิจฺฉตฺตนิยโต นเหตุ ธมฺโม อนิยตสฺส นเหตุสฺส ธมฺมสฺส อารมฺมณปจฺจเยน ปจฺจโย. (1)}

\prose{สมฺมตฺตนิยโต นเหตุ ธมฺโม อนิยตสฺส นเหตุสฺส ธมฺมสฺส อารมฺมณปจฺจเยน ปจฺจโย. (1)}

\prose{อนิยโต นเหตุ ธมฺโม อนิยตสฺส นเหตุสฺส ธมฺมสฺส อารมฺมณปจฺจเยน ปจฺจโย.\csromanspacer{}\csromanspacer{}อนิยโต นเหตุ ธมฺโม มิจฺฉตฺตนิยตสฺส นเหตุสฺส ธมฺมสฺส อารมฺมณปจฺจเยน ปจฺจโย.\csromanspacer{}\csromanspacer{}อนิยโต นเหตุ ธมฺโม สมฺมตฺตนิยตสฺส นเหตุสฺส ธมฺมสฺส อารมฺมณปจฺจเยน ปจฺจโย. (3)}

\proseitem{190}{มิจฺฉตฺตนิยโต นเหตุ ธมฺโม มิจฺฉตฺตนิยตสฺส นเหตุสฺส ธมฺมสฺส อธิปติปจฺจเยน ปจฺจโย.\csromanspacer{} ตีณิ.}

\prose{สมฺมตฺตนิยโต นเหตุ ธมฺโม สมฺมตฺตนิยตสฺส นเหตุสฺส ธมฺมสฺส อธิปติปจฺจเยน ปจฺจโย.\csromanspacer{}\csromanspacer{}สมฺมตฺตนิยโต นเหตุ ธมฺโม อนิยตสฺส นเหตุสฺส ธมฺมสฺส อธิปติปจฺจเยน ปจฺจโย.\csromanspacer{}\csromanspacer{}สมฺมตฺตนิยโต นเหตุ ธมฺโม สมฺมตฺตนิยตสฺส นเหตุสฺส จ อนิยตสฺส นเหตุสฺส จ ธมฺมสฺส อธิปติปจฺจเยน ปจฺจโย. (3)}

\gambhira{ธมฺมานุโลม ติกทุกปฏฺฐานปาฬิ}
\prose{\csromanfolio{534}อนิยโต นเหตุ ธมฺโม อนิยตสฺส นเหตุสฺส ธมฺมสฺส อธิปติปจฺจเยน ปจฺจโย.\csromanspacer{}\csromanspacer{}อนิยโต นเหตุ ธมฺโม สมฺมตฺตนิยตสฺส นเหตุสฺส ธมฺมสฺส อธิปติปจฺจเยน ปจฺจโย. (2)}

\prose{มิจฺฉตฺตนิยโต นเหตุ ธมฺโม อนิยตสฺส นเหตุสฺส ธมฺมสฺส อนนฺตรปจฺจเยน ปจฺจโย. (1)}

\prose{สมฺมตฺตนิยโต นเหตุ ธมฺโม อนิยตสฺส นเหตุสฺส ธมฺมสฺส อนนฺตรปจฺจเยน ปจฺจโย. (1)}

\prose{อนิยโต นเหตุ ธมฺโม อนิยตสฺส นเหตุสฺส ธมฺมสฺส อนนฺตรปจฺจเยน ปจฺจโย.\csromanspacer{}\csromanspacer{}อนิยโต นเหตุ ธมฺโม มิจฺฉตฺตนิยตสฺส นเหตุสฺส ธมฺมสฺส อนนฺตรปจฺจเยน ปจฺจโย.\csromanspacer{}\csromanspacer{}อนิยโต นเหตุ ธมฺโม สมฺมตฺตนิยตสฺส นเหตุสฺส ธมฺมสฺส อนนฺตรปจฺจเยน ปจฺจโย. (3) (สํขิตฺตํ.)}

\proseitem{191}{อารมฺมเณ ปญฺจ, อธิปติยา อฏฺฐ, อนนฺตเร ปญฺจ, สมนนฺตเร ปญฺจ, สหชาเต นว, อญฺญมญฺเญ ตีณิ, นิสฺสเย เตรส, อุปนิสฺสเย สตฺต ฯเปฯ อวิคเต เตรส. (สํขิตฺตํ.)}

\prose{นเหตุยา เตรส, น-อารมฺมเณ เตรส, น-อธิปติยา เตรส. (สํขิตฺตํ.)}

\csromanheader{2}{อารมฺมณปจฺจยา นเหตุยา ปญฺจ. (สํขิตฺตํ.)}
\csromanheader{2}{นเหตุปจฺจยา อารมฺมเณ ปญฺจ. (สํขิตฺตํ.)}
\prose{(ยถา กุสลตฺติเก ปญฺหาวารํ, เอวํ วิตฺถาเรตพฺพํ.)}

\csromanmatikaanchor{mtk.122}
\csromanmatikaanchor{mtk.123}
\csromantocmarknopage{subparagraph}{16. มคฺคารมฺมณตฺติก}{mtk.122}
\bookmark[dest={mtk.122},level=5]{16. มคฺคารมฺมณตฺติก}
\csromantocmarkref{subparagraph}{1. เหตุทุก}{mtk.123}
\bookmark[dest={mtk.123},level=5]{1. เหตุทุก}
\csromanheader{6}{16. มคฺคารมฺมณตฺติก 1. เหตุทุก}
\csromanheader{2}{1-7. ปฏิจฺจาทิวาร ปจฺจยจตุกฺก-เหตุ}
\proseitem{192}{มคฺคารมฺมณํ เหตุํ ธมฺมํ ปฏิจฺจ มคฺคารมฺมโณ เหตุ ธมฺโม อุปฺปชฺชติ เหตุปจฺจยา.\csromanspacer{} ตีณิ.}

\csromanheader{2}{16. มคฺคารมฺมณตฺติก 1. เหตุทุก}
\prose{\csromanfolio{535}มคฺคเหตุกํ เหตุํ ธมฺมํ ปฏิจฺจ มคฺคเหตุโก เหตุ ธมฺโม อุปฺปชฺชติ เหตุปจฺจยา.\csromanspacer{} ตีณิ.}

\prose{มคฺคาธิปติํ เหตุํ ธมฺมํ ปฏิจฺจ มคฺคาธิปติ เหตุ ธมฺโม อุปฺปชฺชติ เหตุปจฺจยา.\csromanspacer{} ปญฺจ.}

\prose{มคฺคารมฺมณํ เหตุญฺจ มคฺคาธิปติํ เหตุญฺจ ธมฺมํ ปฏิจฺจ มคฺคารมฺมโณ เหตุ ธมฺโม อุปฺปชฺชติ เหตุปจฺจยา.\csromanspacer{} ตีณิ.}

\prose{มคฺคเหตุกํ เหตุญฺจ มคฺคาธิปติํ เหตุญฺจ ธมฺมํ ปฏิจฺจ มคฺคเหตุโก เหตุ ธมฺโม อุปฺปชฺชติ เหตุปจฺจยา.\csromanspacer{} ตีณิ. (สํขิตฺตํ.)}

\csromanheader{2}{193. เหตุยา สตฺตรส ฯเปฯ อวิคเต สตฺตรส. (สํขิตฺตํ.)}
\prose{น-อธิปติยา สตฺตรส, นปุเรชาเต สตฺตรส ฯเปฯ น-อาเสวเน นว, นวิปาเก สตฺตรส, นวิปฺปยุตฺเต สตฺตรส. (สํขิตฺตํ.)}

\prose{(สหชาตวารมฺปิ ฯเปฯ สมฺปยุตฺตวารมฺปิ ปฏิจฺจวารสทิสํ วิตฺถาเรตพฺพํ.)}

\csromanheader{2}{\textbf{เหตุ อารมฺมณาทิ}}
\proseitem{194}{มคฺคารมฺมโณ เหตุ ธมฺโม มคฺคารมฺมณสฺส เหตุสฺส ธมฺมสฺส เหตุปจฺจเยน ปจฺจโย.\csromanspacer{} ตีณิ.}

\prose{มคฺคเหตุโก เหตุ ธมฺโม มคฺคเหตุกสฺส เหตุสฺส ธมฺมสฺส เหตุปจฺจเยน ปจฺจโย.\csromanspacer{} ตีณิ.}

\prose{มคฺคาธิปติ เหตุ ธมฺโม มคฺคาธิปติสฺส เหตุสฺส ธมฺมสฺส เหตุปจฺจเยน ปจฺจโย.\csromanspacer{} ปญฺจ.}

\prose{มคฺคารมฺมโณ เหตุ จ มคฺคาธิปติ เหตุ จ ธมฺมา มคฺคารมฺมณสฺส เหตุสฺส ธมฺมสฺส เหตุปจฺจเยน ปจฺจโย.\csromanspacer{} ตีณิ.}

\prose{มคฺคเหตุโก เหตุ จ มคฺคาธิปติ เหตุ จ ธมฺมา มคฺคเหตุกสฺส เหตุสฺส ธมฺมสฺส เหตุปจฺจเยน ปจฺจโย.\csromanspacer{} ตีณิ.}

\prose{มคฺคเหตุโก เหตุ ธมฺโม มคฺคารมฺมณสฺส เหตุสฺส ธมฺมสฺส อารมฺมณปจฺจเยน ปจฺจโย.\csromanspacer{}\csromanspacer{}มคฺคเหตุโก เหตุ ธมฺโม มคฺคาธิปติสฺส เหตุสฺส ธมฺมสฺส อารมฺมณปจฺจเยน ปจฺจโย.\csromanspacer{}\csromanspacer{}มคฺคเหตุโก เหตุ ธมฺโม มคฺคารมฺมณสฺส เหตุสฺส จ มคฺคาธิปติสฺส เหตุสฺส จ ธมฺมสฺส อารมฺมณปจฺจเยน ปจฺจโย. (3)}

\gambhira{ธมฺมานุโลม ติกทุกปฏฺฐานปาฬิ}
\prose{\csromanfolio{536}มคฺคาธิปติ เหตุ ธมฺโม มคฺคาธิปติสฺส เหตุสฺส ธมฺมสฺส อารมฺมณปจฺจเยน ปจฺจโย.\csromanspacer{} ตีณิ.}

\prose{มคฺคเหตุโก เหตุ จ มคฺคาธิปติ เหตุ จ ธมฺมา มคฺคารมฺมณสฺส เหตุสฺส ธมฺมสฺส อารมฺมณปจฺจเยน ปจฺจโย.\csromanspacer{} ตีณิ.}

\proseitem{195}{มคฺคารมฺมโณ เหตุ ธมฺโม มคฺคารมฺมณสฺส เหตุสฺส ธมฺมสฺส อธิปติปจฺจเยน ปจฺจโย.\csromanspacer{} ตีณิ.}

\prose{มคฺคเหตุโก เหตุ ธมฺโม มคฺคเหตุกสฺส เหตุสฺส ธมฺมสฺส อธิปติปจฺจเยน ปจฺจโย.\csromanspacer{} ปญฺจ.}

\prose{มคฺคาธิปติ เหตุ ธมฺโม มคฺคาธิปติสฺส เหตุสฺส ธมฺมสฺส อธิปติปจฺจเยน ปจฺจโย.\csromanspacer{} ปญฺจ.}

\prose{มคฺคารมฺมโณ เหตุ จ มคฺคาธิปติ เหตุ จ ธมฺมา มคฺคารมฺมณสฺส เหตุสฺส ธมฺมสฺส อธิปติปจฺจเยน ปจฺจโย.\csromanspacer{} ตีณิ.}

\prose{มคฺคเหตุโก เหตุ จ มคฺคาธิปติ เหตุ จ ธมฺมา มคฺคารมฺมณสฺส เหตุสฺส ธมฺมสฺส อธิปติปจฺจเยน ปจฺจโย.\csromanspacer{} ปญฺจ.}

\prose{มคฺคารมฺมโณ เหตุ ธมฺโม มคฺคารมฺมณสฺส เหตุสฺส ธมฺมสฺส อนนฺตรปจฺจเยน ปจฺจโย.\csromanspacer{} ตีณิ.}

\prose{มคฺคาธิปติ เหตุ ธมฺโม มคฺคาธิปติสฺส เหตุสฺส ธมฺมสฺส อนนฺตรปจฺจเยน ปจฺจโย.\csromanspacer{} ตีณิ.}

\prose{มคฺคารมฺมโณ เหตุ จ มคฺคาธิปติ เหตุ จ ธมฺมา มคฺคารมฺมณสฺส เหตุสฺส ธมฺมสฺส อนนฺตรปจฺจเยน ปจฺจโย.\csromanspacer{} ตีณิ. (สํขิตฺตํ.)}

\proseitem{196}{เหตุยา สตฺตรส, อารมฺมเณ นว, อธิปติยา เอกวีส, อนนฺตเร นว, สมนนฺตเร นว, สหชาเต สตฺตรส, อญฺญมญฺเญ สตฺตรส, นิสฺสเย สตฺตรส, อุปนิสฺสเย เอกวีส ฯเปฯ อวิคเต สตฺตรส. (สํขิตฺตํ.)}

\prose{นเหตุยา เอกวีส, น-อารมฺมเณ สตฺตรส. (สํขิตฺตํ.)}

\csromanheader{2}{เหตุปจฺจยา น-อารมฺมเณ สตฺตรส. (สํขิตฺตํ.)}
\csromanheader{2}{นเหตุปจฺจยา อารมฺมเณ นว. (สํขิตฺตํ.)}
\prose{(ยถา กุสลตฺติเก ปญฺหาวารํ, เอวํ วิตฺถาเรตพฺพํ.)}

\csromanheader{2}{16. มคฺคารมฺมณตฺติก 1. เหตุทุก \textbf{นเหตุปท 1. ปฏิจฺจวาร}}
\csromanheader{2}{\textbf{ปจฺจยจตุกฺก-เหตุ}}
\proseitem{197}{\csromanfolio{537}มคฺคารมฺมณํ นเหตุํ ธมฺมํ ปฏิจฺจ มคฺคารมฺมโณ นเหตุ ธมฺโม อุปฺปชฺชติ เหตุปจฺจยา.\csromanspacer{}\csromanspacer{}มคฺคารมฺมณํ นเหตุํ ธมฺมํ ปฏิจฺจ มคฺคาธิปติ นเหตุ ธมฺโม อุปฺปชฺชติ เหตุปจฺจยา.\csromanspacer{}\csromanspacer{}มคฺคารมฺมณํ นเหตุํ ธมฺมํ ปฏิจฺจ มคฺคารมฺมโณ นเหตุ จ มคฺคาธิปติ นเหตุ จ ธมฺมา อุปฺปชฺชนฺติ เหตุปจฺจยา. (3)}

\prose{มคฺคเหตุกํ นเหตุํ ธมฺมํ ปฏิจฺจ มคฺคเหตุโก นเหตุ ธมฺโม อุปฺปชฺชติ เหตุปจฺจยา.\csromanspacer{}\csromanspacer{}มคฺคเหตุกํ นเหตุํ ธมฺมํ ปฏิจฺจ มคฺคาธิปติ นเหตุ ธมฺโม อุปฺปชฺชติ เหตุปจฺจยา.\csromanspacer{}\csromanspacer{}มคฺคเหตุกํ นเหตุํ ธมฺมํ ปฏิจฺจ มคฺคเหตุโก นเหตุ จ มคฺคาธิปติ นเหตุ จ ธมฺมา อุปฺปชฺชนฺติ เหตุปจฺจยา. (3)}

\prose{มคฺคาธิปติํ นเหตุํ ธมฺมํ ปฏิจฺจ มคฺคาธิปติ นเหตุ ธมฺโม อุปฺปชฺชติ เหตุปจฺจยา.\csromanspacer{}\csromanspacer{}มคฺคาธิปติํ นเหตุํ ธมฺมํ ปฏิจฺจ มคฺคารมฺมโณ นเหตุ ธมฺโม อุปฺปชฺชติ เหตุปจฺจยา.\csromanspacer{}\csromanspacer{}มคฺคาธิปติํ นเหตุํ ธมฺมํ ปฏิจฺจ มคฺคเหตุโก นเหตุ ธมฺโม อุปฺปชฺชติ เหตุปจฺจยา.\csromanspacer{}\csromanspacer{}มคฺคาธิปติํ นเหตุํ ธมฺมํ ปฏิจฺจ มคฺคารมฺมโณ นเหตุ จ มคฺคาธิปติ นเหตุ จ ธมฺมา อุปฺปชฺชนฺติ เหตุปจฺจยา.\csromanspacer{}\csromanspacer{}มคฺคาธิปติํ นเหตุํ ธมฺมํ ปฏิจฺจ มคฺคเหตุโก นเหตุ จ มคฺคาธิปติ นเหตุ จ ธมฺมา อุปฺปชฺชนฺติ เหตุปจฺจยา. (5)}

\prose{มคฺคารมฺมณํ นเหตุญฺจ มคฺคาธิปติํ นเหตุญฺจ ธมฺมํ ปฏิจฺจ มคฺคารมฺมโณ นเหตุ ธมฺโม อุปฺปชฺชติ เหตุปจฺจยา.\csromanspacer{}\csromanspacer{}มคฺคารมฺมณํ นเหตุญฺจ มคฺคาธิปติํ นเหตุญฺจ ธมฺมํ ปฏิจฺจ มคฺคาธิปติ นเหตุ ธมฺโม อุปฺปชฺชติ เหตุปจฺจยา.\csromanspacer{} มคฺคารมฺมณํ นเหตุญฺจ มคฺคาธิปติํ นเหตุญฺจ ธมฺมํ ปฏิจฺจ มคฺคารมฺมโณ นเหตุ จ มคฺคาธิปติ นเหตุ จ ธมฺมา อุปฺปชฺชนฺติ เหตุปจยา. (3)}

\prose{มคฺคเหตุกํ นเหตุญฺจ มคฺคาธิปติํ นเหตุญฺจ ธมฺมํ ปฏิจฺจ มคฺคเหตุโก นเหตุ ธมฺโม อุปฺปชฺชติ เหตุปจฺจยา.\csromanspacer{}\csromanspacer{}มคฺคเหตุกํ นเหตุญฺจ มคฺคาธิปติํ นเหตุญฺจ ธมฺมํ ปฏิจฺจ มคฺคาธิปติ นเหตุ ธมฺโม อุปฺปชฺชติ เหตุปจฺจยา.\csromanspacer{} มคฺคเหตุกํ นเหตุญฺจ มคฺคาธิปติํ นเหตุญฺจ ธมฺมํ ปฏิจฺจ มคฺคเหตุโก นเหตุ จ มคฺคาธิปติ นเหตุ จ ธมฺมา อุปฺปชฺชนฺติ เหตุปจฺจยา. (3) (สํขิตฺตํ.)}

\gambhira{ธมฺมานุโลม ติกทุกปฏฺฐานปาฬิ}
\proseitem{198}{\csromanfolio{538}เหตุยา สตฺตรส, อารมฺมเณ สตฺตรส ฯเปฯ อวิคเต สตฺตรส. (สํขิตฺตํ.)}

\csromanheader{2}{\textbf{ปจฺจนีย-นเหตุ}}
\proseitem{199}{มคฺคารมฺมณํ นเหตุํ ธมฺมํ ปฏิจฺจ มคฺคารมฺมโณ นเหตุ ธมฺโม อุปฺปชฺชติ นเหตุปจฺจยา. (1) (สํขิตฺตํ.)}

\prose{นเหตุยา เอกํ, น-อธิปติยา สตฺตรส, นปุเรชาเต สตฺตรส, นปจฺฉาชาเต สตฺตรส, น-อาเสวเน นว, นวิปาเก สตฺตรส, นวิปฺปยุตฺเต สตฺตรส. (สํขิตฺตํ.)}

\csromanheader{2}{เหตุปจฺจยา น-อธิปติยา สตฺตรส. (สํขิตฺตํ.)}
\csromanheader{2}{นเหตุปจฺจยา อารมฺมเณ เอกํ. (สํขิตฺตํ.)}
\prose{(สหชาตวารมฺปิ ฯเปฯ สมฺปยุตฺตวารมฺปิ ปฏิจฺจวารสทิสํ วิตฺถาเรตพฺพํ.)}

\csromanheader{2}{\textbf{นเหตุปท 7. ปญฺหาวาร}}
\csromanheader{2}{\textbf{ปจฺจยจตุกฺก-อารมฺมณาทิ}}
\proseitem{200}{มคฺคเหตุโก นเหตุ ธมฺโม มคฺคารมฺมณสฺส นเหตุสฺส ธมฺมสฺส อารมฺมณปจฺจเยน ปจฺจโย.\csromanspacer{}\csromanspacer{}มคฺคเหตุโก นเหตุ ธมฺโม มคฺคาธิปติสฺส นเหตุสฺส ธมฺมสฺส อารมฺมณปจฺจเยน ปจฺจโย.\csromanspacer{}\csromanspacer{}มคฺคเหตุโก นเหตุ ธมฺโม มคฺคารมฺมณสฺส นเหตุสฺส จ มคฺคาธิปติสฺส นเหตุสฺส จ ธมฺมสฺส อารมฺมณปจฺจเยน ปจฺจโย. (3)}

\prose{มคฺคาธิปติ นเหตุ ธมฺโม มคฺคาธิปติสฺส นเหตุสฺส ธมฺมสฺส อารมฺมณปจฺจเยน ปจฺจโย.\csromanspacer{}\csromanspacer{}มคฺคาธิปติ นเหตุ ธมฺโม มคฺคารมฺมณสฺส นเหตุสฺส ธมฺมสฺส อารมฺมณปจฺจเยน ปจฺจโย.\csromanspacer{}\csromanspacer{}มคฺคาธิปติ นเหตุ ธมฺโม มคฺคารมฺมณสฺส นเหตุสฺส จ มคฺคาธิปติสฺส นเหตุสฺส จ ธมฺมสฺส อารมฺมณปจฺจเยน ปจฺจโย. (3)}

\prose{มคฺคเหตุโก นเหตุ จ มคฺคาธิปติ นเหตุ จ ธมฺมา มคฺคารมฺมณสฺส นเหตุสฺส ธมฺมสฺส อารมฺมณปจฺจเยน ปจฺจโย.\csromanspacer{} ตีณิ.}

\proseitem{201}{มคฺคารมฺมโณ นเหตุ ธมฺโม มคฺคารมฺมณสฺส นเหตุสฺส ธมฺมสฺส อธิปติปจฺจเยน ปจฺจโย.\csromanspacer{} ตีณิ.}

\csromanheader{2}{16. มคฺคารมฺมณตฺติก 1. เหตุทุก}
\prose{\csromanfolio{539}มคฺคเหตุโก นเหตุ ธมฺโม มคฺคเหตุกสฺส นเหตุสฺส ธมฺมสฺส อธิปติปจฺจเยน ปจฺจโย.\csromanspacer{}\csromanspacer{}มคฺคเหตุโก นเหตุ ธมฺโม มคฺคารมฺมณสฺส นเหตุสฺส ธมฺมสฺส อธิปติปจฺจเยน ปจฺจโย.\csromanspacer{}\csromanspacer{}มคฺคเหตุโก นเหตุ ธมฺโม มคฺคาธิปติสฺส นเหตุสฺส ธมฺมสฺส อธิปติปจฺจเยน ปจฺจโย.\csromanspacer{}\csromanspacer{}มคฺคเหตุโก นเหตุ ธมฺโม มคฺคารมฺมณสฺส นเหตุสฺส จ มคฺคาธิปติสฺส นเหตุสฺส จ ธมฺมสฺส อธิปติปจฺจเยน ปจฺจโย.\csromanspacer{}\csromanspacer{}มคฺคเหตุโก นเหตุ ธมฺโม มคฺคเหตุกสฺส นเหตุสฺส จ มคฺคาธิปติสฺส นเหตุสฺส จ ธมฺมสฺส อธิปติปจฺจเยน ปจฺจโย. (5)}

\prose{มคฺคาธิปติ นเหตุ ธมฺโม มคฺคาธิปติสฺส นเหตุสฺส ธมฺมสฺส อธิปติปจฺจเยน ปจฺจโย.\csromanspacer{} ปญฺจ.}

\prose{มคฺคารมฺมโณ นเหตุ จ มคฺคาธิปติ นเหตุ จ ธมฺมา มคฺคารมฺมณสฺส นเหตุสฺส ธมฺมสฺส อธิปติปจฺจเยน ปจฺจโย.\csromanspacer{} ตีณิ.}

\prose{มคฺคเหตุโก นเหตุ จ มคฺคาธิปติ นเหตุ จ ธมฺมา มคฺคเหตุกสฺส นเหตุสฺส ธมฺมสฺส อธิปติปจฺจเยน ปจฺจโย.\csromanspacer{} ปญฺจ.}

\csromanheader{2}{\textbf{อนนฺตร อุปนิสฺสย}}
\proseitem{202}{มคฺคารมฺมโณ นเหตุ ธมฺโม มคฺคารมฺมณสฺส นเหตุสฺส ธมฺมสฺส อนนฺตรปจฺจเยน ปจฺจโย.\csromanspacer{}\csromanspacer{}มคฺคารมฺมโณ นเหตุ ธมฺโม มคฺคาธิปติสฺส นเหตุสฺส ธมฺมสฺส อนนฺตรปจฺจเยน ปจฺจโย.\csromanspacer{}\csromanspacer{}มคฺคารมฺมโณ นเหตุ ธมฺโม มคฺคารมฺมณสฺส นเหตุสฺส จ มคฺคาธิปติสฺส นเหตุสฺส จ ธมฺมสฺส อนนฺตรปจฺจเยน ปจฺจโย. (3)}

\prose{มคฺคาธิปติ นเหตุ ธมฺโม มคฺคาธิปติสฺส นเหตุสฺส ธมฺมสฺส อนนฺตรปจฺจเยน ปจฺจโย.\csromanspacer{}\csromanspacer{}มคฺคาธิปติ นเหตุ ธมฺโม มคฺคารมฺมณสฺส นเหตุสฺส ธมฺมสฺส อนนฺตรปจฺจเยน ปจฺจโย.\csromanspacer{}\csromanspacer{}มคฺคาธิปติ นเหตุ ธมฺโม มคฺคารมฺมณสฺส นเหตุสฺส จ มคฺคาธิปติสฺส นเหตุสฺส จ ธมฺมสฺส อนนฺตรปจฺจเยน ปจฺจโย. (3)}

\prose{มคฺคารมฺมโณ นเหตุ จ มคฺคาธิปติ นเหตุ จ ธมฺมา มคฺคารมฺมณสฺส นเหตุสฺส ธมฺมสฺส อนนฺตรปจฺจเยน ปจฺจโย.\csromanspacer{} ตีณิ ฯเปฯ.}

\proseitem{203}{มคฺคารมฺมโณ นเหตุ ธมฺโม มคฺคารมฺมณสฺส นเหตุสฺส ธมฺมสฺส อุปนิสฺสยปจฺจเยน ปจฺจโย.\csromanspacer{} ตีณิ.}

\gambhira{ธมฺมานุโลม ติกทุกปฏฺฐานปาฬิ}
\prose{\csromanfolio{540}มคฺคเหตุโก นเหตุ ธมฺโม มคฺคเนตุกสฺส นเหตุสฺส ธมฺมสฺส อุปนิสฺสยปจฺจเยน ปจฺจโย.\csromanspacer{} ปญฺจ.}

\prose{มคฺคาธิปติ นเหตุ ธมฺโม มคฺคาธิปติสฺส นเหตุสฺส ธมฺมสฺส อุปนิสฺสยปจฺจเยน ปจฺจโย.\csromanspacer{} ปญฺจ.}

\prose{มคฺคารมฺมโณ นเหตุ จ มคฺคาธิปติ นเหตุ จ ธมฺมา มคฺคารมฺมณสฺส นเหตุสฺส ธมฺมสฺส อุปนิสฺสยปจฺจเยน ปจฺจโย.\csromanspacer{} ตีณิ.}

\prose{มคฺคเหตุโก นเหตุ จ มคฺคาธิปติ นเหตุ จ ธมฺมา มคฺคเหตุกสฺส นเหตุสฺส ธมฺมสฺส อุปนิสฺสยปจฺจเยน ปจฺจโย.\csromanspacer{} ปญฺจ. (สํขิตฺตํ.)}

\proseitem{204}{อารมฺมเณ นว, อธิปติยา เอกวีส, อนนฺตเร นว, สมนนฺตเร นว, สหชาเต สตฺตรส, อญฺญมญฺเญ สตฺตรส, นิสฺสเย สตฺตรส, อุปนิสฺสเย เอกวีส ฯเปฯ อวิคเต สตฺตรส. (สํขิตฺตํ.)}

\prose{นเหตุยา เอกวีส, น-อารมฺมเณ สตฺตรส. (สํขิตฺตํ.)}

\csromanheader{2}{อารมฺมณปจฺจยา นเหตุยา นว. (สํขิตฺตํ.)}
\csromanheader{2}{นเหตุปจฺจยา อารมฺมเณ นว. (สํขิตฺตํ.)}
\prose{(ยถา กุสลตฺติเก ปญฺหาวารํ, เอวํ วิตฺถาเรตพฺพํ.)}

\csromanmatikaanchor{mtk.124}
\csromanmatikaanchor{mtk.125}
\csromantocmarknopage{subparagraph}{17. อุปฺปนฺนตฺติก}{mtk.124}
\bookmark[dest={mtk.124},level=5]{17. อุปฺปนฺนตฺติก}
\csromantocmarkref{subparagraph}{1. เหตุทุก}{mtk.125}
\bookmark[dest={mtk.125},level=5]{1. เหตุทุก}
\csromanheader{6}{17. อุปฺปนฺนตฺติก 1. เหตุทุก}
\csromanheader{2}{7. ปญฺหาวาร ปจฺจยจตุกฺก-เหตฺวาทิ}
\proseitem{205}{อุปฺปนฺโน เหตุ ธมฺโม อุปฺปนฺนสฺส เหตุสฺส ธมฺมสฺส เหตุปจฺจเยน ปจฺจโย.\csromanspacer{} อุปฺปนฺนา เหตู สมฺปยุตฺตกานํ เหตูนํ เหตุปจฺจเยน ปจฺจโย.\csromanspacer{} ปฏิสนฺธิกฺขเณ ฯเปฯ. (1)}

\prose{อนุปฺปนฺโน เหตุ ธมฺโม อุปฺปนฺนสฺส เหตุสฺส ธมฺมสฺส อารมฺมณปจฺจเยน ปจฺจโย.\csromanspacer{} อนุปฺปนฺนา เหตู เจโตปริยญาณสฺส, อนาคตํสญาณสฺส อารมฺมณปจฺจเยน ปจฺจโย.\csromanspacer{}\csromanspacer{}อุปฺปาที เหตุ ธมฺโม อุปฺปนฺนสฺส เหตุสฺส ธมฺมสฺส อารมฺมณปจฺจเยน ปจฺจโย.\csromanspacer{} อุปฺปาที เหตู เจโตปริยญาณสฺส อนาคตํสญาณสฺส อารมฺมณปจฺจเยน ปจฺจโย. (2)}

\csromanheader{2}{17. อุปฺปนฺนตฺติก 1. เหตุทุก}
\prose{\csromanfolio{541}อุปฺปนฺโน เหตุ ธมฺโม อุปฺปนฺนสฺส เหตุสฺส ธมฺมสฺส อธิปติปจฺจเยน ปจฺจโย.\csromanspacer{} ตีณิ.}

\prose{สหชาตปจฺจเยน ปจฺจโย.\csromanspacer{} อญฺญมญฺญปจฺจเยน ปจฺจโย.\csromanspacer{} นิสฺสยปจฺจเยน ปจฺจโย.\csromanspacer{} อุปนิสฺสยปจฺจเยน ปจฺจโย.\csromanspacer{} วิปากปจฺจเยน ปจฺจโย.\csromanspacer{} อินฺทฺริยปจฺจเยน ปจฺจโย.\csromanspacer{} มคฺคปจฺจเยน ปจฺจโย.\csromanspacer{} สมฺปยุตฺตปจฺจเยน ปจฺจโย.\csromanspacer{} อตฺถิปจฺจเยน ปจฺจโย ฯเปฯ อวิคตปจฺจเยน ปจฺจโย.}

\proseitem{206}{เหตุยา เอกํ, อารมฺมเณ เทฺว, อธิปติยา ตีณิ, สหชาเต เอกํ, อญฺญมญฺเญ เอกํ, นิสฺสเย เอกํ, อุปนิสฺสเย เทฺว ฯเปฯ อวิคเต เอกํ. (สํขิตฺตํ.)}

\prose{นเหตุยา เทฺว, น-อารมฺมเณ เทฺว. (สํขิตฺตํ.)}

\csromanheader{2}{เหตุปจฺจยา น-อารมฺมเณ เอกํ. (สํขิตฺตํ.)}
\csromanheader{2}{นเหตุปจฺจยา อารมฺมเณ เทฺว. (สํขิตฺตํ.)}
\prose{(ยถา กุสลตฺติเก ปญฺหาวารํ, เอวํ วิตฺถาเรตพฺพํ.)}

\csromanheader{2}{\textbf{นเหตุปท-อารมฺมณาทิ}}
\proseitem{207}{อุปฺปนฺโน นเหตุ ธมฺโม อุปฺปนฺนสฺส นเหตุสฺส ธมฺมสฺส อารมฺมณปจฺจเยน ปจฺจโย. (1)}

\prose{อนุปฺปนฺโน นเหตุ ธมฺโม อุปฺปนฺนสฺส นเหตุสฺส ธมฺมสฺส อารมฺมณปจฺจเยน ปจฺจโย. (1)}

\prose{อุปฺปาที นเหตุ ธมฺโม อุปฺปนฺนสฺส นเหตุสฺส ธมฺมสฺส อารมฺมณปจฺจเยน ปจฺจโย. (1)}

\proseitem{208}{อุปฺปนฺโน นเหตุ ธมฺโม อุปฺปนฺนสฺส นเหตุสฺส ธมฺมสฺส อธิปติปจฺจเยน ปจฺจโย. (1)}

\prose{อนุปฺปนฺโน นเหตุ ธมฺโม อุปฺปนฺนสฺส นเหตุสฺส ธมฺมสฺส อธิปติปจฺจเยน ปจฺจโย. (1)}

\prose{อุปฺปาที นเหตุ ธมฺโม อุปฺปนฺนสฺส นเหตุสฺส ธมฺมสฺส อธิปติปจฺจเยน ปจฺจโย. (1)}

\gambhira{ธมฺมานุโลม ติกทุกปฏฺฐานปาฬิ}
\prose{\csromanfolio{542}อุปฺปนฺโน นเหตุ ธมฺโม อุปฺปนฺนสฺส นเหตุสฺส ธมฺมสฺส สหชาตปจฺจเยน ปจฺจโย.\csromanspacer{} อญฺญมญฺญปจฺจเยน ปจฺจโย.\csromanspacer{} นิสฺสยปจฺจเยน ปจฺจโย.\csromanspacer{} อุปนิสฺสยปจฺจเยน ปจฺจโย.}

\prose{อนุปฺปนฺโน นเหตุ ธมฺโม อุปฺปนฺนสฺส นเหตุสฺส ธมฺมสฺส อุปนิสฺสยปจฺจเยน ปจฺจโย.}

\prose{อุปฺปาที นเหตุ ธมฺโม อุปฺปนฺนสฺส นเหตุสฺส ธมฺมสฺส อุปนิสฺสยปจฺจเยน ปจฺจโย. (สํขิตฺตํ.)}

\proseitem{209}{อารมฺมเณ ตีณิ, อธิปติยา ตีณิ, สหชาเต เอกํ, อญฺญมญฺเญ เอกํ, นิสฺสเย เอกํ, อุปนิสฺสเย ตีณิ ฯเปฯ อวิคเต เอกํ. (สํขิตฺตํ.)}

\prose{นเหตุยา ตีณิ, น-อารมฺมเณ ตีณิ. (สํขิตฺตํ.)}

\csromanheader{2}{อารมฺมณปจฺจยา นเหตุยา ตีณิ. (สํขิตฺตํ.)}
\csromanheader{2}{นเหตุปจฺจยา อารมฺมเณ ตีณิ. (สํขิตฺตํ.)}
\csromanmatikaanchor{mtk.126}
\csromanmatikaanchor{mtk.127}
\csromantocmarknopage{subparagraph}{18. อตีตตฺติก}{mtk.126}
\bookmark[dest={mtk.126},level=5]{18. อตีตตฺติก}
\csromantocmarkref{subparagraph}{1. เหตุทุก}{mtk.127}
\bookmark[dest={mtk.127},level=5]{1. เหตุทุก}
\csromanheader{6}{18. อตีตตฺติก 1. เหตุทุก}
\csromanheader{2}{7. ปญฺหาวาร ปจฺจยจตุกฺก}
\proseitem{210}{ปจฺจุปฺปนฺโน เหตุ ธมฺโม ปจฺจุปฺปนฺนสฺส เหตุสฺส ธมฺมสฺส เหตุปจฺจเยน ปจฺจโย. (1)}

\prose{อตีโต เหตุ ธมฺโม ปจฺจุปฺปนฺนสฺส เหตุสฺส ธมฺมสฺส อารมฺมณปจฺจเยน ปจฺจโย.\csromanspacer{}\csromanspacer{}อนาคโต เหตุ ธมฺโม ปจฺจุปฺปนฺนสฺส เหตุสฺส ธมฺมสฺส อารมฺมณปจฺจเยน ปจฺจโย. (2) (สํขิตฺตํ.)}

\proseitem{211}{เหตุยา เอกํ, อารมฺมเณ เทฺว, อธิปติยา ตีณิ ฯเปฯ สหชาเต เอกํ, อญฺญมญฺเญ เอกํ, นิสฺสเย เอกํ, อุปนิสฺสเย เทฺว ฯเปฯ วิปาเก เอกํ, อินฺทฺริเย เอกํ, มคฺเค เอกํ, สมฺปยุตฺเต เอกํ, อตฺถิยา เอกํ ฯเปฯ อวิคเต เอกํ. (สํขิตฺตํ.)}

\prose{นเหตุยา เทฺว, น-อารมฺมเณ เทฺว. (สํขิตฺตํ.)}

\csromanheader{2}{เหตุปจฺจยา น-อารมฺมเณ เอกํ. (สํขิตฺตํ.)}
\csromanheader{2}{นเหตุปจฺจยา อารมฺมเณ เทฺว. (สํขิตฺตํ.)}
\csromanheader{2}{19. อตีตารมฺมณตฺติก 1. เหตุทุก \textbf{นเหตุปท-อารมฺมณ}}
\proseitem{212}{\csromanfolio{543}อตีโต นเหตุ ธมฺโม ปจฺจุปฺปนฺนสฺส นเหตุสฺส ธมฺมสฺส อารมฺมณปจฺจเยน ปจฺจโย. (1)}

\prose{อนาคโต นเหตุ ธมฺโม ปจฺจุปฺปนฺนสฺส นเหตุสฺส ธมฺมสฺส อารมฺมณปจฺจเยน ปจฺจโย. (1)}

\prose{ปจฺจุปฺปนฺโน นเหตุ ธมฺโม ปจฺจุปฺปนฺนสฺส นเหตุสฺส ธมฺมสฺส อารมฺมณปจฺจเยน ปจฺจโย. (1) (สํขิตฺตํ.)}

\proseitem{213}{อารมฺมเณ ตีณิ, อธิปติยา ตีณิ ฯเปฯ สหชาเต เอกํ, อญฺญมญฺเญ เอกํ, นิสฺสเย เอกํ, อุปนิสฺสเย ตีณิ ฯเปฯ อาเสวเน เอกํ, กมฺเม เทฺว, วิปาเก เอกํ ฯเปฯ อินฺทฺริเย เอกํ, มคฺเค เอกํ, สมฺปยุตฺเต เอกํ, อตฺถิยา เอกํ ฯเปฯ อวิคเต เอกํ. (สํขิตฺตํ.)}

\prose{นเหตุยา ตีณิ, น-อารมฺมเณ ตีณิ. (สํขิตฺตํ.)}

\csromanheader{2}{อารมฺมณปจฺจยา นเหตุยา ตีณิ. (สํขิตฺตํ.)}
\csromanheader{2}{นเหตุปจฺจยา อารมฺมเณ ตีณิ. (สํขิตฺตํ.)}
\csromanmatikaanchor{mtk.128}
\csromanmatikaanchor{mtk.129}
\csromantocmarknopage{subparagraph}{19. อตีตารมฺมณตฺติก}{mtk.128}
\bookmark[dest={mtk.128},level=5]{19. อตีตารมฺมณตฺติก}
\csromantocmarkref{subparagraph}{1. เหตุทุก}{mtk.129}
\bookmark[dest={mtk.129},level=5]{1. เหตุทุก}
\csromanheader{6}{19. อตีตารมฺมณตฺติก 1. เหตุทุก}
\csromanheader{2}{1-7. ปฏิจฺจาทิวาร ปจฺจยจตุกฺก-เหตุ}
\proseitem{214}{อตีตารมฺมณํ เหตุํ ธมฺมํ ปฏิจฺจ อตีตารมฺมโณ เหตุ ธมฺโม อุปฺปชฺชติ เหตุปจฺจยา. (1)}

\prose{อนาคตารมฺมณํ เหตุํ ธมฺมํ ปฏิจฺจ อนาคตารมฺมโณ เหตุ ธมฺโม อุปฺปชฺชติ เหตุปจฺจยา. (1)}

\prose{ปจฺจุปฺปนฺนารมฺมณํ เหตุํ ธมฺมํ ปฏิจฺจ ปจฺจุปฺปนฺนารมฺมโณ เหตุ ธมฺโม อุปฺปชฺชติ เหตุปจฺจยา. (1) (สํขิตฺตํ.)}

\proseitem{215}{เหตุยา ตีณิ, อารมฺมเณ ตีณิ ฯเปฯ อวิคเต ตีณิ. (สํขิตฺตํ.)}

\prose{น-อธิปติยา ตีณิ, นปุเรชาเต ตีณิ, นปจฺฉาชาเต ตีณิ, น-อาเสวเน ตีณิ, นวิปาเก ตีณิ, นวิปฺปยุตฺเต เทฺว. (สํขิตฺตํ.)}

\prose{(สหชาตวารมฺปิ ฯเปฯ สมฺปยุตฺตวารมฺปิ ปฏิจฺจวารสทิสํ วิตฺถาเรตพฺพํ.)}

\csromanheader{2}{ธมฺมานุโลม ติกทุกปฏฺฐานปาฬิ \textbf{เหตุ-อารมฺมณ}}
\proseitem{216}{\csromanfolio{544}อตีตารมฺมโณ เหตุ ธมฺโม อตีตารมฺมณสฺส เหตุสฺส ธมฺมสฺส เหตุปจฺจเยน ปจฺจโย. (1)}

\prose{อนาคตารมฺมโณ เหตุ ธมฺโม อนาคตารมฺมณสฺส เหตุสฺส ธมฺมสฺส เหตุปจฺจเยน ปจฺจโย. (1)}

\prose{ปจฺจุปฺปนฺนารมฺมโณ เหตุ ธมฺโม ปจฺจุปฺปนฺนารมฺมณสฺส เหตุสฺส ธมฺมสฺส เหตุปจฺจเยน ปจฺจโย. (1)}

\prose{อตีตารมฺมโณ เหตุ ธมฺโม อตีตารมฺมณสฺส เหตุสฺส ธมฺมสฺส อารมฺมณปจฺจเยน ปจฺจโย.\csromanspacer{}\csromanspacer{}อตีตารมฺมโณ เหตุ ธมฺโม อนาคตารมฺมณสฺส เหตุสฺส ธมฺมสฺส อารมฺมณปจฺจเยน ปจฺจโย.\csromanspacer{}\csromanspacer{}อตีตารมฺมโณ เหตุ ธมฺโม ปจฺจุปฺปนฺนารมฺมณสฺส เหตุสฺส ธมฺมสฺส อารมฺมณปจฺจเยน ปจฺจโย. (3)}

\prose{อนาคตารมฺมโณ เหตุ ธมฺโม อนาคตารมฺมณสฺส เหตุสฺส ธมฺมสฺส อารมฺมณปจฺจเยน ปจฺจโย.\csromanspacer{} ตีณิ.}

\prose{ปจฺจุปฺปนฺนารมฺมโณ เหตุ ธมฺโม ปจฺจุปฺปนฺนารมฺมณสฺส เหตุสฺส ธมฺมสฺส อารมฺมณปจฺจเยน ปจฺจโย.\csromanspacer{} ตีณิ. (สํขิตฺตํ.)}

\proseitem{217}{เหตุยา ตีณิ, อารมฺมเณ นว, อธิปติยา สตฺต, อนนฺตเร ฉ, สมนนฺตเร ฉ, สหชาเต ตีณิ, อญฺญมญฺเญ ตีณิ, นิสฺสเย ตีณิ, อุปนิสฺสเย นว ฯเปฯ อวิคเต ตีณิ. (สํขิตฺตํ.)}

\prose{นเหตุยา นว, น-อารมฺมเณ นว. (สํขิตฺตํ.)}

\csromanheader{2}{เหตุปจฺจยา น-อารมฺมเณ ตีณิ. (สํขิตฺตํ.)}
\csromanheader{2}{นเหตุปจฺจยา อารมฺมเณ นว. (สํขิตฺตํ.)}
\prose{(ยถา กุสลตฺติเก ปญฺหาวารํ, เอวํ วิตฺถาเรตพฺพํ.)}

\csromanheader{2}{\textbf{นเหตุปท 1-7. ปฏิจฺจาทิวาร}}
\csromanheader{2}{\textbf{ปจฺจยจตุกฺก-เหตุ}}
\proseitem{218}{อตีตารมฺมณํ นเหตุํ ธมฺมํ ปฏิจฺจ อตีตารมฺมโณ นเหตุ ธมฺโม อุปฺปชฺชติ เหตุปจฺจยา. (1)}

\csromanheader{2}{19. อตีตารมฺมณตฺติก 1. เหตุทุก}
\prose{\csromanfolio{545}อนาคตารมฺมณํ นเหตุํ ธมฺมํ ปฏิจฺจ อนาคตารมฺมโณ นเหตุ ธมฺโม อุปฺปชฺชติ เหตุปจฺจยา. (1)}

\prose{ปจฺจุปฺปนฺนารมฺมณํ นเหตุํ ธมฺมํ ปฏิจฺจ ปจฺจุปฺปนฺนารมฺมโณ นเหตุ ธมฺโม อุปฺปชฺชติ เหตุปจฺจยา. (1) (สํขิตฺตํ.)}

\proseitem{219}{เหตุยา ตีณิ, อารมฺมเณ ตีณิ, อธิปติยา ตีณิ ฯเปฯ อวิคเต ตีณิ. (สํขิตฺตํ.)}

\prose{นเหตุยา ตีณิ, น-อธิปติยา ตีณิ, นปุเรชาเต ตีณิ, นปจฺฉาชาเต ตีณิ, น-อาเสวเน ตีณิ, นกมฺเม ตีณิ, นวิปาเก ตีณิ, นฌาเน เอกํ, นมคฺเค ตีณิ, นวิปฺปยุตฺเต เทฺว. (สํขิตฺตํ.)}

\prose{(สหชาตวารมฺปิ ฯเปฯ สมฺปยุตฺตวารมฺปิ ปฏิจฺจวารสทิสํ วิตฺถาเรตพฺพํ.)}

\csromanheader{2}{\textbf{อารมฺมณ}}
\proseitem{220}{อตีตารมฺมโณ นเหตุ ธมฺโม อตีตารมฺมณสฺส นเหตุสฺส ธมฺมสฺส อารมฺมณปจฺจเยน ปจฺจโย.\csromanspacer{} ตีณิ.}

\prose{อนาคตารมฺมโณ นเหตุ ธมฺโม อนาคตารมฺมณสฺส นเหตุสฺส ธมฺมสฺส อารมฺมณปจฺจเยน ปจฺจโย.\csromanspacer{} ตีณิ.}

\prose{ปจฺจุปฺปนฺนารมฺมโณ นเหตุ ธมฺโม ปจฺจุปฺปนฺนารมฺมณสฺส นเหตุสฺส ธมฺมสฺส อารมฺมณปจฺจเยน ปจฺจโย.\csromanspacer{} ตีณิ. (สํขิตฺตํ.)}

\proseitem{221}{อารมฺมเณ นว, อธิปติยา สตฺต ฯเปฯ อวิคเต ตีณิ. (สํขิตฺตํ.)}

\prose{นเหตุยา นว, น-อารมฺมเณ นว. (สํขิตฺตํ.)}

\csromanheader{2}{\textbf{อารมฺมณ}ปจฺจยา นเหตุยา นว. (สํขิตฺตํ.)}
\csromanheader{2}{นเหตุปจฺจยา อารมฺมเณ นว. (สํขิตฺตํ.)}
\prose{(ยถา กุสลตฺติเก ปญฺหาวารํ, เอวํ วิตฺถาเรตพฺพํ.)}

\csromanmatikaanchor{mtk.130}
\csromanmatikaanchor{mtk.131}
\csromantocmarknopage{subparagraph}{20. อชฺฌตฺตตฺติก}{mtk.130}
\bookmark[dest={mtk.130},level=5]{20. อชฺฌตฺตตฺติก}
\csromantocmarkref{subparagraph}{1. เหตุทุก}{mtk.131}
\bookmark[dest={mtk.131},level=5]{1. เหตุทุก}
\csromanheader{6}{ธมฺมานุโลม ติกทุกปฏฺฐานปาฬิ \textbf{20. อชฺฌตฺตตฺติก 1. เหตุทุก}}
\csromanheader{2}{1-7. ปฏิจฺจาทิวาร ปจฺจยจตุกฺก-เหตุ}
\proseitem{222}{\csromanfolio{546}อชฺฌตฺตํ เหตุํ ธมฺมํ ปฏิจฺจ อชฺฌตฺโต เหตุ ธมฺโม อุปฺปชฺชติ เหตุปจฺจยา. (1)}

\prose{พหิทฺธา เหตุํ ธมฺมํ ปฏิจฺจ พหิทฺธา เหตุ ธมฺโม อุปฺปชฺชติ เหตุปจฺจยา. (1) (สํขิตฺตํ.)}

\proseitem{223}{เหตุยา เทฺว, อารมฺมเณ เทฺว ฯเปฯ อวิคเต เทฺว. (สํขิตฺตํ.)}

\prose{น-อธิปติยา เทฺว, นปุเรชาเต เทฺว, นปจฺฉาชาเต เทฺว, น-อาเสวเน เทฺว, นวิปาเก เทฺว, นวิปฺปยุตฺเต เทฺว. (สํขิตฺตํ.)}

\prose{(สหชาตวารมฺปิ ฯเปฯ สมฺปยุตฺตวารมฺปิ ปฏิจฺจวารสทิสํ วิตฺถาเรตพฺพํ.)}

\csromanheader{2}{\textbf{เหตุ อารมฺมณาทิ}}
\proseitem{224}{อชฺฌตฺโต เหตุ ธมฺโม อชฺฌตฺตสฺส เหตุสฺส ธมฺมสฺส เหตุปจฺจเยน ปจฺจโย. (1)}

\prose{พหิทฺธา เหตุ ธมฺโม พหิทฺธา เหตุสฺส ธมฺมสฺส เหตุปจฺจเยน ปจฺจโย.}

\prose{อชฺฌตฺโต เหตุ ธมฺโม อชฺฌตฺตสฺส เหตุสฺส ธมฺมสฺส อารมฺมณปจฺจเยน ปจฺจโย.\csromanspacer{}\csromanspacer{}อชฺฌตฺโต เหตุ ธมฺโม พหิทฺธา เหตุสฺส ธมฺมสฺส อารมฺมณปจฺจเยน ปจฺจโย. (2)}

\prose{พหิทฺธา เหตุ ธมฺโม พหิทฺธา เหตุสฺส ธมฺมสฺส อารมฺมณปจฺจเยน ปจฺจโย.\csromanspacer{}\csromanspacer{}พหิทฺธา เหตุ ธมฺโม อชฺฌตฺตสฺส เหตุสฺส ธมฺมสฺส อารมฺมณปจฺจเยน ปจฺจโย. (2)}

\proseitem{225}{อชฺฌตฺโต เหตุ ธมฺโม อชฺฌตฺตสฺส เหตุสฺส ธมฺมสฺส อธิปติปจฺจเยน ปจฺจโย. (1)}

\prose{พหิทฺธา เหตุ ธมฺโม พหิทฺธา เหตุสฺส ธมฺมสฺส อธิปติปจฺจเยน ปจฺจโย. (1)}

\prose{อชฺฌตฺโต เหตุ ธมฺโม อชฺฌตฺตสฺส เหตุสฺส ธมฺมสฺส อนนฺตรปจฺจเยน ปจฺจโย. (1)}

\csromanheader{2}{20. อชฺฌตฺตตฺติก 1. เหตุทุก}
\prose{\csromanfolio{547}พหิทฺธา เหตุ ธมฺโม พหิทฺธา เหตุสฺส ธมฺมสฺส อนนฺตรปจฺจเยน ปจฺจโย. (1) (สํขิตฺตํ.)}

\proseitem{226}{เหตุยา เทฺว, อารมฺมเณ จตฺตาริ, อธิปติยา เทฺว, อนนฺตเร เทฺว, สมนนฺตเร เทฺว, สหชาเต เทฺว, อญฺญมญฺเญ เทฺว, นิสฺสเย เทฺว, อุปนิสฺสเย จตฺตาริ ฯเปฯ อวิคเต เทฺว. (สํขิตฺตํ.)}

\prose{นเหตุยา จตฺตาริ, น-อารมฺมเณ จตฺตาริ. (สํขิตฺตํ.)}

\csromanheader{2}{เหตุปจฺจยา น-อารมฺมเณ เทฺว. (สํขิตฺตํ.)}
\csromanheader{2}{นเหตุปจฺจยา อารมฺมเณ จตฺตาริ. (สํขิตฺตํ.)}
\prose{(ยถา กุสลตฺติเก ปญฺหาวารํ, เอวํ วิตฺถาเรตพฺพํ.)}

\csromanheader{2}{\textbf{นเหตุปท 1-7. ปฏิจฺจาทิวาร}}
\csromanheader{2}{\textbf{ปจฺจยจตุกฺก-เหตุ}}
\proseitem{227}{อชฺฌตฺตํ นเหตุํ ธมฺมํ ปฏิจฺจ อชฺฌตฺโต นเหตุ ธมฺโม อุปฺปชฺชติ เหตุปจฺจยา. (1)}

\prose{พหิทฺธา นเหตุํ ธมฺมํ ปฏิจฺจ พหิทฺธา นเหตุ ธมฺโม อุปฺปชฺชติ เหตุปจฺจยา. (1) (สํขิตฺตํ.)}

\proseitem{228}{เหตุยา เทฺว, อารมฺมเณ เทฺว ฯเปฯ อวิคเต เทฺว. (สํขิตฺตํ.)}

\prose{นเหตุยา เทฺว, น-อธิปติยา เทฺว ฯเปฯ นวิปฺปยุตฺเต เทฺว. (สํขิตฺตํ.)}

\prose{(สหชาตวารมฺปิ ฯเปฯ สมฺปยุตฺตวารมฺปิ ปฏิจฺจวารสทิสํ วิตฺถาเรตพฺพํ.)}

\csromanheader{2}{\textbf{อารมฺมณ}}
\proseitem{229}{อชฺฌตฺโต นเหตุ ธมฺโม อชฺฌตฺตสฺส นเหตุสฺส ธมฺมสฺส อารมฺมณปจฺจเยน ปจฺจโย.\csromanspacer{}\csromanspacer{}อชฺฌตฺโต นเหตุ ธมฺโม พหิทฺธา นเหตุสฺส ธมฺมสฺส อารมฺมณปจฺจเยน ปจฺจโย. (2)}

\prose{พหิทฺธา นเหตุ ธมฺโม พหิทฺธา นเหตุสฺส ธมฺมสฺส อารมฺมณปจฺจเยน ปจฺจโย.\csromanspacer{}\csromanspacer{}พหิทฺธา นเหตุ ธมฺโม อชฺฌตฺตสฺส นเหตุสฺส ธมฺมสฺส อารมฺมณปจฺจเยน ปจฺจโย. (2) (สํขิตฺตํ.)}

\gambhira{ธมฺมานุโลม ติกทุกปฏฺฐานปาฬิ}
\proseitem{230}{\csromanfolio{548}อารมฺมเณ จตฺตาริ, อธิปติยา จตฺตาริ, อนนฺตเร เทฺว ฯเปฯ อุปนิสฺสเย จตฺตาริ ฯเปฯ อวิคเต ฉ. (สํขิตฺตํ.)}

\prose{นเหตุยา ฉ, น-อารมฺมเณ ฉ. (สํขิตฺตํ.)}

\csromanheader{2}{อารมฺมณปจฺจยา นเหตุยา จตฺตาริ. (สํขิตฺตํ.)}
\csromanheader{2}{นเหตุปจฺจยา อารมฺมเณ จตฺตาริ. (สํขิตฺตํ.)}
\prose{(ยถา กุสลตฺติเก ปญฺหาวารํ, เอวํ วิตฺถาเรตพฺพํ.)}

\csromanmatikaanchor{mtk.132}
\csromanmatikaanchor{mtk.133}
\csromantocmarknopage{subparagraph}{21. อชฺฌตฺตารมฺมณตฺติก}{mtk.132}
\bookmark[dest={mtk.132},level=5]{21. อชฺฌตฺตารมฺมณตฺติก}
\csromantocmarkref{subparagraph}{1. เหตุทุก}{mtk.133}
\bookmark[dest={mtk.133},level=5]{1. เหตุทุก}
\csromanheader{6}{21. อชฺฌตฺตารมฺมณตฺติก 1. เหตุทุก}
\csromanheader{2}{1-7. ปฏิจฺจาทิวาร ปจฺจยจตุกฺก-เหตุ}
\proseitem{231}{อชฺฌตฺตารมฺมณํ เหตุํ ธมฺมํ ปฏิจฺจ อชฺฌตฺตารมฺมโณ เหตุ ธมฺโม อุปฺปชฺชติ เหตุปจฺจยา. (1)}

\prose{พหิทฺธารมฺมณํ เหตุํ ธมฺมํ ปฏิจฺจ พหิทฺธารมฺมโณ เหตุ ธมฺโม อุปฺปชฺชติ เหตุปจฺจยา. (1) (สํขิตฺตํ.)}

\proseitem{232}{เหตุยา เทฺว, อารมฺมเณ เทฺว ฯเปฯ อวิคเต เทฺว. (สํขิตฺตํ.)}

\prose{น-อธิปติยา เทฺว, นปุเรชาเต เทฺว, นปจฺฉาชาเต เทฺว, น-อาเสวเน เทฺว, นวิปาเก เทฺว, นวิปฺปยุตฺเต เทฺว. (สํขิตฺตํ.)}

\prose{(สหชาตวารมฺปิ ฯเปฯ สมฺปยุตฺตวารมฺปิ ปฏิจฺจวารสทิสํ วิตฺถาเรตพฺพํ.)}

\csromanheader{2}{\textbf{เหตุ-อารมฺมณ}}
\proseitem{233}{อชฺฌตฺตารมฺมโณ เหตุ ธมฺโม อชฺฌตฺตารมฺมณสฺส เหตุสฺส ธมฺมสฺส เหตุปจฺจเยน ปจฺจโย. (1)}

\prose{พหิทฺธารมฺมโณ เหตุ ธมฺโม พหิทฺธารมฺมณสฺส เหตุสฺส ธมฺมสฺส เหตุปจฺจเยน ปจฺจโย. (1)}

\csromanheader{2}{21. อชฺฌตฺตารมฺมณตฺติก 1. เหตุทุก}
\prose{\csromanfolio{549}อชฺฌตฺตารมฺมโณ เหตุ ธมฺโม อชฺฌตฺตารมฺมณสฺส เหตุสฺส ธมฺมสฺส อารมฺมณปจฺจเยน ปจฺจโย.\csromanspacer{}\csromanspacer{}อชฺฌตฺตารมฺมโณ เหตุ ธมฺโม พหิทฺธารมฺมณสฺส เหตุสฺส ธมฺมสฺส อารมฺมณปจฺจเยน ปจฺจโย. (2)}

\prose{พหิทฺธารมฺมโณ เหตุ ธมฺโม พหิทฺธารมฺมณสฺส เหตุสฺส ธมฺมสฺส อารมฺมณปจฺจเยน ปจฺจโย.\csromanspacer{}\csromanspacer{}พหิทฺธารมฺมโณ เหตุ ธมฺโม อชฺฌตฺตารมฺมณสฺส เหตุสฺส ธมฺมสฺส อารมฺมณปจฺจเยน ปจฺจโย. (2) (สํขิตฺตํ.)}

\proseitem{234}{เหตุยา เทฺว, อารมฺมเณ จตฺตาริ, อธิปติยา ตีณิ, อนนฺตเร จตฺตาริ, สมนนฺตเร จตฺตาริ, อุปนิสฺสเย จตฺตาริ ฯเปฯ อวิคเต เทฺว. (สํขิตฺตํ.)}

\prose{นเหตุยา จตฺตาริ, น-อารมฺมเณ จตฺตาริ. (สํขิตฺตํ.)}

\csromanheader{2}{เหตุปจฺจยา น-อารมฺมเณ เทฺว. (สํขิตฺตํ.)}
\csromanheader{2}{นเหตุปจฺจยา อารมฺมเณ จตฺตาริ. (สํขิตฺตํ.)}
\prose{(ยถา กุสลตฺติเก ปญฺหาวารํ, เอวํ วิตฺถาเรตพฺพํ.)}

\csromanheader{2}{\textbf{นเหตุปท 1-7. ปฏิจฺจาทิวาร}}
\csromanheader{2}{\textbf{ปจฺจยจตุกฺก-เหตุ}}
\proseitem{235}{อชฺฌตฺตารมฺมณํ นเหตุํ ธมฺมํ ปฏิจฺจ อชฺฌตฺตารมฺมโณ นเหตุ ธมฺโม อุปฺปชฺชติ เหตุปจฺจยา. (1)}

\prose{พหิทฺธารมฺมณํ นเหตุํ ธมฺมํ ปฏิจฺจ พหิทฺธารมฺมโณ นเหตุ ธมฺโม อุปฺปชฺชติ เหตุปจฺจยา. (1) (สํขิตฺตํ.)}

\prose{เหตุยา เทฺว, อารมฺมเณ เทฺว ฯเปฯ อวิคเต เทฺว. (สํขิตฺตํ.)}

\prose{(สหชาตวารมฺปิ ฯเปฯ สมฺปยุตฺตวารมฺปิ ปฏิจฺจวารสทิสํ วิตฺถาเรตพฺพํ.)}

\csromanheader{2}{\textbf{อารมฺมณ}}
\proseitem{236}{อชฺฌตฺตารมฺมโณ นเหตุ ธมฺโม อชฺฌตฺตารมฺมณสฺส นเหตุสฺส ธมฺมสฺส อารมฺมณปจฺจเยน ปจฺจโย. (สํขิตฺตํ.)}

\gambhira{ธมฺมานุโลม ติกทุกปฏฺฐานปาฬิ}
\prose{\csromanfolio{550}อารมฺมเณ จตฺตาริ, อธิปติยา ตีณิ, อนนฺตเร จตฺตาริ, อุปนิสฺสเย จตฺตาริ ฯเปฯ อวิคเต เทฺว. (สํขิตฺตํ.)}

\prose{นเหตุยา จตฺตาริ, น-อารมฺมเณ จตฺตาริ. (สํขิตฺตํ.)}

\csromanheader{2}{อารมฺมณปจฺจยา นเหตุยา จตฺตาริ. (สํขิตฺตํ.)}
\csromanheader{2}{นเหตุปจฺจยา อารมฺมเณ จตฺตาริ. (สํขิตฺตํ.)}
\prose{(ยถา กุสลตฺติเก ปญฺหาวารํ, เอวํ วิตฺถาเรตพฺพํ.)}

\csromanmatikaanchor{mtk.134}
\csromanmatikaanchor{mtk.135}
\csromantocmarknopage{subparagraph}{22. สนิทสฺสนสปฺปฏิฆตฺติก}{mtk.134}
\bookmark[dest={mtk.134},level=5]{22. สนิทสฺสนสปฺปฏิฆตฺติก}
\csromantocmarkref{subparagraph}{1. เหตุทุก}{mtk.135}
\bookmark[dest={mtk.135},level=5]{1. เหตุทุก}
\csromanheader{6}{22. สนิทสฺสนสปฺปฏิฆตฺติก 1. เหตุทุก}
\csromanheader{2}{1-7. ปฏิจฺจาทิวาร ปจฺจยจตุกฺก}
\proseitem{237}{อนิทสฺสน-อปฺปฏิฆํ เหตุํ ธมฺมํ ปฏิจฺจ อนิทสฺสน-อปฺปฏิโฆ เหตุ ธมฺโม อุปฺปชฺชติ เหตุปจฺจยา. (1) (สํขิตฺตํ.)}

\prose{เหตุยา เอกํ, อารมฺมเณ เอกํ ฯเปฯ อวิคเต เอกํ. (สํขิตฺตํ.)}

\prose{(สหชาตวารมฺปิ ฯเปฯ สมฺปยุตฺตวารมฺปิ ปฏิจฺจวารสทิสํ.)}

\proseitem{238}{อนิทสฺสน-อปฺปฏิโฆ เหตุ ธมฺโม อนิทสฺสน-อปฺปฏิฆสฺส เหตุสฺส ธมฺมสฺส เหตุปจฺจเยน ปจฺจโย. (1)}

\prose{อนิทสฺสน-อปฺปฏิโฆ เหตุ ธมฺโม อนิทสฺสน-อปฺปฏิฆสฺส เหตุสฺส ธมฺมสฺส อารมฺมณปจฺจเยน ปจฺจโย. (1) (สํขิตฺตํ.)}

\csromanhangingbegin
\hangingheaditem{238}{เหตุยา เอกํ, อารมฺมเณ เอกํ ฯเปฯ อวิคเต เอกํ. (สพฺพตฺถ เอกํ.)}
\hangingbody{(ยถา กุสลตฺติเก ปญฺหาวารํ, เอวํ วิตฺถาเรตพฺพํ.)}
\csromanhangingend

\csromanheader{2}{\textbf{นเหตุปท 1. ปฏิจฺจวาร}}
\csromanheader{2}{\textbf{ปจฺจยจตุกฺก เหตุ-อารมฺมณ}}
\proseitem{239}{อนิทสฺสนสปฺปฏิฆํ นเหตุํ ธมฺมํ ปฏิจฺจ อนิทสฺสนสปฺปฏิโฆ นเหตุ ธมฺโม อุปฺปชฺชติ เหตุปจฺจยา.\csromanspacer{}\csromanspacer{}อนิทสฺสนสปฺปฏิฆํ นเหตุํ ธมฺมํ ปฏิจฺจ สนิทสฺสนสปฺปฏิโฆ นเหตุ ธมฺโม อุปฺปชฺชติ เหตุปจฺจยา.\csromanspacer{}\csromanspacer{}อนิทสฺสนสปฺปฏิฆํ นเหตุํ ธมฺมํ ปฏิจฺจ อนิทสฺสน-อปฺปฏิโฆ นเหตุ ธมฺโม}

\vspace{\tipitakaparskip}
\csromancenter{สนิทสฺสนสปฺปฏิฆตฺติก 1.\csromanspacer{} เหตุทุก อุปฺปชฺชติ เหตุปจฺจยา.\csromanspacer{}\csromanspacer{}อนิทสฺสนสปฺปฏิฆํ นเหตุํ ธมฺมํ ปฏิจฺจ สนิทสฺสนสปฺปฏิโฆ นเหตุ จ อนิทสฺสน-อปฺปฏิโฆ นเหตุ จ ธมฺมา อุปฺปชฺชนฺติ เหตุปจฺจยา.\csromanspacer{}\csromanspacer{}อนิทสฺสนสปฺปฏิฆํ นเหตุํ ธมฺมํ ปฏิจฺจ อนิทสฺสนสปฺปฏิโฆ นเหตุ จ อนิทสฺสน-อปฺปฏิโฆ นเหตุ จ ธมฺมา อุปฺปชฺชนฺติ เหตุปจฺจยา.\csromanspacer{}\csromanspacer{}อนิทสฺสนสปฺปฏิฆํ นเหตุํ ธมฺมํ ปฏิจฺจ สนิทสฺสนสปฺปฏิโฆ นเหตุ จ อนิทสฺสนสปฺปฏิโฆ นเหตุ จ ธมฺมา อุปฺปชฺชนฺติ เหตุปจฺจยา.\csromanspacer{}\csromanspacer{}อนิทสฺสนสปฺปฏิฆํ นเหตุํ ธมฺมํ ปฏิจฺจ สนิทสฺสนสปฺปฏิโฆ นเหตุ จ อนิทสฺสนสปฺปฏิโฆ นเหตุ จ อนิทสฺสน-อปฺปฏิโฆ นเหตุ จ ธมฺมา อุปฺปชฺชนฺติ เหตุปจฺจยา.\csromanspacer{} สตฺต.}
\proseitem{240}{\csromanfolio{551}อนิทสฺสน-อปฺปฏิฆํ นเหตุํ ธมฺมํ ปฏิจฺจ อนิทสฺสน-อปฺปฏิโฆ นเหตุ ธมฺโม อุปฺปชฺชติ เหตุปจฺจยา.\csromanspacer{}\csromanspacer{}อนิทสฺสน-อปฺปฏิฆํ นเหตุํ ธมฺมํ ปฏิจฺจ สนิทสฺสนสปฺปฏิโฆ นเหตุ ธมฺโม อุปฺปชฺชติ เหตุปจฺจยา.\csromanspacer{}\csromanspacer{}อนิทสฺสน-อปฺปฏิฆํ นเหตุํ ธมฺมํ ปฏิจฺจ อนิทสฺสนสปฺปฏิโฆ นเหตุ ธมฺโม อุปฺปชฺชติ เหตุปจฺจยา.\csromanspacer{}\csromanspacer{}อนิทสฺสน-อปฺปฏิฆํ นเหตุํ ธมฺมํ ปฏิจฺจ สนิทสฺสนสปฺปฏิโฆ นเหตุ จ อนิทสฺสน-อปฺปฏิโฆ นเหตุ จ ธมฺมา อุปฺปชฺชนฺติ เหตุปจฺจยา.\csromanspacer{}\csromanspacer{}อนิทสฺสน-อปฺปฏิฆํ นเหตุํ ธมฺมํ ปฏิจฺจ อนิทสฺสนสปฺปฏิโฆ นเหตุ จ อนิทสฺสน-อปฺปฏิโฆ นเหตุ จ ธมฺมา อุปฺปชฺชนฺติ เหตุปจฺจยา.\csromanspacer{}\csromanspacer{}อนิทสฺสน-อปฺปฏิฆํ นเหตุํ ธมฺมํ ปฏิจฺจ สนิทสฺสนสปฺปฏิโฆ นเหตุ จ อนิทสฺสนสปฺปฏิโฆ นเหตุ จ ธมฺมา อุปฺปชฺชนฺติ เหตุปจฺจยา.\csromanspacer{}\csromanspacer{}อนิทสฺสน-อปฺปฏิฆํ นเหตุํ ธมฺมํ ปฏิจฺจ สนิทสฺสนสปฺปฏิโฆ นเหตุ จ อนิทสฺสนสปฺปฏิโฆ นเหตุ จ อนิทสฺสน-อปฺปฏิโฆ นเหตุ จ ธมฺมา อุปฺปชฺชนฺติ เหตุปจฺจยา.\csromanspacer{} สตฺต. อนิทสฺสนสปฺปฏิฆํ นเหตุญฺจ อนิทสฺสน-อปฺปฏิฆํ นเหตุญฺจ ธมฺมํ ปฏิจฺจ สนิทสฺสนสปฺปฏิโฆ นเหตุ ธมฺโม อุปฺปชฺชติ เหตุปจฺจยา.\csromanspacer{}\csromanspacer{}อนิทสฺสนสปฺปฏิฆํ นเหตุญฺจ อนิทสฺสน-อปฺปฏิฆํ นเหตุญฺจ ธมฺมํ ปฏิจฺจ อนิทสฺสนสปฺปฏิโฆ นเหตุ ธมฺโม อุปฺปชฺชติ เหตุปจฺจยา.\csromanspacer{}\csromanspacer{}อนิทสฺสนสปฺปฏิฆํ นเหตุญฺจ อนิทสฺสน-อปฺปฏิฆํ นเหตุญฺจ ธมฺมํ ปฏิจฺจ อนิทสฺสน-อปฺปฏิโฆ นเหตุ ธมฺโม อุปฺปชฺชติ เหตุปจฺจยา.\csromanspacer{}\csromanspacer{}อนิทสฺสนสปฺปฏิฆํ นเหตุญฺจ อนิทสฺสน-อปฺปฏิฆํ นเหตุญฺจ ธมฺมํ ปฏิจฺจ สนิทสฺสนสปฺปฏิโฆ นเหตุ จ อนิทสฺสน-อปฺปฏิโฆ นเหตุ จ ธมฺมา อุปฺปชฺชนฺติ เหตุปจฺจยา.\csromanspacer{}\csromanspacer{}อนิทสฺสนสปฺปฏิฆํ นเหตุญฺจ อนิทสฺสน-อปฺปฏิฆํ นเหตุญฺจ ธมฺมํ ปฏิจฺจ}

\vspace{\tipitakaparskip}
\csromancenter{ธมฺมานุโลม ติกทุกปฏฺฐานปาฬิ อนิทสฺสนสปฺปฏิโฆ นเหตุ จ อนิทสฺสน-อปฺปฏิโฆ นเหตุ จ ธมฺมา อุปฺปชฺชนฺติ เหตุปจฺจยา.\csromanspacer{}\csromanspacer{}อนิทสฺสนสปฺปฏิฆํ นเหตุญฺจ อนิทสฺสน-อปฺปฏิฆํ นเหตุญฺจ ธมฺมํ ปฏิจฺจ สนิทสฺสนสปฺปฏิโฆ นเหตุ จ อนิทสฺสนสปฺปฏิโฆ นเหตุ จ ธมฺมา อุปฺปชฺชนฺติ เหตุปจฺจยา.\csromanspacer{}\csromanspacer{}อนิทสฺสนสปฺปฏิฆํ นเหตุญฺจ อนิทสฺสน-อปฺปฏิฆํ นเหตุญฺจ ธมฺมํ ปฏิจฺจ สนิทสฺสนสปฺปฏิโฆ นเหตุ จ อนิทสฺสนสปฺปฏิโฆ นเหตุ จ อนิทสฺสน-อปฺปฏิโฆ นเหตุ จ ธมฺมา อุปฺปชฺชนฺติ เหตุปจฺจยา.\csromanspacer{} สตฺต.}
\prose{\csromanfolio{552}อนิทสฺสน-อปฺปฏิฆํ นเหตุํ ธมฺมํ ปฏิจฺจ อนิทสฺสน-อปฺปฏิโฆ นเหตุ ธมฺโม อุปฺปชฺชติ อารมฺมณปจฺจยา. (1) (สํขิตฺตํ.)}

\proseitem{242}{เหตุยา เอกวีส, อารมฺมเณ เอกํ, อธิปติยา เอกวีส ฯเปฯ อวิคเต เอกวีส. (สํขิตฺตํ.)}

\prose{นเหตุยา เอกวีส, น-อารมฺมเณ เอกวีส ฯเปฯ โนวิคเต เอกวีส.}

\csromanheader{2}{เหตุปจฺจยา น-อารมฺมเณ เอกวีส. (สํขิตฺตํ.)}
\csromanheader{2}{นเหตุปจฺจยา อารมฺมเณ เอกํ. (สํขิตฺตํ.)}
\prose{(สหชาตวารมฺปิ ฯเปฯ สมฺปยุตฺตวารมฺปิ ปฏิจฺจวารสทิสํ.)}

\csromanheader{2}{7. ปญฺหาวาร ปจฺจยจตุกฺก-เหตุ}
\proseitem{243}{สนิทสฺสนสปฺปฏิโฆ นเหตุ ธมฺโม อนิทสฺสน-อปฺปฏิฆสฺส นเหตุสฺส ธมฺมสฺส อารมฺมณปจฺจเยน ปจฺจโย. (สํขิตฺตํ.)}

\prose{อารมฺมเณ ตีณิ, อธิปติยา นว, อนนฺตเร เอกํ, สมนนฺตเร เอกํ, สหชาเต เอกวีส, อุปนิสฺสเย ตีณิ ฯเปฯ อวิคเต ปญฺจวีส. (สํขิตฺตํ.)}

\prose{นเหตุยา ปญฺจวีส, น-อารมฺมเณ เอกวีส. (สํขิตฺตํ.)}

\csromanheader{2}{อารมฺมณปจฺจยา นเหตุยา ตีณิ. (สํขิตฺตํ.)}
\prose{นเหตุปจฺจยา อารมฺมเณ ตีณิ. (สํขิตฺตํ.) (ยถา กุสลตฺติเก ปญฺหาวารํ, เอวํ วิตฺถาเรตพฺพํ.)}

\nitthitam{สนิทสฺสนสปฺปฏิฆตฺติกเหตุทุกํ นิฏฺฐิตํ.}
\csromanmatikaanchor{mtk.136}
\csromanmatikaanchor{mtk.137}
\csromantocmarknopage{subparagraph}{1. กุสลตฺติก}{mtk.136}
\bookmark[dest={mtk.136},level=5]{1. กุสลตฺติก}
\csromantocmarkref{subparagraph}{2. สเหตุกทุก}{mtk.137}
\bookmark[dest={mtk.137},level=5]{2. สเหตุกทุก}
\csromanheader{6}{1. กุสลตฺติก 2. สเหตุกทุก \textbf{1. กุสลตฺติก 2. สเหตุกทุก}}
\csromanheader{2}{1-7. ปฏิจฺจาทิวาร ปจฺจยจตุกฺก-เหตุ}
\proseitem{244}{\csromanfolio{553}กุสลํ สเหตุกํ ธมฺมํ ปฏิจฺจ กุสโล สเหตุโก ธมฺโม อุปฺปชฺชติ เหตุปจฺจยา. (1)}

\prose{อกุสลํ สเหตุกํ ธมฺมํ ปฏิจฺจ อกุสโล สเหตุโก ธมฺโม อุปฺปชฺชติ เหตุปจฺจยา. (1)}

\prose{อพฺยากตํ สเหตุกํ ธมฺมํ ปฏิจฺจ อพฺยากโต สเหตุโก ธมฺโม อุปฺปชฺชติ เหตุปจฺจยา. (1) (สํขิตฺตํ.)}

\csromanheader{2}{245. เหตุยา ตีณิ ฯเปฯ อวิคเต ตีณิ. (สํขิตฺตํ.)}
\prose{น-อธิปติยา ตีณิ, นปุเรชาเต ตีณิ, นปจฺฉาชาเต ตีณิ, น-อาเสวเน ตีณิ, นกมฺเม ตีณิ, นวิปาเก ตีณิ, นวิปฺปยุตฺเต ตีณิ.}

\csromanheader{2}{เหตุปจฺจยา น-อธิปติยา ตีณิ. (สํขิตฺตํ.)}
\prose{น-อธิปติปจฺจยา เหตุยา ตีณิ. (สํขิตฺตํ.) (สหชาตวารมฺปิ ฯเปฯ สมฺปยุตฺตวารมฺปิ ปฏิจฺจวารสทิสํ.)}

\csromanheader{2}{\textbf{เหตุ-อารมฺมณ}}
\proseitem{246}{กุสโล สเหตุโก ธมฺโม กุสลสฺส สเหตุกสฺส ธมฺมสฺส เหตุปจฺจเยน ปจฺจโย. (1)}

\prose{อกุสโล สเหตุโก ธมฺโม อกุสลสฺส สเหตุกสฺส ธมฺมสฺส เหตุปจฺจเยน ปจฺจโย. (1)}

\prose{อพฺยากโต สเหตุโก ธมฺโม อพฺยากตสฺส สเหตุกสฺส ธมฺมสฺส เหตุปจฺจเยน ปจฺจโย. (1)}

\prose{กุสโล สเหตุโก ธมฺโม กุสลสฺส สเหตุกสฺส ธมฺมสฺส อารมฺมณปจฺจเยน ปจฺจโย.\csromanspacer{}\csromanspacer{}กุสโล สเหตุโก ธมฺโม อกุสลสฺส สเหตุกสฺส ธมฺมสฺส อารมฺมณปจฺจเยน ปจฺจโย.\csromanspacer{}\csromanspacer{}กุสโล สเหตุโก ธมฺโม อพฺยากตสฺส สเหตุกสฺส ธมฺมสฺส อารมฺมณปจฺจเยน ปจฺจโย. (3)}

\gambhira{ธมฺมานุโลม ติกทุกปฏฺฐานปาฬิ}
\prose{\csromanfolio{554}อกุสโล สเหตุโก ธมฺโม อกุสลสฺส สเหตุกสฺส ธมฺมสฺส อารมฺมณปจฺจเยน ปจฺจโย.\csromanspacer{} ตีณิ.}

\prose{อพฺยากโต สเหตุโก ธมฺโม อพฺยากตสฺส สเหตุกสฺส ธมฺมสฺส อารมฺมณปจฺจเยน ปจฺจโย.\csromanspacer{} ตีณิ.}

\prose{กุสโล สเหตุโก ธมฺโม กุสลสฺส สเหตุกสฺส ธมฺมสฺส อธิปติปจฺจเยน ปจฺจโย. (สํขิตฺตํ.)}

\proseitem{247}{เหตุยา ตีณิ, อารมฺมเณ นว, อธิปติยา สตฺต, อนนฺตเร ปญฺจ, สมนนฺตเร ปญฺจ, สหชาเต ตีณิ, อญฺญมญฺเญ ตีณิ, นิสฺสเย ตีณิ, อุปนิสฺสเย นว ฯเปฯ อวิคเต ตีณิ. (สํขิตฺตํ.)}

\prose{นเหตุยา นว, น-อารมฺมเณ นว. (สํขิตฺตํ.)}

\csromanheader{2}{เหตุปจฺจยา น-อารมฺมเณ ตีณิ. (สํขิตฺตํ.)}
\csromanheader{2}{นเหตุปจฺจยา อารมฺมเณ นว. (สํขิตฺตํ.)}
\prose{(ยถา กุสลตฺติเก ปญฺหาวารํ, เอวํ วิตฺถาเรตพฺพํ.)}

\csromanheader{2}{\textbf{อเหตุกปท 1-7. ปฏิจฺจาทิวาร}}
\csromanheader{2}{\textbf{ปจฺจยจตุกฺก-เหตุ อารมฺมณ}}
\proseitem{248}{อกุสลํ อเหตุกํ ธมฺมํ ปฏิจฺจ อพฺยากโต อเหตุโก ธมฺโม อุปฺปชฺชติ เหตุปจฺจยา. (1)}

\prose{อพฺยากตํ อเหตุกํ ธมฺมํ ปฏิจฺจ อพฺยากโต อเหตุโก ธมฺโม อุปฺปชฺชติ เหตุปจฺจยา. (1)}

\prose{อกุสลํ อเหตุกญฺจ อพฺยากตํ อเหตุกญฺจ ธมฺมํ ปฏิจฺจ อพฺยากโต อเหตุโก ธมฺโม อุปฺปชฺชติ เหตุปจฺจยา. (1)}

\prose{อพฺยากตํ อเหตุกํ ธมฺมํ ปฏิจฺจ อพฺยากโต อเหตุโก ธมฺโม อุปฺปชฺชติ อารมฺมณปจฺจยา. (1) (สํขิตฺตํ.)}

\proseitem{249}{เหตุยา ตีณิ, อารมฺมเณ เอกํ, อธิปติยา เอกํ ฯเปฯ อวิคเต ตีณิ. (สํขิตฺตํ.)}

\csromanheader{2}{1. กุสลตฺติก 2. สเหตุกทุก \textbf{ปจฺจนีย-นเหตุ น-อารมฺมณ}}
\prose{\csromanfolio{555}อพฺยากตํ อเหตุกํ ธมฺมํ ปฏิจฺจ อพฺยากโต อเหตุโก ธมฺโม อุปฺปชฺชติ นเหตุปจฺจยา. (1)}

\prose{อกุสลํ อเหตุกํ ธมฺมํ ปฏิจฺจ อพฺยากโต อเหตุโก ธมฺโม อุปฺปชฺชติ น-อารมฺมณปจฺจยา. (1)}

\prose{อพฺยากตํ อเหตุกํ ธมฺมํ ปฏิจฺจ อพฺยากโต อเหตุโก ธมฺโม อุปฺปชฺชติ น-อารมฺมณปจฺจยา. (1)}

\prose{อกุสลํ อเหตุกญฺจ อพฺยากตํ อเหตุกญฺจ ธมฺมํ ปฏิจฺจ อพฺยากโต อเหตุโก ธมฺโม อุปฺปชฺชติ น-อารมฺมณปจฺจยา. (1) (สํขิตฺตํ.)}

\proseitem{250}{นเหตุยา เอกํ, น-อารมฺมเณ ตีณิ, น-อธิปติยา ตีณิ ฯเปฯ โนวิคเต ตีณิ. (สํขิตฺตํ.)}

\csromanheader{2}{เหตุปจฺจยา น-อารมฺมเณ ตีณิ. (สํขิตฺตํ.)}
\prose{นเหตุปจฺจยา อารมฺมเณ เอกํ. (สํขิตฺตํ.) (สหชาตวารมฺปิ ฯเปฯ สมฺปยุตฺตวารมฺปิ ปฏิจฺจวารสทิสํ.)}

\csromanheader{2}{\textbf{เหตุ อารมฺมณาทิ}}
\proseitem{251}{อกุสโล อเหตุโก ธมฺโม อพฺยากตสฺส อเหตุกสฺส ธมฺมสฺส เหตุปจฺจเยน ปจฺจโย. (1)}

\prose{อกุสโล อเหตุโก ธมฺโม อกุสลสฺส อเหตุกสฺส ธมฺมสฺส อารมฺมณปจฺจเยน ปจฺจโย.\csromanspacer{}\csromanspacer{}อกุสโล อเหตุโก ธมฺโม อพฺยากตสฺส อเหตุกสฺส ธมฺมสฺส อารมฺมณปจฺจเยน ปจฺจโย. (2)}

\prose{อพฺยากโต อเหตุโก ธมฺโม อพฺยากตสฺส อเหตุกสฺส ธมฺมสฺส อารมฺมณปจฺจเยน ปจฺจโย.\csromanspacer{}\csromanspacer{}อพฺยากโต อเหตุโก ธมฺโม อกุสลสฺส อเหตุกสฺส ธมฺมสฺส อารมฺมณปจฺจเยน ปจฺจโย. (2)}

\prose{อกุสโล อเหตุโก ธมฺโม อกุสลสฺส อเหตุกสฺส ธมฺมสฺส อนนฺตรปจฺจเยน ปจฺจโย.\csromanspacer{}\csromanspacer{}อกุสโล อเหตุโก ธมฺโม อพฺยากตสฺส อเหตุกสฺส ธมฺมสฺส อนนฺตรปจฺจเยน ปจฺจโย. (2)}

\gambhira{ธมฺมานุโลม ติกทุกปฏฺฐานปาฬิ}
\prose{\csromanfolio{556}อพฺยากโต อเหตุโก ธมฺโม อพฺยากตสฺส อเหตุกสฺส ธมฺมสฺส อนนฺตรปจฺจเยน ปจฺจโย.\csromanspacer{}\csromanspacer{}อพฺยากโต อเหตุโก ธมฺโม อกุสลสฺส อเหตุกสฺส ธมฺมสฺส อนนฺตรปจฺจเยน ปจฺจโย ฯเปฯ. (2)}

\prose{อกุสโล อเหตุโก ธมฺโม อกุสลสฺส อเหตุกสฺส ธมฺมสฺส อุปนิสฺสยปจฺจเยน ปจฺจโย. (สํขิตฺตํ.)}

\proseitem{252}{เหตุยา เอกํ, อารมฺมเณ จตฺตาริ, อนนฺตเร จตฺตาริ, สมนนฺตเร จตฺตาริ, สหชาเต ตีณิ, อญฺญมญฺเญ เอกํ, นิสฺสเย จตฺตาริ, อุปนิสฺสเย จตฺตาริ. (สํขิตฺตํ.)}

\prose{นเหตุยา ปญฺจ, น-อารมฺมเณ ปญฺจ, น-อธิปติยา ปญฺจ. (สํขิตฺตํ.)}

\csromanheader{2}{เหตุปจฺจยา น-อารมฺมเณ เอกํ. (สํขิตฺตํ.)}
\csromanheader{2}{นเหตุปจฺจยา อารมฺมเณ จตฺตาริ. (สํขิตฺตํ.)}
\prose{(ยถา กุสลตฺติเก ปญฺหาวารํ, เอวํ วิตฺถาเรตพฺพํ.)}

\csromanmatikaanchor{mtk.138}
\csromantocmarkref{subparagraph}{3. เหตุสมฺปยุตฺตทุก}{mtk.138}
\bookmark[dest={mtk.138},level=5]{3. เหตุสมฺปยุตฺตทุก}
\csromanheader{6}{1. กุสลตฺติก 3. เหตุสมฺปยุตฺตทุก}
\csromanheader{2}{1-7. ปฏิจฺจาทิวาร ปจฺจยจตุกฺก-เหตุ}
\proseitem{253}{กุสลํ เหตุสมฺปยุตฺตํ ธมฺมํ ปฏิจฺจ กุสโล เหตุสมฺปยุตฺโต ธมฺโม อุปฺปชฺชติ เหตุปจฺจยา. (1)}

\prose{อกุสลํ เหตุสมฺปยุตฺตํ ธมฺมํ ปฏิจฺจ อกุสโล เหตุสมฺปยุตฺโต ธมฺโม อุปฺปชฺชติ เหตุปจฺจยา. (1)}

\prose{อพฺยากตํ เหตุสมฺปยุตฺตํ ธมฺมํ ปฏิจฺจ อพฺยากโต เหตุสมฺปยุตฺโต ธมฺโม อุปฺปชฺชติ เหตุปจฺจยา. (1) (สํขิตฺตํ.)}

\proseitem{254}{เหตุยา ตีณิ, อารมฺมเณ ตีณิ ฯเปฯ อวิคเต ตีณิ. (สํขิตฺตํ.)}

\csromanheader{2}{น-อธิปติยา ตีณิ ฯเปฯ นวิปฺปยุตฺเต ตีณิ. (สํขิตฺตํ.)}
\prose{(สหชาตวารมฺปิ ฯเปฯ สมฺปยุตฺตวารมฺปิ ปฏิจฺจวารสทิสํ.)}

\csromanheader{2}{1. กุสลตฺติก 3. เหตุสมฺปยุตฺตทุก \textbf{เหตุ-อารมฺมณ}}
\proseitem{255}{\csromanfolio{557}กุสโล เหตุสมฺปยุตฺโต ธมฺโม กุสลสฺส เหตุสมฺปยุตฺตสฺส ธมฺมสฺส เหตุปจฺจเยน ปจฺจโย. (1)}

\prose{อกุสโล เหตุสมฺปยุตฺโต ธมฺโม อกุสลสฺส เหตุสมฺปยุตฺตสฺส ธมฺมสฺส เหตุปจฺจเยน ปจฺจโย. (1)}

\prose{อพฺยากโต เหตุสมฺปยุตฺโต ธมฺโม อพฺยากตสฺส เหตุสมฺปยุตฺตสฺส ธมฺมสฺส เหตุปจฺจเยน ปจฺจโย. (1)}

\prose{กุสโล เหตุสมฺปยุตฺโต ธมฺโม กุสลสฺส เหตุสมฺปยุตฺตสฺส ธมฺมสฺส อารมฺมณปจฺจเยน ปจฺจโย.\csromanspacer{} ตีณิ.}

\prose{อกุสโล เหตุสมฺปยุตฺโต ธมฺโม อกุสลสฺส เหตุสมฺปยุตฺตสฺส ธมฺมสฺส อารมฺมณปจฺจเยน ปจฺจโย.\csromanspacer{} ตีณิ.}

\prose{อพฺยากโต เหตุสมฺปยุตฺโต ธมฺโม อพฺยากตสฺส เหตุสมฺปยุตฺตสฺส ธมฺมสฺส อารมฺมณปจฺจเยน ปจฺจโย.\csromanspacer{} ตีณิ. (สํขิตฺตํ.)}

\proseitem{256}{เหตุยา ตีณิ, อารมฺมเณ นว, อธิปติยา สตฺต, อนนฺตเร ปญฺจ, สมนนฺตเร ปญฺจ, สหชาเต ตีณิ, อุปนิสฺสเย นว ฯเปฯ อวิคเต ตีณิ. (สํขิตฺตํ.)}

\prose{นเหตุยา นว, น-อารมฺมเณ นว. (สํขิตฺตํ.)}

\csromanheader{2}{เหตุปจฺจยา น-อารมฺมเณ ตีณิ. (สํขิตฺตํ.)}
\csromanheader{2}{นเหตุปจฺจยา อารมฺมเณ นว. (สํขิตฺตํ.)}
\prose{(ยถา กสลตฺติเก ปญฺหาวารํ, เอวํ วิตฺถาเรตพฺพํ.)}

\csromanheader{2}{\textbf{เหตุวิปฺปยุตฺตปท 1-7. ปฏิจฺจาทิวาร}}
\csromanheader{2}{\textbf{ปจฺจยจตุกฺก-เหตุ อารมฺมณ}}
\proseitem{257}{อกุสลํ เหตุวิปฺปยุตฺตํ ธมฺมํ ปฏิจฺจ อพฺยากโต เหตุวิปฺปยุตฺโต ธมฺโม อุปฺปชฺชติ เหตุปจฺจยา. (1)}

\gambhira{ธมฺมานุโลม ติกทุกปฏฺฐานปาฬิ}
\prose{\csromanfolio{558}อพฺยากตํ เหตุวิปฺปยุตฺตํ ธมฺมํ ปฏิจฺจ อพฺยากโต เหตุวิปฺปยุตฺโต ธมฺโม อุปฺปชฺชติ เหตุปจฺจยา. (1)}

\prose{อกุสลํ เหตุวิปฺปยุตฺตญฺจ อพฺยากตํ เหตุวิปฺปยุตฺตญฺจ ธมฺมํ ปฏิจฺจ อพฺยากโต เหตุวิปฺปยุตฺโต ธมฺโม อุปฺปชฺชติ เหตุปจฺจยา. (1)}

\prose{อพฺยากตํ เหตุวิปฺปยุตฺตํ ธมฺมํ ปฏิจฺจ อพฺยากโต เหตุวิปฺปยุตฺโต ธมฺโม อุปฺปชฺชติ อารมฺมณปจฺจยา. (1) (สํขิตฺตํ.)}

\proseitem{258}{เหตุยา ตีณิ, อารมฺมเณ เอกํ, อธิปติยา เอกํ ฯเปฯ อวิคเต ตีณิ. (สํขิตฺตํ.)}

\csromanheader{2}{\textbf{นเหตุ น-อารมฺมณ}}
\proseitem{259}{อพฺยากตํ เหตุวิปฺปยุตฺตํ ธมฺมํ ปฏิจฺจ อพฺยากโต เหตุวิปฺปยุตฺโต ธมฺโม อุปฺปชฺชติ นเหตุปจฺจยา. (1)}

\prose{อกุสลํ เหตุวิปฺปยุตฺตํ ธมฺมํ ปฏิจฺจ อพฺยากโต เหตุวิปฺปยุตฺโต ธมฺโม อุปฺปชฺชติ น-อารมฺมณปจฺจยา. (1)}

\prose{อพฺยากตํ เหตุวิปฺปยุตฺตํ ธมฺมํ ปฏิจฺจ อพฺยากโต เหตุวิปฺปยุตฺโต ธมฺโม อุปฺปชฺชติ น-อารมฺมณปจฺจยา. (1)}

\prose{อกุสลํ เหตุวิปฺปยุตฺตญฺจ อพฺยากตํ เหตุวิปฺปยุตฺตญฺจ ธมฺมํ ปฏิจฺจ อพฺยากโต เหตุวิปฺปยุตฺโต ธมฺโม อุปฺปชฺชติ น-อารมฺมณปจฺจยา. (สํขิตฺตํ.)}

\proseitem{260}{นเหตุยา เอกํ, น-อารมฺมเณ ตีณิ, น-อธิปติยา ตีณิ ฯเปฯ โนวิคเต ตีณิ. (สํขิตฺตํ.\csromanspacer{} สหชาตวารมฺปิ ฯเปฯ สมฺปยุตฺตวารมฺปิ ปฏิจฺจวารสทิสํ.)}

\csromanheader{2}{\textbf{เหตุ อารมฺมณ อนนฺตร}}
\proseitem{261}{อกุสโล เหตุวิปฺปยุตฺโต ธมฺโม อพฺยากตสฺส เหตุวิปฺปยุตฺตสฺส ธมฺมสฺส เหตุปจฺจเยน ปจฺจโย. (1)}

\prose{อกุสโล เหตุวิปฺปยุตฺโต ธมฺโม อกุสลสฺส เหตุวิปฺปยุตฺตสฺส ธมฺมสฺส อารมฺมณปจฺจเยน ปจฺจโย.\csromanspacer{}\csromanspacer{}อกุสโล เหตุวิปฺปยุตฺโต ธมฺโม อพฺยากตสฺส เหตุวิปฺปยุตฺตสฺส ธมฺมสฺส อารมฺมณปจฺจเยน ปจฺจโย. (2)}

\csromanmatikaanchor{mtk.139}
\csromantocmarkref{subparagraph}{4. เหตุสเหตุกทุก}{mtk.139}
\bookmark[dest={mtk.139},level=5]{4. เหตุสเหตุกทุก}
\csromanheader{6}{1. กุสลตฺติก 4. เหตุสเหตุกทุก}
\prose{\csromanfolio{559}อพฺยากโต เหตุวิปฺปยุตฺโต ธมฺโม อพฺยากตสฺส เหตุวิปฺปยุตฺตสฺส ธมฺมสฺส อารมฺมณปจฺจเยน ปจฺจโย.\csromanspacer{}\csromanspacer{}อพฺยากโต เหตุวิปฺปยุตฺโต ธมฺโม อกุสลสฺส เหตุวิปฺปยุตฺตสฺส ธมฺมสฺส อารมฺมณปจฺจเยน ปจฺจโย. (2)}

\prose{อกุสโล เหตุวิปฺปยุตฺโต ธมฺโม อกุสลสฺส เหตุวิปฺปยุตฺตสฺส ธมฺมสฺส อนนฺตรปจฺจเยน ปจฺจโย.\csromanspacer{}\csromanspacer{}อกุสโล เหตุวิปฺปยุตฺโต ธมฺโม อพฺยากตสฺส เหตุวิปฺปยุตฺตสฺส ธมฺมสฺส อนนฺตรปจฺจเยน ปจฺจโย. (2)}

\prose{อพฺยากโต เหตุวิปฺปยุตฺโต ธมฺโม อพฺยากตสฺส เหตุวิปฺปยุตฺตสฺส ธมฺมสฺส อนนฺตรปจฺจเยน ปจฺจโย.\csromanspacer{}\csromanspacer{}อพฺยากโต เหตุวิปฺปยุตฺโต ธมฺโม อกุสลสฺส เหตุวิปฺปยุตฺตสฺส ธมฺมสฺส อนนฺตรปจฺจเยน ปจฺจโย. (2) (สํขิตฺตํ.)}

\proseitem{262}{เหตุยา เอกํ, อารมฺมเณ จตฺตาริ, อนนฺตเร จตฺตาริ, สมนนฺตเร จตฺตาริ, สหชาเต ตีณิ, อญฺญมญฺเญ เอกํ, นิสฺสเย จตฺตาริ, อุปนิสฺสเย จตฺตาริ ฯเปฯ อวิคเต จตฺตาริ. (สํขิตฺตํ.)}

\prose{นเหตุยา ปญฺจ, น-อารมฺมเณ ปญฺจ. (สํขิตฺตํ.)}

\csromanheader{2}{เหตุปจฺจยา น-อารมฺมเณ เอกํ. (สํขิตฺตํ.)}
\csromanheader{2}{นเหตุปจฺจยา อารมฺมเณ จตฺตาริ. (สํขิตฺตํ.)}
\prose{(ยถา กุสลตฺติเก ปญฺหาวารํ, เอวํ วิตฺถาเรตพฺพํ.)}

\csromanheader{2}{1. กุสลตฺติก 4. เหตุสเหตุกทุก}
\csromanheader{2}{1-7. ปฏิจฺจาทิวาร ปจฺจยจตุกฺก-เหตุ}
\proseitem{263}{กุสลํ เหตุญฺเจว สเหตุกญฺจ ธมฺมํ ปฏิจฺจ กุสโล เหตุ เจว สเหตุโก จ ธมฺโม อุปฺปชฺชติ เหตุปจฺจยา. (1)}

\prose{อกุสลํ เหตุญฺเจว สเหตุญฺจ ธมฺมํ ปฏิจฺจ อกุสโล เหตุ เจว สเหตุโก จ ธมฺโม อุปฺปชฺชติ เหตุปจฺจยา. (1)}

\gambhira{ธมฺมานุโลม ติกทุกปฏฺฐานปาฬิ}
\prose{\csromanfolio{560}อพฺยากตํ เหตุญฺเจว สเหตุกญฺจ ธมฺมํ ปฏิจฺจ อพฺยากโต เหตุ เจว สเหตุโก จ ธมฺโม อุปฺปชฺชติ เหตุปจฺจยา. (1) (สํขิตฺตํ.)}

\proseitem{264}{เหตุยา ตีณิ, อารมฺมเณ ตีณิ ฯเปฯ อวิคเต ตีณิ. (สํขิตฺตํ.)}

\prose{น-อธิปติยา ตีณิ, นปุเรชาเต ตีณิ, นปจฺฉาชาเต ตีณิ, น-อาเสวเน ตีณิ, นวิปาเก ตีณิ, นวิปฺปยุตฺเต ตีณิ. (สํขิตฺตํ.) (สหชาตวารมฺปิ ฯเปฯ สมฺปยุตฺตวารมฺปิ ปฏิจฺจวารสทิสํ.)}

\csromanheader{2}{\textbf{เหตุ อารมฺมณ อธิปติ}}
\proseitem{265}{กุสโล เหตุ เจว สเหตุโก จ ธมฺโม กุสลสฺส เหตุสฺส เจว สเหตุกสฺส จ ธมฺมสฺส เหตุปจฺจเยน ปจฺจโย. (1)}

\prose{อกุสโล เหตุ เจว สเหตุโก จ ธมฺโม อกุสลสฺส เหตุสฺส เจว สเหตุกสฺส จ ธมฺมสฺส เหตุปจฺจเยน ปจฺจโย. (1)}

\prose{อพฺยากโต เหตุ เจว สเหตุโก จ ธมฺโม อพฺยากตสฺส เหตุสฺส เจว สเหตุกสฺส จ ธมฺมสฺส เหตุปจฺจเยน ปจฺจโย. (1)}

\prose{กุสโล เหตุ เจว สเหตุโก จ ธมฺโม กุสลสฺส เหตุสฺส เจว สเหตุกสฺส จ ธมฺมสฺส อารมฺมณปจฺจเยน ปจฺจโย.\csromanspacer{}\csromanspacer{}กุสโล เหตุ เจว สเหตุโก จ ธมฺโม อกุสลสฺส เหตุสฺส เจว สเหตุกสฺส จ ธมฺมสฺส อารมฺมณปจฺจเยน ปจฺจโย.\csromanspacer{}\csromanspacer{}กุสโล เหตุ เจว สเหตุโก จ ธมฺโม อพฺยากตสฺส เหตุสฺส เจว สเหตุกสฺส จ ธมฺมสฺส อารมฺมณปจฺจเยน ปจฺจโย. (3)}

\prose{อกุสโล เหตุ เจว สเหตุโก จ ธมฺโม อกุสลสฺส เหตุสฺส เจว สเหตุกสฺส จ ธมฺมสฺส อารมฺมณปจฺจเยน ปจฺจโย.\csromanspacer{} ตีณิ.}

\prose{อพฺยากโต เหตุ เจว สเหตุโก จ ธมฺโม อพฺยากตสฺส เหตุสฺส เจว สเหตุกสฺส จ ธมฺมสฺส อารมฺมณปจฺจเยน ปจฺจโย.\csromanspacer{} ตีณิ.}

\prose{กุสโล เหตุ เจว สเหตุโก จ ธมฺโม กุสลสฺส เหตุสฺส เจว สเหตุกสฺส จ ธมฺมสฺส อธิปติปจฺจเยน ปจฺจโย.\csromanspacer{} ตีณิ.}

\csromanheader{2}{1. กุสลตฺติก 4. เหตุสเหตุกทุก}
\prose{\csromanfolio{561}อกุสโล เหตุ เจว สเหตุโก จ ธมฺโม อกุสลสฺส เหตุสฺส เจว สเหตุกสฺส จ ธมฺมสฺส อธิปติปจฺจเยน ปจฺจโย. (1)}

\prose{อพฺยากโต เหตุ เจว สเหตุโก จ ธมฺโม อพฺยากตสฺส เหตุสฺส เจว สเหตุกสฺส จ ธมฺมสฺส อธิปติปจฺจเยน ปจฺจโย.\csromanspacer{} ตีณิ. (สํขิตฺตํ.)}

\proseitem{266}{เหตุยา ตีณิ, อารมฺมเณ นว, อธิปติยา สตฺต, อนนฺตเร ปญฺจ, สมนนฺตเร ปญฺจ, สหชาเต ตีณิ ฯเปฯ อุปนิสฺสเย นว ฯเปฯ อวิคเต ตีณิ. (สํขิตฺตํ.)}

\prose{นเหตุยา นว, น-อารมฺมเณ นว. (สํขิตฺตํ.)}

\csromanheader{2}{เหตุปจฺจยา น-อารมฺมเณ ตีณิ. (สํขิตฺตํ.)}
\prose{นเหตุปจฺจยา อารมฺมเณ นว. (สํขิตฺตํ.) (ยถา กุสลตฺติเก ปญฺหาวารํ, เอวํ วิตฺถาเรตพฺพํ.)}

\csromanheader{2}{\textbf{สเหตุกนเหตุปท 1-7. ปฏิจฺจาทิวาร}}
\csromanheader{2}{\textbf{ปจฺจยจตุกฺก-เหตุ}}
\proseitem{267}{กุสลํ สเหตุกญฺเจว น จ เหตุํ ธมฺมํ ปฏิจฺจ กุสโล สเหตุโก เจว น จ เหตุ ธมฺโม อุปฺปชฺชติ เหตุปจฺจยา. (1)}

\prose{อกุสลํ สเหตุกญฺเจว น จ เหตุํ ธมฺมํ ปฏิจฺจ อกุสโล สเหตุโก เจว น จ เหตุ ธมฺโม อุปฺปชฺชติ เหตุปจฺจยา. (1)}

\prose{อพฺยากตํ สเหตุกญฺเจว น จ เหตุํ ธมฺมํ ปฏิจฺจ อพฺยากโต สเหตุโก เจว น จ เหตุ ธมฺโม อุปฺปชฺชติ เหตุปจฺจยา. (สํขิตฺตํ.)}

\proseitem{268}{เหตุยา ตีณิ, อารมฺมเณ ตีณิ ฯเปฯ อวิคเต ตีณิ. (สํขิตฺตํ.)}

\csromanheader{2}{น-อธิปติยา ตีณิ ฯเปฯ นวิปฺปยุตฺเต ตีณิ. (สํขิตฺตํ.)}
\prose{(สหชาตวารมฺปิ ฯเปฯ สมฺปยุตฺตวารมฺปิ ปฏิจฺจวารสทิสํ.)}

\csromanheader{2}{ธมฺมานุโลม ติกทุกปฏฺฐานปาฬิ \textbf{อารมฺมณ}}
\proseitem{269}{\csromanfolio{562}กุสโล สเหตุโก เจว น จ เหตุ ธมฺโม กุสลสฺส สเหตุกสฺส เจว น จ เหตุสฺส ธมฺมสฺส อารมฺมณปจฺจเยน ปจฺจโย.\csromanspacer{} ตีณิ.}

\prose{อกุสโล สเหตุโก เจว น จ เหตุ ธมฺโม อกุสลสฺส สเหตุกสฺส เจว น จ เหตุสฺส ธมฺมสฺส อารมฺมณปจฺจเยน ปจฺจโย.\csromanspacer{} ตีณิ.}

\prose{อพฺยากโต สเหตุโก เจว น จ เหตุ ธมฺโม อพฺยากตสฺส สเหตุกสฺส เจว น จ เหตุสฺส ธมฺมสฺส อารมฺมณปจฺจเยน ปจฺจโย.\csromanspacer{} ตีณิ. (สํขิตฺตํ.)}

\proseitem{270}{อารมฺมเณ นว, อธิปติยา สตฺต, อนนฺตเร ปญฺจ ฯเปฯ สหชาเต ตีณิ ฯเปฯ อุปนิสฺสเย นว ฯเปฯ อวิคเต ตีณิ. (สํขิตฺตํ.)}

\prose{นเหตุยา นว, น-อารมฺมเณ นว. (สํขิตฺตํ.)}

\csromanheader{2}{\textbf{อารมฺมณ}ปจฺจยา นเหตุยา นว. (สํขิตฺตํ.)}
\csromanheader{2}{นเหตุปจฺจยา อารมฺมเณ นว. (สํขิตฺตํ.)}
\prose{(ยถา กุสลตฺติเก ปญฺหาวารํ, เอวํ วิตฺถาเรตพฺพํ.)}

\csromanmatikaanchor{mtk.140}
\csromantocmarkref{subparagraph}{5. เหตุเหตุสมฺปยุตฺตทุก}{mtk.140}
\bookmark[dest={mtk.140},level=5]{5. เหตุเหตุสมฺปยุตฺตทุก}
\csromanheader{6}{1. กุสลตฺติก 5. เหตุเหตุสมฺปยุตฺตทุก}
\csromanheader{2}{1-7. ปฏิจฺจาทิวาร ปจฺจยจตุกฺก-เหตุ}
\proseitem{271}{กุสลํ เหตุญฺเจว เหตุสมฺปยุตฺตญฺจ ธมฺมํ ปฏิจฺจ กุสโล เหตุ เจว เหตุสมฺปยุตฺโต จ ธมฺโม อุปฺปชฺชติ เหตุปจฺจยา. (1)}

\prose{อกุสลํ เหตุญฺเจว เหตุสมฺปยุตฺตญฺจ ธมฺมํ ปฏิจฺจ อกุสโล เหตุ เจว เหตุสมฺปยุตฺโต จ ธมฺโม อุปฺปชฺชติ เหตุปจฺจยา. (1)}

\prose{อพฺยากตํ เหตุญฺเจว เหตุสมฺปยุตฺตญฺจ ธมฺมํ ปฏิจฺจ อพฺยากโต เหตุ เจว เหตุสมฺปยุตฺโต จ ธมฺโม อุปฺปชฺชติ เหตุปจฺจยา. (สํขิตฺตํ.)}

\csromanheader{2}{1. กุสลตฺติก 5. เหตุเหตุสมฺปยุตฺตทุก}
\csromanheader{2}{272. เหตุยา ตีณิ ฯเปฯ อวิคเต ตีณิ. (สํขิตฺตํ.)}
\csromanheader{2}{น-อธิปติยา ตีณิ. (สํขิตฺตํ.)}
\csromanheader{2}{เหตุปจฺจยา น-อธิปติยา ตีณิ. (สํขิตฺตํ.)}
\prose{\csromanfolio{563}น-อธิปติปจฺจยา เหตุยา ตีณิ. (สํขิตฺตํ.) (สหชาตวารมฺปิ ฯเปฯ สมฺปยุตฺตวารมฺปิ ปฏิจฺจวารสทิสํ.)}

\csromanheader{2}{\textbf{เหตุ อารมฺมณ}}
\proseitem{273}{กุสโล เหตุ เจว เหตุสมฺปยุตฺโต จ ธมฺโม กุสลสฺส เหตุสฺส เจว เหตุสมฺปยุตฺตสฺส จ ธมฺมสฺส เหตุปจฺจเยน ปจฺจโย. (1)}

\prose{อกุสโล เหตุ เจว เหตุสมฺปยุตฺโต จ ธมฺโม อกุสลสฺส เหตุสฺส เจว เหตุสมฺปยุตฺตสฺส จ ธมฺมสฺส เหตุปจฺจเยน ปจฺจโย. (1)}

\prose{อพฺยากโต เหตุ เจว เหตุสมฺปยุตฺโต จ ธมฺโม อพฺยากตสฺส เหตุสฺส เจว เหตุสมฺปยุตฺตสฺส จ ธมฺมสฺส เหตุปจฺจเยน ปจฺจโย. (1)}

\prose{กุสโล เหตุ เจว เหตุสมฺปยุตฺโต จ ธมฺโม กุสลสฺส เหตุสฺส เจว เหตุสมฺปยุตฺตสฺส จ ธมฺมสฺส อารมฺมณปจฺจเยน ปจฺจโย. (สํขิตฺตํ.)}

\proseitem{274}{เหตุยา ตีณิ, อารมฺมเณ นว, อธิปติยา สตฺต, อนนฺตเร ปญฺจ, อุปนิสฺสเย นว ฯเปฯ อวิคเต ตีณิ. (สํขิตฺตํ.)}

\prose{นเหตุยา นว, น-อารมฺมเณ นว. (สํขิตฺตํ.)}

\csromanheader{2}{เหตุปจฺจยา น-อารมฺมเณ ตีณิ. (สํขิตฺตํ.)}
\csromanheader{2}{นเหตุปจฺจยา อารมฺมเณ นว. (สํขิตฺตํ.)}
\prose{(ยถา กุสลตฺติเก ปญฺหาวารํ, เอวํ วิตฺถาเรตพฺพํ.)}

\csromanheader{2}{ธมฺมานุโลม ติกทุกปฏฺฐานปาฬิ \textbf{เหตุสมฺปยุตฺตนเหตุปท 1-7. ปฏิจฺจาทิวาร}}
\csromanheader{2}{\textbf{ปจฺจยจตุกฺก-เหตุ}}
\proseitem{275}{\csromanfolio{564}กุสลํ เหตุสมฺปยุตญฺเจว น จ เหตุํ ธมฺมํ ปฏิจฺจ กุสโล เหตุสมฺปยุตฺโต เจว น จ เหตุ ธมฺโม อุปฺปชฺชติ เหตุปจฺจยา. (1)}

\prose{อกุสลํ เหตุสมฺปยุตฺตญฺเจว น จ เหตุํ ธมฺมํ ปฏิจฺจ อกุสโล เหตุสมฺปยุตฺโต เจว น จ เหตุ ธมฺโม อุปฺปชฺชติ เหตุปจฺจยา. (1)}

\prose{อพฺยากตํ เหตุสมฺปยุตฺตญฺเจว น จ เหตุํ ธมฺมํ ปฏิจฺจ อพฺยากโต เหตุสมฺปยุตฺโต เจว น จ เหตุ ธมฺโม อุปฺปชฺชติ เหตุปจฺจยา. (สํขิตฺตํ.)}

\proseitem{276}{เหตุยา ตีณิ, อารมฺมเณ ตีณิ ฯเปฯ อวิคเต ตีณิ.}

\prose{น-อธิปติยา ตีณิ. (สํขิตฺตํ.) (สหชาตวารมฺปิ ฯเปฯ สมฺปยุตฺตวารมฺปิ ปฏิจฺจวารสทิสํ.)}

\csromanheader{2}{\textbf{อารมฺมณ}}
\proseitem{277}{กุสโล เหตุสมฺปยุตฺโต เจว น จ เหตุ ธมฺโม กุสลสฺส เหตุสมฺปยุตฺตสฺส เจว น จ เหตุสฺส ธมฺมสฺส อารมฺมณปจฺจเยน ปจฺจโย.\csromanspacer{} ตีณิ.}

\prose{อกุสโล เหตุสมฺปยุตฺโต เจว น จ เหตุ ธมฺโม อกุสลสฺส เหตุสมฺปยุตฺตสฺส เจว น จ เหตุสฺส ธมฺมสฺส อารมฺมณปจฺจเยน ปจฺจโย.\csromanspacer{} ตีณิ.}

\prose{อพฺยากโต เหตุสมฺปยุตฺโต เจว น จ เหตุ ธมฺโม อพฺยากตสฺส เหตุสมฺปยุตฺตสฺส เจว น จ เหตุสฺส ธมฺมสฺส อารมฺมณปจฺจเยน ปจฺจโย.\csromanspacer{} ตีณิ. (สํขิตฺตํ.)}

\proseitem{278}{อารมฺมเณ นว, อธิปติยา สตฺต, อนนฺตเร ปญฺจ ฯเปฯ สหชาเต ตีณิ ฯเปฯ อุปนิสฺสเย นว ฯเปฯ อวิคเต ตีณิ. (สํขิตฺตํ.)}

\csromanmatikaanchor{mtk.141}
\csromantocmarkref{subparagraph}{6. นเหตุสเหตุกทุก}{mtk.141}
\bookmark[dest={mtk.141},level=5]{6. นเหตุสเหตุกทุก}
\csromanheader{6}{1. กุสลตฺติก 6. นเหตุสเหตุกทุก}
\prose{\csromanfolio{565}นเหตุยา นว, น-อารมฺมเณ นว. (สํขิตฺตํ.)}

\csromanheader{2}{\textbf{อารมฺมณ}ปจฺจยา นเหตุยา นว. (สํขิตฺตํ.)}
\csromanheader{2}{นเหตุปจฺจยา อารมฺมเณ นว. (สํขิตฺตํ.)}
\prose{(ยถา กุสลตฺติเก ปญฺหาวารํ, เอวํ วิตฺถาเรตพฺพํ.)}

\csromanheader{2}{1. กุสลตฺติก 6. นเหตุสเหตุกทุก}
\csromanheader{2}{1-7. ปฏิจฺจาทิวาร ปจฺจยจตุกฺก-เหตุ}
\proseitem{279}{กุสลํ นเหตุํ สเหตุกํ ธมฺมํ ปฏิจฺจ กุสโล นเหตุ สเหตุโก ธมฺโม อุปฺปชฺชติ เหตุปจฺจยา. (1)}

\prose{อกุสลํ นเหตุํ สเหตุกํ ธมฺมํ ปฏิจฺจ อกุสโล นเหตุ สเหตุโก ธมฺโม อุปฺปชฺชติ เหตุปจฺจยา. (1)}

\prose{อพฺยากตํ นเหตุํ สเหตุกํ ธมฺมํ ปฏิจฺจ อพฺยากโต นเหตุ สเหตุโก ธมฺโม อุปฺปชฺชติ เหตุปจฺจยา. (1) (สํขิตฺตํ.)}

\proseitem{280}{เหตุยา ตีณิ, อารมฺมเณ ตีณิ ฯเปฯ อวิคเต ตีณิ. (สํขิตฺตํ.)}

\prose{น-อธิปติยา ตีณิ. (สํขิตฺตํ.) (สหชาตวารมฺปิ ฯเปฯ สมฺปยุตฺตวารมฺปิ ปฏิจฺจวารสทิสํ.)}

\csromanheader{2}{\textbf{อารมฺมณ}}
\proseitem{281}{กุสโล นเหตุ สเหตุโก ธมฺโม กุสลสฺส นเหตุสฺส สเหตุกสฺส ธมฺมสฺส อารมฺมณปจฺจเยน ปจฺจโย.\csromanspacer{} ตีณิ.}

\prose{อกุสโล นเหตุ สเหตุโก ธมฺโม อกุสลสฺส นเหตุสฺส สเหตุกสฺส ธมฺมสฺส อารมฺมณปจฺจเยน ปจฺจโย.\csromanspacer{} ตีณิ.}

\prose{อพฺยากโต นเหตุ สเหตุโก ธมฺโม อพฺยากตสฺส นเหตุสฺส สเหตุกสฺส ธมฺมสฺส อารมฺมณปจฺจเยน ปจฺจโย.\csromanspacer{} ตีณิ. (สํขิตฺตํ.)}

\gambhira{ธมฺมานุโลม ติกทุกปฏฺฐานปาฬิ}
\csromanmatikaanchor{mtk.142}
\csromantocmarkref{subparagraph}{7-13. สปฺปจฺจยทุกาทิ}{mtk.142}
\bookmark[dest={mtk.142},level=5]{7-13. สปฺปจฺจยทุกาทิ}
\proseitem{282}{\csromanfolio{566}อารมฺมเณ นว, อธิปติยา สตฺต, อนนฺตเร ปญฺจ ฯเปฯ สหชาเต อญฺญมญฺเญ นิสฺสเย ตีณิ, อุปนิสฺสเย นว, อาเสวเน ตีณิ, กมฺเม ปญฺจ, วิปาเก เอกํ, อาหาเร ฯเปฯ สมฺปยุตฺเต ตีณิ ฯเปฯ อวิคเต ตีณิ. (สํขิตฺตํ.)}

\prose{นเหตุยา นว, น-อารมฺมเณ นว. (สํขิตฺตํ.)}

\csromanheader{2}{อารมฺมณปจฺจยา นเหตุยา นว. (สํขิตฺตํ.)}
\prose{นเหตุปจฺจยา อารมฺมเณ นว. (สํขิตฺตํ.) (ยถา กุสลตฺติเก ปญฺหาวารํ, เอวํ วิตฺถาเรตพฺพํ.)}

\csromanheader{2}{\textbf{นเหตุ-อเหตุกปท-เหตุ}}
\proseitem{283}{อพฺยากตํ นเหตุํ อเหตุกํ ธมฺมํ ปฏิจฺจ อพฺยากโต นเหตุ อเหตุโก ธมฺโม อุปฺปชฺชติ เหตุปจฺจยา. (สํขิตฺตํ.)}

\prose{เหตุยา เอกํ, อารมฺมเณ เอกํ ฯเปฯ อวิคเต เอกํ. (สํขิตฺตํ.)}

\prose{นเหตุยา เอกํ ฯเปฯ โนวิคเต เอกํ. (สํขิตฺตํ.) (สหชาตวารมฺปิ ฯเปฯ สมฺปยุตฺตวารมฺปิ ปฏิจฺจวารสทิสํ.)}

\proseitem{284}{อพฺยากโต นเหตุ อเหตุโก ธมฺโม อพฺยากตสฺส นเหตุสฺส อเหตุกสฺส ธมฺมสฺส อารมฺมณปจฺจเยน ปจฺจโย. (สํขิตฺตํ.)}

\prose{อารมฺมเณ เอกํ, อธิปติยา เอกํ ฯเปฯ อวิคเต เอกํ. (สํขิตฺตํ.) (ยถา กุสลตฺติเก ปญฺหาวารํ, เอวํ วิตฺถาเรตพฺพํ.)}

\nitthitam{เหตุโคจฺฉกํ นิฏฺฐิตํ.}
\csromanheader{6}{1. กุสลตฺติก 7. สปฺปจฺจยทุก}
\csromanheader{2}{1-7. ปฏิจฺจาทิวาร ปจฺจยจตุกฺก}
\proseitem{1}{กุสลํ สปฺปจฺจยํ ธมฺมํ ปฏิจฺจ กุสโล สปฺปจฺจโย ธมฺโม อุปฺปชฺชติ เหตุปจฺจยา.\csromanspacer{}\csromanspacer{}กุสลํ สปฺปจฺจยํ ธมฺมํ ปฏิจฺจ อพฺยากโต สปฺปจฺจโย ธมฺโม อุปฺปชฺชติ เหตุปจฺจยา.\csromanspacer{}\csromanspacer{}กุสลํ สปฺปจฺจยํ ธมฺมํ ปฏิจฺจ กุสโล}

\vspace{\tipitakaparskip}
\csromancenter{กุสลตฺติก 7.\csromanspacer{} สปฺปจฺจยทุก สปฺปจฺจโย จ อพฺยากโต สปฺปจฺจโย จ ธมฺมา อุปฺปชฺชนฺติ เหตุปจฺจยา. (3)}
\prose{\csromanfolio{567}อกุสลํ สปฺปจฺจยํ ธมฺมํ ปฏิจฺจ อกุสโล สปฺปจฺจโย ธมฺโม อุปฺปชฺชติ เหตุปจฺจยา.\csromanspacer{} ตีณิ.}

\prose{อพฺยากตํ สปฺปจฺจยํ ธมฺมํ ปฏิจฺจ อพฺยากโต สปฺปจฺจโย ธมฺโม อุปฺปชฺชติ เหตุปจฺจยา. (1)}

\prose{กุสลํ สปฺปจฺจยญฺจ อพฺยากตํ สปฺปจฺจยญฺจ ธมฺมํ ปฏิจฺจ อพฺยากโต สปฺปจฺจโย ธมฺโม อุปฺปชฺชติ เหตุปจฺจยา. (1)}

\prose{อกุสลํ สปฺปจฺจยญฺจ อพฺยากตํ สปฺปจฺจยญฺจ ธมฺมํ ปฏิจฺจ อพฺยากโต สปฺปจฺจโย ธมฺโม อุปฺปชฺชติ เหตุปจฺจยา. (1)}

\proseitem{2}{กุสลํ สปฺปจฺจยํ ธมฺมํ ปฏิจฺจ กุสโล สปฺปจฺจโย ธมฺโม อุปฺปชฺชติ อารมฺมณปจฺจยา. (1)}

\prose{อกุสลํ สปฺปจฺจยํ ธมฺมํ ปฏิจฺจ อกุสโล สปฺปจฺจโย ธมฺโม อุปฺปชฺชติ อารมฺมณปจฺจยา. (1)}

\prose{อพฺยากตํ สปฺปจฺจยํ ธมฺมํ ปฏิจฺจ อพฺยากโต สปฺปจฺจโย ธมฺโม อุปฺปชฺชติ อารมฺมณปจฺจยา. (1) (สํขิตฺตํ.)}

\proseitem{3}{เหตุยา นว, อารมฺมเณ ตีณิ, อธิปติยา นว ฯเปฯ อวิคเต นว. (สํขิตฺตํ.)}

\csromanheader{2}{\textbf{ปจฺจนีย-นเหตุ}}
\proseitem{4}{อกุสลํ สปฺปจฺจยํ ธมฺมํ ปฏิจฺจ อกุสโล สปฺปจฺจโย ธมฺโม อุปฺปชฺชติ นเหตุปจฺจยา. (1)}

\prose{อพฺยากตํ สปฺปจฺจยํ ธมฺมํ ปฏิจฺจ อพฺยากโต สปฺปจฺจโย ธมฺโม อุปฺปชฺชติ นเหตุปจฺจยา. (1) (สํขิตฺตํ.)}

\proseitem{5}{นเหตุยา เทฺว, น-อารมฺมเณ ปญฺจ, น-อธิปติยา นว ฯเปฯ โนวิคเต ปญฺจ. (สํขิตฺตํ.)}

\gambhira{ธมฺมานุโลม ติกทุกปฏฺฐานปาฬิ}
\csromanheader{2}{เหตุปจฺจยา น-อารมฺมเณ ปญฺจ. (สํขิตฺตํ.)}
\prose{\csromanfolio{568}นเหตุปจฺจยา อารมฺมเณ เทฺว. (สํขิตฺตํ.) (สหชาตวารมฺปิ ฯเปฯ สมฺปยุตฺตวารมฺปิ ปฏิจฺจวารสทิสํ.)}

\csromanheader{2}{\textbf{เหตุ อารมฺมณ}}
\proseitem{6}{กุสโล สปฺปจฺจโย ธมฺโม กุสลสฺส สปฺปจฺจยสฺส ธมฺมสฺส เหตุปจฺจเยน ปจฺจโย.\csromanspacer{} ตีณิ.}

\prose{อกุสโล สปฺปจฺจโย ธมฺโม อกุสลสฺส สปฺปจฺจยสฺส ธมฺมสฺส เหตุปจฺจเยน ปจฺจโย.\csromanspacer{} ตีณิ.}

\prose{อพฺยากโต สปฺปจฺจโย ธมฺโม อพฺยากตสฺส สปฺปจฺจยสฺส ธมฺมสฺส เหตุปจฺจเยน ปจฺจโย. (1)}

\proseitem{7}{กุสโล สปฺปจฺจโย ธมฺโม กุสลสฺส สปฺปจฺจยสฺส ธมฺมสฺส อารมฺมณปจฺจเยน ปจฺจโย.\csromanspacer{} ตีณิ.}

\prose{อกุสโล สปฺปจฺจโย ธมฺโม อกุสลสฺส สปฺปจฺจยสฺส ธมฺมสฺส อารมฺมณปจฺจเยน ปจฺจโย.\csromanspacer{} ตีณิ.}

\prose{อพฺยากโต สปฺปจฺจโย ธมฺโม อพฺยากตสฺส สปฺปจฺจยสฺส ธมฺมสฺส อารมฺมณปจฺจเยน ปจฺจโย.\csromanspacer{} ตีณิ. (สํขิตฺตํ.)}

\proseitem{8}{เหตุยา สตฺต, อารมฺมเณ นว, อธิปติยา ทส, อนนฺตเร สตฺต ฯเปฯ สหชาเต นว ฯเปฯ อุปนิสฺสเย นว ฯเปฯ อวิคเต เตรส. (สํขิตฺตํ.)}

\prose{นเหตุยา ปนฺนรส, น-อารมฺมเณ ปนฺนรส. (สํขิตฺตํ.)}

\csromanheader{2}{เหตุปจฺจยา น-อารมฺมเณ สตฺต. (สํขิตฺตํ.)}
\csromanheader{2}{นเหตุปจฺจยา อารมฺมเณ นว. (สํขิตฺตํ.)}
\prose{(ยถา กุสลตฺติเก ปญฺหาวารํ, เอวํ วิตฺถาเรตพฺพํ.)}

\vspace{\tipitakaparskip}
\csromancenter{กุสลตฺติก 9.\csromanspacer{} สนิทสฺสนทุก \textbf{1.}\csromanspacer{}\textbf{ กุสลตฺติก 8.}\csromanspacer{}\textbf{ สงฺขตทุก}}
\csromanheader{2}{1. ปฏิจฺจวาร ปจฺจยจตุกฺก}
\proseitem{9}{\csromanfolio{569}กุสลํ สงฺขตํ ธมฺมํ ปฏิจฺจ กุสโล สงฺขโต ธมฺโม อุปฺปชฺชติ เหตุปจฺจยา. (สปฺปจฺจยทุกสทิสํ.)}
\csromansectionrule

\csromanheader{2}{1. กุสลตฺติก 9. สนิทสฺสนทุก}
\csromanheader{2}{1-7. ปฏิจฺจาทิวาร ปจฺจยจตุกฺก}
\proseitem{10}{กุสลํ อนิทสฺสนํ ธมฺมํ ปฏิจฺจ กุสโล อนิทสฺสโน ธมฺโม อุปฺปชฺชติ เหตุปจฺจยา.\csromanspacer{} ตีณิ.}

\prose{อกุสลํ อนิทสฺสนํ ธมฺมํ ปฏิจฺจ อกุสโล อนิทสฺสโน ธมฺโม อุปฺปชฺชติ เหตุปจฺจยา.\csromanspacer{} ตีณิ.}

\prose{อพฺยากตํ อนิทสฺสนํ ธมฺมํ ปฏิจฺจ อพฺยากโต อนิทสฺสโน ธมฺโม อุปฺปชฺชติ เหตุปจฺจยา. (1)}

\prose{กุสลํ อนิทสฺสนญฺจ อพฺยากตํ อนิทสฺสนญฺจ ธมฺมํ ปฏิจฺจ อพฺยากโต อนิทสฺสโน ธมฺโม อุปฺปชฺชติ เหตุปจฺจยา. (1)}

\prose{อกุสลํ อนิทสฺสนญฺจ อพฺยากตํ อนิทสฺสนญฺจ ธมฺมํ ปฏิจฺจ อพฺยากโต อนิทสฺสโน ธมฺโม อุปฺปชฺชติ เหตุปจฺจยา. (1)}

\proseitem{11}{กุสลํ อนิทสฺสนํ ธมฺมํ ปฏิจฺจ กุสโล อนิทสฺสโน ธมฺโม อุปฺปชฺชติ อารมฺมณปจฺจยา. (1)}

\prose{อกุสลํ อนิทสฺสนํ ธมฺมํ ปฏิจฺจ อกุสโล อนิทสฺสโน ธมฺโม อุปฺปชฺชติ อารมฺมณปจฺจยา. (1)}

\prose{อพฺยากตํ อนิทสฺสนํ ธมฺมํ ปฏิจฺจ อพฺยากโต อนิทสฺสโน ธมฺโม อุปฺปชฺชติ อารมฺมณปจฺจยา. (1) (สํขิตฺตํ.)}

\proseitem{12}{เหตุยา นว, อารมฺมเณ ตีณิ ฯเปฯ อวิคเต นว. (สํขิตฺตํ.)}

\csromanheader{2}{ธมฺมานุโลม ติกทุกปฏฺฐานปาฬิ \textbf{ปจฺจนีย-นเหตุ}}
\proseitem{13}{\csromanfolio{570}อกุสลํ อนิทสฺสนํ ธมฺมํ ปฏิจฺจ อกุสโล อนิทสฺสโน ธมฺโม อุปฺปชฺชติ นเหตุปจฺจยา. (1)}

\prose{อพฺยากตํ อนิทสฺสนํ ธมฺมํ ปฏิจฺจ อพฺยากโต อนิทสฺสโน ธมฺโม อุปฺปชฺชติ นเหตุปจฺจยา. (1) (สํขิตฺตํ.)}

\prose{นเหตุยา เทฺว, น-อารมฺมเณ ปญฺจ, น-อธิปติยา นว ฯเปฯ โนวิคเต ปญฺจ. (สํขิตฺตํ.\csromanspacer{} สหชาตวารมฺปิ ฯเปฯ สมฺปยุตฺตวารมฺปิ ปฏิจฺจวารสทิสํ.)}

\csromanheader{2}{\textbf{เหตุ อารมฺมณ}}
\proseitem{14}{กุสโล อนิทสฺสโน ธมฺโม กุสลสฺส อนิทสฺสนสฺส ธมฺมสฺส เหตุปจฺจเยน ปจฺจโย.\csromanspacer{} ตีณิ.}

\prose{อกุสโล อนิทสฺสโน ธมฺโม อกุสลสฺส อนิทสฺสนสฺส ธมฺมสฺส เหตุปจฺจเยน ปจฺจโย.\csromanspacer{} ตีณิ.}

\prose{อพฺยากโต อนิทสฺสโน ธมฺโม อพฺยากตสฺส อนิทสฺสนสฺส ธมฺมสฺส เหตุปจฺจเยน ปจฺจโย. (1)}

\proseitem{15}{กุสโล อนิทสฺสโน ธมฺโม กุสลสฺส อนิทสฺสนสฺส ธมฺมสฺส อารมฺมณปจฺจเยน ปจฺจโย.\csromanspacer{} ตีณิ.}

\prose{อกุสโล อนิทสฺสโน ธมฺโม อกุสลสฺส อนิทสฺสนสฺส ธมฺมสฺส อารมฺมณปจฺจเยน ปจฺจโย.\csromanspacer{} ตีณิ.}

\prose{อพฺยากโต อนิทสฺสโน ธมฺโม อพฺยากตสฺส อนิทสฺสนสฺส ธมฺมสฺส อารมฺมณปจฺจเยน ปจฺจโย.\csromanspacer{} ตีณิ. (สํขิตฺตํ.)}

\proseitem{16}{เหตุยา สตฺต, อารมฺมเณ นว, อธิปติยา ทส, อนนฺตเร สตฺต ฯเปฯ อุปนิสฺสเย นว ฯเปฯ อวิคเต เตรส. (สํขิตฺตํ.)}

\prose{(ยถา กุสลตฺติเก ปญฺหาวารํ, เอวํ วิตฺถาเรตพฺพํ.)}

\csromanheader{2}{1. กุสลตฺติก 10. สปฺปฏิฆทุก}
\csromanheader{2}{1-7. ปฏิจฺจาทิวาร ปจฺจยจตุกฺก-เหตุ}
\proseitem{17}{อพฺยากตํ สปฺปฏิฆํ ธมฺมํ ปฏิจฺจ อพฺยากโต สปฺปฏิโฆ ธมฺโม อุปฺปชฺชติ เหตุปจฺจยา. (สํขิตฺตํ.)}

\csromanheader{2}{1. กุสลตฺติก 10. สปฺปฏิฆทุก}
\prose{\csromanfolio{571}เหตุยา เอกํ ฯเปฯ อวิคเต เอกํ. (สํขิตฺตํ.) (สหชาตวาเรปิ ฯเปฯ ปญฺหาวาเรปิ สพฺพตฺถ เอกํ.)}

\proseitem{18}{กุสลํ อปฺปฏิฆํ ธมฺมํ ปฏิจฺจ กุสโล อปฺปฏิโฆ ธมฺโม อุปฺปชฺชติ เหตุปจฺจยา.\csromanspacer{} ตีณิ.}

\prose{อกุสลํ อปฺปฏิฆํ ธมฺมํ ปฏิจฺจ อกุสโล อปฺปฏิโฆ ธมฺโม อุปฺปชฺชติ เหตุปจฺจยา.\csromanspacer{} ตีณิ.}

\prose{อพฺยากตํ อปฺปฏิฆํ ธมฺมํ ปฏิจฺจ อพฺยากโต อปฺปฏิโฆ ธมฺโม อุปฺปชฺชติ เหตุปจฺจยา. (1)}

\prose{กุสลํ อปฺปฏิฆญฺจ อพฺยากตํ อปฺปฏิฆญฺจ ธมฺมํ ปฏิจฺจ อพฺยากโต อปฺปฏิโฆ ธมฺโม อุปฺปชฺชติ เหตุปจฺจยา. (1)}

\prose{อกุสลํ อปฺปฏิฆญฺจ อพฺยากตํ อปฺปฏิฆญฺจ ธมฺมํ ปฏิจฺจ อพฺยากโต อปฺปฏิโฆ ธมฺโม อุปฺปชฺชติ เหตุปจฺจยา. (1)}

\csromanheader{2}{\textbf{อารมฺมณ}}
\proseitem{19}{กุสลํ อปฺปฏิฆํ ธมฺมํ ปฏิจฺจ กุสโล อปฺปฏิโฆ ธมฺโม อุปฺปชฺชติ อารมฺมณปจฺจยา. (1)}

\prose{อกุสลํ อปฺปฏิฆํ ธมฺมํ ปฏิจฺจ อกุสโล อปฺปฏิโฆ ธมฺโม อุปฺปชฺชติ อารมฺมณปจฺจยา. (1)}

\prose{อพฺยากตํ อปฺปฏิฆํ ธมฺมํ ปฏิจฺจ อพฺยากโต อปฺปฏิโฆ ธมฺโม อุปฺปชฺชติ อารมฺมณปจฺจยา. (1) (สํขิตฺตํ.)}

\proseitem{20}{เหตุยา นว, อารมฺมเณ ตีณิ, อธิปติยา นว ฯเปฯ อวิคเต นว. (สํขิตฺตํ.)}

\csromanheader{2}{\textbf{ปจฺจนีย-นเหตุ}}
\proseitem{21}{อกุสลํ อปฺปฏิฆํ ธมฺมํ ปฏิจฺจ อกุสโล อปฺปฏิโฆ ธมฺโม อุปฺปชฺชติ นเหตุปจฺจยา. (1)}

\prose{อพฺยากตํ อปฺปฏิฆํ ธมฺมํ ปฏิจฺจ อพฺยากโต อปฺปฏิโฆ ธมฺโม อุปฺปชฺชติ นเหตุปจฺจยา. (1) (สํขิตฺตํ.)}

\gambhira{ธมฺมานุโลม ติกทุกปฏฺฐานปาฬิ}
\proseitem{22}{\csromanfolio{572}นเหตุยา เทฺว, น-อารมฺมเณ ปญฺจ, น-อธิปติยา นว ฯเปฯ โนวิคเต ปญฺจ. (สํขิตฺตํ.\csromanspacer{} สหชาตวารมฺปิ ฯเปฯ สมฺปยุตฺตวารมฺปิ ปฏิจฺจวารสทิสํ.)}

\proseitem{23}{กุสโล อปฺปฏิโฆ ธมฺโม กุสลสฺส อปฺปฏิฆสฺส ธมฺมสฺส เหตุปจฺจเยน ปจฺจโย.\csromanspacer{} ตีณิ.}

\prose{อกุสโล อปฺปฏิโฆ ธมฺโม อกุสลสฺส อปฺปฏิฆสฺส ธมฺมสฺส เหตุปจฺจเยน ปจฺจโย.\csromanspacer{} ตีณิ.}

\prose{อพฺยากโต อปฺปฏิโฆ ธมฺโม อพฺยากตสฺส อปฺปฏิฆสฺส ธมฺมสฺส เหตุปจฺจเยน ปจฺจโย. (1) (สํขิตฺตํ.)}

\vspace{\tipitakaparskip}
\csromancenter{\textbf{เหตุ}ยา สตฺต, อารมฺมเณ นว, อธิปติยา ทส ฯเปฯ อวิคเต เตรส. (สํขิตฺตํ.\csromanspacer{} ยถา กุสลตฺติเก ปญฺหาวารํ, เอวํ วิตฺถาเรตพฺพํ.)}
\csromanheader{2}{1. กุสลตฺติก 11. รูปีทุก}
\csromanheader{2}{1-7. ปฏิจฺจาทิวาร ปจฺจยจตุกฺก-เหตุ}
\proseitem{25}{อพฺยากตํ รูปิํ ธมฺมํ ปฏิจฺจ อพฺยากโต รูปี ธมฺโม อุปฺปชฺชติ เหตุปจฺจยา. (สํขิตฺตํ.)}

\vspace{\tipitakaparskip}
\csromancenter{\textbf{เหตุ}ยา เอกํ, อารมฺมเณ เอกํ ฯเปฯ อวิคเต เอกํ. (สํขิตฺตํ.) (สหชาตวาเรปิ ฯเปฯ ปญฺหาวาเรปิ สพฺพตฺถ เอกํ.)}
\proseitem{26}{กุสลํ อรูปิํ ธมฺมํ ปฏิจฺจ กุสโล อรูปี ธมฺโม อุปฺปชฺชติ เหตุปจฺจยา. (1)}

\prose{อกุสลํ อรูปิํ ธมฺมํ ปฏิจฺจ อกุสโล อรูปี ธมฺโม อุปฺปชฺชติ เหตุปจฺจยา. (1)}

\prose{อพฺยากตํ อรูปิํ ธมฺมํ ปฏิจฺจ อพฺยากโต อรูปี ธมฺโม อุปฺปชฺชติ เหตุปจฺจยา. (1) (สํขิตฺตํ.)}

\csromanheader{2}{1. กุสลตฺติก 12. โลกิยทุก}
\proseitem{27}{\csromanfolio{573}เหตุยา ตีณิ, อารมฺมเณ ตีณิ ฯเปฯ อวิคเต ตีณิ. (สํขิตฺตํ.)}

\prose{นเหตุยา เทฺว, น-อธิปติยา ตีณิ. (สํขิตฺตํ.)}

\prose{(สหชาตวารมฺปิ ฯเปฯ สมฺปยุตฺตวารมฺปิ ปฏิจฺจวารสทิสํ.)}

\csromanheader{2}{\textbf{เหตุ อารมฺมณ}}
\proseitem{28}{กุสโล อรูปี ธมฺโม กุสลสฺส อรูปิสฺส ธมฺมสฺส เหตุปจฺจเยน ปจฺจโย. (1)}

\prose{อกุสโล อรูปี ธมฺโม อกุสลสฺส อรูปิสฺส ธมฺมสฺส เหตุปจฺจเยน ปจฺจโย. (1)}

\prose{อพฺยากโต อรูปี ธมฺโม อพฺยากตสฺส อรูปิสฺส ธมฺมสฺส เหตุปจฺจเยน ปจฺจโย. (1)}

\prose{กุสโล อรูปี ธมฺโม กุสลสฺส อรูปิสฺส ธมฺมสฺส อารมฺมณปจฺจเยน ปจฺจโย.\csromanspacer{} ตีณิ.}

\prose{อกุสโล อรูปี ธมฺโม อกุสลสฺส อรูปิสฺส ธมฺมสฺส อารมฺมณปจฺจเยน ปจฺจโย.\csromanspacer{} ตีณิ.}

\prose{อพฺยากโต อรูปี ธมฺโม อพฺยากตสฺส อรูปิสฺส ธมฺมสฺส อารมฺมณปจฺจเยน ปจฺจโย. (สํขิตฺตํ.)}

\proseitem{29}{เหตุยา ตีณิ, อารมฺมเณ นว, อธิปติยา สตฺต ฯเปฯ อวิคเต ตีณิ. (สํขิตฺตํ.)}

\prose{(ยถา กุสลตฺติเก ปญฺหาวารํ, เอวํ วิตฺถาเรตพฺพํ.)}

\csromanheader{2}{1. กุสลตฺติก 12. โลกิยทุก}
\csromanheader{2}{1-7. ปฏิจฺจาทิวาร ปจฺจยจตุกฺก-เหตุ}
\proseitem{30}{กุสลํ โลกิยํ ธมฺมํ ปฏิจฺจ กุสโล โลกิโย ธมฺโม อุปฺปชฺชติ เหตุปจฺจยา.\csromanspacer{}\csromanspacer{}กุสลํ โลกิยํ ธมฺมํ ปฏิจฺจ อพฺยากโต}

\vspace{\tipitakaparskip}
\csromancenter{ธมฺมานุโลม ติกทุกปฏฺฐานปาฬิ โลกิโย ธมฺโม อุปฺปชฺชติ เหตุปจฺจยา.\csromanspacer{}\csromanspacer{}กุสลํ โลกิยํ ธมฺมํ ปฏิจฺจ กุสโล โลกิโย จ อพฺยากโต โลกิโย จ ธมฺมา อุปฺปชฺชนฺติ เหตุปจฺจยา. (3)}
\prose{\csromanfolio{574}อกุสลํ โลกิยํ ธมฺมํ ปฏิจฺจ อกุสโล โลกิโย ธมฺโม อุปฺปชฺชติ เหตุปจฺจยา.\csromanspacer{} ตีณิ.}

\prose{อพฺยากตํ โลกิยํ ธมฺมํ ปฏิจฺจ อพฺยากโต โลกิโย ธมฺโม อุปฺปชฺชติ เหตุปจฺจยา. (1)}

\prose{กุสลํ โลกิยญฺจ อพฺยากตํ โลกิยญฺจ ธมฺมํ ปฏิจฺจ อพฺยากโต โลกิโย ธมฺโม อุปฺปชฺชติ เหตุปจฺจยา. (1)}

\prose{อกุสโล โลกิยญฺจ อพฺยากตํ โลกิยญฺจ ธมฺมํ ปฏิจฺจ อพฺยากโต โลกิโย ธมฺโม อุปฺปชฺชติ เหตุปจฺจยา. (1) (สํขิตฺตํ.)}

\proseitem{31}{เหตุยา นว, อารมฺมเณ ตีณิ ฯเปฯ อวิคเต นว. (สํขิตฺตํ.)}

\csromanheader{2}{\textbf{ปจฺจนีย-นเหตุ}}
\proseitem{32}{อกุสลํ โลกิยํ ธมฺมํ ปฏิจฺจ อกุสโล โลกิโย ธมฺโม อุปฺปชฺชติ นเหตุปจฺจยา. (1)}

\prose{อพฺยากตํ โลกิยํ ธมฺมํ ปฏิจฺจ อพฺยากโต โลกิโย ธมฺโม อุปฺปชฺชติ นเหตุปจฺจยา. (1) (สํขิตฺตํ.)}

\proseitem{33}{นเหตุยา เทฺว, น-อารมฺมเณ ปญฺจ, น-อธิปติยา นว ฯเปฯ โนวิคเต ปญฺจ. (สํขิตฺตํ.\csromanspacer{} สหชาตวารมฺปิ ฯเปฯ สมฺปยุตฺตวารมฺปิ ปฏิจฺจวารสทิสํ.)}

\csromanheader{2}{\textbf{เหตุ อารมฺมณ}}
\proseitem{34}{กุสโล โลกิโย ธมฺโม กุสลสฺส โลกิยสฺส ธมฺมสฺส เหตุปจฺจเยน ปจฺจโย.\csromanspacer{} ตีณิ.}

\prose{อกุสโล โลกิโย ธมฺโม อกุสลสฺส โลกิยสฺส ธมฺมสฺส เหตุปจฺจเยน ปจฺจโย.\csromanspacer{} ตีณิ.}

\prose{อพฺยากโต โลกิโย ธมฺโม อพฺยากตสฺส โลกิยสฺส ธมฺมสฺส เหตุปจฺจเยน ปจฺจโย. (1)}

\csromanheader{2}{1. กุสลตฺติก 12. โลกิยทุก}
\proseitem{35}{\csromanfolio{575}กุสโล โลกิโย ธมฺโม กุสลสฺส โลกิยสฺส ธมฺมสฺส อารมฺมณปจฺจเยน ปจฺจโย.\csromanspacer{} ตีณิ.}

\prose{อกุสโล โลกิโย ธมฺโม อกุสลสฺส โลกิยสฺส ธมฺมสฺส อารมฺมณปจฺจเยน ปจฺจโย.\csromanspacer{} ตีณิ.}

\prose{อพฺยากโต โลกิโย ธมฺโม อพฺยากตสฺส โลกิยสฺส ธมฺมสฺส อารมฺมณปจฺจเยน ปจฺจโย.\csromanspacer{} ตีณิ. (สํขิตฺตํ.)}

\proseitem{36}{เหตุยา สตฺต, อารมฺมเณ นว, อธิปติยา นว ฯเปฯ อวิคเต เตรส. (สํขิตฺตํ.\csromanspacer{} ยถา กุสลตฺติเก ปญฺหาวารํ, เอวํ วิตฺถาเรตพฺพํ.)}

\csromanheader{2}{\textbf{โลกุตฺตรปท 1-7. ปฏิจฺจาทิวาร}}
\csromanheader{2}{\textbf{ปจฺจยจตุกฺก-เหตุ}}
\proseitem{37}{กุสลํ โลกุตฺตรํ ธมฺมํ ปฏิจฺจ กุสโล โลกุตฺตโร ธมฺโม อุปฺปชฺชติ เหตุปจฺจยา. (1)}

\prose{อพฺยากตํ โลกุตฺตรํ ธมฺมํ ปฏิจฺจ อพฺยากโต โลกุตฺตโร ธมฺโม อุปฺปชฺชติ เหตุปจฺจยา. (1) (สํขิตฺตํ.)}

\proseitem{38}{เหตุยา เทฺว, อารมฺมเณ เทฺว, อธิปติยา เทฺว ฯเปฯ อวิคเต เทฺว. (สํขิตฺตํ.)}

\prose{น-อธิปติยา เทฺว. (สํขิตฺตํ.\csromanspacer{} สหชาตวารมฺปิ ฯเปฯ สมฺปยุตฺตวารมฺปิ ปฏิจฺจวารสทิสํ.)}

\csromanheader{2}{\textbf{เหตุ อารมฺมณาทิ}}
\proseitem{39}{กุสโล โลกุตฺตโร ธมฺโม กุสลสฺส โลกุตฺตรสฺส ธมฺมสฺส เหตุปจฺจเยน ปจฺจโย. (1)}

\prose{อพฺยากโต โลกุตฺตโร ธมฺโม อพฺยากตสฺส โลกุตฺตรสฺส ธมฺมสฺส เหตุปจฺจเยน ปจฺจโย. (1)}

\prose{อพฺยากโต โลกุตฺตโร ธมฺโม อพฺยากตสฺส โลกุตฺตรสฺส ธมฺมสฺส อารมฺมณปจฺจเยน ปจฺจโย.\csromanspacer{}\csromanspacer{}อพฺยากโต โลกุตฺตโร ธมฺโม กุสลสฺส โลกุตฺตรสฺส ธมฺมสฺส อารมฺมณปจฺจเยน ปจฺจโย. (2)}

\gambhira{ธมฺมานุโลม ติกทุกปฏฺฐานปาฬิ}
\prose{\csromanfolio{576}กุสโล โลกุตฺตโร ธมฺโม กุสลสฺส โลกุตฺตรสฺส ธมฺมสฺส อธิปติปจฺจเยน ปจฺจโย. (1)}

\prose{อพฺยากโต โลกุตฺตโร ธมฺโม อพฺยากตสฺส โลกุตฺตรสฺส ธมฺมสฺส อธิปติปจฺจเยน ปจฺจโย.\csromanspacer{}\csromanspacer{}อพฺยากโต โลกุตฺตโร ธมฺโม กุสลสฺส โลกุตฺตรสฺส ธมฺมสฺส อธิปติปจฺจเยน ปจฺจโย. (2)}

\prose{กุสโล โลกุตฺตโร ธมฺโม อพฺยากตสฺส โลกุตฺตรสฺส ธมฺมสฺส อนนฺตรปจฺจเยน ปจฺจโย. (1)}

\prose{อพฺยากโต โลกุตฺตโร ธมฺโม อพฺยากตสฺส โลกุตฺตรสฺส ธมฺมสฺส อนนฺตรปจฺจเยน ปจฺจโย. (1) (สํขิตฺตํ.)}

\proseitem{40}{เหตุยา เทฺว, อารมฺมเณ เทฺว, อธิปติยา ตีณิ, อนนฺตเร เทฺว, สมนนฺตเร เทฺว, สหชาเต เทฺว, อญฺญมญฺเญ เทฺว, นิสฺสเย เทฺว, อุปนิสฺสเย จตฺตาริ ฯเปฯ อวิคเต เทฺว. (สํขิตฺตํ.)}

\prose{(ยถา กุสลตฺติเก ปญฺหาวารํ, เอวํ วิตฺถาเรตพฺพํ.)}

\csromanheader{2}{1. กุสลตฺติก 13. เกนจิวิญฺเญยฺยทุก}
\csromanheader{2}{1-7. ปฏิจฺจาทิวาร ปจฺจยจตุกฺก}
\proseitem{41}{กุสลํ เกนจิ วิญฺเญยฺยํ ธมฺมํ ปฏิจฺจ กุสโล เกนจิ วิญฺเญยฺโย ธมฺโม อุปฺปชฺชติ เหตุปจฺจยา.\csromanspacer{}\csromanspacer{}กุสลํ เกนจิ วิญฺเญยฺยํ ธมฺมํ ปฏิจฺจ อพฺยากโต เกนจิ วิญฺเญยฺโย ธมฺโม อุปฺปชฺชติ เหตุปจฺจยา.\csromanspacer{}\csromanspacer{}กุสลํ เกนจิ วิญฺเญยฺยํ ธมฺมํ ปฏิจฺจ กุสโล เกนจิ วิญฺเญยฺโย จ อพฺยากโต เกนจิ วิญฺเญยฺโย จ ธมฺมา อุปฺปชฺชนฺติ เหตุปจฺจยา. (3)}

\prose{อกุสลํ เกนจิ วิญฺเญยฺยํ ธมฺมํ ปฏิจฺจ อกุสโล เกนจิ วิญฺเญยฺโย ธมฺโม อุปฺปชฺชติ เหตุปจฺจยา.\csromanspacer{} ตีณิ.}

\prose{อพฺยากตํ เกนจิ วิญฺเญยฺยํ ธมฺมํ ปฏิจฺจ อพฺยากโต เกนจิ วิญฺเญยฺโย ธมฺโม อุปฺปชฺชติ เหตุปจฺจยา. (1)}

\prose{กุสลํ เกนจิ วิญฺเญยฺยญฺจ อพฺยากตํ เกนจิ วิญฺเญยฺยญฺจ ธมฺมํ ปฏิจฺจ อพฺยากโต เกนจิ วิญฺเญยฺโย ธมฺโม อุปฺปชฺชติ เหตุปจฺจยา. (1)}

\csromanheader{2}{1. กุสลตฺติก 13. เกนจิวิญฺเญยฺยทุก}
\prose{\csromanfolio{577}อกุสลํ เกนจิ วิญฺเญยฺยญฺจ อพฺยากตํ เกนจิ วิญฺเญยฺยญฺจ ธมฺมํ ปฏิจฺจ อพฺยากโต เกนจิ วิญฺเญยฺโย ธมฺโม อุปฺปชฺชติ เหตุปจฺจยา. (1)}

\proseitem{42}{กุสลํ เกนจิ วิญฺเญยฺยํ ธมฺมํ ปฏิจฺจ กุสโล เกนจิ วิญฺเญยฺโย ธมฺโม อุปฺปชฺชติ อารมฺมณปจฺจยา. (1)}

\prose{อกุสลํ เกนจิ วิญฺเญยฺยํ ธมฺมํ ปฏิจฺจ อกุสโล เกนจิ วิญฺเญยฺโย ธมฺโม อุปฺปชฺชติ อารมฺมณปจฺจยา. (1)}

\prose{อพฺยากตํ เกนจิ วิญฺเญยฺยํ ธมฺมํ ปฏิจฺจ อพฺยากโต เกนจิ วิญฺเญยฺโย ธมฺโม อุปฺปชฺชติ อารมฺมณปจฺจยา. (1) (สํขิตฺตํ.)}

\proseitem{43}{\textbf{เหตุ}ยา นว, อารมฺมเณ ตีณิ ฯเปฯ อวิคเต นว. (สํขิตฺตํ.)}

\csromanheader{2}{\textbf{ปจฺจนีย-นเหตุ}}
\proseitem{44}{อกุสลํ เกนจิ วิญฺเญยฺยํ ธมฺมํ ปฏิจฺจ อกุสโล เกนจิ วิญฺเญยฺโย ธมฺโม อุปฺปชฺชติ นเหตุปจฺจยา. (1)}

\prose{อพฺยากตํ เกนจิ วิญฺเญยฺยํ ธมฺมํ ปฏิจฺจ อพฺยากโต เกนจิ วิญฺเญยฺโย ธมฺโม อุปฺปชฺชติ นเหตุปจฺจยา. (1) (สํขิตฺตํ.)}

\proseitem{45}{นเหตุยา เทฺว, น-อารมฺมเณ ปญฺจ, น-อธิปติยา นว ฯเปฯ โนวิคเต ปญฺจ. (สํขิตฺตํ.\csromanspacer{} สหชาตวารมฺปิ ฯเปฯ สมฺปยุตฺตวารมฺปิ ปฏิจฺจวารสทิสํ.)}

\proseitem{46}{กุสโล เกนจิ วิญฺเญยฺโย ธมฺโม กุสลสฺส เกนจิ วิญฺเญยฺยสฺส ธมฺมสฺส เหตุปจฺจเยน ปจฺจโย.\csromanspacer{} ตีณิ.}

\prose{อกุสโล เกนจิ วิญฺเญยฺโย ธมฺโม อกุสลสฺส เกนจิ วิญฺเญยฺยสฺส ธมฺมสฺส เหตุปจฺจเยน ปจฺจโย.\csromanspacer{} ตีณิ.}

\prose{อพฺยากโต เกนจิ วิญฺเญยฺโย ธมฺโม อพฺยากตสฺส เกนจิ วิญฺเญยฺยสฺส ธมฺมสฺส เหตุปจฺจเยน ปจฺจโย. (1) (สํขิตฺตํ.)}

\vspace{\tipitakaparskip}
\csromancenter{\textbf{เหตุ}ยา สตฺต, อารมฺมเณ นว ฯเปฯ อวิคเต เตรส.}
\prose{(สํขิตฺตํ.\csromanspacer{} ยถา กุสลตฺติเก ปญฺหาวารํ, เอวํ วิตฺถาเรตพฺพํ.)}

\gambhira{ธมฺมานุโลม ติกทุกปฏฺฐานปาฬิ}
\csromanmatikaanchor{mtk.143}
\csromantocmarkref{subparagraph}{14-19. อาสวทุกาทิ}{mtk.143}
\bookmark[dest={mtk.143},level=5]{14-19. อาสวทุกาทิ}
\proseitem{48}{\csromanfolio{578}กุสลํ นเกนจิ วิญฺเญยฺยํ ธมฺมํ ปฏิจฺจ กุสโล นเกนจิ วิญฺเญยฺโย ธมฺโม อุปฺปชฺชติ เหตุปจฺจยา. (เกนจิวิญฺเญยฺยสทิสํ.)}

\nitthitam{กุสลตฺติกจูฬนฺตรทุกํ นิฏฺฐิตํ.}
\csromanheader{6}{1. กุสลตฺติก 14. อาสวทุก}
\csromanheader{2}{1-7. ปฏิจฺจาทิวาร ปจฺจยจตุกฺก-เหตุ}
\proseitem{1}{อกุสลํ อาสวํ ธมฺมํ ปฏิจฺจ อกุสโล อาสโว ธมฺโม อุปฺปชฺชติ เหตุปจฺจยา. (สํขิตฺตํ.)}

\csromanheader{2}{เหตุยา เอกํ ฯเปฯ อวิคเต เอกํ. (สํขิตฺตํ.)}
\prose{(สหชาตวาเรปิ ฯเปฯ สมฺปยุตฺตวาเรปิ ปญฺหาวาเรปิ สพฺพตฺถ เอกํ.)}

\proseitem{2}{กุสลํ โน-อสวํ ธมฺมํ ปฏิจฺจ กุสโล โน-อสโว ธมฺโม อุปฺปชฺชติ เหตุปจฺจยา.\csromanspacer{}\csromanspacer{}กุสลํ โน-อสวํ ธมฺมํ ปฏิจฺจ อพฺยากโต โน-อสโว ธมฺโม อุปฺปชฺชติ เหตุปจฺจยา.\csromanspacer{}\csromanspacer{}กุสลํ โน-อสวํ ธมฺมํ ปฏิจฺจ กุสโล โน-อสโว จ อพฺยากโต โน-อสโว จ ธมฺมา อุปฺปชฺชนฺติ เหตุปจฺจยา. (3)}

\prose{อกุสลํ โน-อสวํ ธมฺมํ ปฏิจฺจ อกุสโล โน-อสโว ธมฺโม อุปฺปชฺชติ เหตุปจฺจยา.\csromanspacer{} ตีณิ.}

\prose{อพฺยากตํ โน-อสวํ ธมฺมํ ปฏิจฺจ อพฺยากโต โน-อสโว ธมฺโม อุปฺปชฺชติ เหตุปจฺจยา. (1)}

\prose{กุสลํ โน-อสวญฺจ อพฺยากตํ โน-อสวญฺจ ธมฺมํ ปฏิจฺจ อพฺยากโต โน-อสโว ธมฺโม อุปฺปชฺชติ เหตุปจฺจยา. (1)}

\prose{อกุสลํ โน-อสวญฺจ อพฺยากตํ โน-อสวญฺจ ธมฺมํ ปฏิจฺจ อพฺยากโต โน-อสโว ธมฺโม อุปฺปชฺชติ เหตุปจฺจยา. (1)}

\csromanheader{2}{\textbf{อารมฺมณ}}
\proseitem{3}{กุสลํ โน-อสวํ ธมฺมํ ปฏิจฺจ กุสโล โน-อสโว ธมฺโม อุปฺปชฺชติ อารมฺมณปจฺจยา. (1)}

\csromanheader{2}{1. กุสลตฺติก 15. สาสวทุก}
\prose{\csromanfolio{579}อกุสลํ โน-อสวํ ธมฺมํ ปฏิจฺจ อกุสโล โน-อสโว ธมฺโม อุปฺปชฺชติ อารมฺมณปจฺจยา. (1)}

\prose{อพฺยากตํ โน-อสวํ ธมฺมํ ปฏิจฺจ อพฺยากโต โน-อสโว ธมฺโม อุปฺปชฺชติ อารมฺมณปจฺจยา. (1) (สํขิตฺตํ.)}

\proseitem{4}{เหตุยา นว, อารมฺมเณ ตีณิ, อธิปติยา นว ฯเปฯ อวิคเต นว. (สํขิตฺตํ.)}

\prose{นเหตุยา เอกํ, น-อารมฺมเณ ปญฺจ, น-อธิปติยา นว ฯเปฯ โนวิคเต ปญฺจ. (สํขิตฺตํ.\csromanspacer{} สหชาตวารมฺปิ ฯเปฯ สมฺปยุตฺตวารมฺปิ ปฏิจฺจวารสทิสํ.)}

\csromanheader{2}{\textbf{เหตุ-อารมฺมณ}}
\proseitem{5}{กุสโล โน-อสโว ธมฺโม กุสลสฺส โน-อสวสฺส ธมฺมสฺส เหตุปจฺจเยน ปจฺจโย.\csromanspacer{} ตีณิ.}

\prose{อกุสโล โน-อสโว ธมฺโม อกุสลสฺส โน-อสวสฺส ธมฺมสฺส เหตุปจฺจเยน ปจฺจโย.\csromanspacer{} ตีณิ.}

\prose{อพฺยากโต โน-อสโว ธมฺโม อพฺยากตสฺส โน-อสวสฺส ธมฺมสฺส เหตุปจฺจเยน ปจฺจโย. (1)}

\prose{กุสโล โน-อสโว ธมฺโม กุสลสฺส โน-อสวสฺส ธมฺมสฺส อารมฺมณปจฺจเยน ปจฺจโย. (สํขิตฺตํ.)}

\proseitem{6}{เหตุยา สตฺต, อารมฺมเณ นว, อธิปติยา ทส ฯเปฯ อวิคเต เตรส.}

\prose{(สํขิตฺตํ.\csromanspacer{} ยถา กุสลตฺติเก ปญฺหาวารํ, เอวํ วิตฺถาเรตพฺพํ.)}

\csromanheader{2}{1. กุสลตฺติก 15. สาสวทุก}
\csromanheader{2}{1-7. ปฏิจฺจาทิวาร ปจฺจยจตุกฺก-เหตุ}
\proseitem{7}{กุสลํ สาสวํ ธมฺมํ ปฏิจฺจ กุสโล สาสโว ธมฺโม อุปฺปชฺชติ เหตุปจฺจยา.\csromanspacer{} ตีณิ.}

\prose{อกุสลํ สาสวํ ธมฺมํ ปฏิจฺจ อกุสโล สาสโว ธมฺโม อุปฺปชฺชติ เหตุปจฺจยา.\csromanspacer{} ตีณิ.}

\gambhira{ธมฺมานุโลม ติกทุกปฏฺฐานปาฬิ}
\prose{\csromanfolio{580}อพฺยากตํ สาสวํ ธมฺมํ ปฏิจฺจ อพฺยากโต สาสโว ธมฺโม อุปฺปชฺชติ เหตุปจฺจยา. (1)}

\prose{กุสลํ สาสวญฺจ อพฺยากตํ สาสวญฺจ ธมฺมํ ปฏิจฺจ อพฺยากโต สาสโว ธมฺโม อุปฺปชฺชติ เหตุปจฺจยา. (1)}

\prose{อกุสลํ สาสวญฺจ อพฺยากตํ สาสวญฺจ ธมฺมํ ปฏิจฺจ อพฺยากโต สาสโว ธมฺโม อุปฺปชฺชติ เหตุปจฺจยา. (1) (สํขิตฺตํ.)}

\proseitem{8}{เหตุยา นว, อารมฺมเณ ตีณิ ฯเปฯ อวิคเต นว. (สํขิตฺตํ.)}

\prose{นเหตุยา เทฺว, น-อารมฺมเณ ปญฺจ, น-อธิปติยา นว. (สํขิตฺตํ.\csromanspacer{} สหชาตวารมฺปิ ฯเปฯ สมฺปยุตฺตวารมฺปิ ปฏิจฺจวารสทิสํ.)}

\proseitem{9}{กุสโล สาสโว ธมฺโม กุสลสฺส สาสวสฺส ธมฺมสฺส เหตุปจฺจเยน ปจฺจโย.\csromanspacer{} ตีณิ.}

\prose{อกุสโล สาสโว ธมฺโม อกุสลสฺส สาสวสฺส ธมฺมสฺส เหตุปจฺจเยน ปจฺจโย.\csromanspacer{} ตีณิ.}

\prose{อพฺยากโต สาสโว ธมฺโม อพฺยากตสฺส สาสวสฺส ธมฺมสฺส เหตุปจฺจเยน ปจฺจโย. (1) (สํขิตฺตํ.)}

\prose{เหตุยา สตฺต, อารมฺมเณ นว ฯเปฯ อวิคเต เตรส. (สํขิตฺตํ.)}

\prose{(ยถา กุสลตฺติเก ปญฺหาวารํ, เอวํ วิตฺถาเรตพฺพํ.)}

\csromanheader{2}{\textbf{อนาสวปท 1-7. ปฏิจฺจาทิวาร}}
\csromanheader{2}{\textbf{ปจฺจยจตุกฺก-เหตุ}}
\proseitem{10}{กุสลํ อนาสวํ ธมฺมํ ปฏิจฺจ กุสโล อนาสโว ธมฺโม อุปฺปชฺชติ เหตุปจฺจยา. (1)}

\prose{อพฺยากตํ อนาสวํ ธมฺมํ ปฏิจฺจ อพฺยากโต อนาสโว ธมฺโม อุปฺปชฺชติ เหตุปจฺจยา. (1) (สํขิตฺตํ.)}

\proseitem{11}{เหตุยา เทฺว ฯเปฯ อวิคเต เทฺว. (สํขิตฺตํ.\csromanspacer{} สหชาตวารมฺปิ ฯเปฯ สมฺปยุตฺตวารมฺปิ ปฏิจฺจวารสทิสํ.)}

\csromanheader{2}{1. กุสลตฺติก 16. อาสวสมฺปยุตฺตทุก \textbf{เหตุ อารมฺมณาทิ}}
\proseitem{12}{\csromanfolio{581}กุสโล อนาสโว ธมฺโม กุสลสฺส อนาสวสฺส ธมฺมสฺส เหตุปจฺจเยน ปจฺจโย. (1)}

\prose{อพฺยากโต อนาสโว ธมฺโม อพฺยากตสฺส อนาสวสฺส ธมฺมสฺส เหตุปจฺจเยน ปจฺจโย. (1)}

\prose{อพฺยากโต อนาสโว ธมฺโม อพฺยากตสฺส อนาสวสฺส ธมฺมสฺส อารมฺมณปจฺจเยน ปจฺจโย.\csromanspacer{}\csromanspacer{}อพฺยากโต อนาสโว ธมฺโม กุสลสฺส อนาสวสฺส ธมฺมสฺส อารมฺมณปจฺจเยน ปจฺจโย. (2)}

\prose{กุสโล อนาสโว ธมฺโม กุสลสฺส อนาสวสฺส ธมฺมสฺส อธิปติปจฺจเยน ปจฺจโย. (1)}

\prose{อพฺยากโต อนาสโว ธมฺโม อพฺยากตสฺส อนาสวสฺส ธมฺมสฺส อธิปติปจฺจเยน ปจฺจโย.\csromanspacer{}\csromanspacer{}อพฺยากโต อนาสโว ธมฺโม กุสลสฺส อนาสวสฺส ธมฺมสฺส อธิปติปจฺจเยน ปจฺจโย. (2)}

\prose{กุสโล อนาสโว ธมฺโม อพฺยากตสฺส อนาสวสฺส ธมฺมสฺส อนนฺตรปจฺจเยน ปจฺจโย. (1)}

\prose{อพฺยากโต อนาสโว ธมฺโม อพฺยากตสฺส อนาสวสฺส ธมฺมสฺส อนนฺตรปจฺจเยน ปจฺจโย. (1) (สํขิตฺตํ.)}

\proseitem{13}{เหตุยา เทฺว, อารมฺมเณ เทฺว, อธิปติยา ตีณิ, อนนฺตเร เทฺว, สมนนฺตเร เทฺว, สหชาเต เทฺว, อญฺญมญฺเญ เทฺว, นิสฺสเย เทฺว, อุปนิสฺสเย จตฺตาริ ฯเปฯ อวิคเต เทฺว. (สํขิตฺตํ.\csromanspacer{} ยถา กุสลตฺติเก ปญฺหาวารํ, เอวํ วิตฺถาเรตพฺพํ.)}

\csromanheader{2}{1. กุสลตฺติก 16. อาสวสมฺปยุตฺตทุก}
\csromanheader{2}{1-7. ปฏิจฺจาทิวาร ปจฺจยจตุกฺก-เหตุ}
\proseitem{14}{อกุสลํ อาสวสมฺปยุตฺตํ ธมฺมํ ปฏิจฺจ อกุสโล อาสวสมฺปยุตฺโต ธมฺโม อุปฺปชฺชติ เหตุปจฺจยา. (สํขิตฺตํ.)}

\prose{เหตุยา เอกํ, อารมฺมเณ เอกํ. (สํขิตฺตํ.)}

\prose{(สหชาตวาเรปิ ฯเปฯ สมฺปยุตฺตวาเรปิ ปญฺหาวาเรปิ สพฺพตฺถ เอกํ.)}

\gambhira{ธมฺมานุโลม ติกทุกปฏฺฐานปาฬิ}
\proseitem{15}{\csromanfolio{582}กุสลํ อาสววิปฺปยุตฺตํ ธมฺมํ ปฏิจฺจ กุสโล อาสววิปฺปยุตฺโต ธมฺโม อุปฺปชฺชติ เหตุปจฺจยา.\csromanspacer{}\csromanspacer{}กุสลํ อาสววิปฺปยุตฺตํ ธมฺมํ ปฏิจฺจ อพฺยากโต อาสววิปฺปยุตฺโต ธมฺโม อุปฺปชฺชติ เหตุปจฺจยา.\csromanspacer{}\csromanspacer{}กุสลํ อาสววิปฺปยุตฺตํ ธมฺมํ ปฏิจฺจ กุสโล อาสววิปฺปยุตฺโต จ อพฺยากโต อาสววิปฺปยุตฺโต จ ธมฺมา อุปฺปชฺชนฺติ เหตุปจฺจยา. (3)}

\prose{อกุสลํ อาสววิปฺปยุตฺตํ ธมฺมํ ปฏิจฺจ อพฺยากโต อาสววิปฺปยุตฺโต ธมฺโม อุปฺปชฺชติ เหตุปจฺจยา. (1)}

\prose{อพฺยากตํ อาสววิปฺปยุตฺตํ ธมฺมํ ปฏิจฺจ อพฺยากโต อาสววิปฺปยุตฺโต ธมฺโม อุปฺปชฺชติ เหตุปจฺจยา. (1)}

\prose{กุสลํ อาสววิปฺปยุตฺตญฺจ อพฺยากตํ อาสววิปฺปยุตฺตญฺจ ธมฺมํ ปฏิจฺจ อพฺยากโต อาสววิปฺปยุตฺโต ธมฺโม อุปฺปชฺชติ เหตุปจฺจยา. (1)}

\prose{อกุสลํ อาสววิปฺปยุตฺตญฺจ อพฺยากตํ อาสววิปฺปยุตฺตญฺจ ธมฺมํ ปฏิจฺจ อพฺยากโต อาสววิปฺปยุตฺโต ธมฺโม อุปฺปชฺชติ เหตุปจฺจยา. (สํขิตฺตํ.)}

\proseitem{16}{เหตุยา สตฺต, อารมฺมเณ เทฺว, อธิปติยา สตฺต ฯเปฯ อวิคเต สตฺต. (สํขิตฺตํ.)}

\prose{(สหชาตวารมฺปิ ฯเปฯ สมฺปยุตฺตวารมฺปิ ปฏิจฺจวารสทิสํ.)}

\proseitem{17}{กุสโล อาสววิปฺปยุตฺโต ธมฺโม กุสลสฺส อาสววิปฺปยุตฺตสฺส ธมฺมสฺส เหตุปจฺจเยน ปจฺจโย.\csromanspacer{} ตีณิ.}

\prose{อกุสโล อาสววิปฺปยุตฺโต ธมฺโม อพฺยากตสฺส อาสววิปฺปยุตฺตสฺส ธมฺมสฺส เหตุปจฺจเยน ปจฺจโย. (1)}

\prose{อพฺยากโต อาสววิปฺปยุตฺโต ธมฺโม อพฺยากตสฺส อาสววิปฺปยุตฺตสฺส ธมฺมสฺส เหตุปจฺจเยน ปจฺจโย. (สํขิตฺตํ.)}

\proseitem{18}{เหตุยา ปญฺจ, อารมฺมเณ นว, อธิปติยา ปญฺจ ฯเปฯ อวิคเต ทส. (สํขิตฺตํ.)}

\prose{(ยถา กุสลตฺติเก ปญฺหาวารํ, เอวํ วิตฺถาเรตพฺพํ.)}

\vspace{\tipitakaparskip}
\csromancenter{กุสลตฺติก 17.\csromanspacer{} อาสวสาสวทุก \textbf{1.}\csromanspacer{}\textbf{ กุสลตฺติก 17.}\csromanspacer{}\textbf{ อาสวสาสวทุก}}
\csromanheader{2}{1-7. ปฏิจฺจาทิวาร ปจฺจยจตุกฺก-เหตุ}
\proseitem{19}{\csromanfolio{583}อกุสลํ อาสวญฺเจว สาสวญฺจ ธมฺมํ ปฏิจฺจ อกุสโล อาสโว เจว สาสโว จ ธมฺโม อุปฺปชฺชติ เหตุปจฺจยา. (สํขิตฺตํ.)}

\prose{เหตุยา เอกํ, อารมฺมเณ เอกํ ฯเปฯ อวิคเต เอกํ. (สํขิตฺตํ.)}

\prose{(สหชาตวาเรปิ ฯเปฯ สมฺปยุตฺตวาเรปิ ปญฺหาวาเรปิ สพฺพตฺถ เอกํ.)}

\proseitem{20}{กุสลํ สาสวญฺเจว โน จ อาสวํ ธมฺมํ ปฏิจฺจ กุสโล สาสโว เจว โน จ อาสโว ธมฺโม อุปฺปชฺชติ เหตุปจฺจยา.\csromanspacer{} ตีณิ.}

\prose{อกุสลํ สาสวญฺเจว โน จ อาสวํ ธมฺมํ ปฏิจฺจ อกุสโล สาสโว เจว โน จ อาสโว ธมฺโม อุปฺปชฺชติ เหตุปจฺจยา.\csromanspacer{} ตีณิ.}

\prose{อพฺยากตํ สาสวญฺเจว โน จ อาสวํ ธมฺมํ ปฏิจฺจ อพฺยากโต สาสโว เจว โน จ อาสโว ธมฺโม อุปฺปชฺชติ เหตุปจฺจยา. (1)}

\prose{กุสลํ สาสวญฺเจว โน จ อาสวญฺจ อพฺยากตํ สาสวญฺเจว โน จ อาสวญฺจ ธมฺมํ ปฏิจฺจ อพฺยากโต สาสโว เจว โน จ อาสโว ธมฺโม อุปฺปชฺชติ เหตุปจฺจยา. (1)}

\prose{อกุสลํ สาสวญฺเจว โน จ อาสวญฺจ อพฺยากตํ สาสวญฺเจว โน จ อาสวญฺจ ธมฺมํ ปฏิจฺจ อพฺยากโต สาสโว เจว โน จ อาสโว ธมฺโม อุปฺปชฺชติ เหตุปจฺจยา. (1) (สํขิตฺตํ.)}

\proseitem{21}{เหตุยา นว, อารมฺมเณ ตีณิ, อธิปติยา นว ฯเปฯ อวิคเต นว. (สํขิตฺตํ.)}

\prose{(สหชาตวารมฺปิ ฯเปฯ สมฺปยุตฺตวารมฺปิ ปฏิจฺจวารสทิสํ.)}

\proseitem{22}{กุสโล สาสโว เจว โน จ อาสโว ธมฺโม กุสลสฺส สาสวสฺส เจว โน จ อาสวสฺส ธมฺมสฺส เหตุปจฺจเยน ปจฺจโย.\csromanspacer{} ตีณิ.}

\prose{อกุสโล สาสโว เจว โน จ อาสโว ธมฺโม อกุสลสฺส สาสวสฺส เจว โน จ อาสวสฺส ธมฺมสฺส เหตุปจฺจเยน ปจฺจโย.\csromanspacer{} ตีณิ.}

\gambhira{ธมฺมานุโลม ติกทุกปฏฺฐานปาฬิ}
\prose{\csromanfolio{584}อพฺยากโต สาสโว เจว โน จ อาสโว ธมฺโม อพฺยากตสฺส สาสวสฺส เจว โน จ อาสวสฺส ธมฺมสฺส เหตุปจฺจเยน ปจฺจโย. (1) (สํขิตฺตํ.)}

\proseitem{23}{เหตุยา สตฺต, อารมฺมเณ นว, อธิปติยา นว ฯเปฯ อวิคเต เตรส. (สํขิตฺตํ.\csromanspacer{} ยถา กุสลตฺติเก ปญฺหาวารํ, เอวํ วิตฺถาเรตพฺพํ.)}

\csromanheader{2}{1. กุสลตฺติก 18. อาสว-อาสวสมฺปยุตฺตทุก}
\csromanheader{2}{1-7. ปฏิจฺจาทิวาร ปจฺจยจตุกฺก}
\proseitem{24}{อกุสลํ อาสวญฺเจว อาสวสมฺปยุตฺตญฺจ ธมฺมํ ปฏิจฺจ อกุสโล อาสโว เจว อาสวสมฺปยุตฺโต จ ธมฺโม อุปฺปชฺชติ เหตุปจฺจยา. (สํขิตฺตํ.)}

\prose{เหตุยา เอกํ, อารมฺมเณ เอกํ ฯเปฯ อวิคเต เอกํ. (สํขิตฺตํ.)}

\prose{(สหชาตวาเรปิ ฯเปฯ สมฺปยุตฺตวาเรปิ ปญฺหาวาเรปิ สพฺพตฺถ เอกํ.)}

\proseitem{25}{อกุสลํ อาสวสมฺปยุตฺตญฺเจว โน จ อาสวํ ธมฺมํ ปฏิจฺจ อกุสโล อาสวสมฺปยุตฺโต เจว โน จ อาสโว ธมฺโม อุปฺปชฺชติ เหตุปจฺจยา. (สํขิตฺตํ.)}

\prose{เหตุยา เอกํ, อารมฺมเณ เอกํ ฯเปฯ อวิคเต เอกํ. (สํขิตฺตํ.)}

\prose{(สหชาตวาเรปิ ฯเปฯ สมฺปยุตฺตวาเรปิ ปญฺหาวาเรปิ สพฺพตฺถ เอกํ.)}

\proseitem{26}{กุสลํ อาสววิปฺปยุตฺตํ สาสวํ ธมฺมํ ปฏิจฺจ กุสโล อาสววิปฺปยุตฺโต สาสโว ธมฺโม อุปฺปชฺชติ เหตุปจฺจยา.\csromanspacer{} ตีณิ.}

\prose{อกุสลํ อาสววิปฺปยุตฺตํ สาสวํ ธมฺมํ ปฏิจฺจ อพฺยากโต อาสววิปฺปยุตฺโต สาสโว ธมฺโม อุปฺปชฺชติ เหตุปจฺจยา. (1)}

\prose{อพฺยากตํ อาสววิปฺปยุตฺตํ สาสวํ ธมฺมํ ปฏิจฺจ อพฺยากโต อาสววิปฺปยุตฺโต สาสโว ธมฺโม อุปฺปชฺชติ เหตุปจฺจยา. (1)}

\prose{กุสลํ อาสววิปฺปยุตฺตํ สาสวญฺจ อพฺยากตํ อาสววิปฺปยุตฺตํ สาสวญฺจ ธมฺมํ ปฏิจฺจ อพฺยากโต อาสววิปฺปยุตฺโต สาสโว ธมฺโม อุปฺปชฺชติ เหตุปจฺจยา. (1)}

\csromanheader{2}{1. กุสลตฺติก 19. อาสววิปฺปยุตฺตสาสวทุก}
\prose{\csromanfolio{585}อกุสลํ อาสววิปฺปยุตฺตํ สาสวญฺจ อพฺยากตํ อาสววิปฺปยุตฺตํ สาสวญฺจ ธมฺมํ ปฏิจฺจ อพฺยากโต อาสววิปฺปยุตฺโต สาสโว ธมฺโม อุปฺปชฺชติ เหตุปจฺจยา. (1) (สํขิตฺตํ.)}

\proseitem{27}{เหตุยา สตฺต, อารมฺมเณ เทฺว ฯเปฯ อวิคเต สตฺต. (สํขิตฺตํ.\csromanspacer{} สหชาตวารมฺปิ ฯเปฯ สมฺปยุตฺตวารมฺปิ ปฏิจฺจวารสทิสํ.)}

\proseitem{28}{กุสโล อาสววิปฺปยุตฺโต สาสโว ธมฺโม กุสลสฺส อาสววิปฺปยุตฺตสฺส สาสวสฺส ธมฺมสฺส เหตุปจฺจเยน ปจฺจโย.\csromanspacer{} ตีณิ.}

\prose{อกุสโล อาสววิปฺปยุตฺโต สาสโว ธมฺโม อพฺยากตสฺส อาสววิปฺปยุตฺตสฺส สาสวสฺส ธมฺมสฺส เหตุปจฺจเยน ปจฺจโย. (1)}

\prose{อพฺยากโต อาสววิปฺปยุตฺโต สาสโว ธมฺโม อพฺยากตสฺส อาสววิปฺปยุตฺตสฺส สาสวสฺส ธมฺมสฺส เหตุปจฺจเยน ปจฺจโย. (สํขิตฺตํ.)}

\proseitem{29}{เหตุยา ปญฺจ, อารมฺมเณ นว, อธิปติยา จตฺตาริ ฯเปฯ อวิคเต ทส. (สํขิตฺตํ.\csromanspacer{} ยถา กุสลตฺติเก ปญฺหาวารํ, เอวํ วิตฺถาเรตพฺพํ.)}

\csromanheader{2}{1. กุสลตฺติก 19. อาสววิปฺปยุตฺตสาสวทุก}
\csromanheader{2}{1-7. ปฏิจฺจาทิวาร ปจฺจยจตุกฺก-เหตุ}
\proseitem{30}{กุสลํ อาสววิปฺปยุตฺตํ อนาสวํ ธมฺมํ ปฏิจฺจ กุสโล อาสววิปฺปยุตฺโต อนาสโว ธมฺโม อุปฺปชฺชติ เหตุปจฺจยา. (1)}

\prose{อพฺยากตํ อาสววิปฺปยุตฺตํ อนาสวํ ธมฺมํ ปฏิจฺจ อพฺยากโต อาสววิปฺปยุตฺโต อนาสโว ธมฺโม อุปฺปชฺชติ เหตุปจฺจยา. (สํขิตฺตํ.)}

\proseitem{31}{เหตุยา เทฺว, อารมฺมเณ เทฺว ฯเปฯ อวิคเต เทฺว. (สํขิตฺตํ.)}

\prose{(สหชาตวารมฺปิ ฯเปฯ สมฺปยุตฺตวารมฺปิ ปฏิจฺจวารสทิสํ.)}

\csromanheader{2}{ธมฺมานุโลม ติกทุกปฏฺฐานปาฬิ \textbf{เหตุ อารมฺมณาทิ}}
\proseitem{32}{\csromanfolio{586}กุสโล อาสววิปฺปยุตฺโต อนาสโว ธมฺโม กุสลสฺส อาสววิปฺปยุตฺตสฺส อนาสวสฺส ธมฺมสฺส เหตุปจฺจเยน ปจฺจโย. (1)}

\prose{อพฺยากโต อาสววิปฺปยุตฺโต อนาสโว ธมฺโม อพฺยากตสฺส อาสววิปฺปยุตฺตสฺส อนาสวสฺส ธมฺมสฺส เหตุปจฺจเยน ปจฺจโย. (1)}

\prose{อพฺยากโต อาสววิปฺปยุตฺโต อนาสโว ธมฺโม อพฺยากตสฺส อาสววิปฺปยุตฺตสฺส อนาสวสฺส ธมฺมสฺส อารมฺมณปจฺจเยน ปจฺจโย. (1)}

\prose{อพฺยากโต อาสววิปฺปยุตฺโต อนาสโว ธมฺโม กุสลสฺส อาสววิปฺปยุตฺตสฺส อนาสวสฺส ธมฺมสฺส อารมฺมณปจฺจเยน ปจฺจโย. (1)}

\prose{กุสโล อาสววิปฺปยุตฺโต อนาสโว ธมฺโม กุสลสฺส อาสววิปฺปยุตฺตสฺส อนาสวสฺส ธมฺมสฺส อธิปติปจฺจเยน ปจฺจโย. (1)}

\prose{อพฺยากโต อาสววิปฺปยุตฺโต อนาสโว ธมฺโม อพฺยากตสฺส อาสววิปฺปยุตฺตสฺส อนาสวสฺส ธมฺมสฺส อธิปติปจฺจเยน ปจฺจโย. (1)}

\prose{อพฺยากโต อาสววิปฺปยุตฺโต อนาสโว ธมฺโม กุสลสฺส อาสววิปฺปยุตฺตสฺส อนาสวสฺส ธมฺมสฺส อธิปติปจฺจเยน ปจฺจโย. (1)}

\prose{กุสโล อาสววิปฺปยุตฺโต อนาสโว ธมฺโม อพฺยากตสฺส อาสววิปฺปยุตฺตสฺส อนาสวสฺส ธมฺมสฺส อนนฺตรปจฺจเยน ปจฺจโย. (1)}

\prose{อพฺยากโต อาสววิปฺปยุตฺโต อนาสโว ธมฺโม อพฺยากตสฺส อาสววิปฺปยุตฺตสฺส อนาสวสฺส ธมฺมสฺส อนนฺตรปจฺจเยน ปจฺจโย. (1) (สํขิตฺตํ.)}

\proseitem{33}{เหตุยา เทฺว, อารมฺมเณ เทฺว, อธิปติยา ตีณิ, อนนฺตเร เทฺว, สมนนฺตเร เทฺว, สหชาเต เทฺว, อญฺญมญฺเญ เทฺว, นิสฺสเย เทฺว, อุปนิสฺสเย จตฺตาริ ฯเปฯ อวิคเต เทฺว. (สํขิตฺตํ.)}

\prose{(ยถา กุสลตฺติเก ปญฺหาวารํ, เอวํ วิตฺถาเรตพฺพํ.)}

\nitthitam{กุสลตฺติก-อาสวโคจฺฉกํ นิฏฺฐิตํ.}
\vspace{\tipitakaparskip}
\csromancenter{กุสลตฺติก 55.\csromanspacer{} สารมฺมณทุก \textbf{1.}\csromanspacer{}\textbf{ กุสลตฺติก 20-54.}\csromanspacer{}\textbf{ สญฺโญชนาทิทุก}}
\csromanmatikaanchor{mtk.144}
\csromantocmarkref{subparagraph}{20-54. สญฺโญชนาทิทุก}{mtk.144}
\bookmark[dest={mtk.144},level=5]{20-54. สญฺโญชนาทิทุก}
\proseitem{34}{\csromanfolio{587}อกุสลํ สญฺโญชนํ ธมฺมํ ปฏิจฺจ ฯเปฯ.\csromanspacer{} อกุสลํ คนฺถํ.\csromanspacer{} โอฆํ.\csromanspacer{} โยคํ.\csromanspacer{} นีวรณํ.}

\prose{กุสลํ โนปรามาสํ ธมฺมํ ปฏิจฺจ กุสโล โนปรามาโส ธมฺโม อุปฺปชฺชติ เหตุปจฺจยา. (อาสวโคจฺฉกสทิสํ.) ทิฏฺฐิํ ปฏิจฺจ ทิฏฺฐิ น อุปฺปชฺชติ.}

\nitthitam{กุสลตฺติเก ฉ โคจฺฉกํ นิฏฺฐิตํ.}
\csromanmatikaanchor{mtk.145}
\csromantocmarkref{subparagraph}{55. สารมฺมณทุก}{mtk.145}
\bookmark[dest={mtk.145},level=5]{55. สารมฺมณทุก}
\csromanheader{6}{1. กุสลตฺติก 55. สารมฺมณทุก}
\csromanheader{2}{1-7. ปฏิจฺจาทิวาร ปจฺจยจตุกฺก}
\proseitem{1}{กุสลํ สารมฺมณํ ธมฺมํ ปฏิจฺจ กุสโล สารมฺมโณ ธมฺโม อุปฺปชฺชติ เหตุปจฺจยา. (1)}

\prose{อกุสลํ สารมฺมณํ ธมฺมํ ปฏิจฺจ อกุสโล สารมฺมโณ ธมฺโม อุปฺปชฺชติ เหตุปจฺจยา. (1)}

\prose{อพฺยากตํ สารมฺมณํ ธมฺมํ ปฏิจฺจ อพฺยากโต สารมฺมโณ ธมฺโม อุปฺปชฺชติ เหตุปจฺจยา. (1) (สํขิตฺตํ.)}

\proseitem{2}{เหตุยา ตีณิ, อารมฺมเณ ตีณิ ฯเปฯ อวิคเต ตีณิ. (สํขิตฺตํ.\csromanspacer{} สหชาตวารมฺปิ ฯเปฯ สมฺปยุตฺตวารมฺปิ ปฏิจฺจวารสทิสํ.)}

\csromanheader{2}{\textbf{เหตุ-อารมฺมณ}}
\proseitem{3}{กุสโล สารมฺมโณ ธมฺโม กุสลสฺส สารมฺมณสฺส ธมฺมสฺส เหตุปจฺจเยน ปจฺจโย. (1)}

\prose{อกุสโล สารมฺมโณ ธมฺโม อกุสลสฺส สารมฺมณสฺส ธมฺมสฺส เหตุปจฺจเยน ปจฺจโย. (1)}

\prose{อพฺยากโต สารมฺมโณ ธมฺโม อพฺยากตสฺส สารมฺมณสฺส ธมฺมสฺส เหตุปจฺจเยน ปจฺจโย. (1)}

\gambhira{ธมฺมานุโลม ติกทุกปฏฺฐานปาฬิ}
\prose{\csromanfolio{588}กุสโล สารมฺมโณ ธมฺโม กุสลสฺส สารมฺมณสฺส ธมฺมสฺส อารมฺมณปจฺจเยน ปจฺจโย. (สํขิตฺตํ.)}

\proseitem{4}{เหตุยา ตีณิ, อารมฺมเณ นว, อธิปติยา สตฺต ฯเปฯ อวิคเต ตีณิ. (สํขิตฺตํ.)}

\prose{(ยถา กุสลตฺติเก ปญฺหาวารํ, เอวํ วิตฺถาเรตพฺพํ.)}

\csromanheader{2}{\textbf{เหตุสหชาตาทิ}}
\proseitem{5}{อพฺยากตํ อนารมฺมณํ ธมฺมํ ปฏิจฺจ อพฺยากโต อนารมฺมโณ ธมฺโม อุปฺปชฺชติ เหตุปจฺจยา. (สํขิตฺตํ.)}

\csromanheader{2}{เหตุยา เอกํ ฯเปฯ อวิคเต เอกํ. (สํขิตฺตํ.)}
\csromanheader{2}{นเหตุยา เอกํ. (สํขิตฺตํ.)}
\prose{(สหชาตวารมฺปิ ฯเปฯ สมฺปยุตฺตวารมฺปิ ปฏิจฺจวารสทิสํ.)}

\proseitem{6}{อพฺยากโต อนารมฺมโณ ธมฺโม อพฺยากตสฺส อนารมฺมณสฺส ธมฺมสฺส สหชาตปจฺจเยน ปจฺจโย.\csromanspacer{} อญฺญมญฺญปจฺจเยน ปจฺจโย.\csromanspacer{} นิสฺสยปจฺจเยน ปจฺจโย.\csromanspacer{} อาหารปจฺจเยน ปจฺจโย.\csromanspacer{} อินฺทฺริยปจฺจเยน ปจฺจโย.\csromanspacer{} อตฺถิปจฺจเยน ปจฺจโย.\csromanspacer{} อวิคตปจฺจเยน ปจฺจโย. (สพฺพตฺถ เอกํ.\csromanspacer{} สํขิตฺตํ.)}

\csromanheader{2}{1. กุสลตฺติก 56. จิตฺตทุก}
\csromanheader{2}{1-7. ปฏิจฺจาทิวาร ปจฺจยจตุกฺก}
\proseitem{7}{กุสโล จิตฺโต ธมฺโม กุสลสฺส จิตฺตสฺส ธมฺมสฺส อารมฺมณปจฺจเยน ปจฺจโย. (สํขิตฺตํ.)}

\prose{อารมฺมเณ นว, อธิปติยา สตฺต. (สํขิตฺตํ.)}

\prose{นเหตุยา นว, น-อารมฺมเณ นว. (สํขิตฺตํ.)}

\prose{(ยถา กุสลตฺติเก ปญฺหาวารํ, เอวํ วิตฺถาเรตพฺพํ.)}

\csromanheader{2}{1. กุสลตฺติก 57. เจตสิกทุก}
\proseitem{8}{\csromanfolio{589}กุสลํ โนจิตฺตํ ธมฺมํ ปฏิจฺจ กุสโล โนจิตฺโต ธมฺโม อุปฺปชฺชติ เหตุปจฺจยา.\csromanspacer{} ตีณิ.}

\prose{อกุสลํ โนจิตฺตํ ธมฺมํ ปฏิจฺจ อกุสโล โนจิตฺโต ธมฺโม อุปฺปชฺชติ เหตุปจฺจยา.\csromanspacer{} ตีณิ.}

\prose{อพฺยากตํ โนจิตฺตํ ธมฺมํ ปฏิจฺจ อพฺยากโต โนจิตฺโต ธมฺโม อุปฺปชฺชติ เหตุปจฺจยา. (1)}

\prose{กุสลํ โนจิตฺตญฺจ อพฺยากตํ โนจิตฺตญฺจ ธมฺมํ ปฏิจฺจ อพฺยากโต โนจิตฺโต ธมฺโม อุปฺปชฺชติ เหตุปจฺจยา. (1)}

\prose{อกุสลํ โนจิตฺตญฺจ อพฺยากตํ โนจิตฺตญฺจ ธมฺมํ ปฏิจฺจ อพฺยากโต โนจิตฺโต ธมฺโม อุปฺปชฺชติ เหตุปจฺจยา. (1)}

\prose{กุสลํ โนจิตฺตํ ธมฺมํ ปฏิจฺจ กุสโล โนจิตฺโต ธมฺโม อุปฺปชฺชติ อารมฺมณปจฺจยา. (สํขิตฺตํ.)}

\proseitem{9}{เหตุยา นว, อารมฺมเณ ตีณิ, อธิปติยา นว ฯเปฯ อวิคเต นว. (สํขิตฺตํ.\csromanspacer{} สหชตวารมฺปิ ฯเปฯ สมฺปยุตฺตวารมฺปิ ปฏิจฺจวารสทิสํ.)}

\proseitem{10}{กุสโล โนจิตฺโต ธมฺโม กุสลสฺส โนจิตฺตสฺส ธมฺมสฺส เหตุปจฺจเยน ปจฺจโย. (สํขิตฺตํ.)}

\prose{เหตุยา สตฺต, อารมฺมเณ นว, อธิปติยา ทส ฯเปฯ อวิคเต เตรส. (สํขิตฺตํ.)}

\csromanheader{2}{1. กุสลตฺติก 57. เจตสิกทุก}
\csromanheader{2}{1-7. ปฏิจฺจาทิวาร ปจฺจยจตุกฺก}
\proseitem{11}{กุสลํ เจตสิกํ ธมฺมํ ปฏิจฺจ กุสโล เจตสิโก ธมฺโม อุปฺปชฺชติ เหตุปจฺจยา. (1)}

\prose{อกุสลํ เจตสิกํ ธมฺมํ ปฏิจฺจ อกุสโล เจตสิโก ธมฺโม อุปฺปชฺชติ เหตุปจฺจยา. (1)}

\gambhira{ธมฺมานุโลม ติกทุกปฏฺฐานปาฬิ}
\prose{\csromanfolio{590}อพฺยากตํ เจตสิกํ ธมฺมํ ปฏิจฺจ อพฺยากโต เจตสิโก ธมฺโม อุปฺปชฺชติ เหตุปจฺจยา. (1) (สํขิตฺตํ.)}

\proseitem{12}{เหตุยา ตีณิ, อารมฺมเณ ตีณิ ฯเปฯ. (สํขิตฺตํ.\csromanspacer{} สหชตวารมฺปิ ฯเปฯ สมฺปยุตฺตวารมฺปิ ปฏิจฺจวารสทิสํ.)}

\csromanheader{2}{\textbf{เหตุ อารมฺมณ}}
\proseitem{13}{กุสโล เจตสิโก ธมฺโม กุสลสฺส เจตสิกสฺส ธมฺมสฺส เหตุปจฺจเยน ปจฺจโย. (1)}

\prose{อกุสโล เจตสิโก ธมฺโม อกุสลสฺส เจตสิกสฺส ธมฺมสฺส เหตุปจฺจเยน ปจฺจโย. (1)}

\prose{อพฺยากโต เจตสิโก ธมฺโม อพฺยากตสฺส เจตสิกสฺส ธมฺมสฺส เหตุปจฺจเยน ปจฺจโย. (1)}

\prose{กุสโล เจตสิโก ธมฺโม กุสลสฺส เจตสิกสฺส ธมฺมสฺส อารมฺมณปจฺจเยน ปจฺจโย.\csromanspacer{} ตีณิ.}

\prose{อกุสโล เจสิตโก ธมฺโม อกุสลสฺส เจตสิกสฺส ธมฺมสฺส อารมฺมณปจฺจเยน ปจฺจโย.\csromanspacer{} ตีณิ.}

\prose{อพฺยากโต เจตสิโก ธมฺโม อพฺยากตสฺส เจตสิกสฺส ธมฺมสฺส อารมฺมณปจฺจเยน ปจฺจโย.\csromanspacer{} ตีณิ. (สํขิตฺตํ.)}

\proseitem{14}{เหตุยา ตีณิ, อารมฺมเณ นว, อธิปติยา สตฺต ฯเปฯ อวิคเต ตีณิ. (สํขิตฺตํ.)}

\prose{(ยถา กุสลตฺติเก ปญฺหาวารํ, เอวํ วิตฺถาเรตพฺพํ.)}

\proseitem{15}{กุสลํ อเจตสิกํ ธมฺมํ ปฏิจฺจ อพฺยากโต อเจตสิโก ธมฺโม อุปฺปชฺชติ เหตุปจฺจยา. (1)}

\prose{อกุสลํ อเจตสิกํ ธมฺมํ ปฏิจฺจ อพฺยากโต อเจตสิโก ธมฺโม อุปฺปชฺชติ เหตุปจฺจยา. (1)}

\prose{อพฺยากตํ อเจตสิกํ ธมฺมํ ปฏิจฺจ อพฺยากโต อเจตสิโก ธมฺโม อุปฺปชฺชติ เหตุปจฺจยา. (1)}

\csromanheader{2}{1. กุสลตฺติก 66. อชฺฌตฺติกทุก}
\prose{\csromanfolio{591}กุสลํ อเจตสิกญฺจ อพฺยากตํ อเจตสิกญฺจ ธมฺมํ ปฏิจฺจ อพฺยากโต อเจตสิโก ธมฺโม อุปฺปชฺชติ เหตุปจฺจยา. (1)}

\prose{อกุสลํ อเจตสิกญฺจ อพฺยากตํ อเจตสิกญฺจ ธมฺมํ ปฏิจฺจ อพฺยากโต อเจตสิโก ธมฺโม อุปฺปชฺชติ เหตุปจฺจยา. (1) (สํขิตฺตํ.)}

\proseitem{16}{เหตุยา ปญฺจ, อารมฺมเณ เอกํ ฯเปฯ อวิคเต ปญฺจ.}

\prose{นเหตุยา เอกํ, น-อารมฺมเณ ปญฺจ, โนวิคเต ปญฺจ.}

\prose{(สหชาตวารมฺปิ ฯเปฯ ปญฺหาวารมฺปิ เอวํ วิตฺถาเรตพฺพํ.)}

\csromanmatikaanchor{mtk.146}
\csromantocmarkref{subparagraph}{58-65. จิตฺตสมฺปยุตฺตาทิทุก}{mtk.146}
\bookmark[dest={mtk.146},level=5]{58-65. จิตฺตสมฺปยุตฺตาทิทุก}
\csromanheader{6}{1. กุสลตฺติก 58-65. จิตฺตสมฺปยุตฺตาทิทุก}
\csromanheader{2}{1-7. ปฏิจฺจาทิวาร}
\proseitem{17}{กุสลํ จิตฺตสมฺปยุตฺตํ ธมฺมํ ปฏิจฺจ ฯเปฯ.\csromanspacer{} กุสลํ จิตฺตสํสฏฺฐํ ธมฺมํ ปฏิจฺจ ฯเปฯ.\csromanspacer{} กุสลํ จิตฺตสมุฏฺฐานํ ธมฺมํ ปฏิจฺจ ฯเปฯ.\csromanspacer{} กุสลํ จิตฺตสหภุํ ธมฺมํ ปฏิจฺจ ฯเปฯ.\csromanspacer{} กุสลํ จิตฺตานุปริวตฺติํ ธมฺมํ ปฏิจฺจ ฯเปฯ.\csromanspacer{} กุสลํ จิตฺตสํสฏฺฐสมุฏฺฐานํ ธมฺมํ ปฏิจฺจ ฯเปฯ.\csromanspacer{} กุสลํ จิตฺตสํสฏฺฐสมุฏฺฐานสหภุํ ธมฺมํ ปฏิจฺจ ฯเปฯ.\csromanspacer{} กุสลํ จิตฺตสํสฏฺฐสมุฏฺฐานานุปริวตฺติํ ธมฺมํ ปฏิจฺจ กุสโล จิตฺตสํสฏฺฐสมุฏฺฐานานุปริวตฺติ ธมฺโม อุปฺปชฺชติ เหตุปจฺจยา. (สํขิตฺตํ.)}

\prose{เหตุยา ตีณิ, อารมฺมเณ ตีณิ ฯเปฯ อวิคเต ตีณิ. (สํขิตฺตํ.\csromanspacer{} สหชาตวารมฺปิ ฯเปฯ สมฺปยุตฺตวารมฺปิ ปญฺหาวารมฺปิ เอวํ วิตฺถาเรตพฺพํ.)}

\csromanheader{2}{1. กุสลตฺติก 66. อชฺฌตฺติกทุก}
\csromanheader{2}{1-7. ปฏิจฺจาทิวาร ปจฺจยจตุกฺก}
\proseitem{18}{อพฺยากตํ อชฺฌตฺติกํ ธมฺมํ ปฏิจฺจ อพฺยากโต อชฺฌตฺติโก ธมฺโม อุปฺปชฺชติ เหตุปจฺจยา. (สํขิตฺตํ.)}

\prose{เหตุยา เอกํ, สหชาเต เอกํ ฯเปฯ อวิคเต เอกํ. (สํขิตฺตํ.)}

\prose{(สหชาตวารมฺปิ ฯเปฯ สมฺปยุตฺตวารมฺปิ ปฏิจฺจวารสทิสํ.)}

\gambhira{ธมฺมานุโลม ติกทุกปฏฺฐานปาฬิ}
\proseitem{19}{\csromanfolio{592}กุสโล อชฺฌตฺติโก ธมฺโม กุสลสฺส อชฺฌตฺติกสฺส ธมฺมสฺส อารมฺมณปจฺจเยน ปจฺจโย.\csromanspacer{}\csromanspacer{}กุสโล อชฺฌตฺติโก ธมฺโม อกุสลสฺส อชฺฌตฺติกสฺส ธมฺมสฺส อารมฺมณปจฺจเยน ปจฺจโย.\csromanspacer{}\csromanspacer{}กุสโล อชฺฌตฺติโก ธมฺโม อพฺยากตสฺส อชฺฌตฺติกสฺส ธมฺมสฺส อารมฺมณปจฺจเยน ปจฺจโย. (3)}

\prose{อกุสโล อชฺฌตฺติโก ธมฺโม อกุสลสฺส อชฺฌตฺติกสฺส ธมฺมสฺส อารมฺมณปจฺจเยน ปจฺจโย.\csromanspacer{} ตีณิ.}

\prose{อพฺยากโต อชฺฌตฺติโก ธมฺโม อพฺยากตสฺส อชฺฌตฺติกสฺส ธมฺมสฺส อารมฺมณปจฺจเยน ปจฺจโย.\csromanspacer{} ตีณิ. (สํขิตฺตํ.)}

\proseitem{20}{อารมฺมเณ นว ฯเปฯ ปุเรชาเต ตีณิ ฯเปฯ อวิคเต ปญฺจ. (สํขิตฺตํ.\csromanspacer{} ยถา กุสลตฺติเก ปญฺหาวารํ, เอวํ วิตฺถาเรตพฺพํ.)}

\proseitem{21}{กุสลํ พาหิรํ ธมฺมํ ปฏิจฺจ กุสโล พาหิโร ธมฺโม อุปฺปชฺชติ เหตุปจฺจยา. (สํขิตฺตํ.)}

\prose{เหตุยา นว, อารมฺมเณ ตีณิ ฯเปฯ วิปาเก เอกํ ฯเปฯ อวิคเต นว. (สํขิตฺตํ.\csromanspacer{} สหชาตวารมฺปิ ฯเปฯ สมฺปยุตฺตวารมฺปิ ปฏิจฺจวารสทิสํ.)}

\proseitem{22}{กุสโล พาหิโร ธมฺโม กุสลสฺส พาหิรสฺส ธมฺมสฺส เหตุปจฺจเยน ปจฺจโย. (สํขิตฺตํ.)}

\csromanhangingbegin
\hangingheaditem{22}{เหตุยา สตฺต, อารมฺมเณ นว ฯเปฯ อวิคเต เตรส. (สํขิตฺตํ.)}
\hangingbody{(ยถา กุสลตฺติเก ปญฺหาวารํ, เอวํ วิตฺถาเรตพฺพํ.)}
\csromanhangingend

\csromanheader{2}{1. กุสลตฺติก 67. อุปาทาทุก}
\csromanheader{2}{1-7. ปฏิจฺจาทิวาร ปจฺจยจตุกฺก}
\proseitem{23}{อพฺยากโต อุปาทา ธมฺโม อพฺยากตสฺส อุปาทา ธมฺมสฺส อาหารปจฺจเยน ปจฺจโย.\csromanspacer{} อินฺทฺริยปจฺจเยน ปจฺจโย.\csromanspacer{} อตฺถิปจฺจเยน ปจฺจโย.\csromanspacer{} อวิคตปจฺจเยน ปจฺจโย.}

\prose{อาหาเร เอกํ, อินฺทฺริเย อตฺถิยา อวิคเต เอกํ. (สํขิตฺตํ.)}

\csromanheader{2}{1. กุสลตฺติก 67. อุปาทาทุก}
\proseitem{24}{\csromanfolio{593}กุสลํ โน-อุปาทา ธมฺมํ ปฏิจฺจ กุสโล โน-อุปาทา ธมฺโม อุปฺปชฺชติ เหตุปจฺจยา.\csromanspacer{} ตีณิ.}

\prose{อกุสลํ โน-อุปาทา ธมฺมํ ปฏิจฺจ อกุสโล โน-อุปาทา ธมฺโม อุปฺปชฺชติ เหตุปจฺจยา.\csromanspacer{} ตีณิ.}

\prose{อพฺยากตํ โน-อุปาทา ธมฺมํ ปฏิจฺจ อพฺยากโต โน-อุปาทา ธมฺโม อุปฺปชฺชติ เหตุปจฺจยา. (1)}

\prose{กุสลํ โน-อุปาทา จ อพฺยากตํ โน-อุปาทา จ ธมฺมํ ปฏิจฺจ อพฺยากโต โน-อุปาทา ธมฺโม อุปฺปชฺชติ เหตุปจฺจยา. (1)}

\prose{อกุสลํ โน-อุปาทา จ อพฺยากตํ โน-อุปาทา จ ธมฺมํ ปฏิจฺจ อพฺยากโต โน-อุปาทา ธมฺโม อุปฺปชฺชติ เหตุปจฺจยา. (1) (สํขิตฺตํ.)}

\proseitem{25}{เหตุยา นว. (สํขิตฺตํ.\csromanspacer{} สหชาตวารมฺปิ ฯเปฯ สมฺปยุตฺตวารมฺปิ ปฏิจฺจวารสทิสํ.)}

\proseitem{26}{กุสโล โน-อุปาทา ธมฺโม กุสลสฺส โน-อุปาทา ธมฺมสฺส เหตุปจฺจเยน ปจฺจโย.\csromanspacer{} ตีณิ.}

\prose{อกุสโล โน-อุปาทา ธมฺโม อกุสลสฺส โน-อุปาทา ธมฺมสฺส เหตุปจฺจเยน ปจฺจโย.\csromanspacer{} ตีณิ.}

\prose{อพฺยากโต โน-อุปาทา ธมฺโม อพฺยากตสฺส โน-อุปาทา ธมฺมสฺส เหตุปจฺจเยน ปจฺจโย. (1)}

\proseitem{27}{กุสโล โน-อุปาทา ธมฺโม กุสลสฺส โน-อุปาทา ธมฺมสฺส อารมฺมณปจฺจเยน ปจฺจโย.\csromanspacer{} ตีณิ.}

\prose{อกุสโล โน-อุปาทา ธมฺโม อกุสลสฺส โน-อุปาทา ธมฺมสฺส อารมฺมณปจฺจเยน ปจฺจโย.\csromanspacer{} ตีณิ.}

\prose{อพฺยากโต โน-อุปาทา ธมฺโม อพฺยากตสฺส โน-อุปาทา ธมฺมสฺส อารมฺมณปจฺจเยน ปจฺจโย.\csromanspacer{} ตีณิ. (สํขิตฺตํ.)}

\proseitem{28}{เหตุยา สตฺต, อารมฺมเณ นว, อธิปติยา ทส ฯเปฯ อวิคเต เอกาทส. (สํขิตฺตํ.)}

\prose{(ยถา กุสลตฺติเก ปญฺหาวารํ, เอวํ วิตฺถาเรตพฺพํ.)}

\csromanheader{2}{ธมฺมานุโลม ติกทุกปฏฺฐานปาฬิ \textbf{1. กุสลตฺติก 68. อุปาทินฺนทุก}}
\csromanheader{2}{1-7. ปฏิจฺจาทิวาร ปจฺจยจตุกฺก}
\proseitem{29}{\csromanfolio{594}อพฺยากตํ อุปาทินฺนํ ธมฺมํ ปฏิจฺจ อพฺยากโต อุปาทินฺโน ธมฺโม อุปฺปชฺชติ เหตุปจฺจยา. (1) (สํขิตฺตํ.)}

\prose{เหตุยา เอกํ, อารมฺมเณ เอกํ ฯเปฯ อวิคเต เอกํ. (สํขิตฺตํ.) (สหชาตวาเรปิ ฯเปฯ ปญฺหาวาเรปิ สพฺพตฺถ เอกํ.)}

\proseitem{30}{กุสลํ อนุปาทินฺนํ ธมฺมํ ปฏิจฺจ กุสโล อนุปาทินฺโน ธมฺโม อุปฺปชฺชติ เหตุปจฺจยา.\csromanspacer{} ตีณิ.}

\prose{อกุสลํ อนุปาทินฺนํ ธมฺมํ ปฏิจฺจ อกุสโล อนุปาทินฺโน ธมฺโม อุปฺปชฺชติ เหตุปจฺจยา.\csromanspacer{} ตีณิ.}

\prose{อพฺยากตํ อนุปาทินฺนํ ธมฺมํ ปฏิจฺจ อพฺยากโต อนุปาทินฺโน ธมฺโม อุปฺปชฺชติ เหตุปจฺจยา. (1)}

\prose{กุสลํ อนุปาทินฺนญฺจ อพฺยากตํ อนุปาทินฺนญฺจ ธมฺมํ ปฏิจฺจ อพฺยากโต อนุปาทินฺโน ธมฺโม อุปฺปชฺชติ เหตุปจฺจยา. (1)}

\prose{อกุสลํ อนุปาทินฺนญฺจ อพฺยากตํ อนุปาทินฺนญฺจ ธมฺมํ ปฏิจฺจ อพฺยากโต อนุปาทินฺโน ธมฺโม อุปฺปชฺชติ เหตุปจฺจยา. (1) (สํขิตฺตํ.)}

\proseitem{31}{เหตุยา นว, อารมฺมเณ ตีณิ ฯเปฯ อวิคเต นว. (สํขิตฺตํ.\csromanspacer{} สหชาตวารมฺปิ ฯเปฯ สมฺปยุตฺตวารมฺปิ ปฏิจฺจวารสทิสํ.)}

\csromanheader{2}{\textbf{เหตุ อารมฺมณ}}
\proseitem{32}{กุสโล อนุปาทินฺโน ธมฺโม กุสลสฺส อนุปาทินฺนสฺส ธมฺมสฺส เหตุปจฺจเยน ปจฺจโย.\csromanspacer{} ตีณิ.}

\prose{อกุสโล อนุปาทินฺโน ธมฺโม อกุสลสฺส อนุปาทินฺนสฺส ธมฺมสฺส เหตุปจฺจเยน ปจฺจโย.\csromanspacer{} ตีณิ.}

\prose{อพฺยากโต อนุปาทินฺโน ธมฺโม อพฺยากตสฺส อนุปาทินฺนสฺส ธมฺมสฺส เหตุปจฺจเยน ปจฺจโย. (1)}

\csromanheader{2}{1. กุสลตฺติก 75. กิเลสทุก}
\csromanmatikaanchor{mtk.147}
\csromanmatikaanchor{mtk.148}
\csromantocmarkref{subparagraph}{69-74. อุปาทานทุก}{mtk.147}
\bookmark[dest={mtk.147},level=5]{69-74. อุปาทานทุก}
\csromantocmarkref{subparagraph}{75-82. กิเลสทุกาทิ}{mtk.148}
\bookmark[dest={mtk.148},level=5]{75-82. กิเลสทุกาทิ}
\proseitem{33}{\csromanfolio{595}กุสโล อนุปาทินฺโน ธมฺโม กุสลสฺส อนุปาทินฺนสฺส ธมฺมสฺส อารมฺมณปจฺจเยน ปจฺจโย.\csromanspacer{} ตีณิ.}

\prose{อกุสโล อนุปาทินฺโน ธมฺโม อกุสลสฺส อนุปาทินฺนสฺส ธมฺมสฺส อารมฺมณปจฺจเยน ปจฺจโย.\csromanspacer{} ตีณิ.}

\prose{อพฺยากโต อนุปาทินฺโน ธมฺโม อพฺยากตสฺส อนุปาทินฺนสฺส ธมฺมสฺส อารมฺมณปจฺจเยน ปจฺจโย.\csromanspacer{}\csromanspacer{}อพฺยากโต อนุปาทินฺโน ธมฺโม กุสลสฺส อนุปาทินฺนสฺส ธมฺมสฺส อารมฺมณปจฺจเยน ปจฺจโย.\csromanspacer{}\csromanspacer{}อพฺยากโต อนุปาทินฺโน ธมฺโม อกุสลสฺส อนุปาทินฺนสฺส ธมฺมสฺส อารมฺมณปจฺจเยน ปจฺจโย.\csromanspacer{} ตีณิ. (สํขิตฺตํ.)}

\proseitem{34}{เหตุยา สตฺต, อารมฺมเณ นว, อธิปติยา ทส, อนนฺตเร ฉ ฯเปฯ สหชาเต นว, นิสฺสเย นว, อุปนิสฺสเย นว ฯเปฯ อวิคเต เอกาทส. (สํขิตฺตํ.)}

\nitthitam{กุสลตฺติกมหนฺตรทุกํ นิฏฺฐิตํ.}
\csromanheader{6}{1. กุสลตฺติก 69-74. อุปาทานทุก}
\proseitem{35}{อกุสลํ อุปาทานํ ธมฺมํ ปฏิจฺจ อกุสโล อุปาทาโน ธมฺโม อุปฺปชฺชติ เหตุปจฺจยา. (สํขิตฺตํ.)}

\prose{เหตุยา เอกํ, อารมฺมเณ เอกํ ฯเปฯ อวิคเต เอกํ. (อุปาทานโคจฺฉกํ วิตฺถาเรตพฺพํ.)}

\nitthitam{กุสลตฺติก-อุปาทานโคจฺฉกํ นิฏฺฐิตํ.}
\csromanheader{2}{1. กุสลตฺติก 75. กิเลสทุก}
\csromanheader{2}{1-7. ปฏิจฺจาทิวาร ปจฺจยจตุกฺก}
\proseitem{1}{อกุสลํ กิเลสํ ธมฺมํ ปฏิจฺจ อกุสโล กิเลโส ธมฺโม อุปฺปชฺชติ เหตุปจฺจยา. (สํขิตฺตํ.)}

\prose{เหตุยา เอกํ, อารมฺมเณ เอกํ ฯเปฯ อวิคเต เอกํ. (สํขิตฺตํ.)}

\vspace{\tipitakaparskip}
\csromancenter{(สหชาตวาเรปิ ฯเปฯ ปญฺหาวาเรปิ สพฺพตฺถ เอกํ.)}
\gambhira{ธมฺมานุโลม ติกทุกปฏฺฐานปาฬิ}
\proseitem{2}{\csromanfolio{596}กุสลํ โนกิเลสํ ธมฺมํ ปฏิจฺจ กุสโล โนกิเลโส ธมฺโม อุปฺปชฺชติ เหตุปจฺจยา.\csromanspacer{} ตีณิ.}

\prose{อกุสลํ โนกิเลสํ ธมฺมํ ปฏิจฺจ อกุสโล โนกิเลโส ธมฺโม อุปฺปชฺชติ เหตุปจฺจยา.\csromanspacer{} ตีณิ.}

\prose{อพฺยากตํ โนกิเลสํ ธมฺมํ ปฏิจฺจ อพฺยากโต โนกิเลโส ธมฺโม อุปฺปชฺชติ เหตุปจฺจยา. (1)}

\prose{กุสลํ โนกิเลสญฺจ อพฺยากตํ โนกิเลสญฺจ ธมฺมํ ปฏิจฺจ อพฺยากโต โนกิเลโส ธมฺโม อุปฺปชฺชติ เหตุปจฺจยา. (1)}

\prose{อกุสลํ โนกิเลสญฺจ อพฺยากตํ โนกิเลสญฺจ ธมฺมํ ปฏิจฺจ อพฺยากโต โนกิเลโส ธมฺโม อุปฺปชฺชติ เหตุปจฺจยา. (1) (สํขิตฺตํ.)}

\proseitem{3}{เหตุยา นว, อารมฺมเณ ตีณิ, อธิปติยา นว ฯเปฯ อวิคเต นว. (สํขิตฺตํ.\csromanspacer{} สหชาตวารมฺปิ ฯเปฯ สมฺปยุตฺตวารมฺปิ ปฏิจฺจวารสทิสํ.)}

\proseitem{4}{กุสโล โนกิเลโส ธมฺโม กุสลสฺส โนกิเลสสฺส ธมฺมสฺส เหตุปจฺจเยน ปจฺจโย.\csromanspacer{} ตีณิ.}

\prose{อพฺยากโต โนกิเลโส ธมฺโม อพฺยากตสฺส โนกิเลสสฺส ธมฺมสฺส เหตุปจฺจเยน ปจฺจโย. (1) (สํขิตฺตํ.)}

\proseitem{5}{เหตุยา จตฺตาริ, อารมฺมเณ นว ฯเปฯ อวิคเต เตรส. (สํขิตฺตํ.\csromanspacer{} ยถา กุสลตฺติเก ปญฺหาวารํ, เอวํ วิตฺถาเรตพฺพํ.)}

\csromanheader{2}{1. กุสลตฺติก 76. สํกิเลสิกทุก}
\csromanheader{2}{1-7. ปฏิจฺจาทิวาร ปจฺจยจตุกฺก}
\proseitem{6}{กุสลํ สํกิเลสิกํ ธมฺมํ ปฏิจฺจ กุสโล สํกิเลสิโก ธมฺโม อุปฺปชฺชติ เหตุปจฺจยา.\csromanspacer{} ตีณิ.}

\prose{อกุสลํ สํกิเลสิกํ ธมฺมํ ปฏิจฺจ อกุสโล สํกิเลสิโก ธมฺโม อุปฺปชฺชติ เหตุปจฺจยา.\csromanspacer{} ตีณิ.}

\csromanheader{2}{1. กุสลตฺติก 76. สํกิเลสิกทุก}
\prose{\csromanfolio{597}อพฺยากตํ สํกิเลสิกํ ธมฺมํ ปฏิจฺจ อพฺยากโต สํกิเลสิโก ธมฺโม อุปฺปชฺชติ เหตุปจฺจยา. (1)}

\prose{กุสลํ สํกิเลสิกญฺจ อพฺยากตํ สํกิเลสิกญฺจ ธมฺมํ ปฏิจฺจ อพฺยากโต สํกิเลสิโก ธมฺโม อุปฺปชฺชติ เหตุปจฺจยา. (1)}

\prose{อกุสลํ สํกิเลสิกญฺจ อพฺยากตํ สํกิเลสิกญฺจ ธมฺมํ ปฏิจฺจ อพฺยากโต สํกิเลสิโก ธมฺโม อุปฺปชฺชติ เหตุปจฺจยา. (1)}

\proseitem{7}{เหตุยา นว, อารมฺมเณ ตีณิ ฯเปฯ อวิคเต นว. (สํขิตฺตํ.\csromanspacer{} สหชาตวารมฺปิ ฯเปฯ สมฺปยุตฺตวารมฺปิ วิตฺถาเรตพฺพํ.)}

\proseitem{8}{กุสโล สํกิเลสิโก ธมฺโม อกุสลสฺส สํกิเลสิกสฺส ธมฺมสฺส เหตุปจฺจเยน ปจฺจโย.\csromanspacer{} ตีณิ.}

\prose{อกุสโล สํกิเลสิโก ธมฺโม อกุสลสฺส สํกิเลสิกสฺส ธมฺมสฺส เหตุปจฺจเยน ปจฺจโย.\csromanspacer{} ตีณิ.}

\prose{อพฺยากโต สํกิเลสิโก ธมฺโม อพฺยากตสฺส สํกิเลสิกสฺส ธมฺมสฺส เหตุปจฺจเยน ปจฺจโย. (1) (สํขิตฺตํ.)}

\proseitem{9}{เหตุยา สตฺต, อารมฺมเณ นว, อธิปติยา นว ฯเปฯ อวิคเต เตรส. (สํขิตฺตํ.\csromanspacer{} ยถา กุสลตฺติเก ปญฺหาวารํ, เอวํ วิตฺถาเรตพฺพํ.)}

\proseitem{10}{กุสลํ อสํกิเลสิกํ ธมฺมํ ปฏิจฺจ กุสโล อสํกิเลสิโก ธมฺโม อุปฺปชฺชติ เหตุปจฺจยา. (1)}

\prose{อพฺยากตํ อสํกิเลสิกํ ธมฺมํ ปฏิจฺจ อพฺยากโต อสํกิเลสิโก ธมฺโม อุปฺปชฺชติ เหตุปจฺจยา. (1) (สํขิตฺตํ.)}

\proseitem{11}{เหตุยา เทฺว, อารมฺมเณ เทฺว ฯเปฯ อวิคเต เทฺว. (สํขิตฺตํ.\csromanspacer{} สหชาตวารมฺปิ ฯเปฯ สมฺปยุตฺตวารมฺปิ ปฏิจฺจวารสทิสํ.)}

\proseitem{12}{กุสโล อสํกิเลสิโก ธมฺโม กุสลสฺส อสํกิเลสิกสฺส ธมฺมสฺส เหตุปจฺจเยน ปจฺจโย. (1)}

\prose{อพฺยากโต อสํกิเลสิโก ธมฺโม อพฺยากตสฺส อสํกิเลสิกสฺส ธมฺมสฺส เหตุปจฺจเยน ปจฺจโย. (1)}

\gambhira{ธมฺมานุโลม ติกทุกปฏฺฐานปาฬิ}
\prose{\csromanfolio{598}อพฺยากโต อสํกิเลสิโก ธมฺโม อพฺยากตสฺส อสํกิเลสิกสฺส ธมฺมสฺส อารมฺมณปจฺจเยน ปจฺจโย.\csromanspacer{}\csromanspacer{}อพฺยากโต อสํกิเลสิโก ธมฺโม กุสลสฺส อสํกิเลสิกสฺส ธมฺมสฺส อารมฺมณปจฺจเยน ปจฺจโย. (2) (สํขิตฺตํ.)}

\proseitem{13}{เหตุยา เทฺว, อารมฺมเณ เทฺว, อธิปติยา ตีณิ, อนนฺตเร เทฺว ฯเปฯ อุปนิสฺสเย จตฺตาริ ฯเปฯ อวิคเต เทฺว. (สํขิตฺตํ.\csromanspacer{} ยถา กุสลตฺติเก ปญฺหาวารํ, เอวํ วิตฺถาเรตพฺพํ.)}

\csromanheader{2}{1. กุสลตฺติก 77. สํกิลิฏฺฐทุก}
\csromanheader{2}{1-7. ปฏิจฺจาทิวาร เหตุ-อารมฺมณ}
\proseitem{14}{อกุสลํ สํกิลิฏฺฐํ ธมฺมํ ปฏิจฺจ อกุสโล สํกิลิฏฺโฐ ธมฺโม อุปฺปชฺชติ เหตุปจฺจยา. (สํขิตฺตํ.)}

\prose{เหตุยา เอกํ, อารมฺมเณ เอกํ ฯเปฯ อวิคเต เอกํ. (สํขิตฺตํ.) (สหชาตวาเรปิ ฯเปฯ ปญฺหาวาเรปิ สพฺพตฺถ เอกํ.)}

\proseitem{15}{กุสลํ อสํกิลิฏฺฐํ ธมฺมํ ปฏิจฺจ กุสโล อสํกิลิฏฺโฐ ธมฺโม อุปฺปชฺชติ เหตุปจฺจยา.\csromanspacer{}\csromanspacer{}กุสลํ อสํกิลิฏฺฐํ ธมฺมํ ปฏิจฺจ อพฺยากโต อสํกิลิฏฺโฐ ธมฺโม อุปฺปชฺชติ เหตุปจฺจยา.\csromanspacer{}\csromanspacer{}กุสลํ อสํกิลิฏฺฐํ ธมฺมํ ปฏิจฺจ กุสโล อสํกิลิฏฺโฐ จ อพฺยากโต อสํกิลิฏฺโฐ จ ธมฺมา อุปฺปชฺชนฺติ เหตุปจฺจยา. (3)}

\prose{อพฺยากตํ อสํกิลิฏฺฐํ ธมฺมํ ปฏิจฺจ อพฺยากโต อสํกิลิฏฺโฐ ธมฺโม อุปฺปชฺชติ เหตุปจฺจยา.\csromanspacer{}\csromanspacer{}กุสลํ อสํกิลิฏฺฐญฺจ อพฺยากตํ อสํกิลิฏฺฐญฺจ ธมฺมํ ปฏิจฺจ อพฺยากโต อสํกิลิฏฺโฐ ธมฺโม อุปฺปชฺชติ เหตุปจฺจยา. (2)}

\proseitem{16}{กุสลํ อสํกิลิฏฺฐํ ธมฺมํ ปฏิจฺจ กุสโล อสํกิลิฏฺโฐ ธมฺโม อุปฺปชฺชติ อารมฺมณปจฺจยา. (1)}

\prose{อพฺยากตํ อสํกิลิฏฺฐํ ธมฺมํ ปฏิจฺจ อพฺยากโต อสํกิลิฏฺโฐ ธมฺโม อุปฺปชฺชติ อารมฺมณปจฺจยา. (1) (สํขิตฺตํ.)}

\csromanheader{2}{1. กุสลตฺติก 79. กิเลสสํกิเลสิกทุก}
\proseitem{17}{\csromanfolio{599}เหตุยา ปญฺจ, อารมฺมเณ เทฺว, อธิปติยา ปญฺจ ฯเปฯ อวิคเต ปญฺจ. (สํขิตฺตํ.\csromanspacer{} สหชาตวารมฺปิ ฯเปฯ สมฺปยุตฺตวารมฺปิ ปฏิจฺจวารสทิสํ.)}

\proseitem{18}{กุสโล อสํกิลิฏฺโฐ ธมฺโม กุสลสฺส อสํกิลิฏฺฐสฺส ธมฺมสฺส เหตุปจฺจเยน ปจฺจโย.\csromanspacer{} ตีณิ.}

\prose{อพฺยากโต อสํกิลิฏฺโฐ ธมฺโม อพฺยากตสฺส อสํกิลิฏฺฐสฺส ธมฺมสฺส เหตุปจฺจเยน ปจฺจโย. (1)}

\prose{กุสโล อสํกิลิฏฺโฐ ธมฺโม กุสลสฺส อสํกิลิฏฺฐสฺส ธมฺมสฺส อารมฺมณปจฺจเยน ปจฺจโย.\csromanspacer{} เทฺว.}

\prose{อพฺยากโต อสํกิลิฏฺโฐ ธมฺโม อพฺยากตสฺส อสํกิลิฏฺฐสฺส ธมฺมสฺส อารมฺมณปจฺจเยน ปจฺจโย.\csromanspacer{} เทฺว. (สํขิตฺตํ.)}

\proseitem{19}{เหตุยา จตฺตาริ, อารมฺมเณ จตฺตาริ, อธิปติยา ปญฺจ, อนนฺตเร จตฺตาริ ฯเปฯ อุปนิสฺสเย จตฺตาริ ฯเปฯ อวิคเต สตฺต. (สํขิตฺตํ.\csromanspacer{} ยถา กุสลตฺติเก ปญฺหาวารํ, เอวํ วิตฺถาเรตพฺพํ.)}

\csromanheader{2}{1. กุสลตฺติก 78. กิเลสสมฺปยุตฺตทุก}
\csromanheader{2}{1-7. ปฏิจฺจาทิวาร-เหตุ}
\proseitem{20}{อกุสลํ กิเลสสมฺปยุตฺตํ ธมฺมํ ปฏิจฺจ อกุสโล กิเลสสมฺปยุตฺโต ธมฺโม อุปฺปชฺชติ เหตุปจฺจยา. (สํกิลิฏฺฐทุกสทิสํ.)}

\csromanheader{2}{1. กุสลตฺติก 79. กิเลสสํกิเลสิกทุก}
\csromanheader{2}{1-7. ปฏิจฺจาทิวาร-เหตุ}
\proseitem{21}{อกุสลํ กิเลสญฺเจว สํกิเลสิกญฺจ ธมฺมํ ปฏิจฺจ อกุสโล กิเลโส เจว สํกิเลสิโก จ ธมฺโม อุปฺปชฺชติ เหตุปจฺจยา. (สํขิตฺตํ.)}

\csromanheader{2}{เหตุยา เอกํ ฯเปฯ อวิคเต เอกํ. (สํขิตฺตํ.)}
\vspace{\tipitakaparskip}
\csromancenter{(สหชาตวาเรปิ ฯเปฯ ปญฺหาวาเรปิ สพฺพตฺถ เอกํ.)}
\gambhira{ธมฺมานุโลม ติกทุกปฏฺฐานปาฬิ}
\proseitem{22}{\csromanfolio{600}กุสลํ สํกิเลสิกญฺเจว โน จ กิเลสํ ธมฺมํ ปฏิจฺจ กุสโล สํกิเลสิโก เจว โน จ กิเลโส ธมฺโม อุปฺปชฺชติ เหตุปจฺจยา.\csromanspacer{} ตีณิ.}

\prose{อกุสลํ สํกิเลสิกญฺเจว โน จ กิเลสํ ธมฺมํ ปฏิจฺจ อกุสโล สํกิเลสิโก เจว โน จ กิเลโส ธมฺโม อุปฺปชฺชติ เหตุปจฺจยา.\csromanspacer{} ตีณิ.}

\prose{อพฺยากตํ สํกิเลสิกญฺเจว โน จ กิเลสํ ธมฺมํ ปฏิจฺจ อพฺยากโต สํกิเลสิโก เจว โน จ กิเลโส ธมฺโม อุปฺปชฺชติ เหตุปจฺจยา.}

\prose{กุสลํ สํกิเลสิกญฺเจว โน จ กิเลสญฺจ อพฺยากตํ สํกิเลสิกญฺเจว โน จ กิเลสญฺจ ธมฺมํ ปฏิจฺจ อพฺยากโต สํกิเลสิโก เจว โน จ กิเลโส ธมฺโม อุปฺปชฺชติ เหตุปจฺจยา. (1)}

\prose{อกุสลํ สํกิเลสิกญฺเจว โน จ กิเลสญฺจ อพฺยกตํ สํกิเลสิกญฺจว โน จ กิเลสญฺจ ธมฺมํ ปฏิจฺจ อพฺยากโต สํกิเลสิโก เจว โน จ กิเลโส ธมฺโม อุปฺปชฺชติ เหตุปจฺจยา. (สํขิตฺตํ.)}

\proseitem{23}{เหตุยา นว, อารมฺมเณ ตีณิ ฯเปฯ อวิคเต นว. (สํขิตฺตํ.\csromanspacer{} สหชาตวารมฺปิ ฯเปฯ สมฺปยุตฺตวารมฺปิ ปฏิจฺจวารสทิสํ.)}

\proseitem{24}{กุสโล สํกิเลสิโก เจว โน จ กิเลโส ธมฺโม กุสลสฺส สํกิเลสิกสฺส เจว โน จ กิเลสสฺส ธมฺมสฺส เหตุปจฺจเยน ปจฺจโย.\csromanspacer{} ตีณิ.}

\prose{อพฺยากโต สํกิเลสิโก เจว โน จ กิเลโส ธมฺโม อพฺยากตสฺส สํกิเลสิกสฺส เจว โน จ กิเลสสฺส ธมฺมสฺส เหตุปจฺจเยน ปจฺจโย. (1) (สํขิตฺตํ.)}

\proseitem{25}{เหตุยา จตฺตาริ, อารมฺมเณ นว, อธิปติยา นว ฯเปฯ อวิคเต เตรส. (สํขิตฺตํ.\csromanspacer{} ยถา กุสลตฺติเก ปญฺหาวารํ, เอวํ วิตฺถาเรตพฺพํ.)}

\vspace{\tipitakaparskip}
\csromancenter{กุสลตฺติก 81.\csromanspacer{} กิเลสกิเลสสมฺปยุตฺตทุก \textbf{1.}\csromanspacer{}\textbf{ กุสลตฺติก 80.}\csromanspacer{}\textbf{ กิเลสสํกิลิฏฺฐทุก}}
\csromanheader{2}{1. ปฏิจฺจวาร-เหตุ}
\proseitem{26}{\csromanfolio{601}อกุสลํ กิเลสญฺเจว สํกิลิฏฺฐญฺจ ธมฺมํ ปฏิจฺจ อกุสโล กิเลโส เจว สํกิลิฏฺโฐ จ ธมฺโม อุปฺปชฺชติ เหตุปจฺจยา. (สํขิตฺตํ.)}

\prose{เหตุยา เอกํ, อารมฺมเณ เอกํ ฯเปฯ อวิคเต เอกํ. (สํขิตฺตํ.) (สหชาตวาเรปิ ฯเปฯ ปญฺหาวาเรปิ สพฺพตฺถ เอกํ.)}

\proseitem{27}{อกุสลํ สํกิลิฏฺฐญฺเจว โน จ กิเลสํ ธมฺมํ ปฏิจฺจ อกุสโล สํกิลิฏฺโฐ เจว โน จ กิเลโส ธมฺโม อุปฺปชฺชติ เหตุปจฺจยา.}

\csromanheader{2}{เหตุยา เอกํ ฯเปฯ อวิคเต เอกํ. (สํขิตฺตํ.)}
\vspace{\tipitakaparskip}
\csromancenter{(สหชาตวาเรปิ ฯเปฯ ปญฺหาวาเรปิ สพฺพตฺถ เอกํ.)}
\csromanheader{2}{1. กุสลตฺติก 81. กิเลสกิเลสสมฺปยุตฺตทุก}
\csromanheader{2}{1. ปฏิจฺจวาร-เหตุ}
\proseitem{28}{อกุสลํ กิเลสญฺเจว กิเลสสมฺปยุตฺตญฺจ ธมฺมํ ปฏิจฺจ อกุสโล กิเลโส เจว กิเลสสมฺปยุตฺโต จ ธมฺโม อุปฺปชฺชติ เหตุปจฺจยา.}

\prose{เหตุยา เอกํ ฯเปฯ อวิคเต เอกํ. (สํขิตฺตํ.) (สหชาตวาเรปิ ฯเปฯ ปญฺหาวาเรปิ สพฺพตฺถ เอกํ.)}

\proseitem{29}{อกุสลํ กิเลสสมฺปยุตฺตญฺเจว โน จ กิเลสํ ธมฺมํ ปฏิจฺจ อกุสโล กิเลสสมฺปยุตฺโต เจว โน จ กิเลโส ธมฺโม อุปฺปชฺชติ เหตุปจฺจยา. (สํขิตฺตํ.)}

\csromanheader{2}{เหตุยา เอกํ ฯเปฯ อวิคเต เอกํ. (สํขิตฺตํ.)}
\vspace{\tipitakaparskip}
\csromancenter{(สหชาตวาเรปิ ฯเปฯ ปญฺหาวาเรปิ สพฺพตฺถ เอกํ.)}
\csromancenter{ธมฺมานุโลม ติกทุกปฏฺฐานปาฬิ \textbf{1.}\csromanspacer{}\textbf{ กุสลตฺติก 82.}\csromanspacer{}\textbf{ กิเลสวิปฺปยุตฺตสํกิเลสิกทุก}}
\csromanheader{2}{1-7. ปฏิจฺจาทิวาร-เหตุ}
\proseitem{30}{\csromanfolio{602}กุสลํ กิเลสวิปฺปยุตฺตํ สํกิเลสิกํ ธมฺมํ ปฏิจฺจ กุสโล กิเลสวิปฺปยุตฺโต สํกิเลสิโก ธมฺโม อุปฺปชฺชติ เหตุปจฺจยา.\csromanspacer{}\csromanspacer{}กุสลํ กิเลสวิปฺปยุตฺตํ สํกิเลสิกํ ธมฺมํ ปฏิจฺจ อพฺยากโต กิเลสวิปฺปยุตฺโต สํกิเลสิโก ธมฺโม อุปฺปชฺชติ เหตุปจฺจยา.\csromanspacer{}\csromanspacer{}กุสลํ กิเลสวิปฺปยุตฺตํ สํกิเลสิกํ ธมฺมํ ปฏิจฺจ กุสโล กิเลสวิปฺปยุตฺโต สํกิเลสิโก จ อพฺยากโต กิเลสวิปฺปยุตฺโต สํกิเลสิโก จ ธมฺมา อุปฺปชฺชนฺติ เหตุปจฺจยา. (3)}

\prose{อพฺยากตํ กิเลสวิปฺปยุตฺตํ สํกิเลสิกํ ธมฺมํ ปฏิจฺจ อพฺยากโต กิเลสวิปฺปยุตฺโต สํกิเลสิโก ธมฺโม อุปฺปชฺชติ เหตุปจฺจยา. (1)}

\prose{กุสลํ กิเลสวิปฺปยุตฺตํ สํกิเลสิกญฺจ อพฺยากตํ กิเลสวิปฺปยุตฺตํ สํกิเลสิกญฺจ ธมฺมํ ปฏิจฺจ อพฺยากโต กิเลสวิปฺปยุตฺโต สํกิเลสิโก ธมฺโม อุปฺปชฺชติ เหตุปจฺจยา. (1) (สํขิตฺตํ.)}

\proseitem{31}{เหตุยา ปญฺจ, อารมฺมเณ เทฺว ฯเปฯ อวิคเต ปญฺจ. (สํขิตฺตํ.\csromanspacer{} สหชาตวารมฺปิ ฯเปฯ สมฺปยุตฺตวารมฺปิ ปฏิจฺจวารสทิสํ.)}

\proseitem{32}{กุสโล กิเลสวิปฺปยุตฺโต สํกิเลสิโก ธมฺโม กุสลสฺส กิเลสวิปฺปยุตฺตสฺส สํกิเลสิกสฺส ธมฺมสฺส เหตุปจฺจเยน ปจฺจโย.\csromanspacer{} ตีณิ.}

\prose{อพฺยากโต กิเลสวิปฺปยุตฺโต สํกิเลสิโก ธมฺโม อพฺยากตสฺส กิเลสวิปฺปยุตฺตสฺส สํกิเลสิกสฺส ธมฺมสฺส เหตุปจฺจเยน ปจฺจโย. (1) (สํขิตฺตํ.)}

\proseitem{33}{เหตุยา จตฺตาริ, อารมฺมเณ จตฺตาริ ฯเปฯ อวิคเต สตฺต. (สํขิตฺตํ.\csromanspacer{} ยถา กุสลตฺติเก ปญฺหาวารํ, เอวํ วิตฺถาเรตพฺพํ.)}

\proseitem{34}{กุสลํ กิเลสวิปฺปยุตฺตํ อสํกิเลสิกํ ธมฺมํ ปฏิจฺจ กุสโล กิเลสวิปฺปยุตฺโต อสํกิเลสิโก ธมฺโม อุปฺปชฺชติ เหตุปจฺจยา.}

\csromanheader{2}{1. กุสลตฺติก 83. ทสฺสเนนปหาตพฺพทุก}
\csromanmatikaanchor{mtk.149}
\csromantocmarkref{subparagraph}{83-99. ทสฺสเนนปหาตพฺพทุกาทิ}{mtk.149}
\bookmark[dest={mtk.149},level=5]{83-99. ทสฺสเนนปหาตพฺพทุกาทิ}
\prose{\csromanfolio{603}อพฺยากตํ กิเลสวิปฺปยุตฺตํ อสํกิเลสิกํ ธมฺมํ ปฏิจฺจ อพฺยากโต กิเลสวิปฺปยุตฺโต อสํกิเลสิโก ธมฺโม อุปฺปชฺชติ เหตุปจฺจยา.}

\prose{เหตุยา เทฺว ฯเปฯ อวิคเต เทฺว. (สํขิตฺตํ.\csromanspacer{} สหชาตวารมฺปิ ฯเปฯ สมฺปยุตฺตวารมฺปิ ปฏิจฺจวารสทิสํ.)}

\proseitem{35}{กุสโล กิเลสวิปฺปยุตฺโต อสํกิเลสิโก ธมฺโม กุสลสฺส กิเลสวิปฺปยุตฺตสฺส อสํกิเลสิกสฺส ธมฺมสฺส เหตุปจฺจเยน ปจฺจโย. (1)}

\prose{อพฺยากโต กิเลสวิปฺปยุตฺโต อสํกิเลสิโก ธมฺโม อพฺยากตสฺส กิเลสวิปฺปยุตฺตสฺส อสํกิเลสิกสฺส ธมฺมสฺส เหตุปจฺจเยน ปจฺจโย. (1) (สํขิตฺตํ.)}

\proseitem{36}{เหตุยา เทฺว, อารมฺมเณ เทฺว, อธิปติยา ตีณิ, อนนฺตเร เทฺว ฯเปฯ อุปนิสฺสเย จตฺตาริ ฯเปฯ อวิคเต เทฺว. (สํขิตฺตํ.\csromanspacer{} ยถา กุสลตฺติเก ปญฺหาวารํ, เอวํ วิตฺถาเรตพฺพํ.)}

\nitthitam{กุสลตฺติกกิเลสโคจฺฉกํ นิฏฺฐิตํ.}
\csromanheader{2}{1. กุสลตฺติก 83. ทสฺสเนนปหาตพฺพทุก}
\csromanheader{2}{1-7. ปฏิจฺจาทิวาร-เหตุ}
\proseitem{1}{อกุสลํ ทสฺสเนน ปหาตพฺพํ ธมฺมํ ปฏิจฺจ อกุสโล ทสฺสเนน ปหาตพฺโพ ธมฺโม อุปฺปชฺชติ เหตุปจฺจยา. (สํขิตฺตํ.)}

\prose{เหตุยา เอกํ, อารมฺมเณ เอกํ ฯเปฯ อวิคเต เอกํ. (สํขิตฺตํ.) (สหชาตวาเรปิ ฯเปฯ ปญฺหาวาเรปิ สพฺพตฺถ เอกํ.)}

\proseitem{2}{กุสลํ นทสฺสเนน ปหาตพฺพํ ธมฺมํ ปฏิจฺจ กุสโล นทสฺสเนน ปหาตพฺโพ ธมฺโม อุปฺปชฺชติ เหตุปจฺจยา.\csromanspacer{} ตีณิ.}

\prose{อกุสลํ นทสฺสเนน ปหาตพฺพํ ธมฺมํ ปฏิจฺจ อกุสโล นทสฺสเนน ปหาตพฺโพ ธมฺโม อุปฺปชฺชติ เหตุปจฺจยา.\csromanspacer{} ตีณิ.}

\prose{อพฺยากตํ นทสฺสเนน ปหาตพฺพํ ธมฺมํ ปฏิจฺจ อพฺยากโต นทสฺสเนน ปหาตพฺโพ ธมฺโม อุปฺปชฺชติ เหตุปจฺจยา. (1)}

\gambhira{ธมฺมานุโลม ติกทุกปฏฺฐานปาฬิ}
\prose{\csromanfolio{604}กุสลํ นทสฺสเนน ปหาตพฺพญฺจ อพฺยากตํ นทสฺสเนน ปหาตพฺพญฺจ ธมฺมํ ปฏิจฺจ อพฺยากโต นทสฺสเนน ปหาตพฺโพ ธมฺโม อุปฺปชฺชติ เหตุปจฺจยา. (1)}

\prose{อกุสลํ นทสฺสเนน ปหาตพฺพญฺจ อพฺยากตํ นทสฺสเนน ปหาตพฺพญฺจ ธมฺมํ ปฏิจฺจ อพฺยากโต นทสฺสเนน ปหาตพฺโพ ธมฺโม อุปฺปชฺชติ เหตุปจฺจยา. (1)}

\proseitem{3}{เหตุยา นว, อารมฺมเณ ตีณิ ฯเปฯ อวิคเต นว. (สํขิตฺตํ.\csromanspacer{} สหชาตวารมฺปิ ฯเปฯ สมฺปยุตฺตวารมฺปิ ปฏิจฺจวารสทิสํ.)}

\proseitem{4}{กุสโล นทสฺสเนน ปหาตพฺโพ ธมฺโม กุสลสฺส นทสฺสเนน ปหาตพฺพสฺส ธมฺมสฺส เหตุปจฺจเยน ปจฺจโย.\csromanspacer{} ตีณิ.}

\prose{อกุสโล นทสฺสเนน ปหาตพฺโพ ธมฺโม อกุสลสฺส นทสฺสเนน ปหาตพฺพสฺส ธมฺมสฺส เหตุปจฺจเยน ปจฺจโย.\csromanspacer{} ตีณิ.}

\prose{อพฺยากโต นทสฺสเนน ปหาตพฺโพ ธมฺโม อพฺยากตสฺส นทสฺสเนน ปหาตพฺพสฺส ธมฺมสฺส เหตุปจฺจเยน ปจฺจโย. (1) (สํขิตฺตํ.)}

\proseitem{5}{เหตุยา สตฺต, อารมฺมเณ นว, อธิปติยา ทส, อนนฺตเร สตฺต ฯเปฯ อวิคเต เตรส. (สํขิตฺตํ.)}

\prose{(ยถา กุสลตฺติเก ปญฺหาวารํ, เอวํ วิตฺถาเรตพฺพํ.)}

\csromanheader{2}{1. กุสลตฺติก 84. ภาวนายปหาตพฺพทุก}
\csromanheader{2}{1. ปฏิจฺจวาร-เหตุ}
\proseitem{6}{อกุสลํ ภาวนาย ปหาตพฺพํ ธมฺมํ ปฏิจฺจ อกุสโล ภาวนาย ปหาตพฺโพ ธมฺโม อุปฺปชฺชติ เหตุปจฺจยา. (สํขิตฺตํ.)}

\prose{เหตุยา เอกํ, อารมฺมเณ เอกํ ฯเปฯ อวิคเต เอกํ. (สํขิตฺตํ.) (สหชาตวาเรปิ ฯเปฯ ปญฺหาวาเรปิ สพฺพตฺถ เอกํ.)}

\proseitem{7}{กุสลํ นภาวนาย ปหาตพฺพํ ธมฺมํ ปฏิจฺจ กุสโล นภาวนาย ปหาตพฺโพ ธมฺโม อุปฺปชฺชติ เหตุปจฺจยา. (สํขิตฺตํ.)}

\prose{เหตุยา นว, อารมฺมเณ ตีณิ ฯเปฯ อวิคเต นว. (สํขิตฺตํ.)}

\vspace{\tipitakaparskip}
\csromancenter{(สหชาตวารมฺปิ ฯเปฯ ปญฺหาวารมฺปิ วิตฺถาเรตพฺพํ.)}
\csromancenter{กุสลตฺติก 87.\csromanspacer{} สวิตกฺกทุก \textbf{1.}\csromanspacer{}\textbf{ กุสลตฺติก 85.}\csromanspacer{}\textbf{ ทสฺสเนนปหาตพฺพเหตุกทุก}}
\csromanheader{2}{1. ปฏิจฺจวาร-เหตุ}
\proseitem{8}{\csromanfolio{605}อกุสลํ ทสฺสเนน ปหาตพฺพเหตุกํ ธมฺมํ ปฏิจฺจ อกุสโล ทสฺสเนน ปหาตพฺพเหตุโก ธมฺโม อุปฺปชฺชติ เหตุปจฺจยา. (สํขิตฺตํ.)}

\prose{เหตุยา เอกํ, อารมฺมเณ เอกํ ฯเปฯ อวิคเต เอกํ. (สํขิตฺตํ.) (สหชาตวาเรปิ ฯเปฯ ปญฺหาวาเรปิ สพฺพตฺถ เอกํ.)}

\proseitem{9}{กุสลํ นทสฺสเนน ปหาตพฺพเหตุกํ ธมฺมํ ปฏิจฺจ กุสโล นทสฺสเนน ปหาตพฺพเหตุโก ธมฺโม อุปฺปชฺชติ เหตุปจฺจยา. (สํขิตฺตํ.)}

\prose{เหตุยา นว ฯเปฯ อวิคเต นว. (สํขิตฺตํ.\csromanspacer{} สหชาตวารมฺปิ ฯเปฯ ปญฺหาวารมฺปิ วิตฺถาเรตพฺพํ.)}

\csromanheader{2}{1. กุสลตฺติก 86. ภาวนายปหาตพฺพเหตุกทุก}
\csromanheader{2}{1. ปฏิจฺจวาร ปจฺจยจตุกฺก}
\proseitem{10}{อกุสลํ ภาวนาย ปหาตพฺพเหตุกํ ธมฺมํ ปฏิจฺจ อกุสโล ภาวนาย ปหาตพฺพเหตุโก ธมฺโม อุปฺปชฺชติ เหตุปจฺจยา. (สํขิตฺตํ.)}

\prose{เหตุยา เอกํ, อารมฺมเณ เอกํ ฯเปฯ อวิคเต เอกํ. (สํขิตฺตํ.) (สหชาตวาเรปิ ฯเปฯ ปญฺหาวาเรปิ สพฺพตฺถ เอกํ.)}

\proseitem{11}{กุสลํ นภาวนาย ปหาตพฺพเหตุกํ ธมฺมํ ปฏิจฺจ กุสโล นภาวนาย ปหาตพฺพเหตุโก ธมฺโม อุปฺปชฺชติ เหตุปจฺจยา.}

\prose{เหตุยา นว, อารมฺมเณ ตีณิ ฯเปฯ อวิคเต นว. (สํขิตฺตํ.)}

\vspace{\tipitakaparskip}
\csromancenter{(สหชาตวาเรปิ ฯเปฯ ปญฺหาวาเรปิ วิตฺถาเรตพฺพํ.)}
\csromanheader{2}{1. กุสลตฺติก 87. สวิตกฺกทุก}
\csromanheader{2}{1-7. ปฏิจฺจาทิวาร เหตุ-อารมฺมณ}
\proseitem{12}{กุสลํ สวิตกฺกํ ธมฺมํ ปฏิจฺจ กุสโล สวิตกฺโก ธมฺโม อุปฺปชฺชติ เหตุปจฺจยา. (1)}

\gambhira{ธมฺมานุโลม ติกทุกปฏฺฐานปาฬิ}
\prose{\csromanfolio{606}อกุสลํ สวิตกฺกํ ธมฺมํ ปฏิจฺจ อกุสโล สวิตกฺโก ธมฺโม อุปฺปชฺชติ เหตุปจฺจยา. (1)}

\prose{อพฺยากตํ สวิตกฺกํ ธมฺมํ ปฏิจฺจ อพฺยากโต สวิตกฺโก ธมฺโม อุปฺปชฺชติ เหตุปจฺจยา. (1) (สํขิตฺตํ.)}

\proseitem{13}{เหตุยา ตีณิ, อารมฺมเณ ตีณิ ฯเปฯ อวิคเต ตีณิ. (สํขิตฺตํ.\csromanspacer{} สหชาตวารมฺปิ ฯเปฯ สมฺปยุตฺตวารมฺปิ ปฏิจฺจวารสทิสํ.)}

\proseitem{14}{กุสโล สวิตกฺโก ธมฺโม กุสลสฺส สวิตกฺกสฺส ธมฺมสฺส เหตุปจฺจเยน ปจฺจโย. (1)}

\prose{อกุสโล สวิตกฺโก ธมฺโม อกุสลสฺส สวิตกฺกสฺส ธมฺมสฺส เหตุปจฺจเยน ปจฺจโย. (1)}

\prose{อพฺยากโต สวิตกฺโก ธมฺโม อพฺยากตสฺส สวิตกฺกสฺส ธมฺมสฺส เหตุปจฺจเยน ปจฺจโย. (1)}

\prose{กุสโล สวิตกฺโก ธมฺโม กุสลสฺส สวิตกฺกสฺส ธมฺมสฺส อารมฺมณปจฺจเยน ปจฺจโย.\csromanspacer{} ตีณิ.}

\prose{อกุสโล สวิตกฺโก ธมฺโม อกุสลสฺส สวิตกฺกสฺส ธมฺมสฺส อารมฺมณปจฺจเยน ปจฺจโย.\csromanspacer{} ตีณิ.}

\prose{อพฺยากโต สวิตกฺโก ธมฺโม อพฺยากตสฺส สวิตกฺกสฺส ธมฺมสฺส อารมฺมณปจฺจเยน ปจฺจโย.\csromanspacer{} ตีณิ. (สํขิตฺตํ.)}

\proseitem{15}{เหตุยา ตีณิ, อารมฺมเณ นว, อธิปติยา สตฺต ฯเปฯ อวิคเต ตีณิ. (สํขิตฺตํ.)}

\prose{(ยถา กุสลตฺติเก ปญฺหาวารํ, เอวํ วิตฺถาเรตพฺพํ.)}

\csromanheader{2}{\textbf{อวิตกฺกปท-เหตุ}}
\proseitem{16}{กุสลํ อวิตกฺกํ ธมฺมํ ปฏิจฺจ กุสโล อวิตกฺโก ธมฺโม อุปฺปชฺชติ เหตุปจฺจยา. (สํขิตฺตํ.)}

\csromanheader{2}{1. กุสลตฺติก 89. สปฺปีติกทุก}
\prose{\csromanfolio{607}เหตุยา สตฺต, อารมฺมเณ เทฺว ฯเปฯ วิปาเก เอกํ, อวิคเต สตฺต. (สํขิตฺตํ.\csromanspacer{} สหชาตวารมฺปิ ฯเปฯ สมฺปยุตฺตวารมฺปิ ปฏิจฺจวารสทิสํ.)}

\proseitem{17}{กุสโล อวิตกฺโก ธมฺโม กุสลสฺส อวิตกฺกสฺส ธมฺมสฺส เหตุปจฺจเยน ปจฺจโย.\csromanspacer{} ตีณิ.}

\prose{อพฺยากโต อวิตกฺโก ธมฺโม อพฺยากตสฺส อวิตกฺกสฺส ธมฺมสฺส เหตุปจฺจเยน ปจฺจโย. (1) (สํขิตฺตํ.)}

\proseitem{18}{เหตุยา จตฺตาริ, อารมฺมเณ นว ฯเปฯ อวิคเต ทส. (สํขิตฺตํ.\csromanspacer{} ยถา กุสลตฺติเก ปญฺหาวารํ, เอวํ วิตฺถาเรตพฺพํ.)}

\csromanheader{2}{1. กุสลตฺติก 88. สวิจารทุก}
\csromanheader{2}{1. ปฏิจฺจวาร ปจฺจยจตุกฺก}
\proseitem{19}{กุสลํ สวิจารํ ธมฺมํ ปฏิจฺจ กุสโล สวิจาโร ธมฺโม อุปฺปชฺชติ เหตุปจฺจยา. (กุสลสวิตกฺกสทิสํ.)}

\csromanheader{2}{1. กุสลตฺติก 89. สปฺปีติกทุก}
\csromanheader{2}{1-7. ปฏิจฺจาทิวาร ปจฺจยจตุกฺก}
\proseitem{20}{กุสลํ สปฺปีติกํ ธมฺมํ ปฏิจฺจ กุสโล สปฺปีติโก ธมฺโม อุปฺปชฺชติ เหตุปจฺจยา. (1)}

\prose{อกุสลํ สปฺปีติกํ ธมฺมํ ปฏิจฺจ อกุสโล สปฺปีติโก ธมฺโม อุปฺปชฺชติ เหตุปจฺจยา. (1)}

\prose{อพฺยากตํ สปฺปีติกํ ธมฺมํ ปฏิจฺจ อพฺยากโต สปฺปีติโก ธมฺโม อุปฺปชฺชติ เหตุปจฺจยา. (1) (สํขิตฺตํ.)}

\gambhira{ธมฺมานุโลม ติกทุกปฏฺฐานปาฬิ}
\proseitem{21}{\csromanfolio{608}เหตุยา ตีณิ, อารมฺมเณ ตีณิ ฯเปฯ อวิคเต ตีณิ. (สํขิตฺตํ.\csromanspacer{} สหชาตวารมฺปิ ฯเปฯ สมฺปยุตฺตวารมฺปิ ปฏิจฺจวารสทิสํ.)}

\proseitem{22}{กุสโล สปฺปีติโก ธมฺโม กุสลสฺส สปฺปีติกสฺส ธมฺมสฺส เหตุปจฺจเยน ปจฺจโย. (1)}

\prose{อกุสโล สปฺปีติโก ธมฺโม อกุสลสฺส สปฺปีติกสฺส ธมฺมสฺส เหตุปจฺจเยน ปจฺจโย. (1)}

\prose{อพฺยากโต สปฺปีติโก ธมฺโม อพฺยากตสฺส สปฺปีติกสฺส ธมฺมสฺส เหตุปจฺจเยน ปจฺจโย. (1) (สํขิตฺตํ.)}

\proseitem{23}{เหตุยา ตีณิ, อารมฺมเณ นว ฯเปฯ อวิคเต ตีณิ. (สํขิตฺตํ.)}

\csromanheader{2}{\textbf{อปฺปีติกปท-เหตุ}}
\proseitem{24}{กุสลํ อปฺปีติกํ ธมฺมํ ปฏิจฺจ กุสโล อปฺปีติโก ธมฺโม อุปฺปชฺชติ เหตุปจฺจยา.\csromanspacer{} ตีณิ.}

\prose{อกุสลํ อปฺปีติกํ ธมฺมํ ปฏิจฺจ อกุสโล อปฺปีติโก ธมฺโม อุปฺปชฺชติ เหตุปจฺจยา.\csromanspacer{} ตีณิ.}

\prose{อพฺยากตํ อปฺปีติกํ ธมฺมํ ปฏิจฺจ อพฺยากโต อปฺปีติโก ธมฺโม อุปฺปชฺชติ เหตุปจฺจยา. (1)}

\prose{กุสลํ อปฺปีติกญฺจ อพฺยากตํ อปฺปีติกญฺจ ธมฺมํ ปฏิจฺจ อพฺยากโต อปฺปีติโก ธมฺโม อุปฺปชฺชติ เหตุปจฺจยา. (1)}

\prose{อกุสลํ อปฺปีติกญฺจ อพฺยากตํ อปฺปีติกญฺจ ธมฺมํ ปฏิจฺจ อพฺยากโต อปฺปีติโก ธมฺโม อุปฺปชฺชติ เหตุปจฺจยา. (1) (สํขิตฺตํ.)}

\proseitem{25}{เหตุยา นว, อารมฺมเณ ตีณิ ฯเปฯ อวิคเต นว. (สํขิตฺตํ.\csromanspacer{} สหชาตวารมฺปิ ฯเปฯ สมฺปยุตฺตวารมฺปิ ปฏิจฺจวารสทิสํ.)}

\csromanheader{2}{1. กุสลตฺติก 93. กามาวจรทุก}
\proseitem{26}{\csromanfolio{609}กุสโล อปฺปีติโก ธมฺโม กุสลสฺส อปฺปีติกสฺส ธมฺมสฺส เหตุปจฺจเยน ปจฺจโย. (สํขิตฺตํ.)}

\prose{เหตุยา สตฺต, อารมฺมเณ นว ฯเปฯ อวิคเต เตรส. (สํขิตฺตํ.\csromanspacer{} ยถา กุสลตฺติเก ปญฺหาวารํ, เอวํ วิตฺถาเรตพฺพํ.)}

\csromanheader{2}{1. กุสลตฺติก 90-92. ปีติสหคตาทิทุก}
\csromanheader{2}{1. ปฏิจฺจวาร ปจฺจยจตุกฺก}
\proseitem{27}{กุสลํ ปีติสหคตํ ธมฺมํ ปฏิจฺจ กุสโล ปีติสหคโต ธมฺโม อุปฺปชฺชติ เหตุปจฺจยา.\csromanspacer{}\csromanspacer{}กุสลํ นปีติสหคตํ ธมฺมํ ปฏิจฺจ ฯเปฯ.\csromanspacer{} กุสลํ สุขสหคตํ ธมฺมํ ปฏิจฺจ ฯเปฯ.\csromanspacer{} กุสลํ นสุขสหคตํ ธมฺมํ ปฏิจฺจ ฯเปฯ.\csromanspacer{} กุสลํ อุเปกฺขาสหคตํ ธมฺมํ ปฏิจฺจ ฯเปฯ.\csromanspacer{} กุสลํ น-อุเปกฺขาสหคตํ ธมฺมํ ปฏิจฺจ. (สํขิตฺตํ.)}

\csromanheader{2}{1. กุสลตฺติก 93. กามาวจรทุก}
\csromanheader{2}{1-7. ปฏิจฺจาทิวาร ปจฺจยจตุกฺก-เหตุ}
\proseitem{28}{กุสลํ กามาวจรํ ธมฺมํ ปฏิจฺจ กุสโล กามาวจโร ธมฺโม อุปฺปชฺชติ เหตุปจฺจยา.\csromanspacer{}\csromanspacer{}กุสลํ กามาวจรํ ธมฺมํ ปฏิจฺจ อพฺยากโต กามาวจโร ธมฺโม อุปฺปชฺชติ เหตุปจฺจยา.\csromanspacer{}\csromanspacer{}กุสลํ กามาวจรํ ธมฺมํ ปฏิจฺจ กุสโล กามาวจโร จ อพฺยากโต กามาวจโร จ ธมฺมา อุปฺปชฺชนฺติ เหตุปจฺจยา. (3)}

\prose{อกุสลํ กามาวจรํ ธมฺมํ ปฏิจฺจ อกุสโล กามาวจโร ธมฺโม อุปฺปชฺชติ เหตุปจฺจยา.\csromanspacer{} ตีณิ.}

\prose{อพฺยากตํ กามาวจรํ ธมฺมํ ปฏิจฺจ อพฺยากโต กามาวจโร ธมฺโม อุปฺปชฺชติ เหตุปจฺจยา. (1)}

\prose{กุสลํ กามาวจรญฺจ อพฺยากตํ กามาวจรญฺจ ธมฺมํ ปฏิจฺจ อพฺยากโต กามาวจโร ธมฺโม อุปฺปชฺชติ เหตุปจฺจยา. (1)}

\gambhira{ธมฺมานุโลม ติกทุกปฏฺฐานปาฬิ}
\prose{\csromanfolio{610}อกุสลํ กามาวจรญฺจ อพฺยากตํ กามาวจรญฺจ ธมฺมํ ปฏิจฺจ อพฺยากโต กามาวจโร ธมฺโม อุปฺปชฺชติ เหตุปจฺจยา. (1) (สํขิตฺตํ.)}

\proseitem{29}{เหตุยา นว, อารมฺมเณ ตีณิ ฯเปฯ อวิคเต นว. (สํขิตฺตํ.\csromanspacer{} สหชาตวารมฺปิ ฯเปฯ สมฺปยุตฺตวารมฺปิ ปฏิจฺจวารสทิสํ.)}

\proseitem{30}{กุสโล กามาวจโร ธมฺโม กุสลสฺส กามาวจรสฺส ธมฺมสฺส เหตุปจฺจเยน ปจฺจโย.\csromanspacer{} ตีณิ.}

\prose{อกุสโล กามาวจโร ธมฺโม กุสลสฺส กามาวจรสฺส ธมฺมสฺส เหตุปจฺจเยน ปจฺจโย.\csromanspacer{} ตีณิ.}

\prose{อพฺยากโต กามาวจโร ธมฺโม อพฺยากตสฺส กามาวจรสฺส ธมฺมสฺส เหตุปจฺจเยน ปจฺจโย. (1) (สํขิตฺตํ.)}

\proseitem{31}{เหตุยา สตฺต, อารมฺมเณ นว, อธิปติยา นว ฯเปฯ อวิคเต เตรส. (สํขิตฺตํ.\csromanspacer{} ยถา กุสลตฺติเก ปญฺหาวารํ, เอวํ วิตฺถาเรตพฺพํ.)}

\csromanheader{2}{\textbf{นกามาวจรปท-เหตุ อารมฺมณ}}
\proseitem{32}{กุสลํ นกามาวจรํ ธมฺมํ ปฏิจฺจ กุสโล นกามาวจโร ธมฺโม อุปฺปชฺชติ เหตุปจฺจยา. (1)}

\prose{อพฺยากตํ นกามาวจรํ ธมฺมํ ปฏิจฺจ อพฺยากโต นกามาวจโร ธมฺโม อุปฺปชฺชติ เหตุปจฺจยา. (1) (สํขิตฺตํ.)}

\prose{เหตุยา เทฺว, อารมฺมเณ เทฺว ฯเปฯ อวิคเต เทฺว. (สํขิตฺตํ.\csromanspacer{} สหชาตวารมฺปิ ฯเปฯ สมฺปยุตฺตวารมฺปิ ปฏิจฺจวารสทิสํ.)}

\proseitem{33}{กุสโล นกามาวจโร ธมฺโม กุสลสฺส นกามาวจรสฺส ธมฺมสฺส เหตุปจฺจเยน ปจฺจโย. (1)}

\prose{อพฺยากโต นกามาวจโร ธมฺโม อพฺยากตสฺส นกามาวจรสฺส ธมฺมสฺส เหตุปจฺจเยน ปจฺจโย. (1)}

\prose{กุสโล นกามาวจโร ธมฺโม กุสลสฺส นกามาวจรสฺส ธมฺมสฺส อารมฺมณปจฺจเยน ปจฺจโย.\csromanspacer{}\csromanspacer{}กุสโล นกามาวจโร ธมฺโม อพฺยากตสฺส นกามาวจรสฺส ธมฺมสฺส อารมฺมณปจฺจเยน ปจฺจโย. (2)}

\csromanheader{2}{1. กุสลตฺติก 94. รูปาวจรทุก}
\prose{\csromanfolio{611}อพฺยากโต นกามาวจโร ธมฺโม อพฺยากตสฺส นกามาวจรสฺส ธมฺมสฺส อารมฺมณปจฺจเยน ปจฺจโย.\csromanspacer{}\csromanspacer{}อพฺยากโต นกามาวจโร ธมฺโม กุสลสฺส นกามาวจรสฺส ธมฺมสฺส อารมฺมณปจฺจเยน ปจฺจโย. (2)}

\proseitem{34}{เหตุยา เทฺว, อารมฺมเณ จตฺตาริ, อธิปติยา ตีณิ ฯเปฯ อวิคเต เทฺว. (สํขิตฺตํ.)}

\prose{(ยถา กุสลตฺติเก ปญฺหาวารํ, เอวํ วิตฺถาเรตพฺพํ.)}

\csromanheader{2}{1. กุสลตฺติก 94. รูปาวจรทุก}
\csromanheader{2}{1-7. ปฏิจฺจาทิวารปจฺจยจตุกฺก-เหตุ}
\proseitem{35}{กุสลํ รูปาวจรํ ธมฺมํ ปฏิจฺจ กุสโล รูปาวจโร ธมฺโม อุปฺปชฺชติ เหตุปจฺจยา. (1)}

\prose{อพฺยากตํ รูปาวจรํ ธมฺมํ ปฏิจฺจ อพฺยากโต รูปาวจโร ธมฺโม อุปฺปชฺชติ เหตุปจฺจยา. (1) (สํขิตฺตํ.)}

\prose{เหตุยา เทฺว ฯเปฯ อวิคเต เทฺว. (สํขิตฺตํ.) (สหชาตวารมฺปิ ฯเปฯ ปญฺหาวารมฺปิ วิตฺถาเรตพฺพํ.)}

\proseitem{36}{กุสลํ นรูปาวจรํ ธมฺมํ ปฏิจฺจ กุสโล นรูปาวจโร ธมฺโม อุปฺปชฺชติ เหตุปจฺจยา.\csromanspacer{} ตีณิ.}

\prose{อกุสลํ นรูปาวจรํ ธมฺมํ ปฏิจฺจ อกุสโล นรูปาวจโร ธมฺโม อุปฺปชฺชติ เหตุปจฺจยา.\csromanspacer{} ตีณิ.}

\prose{อพฺยากตํ นรูปาวจรํ ธมฺมํ ปฏิจฺจ อพฺยากโต นรูปาวจโร ธมฺโม อุปฺปชฺชติ เหตุปจฺจยา. (1)}

\prose{กุสลํ นรูปาวจรญฺจ อพฺยากตํ นรูปาวจรญฺจ ธมฺมํ ปฏิจฺจ อพฺยากโต นรูปาวจโร ธมฺโม อุปฺปชฺชติ เหตุปจฺจยา. (1)}

\prose{อกุสลํ นรูปาวจรญฺจ อพฺยากตํ นรูปาวจรญฺจ ธมฺมํ ปฏิจฺจ อพฺยากโต นรูปาวจโร ธมฺโม อุปฺปชฺชติ เหตุปจฺจยา. (1) (สํขิตฺตํ.)}

\gambhira{ธมฺมานุโลม ติกทุกปฏฺฐานปาฬิ}
\proseitem{37}{\csromanfolio{612}เหตุยา นว, อารมฺมเณ ตีณิ ฯเปฯ อวิคเต นว. (สํขิตฺตํ.\csromanspacer{} สหชาตวารมฺปิ ฯเปฯ สมฺปยุตฺตวารมฺปิ ปฏิจฺจวารสทิสํ.)}

\proseitem{38}{กุสโล นรูปาวจโร ธมฺโม กุสลสฺส นรูปาวจรสฺส ธมฺมสฺส เหตุปจฺจเยน ปจฺจโย.\csromanspacer{} ตีณิ.}

\prose{อกุสโล นรูปาวจโร ธมฺโม อกุสลสฺส นรูปาวจรสฺส ธมฺมสฺส เหตุปจฺจเยน ปจฺจโย.\csromanspacer{} ตีณิ.}

\prose{อพฺยากโต นรูปาวจโร ธมฺโม อพฺยากตสฺส นรูปาวจรสฺส ธมฺมสฺส เหตุปจฺจเยน ปจฺจโย. (1) (สํขิตฺตํ.)}

\proseitem{39}{เหตุยา สตฺต, อารมฺมเณ นว ฯเปฯ อวิคเต เตรส. (สํขิตฺตํ.\csromanspacer{} ยถา กุสลตฺติเก ปญฺหาวารํ, เอวํ วิตฺถาเรตพฺพํ.)}

\csromanheader{2}{1. กุสลตฺติก 95. อรูปาวจรทุก}
\csromanheader{2}{1. ปฏิจฺจวาร ปจฺจยจตุกฺก}
\proseitem{40}{กุสลํ อรูปาวจรํ ธมฺมํ ปฏิจฺจ กุสโล อรูปาวจโร ธมฺโม อุปฺปชฺชติ เหตุปจฺจยา. (1)}

\prose{อพฺยากตํ อรูปาวจรํ ธมฺมํ ปฏิจฺจ อพฺยากโต อรูปาวจโร ธมฺโม อุปฺปชฺชติ เหตุปจฺจยา. (1)}

\prose{กุสลํ น-อรูปาวจรํ ธมฺมํ ปฏิจฺจ กุสโล น-อรูปาวจโร ธมฺโม อุปฺปชฺชติ เหตุปจฺจยา. (รูปาวจรทุกสทิสํ.)}

\csromanheader{2}{1. กุสลตฺติก 96. ปริยาปนฺนทุก}
\csromanheader{2}{1-7. ปฏิจฺจาทิวาร ปจฺจยจตุกฺก-เหตุ}
\proseitem{41}{กุสลํ ปริยาปนฺนํ ธมฺมํ ปฏิจฺจ กุสโล ปริยาปนฺโน ธมฺโม อุปฺปชฺชติ เหตุปจฺจยา.\csromanspacer{} ตีณิ.}

\prose{อกุสลํ ปริยาปนฺนํ ธมฺมํ ปฏิจฺจ อกุสโล ปริยาปนฺโน ธมฺโม อุปฺปชฺชติ เหตุปจฺจยา.\csromanspacer{} ตีณิ.}

\csromanheader{2}{1. กุสลตฺติก 96. ปริยาปนฺนทุก}
\prose{\csromanfolio{613}อพฺยากตํ ปริยาปนฺนํ ธมฺมํ ปฏิจฺจ อพฺยากโต ปริยาปนฺโน ธมฺโม อุปฺปชฺชติ เหตุปจฺจยา. (1)}

\prose{กุสลํ ปริยาปนฺนญฺจ อพฺยากตํ ปริยาปนฺนญฺจ ธมฺมํ ปฏิจฺจ อพฺยากโต ปริยาปนฺโน ธมฺโม อุปฺปชฺชติ เหตุปจฺจยา. (1)}

\prose{อกุสลํ ปริยาปนฺนญฺจ อพฺยากตํ ปริยาปนฺนญฺจ ธมฺมํ ปฏิจฺจ อพฺยากโต ปริยาปนฺโน ธมฺโม อุปฺปชฺชติ เหตุปจฺจยา. (1) (สํขิตฺตํ.)}

\proseitem{42}{เหตุยา นว, อารมฺมเณ ตีณิ ฯเปฯ วิปาเก เอกํ ฯเปฯ อวิคเต นว. (สํขิตฺตํ.\csromanspacer{} สหชาตวารมฺปิ ฯเปฯ สมฺปยุตฺตวารมฺปิ ปฏิจฺจวารสทิสํ.)}

\proseitem{43}{กุสโล ปริยาปนฺโน ธมฺโม กุสลสฺส ปริยาปนฺนสฺส ธมฺมสฺส เหตุปจฺจเยน ปจฺจโย.\csromanspacer{} ตีณิ.}

\prose{อกุสโล ปริยาปนฺโน ธมฺโม อกุสลสฺส ปริยาปนฺนสฺส ธมฺมสฺส เหตุปจฺจเยน ปจฺจโย.\csromanspacer{} ตีณิ.}

\prose{อพฺยากโต ปริยาปนฺโน ธมฺโม อพฺยากตสฺส ปริยาปนฺนสฺส ธมฺมสฺส เหตุปจฺจเยน ปจฺจโย. (1) (สํขิตฺตํ.)}

\proseitem{44}{เหตุยา สตฺต, อารมฺมเณ นว, อธิปติยา นว ฯเปฯ อวิคเต เตรส. (สํขิตฺตํ.\csromanspacer{} ยถา กุสลตฺติเก ปญฺหาวารํ, เอวํ วิตฺถาเรตพฺพํ.)}

\csromanheader{2}{\textbf{อปริยาปนฺนปท-เหตุ อารมฺมณ}}
\proseitem{45}{กุสลํ อปริยาปนฺนํ ธมฺมํ ปฏิจฺจ กุสโล อปริยาปนฺโน ธมฺโม อุปฺปชฺชติ เหตุปจฺจยา. (1)}

\prose{อพฺยากตํ อปริยาปนฺนํ ธมฺมํ ปฏิจฺจ อพฺยากโต อปริยาปนฺโน ธมฺโม อุปฺปชฺชติ เหตุปจฺจยา. (1) (สํขิตฺตํ.)}

\prose{เหตุยา เทฺว, อารมฺมเณ เทฺว ฯเปฯ อวิคเต เทฺว. (สํขิตฺตํ.) (สหชาตวารมฺปิ ฯเปฯ สมฺปยุตฺตวารมฺปิ ปฏิจฺจวารสทิสํ.)}

\proseitem{46}{กุสโล อปริยาปนฺโน ธมฺโม กุสลสฺส อปริยาปนฺนสฺส ธมฺมสฺส เหตุปจฺจเยน ปจฺจโย. (1)}

\prose{อพฺยากโต อปริยาปนฺโน ธมฺโม อพฺยากตสฺส อปริยาปนฺนสฺส ธมฺมสฺส เหตุปจฺจเยน ปจฺจโย. (1)}

\gambhira{ธมฺมานุโลม ติกทุกปฏฺฐานปาฬิ}
\prose{\csromanfolio{614}อพฺยากโต อปริยาปนฺโน ธมฺโม อพฺยากตสฺส อปริยาปนฺนสฺส ธมฺมสฺส อารมฺมณปจฺจเยน ปจฺจโย.\csromanspacer{}\csromanspacer{}อพฺยากโต อปริยาปนฺโน ธมฺโม กุสลสฺส อปริยาปนฺนสฺส ธมฺมสฺส อารมฺมณปจฺจเยน ปจฺจโย. (2)}

\proseitem{47}{เหตุยา เทฺว, อารมฺมเณ เทฺว, อธิปติยา ตีณิ, อนนฺตเร เทฺว ฯเปฯ อุปนิสฺสเย จตฺตาริ ฯเปฯ อวิคเต เทฺว. (สํขิตฺตํ.)}

\prose{(ยถา กุสลตฺติเก ปญฺหาวารํ, เอวํ วิตฺถาเรตพฺพํ.)}

\csromanheader{2}{1. กุสลตฺติก 97. นิยฺยานิกทุก}
\csromanheader{2}{1. ปฏิจฺจวาร ปจฺจยจตุกฺก}
\proseitem{48}{กุสลํ นิยฺยานิกํ ธมฺมํ ปฏิจฺจ กุสโล นิยฺยานิโก ธมฺโม อุปฺปชฺชติ เหตุปจฺจยา. (1) (สํขิตฺตํ.)}

\prose{เหตุยา เอกํ, อารมฺมเณ เอกํ ฯเปฯ อวิคเต เอกํ. (สํขิตฺตํ.) (สหชาตวาเรปิ ฯเปฯ ปญฺหาวาเรปิ สพฺพตฺถ เอกํ.)}

\proseitem{49}{กุสลํ อนิยฺยานิกํ ธมฺมํ ปฏิจฺจ กุสโล อนิยฺยานิโก ธมฺโม อุปฺปชฺชติ เหตุปจฺจยา.\csromanspacer{} ตีณิ.}

\prose{อกุสลํ อนิยฺยานิกํ ธมฺมํ ปฏิจฺจ อกุสโล อนิยฺยานิโก ธมฺโม อุปฺปชฺชติ เหตุปจฺจยา.\csromanspacer{} ตีณิ.}

\prose{อพฺยากตํ อนิยฺยานิกํ ธมฺมํ ปฏิจฺจ อพฺยากโต อนิยฺยานิโก ธมฺโม อุปฺปชฺชติ เหตุปจฺจยา. (1)}

\prose{กุสลํ อนิยฺยานิกญฺจ อพฺยากตํ อนิยฺยานิกญฺจ ธมฺมํ ปฏิจฺจ อพฺยากโต อนิยฺยานิโก ธมฺโม อุปฺปชฺชติ เหตุปจฺจยา. (1)}

\prose{อกุสลํ อนิยฺยานิกญฺจ อพฺยากตํ อนิยฺยานิกญฺจ ธมฺมํ ปฏิจฺจ อพฺยากโต อนิยฺยานิโก ธมฺโม อุปฺปชฺชติ เหตุปจฺจยา. (1) (สํขิตฺตํ.)}

\proseitem{50}{เหตุยา นว, อารมฺมเณ ตีณิ ฯเปฯ อวิคเต นว. (สํขิตฺตํ.)}

\vspace{\tipitakaparskip}
\csromancenter{(สหชาตวารมฺปิ ฯเปฯ ปญฺหาวารมฺปิ วิตฺถาเรตพฺพํ.)}
\csromanheader{2}{1. กุสลตฺติก 99. ส-อุตฺตรทุก \textbf{1. กุสลตฺติก 98. นิยตทุก}}
\csromanheader{2}{1. ปฏิจฺจวาร ปจฺจยจตุกฺก}
\proseitem{51}{\csromanfolio{615}กุสลํ นิยตํ ธมฺมํ ปฏิจฺจ กุสโล นิยโต ธมฺโม อุปฺปชฺชติ เหตุปจฺจยา. (1)}

\prose{อกุสลํ นิยตํ ธมฺมํ ปฏิจฺจ อกุสโล นิยโต ธมฺโม อุปฺปชฺชติ เหตุปจฺจยา. (1) (สํขิตฺตํ.)}

\prose{เหตุยา เทฺว, อารมฺมเณ เทฺว ฯเปฯ อวิคเต เทฺว. (สํขิตฺตํ.) (สหชาตวารมฺปิ ฯเปฯ ปญฺหาวารมฺปิ เอวํ วิตฺถาเรตพฺพํ.)}

\proseitem{52}{กุสลํ อนิยตํ ธมฺมํ ปฏิจฺจ กุสโล อนิยโต ธมฺโม อุปฺปชฺชติ เหตุปจฺจยา. (สํขิตฺตํ.)}

\csromanhangingbegin
\hangingheaditem{52}{เหตุยา นว, อารมฺมเณ ตีณิ ฯเปฯ อวิคเต นว. (สํขิตฺตํ.)}
\hangingbody{(สหชาตวารมฺปิ ฯเปฯ ปญฺหาวารมฺปิ เอวํ วิตฺถาเรตพฺพํ.)}
\csromanhangingend

\csromanheader{2}{1. กุสลตฺติก 99. ส-อุตฺตรทุก}
\csromanheader{2}{1-7. ปฏิจฺจาทิวาร ปจฺจยจตุกฺก}
\proseitem{53}{กุสลํ ส-อุตฺตรํ ธมฺมํ ปฏิจฺจ กุสโล ส-อุตฺตโร ธมฺโม อุปฺปชฺชติ เหตุปจฺจยา.\csromanspacer{} ตีณิ.}

\prose{อกุสลํ ส-อุตฺตรํ ธมฺมํ ปฏิจฺจ อกุสโล ส-อุตฺตโร ธมฺโม อุปฺปชฺชติ เหตุปจฺจยา.\csromanspacer{} ตีณิ.}

\prose{อพฺยากตํ ส-อุตฺตรํ ธมฺมํ ปฏิจฺจ อพฺยากโต ส-อุตฺตโร ธมฺโม อุปฺปชฺชติ เหตุปจฺจยา. (1)}

\prose{กุสลํ ส-อุตฺตรญฺจ อพฺยากตํ ส-อุตฺตรญฺจ ธมฺมํ ปฏิจฺจ อพฺยากโต ส-อุตฺตโร ธมฺโม อุปฺปชฺชติ เหตุปจฺจยา. (1)}

\prose{อกุสลํ ส-อุตฺตรญฺจ อพฺยากตํ ส-อุตฺตรญฺจ ธมฺมํ ปฏิจฺจ อพฺยากโต ส-อุตฺตโร ธมฺโม อุปฺปชฺชติ เหตุปจฺจยา. (1) (สํขิตฺตํ.)}

\gambhira{ธมฺมานุโลม ติกทุกปฏฺฐานปาฬิ}
\proseitem{54}{\csromanfolio{616}เหตุยา นว, อารมฺมเณ ตีณิ ฯเปฯ อวิคเต นว. (สํขิตฺตํ.\csromanspacer{} สหชาตวารมฺปิ ฯเปฯ สมฺปยุตฺตวารมฺปิ ปฏิจฺจวารสทิสํ.)}

\proseitem{55}{กุสโล ส-อุตฺตโร ธมฺโม กุสลสฺส ส-อุตฺตรสฺส ธมฺมสฺส เหตุปจฺจเยน ปจฺจโย.\csromanspacer{} ตีณิ.}

\prose{อกุสโล ส-อุตฺตโร ธมฺโม อกุสลสฺส ส-อุตฺตรสฺส ธมฺมสฺส เหตุปจฺจเยน ปจฺจโย.\csromanspacer{} ตีณิ.}

\prose{อพฺยากโต ส-อุตฺตโร ธมฺโม อพฺยากตสฺส ส-อุตฺตรสฺส ธมฺมสฺส เหตุปจฺจเยน ปจฺจโย. (1) (สํขิตฺตํ.)}

\proseitem{56}{เหตุยา สตฺต, อารมฺมเณ นว, อธิปติยา นว ฯเปฯ อวิคเต เตรส. (สํขิตฺตํ.)}

\csromanheader{2}{\textbf{อนุตฺตรปท-เหตุ อารมฺมณ}}
\proseitem{57}{กุสลํ อนุตฺตรํ ธมฺมํ ปฏิจฺจ กุสโล อนุตฺตโร ธมฺโม อุปฺปชฺชติ เหตุปจฺจยา. (1)}

\prose{อพฺยากตํ อนุตฺตรํ ธมฺมํ ปฏิจฺจ อพฺยากโต อนุตฺตโร ธมฺโม อุปฺปชฺชติ เหตุปจฺจยา. (1) (สํขิตฺตํ.)}

\prose{เหตุยา เทฺว, อารมฺมเณ เทฺว ฯเปฯ อวิคเต เทฺว. (สํขิตฺตํ.\csromanspacer{} สหชาตวารมฺปิ ฯเปฯ สมฺปยุตฺตวารมฺปิ ปฏิจฺจวารสทิสํ.)}

\proseitem{58}{กุสโล อนุตฺตโร ธมฺโม กุสลสฺส อนุตฺตรสฺส ธมฺมสฺส เหตุปจฺจเยน ปจฺจโย. (1)}

\prose{อพฺยากโต อนุตฺตโร ธมฺโม อพฺยากตสฺส อนุตฺตรสฺส ธมฺมสฺส เหตุปจฺจเยน ปจฺจโย. (1)}

\prose{อพฺยากโต อนุตฺตโร ธมฺโม อพฺยากตสฺส อนุตฺตรสฺส ธมฺมสฺส อารมฺมณปจฺจเยน ปจฺจโย. (1)}

\prose{อพฺยากโต อนุตฺตโร ธมฺโม กุสลสฺส อนุตฺตรสฺส ธมฺมสฺส อารมฺมณปจฺจเยน ปจฺจโย. (1) (สํขิตฺตํ.)}

\proseitem{59}{เหตุยา เทฺว, อารมฺมเณ เทฺว, อธิปติยา ตีณิ, อนนฺตเร เทฺว ฯเปฯ อวิคเต เทฺว. (สํขิตฺตํ.)}

\prose{(ยถา กุสลตฺติเก ปญฺหาวารํ, เอวํ วิตฺถาเรตพฺพํ.)}

\csromanmatikaanchor{mtk.150}
\csromantocmarkref{subparagraph}{100. สรณทุก}{mtk.150}
\bookmark[dest={mtk.150},level=5]{100. สรณทุก}
\csromanheader{6}{1. กุสลตฺติก 100. สรณทุก \textbf{1. กุสลตฺติก 100. สรณทุก}}
\csromanheader{2}{1-7. ปฏิจฺจาทิวาร ปจฺจยจตุกฺก}
\proseitem{60}{\csromanfolio{617}อกุสลํ สรณํ ธมฺมํ ปฏิจฺจ อกุสโล สรโณ ธมฺโม อุปฺปชฺชติ เหตุปจฺจยา. (สํขิตฺตํ.)}

\prose{เหตุยา เอกํ ฯเปฯ อวิคเต เอกํ. (สํขิตฺตํ.) (สหชาตวาเรปิ ฯเปฯ ปญฺหาวาเรปิ สพฺพตฺถ เอกํ.)}

\proseitem{61}{กุสลํ อรณํ ธมฺมํ ปฏิจฺจ กุสโล อรโณ ธมฺโม อุปฺปชฺชติ เหตุปจฺจยา.\csromanspacer{} ตีณิ.}

\prose{อพฺยากตํ อรณํ ธมฺมํ ปฏิจฺจ อพฺยากโต อรโณ ธมฺโม อุปฺปชฺชติ เหตุปจฺจยา. (1)}

\prose{กุสลํ อรณญฺจ อพฺยากตํ อรณญฺจ ธมฺมํ ปฏิจฺจ อพฺยากโต อรโณ ธมฺโม อุปฺปชฺชติ เหตุปจฺจยา. (1)}

\prose{กุสลํ อรณํ ธมฺมํ ปฏิจฺจ กุสโล อรโณ ธมฺโม อุปฺปชฺชติ อารมฺมณปจฺจยา. (1)}

\prose{อพฺยากตํ อรณํ ธมฺมํ ปฏิจฺจ อพฺยากโต อรโณ ธมฺโม อุปฺปชฺชติ อารมฺมณปจฺจยา. (1) (สํขิตฺตํ.)}

\proseitem{62}{เหตุยา ปญฺจ, อารมฺมเณ เทฺว ฯเปฯ อวิคเต ปญฺจ. (สํขิตฺตํ.\csromanspacer{} สหชาตวารมฺปิ ฯเปฯ สมฺปยุตฺตวารมฺปิ ปฏิจฺจวารสทิสํ.)}

\proseitem{63}{กุสโล อรโณ ธมฺโม กุสลสฺส อรณสฺส ธมฺมสฺส เหตุปจฺจเยน ปจฺจโย.\csromanspacer{} ตีณิ.}

\prose{อพฺยากโต อรโณ ธมฺโม อพฺยากตสฺส อรณสฺส ธมฺมสฺส เหตุปจฺจเยน ปจฺจโย. (1) (สํขิตฺตํ.)}

\gambhira{ธมฺมานุโลม ติกทุกปฏฺฐานปาฬิ}
\proseitem{64}{\csromanfolio{618}เหตุยา จตฺตาริ, อารมฺมเณ จตฺตาริ ฯเปฯ อวิคเต สตฺต. (สํขิตฺตํ.) (ยถา กุสลตฺติเก ปญฺหาวารํ, เอวํ วิตฺถาเรตพฺพํ.)}

\nitthitam{กุสลตฺติกปิฏฺฐิทุกํ นิฏฺฐิตํ.}
\csromanmatikaanchor{mtk.151}
\csromanmatikaanchor{mtk.152}
\csromantocmarknopage{subparagraph}{2. เวทนาตฺติก}{mtk.151}
\bookmark[dest={mtk.151},level=5]{2. เวทนาตฺติก}
\csromantocmarkref{subparagraph}{100. สรณทุก}{mtk.152}
\bookmark[dest={mtk.152},level=5]{100. สรณทุก}
\csromanheader{6}{2. เวทนาตฺติก 100. สรณทุก}
\csromanheader{2}{1-7. ปฏิจฺจาทิวาร ปจฺจยจตุกฺก}
\proseitem{65}{สุขาย เวทนาย สมฺปยุตฺตํ สรณํ ธมฺมํ ปฏิจฺจ สุขาย เวทนาย สมฺปยุตฺโต สรโณ ธมฺโม อุปฺปชฺชติ เหตุปจฺจยา. (1)}

\prose{ทุกฺขาย เวทนาย สมฺปยุตฺตํ สรณํ ธมฺมํ ปฏิจฺจ ทุกฺขาย เวทนาย สมฺปยุตฺโต สรโณ ธมฺโม อุปฺปชฺชติ เหตุปจฺจยา. (1)}

\prose{อทุกฺขมสุขาย เวทนาย สมฺปยุตฺตํ สรณํ ธมฺมํ ปฏิจฺจ อทุกฺขมสุขาย เวทนาย สมฺปยุตฺโต สรโณ ธมฺโม อุปฺปชฺชติ เหตุปจฺจยา. (1) (สํขิตฺตํ.)}

\proseitem{66}{เหตุยา ตีณิ ฯเปฯ อวิคเต ตีณิ. (สํขิตฺตํ.) (สหชาตวารมฺปิ ฯเปฯ สมฺปยุตฺตวารมฺปิ ปฏิจฺจวารสทิสํ.)}

\proseitem{67}{สุขาย เวทนาย สมฺปยุตฺโต สรโณ ธมฺโม สุขาย เวทนาย สมฺปยุตฺตสฺส สรณสฺส ธมฺมสฺส เหตุปจฺจเยน ปจฺจโย. (1)}

\prose{ทุกฺขาย เวทนาย สมฺปยุตฺโต สรโณ ธมฺโม ทุกฺขาย เวทนาย สมฺปยุตฺตสฺส สรณสฺส ธมฺมสฺส เหตุปจฺจเยน ปจฺจโย. (1)}

\prose{อทุกฺขมสุขาย เวทนาย สมฺปยุตฺโต สรโณ ธมฺโม อทุกฺขมสุขาย เวทนาย สมฺปยุตฺตสฺส สรณสฺส ธมฺมสฺส เหตุปจฺจเยน ปจฺจโย. (1) (สํขิตฺตํ.)}

\csromanheader{2}{3. วิปากตฺติก 100. สรณทุก}
\proseitem{68}{\csromanfolio{619}เหตุยา ตีณิ, อารมฺมเณ นว ฯเปฯ อวิคเต ตีณิ. (สํขิตฺตํ.\csromanspacer{} ยถา กุสลตฺติเก ปญฺหาวารํ, เอวํ วิตฺถาเรตพฺพํ.)}

\csromanheader{2}{\textbf{อรณปท-เหตุ}}
\proseitem{69}{สุขาย เวทนาย สมฺปยุตฺตํ อรณํ ธมฺมํ ปฏิจฺจ สุขาย เวทนาย สมฺปยุตฺโต อรโณ ธมฺโม อุปฺปชฺชติ เหตุปจฺจยา.}

\prose{อทุกฺขมสุขาย เวทนาย สมฺปยุตฺตํ อรณํ ธมฺมํ ปฏิจฺจ อทุกฺขมสุขาย เวทนาย สมฺปยุตฺโต อรโณ ธมฺโม อุปฺปชฺชติ เหตุปจฺจยา. (สํขิตฺตํ.)}

\prose{เหตุยา เทฺว, อารมฺมเณ ตีณิ ฯเปฯ อวิคเต ตีณิ. (สํขิตฺตํ.) (สหชาตวารมฺปิ ฯเปฯ สมฺปยุตฺตวารมฺปิ ปฏิจฺจวารสทิสํ.)}

\proseitem{70}{สุขาย เวทนาย สมฺปยุตฺโต อรโณ ธมฺโม สุขาย เวทนาย สมฺปยุตฺตสฺส อรณสฺส ธมฺมสฺส เหตุปจฺจเยน ปจฺจโย. (1)}

\prose{อทุกฺขมสุขาย เวทนาย สมฺปยุตฺโต อรโณ ธมฺโม อทุกฺขมสุขาย เวทนาย สมฺปยุตฺตสฺส อรณสฺส ธมฺมสฺส เหตุปจฺจเยน ปจฺจโย. (1) (สํขิตฺตํ.)}

\csromanhangingbegin
\hangingheaditem{71}{เหตุยา เทฺว, อารมฺมเณ ฉ ฯเปฯ อวิคเต ตีณิ. (สํขิตฺตํ.)}
\hangingbody{(ยถา กุสลตฺติเก ปญฺหาวารํ, เอวํ วิตฺถาเรตพฺพํ.)}
\csromanhangingend

\csromanheader{2}{3. วิปากตฺติก 100. สรณทุก}
\csromanheader{2}{1-7. ปฏิจฺจาทิวาร ปจฺจยจตุกฺก}
\proseitem{72}{วิปากธมฺมธมฺมํ สรณํ ธมฺมํ ปฏิจฺจ วิปากธมฺมธมฺโม สรโณ ธมฺโม อุปฺปชฺชติ เหตุปจฺจยา. (สํขิตฺตํ.)}

\csromanheader{2}{เหตุยา เอกํ ฯเปฯ อวิคเต เอกํ. (สํขิตฺตํ.)}
\vspace{\tipitakaparskip}
\csromancenter{(สหชาตวาเรปิ ฯเปฯ ปญฺหาวาเรปิ สพฺพตฺถ เอกํ.)}
\gambhira{ธมฺมานุโลม ติกทุกปฏฺฐานปาฬิ}
\proseitem{73}{\csromanfolio{620}วิปากํ อรณํ ธมฺมํ ปฏิจฺจ วิปาโก อรโณ ธมฺโม อุปฺปชฺชติ เหตุปจฺจยา.\csromanspacer{} ตีณิ.}

\prose{วิปากธมฺมธมฺมํ อรณํ ธมฺมํ ปฏิจฺจ วิปากธมฺมธมฺโม อรโณ ธมฺโม อุปฺปชฺชติ เหตุปจฺจยา.\csromanspacer{} ตีณิ.}

\prose{เนววิปากนวิปากธมฺมธมฺมํ อรณํ ธมฺมํ ปฏิจฺจ เนววิปากนวิปากธมฺมธมฺโม อรโณ ธมฺโม อุปฺปชฺชติ เหตุปจฺจยา.\csromanspacer{} ตีณิ.}

\prose{วิปากํ อรณญฺจ เนววิปากนวิปากธมฺมธมฺมํ อรณญฺจ ธมฺมํ ปฏิจฺจ วิปาโก อรโณ ธมฺโม อุปฺปชฺชติ เหตุปจฺจยา.\csromanspacer{} ตีณิ.}

\prose{วิปากธมฺมธมฺมํ อรณญฺจ เนววิปากนวิปากธมฺมธมฺมํ อรณญฺจ ธมฺมํ ปฏิจฺจ เนววิปากนวิปากธมฺมธมฺโม อรโณ ธมฺโม อุปฺปชฺชติ เหตุปจฺจยา. (1) (สํขิตฺตํ.)}

\proseitem{74}{เหตุยา เตรส, อารมฺมเณ ปญฺจ ฯเปฯ อวิคเต เตรส. (สํขิตฺตํ.\csromanspacer{} สหชาตวารมฺปิ ฯเปฯ สมฺปยุตฺตวารมฺปิ ปฏิจฺจวารสทิสํ.)}

\proseitem{75}{วิปาโก อรโณ ธมฺโม วิปากสฺส อรณสฺส ธมฺมสฺส เหตุปจฺจเยน ปจฺจโย.\csromanspacer{} ตีณิ.}

\prose{วิปากธมฺมธมฺโม อรโณ ธมฺโม วิปากธมฺมธมฺมสฺส อรณสฺส ธมฺมสฺส เหตุปจฺจเยน ปจฺจโย.\csromanspacer{} ตีณิ.}

\prose{เนววิปากนวิปากธมฺมธมฺโม อรโณ ธมฺโม เนววิปากนวิปากธมฺมธมฺมสฺส อรณสฺส ธมฺมสฺส เหตุปจฺจเยน ปจฺจโย. (1) (สํขิตฺตํ.)}

\csromanhangingbegin
\hangingheaditem{76}{เหตุยา สตฺต, อรมฺมเณ นว ฯเปฯ อวิคเต เตรส. (สํขิตฺตํ.)}
\hangingbody{(ยถา กุสลตฺติเก ปญฺหาวารํ, เอวํ วิตฺถาเรตพฺพํ.)}
\csromanhangingend

\vspace{\tipitakaparskip}
\csromancenter{อุปาทินฺนตฺติก 100.\csromanspacer{} สรณทุก \textbf{4.}\csromanspacer{}\textbf{ อุปาทินฺนตฺติก 100.}\csromanspacer{}\textbf{ สรณทุก}}
\csromanheader{2}{1-7. ปฏิจฺจาทิวาร ปจฺจยจตุกฺก-เหตุ}
\proseitem{77}{\csromanfolio{621}อนุปาทินฺนุปาทานิยํ สรณํ ธมฺมํ ปฏิจฺจ อนุปาทินฺนุปาทานิโย สรโณ ธมฺโม อุปฺปชฺชติ เหตุปจฺจยา. (สํขิตฺตํ.)}

\prose{เหตุยา เอกํ ฯเปฯ อวิคเต เอกํ. (สํขิตฺตํ.) (สหชาตวาเรปิ ฯเปฯ ปญฺหาวาเรปิ สพฺพตฺถ เอกํ.)}

\proseitem{78}{อุปาทินฺนุปาทานิยํ อรณํ ธมฺมํ ปฏิจฺจ อุปาทินฺนุปาทานิโย อรโณ ธมฺโม อุปฺปชฺชติ เหตุปจฺจยา.\csromanspacer{} ตีณิ.}

\prose{อนุปาทินฺนุปาทานิยํ อรณํ ธมฺมํ ปฏิจฺจ อนุปาทินฺนุปาทานิโย อรโณ ธมฺโม อุปฺปชฺชติ เหตุปจฺจยา. (1)}

\prose{อนุปาทินฺน-อนุปาทานิยํ อรณํ ธมฺมํ ปฏิจฺจ อนุปาทินฺน-อนุปาทานิโย อรโณ ธมฺโม อุปฺปชฺชติ เหตุปจฺจยา.\csromanspacer{} ตีณิ.}

\prose{อุปาทินฺนุปาทานิยํ อรณญฺจ อนุปาทินฺนุปาทานิยํ อรณญฺจ ธมฺมํ ปฏิจฺจ อนุปาทินฺนุปาทานิโย อรโณ ธมฺโม อุปฺปชฺชติ เหตุปจฺจยา. (1)}

\prose{อนุปาทินฺนุปาทานิยํ อรณญฺจ อนุปาทินฺน-อนุปาทานิยํ อรณญฺจ ธมฺมํ ปฏิจฺจ อนุปาทินฺนุปาทานิโย อรโณ ธมฺโม อุปฺปชฺชติ เหตุปจฺจยา. (1) (สํขิตฺตํ.)}

\proseitem{79}{เหตุยา นว, อารมฺมเณ ตีณิ ฯเปฯ อวิคเต นว. (สํขิตฺตํ.) (สหชาตวารมฺปิ ฯเปฯ สมฺปยุตฺตวารมฺปิ ปฏิจฺจวารสทิสํ.)}

\proseitem{80}{อุปาทินฺนุปาทานิโย อรโณ ธมฺโม อุปาทินฺนุปาทานิยสฺส อรณสฺส ธมฺมสฺส เหตุปจฺจเยน ปจฺจโย.\csromanspacer{} ตีณิ.}

\prose{อนุปาทินฺนุปาทานิโย อรโณ ธมฺโม อนุปาทินฺนุปาทานิยสฺส อรณสฺส ธมฺมสฺส เหตุปจฺจเยน ปจฺจโย. (1)}

\prose{อนุปาทินฺน-อนุปาทานิโย อรโณ ธมฺโม อนุปาทินฺน-อนุปาทานิยสฺส อรณสฺส ธมฺมสฺส เหตุปจฺจเยน ปจฺจโย.\csromanspacer{} ตีณิ. (สํขิตฺตํ.)}

\gambhira{ธมฺมานุโลม ติกทุกปฏฺฐานปาฬิ}
\proseitem{81}{\csromanfolio{622}เหตุยา สตฺต, อารมฺมเณ ฉ ฯเปฯ อวิคเต เตวีส. (สํขิตฺตํ.\csromanspacer{} ยถา กุสลตฺติเก ปญฺหาวารํ, เอวํ วิตฺถาเรตพฺพํ.)}

\csromanheader{2}{5. สํกิลิฏฺฐตฺติก 100. สรณทุก}
\csromanheader{2}{1-7. ปฏิจฺจาทิวาร ปจฺจยจตุกฺก-เหตุ}
\proseitem{82}{สํกิลิฏฺฐสํกิเลสิกํ สรณํ ธมฺมํ ปฏิจฺจ สํกิลิฏฺฐสํกิเลสิโก สรโณ ธมฺโม อุปฺปชฺชติ เหตุปจฺจยา. (สํขิตฺตํ.)}

\prose{เหตุยา เอกํ ฯเปฯ อวิคเต เอกํ. (สํขิตฺตํ.) (สหชาตวาเรปิ ฯเปฯ ปญฺหาวาเรปิ สพฺพตฺถ เอกํ.)}

\proseitem{83}{อสํกิลิฏฺฐสํกิเลสิกํ อรณํ ธมฺมํ ปฏิจฺจ อสํกิลิฏฺฐสํกิเลสิโก อรโณ ธมฺโม อุปฺปชฺชติ เหตุปจฺจยา. (1)}

\prose{อสํกิลิฏฺฐ-อสํกิเลสิกํ อรณํ ธมฺมํ ปฏิจฺจ อสํกิลิฏฺฐ-อสํกิเลสิโก อรโณ ธมฺโม อุปฺปชฺชติ เหตุปจฺจยา.\csromanspacer{} ตีณิ.}

\prose{อสํกิลิฏฺฐสํกิเลสิกํ อรณญฺจ อสํกิลิฏฺฐ-อสํกิเลสิกํ อรณญฺจ ธมฺมํ ปฏิจฺจ อสํกิลิฏฺฐสํกิเลสิโก อรโณ ธมฺโม อุปฺปชฺชติ เหตุปจฺจยา. (1) (สํขิตฺตํ.)}

\proseitem{84}{เหตุยา ปญฺจ, อารมฺมเณ เทฺว ฯเปฯ อวิคเต ปญฺจ. (สํขิตฺตํ.) (สหชาตวารมฺปิ ฯเปฯ สมฺปยุตฺตวารมฺปิ ปฏิจฺจวารสทิสํ.)}

\proseitem{85}{อสํกิลิฏฺฐสํกิเลสิโก อรโณ ธมฺโม อสํกิลิฏฺฐสํกิเลสิกสฺส อรณสฺส ธมฺมสฺส เหตุปจฺจเยน ปจฺจโย. (1)}

\prose{อสํกิลิฏฺฐ-อสํกิเลสิโก อรโณ ธมฺโม อสํกิลิฏฺฐ-อสํกิเลสิกสฺส อรณสฺส ธมฺมสฺส เหตุปจฺจเยน ปจฺจโย.\csromanspacer{} ตีณิ. (สํขิตฺตํ.)}

\csromanheader{2}{6. วิตกฺกตฺติก 100. สรณทุก}
\csromanhangingbegin
\hangingheaditem{86}{\csromanfolio{623}เหตุยา จตฺตาริ, อารมฺมเณ ตีณิ ฯเปฯ อวิคเต สตฺต. (สํขิตฺตํ.)}
\hangingbody{(ยถา กุสลตฺติเก ปญฺหาวารํ, เอวํ วิตฺถาเรตพฺพํ.)}
\csromanhangingend

\csromanheader{2}{6. วิตกฺกตฺติก 100. สรณทุก}
\csromanheader{2}{1. ปฏิจฺจวาร ปจฺจยจตุกฺก-เหตุ}
\proseitem{87}{สวิตกฺกสวิจารํ สรณํ ธมฺมํ ปฏิจฺจ สวิตกฺกสวิจาโร สรโณ ธมฺโม อุปฺปชฺชติ เหตุปจฺจยา.\csromanspacer{}\csromanspacer{}สวิตกฺกสวิจารํ สรณํ ธมฺมํ ปฏิจฺจ อวิตกฺกวิจารมตฺโต สรโณ ธมฺโม อุปฺปชฺชติ เหตุปจฺจยา.\csromanspacer{}\csromanspacer{}สวิตกฺกสวิจารํ สรณํ ธมฺมํ ปฏิจฺจ สวิตกฺกสวิจาโร สรโณ จ อวิตกฺกวิจารมตฺโต สรโณ จ ธมฺมา อุปฺปชฺชนฺติ เหตุปจฺจยา. (3)}

\prose{อวิตกฺกวิจารมตฺตํ สรณํ ธมฺมํ ปฏิจฺจ สวิตกฺกสวิจาโร สรโณ ธมฺโม อุปฺปชฺชติ เหตุปจฺจยา. (1)}

\prose{สวิตกฺกสวิจารํ สรณญฺจ อวิตกฺกวิจารมตฺตํ สรณญฺจ ธมฺมํ ปฏิจฺจ สวิตกฺกสวิจาโร สรโณ ธมฺโม อุปฺปชฺชติ เหตุปจฺจยา. (1) (สํขิตฺตํ.)}

\proseitem{88}{เหตุยา ปญฺจ ฯเปฯ อวิคเต ปญฺจ. (สํขิตฺตํ.) (สหชาตวาเรปิ ฯเปฯ ปญฺหาวาเรปิ สพฺพตฺถ วิตฺถาโร.)}

\proseitem{89}{สวิตกฺกสวิจารํ อรณํ ธมฺมํ ปฏิจฺจ สวิตกฺกสวิจาโร อรโณ ธมฺโม อุปฺปชฺชติ เหตุปจฺจยา.\csromanspacer{} สตฺต.}

\prose{อวิตกฺกวิจารมตฺตํ อรณํ ธมฺมํ ปฏิจฺจ อวิตกฺกวิจารมตฺโต อรโณ ธมฺโม อุปฺปชฺชติ เหตุปจฺจยา.\csromanspacer{} ปญฺจ.}

\prose{อวิตกฺก-อวิจารํ อรณํ ธมฺมํ ปฏิจฺจ อวิตกฺก-อวิจาโร อรโณ ธมฺโม อุปฺปชฺชติ เหตุปจฺจยา. (สํขิตฺตํ.)}

\gambhira{ธมฺมานุโลม ติกทุกปฏฺฐานปาฬิ}
\proseitem{90}{\csromanfolio{624}เหตุยา สตฺตติํส ฯเปฯ อวิคเต สตฺตติํส. (สํขิตฺตํ.\csromanspacer{} สหชาตวารมฺปิ ฯเปฯ ปญฺหาวารมฺปิ วิตฺถาเรตพฺพํ.)}

\csromanheader{2}{7. ปีติตฺติก 100. สรณทุก}
\csromanheader{2}{1-7. ปฏิจฺจาทิวาร ปจฺจยจตุกฺก}
\proseitem{91}{ปีติสหคตํ สรณํ ธมฺมํ ปฏิจฺจ ปีติสหคโต สรโณ ธมฺโม อุปฺปชฺชติ เหตุปจฺจยา.\csromanspacer{} ตีณิ.}

\prose{สุขสหคตํ สรณํ ธมฺมํ ปฏิจฺจ สุขสหคโต สรโณ ธมฺโม อุปฺปชฺชติ เหตุปจฺจยา.\csromanspacer{} ตีณิ.}

\prose{อุเปกฺขาสหคตํ สรณํ ธมฺมํ ปฏิจฺจ อุเปกฺขาสหคโต สรโณ ธมฺโม อุปฺปชฺชติ เหตุปจฺจยา. (1)}

\prose{ปีติสหคตํ สรณญฺจ สุขสหคตํ สรณญฺจ ธมฺมํ ปฏิจฺจ ปีติสหคโต สรโณ ธมฺโม อุปฺปชฺชติ เหตุปจฺจยา.\csromanspacer{} ตีณิ. (สํขิตฺตํ.)}

\prose{เหตุยา ทส ฯเปฯ อวิคเต ทส. (สํขิตฺตํ.) (สหชาตวารมฺปิ ฯเปฯ สมฺปยุตฺตวารมฺปิ ปฏิจฺจวารสทิสํ.)}

\proseitem{92}{ปีติสหคโต สรโณ ธมฺโม ปีติสหคตสฺส สรณสฺส ธมฺมสฺส เหตุปจฺจเยน ปจฺจโย.\csromanspacer{} ตีณิ.}

\prose{สุขสหคโต สรโณ ธมฺโม สุขสหคตสฺส สรณสฺส ธมฺมสฺส เหตุปจฺจเยน ปจฺจโย.\csromanspacer{} ตีณิ.}

\prose{อุเปกฺขาสหคโต สรโณ ธมฺโม อุเปกฺขาสหคตสฺส สรณสฺส ธมฺมสฺส เหตุปจฺจเยน ปจฺจโย. (1)}

\prose{ปีติสหคโต สรโณ จ สุขสหคโต สรโณ จ ธมฺมา ปีติสหคตสฺส สรณสฺส ธมฺมสฺส เหตุปจฺจเยน ปจฺจโย.\csromanspacer{} ตีณิ. (สํขิตฺตํ.)}

\csromanheader{2}{9. ทสฺสนเหตุตฺติก 100. สรณทุก}
\proseitem{93}{\csromanfolio{625}เหตุยา ทส, อารมฺมเณ โสฬส ฯเปฯ อวิคเต ทส. (สํขิตฺตํ.) (ยถา กุสลตฺติเก ปญฺหาวารํ, เอวํ วิตฺถาเรตพฺพํ.)}

\proseitem{94}{ปีติสหคตํ อรณํ ธมฺมํ ปฏิจฺจ ปีติสหคโต อรโณ ธมฺโม อุปฺปชฺชติ เหตุปจฺจยา. (สํขิตฺตํ.)}

\csromanheader{2}{เหตุยา ทส ฯเปฯ อวิคเต ทส. (สํขิตฺตํ.)}
\vspace{\tipitakaparskip}
\csromancenter{(สหชาตวาเรปิ ฯเปฯ ปญฺหาวาเรปิ ตตฺตกาว ปญฺหา.)}
\csromanheader{2}{8. ทสฺสนตฺติก 100. สรณทุก}
\csromanheader{2}{1. ปฏิจฺจวาร ปจฺจยจตุกฺก}
\proseitem{95}{ทสฺสเนน ปหาตพฺพํ สรณํ ธมฺมํ ปฏิจฺจ ทสฺสเนน ปหาตพฺโพ สรโณ ธมฺโม อุปฺปชฺชติ เหตุปจฺจยา. (1)}

\prose{ภาวนาย ปหาตพฺพํ สรณํ ธมฺมํ ปฏิจฺจ ภาวนาย ปหาตพฺโพ สรโณ ธมฺโม อุปฺปชฺชติ เหตุปจฺจยา. (1) (สํขิตฺตํ.)}

\prose{เหตุยา เทฺว ฯเปฯ อวิคเต เทฺว. (สํขิตฺตํ.) (สหชาตวาเรปิ ฯเปฯ ปญฺหาวาเรปิ สพฺพตฺถ วิตฺถาโร.)}

\proseitem{96}{เนวทสฺสเนน นภาวนาย ปหาตพฺพํ อรณํ ธมฺมํ ปฏิจฺจ เนวทสฺสเนน นภาวนาย ปหาตพฺโพ อรโณ ธมฺโม อุปฺปชฺชติ เหตุปจฺจยา. (สํขิตฺตํ.)}

\csromanheader{2}{เหตุยา เอกํ ฯเปฯ อวิคเต เอกํ. (สํขิตฺตํ.)}
\vspace{\tipitakaparskip}
\csromancenter{(สหชาตวาเรปิ ฯเปฯ ปญฺหาวาเรปิ สพฺพตฺถ เอกํ.)}
\csromanheader{2}{9. ทสฺสนเหตุตฺติก 100. สรณทุก}
\csromanheader{2}{1. ปฏิจฺจวาร-เหตุ}
\proseitem{97}{ทสฺสเนน ปหาตพฺพเหตุกํ สรณํ ธมฺมํ ปฏิจฺจ ฯเปฯ.\csromanspacer{} ภาวนาย ปหาตพฺพเหตุกํ สรณํ ธมฺมํ ปฏิจฺจ ฯเปฯ.\csromanspacer{} เนวทสฺสเนน นภาวนาย}

\vspace{\tipitakaparskip}
\csromancenter{ธมฺมานุโลม ติกทุกปฏฺฐานปาฬิ ปหาตพฺพเหตุกํ สรณํ ธมฺมํ ปฏิจฺจ เนวทสฺสเนน นภาวนาย ปหาตพฺพเหตุโก สรโณ ธมฺโม อุปฺปชฺชติ เหตุปจฺจยา. (สํขิตฺตํ.)}
\csromanheader{2}{10. อาจยคามิตฺติก 100. สรณทุก}
\csromanheader{2}{1-7. ปฏิจฺจาทิวาร-เหตุ อารมฺมณ}
\proseitem{98}{\csromanfolio{626}อาจยคามิํ สรณํ ธมฺมํ ปฏิจฺจ อาจยคามี สรโณ ธมฺโม อุปฺปชฺชติ เหตุปจฺจยา. (สํขิตฺตํ.)}

\prose{เหตุยา เอกํ ฯเปฯ อวิคเต เอกํ. (สํขิตฺตํ.) (สหชาตวาเรปิ ฯเปฯ ปญฺหาวาเรปิ สพฺพตฺถ เอกํ.)}

\proseitem{99}{อาจยคามิํ อรณํ ธมฺมํ ปฏิจฺจ อาจยคามี อรโณ ธมฺโม อุปฺปชฺชติ เหตุปจฺจยา.\csromanspacer{} ตีณิ.}

\prose{อปจยคามิํ อรณํ ธมฺมํ ปฏิจฺจ อปจยคามี อรโณ ธมฺโม อุปฺปชฺชติ เหตุปจฺจยา.\csromanspacer{} ตีณิ.}

\prose{เนวาจยคามินาปจยคามิํ อรณํ ธมฺมํ ปฏิจฺจ เนวาจยคามินาปจยคามี อรโณ ธมฺโม อุปฺปชฺชติ เหตุปจฺจยา. (1)}

\prose{อาจยคามิํ อรณญฺจ เนวาจยคามินาปจยคามิํ อรณญฺจ ธมฺมํ ปฏิจฺจ เนวาจยคามินาปจยคามี อรโณ ธมฺโม อุปฺปชฺชติ เหตุปจฺจยา. (1)}

\prose{อปจยคามิํ อรณญฺจ เนวาจยคามินาปจยคามิํ อรณญฺจ ธมฺมํ ปฏิจฺจ เนวาจยคามินาปจยคามี อรโณ ธมฺโม อุปฺปชฺชติ เหตุปจฺจยา. (1) (สํขิตฺตํ.)}

\proseitem{100}{เหตุยา นว, อารมฺมเณ ตีณิ, อธิปติยา นว ฯเปฯ อวิคเต นว. (สํขิตฺตํ.\csromanspacer{} สหชาตวารมฺปิ ฯเปฯ สมฺปยุตฺตวารมฺปิ ปฏิจฺจวารสทิสํ.)}

\proseitem{101}{อาจยคามี อรโณ ธมฺโม อาจยคามิสฺส อรณสฺส ธมฺมสฺส เหตุปจฺจเยน ปจฺจโย.\csromanspacer{} ตีณิ.}

\csromanheader{2}{11. เสกฺขตฺติก 100. สรณทุก}
\prose{\csromanfolio{627}อปจยคามี อรโณ ธมฺโม อปจยคามิสฺส อรณสฺส ธมฺมสฺส เหตุปจฺจเยน ปจฺจโย.\csromanspacer{} ตีณิ.}

\prose{เนวาจยคามินาปจยคามี อรโณ ธมฺโม เนวาจยคามินาปจยคามิสฺส อรณสฺส ธมฺมสฺส เหตุปจฺจเยน ปจฺจโย. (1)}

\prose{อาจยคามี อรโณ ธมฺโม อาจยคามิสฺส อรณสฺส ธมฺมสฺส อารมฺมณปจฺจเยน ปจฺจโย. (สํขิตฺตํ.)}

\proseitem{102}{เหตุยา สตฺต, อารมฺมเณ สตฺต, อธิปติยา ทส, อนนฺตเร ฉ ฯเปฯ อวิคเต เตรส. (สํขิตฺตํ.)}

\prose{(ยถา กุสลตฺติเก ปญฺหาวารํ, เอวํ วิตฺถาเรตพฺพํ.)}

\csromanheader{2}{11. เสกฺขตฺติก 100. สรณทุก}
\csromanheader{2}{1-7. ปฏิจฺจาทิวาร-เหตุ}
\proseitem{103}{เนวเสกฺขนาเสกฺขํ สรณํ ธมฺมํ ปฏิจฺจ เนวเสกฺขนาเสกฺโข สรโณ ธมฺโม อุปฺปชฺชติ เหตุปจฺจยา. (สํขิตฺตํ.)}

\prose{เหตุยา เอกํ ฯเปฯ อวิคเต เอกํ. (สํขิตฺตํ.) (สหชาตวาเรปิ ฯเปฯ ปญฺหาวาเรปิ สพฺพตฺถ เอกํ.)}

\proseitem{104}{เสกฺขํ อรณํ ธมฺมํ ปฏิจฺจ เสกฺโข อรโณ ธมฺโม อุปฺปชฺชติ เหตุปจฺจยา.\csromanspacer{}\csromanspacer{}เสกฺขํ อรณํ ธมฺมํ ปฏิจฺจ เนวเสกฺขนาเสกฺโข อรโณ ธมฺโม อุปฺปชฺชติ เหตุปจฺจยา.\csromanspacer{}\csromanspacer{}เสกฺขํ อรณํ ธมฺมํ ปฏิจฺจ เสกฺโข อรโณ จ เนวเสกฺขนาเสกฺโข อรโณ จ ธมฺมา อุปฺปชฺชนฺติ เหตุปจฺจยา. (3)}

\prose{อเสกฺขํ อรณํ ธมฺมํ ปฏิจฺจ อเสกฺโข อรโณ ธมฺโม อุปฺปชฺชติ เหตุปจฺจยา.\csromanspacer{} ตีณิ.}

\prose{เนวเสกฺขนาเสกฺขํ อรณํ ธมฺมํ ปฏิจฺจ เนวเสกฺขนาเสกฺโข อรโณ ธมฺโม อุปฺปชฺชติ เหตุปจฺจยา. (1)}

\prose{เสกฺขํ อรณญฺจ เนวเสกฺขนาเสกฺขํ อรณญฺจ ธมฺมํ ฯเปฯ.\csromanspacer{} เหตุปจฺจยา.}

\gambhira{ธมฺมานุโลม ติกทุกปฏฺฐานปาฬิ}
\prose{\csromanfolio{628}อเสกฺขํ อรณญฺจ เนวเสกฺขนาเสกฺขํ อรณญฺจ ธมฺมํ ปฏิจฺจ เนวเสกฺขนาเสกฺโข อรโณ ธมฺโม อุปฺปชฺชติ เหตุปจฺจยา. (1) (สํขิตฺตํ.)}

\proseitem{105}{เหตุยา นว, อารมฺมเณ ตีณิ ฯเปฯ อวิคเต นว. (สํขิตฺตํ.) (สหชาตวารมฺปิ ฯเปฯ สมฺปยุตฺตวารมฺปิ ปฏิจฺจวารสทิสํ.)}

\proseitem{106}{เสกฺโข อรโณ ธมฺโม เสกฺขสฺส อรณสฺส ธมฺมสฺส เหตุปจฺจเยน ปจฺจโย.\csromanspacer{} ตีณิ.}

\prose{อเสกฺโข อรโณ ธมฺโม อเสกฺขสฺส อรณสฺส ธมฺมสฺส เหตุปจฺจเยน ปจฺจโย.\csromanspacer{} ตีณิ.}

\prose{เนวเสกฺขนาเสกฺโข อรโณ ธมฺโม เนวเสกฺขนาเสกฺขสฺส อรณสฺส ธมฺมสฺส เหตุปจฺจเยน ปจฺจโย. (1) (สํขิตฺตํ.)}

\proseitem{107}{เหตุยา สตฺต, อารมฺมเณ ปญฺจ, อธิปติยา นว, อนนฺตเร อฏฺฐ ฯเปฯ อวิคเต เตรส. (สํขิตฺตํ.)}

\prose{(ยถา กุสลตฺติเก ปญฺหาวารํ, เอวํ วิตฺถาเรตพฺพํ.)}

\csromanheader{2}{12. ปริตฺตตฺติก 100. สรณทุก}
\csromanheader{2}{1. ปฏิจฺจวาร-เหตุ}
\proseitem{108}{ปริตฺตํ สรณํ ธมฺมํ ปฏิจฺจ ปริตฺโต สรโณ ธมฺโม อุปฺปชฺชติ เหตุปจฺจยา. (สํขิตฺตํ.)}

\prose{เหตุยา เอกํ ฯเปฯ อวิคเต เอกํ. (สํขิตฺตํ.) (สหชาตวาเรปิ ฯเปฯ ปญฺหาวาเรปิ สพฺพตฺถ เอกํ.)}

\proseitem{109}{ปริตฺตํ อรณํ ธมฺมํ ปฏิจฺจ ปริตฺโต อรโณ ธมฺโม อุปฺปชฺชติ เหตุปจฺจยา.\csromanspacer{} ตีณิ.}

\prose{มหคฺคตํ อรณํ ธมฺมํ ปฏิจฺจ มหคฺคโต อรโณ ธมฺโม อุปฺปชฺชติ เหตุปจฺจยา.\csromanspacer{} ตีณิ.}

\csromanheader{2}{13. ปริตฺตารมฺมณตฺติก 100. สรณทุก}
\prose{\csromanfolio{629}อปฺปมาณํ อรณํ ธมฺมํ ปฏิจฺจ อปฺปมาโณ อรโณ ธมฺโม อุปฺปชฺชติ เหตุปจฺจยา.\csromanspacer{} ตีณิ.}

\prose{ปริตฺตํ อรณญฺจ มหคฺคตํ อรณญฺจ ธมฺมํ ปฏิจฺจ ปริตฺโต อรโณ ธมฺโม อุปฺปชฺชติ เหตุปจฺจยา.\csromanspacer{} ตีณิ.}

\prose{ปริตฺตํ อรณญฺจ อปฺปมาณํ อรณญฺจ ธมฺมํ ปฏิจฺจ ปริตฺโต อรโณ ธมฺโม อุปฺปชฺชติ เหตุปจฺจยา. (1) (สํขิตฺตํ.)}

\proseitem{110}{เหตุยา เตรส, อารมฺมเณ ปญฺจ ฯเปฯ อวิคเต เตรส. (สํขิตฺตํ.\csromanspacer{} สหชาตวาเรปิ ฯเปฯ ปญฺหาวาเรปิ สพฺพตฺถ วิตฺถาโร.)}

\csromanheader{2}{13. ปริตฺตารมฺมณตฺติก 100. สรณทุก}
\csromanheader{2}{1. ปฏิจฺจวาร-เหตุ}
\proseitem{111}{ปริตฺตารมฺมณํ สรณํ ธมฺมํ ปฏิจฺจ ปริตฺตารมฺมโณ สรโณ ธมฺโม อุปฺปชฺชติ เหตุปจฺจยา. (1)}

\prose{มหคฺคตารมฺมณํ สรณํ ธมฺมํ ปฏิจฺจ มหคฺคตารมฺมโณ สรโณ ธมฺโม อุปฺปชฺชติ เหตุปจฺจยา. (1) (สํขิตฺตํ.)}

\prose{เหตุยา เทฺว ฯเปฯ อวิคเต เทฺว. (สํขิตฺตํ.) (สหชาตวาเรปิ ฯเปฯ ปญฺหาวาเรปิ สพฺพตฺถ วิตฺถาโร.)}

\proseitem{112}{ปริตฺตารมฺมณํ อรณํ ธมฺมํ ปฏิจฺจ ปริตฺตารมฺมโณ อรโณ ธมฺโม อุปฺปชฺชติ เหตุปจฺจยา. (1)}

\prose{มหคฺคตารมฺมณํ อรณํ ธมฺมํ ปฏิจฺจ มหคฺคตารมฺมโณ อรโณ ธมฺโม อุปฺปชฺชติ เหตุปจฺจยา. (1)}

\prose{อปฺปมาณารมฺมณํ อรณํ ธมฺมํ ปฏิจฺจ อปฺปมาณารมฺมโณ อรโณ ธมฺโม อุปฺปชฺชติ เหตุปจฺจยา. (1) (สํขิตฺตํ.)}

\csromanheader{2}{เหตุยา ตีณิ ฯเปฯ อวิคเต ตีณิ. (สํขิตฺตํ.)}
\vspace{\tipitakaparskip}
\csromancenter{(สหชาตวาเรปิ ฯเปฯ ปญฺหาวาเรปิ สพฺพตฺถ วิตฺถาโร.)}
\csromanheader{2}{ธมฺมานุโลม ติกทุกปฏฺฐานปาฬิ \textbf{14. หีนตฺติก 100. สรณทุก}}
\csromanheader{2}{1. ปฏิจฺจวาร-เหตุ}
\proseitem{113}{\csromanfolio{630}หีนํ สรณํ ธมฺมํ ปฏิจฺจ หีโน สรโณ ธมฺโม อุปฺปชฺชติ เหตุปจฺจยา. (สํขิตฺตํ.)}

\csromanheader{2}{เหตุยา เอกํ ฯเปฯ อวิคเต เอกํ. (สํขิตฺตํ.)}
\proseitem{114}{มชฺฌิมํ อรณํ ธมฺมํ ปฏิจฺจ มชฺฌิโม อรโณ ธมฺโม อุปฺปชฺชติ เหตุปจฺจยา. (1)}

\prose{ปณีตํ อรณํ ธมฺมํ ปฏิจฺจ ปณีโต อรโณ ธมฺโม อุปฺปชฺชติ เหตุปจฺจยา.\csromanspacer{} ตีณิ.}

\prose{มชฺฌิมํ อรณญฺจ ปณีตํ อรณญฺจ ธมฺมํ ปฏิจฺจ มชฺฌิโม อรโณ ธมฺโม อุปฺปชฺชติ เหตุปจฺจยา. (1) (สํขิตฺตํ.)}

\csromanheader{2}{เหตุยา ปญฺจ ฯเปฯ อวิคเต ปญฺจ. (สํขิตฺตํ.)}
\vspace{\tipitakaparskip}
\csromancenter{(สหชาตวาเรปิ ฯเปฯ ปญฺหาวาเรปิ สพฺพตฺถ วิตฺถาโร.)}
\csromanheader{2}{15. มิจฺฉตฺตนิยตตฺติก 100. สรณทุก}
\csromanheader{2}{1. ปฏิจฺจวาร-เหตุ}
\proseitem{115}{มิจฺฉตฺตนิยตํ สรณํ ธมฺมํ ปฏิจฺจ มิจฺฉตฺตนิยโต สรโณ ธมฺโม อุปฺปชฺชติ เหตุปจฺจยา. (1)}

\prose{อนิยตํ สรณํ ธมฺมํ ปฏิจฺจ อนิยโต สรโณ ธมฺโม อุปฺปชฺชติ เหตุปจฺจยา. (1) (สํขิตฺตํ.)}

\prose{เหตุยา เทฺว ฯเปฯ อวิคเต เทฺว. (สํขิตฺตํ.) (สหชาตวาเรปิ ฯเปฯ ปญฺหาวาเรปิ สพฺพตฺถ วิตฺถาโร.)}

\proseitem{116}{สมฺมตฺตนิยตํ อรณํ ธมฺมํ ปฏิจฺจ สมฺมตฺตนิยโต อรโณ ธมฺโม อุปฺปชฺชติ เหตุปจฺจยา.\csromanspacer{}\csromanspacer{}สมฺมตฺตนิยตํ อรณํ ธมฺมํ ปฏิจฺจ อนิยโต}

\vspace{\tipitakaparskip}
\csromancenter{มคฺคารมฺมณตฺติก 100.\csromanspacer{} สรณทุก อรโณ ธมฺโม อุปฺปชฺชติ เหตุปจฺจยา.\csromanspacer{}\csromanspacer{}สมฺมตฺตนิยตํ อรณํ ธมฺมํ ปฏิจฺจ สมฺมตฺตนิยโต อรโณ จ อนิยโต อรโณ จ ธมฺมา อุปฺปชฺชนฺติ เหตุปจฺจยา. (3)}
\prose{\csromanfolio{631}อนิยตํ อรณํ ธมฺมํ ปฏิจฺจ อนิยโต อรโณ ธมฺโม อุปฺปชฺชติ เหตุปจฺจยา. (1)}

\prose{สมฺมตฺตนิยตํ อรณญฺจ อนิยตํ อรณญฺจ ธมฺมํ ปฏิจฺจ อนิยโต อรโณ ธมฺโม อุปฺปชฺชติ เหตุปจฺจยา. (1) (สํขิตฺตํ.)}

\prose{เหตุยา ปญฺจ, อารมฺมเณ เทฺว ฯเปฯ อวิคเต ปญฺจ. (สํขิตฺตํ.)}

\vspace{\tipitakaparskip}
\csromancenter{(สหชาตวารมฺปิ ฯเปฯ ปญฺหาวารมฺปิ วิตฺถาเรตพฺพํ.)}
\csromanheader{2}{16. มคฺคารมฺมณตฺติก 100. สรณทุก}
\csromanheader{2}{1-7. ปฏิจฺจาทิวาร-เหตุ}
\proseitem{117}{มคฺคารมฺมณํ อรณํ ธมฺมํ ปฏิจฺจ มคฺคารมฺมโณ อรโณ ธมฺโม อุปฺปชฺชติ เหตุปจฺจยา.\csromanspacer{} ตีณิ.}

\prose{มคฺคเหตุกํ อรณํ ธมฺมํ ปฏิจฺจ มคฺคเหตุโก อรโณ ธมฺโม อุปฺปชฺชติ เหตุปจฺจยา.\csromanspacer{} ตีณิ.}

\prose{มคฺคาธิปติํ อรณํ ธมฺมํ ปฏิจฺจ มคฺคาธิปติ อรโณ ธมฺโม อุปฺปชฺชติ เหตุปจฺจยา.\csromanspacer{} ปญฺจ.}

\prose{มคฺคารมฺมณํ อรณญฺจ มคฺคาธิปติํ อรณญฺจ ธมฺมํ ปฏิจฺจ มคฺคารมฺมโณ อรโณ ธมฺโม อุปฺปชฺชติ เหตุปจฺจยา.\csromanspacer{} ตีณิ.}

\prose{มคฺคเหตุกํ อรณญฺจ มคฺคาธิปติํ อารณญฺจ ธมฺมํ ปฏิจฺจ มคฺคเหตุโก อรโณ ธมฺโม อุปฺปชฺชติ เหตุปจฺจยา.\csromanspacer{} ตีณิ. (สํขิตฺตํ.)}

\csromanheader{2}{เหตุยา สตฺตรส ฯเปฯ อวิคเต สตฺตรส. (สํขิตฺตํ.)}
\prose{(สหชาตวารมฺปิ ฯเปฯ สมฺปยุตฺตวารมฺปิ ปฏิจฺจวารสทิสํ.)}

\gambhira{ธมฺมานุโลม ติกทุกปฏฺฐานปาฬิ}
\proseitem{118}{\csromanfolio{632}มคฺคารมฺมโณ อรโณ ธมฺโม มคฺคารมฺมณสฺส อรณสฺส ธมฺมสฺส เหตุปจฺจเยน ปจฺจโย.\csromanspacer{} ตีณิ.}

\prose{มคฺคเหตุโก อรโณ ธมฺโม มคฺคเหตุกสฺส อรณสฺส ธมฺมสฺส เหตุปจฺจเยน ปจฺจโย.\csromanspacer{} ตีณิ.}

\prose{มคฺคาธิปติ อรโณ ธมฺโม มคฺคาธิปติสฺส อรณสฺส ธมฺมสฺส เหตุปจฺจเยน ปจฺจโย.\csromanspacer{} ปญฺจ.}

\prose{มคฺคารมฺมโณ อรโณ จ มคฺคาธิปติ อรโณ จ ธมฺมา มคฺคารมฺมณสฺส อรณสฺส ธมฺมสฺส เหตุปจฺจเยน ปจฺจโย.\csromanspacer{} ตีณิ.}

\prose{มคฺคเหตุโก อรโณ จ มคฺคาธิปติ อรโณ จ ธมฺมา มคฺคเหตุกสฺส อรณสฺส ธมฺมสฺส เหตุปจฺจเยน ปจฺจโย.\csromanspacer{} ตีณิ. (สํขิตฺตํ.)}

\proseitem{119}{เหตุยา สตฺตรส, อารมฺมเณ นว ฯเปฯ อวิคเต สตฺตรส. (สํขิตฺตํ.)}

\prose{(ยถา กุสลตฺติเก ปญฺหาวารํ, เอวํ วิตฺถาเรตพฺพํ.)}

\csromanheader{2}{17. อุปฺปนฺนตฺติก 100. สรณทุก}
\csromanheader{2}{7. ปญฺหาวาร-เหตุ อารมฺมณ}
\proseitem{120}{อุปฺปนฺโน สรโณ ธมฺโม อุปฺปนฺนสฺส สรณสฺส ธมฺมสฺส เหตุปจฺจเยน ปจฺจโย. (สํขิตฺตํ.)}

\prose{เหตุยา เอกํ, อารมฺมเณ เทฺว, อธิปติยา เทฺว, สหชาเต อญฺญมญฺเญ นิสฺสเย เอกํ, อุปนิสฺสเย เทฺว, กมฺเม ฯเปฯ สมฺปยุตฺเต เอกํ ฯเปฯ อวิคเต เอกํ. (สํขิตฺตํ.)}

\proseitem{121}{อุปฺปนฺโน อรโณ ธมฺโม อุปฺปนฺนสฺส อรณสฺส ธมฺมสฺส เหตุปจฺจเยน ปจฺจโย. (สํขิตฺตํ.)}

\csromanheader{2}{19. อตีตารมฺมณตฺติก 100. สรณทุก}
\prose{\csromanfolio{633}เหตุยา เอกํ, อารมฺมเณ ตีณิ, อธิปติยา ตีณิ ฯเปฯ อุปนิสฺสเย ตีณิ ฯเปฯ อวิคเต เอกํ. (สํขิตฺตํ.)}

\csromanheader{2}{18. อตีตตฺติก 100. สรณทุก}
\csromanheader{2}{7. ปญฺหาวาร-เหตุ}
\proseitem{122}{ปจฺจุปฺปนฺโน สรโณ ธมฺโม ปจฺจุปฺปนฺนสฺส สรณสฺส ธมฺมสฺส เหตุปจฺจเยน ปจฺจโย. (สํขิตฺตํ.)}

\prose{เหตุยา เอกํ, อารมฺมเณ ตีณิ, อธิปติยา ตีณิ ฯเปฯ อุปนิสฺสเย ตีณิ ฯเปฯ อวิคเต เอกํ. (สํขิตฺตํ.)}

\proseitem{123}{ปจฺจุปฺปนฺโน อรโณ ธมฺโม ปจฺจุปฺปนฺนสฺส อรณสฺส ธมฺมสฺส เหตุปจฺจเยน ปจฺจโย. (สํขิตฺตํ.)}

\csromanhangingbegin
\hangingheaditem{123}{เหตุยา เอกํ, อารมฺมเณ ตีณิ ฯเปฯ อวิคเต เอกํ. (สํขิตฺตํ.)}
\hangingbody{(ยถา กุสลตฺติเก ปญฺหาวารํ, เอวํ วิตฺถาเรตพฺพํ.)}
\csromanhangingend

\csromanheader{2}{19. อตีตารมฺมณตฺติก 100. สรณทุก}
\csromanheader{2}{1-7. ปฏิจฺจาทิวาร-เหตุ อารมฺมณ}
\proseitem{124}{อตีตารมฺมณํ สรณํ ธมฺมํ ปฏิจฺจ อตีตารมฺมโณ สรโณ ธมฺโม อุปฺปชฺชติ เหตุปจฺจยา. (1)}

\prose{อนาคตารมฺมณํ สรณํ ธมฺมํ ปฏิจฺจ อนาคตารมฺมโณ สรโณ ธมฺโม อุปฺปชฺชติ เหตุปจฺจยา. (1)}

\prose{ปจฺจุปฺปนฺนารมฺมณํ สรณํ ธมฺมํ ปฏิจฺจ ปจฺจุปฺปนฺนารมฺมโณ สรโณ ธมฺโม อุปฺปชฺชติ เหตุปจฺจยา. (1) (สํขิตฺตํ.)}

\csromanheader{2}{เหตุยา ตีณิ ฯเปฯ อวิคเต ตีณิ. (สํขิตฺตํ.)}
\prose{(สหชาตวารมฺปิ ฯเปฯ ปญฺหาวารมฺปิ สพฺพตฺถ วิตฺถาเรตพฺพํ.)}

\gambhira{ธมฺมานุโลม ติกทุกปฏฺฐานปาฬิ}
\proseitem{125}{\csromanfolio{634}อตีตารมฺมณํ อรณํ ธมฺมํ ปฏิจฺจ อตีตารมฺมโณ อรโณ ธมฺโม อุปฺปชฺชติ เหตุปจฺจยา. (1)}

\prose{อนาคตารมฺมณํ อรณํ ธมฺมํ ปฏิจฺจ อนาคตารมฺมโณ อรโณ ธมฺโม อุปฺปชฺชติ เหตุปจฺจยา. (1)}

\prose{ปจฺจุปฺปนฺนารมฺมณํ อรณํ ธมฺมํ ปฏิจฺจ ปจฺจุปฺปนฺนารมฺมโณ อรโณ ธมฺโม อุปฺปชฺชติ เหตุปจฺจยา. (1) (สํขิตฺตํ.)}

\prose{เหตุยา ตีณิ ฯเปฯ อวิคเต ตีณิ. (สํขิตฺตํ.) (สหชาตวารมฺปิ ฯเปฯ สมฺปยุตฺตวารมฺปิ ปฏิจฺจวารสทิสํ.)}

\proseitem{126}{อตีตารมฺมโณ อรโณ ธมฺโม อตีตารมฺมณสฺส อรณสฺส ธมฺมสฺส เหตุปจฺจเยน ปจฺจโย. (1)}

\prose{อนาคตารมฺมโณ อรโณ ธมฺโม อนาคตารมฺมณสฺส อรณสฺส ธมฺมสฺส เหตุปจฺจเยน ปจฺจโย. (1)}

\prose{ปจฺจุปฺปนฺนารมฺมโณ อรโณ ธมฺโม ปจฺจุปฺปนฺนารมฺมณสฺส อรณสฺส ธมฺมสฺส เหตุปจฺจเยน ปจฺจโย. (1)}

\prose{อตีตารมฺมโณ อรโณ ธมฺโม อตีตารมฺมณสฺส อรณสฺส ธมฺมสฺส อารมฺมณปจฺจเยน ปจฺจโย. (สํขิตฺตํ.)}

\proseitem{127}{เหตุยา ตีณิ, อารมฺมเณ นว ฯเปฯ อวิคเต ตีณิ. (สํขิตฺตํ.\csromanspacer{} ยถา กุสลตฺติเก ปญฺหาวารํ, เอวํ วิตฺถาเรตพฺพํ.)}

\csromanheader{2}{20. อชฺฌตฺตตฺติก 100. สรณทุก}
\csromanheader{2}{1-7. ปฏิจฺจาทิวาร-เหตุ}
\vspace{\tipitakaparskip}
\csromancenter{อชฺฌตฺตํ สรณํ ธมฺมํ ปฏิจฺจ อชฺฌตฺโต สรโณ ธมฺโม อุปฺปชฺชติ เหตุปจฺจยา.}
\csromancenter{อชฺฌตฺตารมฺมณตฺติก 100.\csromanspacer{} สรณทุก พหิทฺธา สรณํ ธมฺมํ ปฏิจฺจ พหิทฺธา สรโณ ธมฺโม อุปฺปชฺชติ เหตุปจฺจยา.}
\prose{\csromanfolio{635}เหตุยา เทฺว ฯเปฯ อวิคเต เทฺว. (สํขิตฺตํ.) (สหชาตวาเรปิ ฯเปฯ ปญฺหาวาเรปิ สพฺพตฺถ วิตฺถาโร.)}

\proseitem{129}{อชฺฌตฺตํ อรณํ ธมฺมํ ปฏิจฺจ อชฺฌตฺโต อรโณ ธมฺโม อุปฺปชฺชติ เหตุปจฺจยา. (1)}

\prose{พหิทฺธา อรณํ ธมฺมํ ปฏิจฺจ พหิทฺธา อรโณ ธมฺโม อุปฺปชฺชติ เหตุปจฺจยา. (1) (สํขิตฺตํ.)}

\csromanheader{2}{เหตุยา เทฺว ฯเปฯ อวิคเต เทฺว. (สํขิตฺตํ.)}
\vspace{\tipitakaparskip}
\csromancenter{(สหชาตวาเรปิ ฯเปฯ ปญฺหาวาเรปิ สพฺพตฺถ วิตฺถาโร.)}
\csromanheader{2}{21. อชฺฌตฺตารมฺมณตฺติก 100. สรณทุก}
\csromanheader{2}{1-7. ปฏิจฺจาทิวาร-เหตุ}
\proseitem{130}{อชฺฌตฺตารมฺมณํ สรณํ ธมฺมํ ปฏิจฺจ อชฺฌตฺตารมฺมโณ สรโณ ธมฺโม อุปฺปชฺชติ เหตุปจฺจยา. (1)}

\prose{พหิทฺธารมฺมณํ สรณํ ธมฺมํ ปฏิจฺจ พหิทฺธารมฺมโณ สรโณ ธมฺโม อุปฺปชฺชติ เหตุปจฺจยา. (1) (สํขิตฺตํ.)}

\prose{เหตุยา เทฺว ฯเปฯ อวิคเต เทฺว. (สํขิตฺตํ.) (สหชาตวาเรปิ ฯเปฯ ปญฺหาวาเรปิ สพฺพตฺถ วิตฺถาโร.)}

\proseitem{131}{อชฺฌตฺตารมฺมณํ อรณํ ธมฺมํ ปฏิจฺจ อชฺฌตฺตารมฺมโณ อรโณ ธมฺโม อุปฺปชฺชติ เหตุปจฺจยา. (1)}

\prose{พหิทฺธารมฺมณํ อรณํ ธมฺมํ ปฏิจฺจ พหิทฺธารมฺมโณ อรโณ ธมฺโม อุปฺปชฺชติ เหตุปจฺจยา. (1) (สํขิตฺตํ.)}

\csromanheader{2}{เหตุยา เทฺว ฯเปฯ อวิคเต เทฺว. (สํขิตฺตํ.)}
\vspace{\tipitakaparskip}
\csromancenter{(สหชาตวาเรปิ ฯเปฯ ปญฺหาวาเรปิ สพฺพตฺถ วิตฺถาโร.)}
\csromanmatikaanchor{mtk.153}
\csromanmatikaanchor{mtk.154}
\csromantocmarknopage{subparagraph}{22. สนิทสฺสนตฺติก}{mtk.153}
\bookmark[dest={mtk.153},level=5]{22. สนิทสฺสนตฺติก}
\csromantocmarkref{subparagraph}{100. สรณทุก}{mtk.154}
\bookmark[dest={mtk.154},level=5]{100. สรณทุก}
\csromanheader{6}{ธมฺมานุโลม ติกทุกปฏฺฐานปาฬิ \textbf{22. สนิทสฺสนตฺติก 100. สรณทุก}}
\csromanheader{2}{1-7. ปฏิจฺจาทิวาร-เหตุ}
\proseitem{132}{\csromanfolio{636}อนิทสฺสน-อปฺปฏิฆํ สรณํ ธมฺมํ ปฏิจฺจ อนิทสฺสน-อปฺปฏิโฆ สรโณ ธมฺโม อุปฺปชฺชติ เหตุปจฺจยา. (สพฺพตฺถ เอกํ.)}

\prose{อนิทสฺสน-อปฺปฏิฆํ อรณํ ธมฺมํ ปฏิจฺจ อนิทสฺสน-อปฺปฏิโฆ อรโณ ธมฺโม อุปฺปชฺชติ เหตุปจฺจยา.\csromanspacer{} สตฺต. (สํขิตฺตํ.)}

\proseitem{133}{เหตุยา เอกวีส ฯเปฯ อวิคเต เอกวีส. (สํขิตฺตํ.) (สหชาตวารมฺปิ ฯเปฯ สมฺปยุตฺตวารมฺปิ ปฏิจฺจวารสทิสํ.)}

\proseitem{134}{อนิทสฺสน-อปฺปฏิโฆ อรโณ ธมฺโม อนิทสฺสน-อปฺปฏิฆสฺส อรณสฺส ธมฺมสฺส เหตุปจฺจเยน ปจฺจโย. (สํขิตฺตํ.)}

\prose{เหตุยา สตฺต, อารมฺมเณ ตีณิ ฯเปฯ อวิคเต ปญฺจวีส. (สํขิตฺตํ.) (ยถา กุสลตฺติเก ปญฺหาวารํ, เอวํ วิตฺถาเรตพฺพํ.)}

\nitthitam{สนิทสฺสนตฺติกสรณทุกํ นิฏฺฐิตํ.}
\nitthitam{\textbf{ธมฺมานุโลเม ติกทุกปฏฺฐานํ นิฏฺฐิตํ.}}
